% 本文件由 LaTeX 排版 Agent 根据有效 translation.json 自动生成。
% author=Augustine of Hippo; work_id=1102; title=希坡的圣奥斯定书信集

\chapter{\Place{希坡}的圣\Person{奥斯定}书信集}

% structure: p0001 source=h1 target=omitted reason=page-title-already-rendered-by-parent

% structure: p0002 source=h2 target=section reason=relative-heading-rank; base=h2

\section{第一部分（公元386–396年）}

\Emph{\Person{奥斯定}成为\Place{希坡}主教前所写的书信}\par

\Emph{书信一}（386年）\Person{奥斯定}致\Person{赫莫革尼安}\newline{}\Emph{书信二}（386年）\Person{奥斯定}致\Person{泽诺比乌斯}\newline{}\Emph{书信三}（387年）\Person{奥斯定}致\Person{内布里迪乌斯}\newline{}\Emph{书信四}（387年）\Person{奥斯定}致\Person{内布里迪乌斯}\newline{}\Emph{书信五}（388年）\Person{内布里迪乌斯}致\Person{奥斯定}\newline{}\Emph{书信六}（389年）\Person{内布里迪乌斯}致\Person{奥斯定}\newline{}\Emph{书信七}（389年）\Person{奥斯定}致\Person{内布里迪乌斯}\newline{}\Emph{书信八}（389年）\Person{内布里迪乌斯}致\Person{奥斯定}\newline{}\Emph{书信九}（389年）\Person{奥斯定}致\Person{内布里迪乌斯}\newline{}\Emph{书信十}（389年）\Person{奥斯定}致\Person{内布里迪乌斯}\newline{}\Emph{书信十一}（389年）\Person{奥斯定}致\Person{内布里迪乌斯}\newline{}\Emph{书信十三}（389年）\Person{奥斯定}致\Person{内布里迪乌斯}\newline{}\Emph{书信十四}（389年）\Person{奥斯定}致\Person{内布里迪乌斯}\newline{}\Emph{书信十五}（390年）\Person{奥斯定}致\Person{罗马尼安}\newline{}\Emph{书信十六}（390年）\Place{马道拉}的\Person{玛息莫}致\Person{奥斯定}\newline{}\Emph{书信十七}（390年）\Person{奥斯定}致\Place{马道拉}的\Person{玛息莫}\newline{}\Emph{书信十八}（390年）\Person{奥斯定}致\Person{塞莱斯廷}\newline{}\Emph{书信十九}（390年）\Person{奥斯定}致\Person{加约}\newline{}\Emph{书信二十}（390年）\Person{奥斯定}致\Person{安多尼努斯}\newline{}\Emph{书信二十一}（391年）\Person{奥斯定}致\Person{瓦莱里}\newline{}\Emph{书信二十二}（392年）\Person{奥斯定}致\Person{奥勒留}\newline{}\Emph{书信二十三}（392年）\Person{奥斯定}致\Person{玛息明}\newline{}\Emph{书信二十五}（394年）\Person{保利努斯}与\Person{泰拉西亚}致\Person{奥斯定}\newline{}\Emph{书信二十六}（395年）\Person{奥斯定}致\Person{利森修斯}\newline{}\Emph{书信二十七}（395年）\Person{奥斯定}致\Person{保利努斯}\newline{}\Emph{书信二十八}（394或395年）\Person{奥斯定}致\Person{热罗尼莫}\newline{}\Emph{书信二十九}（395年）\Person{奥斯定}致\Person{阿利比乌斯}\newline{}\Emph{书信三十}（396年）\Person{保利努斯}与\Person{泰拉西亚}致\Person{奥斯定}\par

% structure: p0005 source=h2 target=section reason=relative-heading-rank; base=h2

\section{第二部分（公元396–410年）}

\Emph{\Person{奥斯定}成为\Place{希坡}主教后，但在与多纳特派举行会议以及白拉奇异端在非洲兴起之前所写的书信}\par

\Emph{第31封信}（396年） \Person{奥斯定}致\Person{保利努斯}与\Person{泰拉西亚}\newline{}\Emph{第33封信}（396年） \Person{奥斯定}致\Person{普罗库莱亚努斯}\newline{}\Emph{第34封信}（396年） \Person{奥斯定}致\Person{优西比乌斯}\newline{}\Emph{第35封信}（396年） \Person{奥斯定}致\Person{优西比乌斯}\newline{}\Emph{第36封信}（396年） \Person{奥斯定}致\Person{卡苏拉努斯}\newline{}\Emph{第37封信}（397年） \Person{奥斯定}致\Person{辛普利齐亚努斯}\newline{}\Emph{第38封信}（397年） \Person{奥斯定}致\Person{普罗富图鲁斯}\newline{}\Emph{第39封信}（397年） \Person{热罗尼莫}致\Person{奥斯定}\newline{}\Emph{第40封信}（397年） \Person{奥斯定}致\Person{热罗尼莫}\newline{}\Emph{第41封信}（397年） \Person{阿里庇乌斯}与\Person{奥斯定}致\Person{奥雷利乌斯}\newline{}\Emph{第42封信}（397年） \Person{奥斯定}致\Person{保利努斯}与\Person{泰拉西亚}\newline{}\Emph{第43封信}（397年） \Person{奥斯定}致\Person{格洛里乌斯}、\Person{埃留西乌斯}、两位\Person{费利克斯}、\Person{格拉玛提库斯}及其他众人\newline{}\Emph{第44封信}（398年） \Person{奥斯定}致\Person{埃留西乌斯}、\Person{格洛里乌斯}与两位\Person{费利克斯}\newline{}\Emph{第46封信}（398年） \Person{普布利科拉}致\Person{奥斯定}\newline{}\Emph{第47封信}（398年） \Person{奥斯定}致\Person{普布利科拉}\newline{}\Emph{第48封信}（398年） \Person{奥斯定}致\Person{优多克西乌斯}\newline{}\Emph{第50封信}（399年） \Person{奥斯定}致\Place{苏费克图姆}殖民地\newline{}\Emph{第51封信}（399年或400年） \Person{奥斯定}致\Person{克里斯皮努斯}\newline{}\Emph{第53封信}（400年） \Person{奥斯定}、\Person{福图纳图斯}与\Person{阿里庇乌斯}致\Person{格内罗苏斯}\newline{}\Emph{第54封信}（400年） \Person{奥斯定}致\Person{雅努阿里乌斯}\newline{}\Emph{第55封信}（400年） \Person{奥斯定}致\Person{雅努阿里乌斯}\newline{}\Emph{第58封信}（401年） \Person{奥斯定}致\Person{帕玛基乌斯}\newline{}\Emph{第59封信}（401年） \Person{奥斯定}致\Person{维克托里努斯}\newline{}\Emph{第60封信}（401年） \Person{奥斯定}致\Person{奥雷利乌斯}\newline{}\Emph{第61封信}（401年） \Person{奥斯定}致\Person{泰奥多鲁斯}\newline{}\Emph{第62封信}（401年） \Person{奥斯定}、\Person{阿里庇乌斯}与\Person{萨姆苏齐乌斯}致\Person{塞维鲁斯}\newline{}\Emph{第63封信}（401年） \Person{奥斯定}致\Person{塞维鲁斯}\newline{}\Emph{第64封信}（401年） \Person{奥斯定}致\Person{昆提亚努斯}\newline{}\Emph{第65封信}（402年） \Person{奥斯定}致\Person{克桑提普斯}\newline{}\Emph{第66封信}（402年） \Person{奥斯定}致\Person{克里斯皮努斯}\newline{}\Emph{第67封信}（402年） \Person{奥斯定}致\Person{热罗尼莫}\newline{}\Emph{第68封信}（402年） \Person{热罗尼莫}致\Person{奥斯定}\newline{}\Emph{第69封信}（402年） \Person{阿里庇乌斯}与\Person{奥斯定}致\Person{卡斯托里乌斯}\newline{}\Emph{第71封信}（403年） \Person{奥斯定}致\Person{热罗尼莫}\newline{}\Emph{第72封信}（404年） \Person{热罗尼莫}致\Person{奥斯定}\newline{}\Emph{第73封信}（404年） \Person{奥斯定}致\Person{热罗尼莫}\newline{}\Emph{第74封信}（404年） \Person{奥斯定}致\Person{普雷西迪乌斯}\newline{}\Emph{第75封信}（404年） \Person{热罗尼莫}致\Person{奥斯定}\newline{}\Emph{第76封信}（402年） \Person{奥斯定}致\OriginalTerm{多纳特派}{Donatists}\newline{}\Emph{第77封信}（404年） \Person{奥斯定}致\Person{费利克斯}与\Person{希拉里努斯}\newline{}\Emph{第78封信}（404年） \Person{奥斯定}致\Place{希波}教会\newline{}\Emph{第79封信}（404年） \Person{奥斯定}致\Person{福图纳图斯}的继任者\newline{}\Emph{第81封信}（405年） \Person{热罗尼莫}致\Person{奥斯定}\newline{}\Emph{第82封信}（405年） \Person{奥斯定}致\Person{热罗尼莫}\newline{}\Emph{第83封信}（405年） \Person{奥斯定}致\Person{阿里庇乌斯}\newline{}\Emph{第84封信}（405年） \Person{奥斯定}致\Person{诺瓦图斯}\newline{}\Emph{第85封信}（405年） \Person{奥斯定}致\Person{保禄}\newline{}\Emph{第86封信}（405年） \Person{奥斯定}致\Person{凯基利安努斯}\newline{}\Emph{第87封信}（405年） \Person{奥斯定}致\Person{埃梅里图斯}\newline{}\Emph{第88封信}（406年） \Place{希波}的神职班致\Person{雅努阿里乌斯}\newline{}\Emph{第89封信}（406年） \Person{奥斯定}致\Person{费斯图斯}\newline{}\Emph{第90封信}（408年） \Person{内克塔里乌斯}致\Person{奥斯定}\newline{}\Emph{第91封信}（408年） \Person{奥斯定}致\Person{内克塔里乌斯}\newline{}\Emph{第92封信}（408年） \Person{奥斯定}致\Person{伊塔利卡}\newline{}\Emph{第93封信}（408年） \Person{奥斯定}致\Person{文琴提乌斯}\newline{}\Emph{第95封信}（408年） \Person{奥斯定}致\Person{保利努斯}与\Person{泰拉西亚}\newline{}\Emph{第96封信}（408年） \Person{奥斯定}致\Person{奥林皮乌斯}\newline{}\Emph{第97封信}（408年） \Person{奥斯定}致\Person{奥林皮乌斯}\newline{}\Emph{第98封信}（408年） \Person{奥斯定}致\Person{波尼法爵}\newline{}\Emph{第99封信}（408年或409年初） \Person{奥斯定}致\Person{伊塔利卡}\newline{}\Emph{第100封信}（409年） \Person{奥斯定}致\Person{多纳图斯}\newline{}\Emph{第101封信}（409年） \Person{奥斯定}致\Person{梅莫尔}\newline{}\Emph{第102封信}（409年） \Person{奥斯定}致\Person{德奥格拉提亚斯}\newline{}\Emph{第103封信}（409年） \Person{内克塔里乌斯}致\Person{奥斯定}\newline{}\Emph{第104封信}（409年） \Person{奥斯定}致\Person{内克塔里乌斯}\newline{}\Emph{第111封信}（409年11月） \Person{奥斯定}致\Person{维克托里亚努斯}\newline{}\Emph{第115封信}（410年） \Person{奥斯定}致\Person{福图纳图斯}\newline{}\Emph{第116封信}（410年） \Person{奥斯定}致\Person{格内罗苏斯}\newline{}\Emph{第117封信}（410年） \Person{狄奥斯库若斯}致\Person{奥斯定}\newline{}\Emph{第118封信}（410年） \Person{奥斯定}致\Person{狄奥斯库若斯}\newline{}\Emph{第122封信}（410年） \Person{奥斯定}致\Place{希波}的民众\newline{}\Emph{第123封信}（410年） \Person{热罗尼莫}致\Person{奥斯定}\par

% structure: p0008 source=h2 target=section reason=relative-heading-rank; base=h2

\section{\Emph{第三部分（主后411--430年）}}

\Emph{书信，写于与\OriginalTerm{多纳特派}{Donatists}的会议之后及\OriginalTerm{白拉奇异端}{Pelagian heresy}兴起期间，即\Person{奥斯定}生命最后二十年间}\par

\Emph{第124封信}（411年）\Person{奥斯定}致\Person{阿尔比娜}、\Person{皮尼安}和\Person{梅拉尼亚}\newline{}\Emph{第125封信}（411年）\Person{奥斯定}致\Person{阿利皮乌斯}\newline{}\Emph{第126封信}（411年）\Person{奥斯定}致\Person{阿尔比娜}\newline{}\Emph{第130封信}（412年）\Person{奥斯定}致\Person{普罗芭}\newline{}\Emph{第131封信}（412年）\Person{奥斯定}致\Person{普罗芭}\newline{}\Emph{第132封信}（412年）\Person{奥斯定}致\Person{沃卢西亚努斯}\newline{}\Emph{第133封信}（412年）\Person{奥斯定}致\Person{玛尔切利努斯}\newline{}\Emph{第135封信}（412年）\Person{沃卢西亚努斯}致\Person{奥斯定}\newline{}\Emph{第136封信}（412年）\Person{玛尔切利努斯}致\Person{奥斯定}\newline{}\Emph{第137封信}（412年）\Person{奥斯定}致\Person{沃卢西亚努斯}\newline{}\Emph{第138封信}（412年）\Person{奥斯定}致\Person{玛尔切利努斯}\newline{}\Emph{第139封信}（412年）\Person{奥斯定}致\Person{玛尔切利努斯}\newline{}\Emph{第143封信}（412年）\Person{奥斯定}致\Person{玛尔切利努斯}\newline{}\Emph{第144封信}（412年）\Person{奥斯定}致\Place{基尔塔}民众\newline{}\Emph{第145封信}（412或413年）\Person{奥斯定}致\Person{安纳斯塔修斯}\newline{}\Emph{第146封信}（413年）\Person{奥斯定}致\Person{白拉奇}\newline{}\Emph{第148封信}（413年）\Person{奥斯定}致\Person{福尔图纳提安}\newline{}\Emph{第150封信}（413年）\Person{奥斯定}致\Person{普罗芭}和\Person{尤利亚纳}\newline{}\Emph{第151封信}（413或414年）\Person{奥斯定}致\Person{凯奇利安}\newline{}\Emph{第158封信}（414年）\Person{埃沃迪乌斯}致\Person{奥斯定}\newline{}\Emph{第159封信}（415年）\Person{奥斯定}致\Person{埃沃迪乌斯}\newline{}\Emph{第163封信}（414年）\Person{埃沃迪乌斯}致\Person{奥斯定}\newline{}\Emph{第164封信}（414年）\Person{奥斯定}致\Person{埃沃迪乌斯}\newline{}\Emph{第165封信}（410年）\Person{热罗尼莫}致\Person{玛尔切利努斯}和\Person{阿纳普西基亚}\newline{}\Emph{第166封信}（415年）\Person{奥斯定}致\Person{热罗尼莫}，论灵魂的起源\newline{}\Emph{第167封信}（415年）\Person{奥斯定}致\Person{热罗尼莫}，论雅各伯书二章十节\newline{}\Emph{第169封信}（415年）\Person{奥斯定}致\Person{埃沃迪乌斯}\newline{}\Emph{第172封信}（416年）\Person{热罗尼莫}致\Person{奥斯定}\newline{}\Emph{第173封信}（416年）\Person{奥斯定}致\Person{多纳特}\newline{}\Emph{第180封信}（416年）\Person{奥斯定}致\Person{奥凯安努斯}\newline{}\Emph{第185封信}（416年）\Person{奥斯定}致\Person{博尼法斯}\newline{}\Emph{第188封信}（416年）\Person{阿利皮乌斯}与\Person{奥斯定}致\Person{尤利亚纳}\newline{}\Emph{第189封信}（418年）\Person{奥斯定}致\Person{博尼法斯}\newline{}\Emph{第191封信}（418年）\Person{奥斯定}致\Person{西斯笃}\newline{}\Emph{第192封信}（418年）\Person{奥斯定}致\Person{塞莱斯廷}\newline{}\Emph{第195封信}（418年）\Person{热罗尼莫}致\Person{奥斯定}\newline{}\Emph{第201封信}（419年）皇帝\Person{霍诺里乌斯}与\Person{狄奥多西}致\Person{奥斯定}\newline{}\Emph{第202封信}（419年）\Person{热罗尼莫}致\Person{阿利皮乌斯}与\Person{奥斯定}\newline{}\Emph{第203封信}（420年）\Person{奥斯定}致\Person{拉尔古斯}\newline{}\Emph{第208封信}（423年）\Person{奥斯定}致\Person{费利西亚}\newline{}\Emph{第209封信}（423年）\Person{奥斯定}致\Person{塞莱斯廷}\newline{}\Emph{第210封信}（423年）\Person{奥斯定}致\Person{费利奇塔斯}和\Person{鲁斯提库斯}\newline{}\Emph{第211封信}（423年）\Person{奥斯定}致一座隐修院\newline{}\Emph{第212封信}（423年）\Person{奥斯定}致\Person{昆蒂利安}\newline{}\Emph{第213封信}（426年9月26日）选择圣\Person{奥斯定}继任者的会议记录\newline{}\Emph{第214封信}（426年）\Person{奥斯定}致\Person{瓦伦提乌斯}\newline{}\Emph{第215封信}（426年）\Person{奥斯定}致\Person{瓦伦提乌斯}\newline{}\Emph{第218封信}（426年）\Person{奥斯定}致\Person{帕拉提努斯}\newline{}\Emph{第219封信}（436年）\Person{奥斯定}、\Person{弗洛伦提乌斯}和\Person{塞孔迪努斯}致\Person{普罗库鲁斯}和\Person{西利努斯}\newline{}\Emph{第220封信}（427年）\Person{奥斯定}致\Person{博尼法斯}\newline{}\Emph{第227封信}（428或429年）\Person{奥斯定}致\Person{阿利皮乌斯}\newline{}\Emph{第228封信}（428或429年）\Person{奥斯定}致\Person{霍诺拉图斯}\newline{}\Emph{第229封信}（429年）\Person{奥斯定}致\Person{达里乌斯}\newline{}\Emph{第231封信}（429年）\Person{奥斯定}致\Person{达里乌斯}\par

% structure: p0011 source=h2 target=section reason=relative-heading-rank; base=h2

\section{第四部分}

\Emph{由于这些书信的写作日期不确定，本笃会编辑将它们归入第四部分，并在此部分下分为两个主要类别：第一类包含一些争议性的书信，第二类则是那些或出于\Person{奥斯定}对个人福祉的关心，或因职分所系而写就的书信。}\par

\Emph{第232封信} \Person{奥斯定}致\Place{马道拉}民众\newline{}\Emph{第245封信} \Person{奥斯定}致\Person{波西迪乌斯}\newline{}\Emph{第246封信} \Person{奥斯定}致\Person{兰帕迪乌斯}\newline{}\Emph{第250封信} \Person{奥斯定}致\Person{奥克西利乌斯}\newline{}\Emph{第254封信} \Person{奥斯定}致\Person{贝内纳图斯}\newline{}\Emph{第263封信} \Person{奥斯定}致\Person{萨皮达}\newline{}\Emph{第269封信} \Person{奥斯定}致\Person{诺比利乌斯}\par

% structure: p0014 source=h2 target=section reason=relative-heading-rank; base=h2

\section{附录}

一封收录于\Person{热罗尼莫}书信集之中的信函，编号即依据该目录。\par

\Emph{\WorkTitle{\Person{热罗尼莫}书信集}第144封} \Person{奥斯定}致\Person{奥普塔特}\par

\SourceRecord{Translated by J.G. Cunningham.；Nicene and Post-Nicene Fathers, First Series；Vol. 1.；Edited by Philip Schaff.；Buffalo, NY: Christian Literature Publishing Co.,；1887.；<http://www.newadvent.org/fathers/1102.htm>.}

% structure: p0001 source=h1 target=section reason=page-title

\section{第一封信（主后386年）}

\Emph{\Person{奥古斯丁}致\Person{赫尔摩根尼阿努斯}问安。}\par

1. 我即使在戏谑的讨论中也不敢攻击\OriginalTerm{学园派}{Academy}的哲学家们；因为，若非我相信他们的观点与通常归于他们的大相径庭，这样杰出人士的权威怎会不能打动我呢？与其驳斥他们——这非我能力所及——我毋宁尽力模仿他们。因为在我看来，这对他们所盛行的时代甚为适宜：凡从柏拉图哲学源泉流出的纯正之物，与其任其流经开阔的草地，在那里无法保持清白纯净以抵御俗众的侵扰，不如引入幽暗多刺的丛林，仅供极少数\Emph{人}清凉之用。我有意使用\Quote{兽群}一词；因为还有什么比认为灵魂是物质的更如禽兽呢？为了防备持此见解的人，在我看来，\OriginalTerm{新学园派}{new Academy}巧妙地设计了这样一门隐藏真理的艺术和方法。但在我们这时代，当我们看不到谁是哲学家——我认为那些只披著哲学家外衣的人不配那尊荣的名号——在我看来（至少对那些因学园派学说的术语精微而不敢试图理解其真义的人来说），应该让他们恢复发现真理的希望，免得那曾适时有效地根除顽固错误的东西，如今竟开始阻碍真理知识的种子的播撒了。\par

2. 在那个时代，各竞争学派的研究被如此热切地从事，以致唯一恐惧的是错误可能被认可。因为每一个被怀疑论哲学家的论证从自以为牢不可破的立场驱赶出来的人，都以与那时代人的更大勤勉和当时流行的坚强信念相应的坚韧和谨慎，着手寻求另一立场——这信念就是：真理虽深奥难解，却确然隐藏在事物和人类心灵的本性之中。然而如今，人们对刻苦努力如此厌恶，对文艺如此冷漠，以致一旦传闻，据最敏锐的哲学家判断，真理不可企及，人们就昏昏欲睡，永蒙心窍。因为他们竟不敢自以为比那些伟人们辨识力更高，能发现\Person{喀涅阿德斯}终其异常漫长的一生，以全部勤勉、才智和闲暇，外加广博多样的学识，都未能发现的东西。而且，如果他们稍奋起克服惰性，去阅读那些仿佛证明了人无法感知真理的书籍，他们就会陷入如此深沉的昏睡，连天国的号角也无法唤醒他们。\par

3. 因此，尽管我以最大的快乐接受您对我那简短论著的坦诚评价，并如此敬重您，以致我依赖您敏锐的判断力不亚于您友谊的真诚，但我恳请您更仔细地注意一点，并就此再次写信给我——即，您是否赞同我在第三卷末尾所提出的意见，其语调或许犹疑多于确定，但陈述，依我想，更可能被发现有裨益，而非被视为不可信。然而，无论那些论著 \WorkTitle{驳学园派} 价值如何，我最欣喜的，不是如您所说的（用的是友好偏爱而非真理的语言）我战胜了学园派，而是我已挣脱并抛弃了那可憎的束缚，它们曾使我因对获得真理——即灵魂的食粮——绝望而被阻于哲学滋养的胸怀之外。\par

\SourceRecord{Translated by J.G. Cunningham.；Nicene and Post-Nicene Fathers, First Series；Vol. 1.；Edited by Philip Schaff.；Buffalo, NY: Christian Literature Publishing Co.,；1887.；<http://www.newadvent.org/fathers/1102001.htm>.}

% structure: p0001 source=h1 target=section reason=page-title

\section{第二封信（公元386年）}

\Emph{致\Person{泽诺比乌斯}，\Person{奥古斯丁}问候。}\par

1. 我想，我们双方都同意，我们身体感官所接触的一切事物，一刻也不能保持不变，相反，它们处于运动之中，不断变迁，没有当下的实在性，也就是说，用拉丁哲学的语言来说，并不存在。因此，真确而神圣的哲学告诫我们，对这些事物的爱恋犹如最危险、最具灾难性的陷阱，必须加以抑制和制伏，好使心灵，即便在使用这个身体时，也能全然专注于并热切关心那些始终如一、其魅力丝毫不靠短暂诱人之处的事物。尽管这完全正确，尽管我的心灵无须借助于感官，便能认识你真实的样子，认识你是可以不受烦扰地去爱的对象，但我仍要承认，当你不在我身边，身体分隔两地时，我会怀念与你相会、相见的快乐，因此，只要这种快乐能够实现，我便热切渴望。我的这个软弱（因为想必就是软弱），如果我对你的了解正确，你是很高兴在我身上看到的；尽管你祝愿你最好、最亲爱的朋友们一切顺利，你却宁愿担心而不是希望他们摆脱这个软弱。然而，如果你的灵魂已达到如此坚定的程度，既能看清这个陷阱，又能嘲笑那些身陷其中的人，那么你确实了不起，与我的境界截然不同。至于我，只要我为别人的离去而惋惜，我也就希望他同样为我的离去而惋惜。同时，我警醒努力，尽量不将自己的爱寄托在任何可能违背我意愿而与我分离的事物上。我将此视为自己的本分，同时也提醒你，无论你此刻心境如何，我们之间已开始的讨论必须完成，只要我们彼此在意。因为我绝不能同意与\Person{阿利比乌斯}一道结束这个讨论，即使他愿意。但他并不愿意；因为他在目前的情况下，绝非这样的人：当我们试图用尽可能多的信件将你留在我们身边时，你却在某种我们不得而知的需要压力下，不愿如此，他便不会与我一同力劝你。\par

\SourceRecord{Translated by J.G. Cunningham.；Nicene and Post-Nicene Fathers, First Series；Vol. 1.；Edited by Philip Schaff.；Buffalo, NY: Christian Literature Publishing Co.,；1887.；<http://www.newadvent.org/fathers/1102002.htm>.}

% structure: p0001 source=h1 target=section reason=page-title

\section{书信三（公元387年）}

\Emph{奥古斯丁向\Person{内布里迪乌斯}致以问候。}\par

我不知道，该将此归因于你那可称之为恭维的言辞之效果，还是事情果真如此，这一点我无法断定。因为这印象来得突然，我尚未确定它究竟在多大程度上值得相信。你好奇这究竟是何事。你怎么看呢？你几乎使我相信——并非相信我幸福，因为幸福乃是唯独智者的产业——而是相信我在某种意义上算得上幸福：正如我们将\Emph{人}这一称谓用于那些仅在与\Person{柏拉图}的理想之人相比较时才算得上人的存在，或者将我们所见之物称为\Emph{圆的}或\Emph{方的}，尽管它们与少数心灵所辨别的完美图形相去甚远。晚餐后，我在灯旁读了你的信；随即我就躺下了，但并未立即入睡；因为我在床上沉思良久，并这样自言自语——\Person{奥古斯丁}自问自答：\Quote{\Person{内布里迪乌斯}断言我幸福，难道这不是真的吗？} \Quote{绝不可能完全真实，因为我自己也承认，我距离智慧尚远。} \Quote{但幸福的人生难道不能也成为那些并非智者之人的命运吗？} \Quote{这几乎不可能；因为若是那样，缺乏智慧便只是小不幸，而非如实际情形那样，是不幸的唯一根源。}那么，\Person{内布里迪乌斯}究竟为何认为我幸福呢？难道是因为读了我这些小书，他便贸然断定我已是智者？诚然，喜悦的冲动不会使他如此轻率，尤其是我深知他的判断极有分量。我现在明白了：他写的是他认为最能让我高兴的话，因为他因我那些论文中的内容而感到愉悦；他是在喜悦的心情中写下这些话的，并未仔细斟酌他欢快的笔端所托付的那些情感。那么，如果他读了我的\WorkTitle{独语录}，他会说什么呢？他会更加欣喜若狂，但却找不到比他已经赐予我的\Quote{幸福}这一称呼更崇高的赞美之词了。于是，他一下子就慷慨地赐予了我可能有的最高称谓，而没有留下一个字，以备将来他因我而更加喜悦时用来增加对我的赞美。看看喜悦会造成什么吧。\par

可是，那真正幸福的生活在哪里？在哪里？啊，在哪里？哦！若是得到了它，人就会鄙视\Person{伊壁鸠鲁}的原子论。哦！若是得到了它，人就会知道，在这下界，除了可见世界便别无他物。哦！若是得到了它，人就会知道，在一个绕轴旋转的球体上，靠近两极点的运动，要比位于两极之间半途之处的点的运动更慢——以及其他我们同样知道的此类事情。然而眼下，我如何、或在何种意义上能被称为幸福呢？我竟不知道这个世界为何具有现在这样的大小，而构成它的那些图形的比例，根本不妨碍它扩大到任何所期望的幅度？或者，难道不会有人对我说——乃至我们不会被迫承认，物质是无限可分的；因此，从任一给定的基点（可以这么说）出发，一定数量的微粒必然达到某个确定而可知的量吗？所以，既然我们不承认有任何微粒小到不能再进一步分割的地步，又有什么能迫使我们承认，任何部分的组合大到不可能再增加呢？或许，我曾在私下对\Person{阿利比乌斯}透露过某个重要的真理，那便是：为理性所认识的数量，可以无限增加，却不能无限减小，因为我们无法将其减至低于单位；而为感官所认识的数量（这当然就是指物质部分或物体的量），则相反地可以无限减小，却有增加上的限度。这或许就是哲学家们公正地宣称，财富在于理性所运用的那些事物中，而贫穷则在于感官所涉及的那些事物中的原因。因为有什么比能够无限地减小更贫穷的呢？又有什么比可以随心所欲地增加、随心所欲地前往、随心所欲地归来，并且所爱的对象宏大而无法使之缩小，更真正地富足呢？因为凡是理解这些数量的人，最爱的莫过于单位；这并不奇怪，因为正是通过单位，他才能爱所有其他的数目。但言归正传：既然世界可能更大或更小，为何它恰恰是现在这样大小呢？我不知道：它的尺寸就是如此，我无法再进一步。还有：为何世界占据它现在的位置，而非另一个位置？在此，最好也不要提问；因为无论答案是什么，其他问题依然会存在。有一件事曾使我大惑不解，那就是物体可以被无限地细分。对此，或许已有一个解答，即提出抽象数目相反的属性[即其可无限倍增的属性]作为对照。\par

3. 但是，且慢：让我们看看那呈现在心灵中的不可定义的对象是什么。这个世界，我们的感官所认识的，无疑是某个由理智把握的世界的影像。现在，我们在镜子反映的影像中观察到一种奇特的现象——无论镜子有多大，它们都不会使影像比置于其前的物体更大，哪怕后者再微小；但在小的镜子中，比如瞳孔，尽管有大平面正对着它们，却只形成极小的影像，与镜子的大小成比例。因此，如果镜子缩小，其中反映的影像也随之缩小；但却不可能通过放大镜子来放大影像。其中确有什么可能值得进一步探究；但与此同时，我必须睡觉了。此外，如果我在 \Person{Nebridius} 看来是幸福的，那不是因为我寻求，而是因为我或许找到了什么。那么，那是什么？是否就是那串推理，我惯于如此珍爱它，仿佛它是我唯一的宝藏，而我或许在其中过于沉溺了？\par

4. \Quote{我们由哪些部分构成？} \Quote{灵魂与身体。} \Quote{其中哪个更高贵？} \Quote{无疑是灵魂。} \Quote{人们在身体中赞扬什么？} \Quote{我所见无非是优美。} \Quote{那么身体的优美是什么？} \Quote{外形上部分的和谐，加上某种悦目的色彩。} \Quote{这种优美在真实时更好，还是在虚幻时更好？} \Quote{毫无疑问，在真实时更好。} \Quote{而它在哪里是真实的？在灵魂中。} \Quote{因此，灵魂比身体更值得爱；但是，这种真实存在于灵魂的哪个部分？} \Quote{在理智和理解力中。} \Quote{理解力必须与什么斗争？} \Quote{与感官。} \Quote{那么我们是否必须竭尽全力抵抗感官？} \Quote{当然。} \Quote{那么，如果感官所认识的事物让我们快乐，该怎么办？} \Quote{我们必须阻止它们这样做。} \Quote{如何阻止？} \Quote{养成没有它们也能生活的习惯，并渴望更好的事物。} \Quote{但如果灵魂死了，那又怎样？} \Quote{那么，真理就死了，或者理智就不是真理，或者理智不是灵魂的一部分，或者某个有不朽部分的东西可能会死：这些结论我很久以前就在我的\WorkTitle{独语录}中证明它们是荒谬且不可能的；我坚定地相信情况就是如此，但不知何故，由于习俗和经历恶事的影响，我们感到恐惧并犹豫。但即便最终承认灵魂会死——我看这也绝无可能——幸福的生活仍不在于那些可感对象所能产生的短暂快乐：对此，我已经深思熟虑并予以证明了。}\par

可能由于我正是基于此类推理，被我的 \Person{Nebridius} 认为，即便不是绝对幸福，至少也算某种意义上的幸福。让我也自我认定为幸福吧：这样做我又会失去什么，或者我为何要吝于对自己的境况抱有良好看法呢？我如此自忖，然后照常祈祷，便入睡了。\par

5. 这些事我思量着写给你是好的。因为当我随意将任何浮现于心头的想法写给你时，你能为此感谢我，这令我愉悦；除了那个我无法令其不悦的人，我还能更乐意地对谁写下荒唐之言呢？但是，如果一个人是否爱另一个人取决于命运，那么，我恳请你注意：当我因命运的恩赐而如此欣喜若狂，并且公然希望我此类美好东西的储备能大量增加时，我又如何能被正当地称作幸福呢？因为那些最为真智慧的人，也是唯一有资格被宣称为幸福的人，都主张命运的恩赐既不应该恐惧，也不应该渴望。\par

在此，我用到了 \Quote{\Emph{cupi}：} 这个词，你能否告诉我，应该是 \Quote{\Emph{cupi}} 还是 \Quote{\Emph{cupiri}}？我很高兴这件事出现了，因为想让你指导我这个动词 \Quote{\Emph{cupio}} 的变位；因为当我将类似的动词与它比较时，我对正确变位的不确定越发增加了。因为 \Quote{\Emph{cupio}} 与 \Quote{\Emph{fugio}}、\Quote{\Emph{sapio}}、\Quote{\Emph{jacio}}、\Quote{\Emph{capio}} 相似；但不定式是 \Quote{\Emph{fugiri}} 还是 \Quote{\Emph{fugi}}，是 \Quote{\Emph{sapiri}} 还是 \Quote{\Emph{sapi}}，我不知道。我原本可以将 \Quote{\Emph{jaci}} 和 \Quote{\Emph{capi}} 视为回答我关于其他动词疑问的平行实例，要不是我担心某个语法家会像玩球一样随意地 \Quote{抓住} 我并 \Quote{抛出} 我，提醒我动名词 \Quote{\Emph{jactum}} 和 \Quote{\Emph{captum}} 的形式与其他动词 \Quote{\Emph{fugitum}}、\Quote{\Emph{cupitum}} 和 \Quote{\Emph{sapitum}} 不同。此外，对于这三个词，我同样不知道倒数第二个音节应该读成长音并带拐折调，还是不带调且读短音。我希望能激起你写一封足够长的信。我恳请你给我一些需要花些时间阅读的东西。因为我阅读你写的文字时的快乐，远超我的表达能力。\par

\SourceRecord{Translated by J.G. Cunningham.；Nicene and Post-Nicene Fathers, First Series；Vol. 1.；Edited by Philip Schaff.；Buffalo, NY: Christian Literature Publishing Co.,；1887.；<http://www.newadvent.org/fathers/1102003.htm>.}

% structure: p0001 source=h1 target=section reason=page-title

\section{信函四（公元387年）}

\Emph{致\Person{内布里狄乌斯}，\Person{奥古斯丁}致意。}\par

1. 令我极为惊讶的是，当我搜索你尚有哪些来信未回复时，发现只有一封信让我成为你的债务人——即你在信中请求我告诉你，在此期间，在我这段你猜想很充裕、并渴望分享的闲暇中，我在学习分辨自然中与感官打交道的事物和那些需运用理智的事物方面，取得了多大进展。但我想你并非不知，若人越来越深地受虚假意见浸染，越充分地、亲密地在这些意见中锻炼自己，那么心灵与真理接触时，就越容易产生相应的效果。然而，我的进步，如同身体的发育一样，是如此渐进，以至难以清晰地界定其步骤，正如同：虽然一个男孩与一个青年之间有极大差异，但若从他幼时起每日询问，谁也不能在任何一天说，现在他不再是男孩，而是青年了。\par

2. 然而，我不愿你如此套用这个比喻，以至于认为，凭借一种更强理解力的旺盛活力，我已可说是步入了灵魂的壮年。因我仍不过是个孩童，尽管也许，如我们所说，是个有希望的孩童，而非一个无用之辈。因为虽然我心灵的眼睛，大都被那些源自可感事物的击打所生的纷扰所困挠和压抑，但它们却被那简短的推理过程所苏醒和提起：\Quote{心灵和理智高于眼睛和通常的视觉能力；若非如此，我们凭理智感知的事物便不会比凭视觉感知的事物更真实。}我恳请你帮我审查，对这一推理是否能提出任何有效的反对。与此同时，我借此得蒙恢复与振奋；当我呼求天主的助佑，开始向他及那些最真实的事物上升时，我有时得以如此把握和享受那永恒常在的事物，以致我有时诧异，竟然需要诸如此类我上述的推理，来说服我那些在我灵魂中确确实实临在于我的事物，正如我临在于我自己一般真实。\par

请你亲自检查你的信函，因为我承认，你为此将比我更为费心，以确保我或许并非无意中仍欠某一封回信：因为简直难以相信我竟如此迅速地卸下了我过去认为如此繁多的债务重担；但同时，我也不能怀疑，你有我的一些信函，至今尚未收到回复。\par

\SourceRecord{Translated by J.G. Cunningham.；Nicene and Post-Nicene Fathers, First Series；Vol. 1.；Edited by Philip Schaff.；Buffalo, NY: Christian Literature Publishing Co.,；1887.；<http://www.newadvent.org/fathers/1102004.htm>.}

% structure: p0001 source=h1 target=section reason=page-title

\section{第五封书信（公元388年）}

\Emph{致\Person{奥斯定}，\Person{内布里迪乌斯}致以问候。}\par

我亲爱的\Person{奥斯定}，你真的把你的精力和耐心花在你在\Place{塔加斯特}的同乡的事务上，而你热切渴望的远离俗务的闲暇仍然被剥夺了吗？我想知道，是谁这样利用你的善良本性，侵占你的时间？我相信他们不知道你最喜爱和最渴望的是什么。你身边难道没有朋友告诉他们你一心所向往的吗？\Person{罗马尼亚努斯}和\Person{卢西尼亚努斯}都做不到吗？至少让他们听我说。我要大声宣告；我要声明，天主是你爱的最高对象，你内心的渴望就是作祂的仆人，紧紧依附着祂。我渴望说服你到我在乡间的家，在这里休息；我不怕被你那些同乡指责为强盗，你对他们爱得太深，而他们也报以过分的爱。\par

\SourceRecord{Translated by J.G. Cunningham.；Nicene and Post-Nicene Fathers, First Series；Vol. 1.；Edited by Philip Schaff.；Buffalo, NY: Christian Literature Publishing Co.,；1887.；<http://www.newadvent.org/fathers/1102005.htm>.}

% structure: p0001 source=h1 target=section reason=page-title

\section{书信第六封（公元389年）}

\Emph{\Person{内布里狄乌斯}向\Person{奥古斯丁}致候。}\par

1. 你的信我极其乐意像保护我的眼睛一样小心珍藏。因为它们是伟大的，不是在篇幅上，而是在它们所讨论的主题的伟大上，以及在展示关于这些主题的真理时表现出的卓越能力上。它们将把\Person{基督}的声音，以及\Person{柏拉图}和\Person{普罗提诺}的教训，带到我耳中。因此，对我而言，它们永远悦耳，因为其优美的文辞；易于阅读，因为其简洁；也因其蕴含的智慧而有益于领会。所以，请不辞辛劳，将凡你判断为神圣或良善的一切都教导我。至于这封信本身，请在你准备好讨论一个关于记忆和想象力所呈现的图像的微妙问题时答复我。我的意见是：虽然可能有独立于记忆的如此图像，但没有任何记忆的运用是独立于如此图像的。你会说：\Quote{那么，当记忆被用来回忆一个领悟或思想的活动时，发生了什么呢？} 我这样回答这个异议：之所以记忆能够回忆起这些行为，是因为在所假设的领悟或思想活动中，我们产生了一个受空间或时间限制的某物，其性质如此，以至于能够被想象力复现：因为我们要么将词语的使用与领悟和思想的运用联结起来，而词语是受时间限制的，因此落入感官或\OriginalTerm{想象力}{imaginative faculty}的领域；要么如果我们没有将词语与思维活动联结，我们的理智无论如何都在思维活动中经历了某种能够通过想象力之助而被记忆所回忆的东西。我照常有些未经深思、略显混乱地陈述了这些。请你细查，并弃绝谬误，写信告诉我你对此问题的真知灼见。\par

2. 请也听听这个问题：\Quote{我很想知道，为什么我们不肯定，\OriginalTerm{想象力}{phantasy [imaginative faculty]}的所有图像都源于其自身，而不是说它从感官接收这些图像？} 因为可能，正如灵魂的理智官能归功于感官，不是为了理智所施用的对象，而是为了激发它看到这些对象的提醒，同样，想象力也可能归功于感官，不是为了作为其施用对象的图像，而是为了激发它观想这些图像的提醒。也许正是通过这种方式，我们才能解释这个事实：想象力觉察到一些感官从未觉察到的对象，由此表明它自身内并源于自身拥有其所有图像。请你也回答我你对这个问题的看法。\par

\SourceRecord{Translated by J.G. Cunningham.；Nicene and Post-Nicene Fathers, First Series；Vol. 1.；Edited by Philip Schaff.；Buffalo, NY: Christian Literature Publishing Co.,；1887.；<http://www.newadvent.org/fathers/1102006.htm>.}

% structure: p0001 source=h1 target=section reason=page-title

\section{书信七（公元389年）}

\Emph{致\Person{内布里狄乌斯}，\Person{奥古斯丁}致以问候。}\par

% structure: p0003 source=h2 target=subsection reason=relative-heading-rank; base=h2

\subsection{第一章 记忆可以不依赖想象所呈现的形像而运作。}

1. 我将免去正式的前言，直接着手讨论你长久以来想听我意见的那个主题；我这样做更是出于乐意，因为阐述需要花费一些时间。\par

在你看来，没有形像——即没有想象所提供的一些对象的把握——便不可能有记忆的活动；你把这种形像乐于称之为\Quote{\OriginalTerm{幻象}{phantasiæ}}。至于我，我持不同的意见。首先，我们必须注意到，我们所记起的事物并不总是那些消逝之物，而多半是那些恒久的事物。因此，既然记忆的功能在于持守属于过去时间的东西，那么它一方面包含那些离我们而去的事物，另一方面包含那些我们离之而去的事物。例如，当我记起我的父亲时，回忆的对象是已离我而去且已不复存在的事物；而当我想起\Place{迦太基}时，对象则是依然存在、只是我已离开的事物。然而在这两种情形下，记忆所保存的，都属于过去的时间。因为我记起这个人或这座城市，不是凭着眼见现在，而是凭着过去曾见过他们。\par

2. 你在这点上或许要问，为何要提出这些事实？而且你可以更轻易地提出这个问题，因为你观察到，在我所引的两个例子中，所记起的事物若不把握你断言总是必需的某种形像，便无论如何也不能进入我的记忆。就我当前的目的而言，我暂时只需以此方式证明，记忆也可被说成是包含那些尚未消逝的事物：现在请留神注意这一点如何支持我的主张。有些人针对\Person{苏格拉底}所创的最著名理论提出了毫无根据的反对，按照该理论，我们所学的事物并非作为新的东西被引入我们的心灵，而是通过\OriginalTerm{回忆}{recollection}的过程被带回记忆；他们用以支持其反对的是断言记忆只与过去之物有关，然而正如\Person{柏拉图}本人所教导的，我们通过运用理智所学到的那些事物是恒久的，因其不朽，所以不能被算在过去之物之列：他们陷入的错误明显来自于此，即他们没有考虑到，我们之所以能辨识这些事物，仅仅是那属于过去的领会这一心智活动；并且正是由于我们在心智活动之流中将这些事物撇在后头，并开始以种种方式关注其他事物，我们才需要借着回忆的努力——即借着记忆——返回于它们。因此，如果我们撇开其他例子，将思想专注于永恒本身——它是永远恒久之物——并且考虑到，一方面，它无需任何由想象所制造的形像作为媒介以引入心灵；另一方面，它若不通过我们记忆它，便绝无可能进入心灵——那么我们就会看到，至少就某些事物而言，记忆可以运作而无须想象呈现出任何有关所记之事的形像。\par

% structure: p0007 source=h2 target=subsection reason=relative-heading-rank; base=h2

\subsection{第二章 心灵若未曾藉由感官获悉外界事物，便缺乏想象所呈现的形像。}

3. 其次，至于你认为心灵有可能不依赖感官的帮助而为自己形成物质事物的形像这一意见，这可由如下的论据加以反驳：——如果心灵在将身体用作感知物质对象的工具之前，就能为自己形成这些对象的形像；并且，如果——正如所有心智正常之人不会怀疑的那样——心灵在卷入感官所产生的错觉之前，获得过更可靠、更正确的印象，那么我们就必须认为，睡梦中的人比清醒的人、以及精神错乱的人比未患此等心智疾病的人，其印象具有更高的价值：因为在这些心智状态下，心灵正是受到同样类型的形像——即在他们尚未凭借这些最具欺骗性的信使即感官获取信息之前所受到的形像——的影响；如此一来，他们在睡梦或疯狂中所见的太阳，就必定比心智健全且清醒之人所见的太阳更为真实，或者虚幻之物必定比真实之物更好。然而我亲爱的\Person{内布里狄乌斯}，如果这些结论正如其明显所是那样全然荒谬，那么这就证明了，你所提到的那种形像无非是感官所施的一击，这些感官对于此类形像的作用并非如你所写的那样，仅仅是促成心灵在自身内形成其的提示或告诫，而是确实将这些错觉——我们通过感官所受到的错觉——带入心灵，或者更确切地说，将其印入心灵。你感到的困难，即为何我们能在思想中构想出我们从未见过的面孔和形状，这正证明了你心智的敏锐。因此我将把本信写得超出通常的篇幅；然而不会超出你所赞同的长度，因为我相信，我写给你的信越是丰满，你便越是欢迎。\par

4. 我注意到，你和许多人所称的\OriginalTerm{幻象}{phantasiæ}，根据其来源是感官、想象或理性能力，可以极其方便而精确地划分为三类。第一类的例子是，心灵在其自身内形成并向我们呈现你的面容、或\Place{迦太基}、或我们已故的朋友\Person{维雷昆杜斯}、或任何现存或曾存且我亲自见过并感知之物的形象。第二类包括我们想象曾存在或如何如何的一切事物：例如，当我们在论述中为了说明而设定并不存在但无害于真理的事物时；或者当我们阅读历史、聆听、创作或拒绝相信虚构叙事时，在心中唤起对所描述事物的生动构想。这样，我随自己的想象，在脑海中描绘出\Person{埃涅阿斯}、或驾著有翼龙群的\Person{美狄亚}、或\Person{克雷梅斯}、或\Person{巴门诺}的模样。此类还包括那些被当作真实而提出的事物，无论是智者将某些真理包裹于这类虚构之中，还是愚人建立了各种类的迷信；例如，受刑者的\Place{弗莱格桑河}、黑暗之族的五个洞穴、支撑苍天的北极、以及诗人与异端者的其他无数怪异之事。此外，在讨论中我们常说：\Quote{假设我们所居住的世界，有三个这样的世界层层相叠；}或\Quote{假设地为四方形所包围，}等等：因为这一切，我们都根据心绪和思想的走向，自行想象描绘。至于第三类形象，主要与数目和度量有关；这些部分见于事物的本性中，例如当整个世界的形象被发现，思考者心中便随之形成相应的形象；部分见于科学，如几何图形、音乐和声及无穷的数字变化：尽管我认为这些作为领悟的对象本身是真实的，但它们仍是幻象练习的肇因，其误导倾向甚至连理性本身也难以抗拒；尽管保持推理科学不受此弊也非易事，因为在逻辑划分和推论中，我们为自己形成所谓\OriginalTerm{小石子}{calculi}或筹码，以便利推理的进行。\par

5. 在这整个形象的密林中，我相信你不认为第一类形象在通过感官进入心灵之前便属于心灵。对此我们无需再辩。至于其他两类，或许可以合理地提出疑问，但若非显而易见，即心灵在尚未受感官及可感事物之欺惑影响时，较少倾向于幻觉，然而谁敢怀疑这些形象远比感官所告知我们的更为虚假？因为我们所假设、相信或想象的事物，在每一点上全然不实；而我们所眼见及其他感官感知的事物，如你所见，远比这些想象的产物更接近真实。至于第三类，无论我凭借此类形象在心灵中构想物体的何种空间广延，尽管看似如思维过程以不容错谬的科学推理产生了此形象，我却证明它具欺瞒性，而这些推理反过来又足以揭露其虚假。因此，我完全无法相信[而若接受你的意见，我便必须相信]灵魂尚未运用肉身的感官，尚未藉这些易谬的工具骤然遭遇那死变易逝之物，却已如此可耻地屈服于错觉。\par

% structure: p0011 source=h2 target=subsection reason=relative-heading-rank; base=h2

\subsection{第三章 答复异议}

6. \Quote{那么，我们从未见过的事物，在思想中构想它们的能力从何而来？}你想，除了心灵与生俱来的某种增减能力，还能有何原因？这一能力，无论心灵转向何处，都不得不随身携带着（这一能力尤其可从数目上观察到）。通过运用此能力，例如，若将眼睛极为熟悉的一只乌鸦的形象，置于心灵之眼前，可以说，通过减去某些特征并增加另一些特征，它几乎可以被变成任何眼睛从未见过的形象。也是凭借此能力，当人心惯于如此思量时，这类形象便仿佛不请自来地涌入思绪。因此，心灵能够如前所述，从感官引入其认知的对象中减去一些事物，并加上一些事物，在想象的运用中产生出整体上从未被任何感官观察到的东西；但其各个部分全都曾在观察之中，尽管分布于种种不同的事物里：例如，当我们还是孩童时，若生于内陆地区并在此长大，我们即使只在一只小杯中见过水，亦能对大海形成某种观念；但草莓与樱桃的滋味，在我们于\Place{意大利}尝到这些果实之前，是绝不会进入我们的思想的。因此，生来即盲者，当被问及光明与色彩时，便不知如何作答。因为他们从未通过感官感知过有色彩的对象，因此不可能在心中拥有此类对象的形象。\par

7. 请勿以为奇怪，尽管心灵寓于并交织于那些按其本性可被我们构形或描绘的一切形象之中，这些形象并非由心灵从自身内部演化而出，而是在它通过感官从外部接收它们之前。因为我们同样发现，伴随着愤怒、喜悦与其他此类情绪，我们在思想官能尚未意识到我们具有产生这类形象（或我们感受的表征）的能力之前，便已在身体相貌与面容上产生变化。这些变化是随着那些奇妙方式中的情绪体验而来的（尤其值得你悉心思考），这些方式在于灵魂中\OriginalTerm{隐藏的数字}{hidden numbers}的反复作用与反作用，而没有任何虚幻物质事物的形象介入。因此，我希望你理解——既然你已察觉心灵有如此多的活动完全独立于所说的那些形象——在心灵可能借以获致物体知识的一切活动中，任何其他方式都比未经辅助的思想创造可感事物形式的过程更为可能，因为我不认为心灵在使用身体与感官之前能够有任何这类构想。\par

因此，我深爱且最可爱的弟兄，凭那联结我们的友谊，以及我们对天主的神圣法律本身的信仰，我劝诫你切勿与那黑暗国度的阴影结为友伴，并即刻断绝你与他们之间可能已开始的任何交谊。抵抗肉体感官的支配，本是我们最神圣的职责，但若我们以宠爱和谄媚对待感官加于我们的打击与创伤，这抵抗便全然废弃了。\par

\SourceRecord{Translated by J.G. Cunningham.；Nicene and Post-Nicene Fathers, First Series；Vol. 1.；Edited by Philip Schaff.；Buffalo, NY: Christian Literature Publishing Co.,；1887.；<http://www.newadvent.org/fathers/1102007.htm>.}

% structure: p0001 source=h1 target=section reason=page-title

\section{书信八（公元389年）}

\Emph{\Person{内布里提乌斯}向\Person{奥古斯丁}致问候。}\par

1. 我急于进入书信主题，故省略前言或引言。当更高（我指的是天上的）力量乐意在睡眠中通过梦向我们启示某事时，这是如何发生的，亲爱的\Person{奥古斯丁}，他们使用的方法是什么？我说，那是什么方法，\Emph{即}，是通过什么技艺或魔法，通过什么中介或法术，他们实现了这一点？他们是否通过自己的思想影响我们的心灵，从而使我们的思想中也出现同样的图像？他们是否将我们梦到的事物呈现在我们面前，仿佛是在他们自己的身体或想象中实际完成的？但如果他们确实在自己的身体中实际做了这些事，那么，为了让我们看到他们所做的事，我们就必须具有别的肉眼，能在睡眠时看到内在发生的事情。然而，如果他们不是借助身体产生所述的效果，而是在自己的想象力中构造这些事物，从而影响我们的想象，由此给予我们梦的内容以可见的形式；那么，我请问，为什么我不能迫使你的想象力再现那些我首先在自己想象中形成的梦呢？毫无疑问，我具有想象力，它能在我心中呈现任何我喜欢的画面；然而，我并未因此在你身上引起任何梦，尽管我看到，我们的身体甚至有能力在我们自己身上引发梦。因为身体通过将它与灵魂联结的共鸣之纽带，以奇异的方式迫使我们通过想象力重复或再现任何它曾经经历的事物。因此，在睡眠中，如果我们口渴，我们常常梦见喝水；如果饥饿，我们就觉得自己在吃东西；还有许多别的例子，可以说，通过某种交换的方式，事物经由想象力从身体传递到灵魂。\par

请不要惊讶于这些问题在此向你陈述时缺乏优雅与精妙；请考虑这主题所涉及的晦涩以及笔者缺乏经验；请你尽己所能弥补他的不足吧。\par

\SourceRecord{Translated by J.G. Cunningham.；Nicene and Post-Nicene Fathers, First Series；Vol. 1.；Edited by Philip Schaff.；Buffalo, NY: Christian Literature Publishing Co.,；1887.；<http://www.newadvent.org/fathers/1102008.htm>.}

% structure: p0001 source=h1 target=section reason=page-title

\section{第九封信（公元389年）}

\Emph{致\Person{内布里迪乌斯}，\Person{奥古斯丁}致以问候。}\par

1. 虽然你了解我的心思，但或许你不知道我多么渴望与你相处。然而，天主有一天会将这巨大的福分赐给我。我读了你的信，信中的话真挚坦诚，你在信中抱怨自己独处，仿佛被朋友们抛弃，而在他们的陪伴中你曾找到生命中最甜美的魅力。但我能向你建议的，无非是我确信你已经在练习的事：与你的灵魂交谈，并尽你所能将它提升至天主。因为在祂内，你以更坚固的纽带也拥有我们，不是借着身体的形象——在此期间我们不得不满足于用这些形象来彼此记念——而是借着那思想的能力，藉此我们认识到我们彼此分离的事实。\par

2. 在我考虑你的信时——回复所有这些信，我确实不得不回答不少困难和重要的问题——我被其中一封信深深震惊，你在信中问我，更高的力量或超人的媒介是通过何种方式将某些思想和梦境放入我们心中的。这问题很大，并且，正如你自己的审慎必能令你信服的，要满意地回答它，需要的不是一封信，而是一场充分的口头讨论或一整篇论文。然而，由于我了解你的才智，我将尝试抛出一些思想的萌芽，或许能为这问题带来启发，以便你或可凭借自己的努力完成对这主题的详尽探讨，或至少不至于对能以满意结果来研究这重要问题的可能性感到绝望。\par

3. 我的意见是，心灵的每一运动都在某种程度上影响身体。我们知道，即使在我们迟钝而呆滞的感官中，当心灵的运动有些剧烈时，例如当我们愤怒、悲伤或喜乐时，这一点也是显而易见的。由此我们可以推测，同样地，当思想忙碌时，尽管心灵活动的身体效果不为我们所察觉，但对于那些具有\OriginalTerm{大气或以太本质}{aërial or etherial essence}的存在——其感知能力极为敏锐，敏锐到相比之下，我们的官能几乎不配被称为感知——可能会有某种可辨别的效果。因此，可以说，心灵在身体上留下的这些它运动的足迹，或许不仅存留，而且带着某种习惯的力量存留下来；并且，当这些足迹被秘密地搅动和拨弄时，它们根据搅动或触动它们者的喜好，将思想和梦境带入我们心中：而这以惊人的容易性完成。因为，显而易见，我们属地而迟钝的身体在操练领域中的成就，例如演奏乐器、走钢丝等，几乎是不可思议的，那么，假设那些以大气或以太身体的能力作用于我们身体、并因其本性构成而能不受阻碍地穿过这些身体的存在，能够以远为敏捷的速度移动它们所欲的任何东西，就绝非不合理；而我们，虽不察觉它们的所作所为，却仍受它们活动结果的影响。我们有一个有些相似的例子在于我们并不察觉胆汁过盛如何驱使我们去更频繁地爆发出激烈情绪；然而它确实产生这种效果，而这胆汁过盛本身又是我们屈服于那种激烈情绪的结果。\par

4. 然而，如果你在接受我如此仓促地提出的这个例子作为一个平行事例时犹豫不决，那就尽你所能地在你的思想中充分思量它。心灵，如果它不断受到某种在实现和完成其所欲之事上的困难的阻碍，便会因此不断地变得愤怒。因为愤怒，就我对其本质所能判断的，在我看来似乎是一种要除去那些限制我们行动自由之物的混乱的热切。因此，我们通常不仅对人生气，而且也对例如我们写字的笔发怒，在激情中将它弄皱或折断；赌徒对他的骰子、画家对他的画笔，以及每个人对他可能正在使用的工具，如果他认为在某种程度上被它所阻挠，也是如此。现在，医生们自己告诉我们，通过这些频繁的愤怒发作，胆汁会增加。然而，另一方面，当胆汁增加时，我们便轻易地、几乎不需要任何刺激就会生气。因此，心灵通过其运动对身体产生的效果，反过来又能再次感动心灵。\par

5. 这些事可以谈论得很长，通过对相关事实的大量归纳，我们对这主题的认知可以被带向更大的确定性和完整性。但请将我之前寄给你的、关于图像和记忆的那封信连同此信一起阅读，并更仔细地研究它；因为从你的回复中，我清楚地看到它没有被完全理解。当你将我在那封信中谈到的一种自然官能的论述，加入你面前的这些话中——藉此官能，心灵在思想中随意地对任何对象进行增减——你就会看出，我们有可能在梦境和清醒的思想中构想出我们从未见过的身体形式的图像。\par

\SourceRecord{Translated by J.G. Cunningham.；Nicene and Post-Nicene Fathers, First Series；Vol. 1.；Edited by Philip Schaff.；Buffalo, NY: Christian Literature Publishing Co.,；1887.；<http://www.newadvent.org/fathers/1102009.htm>.}

% structure: p0001 source=h1 target=section reason=page-title

\section{第十封信（主后389年）}

\Emph{致\Person{内布瑞狄乌斯}，\Person{奥斯定}致以问候。}\par

1. 你提出的问题，从未像我在你上一封信中读到的那段话那样，让我在思考时如此心神不宁；在那段话中，你责备我对安排我们共同生活所需的计划漠不关心。这是一个严重的指责，若非毫无根据，将是极其危险的。但既然有充分的理由似乎证明，我们在这里比在\Place{迦太基}，甚至比在乡间更能如愿地生活，我完全不知所措，我亲爱的\Person{内布瑞狄乌斯}，不知该拿你怎么办。是否该从我们这里派遣一种最适合你健康状况的交通工具到你那里去？我们的朋友\Person{路西尼安}告诉我，你可以乘坐轿子而无伤身体。但另一方面，我考虑到你的母亲，她在你健康时都忍受不了你的离去，在你生病时，就更无法忍受了。那么，我是否该亲自到你那里去？这是我做不到的，因为这里有些人不能陪我前去，而我认为离开他们对我是种罪过。因为你已经能够在依靠自己心灵资源时愉快地度过时光；但对他们而言，当前努力的目标正是要达到这一点。我是否该频繁地往返，时而与你在一起，时而与他们在一起？但这既不是共同生活，也不是按照我们的愿望生活。因为旅程并不短，反而相当漫长，频繁尝试会阻碍我们获得所渴望的闲暇。此外，还加上身体的软弱，正如你所知，这使得我无法完成所想望的事，除非我彻底放弃那些超出我力量范围的愿望。\par

2. 让一个人的思虑终其一生都萦绕于那些无法平静而轻松地完成的旅途，这不是一个思虑专注于那被称为死亡的最终旅途的人所应做的；而正如你所理解的，唯有这最终旅途才真正值得认真思考。天主确实赐予了少数受命治理教会的人这样的能力：他们不仅能平静地等待，甚至热切渴望那最终的旅途，同时又能毫无烦乱地应付其他那些旅途的辛劳；但我不相信，无论是那些因渴望世俗尊荣而被催迫接受这类职务的人，还是那些虽身处平信徒地位却贪求忙碌生活的人，能得到如此大的恩赐，以至于在喧嚣混乱的集会、四处奔波中，能获得我们所寻求的那种对死亡的熟稔：因为这两类人本来都有能力在退隐中寻求建树。若非如此，我虽不敢说是众人中最愚昧的，但至少是最怠惰的，因为我发现，若没有这样一段从忧虑劳苦中解脱出来的间歇，就根本无法品尝并享受那唯一真实的善。请相信我，人需要极大程度地从流逝事物的扰攘中抽身退隐，才能在内心形成那一种能力——不是由于麻木，不是由于傲慢，不是由于虚荣，不是由于迷信的盲目——而能说出：\Quote{我无所畏惧。} 藉此途径，也能获致那种持久的喜乐，任何他处可寻的愉悦兴奋都无法与之相比。\par

3. 但若这样的生活并非人所能得，那么为何我们偶尔会体验到心灵的平静？为何这体验会更频繁，恰与一个人在灵魂深处敬拜天主的虔诚程度成正比？为何当人从那圣所走出、履行生活职责时，这宁静多半仍然常存？为何有些时刻，我们在言谈中不惧怕死亡，在静默中甚至渴望死亡？我对你说——因为我不愿对每一个人说——对你这样一位我深知曾造访上界的人，我要问：你时常感受到，当灵魂对一切纯属肉体的情感死去时，它是何等甜蜜地活着；那么，你是否会否认，人的整个生命最终有可能变得如此无所畏惧，以至于他可被公正地称为智者？或者，你敢断言这种理性所依凭的心灵状态，曾是你的际遇，除非当你闭门自省、与己心交谈之时？既然事已至此，你看，剩下的唯有你与我分担筹划如何安排我们共同生活这一辛劳了。关于你的母亲，你知道得远比我知道得该如何处理，当然，你的兄弟\Person{维克托}不会丢下她独自一人。我不再多写了，以免转移你思考这一提议的心神。\par

\SourceRecord{Translated by J.G. Cunningham.；Nicene and Post-Nicene Fathers, First Series；Vol. 1.；Edited by Philip Schaff.；Buffalo, NY: Christian Literature Publishing Co.,；1887.；<http://www.newadvent.org/fathers/1102010.htm>.}

% structure: p0001 source=h1 target=section reason=page-title

\section{信函十一（公元389年）}

\Emph{\Person{奥古斯丁}向\Person{内布里迪乌斯}致意。}\par

1. 当你以近乎友好的责备长期向我提出的关于我们如何共同生活的问题，严重困扰我心，我已决意写信给你，请求你专就此事答复，不因任何其他关乎我们学业的主题而用笔，好让我们之间的讨论得以结束的时候，最近收到的你信中所述那极为简短而又无可辩驳的结论立时使我摆脱了所有进一步的挂虑；你的说法大意是，关于此事不应再有犹豫，因为一旦我能到你那里，或你能到我这里，我们都会同样感到不得不把握机会。我心如此，正如我所说，既已安歇，我便翻阅了你所有的信，看看还有哪些未予答复。在其中我发现了如此多的问题，即使它们容易解答，单是其数量便足以消耗任何人的时间和才能。然而它们如此困难，以至于若要求我回答其中哪怕一个问题，我也会毫不犹豫地承认自己不堪重负。这段引言的目的是要你暂停提出新问题，直到我无债一身轻；并且在你的回信中只限于陈述你对我答复的看法。同时，我知道，即便只是稍许推迟分享你的神圣思想，也是我自身的损失。\par

2. 因此，请听我所持关于\OriginalTerm{降生成人}{Incarnation}奥迹的见解，我们受教的宗教为成就我们的救恩而将其托付给我们的信德与知识；我选择优先讨论这个问题，尽管它并非最容易回答的。因为你所提出的关于这个世界的问题，在我看来，与获得幸福生活没有足够直接的关系；并且，在探讨时无论享有什么乐趣，都有理由担心它们会占用本应用于更好事情的时间。至于我这次所承担的主题，首先，令我惊讶的是，你曾为为什么不是圣父，而是圣子，被称为降生成人而困惑，却没有同样为关于圣神的同一问题而困惑。因为\OriginalTerm{位格}{Person}在\OriginalTerm{天主圣三}{Trinity}内的合一，在天主教信仰中被提出和相信，并被少数圣洁而有福的人所理解，是如此不可分离，以致天主圣三所做的一切都必须被视为由圣父、圣子和圣神共同完成；圣父所做的，没有不是圣子和圣神也做的；圣神所做的，没有不是圣父和圣子也做的；圣子所做的，没有不是圣父和圣神也做的。由此似乎得出，整个天主圣三都摄取了人性；因为如果是圣子这样做了，但圣父和圣神没有，那么就有一些事是他们分开做的。那么，为什么在我们的奥迹和神圣标记中，\OriginalTerm{降生成人}{Incarnation}只归于圣子呢？这是一个非常重大的问题，如此困难，主题又如此广阔，以致既不可能给出十分清晰的陈述，也不可能用令人满意的证据来支持它。然而，既然我是写信给你，我冒昧地指点而不是解释我的看法，以便你——凭借你的才能和我们的亲密关系，你彻底了解我——能为自己填补这一轮廓。\par

3. \Person{内布里迪乌斯}，没有任何\OriginalTerm{性体}{nature}——事实上，没有任何\OriginalTerm{实体}{substance}——不在自身中包含并展现这三者：第一，它\Emph{存在}；第二，它是\Emph{此或彼}；第三，它尽其可能地\Emph{持存}于其本然状态。这三者中的第一者呈现了万物由以存在的性体的原始因；第二者呈现了万物据以被特定地塑造和形成的型式；第三者呈现了可以说是一种恒存性，万物在此恒存性中持续存在。现在，如果有可能一物能够\Emph{存在}，却不是\Emph{此或彼}，也不\Emph{持存}于其自身的类形式中；或者一物能够是\Emph{此或彼}，却不\Emph{存在}，也不尽其可能地\Emph{持存}于其自身的类形式中；或者一物能够根据其所属的力量\Emph{持存}于其自身的类形式中，却不\Emph{存在}，也不是\Emph{此或彼}——那么，在那一圣三中，一个位格就有可能做某事而其他位格没有参与。但如果你明白，凡是存在者必定立刻是\Emph{此或彼}，并且必定尽其可能地\Emph{持存}于其自身的类形式中，那么你也就明白，这三位不做任何没有所有位格都参与的事情。我明白，至今我只处理了这一问题中使其解答困难的一个部分。但是，我希望向你简要地开启——如果我真地在这方面成功了的话——在天主教真理的体系中，天主圣三位格不可分离性的教义是多么伟大，又是多么难以理解。\par

4. 现在请听，那使你心神不安的事物，如何不再使你不安。存在的方式，即\OriginalTerm{种相}{Species}，前述三者中的第二个，归之于圣子，乃与培育有关，若可就此类事物使用该词，则亦与某种技艺有关，并与理智的操练有关，心灵借此在对事物的思想中得以塑造。因此，既然通过取了人性，所完成的工作是有效地向我们展示某种正道生活的培育，并在某些语句的威严与明晰之下，成为所命令之事的典范，那么将这一切归之于圣子，并非没有理由。因为在我留给你自己的反思和审慎去思考的许多事上，虽然组成要素繁多，但总有某一个突出于其余之上，因此它仿佛并非没有道理地将整体据为己有：\Emph{例如}，在上述三类问题中，虽然提出的问题是事物是否存在，这必然也涉及它\Emph{是什么}（是这个还是那个），因为如果它不是某物，它当然根本无法\Emph{存在}，以及它应被赞同还是反对，因为凡是\Emph{存在}的事物，都是对其\Emph{性质}作出某种意见的合适对象；同样，当提出的问题是事物\Emph{是什么}时，这必然也涉及它\Emph{存在}，并且其性质可依据某种标准来检验；同样，当提出的问题是事物的\Emph{性质}是什么时，这必然也涉及该事物\Emph{存在}，并且是\Emph{某物}，因为万物与自身不可分割地相连——然而，上述每种情形中的问题，并非从所有三者而得名，而是从探询者注意力所指向的特殊方面得名。现在，为人类所需要的是某种培育，藉此他们能依据某种模型而受教导和塑造。然而，关于这种培育在人内所成就者，我们既不能说它不存在，也不能说它不是可欲求的事物【\Emph{i.e.} 除非同时肯定它的\Emph{存在}和\Emph{性质}，否则我们不能说它\Emph{是什么}】；但我们首先寻求知道它\Emph{是什么}，因为知道了这个，我们便知道由此可以推断它是某物，并且我们可以安居于其中。因此，首要的必需之事，是将培育的某种规范和典范清楚显示出来；而这正是由天主所定的降生成人的方式完成的，这方式特归于圣子，好使由此而产生：我们透过圣子认识圣父本身，即认识那万物所由以存在的独一第一原理，以及一种内在而不可言喻的魅力与甘饴，安居于那知识中，并轻看一切必朽事物——此一恩赐与工作，特归于圣神。因此，虽然天主三位在一切事上完全共同行动，且不可能分开，然而由于我们这些从合一堕入繁杂的软弱，他们的运作必须以一种彼此可区分的方式显示出来。因为若不适度俯就他人的水平，没有人能成功将他人提升到他自己所站的高度。\par

在此你有一封书信，它或许确实不能结束你对这端道理的不安，但它可以使你的思想在一个坚固的基础上运作；这样，藉着我很清楚你所拥有的才能，你可以继续探索，并藉着我们尤其必须坚持的虔敬，领悟那尚待发现的。\par

\SourceRecord{Translated by J.G. Cunningham.；Nicene and Post-Nicene Fathers, First Series；Vol. 1.；Edited by Philip Schaff.；Buffalo, NY: Christian Literature Publishing Co.,；1887.；<http://www.newadvent.org/fathers/1102011.htm>.}

% structure: p0001 source=h1 target=section reason=page-title

\section{第13封信（公元389年）}

\Emph{\Person{奥古斯丁}致\Person{内布利丢斯}问安。}\par

1. 我对于写信谈论我们惯常讨论的题目并不感到愉快；我也没有自由写作新的题目。我看到前者不适合你，而后者我没有闲暇。因为，自从我离开你，既没有机会也没有闲暇去拾起和反复思考我们惯常一同探讨的事情。冬夜确实太长，我也并非完全以睡眠度过；但当我闲暇时，其他主题（而非我们惯常讨论的那些）会作为更优先的考虑出现。那么我该怎么办呢？难道我要像一个哑巴，不能说话，或像一个沉默者，不愿说话吗？这都不是你我所愿的。那么，来吧，请忍受这一夜终了时，我在用来构思这封信的时间里所勉强写出的东西。\par

2. 你不可能不记得，我们之间常常激起讨论、并使我们激动、屏息、兴奋的一个问题，是关于一种永远属于灵魂的身体或某种身体，正如你记得的，这东西被一些人称为它的\OriginalTerm{载具}{vehicle}。显然，这东西如果从一处移动到另一处，是不能被理智所认识的。但任何不能被理智认识的东西，是无法被理解的。然而，对于一个超出理智范围的事物，如果它属于感官范围，形成接近真理的意见也并非完全不可能。但当一件事物既超出理智范围又超出感官范围时，由此引发的猜测便过于空洞和琐碎；而我们现在讨论的这件事，如果它确实存在的话，正是如此。那么，我请问，我们为什么不最终放下这个不重要的问题，并藉着对天主的祈祷，将我们自己提升到至高存在者的至高宁静之中呢？\par

3. 也许你可能在此回答说：\Quote{虽然身体不能由理智认识，但我们能用理智认识关于物质对象的许多事；\Emph{例如}，我们知道物质存在。因为有谁会否认这一点，或断言在这件事上我们只能涉及或然性而非真理呢？因此，尽管物质本身属于或然事物之列，但自然界中存在某种类似它的东西，这是最无可辩驳的真理。物质本身因此被宣布为可感知的对象；但断言其存在，则被宣布为理智可以认识的真理，因为除此以外无法认识。那么，我们所探求的这个未知的身体，灵魂依赖于它才能从一处移动到另一处，也许能被比我们所拥有的更强大的感官所认识，尽管不能被我们的感官认识；但无论如何，它是否存在这个问题，是可以通过我们的理智来解决的。}\par

4. 如果你打算这样说，那么让我提醒你，我们称之为理解的心理活动是以两种方式进行的：要么是通过心智和理性在其自身之内，就像我们理解理智本身存在时那样；要么是由于感官的提示，如同在上述情况中，当我们理解物质存在时那样。在这两种方式的第一种里，我们是通过自身理解，即，通过就我们内在的事向天主求问指导；但在第二种中，我们是就身体和感官所给予我们的暗示，向天主求问指导而理解。如果这些道理是真的，那么没有人能凭他的理智发现你所说的那个身体是否存在，除非有某个人的感官曾给他某种暗示。如果存在某种生物，其感官获得了这种暗示，而我们显然不在其中，那么我认为，我方才开始陈述的结论便确立起来了，即关于灵魂\OriginalTerm{载具}{vehicle}的问题，与我们无关。我希望你反复思考此事，并务必让我知道你的思考结果。\par

\SourceRecord{Translated by J.G. Cunningham.；Nicene and Post-Nicene Fathers, First Series；Vol. 1.；Edited by Philip Schaff.；Buffalo, NY: Christian Literature Publishing Co.,；1887.；<http://www.newadvent.org/fathers/1102013.htm>.}

% structure: p0001 source=h1 target=section reason=page-title

\section{第十四封信（公元389年）}

\Emph{\Person{奥斯定}致\Person{内布里迪乌斯}问候。}\par

1. 我宁愿回复你上一封信，并非我看轻你先前的问题，或是对它们兴味索然，而是因为在回答你时，我承担的任务比你所想的更为重大。虽然你吩咐我寄一封极长的信，但我并没有你想象中那么多闲暇，而且你知道我一直向往，至今依然向往那样的闲暇。不要问我为何如此：因为我更容易列举那些阻碍我的事，却难以解释我为何被它们阻碍。\par

2. 你问为什么你和我虽是独立的个体，却能做许多相同的事，而太阳却不能像其他天体那样做相同的事。我必须试着解释其原因。然而，若你我能做相同的事，太阳也做许多其他天体所做的事：若它在某些事上不同于其他天体，这在你我身上也同样属实。我行走，你也行走；它移动，它们也移动：我清醒，你也清醒；它发光，它们也发光：我讨论，你也讨论；它运转，它们也运转。但心灵之事与可见之物之间，却没有可比较的恰当之处。不过，若是合乎理性地以心比心，那么天体若有心灵，其思想或静观，或任何更便于表达其内在活动的用词，必定比人的更均匀一致。再者，关于身体的运动，你若以惯常的专注反思，便会发现两个个体绝不可能做出完全同一的动作。当我们一同行走时，你以为我们必然做同一动作吗？愿如此想法远离你这般智慧之人！因为向北而行的一人，若要迈出与另一人相同的步幅，必会赶在前面，或走得比另一人更慢。这两者都无法为感官所察觉；但假如我未看错，你所关注的，乃是理智所知，而非感官所获。然而，即使我们自北极向南移动，彼此尽可能紧紧相依，踏在平滑如大理石或象牙般光洁的平面上，我们的动作、脉搏跳动、身形面容，绝无可能达致完美的同一。且将我们搁置一旁，换上\Person{格劳科斯}的儿子们，你仍不会因此替换而有所得：因为即便在这对彼此极为肖似的双生子身上，各自动作必然为其所独有的必然性，恰如他们作为独立个体出生的必然性一样大。\par

3. 或许你会说：\Quote{这其中的差异，只有理性才能发现；但太阳与其他天体之间的差异，感官也显而易见。} 倘若你定要我着眼于它们大小的差异，你当知道，就它们距我们的远近有多少可说的话，你所说的显而易见之事会因此被拖入何等的不确定之中。不过，我可以承认天体的实际大小与其表观大小相符，因为我自己也相信这一点；那么，请你指给我看，有谁的感官无法觉察\Person{奈维乌斯}那惊人的身量，他比最高大的人还高出一尺。顺便说一句，我想你之前太过急切地想找到一个能与他匹敌的人；而当你搜寻未果，便决意让我把信拉长，好与他的身量一较高下。因此，若在地上就能见到如此大小的差异，我想我们在天上见到类似情形也不必惊奇。然而，如果令你惊异的是，除了太阳，再没有其他天体的光能充满白昼，那么我要问你，有谁曾向世人显现，比天主使之与自己结合的那人更为伟大呢？天主以全然不同于对待世上一切圣贤的方式，将那人摄取与自己合一。若将他与其他智慧之人比较，他远超他们的优越，远大于太阳超越其他天体的程度。我恳请你务必悉心揣摩这个比较；因为我仅凭这粗略的点拨，或许便已为你那颖悟非凡的心智，暗示了你曾向我提出的有关基督人性的问题的解答。\par

4. 你还问我，那至高的真理与至高的智慧，万物的\OriginalTerm{形式（或原型）}{Form (or Archetype)}，即万物藉祂而受造，我们的信经宣认祂为天主的独生子的那一位，是否在祂内含有一般人类的\OriginalTerm{理念}{idea}，还是也包括我们人类中每一个体的理念。这是一个重大的问题。我的意见是，在创造人时，在祂内只有一般人的理念，而没有你或我作为个体的理念；但在时间的流转中，每一个体的理念，连同人与人之间一切差别的特殊性，都存在于那\OriginalTerm{纯粹真理}{pure Truth}之中。我承认这一点非常隐晦；但我不知能用何种比喻来阐明它，除非我们求助于那些完全存在于我们心灵之内的学问。在几何学中，角的理念是一回事，正方形的理念是另一回事。因此，每逢我欲描写一个角，我心中呈现的只是角的理念，且仅限于此；但我若非要同时关注四个角的理念，就断然无法描写一个正方形。同样，每一个人，作为个体的人，都是依照一个独属于他自己的理念而被造的；但在创造一个民族时，尽管用以创造的也是一个理念，却不是一个人的理念，而是众多人集合的理念。因此，如果\Person{内布里迪乌斯}——正如他那样——是这个宇宙的一部分，而整个宇宙由部分构成，那么创造宇宙的天主，其计划中就不能不包含所有部分的理念。所以，既然这理念涉及极多的人，它本身并不属于人本身；虽然，从另一方面看，所有个体却以奇妙的方式复归于一体。但这点你可自行从容考量。同时，我恳请你满足于我所写的这些，尽管我已经超越了\Person{奈维乌斯}本人。\par

\SourceRecord{Translated by J.G. Cunningham.；Nicene and Post-Nicene Fathers, First Series；Vol. 1.；Edited by Philip Schaff.；Buffalo, NY: Christian Literature Publishing Co.,；1887.；<http://www.newadvent.org/fathers/1102014.htm>.}

% structure: p0001 source=h1 target=section reason=page-title

\section{第十五封信（公元390年）}

\Emph{\Person{奥古斯丁}向\Person{罗马尼亚努斯}致意。}\par

1. 这封信表明纸张短缺，但并非证明这里羊皮纸充裕。我上次寄给你叔叔的信中用了我的象牙书板。你会更容易谅解这张羊皮纸残片，因为写给他的信不能拖延，而我认为因缺乏更好的材料就不给你写信是极其荒谬的。但如果你那里有我任何书板，我请求你送过来以备这样的情形。关于大公信仰，我蒙主俯允，写了一些东西，想在我到来之前寄给你，如果在此期间我的纸还够用的话。因为凡是从与我在一起的弟兄们的书房出来的作品，你都会宽仁地接纳。至于你提到的那些手稿，我完全忘了，除了\WorkTitle{论演说家}那几卷；但我当时信中的意思无非是，你可随意取用，我至今仍持此意；因相距遥远，此事我也别无他法。\par

2. 上封来信中，你愿我分享你家中的喜乐，这令我无比欣慰；但是\par

难道你愿我忘记，这眼下如此平静的深海，\newline{}会何其迅速地变脸，\newline{}并唤醒这沉睡的波涛？\par

然而我知道，你并非要我忘记这一点，你自己也未曾疏忽。因此，如果天主赐予你闲暇以便更深入的默想，就善用这天赐的恩福吧。因为当这些事临到我们时，我们不仅应自庆，还应向赐予者表示感激；因为如果在管理现世的福分时，我们行事公正而仁慈，并以符合这些短暂财物的无常本性所需的节制与清醒的精神去对待它们——如果我们拥有它们却不被它们占有，使之增加却不被其缠累，使用它们却不被束缚，这样的行为配得永福的赏报。因那身为真理的曾说过：\BibleQuote{如果你们在别人的财物上不忠信，谁还把属于你们的交给你们呢？}所以，让我们摆脱对短暂世事之挂虑；让我们寻求那不朽而确实的福分；让我们飞越世俗的财富之上。蜜蜂即使采集了丰盛的花蜜，也丝毫不能缺少翅膀；因为若沉溺于蜜中，它便会死去。\par

\SourceRecord{Translated by J.G. Cunningham.；Nicene and Post-Nicene Fathers, First Series；Vol. 1.；Edited by Philip Schaff.；Buffalo, NY: Christian Literature Publishing Co.,；1887.；<http://www.newadvent.org/fathers/1102015.htm>.}

% structure: p0001 source=h1 target=section reason=page-title

\section{书信 16（公元390年）}

\Emph{来自\Person{马道拉的马克西穆斯}致\Person{奥古斯丁}。}\par

1. 我渴望时常因你的来信而欢欣，也因你最近以极为愉快且不伤和气的方式对我提出的辩驳之激励而欢欣；我不禁以同样的精神回复你，免得你将我的沉默视为认错的表示。但若你认为这些话显出了老年的衰弱，我恳请你宽容慈祥地聆听。\par

希腊神话告诉我们——虽然没有充分的凭据让我们相信这种说法——\Place{奥林波斯山}是诸神的居所。但我们实际看到，我们城镇的广场上挤满了仁慈的神祇；我们对此表示赞同。谁会如此疯狂和愚昧，竟至否认有一位无始无终、无自然所生的至高天主，祂可说是万有的伟大而全能之父？这位天主的大能，弥漫于祂所造的宇宙中，我们以许多名号加以敬拜，因为我们都不知道祂的真名，\OriginalTerm{天主}{God}这个名称通用于各种宗教信仰。由此，当我们在不同的祈求中，仿佛个别地亲近这神性的某些部分时，实际上，我们被视为敬拜那位所有部分合而为一的天主。\par

2. 在另一件事上，你的迷惑如此之大，使我无法掩饰对此感到的不耐烦。因为谁能忍受将\Person{米格多}尊崇于掷雷电的\Person{朱庇特}之上；或将\Person{萨奈}尊崇于\Person{朱诺}、\Person{密涅瓦}、\Person{维纳斯}和\Person{维斯塔}之上；或将那位\OriginalTerm{首席殉道者}{arch-martyr} \Person{南法尼奥}（哦，真可怕！）尊崇于所有不朽神祇之上？在不朽者中，\Person{卢西塔斯}也同样受到不亚于此的宗教敬畏，还有其他数不清的人（其名字为神人所共憎），当他们遭受到了其品行与行为所该当的耻辱结局时，他们竟仍伪称是在为良善事业而高贵地死去，以此为其罪恶生涯加上最后的一笔，却明知自己因所犯的丑事而受判刑。这些人的坟墓（这愚行几乎不值一提）被成群愚昧的人所参拜，他们离弃我们的神庙，蔑视祖先的纪念，以致那位忿怒的吟游诗人的预言显著地应验了：\Quote{\Place{罗马}将在诸神的殿中，指着人的阴魂发誓。} 在我看来，此时仿佛第二次\Place{阿克提姆}战役已经开始，其中那些注定即将灭亡的埃及怪物，竟敢向\Place{罗马}诸神挥舞武器。\par

3. 但是，大智之人啊，我恳求你，暂且放下并丢弃那使你在各处享有盛名的辩才之锋；也请放下你惯于在辩论中使用的\Person{克吕西波}的那些论证；暂时搁置你的逻辑，那逻辑在施其能时，旨在不留任何确定之事予人；请明明切切地指示我，那位你基督徒们声称特别属于你们、并佯称在隐密处与你们同在的天主，究竟是谁？因为我们是在光天化日之下，当着众人的耳目，以虔诚的祈求敬拜我们的神祇，并用悦纳的祭祀取悦他们；我们还竭力使这些事为人所见、为人所认可。\par

4. 然而，我既已衰老年迈，便不再继续这场争辩，并乐意赞同那位\Place{曼图亚}的修辞学家的见解：\Quote{各人被自己所最喜爱的事物所吸引。}\par

此后，卓越的人啊，你已背离了我的信仰，我毫不怀疑这封信将被某个窃贼偷去，并被火或其他方式销毁。若此事发生，纸张会丢失，但我的信不会，我将永远保留一份副本，可供所有虔敬之人取阅。愿你蒙诸神护佑，我们所有生活在地面上的凡夫俗子，正是借着他们，虽表面纷争却真实和谐地，敬畏并敬拜那位诸神与凡人的共同之父。\par

\SourceRecord{Translated by J.G. Cunningham.；Nicene and Post-Nicene Fathers, First Series；Vol. 1.；Edited by Philip Schaff.；Buffalo, NY: Christian Literature Publishing Co.,；1887.；<http://www.newadvent.org/fathers/1102016.htm>.}

% structure: p0001 source=h1 target=section reason=page-title

\section{书信十七（主后390年）}

\Emph{致\Place{马道拉}的\Person{马克西穆斯}。}\par

1. 我们是在进行严肃的辩论，还是您只想彼此戏谑一番？因为，从您来信的语气，我弄不明白，是由于您所持理由薄弱，还是出于您的礼貌，使您在辩论中更愿表现得机智而非有分量。因为，首先，您将\Place{奥林匹斯山}与您的市场相比较，这用意我揣测不透，除非是为了提醒我：在那座山上，\OriginalTerm{朱庇特}{Jupiter}与父亲交战时曾扎营，正如历史所教导的——那些历史，你们的宗教信徒称之为神圣；而在那个市场里，\OriginalTerm{马尔斯}{Mars}以两种形象呈现，一为无武装，一为武装，而在这两尊像对面设置的一尊人像，用三根伸出的手指遏止它们魔性的狂怒，使其不能加害市民。难道我能相信，您提及那个市场，是为了让我想起这些神祇，除非您希望我们以戏谑而非认真的态度继续讨论？然而，关于您所说这些神祇，可以算是那位伟大天主的肢体，我无论如何要劝戒您，既然您持有这种意见，就要非常小心地避免此类亵渎的玩笑。因为倘若您说的是那一位天主，关于祂，无论有学问还是无学问的人，正如古人所言，都意见一致，您怎么断言这些被人像遏止残暴狂怒——或者，若您喜欢，说成权势——的神祇是祂的肢体呢？关于这点，我还可以多说，而您自己的判断或许会向您显明，反驳您观点的门在此是多么敞开。但我克制自己，免得让您觉得我行事更像修辞学家，而非认真捍卫真理的人。\par

2. 至于您收集了一些迦太基死者的名字，以为可用以嘲笑我们的宗教，在您看来想必是机智之举，我不知道我该回答这一嘲讽，还是不予理睬。因为，若以您的明智，这些事如同实际那样微不足道，我无暇顾及此类玩笑。然而，如果这些在您看来很重要，我奇怪您竟没想到，您这位容易为怪异发音的名字所困扰的人，你们宗教信徒的司铎中有名叫\OriginalTerm{欧卡狄雷斯}{Eucaddires}者，其神祇中有名叫\OriginalTerm{阿巴狄雷斯}{Abaddires}者。我想您写信时并非不知道这些名字，而是出于您的礼貌和诙谐，想让我们放松精神，回忆起你们迷信中那些滑稽可笑的事。因为，考虑到您是\Place{阿非利加}人，而我们俩都定居在\Place{阿非利加}，您给阿非利加人写信时不可能忘了，以致认为布匿名字可作指责的话题。因为若我们解释这些词的含义，\OriginalTerm{南法尼奥}{Namphanio}除了\Quote{好脚的人}，即其到来带来好运的人，还能有什么别的意思呢？正如我们惯常说某人到来后便有顺利之事，就说他是踩着幸运的脚来的。而如果您拒绝布匿语，您实际上就否认了大多数学者所承认的事：许多智慧被保存在用布匿语写成的书卷中，免于遗忘。不仅如此，您甚至应为诞生在这片土地上感到羞愧，因为这片语言之摇篮尚有余温，即这语言原本是，而且直到最近仍是人民的语言。然而，如果仅仅因名字的读音而生气是不合理的，且您承认我正确给出了上面那个名字的含义，那么您就有理由对您的朋友\Person{维吉尔}表示不满，因为他向您的大力神\OriginalTerm{赫丘利}{Hercules}发出了参加\Person{埃万德}为敬奉他而举行的神圣仪式的邀请，用的正是这样的词句：\Quote{请以吉祥的脚步光临我们，光临这场为您而设的仪式。} 他祝愿他以\Quote{吉祥的脚步}而来；也就是说，他祝愿\OriginalTerm{赫丘利}{Hercules}作为一位\OriginalTerm{南法尼奥}{Namphanio}而来，而您却乐于以这个名字取笑我们。但如果您偏爱嘲笑，您自己那里就有丰富的笑料——\OriginalTerm{斯忒耳库提乌斯}{Stercutius}神、\OriginalTerm{克洛阿西娜}{Cloacina}女神、秃头\OriginalTerm{维纳斯}{Venus}、\OriginalTerm{恐惧之神}{Fear}、\OriginalTerm{苍白之神}{Pallor}、\OriginalTerm{热病女神}{Fever}，以及其他无数同类神祇，古罗马的偶像崇拜者曾为他们建庙宇，并认为理应崇拜他们；您若忽视这些，便是忽视了罗马的神祇，由此显明您并不通晓罗马的神圣礼仪；然而您却鄙视并嘲弄布匿名字，仿佛您是罗马神祇祭坛前的虔诚崇拜者。\par

3. 然而实际上，我相信您对这些神圣礼仪的看重也许并不比我们更多，只不过在穿越此世的旅途中从中获得某种莫名的乐趣：因为您毫不迟疑地躲在\Person{维吉尔}的羽翼下，并用他的一句诗为自己辩护：\par

\Quote{人人各为其所悦者牵引。}\par

如果，那么，\Person{维吉尔}的权威令您喜欢，如您所表明的那样，您也会喜欢这样的诗句：\Quote{首先，\OriginalTerm{萨图恩}{Saturn}从巍峨的\Place{奥林匹斯山}降临，躲避\OriginalTerm{朱庇特}{Jupiter}的武力，一个被剥夺了王国的流亡者，}——以及其他类似的说法，他借此力图使人理解，\OriginalTerm{萨图恩}{Saturn}与你们其他的类似神祇都是凡人。因为他读过许多由古老权威证实的史书，\Person{西塞罗}也读过，并在他的对话录中更明确地阐述了同样的事，其明确的措辞我们不敢坚持，而他则尽力将这知识带给当时的人们，凡他活着的时代所容许的程度。\par

4. 关于你的说法，即你的宗教仪式应被置于我们之上，因为你们公开敬拜神祇，而我们则使用更隐蔽的聚会场所，我首先要问你，你怎么能忘了你的\Person{巴克斯}呢？你认为只向少数入会者展示\Person{巴克斯}才是正确的。然而，你以为提到你们神圣仪式的公开庆祝，只是为了确保我们会把你们的官员和城中显要人物在街上醉醺醺地狂奔的景象摆在我们眼前；在这庄严的仪式中，如果你们被一位神祇附身，你必然看出那剥夺人理性的神祇是何等性质。然而，如果这种疯狂只是伪装的，那么对于你们吹嘘为公开的仪式中隐藏事物这点，你有何话说？如此卑鄙的欺骗又能达到什么良好目的？再者，如果你们拥有\OriginalTerm{先知之恩}{prophetic gift}，为何不在你们的歌声中预言未来？如果你们神智清醒，又为何要抢夺旁观者？\par

5. 那么，既然你通过书信让我们忆起了这些事情及其他我更愿暂且略过的事，我们为何不能取笑你的神祇呢？凡是了解你心思并读过你书信的人都清楚，你自己已大量地取笑它们。因此，如果你想让我们以一种符合你的年龄与智慧的方式，并且实际上正如我们最亲爱的朋友可以合情合理地要求我们的那样，围绕我们的宗旨讨论这些话题，就请提出值得你我辩论的主题；并在为你的神祇辩护时，留意说出一些能阻止我们认为你故意出卖自己事业的话来，因为我们发现你非但没有提出任何辩护之词，反而让我们忆起了那些可被用来反对它们的事情。然而，最后，为让你不致不知道这一点，并因此无意中陷入亵渎的戏谑，我可以向你保证，\OriginalTerm{天主教基督徒}{Christian Catholics}（他们在你的城市里也建立了一座教会）绝不敬拜任何已故之人，简言之，绝不将任何天主所造所生的东西当作\OriginalTerm{神能}{divine power}来敬拜。他们只敬拜那创造并生成万物的天主本身。\par

这些事，在唯一真天主的助佑下，一旦我得知你愿意认真讨论它们，便会更充分地加以探讨。\par

\SourceRecord{Translated by J.G. Cunningham.；Nicene and Post-Nicene Fathers, First Series；Vol. 1.；Edited by Philip Schaff.；Buffalo, NY: Christian Literature Publishing Co.,；1887.；<http://www.newadvent.org/fathers/1102017.htm>.}

% structure: p0001 source=h1 target=section reason=page-title

\section{第十八封信（公元390年）}

\Emph{致\Person{切莱斯蒂努斯}，\Person{奥斯定}致以问候。}\par

1. 啊，我多么希望能不断地对你说一件事！就是：让我们卸下无用挂虑的重担，只背负那些有益的。因为我不知今生是否可望完全免于挂虑。我曾写信给你，但未收到回复。我将我已能完成的、经修订的反对\OriginalTerm{摩尼教徒}{Manichæans}的著作尽可能多地寄给你，但至今未闻你对它们的判断和感受。现在正是我请求你归还、你也应当归还它们的时候。所以我恳请你不要拖延，连同你本人的一封信一并寄来，我急切地想知道你对它们作何处理，或你认为还需要哪些进一步的帮助，才能成功地攻击那个谬误。\par

2. 因为我深知你，我请求你接受并深思以下关于一个重大主题的简短句子。有一种本性，在时间和空间上都可改变，即有形体的。另一种本性，在空间上全然不可改变，只在时间上可改变，即精神性的。还有第三种本性，无论在空间上还是在时间上都不可改变：那就是天主。我所说在某些方面可改变的那些本性被称为受造物；那不可改变的本性则被称为造物主。然而，既然我们肯定任何事物存在，只在于它继续存在且是一（因此，在一切形式中，统一是美的必要条件），在此本性分类中，你必能分辨出，哪一种以最高方式存在；哪一种居于最低位，却仍在存在范围内；哪一种居于中间位，高于最低者，却不及最高者。那最高者即是本质的真福；最低者则既不能享福也不能受苦；中间的本性，当其趋向下者时，生活中充满痛苦；当其转向最高者时，则生活在真福中。凡信基督的人，不将自己的情感沉溺于最低者中，不在中间者中骄傲自足，如此他便有资格与最高者结合共融；而这便是我们受命、受劝勉、并为圣洁热情所驱使而向往的积极生活的总纲。\par

\SourceRecord{Translated by J.G. Cunningham.；Nicene and Post-Nicene Fathers, First Series；Vol. 1.；Edited by Philip Schaff.；Buffalo, NY: Christian Literature Publishing Co.,；1887.；<http://www.newadvent.org/fathers/1102018.htm>.}

% structure: p0001 source=h1 target=section reason=page-title

\section{书信十九（公元390年）}

\Emph{致\Person{盖尤斯}，\Person{奥斯定}谨致问候。}\par

言语无法表达我与你分别后心中因回忆你而充满的喜悦，且此后也时常充满我心。因为我记得，尽管你在寻求真理时充满惊人的热情，但你从未逾越合宜谦和的辩论界限。我很难找到任何人比你更乐于提问，且同时更耐心聆听回答。因此，我很乐意花很多时间与你交谈；因为这样的时光，不论多久，都不会显得漫长。然而，讨论使我们难以享受如此交谈的障碍又有何用呢？只说极难便够了。也许将来某个时候这会变得很容易；愿天主恩赐如此！同时，情况并非如此。我已托付带此信给你的那位弟兄，让他将我所有的著作呈交给你的杰出智慧和爱德，以便你能阅读。因为我写的任何东西都不会在你那里遇到不情愿的读者；因为我知道你对我所怀的善意。不过，我愿说明，如果你读了这些作品后，赞同它们，并认为它们是真实的，你绝不可视之为我的东西，除非视之为赐予我的；并且你可以转向那同一源泉，那使你能够领悟其真理的能力也出于此。因为没有人能仅仅从手稿中的任何事物，或从作者那里，辨识他所读内容的真理，而是靠着他自己内心的某种东西，若是真理之光，以超越常人常分的清明照耀，且远避肉身的昏暗影响，穿透了他自己的心灵。然而，如果你发现了一些虚假且应被摒弃的东西，我愿你知道，这些东西是如露珠自人性软弱的迷雾中落下，你应将其视为真正出于我的。我本欲劝勉你坚持寻求真理，若非我似乎已看见你心灵之口大大张开，准备啜饮真理。我本也想劝勉你以刚毅的坚韧持守你所领受的真理，若非你已以最清晰的方式表明你拥有心灵的力量和坚定的志向。因为你内里所有的一切，在我们相处的短暂时间里，已向我显明，几乎如同肉身的帷幕被撕裂一般。而确实，我们天主的仁慈眷顾，绝不允许像你这样善良且天赋卓绝的人，成为基督羊群的局外人。\par

\SourceRecord{Translated by J.G. Cunningham.；Nicene and Post-Nicene Fathers, First Series；Vol. 1.；Edited by Philip Schaff.；Buffalo, NY: Christian Literature Publishing Co.,；1887.；<http://www.newadvent.org/fathers/1102019.htm>.}

% structure: p0001 source=h1 target=section reason=page-title

\section{书信二十（公元390年）}

\Emph{致\Person{安多尼努}：\Person{奥思定}谨致问候。}\par

1. 由于您应当收到我们两人致您的信，因此当您亲自见到我们两人中的一人时，我们的一部分债务即以非常丰富的利息偿还了；而由于他的声音，您可以说几乎听到了我的声音，因此，若非那位在我看来已使写信变得不必要的人迫切要求，我或许可以不写了。所以我现在与您的交谈，甚至比亲自与您在一起更满意，因为您既读我的信，又听您知道我心里住着的那一位的话。我怀着极大的喜悦研读并思考了您圣德所寄来的信函，因为它既展现了您的基督徒精神未受这邪恶世纪的诡诈玷污，又展现了您对我充满友善的心。\par

2. 我祝贺您，并感谢我们的天主和主，因为在您内有望德、信德和爱德；我在祂内感谢您，因为您如此看重我，竟相信我是天主的忠仆，并以坦诚的心爱着您在我身上所称赞的；尽管在这件事上，承认您的好意，与其说是感谢的机会，不如说是祝贺的机会。因为您为了那善自身的缘故而爱那善，这为您自己是有益的；凡是因相信他人是善的而爱他人的人，当然也就爱了那善，无论那个人是否真的如他所想的那样。在这件事上只须小心避免一个错误，就是我们不要对真实的善的看法有违真理，不是对个人，而是对人身上真正的善。但，我深爱的弟兄，既然您在相信或是知道为人大善就是欢欣喜乐地、纯洁地事奉天主这一点上，没有丝毫错误，那么当您因相信某人有分于这善而爱他时，即使那人并非如您所想的，您仍能获得赏报。因此，您是有理由因此受到祝贺的；而您所爱的那人之所以应受祝贺，不是因为他因此被爱，而是因为（如果真是这样的话）他确实是那位为此缘故爱他的人所认为的那种人。所以，至于我们真实的品格，以及我们在神圣生命上可能取得的进展，唯有那位对何为人之善及对每个人的品格皆能无误判断的那一位才看得见。就此事而论，为使您获得真福的赏报，您全心接纳我，相信我是天主的忠仆，一如我所应是的，这便足够了。但我仍要为此多多感谢您，因为您称赞我，好像我已达到那样的卓绝，这极大地激励我追求如此卓绝。若您不仅要求分享我的祈祷，而且不停止祈祷，您还将有更多的感谢。因为为弟兄代祷，当作为爱的祭献献上时，便更蒙天主悦纳。\par

3. 我极其亲切地问候您的小儿子，并祈求他能成长起来，走上顺从天主律法的救恩要求的道路。此外，我渴望并祈求，唯一真实的信仰和敬拜，即唯一至公的信仰和敬拜，能在您家中兴旺并增长；如果您认为为促进此事，需要我做任何劳苦，那么请毫不迟疑地要求我服务，仰赖我们共同的主，并仰赖我们必须服从的爱德律法。我尤其要向您虔诚的明智推荐这一点：通过阅读天主的圣言，并与您的伴侣进行严肃的交谈，您或者在她心中栽植，或者培育对天主的明智敬畏。因为任何关心灵魂福祉，并因此毫无偏见决意寻求主旨意的人，几乎不可能在享受一位良善指导者的引导时，辨不清每个形式的裂教与唯一的天主教会之间存在的差别。\par

\SourceRecord{Translated by J.G. Cunningham.；Nicene and Post-Nicene Fathers, First Series；Vol. 1.；Edited by Philip Schaff.；Buffalo, NY: Christian Literature Publishing Co.,；1887.；<http://www.newadvent.org/fathers/1102020.htm>.}

% structure: p0001 source=h1 target=section reason=page-title

\section{信函 21（公元 391 年）}

\Emph{致我最可称颂、最可敬的主教\Person{瓦莱里乌斯}，在主的鉴临下以真诚的爱最热切地珍视的我的父亲，司铎\Person{奥古斯丁}在主内致意。}\par

1. 首先我请求你虔诚的智慧考虑，一方面，如果主教、司铎或执事的职责以敷衍和趋炎附势的方式履行，那么今生没有比这更安逸、更惬意、更能赢得人青睐的工作了，尤其是在我们这个时代，但同时在天主眼中也没有比这更可悲、可叹、应受谴责的了；另一方面，如果遵守我们救恩统帅的命令，那么今生没有比这更艰难、更劳苦、更危险的工作了，尤其是在我们这个时代，但同时在天主眼中也没有比这更受祝福的了。然而，我无论在童年还是成年初期都未学会如何恰当地履行这些职责；而当我开始学习时，我被迫，作为对我罪的公正纠正（因为我不知道还能怎么想），在舵上接受次位，那时我还不知道如何划桨。\par

2. 但我认为，我主以此责备我，乃因我自以为拥有更高的知识和卓越，竟在未经验水手们工作的性质之前，就去指责他们的诸多过失。因此，当我被派到他们中间分担他们的劳苦时，我才开始感受到我指责的轻率；尽管在那之前我已判定这一职务充满许多危险。因此，我的一些弟兄在我\OriginalTerm{晋秩}{ordination}时看见我在城中流泪，他们出于最好的意愿尽力安慰我，但他们的话因为不了解我悲伤的原因，完全未能触及我的处境。然而，我的经验使我更深刻地，无论在程度上还是分量上，比仅仅思考时更清楚地意识到这些：并非我现在看到了任何新的波浪或风暴，是我先前通过观察、传闻、阅读或默想所不知道的；而是我从前不知道自己的技巧或力量以避免或对付它们，并曾估计这技巧和力量有些价值而非毫无价值。然而，主嘲笑我，并乐意通过实际经验向我显明我是什么。\par

3. 但如果祂这样做不是出于审判，而是出于怜悯，正如我至今仍满怀信心地希望，当我已认识自己的软弱时，我的责任便是勤奋学习圣经中为像我这样的情况所包含的一切疗法，并通过祈祷和阅读，致力于使我的灵魂获得履行如此重大职责所必需的健康与活力。这一点我尚未做到，因为我没有时间；因为我正是在想着与其他人一起摆脱一切其他事务，以熟悉圣经，并打算安排好确保为这项伟大工作有不受打扰的闲暇时晋秩的。此外，确实，我之前从未知道自己对于如今使我心神不安、压垮我的艰难工作是多么不胜任。但如果我通过经验了解到一个以神圣的圣事和圣言服事人民的人所需的是什么，却发现自己被阻止不能获得我已知道我所缺少的东西，难道你愿意我丧亡吗，\Person{瓦莱里乌斯}神父？你的爱德何在？你真的爱我吗？你真的爱教会，即你指派这样不胜任的我去服事的那教会吗？我深信你二者都爱；但你认为我胜任，而我更认识自己；然而，若不是我通过经验学到，我本不会认识自己。\par

4. 也许你的圣善回答说：\Quote{我想知道什么缺乏使你能适合你的职务。} 我所缺乏的是如此之多，以至于我更容易列举我现有的，而非我所渴望拥有的。我敢说，我知道并毫无保留地相信关乎我们救恩的教义。但我的难题在于，我如何能用这真理去为他人救恩服务，寻求不仅对我自己有益，而是为许多人有益，使他们得救。或许，不，毫无疑问，在圣经中记载着忠告，通过认识和接受这些忠告，天主的人可以这样在教会中对天主的事务履行他的职责，或至少在不敬虔的人中间常存无亏的良心，无论生死，以确保那唯独谦卑和温顺的基督徒心灵所叹息的生命不致丧失。但如何能做到，除非正如主亲自告诉我们的，通过祈求、寻找、叩门\BibleRef{玛 7:7}，即通过祈祷、阅读和哭泣？为此，我已通过弟兄们提出请求，并在此恳请中再次提出，即恳请你真诚而可敬的爱德，赐我短短的时间，比如直到复活节。\par

5. 我要怎样回答主，我的审判者呢？我能否说：\Quote{我未能获得我所需要的东西，因为我全神贯注于教会事务}？如果祂这样回答呢：\Quote{你这恶仆，倘若属于教会的财产（为收取其果实，已耗费了大量劳苦）正遭受某个压迫者的侵犯，而你又有能力在世俗法官的法庭上为她的权利辩护，那么，你岂不会撇下我以我的血所浇灌的田地，在众人同意、甚至在一些人的严令下前去申诉？倘若判决对教会不利，你岂不会渡海提起上诉；即便离家一年或更久也不会听到召你回家的怨言，因为你的目的是阻止他人夺走并非为灵魂、而是为穷人身体所需的土地——其实，若我的活树得到精心栽培，穷人的饥饿本可以更容易、更令我悦纳的方式得以满足。那么，你为何声称没有时间学习如何耕耘我的田地呢？}请告诉我，我恳求你，我能回答什么？或许你愿意我这样说：\Quote{都怪年迈的\Person{瓦莱里乌斯}；因为他以为我已通晓一切必要之事，出于对我的爱，他执意不准我去学习我尚未掌握的东西}？\par

6. 年迈的\Person{瓦莱里乌斯}，请思考这一切；我恳求你，因着基督的良善与严厉、祂的仁慈与审判，因着那位使你如此爱我，以致我甚至在灵魂利益攸关时也不敢令你不悦的天主，我恳求你思考这一切。此外，你曾呼求天主和基督为我作证，证明你的清白与爱德，以及你对我真诚的爱，仿佛这一切都不是我甘愿为此宣誓的事。因此，我恳求你藉此所宣示的爱与亲情。请怜悯我，为了我所请求的目的，赐予我我所请求的时间；并以你的祈祷帮助我，使我的渴望不致落空，使我的离开不致对基督的教会毫无裨益，也不致对我的弟兄与同仆毫无益处。我知道，主不会轻视你为我代祷的爱，尤其在这类事上；祂将悦纳它，如同馨香的祭献，也许会在比我恳求的更短期限内，以祂所写下的圣经中有益的教导，使我完全装备好从事祂的事工，并将我归还给你。\par

\SourceRecord{Translated by J.G. Cunningham.；Nicene and Post-Nicene Fathers, First Series；Vol. 1.；Edited by Philip Schaff.；Buffalo, NY: Christian Literature Publishing Co.,；1887.；<http://www.newadvent.org/fathers/1102021.htm>.}

% structure: p0001 source=h1 target=section reason=page-title

\section{第二十二封信（主历392年）}

\Emph{致\Person{奥勒留}主教：司铎\Person{奥斯定}谨致问候。}\par

% structure: p0003 source=h2 target=subsection reason=relative-heading-rank; base=h2

\subsection{第一章}

1. 经过长期的犹豫，我不知如何构思一封合适的回信给圣善的您（因为所有试图表达我情感的尝试，都被那自发涌起的强烈友爱之情所挫败，而诵读您的来信，更使这情感猛烈地炽燃起来），我最终俯伏在天主面前，求祂按照我的能力，在我内做工，使我能够写给您这样的回复，既能符合我们共同拥有的、对主的热忱和对祂教会的关切，又能与您的尊严以及我理当对您表示的敬意相称。首先，关于您相信我因我的祈祷而得到助佑，我非但不否认这种确信，反而欣然接纳。因为这样，即使并非透过我的祈祷，也必定透过您的祈祷，我主定会俯听我。至于您以最仁慈的心意，赞许\Person{亚利比乌}弟兄决意继续与我们团聚，成为那些渴望摆脱世俗挂虑的弟兄们的楷模，我心中涌起的感激，远非言语所能表达。愿主以此报偿您的灵魂！因此，在我身边这个已经开始一起成长的全体弟兄团体，都对您的这份大恩深怀感激；您施予这恩惠时，虽然与我们只是在地面上相隔遥远，在心神上却亲近无比，为我们谋取了利益。为此，我们尽己所能地献上祈祷，求主乐意随同您一起扶持托付给您的羊群，永不在任何地方离弃您，而在一切患难时刻，成为您的助佑，并在您履行司铎职分时，向着祂的教会，彰显出灵性之人含泪呻吟所恳求天主施行的慈悯。\par

2. 因此，您知道，至福的尊长，因您爱德的无比丰盛而受人崇敬，我绝不失望，反而怀着活泼的希望，深信借着您所掌握的权威——我们信赖那权威是交托于您的心灵，而非仅仅交托于您的血肉——我们的主天主能借着人们对大公会议以及对您本人的尊敬，治愈非洲教会在许多人的行为中所遭受的、并在少数肢体的悲痛中所哀叹的诸多属肉体的污点与混乱。因为宗徒曾在一处简略列出当受憎恶和远避的三类恶习，由这些恶习生出不可胜数的邪恶行径，而教会唯独严惩其中一类——即他所列的第二类；至于另两类，即第一类和第三类，在世人眼中似乎尚可容忍，由此逐渐演变，甚至不再被视为恶习。蒙拣选的器皿如此说道：\BibleQuote{不可狂宴豪饮，不可淫乱放荡，不可争斗嫉妒；但要穿上主耶稣基督，不应只挂虑肉性的事，以满足私欲。}\BibleRef{罗 13:13}\par

3. 那么，在这三类恶行中，淫乱放荡被视为如此严重的罪行，以致沾染这些罪恶的人，不仅被认为不配在教会中担任圣职，甚至被认为不配领受圣事；这是理所当然的。然而，为什么只把这种谴责局限于这类罪恶呢？因为，狂宴和豪饮受到社会舆论如此的宽容和许可，以至于甚至在为纪念真福殉道者而举行的敬礼中——不仅是在周年庆节（这本身就必须被每一位以灵性眼光看待这些事的人视为可悲的事），而且是每天——都在公开进行。如果这种腐败的做法只是因为可耻而可憎，而不是因为不虔敬，我们或许还可以把它视为一种丑闻，在力所能及的范围内予以容忍。但即便如此，我们该如何应对这样一个事实：当同一位宗徒列出一长串恶行时，其中提到了醉酒，他最后警告说，我们甚至不该与犯这些事的人一同吃饭？\BibleRef{格前 5:11} 不过，如果必须如此的话，就让我们容忍这些事只发生在家庭生活的奢侈与混乱中，发生在私人住宅内举行的宴会上吧；可是，当我们正与那些被禁止与之共享养肉身之粮的人们一起领受基督的圣体时，至少要让这种暴行远离诸圣的坟墓，远离举行圣事的圣地，远离祈祷的殿宇。因为有谁敢在私人生活中禁止那些在圣所内由群众所行、却被冠以尊敬殉道者之名的过分行为呢？\par

4. 如果\Place{非洲}是第一个尝试制止这些事的国家，她的榜样就理应被所有其他国家视为值得效法；然而，在\Place{意大利}的大部分地区，以及海外所有或几乎所有的教会中，这些做法要么——如在有些地方——从未存在过，要么——在有些地方虽然存在过，但无论是新近的还是长久的——都已被那些度圣善生活、对来生持正确见解的主教们，以勤勉和谴责加以根除和制止了——我说，既然情况如此，在已经有许多地方为我们树立了榜样的情形下，我们还怀疑是否能除掉我们道德中这种可怕的缺陷\OriginalTerm{缺陷}{defect}的可能吗？况且，我们拥有一位出身于那些地区的主教，为此我们恳切感谢天主；尽管他是一位如此温和、柔顺，总之，在主内如此明智、热切的人，以至于即使他生于\Place{非洲}，圣经也会说服他，必须为这种放纵无序的习俗所造成的创伤施行治疗。但这一邪恶造成的瘟疫是那样广泛而深重，以致在我看来，若不经由公会议\OriginalTerm{公会议}{council}的权威从中干预，便无法彻底治愈。然而，如果要从一个教会开始，那么在我看来，正如其他教会若试图改变\Place{迦太基}教会仍在坚持的事便是狂妄，同样，其他教会若愿意坚持\Place{迦太基}教会已谴责的做法，那也是极端无耻。为了在\Place{迦太基}进行这样的改革，还能期望有比那位在当执事\OriginalTerm{执事}{deacon}时便郑重谴责过这些做法的主教更好的主教吗？\par

5. 但你所忧伤的那些事，你现在应当抑制，不要严厉，而要如经上所载：\BibleQuote{以温柔的精神}。\BibleRef{迦 6:1} 请原谅我的冒昧，因为你的信向我揭示了你真诚的兄弟之爱，给了我如此大的信心，以致我敢于像对我自己那样畅所欲言地对你说话。这些过犯，至少按我的判断，可用严厉、苛酷和霸道手段之外的其他方法除去——即，与其用命令，不如用教导；与其用谴责，不如用劝告。我们至少应当这样对待群众；至于少数人的罪过，则必须用儆戒性的严厉手段。如果我们确实采用威吓，也要怀着忧伤，以圣经上的话支持我们对将临审判的威吓，好使人们因我们的话而生的恐惧，不是对我们自身权威的恐惧，而是对天主本身的恐惧。这样，首先可使属神\OriginalTerm{属神的}{spiritual}的人，或心态最接近属神的人产生印象；然后再通过最温和却又最恳切的劝告，瓦解其余群众的反抗。\par

6. 然而，由于无知和属血肉的群众习惯于把在墓地举行的这些醉酒狂欢和奢侈宴乐，不仅视为对殉道者的尊敬，而且视为对亡者的慰藉，我认为，要说服他们放弃这些地方的可耻而不合宜的做法，除了指出这些为圣经所禁止之外，还须注意：对于为安眠者灵魂的奉献\OriginalTerm{奉献}{offerings}——我们确实应当相信这些奉献有某种益处——既不可过于奢侈，超出对逝者的怀念所应有的敬重，也不可炫耀，而要慷慨地分施给所有请求分享的人；也不可出售，但若有人愿意出于虔诚而奉献金钱，应当当场施舍给穷人。这样，既可以避免他们因感到亡友的怀念被忽视而产生不小的内心忧伤，又能在教会中恰当地举行那原本是虔诚可敬的宗教服务\OriginalTerm{宗教服务}{religious service}之举。关于纵饮和醉酒，目前说这些或许已足够。\par

% structure: p0010 source=h2 target=subsection reason=relative-heading-rank; base=h2

\subsection{第二章}

7. 至于\Quote{纷争和欺诈}，我有什么资格谈论呢，因为在我们圣职人员\OriginalTerm{圣职人员}{order}内部，这些恶习比在信众\OriginalTerm{信众}{congregations}间更为严重？不过，请允许我说，这些弊端的根源在于骄傲，以及贪求人的称赞，后者也常产生伪善。这只有被对天主的爱和敬畏所充满的人，借着圣经上诸多言明的教导，才能成功地抵抗；然而，这样的人还必须在自己身上展现出忍耐和谦卑的榜样，即当受到称赞和尊敬时，只领受比所献予的更少的份，如同自己应得的一样：同时，对于尊重他的人所给予的一切，既不全盘接受，也不一概拒绝；至于他所接受的部分，他不是为了自己而接受——因为他应度完全生活于天主鉴临之下、轻看人的掌声的生活——而是为了那些人的益处，若他过分自卑以致失去他们的尊敬，便无法促进他们的福祉。因为与这有关的正是这句话：\BibleQuote{不要让人小看你年轻}；\BibleRef{弟前 4:12} 而说这话的同一位在别处也说：\BibleQuote{如果我仍讨人的欢心，我就不是基督的仆人了}。\BibleRef{迦 1:10}\par

8. 不因来自人的尊荣和赞美而欢腾，反而弃绝一切虚荣排场，这是一件大事；倘若其中某些事是必要的，也应使凡这样保留的一切，都能对赐予尊荣者带来裨益与救恩。因为经上并非徒然说：\BibleQuote{天主将粉碎那些寻求取悦于人者的骨头。}因为有什么比那被毁谤者的舌头击倒的人更为脆弱，更缺乏此处以骨头所比喻的坚忍与力量呢？尽管他明知那反对他的言论是虚假的。倘若其灵魂的骨头未曾因贪慕赞美而被打断，这种痛苦绝不会撕裂他灵魂的深处。我确信你心灵的力量：因此，我现在对你说的这些话，也是对我自己说的。不过，我相信，你愿意与我一同思考这些事是何等重要，何等艰难。因为未曾向这仇敌宣战的人，无法了解它的力量；如果当赞美还不易得时，弃绝它还算比较容易，那当赞美被奉上时，要抑制不被它所俘虏则甚为困难。然而，我们心思对天主的依恃应当如此深厚，以致当有人不当赞美我们时，我们若能与他们随意交谈，便应立即纠正他们，以免他们认为我们拥有我们本不具备的，或将那属于天主的当作我们的，或称赞我们那些虽为我们所有，甚或充裕，但按其本性却不值得称赞的事物，例如一切我们与低级动物或恶人共有的善。然而，倘若我们因天主所赐予的而理应受赞美，就让我们为那些因真正善事而感到喜悦的人庆幸；但我们不要因取悦于人而自庆，却要因我们在天主眼中确是（果真如此的话）与他们所推崇相符而自庆，且因赞美并非给予我们，而是给予天主——一切真正并当受赞美之事的赐予者——而自庆。这些事，我自己天天对自己重复，或更确切地说，是由那位所有有益训诲所从出者所重复，无论它们是在诵读圣经时所发现，或是从内心提示于我的；然而，虽与我的仇敌奋力搏斗，当我无法摆脱那献给我的赞美之诱惑力量时，我仍常被他所伤。\par

9. 我写这些事，是为使，倘若它们现在为你的圣德并非必要（因你自己的思想会向你提示此类或其他更有益的思虑，又或你的圣德超乎此类药石之需），至少我的病症可得你知晓，并让你知道那足以推动你屈尊按我的软弱之需，向天主为我恳求的事；我藉着那位命令我们彼此担负重担者的人性恳求你，求你以最迫切的转求为我代祷。关于我的生活和品行，有许多我不愿书写的事，若我们处境允许，只需我的口与你的耳作为心与心沟通之媒介，我必含泪向你坦承。然而，若那年老的\Person{萨图尼努斯}——他为我们所敬重，且在此处为众人以毫无保留、毫不虚伪的爱所爱戴，当我与你同在时，我观察到他对你的兄弟之爱与热忱——若他，我说，乐意在方便时尽早探望我们，那么，我们能与那位圣洁且属灵的人士所享有的任何交谈，都将被我们视同为与您尊威的面晤，即使不完全相同，也相差无几。我以言语无法表达其急迫之情的恳切请求，恳请你屈尊与我们一同从他求得这一恩惠。因为\Place{希波}的人民甚为担忧，且远远超过他们所当担忧的，不愿让我去到离他们如此遥远之处，且绝不肯放我独自前往，以至应允我去视察您以关怀和慷慨赐予弟兄们的那块田地；关于这田地，在你的信来到之前，我们已通过我们的弟兄兼同仆\Person{帕尔特尼乌斯}听闻，从他那里我们也获知了许多其他我们渴望了解的事。主将成就我们尚渴望的其他一切事。\par

\SourceRecord{Translated by J.G. Cunningham.；Nicene and Post-Nicene Fathers, First Series；Vol. 1.；Edited by Philip Schaff.；Buffalo, NY: Christian Literature Publishing Co.,；1887.；<http://www.newadvent.org/fathers/1102022.htm>.}

% structure: p0001 source=h1 target=section reason=page-title

\section{书信二十三（公元392年）}

\Emph{致玛息铭，我深爱的主内弟兄，堪受尊敬者，奥斯定，天主教会的司铎，在主内向您致候。}\par

在进入我决定致函尊驾的主题之前，我将简要说明此信标题所用措辞的原因，以免这些措辞使您或任何其他人感到惊讶。我写下了\Quote{致我主}，因为经上记载：\BibleQuote{弟兄们，你们蒙召是为得自由；只是不可将这自由当作放纵肉欲的机会，却要以爱德彼此服事。}\BibleRef{迦 5:13}因此，看到在给您写信的这项本分中，我实际上是在以爱德服事您，那么我称您为\Quote{我主}，不过是做了合理的事，为的是那一位给了我们这命令的惟一的真主。再者，至于我写下了\Quote{深爱的}，天主知道，我不仅爱您，而且爱您如同爱我自己；因为我深知，我渴望您得到的那些福分，正是我自己亦愿拥有的。至于我加上了\Quote{堪受尊敬}这几个字，我加上它们，并非意在承认我尊敬您的主教职务，因为对我来说，您并不是一位主教；这一点我相信您会视为我并无冒犯之意，而只是出于这样的信念：我们口中应是是，非就非。因为无论您还是任何知道我们的人，都不可能不知道您不是我的主教，我也不是您的司铎。因此，我之所以欣然称您为\Quote{堪受尊敬者}，是基于以下理由：我知道您是一个人；而我知道，人是按天主的肖像和模样受造的，并且由于自然界的秩序和规律本身就处于尊荣之中，只要人能理解他所应当理解的事物，他便能持守其尊荣。因为经上记载：\BibleQuote{人处于尊荣之中而不理解，便与无理性的畜类相似，变得同它们一样。}那么，既然您是人，我为什么不可以称您为堪受尊敬者呢？尤其因为只要您还在此生之中，我就不敢对您的悔改和得救感到绝望。再者，至于我称您为\Quote{弟兄}，您非常熟悉那神圣的诫命，按照这诫命，我们甚至要对那些否认是我们弟兄的人说：\Quote{你们是我们的弟兄}；而这与我决定写信给您——我的弟兄——的原因大有关系。既然我已说明了我在这封信中作如此开场的理由，我请求您静心留意下文。\par

当我在您的地区，正在竭尽全力表达我对那种可悲可叹的恶习的憎恶时——有些人虽以基督徒之名自夸，却毫不犹豫地为基督徒\OriginalTerm{重洗}{rebaptize}——当时不乏有人称赞您，说您并不随从这一恶习。我承认起初我并不相信他们；但后来，我想，对天主的畏惧有可能占据一个习于默想来世生命的人灵，以至阻止人犯下最为昭彰的恶行，于是我信了他们的说法，欢喜地看到，您持守这样的决定，便表明您不愿与天主教会完全隔绝。我甚至期待有机会与您交谈，以便如有可能，消除我们之间尚存的那一点分歧；然而，看，就在几天前，有人向我报告说，您为一位属于我们\Place{穆图根纳}的执事重洗了！我为他可悲的失足，也为您的罪——我的弟兄，这罪既令我惊愕又让我失望——而深感悲痛。因为我知道天主教会是什么。万民是基督的产业，地极是祂的产业。您也知道天主教会是什么；如果您不知道，那就请用心分辨，因为凡愿意学习的人，很容易就能认识它。因此，甚至为一位已在圣洗中领受了\OriginalTerm{圣洁的印记}{seal of holiness}的异端者重洗——这印记是基督教会的惯例传承给我们的——也毫无疑问是罪；但为一位天主教徒重洗，则是最恶劣的罪行之一。然而，因我仍对您持有良好印象，我并未相信那报告，于是亲自去了穆图根纳。我未能找到那不幸的人，但从他的父母那里得知，他已被任命为您的一位执事。尽管如此，我仍对您抱有很大好感，以致我不愿相信他已被重洗。\par

因此，我亲爱的弟兄，我藉着我们主耶稣基督的\OriginalTerm{天主性和人性}{divine and human natures}恳求你，请屈尊答复这封信，告诉我所发生的事，并且你要明白，我打算在教会里向我们的弟兄高声朗读你的信。我写这封信，是为了避免我后来做出你未曾预料的事，以致冒犯你的\OriginalTerm{爱德}{Charity}，并使你有机会在我们共同的朋友面前合理地抱怨我。我看不出有什么理由能合理地阻止你回复这封信。因为如果你确实再施洗，当你写信说你在做他们即使你不情愿也会命令你做的事时，你就不必担心你的同僚们；而且，如果你进一步用你所知的最有力的论证为这事辩护，认为这是应该做的，你的同僚们非但不会因此不悦，反而会称赞你。但如果你不再施洗，我的弟兄\Person{马克西明}，就持守住你的基督徒自由；我恳求你，持守这自由：定睛注视基督，不要因人的指责而惧怕，不要在任何人的权势前颤抖。这世上的荣耀转瞬即逝，尘世野心所追求的一切也都是转瞬即逝的。无论是步步升高的宝座，还是有华盖的讲道台，或是成群\OriginalTerm{献身贞女}{consecrated virgins}的游行和咏唱，都不能在基督的审判台前，为那些现今享有这些荣耀的人提供辩护。到那时，良心将发出控诉的声音，那位审判人良心的主将宣布最后的判决。在此世被尊崇的，那时将成为重担；在此世提升人的，在那一天将重重压在他们身上。那些在此期间为了教会的益处，以示对我们的尊重而做的事，也许能由无亏的良心证实；但它们绝不能为有罪的良心提供遮掩。\par

那么，如果确实是出于虔诚和热心，你避免施行第二次圣洗，反而接受天主教会所施行的圣洗，将它视为那一位真母亲的行为，并视为进入基督那遍及地极的唯一产业的凭证。这位母亲向万民敞开怀抱欢迎他们，好使他们得到重生，并在他们重生后以母亲般的乳汁养育他们。如果，我说，你确实这样做，为什么不迸发出欢欣而独立的信仰宣认呢？为什么把那本来可以如此有益地照亮的灯藏在斗底下呢？为什么不撕掉并扔掉你胆怯奴性之卑污旧衣，穿上基督徒勇敢的全副武装走出去，说：\Quote{我只承认一个以父及子及圣神之名祝圣并盖印的圣洗：无论在哪里找到这圣事，我都应承认并尊重它；我绝不摧毁我认出属于我主的标记；我绝不羞辱我君王的旗帜}呢？就连那些分基督衣服的人，也没有粗暴地将那无缝的长袍撕成碎片；\BibleRef{若 19:24}，而那些人当时还不信基督的复活，他们甚至是在亲睹祂的死亡。那么，如果迫害者在基督挂在十字架上时，尚且忍住不撕裂祂的衣服，为什么基督徒现在，当基督坐在天上宝座之时，却要摧毁祂所建立的圣事呢？如果我生在那古时的子民时代，当一个犹太人，而我也不能成为更好的，我无疑会接受割损礼。那个\Quote{由信德而来的义德的印记}在那个分施中如此重要，直到被主的来临所废除，以至于天使本要杀死\Person{梅瑟}的婴孩，若不是孩子的母亲急忙拿石片给孩子行割损，藉此圣事避免了即将来临的死亡。这圣事也曾使\Place{约旦河}的水停住，使之倒流回源头。主自己在婴孩时也接受了这圣事，尽管祂被钉在十字架上时废除了它。因为这些属灵祝福的标记并没有被定罪，而是让位给更适合后期分施的其他标记。正如割损因主的第一次来临而被废除，圣洗也将因祂的第二次来临而被废除。正如现在，信德的自由已经来到，奴役的轭已被除去，没有基督徒再在肉体上受割损；同样，到那时，当义人与主一同为王，恶人被定罪之后，没有人再受洗，但那两样礼规所预表的实体——即内心的割损和良心的洁净——将永远存留。因此，如果我生活在先前分施的时代，是一个犹太人，有一个\Place{撒玛黎雅}人愿意成为犹太人，放弃主亲自谴责的错误，因为主曾说：\BibleQuote{你们朝拜的，你们不认识；我们朝拜的，我们却认识，因为救恩是出自犹太人}，\BibleRef{若 4:22}——那么，我说，如果一个撒玛黎雅人，已被撒玛黎雅人施了割损，表达了他愿意成为犹太人的意愿，我们就没有理由胆敢坚持重复这礼仪；相反，我们将不得不赞许天主所命令的，尽管它是由异端者施行的。但是，如果在受了割损的人的肉体上，我找不到重复割损的余地，因为只有一个肢体是受割损的，那么在人唯一的心中，就更没有余地重复基督的圣洗了。因此，你们这些想要施洗两次的人，必须寻找那些有双重心灵的人，作为这种双重圣洗的对象。\par

5. 因此，如果你确实不行\OriginalTerm{重洗}{rebaptize}，就请坦诚地公布你是在做正确的事；并将此事写给我，不仅毫无畏惧，更要满怀喜乐。我的弟兄，不要让你那一方的任何会议阻止你迈出这一步：因为如果这让会议不悦，他们就配不上有你列席其中；但如果这让他们喜悦，我们深信，藉著\OriginalTerm{主}{Lord}的仁慈，你与我们之间很快就会享有和平，主从不遗弃那些惧怕惹祂不悦、且竭力在祂眼中做悦纳之事的人；也不要让我们那危险的重担——将来必须交账的尊位——成为障碍，使得我们这些信仰\OriginalTerm{基督}{Christ}、在家中彼此分享日用的粮\BibleRef{玛 6:11}的信众，无法一同在基督的餐桌前坐下。我们岂不悲叹，大多数情况下，当婚姻使夫妻成为一体\BibleRef{创 2:24}时，他们在基督的名内彼此誓许忠贞，却因分属不同的\OriginalTerm{共融团体}{communions}而撕裂基督的身体\BibleRef{格前 12:27}？如果藉著你温和的举措、智慧，以及你对那位为我们倾流宝血者所应有的爱德，这一\OriginalTerm{裂教}{schism}，这一如此可耻的\OriginalTerm{恶表}{scandal}——它使撒殚得意并使许多灵魂丧亡——能被从这个区域除去，那么谁能充分描述主为你预备的赏报是多么光辉灿烂，因为从你身上将传递出一个榜样，这榜样若被效仿（这是极易做到的），就会为基督的所有肢体带来健康，而祂遍布整个\Place{非洲}的肢体目前正悲惨地奄奄一息？我多害怕，因为你不能看到我的内心，会以为我说话是出于讽刺，而非出自真诚的爱！但是，除了将我的话呈现在你眼前，并将我的心呈现在天主面前之外，我还能做什么呢？\par

6. 让我们将那些双方无知之人惯常彼此抛出的虚妄异议，从我们中间除去吧。你那一方不要拿\Person{马卡留斯}的迫害来指责我。我也决不拿\OriginalTerm{割裂派}{Circumcelliones}的暴行来谴责你。你若不为后者负责，我也不为前者担待；这些事与我们无关。主的禾场\BibleRef{玛 3:12}尚未扬净——不可能没有糠秕；让我们只管祈祷，并尽我们所能使自己成为好麦子。我不能对我们的执事被\OriginalTerm{重洗}{rebaptize}一事保持沉默；因为我知道，我的沉默会给我自己带来多大的伤害。我并不打算把自己的时间花费在空洞地享受教会尊位上；我立意行事，时刻铭记：我必须为托付给我的羊群向那位总司牧\BibleRef{伯前 5:4}交账。倘若你宁愿我不如此写信给你，我的弟兄，你必须以我的恐惧为由原谅我；因为我极其恐惧，生怕如果我保持沉默并隐藏我的想法，别人会被你重洗。因此，我决意，凭藉主赐的力量和机会，来安排这场讨论，好让我们所有属于我们共融团体的人，都能知道\OriginalTerm{天主教会}{Catholic Church}与\OriginalTerm{异端}{heresy}和\OriginalTerm{裂教}{schism}相隔多远，并且知道应该何等警惕地防范那等待莠子\BibleRef{玛 13:25}与从主葡萄树上砍下之枝杈\BibleRef{若 15:6}的毁灭。如果你乐意借同意公开宣读双方书信来与我进行这样的商谈，我将喜不自胜。如果这建议令你不悦，我的弟兄，我除了即使没有你的同意，也向天主教信众宣读我们的书信，以便教导他们之外，还能做什么呢？但如果你不肯屈尊给我回信，我至少决意宣读我自己的信，好使当你对自己做法的疑虑被知晓时，其他人可能以被重洗为耻。\par

7. 然而，我不会在士兵的在场下做这件事，免得你们任何人以为我想用暴力方式行事，而非按和平利益的要求；我只有在他们离开后才做，好使所有听我讲的人都能明白，我无意强迫人加入任何一方的共融，而是愿真理能显示给那些在寻求真理时不受恐惧惊扰的人。在我们这一方，绝不利用人们对世俗政权的恐惧；在你们那一方，也不要让\OriginalTerm{割裂派}{Circumcelliones}的暴民进行恐吓。让我们专注于实际辩论的问题，让我们的论证诉诸理智和\OriginalTerm{圣经}{Divine Scriptures}的权威教导，尽我们所能地冷静平和；让我们求、寻、叩门\BibleRef{玛 7:7}，好使我们能获得、寻见，并使门为我们开启，藉著天主对我们共同努力和祈祷的祝福，从而首先彻底清除我们区域中那个使\Place{非洲}蒙羞的不虔敬。如果你不相信我愿意将讨论推迟到士兵离开之后，你可以延迟答复，直到他们离去；而且，如果趁他们还在这里的时候，我就想向民众宣读我自己的信，那么这封信的产生本身就会证明我失信。愿主以仁慈阻止我行任何违背道德，且背离祂所乐意藉著将祂的轭放在我身上而启发我的美善决心的事！\par

8. 我的主教若在这里，也许更愿亲自致函给\OriginalTerm{阁下}{your Grace}；或者我的信，即便不是奉他的命令，至少也会得到他的许可才写。但是，他不在场，又听闻这位执事被\OriginalTerm{重洗}{rebaptize}一事是近日才发生的，我便没有拖延，以免任由此举引发的感受冷却下去，我深为悲痛——我认为这实在是一位弟兄的死亡——之情的催促而行动。我这悲伤，或许藉著我主仁慈与眷顾的帮助，将被我们与你之间和好所带来的补偿之喜乐所治愈。愿我主天主赐予你一颗平静与和解的心，我最亲爱的\OriginalTerm{阁下}{lord}及弟兄！\par

\SourceRecord{Translated by J.G. Cunningham.；Nicene and Post-Nicene Fathers, First Series；Vol. 1.；Edited by Philip Schaff.；Buffalo, NY: Christian Literature Publishing Co.,；1887.；<http://www.newadvent.org/fathers/1102023.htm>.}

% structure: p0001 source=h1 target=section reason=page-title

\section{第25封书信（主历394年）}

\Emph{致奥古斯丁，我们可敬爱的主内弟兄，\Person{保利努斯}和\Person{特辣西亚}，罪人。}\par

1. 那催迫我们的基督的爱，以及那藉着共同\OriginalTerm{信德}{faith}的联系将我们——虽分隔两地——合而为一的爱，促使我放下畏惧，提笔给你写信；并通过你的著作——那些充满知识宝藏、甘美如天上蜜糖、成为我灵魂的良药与食粮的作品——在我内心深处为你安置了一席之地。这些著作目前我有五册，是藉着我们真福且可敬的主教\Person{阿利比乌斯}的善意获得的，这些著作不但用于我个人的训导，也供许多城镇的教会使用。这些著作我正在阅读：我从中得到极大喜乐；我从中找到食粮，不是那可朽坏的，而是那藉着我们的\OriginalTerm{信德}{faith}赋予永生实体的，因为我们在主耶稣基督内，藉着这\OriginalTerm{信德}{faith}成了祂身体的肢体；因为信友的著作与榜样确实大大增强那\OriginalTerm{信德}{faith}，这\OriginalTerm{信德}{faith}不注目于看得见的事物，却以那种对全能天主的真理所包含的一切无不坦然接受的爱，渴慕看不见的事物。啊，真称得上地上的盐，我们的心灵赖你以保存，免遭世俗错谬的败坏！啊，配得教会灯台上位置的光，你将上部圣所七枝灯台的灯油所滋养的火焰光辉，在\OriginalTerm{天主教}{Catholic}城镇中广传，你甚至也驱散了\OriginalTerm{异端}{heresy}的浓雾，并以你那光明的论证之光，从黑暗的混乱中解救出真理的光辉。\par

2. 你已看到，我在基督我主内所亲爱的、敬重的、欢迎的弟兄，我是以何等亲密的情谊自称认识你，以何等的惊奇景仰你，以何等的爱情拥抱你，因为我得以藉着你的著作每日与你交谈，并由你口中的气息得到滋养。因为你的口，我大可称之为输送活水的管道，出自永恒泉源的沟渠；因为基督在你内成了\BibleQuote{涌到永生的活水。}\BibleRef{若 4:14}的泉源。因着对这活水的渴望，我的灵魂渴慕不已，我枯干的心田渴望被你江河的丰沛所淹没。既然你已藉此你的\OriginalTerm{五书}{Pentateuch}，完全武装了我以对抗\OriginalTerm{摩尼教徒}{Manichæans}，如果你还针对其他敌人预备了维护公教\OriginalTerm{信德}{faith}的论著（因为我们的仇敌，以其千般邪恶诡计，必须以种种兵刃——其种类须与攻击我们的诡计相当——方能克胜），我恳请你从你的军械库中为我取出这些，不要拒绝提供给我\Quote{正义的甲胄。}因为即便目前，我在工作中仍背负着沉重负担，作为罪人，在罪人中我是老手，但在为永恒君王效力上，却是一名未经训练的新兵。我先前曾何等不幸地仰慕这世界的智慧，并委身于那些我现在看出无益的文学，和如今被我摒弃的智慧，我在天主面前曾是愚蠢而喑哑的。当我在仇敌的团体中渐趋老迈，并在思想中徒劳无功时，我举目向山岳仰望，仰望法律的规诫和\OriginalTerm{恩宠}{grace}的赠礼，我的救助遂自主而来，祂没有按我的邪恶来报应我，却开启了我这盲目的眼睛，解开了我的桎梏，使我这曾因罪恶而高傲的人谦卑下来，好能于此满怀恩宠的谦卑中提拔我。\par

3. 因此，我虽步履蹒跚，却追随义人的伟大榜样，但愿能藉着你们的祈祷，得到那因天主的怜悯而先已得着我的。所以，请你引导这在地上匍匐的婴儿，以你的步履教他行走。因为我不愿你按我自然出生的年龄来判断我，而应按我属灵新生的年龄来判断。因为就本性生命而言，我的年龄已到宗徒们在\Place{丽门}凭他们话语的能力治好的那位瘸子的岁数。但论及我灵魂的诞生，我的年龄尚如那些婴孩，他们因那瞄准基督的毒手而被祭杀，以无愧于如此尊位的鲜血，作羔羊祭献的先导，并成了主受难的前驱。\BibleRef{玛 2:16}因此，在\OriginalTerm{天主的话}{the word of God}方面我既是婴孩，在属灵年龄上既是哺乳幼儿，请以你的话语——那\OriginalTerm{信德}{faith}、\OriginalTerm{智慧}{wisdom}和\OriginalTerm{爱德}{love}的乳汁，滋养我，以满足我迫切的渴望。若你只考虑我们共有的职务，你是我的弟兄；但若你考虑你领悟力之成熟以及其他能力，你虽在年龄上比我年轻，却是我的父亲；因为可敬的智慧已使你年纪虽轻，却在德业上臻于成熟，并享有属于老年人的尊荣。那么，请培育我，坚固我，因为我正如前所述，在\OriginalTerm{圣经}{the sacred Scriptures}和属灵操练上尚是孩童；而且，在经历了漫长争斗和屡次船难之后，我只学会少许技能，甚至如今仍艰难地从今世的浪涛中探出头来，而你已经安然立于岸上，请接纳我到你胸怀的安全庇护所，若蒙不弃，我们便可一同驶向\OriginalTerm{救恩}{salvation}的港湾。同时，在我努力摆脱此世危难和罪恶的深渊时，请以你的祈祷，如同以一块木板支撑我，使我得以脱离今世，如同人抛下一切，从船难中逃生。\par

因此，我已竭力卸下一切可能导致我拖延的行李和衣裳，以便顺从命令，并靠基督的帮助维持，能不被任何为肉体的衣裳或为明日的忧虑所阻碍，游过今世的生命之海——这海波涛汹涌，回荡着我们罪孽的吠声，如同\OriginalTerm{斯库拉}{Scylla}的犬吠，将我们与天主隔绝。我并不夸耀我已成就此事；即便我可以夸耀，我也只会在主内夸耀，因为成就之事属主，而我们当尽的本分是渴望；但我灵魂所热切追求的，是让主的判断成为她的首要渴望。你可以判断，一个渴望能渴望去承行主旨的人，他在有效履行主旨的道路上，已走了多远。尽管如此，就我而言，我喜爱祂圣所的美，并且，若凭我自己，我会选择在主殿中处于最卑微的位置。但那乐意将我从母腹中分别出来\BibleRef{迦 1:15}，又将我从血肉之谊中引出、归向祂恩宠的天主，祂却乐意从地上和痛苦的深渊提拔我——尽管我毫无功绩——从泥泞和粪土中抬举我\BibleRef{咏 113:7-8}，使我得与祂百姓的权贵同列，并任命我与你处于同等地位；这样，尽管你的功绩远超于我，我却能在职务上与你平起平坐。\par

因此，我之据有这荣誉（我自认不配），并为自己求得与你弟兄情谊的联结，并非出于我自己的擅自，而是按主的旨意和任命；因为我从你品格的神圣深信，你是从真理受教，\BibleQuote{不要心高气傲，倒要俯就卑微的人} \BibleRef{罗 12:16}。因此，我盼望你会欣然仁慈地接纳我们谦卑地向你表达的爱的确信，这爱，我相信，你已经通过最可称颂的司铎\Person{阿利比乌斯}领受了，我们（经他准许）称他为我们的父亲。因为他无疑已亲自给你树立了榜样：既在我们还是陌生人时，且远超我们所配得的，就爱我们；因为他凭着深远而自我扩散的真爱精神，得以藉着亲情看见我们，又藉着书信与我们接触，即便那时我们还不为他所知，且被陆地和海洋的远距所隔绝。他已借上面提及的赠书，向我们提供了他对我们爱的初步证明，以及你的爱的证据。正如他极其关切，要使我们在认识你以后，不只凭他的见证，更全面地借着你著作中所流露的口才和信德，而受到激励，对你怀有热切的爱；同样，我们相信，他也以同等的热忱，使你效法他的榜样，对我们报以极其温暖的爱。啊！在基督内的弟兄，我们所爱的、可敬的、热切期盼的：我们切愿天主的恩宠，既已与你同在，永存不渝。我们以最热诚的弟兄友爱，向你的全家，以及在主内与你圣德为伴、效法你圣德的每一个人，致以问候。我们请求你，在接纳它时，祝福我们送至你爱德的一个饼，以表我们与你同心合一。\par

\SourceRecord{Translated by J.G. Cunningham.；Nicene and Post-Nicene Fathers, First Series；Vol. 1.；Edited by Philip Schaff.；Buffalo, NY: Christian Literature Publishing Co.,；1887.；<http://www.newadvent.org/fathers/1102025.htm>.}

% structure: p0001 source=h1 target=section reason=page-title

\section{第二十六封信（公元395年）}

\Emph{致\Person{利琴提乌斯}} \Emph{\Person{奥古斯丁}寄}\par

我好不容易觅得一个机会给你写信：谁会相信呢？但利琴提乌斯必须相信我的话。我不愿你好奇地探询此事的原因和理由；因为尽管这些理由可以给出，但你对我的信赖便免除了我提供它们的义务。再者，我收到了你的信件，是通过那些无法带回我回信的送信人。至于你托我询问的那件事，我已按我认为当提出的限度以书信加以关注；但结果如何，你或许已看到了。若我尚未成功，及至结果为我所知，或当你亲自提醒我时，我将更迫切地敦促此事。到此为止，我已向你述及那些在其中我们能听见今生锁链声响的事务。我离开这些。现在，请简短地接受我内心对你永生之希望所怀忧虑的倾诉，以及如何为你开启通往\OriginalTerm{天主}{God}之路的问题。\par

我担心，我亲爱的利琴提乌斯，你不断拒绝并畏惧智慧的约束，仿佛它们是枷锁，却正牢牢地、致命地被束缚于必朽的事物上。因为智慧，起初虽以纪律的劳作约束人、制服人，但随即赐予他们真正的自由，并使他们在自由中拥有并享受它本身，且使之富裕；虽然起初它借助暂时的约束来教育他们，但随后便将他们拥入它永恒的怀抱，这怀抱是一切可想枷锁中最甜蜜、最牢固的。我承认，这些最初的约束确实有些难以忍受；但我不能称智慧最终的约束为沉重的，因为它们最甜蜜；也不能称之为容易的，因为它们最牢固：总之，它们拥有一种无法形容的特质，却能成为\OriginalTerm{信德}{faith}、\OriginalTerm{望德}{hope}和\OriginalTerm{爱德}{love}的对象。反之，今世的枷锁，则有真正的严酷与迷惑的诱惑，确定的痛苦与不确定的欢乐，辛苦的劳碌与不安的休息，充满苦难的经历，以及无幸福可言的希望。而你竟将颈项、双手和双脚交给这些锁链，渴望承受此类荣誉的重负，认为若劳苦不以此方式获得回报便是徒然，并且自愿立志固守于那无论是劝导还是强迫都不应使你前往之处？或许你会用\Person{泰伦提乌斯}剧中奴隶的话回答：\par

\Quote{嚯，你在这里倾泻哲言呢。}\par

那么，接受我的话吧，好让我倾泻而不浪费它们。但如果我歌唱，而你宁愿伴另一曲调起舞，即便如此我也不后悔给予劝告的努力；因为歌唱的练习本身便带来愉悦，即使歌曲未能激起那位被怀着爱关顾而为之歌唱的人的回应动作。你的信件中有一些用词错误引起了我的注意，但当我为你行为及整个人生的忧虑所困扰时，我认为讨论这些便显得琐碎了。\par

如果你的诗篇因结构缺陷而受损，或违反韵律法则，或因节奏不完美而刺耳难听，你无疑会感到羞愧，并会毫不迟延、不得安宁，直到你整理、修正、重铸并平衡你的作品，投入任何数量的专注研究及辛劳，以获得并操练诗歌写作的技艺：但当你自己被混乱的生活所毁损，当你违反天主的法律，当你的生活既不符合朋友为你怀有的正当期盼，也不符合你自身学识所赋予的光辉时，你竟以为这是可以抛诸脑后、弃之不顾的小事吗？仿佛你真的认为自己比自己的声音更无价值，并以为因无序的生活得罪天主，比因无序的音节招致文法家的指责更具微末不足道。\par

你如此写道：\Quote{但愿往日的晨光，以其欢悦的战车，带回逝去的明媚时光，那我们在\Place{意大利}的中心及高山之间共同度过的日子，那时我们体验着属于善者的慷慨悠闲与纯洁特权！无论是严寒的冬日携其冻结的冰雪，还是西风粗鲁的吹拂及北风怒号的肆虐，都不能阻止我以热切的步伐追随你的足迹。你只需表达你的愿望。}\par

我若不表达这愿望，不，若不强迫命令，或恳求哀请你来跟随我，我便有祸了。然而，若你的耳朵对我的声音关闭，就让它向自己的声音开启，并留意你自己的诗篇：聆听你自己吧，哦，朋友，你这最为固执、无理且毫无感应之人。你的舌头是金的，内心却是铁的，这于我何益？我将如何不借诗句，而是以哀叹，来充分悲恸你的这些诗章？我从其中发现了一个何等样的\OriginalTerm{灵魂}{soul}，何等样的心智，而我却不被允许将其抓住并作为奉献献给我们天主！你正等着我表达愿望，愿你变好，并享受安息和幸福：仿佛任何一天能比我在天主内享受你天赋心志的那天更令我欢欣照耀，又仿佛你不知道我如何为你而饥渴，又仿佛你没有在这首诗本身中承认这一点。回到你写下这些事时的心境；现在再次对我说：\Quote{你只需表达你的愿望。} 那么，这就是我的愿望，如果我的表达足以感动你遵行：将你自己交给我——将你自己交给我的\OriginalTerm{主}{Lord}，祂是我们双方的主，也是赋予你才智的那一位：因为若非藉着祂，我是什么，不过是你的仆人，并在祂之下是你的同仆？\par

5. 不，难道祂没有表明自己的旨意吗？请聆听福音：它宣告说，\BibleQuote{耶稣站着大声喊说。} \BibleRef{若 7:37} \BibleQuote{凡劳苦和负重担的，你们都到我跟前来，我要使你们安息。你们背起我的轭，跟我学吧！因为我是良善心谦的：这样你们必要找得你们灵魂的安息，因为我的轭是柔和的，我的担子是轻松的。} \BibleRef{玛 11:28} 如果这些话没有被听见，或只是用耳朵听见，那么，\Person{Licentius}啊，你难道期望\Person{Augustine}向他的同仆发出命令，而不是抱怨他主人的旨意被轻视吗？祂命令，不，是邀请，甚至恳求所有劳苦的人到祂那里寻求安息！但对于你那强壮而骄傲的颈项，确实，世俗的轭似乎比基督的轭更容易；然而请想想，祂所加的轭，是谁加的，又是以什么报偿加的。到\Place{Campania}去吧，以\Person{Paulinus}——那位卓越而神圣的天主的仆人——为榜样学习吧：他如何毫不犹豫地摆脱了多么大的世俗尊荣，从他那因谦卑而真正高贵的颈项上卸下，为能如他所行的那样，置于基督的轭下；如今，他心神安宁，温柔地在祂内喜乐，以祂为自己道路的向导。去吧，去学习他怀着何等富饶的心灵向祂献上赞颂的祭献，将他从祂所领受的一切美善归还于祂，免得因为没有将从赐予者那里得到的一切储藏于祂内，而失去一切。\par

6. 你为何如此激动？为何如此摇摆不定？为何你转耳不听从我们，却去听从那致命享乐的虚幻想象？它们是虚假的，它们会朽坏，且引人走向永罚。它们是虚假的，\Person{Licentius}。\Quote{愿真理，}如你所愿，\Quote{通过明证向我们显明，愿它比\Place{Eridanus}更清澈地流出。} 唯有真理宣告何者为真：基督就是真理；让我们到祂那里去，好能脱离劳苦。为使祂治愈我们，让我们背起祂的轭，向祂学习吧，因为祂是良善心谦的，这样我们必能寻得灵魂的安息；因为祂的轭是柔和的，祂的担子是轻松的。魔鬼想要把你当作装饰品佩戴。那么，你若在地上发现一个金圣爵，你会把它献给天主的教会。但你却从天主领受了在灵性上宝贵如黄金的才能；你竟将这些献去伺候你的私欲，将自己交给\Person{撒殚}吗？不要这样做，我恳求你。愿你在某个时候能察觉，我是怀着何等悲伤沉痛的心写了这些话；我祈求你，若你已不再自视珍贵，就请怜悯我。\par

\SourceRecord{Translated by J.G. Cunningham.；Nicene and Post-Nicene Fathers, First Series；Vol. 1.；Edited by Philip Schaff.；Buffalo, NY: Christian Literature Publishing Co.,；1887.；<http://www.newadvent.org/fathers/1102026.htm>.}

% structure: p0001 source=h1 target=section reason=page-title

\section{书信二十七（主历395年）}

\Emph{奥古斯丁向在基督内至圣可敬、堪受至高赞美的我的主、我的弟兄\Person{保利努斯}，在主内致问候。}\par

1. 卓越的人、卓越的弟兄啊，曾有一时，我的心灵尚不认识你；我命令我的心灵耐心忍受你的面容仍不为我双眼所见，但它几乎——不，全然——拒绝服从。它果真在耐心忍受吗？若是这样，那为何对你临在的渴望正煎熬着我至深的灵魂？因为若我身受肉体病痛，而这些痛苦并未扰乱我心灵的宁静，我便堪称为在耐心忍受；但若我不能泰然承受见不到你的缺失，却称我这种心态为耐心，那将是无法容忍的。然而，若我与你分离时竟显为有耐心，因你是如此卓越，那或许更令人无法容忍。因此，我在这种缺失下感到不满足是好的，因为若我对此感到满足，人人都有理由对我不满。我遭遇的事虽怪却真：我因不见你而忧伤，而这忧伤本身却安慰了我；因为对于如你一般的善人之离去，那种能轻易得安慰的坚毅，我既不羡慕也不渴求。因为我们岂不渴望天上的\Place{耶路撒冷}吗？我们越急切地渴望它，岂不也越耐心地为了它而忍受一切吗？谁能如此压抑见你时的喜乐，以致当你不再出现时也不觉痛苦呢？至少我两者都无法做到；既知倘若我能做到，那也只能是践踏了正当而自然的感情，我便欢欣于我不能，且在这欢欣中寻得些许安慰。因此，安慰我的不是忧伤的消除，而是对此忧伤的沉思。我恳求你，不要以你卓越著称的虔敬严肃来责备我；不要说，既然你已向我显示你的心灵——即\OriginalTerm{内在之人}{inner man}，我却因尚未认识你而忧伤，这是错的。因为若我寄居某地，或在你所属的城中，得以认识你为我的弟兄和朋友，且是一位如此卓越的基督徒、如此高尚的人，你怎能认为，若不许我得知你的居所，我竟不会感到失望呢？那么，我既尚未见到你的面容与形体，那藏有我已知如同己有的心灵的居所，又怎能不哀伤呢？\par

2. 因为我已读了你的信，这信流溢着奶与蜜，显露出你在虔敬的引导下寻求主时内心的纯朴，并为主带来光荣和尊荣。弟兄们也读了，且在你那些天主所赋予丰富卓越的恩赐中，获得了不知疲倦、难以言喻的满足。凡读了信的人都将它带走，因为他们读时，信也将他们带走了。言语无法表达你的信所散发的基督的馨香是多么甜美。那封信通过将你呈现在我们眼前，激起了多强烈的愿望，要更充分地认识你！因为它既让我们得以辨认你，又促使我们渴慕你。因为它在某种意义上越有效地使我们意识到你的临在，就越使我们因你的不在而焦躁。所有在那里看见你的人都爱你，并希望得到你的爱。赞美与感谢都归于天主，你是因祂的恩宠而成为现在的你。在你的信中，基督被唤醒，为祂乐意为你平息风浪，引导你的脚步走向祂完美的恒稳。信中，读者看见一位妻子，她不使丈夫变得柔弱，反而因与他的结合，自己得以分享他本性的力量；对于那位在你内，与你完全合一，并因灵性的纯净而坚固的纽带与你相连的她，我们愿以尊敬你的圣德之礼回致敬意。信中，\Place{黎巴嫩}的香柏被砍倒在地，由爱的巧艺造成方舟的形状，无畏腐朽，劈开此世的波浪。信中，光荣被轻看，为要获得它；世界被舍弃，为要赢得它。信中，婴孩，甚至\Place{巴比伦}的更强之子，即暴乱与骄傲的罪，都被摔在磐石上。\par

3. 这些以及其他如此令人愉悦而神圣的景象，都呈现给你信的读者——那封信展现出真实的信德、美善的望德、纯洁的爱德。它向我们呼出你对主的殿宇的渴望、思慕与憔悴！它被何等神圣的爱所启迪！它如何洋溢着一颗真诚心灵的丰盛宝藏！它向天主献上何等的感谢！它从祂那里求得何等的祝福！是你的信的优雅还是热忱，是光明还是赐予生命的能力，在其中最为闪耀？因为它怎能同时安慰我们，又激励我们？它怎能将滋润的甘霖与无云晴空的明亮结合起来？这是怎么回事？我问；我又如何回报你，除非将我全然献给在祂内你完全属于的那一位，使我成为你的？若这算为微小，至少是我所有能给予的。但你使我不再以之为小，因你屈尊在那信中如此称赞我，以致当我以将自己给你作为回报时，若我仍视此礼物为小，就当负有不信你的证言之责。我确实羞愧于相信关于我的如此多善言，但我更不愿拒绝相信你。我有摆脱这窘境的一条出路：我不相信你对我品格的评价，因为在你所描绘的肖像中，我不认识我自己；但我会相信我被你所爱，因为我毫无疑问地感知并体会到这一点。这样，我将既不会在评判自己时轻率，也不会对你的看重忘恩负义。此外，当我将自己呈献给你时，这并非小奉献；因为我呈献的是一位你深切热爱的人，一位虽然并非你所认为的那样，却是你正为他祈祷，使他成为那样的人。如今我更加恳切地请求你的祈祷，以免你因以为我已是那非我，而疏忽祈求补足我所欠缺的。\par

4. 呈递此信与阁下、至卓越的爱德者，是我最亲爱的一位朋友，自幼年便与我极为相知。他的名字见于\WorkTitle{论宗教}一书；您圣善的，正如您在信中所言，曾以极大的愉悦阅读该书，无疑是因为那位如此良善之人——即送书给您者——的推荐，使之更为您所接纳。然而，我并不愿您轻信那位与我如此亲密的朋友可能为赞美我而说的话。因为我常注意到，他并非有意说谎，却因友情的偏私，在对我的评价上出错，竟以为我已拥有许多恩赐，而我的心其实正切望主将这些恩赐赐予我。若他在我面前尚且如此，谁又不会猜想，当他不在我面前时，他心中满溢的情感会说出许多胜于事实的赞美之辞呢？他将呈上供您垂览，并审查我所有的论述；因我不知自己写过什么——无论是写给教会之外的人，还是写给弟兄们的——是他所没有的。然而，我的圣善的\Person{保利努斯}，当您阅读这些时，不要因真理藉我这软弱的工具所说的话，而如此被吸引，以致忽略仔细留意我本人所说的话；免得您在急切地畅饮那些赐予我这仆人的美善与真实之事时，竟忘记为我错误与过失的赦免祈祷。因为在所有若被仔细留意便会公正地使您不悦的事上，您所看到的是我本人；但在所有我的书中被您公正地嘉许的事上，藉着那赐予您的圣神的恩赐，应当受爱慕的、应当受赞美的，乃是那生命之泉源所属的那一位，在祂的光明中，我们将看见光明，不再像如今这般模糊不清，而是面对面。\BibleRef{格前 13:12} 每当我在重读自己的著作时，发现其中有源于我内旧酵母之作用的事物，我便真诚地悲伤自责；但若我所言的任何事，赖天主的恩赐，是出自真诚与真理的无酵饼，我便战战兢兢地以此为乐。因为我们有什么不是领受的呢？然而，或许有人会说，天主赐予更多、更大恩赐的人，其份比领受更少、更小恩赐的人更好。确实如此；但从另一方面看，拥有小的恩赐并向天主献上应当的感谢，胜过拥有大的恩赐却欲将功劳归于自己。我的弟兄，请为我祈祷，使我得以真诚地作出这样的承认，使我的心不与口舌相违。我恳请您祈祷，使我不贪求对自己的赞美，而是向上主呈献赞美，敬拜祂；如此我必脱离我的仇敌而安全。\par

5. 还有一事可能使您更热切地爱这位持我信的弟兄；因他是可敬且真福的主教\Person{阿利比乌斯}的亲戚，您全心爱他，且理应如此：因为凡高度评价此人者，便是高度评价天主所赐予他的宏大仁慈与奇妙恩赐。因此，当他读了您的请求，希望他为您写一篇自传概要，而他虽因您的善意而愿意为之，却因自己的谦卑而不愿为之；我见他无法在爱与谦卑的诉求间作出决定，便将这负担从他的肩头移到我自己的肩上，因他通过信函嘱咐我如此行。我将因此，赖天主的助佑，不久便将\Person{阿利比乌斯}本人真实地呈现于您心中：因为我主要担心的是，他会害怕宣布天主所赐予他的一切，免得（既然他所写下的内容除您之外还会被他人读到）他在那些不太善于分辨的人看来，似乎不是在称赞天主施予人的美善，而是在炫耀自己的功劳；并且，这样一来，您——深知如何正确解读此类陈述——便会因他顾及他人的软弱，而失去那本应与您这位弟兄分享的事。我本已可完成此事，您本已在阅读我对他的描述，若非我的弟兄突然决定比我们预期的更早出发。我恳请您从心中和口中给予他亲切的欢迎，就如您与他并非初识，而是如我一般早已相知。因为他若不怕向您敞开心扉，他将在很大程度上——若非完全——被您的言辞治愈；因为我渴望他能时常听到那些对朋友怀有超越此世之爱的人的言语。\par

6. 即使\Person{罗曼尼安}未计划去拜访您尊贵的爱德，我亦已决定修书向您推荐他的儿子\Person{利琴提乌斯}，他亲如己出（他的名字您也将在我的某些书中找到），以便他不仅因您言语的声响，更因您灵性力量的榜样而受到鼓励、劝勉和教导。我热切渴望，当他的生命尚在青苗之际，莠子能变为麦子，他能相信那些凭经验深知他急于置身其中的危险的人。从我年轻朋友的诗作及我给他的信中，您那至为仁慈和体贴的智慧可察觉出我为他而怀的忧伤、恐惧与挂虑。并非没有希望，赖主的恩宠，我可藉由您而从这种对他的焦虑中解脱。\par

您现今即将阅读许多我所写的东西，若您出于怜悯，秉持公正，纠正和谴责任何令您不悦之处，您的爱将更受我的感激与珍视。因为您不是那以油傅我的头而令我惧怕的人。\par

弟兄们——不仅是那些与我们同住的人，以及那些居住别处以同样方式事奉天主的人，而是几乎一切在基督内成为我们挚友的人——都向您致意，并向您表达他们的敬意与爱慕之情，视您为弟兄、为圣人、为人。我不敢请求；但若您从教会职务中有片刻闲暇，您可目睹全\Place{非洲}与我本人所渴求的恩惠为何。\par

\SourceRecord{Translated by J.G. Cunningham.；Nicene and Post-Nicene Fathers, First Series；Vol. 1.；Edited by Philip Schaff.；Buffalo, NY: Christian Literature Publishing Co.,；1887.；<http://www.newadvent.org/fathers/1102027.htm>.}

% structure: p0001 source=h1 target=section reason=page-title

\section{\Person{奥古斯丁}致\Person{热罗尼莫}书（公元394或395年）}

\Emph{恭致\Person{热罗尼莫}，我至爱之\OriginalTerm{主}{Lord}、兄弟并同侪司铎，堪受至诚孝爱之敬仰与拥抱者，\Person{奥古斯丁}敬问安好。}\par

% structure: p0003 source=h2 target=subsection reason=relative-heading-rank; base=h2

\subsection{第一章}

1. 从未有人的面貌像你在主内那安宁、喜乐且真正崇高的研习热忱一般，对我变得如此熟悉。因为我虽极愿与你相识，却感到在我对你的了解中，除你极少的一部分——即你的容貌——之外，已一无所缺；而即便对此，我亦不能否认，自从我最蒙福的兄弟\Person{亚力庇乌斯}（如今已荣膺主教之职，彼时他确堪当此任）见过你，并在归来后由我重见，他关于你的讲述已几乎将你的形象全然印刻在我心中；我甚且可以说，在他归来之前，当他还在那里见到你时，我已然借着他的眼睛看见了你。因为任何认识我们的人都会说，他与我仅在肉身而非心灵上是两个个体，我们心灵契合之深、亲密友谊之坚，如此之甚；虽然论及德能，我们并不相等，因他远超于我。因此，既然你爱我——既藉由我们彼此相连的心神契合，由来已久，又藉由你从我友人口中得知的关于我的近况——我便觉得，将兄弟\Person{普罗斯佩鲁斯}举荐给你的兄弟之情，并非冒昧（若是我与你素不相识，此举便是冒昧了）；我们深信，他名字的吉兆（\OriginalTerm{快行}{Good-speed}）必将藉我们的努力，并此后蒙你鼎力襄助，而得应验；虽然，或许正因他如此伟大，此举反而显得冒昧，因为由他向你举荐我，远比由我举荐他更为合宜。若我甘愿仅限于一封正式介绍信的格式，或许我便不该再多写；然而，我的心却洋溢著与你交谈的渴望，切望论及我们在主耶稣基督内所从事的研习，祂乐意藉你的爱厚赐我们诸多裨益，并在祂为追随者所指明的道路上，赐予某些沿途的助佑。\par

% structure: p0005 source=h2 target=subsection reason=relative-heading-rank; base=h2

\subsection{第二章}

2. 因此，我们，连同所有在非洲教会中献身于研究的同道，恳请你不要拒绝倾注心力，将那些以希腊文写成、对我们圣经作出最精辟注释者的著作译出。这样，你也可使我们得享这些人的丰盛，尤其那位你似乎乐于在著述中提及其名的\Person{奥力振}。但我恳请你，在将圣经正典经卷译为拉丁文时，除非遵循你翻译\BibleBook{约伯}{书}时所采的方法——即附加注释，让人清楚看到你的译本与极其尊威的\OriginalTerm{七十贤士译本}{LXX}之间的差异，否则切勿枉费心力。就我而言，我无法充分表达我的惊讶：何以时至今日，竟在希伯来手抄本中仍发现某些为众多精通该语言的译者所遗漏的东西。至于\OriginalTerm{七十贤士译本}{LXX}，他们心灵与精神的契合，远超一己独见，对于其权威，我绝不敢以任何方式妄下论断，唯独坚信，在此翻译之工上，无疑应赋予他们极高的尊重。那些在\OriginalTerm{七十贤士译本}{LXX}完成之后进行翻译的译者，据称既精通希伯来词义与句法，却不仅彼此间未能一致，而且留下许多历经如此长时间仍待发掘和阐明之处，这更令我困惑。如今，这些经文或晦涩，或明晰：若为晦涩，则人们相信，你与前人一样可能有误；若为明晰，则人们不会相信他们（\OriginalTerm{七十贤士译本}{LXX}的译者）会有误失。既已陈述我困惑之由，我恳请你的仁爱，就此问题赐我答复。\par

% structure: p0007 source=h2 target=subsection reason=relative-heading-rank; base=h2

\subsection{第三章}

3. 我亦在读一些归于你名下的、论及宗徒\Person{保禄}书信的著述。在读你对\BibleBook{迦拉达}{书}的注释时，我读到一段经文，其中宗徒\Person{伯多禄}因陷于一种危险的佯装而受规劝。在那里，竟然为谎言辩护——无论辩护者是如你这般鼎鼎大名的人物，还是其他作者（若该文出自他人之手）——我不得不坦承，这令我极其痛心，至少，在那些决定我对此事看法的论点被驳倒之前，我将一直如此，如果它们确能被驳倒的话。因为在我看来，若我们相信圣经中存有何种虚假，势将引出最具灾难性的后果：也就是说，这意味着那些将圣经传授给我们并笔之于书的人，在这书中记载了虚假之事。一个善人是否有时有责任行欺骗，这是一个问题；但一位圣经作者是否有责任行欺骗，则是另一个问题：不，这根本不是另一个问题——这根本不成其为问题。因为，一旦你允许在此至高的权威圣所中，存有一条以履行责任为由的虚妄陈述，那么，那些书卷中所记载的每一句话——若有人觉得难以遵行或难以相信——都可能依循同样致命的原则，被解释为作者出于责任感而有意说出不实之言。\par

4. 因为如果宗徒\Person{保禄}在指责宗徒\Person{伯多禄}时没有说真话，他说：\BibleQuote{你既是犹太人，却按照外邦人的方式，而不按照犹太人的方式生活，你为什么强迫外邦人按照犹太人的方式生活呢？}——如果伯多禄在他眼中似乎做了正确的事，而且，尽管如此，他为了安抚那些惹事生非的反对者，既说又写伯多禄做了错事；\BibleRef{迦 2:11}——如果我们这样说，那么，当如他本人所预言的那些败坏的人兴起，禁止婚姻时\BibleRef{弟前 4:3}，如果这些人自辩说：在确认婚姻合法性的所有文字中\BibleRef{格前 7:10}，这位宗徒是考虑到那些可能因爱妻子而成为讨厌的反对者的人，而说假话——这当然不是说因为他相信这些，而是为了转移他们的反对，才说这些——我们就如何作答呢？无需引用许多类似的例子。因为即使是那些有关赞美天主的事，也可以被描绘为怀着虔敬意图说出的谎言，为的是在那些迟钝的人心中燃起对天主的爱；这样，圣经里就再没有一处纯正真理的权威能站得住脚了。我们难道没有看到这位宗徒以何等大的谨慎向我们推举真理吗？他说：\BibleQuote{如果基督没有复活，那么我们的宣讲便是空的，你们的信德也是空的；而且我们也被证明是天主的假证人，因为我们为天主作证说，天主使基督复活了；如果死人真不复活，他也就没有使基督复活。}\BibleRef{格前 15:14}如果有人对他说：\Quote{你为什么因这假话而如此震惊？你所说的即使是假的，也极大地有助于天主的荣耀。}他难道不会憎恶这人的疯狂，用尽一切言辞和表情，敞开心扉的隐秘深处，抗议说：为天主而说的假话唱赞歌，这罪过不亚于，甚至可能大于，毁谤有关天主的真理？因此，为了我们能认识圣经，我们必须小心确保只接受这样的向导：他对圣经充满崇高的敬意，并且深信其真理，从而不致在圣经任何部分自欺，认为某个陈述之所以被说出，不是因为它真实，而是因为它有益；这样的人宁可跳过他不理解的地方，也不愿将自己的感觉凌驾于真理之上。因为，当他说任何事不真实时，他就是在要求别人优先相信他，并且试图动摇我们对圣经权威的信赖。\par

5. 至于我，我愿将主所赐给我的一切力量都用来证明：那些通常被引用来为谎言的合宜性辩护的经文，每一处都应当以别的方式理解，好使这些经文本身的可靠真理在一切地方都能得到一致的维护。因为作为证据提出的陈述，绝不能是虚假的，也不能支持虚假。然而，我把这事留给你的判断。因为如果你更加仔细地查考这章节，或许你会比我更容易看清楚。使你进行这等更仔细研究的，正是那份虔敬之心，由此你分辨出：如果一旦承认——那些将这一切传授给我们的人，可能出于某种责任的考量，在他们的著作中写下一些不真实的事——那么圣经的权威就会动摇（结果每个人都可以相信自己想信的，拒绝自己不想信的）；除非，也许你意在为我们提供某些规则，借以知晓何时谎言可能或不可能成为一项责任。如果你能做到这一点，我恳请你不吝阐明这些规则，并附上既不含糊也不武断的推理；我以我们的主、真理是在他内\OriginalTerm{降生成人}{incarnate}的那位恳求你，不要认为我提出这个请求是冒昧或放肆。因为如果我的错误是为了维护真理，那就不应受到严厉的责备——如果它还值得责备的话——既然你可以为了虚假而利用真理，却不犯错误。\par

% structure: p0011 source=h2 target=subsection reason=relative-heading-rank; base=h2

\subsection{第四章}

6. 还有许多其他事情，我极愿与你那天真纯朴的心彼此交谈，并就基督信仰的研究向你请教；但这个愿望无法在任何信函的篇幅内得到满足。靠着带这封信的弟兄，我或许能更充分地实现这愿望，我因他能去分享你甘饴而有益的交谈并从中获益而感到欣喜。然而，尽管我不认为自己在哪方面比他强，但他从你那里得到的，甚至可能比我希望的还要少；他会原谅我这样说，因为我承认自己比他有更多的空间从你那里领受。我看他的心智已经更为充实，这一点上他无疑胜过我。因此，当他——如我信赖天主恩佑他平安返回——归回时，当我分享他那因你而变得丰富充盈的内心所得时，我仍会感到一种未得满足的空虚，渴望与你当面相交。由此，无论那时还是现在，我们两人中，我将会是较贫穷的，而他则较富足。这位弟兄带去我的一些作品，若你肯屈尊阅读，我恳请你以坦诚和弟兄般的严格来审阅它们。因为圣经上的话：\BibleQuote{义人应当以慈爱指正我，责备我；但不要让罪人的油膏抹我的头，}我理解为那用批评治愈我的朋友，比那用谄媚膏抹我头的人更为真诚。当我重读自己所写的东西时，我极难做出正确的判断，因我若非过于谨慎，便是过于草率。因为有时我看到自己的错误，但我更愿意听到那些比我更能判断的人指正；以免在我或许正当地责备自己错误之后，又开始自我奉承，以为我的自责是出于对自身判断的不当怀疑。\par

\SourceRecord{Translated by J.G. Cunningham.；Nicene and Post-Nicene Fathers, First Series；Vol. 1.；Edited by Philip Schaff.；Buffalo, NY: Christian Literature Publishing Co.,；1887.；<http://www.newadvent.org/fathers/1102028.htm>.}

% structure: p0001 source=h1 target=section reason=page-title

\section{第29封信（公元395年）}

\Emph{\Place{希波}堂区的长老致\Place{塔加斯特}主教\Person{亚吕皮乌}的一封信，关于前\Place{希波}主教\Person{莱昂提乌斯}的诞辰纪念。}\par

1. 由于\Person{玛加略}弟兄不在，我无法就一件令我焦虑的事情写出任何确定的意见：据说他很快就会回来，靠着天主的帮助，这件事上能做的，都必去做。虽然那些曾与我们在一起的、你们城的弟兄们回去时，也许已足以使你们确信我们为他们所怀的关切，但主赐予我的这件事，是值得作为我们彼此间极大安慰的书信往来的主题的；而且，我相信在获得此事的过程中，你对此事的关切曾大大帮助了我，因为你的关切不能不促使你为我们代祷。\par

2. 所以，让我不要忘记将所发生的事告知你的爱德；这样，正如在这项仁慈尚未获赐之前，你曾与我们一同倾注祈祷，如今在领受之后，你也可与我们一同献上感谢。当我得知你们离开后，有人开始公开激烈反对，并宣称他们不能服从（你在此地时曾暗示的）对那个他们称之为\OriginalTerm{欢庆节}{Lætitia}的节日的禁令，他们试图以一种美好的名称来掩饰他们的欢宴，那时，恰巧最适合我的是，由于全能天主隐秘的预定，在第四圣日，那\par

福音的篇章依序在讲道中被讲解，其中记载着这些话：\BibleQuote{不要把圣物给狗，也不要把你们的珍珠扔在猪前}\BibleRef{玛 7:6}。于是我论及狗和猪的方式，迫使那些对神圣规诫顽固咆哮并被弃于肉欲可憎之事的人羞愧脸红；接着我继续说，他们可以清楚看见，若他们在自己家中行那些事，他们必会被禁止领受圣物，以及教会的珍珠般的特权，那么，在宗教的名义下，在教堂的围墙内做这些事，是何等邪恶。\par

3. 尽管这些话很受欢迎，然而因出席聚会的人很少，对于如此大的紧急情况，该做的并未全部完成。不过，当这篇讲道按照每个人的能力和热忱，被听到的人传播开来时，它遇到了许多反对者。但当四旬期的早晨来到，一大群人聚集在讲解圣经的时刻，那一段福音被宣读，其中我们的主论到那些被赶出圣殿的买卖人，以及祂推倒兑换银钱之人的桌子，说\BibleQuote{祂父的家本应是祈祷之所，竟成了贼窝}\BibleRef{玛 21:12}。在提出无节制饮酒的话题以唤醒他们的注意力之后，我自己也宣读了这一章，并加了一个论证，以证明我们的主会以何等更大的愤怒和暴力，将那处处都可耻的醉酒狂欢，从这圣殿中赶出去；祂尚且从这里赶走了那些在其他地方是合法的买卖，尤其那些出售的乃是在那旧约体制下所规定的献祭所需之物；而且我问他们，他们认为一个被出售必需品的人占据的地方，与一个被酗酒无度的人占据的地方，哪一个更像贼窝。\par

4. 此外，因我预备好的圣经章节已准备好被递到我手中，我接着说，犹太民族，尽管在宗教上全无精神性，但他们在圣殿中从未举行过宴乐，哪怕是节制的宴乐，更不用说因无节制而蒙羞的宴乐了，在那时，主的体血尚未被奉献；而且在历史上，他们从未在任何公开的崇拜场合因酒而兴奋，除非是在他们为自己所造的偶像面前举行宴乐的时候\BibleRef{出 32:6}。我说这话时，从助手那里拿过手稿，宣读了那整段经文。我提醒他们宗徒的话，他为了区分基督徒和顽固的犹太人，说\BibleQuote{他们是他的书信，不是写在石版上，而是写在血肉的心版上}\BibleRef{格后 3:3}。我又以最深切的悲痛问道：虽然天主之仆\Person{梅瑟}因这些以色列的首领的罪而摔碎了那石版两块，我为何竟不能打碎这些人的心呢？这些人虽属新约时代，却在庆祝圣人瞻礼时，屡次公然重复旧约时代的子民在拜偶像时才犯过一次的放纵行为。\par

5. 于是，我交还了\BibleBook{出谷}{纪}的手稿后，在时间许可范围内，继续详论醉酒之罪，并拿起宗徒\Person{保禄}的著作，指出醉酒被他列入哪些罪过之中，同时诵读经文：\BibleQuote{若是称为弟兄的，是淫荡的，或贪婪的，或拜偶像的，或辱骂人的，或酗酒的，或勒索人的，就不同这样的人吃饭。}\BibleRef{格前 5:11} 我语重心长地提醒他们，甚至与那些在自己家中就沉湎于无节制的人同食，我们该何等危险。我又诵读了同一书信中稍后所增添的内容：\BibleQuote{你们不要自欺：无论是淫荡的、或拜偶像的、或犯奸淫的、或作娈童的、或好男色的、或偷窃的、或贪婪的、或酗酒的、或辱骂人的、或勒索人的，都不能承继天主的国。你们中从前也有这样的人，但是你们因着我们的主耶稣基督之名，并因我们天主的圣神，已经洗净了，已经圣化了，已经成了义人。}\BibleRef{格前 6:9} 读完这些后，我嘱咐他们思考：信友若仍在自己心中——即在天主的内在圣殿中——容忍这些连天国都为此关闭的恶欲之可憎事物，他们如何能聆听\BibleQuote{但是你们已经洗净了}这句话？然后我继续讲到那一段：\BibleQuote{你们聚集在一处，并不是为吃主的晚餐，因为你们吃的时候，各人先吃自己的晚餐，甚至有的饥饿，有的却醉饱了。难道你们没有家可以吃喝吗？或是你们想鄙视天主的教会？}\BibleRef{格前 11:20} 诵读之后，我更特别恳请他们注意，连清白节制的宴席在教会内也不准许：因为宗徒并没有说：\Quote{难道你们没有自己的家可去醉酒吗？}——仿佛只有醉酒在教会内才算非法；而是说：\Quote{难道你们没有家可以吃喝吗？}——这些事本身是合法的，但在教会内却不合法，因为人们可以回到自己家中通过必要的食物恢复精力。然而，由于时代败坏、道德松弛，我们竟堕落至此，以致我们不再强求人们在家中保持节制，只敢要求他们，将被容忍的放纵限制在他们自己的家里。\par

6. 我也提醒他们记住我前一天所讲解的福音中的一段，论及假先知说：\BibleQuote{你们可凭他们的果实辨别他们。}\BibleRef{玛 7:16} 我还嘱咐他们记住，在那处，\Emph{果实}一词是指我们的行为。接着，我问他们醉酒被列在哪一类果实中，并诵读了\BibleBook{迦拉达}{书}中的那段经文：\BibleQuote{本性私欲的作为是显而易见的：即淫乱、不洁、放荡、崇拜偶像、施行邪术、仇恨、竞争、嫉妒、忿怒、争吵、不睦、分党、妒恨、凶杀、醉酒、宴乐，以及与这些相类似的事。我以前劝戒过你们，如今再说一次：做这种事的人，决不能承受天主的国。}\BibleRef{迦 5:19} 讲完这些话后，我问：当天主命定基督徒应凭他们的果实而被认出时，我们如何能以醉酒这种果实而被认作基督徒呢？我还加上说，我们必须诵读接下来的一段：\BibleQuote{然而圣神的效果却是：仁爱、喜乐、平安、忍耐、良善、温和、忠信、柔和、节制。}\BibleRef{迦 5:22} 我恳求他们想想，如果他们不满足于在家中行这些本性私欲的作为，反而想借着这些作为——说真的——来光荣教会；如果允许的话，竟想使如此广大的敬拜场所，挤满纵饮醉汉的人群；然而，他们却不愿把圣神的果实——那是圣经的权威和我的哀叹所呼召他们去结出的——呈献给天主，而那本是最适宜用来庆祝圣人节日的奉献。这是多么可耻又可悲！\par

7. 讲完这些，我交还了手稿；人们请我继续发言，我便竭尽全力——因为这危险迫使我如此，而主也乐意赐我力量——把托付给我照顾的他们与我共同面临危险的一幕放在他们眼前，我自己必须向至高牧者交代。我以祂的屈尊就卑、祂所忍受的无与伦比的凌辱、掌击和唾面，以祂被刺透的双手和茨冠，以及祂的十字架和宝血，恳求他们：即便他们对自己不满，至少也要怜悯我，并默想那德高望重的长者\Person{瓦莱里乌斯}对我所怀的不可言喻的爱；他为了他们，义无反顾地把向他们宣讲真理之言的危险重担交给了我，且常向他们说，我的到来使他的祈祷得到了应允。他欢欣，当然不是因为我来分担或目睹听众的丧亡，而是因为我能来与他一起为永生而辛劳。最后，我告诉他们，我已决意信赖那不能说谎的、并借着先知的口向我们赐下应许的一位；论到我们的主耶稣基督，祂说：\BibleQuote{倘若他的子孙离弃我的律法，不行我的道路；倘若他们违犯我的规章，不守我的诫命；我必用棍杖惩罚他们的过犯，用鞭笞责罚他们的罪戾，但我绝不由他撤去我的仁慈。}因此，我宣告，我信赖祂；如果他们轻视此时已读给他们听、已对他们宣讲的这些沉重话语，祂必用棍杖与鞭笞惩罚他们，而不会任凭他们与世界一同受罚。在这场恳求中，我倾尽了我们的救主和主宰在这危急万分的时刻乐意赐予的一切心思与言语的力量。我并没有先自己哭泣来诱使他们流泪；但在我说这些话的时候，我承认，被他们开始流下的眼泪所感动，我自己也禁不住效法他们。当我们如此同哭之后，我满心相信，这篇讲道必能约束他们，使他们远离所谴责的恶习，便结束了讲道。\par

然而，次日清晨，当那许多人惯于放纵吃喝的日子天亮时，我得到消息说，有些人，甚至包括那些在我讲道时也在场的人，仍未停止抱怨，可憎的习俗对他们有极大的控制力，他们不用别的论据，只问，\Quote{为什么\Emph{现在}要禁止这习俗呢？难道从前那些不干涉这种做法的人不是基督徒吗？}听到这些话，我不知道还能想出什么更有力的方法影响他们；但我决定，如果他们仍然执意坚持，就仿效\BibleBook{厄则克耳}{先知书}中所说的，守望者若发出警告，即使被警告的人拒绝听从，他已救了自己的灵魂，然后我要抖掉衣服离去。但后来主指示我，他并不撇下我们独自一人，并教导我如何鼓励我们信赖他；因为在我必须登上讲道台之前，那些我听说因干涉这长久以来的习俗而抱怨的人来见我。我友善地接待了他们，只用几句话就使他们回心转意；到了我讲道的时候，我便搁置了已准备好的讲稿，因为不再需要它，就关于上述问题只说了一点，\Quote{为什么\Emph{现在}要禁止这习俗呢？}我说，对可能提出这个问题的人，最简短也最恰当的回答是这样：\Quote{现在，我们终于该把早该禁止的事禁止了。}\par

然而，为了不让人以为我们轻视那些在我们之前，或是容忍或是不敢制止这放荡群氓如此明显放纵的人，我向他们解释这习俗似乎是如何必然在教会中产生的：原来，在经历了那么多次重大的迫害之后到来的和平时期，大批希望接受基督宗教的异教徒却步不前，因为他们惯于以狂欢酗酒来庆祝其偶像崇拜的节日，很难戒除这种既有害又习惯成自然的享乐；因此，我们的先人觉得，对当时的软弱作出让步，允许他们以庆祝圣殉道者的其他节日来取代他们弃绝的那些节日，这些庆祝活动不是像从前那样带着亵渎的意图，而是同样沉湎于自我放纵。我还说，现在，他们既因基督之名而连结一起，且顺服于他威严的权柄之下，就必须接受节制的有益约束——这种约束应激发他们对他应有的敬畏之心，从而顺服——因此，现在时候已到，所有不敢背弃基督名分的人，都该开始按基督的旨意行事；如今既已作了坚定的基督徒，就该弃绝那些只是为使他们成为这样的人而暂时容许软弱的让步。\par

我接着劝他们效法海外教会的榜样；在那些教会中，有些从未容忍过这类做法，另一些则已经因教友们听从良好教会领袖的劝告而将其废止。当人们拿真福宗徒\Person{伯多禄}的教堂中日日酗酒过度的事例来为这种做法辩护时，我首先指出，我听说这些放纵行为已屡次被禁，但因这地方远离主教的管辖，而且在这座城市里，属肉性的人为数众多，尤其是外方人，他们不断涌入，以与无知成正比的固执墨守这种做法，所以还不可能消除如此严重的弊端。但我继续说，如果我们真要尊敬宗徒\Person{伯多禄}，就必须听从他的话，比起我们不了解他的敬拜之处，更应重视那些让我们明白他心思的书信；于是我立刻拿起手稿，宣读他本人的话：\BibleQuote{基督既然在肉身上受了苦难，你们就应该以同样的心志装备自己，因为在肉身上受苦难的，就与罪恶断绝了关系；今后不再顺从人的情欲，而只随从天主的旨意，在肉身内度其余的时日。过去的时候，你们实行外教人的欲望，生活在放荡、情欲、酗酒、宴乐、狂饮和违法的偶像崇拜中，这已经够了！}\BibleRef{伯前 4:1} 之后，当我看到所有人都一致回心转意，放弃了我所反对的习俗时，我便劝他们在正午聚集，聆听天主圣言，咏唱圣咏；我说明我们已决定以更合乎纯洁与虔诚的方式来庆祝节日；而且，从为此聚集的敬拜者人数，就能清楚看出谁是听从理智的，谁又是情欲的奴隶。我就以这些话结束了讲道。\par

十一、下午聚集的人数比上午更多，人们轮流读经和唱赞美诗，直到我和主教一同出去；我们到来后，又读了两篇圣咏。然后那位长者\Person{瓦莱里乌斯}以明确的命令迫使我向民众讲几句话；我本想推辞，因为我渴望结束这一天的焦虑。我作了一篇简短的讲话，以表达我们对天主的感谢。当我们听到宴乐的喧哗——那照常在进行中的异端者的教会里的宴乐——他们还继续欢宴，而我们却从事着截然不同的事情时，我评论说，白日的美丽因黑夜的对比而更显明亮，当黑色近旁时，洁白的纯净更令人愉悦；同样，我们灵性宴席的聚会或许对我们而言，若非因对比那些人放纵肉欲的狂欢，就不那么甘甜了；我且劝勉他们，若已尝过主的恩慈，就当热切渴望我们当时所享受的这种宴席。同时，我说，那些以终将毁灭之物为他们追求的主要目标的人，理应惧怕，因为每一个人都分享他所崇拜之物的命运；宗徒对此曾有明确的警告，当他论及这些人时说，他们的\BibleQuote{神是肚腹}\BibleRef{斐 3:19}，正如他另一处所说的：\BibleQuote{食物是为肚腹，肚腹是为食物；但天主要把这两样都废掉。}\BibleRef{格前 6:13}。我又加上，我们有责任追求那不朽坏的，那远避肉欲、藉着圣神的圣化而获得的；当我就此主题，按当时所需说了主乐意提示我的这些话之后，便举行了每日的晚祷礼仪；当我和主教从教会退出时，弟兄们在那里唱了一首赞美诗，相当多的群众仍留在教会里，赞美天主直到天黑。\par

十二、我已尽我所能简洁地叙述了我确信你们渴望听到的事。请祈求天主乐意保护我们，使我们不致在任何方面得罪人或招致怨恨。当你们在宁静的昌盛中，我们便以热烈的友爱之情与之大大有份，因我们常频繁地听闻你们那灵性高超的\Place{塔加斯特}教会所拥有的恩赐。载着我们弟兄们的船尚未抵达。在\Place{哈斯纳}，我们的弟兄\Person{阿尔根提乌斯}担任司铎的地方，\OriginalTerm{割臂派}{Circumcelliones}闯入我们的教会，摧毁了祭台。此案正在审理中；我们恳切请求你们祈祷，使之能以和平方式并合乎天主教会的体统而得以解决，从而平息那些骚动之异端者的口舌。我已致信\OriginalTerm{亚细亚执政官}{Asiarch}。\par

至为有福的弟兄们，愿你们在主内坚定不移，并记念我们。阿们。\par

\SourceRecord{Translated by J.G. Cunningham.；Nicene and Post-Nicene Fathers, First Series；Vol. 1.；Edited by Philip Schaff.；Buffalo, NY: Christian Literature Publishing Co.,；1887.；<http://www.newadvent.org/fathers/1102029.htm>.}

% structure: p0001 source=h1 target=section reason=page-title

\section{第30封信（公元396年）}

\EditorialNote{此信由\Person{保利努斯}写于收到对其前一封信（第27号，第248页）的回复之前。}\par

\Emph{致我们的主、神圣而亲爱的兄弟\Person{奥古斯丁}，罪人\Person{保利努斯}和\Person{泰拉西亚}致以问候。}\par

1. 我亲爱的在主基督内的兄弟，因着你神圣而虔诚的事工，我在你不知道的情况下认识了你，并在很久以前就看见了不在场的你，我的心灵以全部的感情拥抱你，并急于通过熟悉的兄弟书信交流来获得听到你的满足。我也相信，靠着主的手和恩惠，我的信已送达你处；但是，我们在冬前派去向你和其他同样在天主之名内所爱的人致意的青年尚未返回，我们不能再推迟我们感到应尽的义务，或抑制我们渴望听到你消息的强烈愿望。因此，如果我前一封信有幸到达你手，那么这是第二封；然而，如果它不幸未能到你手中，请接受此信为第一封。\par

2. 但是，我的兄弟，作为一个属神的人判断一切事，不要凭我们所尽的义务或我们通信的频繁程度来估量我们对你的爱。因为那在各处、作为同一的主、将自己的爱运行在他自己之人内的主，可以作证：从我们借着可敬的主教\Person{奥勒留}和\Person{阿利比乌斯}的盛情，通过你反对摩尼教徒的著作而认识你的那时候起，我们对你的爱在我们内占据了如此的位置，以至于我们似乎不是在获得一份新的友谊，而是在恢复一份旧日的感情。如今我们终于以书面形式向你说话；虽然我们在表达上是新手，但在对你的爱上我们不是新手；借着那内在之人的灵的共融，我们好像与你相熟。虽然相隔遥远，我们却彼此亲近；虽然从未相识，我们却彼此熟知——这并不奇怪；因为我们是同一身体的肢体，拥有同一的头，享受同一恩宠的浇灌，以同一食粮为生，行走在同一道路上，居住在同一家园中。总之，在构成我们存在的一切之中——在我们立于今生、或为来世劳作的整个信德和望德之中——我们都在基督的灵和身体内如此合一，以至于如果我们从这合一中堕落，我们就会不复存在。\par

3. 因此，我们身体上的分离所拒绝给我们的，是何等微不足道的事！——因为那无非是那些取悦眼目、只关注暂时之事的果实之一。然而，也许我们不应将我们在身体中所享受的这份乐趣，列入那些只在时间内才属于属神之人的福分中，因为复活将赋予他们的身体不死不灭；而我们，虽然自身不配，却大胆期望借着基督的功绩和天主圣父的仁慈得到它。因此，我祈求天主借我们的主耶稣基督的恩宠，也赐给我们这份恩惠，使我们终能得见你的面。这不但会给我们渴望的心带来巨大的满足，而且还会给我们的心灵带来光照，我们的贫穷也会因你的富足而变得丰富。的确，即使我们不在你身边，你也能赐予我们这一切，尤其是借着我们的儿子\Person{罗玛努斯}和\Person{阿吉利斯}——他们在主内是我们所深爱和至亲的人（我们把他们当作第二个自己托付给你），当他们奉主的名完成那爱的工作后回到我们这里时；我们恳求，在这工作中，他们能特别享有你的爱德的善意。因为你知道，至高者应许给那帮助兄弟的兄弟何等高的赏报。如果你乐意将所赐给你的恩宠中的任何礼物分赐给我，你可以放心地借他们转交；因为，相信我，他们在主内与我们同心同意。愿天主的恩宠常与你同在，如同现在一样，哦，在主基督内蒙爱、可敬、至亲、令人渴慕的兄弟！请代我们问候所有在你那里、在基督内的圣人，因为无疑这样的人会依附于你的团体；请把我们托付给他们所有人，好使他们能与你一起，在祈祷中记念我们。\par

\SourceRecord{Translated by J.G. Cunningham.；Nicene and Post-Nicene Fathers, First Series；Vol. 1.；Edited by Philip Schaff.；Buffalo, NY: Christian Literature Publishing Co.,；1887.；<http://www.newadvent.org/fathers/1102030.htm>.}

% structure: p0001 source=h1 target=section reason=page-title

\section{第三十一封信（公元三九六年）}

\Emph{致\Person{保利诺}弟兄与\Person{特蕾西亚}姊妹，最亲爱的、真诚的，因天主赋予他们极其丰厚的恩宠而真正蒙至福与卓越者，\Person{奥古斯丁}在主内致候。}\par

1. 尽管我渴望在某种意义上迅即靠近你们，在另一种意义上却仍是远离，我非常希望我对你们前一封信的回信（如果我的任何信配得上被称为对你们的回信的话）能以尽可能快的速度送达你们恩宠；然而这延误却给我带来了你们第二封信的益处。主是美善的，祂常扣下我们所渴望的，为能加给我们更喜爱的事物。因为对我来说，你们收到我的信后给我写信是一种喜乐，而你们因未及时收到信而现在写信又是另一种喜乐。如果我这封致你们圣德的信如我所愿并恳切希冀地快速传递给你们，那我就会失去阅读这封信时所感受到的喜悦。而现在，拥有这封信，并期待我自己的信得到回复，倍增了我的满足。此次延误不能归咎于我；主以其更丰厚的仁慈，成就了祂判定更有利于我幸福的事。\par

2. 我们怀着极大的喜乐在主内迎接了圣洁的弟兄\Person{罗曼努斯}和\Person{阿吉利斯}，他们可以说是你们写来的另一封信，能聆听并回应我们的声音；通过他们，我们可悦地部分享受了你们的临在，尽管这只会使我们更切望见到你们。无论何时何地，你们都不可能通过信件向我们提供关于你们自己的那么多信息，如同我们从他们口中得到的一样多，而我们也不应当要求如此。在他们身上也明显地展示出（这是任何纸张所无法传达的）那种讲述你们时所感到的喜乐，以致于从他们说话时的面容和眼神中，我们便能以一种难以言喻的喜悦读出你们写在他们心上的内容。再者，一张信纸，不论它是何种类型，也不论其上所写的内容多么卓越，本身并不能从它所承载的内容获益，尽管展开它可能给他人带来极大的益处；但在阅读你们的这封信时——即这些弟兄的心——当我们与他们交谈时，我们发现，那些被你们写在上面的人所蒙受的真福，显然与他们被你们写上的满全程度相称。因此，为了获得同样的真福，我们通过对你们的一切热切询问，将他们心里所写的转录到我们自己心里。\par

3. 尽管如此，我们还是深表遗憾地同意他们这么快离开我们，即使是回到你们那里。请注意，我恳求你们，那困扰我们的矛盾情感。让他们毫不迟延地离开的义务，因他们渴望服从你们的急切程度而加重；但他们越急切，就越让我们感到你们几乎就在我们中间，因为他们让我们看到你们的感情是何等温柔。因此，我们让他们离开的不情愿，随着他们被允许离开的紧迫性的合理性增加而增加。哦，难以忍受的试炼！若非这种分离终究并未使我们彼此分离——若非我们是\Quote{同一身体的肢体，共有同一个头，享有着同一恩宠的倾注，赖同一饼而生，行于同一道路上，并居于同一家中！}你们认得这话，我想，是引自你们自己的信；我为何不能也使用它们呢？它们何以更像你们的而非我的呢？因为，就它们是真实的而言，它们源于与同一头的相通。且就它们包含着某种特别赋予你们的事物而言，我因此更爱它们，以致它们占据了通过我胸膛的道路，不容任何话由我心传至我舌，除非它们以作为你们之言的优先权先行。我的弟兄姐妹，在主内至圣而可爱的，我们同一身体的肢体，谁能怀疑我们被同一精神所激励呢？除非那些对将我们彼此连结的爱陌生的人。\par

4. 然而，我很想知道，你们是否比我更能耐心且从容地忍受这种身体上的分离。若果真如此，我承认，我不会因此而在你们的这种刚毅中感到任何喜悦，除非或许因其合理性，因为我承认我远不及你们热切渴望我那般配得你们的感情。无论如何，若我发现自己有忍耐你们不在的能力，这会令我不悦，因为它会使我松懈想见你们的努力；还有什么比因忍耐的能力而变得怠惰更荒谬的呢？但我恳求告诉你们的爱德，那使我留在家中的教会职务；因为真福的\Person{瓦莱里}主教（他与我一同问候你们，并以一种你们将从弟兄们那里听到的急切渴望你们），不满足于让我作他的司铎，还坚持加给我与他分担主教职务的更重负担。我不敢拒绝这一职分，因为通过\Person{瓦莱里}的爱德和民众的催促，我确信这是主的旨意，且由于某些类似任命的先例，我无法用其他理由推辞。诚然，基督的轭本身是柔和的，祂的担子是轻松的；\BibleRef{玛 11:30} 然而，由于我的乖谬和软弱，我可能多少会觉得这轭令人烦恼，担子沉重；我无法说出若因你们的来访而获得安慰，我的轭和担子会变得多么轻松轻盈，因为我得知你们的生活更自由，更少此类挂虑。因此，我觉得有理由请求，不，要求并恳请你们屈尊前来\Place{非洲}，这地方对像你们这样的人的渴求，甚至比其众所周知的土壤干旱更为强烈。\par

5. 天主深知我切望你们光临此邦，不仅为悦我心，也不仅为了那些因我，或因公众传言而听闻你们虔诚志向的人；我更是为了那些未曾听闻，或听闻却不信你们虔诚美名的人，他们或许会被迫去爱慕那种美德，因为他们那时不能再无知或怀疑。虽然你们那满有怜悯的善心，其持久与纯洁本是可贵的，但还有更多要求于你们；即：\BibleQuote{你们的光也当在人前照耀，好使他们看见你们的善行，光荣你们在天之父。}\BibleRef{玛 5:16} \Place{加里肋亚}的渔夫们，不仅欢喜地因主的命令捨弃了他们的船和网，而且还宣告他们已经捨弃了一切，跟随了祂。\BibleRef{玛 19:27} 的确，那轻视一切的人，不仅轻视他实际拥有的一切，也轻视他曾渴望拥有的一切。人所渴望拥有的，只有天主的眼目才能看见；实际所拥有的，则连人的眼目也能看见。再者，当我们贪恋世俗的微末之物时，我们对自己已拥有的，不知何故比对自己渴望拥有的，更加执着。因为当初那向主寻求永生之道的富少年，一听到若愿意是成全的，就必须变卖所有，分施给穷人，在天上拥有宝藏，就忧愁地走了——这不正是如福音告诉我们的，因为他拥有许多产业吗？\BibleRef{路 18:22} 因为克制不去取我们所缺乏的，是一回事；割捨那已成为我们自己一部分的，是另一回事：前者有如拒绝食物，后者则如截去肢体。现今，基督徒的爱德满怀惊叹和喜乐，目睹有人为服从基督的福音而欢欣地献上祭献——那正是当初富少年在基督本人的吩咐下，却忧愁地拒绝去做的——这是何等伟大而奇妙的喜乐啊！\par

6. 虽然言语无法表达我内心所构思并竭力想述说的，然而，既然你们凭着自己的明辨和虔诚，认识到这一荣耀并非出于你们，就是说，并非出于人，而是主在你们内的荣耀（因为你们自己极其小心地提防那\OriginalTerm{仇敌}{Adversary}，并且极其虔诚地努力成为基督的学徒，心内良善谦卑；的确，以谦卑持守世财，胜于以骄傲捨弃世财）——我再说，你们既认识到这荣耀并非你们的，而是主的，你们便看出我所说的是多么软弱无力。因为我一直在讲论基督的讚颂，那是超越天使舌头的主题。我们渴望这基督的荣耀能呈现在我们的子民眼前；好使你们，在婚姻的联结中合而为一，能给两性都立下榜样，让人看到如何将骄傲踏在脚下，并满怀希望地追求成全。我想不出你们还能以何等更显著的方式证明你们的善意，莫过于决意让自己的美德被人看见，正如你们热切地追求并保守它一样。\par

7. 我向你们的善意和爱德推荐这孩子\Person{韦图斯蒂努斯}，他的遭遇甚至能引起那些不虔诚者的同情：至于他受苦和离乡的原因，你们将从他本人口中听到。至于他虔诚的志向——即他许愿献身于天主的服事——这志向要等一段时间，待他的力量得到巩固，目前的恐惧消除后，才能更确定地知道。我感觉到你们对我的爱如此热切，并因此相信你们不会吝惜阅读我所写的东西，于是我将三本书寄给你们的尊驾和爱德：但愿这几卷书的篇幅能显示对如此重大主题论述的完整；因为书中讨论的是\OriginalTerm{自由意志}{free-will}问题！我知道这些书，或至少其中一些，我们的弟兄\Person{罗马尼阿努斯}并未持有；但正如我提及的，几乎我所能为任何读者利益而写的全部作品，因你们对我的爱，都可以通过他提供给你们阅读，虽然我并未嘱託他将书带给你们。因为他早已拥有所有这些书，并随身携带：此外，正是通过他，我回覆你们第一封信的回信已经寄出。我想你们的尊驾已凭主赐予的灵性智慧发现，那人在灵魂内承受了多少美善，又因软弱相差多远。在託他转交的信中，我相信你们已读到，我何等恳切地将他和他的儿子託付给你们的同情与爱，以及他们与我所结的纽带何等紧密。愿主藉着你们的手建树他们！这应当向主祈求，而非向你们，因为我知道这早已是你们的渴望。\par

8. 我从弟兄们那里听说，你们正在撰写一篇驳斥\OriginalTerm{异教徒}{Pagans}的论文：如果我们在你心中有任何请求的权利，请立刻将它寄给我们阅读。因为你的心是天主真理的奇妙神谕，我们期待从中得到答案，满意而清晰地解决那些最冗长繁複的辩论。我听说你的尊驾拥有最可敬的教父\Person{盎博罗削}的著作，我极其渴望拜读其中那些他以极大的关注和篇幅，为驳斥某些极其无知而又狂妄之人而写的书，这些人声称我们的主曾受教于\Person{柏拉图}的著作。\par

9. 我们最可敬的弟兄\Person{塞维鲁}，原属我们的团体，现在主持\Place{米莱维斯}教会，并为那城的弟兄们所熟知，他与我一同向你的尊驾致以恭敬的问候。与我一同服事主的弟兄们也热烈地向你问候，他们渴望见到你，如同他们爱你那样深；他们爱你，正如你的卓越善行所应得的。我们寄给你的饼，将通过你的善意所接纳的爱，成为更加丰厚的祝福。愿主永远护佑你们脱离这世代，我最亲爱而真诚的弟兄和姐妹，真正富於善意，并极其卓越地蒙受了主丰厚的恩宠。\par

\SourceRecord{Translated by J.G. Cunningham.；Nicene and Post-Nicene Fathers, First Series；Vol. 1.；Edited by Philip Schaff.；Buffalo, NY: Christian Literature Publishing Co.,；1887.；<http://www.newadvent.org/fathers/1102031.htm>.}

% structure: p0001 source=h1 target=section reason=page-title

\section{第33封信（公元396年）}

\Emph{致\Person{普罗库莱亚努斯}，我的主，尊贵且最亲爱的，\Person{奥思定}致以问候。}\par

1. 那些在此信前冠以的称谓，我无需向你详细辩护或解释，尽管它们可能会冒犯无知之人的虚荣偏见。因为我正确地称你为\Emph{主}，因为我们双方都在寻求使彼此脱离错误，尽管在充分辩论之前，我们中谁陷于错误对某些人来说可能尚不确定；因此，如果我们真诚地努力，为使我们双方都能脱离不和的乖僻，我们便是彼此\Emph{服事}。我以真诚的心，并以基督徒的谦卑的敬畏与战栗致力于此事，这或许不为大多数人所知，但那位洞察所有人心肺的天主却看见。至于我所毫不犹豫地认为你\Emph{尊贵}，你很容易就觉察到。因为我不认为任何由裂教而来的错误配得上任何尊荣，我渴望全人类都脱离它，只要在我力所能及之内；但我自己却一刻也不犹豫地认为你配得尊荣，主要因为你与我在共同人性的纽带中联结，而且在你身上显明了一些较为温和性情的迹象，这鼓励我期望，一旦真理向你显明，你便会欣然接受。至于我对你的\Emph{爱}，我所亏欠的不亚于那位如此爱我们、以至于为了我们承受了十字架的耻辱的主所命令的。\par

2. 然而，请勿惊奇我如此长久地未致信于你的仁善；因为我先前并不认为你的观点如同\Person{埃沃迪乌斯}弟兄满心欢喜地向我宣告的那样，他的见证我不能不信。因为他告诉我，当你们偶然在同一住所相遇，彼此开始谈论我们的希望，也就是基督的嗣业时，你亲切地表示，你愿意在良善之人的见证下与我进行一场会谈。我确实为你屈尊提出此建议而感到欣喜：我决不能错过如此重要的机会，这是你善意所予的，让我运用主所乐意赐给我的任何力量，与你一起思考和辩论，究竟是什么原因、根源或理由，导致在基督的教会中出现了如此可悲可叹的分裂——主曾对那个教会说：\BibleQuote{我将平安赐给你们，我将我的平安留给你们。}\BibleRef{若 14:27}\par

3. 我从上述弟兄那里听闻，你曾抱怨他以近乎凌辱的方式回答了你；但我恳求你，不要视之为凌辱，因为我确信那并非出自傲慢的精神，正如我深知我的弟兄。但是，若他在为自己的信德和教会的爱德进行辩论时，或许言辞稍显激烈，在你看来有些伤及你的尊严，那不应被称为顽固，而是勇敢。因为他渴望辩论并讨论问题，而非仅仅顺服你和奉承你。因为这样的奉承乃是罪人的油，先知不希望自己的头被它所膏抹；因为他曾说：\Quote{义人必以慈爱纠正我，责备我；但罪人的油绝不膏抹我的头。} 因为他宁愿被义人严厉的慈爱纠正，也不愿被奉承的柔和油膏所赞美。因此也有先知的话说：\Quote{那些称你为有福的，反而使你误入迷途。} 所以，论到那些因虚假奉承而变得骄傲的人，常公正地说\Quote{他的头长大了；}因为那是被罪人的油所增大，即不是由严厉真理纠正之人，而是由温顺奉承称赞之人的油。然而，请不要因此认为，我是想说，你已被\Person{埃沃迪乌斯}弟兄如同义人般地纠正了；因为我担心你会认为我此语也带有凌辱的意味，而我正是极度谨防避免此事的。但那位曾说\BibleQuote{我是真理}的，才是义者。\BibleRef{若 14:6} 因此，无论何时，任何真实的话语，即使出自任何人口中，或许稍显粗鲁，我们不是被说话者纠正——他或许和我们一样同为罪人——而是被真理自身纠正，即被那位是义者的基督纠正，以免那柔和却有毒的奉承之油，即罪人的油，膏抹了我们的头。因此，尽管\Person{埃沃迪乌斯}弟兄在维护他所捍卫的共融时，由于过度激动而可能言辞过于激烈，你应因他的年龄和他对此事的重要性估量而谅解他。\par

4. 然而，我恳求你记住你曾乐意应允的事；即与我友好地探讨一件如此重大、且如此涉及共同救恩的事情，在你所选择的听众面前（只要我们的谈论不致于消失无闻，而是用笔记录下来；这样我们便可以更平静有序地进行讨论，并且我们所说的任何可能从记忆中溜走的话，都可以通过阅读记录来回想）。或者，如果你更愿意，我们可以无需任何第三方的介入，通过信件或会面和阅读，在你喜欢的任何地方进行讨论，以免某些不明智地热衷于此的听众，更多地期待我们之间的冲突，而非想到我们的讨论对彼此的益处。然而，在讨论结束后，当通过我们向民众通报结果；或者，如果你更喜欢通过信件交换来讨论，那么让这些信件在双方会众面前宣读，好使他们将来不再分裂，而是合而为一。事实上，我欣然同意你想要的、规定的或偏好的任何条件。至于我最蒙福和可敬的父亲\Person{瓦莱里乌斯}的心意，他此刻不在家中，我满怀信心地担保，他会以极大的喜乐听闻此事；因为我知道他多么热爱和平，而且他丝毫不受任何对空洞尊严炫耀的浅薄顾忌所影响。\par

5.　我问你们，我们与过去世代的纷争有何相干？骄傲之人用苦毒加给我们\OriginalTerm{肢体}{members}的创伤，存留至今，这已经够了；因为时至今日，我们已经不再感到那为消除痛苦通常需寻求医生帮助的那种疼痛。你们看，那破坏基督徒家庭安宁的灾祸，何其大而可悲。夫妇在家里的炉灶边彼此和睦，却在\OriginalTerm{基督的祭台}{altar of Christ}前分裂。他们借着祂彼此立约要和睦相处，却在祂内不能和睦。儿女与亲生父母同住一屋，却不同一天主的殿宇。他们渴望确保从那些因基督的产业而争吵之人获得尘世的产业。仆人与主人分裂了他们共同的主，祂取了\OriginalTerm{奴仆的形体}{form of a servant}，为解救众人脱离奴役。你们一派尊重我们，我们一派也尊重你们。你们的人因我们的主教标志向我们求助，我们的人也向你们表示同样的尊重。我们听取所有人的话，不愿得罪任何人。那么，既然在别处找不出得罪的理由，为什么却在基督内找到，竟然撕裂祂的肢体呢？当有人希望通过我们的帮助了结世俗纷争时，我们能够对他们有益，他们称我们为圣人和天主的仆人，好让我们处置他们的财产问题：现在，让我们终于主动处理一件关系到我们和他们双方\OriginalTerm{救恩}{salvation}的事宜。这事并非关乎金银、土地或牲畜等我们每日都为此被谦恭地致敬、以期我们和平了结纷争的事，而是关乎我们之间的这一分歧，它如此不堪、有害，直接关系到我们的\OriginalTerm{头}{Head}本身。那些向我们致敬、希望我们使他们在此世事务上和睦的人，无论低头多低，我们的头却从天屈尊至十字架，而我们竟不在祂内和睦。\par

6.　我恳求你们，如果你们真有某些人称道的那种弟兄之情，就请以行动表明你们的善良是真诚的，而非为求虚浮的尊荣而假装的：要触动\OriginalTerm{怜悯的心肠}{bowels of compassion}，同意商讨此事；同我一起恒心祈祷，和平地讨论每一点。不要让那些不幸的人民对我们尊严的尊重，在天主的审判中，成为加重我们责罚的因素；反倒要借着我们真诚的爱，带领他们与我们一同脱离错谬和纷争，走上真理与和平的道路。\par

我的主，可敬且至爱者，我祈求你在天主眼中蒙福。\par

\SourceRecord{Translated by J.G. Cunningham.；Nicene and Post-Nicene Fathers, First Series；Vol. 1.；Edited by Philip Schaff.；Buffalo, NY: Christian Literature Publishing Co.,；1887.；<http://www.newadvent.org/fathers/1102033.htm>.}

% structure: p0001 source=h1 target=section reason=page-title

\section{书信三十四（主历三九六年）}

\Emph{致 \Person{欧瑟伯}——我卓越的主内弟兄，堪受爱戴与尊敬，\Person{奥斯定}致候。}\par

1. 天主洞悉人心的隐密；祂知道，正是由于我对基督徒和平的挚爱，才使我为那些人卑鄙而不虔地坚持背离和平的亵渎行为而深感动荡。祂也知道，我这种情感本身即是为了和平，而我渴望的，并非强迫任何人违背意愿地加入\OriginalTerm{天主教会的共融}{Catholic communion}，而是向所有在错误中的人公开宣讲真理，并赖天主的助佑，藉着我的职务将其清晰地显明出来，好使真理以此举荐自身，让他们欣然接受并追随。\par

2. 暂且不论许多其他事情，最近发生的事还有什么是更可恶的呢？一个年轻人因像疯子般经常殴打他的母亲，甚至在那法律对罪大恶极者也施怜悯的日子里，也不克制自己亵渎的双手，去伤害那生育他的人，因而受到他\OriginalTerm{主教}{bishop}的责备。于是他威胁他的母亲，说要转投\OriginalTerm{多纳图斯派}{Donatists}，并要杀死他惯于以难以想象的凶残加以殴打的母亲。他发出这些恐吓，随后便转投了多纳图斯派，在充满邪恶狂怒时\OriginalTerm{重新受洗}{rebaptized}，在燃烧着流他母亲之血的欲望时穿上白衣。他被安置在教堂内栏杆后面一个突出而显眼的位置；在那些悲伤而愤慨的旁观者眼中，这个图谋弑母的人被展示为一个重生的人。\par

3. 我恳求你，作为判断极为成熟的人，这些事能蒙你悦纳吗？我不相信你会如此：我了解你的智慧。母亲被她儿子打伤，伤在那曾生养这忘恩负义之人的身体四肢上；而当教会，他的\OriginalTerm{属灵的母亲}{spiritual mother}，出面干预时，她也在那些\OriginalTerm{圣事}{sacraments}上受了伤，这些圣事正是她曾藉以赐予这同一位忘恩负义的儿子生命与滋养的。难道你不仿佛听见那年轻人为父母的鲜血而咬牙切齿地狂怒，说道：\Quote{我要如何对待那禁止我伤害母亲的教会呢？我已想出办法：让教会本身受她所能够承受的伤害吧；让我接受那能使她的肢体感到痛苦的事。让我投奔那些懂得藐视她赐我\OriginalTerm{属灵的重生}{spiritual birth}的\OriginalTerm{恩宠}{grace}，并损毁我在她母腹中所领受的印记的人。让我以残酷的折磨加害于我的生母与属灵的母亲：让那第二次给我生命的人，先给我带来埋葬；为了她的哀伤，让我寻求属灵的死亡，而为了另一位的死亡，让我延长我肉身的生命。} 欧瑟伯啊！我以一个正直人的身份恳求你，我们现在还能期望什么呢？难道还能期望他作为多纳图斯派，就不会感到自己如此全副武装，以至于毫无畏惧地去攻击那位不幸的、年迈体衰、寡居无依的妇人吗？当他还是一个\OriginalTerm{天主教徒}{Catholic}时，还曾受到约束，不敢伤害她。因为当他向母亲说：\Quote{我要转投多纳图斯派，我要喝你的血}时，他充满激情的心中怀有别的企图吗？看哪，他身穿白衣，良心却因血色而深红，他已部分地实现了他的威胁；剩下的部分，就是喝他母亲的血。因此，如果这些事能蒙你悦纳，那就让他现在的神职人员和祝圣者们催促他，在八天之内完成他誓愿的剩余部分吧。\par

4. 主的手确实强而有力，足以阻止这人的狂暴，使他远离那不幸而无依的寡妇，并以祂的智慧所知道的方法，遏制他这亵渎的图谋；可是，当我的心被这般哀伤刺透时，我又怎能不吐露我的感受呢？他们做出这些事来，难道我还要奉命缄默吗？当祂藉着\OriginalTerm{宗徒}{apostle}的口命令我，说那些教导不当教导之事的人，必须由主教加以斥责\BibleRef{铎 1:9}，难道我还能因害怕他们不快而缄默吗？愿主救我脱离这般的愚昧！至于我渴望将这样一桩不虔的罪行记录在我们的公共档案中，我之所以渴望如此主要是为了这个目的：好叫凡听见我为这些事哀叹的人，尤其在我在其他城镇作此哀叹更为合宜之处，不致认为我是在捏造谎言；更何况，在\Place{希波}本地，人们已经断言，\Person{普罗库莱亚努斯}并未发出官方报告中归咎于他的那道命令。\par

5. 除了通过像您这样一位身居最显赫地位、兼备沉稳与极大的明智及善意之人从中调停之外，我们还能以什么更平和的方式处理这件重要的事呢？因此，正如我先前已借由我派往尊座前的那些善良可敬的弟兄们所请求的，我恳求您屈尊查问：是否\Person{维克托}\OriginalTerm{司铎}{presbyter}确实未从他的主教那里接到官方公开记录所报告的指令；抑或，既然维克托本人另有说法，他们虽与他同属一个共融团体，却在记录中虚捏事实诬陷于他。或者，如果他同意我们平和地讨论我们之间的全部分歧，为使那已经显明的错误更加显露，我乐于接纳此机会。因为我听说他提议，避免民众喧嚷，仅由双方各派十位受人尊敬的人士在场，我们依圣经来查考此事何为真理。至于有人向我报告他提出的另一建议，即我应前往\Place{康斯坦提纳}，因为在那城他的同党人数较多；或我应前往\Place{米莱维斯}，因为据他们说那里即将举行一次会议；——这些事是荒唐的，因我的特别职务只限于\Place{希波}教会。首先，这件事对我的全部重要性，在于它关乎\Person{普洛库莱亚努斯}和我本人；倘若他可能认为不是我对手，就让他恳求任何他愿意的人作为其辩论伙伴施以援助。因为在其他城镇，我们干涉教会事务，只限于那些与我们同享\OriginalTerm{司铎职}{priestly office}的弟兄——即那些城镇的主教们——所允许或命令的程度。\par

6. 然而，我无法理解我一个新手身上有什么，竟令他——一个自称担任主教多年的人——不愿且害怕与我进行讨论。若因我通晓文艺学科，或许他从未修习，或至少不如我研习得多，但这对讨论的问题有何相干？这问题应由圣经或由教会及公开文献来决定，而他多年熟谙这些，应比我更为精通。再者，这里有我的弟兄兼同僚，\Place{图里斯}教会的主教\Person{萨姆苏修斯}，他并未学过任何那些据说他害怕的文科学科：让他代我出席，辩论便在他们之间进行。我将请求他，并且，因我信赖\Person{基督}之名，他必乐意在此事上代我；我也信赖主必在他为真理辩护时施以援助：因为他虽言辞朴拙，却在真信德上受教良深。因此，他并无理由将我推给我不认识的人，而不让我们自行解决我们之间的事。不过，如我所言，如果他本人请求他们协助，我也不会拒绝与他们相见。\par

\SourceRecord{Translated by J.G. Cunningham.；Nicene and Post-Nicene Fathers, First Series；Vol. 1.；Edited by Philip Schaff.；Buffalo, NY: Christian Literature Publishing Co.,；1887.；<http://www.newadvent.org/fathers/1102034.htm>.}

% structure: p0001 source=h1 target=section reason=page-title

\section{第三十五封信（主历三九六年）}

\EditorialNote{另一封致\Person{尤西比乌斯}关于同一主题的信函。}\par

\Signature{\Emph{致\Person{尤西比乌斯}，我卓越的主内兄弟，堪受敬爱与尊重者，\Person{奥古斯丁}谨致问候。}}\par

1. 我并未以强求的劝勉或不顾你推辞的恳求，将你称之为仲裁主教间纠纷的责任强加于你。即使我曾有意推动你承担此任，或许我也不难证明，在这般清晰而简单的事件中你完全有资格在我们之间作出判断；毋宁说，我能指出你其实已经在这样做了，因为你，虽然畏惧审判官的职责，却毫不犹豫地在听取双方陈词之前便作出了偏袒一方的判决。不过，关于这一点，如我所说，此刻暂不多言。因为，除此以外，我从未向你高尚的善意所求过别的——我恳请你在此信中留意这一点，若在前信中未曾留意的话——只求你询问\Person{普罗库莱亚努斯}，他是否亲口对他的司铎\Person{维克托}说过那些经官方派员记录于公共档案中的话，抑或那些被派去的人在公共档案中所写的并非他们从\Person{维克托}听来的实情，而是虚假之辞；此外，还想知道你与他之间就整个问题展开讨论，他的看法如何。我认为，仅仅受一方之请向另一方提问，并俯允写下所获答复之人，并不构成当事双方的仲裁者。如今我再次请求你，不要拒绝做这件事，因为据我经验所知，他不愿收到我的信函，否则我便不会劳烦你的卓越居中调解了。既然他不愿如此，我除了通过你这样一位良善之人、他如此亲密的朋友，来寻求一个关乎我职责之重不容我缄默的问题的答复之外，还有什么更不易得罪人的办法呢？何况，你说（因为儿子殴打母亲的行为受到你健全判断的谴责），\Quote{若\Person{普罗库莱亚努斯}早知此事，他早已禁止那人参与他教派的共融。} 我以一句话回答：\Quote{现在他已知晓，现在就让他禁止那人吧。}\par

2. 让我再提一件事。一个名叫\Person{普里姆斯}的人，曾是\Place{斯帕纳}教会的\OriginalTerm{副执事}{subdeacon}，当他被禁止与修女们进行违反教会法规的交往时，他藐视已定的明令，因而被免去圣职。此人也因受此出于天主之意的惩戒而恼怒，便转入对方教派，并在那里接受了\OriginalTerm{重新受洗}{rebaptized}。同他一起住在天主教会同一村庄里的两名修女，或是被他带去对方教派，或是跟随他而去，同样也接受了重新受洗。如今，置身于\OriginalTerm{游荡派}{Circumcelliones}和一群群无家可归、为逃避约束而拒绝婚姻的妇女之中，他于可憎的放纵狂欢中傲然炫耀，庆幸如今可以毫无阻碍地享有那在天主教会中曾受约束的、最肆无忌惮的恶行。或许\Person{普罗库莱亚努斯}对此事也一无所知。因此，请通过你这位庄重而公正之人，向他陈明此事；让他命令将那人从共融中开除，因为此人之所以选择了他的教派，无非是由于他因不服管束、生活放荡而在天主教会中失去了圣职。\par

3. 就我而言，若主喜悦，我定意遵行这样的准则：凡在他们中间经纪律判决被免职者，若表示愿意转入天主教会，必须接受，条件是需呈明同样的补赎证据，一如他们若留在那里，或许也会被迫呈明的那般。但我恳请你细想，他们对待那些我们因不圣洁生活而以教会惩戒加以约束之人的做法，是何等令人憎恶：他们先劝这些人接受\OriginalTerm{重新受洗}{second baptism}，为使这些人有资格领受，便让他们自称为异教徒（多少殉道者的血已流出，正为免使这样的宣称出自基督徒之口！）；此后，仿佛更新和圣化了一般，实则罪孽愈加深重，竟以新恩宠的伪装，以新的疯狂之亵渎，藐视那他们原本无法顺服的纪律。然而，若我试图通过你善意居间调解以纠正这些弊端而有所错失，莫让人指责我借助公共档案使这些事为\Person{普罗库莱亚努斯}所知——这种通知方式，在这座\Place{罗马城}中，我相信是不能拒绝我的。因为，既然主命令我们去宣讲真理，并在教训中斥责错误，且无论顺境逆境总要勉力，一如我能以主和宗徒们的话语证明的，就莫让人以为可说服我对这些事缄默不言。若他们图谋任何大胆的暴力或暴行，已然将世上一切国度置于其教会怀中的主，其教会广布于整个世界，必不会不护卫她免遭不义。\par

4. 这里教会产业中一个佃农的女儿，她曾是我们的\OriginalTerm{慕道者}{catechumen}之一，违背父母的意愿，被另一方拉走，在他们那里受洗后，立愿当了\OriginalTerm{修女}{nun}。她的父亲想用严厉的手段强迫她回到天主教会；但我不愿我们接纳这位思想如此扭曲的妇女，除非她出于自愿，并在自由判断下选择那更好的。当这个乡下人开始企图用殴打强迫女儿顺服他的权威时，我立刻禁止他使用任何这类手段。尽管如此，后来当我经过\Place{斯帕尼亚地区}时，一位\Person{普罗库莱安}手下的\OriginalTerm{司铎}{presbyter}，站在一位极好的天主教妇女的田里，以极其傲慢的声音向我喊叫，说我是个\OriginalTerm{背教者}{Traditor}和\OriginalTerm{迫害者}{persecutor}；他也对那位妇女——我们教会团体的一员，他正站在她的产业上——发出了同样的责骂。但我听到这话时，不仅没有追究争吵，还拦阻了围在我身边的一大群人。然而，如果我说，让我们调查并确认谁是或曾是真正的背教者和迫害者，他们便回答说：\Quote{我们不要辩论，我们要重新施洗。任凭我们像狼一样，以狡诈残忍的手段掠食你们的羊群吧；如果你们是好牧人，就默默忍受吧。}因为\Person{普罗库莱安}所命令的，除了这个还有什么呢？如果这项命令确实应该归之于他的话：\Quote{如果你是基督徒，就把这事交给天主审判吧；无论我们做什么，你都不要作声。}此外，这位司铎竟敢对一位管理教会庄园的乡下监督发出威胁。\par

5. 我恳请您将这一切告知\Person{普罗库莱安}。请他抑制他手下的神职人员的疯狂；尊敬的\Person{欧瑟伯}，我不得不向您报告这些事。请您写信给我，不要告诉我您对这一切的意见，以免您认为我把我法官的责任加在您身上；而是告诉我他们对我的问题的回答。愿天主的仁慈护佑您免遭凶恶，我最尊贵的大人和弟兄，最值得爱戴与尊敬的人。\par

\SourceRecord{Translated by J.G. Cunningham.；Nicene and Post-Nicene Fathers, First Series；Vol. 1.；Edited by Philip Schaff.；Buffalo, NY: Christian Literature Publishing Co.,；1887.；<http://www.newadvent.org/fathers/1102035.htm>.}

% structure: p0001 source=h1 target=section reason=page-title

\section{第36封信（公元396年）}

\Emph{致我最敬爱且切切怀念的弟兄兼同侪司铎\Person{卡苏拉努斯}，主内\Person{奥斯定}致候。}\par

% structure: p0003 source=h2 target=subsection reason=relative-heading-rank; base=h2

\subsection{第一章}

1. 我不知道为何没有回复您首封来信；但我晓得，这疏忽并非因我对您不存敬意。因为我喜悦您的研学，甚至喜悦您表达思想的言词；我渴望并劝勉您，在早年得以在天主圣言上大有长进，以建树教会。如今收到您的第二封信，您以我们合而为一的手足情谊这一最正义亦最可亲的理由恳请答复，我遂决定不再拖延，满足您爱德所表达的渴望；虽在百忙之中，我仍着手清偿对您的亏欠。\par

2. 关于您希望我发表意见的问题：\Quote{在每周的第七日是否允许守斋}，我回答说：若绝对不许，那么梅瑟、厄里亚，甚至我们的主耶稣本人，就不会连续四十天守斋了。但同样的论证也可证明，主日守斋亦非违法。然而，倘若有人认为主日应当定为斋日，如同某些人遵守第七日一样，这人便会被视为——且不无道理地——给教会带来极大的绊脚石。因为在圣经未立定规的这类事务上，天主子民的习俗，或祖先所立的常规，应被奉为教会的法律。若我们决意为此等事争执，只因一方习俗异于他方便横加指责，其后果必是没完没了的纷争；其间极需谨慎，免得争辩的风暴以乌云遮盖了手足之爱的平静，而我们的力量徒耗于无法提供任何决定性真理证据的纯粹论战。那位作者并未小心避开这危险，他冗长的论述，您认为值得在前信中寄来与我，好让我答复其论据。\par

% structure: p0006 source=h2 target=subsection reason=relative-heading-rank; base=h2

\subsection{第二章}

3. 我眼下没有足够的闲暇逐条驳斥他的意见：我的时间已被其他更重大的工作所占据。但若您稍加用心研读这位匿名罗马作者的大作——您信中证明自己拥有的天赋，我视之为天主恩赐而深爱于您的天赋——便会发现，他竟以最尖刻的言辞，肆无忌惮地诋毁几乎整个基督的教会，从日出之地直到日落之处。岂止几乎，我毋宁说，就是整个教会。因为他甚至没有放过罗马的基督徒，而他自以为在维护他们的习俗；但他不察，他的激烈抨击如何反及己身，因他未曾留意。当他无法提出证明第七日必须守斋的论据时，便厉声咆哮，抨击宴乐的放纵、醉饮的狂欢，以及最可耻的酩酊之习，仿佛斋戒与纵饮之间毫无中庸之道。若然承认这点，星期六守斋于罗马人何益之有？既然根据他的逻辑，在其他不守斋的日子，他们势必被视为贪饕纵酒之辈。因此，若放纵口腹、酗酒醉饮——这常是罪过——与放宽斋戒的严格，而在其他日子恪守节制与平衡，如主日那样不受任何基督徒指摘，这二者果有分别——我若说，这二者果有分别，则他当先区分圣徒的聚餐与那些以肚腹为神者的饕餮，免得他指控罗马人自己在不守斋的日子属于那些\BibleQuote{他们的神就是自己的肚腹}之辈；然后他当研究的，并非在第七日纵醉是否允许——这本在主日亦非允许——而是我们是否有责任在第七日守斋，一如我们在主日向来不守的那样。\par

4. 我愿见他探讨并解答这问题，且不致使自己公然反对那散布于天下、除罗马基督徒及迄今少数几个西方团体外的整个教会。请问，是否应容忍此人，对如此众多、如此杰出的基督仆人，在第七日清醒适度地以食物恢复体力，竟敢说他们\BibleQuote{属于肉身的人，不能悦乐天主}；论及他们时还写着：\BibleQuote{作恶的人，请你们远离我，我不认识你们的路}；又说他们以肚腹为神，宁取犹太礼仪而舍教会礼仪，是为婢女之子；又说他们不受天主正义法律的管辖，而随从自己的私意，只顾私欲，不服从有益的约束；又说他们属肉体，散发着死亡的气息，以及诸如此类的指控。这些话语，若他曾攻击天主的哪怕一位仆人，谁还肯听他？谁不会转身离去？然而，如今他竟以此等斥责辱骂，攻击那在全世界结果增长、几乎处处不在第七日守斋的教会，我便要警告他，不论他是谁，务须小心。因为您不愿向我透露他的名字，这清楚表明您不愿我判断他。\par

% structure: p0009 source=h2 target=subsection reason=relative-heading-rank; base=h2

\subsection{第三章}

5. 他说道：\BibleQuote{人子是安息日的主，在这一天，无论如何，行善总比作恶更合法。}\BibleRef{玛 12:8} 因此，如果我们在破斋时便是作恶，那么就没有任何一个\OriginalTerm{主日}{Lord's day}我们能合乎本分地生活了。至于他承认宗徒们曾在一周的第七日进食，并就此评论说，那时他们\OriginalTerm{禁食}{fasting}的时期还没有来到，因为主亲自说过：\BibleQuote{日子将到，新郎要从他们中被夺去，那时新郎的宾客就要禁食了。}\BibleRef{玛 9:15} 既然有\BibleQuote{欢乐有时，哀悼有时}\BibleRef{训 3:4}，他首先应该注意到，吾主在那里是泛指禁食，而非特指第七日禁食。而且，当他说禁食表示哀恸，进食表示喜乐时，他为何不反思一下，那经上所载\BibleQuote{天主在第七日停止了他所作的一切工程，安息了}——即那安息所彰显的，乃是喜乐而非忧伤——天主究竟要借此指示什么呢？或许，他是想断言，在天主安息并圣化\OriginalTerm{安息日}{Sabbath}时，对犹太人是指示喜乐，对基督徒却是指示哀恸吗？但是，天主无论在起初圣化那一日，因为在那日他安息停止了一切工程时，还是后来他给希伯来民族颁布遵守那日的诫命时，都未曾对一周第七日的禁食或进食立下规条。在那里，对人唯一的要求，就是他自己停止工作，也不要求他的仆人工作。而先前安排\OriginalTerm{先前安排}{former dispensation}中的子民领受了这安息，视之为将来事物的预像，他们遵守了这诫命，所行的正是我们现在看到的犹太人所实践的那种停止劳作；并非如某些人所猜想的那样，是由于他们属血肉，误解了基督徒所正确理解的道理。我们理解这法律，也并不比先知们更好——先知们在此法律仍然有效的时候，所守的安息日安息，正是犹太人相信至今仍当遵守的那种。因此，天主曾命令他们，把在安息日捡柴的人用石头砸死；\BibleRef{户 15:35} 但我们从未读到，有人因为在安息日禁食或进食而被用石头砸死，或被判定应受任何惩罚。这两种行为，哪一种更符合安息，哪一种更属于劳苦，就让我们的作者自己来判断吧——他既把喜乐视为进食者的分，把哀恸视为禁食者的分，或者至少，他认为主在回答有关禁食的问题时，也是这样看待的，因为主说：\BibleQuote{新郎还与宾客在一起的时候，宾客怎能哀恸呢？}\BibleRef{玛 9:15}\par

6. 此外，对于他断言宗徒们在第七日进食（这是长老传统所禁止的），是因为那时他们在那一天禁食的时期还没有来到；我要问：难道那时犹太人停止工作守安息的时期也已经废除了吗？长老的传统岂不一方面禁止禁食，另一方面又命令安息吗？然而，我们读到基督的门徒在安息日进食，而且他们在同一天掐了麦穗，这事当时是不合法的，因为被长老的传统所禁止。因此，让他想一想，是否更有理由这样回答他：主愿意他的门徒在同一天行这两件事——掐麦穗和进食——因为前一件事可以驳斥那些禁止在第七日作任何工的人，后一件事可以驳斥那些命令在第七日禁食的人；因为借着前一件事，他教导说，由于\OriginalTerm{安排}{dispensation}的改变，停止劳作的安息如今已成了迷信行为；借着后一件事，他表明了他的旨意：在这两种安排下，禁食与否的事都留给各人自由选择。我说这话，并不是为了支持我的观点而论证，只是要表明，回答他时，可以提出比他所说过的话有力得多的事。\par

% structure: p0012 source=h2 target=subsection reason=relative-heading-rank; base=h2

\subsection{第四章}

7. 我们的作者说：\BibleQuote{如果我们每周禁食两次，怎能逃脱与法利塞人同受定罪呢？}\BibleRef{路 18:11} 说得好像法利塞人是因为每周禁食两次而被定罪，而不是因为他骄傲地在税吏面前自夸。他同样可以说，那些把自己所有财物的十分之一施舍给穷人的人，也与那法利塞人一同被定罪，因为他也以此自夸于他的其他善工之中；但我却巴不得有许多基督徒这样做，而不是像我们所见的，只有极少数人这样做。或者让他说，凡不是不义的人、奸淫的人、勒索的人，都必须与那法利塞人一同被定罪，因为他自夸不是这些人；但能这样想的人，毫无疑问是心神错乱了。此外，如果法利塞人提到自己所有的这些事，大家公认本身是好的，就不该带着他那种显明的傲慢自夸去持守，而该以他那种所没有的谦卑虔诚去持守；那么，按照同样的原则，每周禁食两次，在像法利塞人那样的人身上是无益的，但在一个谦卑而有信德的人身上，却是一种虔敬的行为。而且，归根结底，圣经并没有说法利塞人被定罪，只说那税吏\BibleQuote{比他更成了义}。\par

8. 再者，当我们的作者在这一点上执意解释主的话：\BibleQuote{除非你们的义德超过经师和法利塞人的义德，你们决进不了天国}\BibleRef{玛 5:21}，并认为除非我们一周禁食超过两次，否则无法履行这一诫命时，他应当清楚知道一周有七天。倘若从这些日子中减去两天，即第七日和主日不禁食，那么便剩下五天，他可以凭着这五天超越那每周只禁食两次的法利塞人。因为我认为，若有人每周禁食三次，他便已超越那每周只禁食两次的法利塞人了。而倘若守了四次禁食，甚至多达五次，只越过第七日与主日不禁食——这是许多人终生奉行的做法，尤其在那些定居在\OriginalTerm{隐修院}{monasteries}里的人中更为常见——这样，不仅在禁食的劳苦上超越了法利塞人，也超越了那些惯常在第四、第六和第七日禁食的基督徒，一如罗马团体在很大程度上所做的那样。然而，你那位匿名的都主教辩论者却称这样的人为\OriginalTerm{属血肉的}{carnal}，即使他一周中有连续五日——除开第七日与主日——都禁绝了身体的一切饮食，仿佛别的日子里的饮食与血肉之躯毫无关系似的；并谴责他以自己的肚腹为神明，好像只有第七日的那一餐才入腹中一般。\par

\EditorialNote{我们毫不愧疚地在此省略了这封信大约八栏的内容，其中奥古斯丁以冗长的细节和浪费的修辞，揭露了那位罗马作者其他软弱且无关紧要的幼稚之谈，而卡苏拉努斯曾将该作者的作品提交给他审阅。我们不必跟随他进入那些浅薄之处，他在追击如此微不足道的对手时被卷入其中；让我们在奥古斯丁就基督徒是否有义务在星期六禁食这一问题给出自己意见的地方，继续翻译。}\par

% structure: p0016 source=h2 target=subsection reason=relative-heading-rank; base=h2

\subsection{第十一章}

25. 至于他论文结尾的那些段落，正如文中其他一些我认为不值一提的内容一样，它们与讨论一周的第七日我们应该禁食还是进食的问题就更加风马牛不相及了。不过，我留给你自己去审察和处理；特别是，如果你从我上述的话中已得到任何帮助的话。现在，我已经尽力对这位作者的论述作了回应，而且我想已经足够了。如果有人问我对此事的见解，我的回答是：在仔细斟酌这个问题之后，我在福音、书信，以及为教导我们而编定的整部称为\OriginalTerm{新约}{New Testament}的书籍中，看到禁食是被吩咐的。但我却没有发现主或宗徒们就我们应当或不应当禁食的日子定下任何明确的规则。因此，我相信，在第七日豁免禁食，更适合于——不是为获得，而是为预表——那永远的安息，在其中真正的\OriginalTerm{安息日}{Sabbath}得以实现，而这安息只能借着信德，并借着那使王的女儿在内里全然荣耀的义德才能获得。\par

26. 然而，在第七日是否禁食这一问题上，在我看来，没有什么比宗徒的以下规则更为稳妥与促进和平的了：\BibleQuote{吃的人不要轻视不吃的人，不吃的人也不要判断吃的人}\BibleRef{罗 14:3}，\BibleQuote{因为我们若吃，并不增加我们的好处；我们若不吃，也不减少我们的好处}\BibleRef{格前 8:8}；我们与同居共处、同在天主内生活的人们之间的共融，不致因这些事而受到扰乱。因为宗徒的话诚然真实：\BibleQuote{人若因食物而叫人跌倒，就有罪了}\BibleRef{罗 14:20}，同样真实的是，人若因禁食而叫人跌倒，也有罪了。因此，我们不要像那些人，见了\OriginalTerm{洗者若翰}{John the Baptist}也不吃也不喝，便说：\BibleQuote{他附了魔}；但我们也同样要避免效法那些人，看见\OriginalTerm{基督}{Christ}也吃也喝，便说：\BibleQuote{看哪，一个贪吃嗜酒的人，税吏和罪人的朋友}\BibleRef{玛 11:19}。主提到了这些话之后，还加了一句极其重要的真理：\BibleQuote{但智慧却藉自己的儿女彰显自己的正义}\BibleRef{玛 11:19}；你若要问这些儿女是谁，请读经上所写的话：\BibleQuote{智慧之子，就是义人的集会}\BibleRef{德 3:1}；他们就是那些人，在吃的时候，不轻视那些不吃的人；在不吃的时候，不判断那些吃的人；但他们确实会轻视并判断那些在吃或不吃上叫人跌倒的人。\par

% structure: p0019 source=h2 target=subsection reason=relative-heading-rank; base=h2

\subsection{第十二章}

27. 对于一周的第七日，实行上述规则困难较少，因为罗马教会以及其他一些教会——尽管为数不多，或近或远——都在那日守禁食；但是在\OriginalTerm{主日}{Lord's day}禁食则是极大的过犯，尤其自从那可憎的\OriginalTerm{摩尼教徒}{Manichaeans}的异端兴起以来，这异端如此明显而严重地违背大公信仰和圣经：因为\OriginalTerm{摩尼教徒}{Manichaeans}曾规定其追随者有义务在这一天禁食；结果，在主日禁食便愈发为人所憎恶。除非，也许有人能够不间断地禁食超过一周，以便尽可能接近四十日的禁食，如同我们曾听说有人做到的那样；我们甚至从最可靠的弟兄那里得到确证，确有一人完成了整整四十日的禁食。正如在旧约圣祖时代，\OriginalTerm{梅瑟}{Moses}和\OriginalTerm{厄里亚}{Elijah}四十日禁食时，并没有做出违反在一周第七日自由进食的事；照样，一个人若能在禁食上超过七天，他也并非刻意选择主日作为禁食之日，只不过在他定为禁食的那些日子当中——因为他已尽力立誓延长禁食——主日恰巧临到罢了。然而，若一连串的禁食要在一周之内结束，则没有比主日更适合作结束的日子了；但如果身体超过一周才得复苏，那么主日在那时并非被选为禁食之日，而是在那些他立意愿守的日子里，主日恰巧在其中出现罢了。\par

28. 不要被\OriginalTerm{普里西利安派}{Priscillianists}（一个与\OriginalTerm{摩尼教徒}{Manichæans}非常相似的教派）惯于引用\WorkTitle{宗徒大事录}中关于宗徒保禄在\Place{特洛阿}所行之事作为论据而动摇。那段经文如下：\BibleQuote{一周的第一日，门徒们聚集擘饼时，保禄向他们讲道，因为他次日要动身离去，就延长讲论直到午夜。}\BibleRef{宗20:7} 后来，当他从他们聚集的楼厅下来，去救治那因沉睡而从窗口跌落、已被人抱起已死的年轻人时，圣经进一步论及宗徒说：\BibleQuote{然后他再次上去，擘了饼，吃了，谈论许久，直到天亮，方才离去。}\BibleRef{宗20:11} 我们绝不接受以此断定宗徒们习惯在主日禁食。因为如今称为主日的日子当时称为一周的第一日，这在福音书中看得更为清楚；因为主复活之日，玛窦称之为\OriginalTerm{一周的第一日}{\GreekText{μία} \GreekText{σαββάτων}}，而其他三位福音作者则称之为\OriginalTerm{一周的第一日}{\GreekText{ἡ} \GreekText{μία} (\GreekText{τῶν}) \GreekText{σαββάτων}}，并且确知这日就是现所称的主日。因此，要么他们的聚会是在第七日结束之后——即在随后的夜间开始之时，而这夜属于主日，或一周的第一日——由此，宗徒在依照基督圣体圣事的方式与他们擘饼之前，延续他的讲论直到午夜，而且在举行圣事后，由于时间紧迫，为能在主日破晓时出发，又继续向聚集的人讲话；要么，若他们的聚会是在一周的第一日，即在主日日落前的一小时内，那么原文中\Quote{保禄向他们讲道，因为他次日要动身离去}这句话本身明确说明了他延长讲论的理由——即他即将离开他们，并希望给予他们充足的教导。因此，这段经文并未证明他们习惯在主日禁食，而只说明，宗徒认为不宜为了进食而中断一篇重要的讲论，当时那些即将与他离别的人正以极大的热情和兴趣聆听，而他由于多次行程，很少探访他们，也许除了这次之外别无其他机会，尤其因为——如后来之事所证——他当时离开他们时，并未预期在此生中还能再见到他们。不仅如此，通过这个实例，反而证明了在主日如此禁食并非惯例，因为记载这事的作者为免使人这样想，特意说明了讲论延长的原因，好叫我们明白，在紧急情况下，晚餐不应妨碍更重要的工作。然而，这些极为热切的听众的榜样更进一步；因为他们极度渴求的，不是水，而是真理之言，当他们想到那活泉即将从他们面前被移开时，不仅晚餐，连一般的饮食也弃之不顾，以不减的热情啜饮从宗徒唇间流出的一切。\par

29. 然而，在那个时代，虽然人们通常并不在主日禁食，但当有如保禄在特洛阿所遇的类似紧急情况时，有人在整个主日，乃至直到次日清晨，都未顾及身体的休养，这对教会来说并不算什么大错。但如今，自从异端者，尤其是这些极不虔敬的\OriginalTerm{摩尼教徒}{Manichæans}，已开始不只是在受环境所迫时偶尔在主日禁食，而是将这种禁食规定为一项受神圣庄严制度约束的义务，且他们的这种做法已为各基督徒团体所熟知；那么，即使出现宗徒所经历的那类紧急情况，我确实认为，现在也不应再行他所做之事，免得因给予冒犯而造成的损害，反而大于所讲言论带来的裨益。无论有何等必要性或充分理由，迫使基督徒在主日禁食——\Emph{例如}，我们在\WorkTitle{宗徒大事录}中看到，在船难危险中，他们在所乘载有宗徒的船上，连续禁食了十四天，期间主日出现了两次，\BibleRef{宗27:33}——我们应当毫不迟疑地相信，主日不应被列入自愿禁食的日子，除非有人许愿要持续禁食超过一周之久。\par

% structure: p0023 source=h2 target=subsection reason=relative-heading-rank; base=h2

\subsection{第十三章}

30. 教会之所以规定一周的第四和第六天进行禁食，其理由可由思考福音叙述而发现。我们在此发现，在一周的第四天，犹太人商议将主处死。相隔一日之后——在那日的晚上，即我们所称的一周第五天结束时，主与祂的门徒一同吃了逾越节晚餐——祂随后在属于一周第六天的夜里被出卖，这一天（众所周知）是祂受难之日。这一天，自晚上开始，是无酵饼的第一天。然而，圣史玛窦却说一周的第五天是无酵饼的第一天，因为在随之而来的晚上，他们将要守逾越节晚餐，届时他们要吃无酵饼和祭献的羔羊。由此可以推断，正是在一周的第四天，主说：\BibleQuote{你们知道，两天之后就是逾越节，人子将被出卖，被钉在十字架上。}\BibleRef{玛 26:2} 为此，这一天被视为适合禁食的日子，因为正如圣史随即补充说：\BibleQuote{那时，司祭长、经师和民间的长老，都聚集在名叫\Person{盖法}的大司祭的庭院里，共同议决要用诡计捉拿耶稣，加以杀害。}\BibleRef{玛 26:3} 中间隔了一天——就是圣史所写的这一天：\BibleQuote{无酵节的第一天，门徒前来对耶稣说：\InnerQuote{祢愿意我们在哪里给祢预备吃逾越节晚餐呢？}} \BibleRef{玛 26:17}——主在一周的第六天受难，这是众所公认的：因此，第六天也同样被正确地视为禁食的一天，因为禁食象征谦卑；因此经上说：\Quote{我以禁食刻苦我的灵魂。}\par

31. 次日是犹太人的安息日，在那一天，基督的身体安息于坟墓，正如在最初创造世界时，天主在那一天停止了祂的一切工程而安息。由此便产生了我们正在思考的祂新娘礼袍上的这种差异：有些人，尤其是东方的团体，宁愿在那一天进食，以象征神圣的安息；另一些人，即\Place{罗马}教会和西方的一些教会，则宁愿在那一天禁食，以纪念主在死亡中的谦卑。一年一次，即在复活节，所有基督徒都通过禁食来遵守一周的第七天，以纪念门徒们如同丧亡之人哀悼主的死亡时所怀有的哀伤（那些在一年中其余所有第七天都进食的人，也以最虔敬之心这样做）；这样便象征性地表达了这两件事——一年中有一个第七天表达门徒的哀伤，而其他所有第七天则表达安息的祝福。有两件事使义人的幸福和一切苦难的终结可被确信地期待，即死亡和死人的复活。在死亡中有那安息，先知论及它说：\BibleQuote{我的百姓，去，进入你的内室，关上门，隐藏片刻，等待盛怒过去。}\BibleRef{依 26:20} 在复活中，全人，包括身体和灵魂，都圆满地享受真福。因此，人们认为这两件事（死亡和复活）应当由食物的欢愉滋养来象征，而不是由禁食的艰苦来象征，唯独复活节星期六除外；在这一天，正如我所说过的，众人决议以更长的禁食来纪念门徒的哀伤，作为应当记念的事件之一。\par

% structure: p0026 source=h2 target=subsection reason=relative-heading-rank; base=h2

\subsection{第十四章}

32. 因此（如我前面所说），我们在福音书以及专属于新约启示的著作中，并未发现任何关于特定日子应守的禁食的律法规定；因此，这事成为许多难以枚举的事情之一，它们构成了君王女儿礼袍（即教会）的多样性——我要告诉你，关于这个问题，我曾请教可敬的\Place{米兰}主教\Person{盎博罗削}，我就是由他施洗的。当我母亲与我在那个城市时，我当时只是慕道者，对这些问题并不关心；但对她来说，是否应当按我们本城的习惯在星期六禁食，还是按\Place{米兰}教会的习惯不禁食，却是一个令她焦虑的问题。为了解除她的困惑，我向这位我刚才提到的人请教。他回答说：\Quote{我除了将我自己的做法推荐给人，还能推荐什么呢？}我以为他这样说是要简单地吩咐我们在星期六进食——因为我知道这是他自己的做法——但他接着又补充了这些话：\Quote{我在此地时，星期六不禁食；但我在\Place{罗马}时，则禁食：无论你到哪个教会，都要随从那教会的习惯，这样你就不致引起或遭受反感。}我把这个答复告诉了我母亲，她感到满意，并对此不再迟疑，我也遵循了同样的规则。然而，特别是在\Place{非洲}，有时会出现这样的情况：一个教会，或同一地区的各教会，可能有些成员在第七日禁食，而另一些则不禁食；在我看来，最好在每个会众中采纳那些在治理上被授予权威之人的习惯。因此，如果你十分愿意听从我的劝告，尤其因为我在这事上说的话已超出了必要的长度，那么不要在这一点上反对你自己的主教，而要毫无顾忌、不加争论地遵从他所奉行的习惯。\par

\SourceRecord{Translated by J.G. Cunningham.；Nicene and Post-Nicene Fathers, First Series；Vol. 1.；Edited by Philip Schaff.；Buffalo, NY: Christian Literature Publishing Co.,；1887.；<http://www.newadvent.org/fathers/1102036.htm>.}

% structure: p0001 source=h1 target=section reason=page-title

\section{书信 37（公元397年）}

\Emph{致\Person{辛普利西}}，\Emph{我最可祝福的主教，我最值得以敬意和真诚爱戴珍视的神父}，\Person{奥古斯丁}在主内向您致意。\par

1. 我收到了您圣善的慈爱所惠寄的信函——这封信带给我许多欢乐的缘由，因为它向我保证您记得我，您像往常一样爱我，并且您对主因怜悯乐意赐给我的每一项恩赐都感到极大的喜悦。在读这封信时，我热切地迎接了从您仁慈的心中流出的父亲般的情感：而这并不是我第一次发现它，好像它是短暂而新奇的，而是早已验证且熟知的，我最可祝福的主，最值得以尊敬和真诚爱戴珍视的主。\par

2. 我花了些文学辛劳撰写了几本书，竟得到如此大的回报，以致您的尊荣竟屈尊阅读它们？这岂不是主——我的灵魂所奉献于祂的那位——定意这样安慰我的忧虑，并减轻我在这种辛劳中不能不经历的恐惧，担心自己尽管在真理的平原上看似平稳，却或因无知或因不慎而失足？因为当我所写的得到您的赞许时，我便知道是谁赞许了它，因为我知道谁居住在您内；而那位所有属神恩赐的赐予者和分施者，正借着您的认可来坚固我对祂的顺服。因为我这些著作中凡配得您赞许的，都来自天主，祂曾借着我作为工具说：\BibleQuote{要有光}，就有了光；而在您的认可中，天主宣告了所成之事是\BibleQuote{好的}。\BibleRef{创 1:3}\par

3. 至于您屈尊命令我解答的那些问题，即使由于我头脑迟钝而不能理解它们，借着您功绩的帮助，我也许会找到答案。我仅请求您，因我的软弱向天主为我转求，并且无论我的何种著作到了您神圣的手中，无论是关于您曾如此慈爱和父亲般地引导我注意的那些题目，还是关于其他主题，您不仅费心阅读它们，而且承担起审阅和纠正它们的职责；因为我承认我自己犯的错误，正如我承认天主赐予我的恩赐一样容易。\par

\SourceRecord{Translated by J.G. Cunningham.；Nicene and Post-Nicene Fathers, First Series；Vol. 1.；Edited by Philip Schaff.；Buffalo, NY: Christian Literature Publishing Co.,；1887.；<http://www.newadvent.org/fathers/1102037.htm>.}

% structure: p0001 source=h1 target=section reason=page-title

\section{书信三十八（公元397年）}

\Emph{\Person{奥古斯丁}向他的兄弟\Person{普洛福图路斯}致以问候。}\par

1. 论到我的灵魂，赖\OriginalTerm{主}{Lord}的喜悦和祂屈尊赐予的力量，我安好无恙；论到我的身体，我却卧病在床。由于一个疮肿的疼痛和肿胀，我既不能行走，不能站立，也不能安坐。但即便如此，这既是\OriginalTerm{主}{Lord}的旨意，除了说我安好之外，我还能说什么呢？因为如果我们不希望祂所乐意做的事，我们就该责备自己，而不应认为祂所做的或容许发生的事可能有任何错误。这一切你都很清楚；但我更愿意对你说什么，除了我对自己的同样的话？因为在我眼中你就像是另一个我。因此，我把我的白昼和黑夜都托付给你虔诚的代祷。请为我祈祷，使我不致因缺乏节制而浪费我的时日，使我能以忍耐承受我的黑夜；请为我祈祷，虽然我行过死荫的幽谷，\OriginalTerm{主}{Lord}与我同在，我就必不惧凶险。\par

2. 无疑，你已听说了年迈的\Person{梅加利乌斯}去世的消息，因为他已卸下这必朽的身体二十四天了。如果可能，我想知道，你是否如你所计划的那样，见到了他的\OriginalTerm{首牧职}{primacy}的继任者。我们并未远离过犯，但同样真实的是，我们并没有被剥夺避难所；我们的忧伤并未止息，但我们的安慰同样持久。我出色的兄弟，你很清楚，在诸如此类的过犯中，我们必须警惕，以免对任何人的憎恨占据我们的心，这样不仅会妨碍我们关上门向\OriginalTerm{天主}{God}祈祷\BibleRef{玛 6:6}，也会向\OriginalTerm{天主}{God}自己关闭心门；因为对别人的憎恨会阴险地蔓延到我们身上，而任何一个发怒的人都不会认为他的愤怒是不义的。因为惯常对人怀有的愤怒会变成憎恨，那看似义怒中所掺杂的甜美，使我们将其留在器皿中的时间超过了我们当做的限度，直到全部变酸，器皿本身也被毁坏。因此，对我们来说，即便有人给了我们正当发怒的理由，我们也远不如克制愤怒更为明智，因为如果从对人发义怒开始，就可能因情绪的这种隐秘倾向而堕入憎恨。我们惯常说，在接待客旅时，宁可忍受接待一个坏人的不便，也不可因为小心提防不接纳任何坏人，而冒将好人拒之门外的风险；但在灵魂的\OriginalTerm{偏情}{passion}上，相反的原则才是正确的。因为为了灵魂的福祉，将心房的深处对愤怒关闭，即使它带着正当的要求来敲门，也远比容许它进入要无可比拟地更有裨益；因为一旦让它进入，将极难把它赶出，它会迅速从一棵小苗长成大树。愤怒敢于比人们想象的更大胆地突然增长，因为当日落之后，它在黑暗中不会感到羞耻\BibleRef{弗 4:26}。如果你想到最近在某次旅行中你对我说的话，你就会明白我是以何等深切的关怀和焦虑写下这些的。\par

3. 我问候我的兄弟\Person{塞维鲁斯}，以及和他在一起的人。如果带信人出发前的有限时间允许，我或许也会给他们写信。我也恳请你协助我说服我们的兄弟\Person{维克托}——我希望能通过你的\OriginalTerm{尊驾}{Holiness}向他表达感谢，因为他告知我他将前往\Place{君士坦丁纳}——不要拒绝从\Place{卡拉马}那条路回来，因为一桩他所知道的事务，在这件事上，我因老\Person{奈克塔里乌斯}对此急切恳求而背负了非常沉重的负担；他曾就此向我作过承诺。再会！\par

\SourceRecord{Translated by J.G. Cunningham.；Nicene and Post-Nicene Fathers, First Series；Vol. 1.；Edited by Philip Schaff.；Buffalo, NY: Christian Literature Publishing Co.,；1887.；<http://www.newadvent.org/fathers/1102038.htm>.}

% structure: p0001 source=h1 target=section reason=page-title

\section{热罗尼莫致奥斯定书（公元397年）}

\Emph{致我主奥斯定，一位父亲} \Emph{真正神圣且至为有福者，热罗尼莫在基督内致问候。}\par

1. 去年，我通过我们的兄弟、\OriginalTerm{副执事}{subdeacon} \Person{阿斯特里乌斯}之手，向阁下寄去一封理应向您问候的信函，我也欣然致意；我想那封信已送达您处。如今，我再通过我的圣善兄弟、执事 \Person{普雷西迪乌斯}写信，首先恳请您不要忘记我，其次请您接待这封信的送信人，我将此人托付给您，请求您视他如我极为亲近的朋友，并在他环境所需时给予鼓励和帮助；并非他缺什么（因为基督已丰厚地赋予了他），而是他极其渴望善人的友谊，并认为得到这份友谊便是获得了最宝贵的祝福。他前往西方的目的，您可以从他亲口得知。\par

2. 至于我们，安居于此隐修院内，却感受着四周波涛的冲击，并背负着尘世旅途的种种忧虑。但我们相信那位曾说过的话：\BibleQuote{你们放心，我已战胜了世界。}\BibleRef{若 16:33}且坚信，藉着祂的恩宠和引导，我们必能战胜我们的仇敌魔鬼。\par

我恳请您向那位圣善而可敬的弟兄、我们的父亲 \Person{阿里皮乌斯}转达我的敬意。在此隐修院中与我一起专心事奉主的弟兄们，向您热烈问候。愿基督我们的全能天主护佑您免受伤害，并使您时常记念我，我主、我父，真正圣善而可敬者。\par

\SourceRecord{Translated by J.G. Cunningham.；Nicene and Post-Nicene Fathers, First Series；Vol. 1.；Edited by Philip Schaff.；Buffalo, NY: Christian Literature Publishing Co.,；1887.；<http://www.newadvent.org/fathers/1102039.htm>.}

% structure: p0001 source=h1 target=section reason=page-title

\section{奥斯定致热罗尼莫（公元397年）}

\Emph{致我至爱之主，并致堪受以最诚挚爱德尊敬与拥抱的弟兄、我的同侪司铎 \Person{热罗尼莫}，\Person{奥斯定} 请安。}\par

% structure: p0003 source=h2 target=subsection reason=relative-heading-rank; base=h2

\subsection{第一章}

1. 我感谢你，因为你没有只给我一封形式上的问候信，而是写了一封信，尽管它远短于我渴望从你那里得到的；因为从你而来的任何东西，都不会使人感到冗长，无论它需要花多少时间。因此，尽管我因他人的事务，而且尤其是世俗事务，而充满巨大的焦虑，我本来难以原谅你书信的简短，若不是我考虑到它是为回复我一封更短的书信而写的。因此，我恳请你致力于书信往来，好使我们得以相通，不让我俩之间的遥远距离将我们完全阻隔；因为我们既在主内，即使搁笔沉默，也借着圣神的合一联结在一起。你在其中辛勤劳碌，从主的府库中取出宝藏的那些书卷，使我对你几乎有了全备的认识。因为，若我因为未曾亲睹你的面容而不能说\Quote{我认识你}，那么同样的道理也可以说你不认识自己，因为你不能看见你自己的面容。然而，如果仅仅因为知道自己的心灵就构成了你对自己的认识，我们借着你的著作对你的心灵也有不小的认识；研读这些著作时，我们称颂天主，因为祂为了你、为了我们、为了所有读你作品的人，赐下了如此这般的你。\par

% structure: p0005 source=h2 target=subsection reason=relative-heading-rank; base=h2

\subsection{第二章}

2. 不久前，在别的事物中，你的某本书到了我手中，我尚不知其名，因为手稿本身，第一页上没有照例写上标题。发现它的那位弟兄说，它的标题是\WorkTitle{Epitaphium}——若我们在这作品中只见到了对已故者生活或著作的记述，我们本会相信这标题是你所认可的。然而，由于书中提到了某些在写作时仍在世、甚至如今仍活着的人的著作，我们便奇怪，你为何给它冠以此名，或让他人相信你如此做了。这本书本身是一部有用的著作，我们完全认可。\par

% structure: p0007 source=h2 target=subsection reason=relative-heading-rank; base=h2

\subsection{第三章}

3. 在你对 \Person{保禄} 的 \BibleBook{迦拉达}{书} 的阐释中，我发现了一处令我深感忧虑的地方。因为，若本身不真实的陈述，却似乎出于责任感和为宗教利益，竟被接纳于圣经之中，那么圣经还剩下什么权威呢？若容许此事，从这些经卷中还能提出什么句子，足以用其分量去摧毁错误的邪恶固执呢？因为一旦你提出这一句子，若与你争辩的人不喜欢它，他便会回答说，在所引用的那段经文中，作者是出于某种可敬的责任感而说了谎。若人们竟能说并相信，宗徒在以这些话——\BibleQuote{我写给你们的，看，我在天主面前并不说谎}\BibleRef{迦 1:20}——引出他的叙述之后，却在对 \Person{伯多禄} 和 \Person{巴尔纳伯} 说\BibleQuote{我一见他们的行为不合乎福音的真理}\BibleRef{迦 2:14}时说了谎，那么谁还能找到不可能逃脱的出路呢？因为，若他们的行为合乎真理，\Person{保禄} 写的便是错误的；若他在 \Emph{这里} 写下了错误，那么他何时说过正确的话呢？难道当他的教导符合读者的偏好时，便应当被认为说的是真理，而凡与读者的印象相抵触的一切，便应算作他出于责任感所说的谎话吗？若我们接受这种解释原则，便不可能阻止人找到理由，认为他不但可能说了谎，而且是必须说谎。毋须多费言辞来继续这一论证，尤其当写信给你的时候，因为你的智慧和审慎，已无须再多言。我绝不敢如此傲慢，竟试图用我微小的铜钱，去充实你那蒙天主恩赐而如金的心灵；并且，没有人比你自己更有能力去修订和纠正我所提及的那部著作了。\par

% structure: p0009 source=h2 target=subsection reason=relative-heading-rank; base=h2

\subsection{第四章}

4. 你无须我教导你，宗徒是在什么意义上说\BibleQuote{对犹太人，我就成为犹太人，为赢得犹太人}\BibleRef{格前 9:20}，以及同一段经文中其它类似的话，这些话应归于怜悯同情的爱，而非蓄意欺骗的手段。因为服事病人的人，就仿佛自己也是病人；他并非假装发烧，而是以真正同情之心思想，若自己处在病人的境况中，会希望别人为自己做什么。\Person{保禄} 确曾是个犹太人；在他成为基督徒之后，他并未放弃那些犹太圣事，那是那一个民族以正确的方式、并在确定的时期内所领受的。因此，尽管他是基督的宗徒，他仍参与遵守这些圣事；但他的目的是要表明，这些圣事对那些即使在相信基督之后，仍愿意保留从祖先律法学来的仪式的人，绝无伤害；只要他们不将得救的希望建立在这些仪式上，因为由这些仪式所预表的救恩，已经由主耶稣带来了。为了同样的理由，他认为这些仪式绝不可强加给外邦的皈依者，因为那会加给他们一个沉重而多余的负担，可能使那些不习惯这些仪式的人偏离信仰。\par

5. 因此，他在伯多禄身上所指责的，并不是他遵守从祖先传下来的习俗——伯多禄若愿意，完全可以这样做，而不致被指控为欺骗或反复无常，因为尽管如今这些习俗已是多余，但对一个习惯了它们的人并无害处——而是他强迫外邦人遵守犹太礼仪，\BibleRef{迦 2:14}，他这样做，只能是使他显得如同甚至在主来临之后，这些礼仪的遵守仍为救恩所必需，而这与保禄藉宗徒职务所宣讲的真理相悖。宗徒伯多禄并非不明白这一点，但他因惧怕那些受割损的人而这样做了。因此，伯多禄显然受到了真正的纠正，保禄对此事的叙述也是真实的，除非我们在此承认一个虚假，就会使得为万代人的信德而赐下的圣经权威变得全然不定和动摇。因为在一封信的篇幅内，既不可能也不适宜说明，作出这样的让步后会带来何等巨大且无法言喻的恶果。然而，如果我们当面交谈，倒可以适时地、以较小的风险指出这一点。\par

6. 保禄已抛弃了属于犹太人且邪恶的一切事，尤其是这一点：\BibleQuote{他们由于不认识天主的正义，企图建立自己的正义，就没有服从天主的正义。}\BibleRef{罗 10:3}此外，他还在这一点上与他们不同：即在基督的受难和复活之后——在基督内，恩宠的奥秘已按照默基瑟德的品位赐下并显明——他们仍旧认为有义务举行旧约体制的圣事，并非仅仅出于对古老习俗的敬意，而是视其为得救所必需，然而这些圣事在某一时期确曾是必需的，否则，玛加伯为坚持这些圣事而遭受的殉道，就会变得无益且徒然了。\BibleRef{加下 7:1}最后，保禄还在这一点上与犹太人不同：他们迫害宣扬恩宠的基督徒宣讲者，视之为法律的仇敌。他宣称，所有这些以及类似的错误和罪恶，\BibleQuote{只当作损失和粪土，为赚得基督；}\BibleRef{斐 3:8}但他这样说，并非贬低犹太法律的礼仪，只要这些礼仪是按照祖先的习俗来遵守，就如他自己遵守它们的方式那样，不视之为得救所必需，而不是按照犹太人所主张的必须遵守的方式，也不是在实行他在伯多禄身上所指责的那种欺骗性的伪装。因为如果保禄是为了赢得犹太人，而假装成犹太人，从而遵守这些圣事，那么当他向那些没有法律的人，成为没有法律的人，为赢得他们的时候，为何不同时参与外邦人的异教祭祀呢？对此的解释是，他之所以参与犹太人的祭献，是因为他生来就是犹太人；并且，当他说我所引用的这一切话时，他的意思并不是他假装成自己所不是的，而是他怀着真诚的怜悯，觉得必须给予他们这样的帮助，如同自己若陷入他们的错误中也会需要的一样。在此，他所行的不是欺骗者的诡诈，而是怜悯的拯救者的同情。在同一段经文中，宗徒更概括地陈述了这一原则：\BibleQuote{对软弱的人，我就成为软弱的，为赢得那软弱的人；对一切人，我就成为一切，为的是总要救些人。}\BibleRef{格前 9:22}——这后半句话引导我们理解前半句的意思，即他表现出怜悯他人的软弱，就如同那是他自己的软弱一样。因为当他说：\BibleQuote{谁软弱，我不软弱呢？}\BibleRef{格后 11:29}时，他不希望别人以为他是假装承受他人的软弱，而是他以同情之心将其表露出来。\par

7. 因此，我恳求你，以一种坦诚而真正属于基督徒的严厉，用于那本书的纠正与修订，并高唱希腊人所谓的\OriginalTerm{帕利诺狄亚}{\GreekText{παλινῴδια}}。因为比起希腊的\Person{海伦}，基督徒的真理更为可爱，无可比拟：为了保卫她，我们的殉道者曾与\Place{索多玛}作战，其勇毅超过了希腊英雄们为\Person{海伦}而攻打\Place{特洛伊}时所表现的。我这样说，并非为让你恢复灵性的视力——我绝无说你已失去这视力的意思！——而是为让你那双既明澈又敏于辨别的眼睛，转视那方，你曾以无从解释的伪装，从那里转离，不肯看见那一旦我们承认一位圣书的作者，竟可以在其著作的任何部分中，正当地、虔诚地说出虚假之事，便会随之而来的灾难性后果。\par

% structure: p0014 source=h2 target=subsection reason=relative-heading-rank; base=h2

\subsection{第5章}

8. 关于此事，我前些时候曾给你写过一封信，那封信没有送达你处，因为受托带信的人未能完成前往你处的行程。从那封信中，我可以引述一个当时口授时浮现于脑海的想法，在这封信中我不应将其省略，好使如果你的意见仍与我的不同，并且更好，你能欣然原谅那促使我写作的焦虑。这想法是：如果你的意见不同，且符合真理（因为只有在这种情况下，它才能比我的更好），你就会承认，\Quote{一个我为了真理的利益而犯的错误，不可能应受重大的指责，假如确有任何指责的话，而既然你能够为了虚假的利益使用真理而不算犯错。}\par

9. 关于你欣然就\Person{奥力振}一事给我的答复，我并不需要被提醒：我们不仅对教会作家，而且对所有其他作家，都应认可并称赞其中正确与真实之处，却要摒弃并谴责其中虚假及有害之点。我从你的智慧与学识所渴望得到的（并且至今仍渴望得到），乃是你应确切地告诉我们，这位卓越的人物在哪些论点上被证实偏离了真理的信仰。此外，在那本你列举了你能忆及的全部教会作家及其作品的书中，我认为，若你在提及你所知道的异端者时（既然你已决定不漏掉他们），能补充说明其教义当在哪些方面加以规避，这样安排或许更为便利。这些异端者中，你亦有所遗漏，我愿知其原由。然而，倘若你之吝于增添对异端者的评论，以及天主教会权威定断其所应受谴责的事项，是出于不欲使该卷书过分臃肿的愿望，那么我恳请你不要吝惜以你部分的文学劳作，——凭我们主天主的恩宠，你已藉此极大地激励并帮助了圣人们学习拉丁语，——对于这一我以谦卑和弟兄之爱提请你关注的主题，略施笔墨，并请（若你其他事务允许）出版一本小书，辑录截至如今所有试图因傲慢、无知或固执己见而颠覆基督信仰之纯朴性的异端者们乖谬的信理之纲要；此书对于那些因缺乏闲暇或因不识希腊文而无法阅读及理解如此繁多事物之人，实属极为必要。我本想更详尽地提出我的请求，若非此举通常被视为对所求恩惠之施予者的善意抱有疑虑的话。同时，我诚恳地将我们的弟兄\Person{保禄}推荐于你在基督内的善意，对于他在此地的崇高声望，我在天主前乐意作证。\par

\SourceRecord{Translated by J.G. Cunningham.；Nicene and Post-Nicene Fathers, First Series；Vol. 1.；Edited by Philip Schaff.；Buffalo, NY: Christian Literature Publishing Co.,；1887.；<http://www.newadvent.org/fathers/1102040.htm>.}

% structure: p0001 source=h1 target=section reason=page-title

\section{书信第41篇（主历397年）}

\Emph{致我们至可颂扬、最堪受尊敬的\Person{奥勒留}主教，我们最亲爱的弟兄和司铎职上的同工，\Person{阿利比乌斯}与\Person{奥古斯丁}在主内问候。}\par

1. \BibleQuote{我们的口满含喜笑，我们的舌充满欢呼}，因你的来信告知，蒙那位以灵感指引你的天主的助佑，你已经实现了你关及我们所有同职弟兄的虔诚计划，尤其是关于司铎们在你面前向民众讲道的规定；借着他们如此投入的口舌，你的爱德在人心中发出的声音，比他们的声音在人耳中更响。感谢天主！我们心中所有、口中所说、笔端所录的，还有比这更好的吗？感谢天主！没有比这更易上口的言辞，没有比这更悦耳、意蕴更深、践行起来更富裨益的话语了。感谢天主！祂赐给你一颗如此真诚顾念子女利益的心，并使你内心隐而未见、无人得窥的愿望显明出来，不但给你行善的意愿，还给了你达成心愿的方法。诚愿如此，诚然如此！愿这些善工在人前照耀，使人们看见，并因此而欢跃，光荣你们在天之父。\BibleRef{玛 5:16} 你该在主内以此类事为乐；也愿你的司铎们因你的祈祷而蒙天主垂允，这位天主，你并不认为经由他们发言而聆听其声是屈尊降贵之事。愿他们前行、奔跑，走在主的道上！愿卑微者与尊贵者一同蒙福，因听到那些人对他们说\BibleQuote{让我们进入主的圣殿！}而欢欣！愿强者引路，愿弱者效仿他们的榜样，跟随他们，如同他们跟随基督一样。愿我们众人如蚂蚁般急切地奔行于神圣勤劳之道途，又如蜜蜂般在圣善职务的芬芳中劳作；并愿靠着恒守到底的拯救恩宠，凭坚忍结出果实！愿主\BibleQuote{不让我们受那超过我们能力的试探，而且在试探中必给我们开一条出路，使我们能够承担}！\BibleRef{格前 10:13}\par

2. 请为我们祈祷：因为你们向天主献上如此真诚的爱德与由善工而来的赞颂，我们便视你们的祈祷为堪蒙垂允的。请祈求，愿这些善工也照耀在我们身上，因你所祈求的那一位知道，我们看见它们照耀在你们身上时，是何等的满心欢喜。这便是我们的心愿；这便是我们内诸多思念中满盈的安慰，使我们的灵魂欢欣。现在如此，因为天主的应许本是如此；正如祂所应许的，将来也必如此实现。我们因那位降福了你，又借着你将祝福赐予你所服事之民众的主，恳求你，命人将你所喜欢的司铎们的讲道稿誊录下来，审校后寄给我们。至于我本人，并未忽略你所要求的；正如我从前屡次写信所说，我仍切望得知你对\Person{蒂科尼乌斯}的\WorkTitle{七条规则或钥匙}的看法。\par

我们诚恳地将我们的弟兄\Person{希拉里努斯}，那位\Place{希波}的首席医生和行政长官，推荐给你。至于我们的弟兄\Person{罗玛努斯}，我们知道你正为他竭力奔走，我们只需祈求天主使你一切顺利即可。\par

\SourceRecord{Translated by J.G. Cunningham.；Nicene and Post-Nicene Fathers, First Series；Vol. 1.；Edited by Philip Schaff.；Buffalo, NY: Christian Literature Publishing Co.,；1887.；<http://www.newadvent.org/fathers/1102041.htm>.}

% structure: p0001 source=h1 target=section reason=page-title

\section{书信第四十二篇（主历397年）}

\Emph{致\Person{保利努斯}与\Person{塞拉西亚}，我在基督内的弟兄与姊妹，堪受尊敬与赞颂，最以虔诚著称者，\Person{奥古斯丁}在主内致意。}\par

难道我们曾盼望或期待，如今竟需藉我们的弟兄\Person{塞维鲁斯}向你们索要那份你们的爱尚未写给我们的回信，而我们如此长久、如此焦急地渴望你们的答复？我们为何注定要经历两个夏季（而且是在干涸的非洲土地上）承受这样的焦渴？我还能说什么呢？哦，慷慨的人，你每天都在施舍你自己的财物，请公正些，偿还所欠我们的债务。或许你长久耽延的原因是你想要完成并寄送给我的\WorkTitle{那本反驳异教崇拜的书}——我曾听说你正致力于此，且已表达了极为恳切的渴望。但愿你能以此丰盛的盛宴满足那因超过一年的斋戒（就你的笔而言）而激起的饥饿！但若此书尚未备妥，我们的抱怨将不会止息，除非你在此期间使我们免于在书成之前便饥饿至极。请问候我们的弟兄，特别是\Person{罗马努斯}和\Person{阿吉利斯}。此处所有与我在一起的人问候你们，若非他们对你们的爱少于实际所怀，他们也不会因你迟延写信而如此被激怒。\par

\SourceRecord{Translated by J.G. Cunningham.；Nicene and Post-Nicene Fathers, First Series；Vol. 1.；Edited by Philip Schaff.；Buffalo, NY: Christian Literature Publishing Co.,；1887.；<http://www.newadvent.org/fathers/1102042.htm>.}

% structure: p0001 source=h1 target=section reason=page-title

\section{第43封信（主历397年）}

\Emph{致\Person{葛洛里}、\Person{厄琉西}、两位\Person{费利克斯}、\Person{格拉玛提库斯}及所有悦纳此书者，我最亲爱且堪受称颂的诸公，\Signature{奥古斯丁}致意。}\par

% structure: p0003 source=h2 target=subsection reason=relative-heading-rank; base=h2

\subsection{第一章}

宗徒保禄曾说：\BibleQuote{对异端人，在谴责一两次以后，就该远离他，因为你知道，这样的人背弃了正道，犯罪作恶，自暴自弃。}\BibleRef{铎 3:10}然而，若有人所持的道理虽然虚妄错谬，却并不以狂热的固执加以捍卫，尤其当他们并非出于自己冒昧的妄断而捏造，乃是从已受误导而陷入错谬的父母那里接受过来，并且正在忧心如焚地寻求真理，一经发现便准备改正，这样的人就不当被视为异端者。若非我深信你们正是如此，或许便不会致书于你们了。但即便如此，对待一个异端者，无论他如何被恶毒的虚荣所膨胀，因顽固地悖逆真理而疯狂，我们虽警告他人要远离他，免得他欺骗软弱无经验的人，却不拒绝尽一切可能的方法来挽救他。因此，我甚至也给\OriginalTerm{多纳徒派}{Donatistae}的一些首领写了信，这些信固然不是共融的书信——由于他们的悖逆，他们早已不再从遍布全世界的、未分裂的天主教会领受——而是私人的函件，一如我们可寄予外教人的那种。然而，他们虽偶尔读了这些信件，却不愿——或者更可能是不能——作出答复。我以为，在这些情况下，我已尽了那爱所要求的本分，这爱是圣神教导我们不仅要施于自己人，也要施于众人，正如圣神借着宗徒所说：\BibleQuote{愿主使你们彼此间的爱情，并对众人的爱情日益增长满溢，正如我们为你们所有的爱情一样。}\BibleRef{得前 3:12}在另一处又劝告我们，对于那些持不同意见的人，当以温柔规劝，\BibleQuote{如果天主开恩，赐他们悔改而认识真理，使他们这些被他活捉去顺从他心意的人，能醒悟过来，脱离魔鬼的罗网。}\BibleRef{弟后 2:25}\par

我说这些话作为前言，以免有人因为你们不在我们的共融之内，便认为我是出于冒昧而非慎重，才寄出这封书信，并切愿如此与你们商谈有关灵魂的福祉；虽然我相信，若果我写信给你们，是为了一宗财产业务，或排解某项银钱纷争，便无人会非难我。世间在人们眼中何其珍贵，而他们对自身的估价又何其微小！所以，这封书信将在天主的审判台前，为我作证，因祂洞悉我写信的心意，且曾说过：\BibleQuote{缔造和平的人是有福的，因为他们要称为天主的子女。}\BibleRef{玛 5:9}\par

% structure: p0006 source=h2 target=subsection reason=relative-heading-rank; base=h2

\subsection{第二章}

因此，我恳请你们回想，当我还在你们城里的时候，曾与你们略略讨论基督徒合一的共融；当时你们拿出了某些案卷，从中当众宣读了一段，声称约有七十位主教，将前\Place{迦太基}主教\Person{凯西利安}，连同他的同僚以及祝圣他的人，一并定罪。同一案卷中又详细陈述了\Place{阿普同加}的\Person{费利克斯}的案件，描绘得尤为可憎恶劣。当这些被全部宣读后，我答说，那些当时制造了这次裂教的人，既不惜窜改案卷，也就难怪会认为，应当定罪那些因受忌妒邪恶之人煽动而成了目标的人，何况那判决是在未曾审慎商议、被定罪者缺席、且未将指控之罪告知他们的情形下作出的。我又补充说，我们另有别的教会案卷，据此，当时任\Place{努米底亚}首席主教的\Place{提吉西斯}的\Person{塞昆都斯}，将那些在场并自承为\OriginalTerm{背教者}{traditor}的人，交付于天主的审判，允许他们继续留任当时所占的主教座堂；我指出，那些人的名字既列在定罪\Person{凯西利安}的名单上，而这位\Person{塞昆都斯}本人更是那次公会议的主席，在会上，他利用那些因当众承认了同一罪行而被他赦免者的表决，确保了那些缺席并被指控为背教者的人被定罪。\par

我又接着说，在\Person{马约里努}被祝圣之后不久——他们以不敬的邪恶，立他抗衡\Person{凯西利安}，竖立一座祭台反对另一座祭台，并以昏聩的争辩撕裂基督的合一——他们便向当时的皇帝\Person{君士坦丁}申诉，请求指派几位主教为法官和仲裁者，审理那些在\Place{阿非利加}发生、扰乱了教会和平的问题。这事既成，\Person{凯西利安}连同那些从\Place{阿非利加}渡海前来控告他的人均已在场，案件由时任\Place{罗马}主教的\Person{默基德}，连同应多纳徒派请求、皇帝所派遣的襄审官一同审理，结果对\Person{凯西利安}提不出任何罪证；于是，他得以确认其主教职位，而作为对方到场的\Person{多纳图}则被定罪。此后，他们全体仍顽固坚持其罪大恶极的裂教行为，皇帝被再次上诉，便特意在\Place{阿尔勒}对此事再次详加审断。然而，他们拒绝教会的裁决，又向\Person{君士坦丁}本人上诉，请求亲审其案。当这次审讯进行时，双方均到场，\Person{凯西利安}被宣告无罪，他们以败诉退下；但仍固执于同样的悖逆。同时，\Place{阿普同加}的\Person{费利克斯}的案件也未被人遗忘，依照同一位皇帝的命令，经总督调查后，他也获判无罪，对其指控的罪行不成立。\par

5. 然而，由于我只是口说这些事，并非宣读记录，你们似乎觉得我的热忱带来的期望未能完全满足。觉察到这一点后，我立即派人去取我答应要宣读的那份文件。在我继续前往\Place{杰利兹}的教会，打算从那里返回你们这里时，所有这些会议录在两天之内就送到了你们那里，并且如你们所知，在一天之内尽可能按时间向你们宣读了。我们先读到，\Place{提吉西斯}的\Person{塞孔杜斯}如何不敢罢免那些自己承认是\OriginalTerm{交付者}{traditores}的同僚；但后来，却正是在这些人的帮助下，竟敢在他们未认罪且缺席的情况下，定\Person{凯基利安}及其同僚的罪。然后我们又读了\OriginalTerm{总督的档案}{Acta Proconsularia}，其中\Person{斐理斯}在经过最彻底的调查后被证明无罪。你们会记得，这些是在上午宣读的。下午我向你们宣读了他们致\Person{君士坦丁}的请愿书，以及他在\Place{罗马}任命法官的教会审理记录，据此多纳徒派被定罪，\Person{凯基利安}的主教尊严得到确认。最后，我宣读了\Person{君士坦丁}皇帝的信函，这些信函中的证据确凿无疑地证实了这一切。\par

% structure: p0010 source=h2 target=subsection reason=relative-heading-rank; base=h2

\subsection{第三章}

6. 先生们，你们还想要求什么呢？你们还想要求什么呢？这里所涉及的不是你们的金银；不是你们的土地、财产，也不是身体的健康。我呼吁你们的灵魂关注获得永生和逃避永死。醒醒吧！我处理的并非晦涩的问题，也不是探究什么隐藏的奥秘——对于这些奥秘的探究，人类理智根本无力承担，或至少只有极少数人能胜任；这事情如同白昼般明晰。还有比这更明显的吗？还有什么能更快看清呢？我断言，无辜且未到场的人被一个非常庞大却行事仓促的会议定罪了。我用总督的档案来证明这一点，你们派别所引用的那次会议的记录，曾宣称那个人是特为有罪的交付者，但在这档案中他却被完全洗清了这一罪名。我进一步断言，对那些被指控为交付者之人的判决，是由一些自己承认犯有同一罪行的人作出的。我用\OriginalTerm{教会会议录}{Acta Ecclesiastica}来证明这一点，其中列出了那些人的名字。\Place{提吉西斯}的\Person{塞孔杜斯}声称渴望维护和平，赦免了这些他知道已犯下罪行的人；后来，尽管破坏了和平，他却又借助这些人的力量，对那些他并无犯罪证据的人作出了判决。这样一来，他就清楚地表明，他先前的决定并非出于对和平的顾及，而是出于对自己的恐惧。因为\Place{利玛塔}的主教\Person{普普里乌斯}曾指控他，说他本人曾因一名督查和士兵的扣押，为迫使他交出圣经而被拘禁，后来却被释放了，毫无疑问是在付出了代价之后，要么是交出了什么东西，要么是命令别人替他这样做。\Person{塞孔杜斯}担心这一嫌疑很容易被证实，于是他听取了\Person{小塞孔杜斯}——他的亲属——的建议，并与所有主教同僚商议后，赦免了那些无需证据、只待天主审判的罪行，从而看似是在维护教会的和平：但这是假的，他只是在保护自己。\par

7. 因为，如果他心中真的顾及和平，那么后来在\Place{迦太基}，他就不会和那些他留给天主审判、当时在场且已认罪的交付者联合起来，对并未在场、且无人证实其罪行的其他人在同一罪行上作出判决。此外，在那个场合他本应更加担心扰乱和平，因为\Place{迦太基}是一座伟大而著名的城市，任何发源于那里的邪恶都可能像从身体的头部蔓延到整个\Place{阿非利加}那样。\Place{迦太基}也靠近海外各地，并以其显赫声望著称，因此它有一位影响力超乎寻常的主教，他能不顾众多敌人的合谋反对，因为他看到自己既通过\OriginalTerm{通功书信}{litterae communicatoriae}与\Place{罗马}教会——在那里\OriginalTerm{宗座}{sedes apostolica}的权威始终盛行——相连，也与\Place{阿非利加}最初接受福音的其他所有地区相通，并且准备在这些教会前为自己辩护，以防对手试图使这些教会与他疏远。因此，既然\Person{凯基利安}拒绝来到他的同僚面前——他察觉或怀疑（或如他们断言，假装怀疑）这些同僚因他的敌人而对案件的实情有偏见——\Person{塞孔杜斯}若想成为真正和平的守护者，就更应避免对那些完全拒绝出席他们法庭的人进行缺席定罪。因为这并不是涉及司铎、执事或低级圣职人员的事，而是涉及同僚主教的事，他们可以将案件完全交给其他主教，特别是那些教会的判决——在这些教会中，对他们作出的缺席判决将毫无效力，因为他们并非在出庭后放弃了法庭，而是一直因自己有怀疑而拒绝出庭。\par

8. 这一考虑本应使当时担任\OriginalTerm{首席主教}{Primate}的\Person{塞康杜斯}深思有加，若其身为\OriginalTerm{公会议}{Council}主席渴望促进和平；因为他若能如此说，或许就能平息或遏止那些对缺席者咆哮的人的口舌：\Quote{弟兄们，你们看，在迫害的巨大浩劫之后，和平已藉着天主的仁慈，由这个世界的君王赐予我们；我们身为基督徒和主教，断不可破坏那甚至连异教仇敌都停止攻击的基督徒合一。因此，要么让我们将那段极其动荡时代的不幸加诸教会的一切事件，留给身为审判者的天主；要么，如果你们中间有人对别人的罪行有如此确凿的认识，足以对他们提出明确的指控，并能在他们申辩无罪时予以证实，且又不愿与此等人\OriginalTerm{共融}{communion}，就让他们速速前往我们海外的弟兄和同侪，即海外教会的主教们那里，首先向他们控告被告的行径和顽抗，指称他们因自知有罪而回避在非洲同侪的审理权，这样，可经由这些海外主教，传唤他们前来出庭，并对指控之事作出答辩。若他们不服从这一传召，其罪行和顽梗将暴露于那些主教；并以他们的名义，发出一封会议通函至现今基督教会延展所及的普世各地，借此将被控各方排除于所有教会的共融之外，以免在迦太基教会的\OriginalTerm{主教座}{see}滋生谬误。既已如此行，他们被隔绝于整个教会之后，我们就可毫不畏惧地给迦太基的社团祝圣另一位主教；但若我们现在祝圣另一位主教，海外教会极可能不与他共融，因为他们不会承认那位主教的罢免有效，而该主教的祝圣曾普世公认，且曾与他们互换共融信函；如此，由于我们急于未经深思便仓促定谳，当教会从外部获得安宁之际，内部却可能出现\OriginalTerm{裂教}{schism}的重大丑闻，而我们竟会胆敢另立祭台，不是反对\Person{凯基利安}，而是反对普世教会，因为它对我们的步骤一无所知，仍将与他保持共融。}\par

9. 若有人立意拒绝如此健全而公允的劝言，他还能做什么呢？他怎能在得不到\OriginalTerm{公会议}{Council}的授权下，获取对其缺席同侪的定罪，更何况首席主教反对他？倘若竟有人如此严重地反叛首席主教的权柄，以致某些人执意要立即定罪那些首席主教希望暂缓处理的人，那么，他通过异议，将自己与这般好斗且结党之徒分开，较之与普世共融分离，岂不更好！但正因在海外主教面前，无法证实任何指控\Person{凯基利安}及参与其祝圣者的罪状，那些判他们罪的人便不愿延缓判决；而宣判之后，他们也不费心向海外教会指明非洲那些应视作已被定罪的\OriginalTerm{背教者}{traditors}，并回避与之共融的名字。因为他们若曾尝试这样做，\Person{凯基利安}及其他人必会自辩，并在海外主教的教会法庭前，通过最彻底的审判，证明自己清白，驳斥诬告者。\par

10. 我们认为，那个乖谬而不公的\OriginalTerm{公会议}{Council}，主要系由\Place{提吉西斯}的\Person{塞康杜斯}在他们供认罪行后予以宽赦的背教者所组成；当有流言传扬，说某些人曾犯有交付圣书的罪行，他们便力图通过对他人横加诽谤，来转移对自己的猜疑，并借着散布层层谎言的迷雾，逃脱对其罪行的揭露；当时全非洲的人听信他们的主教，竟对那些无辜者散布虚言，说他们已在迦太基被定为背教者。由此，亲爱的朋友们，你们便可看出，你们中间某些人所断定为不可能的事，确能发生，即：那些自己承认背教罪行，并获得将其案件发还神圣法庭审理的人，后来竟参与审判和定罪他人，而那些人并未在场自辩，却被控以相同罪行。因为正是他们自身的罪咎，使他们更热切地抓住机会，用无端的指控压倒他人，并借此让众人的口舌忙碌，从而掩护他们自己的劣迹免受调查。再者，若一个人在别人身上断定自己也曾犯过的过失，竟是不可想象的，宗徒\Person{保禄}就不会有机会说：\BibleQuote{因此，人啊！你不论是谁，你判断人，必无可推诿：因为你判断别人，就是定你自己的罪，因为你这判断人的，正作着同样的事。}\BibleRef{罗 2:1} 这正是那些人所作所为，以致宗徒的这话完全可以且恰如其分地应用在他们身上。\par

11. \Person{塞康杜斯}因此，将这些人所坦白的罪行移交天主的法庭，并不是为了促进和平与合一：因为如果是这样，他就会更加谨慎地防止在\Place{迦太基}发生\OriginalTerm{分裂}{schism}，因为当时并没有任何人在场，使他不得不赦免他们坦白的那罪行；相反，维护和平所需要的，只是一道拒绝给不在场者定罪的命令。如果他们甚至决定\Emph{宽恕}这些清白的人，那他们就是以不义对待这些清白者，因为这些人既未经证实有罪，也未曾认罪，实际上根本不在场。因为一旦一个人接受宽恕，他的罪行便无可置疑地成立了。那些以为自己有权凭着自己不知道、因此即使想宽恕也无从宽恕的罪行来定罪的人，是多么放肆和盲目！在前一种情况下，已知的罪行被移交天主仲裁，以免追究其他罪行；在后一种情况下，未知的罪行被当作定罪的依据，以便掩盖已知的罪行。但有人会说，\Person{凯基利安}和其他人的罪行是已知的。即使我承认这一点，他们不在场的事实也应当保护他们免受这样的判决。因为他们不能被指控抛弃了一个他们从未站立其前的法庭；而且，教会并非完全由这些非洲主教代表，以至于他们拒绝在这些主教面前出庭，就可以被视为回避一切\OriginalTerm{教会的管辖权}{ecclesiastical jurisdiction}。因为在海外地区还有成千上万的主教，那些似乎不信任自己在非洲和\Place{努米底亚}的同侪的人，显然可以在他们面前受审。你难道忘记了圣经的诫命：\BibleQuote{未询问以前，不可责斥人；先要考察，然后方可责备}？\BibleRef{德 11:7} 如果圣神禁止我们在询问某人之前责备或纠正他，那么对那些因不在场而无法就被指控的罪行接受讯问的人，不仅责备纠正，甚至连定罪都做了，这罪过该是何等更重呢！\par

12. 再者，至于这些审判者所宣称的：虽然控告的一方不在场，他们并非逃避审判，而是一直表明不信任那个派别，并拒绝出现在他们面前，但他们对这些人定罪所指的罪行是众所周知；我的弟兄们，我问你们，他们怎会知道呢？你们回答说，我们无从得知，因为公共记录中未说明证据。但我要告诉你们他们是如何知道的。请仔细留意\Place{阿普坦加}的\Person{腓理斯}的案件，首先看看他们对他有多么激烈；因为他们对其他人所知，恰恰与对他的情况有着同样的理由，而他却通过彻底而严厉的调查，后来被证明完全清白。既然那个他们怒火最盛、猛烈攻击的人已被证明清白，那么我们更有理由以更大的公义、安全与坦然，相信其他受到较轻控告、定罪较轻的人也是清白的了！\par

% structure: p0018 source=h2 target=subsection reason=relative-heading-rank; base=h2

\subsection{第四章}

13. 也许会有人提出一种反对意见，虽然当它被提出时，你们并不赞成，但我不应略过不提，因为别人也曾提出过它，即：一位主教不该由\OriginalTerm{总督}{proconsul}审讯而获判无罪：仿佛主教是自己设法促成这场审讯，而不是奉皇帝的命令进行的，因为这件事特别关系到他向天主负责的职责。实际上，正是他们自己，就有关交出\OriginalTerm{圣书}{sacred books}和\OriginalTerm{分裂}{schism}的问题，向皇帝呈递请愿书，后来又向他申诉，从而让皇帝成为了此事的仲裁者和审判者；然而他们却拒绝服从他的决定。因此，如果那位并非自己向那个法庭申诉，却被\OriginalTerm{长官}{magistrate}判定无罪的人应当受责备，那么那些巴望一位世俗君王作他们案件的裁判者的人，该当受到何等更重的责备啊！因为如果向皇帝申诉并非不对，那么由皇帝审判也就并非不对，从而由皇帝委派审案的人审判自然也非不对。你的一位朋友曾急于根据以下事实来制造抱怨的口实：在主教\Person{腓理斯}的案件中，一个证人被悬吊在拷问架上，另一个被铁钳折磨。但是，当辩护人正努力发现真相的案件是腓理斯自己的案件时，他有权力阻止审讯的进行——即使是严谨甚至严苛地进行吗？因为这样的拒绝调查，除了被认为是认罪，还能被解释成什么呢？然而那位总督，周围环绕着传令官可畏的叫声和刽子手血污的双手，如果尚有其他地方可以审理此案，他是不会对一位拒绝到他庭前，与自己同级别的缺席者定罪的。或者，如果他是在这种情况下宣判，他自己必定会承受民法规定的公正裁决。\par

% structure: p0020 source=h2 target=subsection reason=relative-heading-rank; base=h2

\subsection{第五章}

14. 然而，如果你们拒绝一位总督的案卷，那么请服从教会的案卷。这些案卷都已按顺序向你们宣读过了。也许你们会说，\Place{罗马}教会的主教\Person{默基德}，连同与他协作的其他海外主教，无权僭取法官之位，去审理一件早已由七十位\Place{非洲}主教（以\Place{蒂吉西斯}的主教兼首席主教为首）裁定的事项。但如果他事实上并没有僭取此位呢？因为皇帝被上告后，派遣主教们与他同席审判，授权他们以他们认为公正的方式裁决整个案件。这一点我们既有\OriginalTerm{多纳图派}{Donatists}的请愿书，也有皇帝本人的话为证；你们记得，这些都曾向你们宣读，现在也提供给你们研究和抄录。请阅读并思考所有这些。你们可以看到，为了保持或恢复和平与合一，每件事是如何被谨慎细致地讨论的；控方的法律资格是如何被审查的，其中某些人在此方面暴露出何种缺陷；以及在场者的证词如何清楚地证明，他们无话可指控\Person{凯基利安}，却想把整个案件移交给属于\Person{马约里努斯}一派的人，即那些反对教会和平的暴乱群众，其目的无非是让\Person{凯基利安}被那伙人控告；他们以为这伙人力量强大，只靠喧嚣吵闹，无需书面证据或事实审查，便足以使审判官的心思偏向他们的意见——除非一个因错谬与邪恶之\Emph{杯}而疯狂痴迷的群众，竟能对\Person{凯基利安}提出真实的指控；在这样一个地方，七十位主教曾以疯狂的轻率，缺席审判了与他们同级的、无辜的人，正如\Place{阿普顿加}的\Person{斐理斯}案件所证明的那样。他们希望\Person{凯基利安}受到那种暴民的指控，正是他们自己向这种暴民让步，对缺席且未经审查的当事人宣判。但可以肯定的是，他们并没有遇到能被说服陷入这种疯狂的审判官。\par

15. 你们的审慎可以使你们在此既注意这些人的顽固，也注意到审判官们的智慧——他们直到最后仍坚持拒绝接受来自\Person{马约里努斯}派群众对\Person{凯基利安}的控告，因为那些人在案件中并无法律资格。你们也会注意到，他们如何被要求提交那些与他们一同从\Place{非洲}前来的人，作为控方或证人，或与案件有其他关联，以及据称他们曾到场，却被\Person{多纳图}撤走。该\Person{多纳图}曾许诺他会把这些人带来，且屡次作此承诺；然而，最终却拒绝再次出现于该法庭之前，而此前他已在庭上承认了那么多，以致看来他拒绝回来只是为了避免在场听到自己被定罪；但是，那些使他定罪的依据却是在他亲自面对且经过审查后，当着他的面被证实的。此外，一些人提交了一份控告\Person{凯基利安}的诉状。随后调查如何重新开始，哪些人提交了诉状，以及最终如何无法证实任何不利于\Person{凯基利安}的事，我无需再说，因为你们都已听过，并可随意反复阅读。\par

16. 至于那次[审判\Person{凯基利安}的]会议中有七十位主教这一事实，你们记得当时是如何借这许多人的可敬权威来指控他的。然而，这些极为可敬的人们（指罗马的审判官）决心不让自己因无望澄清的无尽问题而陷入纠缠，因此他们既不在意那些主教有多少人，也不在意他们从何而来，当他们看见他们被那种鲁莽的假设所蒙蔽，竟对缺席且未经审查的同侪作出草率判决时。然而，那有福的\Person{默基德}本人最终作出的裁决是何等公正、何等周全、何等明智，又多么适宜于促进和平啊！因为他不敢擅自罢免那些未被证明有罪的同侪的职任；他把主要的过失归于\Person{多纳图}，因他查明了此人乃是整场骚乱的根源，于是给予其他所有人悔改回归的机会，并且甚至预备向那些已知由\Person{马约里努斯}祝圣的人发送\OriginalTerm{共融函}{letters of communion}；如此一来，无论何处因这纷争而有两位主教（人数加倍），他决定，按祝圣日期在先的那一位应被确认留任其教座，另一方则另设新的会众。哦，卓越的人啊！哦，基督徒和平之子，基督子民之父！现在请将这一小撮与那一大群主教作一比较，不要只数人数，而要称量他们：一边是节制与审慎，另一边是轻率与盲目。一边是宽仁未曾损害正义，正义也未曾与宽仁相悖；另一边却是恐惧藏于愤怒之下，愤怒被恐惧刺激而过分。在一种情形下，他们聚会是为通过查明真正的罪责所在，使无辜者摆脱诬告；在另一种情形下，他们聚会则是为通过诬告无辜者，来掩护有罪者免受真实控告。\par

% structure: p0024 source=h2 target=subsection reason=relative-heading-rank; base=h2

\subsection{第六章}

\Person{凯基利安}岂能听任这些人审判自己？因为他身边有另一些人，如果在他们面前陈明案情，他便能极轻易地证明自己的清白。即便他是一个新到\Place{迦太基}教会任职的外地人，因而不知道当时一位名叫\Person{路基拉}的极富有的妇人所拥有的蛊惑人心（不论是无赖还是愚人）的能力，他也绝不能将自己交在他们手中；当他还是执事时，曾因执行教会纪律而责备过她，因此得罪了她；因为这一邪恶势力也在作祟，促成了那场不义之举。因为在那个公会议上，那些自称是\OriginalTerm{交付圣经者}{traditor}的人，竟谴责了那些不在场、无辜的人；其中有少数人，希望通过诽谤别人来掩盖自己的罪行，好让人们被无稽的谣言引入歧途，不再追问真相。对这一勾当特别感兴趣的人并不多，尽管优势的权威在他们那边；因为他们有\Person{塞昆德}本人，他因屈服于恐惧而赦免了他们。但据说其余的人都是被\Person{路基拉}的金钱收买和煽动，特别针对\Person{凯基利安}。执政官级人物\Person{泽诺菲洛}手中有一份\WorkTitle{案卷}，根据这份案卷，一个名叫\Person{农迪纳留}的执事（正如我们从同一\WorkTitle{案卷}中得知，他已被\Place{西尔塔}主教\Person{西尔瓦诺}罢免）曾试图通过其他主教的信函举荐自己给那一派，却未能成功，于是盛怒之下揭露了许多秘密，并将其在公开法庭上提出；其中我们在记录中读到这一条：在\Place{非洲}首府\Place{迦太基}的教会里，之所以有人另立祭台，是因为主教们被\Person{路基拉}的金钱收买。我清楚记得我并未向你们宣读过这些\WorkTitle{案卷}，但你们也记得当时没有时间。除了这些影响之外，还有因自尊受损而产生的怨恨，因为他们自己没有祝圣\Person{凯基利安}为\Place{迦太基}主教。\par

当\Person{凯基利安}知道这些人聚集在一起，并非作为公正的法官，而是因这一切而心怀敌意、心存偏见时，他岂能同意，或者他手下的民众岂能允许他，离开教会，进入一处私人住宅，在那里他将不是由同僚公正审判，而是被一个受到女人怨恨驱使的小派别杀害？尤其是，他看出自己的案子本可以在\OriginalTerm{海外教会}{the Church beyond the sea}面前得到公正无私的审理，因为海外教会不受争端双方任何私人恩怨的影响。如果他的对手拒绝在那个法庭辩论，他们便会因此切断自己与整个世界所享有的\OriginalTerm{共融}{communion}，而那正是清白无罪者所享有的。如果他们试图在那里指控他，那么他就会亲自出庭，并针对他们的一切阴谋为自己的清白辩护，正如你们后来所学到的那样：当他们已经犯下\OriginalTerm{裂教}{schism}之罪，并因实际另立祭台而犯下残暴罪行时，才——但为时已晚——请求海外教会裁决。其实，如果他们的理由有真理支持，他们一开始就本该这样做；但他们的策略是，等到虚假谣言随着时间的推移而愈演愈烈，公众的长期传闻可以说已经预先判定了案情之后，再来受审；或者，更可能的是，他们先是随心所欲地定了\Person{凯基利安}的罪，然后依靠他们人多势众以求自保，不敢在其他法官面前展开对这桩如此恶劣案件的讨论，因为那些法官不受贿赂影响，真相可能会被查明。\par

% structure: p0027 source=h2 target=subsection reason=relative-heading-rank; base=h2

\subsection{第七章}

然而，当他们实际发现整个世界与\Person{凯基利安}的共融一如既往，海外教会寄来的\OriginalTerm{共融书信}{letters of communion}是寄给他，而非寄给他们卑鄙地祝圣的那个人时，他们便为一直保持沉默而感到羞耻；因为可能会有人质问他们：为什么他们听任教会在如此众多的国家中继续无知，与已被定罪的人共融？尤其为什么他们因默默容忍他们自己在\Place{迦太基}祝圣的那位主教被排除于共融之外，从而切断了与整个世界的共融，而他们对这个世界却没有提出任何指控？因此，据说他们选择将自己的争端提交给外国的教会，以便达成两件事之一，而这两件事他们都有所准备：一方面，如果他们通过任何诡计成功地使诬告成立，他们就会充分满足自己的复仇欲望；另一方面，如果失败了，他们可以一如既往地顽固，但如今却似乎有了某种借口，声称自己遭受了不公正法庭的对待——这是一切无赖诉讼者的共同喊冤，尽管他们已被真理的最清晰光芒所击败——仿佛不能有人对他们这样说，而且说得极其公正：\Quote{好吧，让我们假设那些在罗马裁决此案的主教们不是好法官；普世教会仍存在一次\OriginalTerm{普世教会的全体会议}{plenary Council of the universal Church}，在这个会议上，这些法官本身也可以为自己辩护；这样，如果他们被证实犯了错误，他们的判决就可能被推翻。}他们是否这样做了，让他们自己去证明：因为我们很容易通过全世界不与他们共融这一事实来证明他们没有这样做；或者，即使这样做了，他们也在那里败诉了，他们与教会分离的状态就是明证。\par

20. 然而，他们后来实际所做的事，在皇帝的信中充分显示了出来。因为他们并不是在别的主教面前，而是在皇帝的法庭上，竟敢对如此崇高权威的教会审判官提出错误判决的指控，正是这些主教们的判决已经宣告了\Person{凯西利安}的无辜，也宣告了他们自己的罪行。皇帝准许了在\Place{阿尔勒}、在其他主教面前进行第二次审判；这并不是因为这是他们应得的，而只是出于对他们顽固的让步，并希望不惜一切抑止如此巨大的无礼。因为这位基督徒皇帝并没有擅自应允他们那无理取闹、毫无根据的投诉，使自己成为对在\Place{罗马}开庭的主教们所颁布判决的审判官；相反，正如我所说的，他任命了其他主教，然而，他们却又宁愿向皇帝本人上诉；你们也已经听到了他对这种行为的措辞。但愿他们当时就停止他们那极其疯狂的争辩，而最终屈服于真理，如同他（后来打算为此做法向尊敬的主教们道歉时）在他们之后审理他们的案件一样，条件是他们若不服从他本人的判决——而这是他们自己提出上诉的——今后就当缄默不语！因为他命令双方当事人在\Place{罗马}会面，进行辩论。当\Person{凯西利安}由于某种原因未能到场时，皇帝应他们的请求，命令所有人随他去\Place{米兰}。那时，他们团体中的一些人开始退出，或许是不满于\Person{君士坦丁}没有仿效他们榜样，立即草率地给缺席的\Person{凯西利安}定罪。当谨慎的皇帝得知此事后，他迫使其余的人在卫兵押送下来到\Place{米兰}。\Person{凯西利安}抵达后，皇帝亲自提出他，正如他在信中所写的；并以他的书信所表明的勤勉、谨慎和审慎，对此事进行了审查，宣告\Person{凯西利安}完全无辜，而他们则是罪大恶极。\par

% structure: p0030 source=h2 target=subsection reason=relative-heading-rank; base=h2

\subsection{第八章}

21. 直到今天，他们还在教会的共融之外施行圣洗，而且，只要有可能，他们就给教会的成员重施圣洗：他们在分裂与不和中祭献，并以和平的名义问候那些他们宣称已处于救恩和平界限之外的团体。基督的合一遭到撕裂，基督的产业受到辱骂，基督的圣洗遭到蔑视；而他们拒绝让已设立的人间权威纠正这些错误，尽管这些人间权威施加暂时的惩罚，正是为了避免他们因如此亵圣而遭永罚。我们谴责那驱使他们走向分裂的狂怒，谴责那使他们重施圣洗的疯狂，也谴责他们脱离遍满全球的基督产业之罪。我们所用的抄本，也和他们手中的一样，其中提及一些教会，这些教会的名称他们现在也在诵读，但他们如今与这些教会已无共融；当在他们的集会中宣读这些书信时，他们对读经者说：\Quote{愿你们平安；}然而，他们与这些书信所写给的人之间却并无和平。他们则把已故之人的罪行归咎于我们，提出要么虚假、要么即便是真实也与我们无关的指控；他们没有觉察到，在我们指摘他们的事情上，他们全都受到了牵连，但在他们指摘我们的事情上，罪责应归于主的庄稼中的糠秕或莠子，而罪行并不属于好籽粒；再者，他们也不考虑，在我们的合一之内，只有那些喜欢恶人的恶行的人，才与恶人相交，而那些厌恶其恶行却又不能纠正它们的人——因为他们不敢在收割之前就拔除莠子，以免连麦子也拔了出来，\BibleRef{玛 13:29}——他们与恶人相交，并不在其行为上，而是在基督的祭台上；这样，他们不仅避免被其玷污，而且按照天主的圣言，他们值得称赞与颂扬，因为，为了防止基督之名因可憎的分裂而受辱骂，他们为了合一而容忍那他们因公义而憎恶的事。\par

22. 如果他们能听，就让他们听圣神向各教会所说的话吧。因为在\BibleBook{若望}{默示录}中我们读到：\BibleQuote{给\Place{厄弗所}教会的天使写：\InnerQuote{那右手握着七颗星，而在那七盏金灯台当中行走的这样说：我知道你的作为、劳苦和坚忍；也知道你不能容忍恶人，并且你也曾查验那些说自己是宗徒而其实不是，发现他们是撒谎的人；你也有坚忍，为了我的名字容忍了他们，而毫不厌倦。}}\BibleRef{默 2:1} 现在，如果主希望这话是写给一位天上的天使，而非写给教会中掌权的人，他就不会接着说：\BibleQuote{可是我有一件事要反对你，就是你抛弃了起初的爱德；所以你该回想你是从那里跌下的，你该悔改，行起初所做的事；不然，我就要临到你那里，把你的灯台从原处挪去，除非你悔改。}\BibleRef{默 2:4} 这番话不可能向天上的天使说，因为他们始终保持爱德不变，因为天使中唯一离弃并失去爱德的，只有魔鬼及其使者。这里所说的\Quote{起初的爱德}，原存在于他们为基督之名而容忍假宗徒的计划中。主命令他们回归这爱德，并做\Quote{起初所做的事。}现在，我们因恶人的罪行而受责难，这些罪行并非我们所为，而是别人所做；而且，其中有些甚至是我们所不知道的。然而，即使这些罪行确实被犯下，而且是在我们眼前犯下，并且我们为了合一的缘故而容忍它们，因着麦子的缘故而任凭莠子不去拔除，凡以敞开的心领受圣经的人都会宣判我们不仅无可指摘，而且值得不小的称赞。\par

23. 亚郎容忍了那些要求、制造并崇拜偶像的群众。梅瑟容忍了成千上万抱怨天主、屡次冒犯祂圣名的人。达味容忍他的迫害者撒乌耳，甚至当他因邪恶的生活离弃上天之事，追求巫术等地下之事时，仍为他报仇，并因他受祝圣的尊贵权利而称他为\Quote{主的受傅者}。撒慕尔容忍了厄里的恶子们，和他自己乖戾的儿子们；人民不肯容忍他们，因此受到天主的警告并被天主的严厉所惩罚。最后，他容忍了这个民族本身，尽管他们骄傲并轻视天主。依撒意亚容忍那些他向他们发出许多应得谴责的人。耶肋米亚容忍那些使他遭受许多痛苦的人。匝加利亚容忍那些经师和法利塞人，关于这些经师和法利塞人当时的品性，圣经告诉了我们。我知道我遗漏了许多例子：让那些愿意且有能力的人自行阅读神圣的纪录：他们会发现，所有圣善的天主仆人和朋友，都常常不得不容忍自己民族中的某些人；虽然与这些人一同领受那时代的圣事，但这样做不但没有使他们受到玷污，反而因他们的容忍精神堪受赞扬，\BibleQuote{竭力保持}，如宗徒所说，\BibleQuote{圣神在和平联系中的合一}。\BibleRef{弗 4:3} 让他们也注意主来临后所发生的事，在那时期，如果我们能把所有实例都写下来并证实，我们会在世界各地找到更多这种容忍的例子：但请留意那些有记录的。主自己容忍犹达斯，他是个魔鬼，是个盗贼，是出卖祂的人；祂允许他与无罪的众门徒一起，领受信友们所知的我们的赎价。宗徒们容忍假宗徒；在那些追求自己的事，而不追求耶稣基督的事的人中间，保禄不追求自己的事，而追求基督的事，活出了极高尚的容忍。最后，正如我不久前提到的那位以\Quote{天使}之名管理一个教会的人，他受到称赞，因为他虽然憎恨恶人，但仍为了主名的缘故容忍他们，即使他们已被查验和揭发。\par

24. 最后，让他们扪心自问：他们难道不也容忍那些由\OriginalTerm{周游派}{Circumcelliones}所行的谋杀与纵火吗？这些周游派尊重那些从危险高处跳下的人的尸体。他们难道不也容忍那使整个\Place{阿非利加}多年呻吟于一个人的惊人暴行之下的苦难吗？这个人就是\Person{奥普塔特}\EditorialNote{塔穆加迪的主教}。我不详述遍布阿非利加各地区、城镇和财产中的暴政和公开掠夺；因为最好还是由你们彼此谈论这些事，无论是低声私语还是公开议论，随你们所愿。因为无论你们望向何处，都会看到我所说的，或更确切地说，我忍住不说的事。而且我们并不因此而谴责你们所爱的那些人，当他们做这些事的时候。我们厌恶那一派人的，不是他们对恶人的容忍，而是他们在分裂的事上不可容忍的邪恶，即祭台对立祭台，以及分离于基督的产业，这产业正如很早所预许的，已遍布全世界。我们悲叹哀悼和平被破坏，合一被撕裂，圣洗被再次施行，圣事遭到轻慢，而这些圣事即使由恶人施行和领受也是神圣的。如果他们视这些为小事，就让他们看看那些已证明天主如何估量这些事的例子。那些制造偶像的人死于普通的死亡，被刀剑杀死：\BibleRef{出 32:27}。但当有人企图在以色列中制造分裂时，领袖们被地裂开口吞没，同伙的群众被火烧灭。从惩罚的差异中，可以看出罪过的不同等级。\par

% structure: p0035 source=h2 target=subsection reason=relative-heading-rank; base=h2

\subsection{第九章}

25. 那么，事实就是：在教难时期，圣书被交给迫害者。那些犯此交付之罪的人承认了，并被移交到天主的法庭；那些无辜者未经审问，就被轻率之人立即定罪。在所有这些被缺席定罪的人中，受指控最激烈的那一位的正直，后来在无可指责的法官面前得到了昭雪。人们从主教们的决定上诉至皇帝；皇帝被选为法官；而当皇帝作出判决时，却被蔑视。当时所发生的事你们已经读过；现在正在发生的事你们亲眼所见。如果在读完所有这些之后，你们仍有怀疑，就被你们所见的事说服吧。我们完全不必再依据古代手稿、公共档案或民事及宗教法庭的记录来争论了。我们有一本更大的书——世界本身。在这本书中，我读到应许的实现，我在天主的书中读到这应许：\BibleQuote{主对我说：你是我的儿子，我今日生了你。你向我请求，我必将外邦人赐你作产业，将地极赐你作领地。} 那不与这产业相通的人，无论他搬出什么书卷，可知自己已被剥夺继承权。那攻击这产业的人，分明被宣示为天主的家庭所弃绝。问题是关于将其中应许了这产业的天主圣书交出的罪犯。那抗拒立遗嘱者意图的人，就应被相信已将遗嘱付诸一炬。多纳特派啊，格林多教会有何得罪你们之处？提到这一个教会时，我希望被理解为也在问远离你们的其他类似教会同样的问题。这些教会对你们做了什么？它们甚至不可能知道你们所做的事，或你们加以定罪者的名字。或者，难道因为\Person{切齐里亚努斯}得罪了\Place{阿非利加}的\Person{卢西拉}，基督的光就离开全世界了吗？\par

26. 让他们最终觉悟到自己的所作所为；因为时光流逝，由于公义的报应，他们的工程反过来击打了自己。请问，是哪个女人的煽动，使\Person{玛克西米安}（据说是\Person{多纳特}的亲属）脱离了\Person{普里米安}的\OriginalTerm{共融}{communion}，又是怎样，在聚集了一群\OriginalTerm{派系}{faction}主教之后，他在\Person{普里米安}缺席时对他宣判，并让自己受\OriginalTerm{祝圣}{ordination}为敌对的主教——正如\Person{马约里诺}受\Person{卢西拉}的影响，聚集了一群主教，在\Person{凯基利安}缺席时定他的罪，并受祝圣为主教以反对他。我假定你们会承认，当\Person{普里米安}被在\Place{阿非利加}与他共融的其他主教们，从\Person{玛克西米安}派系所作宣判中解救出来时，这一决定是有效且充分的吧？那么你们会拒绝承认\Person{凯基利安}被海外同一个教会的主教们，从\Person{马约里诺}派系所作宣判中释放时，情况也是一样吗？我的弟兄们，我请问你们，我向你们所求的是什么大事？理解我摆在你们面前的事，有什么困难呢？\Place{阿非利加}教会，若与世上其他地区的教会相比，与它们大不相同，且在人数和影响上都远远落后；而且即使它保持了合一，若与分布在世界其他各民族中的普世教会相比，也还是远远更小——这差距甚至大于\Person{玛克西米安}派系相比\Person{普里米安}派系时的差距。然而，我所求的只是这一点——而且我相信这是公道的——即：你们不要认为\Place{提吉西斯}的\Person{塞昆杜斯}所召集的会议，就是\Person{卢西拉}煽动起来，反对缺席的\Person{凯基利安}，并反对与\Person{凯基利安}共融的一个\OriginalTerm{宗座}{apostolic see}和全世界的那个会议，比\Person{玛克西米安}的会议更具分量；后一个会议也是以同样的方式，由另一个女人煽动起来，反对缺席的\Person{普里米安}，并反对在\Place{阿非利加}全境与他共融的其余民众。还有什么情况比这更清楚的呢？还有什么要求比这更公道的呢？\par

27. 你们看见并知道这一切，也为此叹息；然而天主同时也看见，没有什么迫使你们停留在这般致命而不虔敬的\OriginalTerm{裂教}{schism}中，只要你们肯克制肉欲以赢得属灵的国，并为了逃避永罚，勇于失去世人的友谊，因为人的好感在天主的审判台前毫无益处。去吧。现在你们一同商议：针对我所写的，找出你们能回答的话来。如果你们拿出支持己方的手抄文件，我们也会这样做；倘若你们的人说我们的文献不可信，那么当我们以同样的话回敬时，也请不要见怪。没有人能从天上抹去天主的法令，也没有人能从地上抹去天主的教会。祂的法令已许给全世界，而教会已充满了世界；它包括善人和恶人。在地上，它只失去恶人；进入天上，它只接纳善人。\par

在我撰写这篇讲论时，天主是我的见证，我是怀着何等真诚的和平之爱与对你们的爱，取用并运用了祂所赐给我的。如果你们愿意，它将成为你们改过的途径；如果不愿意，它也将成为反对你们的见证。\par

\SourceRecord{Translated by J.G. Cunningham.；Nicene and Post-Nicene Fathers, First Series；Vol. 1.；Edited by Philip Schaff.；Buffalo, NY: Christian Literature Publishing Co.,；1887.；<http://www.newadvent.org/fathers/1102043.htm>.}

% structure: p0001 source=h1 target=section reason=page-title

\section{书信四十四（主后398年）}

\Emph{致我最敬爱的主内弟兄，及堪受一切赞美的弟兄们：\Person{埃莱乌修斯}、\Person{格劳里乌斯}与两位\Person{费利克斯}，\Person{奥古斯丁}问候你们。}\par

% structure: p0003 source=h2 target=subsection reason=relative-heading-rank; base=h2

\subsection{第一章}

在前往\Place{基尔塔}的教会的途中经过\Place{图布尔西}时，尽管时间紧迫，我拜访了你们那里的主教\Person{福图尼乌斯}，发现他确实如你们通常亲切地让我期待的那样。当我通知他你们与我关于他的谈话，并表示希望见他时，他并未拒绝这次拜访。因此我去见他，因为我认为由于他的年长，我应该去见他，而不是坚持让他先来见我。于是，我在相当多的人的陪同下前往，这些人当时恰巧在我身边。然而，当我们在他家中坐下后，事情传开，人群大增；但在那整个众人中，在我看来，只有很少的人希望以健全有益的方式讨论此事，或以如此重大问题所要求的深思和庄重。所有其他人都更像是观众，期待看我们辩论的一场戏，而不是怀着基督徒的严肃精神，寻求关于救恩的教导。因此，当我们发言时，他们既不能以沉默表示好感，也不能谨慎发言，甚至不能适当地注意礼节和秩序——除了，如我所言，那些少数虔诚而真诚关心此事的人，对这一点毫无疑问。因此，一切都被人们大声喧哗的声音搅乱，每个人都受自己情绪无节制的冲动支配；尽管\Person{福图尼乌斯}和我都恳求和规劝，我们完全无法说服他们安静地听。\par

尽管如此，问题的讨论还是开始了，我们坚持了几个小时，双方轮流发言，只要偶尔能从喧闹观众的叫嚷声中得到片刻安宁。在辩论之初，意识到说过的话容易被我自己或那些我深切关心其救恩的人遗忘；同时也渴望我们的辩论能以谨慎和自制的方式进行，并能让你们以及其他不在场的弟兄能从记录中得知讨论的内容，我要求由记录员记下我们的话。这一点长时间遭到\Person{福图尼乌斯}或他那一方的反对。不过最终他同意了；但在场的记录员，那些能胜任这项工作的人，却出于某种我不知道的原因拒绝记录。我催促他们，至少让陪同我的弟兄们，尽管不太擅长这项工作，来记录，并承诺我会将记有笔记的写字板交给对方。这一点被同意了。我的一些话先被记录下来，对方的一些陈述也被口述记录。之后，记录员无法忍受对方嘈杂的打断，以及在这种压力下我方辩论变得更加激烈，便放弃了任务。然而，这并没有结束讨论，每个人仍争取机会说了很多。我所挚爱的朋友们，我已决定不让你们错失对整个问题的讨论，或至少我所记得的所有内容；你们可以将这封信读给\Person{福图尼乌斯}听，让他要么确认我的陈述属实，要么他自己毫不犹豫地告诉你们任何他更准确的记忆所提示的内容。\par

% structure: p0006 source=h2 target=subsection reason=relative-heading-rank; base=h2

\subsection{第二章}

他高兴地开始赞扬我的生活方式，他说他是通过你们的陈述（我确信其中善意多于真实）得知的，并补充说，他曾向你们评论道，如果我是在教会内做你们告诉他的关于我的一切事情，我可能已经做得很好。我于是问他，人应当这样生活的教会是哪一个：是那个如圣经早已预言的、遍布全世界的教会，还是那个由少数非洲人或非洲的一小部分所包含的教会。对此，他起初试图回答说，他的共融团体遍布全球。我问他能否向我可能选择的地方颁发我们称为\OriginalTerm{正规的}{regular}的通功信；我肯定，这在所有人看来都是显而易见的，这样问题就能最简单地解决。如果他同意这一点，我的意图是，我们将这样的信函发送给那些我们根据宗徒的权威，都知道在他们的时代就已经建立的教会。\par

4. 然而，他的陈述虚假之本质既已明显，他便在一阵含糊言辞的烟幕中匆忙撤退，其间引用了主的话语：\BibleQuote{你们要提防假先知! 他们来到你们跟前，外披羊毛，内里却是凶残的豺狼。你们可凭他们的果实辨别他们。}\BibleRef{玛 7:15} 当我说这些主的话语也可由我们应用于他们时，他接着夸大了他声称他的党派经常遭受的迫害；意图以此证明他的党派是基督徒，因为他们忍受了迫害。当我准备就此事从福音中回答他时，他却抢先一步，引出了主所说的话：\BibleQuote{为义而受迫害的人是有福的，因为天国是他们的。}\BibleRef{玛 5:10} 感谢他恰当的引述，我立刻补充说，因此需要查问，他们是否确实为义而受了迫害。在继续追问此事时，我希望弄清，尽管事实上众人皆知，这些迫害是在他们属于教会的合一之内时，还是在因\OriginalTerm{裂教}{schism}而脱离之后，才在\OriginalTerm{麦加里教难}{Macarian persecution}中临到他们；以便让那些想要知道他们是否曾为义受迫害的人，可以先考虑先前的问题，即他们脱离全世界的合一是否做得正确。因为如果发现他们在这件事上犯了错，那显然他们是为不义而非为义受迫害，因此不能被算入那些被称为\BibleQuote{为义而受迫害的人是有福的}之列。随后，提到了交付圣书的事，关于此事所说的远多于所证实的事实。我们方面则回答说，\OriginalTerm{交付经书者}{traditors}是他们而非我们的领袖；但如果他们不相信我们用以支持此指控的文件，我们也无法被迫接受他们所提出的那些。\par

% structure: p0009 source=h2 target=subsection reason=relative-heading-rank; base=h2

\subsection{第三章}

5. 因此，我们把那个有争议的问题暂且搁置，然后问道，他们如何能为自己脱离所有其他基督徒的行为辩护；那些基督徒并没有伤害过他们，他们在全世界保存着继承的秩序，建立在最古老的教会中，却完全不知道谁是非洲的\OriginalTerm{交付经书者}{traditors}；而且他们肯定只能与他们听说占据主教座的人保持共融。他回答说，外方教会直到赞同那些在\OriginalTerm{麦加里教难}{Macarian persecution}中受苦者的死亡之时，才伤害了他们。在此，我本可以说，外方教会的清白不可能因麦加里时代所犯的罪行而受损，因为无法证明他所行的是得到了他们的认可。然而，为节省时间，我宁愿问道：即使假设外方教会因麦加里的残酷，从据说他们赞同这些行为之时起便失去了清白，是否也能证明直到那时，\OriginalTerm{多纳特派}{Donatists}仍与东方教会和世界其他地方保持着合一。\par

6. 于是，他拿出某卷书，想要表明\Place{撒底迦}大公会议曾致函属于\Person{多纳特}一派的非洲主教。当这信被朗读时，我听到收信主教的名字中有\Person{多纳特}。因此我坚持要知道，这是否就是他们的派别得名的那位\Person{多纳特}；因为很可能他们写信给的是属于另一宗派〔异端〕、名叫\Person{多纳特}的主教，尤其既然这些名字中并未提及非洲。那么，我问道，我们怎能证明必须相信这里提到的\Person{多纳特}就是\OriginalTerm{多纳特派}{Donatist}主教，既然甚至无法证明这封信是专门寄给非洲主教的？因为虽然\Person{多纳特}是一个常见的非洲名字，但设想在其他国家也可能有人取非洲名字，或者一个非洲出身的人在那里成为主教，这不是不可能。此外，我们在信中没找到执政官的名字或日期，如果比较日期，本可提供某些明确的线索。我确实曾听人说过，\OriginalTerm{亚略派}{Arians}在与公教共融分离之后，曾试图将非洲的多纳特派与他们自己结盟；我的兄弟\Person{阿利比乌斯}在这时低声提醒了我。于是我拿起那卷书，浏览所说的大公会议的谕令，看到\Place{亚历山大里亚}的公教主教\Person{亚大纳削}，他在与亚略派的激烈辩论中声名卓著，以及\Place{罗马}教会的主教\Person{犹利乌斯}，也是一位公教徒，竟被那个撒底迦大公会议定罪；由此我们确信那是亚略派的大公会议，而反对这些异端者的正是这些公教主教，他们曾以非凡的热诚与之抗争。因此，我本想拿起并带走那卷书，以便更仔细地查明该大公会议的日期。然而他拒绝了，说如果我愿意研究其中的任何内容，可以在那里查看。我也请求他允许我在那卷书上作标记；因为我承认，我担心万一将来需要查阅时，会有另一卷替代它。这一点他也拒绝了。\par

% structure: p0012 source=h2 target=subsection reason=relative-heading-rank; base=h2

\subsection{第四章}

7. 随后，他开始坚持要我明确回答这个问题：我认为迫害者与被迫害者哪一方正确？我回答说，这个问题并不公平：可能双方都错了，或者迫害可能来自两方中更为正义的一方；因此，不能总是推断说，受苦的一方就是更好的，尽管事实几乎总是如此。当我察觉他仍然对此极为强调，希望因他们受过迫害，就承认他们一方的正义不容置疑时，我便问他，他是否认为\Place{米兰}教会的\Person{安博罗削}主教是义人和基督徒？他被迫明确否认此人是基督徒和义人；因为如果他承认了，我就会立即反对他说，他认为\OriginalTerm{再洗礼}{rebaptism}是必要的。因此，当他被迫宣布\Person{安博罗削}既非基督徒也非义人时，我便叙述了他在教堂被士兵包围时所遭受的迫害。我还问他，在他们看来，曾在\Place{迦太基}从他们一党中制造分裂的\Person{马克西米安努斯}，是不是义人和基督徒？他不得不否认。于是我就提醒他，那人曾遭受了这样的迫害，以致他的教堂被夷为平地。我极力通过这些例子，如果可能，说服他不要再坚持认为，遭受迫害是基督徒义德最确实无误的标记。\par

8. 他还叙述说，在他们分裂的初期，他的前辈们急于设法遮掩\Person{凯西利安}的过失，以免发生分裂，就在\Person{马约里努斯}被祝圣以对抗\Person{凯西利安}之前，在\Place{迦太基}，为他教会团体中的人任命了一位临时主教。他声称这位临时主教被我们这一派在他自己的聚会所里杀害了。我承认，我以前从未听说过这事，尽管他们对我们的许许多多指控已被驳斥和推翻，而我们对于他们则提出了更大更多的罪行。在叙述了这个故事之后，他又开始坚持要我回答，在这种情况下，我认为杀人者与被害者哪一方更义；好像他已经证实事情就如他所陈述的那样发生了。于是我说道，我们必须先查明这个故事的真实性，因为我们不应该不加审查就相信所听到的一切；而且即使是真的，也可能双方都一样坏，或者一个恶人使另一个比他更坏的人丧了命。因为，事实上，为整个人\OriginalTerm{再洗礼}{rebaptism}者的罪过，可能比仅仅杀害身体者的罪过更为可憎。\par

9. 此后，他没有提出后来向我提出的那个问题。他断言，即使是恶人，也不该被基督徒和义人所杀；好像我们把那些在天主教会中做这些事的人称为义人似的：此外，这个论断，他们宣称起来容易，向我们证明却难，因为，除了少数人以外，他们自己，包括主教、长老和各类圣职人员，不断聚集最狂热之徒，在力所能及之处，制造了许多暴行和破坏，不仅伤害公教徒，有时甚至伤害他们自己的同伙。尽管有这些事实，\Person{福尔图尼乌斯}却假装不知道这些他比我知道得更清楚的恶行，坚持要我举出一个义人甚至处死恶人的例子。这当然与当前问题无关；因为我承认，但凡有基督徒名义的人犯了这种罪行，都不是善人的行为。然而，为了向他指明我们面前的真正问题是什么，我反问他，\Person{厄里亚}在他看来是不是义人；他不得不表示同意。于是我提醒他，\Person{厄里亚}亲手杀了多少假先知。\BibleRef{列上 18:40} 他清楚地看出，事实上他不可能不看出，那时义人做这事是合法的。因为他们行这事，是作为先知，受圣神的引导，并蒙天主的权威所认可；天主无误地知道，对谁来说，被处死甚至可能有益。因此，他要求我向他指出，在新约时代，有哪个义人曾处死过任何人，哪怕是罪犯和不虔敬的人。\par

% structure: p0016 source=h2 target=subsection reason=relative-heading-rank; base=h2

\subsection{第5章}

10. 然后我回到前一封信中使用的论证，其中我努力表明，我们不应以他们中某些人所犯的暴行来指责他们，他们也不应以他们发现我方可能犯下的此类行为来指责我们。因为我承认，在新约中找不到义人处死任何人的例子；但我坚持，借着我主自己的榜样，可以证明恶人曾被无辜者所容忍。因为祂自己的出卖者，就是那已经收了祂血价的人，祂容忍他与同祂在一起的清白者不加分别，直到那最后的平安之吻。祂并未向门徒们隐瞒他们中间有一个能犯如此罪行的人；然而，祂并未排斥这叛徒，而是将祂体血的第一件圣事同样分施给了他们所有人。\BibleRef{玛 26:20} 当几乎所有人都感受到这个论证的力量时，\Person{福尔图尼乌斯}试图用以下说法来应对：在主受难之前，与恶人共融对宗徒们并无害处，因为他们那时还没有\OriginalTerm{基督的圣洗}{baptismus Christi}，只有\OriginalTerm{若翰的洗礼}{baptismus Ioannis}。当他这么说时，我请他解释，经上如何记载耶稣施洗比若翰还多，虽然耶稣本人并不施洗，而是祂的门徒们施洗，也就是说，借着祂的门徒们施洗？\BibleRef{若 4:1} 他们怎能给出他们未曾领受的呢（这是\OriginalTerm{多纳徒派}{Donatistae}自己也常使用的问题）？难道基督是用若翰的洗礼施洗吗？我准备就\Person{福尔图尼乌斯}的这个观点提出许多其他问题；例如——若翰本人如何被问及关于主施洗的事，而他回答说，祂有新娘，并且是新郎？\BibleRef{若 3:29} 那么，新郎能用不过是朋友或仆人的洗礼来施洗吗？再者，如果没有先受洗，他们怎能领受圣体呢？或者，那样的话，主怎能回答\Person{伯多禄}（伯多禄愿意被主完全洗涤）说：\BibleQuote{沐浴过的人，已全身清洁，只需洗脚就够了。}\BibleRef{若 13:10} 因为完全的洁净是出于圣洗，不是若翰的，而是主的，如果领受者配得的话；然而，如果他不配，圣事留在他内，不是为了他的得救，而是为了他的丧亡。当我正要提出这些问题时，\Person{福尔图尼乌斯}自己看出，他不该提起关于主门徒们的洗礼这一话题。\par

11. 从这一点我们转向了别的问题，双方许多人都尽其所能地发言。在其他的言论当中，有人声称我方仍有意迫害他们；\Person{福尔图尼乌斯}则说，他很想看看在这种迫害发生时我将如何行动，我是否会赞同这种残忍，还是完全不予支持。我说，天主鉴察我的心，这是他们看不见的；而且，他们迄今并无理由担心这样的迫害；即便真的发生，那将是坏人所为，然而这些坏人至少不像他们自己派中的某些人那样坏；但是，我们不应当因任何违反我们意愿、甚至不顾我们反对（如果我们有机会公开抗议的话）而做出的事情，就退出与\OriginalTerm{天主教会}{Ecclesia Catholica}的共融，因为我们已学会了宗徒为和平所规定的容忍，如他所说的：\BibleQuote{以爱德彼此担待，尽力以和平的联系，保持心神的合一。}\BibleRef{弗 4:2} 我断言，他们并没有保持这种和平与担待，因为他们制造了分裂，而且在分裂当中，他们中较为温和的人如今容忍了更严重的恶事，以免那已破碎的再被打破，尽管他们为了保持合一，却不愿意在较小的事上行使担待。我也说，在旧约体制下，合一与担待的和平并未像如今借着主的榜样和新约的爱德那样被充分宣示和称扬；然而先知和圣者们常常抗议人民的罪过，却不试图将自己与犹太民族的合一分开，也不脱离与他们一同参与当时所制定的圣事。\par

12. 此后，我不知出于何故，有人怀着敬爱之意追忆已故的\Place{迦太基}宗座\Person{奥勒留}的前任\Person{格内特利乌斯}，因为他压制了某项针对\OriginalTerm{多纳徒派}{Donatists}的法令，并不容许其生效。他们都以极大的善意称赞他、推崇他。我打断了他们的赞美之辞，说尽管如此，倘若\Person{格内特利乌斯}本人落到他们手中，他们定会宣称有必要为他施行第二次洗礼。（此时我们都已站着，因为离去的时间快到了。）对此，那位老人坦言道，现在已经制定了一条规则，凡从我们这边转到他们那边的信徒，都必须受洗；但他说这话时，带着极明显的勉强与真诚的惋惜。当他十分坦率地哀叹他一方许多恶行，并说明——这点也由整个团体的见证进一步证实——他如何远未参与这些事，并告诉我们他惯于向己方人士温和规劝的话；随后我又引用了\Person{厄则克耳}的话——\BibleQuote{父亲的生命怎样属于我，儿子的生命也怎样属于我：谁犯罪，谁就当死亡。}\BibleRef{则 18:4}——根据这节经文所说，儿子的过错不归於父亲，父亲的过错不归於儿子，众人一致同意，在这类讨论中，任何一方都不应将恶人的过分行为提出来反对对方。因此，剩下的只有关于分裂的问题。所以我再三劝勉他，应当以平静无扰的心与我一起努力，藉着勤奋的探究，使如此重要之事的审查得以善终。他友善地回答说，我自己是在以纯真的眼目寻求此事，但我这边的其他人却反对这样探究真理；我便带着这一承诺离开他：我将带给他更多的同工，至少十位，他们愿意以同样的善意、平静与虔诚之心，来仔细审查这个问题，正如我看到他此时在我身上发现并称赞的那种态度。他也给了我类似的承诺，邀请同样数量的己方同工。\par

% structure: p0020 source=h2 target=subsection reason=relative-heading-rank; base=h2

\subsection{第六章}

13. 因此，我奉主的宝血劝勉并恳求你们，要提醒他遵守承诺，并催促他竭尽全力，使那业已开始，且如你们所见，已接近完成的事，得以结束。因为依我看来，你们在你们的主教中很难再找到另一位，其判断与情感都如此健全，如同我们在这位老人身上所见的一样。第二天，他亲自来见我，我们又开始讨论这事。然而，我不能与他久留，因为祝圣一位主教的事要求我离开那地方。我先前已派了一位信使去见\OriginalTerm{崇天派}{Cœlicolæ}的首领——我曾听闻他在他们中间引入了一种新的洗礼，并以此不虔敬之举误导了许多人——我打算在有限的时间内与他会谈。\Person{福图尼乌斯}得知这人要来，察觉到我另有要事，而他自己也有其他职责召唤他离开家，便亲切友好地向我辞别。\par

14. 在我看来，倘若我们要避免喧闹的人群——他们更干扰而非有助于辩论——并渴望藉着主的助佑，完成那业已在真诚的善意与和平精神中开始的伟业，我们就应当在某个小村庄会面，那里双方都没有教会，且居住着属于双方教会的人，比如\Place{提提亚}。就让这个地方，或其它类似的地方，于\Place{图布尔西}或\Place{塔加斯特}地区约定下来；并且我们要备好圣经书卷以备参考。任何一方认为有用的文献也可带至此处；搁下一切其它事务，若蒙天主悦纳，不受其它挂虑干扰，我们将尽可能多的日子专心于这惟一的工作，各自私下祈求主的引导，仰赖那位以基督徒和平为最甘饴者的助佑，使这已在如此良善精神中开启的探究，得以圆满结束。请不要忘记回信，将你或\Person{福图尼乌斯}对此的想法写信告诉我。\par

\SourceRecord{Translated by J.G. Cunningham.；Nicene and Post-Nicene Fathers, First Series；Vol. 1.；Edited by Philip Schaff.；Buffalo, NY: Christian Literature Publishing Co.,；1887.；<http://www.newadvent.org/fathers/1102044.htm>.}

% structure: p0001 source=h1 target=section reason=page-title

\section{第46封信（公元398年）}

\EditorialNote{一封提出若干良心问题的信。}\par

\Emph{致我所敬爱且可敬的父亲、主教\Person{奥古斯丁}，\Person{普布利科拉}敬上。}\par

经上记载：\BibleQuote{询问你的父亲，他必指示你；询问你的长老，他们必告诉你。}\BibleRef{申 32:7}因此我判断理当\Quote{向司铎的口中求问法律}，就我在本信中所要陈述的某个案例而言，同时渴望就其他若干事项蒙受指教。我将各个问题分开陈述，在信中各以单独段落说明，并恳请你就每一问题逐一赐复。\par

一、在\Place{阿祖格}地区，据我所闻，蛮族惯常指着他们的假神发誓，当着驻守边境的\OriginalTerm{十夫长}{decurion}或\OriginalTerm{护民官}{tribune}的面，当他们受雇往某处搬运辎重，或保护庄稼免遭劫掠时；而当十夫长书面证明此誓言已立之后，地主或农夫便雇用他们看守庄稼；或有需要途经其地的旅客也雇用他们，仿佛确信他们现已可靠。如今我心中生疑：如此雇用蛮族的地主，因着那誓言而确信其忠诚，是否使自己以及托付那人管理的庄稼，沾染了那罪孽誓言的污秽呢？而那或许雇用其服务的旅客，是否也是如此呢？然而，我当说明，在这两种情形下，蛮族都因服役而获金钱酬报。不过，在这两项交易中，除了金钱酬报之外，尚有在十夫长或护民官面前所立的这包含大罪的誓言。我所关切的是，这罪是否既污秽了接受蛮族誓言的人，又至少污秽了托付蛮族保管之物。因为无论协议中尚有其他何种条款，哪怕是支付黄金、提供人质作为担保，这罪孽的誓言仍是此项交易的真实部分。敬请你明确而肯定地解除我的疑惑。盖若你的答复表明你自身尚且犹疑，我便可能陷入较前更深的困惑。\par

二、我也听说，我自己的管家从受雇保护庄稼的蛮族那里，接受他们向假神祈求的誓言。这誓言岂不如此污秽了庄稼，以致若基督徒使用它们，或收取其售卖所得之钱，他本人便受玷污吗？请回答此问。\par

三、再者，我听见有人言，蛮族与我的管家立约时并未发誓；又有人告诉我，确曾发过这样的誓。假设后一种说法不实，请告诉我，我是否仅仅因为听闻此说，便应依照经上的规定，放弃使用这些庄稼或其所得之钱：\BibleQuote{若有人向你们说：\InnerQuote{这是祭献过偶像的肉}，你们就不可吃，为了那指证的人。}\BibleRef{格前 10:28}这案例是否与祭过偶像的肉食案例类似？若是，我当如何处置这些庄稼，或其售价呢？\par

四、在这情形下，我是否应当盘问那名说管家面前未立誓的人，以及另一名说确曾立誓的人，并找人作证，以辨明二人中谁所言属实，而在事情未明之前，暂且不动那些庄稼或其售价呢？\par

五、倘那发此罪孽誓言的蛮族要求管家，或驻守边境的护民官（身为基督徒者），以他自己所发、包含大罪的同样誓言，向他保证在看守庄稼的约定上守信，那么这誓言是否只污秽了那名基督徒呢？是否也污秽了他指以立誓的那些事物呢？又或若有异教徒在边境掌权，以此誓言向蛮族表示诚信，他是否使为之立誓的诸方均沾染其罪污呢？若我派人往阿祖格去，他能否合法地从蛮族那里接受那罪孽的誓言？那从蛮族接受此种誓言的基督徒，岂不也因他的罪受玷污吗？\par

六、倘若基督徒明知从打谷场取来的麦子或豆类，或从榨房取来的酒或油，其中一部分已被祭献于假神，他可否合法地使用呢？\par

七、基督徒可否将明知取自偶像丛林的木材用于任何用途呢？\par

八、若基督徒在市场上买了未祭过偶像的肉，心中却疑虑交加，不知是否祭过偶像，但最终认定它未祭过，他若吃了这肉，是否犯罪呢？\par

九、若有人行本身为善的事，却对该事之善恶有所疑虑，他若相信其为善而践行之，尽管先前他可能以为那是恶的，这能算为他的罪吗？\par

十、倘有人妄称某肉曾祭过偶像，随后承认此言不实，且此认错为人所信，那么基督徒可否使用曾听闻此说的肉，或将其出售并使用所得之钱呢？\par

十一、若基督徒旅途中因匮乏而饱受困苦，已禁食一、二或数日，遂至不能复忍此匮乏，值饥饿至极、眼见死亡临近之际，恰巧在偶像庙宇中发现食物，四周无人，亦别无他食可得；他应死去，还是该取食那食物呢？\par

十二、若基督徒即将被蛮族或罗马人杀害，他是否该杀死攻击者以保全己命？或虽不杀死攻击者，是否至少应逐退他、与之相斗？经上说：\BibleQuote{不要抵抗恶人}\BibleRef{玛 5:39}。\par

十三、基督徒可否在其产业周围筑墙以防御敌人？若有人利用那墙作为作战及杀敌之处，基督徒是否成此杀人之因呢？\par

十四. 基督徒可否饮用被投入祭物之泉源或井水？可否饮用荒废庙宇中所发现的井水？若在供奉偶像的庙宇中有一口未被任何东西玷污的井或泉，他可否汲取其水而饮用？\par

十五. 基督徒可否使用那些向偶像献祭之处的浴场？可否使用异教徒在节日所使用的浴场，无论他们还在那里或是已经离去？\par

十六. 基督徒可否使用异教徒在节日从其偶像处下来的同一顶轿子，若他们曾在这轿子中行了其偶像崇拜的任何部分礼仪，而基督徒亦知悉此事？\par

十七. 若一位基督徒在他人处作客，得知摆在他面前的肉是祭献过，便不吃它；但以后无意中发现同样的肉在出售而买来，或在别人的餐桌上又摆给他，他因不知道是同样的肉而吃了，他是否有罪？\par

十八. 基督徒可否购买并食用他知道是从庙宇或偶像的司祭的园子里取来的蔬菜或水果？为了不使你费心翻查圣经中关乎我所说之誓约与偶像的经文，我决意将靠主助所找到的那些经文摆在你面前；但若你在圣经中发现了更好或更切题的，敬请告知。例如，\Person{拉班}对\Person{雅各伯}说：\BibleQuote{\Person{亚巴郎}的天主和\Person{纳曷尔}的天主在我们之间裁断，} \BibleRef{创 31:53}圣经并未说明指的是哪位天主。再者，\Person{阿彼默肋客}来到\Person{依撒格}那里，他和与他同行的人向\Person{依撒格}起誓，我们并未被告知那是什么样的誓约。\BibleRef{创 26:31}至于偶像，\Person{基德红}奉主之命，将他所宰杀的公牛献作全燔祭。\BibleRef{民 6:26}又在\Person{农}的儿子\BibleBook{若苏厄}{书}中，论及\Place{耶里哥}说，所有的金银铜器都要送入主的府库，这应毁灭之城中所发现的物件都被称为圣物。\BibleRef{苏 6:19}我们在\BibleBook{}{申命纪}中也读到：\BibleRef{申 7:26}\BibleQuote{你不可将可憎之物带入你的家，免得你像它一样成了应受诅咒的。}\par

愿主保佑你。我问候你。请为我祈祷。\par

\SourceRecord{Translated by J.G. Cunningham.；Nicene and Post-Nicene Fathers, First Series；Vol. 1.；Edited by Philip Schaff.；Buffalo, NY: Christian Literature Publishing Co.,；1887.；<http://www.newadvent.org/fathers/1102046.htm>.}

% structure: p0001 source=h1 target=section reason=page-title

\section{书信四十七（主历三九八年）}

\Emph{致尊敬的\Person{普布利科拉}，我深爱的儿子，\Person{奥古斯丁}在主内致以问候。}\par

1. 自从我从你的来信得知你的困惑后，这些困惑也成了我的困惑，不是因为那些你告诉我使你困扰的事情搅扰了我的心灵；而是，我承认，我对于如何消除你的困惑感到极为困惑；特别是因为你要求我给出一个决定性的答复，以免你陷入比我未为你解惑之前更大的疑虑。因为我明白我无法给出这样的答复，尽管我可能写下在我看来最为确定之事，若我不能说服你，你无疑会比以前更加茫然；虽然我有能力运用那些我个人认为有力的论证，我却可能无法用这些论证说服他人。然而，为了不拒绝你的爱所要求的这小小服侍，我经过一番考虑后，决定写信答复。\par

2. 你疑虑之一，是关于使用一个以其假神发誓来保证其信实的人之服务。在这件事上，我恳请你思考：若一个人以这样的誓言作担保后食言，你是否不视为他犯了两重罪？因为若他遵守了以此誓言所确认的约定，他仅在以这些神明发誓这一点上被定罪；但没有人会正当地因他遵守约定而责备他。但在所假定的事例中，他既以那不应敬拜的神明发誓，又食言做了不当做的事，他便犯了两重罪：由此显而易见，若使用一个人的服务，并非为恶事，而是为某种良善且合法之目的，且已知此人之信实曾以假神之名起誓而得以确保，那么参与其中者，所分享的不是以假神发誓的罪，而是他遵守诺言的信实。我在此所说的信实，并非那使在基督内受洗者被称为信友的信实：那是全然不同、与人类契约和协定所要求的信实相去甚远的。然而，毫无疑问，以真天主之名发虚誓，比以假神之名发真誓更恶劣；因为我们凭之起誓者越神圣，发虚誓的罪就越大。因此，另一个问题是：那要求他人以其神明之名起誓以作担保的人，是否因其敬拜假神而有罪。在回答这个问题时，我们可接受你本人所引用的\Person{拉班}和\Person{阿彼默肋客}的例证为决定性的（如果\Person{阿彼默肋客}确曾以其神明发誓，如同\Person{拉班}以\Person{纳曷尔}之神发誓）。正如我所说的，这是另一个问题，若非有\Person{依撒格}和\Person{雅各伯}的例证——就我所知，或许还可加上其他例证——这个问题可能会使我困惑。可能有人仍会因新约中的诫命，\BibleQuote{什么誓都不可发}；\BibleRef{玛 5:34}这话，而心生疑虑；在我看来，这话之所以说，不是因为发真誓是罪，而是因为发虚誓是极重的罪：我们的主命令我们完全不可发誓，正是要我们远远躲避这罪行。然而我知道，我们的看法不同：因此此时不宜讨论；我们还是来探讨你来信请求我提供意见的那件事吧。基于你自己避免发誓的同样立场，若你持此见解，你或许可以视强求他人发誓为被禁止的，尽管明确说过，\Quote{不可发誓}；但我不记得在圣经中哪里读到我们不可接受他人的誓言。至于我们是否应利用他人之间通过交换誓言所建立的和谐，则完全是另一个问题。若我们对此作否定回答，我不知道我们在地上哪里还能生存。因为不仅是在边境，在所有行省，和平的保障都依赖于蛮族的誓言。由此推论，不仅那些由以假神之名发誓忠信者所守护的庄稼，而且一切受由类似誓言所确认的和平庇护的事物，都被玷污了。若你承认这纯属荒谬，便请连带你对所提名事例的疑虑一并抛开。\par

3. 再次，若有基督徒明知有人从他的打谷场或榨酒池中取东西去献给假神，而他有能力阻止却容许，他就有罪了。如果他事后发现此事已行，或无力阻止，他便坦然使用其余的谷物或酒，视为未受玷污，正如我们使用水泉，虽知道有人从中取水用于偶像崇拜一样。同样的原则也解决了浴场的问题。因为我们坦然吸入空气，虽知道所有偶像崇拜者的祭坛和香火的烟都升腾其中。由此显然，所禁止的是我们将任何东西归于假神的荣耀，或是行为上看似如此，以至于鼓动那些不知我们心意的人参与这样的崇拜，尽管我们内心鄙视那些偶像。当庙宇、偶像、丛林等奉官长的准许被拆毁时，虽然我们参与此工清楚证明我们并不尊崇、反而憎恶这些事物，但我们仍须克制，不可将其中任何物品据为己有、私自使用；好叫人人看出，我们拆毁这些是出于虔诚，而非贪婪。然而，当这些场所的战利品被用于公益，或奉献于事奉天主时，对待它们的方式，就与那些从不虔和亵圣中转向真宗教的人一样。我们从你自己引述的例子中领会到，这正是天主的旨意：祂命人从假神的丛林中取木料作全燔祭（由\Person{基德红}）；又命将\Place{耶里哥}的一切金银铜都带入主的府库。由此，在\BibleBook{申}{命纪}中也有一句诫命：\BibleQuote{你不可贪恋其上的金银，也不可收取，免得你陷于罗网，因为这是主你的天主所憎恶的。可憎之物，你不可带进你的家，免得你也成了当毁灭的，如同它一样；你应憎恶它，憎恨它，因为那是当毁灭之物。}\BibleRef{申 7:25}由此显然看出，要么是绝对禁止将这类战利品据为己有、私自使用，要么是禁止将任何这类物品带回家中，意图对它表示尊崇；因为那时它在天主眼中便是可憎和当诅咒的；相反，人们悖逆地给予那些偶像的尊崇，因着公开的毁灭而被彻底废除。\par

4. 关于祭偶像的肉，我向你保证，除了谨守宗徒对此的教导外，我们并无其他义务。因此，请研读他关于此事的言论，若它们对你而言隐晦不明，我愿尽力解释。若有人无意中吃了先前因曾祭过偶像而拒绝的食物，他并没有罪。一棵菜蔬，或地上所出的任何其他果实，都属于那创造它的主；因为\BibleQuote{大地和其中的万物，属于主}\BibleRef{咏 24:1}，并且\BibleQuote{凡天主所造的样样都好}\BibleRef{弟前 4:4}。但若地上所出产之物已被奉献或祭献给偶像，那么我们就必须将它算作祭偶像之物。我们必须小心，以免在宣称不应吃属于偶像庙宇园子里的果实时，被卷入这样的推论：宗徒在\Place{雅典}进食是不对的，因为那城属于\Person{密涅瓦}，且作为守护神而被奉献于她。对于庙内所围的水井或水泉，我也给出同样的回答，尽管若祭祀的一部分被投入这井或泉中，我的顾虑会稍增。但情形，正如我先前所说，完全与我们使用那接纳这些祭祀之烟的空气一样；或者，若有人以为这也有分别，即那与空气混合的烟所出的祭祀，并非献给空气本身，而是献给某些偶像或假神，而有时供物被投入水中，意图便是祭祀水本身，那么只需说，同样的原则也会阻止我们使用太阳的光，因为恶人在各处只要被容许，便不断崇拜那发光体。祭祀也献给风，但我们为自己方便仍使用风，虽然风似乎贪婪地吸入并吞下这些祭祀的烟。若有人对肉是否祭过偶像存疑，而事实是并未祭过，当他以这肉未曾祭过偶像的想法吃它时，他绝没有做错；因为事实上，且按他现时的判断，这也不是祭偶像的食物，尽管他先前以为是。因为无疑，用真实的印象纠正虚假的印象是合法的。但若有人以恶为善，并照着行，他便因怀着那信念而犯了罪；这些都是出于无知的罪，即人自以为做得对，其实却是他所不当做的。\par

5. 至于为保护自己的生命而杀害他人，我不赞同此事，除非这人恰好是士兵或公务人员，并非为自身，而是为保卫他人或他所居住的城市而行动，且是依照合法授予的委任，并以符合其职务的方式行事。然而，当人们因受到惊吓而被阻止作恶时，可以说，这对他们自己而言是一桩真正的服务。诫命\BibleQuote{勿抵抗恶}\BibleRef{玛 5:39}的颁布，是为防止我们以报复为乐，在这种心态中，心灵因他人的受苦而获得满足；并非叫我们忽略阻止人犯罪的义务。由此推知，一个人若为他的庄园筑了一道墙，即使有人因推倒墙而被倒塌的墙砸死，他并不犯有杀人罪。因为一个基督徒并不犯有杀人罪，即便他的牛顶了人，或他的马踢了人，以致那人死亡。按此原则，基督徒的牛就不应有角，马不应有蹄，狗不应有牙。按此原则，当宗徒\Person{保禄}特意通知千夫长，有某些亡命之徒为他设下伏击，并因此接受武装护送时\BibleRef{宗 23:17}，若那些图谋杀害他的歹徒撞上士兵的武器而死，\Person{保禄}就得承认他们的流血是由他所犯的罪。愿天主不许我们因某些意外而受责备，这类意外并非出于我们的意愿，而是借着我们所行或所拥有的、本身是善且合法的事物，发生在他人身上。若是那样，我们就不该在家中或田地里拥有铁器，免得有人用它自杀或杀害他人；我们的宅院内不该有树或绳索，免得有人上吊；我们的房屋不该有窗户，免得有人跳下。但我为何还要再举更多例子呢，这样的清单必是无尽的。因为人间所采用的善而合法的事物，有哪一样不会成为毁灭的工具而受指责呢？\par

6. 我现在只需注意（除非我记错了）你提到的一个事例：一个基督徒在旅途中，因极度饥饿而濒临绝境；若他找不到任何食物，只有放在偶像庙中的肉，并且附近无人能救济他，那么他宁愿饿死，也不取那食物充饥，是否更好？由于在这一问题中，并不假定如此找到的食物是祭献给偶像的（因为它可能是由在旅途中转往那里进食的人，因误放或故意留下的；或是为其他目的而放置的），我简要回答如下：要么确定这食物是祭献给偶像的，要么确定它不是，要么两者皆不知。如果确定，那么以基督徒的坚毅拒绝它更好。至于其他两种情形，他可以毫无良心不安地用来解决急需。\par

\SourceRecord{Translated by J.G. Cunningham.；Nicene and Post-Nicene Fathers, First Series；Vol. 1.；Edited by Philip Schaff.；Buffalo, NY: Christian Literature Publishing Co.,；1887.；<http://www.newadvent.org/fathers/1102047.htm>.}

% structure: p0001 source=h1 target=section reason=page-title

\section{书信48（主历398年）}

\Emph{致我的主\Person{欧多西乌斯}，我的弟兄及同侪司铎，我所热爱并渴望的，以及与他同在的弟兄们：} \Emph{\Person{奥古斯丁}与在此的弟兄们致以问候。}\par

1. 当我们思念你们在基督内享有的安宁休息时，我们虽身处多种劳苦艰辛中，也同享你们的安息，亲爱的弟兄们。我们是在同一元首下的同一身体，因此你们分担我们的劳苦，我们也分享你们的安宁：因为\BibleQuote{若是一个肢体受苦，所有的肢体都一同受苦；若是一个肢体受尊荣，所有的肢体都一同欢乐。}\BibleRef{格前 12:26} 因此，我们藉着基督的至深谦卑和至慈威严，恳切地劝勉并恳求你们，在你们圣洁的代祷中记念我们；因为我们相信你们的祈祷比我们的更为活泼、不受干扰，我们的祈祷常因世俗事务带来的黑暗与混乱而受损削弱——这不是我们自己要有的这些事务，而是由于人们强加给我们的重担，我们几乎无法喘息，他们仿佛强迫我们同他们走一里路，而我们的主却命令我们同他们走更远的路。\BibleRef{玛 5:41} 然而我们相信，那位垂听囚徒叹息的主，必眷顾我们，因祂乐意将我们安置在这职分上，并以应许的赏报鼓励我们坚守，也必因你们的祈祷援助，拯救我们脱离一切困苦。\par

2. 我们在主内劝勉你们，弟兄们，要坚定你们的志向，并坚持到底；如果教会——你们的母亲——召叫你们从事活跃的服事，一方面要提防因过分的喜乐而接受，另一方面也要提防因懒散怂恿而推辞；并以谦卑的心服从天主，温顺地顺服于那位治理你们的，祂将引领温顺的人公正判断，将教给他们自己的道路。不要贪图自己的安逸，而忽视教会的需要；因为如果当时没有善人愿意服事她，帮助她生育属灵的子女，你们自己属灵生命的开始就不可能实现。正如人们在火与水之间行走时必须小心道路，以免被烧或被淹，我们同样必须谨慎行走在骄傲的尖顶与懒散的漩涡之间；正如经上所载：\BibleQuote{不偏右也不偏左。}\BibleRef{申 17:11} 因为有些人过于谨慎地提防被高举，仿佛被提升到右边危险的顶峰，却掉落并沉入了左边的深渊。又有一些人，过于急切地逃避左边陷入怠惰柔弱的危险，却反而被自负的狂妄所毁灭焚烧，化为灰烬和烟雾。所以，亲爱的弟兄们，要在你们爱好安静中克制一切世俗的享乐，并要记住，没有一个地方是那惧怕我们飞向天主的捕鸟者不会为我们设下陷阱的；我们要视那曾经俘虏我们的，为一切善人的死敌；切莫以为我们已经享有完全的平安，直到罪恶止息，\Quote{审判回归公义}为止。\par

3. 当你们强健而热切地努力时，无论做什么，或勤于祈祷、禁食、施舍，或分施于穷人，或宽恕过犯，\BibleQuote{如同天主在基督内宽恕了我们一样}\BibleRef{弗 4:32}，或克制恶习，痛击身体，使其为奴\BibleRef{格前 9:27}，或忍受苦难，尤其要彼此相爱忍耐（因为一个人若不宽容弟兄，还能忍受什么呢？），或提防诱惑者的诡计和伎俩，用信德的盾牌抵挡并熄灭他的火箭\BibleRef{弗 6:16}，或\BibleQuote{以圣咏、颂词、属神歌曲，从你们心中歌颂赞美主}\BibleRef{弗 5:19}，或用声音与心灵和谐一致地歌唱——总之，我再说，\BibleQuote{一切都要为光荣天主而作}\BibleRef{格前 10:31}，那位\BibleQuote{在一切内工作一切}的天主\BibleRef{格前 12:6}，且要如此\BibleQuote{在圣神内热忱}\BibleRef{罗 12:11}，使\Quote{你们的灵魂在主内欢躍}。这就是那些行走在\Quote{正直道路}上的人的行径，他们的\Quote{眼睛不断仰望上主，因为祂将把他们的脚从网罗中拔出}。这样的行径既不被俗务打断，也不因闲散而麻木，既不喧嚷也不懈怠，既不傲慢也不沮丧，既不鲁莽也不懒散。\BibleQuote{实行这些事吧，这样赐平安的天主必与你们同在。}\BibleRef{斐 4:9}\par

4. 愿你们的爱德使你们不要认为我冒昧地想要给你们写信。我提醒你们这些事，并不是因为我认为你们在这些方面有欠缺，而是因为我想，如果你们在向天主尽职时，能记起我的劝勉，那么我在天主面前就因你们而倍受举荐。因为美好的报告，甚至在弟兄\Person{欧斯塔西乌斯}和\Person{安德肋亚斯}从你们那里到来之前，就已经如同他们所带来的那样，传来了你们圣善交往所散发的基督的馨香之气。其中，欧斯塔西乌斯已经先我们去了那片安息之地，那里的岸边没有像围绕你们海岛家园那样的狂风巨浪拍打；在那地方，他不再为\Place{卡普雷拉岛}遗憾，因为那里提供给他的粗布衣裳，他已不再穿戴。\par

\SourceRecord{Translated by J.G. Cunningham.；Nicene and Post-Nicene Fathers, First Series；Vol. 1.；Edited by Philip Schaff.；Buffalo, NY: Christian Literature Publishing Co.,；1887.；<http://www.newadvent.org/fathers/1102048.htm>.}

% structure: p0001 source=h1 target=section reason=page-title

\section{第50封信（公元399年）}

\Emph{致\Place{苏菲克托姆}殖民地的行政官与领袖，或称长老，\Person{奥思定}主教致意。}\par

大地震撼，苍天战栗，皆因听闻那骇人听闻的罪行与前所未有的残暴，竟使你们的街道和庙宇血流成河，回荡着凶手的呼喊。你们已将罗马法律埋葬于耻辱的坟墓，轻蔑地践踏了对公正法令应有的敬畏。你们既不尊重也不畏惧皇帝的权威。在你们的城中，我们六十位弟兄的无辜之血已被倾流；凡在这次屠杀中表现最为积极的人，都得到了你们的喝彩，并在你们的议事会中获得高位。来吧，我们现在来谈谈这次暴行的主要借口。如果你们说\Person{赫拉克勒斯}是属于你们的，我们必定会弥补你们的损失：我们有金属在手，石头也不缺乏；不仅如此，我们还有各种大理石和许多工匠。不要害怕，你们的神正在工匠手中，将被精心雕凿、打磨和装饰。我们还会额外加上一些红赭石，好让他脸红得与你们的虔敬相称。或者，如果你们说\Person{赫拉克勒斯}必须由你们自己制造，我们就发起一分钱的募捐，从你们自己的工匠那里为你们买一个神。只是请你们同时向我们作出赔偿；正如你们的神\Person{赫拉克勒斯}归还给你们一样，请把你们暴力所毁灭的许多人的生命归还给我们。\par

\SourceRecord{Translated by J.G. Cunningham.；Nicene and Post-Nicene Fathers, First Series；Vol. 1.；Edited by Philip Schaff.；Buffalo, NY: Christian Literature Publishing Co.,；1887.；<http://www.newadvent.org/fathers/1102050.htm>.}

% structure: p0001 source=h1 target=section reason=page-title

\section{第五十一封书信（主后399年或400年）}

邀请\Place{加辣玛}城的\Person{克里司皮努斯}多纳图派主教，就多纳图派分裂的全部问题展开讨论。\par

（信首无问候语。）\par

1. 我之所以采用这样的信首格式，是因为你们一派对于我在其他信函开头所使用的谦卑问候语感到不悦。若非我相信你在回信时也会同样如此，别人可能会以为我这样做是对你的侮辱。关于你在\Place{迦太基}的承诺，以及我催你履行承诺之事，我何必多言呢？让我们把过去彼此对待的方式连同过去一起遗忘，免得它阻碍将来的交谈。现在，除非我弄错了，靠着主的帮助，路上已无障碍：我们两人都在\Place{努米底亚}，彼此相距不远。我听人说，你仍愿意与我辩论那使我们无法共融的问题。你看，一切疑虑可以多么迅速地得到澄清：如果你愿意，请给我回信答复，也许这样不仅对我们俩，就是对那些想听我们辩论的人而言也就够了；倘若不够，就让我们一再通信，直到讨论透彻为止。因为我们所居住的城市相距较近，这岂不是一个很大的便利吗？我决心只以书信向你展开辩论，这一方面是为了防止任何说过的话从记忆中溜走，另一方面也是为了让那些对此问题感兴趣却无法亲临辩论的人不至于失去学习的机会。你常常不是出于故意撒谎，而是由于误解，拿一些适合你目的的指控来责难我们，这些指控是关于一些我们斥为不实的往事。因此，如果你愿意，就让我们依照目前的状况来衡量这个问题，而把过去的事搁在一边吧。你无疑知道，在犹太教的制度下，百姓犯了拜偶像的罪，又有一次，天主的先知的圣书被一位悖逆的君王焚烧了；\BibleRef{耶 36:23} 对分裂之罪的惩罚，若不是因其罪过更大，就不会比这两个罪所受的更严厉了。你当然记得，大地如何裂开，活活吞下了为首的分裂者，又有火从天降下，摧毁了他们的同党。\BibleRef{户 16:31} 制造并崇拜偶像，或焚烧圣书，都不配受这样的惩罚。\par

2. 你们惯常以一项罪行来责备我们，这项罪行虽然确凿无疑地证实于你们自己一派中的某些人，却从未在我们身上得到证实——这项罪行就是在迫害的恐惧下交出圣经任其焚烧。因此，我要请问，你们为何重新接纳了那些被你们以\Quote{你们那\InnerQuote{全体会议无误的声音}}（我是根据记录引述的）定为分裂之罪的人，并将他们重新安置在你们对他们宣判时他们所占据的同一主教教座上？我指的是\Place{穆斯}的\Person{费里西安}和\Place{阿苏里}的\Person{普雷特克塔图斯}。这两个人并不像你们想让无知者相信的那样，属于你们的大公会议指定并告知了一个期限的那些人，过了那期限，如果他们还没有回到你们中，判决就会成为定局；相反，他们属于你们当天就毫不迟延地定了罪的那些人——而在同一天，正如我所说的，你们给了某些人缓期。这一点，如果你否认，我可以证明。你们自己的大公会议就是证人。我们也有\WorkTitle{总督纪事}，其中你们对此不只一次，而是反复确认。因此，如果你能，就另找辩护之辞吧，免得你否认我能证明的事，白白浪费时间。那么，如果费里西安和普雷特克塔图斯是无辜的，他们为什么被这样定罪？如果他们有罪，又为什么被这样恢复职位？如果你证明他们无辜，那么当存在以下情况时——即一些无辜者被诬陷为\OriginalTerm{交付圣经者}{traditors}，却被比你们先辈人数少得多的你们的先辈们定罪；而另一些无辜者被诬陷为\OriginalTerm{分裂者}{schismatics}，却被他们的后继者三百一十人定罪，且那次定罪的宣判还被夸张地描述为出自\Quote{\InnerQuote{全体会议无误的声音}}——你还能反对我们相信无辜者可能被诬定罪吗？然而，如果你证明他们被定罪是公正的，那么你又能提出什么理由来为他们恢复原职、重新占据同一主教教座辩护呢？除非你因强调和平的重要性与益处而主张：为了保持合一纽带不受破坏，甚至这样的事情也当被容忍。但愿你不只是口头上，而是全心地说出这个理由！那么你不会看不到，任何诽谤都不能成为破坏遍及世界的基督和平的理由，既然在非洲，允许人们宁可恢复那些曾因不虔敬的分裂而被定罪的人原职，也不愿破坏\Person{多纳图}及其一派的和平。\par

3. 你们还惯常责备我们借助世俗权力迫害你们。关于这一点，我既不从如此大不虔敬所包含的严重过犯来立论，也不从调节我们行动严厉性的基督温和来立论。我持以下立场：如果这算一种罪，那么你们为何借助那些皇帝所任命的法官，残酷迫害\OriginalTerm{马西米安派}{Maximianists}呢？那些皇帝藉福音获得属灵新生，本是归功于我们的教会。你们为何以喧嚣的争辩、谕令的权势和士兵的暴力，把他们从他们原本拥有、在退出你们时还占有的礼拜场所赶走？他们在各处那场斗争中所受的冤屈，有不久前事件的痕迹可以作为见证。文件宣告了所下的命令。那些所作所为在整个地区都声名狼藉，而在这同一地区，你们领袖\Person{奥普塔图斯}的神圣记忆正受到尊崇。\par

4. 再者，你们惯常说我们没有基督的圣洗，并且说在你们的共融之外找不到圣洗。关于这一点，我本要进入更长的辩论；但对付你们时这并不必要，因为你们连同 \Person{费利西安} 和 \Person{普雷特克斯塔特}，也承认了马克西米安派的圣洗为有效。因为这些主教在与 \Person{马克西米安} 共融期间所施洗的所有人，尽管你们当时正全力在世俗法庭上进行旷日持久的争斗以将他们（即 \Person{费利西安} 和 \Person{普雷特克斯塔特}）逐出他们的教会，正如案卷所证实的——我说，他们在那一时期所施洗的所有人，现在都与他们和你们共融在一起。尽管这些人是在他们受绝罚且身负分裂之罪时由他们所施洗的，不仅是在因危险疾病而临危的情况下，而且是在复活节礼仪中，在他们所属城市的众多教堂里，就在这些重要城市本身——对于所有这些人的情况，圣洗的仪式一次也没有被重复过。我希望你们能证明，当 \Person{费利西安} 和 \Person{普雷特克斯塔特} 受绝罚且身负分裂之罪时，那些可说是徒然被他们施洗的人，在他们恢复身份后，又被他们恰当地重新施洗了。因为如果为子民重新施洗是必要的，那么为主教们重新祝圣也并不比这更不必要。因为他们若离开你们后不能在你们的共融之外施洗，便丧失了主教职权；因为若他们离开你们时尚未丧失主教职权，他们就仍能施洗。但如果他们丧失了主教职权，在回来时就应重新领受祝圣，好使他们所失去的得以恢复。然而，不要让此事令你们惊慌。正如他们返回时确实持有与他们当初离开你们时相同的主教身份，同样确实的是，他们带回了那些他们在 \Person{马克西米安} 的分裂中所施洗的所有人，并且没有重复他们的圣洗，一同回到你们的共融中。\par

5. 当我们看到马克西米安派的圣洗被你们承认，而普世教会的圣洗却被藐视时，我们怎能不极度哀哭呢？你们对 \Person{费利西安} 和 \Person{普雷特克斯塔特} 的定罪，无论是听取了他们的辩护还是未听取，无论是公正还是不公，我不问这些；但请告诉我：\Place{格林多}教会的主教有哪一位曾在你们的法庭上为自己辩护过，或从你们那里接受过判决？或者说，\Place{迦拉达}、\Place{厄弗所}、\Place{哥罗森}、\Place{斐理伯}、\Place{得撒洛尼}的主教，或任何其他被包含在如下应许中的城市的主教，有谁这样做过？那应许说：\Quote{地上万民都要朝拜祢。}然而，你们接纳前者的圣洗，却藐视后者的圣洗；而圣洗本不属于前者也不属于后者，而是属于那位论及祂时说：\BibleQuote{这人就是要用圣神施洗的那一位。} \BibleRef{若 1:33}的那一位。然而，我暂时不细讲这一点：请留意我们近在眼前的事——看哪，甚至对瞎子也能产生深刻印象的事！我们在哪里找到你们所承认的圣洗呢？诚然，是在那些你们已定罪的人那里，而不是在那些从未在你们的法庭上受审的人那里！——是在那些因分裂之罪被指名谴责并被你们驱逐的人那里，而不是在那些你们不认识、居住在遥远地域、从未被你们指控或定罪的人那里！——是在那些不过是 \Place{阿非利加} 一隅居民中的一小部分人那里，而不是在那些从他们的国家福音最初传到 \Place{阿非利加} 的人那里！我何必增添你们的负担呢？请就这些事给我一个答复。看看你们的公会议对马克西米安派的指控，指控他们犯了邪恶的分裂之罪；看看你们向他们所借助的世俗法庭的迫害；看看你们恢复了一些人而并未重新祝圣他们，并且认可了他们的圣洗为有效这一事实；然后，如果你们能，请回答：你们是否有能力向那些无知的人隐瞒，你们在一个比你们曾炫耀已定罪于马克西米安派的分裂更加可憎的分裂中，使自己与普世分离的原因？愿基督的平安在你们心中得胜！那么一切都会好的。\par

\SourceRecord{Translated by J.G. Cunningham.；Nicene and Post-Nicene Fathers, First Series；Vol. 1.；Edited by Philip Schaff.；Buffalo, NY: Christian Literature Publishing Co.,；1887.；<http://www.newadvent.org/fathers/1102051.htm>.}

% structure: p0001 source=h1 target=section reason=page-title

\section{第 53 封信（公元 400 年）}

\Emph{致我们至爱而尊贵的弟兄\Person{格内罗苏斯}，\Signature{福图纳图斯、阿利皮乌斯和奥古斯丁}在主内致意。}\par

% structure: p0003 source=h2 target=subsection reason=relative-heading-rank; base=h2

\subsection{第一章}

1. 既然你愿意把那封由一名多纳特派司铎寄给你的信向我们告知，虽然你本着真正公教徒的精神轻视了它，但为了帮助你追求他的福祉——如果他的愚妄并非不可救药——我们恳请你把下面的答复转交给他。他写道，有一位天使命令他向你宣告你城里的基督信仰的主教继承体系；向你，诚然，你所持守的基督信仰不仅属于你自己的城，也不仅属于\Place{非洲}，而是属于全世界，那已经传给了万民，并正在传给万民的基督信仰。这证明他们认为自己被切割出去却不感到羞愧是小事，并且只要还有可能，也不设法重新被嫁接；他们不满足，除非他们尽力切割别人，使他们与自己同命运，如同适合焚烧的枯枝。因此，即使那位天使亲自显现给你——他声称天使向他显现了，我们视此为狡猾的虚构——即使天使对你说了那些话，就是他以所谓命令为依据向你复述的——即使那样，你也应记起宗徒的话：\BibleQuote{即使我们，或是从天上来的天使，若给你们宣讲的福音与我们曾经给你们宣讲的不同，他就该受诅咒。}\BibleRef{迦 1:8} 因为主耶稣基督亲自的声音已向你们宣告，他的\BibleQuote{福音将传给万民，然后末日才来到。}\BibleRef{玛 24:14} 此外，众先知和众宗徒的著作也已向你们宣告，那应许是赐给亚巴郎和他的后裔，即基督的，\BibleRef{迦 3:16} 当天主对他说：\BibleQuote{地上万民都要因你的后裔蒙受祝福。} 既然你们有了这些应许，若有一位从天上来的天使对你说：\Quote{放弃全地的基督信仰，依附\Person{多纳图斯}派，其主教继承已由你们城里的他们的主教在一封信中列出}，那么他在你们眼中就该受诅咒；因为他正试图将你们从整个教会割离，推入一个小的党派，使你们失去对天主应许的份。\par

2. 因为如果应该考虑主教的血统继承，那么我们向上追溯直到\Person{伯多禄}本人，对教会来说，将会多么更加确定和有益啊；主向他说——他以预象的方式代表整个教会——\BibleQuote{在这磐石上，我将建立我的教会，阴间的门不能胜过它！}\BibleRef{玛 16:18} \Person{伯多禄}的继承人是\Person{理诺}，他绵延不绝的继承人如下：——\Person{克肋孟}、\Person{阿纳克肋图}、\Person{埃瓦里斯图}、\Person{亚历山大}、\Person{西斯笃}、\Person{泰莱斯福鲁斯}、\Person{伊吉诺}、\Person{阿尼切图}、\Person{庇护}、\Person{索特尔}、\Person{埃留特里乌斯}、\Person{维克托}、\Person{则斐林}、\Person{卡利斯图}、\Person{乌尔班}、\Person{彭提安}、\Person{安特鲁}、\Person{法比安}、\Person{科尔内略}、\Person{路基}、\Person{斯德望}、\Person{西斯笃}、\Person{狄奥尼修}、\Person{菲利克斯}、\Person{欧提基安}、\Person{凯尤}、\Person{马尔切林}、\Person{马尔切卢斯}、\Person{欧瑟比乌}、\Person{米迪亚德}、\Person{西尔维斯特}、\Person{马尔谷}、\Person{尤里乌}、\Person{利贝里乌}、\Person{达玛苏}、\Person{西里奇}，其继任者是现任主教\Person{安纳斯塔修}。在这个继承序列中，找不到任何多纳特派主教。但是，逆转事物的自然进程，多纳特派从非洲派了一位已祝圣的主教到\Place{罗马}，他居于大都会中少数非洲人的首领地位，使他们以\Quote{山民}或\OriginalTerm{库特祖皮特}{Cutzupits}闻名，世人即以此名称呼他们。\par

3. 现在，即使在这些世纪中，由于疏忽，某个\OriginalTerm{背教者}{traditor}获得了从\Person{伯多禄}本人到现今占据该宗座的\Person{安纳斯塔修}之间的主教序列中的位置——这一事实也对教会和无辜的基督徒没有害处；因为主预见到这种情况，论到教会中的管理者说：\BibleQuote{凡他们所吩咐你们遵守的，你们都要遵守遵行；但不要效法他们的行为，因为他们只说，却不做。}\BibleRef{玛 23:3} 这样，信友们的希望就稳固了，因为它建立于主，而非人，绝不会被不虔的分裂的狂怒冲走；而那些在读圣经中提到宗徒们曾致书的教会名字，却在这些教会中没有主教的人，他们自己反而被冲走了。还有什么比这更能证明他们的悖逆和愚昧呢？当他们读出这些信时，对自己的圣职人员说：\Quote{愿平安与你们同在}，而他们自己却与这些信原本写给的那些教会的平安分离。\par

% structure: p0007 source=h2 target=subsection reason=relative-heading-rank; base=h2

\subsection{第二章}

然而，为免他因\Place{君士坦丁纳}——你自己的城市——的主教继承而过于自鸣得意，你要读给他听在曾任执政官的\Person{穆纳提乌斯·菲利克斯}面前的庭审记录；菲利克斯是\OriginalTerm{祭司}{Flamen}（异教祭司），在\Person{戴克里先}第八次与\Person{马克西米安}第七次执政官任期之年，担任你们城市的行政长官，日期是六月朔日前十一天。这些记录证明，主教\Person{保禄}曾是一名\OriginalTerm{背教者}{traditor}；事实是，那时\Person{西尔瓦努斯}是他的副执事之一，与他一起拿出并交出了一些属于主之圣殿的物件，这些物件曾被精心藏匿，即一只银匣和一盏银灯，看见这些东西时，据称一个名叫\Person{维克托}的人说道：\Quote{你若没有找到这些，早就被处死了。} 你们的多纳图派司铎在他写给你的信中，对这位已明确被定罪的背教者\Person{西尔瓦努斯}极为推崇，提到他由\Place{提吉西斯}的\OriginalTerm{首席主教}{Primate}\Person{塞昆杜斯}祝圣为主教。让他们闭紧骄傲的嘴，承认真正可以加于他们身上的指控，而不是对他人发出愚蠢的诽谤。如果允许，你也要读给他听这位提吉西斯的\Person{塞昆杜斯}在\Person{乌尔巴努斯·多纳图斯}宅邸中所进行的教会庭审记录，在那次庭审中，他将那些自认是背教者的人交给天主作审判者——包括\Place{马斯库利}的\Person{道纳图斯}、\Place{阿奎·提比利坦奈}的\Person{马里努斯}、\Place{卡拉马}的\Person{道纳图斯}——并且他竟与这些人作为同僚，虽然他们是自认的背教者，却祝圣了他们的主教\Person{西尔瓦努斯}，关于他在这同一件事上的罪过，我在前面已经记述了。你还要读给他听在曾任执政官的\Person{泽诺菲鲁斯}面前的庭审记录；在那次庭审中，他们的一个名叫\Person{农迪纳里乌斯}的执事，因\Person{西尔瓦努斯}将他绝罚而恼怒，将所有事实诉诸法庭，以真实文件、讯问证人、宣读官方记录和多份书信，不容置辩地证明了这些事实。\par

还有许多其他事情，你可以读给他听，如果他愿意不怒气冲冲地辩论，而是审慎地倾听。比如：多纳图派向\Person{君士坦丁}的请愿，恳求他从\Place{高卢}派遣主教来解决这场分裂非洲主教的争端；记录在\Place{罗马}发生的事情的\WorkTitle{纪事}，当时由他派往那里的主教们受理并裁决了此案；你还可以从其他书信中读到，上述皇帝如何陈述，他们曾向他申诉，反对他们的同僚——即他派往罗马的主教们——的裁定；他又如何指定其他主教在\Place{阿尔勒}重新审理此案；他们又如何从那个法庭再次上诉到皇帝那里；最终他如何亲自调查此事；以及他如何极其强调地宣布，\Person{凯西利安}的无辜战胜了他们。如果他愿意，就让他听听这些事，他就会沉默，不再密谋反对真理。\par

% structure: p0010 source=h2 target=subsection reason=relative-heading-rank; base=h2

\subsection{第三章}

然而，我们所依靠的，并不怎么是这些文件，而是\WorkTitle{圣经}，其中应许给基督的产业是一个伸展到地极、遍布万国的权柄；他们因自己罪恶的裂教而与此分离，却将属于主\OriginalTerm{禾场}{threshing-floor}上糠秕的罪过归咎于我们，而糠秕必被允许与好麦子混在一起，直到终结之时，直到末日的审判将一切扬净。由此显明，无论这些指控是真是假，它们都不属于主的麦子，\BibleRef{玛 13:30}这麦子必须在整个田地，\Emph{即}全地，长到世界末日；我们如此知道，不是藉着那位印证了你通信者错误的假天使的见证，而是藉着主在福音中的话语。而且，因为这些不幸的多纳图派，将许多虚假空洞的指控的耻辱，加给那些本就无可指摘、却在全世界与糠秕或莠子，\Emph{即}与名不副实的基督徒，混杂在一起的基督徒，因此天主以公义的报应，命定他们借其普世大公会议，谴责马克西米安派为裂教者，因为他们曾谴责\Person{普里米安}，并在未与\Person{普里米安}共融的情况下施洗，且给那些由\Person{普里米安}施洗的人重洗，然后，在短暂间隔之后，在\Person{吉尔多}的走狗\Person{奥普塔图斯}的胁迫下，恢复其中两位——\Place{穆斯提}的主教\Person{菲利西安}和\Place{阿苏里}的主教\Person{普雷特克塔图斯}——的原有职务之荣誉，并承认所有由这些受绝罚的人施洗者的洗礼。因此，如果他们没有因与这些恢复原职的人共融而被玷污——这些人，他们曾亲口谴责为邪恶不虔，并将他们比作被大地活活吞下的第一批异端者——就让他们最终醒悟，思考一下，他们认为全世界因非洲人未知的罪恶而被玷污，并认为基督的产业（按照应许已向万民彰显）因与这些非洲人保持共融而被他们的罪所摧毁，这该是何等的盲目和愚昧；同时，他们却拒不承认，自己因与那些他们既已知道又已谴责其罪行的人共融，而遭到败坏和玷污。\par

7. 因此，由于宗徒\Person{保禄}在另一处说，连撒殚也常冒充光明的天使，所以他的仆役若冒充正义的仆役，并不足为奇。\BibleRef{格后 11:13}如果你的通信者果真看见一位天使教导他错谬，想要使基督徒远离教会的合一，他遇到的便是那冒充光明天使的撒殚使者。但如果他对你撒了谎，并没有见到这样的异象，那么他自己就是撒殚的仆役，冒充正义的仆役。然而，如果他不是顽固不化、执迷不悟，他或许能藉着思考上述一切，既免于误导他人，也免于自己受误导。因为，我们把握你提供的机会，毫无怨恨地与他交锋，并记起宗徒论到他的话：\BibleQuote{主的仆人不可争吵，但要和气对待众人，善于教导，凡事忍耐；以温和规劝反对的人，或许天主会赐给他们悔改的心，得以认识真理，使这些被魔鬼活捉去顺从他心意的人，能脱离他的罗网。}\BibleRef{弟后 2:24}因此，如果我们说了什么严厉的话，要他明白，那并非出于争辩的尖刻，而是出于热切渴望他回归正途的爱。愿你在基督内安然无恙，最亲爱的、可敬的弟兄！\par

\SourceRecord{Translated by J.G. Cunningham.；Nicene and Post-Nicene Fathers, First Series；Vol. 1.；Edited by Philip Schaff.；Buffalo, NY: Christian Literature Publishing Co.,；1887.；<http://www.newadvent.org/fathers/1102053.htm>.}

% structure: p0001 source=h1 target=section reason=page-title

\section{第54封信（公元400年）}

\EditorialNote{也称为《\WorkTitle{对雅努阿里乌斯问题的回答}》卷一。}\par

\Emph{致他亲爱的儿子\Person{雅努阿里乌斯}，\Person{奥古斯丁}在主内致以问候。}\par

% structure: p0004 source=h2 target=subsection reason=relative-heading-rank; base=h2

\subsection{第1章}

1. 关于你问我的问题，我本愿知道你自己会作何回答；这样我就能以更少的话答复你，也最容易通过赞同或修正你的答案来确认或纠正你的意见。我本来会很乐意的。但为了立即答复你，我想写一封长信比浪费时间要好。因此，首先，我愿你在目前的讨论中持守一个基本原则，即我们的主耶稣基督，正如祂在福音中所宣布的：\BibleRef{玛 11:30}，已为我们指定了一个\Quote{柔和的轭}和一个\Quote{轻松的担子}。祂依此在新律法下，借着数极少、最易遵行且意义极其卓越的圣事，将祂的子民结合在共融之中：例如以天主圣三之名举行的圣洗，祂圣体圣血的共融，以及正典圣经中所规定的其他事项——但那些对天主的古代子民而言是奴役之轭的律令除外，那些律令适合他们当时的心理状态和先知时代，见于\WorkTitle{梅瑟五书}中。至于我们根据传统而非圣经权威所持守、并在全世界遵守的其他事物，可以理解为它们是由宗徒们自己，或由大公会议所批准和设立的——大公会议在教会中的权威极为有益。\Emph{例如}，以特殊庆典纪念主的受难、复活与升天，以及圣神自天降临，以及普世教会在任何地方建立后以同样方式遵守的其他一切。\par

% structure: p0006 source=h2 target=subsection reason=relative-heading-rank; base=h2

\subsection{第2章}

2. 然而，有其他事物在不同的地方和国家有所不同：\Emph{例如}，有些人在星期六守斋，其他人则不；有些人每日领受基督的圣体圣血，其他人则在固定的日子领受：有些地方没有一天不举行祭献；其他地方则只在星期六和主日，或者只在主日举行。至于这些以及任何地方可能遇到的各种可变惯例，每个人有自由选择遵守或不遵守；在这事上，对明智而严谨的基督徒而言，没有比适应他可能来到的教会所盛行的做法更好的规则了。因为这种习俗，只要明显不违背信德或健全的道德，就应视为无关紧要之事，并应为了与我们所共处的人保持共融而遵守。\par

3. 我想你可能以前听我说过，但我现在还是要提一下。当我的母亲随我来到\Place{米兰}时，她发现那里的教会在星期六不守斋。她开始感到不安，不知如何是好；于此，我当时虽然对这类事并不特别在意，还是替她请求了至为可怀念的\Person{安博罗}的建议。他回答说，他只能教他自己所实行的，因为如果他知道有更好的规则，他自己就会遵守了。当我以为他单凭自己的权威、不提供任何论证，意在建议我们放弃星期六守斋时，他随即对我说：\Quote{当我访问\Place{罗马}时，我在星期六守斋；当我在这里时，我不守斋。基于同样的原则，无论你来到哪个教会，你都应遵守那教会通行的惯例，如果你既不愿因自己的行为使人跌倒，也不愿因别人的行为而跌倒。}当我将这话告诉我的母亲时，她欣然接受了；而我自己，在反复思考他的决定之后，一直珍视它，仿佛是从天上得来的神谕。因为我常常极为难过地看到，许多软弱的弟兄因某些人好辩的固执或迷信的摇摆不定而陷入不安；这些人在此类问题上——这类问题既不能由圣经的权威、普世教会的传统、也不能由它们对品行的明显良好影响来作出最终决定——或许出于自己某种怪僻的念头，或许出于对本乡惯例的执着，或许出于对在国外所见的偏爱，以为智慧与离家远行的距离成正比，便以如此激烈的态度鼓动这些问题，以致认为除了他们自己所行的，一切都不对。\par

% structure: p0009 source=h2 target=subsection reason=relative-heading-rank; base=h2

\subsection{第3章}

4. 也许有人会说：\Quote{圣体圣事不应每日领受。}你问：\Quote{有何根据？}他回答：\Quote{因为，为使一个人能相称地接近如此伟大的\OriginalTerm{圣事}{sacrament}，他应选择那些自己生活在更特别的纯洁与克己之中的日子；因为\BibleQuote{谁若不相称地吃喝，就是吃喝自己的罪案。}} \BibleRef{格前 11:29}另有一人回答：\Quote{的确；若罪恶所导致的创伤和灵魂疾病的猛烈，使得必须暂时推迟使用这些良药，那么每个人在此情况下，应由主教的权威禁止接近祭台，并被指派做\OriginalTerm{补赎}{penance}，之后再由同一权威恢复特权；因为如果在应当作补赎的时候领受，那便是不相称地领受；而且，人不可凭自己的判断随意退出教会的共融或自行恢复。然而，若他的罪尚未大到足以公正地受到\OriginalTerm{绝罚}{excommunication}的判决，他便不应为了灵魂的治愈而避开每日领受主的身体。}或许有第三位提出对此问题更为公正的裁决，提醒他们，首要之事是存留在基督的平安中保持合一，并且各人应自由地按照自己的信仰，凭良心视其为本分而行。因为他们二人皆没有轻看主的身体和血；相反，两者都在争相以最高的尊崇来敬重这充满祝福的圣事。福音中提及的那两位——匝凯和百夫长——之间并无争论；两人都不认为对方比自己更好，尽管前者欢欢喜喜地将主接到自己家中，\BibleRef{路 19:6}后者却说：\BibleQuote{主，我当不起祢到我舍下来，}\BibleRef{玛 8:8}——两人都以不同的、看似彼此相反的方式尊敬救主；两人都因罪而可怜，也都获得了他们所需的怜悯。我们还可以进一步借用一则比喻：昔日赐给天主古代子民的玛纳，在每个人口中尝起来的味道恰如他所愿。这征服世界的圣事在每位基督徒心中也是如此。那些不敢每日领受的人，和那些不敢任一天漏掉的人，同样都是出于敬意。这神圣食粮不容人轻看，正如玛纳不能无罪地厌恶。因此，宗徒说，那些不相称地分享此圣事的人，是因为他们没有将它与所有其他食物区分开来，没有给予它所当受的特殊尊敬；因为在已引用的经文\BibleQuote{吃喝自己的罪案}之后，他又加上了这句话：\BibleQuote{没有辨别主的身体；}从\BibleBook{格林多}{前书}的整段经文来看，如果仔细研读，这一点是显而易见的。\par

% structure: p0011 source=h2 target=subsection reason=relative-heading-rank; base=h2

\subsection{4}

5. 假设有位来自外地的人，访问一个在\OriginalTerm{四旬期}{Lent}期间习惯禁用沐浴、并在周四继续禁食的地方。他说：\Quote{我今天不守这禁食。}当被问及理由时，他说：\Quote{在我的家乡不是这样的风俗。}这样的行为，难道不是试图宣称自己的风俗优于他们的吗？因为他无法从\WorkTitle{天主的书}中引用明确的经文作为依据；也无法借着普世教会无论散布何处的共识来证明自己的意见正确；更无法证明他们行事违背信仰，而他按信仰行事，或是他们正在进行有损健全道德的事，而他则在维护其利益。那些人因在一个不必要的问题上争吵，而损害了自己的安宁与和平。我更愿意建议，在这类事情上，当一个人旅居他乡，发现当地的风俗与自己不同时，应同意随从当地人的做法。另一方面，倘若一位基督徒在海外旅行时，到了某个地方，那里的天主子民人数众多、更容易聚集、且宗教热情更高，他看见，\Emph{例如}，在四旬期最后一周的周四，祭献举行了两次，分别在早晨和傍晚，于是当他回到自己的家乡，那里的祭献只在傍晚举行时，他便抗议说这是错误且不合法的事，因为他曾在异地见过不同的风俗，这便显出一种幼稚的判断力软弱，对此我们应加以防范，也必须容忍他人如此，但要纠正所有受我们影响的人。\par

% structure: p0013 source=h2 target=subsection reason=relative-heading-rank; base=h2

\subsection{5}

6. 现在来看你的来信中第一个问题属于这三类中的哪一类。你问：\Quote{在四旬期最后一周的星期四，我们应当做什么？我们应当在早晨举行祭献，然后再在晚餐后举行祭献，因为福音中有\BibleQuote{同样地……晚餐后？}的话吗？\BibleRef{路 22:20} 还是我们应当禁食，只在晚餐后举行祭献？或者我们应当禁食直到祭献完成，然后照常吃晚餐？} 因此我回答说，如果圣经的权威决定了这些做法中哪一个是正确的，我们就毫无疑问应当按经上所记的去行；我们的讨论便当专注于一个问题，不是关于责任，而是关于解释这神圣建立的含义了。同样地，若普世教会遵循其中任何一种做法，我们就毫无疑问应当遵行；因为讨论是否应当遵行，会是极其狂妄的疯狂之举。但你提出的问题既非由圣经，亦非由普世习俗所决定。因此它必须被归入第三类——即那些在不同地方、不同国家做法各异的事。所以，每个人当顺应他所往之地的教会的惯例。因为这些做法，无一与基督信仰或道德利益相悖，不会因采纳某一种习俗更受青睐。若信仰或健全的道德受到威胁，那就必须要么改正错谬的做法，要么确立被忽略的当行之事。但仅仅改变习俗，即便在某些方面可能有益，也会因其新颖而扰乱人心：因此，若改变并无益处，便因无益地扰乱教会而大有害处。\par

7. 我还要补充说，那种认为在许多地方盛行的、在那一天进食后举行\OriginalTerm{祭献}{sacrifice}的习俗，可追溯至\Quote{同样地晚餐后}等语，是错误的。因为主或许将祂的身体称为晚餐，他们既已领受了祂的身体，便是在这之后才领了杯：正如宗徒在另一处所说：\BibleQuote{你们聚集在一处的时候，这并不是\Emph{吃}主的晚餐，}\BibleRef{格前 11:20} 从而将领受\OriginalTerm{圣体圣事}{Eucharist}（亦即吃饼）称为主的晚餐。\par

% structure: p0016 source=h2 target=subsection reason=relative-heading-rank; base=h2

\subsection{第六章}

至于在那一天，在举行祭献或领受圣体圣事之前进食是否合宜的问题，福音中的这些话或许足以让我们下定决心：\BibleQuote{他们正吃的时候，耶稣拿起饼来，祝福了，} 与上文所说的相连：\BibleQuote{到了晚上，祂与十二宗徒坐席。他们吃的时候，祂说：我实在告诉你们：你们中有一个人要出卖我。} 因为正是在那之后，祂建立了圣事；显然，门徒们初次领受主的体血时，并未禁食。\par

8. 那么，我们是否因此要谴责普世教会，因为圣事在各处都是由禁食的人领受的？不，断然不是，因为从那时起，圣神乐意规定，为如此伟大圣事的尊荣，主的身体应先于所有进入基督徒口中的其他食物；正是为此，这一习俗才被普遍遵行。因为主在进用其他食物之后建立了圣事，这并不证明弟兄们应当在进餐或晚餐后聚集领受圣事，或是效法那些被宗徒责备并纠正的人，他们不区分主的晚餐与普通餐食。的确，救主为将这奥迹的深邃更动人地托付给门徒，乐意借着在离开他们受难前的最后行动，将建立此圣事铭刻在他们心中和记忆里。因此，祂并未规定当如何遵行的次序，而将此事留给宗徒们，祂意图通过他们安排一切有关教会的事宜。倘若祂规定圣事必须总是在其他食物之后领受，我相信无人会偏离此做法。然而，当宗徒论及此圣事时说：\BibleQuote{所以，我的弟兄们，当你们聚集吃的时候，要彼此等待；若有人饿了，当在家里吃，免得你们聚集自取判决。} 他立即补充道：\BibleQuote{其余的事，等我来时再安排。}\BibleRef{格前 11:33} 由此我们明白，既然在一封书信中全面规定那为普世教会所遵循的方式，对他而言过于繁重，这便属于他亲自安排的事之一；因为我们在各种其他习俗的差异中，发现此惯例是统一的。\par

% structure: p0019 source=h2 target=subsection reason=relative-heading-rank; base=h2

\subsection{第七章}

9. 确实，有些人认为（他们的看法并非不合理），在一年之中一个固定的日子，即主建立晚餐的那一天，在其他食物用后，祭献并领受主的体血是合法的，为的是赋予该纪念日的礼仪特殊的庄严性。我认为，在这种情况下，将庆祝安排在这样一个时辰更为适宜，即能让任何禁食的人在第九时辰惯常的餐前参与礼仪。因此，我们既不强迫，也不敢禁止任何人在那一天主的晚餐前开斋。然而我相信，这一习俗赖以存在的真正理由是，在许多地方，许多人，几乎是所有人，习惯在那一天沐浴。又因为有些人继续禁食，所以祭献便在早晨为那些进食的人举行，因为他们无法同时承受禁食和沐浴；而在晚上，则为那些全天禁食的人举行。\par

10. 若你问我，在那一天沐浴的习俗由何而来，我细细思索后觉得，最可能的原因是为了避免在圣洗池边发生不体面之事：那些即将领受圣洗的人，因在四旬期内严格禁绝沐浴，身体不免污秽，若不在之前一日得到清洗，难免有碍观瞻；而特意选定这一天，则是因为它乃是\OriginalTerm{晚餐}{the Supper}建立的周年纪念。那些即将领受圣洗的人既蒙准许享受沐浴，许多其他人便也想加入其中，同享沐浴之乐，并放松他们的斋戒。\par

我已尽己所能讨论了这些问题。我劝勉你，在能力所及之处，遵守我所陈述的，一如教会明智而爱好和平的子女所当作的。至于你其余的问题，若主愿意，我打算另找时间回答。\par

\SourceRecord{Translated by J.G. Cunningham.；Nicene and Post-Nicene Fathers, First Series；Vol. 1.；Edited by Philip Schaff.；Buffalo, NY: Christian Literature Publishing Co.,；1887.；<http://www.newadvent.org/fathers/1102054.htm>.}

% structure: p0001 source=h1 target=section reason=page-title

\section{书信第五十五篇（公元400年）}

% structure: p0002 source=h2 target=subsection reason=relative-heading-rank; base=h2

\subsection{第一章}

1. 读过来信，其中你提醒我，我尚有义务答复你很久以前向我提出的其余问题，因此我不能继续拖延满足你对教导的渴望——看到你如此渴望，令我心甚喜甚慰。虽然被诸多事务缠身，我已将为你提供所需答复的工作置于首位。我不再对你的来信内容多作评论，以免因此耽误我清偿所欠。\par

2. 你问道：\Quote{为什么我们纪念主受难的日子，不像传统所传下来的主诞生日那样，每年都在同一天庆祝？}你又补充说：\Quote{如果这个原因与星期和月份有关，那么在这一庆典中，我们与星期几或月相状态有何相干？}你首先必须知道并记住的是，守主诞生日并非圣事性的，而仅仅是纪念祂的诞生，因此只需通过一个年度宗教节日来标记该事件发生的日子即可。只有在纪念事件被如此安排，以至于被理解为象征某件应当受虔敬尊崇的神圣事物时，庆祝活动才在本质上成为圣事性的。因此，我们庆祝复活节的方式，不仅是回忆基督的死亡与复活的史实，也要容纳所有与这些事件相关的、证明该圣事意义的事物。因为正如宗徒所写的：\BibleQuote{祂为我们的过犯被交付，又为使我们成义而复活，}\BibleRef{罗 4:25}在主的那次受难与复活中，一种从死亡到生命的过渡被祝圣了。因为\OriginalTerm{巴斯挂}{Pascha}这个词本身并非如通常所想的那样是希腊词：通晓两种语言的人断言它是一个希伯来词。因此，这个词并非源自受难，因为希腊词\OriginalTerm{受苦}{\GreekText{πάσχειν}}，而是取意于我所提到的从死亡到生命的过渡；希伯来词\OriginalTerm{巴斯挂}{Pascha}的意思是逾越或过渡。主自己曾有意指涉这一点，当祂说：\BibleQuote{那信从我的，已出死入生。}\BibleRef{若 5:24}记录这话的同一位圣史，在叙述主即将与门徒庆祝逾越节并在那时设立圣体圣事晚餐时，有意对此予以强调，他说：\BibleQuote{耶稣知道自己离世归父的时候到了。}\BibleRef{若 13:1}这种从有死生命到另一不死生命的过渡，即从死亡到生命，在主受难和复活中被彰显出来。\par

% structure: p0005 source=h2 target=subsection reason=relative-heading-rank; base=h2

\subsection{第二章}

3. 这种从死亡到生命的过渡，在当前是由我们的信德在我们内成就的，这信德是我们为得罪赦和永生盼望而拥有的，当我们爱天主和近人时；\BibleQuote{因为信德由爱德行事，}\BibleRef{迦 5:6}并且\BibleQuote{义人因信德而生活；}\BibleRef{哈 2:4}\BibleQuote{所见的盼望不是盼望：因为人所看见的，为什么还要盼望呢？但我们若盼望那看不见的，就必忍耐等候。}\BibleRef{罗 8:24}根据这信德、望德和爱德，藉此我们开始处于\Quote{恩宠之下}，我们已经与基督同死，并藉着圣洗与祂同葬于死亡之中；正如宗徒所说：\BibleQuote{我们的旧人已与祂同钉在十字架上；}\BibleRef{罗 6:6}并且我们与祂一同复活了，因为\BibleQuote{祂使我们一同复活，并使我们在天上的地方一同坐在基督内。}\BibleRef{弗 2:6}因此他也给予这劝勉：\BibleQuote{你们既然与基督一同复活了，就该追求天上的事，在那里有基督坐在天主的右边。你们该思念天上的事，不该思念地上的事。}\BibleRef{哥 3:1}在接下来的话中，\BibleQuote{因为你们已经死了，你们的生命已与基督一同藏在天主内；当基督——我们的生命——显现时，那时你们也要与祂一同出现在荣耀中，}\BibleRef{哥 3:3}他清楚地使我们明白，我们如今藉着信德由死亡到生命的过渡，是在对未来最终复活和荣耀的盼望中完成的，那时\BibleQuote{这必朽坏的}，即我们如今在其中呻吟的这肉躯，\BibleQuote{必须穿上不朽坏的，这必死的必须穿上不死的。}\BibleRef{格前 15:53}因为现在，我们确实藉着信德拥有\Quote{圣神的初果}；但我们仍\Quote{在自己内叹息，等候义子期望的实现，即我们肉身的救赎：因为我们得救在于盼望。}我们在这盼望中时，\Quote{肉身固然因罪而死，心神却因义而生活。}现在注意接下来的话：\Quote{然而，如果那位使耶稣从死者中复活者的圣神住在你们内，那么那位使基督从死者中复活的，也必要藉着住在你们内的圣神，使你们有死的身体复活。}因此，整个教会，在此尘世旅途和死亡状态下，期盼在世界终结时，那首先在我主耶稣基督的身体上展示的，也要在她身上实现，祂是\BibleQuote{死者中的首生者，}\BibleRef{哥 1:18}因为祂作为元首的身体正是教会。\par

% structure: p0007 source=h2 target=subsection reason=relative-heading-rank; base=h2

\subsection{第三章}

4. 的确，有些人仔细研究宗徒经常使用的那些话——关于我们与基督同死同复活——却误解了那些话的意思，便以为复活已是过去的事，在终结时再无别的可盼。他说：\BibleQuote{其中有\Person{依默纳约}和\Person{菲肋托}；他们离开了真理，说复活已是过去的事，颠覆了一些人的信德。} \BibleRef{弟后 2:17} 这位宗徒虽然如此谴责并指证他们，却也教导我们与基督一同复活了。这表面的矛盾如何消除呢？难道他不是指这事借着信德、望德和爱德，按着圣神的初果，在我们身上完成吗？但正因为\BibleQuote{看得见的希望，不算是希望}，\BibleQuote{我们若希望那看不见的，就必须坚忍等待。} 毫无疑问，肉身的救赎尚在未来，我们为此而\BibleQuote{心中叹息}。因此又有话说：\BibleQuote{论望德，要喜乐；在困苦中，要忍耐。} \BibleRef{罗 12:12}\par

5. 所以，我们生命的这种更新，是一种由死入生的过渡，首先借着信德完成，使我们论望德而喜乐，在困苦中忍耐，同时我们的\BibleQuote{外在的人日渐损坏，但内在的人却日日更新。} \BibleRef{格后 4:16} 正是由于这新生命的开始，由于我们奉命脱去旧人，穿上新人 \BibleRef{哥 3:9}，\BibleQuote{除净旧酵，好使我们成为新团，因为我们逾越节的羔羊基督，已被祭献作了牺牲。} \BibleRef{格前 5:7} 我说，正是由于我们内这生命的新，才选定一年之首月，作为这庆节的时期。这月也的确被称为\OriginalTerm{阿彼布月}{Abib}，即月份之首 \BibleRef{出 23:15}。再者，主在第三日复活，因为世界第三时代也由此开始。第一时代是在法律以前，第二时代是在法律之下，第三时代是在恩宠之下，如今在此时代中，那从前隐藏在晦暗预言里的奥秘已经彰显出来。因此，这也在庆祝所定的那部分月份中表征出来；因为数字七在圣经中通常作为奥秘的数字，用来表示某种圆满，所以复活节的庆祝日落在该月的第三周内，即在十四日到二十一日之间。\par

% structure: p0010 source=h2 target=subsection reason=relative-heading-rank; base=h2

\subsection{第4章}

6. 这里还有另一个奥秘，如果你因对这类学问不太熟悉而不能轻易领悟，也无须苦恼；也不要因为我自幼学得这些事，便以为我比你强：因为主说：\BibleQuote{凡要夸耀的，应当因这事夸耀：就是明白认识我是上主。} \BibleRef{耶 9:24} 有些留意这类学问的人，曾对天体的数字和运行作过许多探讨。其中最有才能的人发现，月亮的盈亏是由于其球体的转动，而不是像愚昧的\OriginalTerm{摩尼教徒}{Manichæans}所设想的那样，其本体有实际的增加或减少。他们说，如同船被装满一样，月亮被逃亡的天主本体的一部分所充满；他们以不虔的心和口，毫不迟疑地相信并宣称这部分与黑暗的掌权者相混合，并为其污秽所玷染。他们解释月亮的渐盈说，当那失落的天主本体的部分，经过巨大的劳苦，从污染中被净化，逃出整个世界和一切污秽不堪之物，回归于那位哀悼它直至归来的天主之时，便发生这种渐盈；由于这样，月亮直到月半就被充满，而在月半之后，这一部分又被倒回太阳，如同倒进另一条船。在这些可咒的亵渎中，他们却从未能设法解释，为何月亮在发光之初或末尾，以角状发光，或者为何它在月中开始亏缺，而不一直满盈，直至将其增长倒回太阳为止。\par

7. 然而，我所指的那些人，曾以可靠的计算探究了这些事，以致他们不仅能够说明日食和月食的原因，还能在许久之前预言它们的发生，并能通过数学运算确定它们必然发生的精确时刻，并将结果载于论著中；任何人通过阅读和理解这些论著，也能同他们一样预言这些食相的到来，并发现他们的预言得到事实验证。这些人——正如\WorkTitle{圣经}所教导的，他们应受谴责，因为\BibleQuote{他们虽有了智慧，能测量世界的周期，却没有更容易地——凭着谦卑的孝爱，他们本可以——认识它的主}\BibleRef{智 13:9}——我说，这些人从\OriginalTerm{月牙}{horns of the moon}推论出，月牙在盈缺时都背向太阳，要么月球是由太阳照亮的，并且月球离太阳越远，在从地球可见的一面上就越充分地暴露于太阳光下；但在月中之后，当月球沿轨道的另一半靠近太阳时，它的上半部分就更加明亮，而朝向地球的一面接收阳光越来越少，因此我们看起来月球在逐渐亏缺：要么，如果月球自身有光，那么这光只存在于它的一侧半球上，随着月球远离太阳，它逐渐将这一面更多地转向地球，直到完全展现出来，从而显示出表面的增大，不是由于补充了缺失的部分，而是由于显露了原本就在那里的部分；同样，在向太阳运行的过程中，月球又逐渐将已显露的部分从我们视线中移开，因此看起来在减小。无论这两种理论哪一种正确，至少这一点是明显的，任何细心的观察者都很容易发现：月球在我们的眼中，只有当它远离太阳时才会增大，只有当它返回太阳时才会减小。\par

% structure: p0013 source=h2 target=subsection reason=relative-heading-rank; base=h2

\subsection{第五章}

8. 现在，请注意\BibleBook{}{德训}中所说的：\BibleQuote{智慧人如太阳坚定不移；愚昧人却如月亮变幻无常。}\BibleRef{德 27:12} 而那没有变化的智者是谁呢？岂不是那\OriginalTerm{义德的太阳}{Sun of Righteousness}，论到祂有话说：\BibleQuote{义德的太阳已为我升起，} 而恶人将在审判之日哀哭这太阳没有为他们升起时，要说：\BibleQuote{正义的光没有照耀我们，太阳也没有为我们升起。}\BibleRef{智 5:6}？因为那肉眼可见的太阳，天主使之同样升起在恶人与善人之上，正如祂降雨给义人和不义的人；\BibleRef{玛 5:45} 然而，恰当的比喻常取自可见之物，以解释不可见之事。再问，那\Quote{愚昧人}，那\Quote{如月亮变化}的，若不是亚当——众人都因他犯了罪——又是谁呢？因为人的灵魂，离开了\OriginalTerm{义德的太阳}{Sun of Righteousness}，即离开了对不变真理的内在默观，便将全力转向外在事物，在其更深更高贵的能力上越发昏暗；但当灵魂开始回向那不变的智慧，越是怀着虔敬的渴望亲近它，\OriginalTerm{外在的人}{outward man}便越趋朽坏，而\OriginalTerm{内在的人}{inward man}却日复一日地更新，灵魂中凡倾向于下地事物的光，都转向了上界的事物，因此从世间事物中抽离出来；从而灵魂对这个世界死去得越来越多，而其生命与基督一同藏在天主内。\par

9. 因此，当灵魂转向外在事物，抛弃与内修生活相关的一切时，它是变得更坏了；而对世界，即对思念世事的人而言，灵魂在这种情况下反显得更好了，因为恶人因心中的贪欲而受赞扬，不义者竟受祝福。但当灵魂逐渐将目标与抱负从看似重要的世间事物转离，而指向更高尚、不可见的事物时，它便是变得更好；而对世界，即对思念世事的人而言，灵魂在这种情况下似乎变坏了。因此，那些最终必将徒然悔罪的恶人，在别的怨言中会这样说：\BibleQuote{这些人就是我们曾嘲笑侮辱过的；我们在愚蠢中曾以为他们的生活方式是疯狂。}\BibleRef{智 5:3} 现在，\OriginalTerm{圣神}{Holy Spirit}借着从可见之物到不可见之物、从有形之物到属灵奥迹的比拟，乐意命定那象征旧生命进入新生命的节日——即\OriginalTerm{逾越节}{Pascha}——应在月份的十四日至二十一日之间庆祝：在十四日之后，是为获得关于属灵事实的双重阐释，即关于世界的第三个时代（此即我说它在第三周发生的原因），以及关于灵魂从外在转向内在的变化——这一变化与月亮亏缺时的变化相应；不晚于二十一日，是因为数字七本身常被用来表示宇宙的概念，也因教会与宇宙的相似而应用于她。\par

% structure: p0016 source=h2 target=subsection reason=relative-heading-rank; base=h2

\subsection{第六章}

10. 为此，宗徒若望在\BibleBook{若望}{默示录}中写信给\Emph{七个}教会。此外，当教会仍处于我们现世肉身生命的状况下时，由于她易于变化，\WorkTitle{圣经}以月亮之名称呼她；例如，\BibleQuote{他们已在弓囊中备好了箭，要在月亮昏暗时，射伤那些心地正直的人。}因为在那宗徒所说的\BibleQuote{当基督，我们的生命显现时，你们也要与祂一同在光荣中显现}\BibleRef{哥 3:4}的事发生之前，教会在她旅居尘世期间似乎被遮蔽，在许多罪恶中呻吟；这时，那些欺骗和误导人的阴谋诡计是当畏惧的，这些就是经文中\Quote{箭}这个词所指的。再者，我们在\BibleBook{}{圣咏}第八十九篇中另有实例，在那里，由于她到处为真理带来忠实的见证，教会被称作\BibleQuote{月亮，天上的忠实见证。}当圣咏作者歌颂主的王国时，他说：\BibleQuote{在他的日子里，正义与和平的丰盛，直到月亮被消灭；}\Emph{即}和平的丰盛将如此增长，直到祂最终除去这现世状况固有的一切变化。那时，死亡——最后的仇敌——将被消灭；凡由于我们肉身的软弱而阻碍我们获得完美平安的一切障碍，当这可朽坏的穿上不可朽坏，这可死的穿上不死的时候，都要被彻底消灭。我们还有另一个实例：耶里哥城的城墙——在希伯来语中据说意指\Quote{月亮}——当约柜绕城七次时，它们倒塌了。因为绕耶里哥的约柜所象征的天国来临的预许，除了表示这现世生命的一切堡垒，即一切属于此世的、抵抗来世希望的希望，必须藉着灵魂的自由同意，由圣神的七恩予以摧毁之外，还传达什么别的呢？因此，当约柜绕行时，那些城墙自己倒塌了，并非由于猛烈的攻击。除此之外，\WorkTitle{圣经}中还有其他论及月亮的章节，以这个形象向我们讲明教会在此世的光景：在忧虑和劳作中，她处于必死的命运下，是远离那个圣天使为公民的耶路撒冷的旅客。\par

11. 这些拒绝改过迁善的愚妄人，并没有理由认为那些天上的光体应受崇拜，因为有时借用它们的形象来象征神圣的奥迹；因为这样的借喻是从一切受造物中汲取的。也没有理由让我们招致宗徒对某些人宣判的定罪之词，这些人敬拜并事奉受造物，而不是那永远受赞颂的造物主。\BibleRef{罗 1:25}我们并不朝拜羊或牛，尽管基督被称为羔羊\BibleRef{若 1:29}，先知也称祂为牛犊\BibleRef{则 43:19}；也不朝拜任何猛兽，尽管祂被称作犹大支派的狮子\BibleRef{默 5:5}；也不朝拜石头，尽管基督被称作磐石\BibleRef{格前 10:4}；也不朝拜熙雍山，尽管在它内有教会的预像\BibleRef{伯前 2:4}。同样，我们也不朝拜太阳或月亮，尽管为了传达神圣奥迹的教导，从造物主的这些天体化工中借用了圣事的形象，正如从祂在地上创造的许多事物中借用一样。\par

% structure: p0019 source=h2 target=subsection reason=relative-heading-rank; base=h2

\subsection{第七章}

12. 因此，我们有责任深恶痛绝地谴责那些占星者的胡言乱语。他们自己陷入了虚妄的空想，还以此拖累他人，当我们指责他们这些空洞的虚构时，他们自以为可以振振有辞地回答：\Quote{那么，你们为什么根据日月的位置计算来规定庆祝复活节的日期呢？}——好像我们所指责的是天体的秩序，或是季节的更替，而这些是全能全善的天主所安排的，而不是指责他们扭曲地滥用这些天主以完美智慧所定的事物，来支持最荒谬的见解。如果占星者可以据此禁止我们从天体借用比较来象征圣事的实体，那么占鸟兆者也可以同样理由阻止我们使用\WorkTitle{圣经}中的话：\BibleQuote{你们应如同鸽子般纯朴}，以及\BibleQuote{你们应如同蛇一般机警}\BibleRef{玛 10:16}；而演戏者也可以干扰我们提及\BibleBook{}{圣咏}中的琴瑟。因此，如果他们愿意，就让他们说吧：因为表现天主圣言奥迹的比喻取自我们所提及的这些事物，我们就该被指控为或是依鸟的飞行问卜，或是调制耍蛇者的毒药，或是沉溺于剧场的放荡——这种说法真是荒谬绝伦。\par

13. 我们并不通过研究日月、年月的时节来预测我们事业的结果，以免在人生最艰难的紧急关头，我们撞上可悲奴役的礁石，使我们的自由意志遭受船难；而是以最虔诚的虔诚敬神之心，我们接受这些天体所提供的适合说明神圣事物的比喻，正如我们从所有其他受造物，如风、海、陆、鸟、鱼、牛、树木、人等，在我们的谈话中借用多种形象；在举行圣事时，我们采用基督宗教的较为自由的安排所规定的极少事物，如水、饼、酒、油。然而，在旧约的安排之下，规定了许多礼仪，这些礼仪只为了教导我们其意义而让我们知晓。我们现在不再遵守年、月、季节，以免\Person{保禄}宗徒的话适用于我们：\BibleQuote{我为你们担心，恐怕我在你们身上所费的劳苦是白费了。}\BibleRef{迦 4:11} 因为他责备那些说：\Quote{我今天不出发，因为今天是凶日，或因为月亮是如此这般；} 或 \Quote{我今天要去，以便万事顺利，因为星辰的位置是这样或那样；这个月我不做生意，因为某星掌管；} 或 \Quote{我要做生意，因为另一星辰已取而代之；今年我不种葡萄园，因为今年是闰年。} 然而，任何有常识的人都不会认为那些研究时节的人该受责备，\Emph{例如} 谁说：\Quote{我今天不出发，因为暴风雨已经开始了；} 或 \Quote{我不会出海，因为冬天还没有过去；} 或 \Quote{现在是播种的时候了，因为大地已被秋雨浸透；} 等等，关于那些与天体的运动和湿度有关的、经观察得到的其他自然效应，这些天体是被完美安排运行的，当它们被造时，经上说：\BibleQuote{让它们作为记号、时节、日子和年岁。}\BibleRef{创 1:14} 同样，每当从受造物中借用说明性的象征来宣扬属灵的奥秘时，不仅从天体和星球，也从更低微的生物中借用，这是为了使救恩的教义具有一种雄辩，适合将接受者的情感从可见的、有形体的、短暂的事物，提升到不可见的、属灵的、永恒的事物。\par

% structure: p0022 source=h2 target=subsection reason=relative-heading-rank; base=h2

\subsection{第八章}

14. 我们中没有人关注这样一件事：在我们庆祝复活节的时候，太阳位于\Quote{白羊宫}——人们如此称呼天体的某个区域，事实上，在每个月的开始，太阳就在那里。但无论他们选择称那部分天空为白羊宫还是别的什么，我们从圣经中得知，天主造了所有天体，并按自己的旨意安置它们的位置；无论天文学家怎样划分这些为不同星座所指定和安排的区域，无论他们用什么名称区分它们，太阳在第一个月所处的位置，正是这个圣事庆典应当找到这颗星球的位置，因为它圣化了生命更新的一个神圣奥秘，对此我已经充分论述过了。然而，如果称那一部分天体为\Quote{白羊宫}是因为其形状与名称有些对应，那么天主的话不会犹豫从这类事中借用来说明一个神圣奥秘，正如它不但从其他天体，也从地上之物借用，\Emph{例如} 从\Place{猎户座}和\Place{昴宿}、\Place{熙雍山}、\Place{西乃山}，以及那些被命名的河流：\Place{基红河}、\Place{丕雄河}、\Place{底格里斯河}、\Place{幼发拉底河}，尤其是\Place{约旦河}，它在神圣奥秘中被如此频繁地提到。\par

15. 然而，谁察觉不到以下两者之间的巨大差异呢？一方面是对天体进行有益的观察，以此预测天气，像农夫或水手所做的那样；或者为了确定他们所在的方位以及应遵循的路线，像船只领航员或穿过南部非洲无路可循的沙地的人所做的那样；或者为了借用关于天体的一些事实来以比喻的方式呈现某种有益的教义——另一方面，人们徒劳地迷信，观察天空不是为了知道天气、他们的路线，也不是为了进行科学计算，或寻找属灵事物的说明，而只是为了窥探未来，获知命运已决定了什么？\par

% structure: p0025 source=h2 target=subsection reason=relative-heading-rank; base=h2

\subsection{第九章}

16. 现在让我们用心思考：为什么在庆祝复活节时，要小心安排日子，使星期六在其之前？因为这是基督宗教所特有的。犹太人从尼散月十四日到二十一日守逾越节，无论那个星期从哪一天开始。但是，在主受难的那次逾越节，事情是这样：主的死亡与复活之间，有犹太人的安息日。因此，我们的教父们认为，在他们的复活节庆典中加上这一特殊规定是适当的，一方面为了使我们的节日区别于犹太人的逾越节，另一方面为了使后代在他们每年纪念主的受难时，能够持守那件事，我们必须相信，那位在时间之前就存在、时间也是藉着他而造、并在时间圆满时来临、当他说\Quote{我的时辰还没有到}时有权舍去生命并再取回来的那位，是为了某种好的理由而行此事，因此他等候的不是由盲目的命运所定的时辰，而是适合他所决意要托付给我们观察的神圣奥秘的时辰。\par

17.\ 我们在信德与望德中所拥有的，并由爱德努力奔赴的，正如我以上所说，是一种摆脱所有挂虑重负的某种圣洁永恒的安息；从今生过渡到那安息，我们的主耶稣基督甘愿在自己受难中予作示范并祝圣。然而，这安息并非怠惰的无所作为，而是一种因工作而生的不可言喻的宁静，其中没有痛苦的努力。因为在此生劳苦之后所进入的憩息，其性质是与更好生命中欢快的活动相结合的。然而，由于这活动是在颂扬天主中进行的，没有身体的劳苦或精神的焦虑，因此过渡到那活动并非通过一种随后仍有劳苦的憩息，\Emph{即}那种一旦活动开始便不再是憩息的憩息：因为在这些活动中，再无劳苦与挂虑的反复；而构成安息者——免除工作疲劳和思想疑虑——总是存在于其中。如今，既然借着安息我们恢复了灵魂因罪丧失的原初生命，这安息的预表便是每周的第七日。但那原初生命本身，恢复给那些从流浪归来、并作为欢迎而领受他们最初袍子的人，则由每周的第一日代表，我们称之为\OriginalTerm{主日}{Lord's day}。如果你读\BibleBook{创世}{纪}，查考七日的记载，你会发现第七日没有夜晚，这表示它所预表的安息是永恒的。最初赋予的生命并非永恒，因人犯了罪；但最终的安息，即以第七日为预表的安息，是永恒的，因此第八日也将拥有永恒的福乐，因为那永恒的安息被第八日吸收，而非被它毁灭；因为如果被毁灭，它就不是永恒的了。因此，第八日，即每周的第一日，向我们代表那原初生命，不是被拿走，而是被转化为永恒的。\par

% structure: p0028 source=h2 target=subsection reason=relative-heading-rank; base=h2

\subsection{第十章}

18.\ 然而，第七日被指定给犹太民族，作为身体安息应遵守的日子，以便成为那圣化的预表，人是借着在圣神内的安息而达到圣化的。在\BibleBook{创世}{纪}所载之前各日的历史中，我们并没有读到圣化的事：唯独论及安息日时，才说\BibleQuote{天主祝福了第七日，并定为圣日。}\BibleRef{创 2:3}。人的灵魂，无论善恶，都喜爱安息，但如何达到他们所喜爱的，对大多数人来说却是未知的：身体按其重量所寻求的，正是灵魂按其爱所寻求的，即一个安息之所。因为就如物体依其比重下降或上升，直到抵达能安息之处——例如油若倒入空气中便下降，若倒入水中则上升——同样，人的灵魂奋力趋向它所爱的事物，为能借着抵达它们而得到安息。固然有很多通过身体而使灵魂愉悦的事物，但灵魂在这些事物中的安息并非永恒，甚至不能长久延续；因此它们反而贬低灵魂，拖累它，使之牵制那股趋向更高事物的纯粹无重量之势。当灵魂从自身找到愉悦时，它尚未在不变之物中寻求快乐；因此它仍是骄傲的，因为它将自己置于最高位，而天主更高。在这样的罪中，灵魂不会不受惩罚，因为\BibleQuote{天主拒绝骄傲人，却赐恩宠于谦逊人。}\BibleRef{雅 4:6}。然而，当灵魂以天主为乐时，就在那里找到了真实、稳固而永恒的安息，这是在一切其他事物中徒然寻求的。因此\BibleBook{圣咏}{集}中给出了这样的劝勉：\BibleQuote{你要在主内寻乐，祂必赐给你心中所渴望的。}\par

19.\ 因此，\BibleQuote{天主的爱，借着所赐与我们的圣神，已倾注在我们心中了，}\BibleRef{罗 5:5}圣化便与第七日相连，即那规定安息的日子。但我们既不能行任何善工，除非蒙天主的恩赐而得到帮助，正如宗徒所说的：\BibleQuote{因为是天主在你们内工作，使你们愿意，并使你们力行，为成就他的善意，}\BibleRef{斐 2:13}同样，在我们完成此生所有善工之后，也不能得以安息，除非借着同一恩赐被圣化和成全，直到永恒；因此，论到天主自己，说祂造齐了一切，认为\BibleQuote{很好，}\BibleRef{创 1:31}之后，就在\BibleQuote{第七天休息，停止了祂所作的一切工程。}\BibleRef{创 2:2}因为祂这样做，正预表了将来祂要在我们善工完成后，赐予我们人类的安息。因为正如在我们行善时，是祂在我们内工作，借着祂的恩赐我们才能行善；同样在我们的安息中，祂也被说成是安息者，我们借着祂的恩赐而安息。\par

% structure: p0031 source=h2 target=subsection reason=relative-heading-rank; base=h2

\subsection{第十一章}

20. 此外，这正是为什么在十诫中，三条指出我们对天主本分的诫命中，安息日的法律被置于第三位（因为其余七条涉及我们的近人，即人；全部法律都系于这两条诫命）\BibleRef{玛 22:10}。第一条诫命禁止我们崇拜任何人类技艺所制造的天主肖像，我们应理解这条诫命是指向圣父：这一禁令的颁布，不是因为天主没有肖像，而是因为除了那位与祂自己相同的肖像之外，不该崇拜祂的其他任何肖像；而且这同一肖像，不是代替祂，而是与祂一同受人崇拜。然后，由于受造物是变化无常的，因此经上说：\BibleQuote{一切受造物都屈服于虚妄，}\BibleRef{罗 8:20} 而且整体的本性也在其任何一部分中彰显出来，免得有人将天主圣子，即创造万物的圣言，视为受造物，所以第二条诫命是：\BibleQuote{不可妄称上主你天主的名。}\BibleRef{出 20:7} 又因为天主圣化了祂休息的第七日，圣神——在祂内，我们获得了那各处所寻求，但只在爱天主时才寻获的安息，正如经上所说：\BibleQuote{天主的爱，藉着赐给我们的圣神，已倾注在我们心中了。}\BibleRef{罗 5:5} ——便借着关于遵守安息日的第三条诫命呈现于我们的心思中，这并非使我们以为在今生就能获得安息，而是要使我们一切善工的努力都指向那永恒的安息。因为我特别提醒你要记住上面所引用的经文：\BibleQuote{我们得救，还是在于希望；所见的希望不是希望。}\BibleRef{罗 8:24}\par

21. 为了滋养并煽旺那热烈的爱德——在这种爱德中，我们如同受到一种类似重力的法则，被向上或向内引向安息——透过\OriginalTerm{象征}{emblems}来呈现真理，具有极大的力量：因为如此呈现的事理，比那些不加\OriginalTerm{圣事象征}{sacramental symbols}的、赤裸裸的陈述，更能感动并点燃我们的情感。何以如此，难以说明；但事实是，我们透过\OriginalTerm{寓意}{allegory}或象征所学到的一切，比用最清晰、最平白的词语陈述出来的，更能影响我们、取悦我们，也更受我们的重视。我相信，当灵魂完全沉浸于尘世事物时，情感不易被燃起；但如果它先被带向那些作为灵性事物象征的有形事物，然后从这些象征被引向它们所代表的灵性实体，那么单凭从一者过渡到另一者的行动，灵魂便聚集力量，如同一支点燃的火炬，因着运动而燃烧得更加明亮，并被一种更为强烈炽热的爱德带往安息。\par

% structure: p0034 source=h2 target=subsection reason=relative-heading-rank; base=h2

\subsection{第十二章}

22. 也正因如此，在全部十诫中，与安息日相关的那一条，是唯一所命令的事具有\OriginalTerm{预像}{type}意义的；所规定的身体安息乃是一种预像，我们接受它作为教导的方式，却不是我们也必须遵守的义务。因为安息日显示的是灵性安息的预像，关于这安息，圣咏上说：\BibleQuote{你们要停手，应承认我是天主，} 而且主自己也用这些话邀请人进入这安息：\BibleQuote{凡劳苦和负重担的，你们都到我跟前来吧！我要使你们安息。你们背起我的轭，跟我学吧！因为我是良善心谦的：这样你们必要找得你们灵魂的安息；}\BibleRef{玛 11:28} 至于其他诫命所吩咐的一切事，我们应遵守它们，其中毫无预像。因为我们已经按字面被教导不可崇拜偶像；而且那些吩咐我们不可妄称天主的名、孝敬父母、不可奸淫、不可杀人、不可偷盗、不可作假见证、不可贪恋近人的妻子、不可贪图近人一切所有的，\BibleRef{出 20:1} 都毫无预像或奥秘的意义，必须按字面遵守。但我们没有被命令按字面遵守安息日，像犹太人那样停止身体劳动而休息；甚至他们按照规定遵守的安息，除了作为另一种安息，即灵性安息的象征外，都应被视为可轻视的。由此我们可以合理地推论，圣经中一切用形象表述的事理，在激励那引导我们趋向安息的爱德上，都具有强大的力量；因为在十诫中唯一的形象性或预像性规条，就是那向我们举荐安息的规条，这安息虽为人各处所寻求，却唯独在天主内才确然、神圣地寻获。\par

% structure: p0036 source=h2 target=subsection reason=relative-heading-rank; base=h2

\subsection{第十三章}

23. 然而，主日并非启示给犹太人，而是启示给基督徒，借着主的复活，并从祂开始具有了那专属于它的节庆性质。因为虔敬死者的灵魂，确实在肉身复活之前处于安息的状态，但他们并不从事那种在身体复原后将要运用其身体力量的活动。而这种活动状态，由第八日（也是每周的第一日）作了预象，因为它并不结束那种安息，而是使之受荣显。因为随着身体重新结合，灵魂安息的任何障碍都不会重来，因为在复原的身体里没有腐朽：因为\BibleQuote{这可朽坏的，必须穿上不可朽坏的；这可死的，必须穿上不可死的。} \BibleRef{格前 15:53} 因此，虽然数字八作为复活象征的\OriginalTerm{圣事性}{sacramental}含义，对于那些充满先知精神的古圣先贤来说绝非隐藏（因为在\WorkTitle{圣咏集}的标题[第六篇和第十二篇]中，我们找到\BibleQuote{为第八}的话；婴孩在第八天受割损；在\WorkTitle{训道篇}中，联系两个盟约说道：\BibleQuote{把一份分给七，也分给八}）；然而在主的复活之前，这事被保留隐密着，惟有安息日被指定当守，因为在那个事件之前，确实有死者的安息（安息日的休息是其预象），但没有一个人从死者中复活、不再死亡、死亡不再统治他的实例；这样做乃是为了，从主的身体果真发生这种复活之时起（教会的\OriginalTerm{元首}{Head}先经历那祂的身体——教会——在时间终结时所期待的事），祂复活的那日，即第八日（也就是每周的第一日），应当开始被遵守为主日。同样的理由使我们能够理解，为什么关于守逾越节的日子——在那日子犹太人受命宰杀并吃羔羊，这显然预表了主的受难——并没有命令他们考虑每周的日子，等待安息日过去，并使月亮第三周的开始与正月第三周的开始相一致；原因是，主更愿意在祂自己的受难中，宣示那个日子的意义，正如祂来也是要宣示如今被称为主日的那个日子的奥迹，即第八日，也就是每周的第一日。\par

% structure: p0038 source=h2 target=subsection reason=relative-heading-rank; base=h2

\subsection{第十四章}

24. 现在请仔细思量这三日最神圣的日子，即因主的被钉十字架、坟墓中的安息和复活而著称的日子。这三者中，以十字架为象征的是我们现世生活的事务；而由祂的坟墓安息与复活所象征的，我们则以信德和望德持守。因为\Emph{现在}对每一个人所下的命令是：\BibleQuote{背着自己的十字架，跟随我。} \BibleRef{玛 16:24} 然而，当我们在世的肢体被克制，诸如邪淫、不洁、放荡、贪婪等，肉身便被钉在十字架上了。关于这些，宗徒在另一处说：\BibleQuote{如果你们随从肉性生活，必要死亡；然而如果你们依赖圣神，去致死肉性的妄动，必能生活。} \BibleRef{罗 8:13} 因此他也论到自己说：\BibleQuote{世界于我已被钉在十字架上了；我于世界也被钉在十字架上了。} \BibleRef{迦 6:14} 又说：\BibleQuote{知道我们的旧人已与他同钉在十字架上了，使那属罪恶的身体消逝，叫我们不再作罪恶的奴隶。} \BibleRef{罗 6:6} 我们的工作旨在削弱并摧毁罪恶之身的时期，在此期间，我们外在的人日渐消损，内在的人却日日更新——那就是十字架的时期。\par

25. 这些确实是善工，以安息为酬报，但它们同时又是劳苦痛苦的：因此我们受命\BibleQuote{在希望中要喜乐}，好使我们在默观未来的安息时，能以喜乐之心从事目前的劳苦。这种喜乐，由钉上双手的横木所代表的十字架广度是一个象征：因为我们认为双手象征工作，而广度象征工作者内心的喜乐，因为忧愁会使心神窘迫。在十字架的高处，就是头所靠的地方，我们得到一个象征，表示期待从天主至尊公义而来的酬报，祂\BibleQuote{要照每人的行为予以报应：凡恒心行善，寻求真荣、尊贵和不朽的人，赐以永生。} \BibleRef{罗 2:6} 因此，十字架的长度，就是整个身体伸展于其上的，是那种在天主圣意内恒忍耐的象征，为此那些忍耐的人被称为\OriginalTerm{坚忍}{long-suffering}者。还有深度，\Emph{即}那插入地下的部分，代表着神圣奥迹的隐秘性质。我想你们记得宗徒的话，在这十字架的描述中，我正试图阐释的：\BibleQuote{使你们在爱德上根深蒂固，建定基础，为能同众圣徒领悟基督的爱是怎样的宽、广、高、深。} \BibleRef{弗 3:17}\par

那些我们尚未看见或拥有，却凭\OriginalTerm{信德}{faith}和\OriginalTerm{望德}{hope}持守的事物，正是上述三个纪念日中第二日和第三日所表征的事件（即主的安息于坟墓和祂的\OriginalTerm{复活}{resurrection}）。但那些在此尘世生活中使我们忙碌，我们因敬畏天主而受诫命约束，如同被钉子钉透肉身（经上记载：\BibleQuote{求祢以敬畏祢的钉子钉牢我的肉身}）的事物，应被视为必要之事，而非为其本身值得渴求和贪恋之事。因此保禄说他切望离世与基督同在，那是好得无比的：\BibleQuote{然而，}他接着说，\BibleQuote{我活在肉身内却是为了你们的缘故。}\BibleRef{斐 1:23}——为你们的福祉是必要的。这离世与基督同在乃是不被中断之安息的开始，却因复活而受荣耀；这安息如今凭信德享受，\BibleQuote{因为义人必因信德而生活。}\BibleRef{哈 2:4}同一位宗徒说：\BibleQuote{你们岂不知道：我们受过洗归于基督耶稣的人，就是受洗归于祂的死亡吗？所以我们藉着\OriginalTerm{圣洗}{baptism}归入死亡，与祂同葬了。}\BibleRef{罗 6:3}如何成就？凭信德。因为只要我们还\BibleQuote{在心里叹息，等待义子之期望，即我们身体的救赎：因为我们得救是藉着望德；但所见的希望不是希望：因为人所看见的为什么还要盼望呢？但我们若盼望那未见的，就必坚忍等待。}这在我们身上尚未实际完成。\par

26. 你要记得我屡次向你重申：我们不要以为应该在此尘世生活中就获得幸福，脱离一切艰难，因而在遭受世事窘迫时，就可以亵渎地抱怨天主，好像祂没有履行自己的许诺。祂确实许诺了今生所需之物，但那些减轻我们当前困苦的安慰，与圆满真福者的喜乐大不相同。信徒说：\BibleQuote{我心里多忧多虑，上主，祢的安慰使我的灵魂欢欣。}因此，我们不要因艰难而抱怨；我们不要丧失喜乐的开阔心胸，因为经上记载：\BibleQuote{在希望中要喜乐，}其后续是：\BibleQuote{在患难中要忍耐。}\BibleRef{罗 12:12}因此，新生命目前是在信德中开始，并靠望德维持：因为只有等到这必死的被生命吞灭，死亡被胜利吞没；当最后的仇敌——死亡——被摧毁；当我们改变，肖似天使的时候，才得圆满：因为\BibleQuote{我们众人都会复活，但并非我们全部改变。}主又说：\BibleQuote{他们将相似天使。}\BibleRef{路 20:36}如今，我们凭信德怀着敬畏被祂所抓住；将来，我们要凭眼见在爱中抓住祂。因为\BibleQuote{当我们还存留在这肉身内，就是与主远离，因为我们行事为人是凭信德，不是凭眼见。}\BibleRef{格后 5:6}因此，保禄宗徒自己虽说过：\BibleQuote{我追求，是否可以抓住那曾抓住我的基督耶稣，}却也坦白承认自己尚未达到。他说：\BibleQuote{弟兄们，我并不以为我已经抓住了。}\BibleRef{斐 3:12}然而，我们的望德因应许的真实而稳固，所以他在别处说：\BibleQuote{我们藉着圣洗归入死亡，与祂同葬了，}接着又加上：\BibleQuote{为的是，正如基督藉着圣父的光荣从死者中复活，我们也要在新生活中度生。}\BibleRef{罗 6:4}因此，我们是在实际的劳作中行走，但在安息的盼望中；在旧生命的肉躯内，却在新生命的信德中。因为他又说：\BibleQuote{身体因罪而死，心神却因义德而生活。而且，如果那使耶稣从死者中复活者的\OriginalTerm{圣神}{Holy Spirit}住在你们内，那使基督从死者中复活者，也要藉着住在你们内的圣神，使你们有死的身体复活。}\par

27. 圣经的权威与普世教会的一致同意，共同规定了每年在\OriginalTerm{复活节}{Easter}以礼仪纪念这些事件，这些礼仪正如你现在所见，充满了灵性意义。从旧约圣经中，我们未获教导复活节的具体日期，仅知应在首月的第十四日至二十一日之间；但由于我们从福音书中确知，一周之中哪些日子依次因主的被钉十字架、安眠于坟墓及复活而成为圣日，所以\OriginalTerm{教父}{Fathers}们的\OriginalTerm{大公会议}{Councils}更附加规定，要求遵守这些日子，而整个基督教世界也一致相信，这是正确举行复活节的方式。\par

% structure: p0044 source=h2 target=subsection reason=relative-heading-rank; base=h2

\subsection{第十五章}

28. 四旬期的根据，在旧约中有梅瑟的禁食\BibleRef{出 34:28}和厄里亚的禁食\BibleRef{列上 19:8}，在福音中则有我主同样禁食四十天的事实\BibleRef{玛 4:2}，由此证明福音并非与法律和先知相悖。因为法律和先知分别由梅瑟和厄里亚的人物代表；主也曾荣耀地显现在他们中间于山上，为使宗徒所说关于他的话——\BibleQuote{他受到法律和先知的共同见证}\BibleRef{罗 3:21}——得以更清楚地显明。那么，一年之中，哪段时间更适合遵守四旬期的斋戒呢？莫过于紧接并邻近主受难时期的时段了。因为藉此斋戒，正象征着这劳苦的生命，其主要工作在于操练自制，弃绝世俗的友谊，这友谊永不停止地虚假地抚慰我们，并大量散布其迷人的诱惑。至于为何用数字40来象征这劳苦和自制的生活，我似乎看到数字10（其中包含我们真福的完满，如同数字8，因为它回归于单位）在这个数字中具有相似的地位[正如单位赋予数字8意义一样]；因此，我认为40这个数字是今生的合适象征，因为在其中，受造物（其象征数字为7）依附于造物主，在造物主内揭示了天主圣三的合一，这合一要在时间终结之前，被传扬于普世——这世界由四风鼓荡，由四元素构成，并经历一年四季的变化。现在，4乘10[7加3]等于40；但40这个数字若加上它的一个部分，就加上了数字10，[如同7加上它的一个部分，就加上了单位，]总和成为50——可以说，是劳苦和自制的酬报的象征。因为并非无缘无故，主自己在复活之后，曾在这地上、在今生中与门徒们共处四十天，并在升天之后，又过了十天，即在五旬节完全来到的日子，派遣了所恩许的圣神。此外，这第五十天还蕴含着另一个神圣奥迹：因为7个7天是49天。当我们回到另一个七天的开端，并加上第八天，这第八天也是一周的第一天，我们就完成了50天；这五十天的一段时期，我们在主复活之后庆祝，所代表的不是劳苦，而是安息和喜乐。因此，我们在此期间不禁食；祈祷时也直立，这姿势是复活的象征。因此，在五十天内的每一个主日，祭台上也遵守此惯例，并咏唱阿肋路亚，这表示我们未来的活动将完全在于赞美天主，如经上所说：\BibleQuote{上主，住在你殿宇内的人真有福！他们将永远（\Emph{即}无休止地）赞美你。}\par

% structure: p0046 source=h2 target=subsection reason=relative-heading-rank; base=h2

\subsection{第十六章}

29. 第五十天在圣经中也受到推荐；不仅在福音中，因为圣神在那一天降临，而且在旧约的书中也有推荐。因为在旧约中我们得知，当犹太人按照规定宰杀羔羊举行了第一个逾越节后，从那天起直到在西乃山上天主用\BibleQuote{天主的手指}将法律赐给梅瑟的日子，中间相隔了五十天；而这\BibleQuote{天主的手指}在福音中极其明确地被宣示为指圣神：因为一位福音作者引用主的话说：\BibleQuote{我若以天主的手指驱魔}\BibleRef{路 11:20}，另一位福音作者则引用说：\BibleQuote{我以天主的神驱魔}\BibleRef{玛 12:28}。谁不愿选择这些神圣奥迹所带来的喜乐，当它们将治疗的真理之光照射到灵魂上时，胜过这世间万国，哪怕这些国家处于完全的繁荣与和平之中？我们岂不能说，正如两个色辣芬在歌颂至高者时彼此应和：\BibleQuote{圣！圣！圣！万军的上主！}\BibleRef{依 6:3}，同样，旧约和新约以完美的和谐，一同为神圣真理作证？羔羊被宰杀，逾越节被庆祝，五十天后，那使人敬畏的法律被赐下，是由\BibleQuote{天主的手指}书写的。基督被宰杀，如依撒意亚所见证的，他像羔羊被牵去宰杀\BibleRef{依 53:7}；真正的逾越节被庆祝；五十天后，圣神被赐下，他就是\BibleQuote{天主的手指}，其果实是爱，因此他与那些只求自己利益的人相对立，那些人肩上荷着沉重的轭和重担，灵魂得不到安息；因为爱\BibleQuote{不求自己的益处}\BibleRef{格前 13:5}。因此，在异端者冷酷无爱的精神中是没有安息的，宗徒宣称他们的行为如同法郎的术士，说：\BibleQuote{正如雅乃斯和杨布勒反对梅瑟，这些人也反对真理；他们心思败坏，在信德上已被废弃。但他们不能再有所进展，因为他们的愚昧必将暴露在众人前，如同那两人一样}\BibleRef{弟后 3:8}。正是由于这种心思的败坏，他们极其不安，在第三个奇迹前失败了，承认在梅瑟内的天主之神是反对他们的：他们承认失败时说：\BibleQuote{这是天主的手指}\BibleRef{出 8:19}。圣神对温良和心里谦逊的人是和好与施恩的，并赐予他们安息；而对骄傲和高慢的人则是毫不留情的对头，使他们不得安宁。这种不安正由那些可憎的昆虫所预表，法郎的术士们承认自己被击败时，便说：\BibleQuote{这是天主的手指}。\par

30. 请阅读\WorkTitle{出谷纪}，并留意第一个逾越节与颁布法律之间的天数。\Place{西乃}旷野中，天主在第三个月的第一天对\Person{梅瑟}说话。那么，记下这个月的一天，再注意主在那一天所说的（除了其它事之外）：\BibleQuote{你到百姓那里，叫他们今天明天圣洁自己，洗净自己的衣服，第三天都应准备停当，因为第三天主要在众百姓前降到西乃山上。}\BibleRef{出 19:10} 律法于是在该月的第三日被颁布。现在计算从第一个月的第十四天，即逾越节那天，到第三个月的第三天之间的天数：这样你就得到第一个月的十七天，第二个月的三十天，第三个月的三天——总共是五十天。约柜中的律法象征主身体内的圣洁，藉着他的复活，应许了我们未来的安息；为了使我们能领受这安息，圣神将爱德倾注在我们心中。然而那时圣神还没有赐下，因为耶稣还没有受到光荣。\BibleRef{若 7:39} 因此那首预言诗篇说：\BibleQuote{上主，请起来，驾临你的安息之所，你与你的约柜！}[七十士译本作\Quote{圣洁}]。哪里有安息，那里就有圣洁。因此，我们现在已领受了这安息的凭据，好使我们能爱慕它、渴望它。因为那属于来世的安息，是我们因逾越节所象征的、脱离此世的过渡而得以进入的；如今，一切人都因圣父、圣子及圣神之名而被邀请进入。\par

% structure: p0049 source=h2 target=subsection reason=relative-heading-rank; base=h2

\subsection{第十七章}

31. 因此，我主复活后，为显示这新生命，命人在船的右边捕到的大鱼数目中，也出现了五十的三倍再加三[天主圣三的象征]，使这神圣的奥迹更显明；而门徒的网却没有破，\BibleRef{若 21:6} 因为在那种新生命中，不再有因异端者的扰乱而引起的分裂。于是，在这种新生命中，人既达到成全和安息，身体和灵魂都被纯净的天主圣言所净化——这圣言如同炼净渣滓、七次提炼的白银——必将领受他的赏报，就是\OriginalTerm{德纳}{denarius}；\BibleRef{玛 20:9} 因此，藉着那个赏报，数字10与7便在他身上结合。因为在数字17里，就像在其他具有象征意义的组合数字里一样，蕴含着一个奇妙的奥迹。第十七首圣咏（\WorkTitle{列王纪}中唯一完整记载的圣咏）之所以被完整保存，也不是没有理由的，因为它象征那个我们将不再有仇敌的王国。它的标题是：\Quote{达味的圣咏，在上主救他脱离一切仇敌和撒乌耳手的那一日。} 达味是谁的预像呢？岂不是那按肉身出自达味后裔的那一位吗？\BibleRef{罗 1:3} 他在他的教会内，即在他的身体内，仍然忍受着仇敌的恶意。因此，那从天上临到迫害者的耳朵的话——耶稣用他的声音将他击倒，并把他转变成他身体的一部分（就如我们所用的食物成为我们身体的一部分一样）——就是这些话：\BibleQuote{扫禄，扫禄，你为什么迫害我？}\BibleRef{宗 9:4} 而他的这个身体何时才能最终脱离仇敌呢？岂不是当最后的仇敌——死亡——被毁灭的时候吗？那153条鱼的数字正是与那个时刻相关的。因为如果数字17本身是一个算术三角形的边长，该三角形由上面一行行单位数构成，数目从1逐行增加到17，那么这些单位数的总和就是153：因为1加2得3；3加3得6；6加4得10；10加5得15；15加6得21；如此继续，一直加到17，总和是153。\par

32. 因此，复活节和五旬节的庆祝，是以圣经为最稳固基础的。至于复活节前四十天的守斋，已由教会的实践所确认；同样，新教友的八日庆期也按此顺序安排，即第八日与第一日相合。在教会内仅于这五十天歌唱\Quote{阿肋路亚}的习俗，并非普遍遵守；因为在其他地方，其他时间也唱，但在这几天却到处都唱。至于在这些日子以及所有的\Quote{主日}站着祈祷的习俗是否处处遵守，我不得而知；但无论如何，我已经告诉你们教会在这方面的做法所依据的准则，而且我认为足够明显。\par

% structure: p0052 source=h2 target=subsection reason=relative-heading-rank; base=h2

\subsection{第十八章}

33. 关于洗脚礼，既然主亲自举荐它，作为他来教导的那种谦卑的榜样，正如他后来所解释的，便产生了这样的问题：究竟在什么时候，按字面实行此事，来公开教导它所阐释的重要义务最为适宜；而这段时间 [四旬期] 被提出，为的是使其中所教导的教训能产生更深刻、更严肃的印象。然而，许多人并不接受这种习俗，恐怕它被误会为属于圣洗的仪式；有些人甚至毫不迟疑地完全否认它在我们仪式中的位置。不过，另有些人，为将这礼仪与这四旬期的更神圣的联想结合起来，同时又无论如何避免与圣洗混淆，便选择了在第八日，或第三个第八日举行这礼仪，因为在许多神圣奥迹中，数字三具有重大的意义。\par

34. 你表示希望我就那些在不同国家有不同规定的仪式写些东西，这令我惊讶，因为此事并无必要；而且，对于这些习俗，当它们绝不与真道或健全道德相悖，反倒包含某些美好生活的激励时，就有一项极好的规则当遵守：即，无论我们在哪里看到它们被遵守，或知道它们被确立，我们便不仅不该指摘它们，更该以赞同和效仿来举荐它们，除非因他人的软弱，怕这样做弊大于利而不敢推荐。然而，若赞同那些热心此仪式的人所预期的利益，甚于因反对者不悦而须顾虑的损失，我们就不该因而却步。在这种情况下，我们务必采用它，尤其当有圣经可引以为证时：例如唱圣咏和赞美诗，对此，我们在圣经上存有主及其宗徒的榜样和训令。在这项大有助于引导虔敬心境、炽热对天主之爱的圣操上，各地用法不同，而在\Place{阿非利加}，教会的成员对此多少有些冷漠；正因如此，多纳图派责备我们在教堂中，用庄重的吟唱来咏唱先知们的圣歌，他们却在狂欢中，用人为的圣咏来挑动激情，如同战场上激荡的号角声。但当弟兄们聚集在教堂中时，为什么不应将时间用于歌唱圣歌呢？当然，在读经或讲道，或主礼者高声祈祷，或执事的声音带领会众同祷时除外。在其他不占用的间歇，我看不出对一个基督徒会众有什么比这更卓越、更有益、更神圣的操练。\par

% structure: p0055 source=h2 target=subsection reason=relative-heading-rank; base=h2

\subsection{第十九章}

35. 然而，我不能赞成那些背离教会习俗、并以象征某种神圣奥迹为借口而设立的仪式；尽管为了避免冒犯某些人的虔诚和另一些人的好斗，我不敢严厉谴责许多此类事情。但我痛惜（且太常有机会痛惜）的是，对圣经所命令的许多最有益处的礼仪，人们关注甚少；而各处又盛行着如此多的错误观念，以至于在\OriginalTerm{八日庆期}{octaves}（洗礼前）赤脚触地的人，会比一个沉溺于醉乡的人受到更严厉的责斥。因此，我的意见是：但凡可能，便应毫不迟疑地废除这一切；它们既无圣经的凭据，也不见由主教会议规定，更未被普世教会的实践所确认，却随着各地不同的习俗而千差万别，以至于即使有可能，也很难发现当初设立它们的缘由。因为即使或许找不到什么违反真信德之处，然而基督的宗教——天主在仁慈中使其自由，并赋予数目极少且极易遵守的圣事——竟被这些累人的仪式所压制，以致犹太教会本身的处境反倒更可取了：因为他们虽未认识自由的时代，却臣服于天主法律所加的重担，而非人的虚妄想法。不过，天主的教会既在目前情况下容纳着大量糠秕和莠子，便容忍许多事；但若有任何事违反信德或圣洁生活，她既不默许，也不实行。\par

% structure: p0057 source=h2 target=subsection reason=relative-heading-rank; base=h2

\subsection{第二十章}

36. 因此，你写到有关某些弟兄以礼仪上不洁为由，戒用动物性食物一事，这极其相反信德和健全道理。我若对此事进行任何形式的详尽讨论，或有人会以为宗徒在这事上的训令有模糊之处；其实，在众多相关言论中，他用以下的话表达了对异端者这种意见的憎恶：\BibleQuote{圣神明明地说：在最后的时期，有些人要背弃信德，听信欺诈的神和魔鬼的训言，这训言是出于那些伪善的说谎者，他们的良心已烙上了烙印。他们禁止嫁娶，戒避一些食物；这些食物都是天主所造，叫那信仰而认识真理的人，以感恩的心所享用的。因为天主所造的样样都好，如以感恩的心领受，没有一样是可摈弃的；因为样样都是藉天主的话和祈祷祝圣了的。} \BibleRef{弟前 4:1} 在另一处，他又论及这些事说：\BibleQuote{为洁净人一切都是洁净的，但为败坏的人和不信的人没有一样是洁净的，就连他们的理智和良心都是污秽的。} \BibleRef{铎 1:15} 你自己去读其余的，并把这些经文读给别人听——能读多少就读多少——好叫他们既蒙召获得自由，就不使天主的恩宠在他们身上落空；只让他们不要以自由作为放纵肉欲的机会；更不要托词说在迷信或不信的驱使下这样做是不合法的，就拒绝实行克制肉欲的目的，即戒避某些食物。\par

37. 至于那些随手翻阅福音书的一章一节以预卜未来的人，尽管这样做比去求问占卜的神灵要好，然而在我看来，试图将那些旨在教导我们更高生命的圣言，转用于世俗事务和此生之虚幻，仍是一种应受谴责的做法。\par

% structure: p0060 source=h2 target=subsection reason=relative-heading-rank; base=h2

\subsection{第二十一章}

38. 若你认为我此刻对你的问题答复得尚不够多，那你必是对我的能力或职责所知甚少。因为我远非如你所想的那样通晓万事；我在你的信中读到那段话时，心中最为忧伤，这不仅因为它显然绝非事实，也因为我惊讶于你竟不知道，不仅在其他无数领域中有许多事为我所不知，就连在圣经中，我所不知的事也远比我所知的为多。然而，我在基督的名内怀着希望，这希望并非没有酬报，因为我不仅相信了我的天主的证言：\BibleQuote{这两条诫命是全部法律和先知的总纲}；\BibleRef{玛 22:40}，而且我还借着经验亲自证明了这一点，并且日日都在证明。因为没有任何一个神圣的奥迹，也没有任何一段天主圣言的艰深章节，在向我阐明之时，我不从中发现同样的这些诫命：因为\BibleQuote{这诫命的终向就是爱德，即由纯洁的心、无愧的良心及无伪的信德所发出的爱德}；\BibleRef{弟前 1:5}，而\BibleQuote{爱就是法律的圆满}。\BibleRef{罗 13:10}\par

39. 因此，我恳求你，我至亲爱的，无论你在研习这些著作还是其他著作，都要这样阅读和学习，好能铭记那句非常真实的话：\BibleQuote{知识只会使人傲慢自大，爱德才能立人}；\BibleRef{格前 8:1} \BibleQuote{爱德不誇张，不自大}。因此，让知识被用作一种脚手架，借着它来建造爱德的殿宇，这殿宇将在知识废去时永远长存。\BibleRef{格前 13:4}知识，若作为达到爱德的方法来使用，便是大有裨益的；但若离开这个崇高的目的，它已被证实不但是多余的，甚至是有害的。我知道，圣善的默想如何保守你安息在我们天主的翅膀的荫下。我之所以讲述这些，尽管很简短，是因为我知道你那\Quote{不誇张}的爱德，会促使你将这封信借给许多人传阅。\par

\SourceRecord{Translated by J.G. Cunningham.；Nicene and Post-Nicene Fathers, First Series；Vol. 1.；Edited by Philip Schaff.；Buffalo, NY: Christian Literature Publishing Co.,；1887.；<http://www.newadvent.org/fathers/1102055.htm>.}

% structure: p0001 source=h1 target=section reason=page-title

\section{信函58（主历401年）}

\Emph{致我尊贵可敬的阁下\Person{帕玛基乌斯}，我在基督的心肠内所挚爱的儿子，\Person{奥古斯丁}在主内致意。}\par

1. 源于基督恩宠的善工，使你在我们这些祂的肢体面前获得了受尊敬的名分，并让你真正为我们所认识，也为你所能地为我们所挚爱。因为即使我天天见到你的面容，这也不能使我对你的认识比现在更完全，因在你某一件行为的光辉中，我已看见了你的内心，因和平之可爱而美好，因真理之光辉而闪耀。看见这些，使我认识了你；认识了你，便使我爱你；因此，在这封信中，我是写给一个虽彼此相隔遥远，却已为我所认识，并成为我心爱朋友的你。那维系我们的纽带，其实日期更早，我们曾生活在同一元首之下：因为如果你没有植根于祂的爱内，\OriginalTerm{天主公教}{Catholic}的合一就不会对你如此宝贵，你就不会如同己任般对待那些定居在\Place{努米底亚}执政官行省中心的你的非洲佃户——恰恰就是\OriginalTerm{多纳特派}{Donatists}愚妄肇始之处——用那样的言语劝勉他们，用那样的热忱鼓励他们，以致说服了他们毫不犹豫地献身于他们所信、一个有你这样品格和地位的人，除非出于已查明和承认的真理，决不会选择的道路，并且使他们虽离你如此遥远，仍服从了同一元首；这样他们就得以和你一起，永远被视为祂的肢体，而他们此时因祂的安排正暂时依赖于你。\par

2. 因此，因这一事件而认识了你，我便怀着喜悦之情，在我的主耶稣基督内祝贺你，并给你寄这封信，作为我心灵对你的爱的一个证明；因为我所能做的不能更多。不过我恳求你，不要用这封信来衡量我的爱之多寡；而要借着这封信，当你读完之后，循着思想开辟的无形通道进入我的内心，看看那里对你怀着怎样的感情。因为在爱的目光前，那爱的圣所将被揭示，当我们在这圣所中朝拜天主时，我们关闭了它，使之不受这世上纷扰琐事的打扰；在那里你将看到我因你的善工而洋溢的狂喜，这狂喜我无法用舌和笔形容，在对那以启发使你愿意、并以助佑使你能如此事奉祂的天主的赞美祭献中，闪耀燃烧。\BibleQuote{感谢天主，为祂无可言喻的恩赐！} \BibleRef{格前 9:15}\par

3. 哦，我们在非洲多么渴望看到，有许多像你一样身为国家元老和圣教会之子的人，能做出你这使我们快慰的事啊！然而，给他们这样的劝勉是危险的：他们可能拒绝听从，而教会的敌人就会利用这一点欺骗弱小者，仿佛他们在那些漠视我们劝告之人的心里，胜过了我们似的。但对你表达感激却是安全的；因为你已经做了，在解救弱小者的过程中，教会的敌人已被击败。因此，我想这就足够了：请你以友善的勇气，向任何因他们的基督徒身份而你能够这样读给他们的人读这封信。因着得知你所成就的，他们就会相信，那他们如今因视为不可能而漠不关心的事，在非洲是能够做到的。至于这些异端者以乖谬心思所设下的陷阱，我已决意不在这封信中提及，因为我只是感到好笑，他们竟以为能占得上风胜过你那基督据为己有的心灵。然而，你会从我恳切推荐给你的弟兄们那里听到他们的言论：他们担心你会不屑于某些事，那些事在你看来，可能为那些因你的工作而使其天主教慈母欢欣的人们的伟大且意想不到的救恩，似乎是多余的。\par

\SourceRecord{Translated by J.G. Cunningham.；Nicene and Post-Nicene Fathers, First Series；Vol. 1.；Edited by Philip Schaff.；Buffalo, NY: Christian Literature Publishing Co.,；1887.；<http://www.newadvent.org/fathers/1102058.htm>.}

% structure: p0001 source=h1 target=section reason=page-title

\section{书信五十九（公元401年）}

\Emph{致我最蒙福的主、可敬的父亲\Person{维克托里努斯}——我在司铎职内的兄弟，\Person{奥斯定}在主内致候。}\par

1. 你的\OriginalTerm{会议}{Council}召集令于\OriginalTerm{十一月望日前五天}{the fifth day before the Ides of November}的傍晚送达我处，当时我身体极为不适，因此无法出席。不过，我请你以虔敬而明智的判断考虑：是召集令引起的某些困惑是由于我自己的无知，还是有充分的理由。我在召集令中读到，该信也发给了\Place{毛里塔尼亚}的教区，而据我们所知，那些教区有自己的首席主教。那么，如果这些省份要参加在\Place{努米底亚}举行的\OriginalTerm{会议}{Council}，在通函上附上在\Place{毛里塔尼亚}的一些较显赫主教的名字，是绝对合宜的；而我没有看到这一点，感到非常惊讶。而且，该信发给\Place{努米底亚}主教们的方式是如此混乱和草率，以至我发现我的名字排在第三位，尽管我知道我在主教名册上的顺序理应更靠后得多。这既错了待了别人，也让我痛心。再者，我们可敬的父亲和同僚，\Place{塔戈萨}的\Person{克桑提普斯}，声称\OriginalTerm{首席权}{primacy}属于他，并且很多人尊他为首席主教，他也发出如你所发的此类信函。即使这是错误，你\OriginalTerm{尊座}{Holiness}可以轻易发现并纠正，但他的名字当然不应该在你发出的召集令中被遗漏。如果他的名字被放在名单中间，而不是第一行，我会非常惊讶；那么，当我发现召集令中根本没有提及他时，我的惊讶又大了多少呢？他比任何人都更应该出席这次\OriginalTerm{会议}{Council}，以便在所有\Place{努米底亚}各教会的主教们面前，首先讨论这个\OriginalTerm{首席权}{primacy}顺序的问题！\par

2. 鉴于这些理由，我甚至可能会犹豫是否来参加\OriginalTerm{会议}{Council}，唯恐这份充满诸多明显错误的召集令是伪造的；即使我没有因通知仓促以及许多其他重要事务的阻碍而无法成行。因此，我恳请你，最蒙福的教长，原谅我，并请首先注意促成你\OriginalTerm{尊座}{Holiness}与年长的\Person{克桑提普斯}之间就你们中谁应当召集\OriginalTerm{会议}{Council}的问题达成真诚的相互理解；或者，至少——我认为这甚至更好——让你们两位，在不预判任何一方主张的情况下，联合召集我们的同僚们，特别是那些几乎和你们一样长期担任\OriginalTerm{主教职}{episcopate}的人，他们很容易发现并决定你们中谁的主张符合真理，以便这个问题先在你们少数几位中解决；然后，错误纠正之后，再让较年轻的主教们聚集起来，他们目前除了你们之外没有其他可以或应当让他们作为此事见证的人，因而暂时无法得知应当优先考虑你们中的哪一位。\par

我已寄出此信，封缄所用戒指上刻有一男子侧面像。\par

\SourceRecord{Translated by J.G. Cunningham.；Nicene and Post-Nicene Fathers, First Series；Vol. 1.；Edited by Philip Schaff.；Buffalo, NY: Christian Literature Publishing Co.,；1887.；<http://www.newadvent.org/fathers/1102059.htm>.}

% structure: p0001 source=h1 target=section reason=page-title

\section{第60封书信（公元401年）}

\Emph{致\Person{奥勒留}神父，我至福的主教，因着极应得的尊敬备受敬重，我在司铎职分上的至亲兄弟，\Person{奥古斯丁}在主内致问候。}\par

1. 自从我们分别后，我不曾收到过主教阁下的任何信函；但现今我已读到了主教阁下关于\Person{多纳图斯}及其兄弟的信，长久以来犹豫着应该给予怎样的答复。我屡次省思，在此等情况下，何种做法有助于我们在基督内所服事并力求在祂内培育的那些人的利益，却从未想出任何能使我改变主意的理由，即我不认为应该给天主的仆人们以口实，使他们以为，晋升到更好的职位更轻易地赐给了那些变得更坏的人。这样的规则会使隐修士们对失足少存戒心，而如果让那些背弃了自己作为隐修士职责的人被拣选去担任神职，那对圣秩便构成了极大的不公，因为我们通常只从那些继续忠于隐修圣召的人中，拣选更为久经考验且更为优越的人来充任此职——除非，也许要教导普通民众来取笑我们，说什么\Quote{一个坏的隐修士造就一个好的神职人员}，正如他们惯常说\Quote{一个拙劣的吹笛者造就一个好的歌手}。倘若我们助长隐修士们滋生出如此致命的骄傲，并同意使我们自身所属的圣秩蒙受如此严重的耻辱，那将是难以承受的灾祸：因为就连一个好的隐修士有时也难于成为一个好的神职人员，他虽善于克己，却可能缺乏必要的学识，或因个人的某种缺陷而不合格。\par

2. 然而，我以为主教阁下当初认为这些隐修士是在我的同意下离开隐修院的，为使他们能更有效地服务于自己地区的人民；但这并非事实：他们是擅自离开的，擅自离弃了我们，尽管我为了他们的益处而尽力加以反对。至于\Person{多纳图斯}，既然他在我们能够在公会议上就他的案子作出任何决定之前，便已获得了圣秩，就请主教阁下按您的智慧所指引的去做吧；或许他那骄傲的顽固已被克服了。但至于他的兄弟，就是导致\Person{多纳图斯}离开隐修院的主要原因，我不知道该写些什么，因为主教阁下知道我对他的看法。我绝不敢反对像您这样具有智慧、地位和虔诚的人所认为最好的做法；而且我全心希望，您会做您所判断为对教会的肢体最有益的事。\par

\SourceRecord{Translated by J.G. Cunningham.；Nicene and Post-Nicene Fathers, First Series；Vol. 1.；Edited by Philip Schaff.；Buffalo, NY: Christian Literature Publishing Co.,；1887.；<http://www.newadvent.org/fathers/1102060.htm>.}

% structure: p0001 source=h1 target=section reason=page-title

\section{第61封信（公元401年）}

\Emph{致他至爱且可敬的兄弟\Person{西奥多鲁斯}，主教\Person{奥古斯丁}在主内致意。}\par

1. 我决定在这封信中写下当你我一起交谈时所说的话，关于我们以何种条件接纳希望成为天主教徒的\OriginalTerm{多纳图派}{Donatist}神职人员，以便如果有人问起你我们在这方面的看法和做法，你可以展示出我亲手写下的这些话。因此，你要确信，在多纳图派神职人员中，我们只憎恶那使他们成为\OriginalTerm{裂教者}{schismatics}和\OriginalTerm{异端者}{heretics}的东西，即他们与\OriginalTerm{天主教会}{Catholic Church}的合一和真理分离，因为他们不与遍布世界的天主子民保持和平，并且拒绝承认在领受过的人中所领受的基督的\OriginalTerm{圣洗}{baptism}。因此，我们弃绝他们这个严重的错误；但我们承认在他们身上所拥有的天主的美名，以及他们所领受的\OriginalTerm{圣事}{sacrament}，并以敬畏和爱慕接纳之。正因如此，我们为他们的迷途而痛心，并渴望以基督的爱为他们赢得天主，使他们能在教会的和平内拥有那为他们\OriginalTerm{救恩}{salvation}的圣事，这圣事他们目前却在教会之外拥有，导致他们的毁灭。因此，如果从我们之间除去那些出于人的恶事，并适当尊崇那来自天主、属于双方共有的美善，就会产生兄弟般的和谐，如此可爱的和平，以致基督的爱将在人们心中战胜魔鬼的诱惑。\par

2. 因此，当任何人从多纳图派来到我们这里时，我们不欢迎属于他们的恶事，即他们的错误和\OriginalTerm{裂教}{schism}：这些，我们之间唯一的障碍，从我们之间除去，我们就拥抱我们的兄弟，与他们一同站立，正如宗徒所说，在\BibleQuote{圣神的合一，在和平的联系中}\BibleRef{弗 4:3}，并承认在他们身上属于神的美善，如他们的\OriginalTerm{圣洗}{baptism}、\OriginalTerm{祝圣}{ordination}所赋予的祝福、他们\OriginalTerm{克己}{self-denial}的宣认、他们\OriginalTerm{独身}{celibacy}的誓愿、他们对\OriginalTerm{天主圣三}{Trinity}的\OriginalTerm{信德}{faith}，以及类似的一切；所有这些事情，他们从前也有，但\Quote{对他们毫无益处，因为他们没有\OriginalTerm{爱德}{charity}}。因为，如果不拥抱基督徒的合一，一个宣认基督徒爱德的人，又有什么真理呢？因此，当他们来到天主教会时，他们获得的不是他们原本拥有的，而是他们以前所没有的——也就是说，那些他们原本拥有的东西，从那时起开始对他们有益。因为在\OriginalTerm{天主教会}{Catholic Church}中，他们获得了\OriginalTerm{爱德}{charity}的根，在和平的联系和合一的共融中：这样，他们所拥有的一切真理的\OriginalTerm{圣事}{sacraments}，不再用来定他们的罪，而是使他们得救。枝子不应夸耀自己的木料是葡萄树的木料，不是荆棘的木料；因为如果他们不借着与根相连而生活，他们虽然外表如此，也必将被投入火中。但关于一些被折断的枝子，宗徒说，\BibleQuote{天主能重新把他们接上去}\BibleRef{罗 11:23}。因此，亲爱的兄弟，如果你看见多纳图派的任何人怀疑我们将以何种地位接纳他们，就把我亲笔写的这封信给他们看，这是你所熟悉的，如果他们愿意，就让他们阅读；因为\Quote{我呼求天主为我的灵魂作证}，我将以这样的条件接纳他们：他们不仅保留所领受的基督的\OriginalTerm{圣洗}{baptism}，也保留对他们圣洁的誓愿和\OriginalTerm{克己}{self-denial}的美德应有的尊重。\par

\SourceRecord{Translated by J.G. Cunningham.；Nicene and Post-Nicene Fathers, First Series；Vol. 1.；Edited by Philip Schaff.；Buffalo, NY: Christian Literature Publishing Co.,；1887.；<http://www.newadvent.org/fathers/1102061.htm>.}

% structure: p0001 source=h1 target=section reason=page-title

\section{第62封信（公元401年）}

\Emph{\Person{阿利比乌斯}、\Person{奥斯定}、\Person{撒姆苏奇乌斯}，以及与他们在一起的弟兄们，在主内向\Person{塞维鲁}致候}，\Emph{我们的主，最可称颂的，并以一切崇敬最可爱的，真理上的兄弟，司铎职上的同伴，以及所有与他在一起的弟兄们。}\par

1. 当我们来到\Place{苏布萨纳}，并调查在我们离开期间、违背我们意愿所发生的一切时，我们发现有些事情正如我们听说的，有些事情却不然，但所有事情都令人悲痛，需要忍耐；我们尽力，靠着主的帮助，藉着责备、劝戒和祈祷，予以纠正。最令我们难过的是，自从你离开那里之后，那些从那里去你那里的弟兄们竟被允许在没有向导的情况下出发，为此我们请求你原谅，因为此事并非出于恶意，而是出于过分的谨慎。因为他们相信这些人是我们的儿子\Person{弟茂德}派来，为要使你不悦，于是他们希望将整个事情保留原状直到我们到来（到时候他们指望能与你一同见我们），他们认为如果不给他们配备向导，这些人的离开就会被阻止。他们这样试图扣留弟兄们，我们承认是做错了——的确，谁能怀疑呢？由此也产生了讲给\Person{福索尔}的那个说法，说\Person{弟茂德}已经和这些弟兄们一同到你那里去了。这完全是假的，但那陈述并非由那位长老所出；而我们的兄弟\Person{卡尔凯多纽斯}对这一切毫不知情，这一点已由所有能证明此类事情的方式，向我们最清晰地证明了。\par

2. 但为什么在这些事情上花费更多时间呢！我们的儿子\Person{弟茂德}，因为发现自己完全违背意愿地陷入了如此出乎意料的困惑而极为不安，他告诉我们，当你催促他在\Place{苏布萨纳}事奉天主时，他激烈地脱口而出，并发誓他绝不会在任何情况下离开你。当我们问及他目前的愿望时，他回答说，由于这个誓言，他不能去我们先前希望他去的那个地方，即便有证据表明他并非出于勉强，他的心也得安定了。我们向他指出，如果通向与你在一起的途径，并非由他，而是由你加以阻拦，为的是避免恶表，那么他并不会犯违背誓言的罪；因为他凭誓言只能约束他自己的意志，而非你的，他也承认你并未以誓言相互约束；最后，他说，正如一位天主的仆人和教会之子所当说的那样，他将毫无犹豫地同意，由我们，连同你的圣善，就他所作出的任何决定。因此，我们请求，并藉着基督的爱恳求你，凭你的明智，记念我们在这件事上彼此谈论的一切，并以你对此信的回复使我们欢欣。因为\BibleQuote{我们强壮的人}（倘若我们在这许多且危险的诱惑中，还真敢妄用此名）按宗徒所说的，有责任\BibleQuote{担待不强壮人的软弱}。\BibleRef{罗 15:1} 我们的兄弟\Person{弟茂德}没有写信给您的圣善，因为您可敬的兄弟已向您全体禀报。愿您在主内喜乐，并记念我们，我们的主，最可称颂的，并以一切崇敬最可爱的，真诚上的兄弟。\par

\SourceRecord{Translated by J.G. Cunningham.；Nicene and Post-Nicene Fathers, First Series；Vol. 1.；Edited by Philip Schaff.；Buffalo, NY: Christian Literature Publishing Co.,；1887.；<http://www.newadvent.org/fathers/1102062.htm>.}

% structure: p0001 source=h1 target=section reason=page-title

\section{第63封书信（主历401年）}

\Emph{致\Person{塞维鲁}——我至为真福、可敬的主，一位配得以诚挚之爱拥抱的弟兄，司铎职的同工，并致与他同在的弟兄们：\Person{奥斯定}及与他同在的弟兄们在主内致候。}\par

1. 如果我坦率地说出此案迫使我必须说的一切，你或许会问我，我顾念维护爱德之心何在；但如果我不能说出此案所要求的一切，我难道不能问你，友谊所赋予的自由又在哪里？踌躇于此两者之间，我选择写下那些足以自辩而不致控诉你的话。你曾写道，你对我们深为已行之事悲痛，却仍默许之感到惊讶，因那原本可由我们的纠正来补救；就好像事后再尽可能加以纠正之后，那些错误就不应再被痛惜了；更仿佛我们默许一件虽已错误施行却无法挽回之事，是荒谬的。因此，我真诚敬爱的弟兄，你的惊讶可以休矣。因为\Person{弟茂德}是在\Place{苏布撒纳}被祝圣为\OriginalTerm{五品}{subdeacon}的，此事既违我劝亦非我愿，且当时他的案件作为我们之间商议与讨论的主题，其裁决尚未决定。看哪，我至今仍为此悲痛，尽管他已回到你那里；而我们并不后悔在同意他回去时服从了你的意愿。\par

2. 愿你喜欢听闻，早在他离开此地之前，我们便已通过责备、劝诫和祈祷，纠正了所发生的错误，免得你以为那时尚未纠正任何错误，只因他尚未回到你那里。通过\Emph{责备}，我们首先针对\Person{弟茂德}本人进行训诫，因为他没有服从你，而是在没有征询我们弟兄\Person{卡尔切多尼乌斯}的情况下，就直接去见你的尊前，而他的这一举动正是这苦恼的根源；随后我们指责了\OriginalTerm{司铎}{presbyter}\Person{卡尔切多尼乌斯}和\Person{维里努斯}，因为我们发现他们策划了\Person{弟茂德}的祝圣。当所有这些人在我们的责备下承认了所有被指控的错事并请求宽恕时，我们若拒绝相信他们已充分悔改，那便是过于傲慢了。因为他们无法使已行之事变得未曾发生；而我们通过责备所期望或渴望的，无非是让他们承认错误并为之悲伤。通过\Emph{劝诫}：首先，警告所有人永远不再敢做这样的事，以免招致天主的义怒；然后特别吩咐\Person{弟茂德}，他曾说他只是因誓言所缚才去你尊前，若你尊驾考虑到我们共同商讨此事的全部内容，可能决定不将他留在身边——为顾及\BibleQuote{基督曾为之而死的那软弱的人}，他们可能因之跌倒，也为着教会的纪律，忽视它是危险的，并且他已在本地教区开始担任读经员——那么，他就应摆脱誓言的束缚，以不受干扰的心志献身于天主的服务，我们都必须为自己的作为向天主交账。通过这些我们所能给予的劝诫，我们也说服了弟兄\Person{卡尔切多尼乌斯}，完全顺从地接受为维护教会纪律而必须对他采取的任何处置。通过\Emph{祈祷}，我们更努力纠正我们自己，将我们计划的指引和结果交托于天主的仁慈，并祈求若有任何罪性的愤怒伤害了我们，我们可藉投奔祂治愈的右手而得痊愈。看，我们通过责备、劝诫和祈祷，已纠正了多少！\par

3. 而今，念及爱德的纽带，\BibleQuote{免得我们为撒殚所利用——因为我们并非不知道他的诡计}\BibleRef{格后 2:11}——除了服从你的愿望，我们还能做什么呢？因为你认为，已行之事除了我们将此人交还于你的权下，便无法补救；你正是在他身上抱怨自己受了委屈。即使是我们的弟兄\Person{卡尔切多尼乌斯}本人也同意了这一做法，虽并非毫无心灵上的极大痛苦——为此我恳请你为他祈祷——但最终还是毫无反对，相信他在顺服你时便顺服了基督。更有甚者，当我还在考虑是否应给你写第二封信时，我的弟兄，在\Person{弟茂德}还在此地的期间，他本人出于子女般的敬畏，害怕惹你不悦，便终止了我的思量，不仅同意，甚至请求，将\Person{弟茂德}归还于你。\par

4. 因此，\Person{塞维鲁}弟兄，我将我的事态交由你来裁定。因为我确信基督居在你心中，我藉着祂恳求你，向祂请训，并将你的心意顺服于祂的指导，去思想这个问题：当一个人已开始在托付给我照管的教会中担任\OriginalTerm{读经员}{reader}，他不仅在\Place{苏布撒纳}读过一次经，而是第二次、第三次，并且与\Place{苏布撒纳}教会的司铎一起，在\Place{图雷斯}、\Place{基萨}和\Place{韦尔巴利斯}也做了同样的事，那么，宣称他从未作过读经员，这是可能的，或是正当的吗？我们既已顺从天主，纠正了后来违背我们意愿所行的事，那么，请你同样顺从祂，纠正先前因你不了解案情而错误施行的那件事。因为我不担心你会不明白，若某一教会的圣职人员向另一教会的圣职人员发誓不会离开他，而后者却鼓励他留在自己身边，声称这样做是为了不使誓言遭破坏，那么这将为破坏教会纪律打开多大的破口；其实，禁止此事并拒绝让那人留在自己身边的人（因为那人只能通过其誓言约束他自己的良心），无疑是以一种无人能公正指责的方式，维护了和平所需的秩序。\par

\SourceRecord{Translated by J.G. Cunningham.；Nicene and Post-Nicene Fathers, First Series；Vol. 1.；Edited by Philip Schaff.；Buffalo, NY: Christian Literature Publishing Co.,；1887.；<http://www.newadvent.org/fathers/1102063.htm>.}

% structure: p0001 source=h1 target=section reason=page-title

\section{书信六四（主历四〇一年）}

\Emph{致我主 \Person{昆提安努斯}，我最亲爱的兄弟及同铎，\Person{奥斯定} 在主内致以问候。}\par

我们不轻看那些在美丽上有缺陷的身体，尤其因为我们的灵魂本身也尚不如我们期望的那么美丽，我们期望当那位拥有不可言喻之美者显现时，我们将变得美丽；虽然现在我们看不见祂，但我们信祂；因为那时，\BibleQuote{我们必要相似祂，因为我们要看见祂实在怎样。}\BibleRef{若一 3:2} 你若以友善和兄弟般的精神接纳我的劝告，我就劝你如此看待你的灵魂，如同我们看待自己的一样，不要狂妄地想象它已是完美无瑕；但如宗徒所命，\BibleQuote{论望德，要喜乐；在困苦中，要忍耐；}\BibleRef{罗 12:12} 他又说：\BibleQuote{因为我们得救，还是在于望德；所希望的若已看见，就不是望德了；那有人还望那所见的事物呢？但我们若盼望那未看见的，必须坚忍等待。}\BibleRef{罗 8:24} 你不可缺乏这忍耐，反要以无愧的良心，\Quote{等待主；鼓起勇气，主必坚固你的心；我再説，等待主。}\par

2. 当然，如果你在被禁与可敬的\Person{奥勒留}主教共融的情况下来到我们这里，你显然不能获准与我们共融；但我们对待你的态度，将如同我们相信他会采取的态度一样，充满同样的爱德。然而，你不必因此不好意思前来我们这里，因为你应当出于对教会纪律的尊重，感受到顺服此事是你的本分，尤其若你有无愧的良心，这良心为你和天主所知。\Person{奥勒留}若延迟审理你的案件，并非出于对你的厌恶，而是由于其他事务的紧迫；你若像了解自己那样了解他的处境，你对延迟就不会惊讶或忧愁。我恳请你相信我的话，因为你也同样无从知道我的忙碌情况。但还有其他主教比我年长，无论权威更值得尊敬，还是地点更便捷，借助他们的帮助，你可以更容易地解决你负责的教会中悬而未决的事务。然而，我并非没有将你的苦楚以及你信中抱怨的事，告知我可敬的兄弟兼同僚年长的\Person{奥勒留}；我敬重他，配得上他的美德；我特意将你的信件的副本寄给他，让他知晓你所受的指控是无辜的。我收到你的信，是在圣诞节前一天，至多前两天；你在信中告知我，他有意前往\Place{巴德西尔}的教会，你担心民众会因此受扰动并转而反对你。因此，我不敢写信给你的民众；因为我可以回复任何写信给我的人，但我怎能未经请求擅自写信给并非我负责的民众呢？\par

3. 然而，我现在对你、唯一写信给我的人所说的话，可以通过你传达给那些应该听到的人。那么，我首先嘱咐你，不要在公开集会中宣读教会正典未认可的著作，以免给教会带来耻辱；因为异端者，尤其是摩尼教徒（我听说他们中的一些人潜伏在你所在的地区，并非没有受到鼓励），常常利用这些著作颠覆缺乏经验者的心灵。令我惊讶的是，像你这样有智慧的人，竟然劝告我禁止接纳那些从你那里来到我们这里的人进入隐修院，为的是遵守大公会议的某项教规，同时却忘记了同一个大公会议的另一项教规，即宣布哪些是应该向民众宣读的正典圣经。请再次阅读大公会议的文献，并牢记在心：你将在那里发现，你所提及的禁止不加区别地接纳申请进入隐修院的人的教规，并非针对平信徒制定，而仅适用于圣职者。的确，该教规中并未提及隐修院；但总体上规定，任何人不得接纳属于其他教区的圣职者〔除非是以维护教会纪律的方式〕。而且，在最近的一次大公会议上还制定了教规，任何擅离隐修院或被逐出隐修院的人，不得在其他地方被接纳担任圣职或管理隐修院。因此，如果你对\Person{普里瓦提奥}有所担忧，让我告诉你，他尚未被我们接纳进隐修院；我已将他的案件提交给年长的\Person{奥勒留}，并将依照他的决定行事。我觉得奇怪的是，一个人如果仅仅在公开场合读过一次，并且那次读的是非正典的著作，就能被视为读经员。如果因此他被视为教会的读经员，那么他读的著作也必然被视为教会所认可的。但如果该著作未被教会认可为正典，那么，尽管一个人可能已经向会众读过它，他并不能因此成为教会的读经员〔而是仍然如从前一样，是个平信徒〕。然而，对于这位年轻人，我必须遵守我所指派的仲裁者的决定。\par

4. 至于\Place{维盖西莱}的民众，他们在\OriginalTerm{基督的心肠}{bowels of Christ}里为我们也为你们所爱；如果他们拒绝接受一位已被非洲全体公会议罢免的主教，那他们行事明智，既不能被迫屈服，也不应该被迫屈服。凡试图以暴力强迫他们接受那主教的人，必将清楚地显明他的品格如何，也会让人清楚明白他早先让人相信他毫无邪恶时的真实品格是怎样的。因为，任何人若妄用世俗权力或任何其他暴力的手段，企图通过煽动和投诉来恢复他所丧失的教会品级，那么他比任何人都更有效地暴露了自己事业的毫无价值。因为他所渴望的，并不是把基督所要求的服务献给祂，而是夺取基督徒所不承认的权威。弟兄们，要谨慎；魔鬼的诡计极大，但基督是天主的智慧。\par

\SourceRecord{Translated by J.G. Cunningham.；Nicene and Post-Nicene Fathers, First Series；Vol. 1.；Edited by Philip Schaff.；Buffalo, NY: Christian Literature Publishing Co.,；1887.；<http://www.newadvent.org/fathers/1102064.htm>.}

% structure: p0001 source=h1 target=section reason=page-title

\section{第65封信（公元402年）}

\Emph{致可敬的\Person{克桑蒂普斯}，我最可称颂、堪受尊敬的主，我在司铎职务上的父亲与同僚，\Person{奥古斯丁}在主内致候。}\par

1. 向阁下致以与您尊贵相称的敬意，并恳求您在祈祷中纪念我，我愿将一位名叫\Person{阿邦丹提乌斯}的司铎的案件呈请您的睿智裁夺。此人是在我管辖的\Place{斯特拉博尼亚}地区被祝圣为司铎的。他开始因行事不符合天主仆人的正道而招致不好的传闻；我为此感到不安，虽不经查核不信谣言，却由此更密切注意他的行为，并尽力获取（若可能）他受到指控的那些恶行无可辩驳的证据。我首先确定的是，他挪用了一位乡民托他用于宗教目的的钱款，且对他的管家职分不能给出令人满意的交代。接着被证实，且由他自己认罪的事情是：在圣诞节那天，\Place{吉佩}的教会与其他所有教会一样守斋，他在上午约十一时向他在\Place{吉佩}的司铎同僚告别，佯装要回自己的教堂，却未带任何圣职人员作伴，逗留在同一堂区，在一名声名狼藉的妇人家里用了午膳、晚膳，并过了夜。恰巧有一位我们\Place{希波}的圣职人员也在那地方投宿，他曾到过那里；这些事实对这位证人而言是确凿无疑的，因此\Person{阿邦丹提乌斯}无法否认这项指控。至于他所否认的事，我将之留待天主的法庭判断，只就那些他未能隐瞒之事对他宣判。我不敢让他继续管理一个教会，尤其是像他那种处于狂暴咆哮的异端者包围中的教会。当他恳求我写一封信，向\Place{布拉}地区的\Place{阿梅玛}堂区的司铎说明他的案件，他原是从那里来到我们这里的，以免那里对他品行产生过分的猜疑，并使他能在那里，若可能，过一个更规矩的生活，不再承担司铎职务，我出于怜悯便照他所愿的做了。同时，我有特别的责任将这些事实呈报给睿智的您，以免您受到任何欺骗。\par

2. 我于复活主日前一百天，即今年四月七日，对此案作出宣判。我特地将日期告知您，因为这是公会议法令的规定；我也没有向他隐瞒，而是向他解释了教会的律法：若他认为可以通过任何方式推翻我的判决，除非在一年内为此启动程序，否则期限一过，无人会再听取他的申诉。就我而言，我最可称颂的主，堪受一切尊敬的父亲，我向您保证：若不是我认为一个圣职人员这些邪恶行为的实例，尤其伴有坏名声时，理应受到公会议所规定的惩罚，那么我如今将不得不尝试去审查那些无法确知的事，并要么根据可疑的证据定他有罪，要么因证据不足而宣告他无罪。当一位司铎，在众人守斋之日，而该地也如此守斋，他却向他在那地的司铎同僚告别后，未带任何圣职人员作伴，竟敢逗留在一名声名狼藉的妇人家里，并在那里用膳、晚餐和过夜，在我看来，无论别人如何想，他都应被革除职务，因为我不敢再将天主的教会托付给他。假若他可能上诉的教会法官们持有不同意见，鉴于公会议规定，对司铎的案件须由六位主教作出终裁，那么无论谁愿意在他管辖范围内将教会托付给他，就我而言，我承认，我害怕将任何团体托付给像他这样的人，尤其当不能举出一般良好品德的理由来为其过失开脱时；以免一旦他爆发出更具毁灭性的恶行，我不得不痛心地为他的罪行所造成的伤害自责。\par

\SourceRecord{Translated by J.G. Cunningham.；Nicene and Post-Nicene Fathers, First Series；Vol. 1.；Edited by Philip Schaff.；Buffalo, NY: Christian Literature Publishing Co.,；1887.；<http://www.newadvent.org/fathers/1102065.htm>.}

% structure: p0001 source=h1 target=section reason=page-title

\section{第66封信函（公元402年）}

\Emph{无问候语，致\Place{卡拉马}的\OriginalTerm{多纳徒派}{Donatist}主教\Person{克里斯皮努斯}。}\par

1. 你本该因敬畏天主而受影响；然而，既然你在为\OriginalTerm{马帕利安人}{Mappalians}\OriginalTerm{再施洗}{rebaptizare}的过程中，选择利用你作为人所能引发的恐惧，那么让我问你，当省长的命令在一个村庄里已被完全执行时，君主的命令又何以不能在全省执行呢？如果你比较所涉及的人，你不过是一个占据封地的臣属；他是皇帝。如果你比较两者的地位，你是在一份产业上，他是在宝座上；如果你比较两者所持的动机，他的目标是治愈分裂，而你的目标则是撕裂合一。但我们并不让你畏惧人：尽管我们可以采取措施，按照帝国的法令，迫使你因暴行而支付十磅黄金的罚金。或许你无法缴纳对那些为教会成员再施洗者所处的罚金，因为你已花费巨资购买那些你可能逼迫其接受那礼仪的人。但是，正如我所说，我们并不让你畏惧人：倒不如让\Person{基督}使你充满敬畏。我想知道，如果基督这样对你说，你能如何作答：\Quote{克里斯皮努斯，你为收买马帕利安乡民们的恐惧付出了高昂代价；而我的死亡，就是我为换取万民的爱所付出的代价，在你眼中竟如此微不足道吗？你为获得这些农奴以给他们再施洗而从荷包里掏出的金钱，难道比我为救赎万民以给他们施洗而从我肋旁流出的血更为昂贵的祭献吗？}我知道，如果你肯聆听\Person{基督}，就能听到更多这样的呼召，甚至能因你已获得的产业而获得警告，认识到你所讲论反对\Person{基督}的事是何等不虔。因为如果你认为，按照人间法律，你对你用金钱买来的东西的所有权是确定的，那么，按照神圣法律，\Person{基督}对他用自己的宝血所买来的东西的所有权该是何等更加确定呢！那经上论及他所说：\BibleQuote{他要执掌权柄，从这海直到那海，从大河直到地极}的话是真的，他要以不可战胜的大能持守他所买来的一切；但你却断言\Person{基督}已失去整个世界，只留下\Place{非洲}作为他的份，你又怎能指望确实保住你认为在\Place{非洲}用金钱买来的东西呢？\par

2. 但何必多言？若这些\OriginalTerm{马帕利安人}{Mappalians}是出于自由意志归入你的\OriginalTerm{共融}{communio}，就让他们听取你我双方关于使我们分裂的问题的意见——我们各自的话被记录下来，经我们签署证明后译成\OriginalTerm{布匿语}{Punic}；然后，在消除了一切因惧怕其上司而带来的压力之后，让这些农奴选择他们所喜欢的。因为我们所说的话将显明，他们是在强迫下持守错误，还是因同意而持守他们所相信的真理。他们要么明白这些事，要么不明白：如果不明白，你怎敢在他们无知的情况下把他们转入你的\OriginalTerm{共融}{communio}？而如果他们明白，正如我所说的，就让他们听取双方的意见，并为自己自由抉择。如果有任何团体从你那里转到我们这里，而你认为他们是屈服于其上司的压力，那么在他们那里也照此办理；让他们听取双方的意见，并为自己选择。现在，如果你拒绝这个提议，谁还能不相信，你所依靠的不是真理的力量呢？但你应在此世和来世都畏惧天主的义怒。我因\Person{基督}恳求你，回复我所写的。\par

\SourceRecord{Translated by J.G. Cunningham.；Nicene and Post-Nicene Fathers, First Series；Vol. 1.；Edited by Philip Schaff.；Buffalo, NY: Christian Literature Publishing Co.,；1887.；<http://www.newadvent.org/fathers/1102066.htm>.}

% structure: p0001 source=h1 target=section reason=page-title

\section{\Person{奥古斯丁}致\Person{热罗尼莫}书（公元402年）}

\Emph{致我至爱且切望之主、我在基督内可敬的弟兄、\OriginalTerm{同作司铎的}{fellow-presbyter}\Person{热罗尼莫}，\Person{奥古斯丁}在主内向您致意。}\par

% structure: p0003 source=h2 target=subsection reason=relative-heading-rank; base=h2

\subsection{第一章}

1. 我听说我的信已到你手中。我尚未收到回信，但我并不因此怀疑你的情谊；无疑是有事耽搁了。因此，我深知并承认，我应当祈求天主使你能够寄出回信，因为祂早已赐你能力撰写回信，只要你愿意，这原是极容易办到的。\par

% structure: p0005 source=h2 target=subsection reason=relative-heading-rank; base=h2

\subsection{第二章}

2. 我曾犹豫是否相信我所听闻的某件事；但我认为，关于此事，我应当毫不犹豫地给你写上几句。简言之，我听说有几位弟兄曾告诉你的\OriginalTerm{爱德}{Charity}，说我写了一本反对你的书，并将它寄到了\Place{罗马}。请确信这是假的：我呼求天主作证，我并未做过此事。然而，倘若在我的一些著作中，有些地方我持与你不同的见解，我认为你应当知道——或者，即便不能确知，至少应当相信——这些文字的用意并非为了反对你，而仅是阐述我自己的观点。尽管如此，请容我向你保证，我不仅极其乐意以弟兄之情听取你对我的著作中任何令你不快之处所持的异议，而且还恳求你，乃至央求你将这些意见告知我；这样，无论是我纠正错误，还是至少得悉你的善意，我都将欣喜不已。\par

3. 但愿我能有此幸运，得以与你相邻而居，即便不在同一屋檐下，也能时常在主内与你进行愉快而亲密的交谈！然而，既然此事未能获允，我便恳请你费心，使这一条能让我们在主内相聚的途径得以维系；不，更要有所增进和完善。不要拒绝给我回信，哪怕只是偶尔写来。\par

请代我问候我们的圣善弟兄\Person{保利尼亚努斯}，以及所有与你同在、并因你而在主内欢欣的众弟兄。愿你记念我们，使你的所有圣善愿望在主前蒙受应允，我至爱且切望之主、我在基督内可敬的弟兄。\par

\SourceRecord{Translated by J.G. Cunningham.；Nicene and Post-Nicene Fathers, First Series；Vol. 1.；Edited by Philip Schaff.；Buffalo, NY: Christian Literature Publishing Co.,；1887.；<http://www.newadvent.org/fathers/1102067.htm>.}

% structure: p0001 source=h1 target=section reason=page-title

\section{\Person{热罗尼莫}致\Person{奥古斯丁}（主历402年）}

\Emph{致\Person{奥古斯丁}，我主，至圣至福之父}，\Emph{\Person{热罗尼莫}在\Person{基督}内致候。}\par

1. 当我的亲戚、我们的圣子\Person{阿斯忒里乌斯}副执事正要启程时，尊驾的信函到了；在信中，你澄清了关于你曾向罗马寄送一本攻击你卑仆之书的指控。我未曾听闻此指控；但通过我们的兄弟、执事\Person{西辛尼乌斯}，有人寄给我的一封信的抄本传到了这里。在那封信中，我被劝诫撰写\OriginalTerm{翻案诗}{\GreekText{παλινωδία}}，承认在宗徒著作的一段中犯了错误，并效法\Person{斯忒西科罗斯}；他曾在责备和赞美\Person{海伦}之间摇摆，通过赞美她，重获了因出言反对她而丧失的视力。虽然信的文笔和论证方式看似是你的，但我必须坦诚地向尊驾承认，我认为，在仅见到抄本的情况下，不宜未经查证就认定一封信的真实性，以免万一你因我的回复而恼怒，你有理由抱怨说，我理应先确证你是作者，待查明后再回复你。我延迟的另一个原因是虔诚可敬的\Person{保拉}久病不愈。因为，在长时间忙于照料她重病期间，我几乎忘记了你的信，更准确地说，是以你的名义写来的信，此时我想起那句经上的话：\BibleQuote{不合时宜的言论，有如居丧时的音乐；}\BibleRef{德 22:6}因此，若那确是你的信，请坦率地书面告诉我确是如此，或寄给我一份更准确的抄本，以便我们毫无激愤地专心讨论圣经真理；而我可以或是改正自己的错误，或是表明别人无端指责了我。\par

2. 我绝不敢擅自攻击尊驾所写的任何东西。因为对我来说，证明自己的观点已足够，无需反驳他人所持的。然而，以你的智慧，深知每个人都满足于自己的意见，而攻击已成名的前辈以图为自己扬名，那是青年人旧时的惯技，是一种幼稚的自满。我并非如此愚蠢，以至于认为你给出与我不同的解释是对我的侮辱；同样，你也并未因我的观点与你所持的相反而受到伤害。但那正是朋友能真正互益的一种责备：各人看不见自己背上装满过错的包袱，却如\Person{珀尔修斯}所言，注意别人背负的袋子。容我再说，爱那爱你的人，不要因为年轻，就挑战圣经领域中的老兵。我已有过我的时刻，已尽我之力驰骋了全程。理应我该休息，而你接替奔跑，去完成长远的距离；与此同时（请容我说，并无不敬之意），免得似乎只有你能引用诗人的作品，让我提醒你\Person{达勒斯}与\Person{恩特鲁斯}的较量，以及那句谚语：\Quote{疲倦的牛迈步更稳当。}我是怀着忧伤口授这些话的。但愿我能得到你的拥抱，并借着交谈在学问上互相增益！\par

3. \Person{卡尔普尼乌斯}，别号\Emph{拉纳里乌斯}，以他一贯的厚颜无耻，给我寄来了他可憎的作品；我得知他还不遗余力地在非洲散布它们。对此，我先前已简短回复；并已寄给你我的一篇论著的抄本，打算一有机会，当我有暇备稿时，再寄给你一份更详细的论著。在这篇论著中，我注意不伤害任何基督徒的情感，而只是驳斥出于他的无知与疯狂而编造的谎言与幻觉。\par

请记我一念，圣洁而可敬的父亲。看，我多么真诚地爱你，因为我即使被挑衅，也不愿回应，并拒绝相信那封在别人那里我会厉加斥责的信，是出于你手。我们的兄弟\Person{康穆尼斯}向你致以问候。\par

\SourceRecord{Translated by J.G. Cunningham.；Nicene and Post-Nicene Fathers, First Series；Vol. 1.；Edited by Philip Schaff.；Buffalo, NY: Christian Literature Publishing Co.,；1887.；<http://www.newadvent.org/fathers/1102068.htm>.}

% structure: p0001 source=h1 target=section reason=page-title

\section{书信六十九（主后402年）}

\Emph{致我们应受尊爱的主\Person{卡斯托里乌斯}，我们真诚欢迎且当受尊敬的儿子，\Person{亚利比乌斯}和\Person{奥斯定}在主内致候。}\par

基督徒的仇敌试图借着我们至亲至爱的儿子、你的兄弟，在天主教会内煽动一场最危险的丑闻；这慈母教会，当你从一份被剥夺继承权而分离的碎片中逃入基督的产业时，曾欢迎你进入她慈爱的怀抱。那仇敌显然愿望用不合宜的忧郁，遮蔽你皈依的恩赐所带给我们喜乐的宁静之美。但是我们的主天主，慈悲而怜悯，安慰沮丧者，哺育幼儿，眷顾弱者，却容许他在这图谋上略有所成，只为叫我们因灾难得以避免而欢欣，更胜于因危险而忧愁。因为辞去主教尊严的重担以拯救教会脱离危险，远比接受这些责任以参与教会治理更为大度。若有人不能在不危及教会和平的情况下保留主教职位，就不坚持保留，他便证明了，假如和平的利益允许，他原是配得上担任此职的。因此天主乐意藉着你的兄弟，我们可爱的儿子\Person{马克西米安}，向祂教会的仇敌展示，在教会内有那些不求己事，只求耶稣基督之事的人。因为他在放下那天主奥迹管家之职务时，并非在世俗私欲的压力下放弃职责，而是出于对和平的虔诚之爱的激励，以免因授予他的尊位，在基督的肢体间引发一场不体面而危险、甚至可能致命的争论。试想，有什么比为了天主教会的和平而离弃\OriginalTerm{裂教者}{schismatics}，然后又因自己的地位和晋升问题扰乱同一公教和平，更疯狂而完全应受谴责的事呢？另一方面，又有什么比他在脱离了\OriginalTerm{多纳图派}{Donatists}的狂热骄傲之后，在依附基督产业的方式中，于爱的力量下为教会合一而表现出如此显著的谦逊明证，更值得称赞的呢？所以关于他，我们实在欢喜，因为他已被证实如此坚固，以致这诱惑的风暴并未推倒天主真理在他心内所建造的；因此我们愿望并祈求主恩赐，藉着他今后的生活和行事，越来越显明他本会如何善尽那职务的责任，倘若接受那职务是他的本分。愿那应许给教会的永远的和平赏报他，因他看明，凡与教会和平不相容的事，于他都是无益的！\par

至于你，我们亲爱的儿子，我们因你而大感喜乐；既然你并不像你兄弟那样，因任何考虑而不敢接受主教职务，那些考虑引导他拒绝了它，那么合乎你性情的事，就是把你身上因基督的恩赐而有的，都奉献给基督。你的才能、审慎、口才、庄重、自制，以及所有装饰你言行的其它品质，都是天主的恩赐。它们该当用于什么服务，比用于那位赐予它们的主更合适呢？好让祂保守、增进、成全和赏报它们。不要把它们用于服侍此世，免得随此世一同消逝灭亡。我们知道，与你打交道时，无需多加强调，你就很容易反思到：虚荣者的希望、他们无厌的欲望，及生命的无常。所以，丢弃你心灵曾抓住的一切欺骗性世间福乐的期望吧；劳作在天主的葡萄园中，那里的果实是确实的，许多应许已经如此大地应验，以致对剩余的应许失望就真是疯狂之极了。我们以基督的天主性和人性，并以那属天之城——我们在尘世客旅劳苦之后获得永息之所——的和平，恳请你接受\Place{瓦吉纳}教会的主教之位，那是你兄弟所辞去的，不是因不名誉的免职，而是因宽宏大量的出让。叫我们期望因你那天主恩赐所装备和装饰的心智与口才而获最丰沛恩福的子民——我们说，叫那子民通过你认识到，你兄弟的所做所为，并非图谋自己的安逸，而是顾及他们的和平。\par

我们已吩咐，直到那些需要你的人将你实际留下，这信才读给你听。因为我们以灵性之爱的绳索紧系你，因为你对我们也是极需要的同僚。我们未亲自去你那里的原因，你日后必会明白。\par

\SourceRecord{Translated by J.G. Cunningham.；Nicene and Post-Nicene Fathers, First Series；Vol. 1.；Edited by Philip Schaff.；Buffalo, NY: Christian Literature Publishing Co.,；1887.；<http://www.newadvent.org/fathers/1102069.htm>.}

% structure: p0001 source=h1 target=section reason=page-title

\section{\Person{奥斯定}致\Person{热罗尼莫}（公元403年）}

\Emph{致我可敬的主\Person{热罗尼莫}，我尊敬且圣善的弟兄与同司铎，\Person{奥斯定}在主内向您致意。}\par

% structure: p0003 source=h2 target=subsection reason=relative-heading-rank; base=h2

\subsection{第一章}

1. 自从我开始给您写信，并渴望您的回信以来，我从未有过比现在更好的机会来互通音信，因为我的信将由天主最忠信的仆人与事奉者送去，他也是一位我极其亲爱的朋友，即我们的孩子\Person{西彼廉}\OriginalTerm{执事}{deacon}。通过他，我期盼收到您的信，其确定程度正是此类事情所能保证的。因为我所提到的这位孩子，在热切请求或有力劝说以获取您的回信方面，绝不会令人失望；在勤勉保管、迅速传递和忠实交付回信方面，他也不会失职。我只祈求，倘若我配得上如此，愿主赐予您的心灵和我的心愿以助佑和恩惠，好使任何更高的意志都无法阻挠您弟兄般的善意所倾向去行的事。\par

2. 既然我先前给您寄了两封信，都未收到回音，我决定这次将这两封信的抄本一并发给您，因为猜想它们从未送达您手中。如果它们确实送达了，而您的回信——也可能如此——未能送到我这里，那么请您将之前的回信再寄一次，如果您保存了副本的话；如果没有，请再口述重写一遍，供我阅读，不要拒绝为这些先前的信赐予我等候已久的答复。我给您写的第一封信，是在我任\OriginalTerm{司铎}{presbyter}期间准备的，本来要由我们的一位弟兄\Person{普罗福图鲁斯}带给您——他后来成为我在\OriginalTerm{主教职}{episcopate}上的同僚，并已在此后离世；但他当时未能亲自将信带给您，因为就在他打算启程之时，他因被祝圣为主教这一重职而受阻，不久之后他就去世了。我也决定在此时将这封信寄出，好让您知道，我心中向往与您交谈的热望已珍藏了多久，并且我是何等地勉强接受这遥远的分离——它使我的心灵无法透过肉身感官而触及您的内心，我最可爱且尊荣的主内弟兄。\par

% structure: p0006 source=h2 target=subsection reason=relative-heading-rank; base=h2

\subsection{第二章}

3. 在此信中，我还要补充说，我后来听说您已从\OriginalTerm{希伯来文}{Hebrew}原文翻译了\BibleBook{约伯}{传}，尽管在您从\OriginalTerm{希腊文}{Greek}翻译的同一先知书中，我们已经有了该书的译本。在您早先的译本中，您用\OriginalTerm{星号}{asterisk}标出了希伯来文中有而希腊文中无的词，用\OriginalTerm{短剑号}{obelisk}标出了希腊文中有而希伯来文中无的词；这项工作做得如此令人惊异的精确，以至于在某些地方，每个词都分别用一个星号标出，以表示这些词在希伯来文中存在，而在希腊文中没有。然而，在这次的希伯来文新译本中，对词句的忠实却不如先前那般一丝不苟；这使得任何细心的读者都感到困惑，既要揣测先前译本如此慎重地使用星号、甚至标出希腊译本中缺失的、从希伯来文未译出的最小虚词的原因，又要揣测这次的希伯来文新译本为何没有同样细心的确保这些虚词出现在其应有的位置上。我本想在这里摘引一两段以说明此批评；但目前我无法查阅希伯来文译本的抄本。不过，您那敏锐的辨识力早已预见并超越了我所说的，甚至超越了我意欲说的，我想您已经足够明白，仅需解释您所采用之方案的理由，便能释我之惑。\par

4. 就我而言，我更希望您能为我们提供那被称为\OriginalTerm{七十贤士译本}{Seventy translators}的希腊文\OriginalTerm{正典圣经}{canonical Scriptures}的译本。因为，如果您的译本开始在许多教会中更普遍地被诵读，那么，在读经时，拉丁教会与希腊教会之间将势必产生分歧，这实为可悲之事；尤其是，这种歧异在拉丁译本中，很容易因出示希腊原文而受到指摘，因为希腊文是一门广为人知的语言；相反，如果有人因读到他不习惯的、来自希伯来文的译本而感到不安，并声称新译本有误，那么，要找到希伯来文原典以维护那遭受非议的译本，若非不可能，也必极其困难。纵使找到了这些原典，又有谁愿意承认如此众多的拉丁和希腊权威都是错误的呢？除此之外，倘若就希伯来文本的含义咨询犹太人，他们可能会提出与您不同的见解；这样一来，您的亲自出场就显得不可或缺，因为唯有您能反驳他们的观点；而要在您与他们之间找到一位能担任仲裁者的人，简直是个奇迹。\par

% structure: p0009 source=h2 target=subsection reason=relative-heading-rank; base=h2

\subsection{第三章}

5. 我们的一位弟兄，某主教，在他所主持的教会中推行诵读你的译本，当读到\BibleBook{约纳}{先知书}时，遇到一个词，你译得与素来所有崇拜者感官和记忆所熟悉的、在教会中传唱了许多代的译文大不相同\BibleRef{纳 4:6}。于是会众中哗然骚动，特别是希腊人，他们纠正所读的，并指责该译文为谬误，以致主教被迫去征求当地犹太居民的意见（这事发生在\Place{俄亚}城）。这些人，无论出于无知还是恶意，回答说那些希伯来文手稿中的词语在希腊文译本及由此而来的拉丁译本中，都被正确地译出了。我还需要说什么呢？这人被迫在你那段译文上做出修正，仿佛它是错译一般，因为他不想失去会众——他险些遭遇的一桩祸事。从这个事例，我们也认为你偶尔可能有误。你也会看到，若此事发生在那些无法通过比较现今使用语言的证据来阐释的经卷中，困难该有多大。\par

% structure: p0011 source=h2 target=subsection reason=relative-heading-rank; base=h2

\subsection{第4章}

6. 同时，我们对天主满怀感激，因你从希腊原文翻译福音的工作，因为当我们与希腊文圣经对照时，几乎在所有经节中都未发现可指摘之处。借着这项工作，任何支持旧有错误译本的好辩者，只要出示和校对手稿，就可轻易地被说服或驳倒。而且，如果——确实极少发生——发现一些可指摘之处，谁会如此不通情理，不肯轻易原谅这么一部如此有用、再怎么高度赞扬也不为过的工作呢？希望你能惠然向我阐明，你认为希伯来抄本所支持的文本与\OriginalTerm{希腊文七十贤士译本}{Greek Septuagint version}之间频繁出现差异的原因是什么。因为后者具有不小的权威，不仅因为它流传甚广，而且是\OriginalTerm{宗徒们}{apostles}所用的，这一点既可通过审视文本本身得到证实，而且，我记得，你本人也予以肯定。因此，你若能将希腊文七十贤士译本精确地译成拉丁文，就是给我们莫大的恩惠：因为拉丁文不同抄本中发现的差异多得令人难以容忍；并且它极有理由被怀疑可能与希腊文中所见者不同，以至于人们在引用它或借助它证明任何事时，都毫无信心。\par

我原以为这封信会很短，但不知怎的，继续写下去竟如同与你交谈一般令人愉快。最后，我恳切地因着主祈求你，请善意地给我一个全面的答复，从而只要在你权能范围内，让我享受你亲临般的愉悦。\par

\SourceRecord{Translated by J.G. Cunningham.；Nicene and Post-Nicene Fathers, First Series；Vol. 1.；Edited by Philip Schaff.；Buffalo, NY: Christian Literature Publishing Co.,；1887.；<http://www.newadvent.org/fathers/1102071.htm>.}

% structure: p0001 source=h1 target=section reason=page-title

\section{\Person{热罗尼莫}致\Person{奥斯定}书（主历404年）}

\Emph{致\Person{奥斯定}，我真正神圣的主，最可称颂的父，\Person{热罗尼莫}在主内致意。}\par

% structure: p0003 source=h2 target=subsection reason=relative-heading-rank; base=h2

\subsection{第一章}

你一封接一封地寄信给我，频频催促我答复你某一封信，该信的抄本，如我已写过的，是通过我们的执事兄弟\Person{西斯尼乌斯}转交我的，上面并未署名；你告诉我，你最初将这封信托付给我们的兄弟\Person{普洛福图鲁斯}，后来又交给别人；但\Person{普洛福图鲁斯}未能完成他已首途的旅行，并被祝圣为主教后，猝然辞世；而第二位送信人，你未提其名，因害怕大海的危险，放弃了预期的航程。事既如此，我不知如何表达我的惊愕，因为同一封信据说在\Place{罗马}以及整个\Place{意大利}的大多数基督徒中流传，除了我之外每个人都收到了，而它表面上是唯独寄给我的。这令我格外惊讶，因为前述兄弟\Person{西斯尼乌斯}告诉我，他是在约五年前于\Place{亚得里亚海}的一个岛上，在你的其他已发表著作中发现了这封信，既不在\Place{非洲}，也不在你手中。\par

2. 真正的友谊不能存有猜疑；朋友必须像对自己的第二个自我一样，坦率地对朋友说话。我的一些相识，乃基督的器皿，在\Place{耶路撒冷}和圣地人数众多，向我暗示说，这事并非你以真诚之心所为，而是出于追求赞颂与名声，以及在众人眼前\Emph{显耀}，意图藉我身败名裂而扬名；使众人知道你挑战我，我却不敢应战；你以学者身份写信，而我则以沉默承认无知，并终于找到能止住我饶舌的人。然而，容我坦率地说，我起初拒绝答复阁下，因为我不相信这封信，或如我可称之（借用一个谚语说）的蜜糖之剑，是由你发出的。此外，我谨慎以免显得无礼回复一位与我共融的主教，并挑剔一位挑剔我者的信中的任何内容，尤其因为我断定其中某些陈述沾染了异端的色彩。最后，我害怕你可能有理由责备我说，\Quote{什么！你若已看到那信是我的——你若在附于其上的签名中认出了你所熟悉的笔迹的亲笔签名，却如此粗心地伤害了朋友的感情，并因他人恶意构陷之事责备我？}\par

% structure: p0006 source=h2 target=subsection reason=relative-heading-rank; base=h2

\subsection{第二章}

3. 因此，如我早已写过的，要么将有问题的亲笔签名的原信寄给我，要么别再骚扰一个在隐修小室中寻求退隐的老人。你若要施展或炫耀学识，请选择年轻、雄辩、杰出的人作对手，据说在\Place{罗马}有很多这样的才俊，他们既不会不能、也不敢与你交锋，在关于圣经的辩论中与一位主教比武。至于我，曾是士兵，今已退役的老兵，对我而言，更适合为你们和其他人赢得的胜利喝彩，而不是以我衰退的身体参与争战；要小心，若你坚持索要答复，我恐怕会想起\Person{昆图斯·马克西穆斯}以忍耐击败骄傲自大、满有胜算的年轻\Person{汉尼拔}的那段历史。\par

岁月带走一切，也带走心神。我常记得漫长的\newline{}歌唱着少年我送走太阳；\newline{}如今我已忘却那许多的诗歌：\Person{莫伊里斯}的歌喉也\newline{}已弃我而逝。\par

或者更确切地说，引用圣经中的例子：\Place{基肋阿得}的\Person{巴尔齐来}，当他为了年轻的儿子而谢绝了\Person{达味}王的恩待以及他宫廷的一切魅力时，教导我们老年既不该贪图这些，也不该接受给予。\par

4. 至于你呼天主作证，说你没有写过攻击我的书，当然也没有将你从未写过的东西送到\Place{罗马}，并且补充说，即便在你的著作中发现一些与我的观点不同的地方，也并未因此对我构成任何冒犯，因为你并非有意得罪我，而只是写下了你确信为正确的见解；我恳请耐心倾听。你从未写过攻击我的书：那么为什么有人带给我一份由他人抄写、包含着你对我的责备的论著？你未曾写过的书，\Place{意大利}又如何得到它？而且，你如何合理地要求我答复你郑重保证非你所写的信件？我也并非愚蠢到认为，如果有什么地方你的意见与我不合，我就是受你侮辱了。但是，如果你要求与我对质一般，挑剔我的见解，要求我说明我的写作理由，坚持要我改正你判为错误的地方，并要我卑微地写一篇\OriginalTerm{翻案诗}{\GreekText{παλινῳδία}}以撤回前言……\par

\SourceRecord{Translated by J.G. Cunningham.；Nicene and Post-Nicene Fathers, First Series；Vol. 1.；Edited by Philip Schaff.；Buffalo, NY: Christian Literature Publishing Co.,；1887.；<http://www.newadvent.org/fathers/1102072.htm>.}

% structure: p0001 source=h1 target=section reason=page-title

\section{\Person{奥古斯丁}致\Person{热罗尼莫}（主后404年）}

\Emph{致\Person{热罗尼莫}，我可敬且最尊敬的弟兄及同工司铎：\Person{奥古斯丁}在主内致意。}\par

% structure: p0003 source=h2 target=subsection reason=relative-heading-rank; base=h2

\subsection{第一章}

1. 我虽然以为，在这封信到达你之前，你已经透过我们的儿子、天主的仆人、执事\Person{西彼廉}收到了我托他带去的那封信；从该信你定能确知，我写了你提到有人带给你抄本的那封信；因此我设想，我已开始像鲁莽的\Person{达列斯}一样，在你所写的回复中，被第二位\Person{恩特鲁斯}的投掷物与无情重拳连连击中，饱受痛打；然而我同时也要回复你透过我们圣善的儿子\Person{阿斯特里乌斯}惠然寄给我的信，我在其中发现许多你对我极为友善的证据，同时也有一些迹象表明你在一定程度上因我而感到受了冒犯。因此，在阅读此信时，我时而因一句话而得到安慰，时而又被另一句话打击；尤其令我惊异的是：你声称之所以不愿轻率接受你手头那封声称出自我手的信为真迹，乃是因为我若受你的回复冒犯，便可能会理直气壮地向你抗议，因为你本应先确定那信是我的，然后再回复；但你接着却吩咐我，如果那信真是我的，就坦率承认，或者寄出一份更可靠的抄本，以便我们能不带任何苦毒情绪，专心讨论圣经教义。因为，如果你已决意要冒犯我，我们又怎能不带苦毒情绪地进行这样的讨论呢？或者，如果你并无此意，那么，在你并未冒犯我的情况下，我又有何理由，因受你冒犯而理直气壮地抱怨你应先确定信是我的，然后再回复——也就是说，然后再冒犯我呢？因为如果你的回复中并无冒犯我的内容，我就没有正当的抱怨理由。因此，当你对那样一封信写出必然会冒犯我的回复时，我们还有何希望能在不带任何苦毒的情况下讨论圣经教义呢？我绝不敢因你愿意且能够以无可辩驳的论证证明你比我更正确地领会了\Person{保禄}宗徒\BibleBook{迦拉达}{书}中那段经文，或\WorkTitle{圣经}中任何其他经文的意义，而感到冒犯；不但如此，我绝不敢将以任何方式从你的教导中获得启发，或经你纠正而改正，视为我自身的损失，而非向你感恩的理由，反而要视之为益。\par

2. 但是，我最亲爱的弟兄，若非你觉得自己因我所写的内容而受了冒犯，你便不可能认为我会因你的回复而受冒犯。因为我若从未如此看待你，设想你并未觉得自己因我而受冒犯，而你竟如此措辞回复，以致反过来冒犯我，我是不可能如此设想你的。另一方面，如果我被你视为那种荒谬愚蠢之人，竟会在你并未措辞冒犯我的情况下仍感到受冒犯，那么你在此事上已经错待了我，因为你竟对我抱持这种看法。然而你如此谨慎，虽然你在那封你持有抄本的信中认出了我的文风，却仍不肯相信其真实性，你必定不会不加思索地认为，我竟与你经验中所认识的我判若两人。因为，如果你有充分理由看出，我可能会理直气壮地抱怨，倘若你轻率地将一封并非我所写的信归为我所写，那么，倘若你不加思索地对我形成一种与你自身经验相矛盾的评价，我岂非更有理由抱怨吗？因此，你在判断上不至于如此偏差，竟相信当你并未写出任何可能冒犯我的内容时，我竟会如此愚拙，竟能被那样一封回复所冒犯。\par

% structure: p0006 source=h2 target=subsection reason=relative-heading-rank; base=h2

\subsection{第二章}

3. 因此毫无疑问，只要你拥有那封信确为我所写的确凿证据，你便打算以冒犯我的方式来回复。因此，既然我不相信你在没有正当理由的情况下会认为冒犯我是对的，我便只能承认我的过错，正如我现在所做的，承认我因写了那封我无法否认出自我手的信，而首先冒犯了你。我何必逆流挣扎，而不请求宽恕呢？因此，我靠着\Person{基督}的慈爱恳求你，宽恕我让你受伤之处，不要以伤害我来以恶报恶。因为你若对我言行中的任何错误保持沉默，那将是对我的伤害。其实，如果你指责我本不该指责之处，那是伤害你自己，而非伤害我；因为，以你的生活和神圣誓愿，断乎不会出于冒犯之心而指责，用恶毒的舌头来谴责你借由真理启示之理性之光明知不配谴责的事。所以，要么你以善意责备那位你虽认为他应受责备、实则并无过错的人；要么以慈父般的温情抚慰那位你无法使他同意你观点的人。因为有可能你的看法与真理不符，而你的行为却与基督徒的爱德和谐一致：因为我也将满怀感激地接受你的责备，视之为最友善的行动，即使你所谴责的事本身可以辩护，因而本不应受谴责；否则我便会承认你的善意和我的过失，并要在主所许可的范围内，为前者感恩，为后者改正。\par

4. 那么，我为何要惧怕你的话语呢？它们或许严厉，就像\Person{恩泰路斯}的拳击手套，但无疑是为使我获益。\Person{恩泰路斯}的击打本意不是医治，而是伤害，因此他的对手是被征服，而非被治愈。但我，如果我平静地接受你的指正，如同接受必需的良药，就不会因它感到痛苦。不过，倘若由于人性共有的软弱，或我独有的弱点，我在受到谴责时，即便那是公正的责备，也难免感到一些痛苦，那么，宁愿让头上的肿瘤被治愈，即使长矛刺入会带来疼痛，也胜过为了逃避痛苦而任由疾病蔓延。这一点，那位说过下面这番话的人看得分明：通常，揭示我们过错的敌人，比害怕责备我们的朋友更为有益。因为前者在愤怒的指责中，有时确实指控了我们应当纠正之处；而后者，为了怕破坏友谊的温馨，在为正义辩护时，往往缺少应有的勇气。因此，既然用你自己的比喻来说，你是一头牛，或许因年岁体力已衰，但心智活力不减，仍在主的打谷场上勤勉而有益地劳作；那么，我在此，无论我说错了什么，就请重重地踩踏我：你那可敬年岁的分量不应令我难受，只要我过错的糠秕被如此踩碎，从而与我分离。\par

5. 再者，我怀着最深切的渴望之情，诵读并忆起你信末的这些话：\Quote{巴不得我能接受你的拥抱，并借着谈话互相帮助学习。} 就我而言，我说——巴不得我们哪怕居住在地球上相隔较近的地方；这样，即使我们不能见面交谈，至少能更频繁地互通书信。因为实际上，我们相隔如此遥远，以致无论通过眼目还是耳朵都无从接触，我记得自己年轻时曾就\BibleBook{迦拉达}{书}中的这些话致函你的圣德；看哪，如今我已年迈，却仍未收到回复，而我那封信的抄本竟因某种奇特的意外，比原信更早到达你手中，我为传递那封信可是费了不少心思。因为受托之人既未将信交给你，也未退还给我。在我们所能获得的你的著作中，我极其珍视它们的价值，以至于若有可能，我宁愿放弃其他一切研究，也要获得坐在你身旁向你学习的特权。既然我无法亲自做到，我打算派我在主内的一个儿子到你那里去，使他可以为我而受教于你，倘若我能就此事收到你善意的答复。因为我如今没有，也永远不敢指望拥有，像我所见你那样对圣经的知识。我无论有何等研究能力，都已全然用于教导天主托付给我的子民；教会的职务完全占据了我，使我除了为尽职公开教导所需之外，别无余暇进一步钻研。\par

% structure: p0010 source=h2 target=subsection reason=relative-heading-rank; base=h2

\subsection{第三章}

6. 我不曾读过那些攻击你的作品，就是你告诉我已传入非洲的那些。不过，我已收到你惠然寄来的答辩。坦率地说，读过之后，我极其痛心，因为在曾经如此相亲相善、昔日因几乎为所有教会所共知的友谊联系在一起的两人之间，竟发生了这种令人痛苦的纷争，造成了如此的不幸。在你的那篇论著里，任何人都能看到，你是如何自我节制，如何强压胸中刺骨的义怒，以免以辱骂还辱骂。然而，即便在读你这篇答辩时，我已因悲痛而昏厥，因恐惧而战栗，那么，如果他写来攻击你的那些东西落到我手中，又将在我的内心产生何等的影响呢！\BibleQuote{世界因了恶表是有祸的！}\BibleRef{玛 18:7} 看哪，那身为真理者所预言的话完全应验了：\BibleQuote{由于罪恶增多，许多人的爱德必要冷淡。}\BibleRef{玛 24:12} 因为如今，有谁信赖的心还能坦然倾注，确知自己的信任会得到回报？如今，信赖的爱还能毫无保留地投向谁的怀抱？简言之，如果在\Person{热罗尼莫}与\Person{鲁菲诺}之间都可能发生我们所痛惜的破裂，那么，哪里还有朋友可以不怕他日后成为敌人呢？哦，我们的处境何等悲惨可怜！当一个人对朋友现在的为人了然于心，却无法预知他们将来会变成什么样子时，谁能倚靠朋友的友爱呢？然而，既然一个人就连自己将来会变成什么都不知道，我又怎能因某人不了解他人将来的情形而感到悲哀呢？他最多只能知道自己目前的状态，而且这种了解也是不完整的；至于他日后会成为怎样的人，他是一无所知的。\par

7. 神圣而有福的天使是否不仅拥有对他们实际品性的这种认识，还拥有对他们将来会成为什么状态的预知？如果他们拥有，我就无法看出撒殚怎么可能曾是幸福的，甚至在他还是善良天使时，因为在这种情况下他必定已经知道他将来的悖逆和永罚。关于这个问题，如果它确实是一个我们能有益地作出回答的问题，我希望听听你的看法。但请注意我因使我们身体彼此远离的陆地和海洋所遭受的痛苦。如果我自己就是你现在正在读的那封信，你或许已经告诉了我刚才所问的；但现在，你何时才会写信回复我？你何时才会把信发出？它何时才会到达这里？我何时才能收到它？然而，但愿我至少能确知它终会到达，虽然在这期间我必须鼓起我能掌控的一切忍耐，来忍受这不受欢迎却无法避免的迟延！因此，我回到你信中那些最令人愉悦的话，它们充满了你圣善的渴望，我也同样把它们变成我自己的渴望：\Quote{但愿我能得到你的拥抱，借由交谈在学问上互相帮助，}——如果确实有某种我能向你传授教导的意义的话。\par

8. 当这些话，现在不仅是你的，也是我的，让我欢喜并振作，当我因我们彼此相见的渴望存在却未得满足这件事而大得安慰时，没过多久，我便被最尖锐的悲伤之箭刺透，当我想到\Person{鲁菲努斯}和你——天主曾以最丰盛的程度并长久地赐予你们我们两人所渴望的东西，使你们能在最亲密可爱的友谊中一同以圣经的甜蜜为盛宴——并且想到你们之间竟生出了如此剧烈的苦毒，迫使我们不得不问，同样的灾祸何时、何地以及在何人身上不会令人合理地恐惧？这灾祸临到你们的时候，正是你们卸下世俗重担，追随主，一同生活在我们的主曾用脚行走过的那片土地上时，那时祂说：\BibleQuote{我把平安留给你们，我将我的平安赐给你们；}\BibleRef{若 14:27} 况且你们是成熟的人，你们的生活都专务于学习天主的圣言。的确，\BibleQuote{人生在世是一场考验。} 如果我能在何处见到你们两人在一起——唉，我无法指望做到——我的激动、悲伤和恐惧是那么强烈，以致我想我要俯伏在你们脚前，在那里哭泣直到不能再哭，然后以爱的全部热诚，首先为你们各自的缘故，再为你们彼此的缘求你们，同时也为那些人的缘故，尤其是软弱的，\BibleQuote{基督为他们死了}\BibleRef{格前 8:11} 的人，他们的救恩正处在危险中，因为他们看着你们，你们在时间的舞台上占据了如此显著的位置；恳求你们不要写下并散播这些针对彼此的刻薄言语，这些言语如果在你们现在不和而将来又和好时，你们都无法销毁，那时你们也不敢去读，以免再次点燃纷争。\par

9. 但我对你的爱德说，当我读到你的信中一些表明你对我不满的迹象时，没有什么比你和\Person{鲁菲努斯}之间的疏远更令我恐惧战栗。我指的与其说是你关于\Person{恩特鲁斯}和疲乏的牛的言论——在我看来，你在那里使用了温和的戏谑而非愤怒的威胁——不如说是我已经谈及的字句，或许我所讲得超出了适当，但没有超出我的恐惧所迫使我讲的：就是那句话，\Quote{免得你或许因被触怒而有理由责备我。} 若我们有可能，在没有任何不和的苦毒下，审视和讨论那能滋养我们心灵的东西，我恳求你，让我们致力于此。但如果你我中的任何一个在指出对方著作中他认为需要纠正之处时，不可能不受到猜忌并被视作伤害友谊，那么，为了我们灵性的生命和健康，就让我们放弃这样的讨论吧。让我们满足于那能使人自高自大的\OriginalTerm{知识}{knowledge}上较少的程度，只要我们能借此保持那能建树的\OriginalTerm{爱德}{charity}不受伤害。\BibleRef{格前 8:1} 我觉得我远未达到经上所写的那种完美：\BibleQuote{谁若在言语上不犯过失，他便是个完人；}\BibleRef{雅 3:2} 但靠着天主的仁慈，我确实相信自己能够请求你宽恕我在什么事上得罪了你：而你应将这事向我指明，好使通过听从你，你得以赢取你的兄弟。\BibleRef{玛 18:18} 你也不要因为我们在地表上的距离使我们无法面对面相见，就以此为理由让我留在错误之中。至于我们所探讨的问题，如果我在某件事上，得知、相信或认为已把握了真理，而你的意见与我不同，我会尽一切可能，按主所赐的能力，在不伤害你的情况下维护我的观点。至于我可能冒犯你的任何事，一旦我觉察到你的不悦，我定会毫不保留地恳求你的宽恕。\par

10. 此外，我想，你对我感到不悦的唯一原因，只能是我要么说了不该说的话，要么未能以应有的方式表达：因为我不奇怪我们彼此的了解，不如我们对自己最亲密的朋友那样透彻。对于这样的朋友的爱，我毫不犹豫地全心信赖，尤其当我被这世界的丑闻所激怒和厌倦时；我在他们的爱中安歇，没有任何纷扰的忧虑：因为我感觉到天主就在那里，我满怀信赖地投向祂，并在祂内安然憩息。在这信赖中，我也不被任何对明日的不确定的恐惧所困扰，这种恐惧当我们依靠人的软弱时必然存在，也是我在前段中所哀叹的。因为当我察觉到一个人燃烧着基督徒的爱德，并感到他由此成为我的忠实朋友时，我将自己的一切计划或思想托付给他，并不认为托付给了人，而是托付给了那位因他的品行而明显居住在他内的那位：因为\BibleQuote{天主是爱，那住在爱内的，就住在天主内，天主也住在他内；}\BibleRef{若一 4:16}如果他不再住在爱内，他那离弃爱所带来的痛苦，必然不亚于他住在爱内所带来的喜乐。然而，在这种情况下，当一位密友成为敌人时，他费尽心机编造对我们不利的谎言，胜过让愤怒激使他揭露真相。每个人都最容易通过以下方式确保这一点：不是隐藏自己所行之事，而是不做任何他想要隐藏的事。天主的仁慈赐予善良虔敬的人：他们无论朋友日后变成怎样，都在朋友中间自由出入，无恐无惧；他们不揭露自己所知的别人所犯的罪，自己也避免去做那些害怕被揭露的事。因为当诽谤者捏造虚假指控时，要么它不被相信，要么，即使被相信，也只是我们的声誉受损，我们的灵性福祉不受影响。但是，当犯下任何罪过时，那行动本身就成为一个隐秘的敌人，即使它未被那些知道我们秘密的人轻率或恶意的言谈揭露。因此，任何有洞察力的人都可以在你自身的例子中看到，你如何凭着良心的安慰，忍受了那原本难以承受的事——就是你从前最亲密、最心爱朋友的难以置信的敌意；以及你如何甚至将他所说反对你的话，甚至被一些人信以为真而对你不利的话，都转化为有益之用，如同左手持有正义的盔甲，这盔甲在与魔鬼的战争中，并不逊于右手的盔甲\BibleRef{格后 6:7}。但说真的，我宁愿看到他的指责不那么刻薄，也不愿看到你因此更全副武装。从如此的友谊转为如此的敌意，这是一件可悲可叹的奇事：若你能从这样的敌意回归昔日的友情，那将是一件令人欣喜、更为伟大的盛事。\par

\SourceRecord{Translated by J.G. Cunningham.；Nicene and Post-Nicene Fathers, First Series；Vol. 1.；Edited by Philip Schaff.；Buffalo, NY: Christian Literature Publishing Co.,；1887.；<http://www.newadvent.org/fathers/1102073.htm>.}

% structure: p0001 source=h1 target=section reason=page-title

\section{书信74（公元404年）}

\Emph{致我的主\Person{普雷西迪乌斯}，最可尊敬的、我的兄弟与\OriginalTerm{司铎职}{Priestly Office}的同僚，由衷敬爱者，\Person{奥斯定}在主内向您致意。}\par

1. 我写这封信是为提醒你，当你在这里时，我作为真诚的朋友向你提出的请求，即请你不要拒绝将我的一封信转交给我们的圣善兄弟与司铎同僚\Person{热罗尼莫}；此外，也让你的\OriginalTerm{爱德}{Charity}知晓，你应以何种措辞代表我写信给他。我已寄出我给他的信以及他给我的信的一份抄本，通过阅读它，你虔敬的智慧不难看出我谨慎保持的节制语气，以及他使我相当恐惧的那种激烈态度，这并非毫无理由。然而，如果我写了任何本不该写的东西，或以不得体的方式表达了自己，那么请你不要以兄弟之爱向他，而是向我发送你对我所为的看法，以便若经你的责备使我确信己过，我可请求他的宽恕。\par

\SourceRecord{Translated by J.G. Cunningham.；Nicene and Post-Nicene Fathers, First Series；Vol. 1.；Edited by Philip Schaff.；Buffalo, NY: Christian Literature Publishing Co.,；1887.；<http://www.newadvent.org/fathers/1102074.htm>.}

% structure: p0001 source=h1 target=section reason=page-title

\section{\Person{热罗尼莫}致\Person{奥斯定}（公元404年）}

\EditorialNote{热罗尼莫对第28、40和71封信的答复。}\par

\Emph{致\Person{奥斯定}，我真正神圣的主，并最蒙福的父亲，\Person{热罗尼莫}在基督内致以问候。}\par

% structure: p0004 source=h2 target=subsection reason=relative-heading-rank; base=h2

\subsection{第一章}

1. 我由执事\Person{西彼廉}收到您阁下的三封信，或不如说是三本小书，同时，包括您所说的各种问题，但我觉得是对我已发表的意见的批评，若要回答，如果我打算做，将需要一本相当大的书。不过我将尝试在不超出适当长度信件的限度内回复，并且不耽搁我们那位急于离开的弟兄，他在预定行程日期的前三天才恳切要求带一封信，因此我不得不几乎未经预谋地倾泻出这些句子，以散漫的流露回答您，不是在从容构思的闲暇中准备，而是在即席口授的匆忙中，这通常产生更多是偶然的产物而非教导的父亲的话语；正如突然袭击使最勇敢的士兵也陷入混乱，他们被迫在披挂铠甲之前就逃跑。\par

2. 但是我们的铠甲是基督；这就是宗徒保禄在致厄弗所人的书信中所吩咐的，他说：\BibleQuote{你们应拿起天主的全副武装，好在邪恶的日子能够抵抗}；又说：\BibleQuote{所以要站稳，用真理作腰带，穿上正义的胸甲，以和平的福音作走路的鞋穿上；此外，再拿起信德的盾牌，用以熄灭恶者的一切火箭；戴上救恩的头盔，并拿起圣神的利剑，即是天主的话。} \BibleRef{弗 6:13} 达味王当年就是佩戴这些武器出战的；他从溪流中取了五块光滑圆石，证明即使在世界各种涡流之中，他的情感既无粗糙也无玷污；在路边喝了溪水，因此心神振奋，他用那傲慢敌人的剑作为最合适的死亡工具，砍下了哥肋雅的头，击打那亵渎的狂妄者的额头，伤在与他僭越司祭职务而受癞病打击的乌齐雅相同的位置；\BibleRef{编下 16:19} 同样这地方，闪耀着令圣人在主内喜乐的荣光，说：\BibleQuote{上主，祢的慈颜已印在我们身上。} 所以我们也当说： \BibleQuote{天主，我的心已踏实，我的心已踏实：我要歌唱赞美；我的荣耀，醒起；琴瑟啊，醒起；我自己要及早醒起；} 好使我们内应验那句话： \BibleQuote{你大张开口，我要使它充满；} 以及： \BibleQuote{上主必以大力赐给传报者的话。} 我深信，你的祈祷与我的祈祷一样，就是我们的争辩中，胜利归于真理。因为你寻求基督的光荣，而非你自己的：若你获胜，我若发现自己的错误，我也同样获得胜利。相反，若我占上风，获益的是你；因为子女不该为父母积蓄，而是父母为子女积蓄。\BibleRef{格后 12:14} 我们又在\WorkTitle{编年纪}中读到，以色列子民怀揣和平之心前去作战， \BibleRef{编上 12:17} 即便在刀剑流血、尸横遍野中，他们不为自己寻求胜利，而是为和平。因此，若合基督之意，让我答复你写来的一切，并尝试在一篇短论中解答你众多的问题。我不提你礼貌问候中的委婉措辞；也不谈你试图减弱批评的恭维话：让我直接进入辩论的问题。\par

% structure: p0007 source=h2 target=subsection reason=relative-heading-rank; base=h2

\subsection{第二章}

3. 你说你从某位弟兄那里收到我的一本书，其中列有希腊和拉丁教会作家的名录，但没有书名；当你问那位弟兄（我引用你的原话）为何题页无题词，或者说这书以何名而为人所知时，他回答说它被称为\Quote{\OriginalTerm{墓志铭}{Epitaphium}}，\Emph{即} \Quote{讣闻}：对此你发挥推理，评论说如果读者在其中找到的是已故作者的生平和著作，那么\WorkTitle{墓志铭}这个名字本来很合适，但因为书中提到许多在写作时仍在世、且如今依然健在的作者的作品，你感到奇怪我为何给这本书取了如此不相称的名。我认为，凭你的常识，你应当不难从书的内容发现它的书名，而无须他人帮助。因为你读过希腊和拉丁的杰出人物传记，你知道这类作品不叫\WorkTitle{墓志铭}，而简单地称为\Quote{\WorkTitle{杰出人物}}，\Emph{例如} \Quote{\WorkTitle{著名将领}}，或\Quote{哲学家、演说家、史家、诗人}，\Emph{等等}，视情况而定。所谓\WorkTitle{墓志铭}是为死者写的著作；正如我记得很久以前在令人怀念的司铎\Person{内波提亚努斯}去世后写过的那样。因此，你所谈的那本书应名为\Quote{\WorkTitle{论杰出人物}}，或更精确地\Quote{\WorkTitle{论教会作家}}，尽管据说由于许多人不具备修正书名之资格，它已被称作\Quote{\WorkTitle{论作家}}。\par

% structure: p0009 source=h2 target=subsection reason=relative-heading-rank; base=h2

\subsection{第三章}

4. 第二，你问道，我在\BibleBook{迦拉达}{书}注释中为何说保禄不可能因自己曾做过的事斥责伯多禄，\BibleRef{迦 2:14} 也不可能指责别人虚伪，而他自己也承认犯过这种虚伪；你确认那宗徒的斥责并非虔敬的权宜之计，而是真实的；你还说我不可教授虚假，而圣经中的一切都应按字面意义接受。\par

对此我首先要回答说，您的智慧本应提醒您记起我为注释所写的简短序言，在那里论到我自己说：\Quote{那么怎么样呢？我岂是这样愚昧狂妄，竟敢答应那他自己也不能完成的事呢？绝无其事；然而，依我看来，我反倒更加保留和犹豫，因为感到自己力量不足，而在这一事上依从了\OriginalTerm{\Person{奥利金}}{Origen}的注释。}因为那位杰出的人士就\BibleBook{迦拉达}{书}撰写了五卷书，并在他的\Emph{\OriginalTerm{\WorkTitle{杂集}}{Stromata}}第十卷中，以一篇简短的论文阐述了他对这封书信的解释。他还撰写了数篇专题论文和断片，即便它们单独存在，也足够了。我且略过我所尊敬的老师\OriginalTerm{\Person{迪迪莫斯}}{Didymus}（他确是眼瞎，但在辨识灵性事物上却目光锐利），以及最近离开了教会的\OriginalTerm{\Place{劳狄刻雅}}{Laodicea}的主教，还有早期的异端者\OriginalTerm{\Person{亚历山大}}{Alexander}，以及\OriginalTerm{\Person{埃梅萨的优西比乌}}{Eusebius of Emesa}和\OriginalTerm{\Person{赫拉克利亚的狄奥多鲁斯}}{Theodorus of Heraclea}，他们也留下了一些关于这个主题的简短探讨。如果从这些作品中哪怕摘录出几段，也能产生一部不可轻视的著作。因此，容我坦率地说，所有这些我都读过了；并将其中包含的许多内容储存在心中，我有时向我的抄写员口述借用自其他作者的内容，有时则是我自己的，却未能清楚地记得哪些方法、哪些措辞或哪些见解各属于何人。我如今仰望主的仁慈，求祂赏赐，不要因我缺乏技巧和经验，而使别人讲得好的东西失传，或在外国读者中得不到它们最初写成时在那语言中所获得的接纳。因此，如果我的解释中有任何地方您认为需要纠正，那么以您的学识，理应首先探究我所写的，是否见于我所引证的希腊作家之中；倘若他们并未提出您所谴责的见解，那么您就可以合宜地谴责我所提供的我自己的观点，尤其是我在序言中已经公开承认，我曾依从\Person{奥利金}的注释，有时口述他人的观点，有时口述我自己的观点，并且我在...末尾写上了\par

您所挑剔的那一章：\Quote{若有人不满意此处给出的解释，即藉此表明\OriginalTerm{\Person{伯多禄}}{Peter}既未犯罪，\OriginalTerm{\Person{保禄}}{Paul}也未僭越地斥责一位比他自己更大的，那么他有义务说明\Person{保禄}如何能一贯地责备别人所做的，正是他自己所做的事。}藉此我已表明，我并没有最终且不可更改地采纳我从这些希腊作者所读到的内容，而是提出我所读到的，让读者自行判断是当拒绝还是接纳。\par

5. 然而，您为了避免照我所请求的去做，便针对所提出的观点设计了一种新的论证；主张那信了基督的外邦人不受礼仪法律的束缚，但犹太的归信者仍在法律之下，而\Person{保禄}作为外邦人的教师，正确地斥责了那些遵守法律的人；而\Person{伯多禄}，他是\OriginalTerm{\Quote{受割损者}}{circumcision}的首领\BibleRef{迦 2:8}，却因命令外邦归信者去做那只有出自犹太的归信者才有义务遵守的事，受到了正义的斥责。如果这是您的意见，或者更确切地说，既然这是您的意见，即所有信主的犹太人都负有遵守全法律的义务，那么您作为举世闻名的主教，理应公布您的教义，并努力说服所有其他主教与您一致。至于我，在我卑微的斗室里，与我的同是罪人的隐修士同伴们一起，我不敢在如此重大的事上妄断教义；我只坦率地承认，我阅读教父们的著作，并遵照普遍的惯例，在我的注释中列出各种解释，以便每个人能从所提供的解释中选取那令他满意的一种。我想，您在阅读中也发现并赞同这种方法，无论是在世俗文学还是\WorkTitle{圣经}方面。\par

6. 此外，至于\Person{奥利金}最先提出、且在他之后所有其他注释者都采纳的这种解释，他们提出来，主要是为了回应\OriginalTerm{\Person{波斐留}}{Porphyry}的亵渎之词；\Person{波斐留}指责\Person{保禄}僭妄，因为他竟敢当面指责\Person{伯多禄}并斥责他，且通过推理证明他做了错事；也就是说，犯了那个他自己——这位责备别人犯过的人——所犯的同样的过失。关于那位长期治理\Place{君士坦丁堡}教会、居有主教尊位的\OriginalTerm{\Person{若望}}{John}，我又该说什么呢？他针对这一段写了一本篇幅很大的书，并遵循了\Person{奥利金}和古代解经家的意见。因此，如果您指责我为错误，请容我，我恳求您，与这些人一同犯错吧；当您看到我在错误中有如此多的同伴时，您至少需要提出一位支持者来为您的真理辩护。关于\BibleBook{迦拉达}{书}中这一段解释，就说到这里。\par

7. 然而，为避免让人觉得我对于您的推理所作的答复，全赖我这边证人的数目，并利用杰出人士的名字作为逃避真理的手段，而不敢与您在论据上交锋，我将简短地从\WorkTitle{圣经}中举出一些例证。\par

\WorkTitle{宗徒大事录}中，有一个声音对\Person{伯多禄}说：\BibleQuote{起来，\Person{伯多禄}，宰了吃！}当时有各种四足兽、爬虫和飞鸟呈现给他；这话证明了没有人本性上是（礼仪上）不洁的，众人都同样蒙召迎接基督的福音。\Person{伯多禄}却回答说：\BibleQuote{主，绝对不可！因为我从来没有吃过什么俗物和不洁之物。}那声音第二次向他说：\BibleQuote{天主所洁净的，你不可称为俗物。}于是他去了\Place{凯撒勒雅}，进了\Person{科尔乃略}的家，\BibleQuote{便开口说：我真正明白了：天主是不看情面的，凡在各民族中，敬畏他而又履行正义的人，都是他所中悦的。}此后，\BibleQuote{圣神降在所有听道的人身上。那些受过割损与\Person{伯多禄}同来的信徒，都惊讶圣神的恩惠也倾注在外邦人身上，因为听见他们说各种语言，并颂扬天主。那时\Person{伯多禄}就发言说：这些人既领受了圣神，和我们一样，谁能阻止他们受水洗呢？遂吩咐人以主的名给他们施洗。}\BibleRef{宗 10:13} \BibleQuote{宗徒和在\Place{犹太}的弟兄们听说外邦人也接受了天主的话。及至\Person{伯多禄}上到\Place{耶路撒冷}，那些受割损的人非难他说：你竟进了未受割损人的家，且同他们吃了饭！}于是他对他们详细解释了自己行事的缘由，最后说：\BibleQuote{那么，既然天主给了他们同样的恩惠，如同我们信主耶稣基督时，给了我们一样；我是谁，岂能拦阻天主呢？他们听了这话，才平静下来，并光荣天主说：原来天主也赐恩给外邦人，叫他们悔改而获得生命。}\BibleRef{宗 11:1}后来又过了很久，\Person{保禄}和\Person{巴尔纳伯}到了\Place{安提约基雅}，\BibleQuote{他们召集了教会，报告天主偕同他们所行的一切，以及他怎样给外邦人打开了信德的门。有几个从\Place{犹太}下来的人教训弟兄们说：你们若不按\Person{梅瑟}的规例受割损，不能得救。当\Person{保禄}和\Person{巴尔纳伯}同他们发生不少争执和辩论时，大家决定派\Person{保禄}和\Person{巴尔纳伯}，及他们中的几个人，上\Place{耶路撒冷}去见宗徒和长老，讨论这问题。他们到了\Place{耶路撒冷}，有几个信教的法利塞党人起来说：必须叫外邦人受割损，又应该命他们遵守\Person{梅瑟}法律。}经过许多辩论后，\Person{伯多禄}起来，以他一贯的爽快态度，\BibleQuote{说：诸位弟兄！你们深知，多时以前，天主就在你们中选拔了我，为叫外邦人借我的口听到福音的道理而信从。洞悉人心的天主，为他们作了证，因为赐给了他们圣神，如同赐给了我们一样；在我们和他们中间没有作任何区别，因他以信德净化了他们的心。既然如此，现今你们为什么试探天主，在门徒的颈项上，放上连你们的祖先和我们自己都不能负荷的轭呢？但是，我们信我们得救是借着主耶稣基督的恩宠，正和他们一样。于是众人都缄默不语；}宗徒\Person{雅各伯}和众长老都赞同了他的意见。\par

8. 这些引证不应使读者厌烦，而应对他和我都有益处，因为它们证明，早在宗徒\Person{保禄}之前，\Person{伯多禄}已经知道福音赐下后，法律就不再生效；不仅如此，\Person{伯多禄}还是颁布命令肯定此事的主要推动者。再者，\Person{伯多禄}的权威如此之大，\Person{保禄}在自己的书信中记载道：\BibleQuote{三年后，我才上\Place{耶路撒冷}去拜见\Person{刻法}，和他同住了十五天。}\BibleRef{迦 1:18}在下文，他又补充说：\BibleQuote{过了十四年，我同\Person{巴尔纳伯}再上\Place{耶路撒冷}去，还带了\Person{弟铎}同去。我是奉了启示上去的，把我在这外邦人中所宣讲的福音，陈说给他们；}以此证明，若没有\Person{伯多禄}及同他在一起的人的认可，他对自己所宣讲的福音就缺乏信心。接下来的话是：\BibleQuote{但只私下向那些有众望的人陈述，免得我从前或是如今都白跑了。}他为何私下而不是公开去？为避免伤害那些从犹太人中归信之人的信德，因为他们以为法律仍在生效，他们应把守法和信靠主救主结合起来。因此，那时\Person{伯多禄}到了\Place{安提约基雅}（虽然\WorkTitle{宗徒大事录}没有提及此事，但我们必须相信\Person{保禄}的话），\Person{保禄}断言，他\BibleQuote{当面对立了他，因为他有可责之处。原来在和某些从\Person{雅各伯}那里来的人未到之前，他惯常同外邦人一起吃饭；可是他们一来到，他因怕那些受割损的人，就退避了，自己躲开。其余的犹太人也都跟他一起装假，以致\Person{巴尔纳伯}也因他们的假善而误入歧途。但我一看出，}他说，\BibleQuote{他们行事不合福音的真理，就当众对\Person{刻法}说：如果你这犹太人，按外邦人的方式，而不按犹太人的方式生活，你怎么竟敢强迫外邦人按犹太人的方式生活呢？}等等。因此，无人能怀疑，宗徒\Person{伯多禄}本人正是他所被指控偏离的那条规则的制定者。而且，他偏离的原因，是出于对犹太人的惧怕。因为圣经说，\BibleQuote{起初他惯常同外邦人一起吃饭，但某些从\Person{雅各伯}那里来的人一到，他因怕那些受割损的人，就退避了，自己躲开。}他惧怕那些他曾受命作宗徒的犹太人，怕他们因外邦人的缘故而偏离对基督的信德；他效法善牧，深怕失去托付给他的羊群。\par

9. 因此，正如我已指出的，\Person{伯多禄}深知梅瑟的法律已被废除，却因恐惧而假装遵守；现在让我们也来看看，那指责别人的\Person{保禄}自己是否也做过同样的事。我们在同一卷书中读到：\BibleQuote{保禄走遍了叙利亚和基里基雅，坚固各教会。然后来到了德尔贝和吕斯特辣。在那里，有一个门徒，名叫弟茂德，是一个信主的犹太妇人的儿子，父亲却是希腊人，在吕斯特辣和依科尼雍的弟兄们都称赞他。保禄愿意带他一同去；但为了那些地方的犹太人，就给他行了割损礼，因为他们都知道他父亲是希腊人。}啊，有福的宗徒保禄！你曾因\Person{伯多禄}在那些从\Person{雅各伯}那里来的犹太人面前，因恐惧而退避外邦人，责备他装假；为什么你自己不顾自己的教训，竟被迫给一个外邦人的儿子——不，他自己就是外邦人（因为他不是犹太人，未曾受割损）——行割损礼呢？你会回答说：\Quote{为了那些地方的犹太人。}那么，你若原谅自己给这位出自外邦人的门徒行割损礼，也请原谅那地位在你之上的\Person{伯多禄}，他因惧怕信主的犹太人而做了类似的事。经上又记载：\BibleQuote{此后，保禄又住了多日，然后辞别弟兄们，乘船往叙利亚去；同他一起的有普黎斯加和阿桂拉，因为他在耕格勒许过愿，就在那里剃了头。}\BibleRef{宗 18:18} 就算承认在其他情况下他是因惧怕犹太人而被迫做了他不愿做的事，那么，他为什么又按照自己许的愿而留起头发，随后在耕格勒又依法剃头，正如那些按照梅瑟法律许愿献身于天主的\OriginalTerm{纳齐尔人}{Nazarites}通常所做的那样呢？\par

10. 但这些事若与随后的事相比，还算小事。神圣的历史学家\Person{路加}进一步叙述：\BibleQuote{我们到了\Place{耶路撒冷}，弟兄们欢欢喜喜地接待了我们。}次日，\Person{雅各伯}及所有长老，在表示赞同他的福音后，对保禄说：\BibleQuote{兄弟，你看，在犹太人中信主的有多少万，他们都是热爱法律的人；他们听说你教训所有在外邦人中的犹太人背弃\Person{梅瑟}，说他们不必给自己的孩子行割损礼，也不必遵守规例。那么这怎么办？群众必定会聚集，因为他们听说你来了。所以你就照我们的话做吧：我们这里有四个人都有愿在身；你带这些人去，同他们一起行取洁礼，并替他们出费，使他们得剃头；这样，众人就会知道先前所听到关于你的事都是假的，而你本身仍是循规蹈矩，遵守法律的。}于是\Person{保禄}带着这些人，第二天就同他们一起行了取洁礼，进入圣殿，报告取洁的日子满了，等待为他们每一个人献祭。\BibleRef{宗 21:17} 啊，\Person{保禄}！在这里容我再问你：你为什么剃头？为什么按照犹太礼仪赤足行走？为什么献祭？为什么依照法律为你宰杀牺牲？你无疑会回答说：\Quote{为了避免冒犯那些信主的犹太人。}你为争取犹太人，就装成犹太人；\Person{雅各伯}和其他所有长老教给了你这种装假。然而你终究未能逃脱。因为你在引发的骚乱中险些被杀，幸为千夫长所救，并由他派人严密护送至\Place{凯撒勒雅}，以免犹太人将你这个\OriginalTerm{假善人}{dissembler}和法律的破坏者杀死；你从\Place{凯撒勒雅}来到\Place{罗马}，在你所租的房子里，向犹太人和外邦人宣讲\Person{基督}，最后你在\Person{尼禄}的刀下，用自己的血给\OriginalTerm{证道}{testimony}盖上了印。\par

11. 因此，我们认识到，由于惧怕犹太人，\Person{伯多禄}和\Person{保禄}都曾假装遵守法律的规条。保禄怎能厚颜且放肆地指责别人做了他自己也做过的事呢？至少我，或者不如说，在我之前的其他人，已经按照他们所能找到的最好方式解释了这件事，他们并未如你所指责的那样，为了宗教的利益而辩护使用虚假，而是教导一种明智的分辨所带来的可敬的操练；他们力求既显示宗徒们的智慧，又遏制\Person{波尔菲里}那些无耻的亵渎。波尔菲里说，伯多禄和保禄像孩子般互相争吵，并声称保禄因伯多禄的卓越而心生嫉妒，自夸地写了一些事，这些事他要么没有做过，要么如果做了，也是以不可原谅的鲁莽做的，却指责别人做了他自己做过的事。他们在回答他时，给出了他们所能找到的关于那段经文的最好解释；那么你打算提出什么解释呢？既然你拒绝了古代注释家的意见，你必定是打算说出一些比他们更好的话。\par

% structure: p0021 source=h2 target=subsection reason=relative-heading-rank; base=h2

\subsection{第4章}

12. 你在信中说：\Quote{你不要求我教导你，宗徒所言\BibleQuote{向犹太人，我就成为犹太人，为赢得犹太人；}\BibleRef{格前 9:20} 以及同段经文中其他类似的话，是应归于怜悯之爱的同情心，而非故意欺骗的诡计。因为服事病人的人，仿佛自己成了病人；他并非假装发热，而是以真诚同情者的心境，设想自己若处在病人的位置，希望别人怎样对待自己。\Person{保禄}确实是犹太人；在他成为基督徒之后，也未抛弃那些犹太礼仪，那是该民族以正当方式为某一特定时期所领受的。所以，即使身为基督的宗徒，他也参与遵行这些礼仪，但目的在于表明，这些礼仪对那些在信了基督之后仍愿持守从祖先学来的法律礼节之人，决无害处；只须他们不将救恩的希望建立在这些礼节之上，因为这些礼节所预表的救恩，如今已由主耶稣带来了。} 你将你的整个论证铺陈为一篇最冗长的论述，其要旨是：\Person{伯多禄}认为法律对信奉基督的犹太人仍有约束力，这并没有错；他偏离正轨之处在于，强迫外邦归正者遵守犹太习俗。然而，倘若他强迫了他们，并非以导师的权威，而是以自己实践的榜样。而\Person{保禄}，按你的看法，并非抗议\Person{伯多禄}本人所行之事，而是责问\Person{伯多禄}为何强迫外邦人遵守犹太习俗。\par

13. 因此，所辩论的问题，或者更确切地说，你对此的意见，可归结为：自从宣讲基督的福音以来，信主的犹太人遵守法律的规条，\Emph{即}如\Person{保禄}那样奉献祭献，如\Person{保禄}给\Person{弟茂德}行割损那样给自己的孩子行割损，以及遵守犹太安息日，如同所有犹太人向来所行的那样，这样做是好的。如果这是真的，我们便陷入了\Person{克林妥}和\Person{艾比恩}的异端，他们虽然信基督，却因这一个错误——将法律的礼仪与基督的福音混杂，既宣认新信仰，又不放弃旧事物——而被教父们予以绝罚。我何必再提那些妄自称基督徒的\Person{艾比恩}派呢？在我们这个时代，东方所有会堂的犹太人中，存在一个称为米尼派的教派，至今仍受法利塞人谴责。此派的信徒通常被称为纳匝肋派；他们信基督，天主子，由童贞玛利亚所生；并说那位在\Person{般雀·比拉多}手下受难并复活者，正是我们所信的那一位。但他们既想当犹太人，又想当基督徒，结果既不是犹太人，也不是基督徒。因此，我恳求你，你自以为蒙召来医治我这不过是针尖刺戳般的小伤口，却当将你的医术用于那由标枪猛力刺入而造成的重创。因为在圣经注释中展示教父们所持的各种不同意见的过失，与将最致命的异端重新引回教会的罪责相比，实在不可同日而语。然而，如果我们别无选择，必须让犹太人连同其法律规定的习俗一起进入教会；简言之，若宣布他们可以在基督的教会中继续行他们在撒殚会堂中惯常所行之事，那么我告诉你我对这事的看法：他们不会成为基督徒，反倒会使我们变成犹太人。\par

14. 有哪个基督徒愿意听信你信中所说的呢？\Quote{保禄确实曾是犹太人；他在成为基督徒后，并未放弃那个民族以正确的方式、并在一定时期内所领受的那些犹太礼仪。因此，即使在他作了基督的宗徒之后，他仍参与遵行这些礼仪；但目的只是为表明，对于那些即使已信了基督后仍愿保留他们由祖先通过法律所学得的礼仪的人，这些礼仪本来是毫无害处的。}\newline{}现在我恳求你耐心听我的抱怨。保禄即使在他作基督的宗徒时，也遵行犹太礼仪；而你断言，对于那些愿意保留他们由祖先通过法律所学得的礼仪的人，这些礼仪毫无害处。我则相反，我要坚持，并且即使全世界都反对我的观点，我仍要公然宣告：犹太礼仪对基督徒既是有害的，也是致命的；而且无论谁遵行它们，不管他原本是犹太人还是外邦人，都被抛入丧亡的深渊。\BibleQuote{因为基督就是法律的终向，使凡相信的人获得正义，}\BibleRef{罗 10:4} 就是指犹太人和外邦人两者；因为若将犹太人排除在外，祂就不是使每一个相信的人获得正义的法律的终向了。\newline{}而且，我们在\WorkTitle{福音}中读到：\BibleQuote{法律及先知到若翰为止。} 另外，还有一处：\BibleQuote{为此犹太人越发想要杀他，因为他不单犯了安息日，并且又称天主是自己的父，使自己与天主平等。}\BibleRef{若 5:18} 又说：\BibleQuote{从他的满盈中，我们都领受了恩宠，而且恩宠上加恩宠。因为法律是借梅瑟传授的，但恩宠和真理却是由耶稣基督来的。}\BibleRef{罗 1:16} 取代已逝去之法律的恩宠，我们领受了常存的福音的恩宠；取代旧约时代的影像和预像，真理却由耶稣基督来了。\newline{}\Person{耶肋米亚}也曾奉天主的名这样预言说：\BibleQuote{时日将到——上主的断语——我必要与以色列家和犹大家订立新盟约，不像我昔日握住他们的手，引他们出离埃及时，与他们的祖先订立的盟约。}\BibleRef{耶 3:31} 你要注意先知说的，不是对外邦人，因为他们未曾参与过先前的盟约，而是对犹太民族说的。那位通过梅瑟赐给他们法律的，应许以福音的新盟约取而代之，好使他们不再生活在文字的古旧中，而是生活在精神的新意中。\newline{}此外，引出这个个问题的保禄本人，也多次表达了这样的见解，我因愿简略，仅摘录少数于此：\BibleQuote{请注意\Person{保禄}对你们说：若你们还受割损，基督对你们便毫无益处。} 又说：\BibleQuote{凡是愿意因\WorkTitle{法律}成义的，你们都与基督断绝了关系，由恩宠上跌了下来。} 又说：\BibleQuote{如果你们随从他的引导，就不在法律权下。} 由此可见，那顺从法律的，并没有圣神；这种顺从，并不像我们的祖先所断言宗徒们所做的那样，是出于明智的考虑而假意为之，而是如你所设想的那样，真心实意地顺从。\newline{}至于这些法律规定的性质，让我们从天主的亲自教训中学习吧：\BibleQuote{我给了他们，}祂说，\BibleQuote{不美好的法度，和不能赖以生存的律例。}\BibleRef{则 20:25} 我说这话，并非要像\Person{摩尼}和\Person{马西}那样毁坏法律，我从宗徒言知道法律是圣洁而属神的；而是因为当\BibleQuote{信德来到，} 且时期一满\BibleQuote{天主就派遣了自己的儿子，生于人，生于法律权下，为把在法律权下的人赎出来，使我们获得义子的地位，}\BibleRef{迦 4:4} 并且不再生活在作为启蒙师的法律之下，而是生活在如今已达成年且为主的后嗣之下。\par

15. 您信中进一步说：\Quote{因此，他在伯多禄身上所斥责的，并不是他遵守从祖先传下来的习俗；如果伯多禄愿意的话，他可以这样做，而不必被指责为欺骗或言行不一。} 我再次说：您既然是主教，是基督的众教会中的教师，如果愿意证明自己的主张，就接纳一个犹太人，他在成为基督徒后，若生了儿子，就给儿子行割损礼，遵守犹太的安息日，禁戒天主所造本应怀着感恩享用的食物，并在正月十四日晚宰杀逾越节羔羊；当您这样做时，或更确切地说，拒绝这样做时（因为我知道您是基督徒，不会犯亵渎的行为），您就会被迫，无论情愿与否，放弃您的意见；那时您就知道，拒绝别人的意见比自己确立主张更为困难。再者，我们恐怕不信您的话，更确切地说，恐怕不理解您的话（因为往往一篇过分冗长的论述不易理解，在辩论中不熟练的人较少指责它，因为其弱点不容易被察觉），您便反复重申自己的意见，用以下的话来灌输：\Quote{保禄舍弃了犹太人一切邪恶的特色，尤其是这一点，即\BibleQuote{他们不认识天主的正义，企图建立自己的正义，而不服从天主的正义}。\BibleRef{罗 10:3} 此外，他在这一点上也与犹太人不同：在基督受难与复活之后——在祂内，按照\Person{默基瑟德}的品位，恩宠的奥迹已被赐予并显明——他们仍认为必须遵守，不是出于对旧习俗的敬重，而是视为得救所必需的旧约的圣事；这些圣事确曾一度是必需的，否则，\Person{玛加伯}为坚持遵守它们而殉道，就无益且徒然了。\BibleRef{加下 7:1} 最后，在这一点上，保禄也与犹太人不同：犹太人把宣扬恩宠的基督徒宣道者当作法律的仇敌来迫害。这一切以及所有类似的错误和罪恶，他声明都视之为损失和粪土，为获得基督。}\par

16. 我们从您那里得知保禄舍弃了犹太人哪些邪恶的特色；现在让我们从您的教导中得知他保留了犹太人哪些好的东西。您会回答说：\Quote{他们在礼仪方面的遵守，继续遵循他们祖先的做法，保禄自己也同样遵守这些礼仪，但根本不认为它们为得救所必需。} 我不完全理解您所说的\Quote{根本不认为它们为得救所必需}是什么意思。因为如果它们无助于得救，为什么要遵守呢？如果必须遵守，它们就必定有助于得救；尤其看到，因着遵守它们，有些人成了殉道者：若非有助于得救，它们就不会被遵守。因为它们不是无善无恶、中性的东西，如哲学家所说的那种既非善亦非恶。节制是善，放纵是恶：介于两者之间，作为没有道德品质的中性事物，就如走路、擤鼻涕、吐痰之类。这种行为既非善亦非恶；因为无论你做或不做，都不影响你是否成为义人或罪人。但遵守法律礼仪并非中性的；它要么是善，要么是恶。您说它是善，我断定它是恶，而且不仅当由外邦皈依者遵行时为恶，就是当由已信主的犹太人遵行时也是恶。在这一点上，如果我没弄错，您是在避免一个错误时陷入了另一个错误。因为您在提防\Person{波尔费留}的亵渎言论时，却陷入\Person{厄俾翁}的圈套；宣称法律对出自犹太人的信者仍有约束力。您又意识到所说是危险的教义，便试图用多余的话来修饰：即，\Quote{法律的遵守不应出于任何信仰，例如像犹太人那样认为这是得救所必需的，也不应出于任何误导人的伪善，如保禄斥责伯多禄的那种。}\par

因此，\Person{伯多禄}假装遵守法律；但这位批评伯多禄的人却勇敢地遵行了法律所规定的。你信中接下来的话是：\Quote{因为，如果\Person{保禄}为了争取犹太人，而假装为犹太人，按部就班地遵守这些\OriginalTerm{礼仪规定}{sacraments}，那么，当他为了争取那些没有法律的人，而成为如同没有法律的人，为什么也不同外邦人一起参与异教祭献呢？对此的解释是：他参与犹太礼仪，因为他自己就是犹太人；而当他讲述这一切我引述的话时，他的意思并非假装自己是某种不是的人，而是出于真实的同情，感到自己必须给予他们这样的帮助，就像如果他自己陷入歧途，自己也需要这样的帮助一样。在这里，他所运用的，不是欺骗者的狡诈，而是慈悲拯救者的同情。}这是对保禄的辉煌辩护！你证明他并非假装分担犹太人的错谬，而是实际卷入其中；并且他拒绝模仿伯多禄的欺骗行径——出于对犹太人的畏惧而掩饰自己的真实身份，反而毫不保留地公然承认自己是犹太人。哦，宗徒那闻所未闻的怜悯！为求使犹太人成为基督徒，他自己却成了犹太人！因为他若未曾像他们一样奢侈，便无法说服奢侈者变得节制；而他若不以亲身经历他们的可怜处境，便无法如你所说，以他的怜悯帮助那些可怜的人！那些被争辩的激烈冲动所裹挟，且因爱那已被废除的法律，而将\Person{基督}的宗徒变成犹太人的人，真是可悲，值得最深切的怜悯！终究，我的意见与你的并无太大差别：因为我说，\Person{伯多禄}和\Person{保禄}都因惧怕信了主的犹太人，而遵行——或不如说假装遵行——犹太法律的规定；而你却主张他们是出于怜悯，\Quote{不是用欺骗者的狡诈，而是用慈悲拯救者的同情。}但双方都同样承认，他们（无论出于恐惧还是出于怜悯）假装成自己并不是的人。至于你反对我们观点的论证：如果他为了犹太人而成为犹太人，那么他就应为了外邦人而成为外邦人，这论点反而有利于我们的看法而不是你的：因为他既未真正成为犹太人，也未真正成为外邦人；正如他未真正成为外邦人，也未真正成为犹太人。他与外邦人的认同在于：他接纳了信仰\Person{基督}而未受割损者为基督徒，并允许他们毫不顾忌地享用犹太法律所禁止的食物；而不像你所设想的那样，参与他们崇拜偶像的活动。因为\BibleQuote{在基督耶稣内，受割损算不得什么，不受割损也算不得什么，只该遵守天主的诫命。}\par

因此，我恳求你，并且极其迫切地请求，宽恕我这番试图讨论此事的卑微尝试；并且我若有任何逾越之处，请你归咎于你自己，因为你迫使我回信作答，又将我说成如\Person{斯泰西科鲁斯}那般盲目。不要将教导说谎的污名加在我这个\Person{基督}的追随者身上，他曾说：\BibleQuote{我是道路、真理、生命。}\BibleRef{若 14:6} 我身为真理的朝拜者，绝不可能屈服于虚假的轭下。此外，请不要煽动那些不学无术的民众来反对我；他们尊你为他们的主教，并以对待司铎职务的敬意重视你在教会中的讲论，却轻视像我这样衰老残躯、远离尘嚣、栖身隐修院的老者；你最好去找其他更适于接受你教导或纠正的人。因为你的声音几乎无法传到我这里，我们之间被海陆相隔如此遥远。而且，你若偶然写信给我，\Place{意大利}和\Place{罗马}必定早在信件送达我这唯一应收之人以前，便已知悉其内容。\par

% structure: p0029 source=h2 target=subsection reason=relative-heading-rank; base=h2

\subsection{第五章}

在另一封信中，你问及，我先前对一些\OriginalTerm{正典书卷}{canonical books}所做的译本为何仔细标注了\OriginalTerm{星号和短剑号}{asterisks and obelisks}，而后我发表的译本却没有这些标记。你必须原谅我这么说，因为在我看来，你似乎并不明白此事：因为先前的译本是依据\WorkTitle{七十贤士译本}；标注短剑号的地方，意在指明七十子译者所说的话比希伯来文中所见的要多。而星号则指明\Person{奥力振}从\Person{德敖多齐奥}的译本中所增添的部分。在那个译本中，我是从希腊文翻译；但在后来的译本中，我是直接从希伯来文本翻译，我表达了我所理解的意思，注意保留确切的意义甚于字词的次序。我感到惊讶的是，你并不阅读七十子译者最初流传于世的原貌，反而阅读那经过奥力振以短剑号和星号修正——或更确切地说，败坏——的文本；并且你拒绝接受一个基督徒人士所作的任何翻译，无论其如何拙陋，尤其是鉴于奥力振所增添的内容，乃借自一个在\Person{基督}受难后，身为犹太人和亵渎者的人的版本。你希望成为七十子译者真正的仰慕者和追随者吗？那么就不要阅读星号下所发现的内容；宁可将其从经卷中删去，这样你才能证明自己确实是古人的追随者。然而，你若这样做，就不得不指责所有教会图书馆；因为你几乎找不到几本手抄本没有这些添加的内容。\par

% structure: p0031 source=h2 target=subsection reason=relative-heading-rank; base=h2

\subsection{第六章}

现在，就你所说的我不应在古人已经完成翻译后再提供译本，以及你所使用的新奇\OriginalTerm{三段论}{syllogism}，略作回应：\Quote{七十子译者给予解释的经文，要么晦涩，要么明白。如果它们晦涩，那么人们相信你与其他人一样可能犯错；如果它们明白，人们不会相信七十子译者可能犯错。}\par

所有在我们之前的，在主内作我们前辈、从事解释圣经工作的注释家，都已解释了那些隐晦的或明显的经文。如果某些段落是隐晦的，在他们之后，你怎能擅自讨论他们无法解释的内容呢？如果那些段落是明显的，你着手处理那些他们不可能遗漏的内容，便是浪费时间。这个三段论特别适用于 \WorkTitle{圣咏集}，对此，希腊注释家们写了许多卷：即，1\Emph{st}.，\Person{奥力振}；2\Emph{nd}.，\Place{凯撒勒亚}的\Person{欧瑟比}；3\Emph{rd}.，\Place{赫拉克勒亚}的\Person{狄奥多若}；4\Emph{th}.，\Place{斯齐托颇利斯}的\Person{阿斯特里}；5\Emph{th}.，\Place{劳狄刻雅}的\Person{阿波林纳}；及6\Emph{th}.，\Place{亚历山大}的\Person{迪狄默}。据说还有关于 \WorkTitle{圣咏集} 选段的较小著作，但目前我谈论整部书。此外，在拉丁作家中，\Place{普瓦提埃}的主教 \Person{依拉略} 和 \Place{韦尔切利} 的主教 \Person{欧瑟比} 曾翻译了 \Person{奥力振} 和 \Place{凯撒勒亚}的\Person{欧瑟比}，前者在某些方面被我们自己的 \Person{盎博罗削} 所遵循。现在，我请教你的明智，为何在这么多且如此胜任的诠释者的辛劳之后，你在某些段落的解释上与他们不同？如果 \WorkTitle{圣咏集} 是隐晦的，那么必须相信你同样可能犯错，就像其他人一样；如果它们是明显的，令人难以置信的是，其他人竟会犯错。在这两种情况下，按你自己的说法，你的解释是多余的辛劳；本着同一原则，在别人发表了意见之后，没有人会再冒险谈论任何主题，也没有人有自由写任何关于他人已处理过的事，无论该事多么重要。\par

然而，更符合你明达判断的做法是，将你对自己所容忍的自由同样赋予所有其他人。因为在我试图将圣经的校订 \OriginalTerm{希腊文}{Greek} 译本翻译为 \OriginalTerm{拉丁文}{Latin}，以造福与我操同一语言的人时，我并非致力于取代久已为人所珍视的译本，而是仅将那些被犹太人删节或篡改之处明显地凸显出来，以便 \OriginalTerm{拉丁文}{Latin} 读者能知悉在 \OriginalTerm{希伯来文}{Hebrew} 原文中所发现的内容。如果有人不愿读它，没有人强迫他。让他满意地喝旧酒，并藐视我的新酒，\Emph{即}，我为了解释先前的作家而发表的那些语句，其用意在于通过我的评注，使他们在其中看来隐晦的地方变得更为明显。\par

至于解释圣经时应遵循的原则，已陈述于我所写的书中，以及在我所编版本中附于每卷圣经前的所有引言中；我认为，引荐这些给审慎的读者便足够了。既然如你所说，你赞同我校订 \WorkTitle{新约} 译本的辛劳——同时你也以此作为理由，即很多人通晓 \OriginalTerm{希腊文}{Greek}，因此是我工作的胜任评判者——但若以同样的信任来对待我在 \WorkTitle{旧约} 中的忠实，原本才是公平的；因为我并非随从自己的想象，而是按照我发现那些讲 \OriginalTerm{希伯来文}{Hebrew} 的人所理解的意义，来翻译了神圣的言语。如果你对任何段落的有疑问，去问犹太人原文的意思是什么。\par

21. 或许你会说：\Quote{如果犹太人拒绝回答，或故意欺骗我们，那怎么办？} 能想象所有的犹太众人，在被问及我的译本时，一致同意沉默，且没有一个通晓 \OriginalTerm{希伯来文}{Hebrew} 的人被找到吗？或者他们都会效法你提到的、在某个小镇上合谋损害我声誉的那些犹太人吗？因为你在信中拼凑了如下的故事：— \Quote{某位主教，我们的一位弟兄，在他所主持的教会中采用了你的译本诵读，读到 \BibleBook{约纳}{先知书} 中的一个词，你的译法与此前历代所有敬拜者的感官和记忆所熟悉的，并在教会中传颂了如此多代的译法大相径庭。于是，在会众中，尤其是在希腊人中，引起了如此大的骚动，他们纠正所读的内容，并指责译本为错误的，以致主教被迫去寻求犹太居民（那是在 \Place{奥亚} 镇）的见证。这些人，无论是出于无知还是恶意，回答说 \OriginalTerm{希伯来文}{Hebrew} 抄本中的词语在 \OriginalTerm{希腊文}{Greek} 译本中，以及在出自它的 \OriginalTerm{拉丁文}{Latin} 译本中，都正确地译出了。我还有什么需要多说的呢？那人被迫在所涉段落中纠正你的译本，仿佛它被错译了一般，因为他不想失去全体会众——一场他勉强逃脱的灾难。由此，我们也认为你可能偶尔会出错。}\par

% structure: p0037 source=h2 target=subsection reason=relative-heading-rank; base=h2

\subsection{第七章}

22. 你告诉我，我在\BibleBook{约纳}{书}中某个词的翻译错了，并且一位可敬的主教因百姓群情激愤，险些失去他的职务，而这全因这一个词的不同译法所致。同时，你却对我隐瞒了我误译的是哪个词；从而使我无法为自己作任何辩护，免得我的答复会对你的反对意见构成致命反驳。也许，那是关于\Quote{葫芦}的旧争论又被重新提起，自那位著名人物——他在当时集科尔乃略家族与\Person{阿西尼乌斯·波利奥}的祖先荣耀于一身——指控我在译文中用了\Quote{常春藤}而非\Quote{葫芦}以来，这场争论已沉寂了漫长岁月。我已在\BibleBook{约纳}{书}注释中对此作了充分答复。目前，我认为只需指出：在那段经文中，\WorkTitle{七十贤士译本}作\Quote{葫芦}，而\Person{阿基拉}等人则把这词译作\Quote{常春藤}（\OriginalTerm{常春藤}{\GreekText{κίσσος}}），希伯来文抄本则作\Quote{ciceion}，按现今的叙利亚语，则说\Quote{ciceia}。这是一种灌木，叶大如葡萄树，一经种植，便迅速长成一棵小树般大小，靠自己的茎直立，无需像葫芦和常春藤那样依赖芦苇或竿子支撑。因此，如果我逐字翻译，用了\Quote{ciceia}这个词，没有人会懂它的意思；如果我用了\Quote{葫芦}，我说的就不是希伯来文所载的内容。因此，我写下了\Quote{常春藤}，以便与其他所有译者保持一致，不致立异。但是，如果你那里的犹太人——正如你本人所暗示的——出于恶意或无知，声称希腊文和拉丁文译本中所见的那个词也见于希伯来原文，那么显然，他们要么不通希伯来文，要么乐于说谎，以取笑那些种植葫芦的人。\par

结束这封信时，我恳请你体谅一位现已年老、久已退出战场的士兵，不要强迫他重上战阵，再让他的生命冒战争的危险。你尚年轻，又被擢升到主教尊位这显要职位，理应专心教导子民，并用来自丰饶非洲的新储备充实\Place{罗马}。我则满足于在隐修院的一个幽暗角落里，不发大声，只有一人听我说话或为我诵读。\par

\SourceRecord{Translated by J.G. Cunningham.；Nicene and Post-Nicene Fathers, First Series；Vol. 1.；Edited by Philip Schaff.；Buffalo, NY: Christian Literature Publishing Co.,；1887.；<http://www.newadvent.org/fathers/1102075.htm>.}

% structure: p0001 source=h1 target=section reason=page-title

\section{书信第七十六篇（公元402年）}

听，多纳特派啊，天主教会对你们说：\BibleQuote{你们众人，你们这些人的子孙，你们的心硬要到何时？你们爱慕虚幻，追求虚伪究竟何为？} 你们为什么以可憎的\OriginalTerm{裂教}{schism}不虔，将自己从整个世界的合一中分裂出去？你们听从那些关于将圣经交给迫害者的谎言，这些谎言是由那些欺骗你们或自己受骗的人编造的，为要使你们死于异端分裂的状态；而你们却不听从这些圣经本身所宣告的，为使你们能在天主教会的平安中生活。因此，你们为何侧耳倾听那些人告诉你们从未能证实的事，却对天主如此说话的圣言充耳不闻：\BibleQuote{上主对我说：你是我的儿子，我今日生了你。你向我请求，我必将万民赐你作产业，以地极作你的领地。} \BibleQuote{所恩许原是向\Person{亚巴郎}和他的后裔所许诺的。并没有说\InnerQuote{后裔们}，好像是向许多人说的，而是向一个人说的，且说\InnerQuote{你的后裔}，就是指基督。}\BibleRef{迦 3:16} 而宗徒所指的恩许是：\BibleQuote{地上万民都要因你的后裔蒙受祝福。}\BibleRef{创 22:18} 因此，抬起你们灵魂的眼睛，看看在全世界，万民如何在\Person{亚巴郎}的后裔内蒙受祝福。\Person{亚巴郎}在他那时代相信了那未见之事；而你们看见了却拒绝相信那已实现的事。主的圣死是世界的赎价；他支付了赎价，为的是全世界；你们却不与全世界和谐共处，这本是对你们有益的，反而分裂出去，争论着试图破坏全世界，损害的是你们自己。现在请听\BibleBook{}{圣咏}中关于这赎价所说的话：\BibleQuote{他们穿透了我的手脚，我竟能数清我的骨骼；他们却冷眼观望着我。他们瓜分了我的衣衫，为我的长衣抽签。} 为什么你们竟犯下分裂主衣裳的罪，却不与全世界共同持守那件从上到下编织而成的\OriginalTerm{爱德之衣}{coat of charity}，连他的行刑者都没有撕裂它呢？同一篇\BibleBook{}{圣咏}中，我们读到全世界持守这一点，因为他说：\BibleQuote{整个世界将觉醒而归顺上主，天下万民将在你面前屈膝叩首；因为君王权属于上主，他是统治万民的那位。} 打开你们灵魂的耳朵，听：\BibleQuote{上主，全能者天主，出令传召大地，从日出之处直到日落之地。由美丽绝伦的\Place{熙雍}，天主发出光明。} 如果你们不愿理解这话，那么就请听主亲口所说的福音，他如何说：\BibleQuote{诸凡\Person{梅瑟}法律、先知并\BibleBook{}{圣咏}上指着我所记载的话，都必须应验；并且必须从\Place{耶路撒冷}开始，因他的名向万邦宣讲悔改，以得罪之赦。}\BibleRef{路 24:44} \BibleBook{}{圣咏}中\Quote{从日出之处直到日落之地}这句话，对应福音中\Quote{向万邦}；正如他在\BibleBook{}{圣咏}中说\Quote{由美丽绝伦的\Place{熙雍}}，在福音中则说\Quote{从\Place{耶路撒冷}开始}。\par

你们想象自己在收割之前，从那与麦子混在一起的莠子中将自己分离出来，这恰恰证明你们只是莠子。因为如果你们是麦子，你们就会容忍莠子，而不从基督田地中生长的分离出来。关于莠子，经上说：\BibleQuote{因为罪恶的增多，许多人的爱德必要冷淡；}但关于麦子，经上说：\BibleQuote{唯独坚持到底的，才可得救。}\BibleRef{玛 24:12} 你们有什么理由相信莠子增多并充满了世界，而麦子减少了，如今只存在于\Place{非洲}呢？你们自称是基督徒，却否认\Person{基督}的权威。他\BibleQuote{让两样一起长到收割的时候}；他没有说：\Quote{让麦子减少，让莠子增多}。他\BibleQuote{田地就是世界}；他没有说：\Quote{田地是\Place{非洲}}。他\BibleQuote{收割的时候是今世的终结}；他没有说：\Quote{收割的时候是\Person{多纳特}的时代}。他\BibleQuote{收割者是天使}；他没有说：\Quote{收割者是巡回苦行派的首领}。\BibleRef{玛 13:30} 但你们指控好麦子为莠子，正好证明你们自己是莠子；更糟的是，你们过早地将自己与麦子分离了。因为你们的一些前辈，你们仍顽固地留在他们的亵圣\OriginalTerm{裂教}{schism}中，他们将圣经手稿和教会的圣器交给了迫害者（正如城市档案中所见）；他们中的另一些人则轻忽了这些人所承认的过错，仍与他们保持共融；两派作为受迷惑的派别一同聚集在\Place{迦太基}，不经审讯便以那过错（他们之间却已同意宽恕）谴责他人，然后设立一位主教以对抗已祝圣的主教，并树立一座祭台以对抗已被认可的祭台。后来，他们差人送信给皇帝\Person{君士坦丁}，恳求指定海外教会的主教们来仲裁\Place{非洲}的主教们之间的事。当他们所寻求的裁判官得到允准，并在\Place{罗马}作出判决后，他们却拒绝服从，反而向皇帝控告主教们审判不公。从另一批被派往\Place{阿尔勒}审理此案的主教团的判决，他们又上诉到皇帝本人。当他听了他们，他们被证实犯有诬告之罪后，他们仍执迷不悟。醒悟吧，为了你们的救恩！热爱和平，回归合一！无论何时你们愿意，我们随时准备详细叙述我们提及的这些事件。\par

赞同恶人行为的人，就是恶人的同伙；却不是那让稗子在主的田里生长直到收割，或让糠秕存留直到最终簸扬时刻的人。如果你们憎恶作恶的人，就应使自己摆脱\OriginalTerm{裂教}{schism}的罪。如果你们果真害怕与恶人交往，你们就不会多年来容许\Person{奥普塔图斯}留在你们中间，当他生活在最昭彰的罪恶中时。如今你们称他为殉道者，那么，你们若前后一致，就必须称他为之而死的那一位为基督。最后，基督徒世界何事得罪了你们，你们竟疯狂而邪恶地从中分裂出去？而那些\Person{马克西米安}的追随者，你们在谴责了他们之后，又借民政当局的命令将他们从各教会中暴力驱逐，随后却又尊荣地将他们接回，他们对你们的敬重又有何理由呢？基督的和平何事得罪了你们，使得你们通过脱离自己所诽谤的人来抗拒它？而\Person{多纳图斯}的和平又怎样赢得了你们的欢心，使你们为了促进它而接回那些你们所谴责的人？\Place{穆斯蒂}的\Person{费利齐安}如今是你们中的一分子。关于他，我们曾读到，他先被你们的\OriginalTerm{公会议}{council}谴责，随后在\OriginalTerm{总督}{proconsul}的法庭上被你们控告，并在\Place{穆斯蒂}城中，据市政记录所载，遭到了攻击。\par

如果交出\OriginalTerm{圣书}{sacred books}加以销毁是一桩罪行——在那位烧毁\BibleBook{耶肋米亚}{书}的国王的事例中，天主以使其被掳死于战争来惩罚他——那么\OriginalTerm{裂教}{schism}的罪又该是多么大得多呢！因为对于你们将\Person{马克西米安}的追随者与之相比的那些裂教的始作俑者，地便裂开口，将他们活活吞下。\BibleRef{户 16:31} 那么，你们为什么用你们无法证实的交出圣书的指控来反对我们，同时又谴责并重新接纳你们当中的裂教者呢？如果你们因遭受皇帝的迫害而声称自己是正确的，那么\Person{马克西米安}的追随者必定比你们有更强有力的理由，因为你们自己就在\OriginalTerm{公教皇帝}{Catholic emperors}差遣给你们的法官的协助下迫害了他们。如果只有你们拥有\OriginalTerm{圣洗}{baptism}，那么你们认为，在那些由\Person{费利齐安}于被你们判处谴责期间施洗、后来又随他一同过来的人身上，\Person{马克西米安}追随者所施行的圣洗有何分量呢？让你们的主教至少向你们的平信徒回答这些问题吧，如果他们不愿与我们辩论；而你们，若珍惜你们的\OriginalTerm{救恩}{salvation}，就当思考他们拒绝与我们讨论的那种\OriginalTerm{教义}{doctrine}必定是什么样的。如果狼有足够的智慧避开牧羊人的路，为什么羊群却如此失去了智慧，竟走进狼的洞穴里去呢？\par

\SourceRecord{Translated by J.G. Cunningham.；Nicene and Post-Nicene Fathers, First Series；Vol. 1.；Edited by Philip Schaff.；Buffalo, NY: Christian Literature Publishing Co.,；1887.；<http://www.newadvent.org/fathers/1102076.htm>.}

% structure: p0001 source=h1 target=section reason=page-title

\section{书信七十七（公元404年）}

\Emph{致\Person{斐理斯}和\Person{希拉里努斯}，我最亲爱的尊主，及堪受所有尊敬的弟兄们，\Person{奥思定}在主内致候。}\par

我毫不惊奇看到信徒的心灵被\OriginalTerm{撒殚}{Satan}扰乱；你们要抵抗他，持续坚持那寄托于天主许诺的\OriginalTerm{望德}{hope}，天主是绝不会撒谎的。祂不仅屈尊许诺，将永远的赏赐赐予我们这些相信祂、寄望于祂、并在\OriginalTerm{爱德}{love}中坚持到底的人，而且还预言，在时间中，那些试探和考验我们\OriginalTerm{信德}{faith}的恶表必不缺乏；因为祂说过：\BibleQuote{由于罪恶增多，许多人的爱德必要冷淡；}但祂立刻又加上说：\BibleQuote{唯独坚持到底的，才可得救。} \BibleRef{玛 24:12}。那么，为什么人会感到奇怪，有人诽谤天主的仆人，既然无法使他们偏离正直的生活，就竭力玷污他们的名声？要知道，这些人如果对天主——这些仆人的主——任何公义而奥妙的计划的执行感到不满，而此执行与他们的愿望相违，他们便天天不停地亵渎天主。因此，我最亲爱的尊主和堪受尊敬的弟兄们，我呼吁你们的智慧，并劝勉你们以最合宜基督徒的方式锻炼你们的心志，用圣经来对抗人们空洞的诽谤和无根据的猜疑；圣经早已预言这些事必会发生，并告诫我们要勇敢面对。\par

因此，我简短地对你们的爱德说几句：\OriginalTerm{司铎}{presbyter}\Person{鲍尼法斯}并未被我查出犯有任何罪行，我从未相信，至今也不相信任何针对他的指控。那么，我怎能因我们主在福音中的这句话而充满畏惧，就命令从司铎名册上删除他的名字呢？那话说：\BibleQuote{你们用什么尺度量给人，也要用什么尺度量给你们。} \BibleRef{玛 7:2}。因为，他与\Person{斯佩}之间的争执，经双方同意，已呈交天主的仲裁；其方式，如果你们愿意，可以告知你们。那么，我是谁，竟敢擅自抢先天主的裁决，删除或略去他的名字？作为\OriginalTerm{主教}{bishop}，我不应草率地怀疑他；而作为一个人，我无法无误地判决那隐藏于我的事。即使在世俗事务中，当案件上诉至更高权威后，在等待那不允再上诉的裁决期间，一切程序都暂停；因为如果在案件尚待仲裁时有所更改，那将是对更高法庭的侮辱。那么，甚至最高的人间权威与天主的权威之间，又有着何等巨大的差距！\par

愿我们主天主的仁慈永不离开你们，我最亲爱的尊主，及堪受所有尊敬的弟兄们。\par

\SourceRecord{Translated by J.G. Cunningham.；Nicene and Post-Nicene Fathers, First Series；Vol. 1.；Edited by Philip Schaff.；Buffalo, NY: Christian Literature Publishing Co.,；1887.；<http://www.newadvent.org/fathers/1102077.htm>.}

% structure: p0001 source=h1 target=section reason=page-title

\section{第78封书信（公元404年）}

\Emph{致我在基督的爱内所服侍的、我所挚爱的\Place{希波}教会的弟兄们：神职班、长老及信众，我，\Person{奥思定}，在主内致候。}\par

1. 但愿你们，专注于天主的话，就不需要我的劝告来支持你们面对可能发生的任何恶表！但愿你们的安慰更来自那位，我们也由祂获得安慰的主；祂不仅预告了祂将为圣洁忠信的人赐下的美善，也预告了这世界将充满的邪恶；祂使这些事被记载下来，好让我们能怀着绝不逊于如今亲身经历世界终结之前所预告的邪恶的确定性，去期待那紧随世界终穷之后的祝福！故此，宗徒也说：\BibleQuote{凡从前所写的，都是为教训我们而写的，好使我们因圣经的忍耐和安慰，获得希望。}\BibleRef{罗 15:4}我们主自己又为何断定必须不仅说：\BibleQuote{那时，义人将在他们父的国里，发光如同太阳。}\BibleRef{玛 13:43}（这事将在世界终穷后实现），而且还呼喊说：\BibleQuote{世界因恶表而遭祸了！}\BibleRef{玛 18:7}若不为了不让我们自欺，以为能不经过在尘世磨难中的试探而得胜，便可抵达永恒真福的居所？祂又为何必须说：\BibleQuote{因罪恶增多，许多人的爱心就冷淡了，}若不为了使那些祂在下一句所说：\BibleQuote{但坚持到底的，必能得救}\BibleRef{玛 24:12}的人，在看到因罪恶增多而爱心冷淡时，能免于因这些事而感到羞愧、恐惧或被忧愁压倒，仿佛它们是怪异意外之事，反而因见证这些预定了在世界终穷前发生的事件，得助以坚忍到底，从而在终穷之后，获得在那无尽生命中和平为王的赏报？\par

2. 因此，弟兄们，关于因司铎\Person{卜尼法斯}而使某些人困惑的那个恶表，我并不对你们说：你们不应对此感到忧伤；因为凡不为此类事忧伤的人，就没有基督的爱，而那些喜好此类事的人则充满了魔鬼的恶意。然而，并非上述司铎有什么该受谴责的过错已为我们所知，而是我们院内有两个人处于这种状况，以致其中一人应被毫无疑问地视为邪恶；虽然另一人的良心并未受玷污，但他的好名声在一些人眼中已经丧失，并为其他人所猜疑。你们应为这些事忧伤，因为这值得哀悼；但不要忧伤到使你们的爱德冷淡，或对圣善生活漠不关心。倒要使其更炽热地燃烧于向天主的祈祷之中，若你们的司铎是无辜的，（我更倾向于这样相信，因为当他发现另一人的邪恶下流建议时，他既不肯同意，也不隐匿它，）愿天主的裁断能迅速恢复他行使职权，并洗清他的无辜；另一方面，如果他自知有罪（我绝不敢如此猜疑），并故意在未能败坏对方道德时，企图毁谤他的好名声，正如他指控其控告者所行的那样，愿天主不容许他掩藏自己的邪恶，以致人所不能发现的，能因天主的审判而显露出来，使这人或那人得以信服。\par

3. 因为这件案子已使我长久不得安宁，而我无法确定两人中谁是有罪的，尽管我较倾向于相信那位\OriginalTerm{司铎}{presbyter}是无辜的，起初我决意将他们二人一起交托在天主手中，暂且不下定案，直到我所怀疑的那个人做出什么事，为将他从我们中间逐出提供公正而无可置疑的理由。但当他竭力谋求晋升到圣职的行列，要么在本处从我这里获得晋升，要么在其他地方藉着我的推荐信获得晋升，而我绝对不能受到诱惑，既不愿对这样一个我所鄙弃的人施行覆手，授予圣秩，也不愿同意通过我的推荐将其介绍给任何弟兄，以达到同样目的，他便开始更加猛烈地要求：如果他不被提升为圣职人员，就不应允许\Person{鲍尼法斯}保留其司铎的身份。他提出这个要求之后，我看出\Person{鲍尼法斯}不愿意，因为别人对他圣善生活的怀疑，而使那些软弱和易起疑心的人跌倒，并且他情愿在人间承受名誉的损失，也不愿徒然坚持——甚至扰乱教会——在一场论辩中，此事本身就使他无法向那些不认识他的人、或对他疑惑或易于猜疑的人满意地证明他的无辜（他自己良心是清楚的），于是我决定用以下方法来查明真相。两人庄严约定前往一处圣地，在那里天主的更可敬畏的工程或许能更容易地显明他们二人中任何一个良心上的邪恶，并迫使有罪的认罪，无论是通过天主的判断，还是由于对判决的恐惧。诚然，天主无所不在，创造万物的那位不受任何地方的容纳或局限；祂的真正朝拜者应以心神和真理朝拜祂\BibleRef{若 4:24}，好使天主在暗中垂听，也在暗中使人成义并予以赏报。但至于人前可见的祈祷答复，谁能探究祂选定某些地方而非别处来施行奇迹干预的理由呢？许多人熟知真福\Person{斐理斯}遗体安葬之地的神圣，我愿他们前往那里；因为从那里，我们可以更容易、更可靠地获得一份书面记录，记载天主对他们二人中任何一个所作的启示。因为我亲知，在\Place{米兰}，在圣人们的墓前，那里邪魔以极为可奇可畏的方式被带来供认其恶行，有一个小偷来到那里想发假誓进行欺骗，却被强制认罪，归还了他所偷的东西；难道\Place{阿非利加}不也是遍地是圣殉道者的遗体吗？然而我们却不知在此地有任何类似的事发生。正如宗徒所宣讲的，治病的奇恩和辨别神恩的神恩并非赐予所有的圣人；同样，祂按己意分施每人恩宠的那位，乐意在圣人们的墓前施行这些奇迹的，也并非在所有墓前都施行。\par

4. 因此，虽然我原打算不让我心中这极重的负担为你们所知，免得我以痛苦而徒然的烦恼使你们不安，但天主却乐意让你们知道，这或许是为了使你们可以和我一起专心祈祷，恳求祂俯允，启示那祂知道而我们对此事无法得知的。因为我不敢擅自将这位司铎的名字从同僚的名册上删去，免得我似乎侮辱了天主的尊威——此案现在正系于祂的判断——如果我以任何仓促的决定抢先于祂的判决：即使在世俗事务中，当一件棘手的案件提交给更高的权威时，在下级的法官也不得在案件审理期间作出任何更改。况且，在一次主教会议上曾决定：任何神职人员，只要尚未证实有罪，不得暂停其共融，除非他未能亲自到场接受对他的控告的审查。然而，\Person{鲍尼法斯}谦逊地同意放弃他对推荐信的要求——使用这封信可使他在旅途上确保其身份得到承认——宁愿两人在双方都无人知晓的地方处于平等地位。现在，如果你们宁愿不宣读他的名字，以便我们\Quote{绝去机会}，正如宗徒所说的，不叫那些想找机会的人找到机会\BibleRef{格后 11:12}，好使他们不愿来教会的借口不成立，那么对他的名字的省略将不是我们的作为，而是那些为此事而这样做的人的作为。因为别人无知而拒绝在牌上宣读他的名字，这对他有什么害处呢？只要愧疚的良心没有将他的名字从生命册上抹去就好了。\par

5. 所以，敬畏天主的弟兄们，要记得宗徒\Person{伯多禄}的话：\BibleQuote{你们的仇敌魔鬼，如同咆哮的狮子巡游，寻找可吞食的人。}\BibleRef{伯前 5:8}当他不能通过引诱人陷入罪恶来吞噬人时，他就试图损害人的好名声，如果可能，使人屈服于人的羞辱和毁谤的舌头，从而落入他的口中。然而，如果他甚至不能玷污一个无辜者的好名声，他就试图说服人对自己的弟兄怀有恶意的猜疑，严厉地判断他，从而陷入网罗，成为易得的猎物。谁能知道或述说他的所有陷阱和诡计呢？然而，对于与我们目前情况更特别相关的三件事：首先，为了使你们不至因效法坏榜样而偏向邪恶，天主借着宗徒给你们这些警告：\BibleQuote{你们不要与不信的人共负一轭，因为正义与不法之间，那能有什么相通？或者说，光明与黑暗之间，那能有什么联系？}\BibleRef{格后 6:14}另一处说：\BibleQuote{你们不要为人所误：\InnerQuote{滥交败坏善行。}你们当醒寤，别再犯罪了。}\BibleRef{格前 15:33}第二，为了使你们不至屈服于毁谤者的舌头，祂借着先知说：\BibleQuote{你们认识正义，心怀我法律的百姓，不要怕人的侮辱，不要因人的毁谤而惊惶，因为他们必像衣服一样，为蠹虫所蛀，必像羊毛一样，为衣鱼所蚀；但我的正义永远长存。}\BibleRef{依 51:7}第三，为了使你们不至因对天主的任何仆人产生毫无根据的和恶意的猜疑而败坏，要记得宗徒的话：\BibleQuote{所以，时候未到，你们什么也不要判断，只等主来，祂要揭发暗中的隐情，且要显露人心的计谋：那时，各人才能由天主那里获得称誉。}\BibleRef{格前 4:5}也要记住这句话：\Quote{已启示的事是属于你们的，隐秘的事是属于上主你们的天主的。}\par

6. 确实明显的是，这样的事发生在教会中，没有不引起圣人和信徒的极大悲伤的；但让那预先讲了这一切事，警告我们不要因罪恶增多而爱心冷淡，要忍耐到底才能得救的天主作我们的安慰者。因为，就我而言，如果我内有基督之爱的火花，你们中谁软弱，我不软弱呢？谁跌倒，我不心焦呢？\BibleRef{格后 11:29}所以不要因你们屈服于无端的猜疑或受他人之罪的影响而加重我的痛苦。我恳求你们不要这样，免得我对你们说：\Quote{他们增加了我的伤痛。}因为更容易忍受那些公开以我们这些痛苦为乐的人的辱骂，论到他们，在论及基督自己时已有预言：\BibleQuote{坐在城门口的人对我议论纷纷，醉汉以我为歌曲取乐。}我们已受教为他们祈祷并寻求他们的福祉。因为他们为什么坐在城门口，守望什么呢？难道不是为这个：一旦有某位主教、司铎、隐修士或修女跌倒，他们就有理由相信、夸耀并坚称所有人都与那跌倒的人一样，但所有人都不能被定罪和揭露？然而这些同样的人并不因某位已婚妇女被发现犯奸淫，就立即休弃自己的妻子或控告自己的母亲。但一旦有任何虔诚信仰告白的人被诬告或确证犯了罪，这些人就不断不倦地努力使这指控被所有人谴责，责备所有虔诚的人。因此，那些在我们悲痛之事中急切地寻找他们恶毒舌头所喜之物的家伙，我们可以比作那些狗，如果它们确实被理解为增加他的痛苦的话，它们舔躺在富翁大门前的乞丐的疮痍，乞丐耐心忍受一切艰辛与侮辱，直到他在\Person{亚巴郎}怀中得到安息。\BibleRef{路 16:21}\par

7. 你们这些对天主怀有某种希望的人啊，不要加增我的忧伤。不要让这些所舔的伤疮因你们而增多，我们为你们时刻冒着危险，外有争斗，内有恐惧，经历城里的危险、旷野的危险、外邦人的危险、假弟兄的危险。我知道你们忧伤，但你们的忧伤比我的更痛切吗？我知道你们不安，我害怕因毁谤者的舌头，那些基督为之而死的人中有些软弱者会丧亡。不要让你们加增我的忧伤，因为这忧伤成为你们的并非由于我的过错。因为我曾尽可能谨慎，为能不让防止这邪恶的必要措施被忽略，也不让你们知晓这事，因为这只能给强壮者带来无益的烦恼，给软弱者带来危险的不安。但愿那位容许你们因知道这事受试探的天主，赐给你们力量忍受试炼，并\BibleQuote{照祂的法律指教你们，使你们在患难之日得享安宁，直到为恶人掘下陷阱。}\par

8. 我听说你们中有些人因这事比因那两位由多纳特派加入我们中的执事的堕落更为忧伤，好像他们使\Person{普罗库莱亚努斯}的规训蒙羞似的；而这事却阻挠了你们对我的夸耀，认为在我的规训下绝不会在圣职人员中出现这样的不一致。我要坦率地对你们说，无论做了这事的是谁，你做得不好。看，天主指教了你们：\BibleQuote{凡要夸耀的，当在主内夸耀；} \BibleRef{格前 1:31} 你们只应以这一件事指责异端者，即他们不是天主教徒。你们不要像这些异端者，他们因为在为自己的分裂辩护时无可置辩，便只知堆积对与他们分裂之人的指控，并且极其虚假地夸耀说，我们在这些方面具有一种令人不快的优越性，目的就是，既然他们既不能攻击也不能遮蔽圣经的真理——那散播各地的基督教会由之获得明证的真理——他们便使宣讲这真理的人失去人心，他们能随意指控这些人的任何事——只要他们想得出来的。\BibleQuote{你们却不是这样学了基督，如果你们真听过他，并受了他的教导。} \BibleRef{弗 4:20} 因为他亲自保护了他信从的子民，免得他们因邪恶——即使是在天主奥秘的管家身上的邪恶——而过度不安，因为他们行的是自己的恶，而说的是他的善。\BibleQuote{所以，凡他们对你们所说的，你们要行要守；但不要照他们的行为去做，因为他们只说不做。} \BibleRef{玛 23:3} 请务必为我祈祷，免得我\BibleQuote{给别人报捷，自己反而落选；} \BibleRef{格前 9:27} 但当你们夸耀时，不要因我，而要因主而夸耀。因为，无论我家的规训多么警醒，我不过是一个人，并且生活在人群当中；我决不敢妄想我的家比\Person{诺厄}的方舟更好，在方舟里的八个人中，竟有一人被弃绝；\BibleRef{创 9:27} 或比\Person{亚巴郎}的家更好，论到他的家，有话说：\BibleQuote{你该将这婢女和她的儿子赶出去；} \BibleRef{创 21:10} 或比\Person{依撒格}的家更好，论到他的孪生子，有话说：\BibleQuote{我爱了\Person{雅各伯}，而恨了\Person{厄撒乌}。} \BibleRef{拉 1:2} 或比\Person{雅各伯}本人的家更好，在他的家里，\Person{勒乌本}玷污了他父亲的床；\BibleRef{创 49:4} 或比\Person{达味}的家更好，在他的家里，一个儿子与他的姊妹行了丑事，\BibleRef{撒下 13:14} 另一个儿子竟反叛一位如此圣善仁慈的父亲；或比宗徒\Person{保禄}的同伴团体更好，如果保禄只与好人同住，他决不会如前面所引的那样说：\Quote{外有争斗，内有恐惧；} 也不会在论及\Person{弟茂德}的圣洁与忠诚时说：\BibleQuote{我没有一个像他那样诚心关照你们的人，因为其他的人都求自己的事，而不求耶稣基督的事；} \BibleRef{斐 2:20} 或比主\Person{基督}本人的门徒团体更好，在当中，十一位善人容忍了\Person{犹达斯}，那个贼与叛徒；最后，或比天堂本身更好，天使就是从那里堕落了的。\par

9. 我在我们的主天主面前，向你们的爱德坦率承认——自从我开始事奉他起，我就请他作我灵魂的见证——正如我几乎找不到比隐修院中行善之人更好的人，我也找不到比堕落的隐修士更坏的人；因此我想，\WorkTitle{默示录}上的话正适用于他们：\BibleQuote{让行义的，仍行义罢！让污秽的，仍受污秽罢！} \BibleRef{默 22:11} 所以，如果我们因一些污秽的瑕疵而忧伤，我们便因更多相反的例子而得安慰。因此，不要让那冒犯你们眼睛的渣滓，使你们憎恨榨油坊，正是从那里，主的仓库被供应了更为光明照耀的油，以使他们获益。\par

愿我们的主的仁慈保守你们在他的平安中，远离仇敌的一切罗网，我至亲爱的弟兄们。\par

\SourceRecord{Translated by J.G. Cunningham.；Nicene and Post-Nicene Fathers, First Series；Vol. 1.；Edited by Philip Schaff.；Buffalo, NY: Christian Literature Publishing Co.,；1887.；<http://www.newadvent.org/fathers/1102078.htm>.}

% structure: p0001 source=h1 target=section reason=page-title

\section{第79封书信（公元404年）}

\EditorialNote{对某位摩尼教导师的简短而严厉的挑战，他接替了\Person{福图纳图斯}（据信是\Person{费利克斯}）。}\par

你的回避企图毫无用处：你的真实品性即便远在千里之外也昭然若揭。我的弟兄们已向我报告了他们与你的交谈。你说你不怕死；这很好：但你应当惧怕那种因你对天主的亵渎言论而给自己招来的死亡。至于你理解到众人所知的可见死亡是灵魂与肉身的分离，这是一个无需高深智识便可明白的真理。但你所附加的说法，即死亡是善与恶的分离，难道你看不出，如果灵魂是善的而肉身是恶的，那么将二者结合起来的祂，便不是善的吗？然而你断言善的天主将二者结合了起来；由此可推论，祂要么是恶的，要么是受制于对恶者的恐惧。你却自夸你无所畏惧于人，同时你又臆想天主是这样的：由于对黑暗的恐惧，祂竟将善与恶结合起来。不要像你的文字所显示的那样趾高气扬，以为我决心遏止你毒液的流溢，免得其阴险有害的力量造成伤害，是在抬举你；因为宗徒并没有抬举他称之为\Quote{狗}的人，他曾对斐理伯人说：\BibleQuote{你们要提防狗}\BibleRef{斐 3:2}；也没有抬举那些他说\BibleQuote{这些人的言论如同毒癌，愈烂愈大}的人\BibleRef{弟后 2:17}。因此，我因基督之名，要求你，如果你有能力的话，回答那个难倒了你的前任\Person{福图纳图斯}的问题。因为他离开我们辩论的现场时声称，除非与他的同党商议后找到可以反驳我们弟兄所用论证的东西，否则他就不再回来。如果你不打算这样做，那么就离开此地，不要歪曲主的正道，用你的毒液诱捕并感染软弱者的心灵，以免凭主右手相助于我，你以一种出乎你意料的方式遭到挫败。\par

\SourceRecord{Translated by J.G. Cunningham.；Nicene and Post-Nicene Fathers, First Series；Vol. 1.；Edited by Philip Schaff.；Buffalo, NY: Christian Literature Publishing Co.,；1887.；<http://www.newadvent.org/fathers/1102079.htm>.}

% structure: p0001 source=h1 target=section reason=page-title

\section{\Person{热罗尼莫}致\Person{奥斯定}（公元405年）}

\Emph{致\Person{奥斯定}，我真正圣洁的\OriginalTerm{主}{my lord}，至为有福的教父，\Person{热罗尼莫}谨在主内致候。}\par

我急切地向我们的圣洁弟兄\Person{菲尔穆斯}询问您的近况，得知您安好，我甚为欣慰。我本期待他带来，更确切地说，我坚持要他带给我一封您的信；他却告诉我，他动身离开\Place{非洲}时并没有告知您他的打算。因此，我透过这位以特别温暖的感情依恋您的兄弟，向您致以我的敬意；同时，关于我的上一封信，我恳请您原谅我因谦逊，在您长期要求我回信时，竟未能拒绝您。再者，那封信并非我对您的答复，而是把我的论据与您的论据加以对照。如果说，寄出回信是我的过错（我恳请您耐心听我说），那么坚持要我回信之人的过错就更大了。但让我们摒弃这些争论吧；让我们之间保持真诚的兄弟情谊；从此以后，我们彼此往来的书信，不再是争辩，而是出于相互\OriginalTerm{爱德}{charity}的书信。与我一同侍奉主的圣洁弟兄们，向您致以衷心的问候。请代我们问候那些与您一同背负基督轻轭的圣洁弟兄们；尤其我恳请您，向极堪尊敬的圣洁的教父\Person{阿利比乌斯}转致我的敬意。愿基督，我们全能的天主，保护您安康，并记念我，我真正圣洁的\OriginalTerm{主}{my lord}，至为有福的教父。如果您读过我的\WorkTitle{约纳注释}，我想您就不会再提起那可笑的葫芦之争了。再者，如果那位率先挥剑攻击我的朋友，已被我的笔击退，我仰赖您的善意与公正，会将指责归咎于那挑起事端者，而非那予以回击者。如果您愿意，让我们在圣经的园地中锻炼自己，彼此不加伤害。\par

\SourceRecord{Translated by J.G. Cunningham.；Nicene and Post-Nicene Fathers, First Series；Vol. 1.；Edited by Philip Schaff.；Buffalo, NY: Christian Literature Publishing Co.,；1887.；<http://www.newadvent.org/fathers/1102081.htm>.}

% structure: p0001 source=h1 target=section reason=page-title

\section{\Person{奥斯定}致\Person{热罗尼莫}（公元405年）}

\EditorialNote{答复信函72、75、81}\par

\Emph{致\Person{热罗尼莫}，我在基督的爱中深爱和尊敬的、我的圣善弟兄与同为司铎者，\Person{奥斯定}在主内致问候。}\par

% structure: p0004 source=h2 target=subsection reason=relative-heading-rank; base=h2

\subsection{第一章}

1. 很久以前，我曾给您的爱德寄去一封长信，答复您记得由您的圣善子\Person{阿斯特里乌斯}转来的那封信；如今他不仅是我主内的弟兄，也是我的同工。那封回信您是否收到，我不得而知；除非我从您由我们最真诚的朋友\Person{菲尔穆斯}转来的信中推断出来——您说，如果那首先用剑攻击您的人已被您的笔击退，您信赖我的善意与公正，会将责备加于那提出控告者，而非那驳斥控告者。从这极微弱的迹象，我推断您已读过了我的信。在那封信中，我确实表达了痛心，因您与\Person{鲁菲努斯}之间竟兴起了如此大的纷争；昔日的友谊曾那样深厚，凡其名声所及之处，兄弟之爱都为之欢欣。但我并无意以此责备您，我的弟兄，我绝不敢说在此事上发觉您有可责之处。我只是哀叹世人的可悲命运，他们的友谊，既依赖于彼此相敬的延续，便无稳固之时，无论那份相敬有时是何等深厚。然而，我更愿您来信告知，您是否已宽恕我恳求的过犯，我现在盼望您给我更明确的保证；虽然您信中更和蔼愉快的语气似乎表明，我已从我信中的请求得到了赦免——倘若那确是在您读了我的信之后发出的，但正如我所说，这一点并不明确。\par

2. 您问，更确切地说，您是凭着爱德充满信赖的大胆命令，要我们在圣经的园地中彼此取乐，而不彼此伤害。至于我，我完全情愿在这些主题上认真操练，远胜于仅作取乐。然而，如果您选用这个词，是因为它暗示轻松的练习，那么请容我直说，我期望从您那里得到更多；您既拥有伟大才能，又受温和性情的约束，更兼有学问光照的智慧，且专心研究，不受其他事务分心，年复一年，在热情的推动与天赋的引导下奋进：圣神不仅赐予您这一切优长，更明确嘱托您去帮助那些正从事重大而艰难研究的人；他们研究圣经，不是如在平原上取乐，而是像人正喘着气、奋力攀爬陡峭的山崖。不过，或许您选用\Quote{ludamus}（让我们彼此取乐）这个词，是出于爱友之间讨论所宜有的友善温和，无论所辩论的问题浅显容易，还是错综复杂，我恳请您教导我如何能达致这种境地；好使当我对于任何事感到不满时——那事并非因我疏忽，而也许是因我理解迟钝，尚未向我显明——倘若我试图证明相反的意见，而以某种不够谨慎的坦率表达自己时，我不会蒙上幼稚自负与冒昧的嫌疑，好像我企图通过抨击名人来宣扬自己的名声；同样，当论证的严苛要求说出一些刺耳的话时，我若试图用温和有礼的措辞使其不那么难以忍受，我不会被宣判为挥舞着一把\Quote{蜜剑}。依我看，要避免这两种过错，或避免这两者中任何一种的嫌疑，唯一的办法是同意：当我这样与一位比我更有学问的朋友辩论时，我必须赞许他所说的一切，并且，即使是为了更准确地获取信息，也不可在接受他的决定之前有丝毫犹豫。\par

3. 若照这样的条件，我们或许可以在圣经的园地中彼此取乐，而不必担心冒犯对方；但我很可能要奇怪，这取乐是否由我付出代价。因为我向您的爱德承认，我已学会将这种尊崇与敬意只归于圣经的正典书卷：唯独对这些书卷，我最坚定地相信，它们的作者完全免于错误。如果在这些著作中，我被任何看来与真理相悖的事所困惑，我会毫不犹豫地设想，或是抄本有误，或是译者未能领悟所言之义，或是我自己未能理解。至于所有其他著作，在阅读它们时，无论作者们在圣德与学问上比我超越多少，我不会仅仅因为他们持有那意见，就接受他们的教导为真理；我接受，只是因为他们借着这些正典著作本身，或是借着针对我理性提出的论证，成功地使我的判断信服其真理。我的弟兄，我相信这是您我共同的看法。我无需说，我不认为您愿意您的著作被读如先知或宗徒的著作，对于后者，若怀疑它们免于错误，便是错误的。但愿您那谦卑的虔诚和对自己正确的认识远远摒弃这种傲慢，我知道您拥有这样的谦卑和认识，若没有这些，您肯定就不会说：\Quote{但愿我能得到您的拥抱，并能通过交谈在学问上彼此帮助！}\par

% structure: p0008 source=h2 target=subsection reason=relative-heading-rank; base=h2

\subsection{第二章}

4. 我既了解你的生活与品行，就不相信你曾出于虚伪和欺骗而说话；那么，对于宗徒\Person{保禄}，我若相信他论及\Person{伯多禄}和\Person{巴尔纳伯}时所写的并非想一套说一套，岂不更为合理？他说：\BibleQuote{我一见他们的行为与福音的真理不合，就在众人面前对\Person{伯多禄}说：\InnerQuote{你既是犹太人，竟按照外邦人的方式，而不按照犹太人的方式生活，你怎么敢强迫外邦人犹太化呢？}}\BibleRef{迦 2:14} 因为，如果宗徒也欺骗了他自己的\BibleQuote{儿女们}，他曾为这些儿女\BibleQuote{再受产痛，直到基督（即真理）在他们内形成}\BibleRef{迦 4:19}，那么，我还能信赖谁不会在言语或文字上欺骗我呢？他曾事先对他们说：\BibleQuote{我给你们写的这些话，看，我在天主前，绝不说谎}，难道他能对同一些人写出不实之言，并以某种权宜之计认可的欺骗来蒙蔽他们，声称他看见\Person{伯多禄}和\Person{巴尔纳伯}没有按福音的真理正直行事，并为此当面反对了\Person{伯多禄}，因为他强迫外邦人按犹太人的方式生活？\par

5. 但你会说，宁可信宗徒\Person{保禄}写了不实之言，也不信宗徒\Person{伯多禄}行了不义之事。照此原则，我们必将说（但愿我们绝不这样说），宁可信福音史是虚假的，也不信\Person{伯多禄}曾否认基督；\BibleRef{玛 26:75} 且宁可指责\BibleBook{列王}{纪}（即\BibleBook{撒慕尔}{纪下}）记载不实，也不信像\Person{达味}这样伟大的先知、如此特蒙主天主拣选的人，竟因贪恋他人之妻而犯奸淫，并设计杀害其夫，行如此可憎的凶杀。\BibleRef{撒下 11:4} 我更愿意带着确信和对其真理的信服，研读被置于最高（甚至属天）权威之巅的圣经，而不质疑其记述的可信性，从中得知有人受称赞、被纠正或遭责罚；而不愿因担心相信那些虽具备受称扬的美德、却仍是凡人的人有时会做出应受谴责的事，就让我对整个\Quote{\OriginalTerm{天主的圣言}{oracles of God}}的可信性产生怀疑。\par

6. 摩尼教徒主张，圣经的大部分内容——正是这些内容以最明确的措辞驳斥了他们的邪恶谬误——是不值得采信的，因为他们无法曲解经文的语言以支持己见；然而他们却将所谓错误的罪责不归于最初写下这些文字的宗徒，而归咎于一些不为人知的抄本篡改者。但是，由于他们从未成功地通过更多更早的抄本，或借助拉丁译本所据的原文语言来证明这一点，便在辩论场上退却，被众所周知的真理击败而羞愧。难道你圣善的明智没有看出，如果我们不是说（如他们所说）宗徒著作被他人篡改，而是说宗徒自己写了他们明知不实的内容，这将给他们对真理的恶意攻击留下多大的余地？\par

7. 你说，\Person{保禄}竟责备\Person{伯多禄}做了保禄自己也做过的事，这令人难以置信。我此刻并不探究保禄做了什么，而是探究他写了什么。这与我手头的事——即确证交付给我们的圣经的普遍且无可置疑的真理——至为相关；这些经文是由宗徒们，而非无名之辈，交付给我们，以建树我们的信德的，并因此被接纳为权威的\OriginalTerm{正典标准}{canonical standard}。因为，如果伯多禄那次做了他该做的事，那么保禄声称看见他没有按福音的真理正直行事，便是虚假的。因为凡做他该做的事，就是正直行事。因此，明知他人做了该做的事，却说他没有正直行事的人，便是说谎。那么，如果保禄所写的是真的，伯多禄那时确实没有按福音的真理正直行事。他因此做了不该做的事；而如果保禄自己先前也做过类似的事，我宁愿相信，他自己既已受纠正，就不能忽略指出他的一位宗徒弟兄的错误，也不愿相信他在书信中写下任何虚假的话；如果在保禄的任何书信中出现虚假会令人惊讶，那么在这封书信中——他在其序言中说：\BibleQuote{我给你们写的这些话，看，我在天主前，绝不说谎}——又岂非更加令人惊讶呢！\par

8. 就我而论，我相信\Person{伯多禄}在这件事上的作为是要强迫外邦人如同犹太人一样生活：因为我读到\Person{保禄}写了这事，我不信他是在撒谎。因此\Person{伯多禄}行事并不正直。因为这违背了福音的真理，就是那些信了\Person{基督}的人，竟认为没有那些旧时的礼仪，他们就不能得救。这正是那些信了主的割损派在\Place{安提约基雅}所持的立场；\Person{保禄}曾一贯并强烈地反对他们。至于\Person{保禄}给\Person{弟茂德}行割损\BibleRef{宗 16:3}，在\Place{耕格勒}许愿\BibleRef{宗 18:18}，以及在\Place{耶路撒冷}依\Person{雅各伯}的建议，与几个发了愿的人一同遵行了规定的礼仪\BibleRef{宗 21:26}，显然\Person{保禄}在这些事上的用意，并不是要让别人以为，他认为在基督信仰的时代，救恩可由这些仪式获得；而是要避免人们相信，他谴责天主在先前时代所规定的那礼仪——这礼仪是与那时代相宜的，也是未来事物的预像——与异教偶像崇拜无异。因为\Person{雅各伯}正是这样对他说的，即有传闻说他教导人\BibleQuote{背弃\Person{梅瑟}}\BibleRef{宗 21:21}。那些信了\Person{基督}的人，若背弃那位预言了\Person{基督}的人，仿佛憎恶并谴责那位曾说\BibleQuote{如果你们信\Person{梅瑟}，也必信我，因为他是关于我而写的}之人的教训，这无论如何都是错的。\par

9. 我恳请你注意\Person{雅各伯}的话：\BibleQuote{你看见，弟兄，在信主的犹太人中，数以万计，他们都是热爱法律的人；关于你，他们听说你教训在外邦人中所有的犹太人背弃\Person{梅瑟}，说不要给自己的孩子行割损，也不要按习俗生活。那么怎么办呢？众人必要聚集，因为他们听说你来了。你就照我们给你说的去行吧：我们这里有四个人，他们都有誓愿在身；你带他们去，同他们一起行取洁礼，并替他们出费，叫他们剃头；这样，众人便知道，关于你所听到的事，都是虚的，而你却循规蹈矩，遵守法律。至于信主的外邦人，我们已经写信决定，叫他们戒避偶像的玷污、戒避血、戒避窒死之物，并戒避奸淫。}\BibleRef{宗 21:20}依我看来，\Person{雅各伯}之所以提出这个建议，原因非常清楚，就是要让那些犹太人——他们虽然信了\Person{基督}，却热衷于维护法律的荣誉，不愿意让人以为\Person{梅瑟}传给祖宗的制度，被\Person{基督}的教义视为亵渎而加以谴责，好像那些制度原本不是由神圣权威所颁布的那样——知道他们所听到关于\Person{保禄}的传闻是虚假的。因为那些给\Person{保禄}带来这种非难的人，并不是那些理解信主的犹太人在基督信仰时代应如何正确遵守这些礼仪的人——也就是说，他们明白遵守礼仪是为了表明这些仪式的神圣权威，以及它们在先知时代的圣善用途，而不是作为获得救恩的方法，因为救恩已经在\Person{基督}内启示给他们，并由圣洗施行。相反，那些散布反对宗徒的报告的人，乃是主张必须遵守这些礼仪的人，仿佛不遵守这些礼仪，那些信从福音的人便不能得救。因为这些假教师发现\Person{保禄}是一位白白恩宠最热切的宣讲者，也是他们观点最坚决的反对者，他所教导的是，人不是靠这些事成义，而是靠\Person{耶稣基督}的恩宠，这些法律的礼仪正是预定为这恩宠做预像的。因此，这一派的人竭力煽动对\Person{保禄}的憎恨和迫害，指控他反对法律和天主的制度；而要消除这种诬蔑所引起的憎恨，最合适的方法莫过于\Person{保禄}亲自举行那些他被认为谴责为亵渎的礼仪，从而表明，一方面，不应禁止犹太人遵守这些礼仪，仿佛它们是不合法的，另一方面，也不应强迫外邦人遵守它们，仿佛它们是必要的。\par

10. 因为如果他果真如传闻所说的那样谴责这些事，并且他履行这些礼仪是为了掩饰、伪装他的真实感受，那么\Person{雅各伯}就不会对他说：\BibleQuote{众人便知道}，而是说：\Quote{众人便\Emph{以为}那些关于你所听到的事都是虚的}\BibleRef{宗 21:24}；尤其是看到在\Place{耶路撒冷}，宗徒们早已颁布法令，不得强迫外邦人遵守犹太礼仪，但并没有颁布法令，说任何人因此就该阻止犹太人按自己的习俗生活，尽管基督教义也没有将这种义务强加于犹太人。因此，如果\Person{伯多禄}在\Place{安提约基雅}的伪装，是在宗徒的法令颁布之后发生的，而且他这样做是强迫外邦人按犹太人的方式生活——他本人并没有被强迫这样做，尽管他也没有被禁止使用犹太礼仪，以彰显交托给犹太人的\OriginalTerm{天主的谕令}{oracles of God}的荣耀——那么，如果情况是这样，\Person{保禄}劝勉他公开宣告那项法令，就是他和别的宗徒在\Place{耶路撒冷}一起制定的，这有什么奇怪呢？\par

11. 然而，如果如我更倾向于认为的那样，伯多禄是在耶路撒冷那次会议之前这样做的，那么保禄希望他不要怯懦地隐藏，而应大胆地宣布一项实践规则，关于这规则他已知他们二人心意相同，这就不奇怪了；保禄之所以知道这一点，也许是因为他与他商量过他们所传的福音，也许是因为他听说在百夫长科尔乃略蒙召时，伯多禄曾受到天主的指教关于此事，也许是因为他曾看见伯多禄在他所怕得罪的那些人来到安提约基雅之前，与外邦皈依者一同吃饭。因为我们并不否认伯多禄在此问题上已经与保禄本人持相同意见。所以，保禄不是教导伯多禄关于此事的真理，而是责备他的\OriginalTerm{假装}{simulatio}，因为这种假装迫使外邦人照犹太人的方式行事；这样做的唯一原因是，所有这类假装的倾向是要传达或加强一种印象，即那些认为没有\OriginalTerm{割损}{circumcision}和其他礼仪（这些是将来的事的预像）信者便不能得救的人所教导的是真的。\par

12. 为此，他也给弟茂德行了割损礼，免得在犹太人，尤其是他母系的亲属看来，信了基督的外邦人憎恶割损，如同憎恶偶像崇拜一般；而割损是由天主所立，偶像崇拜却由撒殚发明。再者，他没有给弟铎行割损，免得给那些说没有割损信者便不能得救的人留下口实，那些人为了欺骗外邦人，公然声称这就是保禄所持的见解。他本人十分清楚地说：\BibleQuote{但是，同我在一起的弟铎，虽是希腊人，也没有被迫受割损；因为有些潜入的假弟兄，他们暗中来窥探我们在基督耶稣内所享有的自由，为使我们成为奴隶；可是我们没有一刻顺从让步，为使福音的真理在你们中保持不变。}\BibleRef{迦 2:3} 这里我们清楚地看出，他察觉到他们急欲窥探什么，以及为什么他对弟铎没有像对弟茂德那样做，而这原本是他凭那自由可以做的，他藉这自由已经表明，这些规例既非须被要求为得救所必需，也非可被指为不法。\par

13. 可是，你说，在这讨论中我们必须避免与哲学家一同论断，某些人的一些行为处于对错之间的地带，因此既不算善行也不算恶行；并且你在论证中强调，遵守法律礼仪不能是一件\OriginalTerm{中性事物}{adiaphora}，而是非善即恶的；所以，如果我认为它是善的，我就等于承认我们也有义务遵守这些礼仪；但如果我认为它是恶的，我就必须相信宗徒们并非诚心遵守，而是以假装的方式。就我而言，我既不会害怕借用哲学家的权威来为宗徒们辩护，因为毕竟他们在论述中也教导一定程度的真理，也不会害怕模仿法庭上律师的做法，有时为了当事人的利益而牺牲真理。如果像你在\WorkTitle{迦拉达书注释}中所说，律师的这种做法可以被你恰当地引用来比拟和证明伯多禄与保禄的假装，那我为什么要害怕向你引述哲学家的权威呢？我们认为他们的教导毫无价值，并非因为他们所说的一切都是虚假的，而是因为他们大部分事情上都错了，而且即使他们偶尔肯定了真理，他们也与那位是真理的基督的恩宠无分。\par

14. 可是，关于旧约制度中的这些规定，我为什么不能说它们既非善也非恶呢：非善，因为人不能因它们而成义，它们只不过是预像，预告那使我们成义的恩宠；非恶，因为它们是天主按时机为这人民所定的。我为什么不能这样说，既然有先知的话支持，他说道，天主赐给祂的子民\BibleQuote{不美好的律例}？\BibleRef{则 20:25} 或许这正好解释了他为何没有称它们为\Quote{邪恶的}，而是称它们为\Quote{不美好的}，\Emph{即}，它们不是那种能使人成为善的，也不是那种没有它们人就无法成为善的。我愿有幸得知你的至诚之见：若有圣人从东方来到罗马，在每周的第七天守斋，唯独复活节前的星期六例外，他是否算是假装呢？因为若说在第七天守斋是错的，我们不但要定罪罗马教会，也要定罪许多其他教会，包括邻近的和遥远的，在那里这习俗仍被遵守。反之，若说不守第七天的斋是错的，那么我们谴责那么多东方教会，甚至基督教世界的大多数，该是多么狂妄！或者，你宁愿说这行为本身是一件中性之事，但遵守它的人值得称赞，不是作为假装者，而是出于合宜的团结合一及体恤他人感受的愿望？然而，正典圣经并未给基督徒任何关于此事的命令。那么，我更当如何，怎敢将我无法否认是天主所设立的定为恶事！——我之所以承认这一点，乃是凭着同样的信德，我藉这信德知道，并非通过这些规例，而是靠天主的恩宠，通过我们的主耶稣基督，我得以成义。\par

我因此主张，\OriginalTerm{割损}{circumcision}，以及此类其它事物，是藉着所谓\OriginalTerm{旧约}{Old Testament}，以天主的权威赐给犹太人的，作为将来要在基督内实现之事的标记；如今这些事既已实现，关于这些仪式的法律，留给基督徒去阅读，只是为了让他们能理解先前所赐的预言，而不是必须去遵行，仿佛它们所预表的\OriginalTerm{信德}{faith}启示尚属未来一样。虽然如此，这些仪式不应强加给外邦人，但也不应以一种令人以为它们该受憎恶和谴责的方式去禁止犹太人素常的遵行。因此，要藉着正确宣讲基督\OriginalTerm{恩宠}{grace}真理的大能，使所有这些\OriginalTerm{预像}{types}的遵守逐渐消逝；唯独这恩宠，信徒们将受教去认为，他们的\OriginalTerm{成义}{justification}和\OriginalTerm{救恩}{salvation}只应归于它，而不是归于那些预像和\OriginalTerm{阴影}{shadows}——那些当时尚属未来，但如今已新近到来并临在的预像和阴影，正如在我们的\Person{主}亲身降来和宗徒时代，那些犹太人的蒙召之时，他们被发觉仍惯于遵守这些礼仪制度。暂时容忍他们继续遵守这些，就足以表明它们的卓越：它们虽要被放弃，却不像偶像崇拜那样该受憎恶；但也不可强加给别人，免得它们被认为是救恩不可或缺的方法或条件。这就是那些异端者的意见，他们既急于作犹太人，又急于作基督徒，结果两者都作不成。针对这意见，您曾极其仁慈地屈尊警告我，尽管我从未持有这意见。这也是\Person{伯多禄}出于畏惧而陷入假装同意的过失的那个意见，虽然他内心并不同意；为此，\Person{保禄}极其真实地写到他，说他行事不合乎\WorkTitle{福音}的真理，并且极其中肯地说，他在强迫外邦人按犹太人的方式生活。\Person{保禄}并非藉着他在必要时真诚遵守这些礼仪，来将这重担加给外邦人，其用意是要证明这些礼仪不该被彻底弃绝（如同偶像崇拜应被弃绝那样）；因为他不断宣讲：信徒获得救恩不是靠这些事，而是靠向\OriginalTerm{信德}{faith}启示的恩宠，免得引导任何人将遵守这些犹太礼仪视为得救的必要条件。因此，我相信宗徒\Person{保禄}诚诚实实地做了这一切，毫无伪装；然而，现今若有任何人离开犹太教成为基督徒，我既不强逼也不允许他效法\Person{保禄}的榜样，继续真诚地遵守犹太礼仪，正如您——既然认为\Person{保禄}遵行这些礼仪时是在伪装——也不会强逼或允许那样的人在那种伪装中效法宗徒一样。\par

我还要总结一下\Quote{争论的要点，或者更确切地说，您对此的看法}（用您自己的话说）吗？在我看来是这样的：在基督的\WorkTitle{福音}被传扬之后，那些信主的犹太人若是像\Person{保禄}那样献上\OriginalTerm{祭献}{sacrifices}，像\Person{保禄}给\Person{弟茂德}行割损那样给他们的孩子行割损，并且守\Quote{一周的第七日，正如犹太人素常所行的，只要他们做这一切时是伪装者和欺骗者}，那他们做得正确。如果这是您的教义，我们如今不是坠入了\OriginalTerm{厄俾奥尼派}{Ebion}的异端，也不是那些通常被称为\OriginalTerm{纳匝肋派}{Nazarenes}者的异端，或任何其它已知的异端，而是陷入某种新的错误，这错误更为有害，因为它不是出于误解，而是出于故意且精心设计的欺骗企图。如果您为了洗脱怀有这类想法的嫌疑而回答说，宗徒们在这些事例中值得称赞的伪装，是为了避免触犯许多软弱的犹太信徒，因为那些人还不明白这些事应当被摒弃；但如今，当基督\OriginalTerm{恩宠}{grace}的教义已在如此众多的民族中稳固建立，并且藉着在基督的所有教会中诵读\WorkTitle{法律}与\WorkTitle{先知}，人们清楚知道诵读这些不是为了遵行，而是为了教导，那么任何人若是提议假装遵行这些礼仪，都会被视为疯子——那么，我坚称宗徒\Person{保禄}以及其他纯正忠信的基督徒，有责任藉着偶尔尊崇这些古礼来真诚地表明其价值，免得有人以为这些制度——原本充满先知性寓意，且被他们最虔诚的祖先圣善地珍视——会被他们的后代当作魔鬼的亵渎发明而憎恶，对此能有什么反对意见呢？因为如今，当信仰已经来到——这信仰曾被这些礼仪所预表，并在\Person{主}的死亡与复活之后被启示出来——就它们的职任而言，它们已经失效了。然而，正如将逝者的遗体由亲属恭敬地抬到坟墓是合宜的，同样，这些礼仪也宜以一种与其起源和历史相称的方式被除去，这不是出于假装的敬意，而是作为一种宗教义务，而不是被立刻抛弃，或被丢出去，任仇敌的辱骂如同狗牙一般将其撕碎。再进一步引申这个比喻：如今若有任何基督徒（尽管他可能是从犹太教皈依的）打算效法宗徒们遵行这些礼仪，就如同那打扰安息者骨灰的人一样，他不是在虔诚地履行葬礼中的职责，而是在亵渎地侵犯坟墓。\par

我承认，我在那封信中的陈述——即\Person{保禄}（他虽是基督的宗徒）之所以履行某些礼仪，是为了显示这些礼仪对那些愿意继续持守他们从祖先所领受法律之礼的人并非有害——我没有足够明确地限定这一陈述，即加上这仅是\Emph{在信德的恩宠初次揭示的那一时代}才如此；因为在那时，这并不有害。这些礼仪应由所有基督徒随着时间的推移逐步放弃，而不是突然全部放弃；以免如果这样做，人们就无法察觉天主藉\Person{梅瑟}为祂的古代子民所规定的礼仪，与不洁的神教导人们在异教殿宇中所行的仪式之间的区别。因此，我承认，在这一点上你的指责是合理的，我的疏漏理应受责备。然而，早在我收到你的信之前，当我撰写\WorkTitle{驳摩尼教徒福斯图斯}论著时，我在对该相同段落的简短阐释中没有忽略加入我刚刚陈述的限定条件，如你肯屈尊阅读那部论著，你自己便可看到；或者你可以通过此信的送信人以你喜欢的任何其他方式确认，我早已公开发表了这一普遍陈述的限制。现在，我犹如在天主面前说话，因爱德的定律恳求你，相信我全心所说的话：我从未认为在我们这时代，成为基督徒的犹太人应当或以任何方式、出于任何动机自由地遵守旧约时代的礼仪；尽管我自开始阅读宗徒\Person{保禄}的著作以来，一直对他持有你质疑的观点。你也不会认为这两件事不相容；因为你不认为在我们今天任何人有义务假意遵守这些犹太礼仪，尽管你相信宗徒们曾这样做。\par

相应地，正如你在反驳我时所肯定的，并且引用你的原话，\Quote{即使全世界抗议，仍要大胆宣布：犹太礼仪对基督徒既是有害的，又是致命的；凡遵守这些礼仪的，不论原本是犹太人还是外邦人，都正在走向地狱的深渊，}我完全赞同这一陈述，并补充说：\Quote{凡遵守这些礼仪的，不论原本是犹太人还是外邦人，都正在走向地狱的深渊，不仅真诚地遵守，而且伪装地遵守也是如此。}你还要什么呢？但正如你区分了你认为宗徒们所实践的伪装与适合现今时代的行事准则，我也同样区分了\Person{保禄}——依我看来——当时在所有这些例子中真诚遵循的路线，与在我们今日遵守这些犹太礼仪之事，尽管这事或许如他所做，毫无伪装，因为那时是应予赞同的，如今却应憎恶。因此，尽管我们读到\BibleQuote{法律和先知到若翰为止}\BibleRef{路 16:16}，且读到\BibleQuote{为此犹太人越想要杀害祂，因为祂不但犯了安息日，又称天主是自己的父，使自己与天主平等}\BibleRef{若 5:18}，又读到\BibleQuote{从他的满盈中，我们都领受了恩宠，而且恩宠上加恩宠，因为法律是藉\Person{梅瑟}传授的，恩宠和真理却是由\Person{耶稣基督}而来的}\BibleRef{若 1:16}；并且尽管\Person{耶肋米亚}曾预言，天主将与犹大家订立新约，不像与他们的祖先订立的盟约\BibleRef{耶 31:31}；然而我不认为我们的主受割损是伪装之举。如果有人反对说，祂之所以没有禁止这事是因为祂只是一个婴儿，我则接着说，我不认为祂对癞病人说\BibleQuote{但洁净了的，就去献上\Person{梅瑟}所吩咐的礼物，给他们当作证据}\BibleRef{谷 1:44}是出于骗人的意图——从而祂将自己的命令加于\Person{梅瑟}法律中关于那礼仪的权威之上。祂上\Place{耶路撒冷}过节\BibleRef{若 7:10}也不是伪装，同样祂也不想引人注意；因为祂是暗中去的，并非公开。\par

19. 然而，宗徒自己的话可能会被用来反对我：\BibleQuote{请注意，我\Person{保禄}告诉你们：如果你们受割损，基督对你们毫无益处。} \BibleRef{迦 5:2} 据此推论，他欺骗了\Person{弟茂德}，并使基督对他毫无益处，因为他给\Person{弟茂德}行割损。你会回答说，这次割损并未伤害\Person{弟茂德}，因为它是出于欺骗的意图？我答复说，宗徒并未提出任何这种例外。他没有说，如果你不受伪装地受割损，就如他没有说，如果你伪装地受割损。他毫无保留地说：\BibleQuote{如果你们受割损，基督对你们毫无益处。} 因此，既然你坚持要为你的解释找到空间，提议补充\Quote{除非是作为伪装的行为}这些词，我并非不近情理地要求你允许我理解\Quote{如果你们受割损}这句话，在那段经文中是针对那些要求割损的人说的，因为他们认为，除非如此，他们便不能藉着基督得救。凡因这种信念和意愿，并怀着这种目的而受割损的人，基督确实对他毫无益处，正如宗徒在别处明确肯定：\BibleQuote{如果成义是藉着法律而来，基督就白白死了。} \BibleRef{迦 2:21} 你引用的那些话也肯定了同样的事：\BibleQuote{你们这些要靠法律成义的人，已与基督断绝了关系，由恩宠上跌了下来。} \BibleRef{迦 5:4} 因此，他的谴责是针对那些相信自己要靠法律成义的人——而不是针对那些深知法律仪式作为预表真理而被设立的目的，以及它们预定生效的时期，为了尊崇那起初设立这些仪式的主而遵守它们的人。因此他在别处也说：\BibleQuote{如果你们受圣神引导，就不在法律之下，} \BibleRef{迦 5:18} ——从这节经文你推断出，显然 \Quote{凡服从法律的，便没有圣神，并非如我们的教父宣称宗徒们所做的那样，出于明智的审慎而伪装服从，而是}——正如我所设想的情形——\Quote{真诚地服从。}\par

20. 在我看来，明确宗徒在此所谴责的对法律的服从究竟指什么，至关重要；因为我不认为他这里仅仅指割损，或我们的祖先当日所献、但基督徒如今不再奉献的祭献，以及其他同类的礼仪。我更认为，他也包括了那条法律诫命：\BibleQuote{不可贪婪}，我们承认基督徒无疑有义务遵守这诫命，并且我们发现它在福音所投射的光照下被最充分地宣示出来。他说：\BibleQuote{法律是圣的，诫命也是圣的、正义和美善的；}随后又说：\BibleQuote{那么，这美善的，竟成了我的死亡吗？断然不是！}\BibleQuote{但罪藉着那美善的，在我内造成了死亡，好显示出它真是罪，并可使罪因着诫命，变得极端邪恶。} \BibleRef{罗 7:13} 正如他在这里说：\BibleQuote{可使罪因着诫命，变得极端邪恶，}同样在别处说：\BibleQuote{法律渗入了，为使过犯增多；但罪恶在哪里增多，恩宠在那里也格外丰富。} \BibleRef{罗 5:20} 又在另一处，论及恩宠的安排，肯定唯独恩宠使人成义之后，他问道：\BibleQuote{那么，法律是为什么？}并即刻回答：\BibleQuote{它因过犯的缘故而深入，直到蒙受应许的后裔来到。} \BibleRef{迦 3:19} 因此，宗徒在此宣称其服从法律成为自身定罪原因的人，乃是那些被法律定为有罪者，因为他们不履行法律的要求，并且不明白白白恩宠的功效，反而自负地依赖自己的力量去遵行天主的法律；因为\BibleQuote{爱就是法律的满全。} \BibleRef{罗 13:10} 然而，\BibleQuote{天主的爱，藉着赐给我们的圣神，已倾注在我们心中了，}不是靠我们自己的力量，而是\BibleQuote{藉着赐给我们的圣神。} \BibleRef{罗 5:5} 不过，要圆满地讨论这一点，如果不另写一卷书，至少也需要过长的离题。那么，如果那条法律诫命，\BibleQuote{不可贪婪}，将那些人性软弱未获天主恩宠扶助的人定为有罪，不但不开释罪人，反而判定他为犯法者，那么，那些仅仅是预象的礼仪，如割损等，它们预定在恩宠的启示更广为人知时被废除，又如何能成为使任何人成义的方法呢！尽管如此，这些礼仪并不因此就应如同外教人那魔鬼般的邪恶一样，从恩宠启示开始（这恩宠原由这些影像所预表）便立即被厌恶地摒弃；而应在一小段时间内容忍，尤其是在那些加入基督教会的、来自曾受这些礼仪的民族的人中。然而，一旦它们仿佛被光荣地安葬之后，从此便当由一切基督徒彻底弃绝。\par

21. 现在，关于你使用的言辞——\Quote{非出于权宜，如我们前辈所愿}，——\Quote{非以权宜之计为正当，此即我们教父们原用以解释宗徒行事之依据}，——请问这些话是什么意思？这无疑就是我所称的\OriginalTerm{义务性的谎言}{officiosum mendacium}，即权宜所容许的自由相当于以虔敬的意向说出虚谎的义务。至少我看不出其他解释，除非仅仅加上\Quote{权宜所容许}这几个字就足以使谎言不再是谎言；若这是荒谬的，你何不公开说，以义务为理由说的谎言是可以辩护的呢？或许\OriginalTerm{义务}{officium}一词冒犯了你，因这词在教会书籍中不常见；但这并未阻止我们的圣\Person{安博罗削}使用它，因为他为自己的某些充满有益道理的书籍选用了\WorkTitle{论义务}这一书名。你的意思是说，凡出于义务感而说谎的应受指责，而基于权宜说同样谎言的应被赞许吗？我恳求你思想，凡持此见解的人，可能在他认为合适时便说谎，因为这涉及一个全然重要的问题：说虚谎的话是否在任何时候都可算为善人的义务，尤其是一个基督徒的义务，因为对这样的人曾说，\BibleQuote{你们的话，是就说是，非就说非，免得你们遭受审判。}\BibleRef{雅 5:12} 他也相信\OriginalTerm{圣咏作者}{Psalmist}的话，\BibleQuote{你必消灭一切说谎言的人。}\par

22. 然而，如我所言，这是另一个重大的问题；我让持此见解的人自己去判断他在何种情况下有自由说谎：但有一件事必须最坚定地相信并坚持，即这种说谎方式与那些被用来撰写神圣著作，特别是正经圣经的作者们相去甚远；免得基督的管家们，圣经论到他们说，\BibleQuote{要求于管家的是要人表现忠信。}\BibleRef{格前 4:2}让人以为他们是以学会在必要时为真理的缘故说虚谎为重要的功课，来证明自己的忠信，尽管在拉丁语中，\Quote{忠信}一词原意即言行之间的真实相符。然而，凡所言确实实行的，那里就绝无虚谎的余地。因此，保禄无疑作为\Quote{忠信的管家}应被视为在其著作中证实了自己的忠信；因为他是真理的管家，而非虚谎的管家。所以，当他写道，他看见伯多禄没有按福音的真理正直行事，而且因他强迫外邦人按犹太人的方式生活，便当面反对他时，他所写的是真理。伯多禄自己则以相称于他的圣洁而仁爱的谦逊，领受了保禄为维护真理并以爱的勇敢所施加的责备。伯多禄借此给后来的人留下榜样：倘若他们随时偏离正道，也不应以为接受年轻于己者的纠正乃是降格——这个榜样比保禄在同一事件上留给我们的更罕见，也需要更大的虔诚：即较年轻的人应有勇气甚至违抗他们的长者，若为福音真理的辩护所需，但仍要以保持兄弟友爱不受损害的方式进行。因为虽然圆满地持守正道更为可取，但温良地接受劝戒，比大胆地规劝有过者，远远更值得赞赏和称许。因此，在那件事上，保禄因正直的勇气应受赞扬，伯多禄因圣善的谦逊应受赞扬；就我的判断所能形成的见解而言，这一点本应被用来回答\Person{波尔菲里}的诽谤，而不应给他进一步挑剔的机会，让他能够对基督徒提出更严重的指控，即他们在撰写书信或遵从\OriginalTerm{天主的诫命}{the ordinances of God}时行了欺骗。\par

% structure: p0028 source=h2 target=subsection reason=relative-heading-rank; base=h2

\subsection{第三章}

23. 你要求我至少举出一个人，在这事上我所遵从的是他的意见，而同时你却引用了许多在你之前持守你所辩护之意见者的名字。你又说，如果我因你在这事上的错误而谴责你，你恳求允许你与这些伟人一同留在错误中。我并未读过他们的著作；然而，虽然他们总共不过六、七人，你自己却已驳斥了其中四人的权威。因为关于那位未提名字的\Place{劳狄刻亚}的作者，你说他近来已经离弃了教会；你描述\Person{亚历山大}为旧时的\OriginalTerm{异端者}{heretic}；至于\Person{奥力振}和\Person{狄迪慕斯}，我在你某些更近期的著作中读到，你对他们的见解给予了严厉的批评，而且并非针对不重要的问题，尽管你从前曾给奥力振极高的赞扬。因此我料想，你自己也不会甘心与这些人一同处于错误中；虽然我所指的那些话语，等于是断言在这件事上他们并未犯错。因为谁愿意明知故错，不论与他同错的人何等之多呢？因此，七人中只剩下三位：\Person{埃梅萨的优西比乌}、\Person{黑拉克莱亚的狄奥多若}，以及你后来提到的从前在\Place{君士坦丁堡}教会作过主教的\Person{若望}。\par

24. 然而，如果你询问或回忆起我们的\Person{盎博罗削}和我们的\Person{西彼廉}在这个问题上的意见，或许你会发现，我并非没有遵循一些人的脚步，在我所坚持的观点上。同时，正如我已经说过的，我只对\OriginalTerm{正典圣经}{canonical Scriptures}负有完全服从的义务，要遵循它们的教导，不容有丝毫怀疑其中可能存在任何错误或有意误导的言论。因此，当我环顾四周，想找出第三个名字，以便能用我方三个人名来对抗你的三个时，我确实可以轻易找到一个，我想，如果我的阅读范围广泛的话；但我想起了一个人，他的名字与所有这些人一样好，不，更具有权威——我指的是宗徒\Person{保禄}本人。我投奔他；我向他本人申诉，反对所有对他的著作提出与我不同意见的注释者的判断。我质问他，要求知道他是否写了真实的事，还是出于方便的借口，写了他明知为假的事，当他写道他看见\Person{伯多禄}没有按福音的真理正直行事，就当面抵挡他，因为他的假装迫使外邦人按犹太人的方式生活。我听到他在书信的前面部分庄严地宣告，他叙述这件事时，\BibleQuote{我给你们写的都是真的，我在天主前作证，我没有撒谎。}\BibleRef{迦 1:20}\par

25. 请允许我与那些持不同意见者，不论他们的名声多大，有所分歧。这位如此伟大的宗徒，在他自己的书信中使用誓言来证实其真实性，这对我而言比任何学者在讨论他人著作时提出的意见更有分量。我也不害怕人们会说，在为宗徒\Person{保禄}辩护，反驳他假装随从犹太偏见的错误这一指控时，我反而坚称他确实随从了。因为一方面，当他以宗徒的自由（这种自由适合那个过渡时期），通过实践那些古老的神圣礼仪（当有必要说明它们最初的美善，因为它们不是由\Person{撒殚}的诡计所设立来欺骗人，而是由天主的智慧所设立，为预表将来的事），他这样做时，他并没有假装随从错误；另一方面，他也没有真正随从犹太教的错误，因为他不仅知道，而且不断并强烈地宣讲，那些认为这些礼仪应强加给外邦归依者，或认为相信者得因这些礼仪成义的人，是错误的。\par

26. 此外，至于我说他对于犹太人，就当作犹太人，对于外邦人，就当作外邦人，不是出于蓄意欺骗的狡诈，而是出于怜悯之爱的同情，在我看来，你没有充分思考我这些话的意思；或者，也许是我没有把它表达清楚。因为我这样说，并不是我认为他在任何情况下都实行了伪装的随从；我这样说，是因为他的随从是真诚的，他在成为犹太人中的犹太人方面，不亚于他在成为外邦人中的外邦人方面——这种类比是你自己提出的，我感谢地承认，你由此极大地帮助了我的论证。因为当我在信中请你解释，怎么能认为宗徒\Person{保禄}对于犹太人成为犹太人，就意味着他在遵守犹太规矩时必定行动诡诈，而他在对于外邦人成为外邦人时，却不涉及这样诡诈地随从异教习俗；你的回答是，他对于外邦人成为外邦人，无非是指他接纳未受割损的人，并容许自由取用那些被犹太法律定为不洁的食物。那么，如果我问，他在这件事上是否也行了伪装，这个想法会被视为显然荒谬和虚假而遭到摒弃：那么明确的推断是，在他实行那些使他对于犹太人成为犹太人的事上，他运用了明智的自由，而不是屈服于卑劣的强迫，也没有做更与他身份不符的事，即为了权宜之计而从正直堕落为欺诈。\par

27. 因为对于信徒，以及那些认识真理的人，一如宗徒所证实的（除非在这里他也可能欺骗他的读者），\BibleQuote{因为天主所造的样样都好，如以感恩的心领受，没有一样是可摈弃的；}\BibleRef{弟前 4:4} 因此，对\Person{保禄}本人而言，不仅作为一个人，更作为一位极其忠信的管家，不仅作为认识真理者，也作为真理的导师，凡用作食物的天主受造物，并非假装的，而是真正美好的。那么，如果他对\Place{外邦人}成了外邦人，是通过宣讲和教导关于食物和割损的真理，尽管他并未佯装遵从外邦人的礼仪和规条，为什么说他不可能对犹太人成为犹太人，除非他以实行他们宗教的仪式来施行掩饰呢？为什么他对那被接上的野橄榄枝，保持了一个管家的真正忠信，却在对那些橄榄树本来的、原有的枝子的人面前，以权宜之计为借口，举起了那陌生的掩饰之幕呢？为什么他对外邦人成为外邦人时，他的教导和行为与他的真实情感一致；但对犹太人成为犹太人时，他却将一件事藏在心里，而在言语、行为和著作中说出完全不同的东西呢？但我们对宗徒决不可怀有这样的想法。他向犹太人和外邦人都欠着\BibleQuote{这训令的目的就是爱德，这爱德出于纯洁的心、光明磊落的良心和真诚的信仰；}\BibleRef{弟前 1:5} 因此他向什么样的人，就成了什么样的人，为要赢得所有的人，\BibleRef{格前 9:22} 不是用骗子的诡计，而是用充满怜悯者的爱心；这就是说，不是藉着假装去做其他一切人所做的恶事，而是尽心竭力，以全部怜悯之心施予他人因所受的恶疾而需的救治，仿佛那些恶疾就是他自己的一般。\par

28. 因此，当他并不拒绝实行某些旧约的礼仪时，并非出于对犹太人的怜悯而佯装遵从，而毫无疑问是真诚地实行；并且藉着这一行动，表明他对那些在先前安排中曾一度由天主所命令的事物的尊重，他将这些事物与异教的不虔敬礼仪区分开来。此外，在那时，不是用骗子的诡计，而是用被怜悯所感动的爱心，他向犹太人成为犹太人，当他看到他们陷于错误之中，这错误要么使他们不愿相信\Person{基督}，要么使他们认为藉着这些古老的祭献和礼仪规矩，就能洗净罪恶并得享救恩；他渴望将他们从那错误中解救出来，仿佛看到的不是他们，而是他自己陷于其中；如此真心爱近人如己，并且以己之所欲，施之于人，倘若自己需要他人的帮助——对于这一责任，我们的主曾加上说：\BibleQuote{法律和先知即在于此。}\BibleRef{玛 7:12}\par

29. \Person{保禄}在\BibleBook{迦拉达}{书}中推荐了这种怜悯之情，说：\BibleQuote{弟兄们，如果见一个人陷于某种过犯，你们既是属神的人，就该以柔和的心神矫正他；但你们自己要小心，免得也受试探。}\BibleRef{迦 6:1} 请看，他是否没有说：\Quote{你就和他一样，好能赢得他。} 当然，不是说要犯下或假装犯下那被过犯所胜之人同样的过失，而是在他人的过失中，他应考虑自己可能会遇到什么事，并如此以怜悯之心帮助那人，如同他愿意那人在同样情况下帮助自己一样；即，不是用骗子的诡计，而是用充满怜悯者的爱心。因此，无论\Person{保禄}在犹太人或外邦人，或任何人身上发现什么错误或过失，他向什么样的人，就成了什么样的人——不是通过假装不真实的事，而是通过感受，因为这事本可能临于自身，怀着设身处地的怜悯之情——为要赢得所有的人。\par

% structure: p0036 source=h2 target=subsection reason=relative-heading-rank; base=h2

\subsection{第四章}

30. 我恳求你，如果愿意，请稍微省察一下你自己的内心——我指的是你对我所怀有的内心——并回忆，或如果你手边有那封信的副本，重新阅读，你通过我们的弟兄\Person{西彼连}，如今是我的同工，带给我的那封（未免太过简短的）信中的话语。请读一读你是以怎样真诚的兄弟情谊和热切爱心，在对我所做的事提出严重抱怨之后，加上了这些话：\Quote{在这事上，友谊受了伤，弟兄联合的法则已被轻视。不要让世人看到我们像小孩子一样争吵，并为那些可能成为我们各自支持者或反对者的人提供愤怒争斗的材料。} 我看出这些话是你从心底说出的，出自一颗善意想给我忠告的心。然后你又加上，即使你不说我也明白的话：\Quote{我写下现在所写的，因为我渴望对你保持纯洁的基督徒之爱，而不愿在心里隐藏任何与口中之言不相符的事。} 啊，虔敬的人，我所亲爱的，洞察我灵魂的天主作证，我真心相信你的话；正如我不怀疑你在信中向我作出的这番表白的真诚，我也完全相信宗徒\Person{保禄}在书信中所作的完全相同的表白——这信不是写给某一个体，而是写给犹太人和希腊人，以及所有那些在福音中作他儿女的外邦人——他为了他们在灵性上的诞生而受产痛，之后又为了那成千上万信仰\Person{基督}的信徒，这书信因他们的缘故而被保存下来。我说，我相信，他并未\Quote{在心里隐藏任何与口中之言不相符的事。}\par

31. 你对我所做的正是这事——你成为我这样的人——\Quote{不是出于诡诈的狡猾，而是出于怜悯的爱，}当你认为，你有责任尽力避免我陷入你所以为的错误中，正如如果情况发生在你自己身上，你希望别人为你解救时那样。因此，我感激地承认你对我善意的这一证据，我也主张，如果我对你著作中的某些内容感到不安，并向你倾诉我的苦恼时，你也同样不要对我不满；因为我希望所有人都能以我对你所采取的态度来对待我，即，如果他们认为我的著作中有值得谴责之处，他们既不应以虚假的赞扬来奉承我，也不应在别人面前指责那些他们对我本人保持沉默的事；在我看来，这样做会更严重地\Quote{伤害友谊，轻视兄弟团结的法则。}因为，对于那些友谊，我不敢称之为基督徒的友谊，其中流行的谚语\Quote{谄媚结交朋友，真理招致敌人}比圣经中的谚语更受重视，圣经说：\BibleQuote{朋友的创伤是忠诚，敌人的亲吻是欺骗。} \BibleRef{箴 27:6}\par

32. 因此，我们更应尽力在我们挚爱的朋友们面前立下榜样，他们衷心希望我们工作顺利，使他们知道，最亲密的朋友之间也可能意见相左，以至于一方的观点为另一方所反对，却不会减少彼此的友爱，也不会因真诚友谊所应有的真理而点燃仇恨，不论这反对本身是否符合真理，或者至少，无论其内在价值如何，都是由一颗决心不\Quote{在心里隐藏任何与口中所言不一致的事}的真诚之心说出的。因此，愿我们的弟兄，你的朋友们——你见证他们是基督的器皿——相信我，我那封信在到达你手中之前，就已落在许多别人手里，这完全违背我的意愿，我因这一失误而心中充满不小的忧伤。这事的经过说来话长，而且如果我未弄错的话，现在也无此必要；因为，如果关于此事我的话尚可采信，我只消说，这事并非出于某些人猜想的恶意意图，也非出于我的愿望、安排、同意或图谋。如果他们对我在天主面前所宣示的不信，我也无法再令他们满意。然而，我绝不相信他们向你的圣德提出这种暗示，是出于恶意，企图在你我之间燃起敌意，愿天主以祂的仁慈保护我们免陷于此！无疑地，他们并无故意伤害我的意图，只是轻易怀疑我，作为一个人，可能犯下人性共有的缺失。因为如果他们是基督的器皿，被立定不是为羞辱，而是为尊荣，并由天主装备，适于在祂的大户人家中行各种善工，那么我理当如此相信他们。\BibleRef{弟后 2:20}然而，倘若我这庄重的声明为他们所知，而他们仍然持守对我行为的同样看法，你自己将看出他们这样做是错了。\par

33. 至于我曾写道，我从未寄给罗马一本反对你的书，我之所以这样写，首先是因为我不认为\Quote{书}这个名称适用于我的信函，因此我以为你听说的是一件与此截然不同的事；其次，我并未将所说的那封信寄往罗马，而是寄给了你；再者，我并不认为它是反对你的，因为我知道我是出于友谊的真诚，这友谊应当给予我们彼此自由地提出建议和纠正。暂且略去我方才提到的你的朋友们，我恳求你，因着我们借以蒙救赎的恩宠，不要以为，在我信函中论及因主的仁慈而赐予你的那些美好恩赐时，我犯了任何巧言奉承的过错。然而，如果我有什么地方得罪了你，请原谅我。至于那位已被人遗忘的诗人生平中的那件事，我或许以过于炫学而非高雅的方式从古典文学中引用而来，我恳请你不要将其应用到你身上，超出我所允许的范围；因为我紧接着补充说：\Quote{我这样说，并非为了让你恢复属灵的视力——我绝无冒失说你已失去它！——而是为了，你既有清晰敏锐的眼目，可将它们转向这件事。}我以为引用那件事，完全只与\OriginalTerm{翻案诗}{\GreekText{παλινῳδία}}有关，在那件事上，我们都应效法\Person{斯特西科鲁斯}，如果我们写了什么，便有责任在日后的著作中纠正它；而绝非与斯特西科鲁斯的失明有关，我既未将其归咎于你的心智，也不担心你会遭遇它。我再次恳求你，凡你认为我的著作中有必要谴责之处，就大胆地纠正。因为，虽然在教会中盛行的尊荣名衔上，主教的位分高于司铎，但\Person{奥古斯丁}在许多事上不及\Person{热罗尼莫}；尽管即使纠正来自一个在各方面可能都不如自己的人，也不应拒绝或轻视。\par

% structure: p0041 source=h2 target=subsection reason=relative-heading-rank; base=h2

\subsection{第五章}

34. 关于你的翻译工作，你现在已使我信服，你提议从\OriginalTerm{希伯来原文}{Hebrew}翻译\WorkTitle{圣经}，以便揭露那些被犹太人遗漏或歪曲之处，这一计划确有益处。但我恳请你说明，这是由哪些犹太人所为：是那些在主降临前翻译\WorkTitle{旧约}的犹太人——若是如此，是他们中的哪一个或哪几个——还是后来的犹太人，他们可能被认为为了无法回答这些抄本所提供的有关\OriginalTerm{基督信仰}{Christian faith}的证据，而残害或败坏了\OriginalTerm{希腊文抄本}{Greek manuscripts}？我找不到任何理由，能促使早期的犹太译者们如此不忠。此外，我也恳请你，将你的\OriginalTerm{七十贤士译本}{Septuagint}寄给我们，我此前不知你已出版。我也切望拜读你提到的那本名为\WorkTitle{\OriginalTerm{论最好的翻译方式}{De optimo genere interpretandi}}的书，并从中了解，如何权衡译者对原文语言的熟识所产生的成果，与那些称职的\WorkTitle{圣经}评注家们的推测——这些评注家虽然共同忠于同一真信德，但由于许多经文晦涩难解，常常不得不提出各种不同见解；然而这种解释上的差异绝不意味着背离信德的统一；正如一位评注家本人，可以持守他所拥有的信德，对同一段经文给出两种不同的解释，因为经文的晦涩使两者都同样可以接受。\par

35. 此外，我也切望你的\OriginalTerm{七十贤士译本}{Septuagint}，为使我们尽可能摆脱那些不论称职与否、试图从事拉丁翻译之人显著无能所带来的后果；也为了使那些以为我对你的辛劳工作怀有嫉妒的人，若有可能，终能理解：我之所以反对在我们的教会中公开宣读你译自\OriginalTerm{希伯来文}{Hebrew}的译本，唯一的理由是，生怕提出任何新的、有违\OriginalTerm{七十贤士译本}{Septuagint}权威之处，会给基督的羊群造成严重的困扰——这些羊群的耳朵和心灵，早已习惯聆听那个由宗徒们亲自给予认可印鉴的译本。因此，关于\BibleBook{约纳}{先知书}中的那棵植物，如果在希伯来文中既非\Quote{\OriginalTerm{蓖麻}{gourd}}，也非\Quote{\OriginalTerm{常春藤}{ivy}}，而是某种挺直而立、靠自己的茎支撑、无需其他扶持的植物，我更愿在我们所有的拉丁译本中沿用\Quote{\OriginalTerm{蓖麻}{gourd}}；因为我不认为\OriginalTerm{七十贤士}{Seventy}会随意这样翻译，除非他们知道那植物与蓖麻相似。\par

36. 我想我现已对你的三封信给出了足够的答复（或许过于足够）；其中两封由\Person{西彼廉}带来，一封由\Person{菲尔穆斯}带来。在回信中，请将你认为有助于教导我及他人的一切内容寄来。我也将更加留意，在主的助佑下，凡我写给你的信，必先抵达你本人，而不致落入他人之手，导致内容外泄；因为我承认，我不愿你的任何来信遭遇到你合理抱怨我的去信所遭遇的那般命运。然而，让我们决心在我们之间保持朋友的自由与友爱；如此，在我们往来的书信中，任何一方都可不加拘束地向对方坦诚指出自己认为应予纠正之处，只要这总是本着不违背兄弟友爱、不致招惹天主不悦的精神去做。然而，如果你认为在我们之间这样做，难免会危及那兄弟之爱，那么我们就不这样做：因为，我所希望在我们之间保持的爱，固然是那能使这种相互自由成为可能的更伟大的爱；但若有所不及，也总比完全没有好。\par

\SourceRecord{Translated by J.G. Cunningham.；Nicene and Post-Nicene Fathers, First Series；Vol. 1.；Edited by Philip Schaff.；Buffalo, NY: Christian Literature Publishing Co.,；1887.；<http://www.newadvent.org/fathers/1102082.htm>.}

% structure: p0001 source=h1 target=section reason=page-title

\section{书信 83（公元405年）}

\Emph{致我至福的主\Person{Alypius}，我最亲爱的兄弟与同工，可亲可爱的；并致真诚的敬意，向与他同在的弟兄们，\Person{Augustine}及与他同在的弟兄们在主内问候。}\par

1. \Place{Thiave}教会众成员的忧伤使我心无法安息，直到我听闻他们恢复了对你们往日的心意；此事必须尽快完成。因为宗徒曾为一个人担心，\BibleQuote{免得这样的人因过度忧伤而沉溺}，在同一上下文中又说，\BibleQuote{免得撒殚乘机胜过我们，因为我们并非不知道他的诡计}\BibleRef{格后 2:7}。何况我们更应谨慎从事，免得让整个羊群，尤其是新近与天主教会和好的羊群，遭受同样的忧伤，而我绝不能离弃他们！然而，短暂的时间不容我们充分商议以达成成熟满意的决定，愿你的尊威在此信中接受我自离别后长久深思而最赞成的意见；若你赞同，请立即将我这以你我共同名义写给他们的附函送交他们。\par

2. 你提议他们应得一半（\Person{Honoratus}留下的财产），另一半由我从可支配的资源中补足给他们。然而我认为，如果全部财产被取走，人们或许有理由说我们在这事上费心尽力是为了正义，而非为金钱利益。但当我们给他们一半，以妥协方式解决时，便显明我们只关心金钱；你明白由此而来的损害。因为一方面，他们会认为我们取走了一半原无权要求的财产；另一方面，我们会认为他们不正当地、不义地同意接受本应全归穷人的一半财产的援助。因为你所说的，\Quote{我们必须谨慎，免得在试图公正解决难题时，造成更严重的伤害}，若只给他们一半，这话同样有力。因为那些我们希望确保他们弃俗修道的人，会为了这一半私人产业，设法找借口推迟出售，以便按此先例处理他们的事。而且，在这类双方都有话可说的问题上，若整个团体因我们不避恶表，因认为他们高度敬重的主教们被卑鄙的贪婪所击中而跌倒，岂不奇怪？\par

3. 因为任何人归向隐修士生活时，若真心如此，便不会思想我刚才所提的事，尤其若他受劝戒明白这种行为的罪恶性。但若他是骗子，寻求\BibleQuote{自己的事，而不是耶稣基督的事}\BibleRef{斐 2:21}，便没有爱德；没有爱德，即便\BibleQuote{将全部财产施舍给穷人，且舍身被焚烧}\BibleRef{格前 13:3}，为他有何益处？再者，如我们交谈时所同意的，今后可避免此事，而与每位愿意进隐修院的人作出安排：除非他先解除所有这类缠累，轻松自如，因其产业已不再属他，否则不得加入弟兄团体。因为除了清楚表明我们在这些事上绝非关心金钱以外，别无他法能避免软弱弟兄的这种灵性死亡，并除去我们所切愿领回与天主教会和好之人得救的严重障碍。而唯有让他们使用那一直以为属于已故司铎的产业，他们才能明白这一点；因为即使产业不属于他，他们也本应从起初便知道这一点。\par

4. 因此，在我看来，对于这类事情，应遵循的规则是：按照关于财产的普通民法，凡属于在任何地方受祝圣的圣职人员的财产，在他死后，应归于他受祝圣所归属的那个教会。而根据民法，该财产属于司铎\Person{Honoratus}；所以，不仅因为他是在别处受祝圣，而且即使他留在\Place{Thagaste}的隐修院，如果他在未出售其地产或通过明示的赠与契约将其移交给任何人之前去世，继承权也仅属于他的继承人：就如\Person{Æmilianus}弟兄继承了\Person{Privatus}弟兄留下的那三十个银币一样。因此，这事本应及时考虑和安排；但如果没有作出安排，在处理此地产时，我们必须遵守那些为规范公民社会中财产的持有或放弃而制定的法律；好使我们能尽己所能，不仅避免恶事本身，也避免一切恶表，并保全我们作为管家职务所必需的美名。我恳请您的智慧留意，这做法是多么具有恶表。因为我听到他们的忧伤（我们自己在\Place{Thiave}曾亲眼目睹），担心如常发生的那样，我因偏袒己见而弄错，于是我将此事的实情告诉我们的弟兄兼同僚\Person{Samsucius}，当时我没有告诉他我对此事的当前看法，反而陈述了我们在抵制他们要求时两人共持的看法。他极为震惊，且惊讶我们竟持有这种看法；令他不安的不是别的，正是这交易的可憎面目，这面目不仅完全不配我们，而且不配任何人。\par

5. 因此，我恳求你签署并毫不迟延地转交那封我以我们两人名义写给他们的书信。即便你可能认为另一做法在此事上是正义的，也不要让那些软弱的弟兄被迫立刻学习连我本人也无法理解的事；反而在与他们交往时应记起主所说的这句话：\BibleQuote{我本来还有许多事要告诉你们，然而你们现在不能承当。}\BibleRef{若 16:12} 因为主自己，出于对这种软弱的俯就，曾在另一场合（关于纳税的事）说过：\BibleQuote{所以儿子就免税了。但为避免使他们疑怪}等等；并派遣\Person{伯多禄}去缴纳当时所征收的\OriginalTerm{殿税}{didrachmæ}。\BibleRef{玛 17:26} 因为祂知道另一种法律，依照那法律祂并不负有缴纳此税的义务；但祂却按照那法律（如我所言，\Person{Honoratus}的遗产继承亦应依此法律调整）所课于祂的，缴纳了税。不但如此，即使在有关教会法律的事上，\Person{保禄}也对软弱者显示了宽容，并没有坚持他应得的报酬，尽管他在良心中完全确信自己完全有权要求；然而他放弃自己的权利，仅因为他借此避免对其动机的猜疑，以免那猜疑损坏他在他们中基督的馨香之气，并避免在某地区出现恶表，他知道这是他的义务，而且很可能在凭经验知道此事将引起的忧伤之前，他就已经这样做了。我们如今虽已有些迟延，并已由经验得到告诫，就应当纠正那本应预见的事。\par

6. 我记得，我们分别时你曾提议，应由\Place{Thagaste}的弟兄们要求我负责补足所要求款项的一半；作为结论，我愿说，既然我害怕一切可能使我的努力失败的事，如果你清楚地认为那提议是公正的，我将不拒绝遵行此事，然而，附此条件：我只在我有能力支付时才支付此数额，\Emph{即}当有某笔相当可观的捐赠归于我们\Place{Hippo}的隐修院，使得这样做不会过分窘迫我们---在扣除如此大额款项后，剩余数额仍足以按居住在此隐修院的弟兄人数比例，为我们的隐修院提供同等的份额。\par

\SourceRecord{Translated by J.G. Cunningham.；Nicene and Post-Nicene Fathers, First Series；Vol. 1.；Edited by Philip Schaff.；Buffalo, NY: Christian Literature Publishing Co.,；1887.；<http://www.newadvent.org/fathers/1102083.htm>.}

% structure: p0001 source=h1 target=section reason=page-title

\section{第八十四封信（主后405年）}

\Emph{致我主\Person{诺瓦图斯}——最可敬者，我在司铎职上的弟兄与同侪，我所尊敬并渴望的——并致与他同在的弟兄们：奥古斯丁及与他同在的弟兄们在主内致候。}\par

1. 我自知在你看来必定显得心硬，而且我几乎无法为自己开脱，因为我不肯答应将我的神子执事\Person{卢西卢斯}——你的亲兄弟——遣至你的圣善那里。但等你自己必须为了远方教会的要求，而放弃那些你所培育、为你所至爱而亲昵的人时，那时，只有那时，你才会明白那种渴望的痛苦如何将我刺透，它来自那些曾与我亲密无间、极其融洽相处的人，如今却不再在我身边了。让我将一位远方之人的境况呈于你的思念中。弟兄之间的血亲关系无论多么密切，也不及我所弟兄\Person{塞维鲁}与我之间的连结；然而你知道，我何等难得有与他相见的福分。如此情形，既非出于他的意愿，也非出于我的意愿，而是因为我们为了那将来永远不再分离的永恒，甘愿将我们母亲教会的要求，置于现世的要求之上。因此，你为了教会的益处而忍受那位弟兄的离别，不是更为合理吗？与他一同分享主我们的牧者所赐食粮的时日，远不及我与最可爱的同乡\Person{塞维鲁}一起的时日长久；而如今，我与他只能费力地、隔很久才借着简短的信函交谈——而且这些信函，多半还充满了他人的挂虑和事务，却不为我带来丝毫关于那些青绿草场的回忆，曾经我们在基督的爱护下，一起躺卧在那些草场上！\par

5. 你也许会答复说：\Quote{那又如何？我的弟兄不也可以在此地有益于教会吗？我想要他与我在一起，岂不正是为了教会的益处？}的确，如果他在你身边对我来说，在聚集或管理主的羊群方面，与他在此地的用处一样重大，那么每个人都会理所当然地责备我不仅心硬，而且不公。然而他通晓布匿语，而这里正因为缺乏此语，宣讲福音的工作大受阻碍；你们那里则普遍使用该语。你认为，如果我们为了主之羊群的益处尽心尽力，却将这特有的才能差派到并不特别需要的地方，而将它从我们如此焦渴盼望它的地区撤走，这样算是尽到了职责吗？因此，请你原谅我，当我违拗你的意愿并且也违反自己的情感，这样做实乃我身上的负担的关切迫使我如此。你将心灵献给了他，主必赐给你在工作中所需的助佑，使你因这份善意而得到酬报；因为我们承认，你是出于甘心，而非迫不得已，才将执事\Person{卢西卢斯}让给这我们命运所在的地区之焦渴者。你若不再因这事对我提出任何请求，便是给了我莫大的恩惠，免得我显得对你不仅是稍微心硬了些；我是因你那圣善的性情而尊敬你的。\par

\SourceRecord{Translated by J.G. Cunningham.；Nicene and Post-Nicene Fathers, First Series；Vol. 1.；Edited by Philip Schaff.；Buffalo, NY: Christian Literature Publishing Co.,；1887.；<http://www.newadvent.org/fathers/1102084.htm>.}

% structure: p0001 source=h1 target=section reason=page-title

\section{第八十五封书信（主历405年）}

\Emph{致我最亲爱的、在\OriginalTerm{司铎职}{priesthood}上的弟兄与同僚我主\Person{保禄}，愿我的一切祈祷为他祈求至高福祉，\Person{奥思定}在\OriginalTerm{主}{Lord}内致以问候。}\par

1. 若非你也认为我是个伪君子，你便不会称我如此不可通融。因为若凭你所写来论断我的心意，你还能相信什么别的呢？无非是我对你怀有应受责备和憎恶的反感与敌意；仿佛在这只能有一种意见的事上，我竟不小心谨慎，免得在规劝别人时，我自己却应受责斥，\BibleRef{格前 9:27}，或者想要去掉你眼中的木屑，却保留和助长自己眼中的大梁，\BibleRef{玛 7:4}。绝非如你所想。看！我重申此事，并呼求\OriginalTerm{天主}{God}作证：只要你为自己渴求我所为你渴求的，你现在就会在\Person{基督}内生活，免于一切不安，并使整个教会因你带给祂圣名的光荣而欢欣。请注意，我恳请你，我不仅称你为我的弟兄，也称你为我的\OriginalTerm{同僚}{colleague}。因为\OriginalTerm{天主教会}{Catholic Church}的任何一位\OriginalTerm{主教}{bishop}，只要未被任何一个教会法庭定罪，就不可能不再是我的\OriginalTerm{同僚}{colleague}。至于我拒绝与你共融，唯一的理由是我不能奉承你。因为我在\Person{基督}内生了你，便有特别的责任，以忠信的劝诫和谴责，向你施以有益于得救的爱的严厉。诚然，我欣喜于那些由于天主祝福你的工作而被聚集到\OriginalTerm{天主教会}{Catholic Church}中的人；但这并未减轻我的哀恸，因为更多的人正因你而四散离开教会。因为你如此伤害了\Place{希波}教会，以致除非\OriginalTerm{主}{Lord}使你摆脱一切世俗的忧虑和负担，并召叫你回到符合真\OriginalTerm{主教}{bishop}的身份的生活方式和品行，否则这创伤不久便会无法医治。\par

2. 然而，看到你仍然越来越深地卷入这些事务，并且不顾你的弃绝誓愿，再次将自己纠缠于你曾庄重搁置的事物中——这一步即使在普通人事的法律下也无法自圆其说；又听说你以一种奢侈的方式生活，而这单凭你教会的微薄收入是难以维持的——你为何还坚持要与我共融，却拒绝听取我对你过错的谴责呢？难道是因为我无法忍受别人的怨言，他们便可以因你所做的一切而公正地指责我吗？此外，你误以为那些指责你的人，就是很早以前甚至在你早年生活中便一直反对你的人。并非如此：你无需惊讶有许多事逃避了你的观察。但即使情况如此，你也有责任确保，你的行为中找不出任何他们可以合理责备并因此毁谤教会的事。或许你以为我说这些话的原因，是我没有接受你在辩护中所提出的理由。其实不然，我的理由乃是：对于这些事，我若缄口不言，在天主面前我将犯下无可推诿的罪过。我知道你的才能；但即使心智迟钝的人，若将自己的情感放在天上的事上，也能免于不安；而心智敏锐的人若将自己的情感放在地上的事上，这恩赐也徒然无用。\OriginalTerm{主教}{bishop}的职分并非为了让人可以度过虚荣的一生。\OriginalTerm{主}{Lord}天主已向你关闭了所有你曾想利用祂为自己谋利的道路，为的是倘若你明白祂，祂便可引导你走上那条道路，正是为了追求那条道路，那神圣的责任才被放在你身上，祂将亲自教导你我现在所说的。\par

\SourceRecord{Translated by J.G. Cunningham.；Nicene and Post-Nicene Fathers, First Series；Vol. 1.；Edited by Philip Schaff.；Buffalo, NY: Christian Literature Publishing Co.,；1887.；<http://www.newadvent.org/fathers/1102085.htm>.}

% structure: p0001 source=h1 target=section reason=page-title

\section{书信86（公元405年）}

\Emph{致我高贵的主\Person{切齐利亚努斯}，我真正且应受尊敬并在基督之爱内蒙我珍爱的儿子，\Person{奥古斯丁}主教在主内致意。}\par

您行政的声誉、您德行的美名，以及您基督徒虔敬之可嘉的热忱和忠信的真挚——这些是天主的恩赐，使您因那位赐恩者而喜乐，并从祂那里盼望领受更大的事——都促使我借这封信向尊贵的您述说那搅扰我心灵的挂虑。我们在极大的喜乐中，因为您在\Place{非洲}其余地方为\OriginalTerm{公教的合一}{Catholic unity}采取了措施，并取得了非凡的成就；同样，我们也感到极大的忧伤，因为\Place{希波}地区及\Place{努米底亚}边境邻近的地带，尚未享受到您作为执法官强力执行法令所带来的效益——我高贵的主，我真正且应受尊敬并在基督之爱内蒙我珍爱的儿子。为了避免此事被视作我——在\Place{希波}肩负\OriginalTerm{主教职务}{episcopal office}者——玩忽职守所致，我认为有责任向尊贵的您提及。若您屈尊了解\Place{希波}地区内\OriginalTerm{异端者}{heretics}的厚颜无耻已到了何等极端的地步——您可以询问我的司铎弟兄和同事，他们能向尊贵的您提供信息，或询问我随此信派去的\OriginalTerm{司铎}{presbyter}——我相信您定会如此处理这不虔傲慢的毒瘤，使之借着警告而得医治，而非日后以惩罚痛苦地割除。\par

\SourceRecord{Translated by J.G. Cunningham.；Nicene and Post-Nicene Fathers, First Series；Vol. 1.；Edited by Philip Schaff.；Buffalo, NY: Christian Literature Publishing Co.,；1887.；<http://www.newadvent.org/fathers/1102086.htm>.}

% structure: p0001 source=h1 target=section reason=page-title

\section{书信 87（公元 405 年）}

\Emph{致他挚爱且思念的兄弟 \Person{埃梅里图斯}，\Person{奥古斯丁}致候。}\par

1. 我知道，灵魂的救恩并不取决于拥有良好的天赋和博雅的教育；但当我听说有如此天赋的人，在一个极易考察的问题上，却持守与真理所强求的不同的观点时，我对这样的人越感惊讶，就越燃烧着渴望去结识他、与他交谈；如果做不到，就切望通过书信往来使他的思想与我交流，书信即使远隔千山万水也能飞越。我听说你正是我所描述的这样一个人，便难过你竟被割绝于天主教会之外，不能相共融通，正如圣神所预言的，她遍布整个世界。我不知你为何这样分离。因为毫无疑问，在罗马世界的大部分地区，\OriginalTerm{多纳图斯派}{Donatists}是完全不为人知的；更不用说那些蛮族国家（宗徒也曾说，他欠他们的债 \BibleRef{罗 1:14}），他们在基督信仰上与我们共融，事实上他们甚至根本不知道这分裂是何时、因何而起的。如今，除非你承认这些基督徒并没有你们指控非洲基督徒的那些罪行，否则你就必须承认，只要那些无赖之辈（说得温和些）在你们中间未被察觉，你们所有人就因参与这些无赖的恶行而被玷污了。因为你们偶尔会将某个成员从你们的共融中开除，而开除总是在他犯下应被开除的罪行之后才发生的。在他被发现、定罪、并被你们判罚之前，难道不是有一段他未被察觉的中间时间吗？所以我问，在他未被你们发现的那段时间里，他是否使你们沾染了他的污秽？你回答说，\Quote{绝不可能。}那么，如果他根本未被发现，他就绝不会使你们沾染他的污秽；因为有时人们的罪行在死后才暴露出来，但这并不会使那些在他们生前与他们共融的基督徒负有罪责。那么，你们为何以如此轻率而不敬的\OriginalTerm{裂教}{schism}，割裂自己，不与无数东方教会共融呢？在那些教会中，你们或真或假地坚称在非洲所发生的一切，始终是完全不为人知的。\par

2. 至于你们的主张是否真实，那完全是另一个问题。我们以远比你们所提出的更可信的文献驳斥这些主张，而且我们还在你们自己的文献中找到了更多证明你们所攻击的那些立场的证据。但如我所言，这完全是另一个问题，必要时再予以讨论。与此同时，请你的心思特别留意这一点：没有人会因不知名之人所犯的不知之罪而分担罪责。由此显明，你们犯下了不虔敬的裂教之罪，因为你们割裂自己，不与整个世界共融；而你们无论真伪地指控非洲某些人的那些事，对这个世界来说，始终是完全不为人知的。然而也不应忘记这一点：即使恶人被知晓并发现，若缺乏阻止他们参与共融的权力，或者教会的和平利益不允许这样做，他们也不会伤害教会中的善人。因为按照先知\Person{厄则克耳}的预言，那些在恶人毁灭前获得被标记的赏报、并在他们毁灭时未受伤害的人，不就是那些因天主子民中发生的罪恶与不义而叹息哀号的人吗？有谁会对不知之事叹息哀号呢？基于同一原则，宗徒\Person{保禄}容忍了假弟兄。因为他所说的\Quote{各人只顾自己的事，而不顾耶稣基督的事；}并非指他不知之人；而他明确表示这些人就在他身旁。那些宁愿向偶像焚香或交出圣经而不愿受死的人，如果不属于那些\Quote{只顾自己的事，而不顾耶稣基督的事}的人，又属于哪一类呢？\par

3. 我略去了许多本可从圣经中引用的证据，以免使这封信超出必要的长度；并且我还留下许多事情，让你凭自己的学识去思考。但我恳请你注意这一点，这已足以决定整个问题：如果在一个民族——当时的天主教会——中，如此众多的违法者并没有使那些与他们往来的人也像他们一样有罪；如果那群假弟兄并没有使与他们同属一个教会的宗徒\Person{保禄}，成为一个不顾耶稣基督的事而只顾自己事的人——那么显而易见，一个人不会因与他同赴基督祭台之人的邪恶而变得邪恶，即使那人并非不为他所知，只要他不怂恿那人的邪恶，而是凭着良心不赞成其行为，使自己与他隔开。因此显然，要成为盗贼的同伙，就必须要么帮他行窃，要么欣然收受他所偷窃之物。我这样说，是为了去除那些有关人类行为的无尽而不必要的问题，这些问题在用来反对我们的立场时，是完全不相干的。\par

4. 然而，如果你不同意我所说的，你就使你们整个党派都蒙受了像\Person{奥普塔特}那样的人的羞辱，尽管你们明知他的罪行，却容忍他与你们共融；我绝无将此话用于像\Person{埃梅里特}这样的人，以及你们中间其他正直之士，我确信他们都完全厌恶那些玷污了他的恶行。因为我们并未指控你们别的，只是指控你们\OriginalTerm{裂教}{schisma}，而你们由于顽固坚持，如今已使它成为\OriginalTerm{异端}{haeresis}。你们可以借着阅读你们无疑早已读过的那段经文，看出这罪行在天主的审判中是多大。你们会发现\Person{达堂}和\Person{阿彼兰}被地裂开的口吞噬，所有与他们同谋的人被从他们中间喷出的火烧死\BibleRef{户 16:31}。上主天主为警告人逃避此罪，以这一立即的惩罚昭示了犯罪者的下场，为要表明祂为犯有同样过犯的人在最后审判时所保留的报应，祂现在仍以极大的宽容暂时容忍他们。的确，我们不苛责你们容忍\Person{奥普塔特}留在你们中间的理由。我们不怪责你们，因为当\Person{奥普塔特}因疯狂滥权而被指控，遭受整个\Place{阿非利加}的悲叹攻击时——如果你们真是那位良好名声所传说的那样（天主知道，我极愿意相信这个传闻），你们也一同悲叹——你们没有将他\OriginalTerm{绝罚}{excommunicatio}，是怕他在这种判决下拉走许多人，以裂教的疯狂撕裂你们的\OriginalTerm{共融}{communio}。然而，这件事本身就是在天主的法庭上对你们的控诉，\Person{埃梅里特}弟兄啊：你们虽然看到\OriginalTerm{多纳图派}{Donatistae}的分裂是如此大恶，以至于认为宁可容忍\Person{奥普塔特}留在你们的共融中，也比在你们中间引入分裂更好，但你们却依然延续着由你们祖先造成的基督教会分裂的祸害。\par

5. 这里，你可能会因辩论的需要而倾向于为\Person{奥普塔特}辩护。我恳求你，不要这样；不要这样，我的弟兄：这不适合你；而且即使有人适宜这样做（尽管事实上，任何不正当的事都是不宜的），但\Person{埃梅里特}为\Person{奥普塔特}辩护肯定是不宜的。或许你会回答说，你同样不宜指控他。完全同意。那么，采取介于辩护和指控之间的态度吧。说：\BibleQuote{每人要背负自己的重担} \BibleRef{迦 6:5} \BibleQuote{你是谁，竟敢判断别人的仆人？}\BibleRef{罗 14:4} 那么，尽管有整个\Place{阿非利加}的见证——不，甚至所有听说过\Person{吉尔多}名字的地方，因为\Person{奥普塔特}的臭名并不亚于他——你们却不敢对\Person{奥普塔特}作出判断，唯恐鲁莽地对你们不了解的人下结论。我问，我们若仅仅根据你们的证词，就轻率地对我们不了解的人宣判，这难道是可能的或是正当的吗？你们以自己都没有把握的事情来指控他们，难道还不够吗？非要我们也对和你们一样无知的事情定他们的罪吗？因为即使\Person{奥普塔特}因诽谤者的谎言而身处险境，当你说\Quote{我不知道他品行如何}时，你并不是在为他辩护，而是在为自己辩护。那么，东方世界对那些你们虽远不如了解\Person{奥普塔特}那样了解、却被你们指责的阿非利加人的品行一无所知，岂不更明显吗？然而，你们却因可耻的裂教而与东方的各教会断绝了，你们拥有这些教会的名字，并在圣经中读到它们。如果你们那位最有名、最臭名昭著的\Place{塔穆加迪}主教，在那时甚至不为他的同事所知，我就不说\Place{凯撒勒亚}了，单说近在咫尺的\Place{西蒂菲斯}，那么\Place{格林多}、\Place{厄弗所}、\Place{哥罗森}、\Place{斐理伯}、\Place{得撒洛尼}、\Place{安提约基}、\Place{本都}、\Place{迦拉达}、\Place{卡帕多细亚}的教会，以及其他由宗徒们在基督内建立的教会，怎么可能知道这些阿非利加的\OriginalTerm{背教者}{traditor}——无论他们是谁——的事呢？你们又怎么能以他们不知道为由定他们的罪呢？然而，你们与这些教会并不共融。你们说他们不是基督徒，并致力于为他们的成员\OriginalTerm{重新施洗}{rebaptizare}。我还需要说什么？这里还需要什么控诉、什么抗议？如果我是在对一个正直的人说话，我知道你与我有同感共愤。因为你无疑立刻就能看出我若愿意还会说什么。\par

6. 然而，也许你们的祖先自己组织了一个会议，将除他们自己以外的整个基督教世界都处于绝罚的判决之下。你们竟如此判断事理，声称\Person{马克西米安努斯}的追随者（他们从你们中分裂出去，正如你们从教会分裂出去）所开的会议对你们没有权威，因为他们的数目与你们相比是少数；然而，你们却又为你们的会议要求对万民（他们是基督的产业）和地极（这是祂的基业）的权威吗？我怀疑那个不以这样的自诩为耻的人，身上是否还有一点血性。我恳求你写回信给我；因为我从某些我不得不信赖的人那里听说，如果我写信给你，你会回信给我。此外，前些时候我曾寄给你一封信；但我不知道你是否收到或回复了，也许你的回复未能送达我这里。但现在，我恳求你切勿拒绝回复这封信，并说明你的看法。但请不要纠缠于我提出的问题以外的其他问题，因为这是对裂教起源进行有序讨论的关键所在。\par

7. 世俗权力为其迫害分裂者的行为辩护，依据宗徒所定的规则：\BibleQuote{所以谁反抗权柄，就是反抗天主的规定，而反抗的人就是自取惩罚。因为长官为行善的人，不是可怕的；为行恶的人，才是可怕的。你愿意不怕掌权的吗？你行善罢！那就可由他得到称赞，因为他是天主的仆役，是为相帮你行善；你若作恶，你就该害怕，因为他不是无故带剑；他既是天主的仆役，就负责惩罚作恶的人，向作恶的人，发泄义怒。}\BibleRef{罗 13:2} 因此，整个问题在于，分裂是否不是一件恶事，或者你是否未曾造成分裂，以至于你对现存的权柄的抵抗，是出于良好的动机，而非恶行，从而给自己招致惩罚。因此，主以无限的智慧，不仅说：\BibleQuote{为义而受迫害的人是有福的，}而且加上：\BibleQuote{为义。}\BibleRef{玛 5:10} 所以，根据我前面所说的，我愿意从你那里知道：发起并使你们与教会分离的状态延续下去，究竟是不是一件义行？我也愿意知道：在未听申诉的情况下，就谴责整个基督徒世界，——要么因为它没有听到你们所听到的，要么因为没有向它提供任何证据，证明那些轻易相信的，或你们在缺乏充分证据的情况下提出的指控，——并据此提议为如此众多由主在世时的亲自宣讲和辛劳，或由祂的宗徒们所建立的教会的成员重行圣洗，这一切难道不是一件不义之举吗？而且这一切都基于这样的假设：你们要么不知道与你们比邻而居、与你们共同使用同样圣事的非洲同伙的邪恶，要么即使知道了也容忍他们的恶行，以免多纳图派分裂，这是可以谅解的；但对他们而言，尽管他们居住在最遥远的地区，却不可原谅他们不知道你们关于非洲人的所知、所信、所闻或所想象的事。那些固守自己的不义，却指责世俗权力的严厉的人，是何等乖戾啊！\par

8. 你也许会回答说，基督徒甚至不应该迫害恶人。就算是这样吧；让我们承认他们不应该：但是，难道可以以此反对那正是为此目的而设立的权柄吗？难道我们要删除宗徒的话吗？或者，你们的抄本不包含我刚才提到的那些话吗？但你会说，我们不应该与这样的人相通。那又怎样？你们前些时候，难道不是因为副巡抚\Person{弗拉维安}在执法时处死了他认定有罪的人，就与他断绝了相通吗？再者，你会说，罗马皇帝是被我们煽动起来反对你们的。不，倒要为此责备你们自己，因为正如很久以前关于基督的应许所预言的：\BibleQuote{众王都要朝拜他，}他们如今成了教会的成员；而你们竟敢以分裂伤害教会，并且仍擅自坚持为其成员重行圣洗。我们的弟兄确实寻求那设立的权柄的援助，不是为了迫害你们，而是为了保护自己，免受你们派别中某些个人的不法暴力行为的侵害——你们自己不参与这种行为，但也为此哀叹悲愤；正如在罗马帝国成为基督徒之前，宗徒\Person{保禄}采取措施，要求罗马士兵武装保护，以对付那些密谋杀害他的犹太人。然而，无论这些皇帝因何契机得知你们分裂的罪行，他们都针对你们制定了他们的热忱和职责所要求的法令。因为他们不是无故带剑；他们是天主的仆役，要向作恶的人执行义怒。最后，如果我们的某些人在此问题上超越了基督徒节制的界限，我们是不悦的；尽管如此，我们并不因此就离弃天主教会，因为我们无法在最后的簸扬之前，将麦粒与糠秕分开，特别是你们自己也没有因为\Person{奥普塔图斯}的缘故而离弃多纳图派，因为你们没有勇气因他的罪行将他开除教籍。\par

9. 然而，你说：\Quote{既然我们如此被罪污所玷染，为什么还要设法让我们与你们联合呢？} 我答复说：因为你们仍然活着，如果你们愿意，仍能恢复。当你们与我们，\Emph{即}，与天主的教会——基督的产业，祂拥有地极作为自己的产业——相连时，你们便得以恢复，从而与\OriginalTerm{根}{Root}在生命中联合。因为宗徒论及那些被折下的枝条说：\BibleQuote{天主能重新接上他们。}\BibleRef{罗 11:23} 我们劝勉你们，在背离教会这一点上作出改变；尽管至于你们所领受的圣事，我们承认它们是神圣的，因为它们在一切人中是相同的。因此，我们渴望见到你们改变你们的固执，也就是说，使你们这些已被切断的人能重新与根在生命中相连。因为你们未曾改变的圣事，正如你们所拥有的一样，我们予以赞同；否则，我们若试图纠正你们的罪，就会对基督的奥迹做出不虔的伤害，这些奥迹并未因你们的不配而失去其价值。因为即使是撒乌耳，虽犯有诸多罪恶，也没有摧毁他所领受的傅油的效力；对这番傅油，天主的虔敬仆人达味表示了极大的尊重。因此，我们并不坚持为你们重新施行圣洗，因为我们只希望恢复你们与根的连接：那已被切断的枝条的形状，若未曾改变，我们欣然接受；但枝条，无论其形状多么完美，除非与根联合，便不能结果实。至于你们所说的，从我们这一方所受的、温和且充满仁慈的迫害——而毋庸置疑，你们那一方却以不法且无序的方式，对我们施加了更大的伤害——这是一个问题；与圣洗有关的问题，则完全与此不同；在圣洗的问题上，我们询问的，不是它在哪里，而是它在何处产生效益。因为无论圣洗在哪里，它都是一样的；但不能说，领受圣洗的人，无论他身处何处，他也是一样的。因此，我们憎恶人们在裂教状态下，作为个人所犯的不虔罪行；但我们随处尊敬基督的圣洗。如果逃兵携带着帝国的军旗，当逃兵要么受到严厉判决的惩罚，要么在仁慈的运作下被恢复时，这些军旗若仍保持完好，会以原样被重新接纳。关于此事，若要进行更详尽的探讨，那是——正如我说过的——另一个问题；因为在这些事情上，天主的教会的惯例是我们行动的准则。\par

10. 然而，我们之间的问题是，究竟是你们的教会还是我们的教会才是天主的教会。要解决这个问题，我们必须从最初的调查开始，即你们为何成为裂教者。如果你们不给我回信，我相信，在天主的审判台前，我因已在这件事上尽了自己的本分，将很容易得到辩白；因为我给我所听说的一位善良且受过良好教育的人——仅因他依附裂教者——寄去了一封为和平利益而写的信。你们应当考虑，该如何答复那位，祂如今的宽容要求你们的赞美，而祂的审判终将要求你们的畏惧。然而，如果你们给我回信，像我写这封信一样用心，我相信，藉着天主的仁慈，如今使我们分裂的谬误，将在爱和平与真理的逻辑面前消逝。请注意，对于\Person{罗加图斯}的追随者，我绝口未提，他们称你们为\OriginalTerm{斐尔米安派}{Firmiani}，正如你们称我们为\OriginalTerm{玛加里安派}{Macariani}。我也未提及你们在\Place{鲁卡塔}（或称\Place{路西卡达}）的主教，据说他曾与\Person{斐尔姆斯}达成协议，许诺，在他所有追随者安全的条件下，为他打开城门，并将天主教徒交出去受屠杀和抢劫。还有许多其他类似的事，我略过不提。因此，你们也应同样停止对你们所听到或知道的人的行为，作老套的修辞夸张；因为你们看到，我对你们一方的行为保持缄默，是为了将辩论限定在整个问题的关键上，即裂教的根源。\par

我亲爱的、思念的兄弟，愿主我们的天主在你内吹入趋向和解的意念。\par

\SourceRecord{Translated by J.G. Cunningham.；Nicene and Post-Nicene Fathers, First Series；Vol. 1.；Edited by Philip Schaff.；Buffalo, NY: Christian Literature Publishing Co.,；1887.；<http://www.newadvent.org/fathers/1102087.htm>.}

% structure: p0001 source=h1 target=section reason=page-title

\section{书信八十八（公元406年）}

\Emph{致\Person{雅努阿里乌斯}，\Place{希波}地区的天主教圣职人员} \Emph{寄送如下：}\par

1. 你的圣职人员和你的\OriginalTerm{Circumcelliones派}{Circumcelliones}正以一种新颖且空前凶残的迫害向我们发泄怒火。倘若我们以恶报恶，便违背了基督的律法。但现在，当公正地审视你方与我所做的一切，便发现我们正在遭受经上所记载的：\BibleQuote{他们以恶报善}\BibleRef{咏 35:12}；又（另一篇圣咏）：\BibleQuote{我的灵魂久居在憎恨和平的人中。我言和平，他们却主张战争。}\BibleRef{咏 120:6-7} 因为，鉴于你已届如此高龄，我们料想你完全知道，\Person{多纳图}派——最初在\Place{迦太基}被称为\Person{玛约里努}派——曾主动在当时著名的\Person{君士坦丁}皇帝面前控告当时\Place{迦太基}主教\Person{凯奇利安}。然而，可敬的先生，恐怕你已忘记此事，或佯作不知，或可能（我们几乎认为不可能）从来不知，我们兹将当时总督\Person{阿努利努}的叙述副本插入于此；\Person{玛约里努}派曾向他申诉，请求他作为总督将他们控告\Person{凯奇利安}的陈述呈送前述皇帝：——\par

\Emph{致\Person{君士坦丁}·奥古斯都：\Person{阿努利努}，一位具有执政官等级的\Place{非洲}总督，谨呈如下：}\par

至尊陛下致\Person{凯奇利安}及其所管被称为圣职人员者之受崇敬属天诏书，已由陛下最卑微的仆人录入档案；并已敦促这些人士，既然他们因陛下宽仁而免除其他所有负担，便当彼此同心合意，保持天主教合一，以法律与神圣事务应受之敬礼，专注于自身职责。然而，不数日后，出现某些人，并有大批民众附从，他们认为应对\Person{凯奇利安}提起诉讼，遂向我呈交一个以皮革包裹的密封包裹及一小份未封缄之文件，恳切请求我将它们转呈陛下神圣庄严的朝廷。陛下最卑微的仆人已遵嘱办理，其间\Person{凯奇利安}仍维持原状。有关此案之卷宗附呈于后，以便陛下就全部事宜作出裁决。所递文件有二：其一置于皮封内，标题为\Quote{天主教会控告\Person{凯奇利安}之文件，由\Person{玛约里努}派呈交}；另一份未封缄，附于同一皮封之上。\par

写于五月朔日前十七日，即吾主\Person{君士坦丁}·奥古斯都第三次执政官任期（\Emph{即}公元313年4月15日）。\par

3. 此报告上呈皇帝后，皇帝遂召集各方至在\Place{罗马}设立的主教法庭。教会档案显示该案如何辩论及裁决，并宣布\Person{凯奇利安}无罪。至此，在主教法庭作出带来和平的裁决之后，一切顽梗的争竞与苦毒本应平息。然而，你们的先祖再次上诉至皇帝，抱怨裁决不公，且案件未获充分审理。因此，皇帝又在\Place{高卢}城镇\Place{阿尔勒}设立第二个主教法庭；在判决谴责你们毫无价值且属魔鬼的裂教后，你们中许多人回归与\Person{凯奇利安}修好；但有些最顽固好争之徒，又向皇帝上诉。其后，皇帝对他们的强求让步，亲自介入这原属主教裁决的争端；审案之后，他判决反对你们一派，并率先立法没收你们团体的财产；这一切我们本可在此插入文件证据，惟恐封信过长。然而，我们万不可省略对\Person{阿普图加的费利克斯}一案的调查与公开判决。在\Place{提吉西斯}主教\Person{塞孔杜斯}担任首席主教期间的\Place{迦太基}会议上，你们的教父们断言他是所有这一切祸患的原始起因。因为前述皇帝在我们将附上副本的信函中证实，在那次审判中，你们一派曾在他面前作为原告和最激烈的起诉人：——\par

\Emph{皇帝\Person{弗拉维乌斯·君士坦丁}、\Person{马克西穆斯}·凯撒，及\Person{瓦莱里乌斯·李锡尼}·凯撒，致\Place{非洲}总督\Person{普罗比安努斯}：}\par

你的前任\Person{埃利亚努斯}，他在那位最卓越的长官\Person{维鲁斯}因病重而无法理事期间，代理了监察长之职；在他代理期间，于我们管辖下的非洲地区，他认为有责任——也是最公正地——重新调查并裁决\Person{凯基利安努斯}的事件，更确切地说，是针对这位天主教会主教的仇恨情绪。因此，他传唤了百夫长\Person{苏佩里乌斯}、\Place{阿普同加}的长官\Person{凯基利安努斯}、前治安官\Person{撒图尔尼努斯}及其继任者\Person{小卡利比乌斯}，以及\Place{阿普同加}的一名官员\Person{索隆}，听取了这些证人的证词；结果表明，尽管有人因\Person{凯基利安努斯}的主教祝圣礼是由\Person{斐理斯}所执行而对他提出反对，因为\Person{斐理斯}似乎被指控交出了并焚毁了圣书，但\Person{斐理斯}在此事上的清白已被明确证实。此外，当\Person{马克西穆斯}声称\Place{齐库阿}镇的市议员\Person{英根提乌斯}伪造了一封前长官\Person{凯基利安努斯}的信件时，我们在查阅了面前的案卷后发现，这同一位\Person{英根提乌斯}曾因该罪行而受到拷问刑架之刑，而他遭受酷刑，并非如所言是因为他声称自己是\Place{齐库阿}的市议员。因此，我们希望你派遣适当看押的上述\Person{英根提乌斯}到奥古斯都\Person{君士坦丁}的朝廷，以便在那些如今为此案辩护的人、以及日复一日坚持控诉的人面前并当庭聆听时，可以显明并完全知晓：他们为煽动对主教\Person{凯基利安努斯}的仇恨而徒劳地操劳，并对他猛烈叫嚣。我们希望，这将使民众终止他们本应停止的这类纷争，转而以应有的敬畏专心于他们的宗教本分，不再因彼此的争执而分心。\par

因此，既然你看到事情是如此，你为何以针对你们的帝国法令为由，激起对我们的仇恨呢？而你们自己却在我们效仿你们之前，早已行过这一切。如果皇帝不应在此类事务上运用其权威，如果这些事务超出基督徒皇帝管辖的范围，那么是谁敦促你们的先人，通过总督，将\Person{凯基利安努斯}的案件呈报皇帝，并再次将一个你们在缺席情况下不知怎样就定罪了的主教的指控呈于皇帝面前，且在他被宣告无罪后，又编造关于\Person{斐理斯}的其他诽谤之词，呈于同一位皇帝面前，而\Person{斐理斯}正是那位上述主教的祝圣者？如今，除了那个\Person{君士坦丁}的裁决之外，还有什么法律在强制你们这一派呢？你们的先人曾出于自愿而诉诸于它，他们通过不休的控诉从他那里强索了它，并且他们宁愿要它，也不要主教会议的裁决。如果你们对皇帝的法令不满，那么是谁首先迫使皇帝将这些法令加诸你们呢？因为你们没有理由因针对你们的皇帝法令而向天主教会叫嚣，正如那些在\Person{达尼尔}获救后，被扔进去被同样的狮子吞食的人，没有理由向\Person{达尼尔}叫嚣一样；正是他们先试图用这些狮子毁灭\Person{达尼尔}。正如经上所说：\BibleQuote{君王震怒，有如狮子的咆哮。}\BibleRef{箴 19:12} 这些诽谤的敌人执意要将\Person{达尼尔}投入狮子坑中；他的清白胜过了他们的恶意；他毫发无伤地被从坑中取出，而他们被扔进去后便灭亡了。同样地，你们的先人将\Person{凯基利安努斯}及其同伴们抛入，让他们在国王的震怒下被毁灭；而当他们因清白而从中得以脱身时，如今你们自己却因这些国王而遭受了你们那一派想让他们遭受的；正如经上所说：\BibleQuote{谁为邻人挖掘陷阱，自己必陷在其中。}\par

因此，你们没有理由抱怨我们：不仅如此，天主教会的慈爱本会使我们甚至放弃执行这些皇帝法令，若不是你们的圣职人员和\OriginalTerm{巡游派}{Circumcelliones}，以他们最骇人的罪行和疯狂的暴力行为，扰乱我们的和平，毁灭我们，迫使我们重新恢复并执行这些法令。因为在你们所抱怨的这些较近期谕令传到非洲之前，这些亡命之徒就在路上伏击我们的主教，以野蛮的殴打虐待我们的圣职人员，并以同样极端残忍的方式袭击我们的平信徒，纵火焚烧他们的住所。有一位长老，他出于自己的自由选择而偏爱我们教会的合一，就因此被人从他的房屋中拖出，未经法律程序而被残酷殴打，在泥坑中被来回滚，盖上一层灯心草编的席子，被展出，让一些人同情，另一些人嘲笑，而迫害他的人则以罪行为荣；之后，他们将他随意带走，十二天后才不情愿地将他释放。当我们的主教因这次暴行而向\Person{普罗库莱亚努斯}提出质问时，在一次市政法庭会议上，他起初试图通过假装对此一无所知来逃避调查；而当要求立刻被重申时，他公开宣称对此事不再多说。而那暴行的实施者，如今仍在你们的长老之中，并且继续恐吓我们，极尽所能地迫害我们。\par

7. 然而，我们的主教并未向皇帝们申诉那些日子我们地区天主教会所遭受的冤屈和迫害。但是，当一次会议召开后，大家一致同意邀请你们与我们会面，以和平的方式进行，如果可能的话，你们[\Emph{即}双方的主教，因为此信是由\Place{希波}的圣职人员所写]可以举行一次会谈，从而消除错误，使兄弟之爱在我们中间因和平的联结而欢欣。你们可以从自己的记录中得知\Person{普罗库雷亚努斯}最初在那时机所作的回答，即你们会召集一次会议，并在会上决定如何答复；以及后来，当他再次被公开提醒他的承诺时，正如记录所证实的，他说他拒绝举行任何旨在和平的会谈。此后，当你们圣职人员和\OriginalTerm{西库姆塞利翁派}{Circumcelliones}恶名昭彰的暴行仍在继续，一个案件被送交审判；\Person{克里斯皮努斯}被定为异端者，虽然由于天主教徒的宽容，他被免除了皇帝谕令对异端者所课的十磅黄金的罚款，但他仍认为自己有理由向皇帝们上诉。至于对那次上诉所作的答复，难道不是由于你们一派先前的邪恶和他自己的上诉而逼出来的吗？然而，即使在那个答复作出之后，他还是被允许免于缴纳那笔罚款，这是由于我们的主教们为他向皇帝求情。然而，从那次会议后，我们的主教们派遣代表前往宫廷，他们获得一项法令，规定并非你们所有的主教和圣职人员都须承担这笔对所有异端者所课的十磅黄金的罚款，而只是那些在其辖区内天主教会遭受你们一派暴行的人。但当代表们抵达\Place{罗马}时，\Place{巴加}的天主教主教那时刚刚受到可怕的伤害，他的伤口感动了皇帝，使他发出了实际发出的那些谕令。当这些谕令送达\Place{非洲}时，看到特别是开始对你们施加强大压力，并非为了任何邪恶目的，而是为了你们的益处，你们所当做的是不是邀请我们的主教与你们会面，正如他们曾邀请你们的主教与他们相会一样，好经由会谈使真理显明出来呢？\par

8. 然而，你们不仅未能这样做，而且你们一派还不断对我们加以更大的伤害。他们不满足于用棍棒殴打我们，用刀剑杀害一些人，甚至以难以置信的罪恶巧计，把混有酸\EditorialNote{?硫酸}的石灰投到我们的人眼中，要使他们失明。此外，为了劫掠我们的家宅，他们还制造了巨大而可怕的器具，带着这些器具到处游荡，口中发出威胁，要杀戮、抢劫、焚烧房屋、弄瞎我们的眼睛；因着这些事，我们首先不得不向你申诉，尊敬的大人，恳求你考虑，在所谓天主教皇帝的严厉法律之下，你们中许多人、其实所有人，尽管声称自己是受迫害者，却平安地定居在你们自己的产业上，或你们从别人那里夺得的产业上，而我们却在你们一派手下遭受如此闻所未闻的冤屈。你们说你们受迫害，而我们却被你们的武装人员用棍棒刀剑杀害。你们说你们受迫害，而我们的家宅却被你们的武装强盗劫掠。你们说你们受迫害，而我们中许多人却被你们的成员为此目的而装备的石灰和酸毁坏了视力。而且，如果他们犯罪的行径导致他们中一些人死亡，他们就声称这些死亡理应引起对我们的憎恨，而为他们带来荣耀。他们不为他们对我们的伤害承担责任，反而把他们自己招致的伤害归咎于我们。他们活着像强盗，死得像\OriginalTerm{西库姆塞利翁派}{Circumcelliones}，却被尊为殉道者！不，我这样比较是对强盗不公平的；因为我们从未听说过强盗会毁坏他们所劫掠者的视力：强盗确实把所杀之人从光明中除去，却不夺走那些被留下活命之人的光明。\par

9. 另一方面，若我们任何时候将你们一派的人拿在手中，我们便使他们毫发无损，向他们显示极大的爱；并告诉他们一切，以此驳斥那使弟兄与弟兄分离的谬误。我们正是遵照主亲自藉着先知\Person{依撒意亚}的话所吩咐我们的去行：\BibleQuote{敬畏主言语的人，你们听主的话；对那憎恨你们、驱逐你们的人说：你们是我们的弟兄，好使主的名受光荣，并使他们看见而喜悦；但让他们蒙羞。}\BibleRef{依 66:5} 这样，我们便说服了他们中的一些人，因为他们思考了真理的证据及和平的美好，不再重新领受圣洗以表示效忠我们君王的记号（尽管他们曾是叛逃者），而是接受他们先前未曾有的那信德、对圣神的爱德，以及与基督身体的合一。因为经上记载：\BibleQuote{藉信德洁净了他们的心；}\BibleRef{宗 15:9} 又说：\BibleQuote{爱德遮盖许多罪过。}\BibleRef{伯前 4:8} 然而，倘若由于过分的顽固，或由于羞耻使他们无法忍受那些惯常与他们一同频繁地虚假毁谤我们、图谋加害我们之人的讥讽，或者更由于害怕会与我们一同遭受他们从前惯于加给我们的那种伤害——倘若，我说，出于以上任何一种原因，他们拒绝与基督的合一修和，我们就让他们离去，正如他们被拘留时一样，不受任何伤害。我们也尽力劝勉我们的平信徒，要拘留他们而不加任何伤害，并带他们到我们这里来接受劝告和教导。他们中的一些人听从了我们并如此行了，只要他们力所能及；另有人则待他们如同强盗，因为他们实际遭受了此类人常行的强盗行径。有些人为了保护自己的身体免受其殴打，便反击那些攻击者；还有些人将他们捉拿并送交官长；尽管我们为他们求情，官长也不放过他们，因为他们极其惧怕他们的野蛮暴行。然而，与此同时，这些人尽管坚持强盗的勾当，却在自己行为的应有报应临到时，声称应被尊为殉道者！\par

10. 因此，可敬的先生，我们藉着这封信和我们派遣的弟兄们，向您提出的愿望如下：首先，如果可能，请与我们的主教们举行一次和平的会议，好使谬误本身得以终结，而非使持守它的人被终结，使人得以纠正而非受惩罚；并且，如同您先前曾轻看他们和解的提议，如今请让他们从您这一方提出。你们的主教与我们的主教之间举行这样一次会议，将会议记录书写下来，经双方签署后呈送皇帝，这比你们与民事长官商谈要好得多，因为那些长官除了执行那已经制定以反对你们的法律之外，别无他法！因为你们那些从这地方航行出去的同僚曾说，他们是来向长官们呈诉案件的。他们还提到了我们圣善的父亲、天主教主教\Person{瓦伦提诺}，他当时正在宫廷，说他们愿意与他一同受审。法官不能准许这点，因为他在司法职务上受制于那已制定以反对你们的法律；而且，这位主教也并非由此途径而来，亦未奉其同僚的任何此类指示。因此，当那次会议的记录呈递皇帝面前时，皇帝本人将何等更有资格对你们的案件做出判决，因为他并不受这些法律的约束，有权制定别的法律取而代之；尽管可以说，早先已对此案作出了终审判决！然而，我们希望与你们举行这次会议，并非寻求再一次的终审判决，而是要使那些尚不知情的人知晓该案业已定谳。如果你们的主教愿意这样做，你们会损失什么？岂不反得益处，因为你们对此会议的热忱将会显明，而先前理所当然加给你们的、你们不信任自己理由的谴责，也将被消除？或许，你们认为这样的会议是不合法的？你们当然知道，我们的主基督也曾对魔鬼论及法律，\BibleRef{玛 4:4}；而宗徒\Person{保禄}不但与犹太人，也与属于斯多葛派和伊壁鸠鲁派的外教哲学士辩论。\BibleRef{宗 17:18} 或许，是皇帝的法律不许你们与我们的主教会晤？如果这样，那么请将你们在\Place{希波}地区的主教们暂时聚集起来，我们正在该地区遭受你们一派的人如此多的伤害。因为，你们可能寄给我们的信函，比他们用来攻击我们的武器，其通道要合法、公开得多！\par

11. 最后，我们恳请你们，通过我们派去的弟兄，将这类著作交回。但如果你们不愿这样做，至少也请听我们一方的话，正如你们听你们自己一派人说话一样；我们从他们手中蒙受了如此不义之事。请你们将那据你们声称因此受迫害的真理指示给我们，而此时我们正从你们一方遭受如此巨大的残暴对待。因为，如果你们定我们有错，或许你们会允许我们免受你们的再洗，因为我们是由你们未曾定罪的人施洗的；而你们曾将这种豁免给予那些由\Person{穆斯蒂的费利奇安努斯}和\Person{阿苏里的普雷特克斯塔图斯}施洗的人，那是在你们试图通过法律禁令将他们逐出自己教堂的漫长时期中，因为他们与\Person{马克西米安}共融，而他们连同马克西米安一起，在\OriginalTerm{巴盖会议}{Council of Bagæ}上被点名明确地定罪。这一切我们都能通过司法和市政记录来证明，在这些记录中，你们出示了你们同一会议的决议，那时你们想向法官表明，你们正从教会建筑中驱逐的人，是那些因分裂而与你们分离的人。然而，你们藉着分裂，已将自己与\Person{亚巴郎}的子孙隔绝，\BibleQuote{地上万民都要因他蒙受祝福}\BibleRef{创 22:18}，却拒绝被逐出我们的教会建筑，而此项命令并非出自你们在处理自己教派分裂分子时所用的那种法官，而是出自地上的君王本人，他们正如预言所预言的敬拜基督；当你们控告\Person{凯基利安}时，你们曾从他们的法庭败退而归。\par

12. 然而，如果你们既不愿教导我们，也不愿听我们，那就请你们亲自来，或派你们中的一些人到\Place{希波}地区，由我们中的一些人为他们作向导，好让他们看看你们那支装备着武器的军队；不，比以往任何军队装备得更充分的军队，因为从未听说过有士兵在与蛮族作战时，会在其他武器之外加上石灰和酸液，以毁坏敌人的眼睛。如果你们连这也拒绝，我们至少恳请你们写信给他们，要他们现在停止这些事，不要再杀害、抢劫和弄瞎我们的人民。我们不会说：定他们的罪吧；因为，是否因容忍那些我们证明为强盗的人留在你们的共融内，而你们丝毫不受玷污，这是你们自己应当看明的事；而我们却因拥有你们从未能证明是\OriginalTerm{叛教者}{traditors}的成员，就受到玷污。然而，如果你们藐视我们这一切规劝，我们绝不后悔我们曾渴望以和平有序的方式行事。主必要这样为祂的教会辩护，使得你们反倒要后悔藐视了我们谦卑的和解尝试。\par

\SourceRecord{Translated by J.G. Cunningham.；Nicene and Post-Nicene Fathers, First Series；Vol. 1.；Edited by Philip Schaff.；Buffalo, NY: Christian Literature Publishing Co.,；1887.；<http://www.newadvent.org/fathers/1102088.htm>.}

% structure: p0001 source=h1 target=section reason=page-title

\section{第89号信函（公元406年）}

\Emph{致\Person{费斯图斯}，我深爱的阁下，我可敬而值得赞誉的儿子，\Person{奥斯定}在主内致候。}\par

1. 如果人们为了错误、不可原谅的分裂以及在各方面已被充分反驳的虚妄，竟如此胆大妄为，坚持不懈地公然攻击并威胁那寻求他们救恩的天主教会，那么，那些持守基督徒和平与合一的真理——这真理甚至令那些声称否认或试图抗拒它的人也折服——的人，岂不更应理直气壮地、不知疲倦地、用力地辛勤工作，不仅为保卫那些已经是天主教徒的人，也为纠正那些尚未在教会内的人吗？因为，如果顽固不化旨在拥有并施展不屈不挠的力量，那么，那将持久不懈、不辞辛劳的努力奉献于一项既蒙天主悦纳又无疑为智者判断所嘉许的事业的恒心，该有何等巨大的力量啊！\par

2. 再者，还有什么比这更为可悲的执拗例证呢？就是有人不但不肯因纠正其恶行而谦卑自下，反而为自己的行为要求褒奖，正如\OriginalTerm{多纳图派}{Donatists}所做的那样，他们自诩遭受迫害；这若不是出于不可思议的盲目，不认识真理，便是出于不可原谅的冲动，佯装不认识真理——他们不知道，人成为\OriginalTerm{殉道者}{martyrs}，并非因其受苦受难的程度，而是因为他们为之受苦的原因。即便我只是在反驳那些仅陷入错误黑暗、并因此遭受全然应得的惩罚、且未曾以疯狂暴力攻击任何人的人，我也会这样说。但是，对于那些致命顽固到只有因害怕损失才受约束、只有通过流放才得知教会（正如预言的那样）何其广布四方——他们宁愿攻击教会也不愿承认它——的人，我又能说些什么呢？而如果在他们以这种最温和的纪律所遭受的事与他们在肆意狂暴中所做的事之间作一比较，谁看不出哪一方更真正配称为迫害者呢？不，即使是邪恶的子女不施暴力，他们那种放纵的生活方式，也比父母用与爱心深厚的程度相称的严厉、毫无掩饰地设法强迫他们正直为人，对慈爱的父母所造成的伤害严重得多。\par

3. 在公开档案中存在着最有力的证据，您若愿意，可以阅读；我更恳切地请求并劝勉您阅读这些档案，它们证明，他们那些起初脱离教会和平的前辈，确实曾自行经由当时的行省总督\Person{阿努利努斯}，向皇帝控告\Person{凯西里安}。倘若他们在最初的审判中得胜，\Person{凯西里安}在皇帝手中将遭受的，岂不就是当他们败诉时皇帝加于他们的同样待遇吗？但实际上，如果他们这些控告者获胜，\Person{凯西里安}及其同僚被逐出教区，或者，由于坚持其阴谋而在民事法庭上败诉后，抗拒不从而招致更严厉的惩罚（因为皇帝的责罚不可能不惩罚那些在民事法庭败诉者的抗拒），那么，他们当时就会宣扬，皇帝的明智措施及其为教会福祉的殷切关怀，是值得一切赞美的。但现在，他们自己败诉了，完全不能证明他们提出的指控，便因他们的不义而遭受任何惩罚，就称之为迫害；不仅毫不约束其邪恶暴力，还要求被尊为殉道者：仿佛天主教基督徒皇帝们在处置他们这些极其顽固的邪恶时，除了遵循\Person{君士坦丁}的先例之外，还遵循了别的先例似的——他们当初就是自愿向\Person{君士坦丁}提出上诉，作为控告\Person{凯西里安}的人，而且他们如此看重他的权威，胜于所有海外的众主教，甚至宁可将这件教会争端提交给他，而非那些主教。他们又曾向\Person{君士坦丁}抗议，反对他指派在\Place{罗马}审理此案的主教们所作的第一次判决；并且还向\Person{君士坦丁}上诉，反对在\Place{阿尔勒}的主教们所作的第二次判决：然而，最终当他们被\Person{君士坦丁}本人的裁决击败时，他们却顽固到底，不肯改变。我想，即使是魔鬼本人，若在一个他自愿选择向其上诉的法官的权威之下一再被驳倒，也不致厚颜无耻地坚持这种论据。\par

4. 可能有人会说，这些是人间的法庭，它们可能被哄骗、被误导或被贿赂。那么，为什么基督徒世界被诽谤和烙上某种罪行的印记，而这罪行乃是某些据称向迫害者交出圣书的人所被指控的呢？因为基督徒世界既不可能、也无义务不去相信原告所选择的法官，反倒去相信那些被这些法官判其败诉的原告。这些法官无论公正与否，都要为他们的裁决向天主负责；但遍布全世界的教会究竟做了什么，竟使\OriginalTerm{多纳徒派}{Donatists}认为必须给教会重新受洗，而其唯一理由不过是因为教会在这件自己无法断明真相的案件中，认为自己理应相信那些能够正确判断的人，而非那些虽在世俗法庭败诉却仍不肯屈服的人？啊，这是对所有民族的严厉控诉！天主曾应许他们要在\Person{亚巴郎}的后裔中蒙受祝福，如今祂已实现了这应许！当他们同声质问：\Quote{你们为什么想要使我们重新受洗？}时，得到的回答是：\Quote{因为你们不知道在非洲有哪些人犯了交出圣书的罪；并且你们既然无知，便接受了对案件作出裁决的法官的证词，认为它比提出指控者的证词更可信。}没有人该为他人的罪行受责备；那么，全世界与某个在非洲可能犯下的罪有何相干呢？没有人该为他毫不知情的罪行受责备；而在这件事上，全世界怎能知道究竟有罪的是法官还是被定罪的一方呢？你们明智的人，请判断我所说的。这就是异端者的正义：多纳徒派未经审讯就谴责全世界，因为全世界不谴责一件自己所不知道的罪行。但对于世界来说，有天主的应许就足够了，并在自己身上看到先知很久以前所预言的事得以应验，并借着那识别其君王\Person{基督}的同一\WorkTitle{圣经}来识别教会。因为正如在这些经书中论及\Person{基督}的事已预先记录，而我们读到在福音历史中已在祂身上应验，同样，论及教会的事也已预先记录，而我们现在看到在全世界已应验。\par

5. 或许有些有思想的人会因他们惯常对圣洗的说法而感到不安，即圣洗只有由一位义人施行时才是\Person{基督}的真正圣洗，然而，有关这个问题，基督徒世界持守的是最明显不过的福音真理，正如\Person{若望}的话所教导的：\BibleQuote{那派遣我来以水施洗的，给我说：你看见圣神降下，停在谁身上，谁就是那要以圣神施洗的人。}\BibleRef{若 1:33} 因此，教会平静地拒绝将希望寄托于人，免得落入\WorkTitle{圣经}所宣判的诅咒中：\BibleQuote{凡信赖世人的人，是可咒骂的，}\BibleRef{耶 17:5} 却将希望寄托于\Person{基督}，祂虽取了奴仆的形体，却不失天主的形体，论到祂说：\BibleQuote{就是那施洗的一位。} 因此，无论那施行此圣事的人是谁，无论他担当什么职务，施洗的并不是他——那是鸽子降在其身上的那位的工作。这些愚蠢的见解给多纳徒派带来的荒谬是如此之大，以致他们无法逃脱。因为当他们承认，若属他们一派的一位有罪的人施洗，但罪隐藏不露，圣洗仍然有效且真实时，我们问他们：在这情况下，是谁在施洗？他们只能回答说：天主；因为他们不能肯定一个有罪的人（比如犯奸淫者）能圣化任何人。那么，如果圣洗由一位公认的义人施行时，他圣化那受洗者；但若由一位邪恶的人施行，而其邪恶被掩盖时，圣化人的不是他，而是天主。那么，准备受洗的人应当宁愿由暗中邪恶的人施洗，而不愿由明显良善的人施洗，因为天主远比任何义人更能有效地圣化人。如果说准备受洗的人应当希望由伪善的奸淫者施洗，而不愿由公认贞洁的人施洗，这显然是荒谬的，那么清楚可见，无论那施行礼节的人是谁，圣洗都是有效的，因为是那位鸽子降在其身上的祂亲自施洗。\par

6. 虽然如此明显的真理应当在人们的耳中和心中产生印象，但邪恶习俗的漩涡吞没了一些人，以致他们宁愿抗拒一切权威和论证，也不愿屈服。他们的抗拒有两种——或是激烈的愤怒，或是被动的麻木。那么，当教会怀着母亲的焦虑寻求所有人的救恩，而又被一些人的狂躁和另一些人的昏睡所困扰时，她必须采用什么治疗方法呢？她轻视或放弃任何可能促进他们康复的方法，是正当的吗？是可能的吗？她必然被双方都视为烦扰，这正是因为她不是任何一方的敌人。因为狂躁的人不喜欢被束缚，昏睡的人不喜欢被惊醒；然而\OriginalTerm{爱德}{charity}的勤勉坚持不懈地约束前者，激励后者，出于对两者的爱。两者都被激怒，但两者都被爱着；两者在病中时，都怨恨这治疗是令人烦恼的；而当他们痊愈时，都对此表示感谢。\par

7. 而且，他们以为并夸耀我们按照他们原来的样子接纳他们进入教会，但事实并非如此。我们接纳的是已然彻底改变了的他们，因为他们直到不再是\OriginalTerm{异端者}{heretics}才开始成为\OriginalTerm{天主教徒}{Catholics}。因为他们与我们共有的\OriginalTerm{圣事}{sacraments}，我们并不厌恶，因为这些圣事不是来自人，而是来自天主。必须从他们那里去除的是他们自己的错误，是他们邪恶地吸收的；不是圣事，这些圣事他们像我们一样领受了，他们持有并拥有——确实，他们拥有圣事是为他们自己的定罪，因为他们如此不相称地使用它们；然而，他们确实拥有它们。因此，当他们离弃错误，在他们身上纠正了\OriginalTerm{分裂}{schism}的乖僻时，他们就从异端过渡到教会的平安，这平安他们以前未曾拥有，而没有这平安，他们所有的一切只是对他们有害。然而，如果在这种过渡中他们并不真诚，这不是我们，而是天主来判断的事。然而，有些人因为出于恐惧而归向我们，被怀疑不真诚，但他们在后来的一些考验中却显得如此\OriginalTerm{信德}{faith}坚固，以至于超过了那些原本就是天主教徒的人。因此，不要说采用强硬手段毫无成效。因为当顽固习惯的壁垒被对人的权威的恐惧所攻破时，这还不是全部，因为同时，信德得以坚固，理智得以说服，这是藉着天主的权威和论证。\par

8. 既然如此，请主教阁下知晓，您在\Place{希波}地区的人仍是\OriginalTerm{多纳徒派}{Donatists}，您的信件对他们没有产生任何影响。信件未能打动他们的原因，我无需多写；但请派遣一个人，或是您的仆人，或是您的朋友，其忠诚您可以托付此项使命，让他不要先去找他们，而是先到我们这里来，不让他们知道；当他与我们仔细商议当行何事为佳之后，便靠着\OriginalTerm{主}{Lord}的助祐去行。因为我们采取这些措施，不仅是为了他们的福祉，也是为了我们自己那些已成为\OriginalTerm{天主教徒}{Catholics}的人，这些多纳徒派的邻近对他们极其危险，以至于我们不能视之为小事。\par

我本可以写得更简短；但我希望您能收到我的一封信，藉此您不仅可以了解我关切的原因，而且如果有人劝您不要热心致力于纠正属于您的人，并反对我们希望您这样做，您也可以以这封信作为答复。如果我这样做是多此一举，因为您自己已经领会或思考了这些原则，或者如果我给如此忙碌于公共事务的您强加如此长的一封信而成为您的负担，我请求您宽恕我。我只恳求您不要轻视我向您提出的请求和对您的要求。愿\OriginalTerm{天主}{God}的仁慈护佑您！\par

\SourceRecord{Translated by J.G. Cunningham.；Nicene and Post-Nicene Fathers, First Series；Vol. 1.；Edited by Philip Schaff.；Buffalo, NY: Christian Literature Publishing Co.,；1887.；<http://www.newadvent.org/fathers/1102089.htm>.}

% structure: p0001 source=h1 target=section reason=page-title

\section{信函九十（主后四〇八年）}

\Emph{致我尊贵的主与兄弟、最堪尊敬的\Person{奥古斯丁}主教，\Person{内克塔里乌斯}敬祝安好。}\par

我不会多谈人们对自己故土的热爱之情，因为您深知此情。它是唯一比亲情更强烈的感情。倘若正直的心灵可以不受其约束，在一定的期限或时候之后便可解脱，那么我如今早就该从它所加给我的负担中获得了豁免。然而，因着对我们的国家所怀的爱戴与感激之情与日俱增，而且人越接近生命的终点，就越发热烈地渴望使自己的祖国处于安全与繁荣的境地，所以当我开始写此信时，我感到欣慰，因我正在向一位精通各类学问、因而能够体会我的感受的人说话。\par

在\Place{卡拉马}殖民地，有许多事物使我的爱理所当然地系于此地。我生于此地，并且（在别人看来）我为这个社群立过重大的功劳。然而，我至尊贵、最堪尊敬的主，如今这城因着其公民一次严重的过犯而遭了殃，倘若按国家法律的严苛待我们，此过犯必须受到极为严厉的惩处。但一位主教受着另一种法律的引导。他的职责是促进人们的福祉，在任何案件中只关心当事人的利益，并从全能的天主手中为他人求得罪过的宽赦。因此，我以所有可能的迫切之情恳求您，务必做到：若此事要被起诉，要使无罪者受到保护，并在无辜者与行恶者之间加以区分。这正是您自己本性的情感所要求的，我祈求您应允。对骚乱造成的损失进行估定赔偿，可以轻易征收。我们仅恳请不要施行严厉的报复。愿您更充分地享受天主的恩宠，我尊贵的主、最堪尊敬的兄弟。\par

\SourceRecord{Translated by J.G. Cunningham.；Nicene and Post-Nicene Fathers, First Series；Vol. 1.；Edited by Philip Schaff.；Buffalo, NY: Christian Literature Publishing Co.,；1887.；<http://www.newadvent.org/fathers/1102090.htm>.}

% structure: p0001 source=h1 target=section reason=page-title

\section{第九十一封信（公元408年）}

\Emph{致我尊贵的阁下及理当尊敬的弟兄\Person{内克塔里乌斯}，\Person{奥斯定}致以问候。}\par

1. 我并不惊讶，虽然您的四肢因年迈而冰冷，您心中仍燃烧着爱国热忱。我钦佩这一点，并且非但不悲伤，反而欣喜地得知您不仅铭记，而且以您的生活和实践体现了这一格言：正直之人对国家所负有的关怀与服务的责任，无论从程度还是时间上都是没有限度的。因此，我们渴望通过神圣的爱，使您列入一个更高贵更崇高的国家的事务中；我们按自己的能力，在那些我们希望敦促他们以那国家为自己家乡的人们中间，遭遇危险与劳苦时，赖以得到支撑的正是这爱。啊，倘若您成了那国家的公民，您便会认为，为了她那些尚在地上作旅客的一小部分公民的利益，您的努力在程度和时间上都应是没有限度的！这才是更好的忠诚，因为您是在回应一个更好的国家的要求；而且，如果您决定在有生之年为她福祉的辛劳永不止息，那么您必将在她永恒的平安中得享无尽喜乐的赏报。\par

2. 但在实现这点之前——您能够获得，或者甚至现在就在极明智地考虑应当获得，那您父亲先您而去的国家，并非毫无希望——我说，在实现这点之前，请您原谅我们，若我们为了那渴望永远不离弃的国家，而给您渴望在荣誉与繁荣的盛放中离开的国家，带来一些痛苦。至于在您国家如此盛开的鲜花，如果我们与如您般智慧的人讨论这主题，我们毫不怀疑您会轻易被说服，或者更确切地说，您自己会立即领悟，一个共和国当以何种方式繁荣。你们最杰出的诗人曾歌唱意大利的某几种鲜花；但在您的国家，我们凭经验得知的，与其说是英雄们如何使之繁花盛开，不如说是战争的武器如何闪耀其中。不，我应当写它如何燃烧而非如何闪耀；而且与其说战争的武器，不如说纵火者的火焰。如果如此重大的罪行竟不受惩罚，没有恶棍们当受的任何谴责，您认为您将让您的国家停留于荣誉与繁荣的盛放中吗？哦盛开的鲜花，不结果实，反生荆棘！现在请想想，您是更愿看到您的国家因其居民的虔诚而繁荣，还是因其逃避罪恶的惩罚而繁荣；因其风气的纠正而繁荣，还是因有罪不罚所助长的恶行而繁荣。我说，请比较这些事，并判断，您是否比我们更爱您的国家；它的繁荣与荣誉是更真实且更热切地由您还是由我们所寻求。\par

3. 请稍思考一下那些书，\WorkTitle{论共和国}，您从中汲取了那最忠诚公民的情操，即正直之人对国家所负有的关怀与服务的责任，无论从程度还是时间上都是没有限度的。我请求您思考它们，并注意它们何等大力赞扬了节俭、自控、婚姻忠实，以及那些贞洁、可敬且正直的风尚，这些风尚在任何城邦的普及，使该城邦堪称为繁荣。如今，遍布全世界的各教会，犹如神圣的公共教育场所，在其中谆谆教诲并学习这健全的道德，并且首要的是，在其中教导人们当敬拜真实信实的天主，祂不仅命令人们践行这一切，更赐予恩宠去履行，藉此，人的灵魂得以装备并适于与天主共融，适于居住于永恒的天国。为此，祂既预言又命令了推倒世上众多假神的偶像。因为没有什么比效法异教著作中所描述和颂扬的那些神灵，更有效地使人在行为上堕落，并且不适合作社会的好成员。\par

4. 事实上，那些最博学的人（顺便说一句，他们对于此世一个共和国或共同体的\OriginalTerm{理想典范}{beau ideal}，更多是在私下讨论中探究或描述，而非在公共措施中建立和实施），习惯于为教育青年提出他们所尊崇和赞许的人作为模范，而非他们的神所立的榜样。毫无疑问，\Person{泰伦斯}笔下的年轻人，他看到墙上有一幅描绘神王\Person{朱庇特}淫荡行为的图画，便因这如此高的权威给邪恶的鼓励，燃起主导他的情欲之火；如果他选择效法\Person{加图}而非\Person{朱庇特}，他就不会陷入这样的欲望，也不会犯下如此可耻的罪行；但在被迫于庙宇中敬拜\Person{朱庇特}而非\Person{加图}的情况下，他又如何能作出这样的选择呢？或许有人会说，我们不该从一出喜剧中引证，以使不敬神者的淫荡和亵渎的迷信蒙羞。但请阅读或回想，上面提到的那些书里何等明智地论证说，喜剧的风格和情节若非与欣赏它们的人的风尚和谐一致，便绝不会得到公众的认可；因此，借由这些在他们所属的共和国中最显赫尊贵、且正辩论完美共和国条件之人的权威，我们的论点得以确立，即最堕落的人若效法他们的神，就可能变得更坏——这些神当然不是真的，而是虚假和捏造的。\par

5. 你或许会回答说，那些古时所写、关于诸神生活与品行的种种，智者应当理解并解释为与字面意义大相径庭。其实，就在几天之前，我们听说这样有益的阐释如今正在庙宇中向聚集的民众宣读。请告诉我，人类竟如此盲目于真理，以致不能觉察如此明白可触的事实吗？画家、铸工、锻工、雕刻师、作家、演员、歌者和舞者，以他们的技艺在那么多地方将\Person{朱庇特}公开展示其丑恶可耻的淫行，而在他的\Place{卡比托利欧}，至少崇拜他的人若能读到他亲自颁布的、禁止此类罪行的法令，那该是何等重要！正因为缺少了这样的禁令，那些集羞耻与亵渎之大成的怪物，才会被民众热切推崇，在庙宇中受崇拜，在剧场中受人嘲笑；为了给它们献祭，甚至连穷人都被夺去牲畜；为了让演员能以姿态与舞蹈再现其臭名昭著的劣迹，富人们挥霍家财——实行这些事的城邦，还能说是繁荣吗？它们繁荣所绽放的花朵并非生自肥沃的土壤，也非生自丰裕的德行；它们那位相称的母亲乃是女神\Person{福罗拉}，庆祝她的戏剧节庆时，那种放荡是如此极度放纵、不知羞耻，以至任何人都能由此推断出，那需要以——不是鸟，不是四足兽，甚至也不是人命——而是（哦，何等更大的恶德！）人的贞洁与德性在她的祭台上沦为祭品，方能平息的那位，到底是个怎样的\OriginalTerm{邪魔}{demon}。\par

6. 我说这些话，是因为你曾写道，你愈近生命的终点，就愈渴望让你的国家处于安全与繁荣之中。祛除这一切虚妄与愚行，使人们归向天主的真实崇拜，归向贞洁与虔敬的生活吧！那时，你将看见你的国家繁荣，不是在愚人虚妄的见解中，而是在智者健全的判断里；那时，你在地上的家园将成为那\OriginalTerm{天乡}{Fatherland}的一部分，我们进入那天乡，不是藉着血肉，而是藉着信德；在那里，所有天主圣洁忠信的仆人，在时间的寒冬中辛劳已毕，将要在永恒的盛夏中绽放。因此，我们决心，一方面不放弃基督徒的温和，另一方面也不在你们的城市中留下那将成为他人竞相效尤的、极为有害的榜样。我们信赖天主的助佑，能使此事务成功，只要祂对作恶者的义怒不致使祂收回祝福。因为显然，我们所愿保持的温和，以及我们将努力不带忿怒而施行的纪律，都可能受阻，倘若天主以祂隐秘的上智安排另作决定，或定意以更严厉的惩罚来惩治此番大恶，或出于更大的不悦，让这罪在世上不受惩罚，使犯罪者既不受责备，也不得改正。\par

7. 你依自己的判断，列出了主教当遵循的原则；你述说你的城市因公民的严重恶行而惨遭灾祸，若按俗世法律的严厉处治，这些行为必受极严厉的惩罚，之后你又补充说：\Quote{但主教受另一种法律引导；他的职责是促进人的福利，在一切案件中只关心相关者的利益，并为他人向全能的天主求得罪过的宽赦。}我们竭力确保的正是这一点：没有人因我们或任何因我们的介入而受影响的人，遭受过度的惩罚；我们寻求促进人类真正的福利，这福利在于行善的福祉，而非在于作恶而免罚的保证。我们也热切地，不仅为自己，也为他人寻求罪的宽赦：然而，我们只能为那些因受纠正而弃绝罪行的人，求得此宽赦。你还说：\Quote{我以一切可能的急切恳求你，务要确保，若此事须提出诉讼，要保护无辜，并区分清白者与作恶者。}\par

8. 请听我简述所发生之事，并由你自己划清无辜者与有罪者的界限。有人公然藐视最新颁布的法律，在异教节日\OriginalTerm{六月朔日}{calends of June}举行某些不敬的仪式，无人干涉加以禁止，且大胆放肆至极，以至于一伙极为无礼的暴民沿着圣堂所在的街道游行，并在圣堂前跳舞——这一暴行甚至在\Person{尤利安}时代也未曾发生。当圣职人员力图阻止这种最不合法且侮辱性的行为时，圣堂遭到投石攻击。大约八天后，主教已提请当局注意有关此事的众所周知的法律，而当局正准备依法执行时，圣堂第二次遭到投石攻击。次日，我们的民众希望公开在法庭上提出他们认为必要的控告以震慑这些恶棍，却被剥夺了作为公民的权利。当天，下了一场冰雹，好让他们若非因人的缘故，至少因天主的大能而惧怕，以此报应他们对圣堂的投石如雨；但这件事一过去，他们第三次用石头攻击，最后企图用火烧毁建筑和里面的人：一个天主的仆人因迷路遇见他们，被当场杀死，其余的人尽己所能逃跑或躲藏起来；主教则挤入一个难以容身的缝隙中，躬身蜷缩，同时听见那些暴徒搜寻他、要杀害他的声音，并听到他们因他逃脱而在这次严重暴行中主要目的未能得逞的懊恼。这些事情从大约第十时辰持续到深夜。那些有权势可能影响暴民的人，既未试图抵抗，也未设法营救。唯一出手干预的是一位外乡人，因他的努力，许多天主的仆人从试图杀害他们的人手中被救出，大量财产也被强行从劫掠者那里夺回：由此可见，如果市民，尤其是领头人物，从一开始或在事态开始后加以禁止，这些暴乱行为原本极易完全防止或制止。\par

9. 因此，你在那个团体中无法区分无辜者与有罪者，因为所有人都有罪；但你或可区分罪责的轻重。那些因恐惧，尤其是害怕冒犯那些他们知道在该团体中具有领导影响力且仇视教会的人，而不敢援助基督徒的人，其过错相对较小；所有同意这些暴行的人都有罪，即便他们既未施行暴行，也未教唆他人犯罪；施行不义者更有罪，而教唆他们犯罪者罪责最深。然而，我们假定教唆他人犯下这些罪行仅是疑点，而非确知，我们也不去查究那些除非拷问证人否则无从查明的事。我们也宽恕那些因恐惧而认为与其公开得罪教会的有力仇敌，不如暗中向天主祈求保护主教和祂的其他仆人的人。你能给出什么理由，主张那些剩下的人不应受到任何纠正和约束？你真的认为，如此残忍暴怒的事件应该昭告世人，且不受惩罚吗？我们并不想通过报复性的惩罚来满足我们的愤怒，而是以真正慈悲的精神为将来作安排。恶人有一些方面可受惩罚，且可由基督徒以慈悲的方式惩罚，以促进他们自身的利益和福祉。因为他们有三样东西：身体的生命与健康、维持生命的手段，以及过邪恶生活的手段和机会。让前两样在那些悔改其罪孽的人身上保持不受触动：这是我们渴望的，我们也为此不遗余力去确保。至于第三样，天主若乐意，将以祂的大慈大悲施以惩罚，把它像腐烂或患病部分一样，必须用修枝刀切除。然而，如果祂乐意超出此限，或不许惩罚达到如此地步，这一更高且无疑更正义的谋虑的原因在于祂：我们的本分是依照所赐予我们的光明，尽心竭力并运用影响力，恳求祂悦纳我们为众人最大益处所做的努力，并祈求祂不许我们做任何事，即那比我们更洞悉万有的祂所知道对我们及祂的教会均无益的事。\par

10. 最近我到\Place{卡拉马}去，为的是在这般沉重的忧伤下，可以安慰我们当中沮丧的人，或平息义愤者的怒火。我运用了自己在基督徒中的全部影响力，劝说他们去做我当时认为他们应尽的本分。之后，应他们自己的请求，我也接见了\OriginalTerm{外教人}{Pagans}——他们是这一切祸患的根源和起因——为要劝诫他们：若是有智慧，应当做什么，这不仅为消除当前的忧虑，更是为获得\OriginalTerm{永远的救恩}{everlasting salvation}。他们听了我许多话，也向我提出许多请求；但是，我绝不愿作这样的仆人：那些不在我主面前谦卑自下、向祂祈求的人，竟来向我请愿，我却以此为乐。以你的敏捷才智，你必立刻看出，我们的目标必定是：在持守基督徒的温和与节制的同时，要如此行事，好叫我们或是使人不敢效法他们的乖僻，或是使人因他们的受教而获益，从而渴望效法。至于所遭受的损失，这损失或由基督徒承受，或藉着弟兄们的帮助而得到补救。我们所关心的是争取灵魂，即使冒生命危险，我们也迫不及待地要确保得到他们；我们的愿望是，在你的地区我们能取得更大的成功，而在其他地区，我们也不会因你的榜样所及的影响而受到阻碍。愿天主因祂的仁慈，赐予我们为你的救恩而欢欣！\par

\SourceRecord{Translated by J.G. Cunningham.；Nicene and Post-Nicene Fathers, First Series；Vol. 1.；Edited by Philip Schaff.；Buffalo, NY: Christian Literature Publishing Co.,；1887.；<http://www.newadvent.org/fathers/1102091.htm>.}

% structure: p0001 source=h1 target=section reason=page-title

\section{第九十二封书信（主后408年）}

\Emph{致高贵而正直的夫人\Person{意塔利卡}，一位在基督的爱内堪受尊敬的女士，主教\Person{奥斯定}在主内致以问候。}\par

1. 我从你的信以及送信人的陈述中得知，你恳切地请求我回一封信，相信从其中可以获得极大的安慰。你从我的信中得到什么，应由你自己判断；我至少认为，我不应拒绝也不应延迟满足你的请求。但愿你自己信德和望德安慰你，以及那借著圣神倾注在虔诚人心中的爱德，\BibleRef{罗 5:5}，我们如今已得著一部分作为全体的凭据，好使我们学会渴望它的圆满成全。因为你不应以为自己是孤独的，因为你借着信德有基督居住在你心中；你也不应像那些没有望德的外邦人一样忧伤，因为我们对于那些朋友——他们并未丧失，只是比我们先蒙召回到那我们将要追随他们而去的家乡——怀有望德，这望德建立在最确实的应许上，即我们将从今生过渡到那另一生命中；在那生命中，他们对我们将更加亲爱，因为我们将更好地认识他们；而在那种生命中，我们喜爱他们的快乐将不会混杂任何分离的恐惧。\par

2. 你已故的丈夫——因他的逝世你现在成了寡妇——你确实很了解他，但他自己比你自己更了解他。怎么可能这样呢？因为你看见了他的面容，而他自己并没有看到；除非我们对自己内在的认识更为确实，因为没有人\BibleQuote{除了在人内的心神，没有人知道人的事}？\BibleRef{格前 2:11}，但当主来临时，\BibleQuote{祂要照亮黑暗中的隐秘，并要显明人心的意念}，\BibleRef{格前 4:5}，那时，任何人在近人面前将没有任何隐藏的事；也没有任何事是某人可以告诉朋友却要向陌生人隐瞒的，因为那里不再有陌生人。有谁能用言语描述那光明的性质和伟大，借著那光明，所有现在隐藏在人心里的事物都将显露出来？我们微弱的官能又怎能接近它呢？诚然，那光明就是天主自己，因为\BibleQuote{天主是光，在祂内没有一丝黑暗}；\BibleRef{若一 1:5}，但祂是净化了的心灵之光，而不是这肉眼的光。那时的心灵将能看见那光，而现在却不能。\par

3. 但这肉眼现在不能，将来也不能看见那光。因为凡能被肉眼看见的，必定存在于某个地方，不能同时全部地存在于各处，而是以较小部分占据较小空间，以较大部分占据较大空间。天主却并非如此，祂是无形的、不朽的，\BibleQuote{独享不死不灭，住於人不能接近的光中，没有人见过，也不能见}；\BibleRef{弟前 6:16}，因为人不能通过那用来观看有形事物的肉体的器官看见祂。因为如果祂对圣人们的心灵也是不可接近的，那就不会说：\BibleQuote{他们瞻仰祂，就蒙光照}\EditorialNote{奥斯定译作\Quote{\InnerQuote{靠近祂，并蒙光照}}}；倘若祂对圣人们的心灵是不可见的，也就不会说：\BibleQuote{我们要看见祂实在怎样}：因为请考虑那部\BibleBook{若望}{一书}中的整个上下文：\BibleQuote{可爱的，现在我们是天主的子女，但我们将来如何，还没有显明；可是我们知道：当祂显现时，我们必要相似祂，因为我们要看见祂实在怎样。}\BibleRef{若一 3:2}，因此，我们将按照我们相似祂的尺度看见祂；因为现今我们不看见祂的尺度，正是我们与祂不相像的尺度。所以我们将通过我们与祂相似之处看见祂。但是谁会如此愚蠢，竟敢断言我们现今或将来在身体上相似天主呢？因此所说的相似是在内在的人身上，\BibleQuote{这内在的人正在知识上更新，依照创造他者的肖像}；\BibleRef{哥 3:10}，我们越在知识和对祂的爱上长进，就会越相似祂；因为\BibleQuote{纵然我们外在的人日渐损坏，但我们内在的人却日日更新}，\BibleRef{格后 4:6}，然而，无论人在今生能有多大的长进，他仍远不及那为看见天主所需的完美相似，正如宗徒所说，\BibleQuote{面对面}地看见\BibleRef{格后 13:12}。如果这些话我们理解为肉体的脸，那么就可以推论说天主有像我们一样的脸，在我们将来面对面看见祂时，我们的脸和祂的脸之间必定有一个空间间隔。如果有一个空间间隔，那就预设了肢体的界限和一定形状，以及其他一些荒谬可笑、甚至连想都是亵渎的事。那些属血气的人，\BibleQuote{他们不接受天主圣神的事}，\BibleRef{格前 2:14}，就被这些极端空洞的妄想所欺骗。\par

4. 因为我听说，那些如此愚拙说话的人中，有人断定，我们现今是以心灵看见天主，将来却要以身体看见他；他们甚至说，恶人也要以同样方式看见他。请注意他们如何每况愈下，因他们的愚拙言论未受惩罚，便肆无忌惮地信口开河，既不惧怕，也不羞愧。起初他们说，基督只将他自己的肉身赋予了以肉眼看见天主的这种能力；随后他们又加上了，众圣人在复活中重新获得身体后，也都将以同样方式看见天主；如今他们竟承认，恶人同样能做到这一点。好吧，就让他们按自己的意愿将能力赐予任何人吧：因为谁能反对人把自己所有的东西送出去呢？\BibleQuote{他说谎话，是出于他自己。} \BibleRef{若 8:44} 然而，你们与所有持守健全教义的人，切不可如此擅自从你们自己那里拿出这些错误；但你们读到：\BibleQuote{心里洁净的人是有福的，因为他们要看见天主，} \BibleRef{玛 5:8} 就要从中学得，不敬虔的人决不能看见他：因为不敬虔的人既不蒙福，心里也不洁净。再者，你们读到：\BibleQuote{我们现在是借着镜子观看，模糊不清，到那时，就要面对面地观看了，} \BibleRef{格前 13:12} 就要从中学得，我们将要面对面看见他，就如同现今我们借着镜子模糊不清地看见他一样。在这两种情形中，看见天主都属于内在的人，无论是当我们在这今世的旅途中，仍凭信德行事，此时我们使用镜子和\OriginalTerm{谜语}{\GreekText{αἴνιγμα}}，还是当我们在那家乡的国度里，我们将凭着直观得以看见，而\Quote{面对面}一词所指的正是这种观看。\par

5. 愿那沉溺于肉欲幻想的肉躯听听这话：\BibleQuote{天主是神，朝拜他的人，应当以心神以真理去朝拜他。} \BibleRef{若 4:24} 若朝拜他应当如此，更何况看见他呢！因为当天主不许可人以肉身的方式朝拜他时，谁敢断言能以肉身的方式看见\OriginalTerm{天主的性体}{divine essence}呢？然而，他们自以为十分敏锐，提出这样一个问题迫使我们回答：基督能否赋予他的肉身，使他能以肉眼看见圣父？如果我们回答不能，他们便会四处宣扬，说我们否定了天主的全能；反之，如果我们承认他能，他们便会声称我们的回答证实了他们的论点。那些主张肉身将被转化为天主性体，并将成为天主自身的人，他们的愚妄更值得原谅，因为他们想以此使那因本性与天主迥异而不肖似的肉身，获得看见天主的能力！然而，我相信他们已从信条中摒弃了这一错误，或许连听都不愿听。不过，倘若同样以刚才那个问题质问他们——即天主能否做到（即将我们的肉身转化为天主性体）这件事，他们会选择哪个选项？若他们回答\Quote{不能}，是否限制了天主的权能？若他们承认\Quote{能}，是否就要因此让步，承认这事必将发生？让他们用自己摆脱这个质问困境的方法，同样摆脱他们提出的另一个困境吧。再者，他们为何主张这种恩赐只应归于眼睛，而不归于基督的其他所有感官呢？难道天主将成为一种声音，以便被耳朵听见？成为一种气味，以便被嗅觉分辨？成为某种液体，以便被饮用？成为固体，以便被触摸？不，他们说。那么？我们回答：天主能否成为这样，还是不能？若他们说不能，为何要贬损天主的全能？若他们说能，但不愿如此，为何只偏爱眼睛，却吝于将同样的荣耀赐给基督的其他感官？难道他们可以随心所欲地扩展自己的愚妄吗？我们的做法要好得多：我们不为他们的愚妄设限，而是尽力防止他们陷入其中！\par

6. 为驳斥那种疯狂，可以提出许多论点。然而，若他们再侵扰你的耳鼓，请将这封信读给那些支持此谬误的人听，并不要以为写信尽可能回复我他们的答辩是过于繁重的工作。我还要补充说：我们的心是因信德得以净化的，因为看见天主的异象乃是应许给信德的酬报。如果这种看见天主是通过肉眼进行，那么，为领受它而操练的圣人们的灵魂便是徒劳了；不仅如此，一个怀有这种想法的灵魂，并非在自身内受操练，而是完全沉溺于肉身。因为，灵魂在何处能比在它期望借以看见天主的那种器官中，更坚决、更牢固地居住呢？这事有多大的害处，我宁愿让你自己的聪慧去察悟，而不费长篇大论去证明。\par

愿你的心常在主内得享安宁，尊贵而实至名归的女士，在基督的爱内堪受尊敬的女儿！请代我以你应得的敬意，问候你的儿子们，他们与你同为尊贵，并是我在主内深切爱慕的。\par

\SourceRecord{Translated by J.G. Cunningham.；Nicene and Post-Nicene Fathers, First Series；Vol. 1.；Edited by Philip Schaff.；Buffalo, NY: Christian Literature Publishing Co.,；1887.；<http://www.newadvent.org/fathers/1102092.htm>.}

% structure: p0001 source=h1 target=section reason=page-title

\section{第九十三封信（主历四〇八年）}

\Emph{致\Person{文森提乌斯}，我亲爱的弟兄，\Person{奥古斯丁}致以问候。}\par

% structure: p0003 source=h2 target=subsection reason=relative-heading-rank; base=h2

\subsection{第一章}

我收到了一封信，我相信是您写给我的：至少我并未认为这是不可信的，因为带信的人是我所认识的一位天主教基督徒，我想他不至于欺骗我。但即使这封信或许并非出自您手，我认为仍有必要给写信的人——无论他是谁——写一封回信。您知道我现在比起早年在\Place{迦太基}时，就是您的直接前任\Person{罗伽图斯}还在世的时候，更加渴望安息，并且热切地寻求它。然而，我们却因\OriginalTerm{多纳徒派}{Donatists}而无法得到这种安息；在我看来，藉由天主所设立的权力来压制和纠正他们，并非徒劳之举。因为我们已经为许多人的纠正而欢欣，他们以极大的真诚持守并捍卫着天主教会的合一，且为从先前的错误中被解救出来而如此喜悦，以至于我们满怀感恩和喜乐地赞赏他们。然而，同样是这些人，在某种难以言喻的习惯束缚下，若非受到这种恐惧的震动，将心思转向对真理的认真研究，是绝不会想到转向更好的状态的；他们担心，如果徒劳无益地、在人手下遭受苦难，不是为了正义，而是为了自己的固执和傲慢，那么日后从天主手中所得的，除了惩罚——就是那些藐视祂如此温和的劝诫和慈父般的纠正的恶人应得的惩罚——之外，再无其他；他们经过这样的反思，变得易于受教，于是发现，并非在有害的或轻浮的人间寓言中，而是在神圣经卷的许诺中，看到了那个按着许诺扩展到万邦的普世教会：正如基于同一圣经中先知的见证，他们毫不犹豫地相信基督已被高举于诸天之上，尽管他们并未看见祂的荣耀。难道我应该对这些人的得救感到不悦，并将我的同工们从这种慈父般的勤勉中唤回吗？这种勤勉的结果是，我们看到许多人责备自己从前的盲目：因为他们看到，那些相信基督已被高举于诸天之上、虽然看不见祂，却否认祂的荣耀遍及全地、虽然可以看见的人，其实是盲目的；然而，先知却如此清晰地在一句话中同时包含了这两者：\BibleQuote{天主，愿祢崇高，遍及诸天，愿祢的荣耀普照尘寰！}\BibleRef{咏 57:5}。\par

因此，如果我们对那些以多种严重暴力与诡计严重扰乱我们和平与安宁的残忍敌人如此置之不理、一味容忍，以至于我们完全不谋划、不采取任何旨在惊慑并纠正他们的措施，那么我们就是以恶报恶了。因为如果有人看见自己的仇敌因高烧而精神错乱，正一头冲向毁灭，那么在此情形下，如果任由他继续冲下去，岂不是远比设法将他捉住捆绑起来，更加是以恶报恶吗？然而在那一刻，他在对方看来却似乎是极其讨厌、极像仇敌的，而事实上他已证明自己是最有益、最富同情心的；尽管毫无疑问，当健康恢复后，对方会以何等热烈的感激之情来回报他，其程度正与他当初在自己狂乱之时拒绝纵容他的感受相对应。哦，但愿我能向您展示，我们甚至从\OriginalTerm{西康塞里昂派}{Circumcelliones}中得到了多少人，他们如今已是公认的天主教徒，谴责自己从前的生活，以及那可悲的妄念；在那妄念下，他们以为自己在为天主的教会做事，实际上却是在一种不安分的鲁莽冲动下行事；然而，倘若他们未曾如同心神错乱的人一般，被你们所不喜的那些法律的绳索所捆绑，他们就绝不会达到这种健全的判断力！至于另一种极其严重的病症——即那些虽然并无导致暴力行为的胆量，却被一种根深蒂固的迟钝所压制，会对我们说：\Quote{你们所肯定的都是真的，无可反驳；但要我们放弃从祖先那里因传统而领受下来的，实在很难}的人——为什么不应当藉由一种给他们在世俗事物上带来不便的法律，以有益的方式震动他们，好使他们从昏睡中醒来，醒悟到在教会的合一中才能找到的救恩呢？他们中有多少人，如今与我们一同喜乐，痛切地述说过去那毁灭性的道路曾如何压迫他们，并承认我们本就有责任令他们烦恼，以免他们在昏睡习惯的疾病中，如同在致命的沉睡中灭亡！\par

3. 你会说，这些补救对某些人毫无用处。难道就因为有些人的疾病无法治愈，便应放弃治疗之术吗？你所顾及的，只是那些刚愎自用、连这样的纠正都不肯接受的人。论到这种人，经上记载说：\BibleQuote{我徒然责打了你们的儿女，他们不肯受教：}\BibleRef{耶 2:30} 不过我仍相信，先知所提及的那些人，是出于爱而受责打，而非出于恨。然而，你也应当考虑那些为数极多、我们因他们的救恩而喜乐的人。因为假使他们只是受惊惧，而未得教导，这或许会被视为一种不可原谅的苛政。再者，假使他们只受教导，而未受惊惧，他们便会因积习日久而心硬，更难被说服去接受救恩之道：我们认识的许多人，当我们向他们提出论证，并以天主的话为证使之看清真理后，他们便回答我们说，他们渴望进入天主教会的共融，但畏惧那些无赖之徒的暴力，怕招致他们的敌意；其实，为了正义和永生的缘故，他们本当全然蔑视这种暴力。然而，我们决不可将这些人的软弱视为无望，而应扶持他们，直到他们获得更大的力量。我们也不可忘记，主曾对那时尚软弱的伯多禄说：\BibleQuote{你现在不能跟随我，但以后你要跟随我。}\BibleRef{若 13:36} 然而，当有益的教导与引发救恩恐惧的方法相结合，以致不仅真理之光驱散错谬的黑暗，同时畏惧的力量也击碎恶习的桎梏，那时，正如我所说，我们便为许多人的得救而欢欣；他们与我们一同赞颂天主，感谢祂，因为祂藉着祂盟约的实现——其中祂应许世上的君王要服事基督——如此医治了有病的，使软弱的恢复健康。\par

% structure: p0007 source=h2 target=subsection reason=relative-heading-rank; base=h2

\subsection{第二章}

4. 并非每个宽容的人都是朋友，也并非每个责打的人都是仇敌。\BibleQuote{朋友的创伤，胜过仇敌的假意亲吻。}\BibleRef{箴 27:6} 宁可以严厉去爱，而不要以温和去欺骗。让饥饿的人停止进食，如果这样能使他不因饮食无忧而忘记正义，比把饼给饥饿的人，让他因而受贿同意不义，更能成就善事。束缚疯狂者的人，和唤醒昏睡者的人，同样地使两人都感到烦恼，也同样地都是出于对患者的爱。谁能比天主更爱我们呢？然而，祂不仅赐予我们甘饴的教导，而且也不断地用救恩的畏惧来激励我们。祂时常在用以安慰人的柔和药剂之外，加上磨难这剂苦药；祂甚至以饥荒磨难那些虔敬热心的先祖，以更严厉的惩戒动摇叛逆的百姓，并且虽经宗徒再三恳求，仍不除去他身上的刺，好叫祂的能力在软弱中显得完全。\BibleRef{格后 12:7} 我们无论如何要爱我们的仇敌，因为这是正当的，天主也这样命令我们，为使我们成为在天之父的子女，\BibleQuote{祂使太阳上升，光照恶人，也光照善人；降雨给义人，也给不义的人。}\BibleRef{玛 5:45} 然而，当我们赞颂祂的这些恩赐时，同样地，也让我们深思祂对所爱之人的纠正。\par

5. 你认为，不应强迫任何人去追随正义；然而，你却读到主人对仆人说的话：\BibleQuote{无论你们遇着谁，都强迫他们进来。}\BibleRef{路 14:23} 你也读到，起初叫扫禄、后来叫保禄的人，怎样为了认识并拥抱真理，而受到基督强迫他的大力；因为你必定认为，你们目中所享有的光，比金钱或任何其他财物更宝贵。这光，当他因从天上来的声音跌倒在地时，忽然丧失，直到他成为圣教会的肢体后才恢复。你也认为，不应使用强迫手段，以使任何人脱离错误所带来的致命后果；然而，你却看到，在一些无可争论的事例中，天主就是这样做的，祂对我们的真正关切，胜过其他任何人；而你听到基督说：\BibleQuote{凡圣父不吸引的人，就不能到我这里来，}\BibleRef{若 6:44} 这件事，发生于所有因畏惧天主的义怒而投奔祂的人心中。你也知道，有时盗贼会把食物撒在羊群前，为要引它们走入迷途；而有时牧者会用他的棍杖，把迷路的羊领回羊群。\par

6. 撒辣当有能力时，岂不宁愿折磨那傲慢的婢女吗？她确实并非残忍地憎恨她——她曾因自己的善意而使后者成为母亲；她只是以有益的约束抑制了她的骄傲。\BibleRef{创 16:5} 而且，正如你所熟知的，这两个女人，撒辣与哈加尔，以及她们的两个儿子依撒格和依市玛耳，乃是预像，代表属神的人与属血肉的人。尽管我们读到婢女和她的儿子从撒辣受了极大的苦楚，但宗徒保禄却说，依撒格受了依市玛耳的迫害：\BibleQuote{可是，正如那时按血肉生的人迫害了按神生的人，如今仍是如此；}\BibleRef{迦 4:29} 由此，有悟性的人可以明白，那遭受迫害的，其实是\OriginalTerm{天主教会}{Catholic Church}，她透过那些属血肉之人的骄傲与不虔，承受迫害，而她正努力以现世的苦难与恐吓纠正这些人。因此，无论那真正而合法的母亲做什么，即使她的子女从她手中感受到严厉与痛苦，她也不是以恶报恶，而是运用\OriginalTerm{惩戒}{discipline}的益处来对抗罪的邪恶，并非出于怀恨而欲加害，而是出于\OriginalTerm{爱}{love}而欲治愈。当善人与恶人做同样的事，受同样的苦时，区分他们的，不是他们的行为或遭遇，而是各自的原因：\Emph{例如}，法郎以严酷的奴役压迫天主的子民；梅瑟在子民犯下不虔之罪时，以严厉的纠正磨炼他们：他们的行为相似，但他们对子民福祉的动机却不同——一个因\OriginalTerm{统治欲}{lust of power}而膨胀，另一个则因爱而炽热。依则贝耳杀害先知，厄里亚也杀害假先知；\BibleRef{列上 18:4} 我想，在这两种情形中，行此事者与受此事者各自的功过，是完全不同的。\par

7. 再请看新约时代，那时\OriginalTerm{爱}{love}本质上的温柔不仅应存于心中，也应彰显于外：在这时代，伯多禄的剑被基督命其收回鞘内，我们被教导，即便为护卫基督，也不应拔出剑来。\BibleRef{玛 26:52} 然而我们读到，不但犹太人鞭打宗徒保禄，而且希腊人为了宗徒保禄的缘故，也鞭打了犹太人索斯特乃。表面上，事件的相似性岂不将两者并列？同时，其原因的不同岂不造成了真正的差异？再者，天主既然没有怜惜自己的儿子，反而为众人把祂交出了。\BibleRef{罗 8:32} 论到子，也说：\BibleQuote{祂爱了我，且为我舍弃了自己；}\BibleRef{迦 2:20} 论到犹达斯，则说撒殚进入了他，为使他出卖基督。\BibleRef{若 13:2} 因此，既然圣父交出了自己的圣子，基督交出了自己的身体，而犹达斯交出了自己的师傅，在这交出基督的事上，为何天主是圣善的，而人是罪恶的，若非因为他们虽然做了同一件事，但做这事的原因却不同？三座\OriginalTerm{十字架}{crosses}立于一处：一座上是那将得救的强盗；第二座上是那将受罚的强盗；第三座，在二者中间的，是基督，祂将拯救一个强盗而惩罚另一个。有什么比这些十字架更相似的呢？又有什么比悬于其上的人更不同的呢？保禄被交出去受监禁和捆锁，\BibleRef{宗 21:23} 而撒殚无疑比任何狱卒更残忍：但保禄却亲自将一个人交给他，为使他那属血肉的部分被摧毁，好让他的灵魂在主耶稣的日子可以得救。\BibleRef{格前 5:5} 对此我们该说什么呢？看，两人都将人交给捆绑；但那残忍的人将囚犯交给较温和者，而那慈悲的人却将囚犯交给更残忍者。我的弟兄，让我们学会在相似的行为中分辨行事者的意图；不要闭眼不看，无端指责，将那些寻求人灵福祉的人，当作是在加害他们。同样，当这同一位宗徒说他已将某些人交给撒殚，为叫他们学得不再亵渎，\BibleRef{弟前 1:20} 难道他是以恶报恶呢？还是他其实将此视为善工，藉着恶者来纠正恶人呢？\par

8. 如果受迫害在任何情况下都是一件值得称赞的事，那么主只说：\BibleQuote{受迫害的人是有福的}，无需加上\BibleQuote{为义}。\BibleRef{玛 5:10} 再者，如果施加迫害在任何情况下都是该受谴责的，那么圣经上就不会写：\BibleQuote{暗中诽谤邻人的，我要迫害他（英译本作\InnerQuote{剪除}）。}因此，在某些情形中，受迫害的人有过错，而施加迫害的人却正义。但真相是：恶人常常迫害善人，善人也常常迫害恶人；前者以不义造成伤害，后者则以施行\OriginalTerm{惩戒}{discipline}谋求善果；前者出于残忍，后者出于节制；前者受\OriginalTerm{贪欲}{lust}驱使，后者受\OriginalTerm{爱}{love}约束。因为那旨在杀人的，不在意如何伤人；但那旨在医治的，却谨慎使用手术刀；因为一个要摧毁健全之物，另一个则要摧毁腐烂之物。恶人杀害先知；先知也曾杀害恶人。犹太人鞭打基督；基督也曾鞭打犹太人。宗徒们被世人交给世俗权柄；宗徒们自己也把人交给撒殚的权下。在所有这些事例中，唯一重要的是注意：谁站在真理一边，谁站在不义一边；谁是出于伤害的欲望，谁是出于纠正过失的渴望？\par

% structure: p0013 source=h2 target=subsection reason=relative-heading-rank; base=h2

\subsection{第三章}

9. 你们说，在福音作者和宗徒的著作中，找不到任何代教会向世上的君王祈求攻击教会的敌人的例子。谁否认这点？这样的例子确实找不到。但那时，那预言，\BibleQuote{因此，你们这些君王，现在要明智；你们这些世上的判官，要受教：应以敬畏服事主，}尚未应验。直到那时，我们在这同一首圣咏开头所读到的这些话正逐渐应验，\BibleQuote{外邦为什么嚣张，万民为什么妄想？世上的君王起来，首领们联合商议，反抗主，反抗他的受傅者。}的确，如果先知书中记载的过去事件是未来的预像，那么在\Person{拿步高}王时代，就提供了一个预像，既预表教会处于宗徒时代的时期，也预表她现在的时期。在宗徒和殉道者的时代，那预表的事情应验了：那时，前述的君王强迫虔诚而正义的人朝拜他的金像，并将拒绝朝拜的人投入火中。然而，现在应验的，是那同一君王不久后的另一预表事件：当他转而敬拜真天主后，便在其帝国境内颁布谕令，凡亵渎\Person{沙得辣客}、\Person{默沙客}与\Person{阿贝得乃哥}的天主者，应受其罪行所当得的刑罚。那君王的前期代表了那些不信基督的皇帝们的早期时代，在他们的治下，基督徒因恶人而受苦；但那君王的后期代表了继承皇位者、现今已信基督的时代，在他们的治下，恶人因基督徒而受苦。\par

10. 然而，很明显地，对那些在基督徒名号下被邪僻之人引入歧途的人，我们小心谨慎地施加一种适度的严厉，或更确切地说，一种宽仁；我们所用的手段旨在阻止基督的羊走失，并将他们带回羊栈，通过诸如流放与罚款之类的处罚，劝告他们思考自己所受的是什么、原因何在，并教导他们珍视自己所读的圣经，胜过人的传说与毁谤。我们中谁，甚至你们中谁，不赞许皇帝们颁布的反对异教祭献的法律呢？在这些法律中，确实规定了更为严厉的刑罚，因为那等不敬之罪的惩罚乃是死刑。然而在镇压和约束你们时，所企图的毋宁是劝告你们远离邪恶，而非因罪行惩罚你们。因为或许宗徒论及犹太人的话也可用于你们：\BibleQuote{我可以为他们作证，他们对天主有热心，但不按照真知，因为他们不认识由天主而来的正义，企图建立自己的正义，而不服从天主的正义。}\BibleRef{罗 10:2} 因为，当你们说除了那些有机会由你们施洗的人之外，无人能成义，这岂非企图建立自己的正义，还能是什么呢？关于宗徒论及犹太人的这言论，你们与这言论原初所指的那些人之间的区别，在于你们拥有基督的圣事，而他们至今仍缺乏这些圣事。然而，就这些话而言：\Quote{不认识天主的正义，企图建立自己的正义}，以及\Quote{他们对天主有热心，但不按照真知}，你们与他们完全一样，只有你们中间那些知道什么是真理，却在自己乖戾的顽固中继续攻击自己完全清楚之真理的人除外。这些人的不虔敬，或许甚至是比偶像崇拜更大的罪。然而，由于他们不易被定此罪（因为这罪隐藏于内心），你们所有人同等地受到一种相对温和的严厉约束，因为你们与我们并非如此疏远。这一点我可以论及所有不分青红皂白的异端者，他们虽保留了基督的圣事，却背离了基督的真理与合一；也可以论及所有\OriginalTerm{多纳徒派}{Donatists}，无一例外。\par

11. 然而，至于你们，不仅与上述那些人同被称为\OriginalTerm{多纳徒派}{Donatists}（由\Person{多纳徒}得名），还特别被称为\OriginalTerm{罗加特派}{Rogatists}（由\Person{罗加特}得名），你们似乎性情较为温和，因为你们不随同那些凶残的\OriginalTerm{西康瑟里翁派}{Circumcelliones}团伙四处猖狂；但若一只野兽因为没有牙齿和利爪而不伤人，决不能因此就被称为温和。你们说自己无意残忍：我认为你们所缺少的是能力，而非意愿。因为你们人数如此之少，即使你们有此意欲，也不敢对抗那些敌视你们的群众。然而，我们姑且假定，你们并不想去做你们无力做到的事；我们姑且假定，你们如此理解并遵守福音的规诫：\BibleQuote{若有人想要控告你，拿你的里衣，你连外衣也让他拿去，}\BibleRef{玛 5:40} 以至于抵抗迫害你们的人便是不合法的，无论迫害你们的人有理或无理。你们教派的创立者\Person{罗加特}，要么并不持守此观点，要么便犯了言行不一之罪；因为他曾以极其坚决的态度，在一桩据你们声称属于你们的某些财物的诉讼中进行争辩。如果当时有人对他说：有哪一个宗徒曾因信仰之事向世俗法庭申诉，以维护自己的财产？正如你在信中所提出的问题：\Quote{有哪一个宗徒曾在信仰之事上侵占他人的财产？} 他在圣经中找不到任何此类例子；但他或许能找到某种真正的辩护，倘若他没有背离真教会，随后又大胆地以真教会的名义要求拥有那争议的财产。\par

% structure: p0017 source=h2 target=subsection reason=relative-heading-rank; base=h2

\subsection{第四章}

关于获取或执行这世界之权势针对裂教者和异端者的法令，那些你们离开的人，据我们所闻，积极从事此事，既反对你们，也反对\Person{马克西米安}的追随者，我们以他们自己无可否认的纪录为证；但你们尚未离开他们的时候，他们在请愿中对\Person{尤里安}皇帝说：\Quote{除了正义，在他那里找不到别的，}——然而他们始终都知道他是一个背教者，且看到他如此沉溺于偶像崇拜，以至于他们要么承认偶像崇拜就是正义，要么无法否认当他们说除了正义在他那里找不到别的时，是邪恶地撒谎，因为他们看到偶像崇拜在他那里如同正义般占据如此大的位置。不过，姑且说那只是用词上的错误，对于行为本身你们又怎么说？如果连正当的事都不该通过向皇帝恳求来寻求，那为什么你们却从\Person{尤里安}那里寻求你们认为是正当的事呢？\par

你是否回答说，为了追回自己的财产而向皇帝请愿是合法的，但为了迫使皇帝惩罚某人而控告他则不合法？我可以顺便指出，即使是为了追回自己的财产而请愿，也放弃了宗徒的榜样所涵盖的立场，因为我们从未发现有宗徒这样做过。但撇开这一点不谈，当你们的前辈们通过总督\Person{阿努利努斯}，向\Place{君士坦丁}皇帝控告当时的\Place{迦太基}主教\Person{凯齐利安}，他们拒绝与这个他们视为有罪的人共融，他们并非试图追回自己所失去的东西，而是以诬告攻击一个我们认为是——而且司法程序的结果也表明——无辜的人；还有什么比这更令人发指的罪过呢？然而，如果像你们错误地认为的那样，他们确实把一个真正有罪的人交给了世俗政权的审判，那么，你们为什么反对我们做你们自己的人首先擅自做的事呢？如果他们做这件事不是为了嫉妒地伤害人，而是为了谴责和纠正错误，我们也不会责怪他们。但我们毫不犹豫地责备你们，因为你们认为，我们向一位基督徒皇帝控告与我们共融为敌的人是犯罪，然而，你们的前辈们交给总督\Person{阿努利努斯}、并由他转交给\Person{君士坦丁}皇帝的一份文件，其题目是：\Quote{\WorkTitle{Libellus Ecclesiæ Catholicæ, criminum Cæciliani, traditus a parte Majorini.}}。我们更特别地责备他们，因为当他们主动向皇帝控告\Person{凯齐利安}时，他们本应该首先在海外那些他的同僚主教们面前证明这些指控，而当皇帝以远比他们更为有序的方式，把交到他面前的这桩主教案情交给主教们裁决时，他们即使在裁决对他们不利、被击败之后，也不愿意与弟兄们和好。相反，他们接着在世俗君王的法庭上，不仅控告\Person{凯齐利安}，还控告那些被任命为审判官的主教们；最后，他们又从第二次主教法庭向皇帝上诉。当皇帝亲自调查并作出判决时，他们也不认为有责任顺服真理或和平。\par

如果当你们的前辈们控告\Person{凯齐利安}时，\Person{凯齐利安}和他的朋友们败诉了，\Person{君士坦丁}除了对那些主动提出控告、却未能证明其指控、且在败诉后仍拒绝顺服真理的人所颁布的法令之外，还能作出什么别的裁决呢？如你们所知，皇帝在那个案件中首次下令，那些被定为有罪、顽固抵抗教会合一者的财产应予没收。然而，如果事情的结果是，你们提出控告的前辈们赢了官司，而皇帝对\Person{凯齐利安}所属的共融团体颁布了类似的法令，你们就会希望皇帝们被称为教会利益的友人、教会和平与合一的守护者。但当皇帝们对那些主动提出控告、却无法证实其控诉、并在获得悔改条件即可重返和平怀抱时却拒绝接受的人颁布这类法令时，却有人大声抗议，说这是不当的伤害，并主张不应强迫任何人合一，也不应以恶报恶。这岂不是你们自己人曾写过的：\Quote{我们所愿望的就是圣的}吗？鉴于这些事，你们本不难反思并发现，由于你们的前辈们不断向皇帝提出他们无法证实的指控，\Person{君士坦丁}针对你们颁布的法令和判决就应依然有效；而所有继任的皇帝，尤其是那些公教基督徒皇帝，在你们的顽固使得他们有必要对你们采取任何措施时，都自然会照此行事。\par

15. 你们本可轻易反省这些事，或许有时会对彼此说：既然\Person{凯基利安}要么无辜，要么至少无法证明有罪，这遍布全世界的基督教会，在这件事上犯了什么罪呢？对于连那些以此控告他人的人都无法证明的事，整个基督徒世界仍不知情，这又有何不可呢？基督把祂的种子撒在祂的田里，也就是这世界，并吩咐让它们与莠子一起生长，直到收割的时候，\BibleRef{玛 13:24}——这些万国中成千上万的信徒，主曾将他们的数目比作天上的星和海边的沙，祂在古时应许给他们、如今已赐给他们在\Person{亚巴郎}后裔中的祝福——请问，为什么要否认所有这些人享有基督徒之名呢？因为，诚然，对于这件他们从未参与讨论的案件，他们宁愿相信那些在重大责任下作出裁决的法官，而不相信那些裁决对其不利的原告。确实，一个人的罪行不可能玷污另一个对此毫不知情的人。散居全世的信徒，怎么可能知道那些人犯下交出圣书的罪行，而连他们的指控者，即使知道这罪行，也至少无法证明？毫无疑问，他们不知情的这一事实，最轻易地证明他们与这罪行毫无干系。那么，为什么要因不知他人被公正或不公正地指控的罪行，而使无辜者被控以他们从未犯过的罪行呢？如果单因不知他人的罪行就构成犯罪，那无辜还有什么余地可言？此外，正如前述，单单不知情就证明如此多民族的信众是无辜的，那么与这些无辜者断绝共融的罪过，该是多么重大啊！因为，那些既无法向无辜者证明、也无法为他们所相信的有罪之人的行为，绝不会玷污任何人，因为即使知道了这些行为，为了保持与无辜者的共融，也会容忍它们。因为，不应为了恶人而离弃善人，反而要为善人而容忍恶人；正如先知们容忍了那些他们曾作证反对的人，并且没有停止参与犹太民族的圣事；又如我们的主容忍了有罪的\Person{犹达斯}，直到他遭受应得的结局，并允许他与无辜的门徒一起参与神圣的晚餐；又如宗徒们容忍了那些出于嫉妒——一种尤为魔鬼般的罪——而宣讲基督的人；又如\Person{西彼廉}容忍了犯有贪婪之罪的同僚，按照宗徒的榜样，\BibleRef{哥 3:5}，他称贪婪为偶像崇拜。总之，无论当时在这些主教中间发生了什么，尽管或许其中有些人知道，但只要在审判时不偏待人，对所有人都是未知的：那么，为什么众人不爱慕和平呢？这些想法你们本容易想到；或许你们已经怀有这些想法。但是，你们与其执着于无价值的虚荣，担心因赞同真理而在人前失去它，倒不如执着于世间的财产，因害怕失去财产，你们或许会被显明是赞同已知真理的。\par

% structure: p0022 source=h2 target=subsection reason=relative-heading-rank; base=h2

\subsection{第五章}

16. 因此，我想你们现在明白了，当任何人被强迫时，要考虑的不仅仅是强迫这一事实，而是他被强迫去接受的事物的性质，无论是善是恶：这并不是说任何人能违背自己的意愿而成为善的，而是说，通过害怕遭受他不愿受的痛苦，他要么放弃了敌对的偏见，要么被迫去审视他先前满足于无知的那真理；并在这恐惧的影响下，摒弃了他习惯性辩护的谬误，或寻求他原先毫无所知的真理，现在心甘情愿地持守他先前所拒绝的。也许，若没有这么多事例的证明，用言语断言这一点会完全无用。我们看到，不是这里那里少数几个人，而是许多城市，从前是多纳图派，如今是天主教，极度憎恶那魔鬼般的分裂，并热切爱慕教会的合一；而他们成为天主教徒，是受了那种令你们如此反感的恐惧的影响，即通过皇帝们的法律，从\Person{君士坦丁}开始（你们的党派曾在他面前自行控告\Person{凯基利安}），直到我们这个时代的皇帝们，他们最公义地裁定，你们自己的党派所选择、且他们优先于主教法庭的那位法官的判决，应当对你们保持效力。\par

17. 因此，我已被我的同工们摆在我面前的这些实例所给出的证据所说服。因为起初我的意见是，任何人都不应被强迫进入基督的合一之中，我们只应藉言语行事，只应藉辩论战斗，并应靠理性的力量得胜，免得我们得到那些我们明知是公然异端者的人，却假装成公教徒。但我的这个意见，不是被那些反对它的人的话所胜过，而是被他们所能指出的确凿实例所胜过。首先，与我的意见对立的是我自己的城镇。这城镇虽然曾经完全站在\Person{多纳图斯}一边，却因畏惧帝国敕令而被带入公教的合一之中，如今我们看到它充满了对你那可悲悖逆的憎恶，以至于几乎难以相信它曾陷入过你的错谬之中。还有许多其他城镇被指名道姓地告诉了我，以致从事实本身，我不得不承认圣经的话可以被理解为适用于这件事：\BibleQuote{给智慧人机会，他就会更有智慧。}\BibleRef{箴 9:9} 因为正如我们确知的，有多少人早已愿意成为公教徒！他们被不容置辩的真理的明晰所打动，却日复一日因害怕冒犯自己的党派而延迟公开宣认。有多少人不是被真理所束缚——因为你们从不曾声称真理是你们的——而是被积习的沉重锁链所束缚，以致在他们身上应验了天主的这句话：\BibleQuote{仆人（刚硬的）不受言语的纠正；他虽然明白，却不回答}！\BibleRef{箴 29:19} 有多少人仅仅因为安逸使他们过分懒散、自负或怠惰，不肯费力去考察公教的真理，就认为\Person{多纳图斯}派才是真教会！若非那些诽谤者的谗言将他们拒之门外——这些人声称我们在天主的祭台上所献的并非我们实际所献的——有多少人早就进来了！有多少人，认为一个基督徒属于哪一派无关紧要，仅仅因为他们出生在其中，且无人强迫他们离开它而转入天主教会，就继续留在\Person{多纳图斯}的\OriginalTerm{裂教}{schism}之中！\par

18. 对于所有这些类别的人，对法律的畏惧——君王在敬畏中事奉主，正是以颁布这些法律来事奉——已如此有益，以致如今有些人说：我们很早以前就愿意如此；但感谢天主，祂给了我们立刻这样做的机会，并斩断了拖延的犹豫！另有些人说：我们早已知道这是真的，却被旧习惯的力量所囚禁：感谢主，祂已将这些捆绑折断，把我们带入和平的联结之中！又有些人说：我们不知真理就在这里，也无意愿去学习；但恐惧使我们变得认真，我们在惊慌中开始考察它，免得在永恒之事上毫无所得，却在今世之事上遭受损失：感谢主，祂以恐惧的刺激惊醒了我们的懈怠，使我们如今在不安中去探究那些安逸时我们无心去知道的事！还有些人说：我们被虚假的谣言阻挡，无法进入教会，这些谣言我们若不进入教会，就不能知道是假的；我们本不愿进入，除非受到强迫：感谢主，祂以祂的鞭笞除去了我们的胆怯犹豫，教导我们亲自发现，谣言对祂的教会所散布的谎言是何等虚妄而荒谬：藉此我们确信，这异端的创立者所发出的指控毫无真实可言，因为他们的追随者所捏造的更严重的罪名也毫无根据。又有些人说：我们确实曾以为，在哪个\OriginalTerm{共融}{communion}中持守基督的信德并不重要；但感谢主，祂将我们从\OriginalTerm{裂教}{schism}的状态中聚集起来，并教导我们，在合一中敬拜唯一的天主原是合宜的。\par

19. 那么，我怎能坚持反对我的同工们，且因抗拒他们而阻挡主的这些胜利，并阻止基督的羊——它们曾在你们的山岭和冈陵上，即在你们的骄傲的膨胀之处——被聚集到和平的羊栈中，在那里只有一个羊群、一个牧者？\BibleRef{若 10:16} 难道我有责任阻挠这些措施，为的是让你们不致失去你们所谓的自己的东西，并让你们能无忧无惧地抢夺属于基督的东西吗？让你们能按罗马法拟定你们的遗嘱，并能以诽谤性的指控破坏那以天主法律之认可、针对列祖所立的\OriginalTerm{约}{Testament}，其中写着：\BibleQuote{地上万民都要因你的后裔蒙受祝福}：\BibleRef{创 26:4} 吗？让你们能在买卖交易中享有自由，并能壮胆去分割并声称基督以舍己为代价所买来的东西是你们的吗？让你们中一人赠予另一人的任何礼物不受质疑，而万神之神赐予其子女的礼物——他们从日出之地蒙召到日落之处——却变成无效吗？让你们不致从你们天然出生的土地被流放，却让你们竭力将基督从祂以宝血买来的国度中驱逐出去，这国度从这海伸展到那海，从大河直到地极吗？不，坚决不；愿地上的君王通过为基督和祂的事业制定法律来事奉基督。你们的先辈曾将\Person{凯基利安}和他的同伴暴露于世上的君王面前，让他们因诬告的罪名而受惩罚：现在就任狮子转身，去粉碎那些诽谤者的骨头吧，当\Person{达尼尔}已被证明无辜、并从他们将遭遇丧命的狮坑中被解救出来时，不要为他们求情\BibleRef{达 6:23}；因为那为邻舍掘坑的，自己反倒在其中掉得最为公正\BibleRef{箴 26:27}。\par

% structure: p0027 source=h2 target=subsection reason=relative-heading-rank; base=h2

\subsection{第六章}

20. 因此，我的弟兄，趁你还在现世生命，你要救自己脱离那要临于顽固和骄傲者的愤怒。这世界的官长们的可怕权力，当它攻击真理时，为强者提供了光荣的考验机会，却为软弱的义人设下了危险的诱惑；但当它协助宣扬真理时，它为有智慧者带来有益的劝诫，却为迷途中愚昧的人带来无益的烦恼。\BibleQuote{因为没有权柄不是出于天主的：所以谁反抗权柄，就是反抗天主的规定；因为官长对于善行不是可怕的，而是对于恶行。你愿意不惧怕权柄吗？你行善罢！就可由他得到称赞。}\BibleRef{罗 13:1} 因为如果权力站在真理一边，纠正某个陷于错谬的人，那因纠正而改过的人便会从权力获得称赞。相反，如果权力与真理为敌，残酷地迫害某人，那位在斗争中获胜并戴上荣冠的人，也会从他抵抗的权力那里获得称赞。但你却不行善，好叫你不惧怕权柄；除非这竟是善：坐在那里发言反对不是一个弟兄，而是反对散居万民中的你所有的弟兄们，而先知们、基督和宗徒们，在圣经的话语中为这些弟兄作证说：\BibleQuote{地上万民都要因你的后裔蒙受祝福；} 又说：\BibleQuote{从日出直到日落之处，必有纯洁的祭品奉献于我的名；因为我的名在外邦中必受尊崇，主说。}\BibleRef{拉 1:11} 要注意这话：\BibleQuote{主说}；不是说\Person{多纳图斯}，也不是\Person{罗加图斯}，也不是\Person{文森提乌斯}，也不是\Person{安博罗削}，也不是\Person{奥斯定}，而是\BibleQuote{主说}；再有：\BibleQuote{世上万邦都要因他蒙受祝福，万民都要称他为有福。愿主天主，以色列的天主，受赞颂，因为唯有他行奇事；愿他的光荣名号永受赞颂，愿他的光荣充满整个大地！但愿如此，但愿如此！} 而你却在\Place{卡尔滕纳}坐着，和十几个残存的罗加图斯派人士一起说：\Quote{但愿不如此！但愿不如此！}\par

21. 你在福音中听到基督这样说话：\BibleQuote{凡梅瑟法律、先知和圣咏上指着我所记载的话，都必须应验。耶稣遂开启他们的明悟，叫他们理解经书；又向他们说：\InnerQuote{经上曾这样记载：默西亚必须受苦，第三天要从死者中复活；并且必须从耶路撒冷开始，因他的名向万邦宣讲悔改，以得罪之赦。}}\BibleRef{路 24:44} 你也在\WorkTitle{宗徒大事录}中读到这福音怎样从\Place{耶路撒冷}开始，在那里圣神先充满了那一百二十人，然后从那里传到\Place{犹太}和\Place{撒玛黎雅}，并传到万民中，正如他即将升天时对他们所说的：\BibleQuote{你们要在耶路撒冷、犹太全境和撒玛黎雅，直到地极，为我作证人。} 因为\BibleQuote{它们的声音传遍普世，它们的言语达于地极。} 而你却抗拒如此坚定确立又如此清晰启示的天主的明证，并企图如此彻底地剥夺基督的产业，以致虽然悔改如他所说的，因他的名已向万民宣讲，但任何地方的人，哪怕被这宣讲所感动，也没有可能获得罪赦，除非他寻找并发现\Place{卡尔滕纳}的\Person{文森提乌斯}，或他十个同伴中的某一位，隐没于帝国殖民地\Place{茅利塔尼亚}的偏僻之地。卑微之人的骄傲还有什么不敢做的呢？血肉之躯的妄自尊大会把人驱逼到何等极端！这就是你所行的那善，让你因此就不惧怕权柄吗？你在你母亲的儿子（基督为他死了）的路上放置了这严重的绊脚石，\BibleRef{格前 8:11} 而他还处在软弱的婴儿期，不能吃硬食，还需要喂以母乳；\BibleRef{格前 3:2} 你为了否认教会在万民中增长的事实，竟引用\Person{依拉略}的著作来反驳我；这增长直到世界终结，正是天主的许诺，为制服你的不信，他曾以誓言确认！当初这许诺被赐下时，你若反对它，必定极为不幸；如今这许诺已经应验，你却仍然反对它。\par

% structure: p0030 source=h2 target=subsection reason=relative-heading-rank; base=h2

\subsection{第七章}

22. 然而，你凭着你深湛的学识，发现了一些你认为足以作为反对天主见证的重大异议的东西。因为你说：\Quote{如果我们考虑整个世界所包括的部分，基督信仰为世人所知的地方所占的比例是相当小的。} 你还拒绝看见，或假装不知道，在如此短暂的时间里，福音已经传入了多少野蛮民族，甚至连基督的仇敌也不能怀疑，再过不多时，我们主所预言的那事就要成就：当门徒们问及世界末日时，他回答说：\BibleQuote{这天国的福音必在全世界宣讲，给万民作证；然后结局才会来到。}\BibleRef{玛 24:14} 同时，你尽你所能地宣称并坚持，即使福音传到了\Place{波斯}和\Place{印度}——事实上早已如此——但凡听到的人，若不来到\Place{卡尔滕纳}或\Place{卡尔滕纳}附近，都丝毫不能洗净自己的罪！如果你没有明确这样说，那无疑是因为怕人嘲笑你；然而，当你果真这样说时，你难道拒绝人为你哭泣吗？\par

23. 你以为你提出了一个非常精辟的论点，当你断言\OriginalTerm{大公}{Catholic}之名意指普世，并非指包容整个世界的共融，而是指遵守一切神圣诫命和一切圣事时，好像我们（即使接受这立场：教会之所以称为\OriginalTerm{大公}{Catholic}，是因为它诚实地持守全部真理，而某些异端中处处可见这真理的碎片）依据的是这个词意义的证据，而不是天主的应许，以及如此多无可辩驳的真理本身的见证，来证明天主的教会存在于万国之中。然而，实际上，这就是你们试图让我们相信的全部：只有\OriginalTerm{罗加特派}{Rogatists}才配得上\OriginalTerm{大公信徒}{Catholics}之名，因为你们遵守了一切神圣诫命和一切圣事；并且只有你们才是\BibleQuote{人子来临时，能遇见信德}的人。\BibleRef{路 18:8} 请原谅我这样说，我们对此一个字也不相信。因为，虽然为了使主说过在地上找不到的那种信德能在你们里面找到，你们或许甚至擅自声称你们应该被视为在天上，而非在地上，至少我们从宗徒的警告中获益，他教导我们，\BibleQuote{即使是从天上来的天使，若给我们宣讲的福音，与我们领受的不同，就该受诅咒}。\BibleRef{迦 1:8} 但是，如果我们不接受同一天主圣言为教会所作的无可辩驳的见证，我们又怎能确信在圣言中我们有关于基督的无可辩驳的见证呢？因为，不论人想出何等精巧复杂的诡辩来反对这单纯的真理，也不论用何等浓重的诡诈迷雾来遮蔽它，凡宣讲基督没有受难、第三天没有从死者中复活的人，都该受诅咒——因为我们在福音的真理中学到，\BibleQuote{基督必须受苦，第三天从死者中复活；}\BibleRef{路 24:46}——基于完全相同的理由，凡宣讲教会不在包容万国的共融中的人，也必该受诅咒：因为在同一段经文的接下来的话中，我们还学到，\BibleQuote{必须从耶路撒冷开始，因他的名向万邦宣讲悔改，以得罪之赦。}\BibleRef{路 24:47} 我们应当坚守这条准则：\BibleQuote{如果任何人给你们宣讲的福音，与你们领受的不同，他该受诅咒。}\BibleRef{迦 1:9}\par

% structure: p0033 source=h2 target=subsection reason=relative-heading-rank; base=h2

\subsection{第八章}

24. 此外，即便我们不听从整个\OriginalTerm{多纳特派}{Donatists}声称他们是基督的教会，因为他们在这方面并不能从圣经中提出任何证据；那么请问，我们岂不更不应当听从\OriginalTerm{罗加特派}{Rogatists}吗？他们竟为了自己党派的利益，试图解释圣经的话：\Quote{你在何处放牧？你在何处午间休息？}！因为，如果以此指非洲的南部——也就是多纳特派所占据的地区，因它处于天空更炽热的部分之下——那么，\OriginalTerm{马克西米安派}{Maximianists}必然胜过你们党派的其他一切，因为他们的分裂之火在\Place{拜占庭}和\Place{的黎波里}爆发。让\OriginalTerm{阿尔祖格人}{Arzuges}，如果他们愿意，与他们就这一点争论，并争辩这节经文更适用于他们；但是，属帝国的\Place{毛里塔尼亚}行省，位置更偏西而非偏南，既然它拒绝被称为非洲——我要说，它怎能从\Quote{南方}一词中找到夸耀的理由？我不说对抗世界，而是对抗连多纳特派本身，而罗加特派正是从那另一更大的碎片中脱离出来的极微小碎片。一个人若没有明白的见证，借其光照可以使隐晦的寓意显明，就擅自按自己的利益解释任何寓意性说法，这岂非极度的厚颜无耻？\par

我们常对所有多纳徒派说的话，现在可以更有力地对你们说：如果有人有充分的理由（其实没有人能有）使自己与全世界的共融分离，并将自己的共融称为基督的教会，理由是他们有理由地从所有民族的共融中退出了——那么，你们怎么知道，在基督教的社会里，这社会如此广阔地分布，可能在某个非常遥远的地方，那里的义人的名声无法传到你们耳中，在你们分裂之前，他们就已经为了某个正义的原因而脱离了全世界的共融？这样一来，教会怎么可能在你们一派之中，而不是在那些比你们更早分离的人之中呢？因此，只要你们对此无知，你们就无法确定地提出任何主张：这必然是所有在捍卫自己党派的事业时，诉诸自己的见证而非天主的见证的人的命运。因为你们不能说：如果发生了这样的事，不可能不为我们所知；因为，即使不超出\Place{非洲}本身，当你们被质问时，你们也说不出多纳徒派内部发生了多少分裂：对此我们必须特别记住，在你们看来，分离的人数越少，他们事业的正义性就越大，而人数的稀少无疑使他们更容易不被注意到。因此，你们也根本无法确定，在\Place{非洲}南部很遥远的某个地方，是否有一些义人，人数很少，因而无人知晓，早在多纳徒派将自己的义从与所有其他人的不义交往中撤回之前，就在他们遥远的北方地区，以同样的方式，为了一些极其充分的理由而分离出来，并且现在，凭借比你们更优越的权利，成为天主的教会，作为在正当分离的事上先于你们所有派别的属灵的\Place{熙雍}，而且他们解释圣咏中的话：\BibleQuote{熙雍山，在极北的边界，是大王的城邑}，以此为自己辩护，这比多纳徒派解释另一句话：\BibleQuote{你在何处牧羊？正午你让他们在何处歇息？}\BibleRef{歌 1:7}，以此为自己辩护要有理由得多。\par

然而你们却声称害怕，当你们被皇帝的法令强迫接受合一时，天主的名会更长久地被犹太人和外邦人亵渎：好像犹太人不知道，他们的民族以色列在历史上起初的时候，曾想用战争消灭那在\Place{约但河}东承受了产业的二支派半，因为他们认为这些支派已经脱离了民族的合一。\BibleRef{苏 22:9} 至于外邦人，他们确实可以更合理地指责我们，因为基督徒皇帝制定了反对偶像崇拜者的法律；然而，通过这些法律，他们中许多人已经，并且现在每天都有，从偶像转向永生的真实的天主。但实际上，如果犹太人和外邦人认为基督徒的人数像你们一样微不足道——你们居然主张只有你们才是基督徒——他们就不会屈尊说任何反对我们的话，而是会不停地嘲笑和藐视我们。难道你们不怕犹太人对你们说：\Quote{如果你们这一小撮人就是基督的教会，那么你们的宗徒\Person{保禄}所说的，你们的教会由经上的话所描写的那……岂不是落空了吗？\InnerQuote{不生育的石女，欢乐罢！未经产痛的女人，欢呼高唱罢！因为被弃者的子女比有夫者的子女还多。}\BibleRef{迦 4:27}}？你们会对他们说：\Quote{我们人数虽少，但我们更义；} 你们难道期望他们不会回答：\Quote{无论你们自称什么，如果你们被缩减到如此小的数目，你们就不是经上所说的那个\InnerQuote{被弃者却有\Emph{许多}子女}的}？\par

也许你们会拿那位义人的例子来反驳这一切，在洪水来临时，唯独他和他全家被认为配得拯救。那么，你们看出你们离义还有多远吗？我们当然不会根据这个例子肯定你们是义人，除非你们的同伴减少到七人，而你们自己是第八人：然而始终要假设，正如我所说的，没有其他人在\Person{多纳徒}派之前抢先夺取了那义，因为他在某个遥远的地方，出于某些充分的理由，同另外七人一起脱离了全世界的共融，从而从淹没它的洪水中拯救了自己。因此，由于你们不知道是否可能有人这样做了，而且你们对此全然不知，就像远方许多基督徒民族从未听闻\Person{多纳徒}之名一样，你们就无法确定地说教会在哪里可以找到。因为教会必定是在那个地方，在那里，你们现在所做的事可能早已被别人做过了，如果有任何可能的正义理由支持你们脱离全世界的共融的话。\par

% structure: p0038 source=h2 target=subsection reason=relative-heading-rank; base=h2

\subsection{第九章}

28. 然而，我们确信，没有人能够有理由使自己脱离所有民族的共融，因为我们每一个人在自身义德中找不到教会的标记，而是在圣经中寻找，并按照应许看见它实际存在。因为论到教会，经上说：\BibleQuote{我的爱卿在少女中，有如荆棘中的百合花；}\BibleRef{歌 2:2} 这教会一方面因人的品行邪恶而可称为\Quote{荆棘}，另一方面因共享同一圣事而可称为\Quote{少女}。再者，教会说：\BibleQuote{当我心灵忧戚时，我从地极呼号你；} 又在另一篇\BibleBook{}{圣咏}中说：\BibleQuote{为了那些背弃你法律的罪人，我常有恐怖攻心；} 又说：\BibleQuote{我一见奸党，就心生厌烦。} 同一位教会又对她的新郎说：\BibleQuote{我心爱的！你告诉我：你在何处牧羊？正午在何处卧卧？为什么我在你同伴的羊群边，独自徘徊？}\BibleRef{歌 1:7} 这也就是在另一处所说：\BibleQuote{求你使我知道你的右手，和那些心中受教而有智慧的人；} 在他们当中，由于他们闪耀着光明、燃烧着爱德，你如同在正午安息；免得我像蒙纱的，即隐藏而不为人知的，跑往别处——不是你的羊群，而是你同伴的羊群，\Emph{即}异端者的羊群，在这里新娘称他们为同伴，正如祂将\BibleRef{歌 2:2}中的荆棘也称为\Quote{少女}，是因为共同参与圣事的缘故：对于这些人，在别处经上说：\BibleQuote{你本是我的同伴，我的知交，我的知心好友！我们彼此畅谈，在人群中一同进入天主的殿宇。愿死亡突袭他们，使他们活活堕入阴府！} 就像\Person{达堂}和\Person{阿彼兰}，是不虔敬裂教的始作俑者。\par

29. 这也是对教会所作的答复，就在上文所引经文的后面：\BibleQuote{你这女中极美丽的啊！你若不知，出去追踪羊群，在牧人的牧场边，牧放你的小羊。}\BibleRef{歌 1:8} 啊，新郎无可比拟的甘饴，祂这样回答她的问题：\Quote{你若不知，}祂说；仿佛祂说，\BibleQuote{城造在山上是不能隐藏的；\BibleRef{玛 5:14} 因此，\InnerQuote{你并非像蒙纱的，要跑到我同伴的羊群那里去。} 因为我是在诸山之巅建立的山，万民都要向我涌来。\BibleRef{依 2:2} \InnerQuote{你若不知，}凭你从我的书卷的见证中，而非假见证人的话中，所能获得的知识；\InnerQuote{你若不知，}从有关你的这种见证：\InnerQuote{扩张你帐幕的布幔，伸张你居所的帷幔，不要吝惜！放长你的绳子，坚固你的木橛！因为你要向左右拓展，你的后裔将继承外邦人，使废弃的城邑再有人居住。不要害怕！你再不会蒙羞；不要羞惭！你再不会受辱；你几乎忘了你青春的耻辱，不再记得你寡居的羞辱；因为你的夫君是你的造主，万军的上主是他的名字，你的救主是以色列的圣者，他将称为全世界的天主。}\InnerQuote{你若不知，}你这女中极美丽的啊，从这句论到你所说的，\InnerQuote{君王爱恋你的美丽，}和\InnerQuote{你的子孙要继承你的列祖，你将立他们为全地的统治者。}因此，如果你真的\InnerQuote{不知自己，}出去吧：我不是驱逐你，而是\InnerQuote{出去吧，}好让经上论到这样的人说，\InnerQuote{他们出于我们，但不是属于我们。}\BibleRef{若一 2:19} \InnerQuote{出去吧，}沿着羊群的足迹，不是我的足迹，而是羊群的足迹；不是那唯一的羊群，而是分裂走散的羊群。\InnerQuote{去牧放你的小羊，}不要像伯多禄，他受命说，\InnerQuote{你喂养我的羊群；}\BibleRef{若 21:17} 而是，\InnerQuote{在牧人的帐篷边牧放你的小羊，}不是在牧者的帐篷边，那里有\InnerQuote{一个羊群，一个牧人。}\BibleRef{若 10:16}} 但教会认识她自己，从而避免了那些不认识自己在教会内的人所遭遇的命运。\par

30. 当论到她的成员数目与恶人的众多相比是稀少时，圣经论到同一个教会说：\BibleQuote{那导入生命的门是多么窄，路是多么狭！找到它的人果然不多。}\BibleRef{玛 7:14} 又论到同一个教会成员的众多说：\BibleQuote{我要使你的后裔繁多，如天上的星辰，如海边的沙粒。}\BibleRef{创 22:14} 因为那同一由圣洁善良的信者组成的教会，若与为数更多的恶人相比，则是小的；若单独考虑，则是大的；\Quote{因为被遗弃者比有丈夫者更有众多子女，}并且\BibleQuote{将有许多人从东方和西方来，同\Person{亚巴郎}、\Person{以撒格}和\Person{雅各伯}在天国里一起坐席。}\BibleRef{创 8:11} 天主还为自己准备了一个\Quote{众多的，热心行善。}而在\BibleBook{默示录}中，有无数群众，\BibleQuote{没有人能够数清}，从各支派、各语言来，他们身穿白衣，手持胜利的棕榈枝。\BibleRef{默 7} 这同一个教会有时被众多的过犯所遮蔽，如同乌云笼罩，那时罪人们弯弓，要在暗月下射击心灵正直的人。但即使在这样的时刻，教会仍在那些对她最坚忠的人身上发显光辉。如果对于上面引用的神圣许诺，要对其两个分句作任何特定的应用，那么将\Person{亚巴郎}的后比作\BibleQuote{天上的星辰}和\BibleQuote{海边的沙粒}，或许并非没有道理：因此，\BibleQuote{星辰}可被理解为那些在数目上较少、但更为固定和更为光辉的人；而\BibleQuote{海边的沙粒}被理解为教会内那众多的软弱和属血肉的人，他们在风平浪静时显得安详自由，但在磨难和诱惑的风浪下却被淹没和摇动。\par

31. 这样，\Person{依拉略}所写那段经文的时候，正是这样动乱的时代。你们却认为适宜狡猾地引证那段经文，来反对如此多的神圣证据，好像借此可以证明教会已从地上灭绝。你们倒不如说：当宗徒写信给\Place{迦拉达}的各教会时，那些教会也不复存在了，因为他说：\BibleQuote{无知的迦拉达人啊！谁迷惑了你们呢？}又说：\BibleQuote{你们以圣神开始了，如今却愿以肉身结束吗？}因为你们会这样曲解那位博学的人，他（如同宗徒一样）严厉地责斥那些心灵迟钝和胆怯的人，他为他们再次受产痛，直到基督在他们内形成。\BibleRef{迦 4:19} 谁不知道当时许多判断力薄弱的人，被模棱两可的言辞所欺骗，以为\OriginalTerm{亚略派}{Arians}和他们自己所信的道理相同？谁不知道另有些人出于畏惧而屈服，假装同意，没有按福音的真理正直行事？当他们从错误中回转时，你们竟不愿给他们那后来所给予的宽恕。但你们拒绝那样的宽恕，就证明你们完全不认识天主的圣言。因为，请读一读\Person{保禄}关于\Person{伯多禄}所记载的事，\BibleRef{迦 2:11} 并读一读\Person{西}根据那段陈述所表达的观点，然后不要责怪教会的怜悯，这怜悯不会在基督的肢体聚集时驱散他们，反而努力聚集他四散的肢体合而为一。对那些当时立场最坚定、能够识破异端者诡诈言辞的人，确实，他们与其余人相比为数不多；但应当记得：其中有些人当时勇敢地承受着流放的判决，另有些人则为了安全躲藏在世界各地。这样，教会，就像主的麦子，在万民中增长，一直保全至今，并且将要保全至终，直到万民，甚至蛮族，也都怀抱在教会内。因为这就是人子所撒下好种子的那教会，他所预言将要与莠子一同生长，直到收割。田地就是世界，收割就是今世的终结。\BibleRef{13:24}\par

32. 因此，\Person{希拉里}要么不是在责备麦子，而是在责备\Place{亚细亚的十个省区}里的莠子，要么是在针对麦子说话，因为麦子由于某种不忠信而陷入危险，并且他的言论如同一个认为责备越严厉就越有益处的人。因为正典圣经中包含着以同样方式责备的例子，其中针对某些人的话却好像是对所有人说的。因此，当宗徒对格林多人说：\BibleQuote{怎么你们中间有人说没有死人的复活呢？}\BibleRef{格前 15:12}时，他清楚地表明他们并非所有人都是那样的；但他作证说，那些那样的人并没有脱离他们的共融，而是在他们中间。稍后，为了不让那些持不同意见的人被他们引入歧途，他给出了这警告：\BibleQuote{你们不要为人所欺瞒，滥交败坏善行。你们当醒悟，不为罪；因为有些人不认识天主：我说这话是为叫你们羞愧。}\BibleRef{格前 15:33}但当他说：\BibleQuote{你们中间既有嫉妒、纷争和分裂，你们岂不还是属血肉的人，按照人的方式行事吗？}\BibleRef{格前 3:3}时，他仿佛针对所有人说的，而且你看他提出的指控是何等严重。因此，若不是我们在同一书信中读到：\BibleQuote{我常为你们感谢我的天主，因天主的恩宠已在耶稣基督中赐给了你们；你们在祂内事事富足，在一切言语和一切知识上也是如此；正如我为基督所作的见证在你们中得以坚固，以致你们在任何恩赐上都不落后。}\BibleRef{格前 1:4}，我们会以为所有格林多人都属于血肉，属血气的人，不能领受天主圣神的事，\BibleRef{格前 2:14}喜爱纷争，充满嫉妒，并\Quote{按照人的方式行事}。同样，一方面经上说：\BibleQuote{全世界都卧在那恶者手下。}\BibleRef{若一 5:19}，这是因为全世界的莠子；另一方面，基督\BibleQuote{就是赎罪祭，为我们的罪过，不但为我们的，也是为全世界的罪过。}\BibleRef{若一 2:2}，这是因为全世界的麦子。\par

33. 然而，由于跌倒的事，许多人的爱情会冷淡；这些跌倒的事随着那些在基督圣事的共融内，为光荣祂的名而聚集起来的人中间，也有邪恶者以及坚执错误顽固者的增多而愈加普遍；不过，如同糠秕与麦子的分离，他们的分离只有在主的禾场最终清理时才能完成。\BibleRef{玛 3:12}这些人不能毁灭那些主的麦子——尽管与其余的人相比，人数确实少，但本身却是一大群人；他们不能毁灭天主的选民，这些选民要在世界末日时，从四方，从天这边到天那边，被聚集起来。\BibleRef{玛 24:31}因为这呼声正是出自选民：\BibleQuote{上主，求你拯救！因为虔敬的人已不存在，因为在人们的子孙中，忠实的人已绝迹。}并且主论到他们说：\BibleQuote{唯有坚持到底的（那时不法的事增多），必然得救。}\BibleRef{玛 24:12}此外，以下经文表明所引用的圣咏并非一个人的言语，而是多人的言语：\BibleQuote{上主，求你保护我们，求你从这一代永远护佑我们。}由于主所预言的不法之事的增多，别处也说：\BibleQuote{人子来临时，能在世上找到信德吗？}全知者的这一疑问预表了我们在祂内所经受的疑虑：当教会常常在许多曾被寄予厚望的人身上失望，而他们表现得与所期望的截然不同时，教会对自己的成员如此惊恐，以致迟迟不愿相信任何人的善。然而，怀疑那些祂将在世上找到信德的人正与莠子一起在整个田地中生长，是错误的。\par

34. 因此，这同一教会也在主的网内与坏鱼一起游泳，但在心灵和生活上却已与它们分离，并远离它们，好能作为\BibleQuote{光耀的教会，没有瑕疵，没有皱纹的}被呈献给她的主。\BibleRef{弗 5:27}但她所期待的、实际可见的分离只在海边，\Emph{即}在世界终穷时——其间她尽力纠正所能纠正的人，容忍她无法纠正的人；但她并不因发现那些无可救药者的邪恶而放弃善人的合一。\par

% structure: p0046 source=h2 target=subsection reason=relative-heading-rank; base=h2

\subsection{第10章}

35. 所以，我的弟兄，请勿再收集如此多、如此明显、如此无懈可击的神圣证据，来反对那些可能见于主教著作中的诽谤——无论是我们共融内的主教，如\Person{希拉里}，还是\Person{多纳图斯}裂教之前那未分裂的教会本身时代的主教，如\Person{西彼廉}或\Person{阿格里皮努斯}；因为，首先，就权威而言，这类著作必须与圣经正典区分开来。因为我们阅读它们，并非如同从其中提出的见证不可有任何异议，因为他们所持的观点，可能与真理要求我们赞同的并不相同。因为我们属于那些不拒绝，甚至也不拒绝宗徒所教导的：\BibleQuote{如果你们在什么事上意见不同，天主也要将这事启示给你们；然而，我们无论达到什么程度，都应该按同一规则行事。}\BibleRef{斐 3:15}——即，按基督的道路；关于这道路，圣咏作者如此说：\BibleQuote{愿天主怜悯我们，并降福我们；以自己的慈容光照我们！好让世人认识你的道路，万民认识你的救恩。}\par

36. 其次，如果那位主教和光荣的殉道者\Person{圣西彼廉}的权威吸引了你——我们确实，正如我已经说过的，认为这权威与\OriginalTerm{正典圣经}{canonical Scripture}的权威截然不同——那么你为什么不被他的如下事迹所吸引：他忠诚地维护并辩论捍卫了教会在这世上及万民中的合一；他谴责那些想要自命为义而脱离教会的人，指责他们充满自负和骄傲，嘲笑他们擅自取得连宗徒们也没有从主那里得到许可的事——即在收割之前收集莠子——并试图把糠秕从麦子中分离出来，仿佛清扫糠秕和清理打谷场的职责已经交给了他们；他证明没有人能被别人的罪所玷污，从而扫除了\OriginalTerm{分裂}{schism}制造者为其分离所提出的唯一理由；就在那教义问题上，当他与同僚意见不同时，他并没有下令谴责或\OriginalTerm{绝罚}{excommunicated}那些持有与他不同意见的人；甚至在他致\Person{犹巴亚努斯}的信中（这封信首次在那次公会议上宣读，而你们惯于援引该公会议的权威来辩护\OriginalTerm{再洗礼}{rebaptizing}的做法），尽管他承认过去在其他团体中受洗的人曾被接纳进入教会而未经第二次洗礼，并因此被他视为未曾受洗，但他仍然认为教会内和平的好处与利益是如此之大，以致为了和平，他主张这些人（虽然按他的判断他们是未曾受洗的）不应被排除在教会的职务之外？\par

37. 由此你将很容易明白（因为我知道你心思敏锐），你的主张完全被颠覆和消灭了。因为，如果真如你所假设，那曾传遍世界的教会因接纳罪人领受她的圣事而灭亡（而这是你们离去的理由），那么它早就在此之前完全灭亡了——即如西彼廉所说，当未受洗者被接纳进入教会之时——于是西彼廉本人也就没有在其中诞生的教会了；若是如此，那么对于像\Person{多纳图斯}那样属于较晚时代的人，即你们\OriginalTerm{分裂}{schism}的始作俑者和你们教派的创始人，岂非更不可能有教会？然而，如果在那个时期，尽管未受洗者被接纳进教会，教会却依然存在，从而诞生了西彼廉和后来的多纳图斯，那么显然义人并不会因参与圣事而与其他人一同被别人的罪所玷污。因此，你们没有任何借口可以清洗这分裂的罪，正是由于这分裂你们脱离了教会的合一；并且圣经上的这句话应验在你们身上了：\BibleQuote{有一类人自以为纯净，却没有洗除自己的污秽}\BibleRef{箴 30:12}。\par

38. 一个人出于对圣事相同性的考虑，即使对异端者也不敢坚持施行第二次洗礼，他因如此避免西彼廉的错误，并不因此在功绩上与西彼廉并驾齐驱，正如一个人不坚持外邦人遵守犹太礼规，并不因此而在功绩上与宗徒\Person{伯多禄}并驾齐驱一样。然而，在伯多禄的事上，\OriginalTerm{正典圣经}{canonical Scriptures}不仅记载了他的失足，也记载了他的纠正；而关于西彼廉曾持有与教会宪章和实践认可的见解相左的观点这一陈述，却不是见于正典圣经，而是见于他自己的著作以及一次大公会议的文献中；并且，尽管在同样的纪录中找不到他纠正了那个观点的证据，但推测他确实纠正了，以及这一事实或许被那些过于喜爱他所陷入的错误、不愿失去如此伟大人物的支持的人所隐瞒，这并非完全不合情理。同时，也不乏有人坚称西彼廉从未持有归于他的那种观点，而是说谎者假托他的名字进行的无法无天的\OriginalTerm{伪造}{forgery}。因为任何一位主教，无论多么著名，其著作的完整性和真实性都无法像正典圣经那样，通过被翻译成如此多的语言，以及通过教会公开敬礼中定期和持续使用的方式，得到保障和保存。即便如此，还是发现有人假托宗徒们的名字伪造了许多东西。的确，他们的这些企图是徒劳的，因为正典圣经的文本已如此受到证实，并被如此普遍地使用和认识；但这种邪恶的胆大妄为之举，竟敢攻击那些凭借如此广为人知的力度而受到保护的著作，这证明了它能够对那些不具备正典权威的著作做出何种尝试。\par

39. 然而，我们并不否认西彼廉持有归于他的那些观点：首先，因为他的文风具有某种可以辨认的表达特质；其次，因为在这件事上，我们的主张而非你们的主张被证明为得胜，而你们\OriginalTerm{分裂}{schism}所依据的借口——即你们可能不被别人的罪所玷污——以最简单的方式被推翻了；因为从西彼廉的书信中明显可知，那些按你们判断（且如你们所愿，也是按他的判断）是未曾受洗的人，就那样被接纳进了教会，容许罪人参与圣事，而教会却并未灭亡，反而保持了其本性所应有的尊严，即遍及世界的上主之麦。因此，如果你们在惊恐中这样投靠西彼廉的权威，如同投靠避难港，你们会看到在此航程中你们的错误所撞上的礁石；另一方面，如果你们不敢逃往那里，你们就连挣扎逃生的机会都没有就船毁沉没了。\par

40. 此外，\Person{西彼廉}要么根本没有持守你归咎于他的那些意见，要么后来用真理的规则纠正了自己的错误，要么用他丰厚的爱德遮盖了这个——我们可以称之为——他那无瑕心灵上的瑕疵，因为他最充分地捍卫了在全天下不断增长的教会的合一，最坚定地持守着和平的纽带；因为经上写道：\BibleQuote{爱德（爱）遮盖许多罪过。}\BibleRef{伯前 4:8}此外，还要加上：他像一根结果丰硕的枝条，父用苦难的修剪刀除去了他里面可能需要纠正的一切；主说：\BibleQuote{凡在我身上结果实的枝条，祂就修剪，使它结更多的果实。}\BibleRef{若 15:2}而父对他的这种关怀，岂不是因为他继续作为广阔葡萄树上的枝条，没有离弃合一的根源吗？\BibleQuote{因为，即使他舍身被火焚烧，若没有爱德，为他也没有益处。}\BibleRef{格前 13:3}\par

41. 现在请稍加注意\Person{西彼廉}的书信，好让你看出，他如何证明那些人不可原谅：他们表面上基于自己的义德，想要从教会的合一（天主所预许并在万民中成就的）中退出；也让你更清楚地理解我稍前所引经文的真义：\Quote{有一代人心目中自以为义，却未曾清洁自己离去的罪愆。}在他致\Person{安东尼安}的一封信中，他讨论了一个与我们正在辩论的极为相似的问题；但最好还是引用他本人的话。他说：\Quote{在我们此省的主教职务中，我们的一些前任认为，不应将教会的平安给予淫乱者，并最终向犯奸淫的人关闭了悔改之门。然而，他们并未撤出与同僚主教的共融；也未因如此严厉和固执地坚持批评而撕裂天主教会的合一，以致于因他人将教会的平安给予奸淫者而他们自己拒绝给予，从而脱离教会。和好的纽带未断，教会的圣事保持不分，每位主教安排和管理自己的行为，如同一个要为自己的行事向主交账的人。}对此，\Person{文森修}弟兄，你有何言？你必定看出，这位伟人，这位热爱和平的主教和无畏的殉道者，最关切者莫过于防止合一的维系被撕裂。你看到他如同为人的灵魂受产痛，不仅使他们在受孕时生于基督，而且也使他们在出生后不致因被摇离母怀而灭亡。\par

42. 现在，我请你进一步留意他在抗议不虔敬的分裂者时所提及的这一点。如果那些将平安给予悔改其罪的奸淫者的人分担了奸淫者的罪，那么那些不这样做的人是否因与同僚的共融而受玷污呢？再者，如果如真理所断言、教会所坚持的，将平安给予悔改其罪的奸淫者是正确的事，那么那些彻底向奸淫者关闭经由悔改的和解之门的人，无疑犯下了不虔敬的罪，因他们拒绝医治基督的肢体，夺走那些叩门求进者的教会的钥匙，并以无情的残酷反对天主的至慈容忍——天主容忍他们存活，好让他们借着悔改，以痛悔和谦卑的心的祭献而得医治。然而，他们这种无情的错谬和不虔敬并未玷污其他那些充满同情、热爱和平的人，因为后者与他们共享基督徒的圣事，并在合一的网中容忍他们，直到他们被带到岸边、彼此分开之时；或者说，如果别人的这种错谬和不虔敬果真玷污了他们，那么教会那时便早已毁灭，也就不存在一个能诞生\Person{西彼廉}的教会了。但是，如果无可置疑地，教会继续存在，那么同样无可置疑的是，在基督的合一内，任何人若不同意恶人的行为，并因此实际参与其罪行而受玷污，而只是为着与善人共融的缘故容忍恶人，如同容忍那躺在主的禾场上直待最终扬净的糠秕，他便不会被他人罪恶的罪责所玷污。事情既是如此，你们分裂的借口何在？你们岂不正是那一代\Quote{恶人，自以为义，却未曾洗净你们离去（离开教会）的罪愆}吗？\par

43. 倘若我现在有意引用\Person{提科纽斯}——你们共融中的一员，他著作的主旨是维护教会并反对你们，而非相反——的著作来反驳你们，他徒劳地否认非洲基督徒的共融为\OriginalTerm{背教者}{traditores}（\Person{帕梅尼安努斯}正是以此事使他缄默），那么除了\Person{提科纽斯}本人论及你们时所说、我方才提醒你们的那句话之外，你们还能作何回答呢：\Quote{凡合乎我们意愿的，便是圣的}？因为这位\Person{提科纽斯}——如我所言，是你们共融中的一员——写道，由你们二百七十位主教在\Place{迦太基}举行了一次会议；在此会议中，经过七十五天的审议，搁置了先前对此事的一切决议，公布了一项精心修订的决议，大意是：对那些犯有可怕罪行的背教者，若他们拒绝受洗，则应如同无罪之人一样，授予共融的特权。他还说，你们党派的\Place{马克里亚纳}的\Person{德乌特里乌斯}主教，将一大群背教者纳入教会，对他们与他人一视同仁，依照你们那二百七十位主教所举行的会议的决议，向这些背教者敞开了教会的合一之门，并且在那次行动之后，\Person{多纳图斯}未曾中断与前述\Person{德乌特里乌斯}的共融，不仅与他，也与所有\Place{毛里塔尼亚}的主教们维持了四十年的共融，这些主教按照\Person{提科纽斯}的陈述，接纳背教者共融而不要求他们重新受洗，直至\Person{马卡里乌斯}发动的迫害时期。\par

44. 你们会说：\Quote{那\Person{提科纽斯}与我何干？}诚然，\Person{提科纽斯}正是\Person{帕梅尼安努斯}以其回应所钳制，并有效警告其不得书写此类言论的人；但他并未驳倒那些陈述本身，而是如前所述，仅以此一事使其缄默：当\Person{提科纽斯}就那散布于全世界的教会说出这些话，并承认在其合一之内他人的过失无法玷污无辜者时，他自己却因\Place{非洲}基督徒是背教者而退出了与他们共融的交往，并成为\Person{多纳图斯}一派的支持者。\Person{帕梅尼安努斯}固然可以说，\Person{提科纽斯}在所有这些事上说的都是虚假的；但正如\Person{提科纽斯}本人所指出的，当时尚有许多人健在，由他们能证明这些事是毫无疑问地真实且众所周知的。\par

45. 然而，我不再就此多说：如果你们愿意，就坚持说\Person{提科纽斯}说的是假话吧；我把你们引回你们自己已引证的权威\Person{西普里安}。如果根据他的著作，教会合一中的每一个人都被其他成员的罪所玷污，那么教会在\Person{西普里安}时代之前早已彻底灭亡，而\Person{西普里安}自己（作为教会的一员）存在的任何可能性也被取消。然而，若连这样想都是不虔敬的，且教会一直存续无疑，则推知在公教合一之内，没有人会被他人之罪的罪责所玷污；因此，你们这\Quote{邪恶的世代}，声称自己是义人，而你们却\Quote{未洗去你们出离之罪咎}，乃是徒然的。\par

% structure: p0058 source=h2 target=subsection reason=relative-heading-rank; base=h2

\subsection{第十一章}

46. 你们会说：\Quote{那么你们为何寻找我们？为何接纳你们所称的异端者？}请留意我的回答是何等简单而简短。我们寻找你们，因为你们迷失了，好使我们在找到你们时欢喜，正如我们曾为迷失的你们忧伤。我们又称你们为异端者；但这名称只适用于你们尚未转向天主教会之和平、尚未从你们被其诱陷的错误中解脱之时。因为当你们来到我们这里时，你们便完全放弃了你们先前所持的立场，从而不再是异端者，而来到我们这里。你们会说：\Quote{那么给我施洗吧。}如果你们尚未受洗，或是领受了\Person{多纳图斯}或\Person{罗加图斯}的洗，而非基督的洗，我就会给你们施洗。使你们成为异端者的，并非基督徒的圣事，而是裂教之罪。由你们自身而来的恶，不是我们否认那存留在你们内的善的理由；那善若不存在于那赋予它行善效能的教会内，你们虽拥有它，却于己有害。因为主的一切圣事皆出自天主教会，你们持有并施行这些圣事，正如同你们从她分离之前所持有和施行的一样。然而，你们虽不再处于那赋予你们所拥有之圣事的教会内，但这并不减损你们依然拥有它们的事实。因此，我们不改变你们与我们一致之处，因为在许多事上你们是与我们一致的；论到这样的人，经上说：\Quote{他们曾在许多事上与我同心。}但我们要纠正你们与我们不一致之处，并愿你们领受那些你们在现今所在之处所没有的事物。在圣洗、信经和主的其他圣事上，你们与我们一致。但在合一精神与和平联络中，总之，在天主教会本身内，你们却与我们不一致。如果你们领受了这些，那么你们所已有的其他事物就不至于开始为你们所有，而将开始为你们所用。因此，我们并非如你们所想，把你们党派中的人当作仍属你们的人来接纳，而是在接纳他们的行动中，将那些离弃你们以便被我们接纳的人融入我们自身；为使这些人属于我们，他们的第一步就是断绝与你们的联系。我们也不是强迫那些殷勤服侍我们所憎恶之错误的人与我们合一；我们希望那些人合一于我们的理由，乃是使他们不再配受我们的憎恶。\par

47. 但你将会说：\BibleQuote{宗徒保禄在若翰之后施洗。}\BibleRef{宗 19:5} 难道他是在异端者之后施洗吗？如果你竟敢称这位新郎的朋友为异端者，并说他不在教会的合一之中，我请你将这话写下来。但如果你相信这样想或这样说乃是极度的愚妄，那么凭你的智慧，你仍须解决为何宗徒保禄在若翰之后施洗的问题。因为，如果他是在一个与他同等的人之后施洗，你们众人就应当彼此之后施洗。如果是在一位比他大的人之后，你们就应当在\Person{罗加图斯}之后施洗；如果是在一位比他小的人之后，\Person{罗加图斯}就应当在你之后，给你作为长老所施洗的人付洗。然而，倘若如今所施行的圣洗，对领受者而言，无论施洗者功德如何不等，都具有同等的价值，因为这是基督的圣洗，而非施洗者的权利，那么我想你必定已经明白，保禄之所以给某些人施以基督的圣洗，乃是因为他们只领受了若翰的洗礼，而非基督的圣洗；因为圣经明确称之为若翰的洗礼，正如圣经多处所证实的，也如主亲口所说：\BibleQuote{若翰的洗礼是从哪里来的？是从天上来的，还是从人来的？}\BibleRef{玛 21:25} 但伯多禄所施行的圣洗，不是伯多禄的，而是基督的；保禄所施行的圣洗，不是保禄的，而是基督的；那些在宗徒时代传播基督并非出于真诚，而是出于纷争的人，他们所施行的也不是他们自己的圣洗，而是基督的圣洗；那些在\Person{西彼廉}时代，或用狡诈的不义手段取得财产，或用重利盘剥的人，他们所施行的也不是他们自己的圣洗，而是基督的圣洗。正因为这圣洗属于基督，所以尽管施洗者的身份极不相称，对领受者而言却同等有效。因为，倘若每一次圣洗的卓越程度取决于施洗者本人的卓越程度，那么宗徒感谢天主，说他除了\Person{克黎斯颇}、\Person{加约}和\Person{斯德法纳}一家，没有给格林多的其他人施洗，\BibleRef{格前 1:14} 这就错了；因为如果由他自己施洗，格林多归信者的圣洗就会更为卓越，正如保禄本人比其他人更为卓越一样。最后，当他说：\BibleQuote{我栽种，\Person{阿颇罗}浇灌，}\BibleRef{格前 3:6} 他似乎暗示他宣讲了福音，而\Person{阿颇罗}施了洗。难道\Person{阿颇罗}比若翰更优秀吗？那么，为什么那位在若翰之后施洗的保禄，没有在\Person{阿颇罗}之后施洗呢？显然，因为前一种情况，无论由谁施洗，都是基督的圣洗；而后一种情况，无论由谁施洗，都只是若翰的洗礼，尽管它为基督预备了道路。\par

48. 在你看来，说有些人在若翰施洗之后还领受了圣洗，却又说在异端者施洗之后不应再为人施洗，这是一件可憎的事；但同样可以公正地说，有些人在若翰施洗之后领受了圣洗，而在不节制者施洗之后却不应再为人施洗，这也是一件可憎的事。我提这不节制的罪，而不提其他罪，是因为沉溺于此罪的人无法隐藏它；然而，即使是瞎子，谁又不知道到处有多少沉溺于这恶习的人呢？可是，在肉身的作为中——经上说，行这些事的人，必不能承受天主的国——宗徒列举这些事时，也将异端包括在内：他说：\BibleQuote{肉身的作为是显而易见的：即淫乱、不洁、放荡、崇拜偶像、施行邪法、仇恨、竞争、嫉妒、忿怒、争吵、不睦、分党、妒恨、凶杀、醉酒、宴乐，以及与这些相类似的事。我以前劝戒过你们，如今再说一次：做这种事的人，决不能承受天主的国。}\BibleRef{迦 5:19} 因此，尽管在若翰之后能施洗，却不能在一个异端者之后施洗，其原则完全相同：尽管在若翰之后能施洗，却不能在一个不节制者之后施洗；因为异端和醉酒都属于那些使人不能承受天主的国的作为。在你看来，难道这不是一件极其荒谬的事吗？即，圣洗在由那位连酒也不适度喝、而是完全戒绝、为天主的国预备了道路的人施洗之后，竟能重行，而在由一个不能承受天主的国的不节制者施洗之后，却又不能重行？对此能回答什么呢？只能说，前者是若翰的洗礼，宗徒在其后施行的是基督的圣洗；而后者，由不节制者施行的，乃是基督的圣洗。洗者若翰与不节制者之间，有极大的差异，如同对立；基督的圣洗与若翰的洗礼之间，并无对立，但有极大的差异。宗徒与不节制者之间，有极大的差异；但由宗徒施行的基督圣洗，与由不节制者施行的基督圣洗，却毫无差异。同样，若翰与异端者之间，有极大的差异，如同对立；而若翰的洗礼与异端者所施行的基督圣洗之间，并无对立，但有极大的差异。但是，由宗徒施行的基督圣洗，与由异端者施行的基督圣洗，却是毫无差异。因为圣事的形式被公认是相同的，即使施行的诸人在功德方面有极大的差别。\par

49. 但是，请宽恕我，因为我想通过一个放纵之人施洗的例子来说服你，这错了；因为我忘了，我所面对的是罗加提派，而非更广义的多纳特派。在你们人数稀少的同僚中间，以及你们全部圣职人员中，或许找不到一个沾染此恶习的。因为你们这些人认为，\OriginalTerm{公教}{Catholic}之名赋予信仰，并非因为持守此信仰者的共融遍及全世界，而是因为他们遵守全部神圣诫命与全部圣事；你们就是这样的人，人子来到时，唯独在你们身上会找到信德，而在地上却找不到信德，因为你们不是地，也不属地，而是属天的，居于天上！难道你们不畏惧，或不留意到：\BibleQuote{天主拒绝骄傲人，却赐恩宠于谦卑人}吗？\BibleRef{雅 4:6} 福音中那段话难道不使你们惊觉吗？主在那里说：\BibleQuote{人子来临时，能在世上找到信德吗？}\BibleRef{路 18:8} 紧接着，仿佛预见有人会骄傲地自诩拥有这信德，祂对那些自以为义而轻视别人的人，讲了两个人上圣殿祈祷的比喻：一个是法利塞人，另一个是税吏。此后的话语，我留给你自己去思量并回答。然而，请更仔细地审查你们那小教派，看看是否连一个施洗者都不是放纵之人。因为此一瘟疫对灵魂造成的破坏如此广泛，令我极其惊讶，它竟然没有蔓延到你们那极微小的羊群，尽管你们夸口说，早在基督——那位唯一的善牧——来临之前，你们就已经将绵羊和山羊分开了。\par

% structure: p0063 source=h2 target=subsection reason=relative-heading-rank; base=h2

\subsection{第十二章}

50. 请聆听那些作为主的麦子之人，通过我向你们所作的见证；他们目前在主的打谷场上，混杂于糠秕之中，忍受痛苦直到最终的簸扬，\BibleRef{玛 3:12} 即在全世界，因为\BibleQuote{天主从日出之地直到日落之处召唤了大地}，而在这同一广阔田野中，\BibleQuote{孩童们赞美祂}。我们不赞成任何人利用这道帝国法令来迫害你们，并非出于对你们改过的慈爱关怀，而是出于敌人的恶意。再者，虽然每一份世上的财产，只有依据神权——据此一切属于义人——或依据人权——此权在地上有君王的司法权内——才能正当保留，但你们错误地称那些你们并未作为义人拥有的、且已因地上君主的法律而丧失的财物为你们的，并枉然辩称：\Quote{我们辛苦积聚了它们}，因为你们可以读到经上的话：\BibleQuote{罪人的财富，是为义人积蓄的。}\BibleRef{箴 13:22} 尽管如此，我们不赞成任何人，利用地上君王为向基督表示敬奉而颁布的、旨在纠正你们不虔的这条法律，贪婪地图谋占有你们的财产。我们也不赞成任何人，不以正义而以贪婪为由，夺取并扣留原本属于穷人的供应，或你们聚会敬拜的礼拜堂——这些曾是你们以教会名义拥有的，并且无论如何只属于基督的真教会，那才是合法的财产。我们不赞成任何人，以与接纳那些在你们中间生活时除了与你们一起背离我们的错误外别无其他罪过的人相同的条件，来接纳因某些可耻行为或罪行被你们驱逐的人。但这些都是你们难以轻易证明的事；即使你们能证明，我们也会容忍一些我们无法纠正甚或无法惩罚的人；我们不会因为打谷场上有糠秕就离开主的打谷场，不会因为网中有坏鱼就撕裂主的网，不会因为羊群中有终将被分开的山羊就离弃主的羊群，不会因为主的屋内有遭耻辱的器皿就走出主的家。\par

% structure: p0065 source=h2 target=subsection reason=relative-heading-rank; base=h2

\subsection{第十三章}

51. 然而，我的弟兄，如果你放弃寻求来自人的虚浮荣耀，并轻看愚妄人的嘲骂——他们随时会说：\Quote{你为何现在毁掉你从前辛苦建造的？}——那么在我看来，毫无疑问，你现在将会转入那教会，我认为你已承认那就是真教会：因为我在手边找到了你这种想法的证据。在我此刻回复的你的信函开头，你有这样一段话：\Quote{我曾认识你，我卓越的朋友，在你仍远离基督信仰、早年从事文学研究时，已是一个致力于和平与正直的人；但自从你较近时归向基督信仰后，据许多人告知，你即把你的时间和辛劳投入了神学争论。} 如果你正是寄给我那封信的人，这些话无疑出自你口。因此，既然你承认我已归向基督信仰，尽管我并非归向多纳特派或罗加提派，你就无可辩驳地确认了：在罗加提派和多纳特派的界限之外，存在着基督信仰。因此，正如我们所说，这信仰已传遍万国，即那些按照天主的见证因亚巴郎的后裔而蒙福的万民。\BibleRef{创 22:18} 那么，为什么你在认识到那是真理之后，还犹豫不采纳它呢？除非因为你在过去未曾看清，或持别样观点，如今自觉羞愧——但更合理的羞愧应是——羞于纠正错误，却竟不以执迷不悟为耻（而此时羞耻本该更合情理）？\par

52. 这样的行为，圣经并没有保持沉默；因为我们读到：\BibleQuote{有一种羞耻带来罪，又有一种羞耻带来光荣和恩宠。}\BibleRef{德 4:21} 羞耻带来罪，当一个人因羞耻而不肯改正错误的观念，免得别人以为他反复无常，或以为他是根据自己判断承认自己长期错误时：这样的人活活地坠入阴府，也就是说，意识到自己的丧亡；他们的未来命运早已由\Person{达堂}、\Person{阿彼兰}和\Person{科辣黑}被地裂吞噬的死亡所预表。\BibleRef{户 16:31} 但是，当一个人为自己的罪而脸红，并藉着悔改转向更好的事物时，羞耻便是光荣和恩宠的；你却不情愿这样做，因为你被那种虚假而致死的羞耻所压服，生怕那些不知自己所说为何事的人，会用宗徒的这句话来反对你：\BibleQuote{若我再重建我所拆毁的，我就证明自己是罪犯。}\BibleRef{迦 2:18} 然而，如果这句话适用于那些改正之后宣讲他们曾因乖谬而反对的真理的人，那么当初它就可能被用来反对\Person{保禄}本人，因为基督的众教会曾因听见他现在\BibleQuote{传扬他从前所迫害的信仰}\BibleRef{迦 1:23} 而光荣天主。\par

53. 但是，不要以为一个人能从错误转向真理，或从任何罪过，无论大小，转到改正他的罪过，而不对他的悔改给出某种证明。然而，人若责备教会——教会已由如此多的天主的见证证明为基督的教会——说她以一种方式对待那些离弃她的人，即在接受他们回来时，条件是通过某种悔改的承认来纠正这过错，而以另一种方式对待那些从未在她怀抱中、并首次受到欢迎进入她平安的人，这是一种不可容忍的狂妄错误；她的做法是：更充分地谦抑前者，而以更宽松的条件接纳后者，对两者都怀有慈爱，并以母亲的爱心服侍两者的健康。\par

你在这里也许有比你期望更长的信。如果我在回复时只想到你一人，它会短得多；但即便如此，即使它对你没有用处，我也不认为它对那些在天主敬畏中、不凭外貌待人而用心阅读的人会毫无益处。阿们。\par

\SourceRecord{Translated by J.G. Cunningham.；Nicene and Post-Nicene Fathers, First Series；Vol. 1.；Edited by Philip Schaff.；Buffalo, NY: Christian Literature Publishing Co.,；1887.；<http://www.newadvent.org/fathers/1102093.htm>.}

% structure: p0001 source=h1 target=section reason=page-title

\section{书信第九十五封（主历408年）}

\Emph{致\Person{保利努斯}弟兄和\Person{特拉西亚}姊妹，最亲爱、真诚、堪受爱戴和尊敬的圣人，同在吾主耶稣门下同为门徒，} \Signature{奥古斯丁} \Emph{在主内致问候。}\par

1. 当那些与我们极为亲密的弟兄——您我和他们惯常彼此珍惜并表达大家由衷回应的亲切情怀——经常有机会拜访您时，这种神益，与其说令我们因善的增加而喜乐，不如说是在患难中给我们安慰。因为我们竭尽全力避免那些迫使他们奔波的原因和紧急情况，然而——我不知道怎么，除非是公正的报应——这些事却无法免除：可是当他们回到我们身边见到我们时，圣经上的话便在我们的经验中应验了：\BibleQuote{当我心中忧虑丛生时，祢的安慰令我的灵魂喜悦。}因此，当您从我们的弟兄\Person{波西迪乌斯}本人那里得知，迫使他去\Place{意大利}的缘由是多么令人悲伤时，您就会明白我才说的话多么真实，关乎他见到您的喜乐；然而，如果我们中任何人仅仅为了享受与您会面的快乐而漂洋过海，还有什么比这更迫切或更好的理由呢？但这样做，与我们因着职责所系、必须服侍那些因疾病而虚弱的人，而不能从他们那里抽身离开，是相违背的，除非他们的病情变得危险，使得离去成为必要。在这些事上，我不知道自己是在接受惩戒还是受审判；但我知道，祂并没有按我们的罪过对待我们，也没有照我们的恶行报复我们，却在我们受苦时融入如此大的安慰，并以令人惊叹的补救措施，确保我们不爱世界，也不致因世界而堕落。\par

2. 我在前一封信中曾询问您对圣人未来生命本质的看法；但您在回信中说，我们对于现世生活的状况仍需多加研究，您这样做很对，只是有一点不妥：您表示愿意从我这里学得那一问题的解答，而这事您要么与我同样无知，要么同样明了，甚至您比我知道的要多得多；因为您说得完全正确，在我们遭遇这必死之躯的消亡之前，我们必须按福音的意义，经由自愿的离去而死亡，不是藉着肉身的死，而是藉深思熟虑的决心，从今世的生活中退隐。这条道路是单纯的，没有疑虑的风浪困扰；因为我们主张，我们在此必死之生活中，应当以某种方式生活，好使自己多少能适于永生。然而，整个问题——当讨论和探究时，令我这样的人困惑不已——乃是：对于那些尚未学会经由死去而生活的人，我们应当如何生活在他们中间或为了他们的福祉而生活；那不是经由肉身的消亡，而是以某种属灵的决心，从纯自然事物的诱惑中转离。因为在大多数情况下，看来除非我们在那些我们切愿见他们脱离的事物上稍稍迁就他们，否则我们便无法为他们带来任何益处。而当我们这样迁就时，对这些事物的兴趣便偷偷潜入我们内心，以致我们常常乐于谈论和聆听轻浮的事，不仅对其微笑，甚至被大笑完全征服：如此，我们的灵魂便背负了依恋尘世尘埃、甚或世间泥淖的情感，我们感到更难以且更不情愿提升自己归向天主，好能藉福音之死度福音之生。每当这种心态达成时，随即就有称赞声说：\Quote{做得好！做得好！}这不是来自人，因为没有人能察觉他人内心领悟神圣事物的行动，而是在某种内在的静默中，不知从哪里传来：\Quote{做得好！做得好！}由于这类试探，伟大的宗徒承认，他受到天使的打击。\BibleRef{格后 12:7}看，这便是为何我们整个世间生活都是一个试探；因为人即便在尽可能效法天上生活的形态时，也难免在这事上受试探。\par

3. 对于惩罚的施加或免除，若是我们唯一的愿望只是促进那些我们判定应受惩罚或不应受惩罚的人的灵性福祉，我又能说什么呢？并且，如果我们不仅考虑过犯的性质和严重程度，也考虑每个人按其心灵的力量所能或不能承受的，那么调整惩罚的程度，以使受罚者不仅得不到好处，反而因此遭受损失，这是多么深奥难明的问题啊！此外，我不知道，当人们因畏惧来自人的这种惩罚的威胁而恐慌时，更多的人是改进了还是变得更坏了。那么，责任之路何在呢？因为常发生这样的事：你若惩罚一个人，他就毁灭了；而你若放过他，不施惩罚，另一个人却毁灭了。我承认，我在这方面天天犯错，不知道何时以及怎样遵守圣经的准则：\BibleQuote{在众人前斥责犯罪的人，为使其余的人有所警惕；}\BibleRef{弟前 5:20} 以及另一准则：\BibleQuote{在他和你独处时，指出他的过失；}\BibleRef{玛 18:15} 以及准则：\BibleQuote{时候未到，你们什么也不要判断；}\BibleRef{格前 4:5} \BibleQuote{你们不要判断，免得你们受判断}\BibleRef{玛 7:1}（在这条命令中，主没有加上\InnerQuote{时候未到}的话）；以及圣经上的这句话：\BibleQuote{你是谁，竟敢判断别人的仆人？他或站立，或跌倒，都由他自己的主人管；而且他必站得住，因为天主能够使他站得住。}\BibleRef{罗 14:4} 这些话清楚地表明，他指的是教会内的人；然而，另一方面，他又命令要判断他们，当他说：\BibleQuote{审断教外的人，与我何干？你们岂不是审断教内的人吗？所以，你们该把那个恶人从你们中间除去。}\BibleRef{格前 5:12} 但当这成为必要之时，对于应做到何种程度的问题，会引发何等的谨慎与恐惧，以免发生那事，就是宗徒在给他们的第二封书信中劝诫他们提防的，并以那个例子为证，说：\BibleQuote{免得这样的人为过度的忧愁所吞噬；} 他又加上说，为免人们认为这事无需焦虑挂心：\BibleQuote{免得撒殚占了我们的便宜，因为我们不是不知道他的诡计。}\BibleRef{格后 2:7} 面对这一切，我们感到何等战栗，我的兄弟\Person{保利诺}，天主的圣者啊！何等战栗，何等黑暗！我们岂不是可以认为，针对这些事，经上说了这话：\BibleQuote{惊惧与战栗侵袭了我，恐怖笼罩了我。我便说：但愿我有鸽子般的翅膀，那样我就能飞翔而去，得以安息。看，我必逃往远处，在旷野里暂住。} 然而，即便在旷野里，也许他仍然经历这事；因为他接着说：\BibleQuote{我期待那位救我脱离软弱和风暴的主。} 因此，的确，人在世上的生活是一场试探。\BibleRef{约 7:1}\par

4. 再者，关于\OriginalTerm{天主的神谕}{oracles of God}，我们岂不是只轻轻触及，而非紧紧把握并透彻理解吗？因为其中绝大部分，我们还没有形成明确而确定的见解，而是仍在探究我们应有何种见解。这种谨慎，虽然伴随着极大的不安，却远胜于武断的鲁莽。再者，如果一个人不属血肉之见（宗徒说这是死亡），那么他岂不会在许多圣经章节的阐释中，大大得罪那些仍属血肉之见的人吗？在这些阐释中，说出你所信的，极其危险；闭口不言，又极其痛苦；说出与你所信不同的话，更是极其有害。更有甚者，当我们在教会内人士的讲论或著作中发现一些应受指责的事，并且不掩饰我们的不赞同（假定这样的纠正符合兄弟友爱的自由），我们若是被人怀疑这样做是出于嫉妒而非善意，那又是何等的罪过加在我们身上！而当我们同样地把那些批评我们意见的人，看作是意在伤害而非纠正我们时，我们又是何其得罪他人啊！确实，通常由此便生出苦毒的仇视，即便在彼此以最深感情和亲密相连的人之间也是如此，当 \BibleQuote{思想人超过所记载的，就有人为了某一人的缘故而自高自大，反对另一人。}\BibleRef{格前 4:6} 当他们彼此相咬相吞时，\BibleQuote{就有理由害怕，免得彼此消灭。}\BibleRef{迦 5:15} 因此，\BibleQuote{但愿我有鸽子般的翅膀，那样我就能飞翔而去，得以安息。} 因为，无论是由于人身处其中的危险在他看来，比他未曾经历的危险更大，还是我的印象确实正确，我都不禁认为，旷野中的任何软弱与风暴，都比我们在繁忙世界中所感受或畏惧的事更容易忍受。\par

因此，我极表赞同你的说法：我们应以人们在此生所处的状态，或更确切地说，他们所跑的赛程，作为讨论的主题。我之所以主张优先讨论此事，还有另一个理由：必须先寻见并遵循那赛程本身，然后才能寻见并拥有它所导向的目标。因此，当我征询你对此事的看法时，我的态度就好像我们已经谨慎地持守、奉行了今生的正道，因而对它的行程已无忧虑；虽然我自觉在许多方面，尤其在我所提及的那些事上，我仍在极大的危险中劳苦。不过，就我看来，这整个无知与困惑的原因在于：在极为不同的生活方式与心性之中——这些心性的倾向与软弱完全隐藏在我们的视线之外——我们寻求的乃是那些公民与臣民的利益，他们并非属于地上的\Place{罗马}，而是属于天上的\Place{耶路撒冷}。因此，我觉得与你谈论我们将来的状态，比谈论我们现在的状态，更为惬意。因为，虽然我们不知道那里将要享受的福分是什么，但至少有一件事我们可以确信，且此事非同小可：那里绝不会有我们在此地所经历的种种祸患。\par

因此，就安排此生生活的途径——我们必须遵循此途径以获得永生——而言，我知道我们必须克制肉欲，只容许满足身体感官的程度足以维持此生并积极履行它的义务；而且，凡与此生联系在一起的、因天主的真理以及我们或邻人永恒福祉而临于我们的种种烦恼，都必须以忍耐和坚毅来承受。我也知道，我们应当以全部的爱德热忱寻求邻人的利益，使他能善度此生，以获得永生。我还知道，我们应喜爱属神的事物胜于属血肉的事物，喜爱不变的事物胜于可变的事物；一个人能做到这一切的程度，全看他得到多少恩宠的助佑——这恩宠是藉着我们的主\Person{耶稣基督}而来的。但我不知道为何这人或那人或多或少地得到那恩宠的助佑，或不受助佑；我只知道，天主这样做是至公义的，并且祂有自己认为充分的理由。至于我上面所提到的那些事，即我们应当在人群中如何生活，如果你透过经验或默想有所领悟，我恳请你指教。如果这些事也让你如我一样困惑，那么请你寻找一位贤明的神修良医（或许你可以在你居住的地方找到，或当每年访问罗马城时在\Place{罗马}找到），与他一起研讨此事，之后写信给我，无论主藉着他的教导，或当你们就此课题交谈时向你们两人所启示的，都请告诉我。\par

关于肉身的复活，及其肢体在不可朽坏、不死不灭状态中的未来功用，既然你在回覆我的问题时也询问了我对这些事的看法，那么请听我简短的陈述；如不充分，日后可仰赖主的助佑，通过更充分的讨论予以扩充。我们必须最坚定地持守这一教义：圣经对此的见证真实无误，即这些可见的、属土的身体，现今称为\OriginalTerm{属生灵的身体}{natural body}的，在信者和义人复活时，将成为\OriginalTerm{属神的身体}{spiritual body}。与此同时，我们既不知如何理解或描述属神身体的特性，因为这事超出我们的经验范围。毫无疑问，这样的身体将没有朽坏，因此它们不再需要现今所必需的可朽坏的食粮；然而，虽然不需要，它们若要随意取用并实际进食，也并非不可能；否则，复活后的主就不会进食了，祂既给了我们肉身复活的榜样，以致宗徒据此论证说：\BibleQuote{如果死人没有复活，基督也就没有复活。} \BibleRef{格前 15:16}，但祂显现给门徒时，拥有全部肢体，并照其功用使用它们，又指给他们看祂的伤痕所在之处——关于这些伤痕，我一向认为那只是疤痕，并非伤口本身，而且它们存在于那里，不是出于必要，而是按祂自由施展的大能。祂当时极其清楚地显明了自己施展这大能何等容易：一方面，祂显现给两位门徒时以另一个形像出现；另一方面，当祂在紧闭房门的楼上向门徒们显现时，不是以幽灵的形态，而是以自己真实的身体。\par

8. 由此引出了关于天使的问题：他们是否具有适应其职务和自一处迅速移至另一处的身体，或者仅是\OriginalTerm{神体}{spirits}？因为如果我们说他们有身体，我们就会面对这一段经文：\BibleQuote{祂使自己的天使成为神体；}而如果我们说他们没有身体，那么解释如下事情时就会遇到更大的困难：如果他们没有形体，那么经上记载他们向人的身体感官显现、接受款待、任人洗脚，并食用为他们预备的饮食，这该如何解释。因为，如果我们设想，天使被称为神体，与人被称为灵魂的方式相同（\Emph{例如}，圣经说，那些与雅各伯一同下到埃及去的有那么多灵魂，这并非意指他们也没有身体，见：\BibleRef{创 46:27}），这似乎比设想天使没有形体却做了这一切事情所遇到的困难要少。再者，在\BibleBook{若望}{默示录}中，明确提到了一个天使的确定高度，其尺寸只能适用于身体，这证明那显现于人眼前的事，不应解释为幻觉，而应解释为源于我们所说的\OriginalTerm{属神的身体}{spiritual body}易于发挥的那种能力。然而，无论天使是否有身体，也无论是否有人能够说明没有身体他们如何能做这一切事情，有一点是确定的：在圣者的城中，我们人类中那些蒙基督救赎的人将永远与成千上万的天使联合在一起；在那里，由言语器官发出的声音将表达胸怀之中毫无隐藏的思想；因为在那种神圣的共融中，一个人心中的任何意念都不可能对另一个人隐藏，而是将在赞颂天主时完全和谐与同心，这不仅发自精神，而且通过\OriginalTerm{属神的身体}{spiritual body}作为其工具；至少，这是我所相信的。\par

9. 与此同时，如果你已经找到或能从其他老师那里学到任何比这更完全符合真理的见解，我极渴望从你那里得到教诲。请你察阅我的信；对于这封信，你曾以带信给我的执事行程匆忙为理由，解释你的答复非常简短，我对此并无抱怨，而是提醒你，以便你将在那次答复中省略的内容，现在予以补充。请再查阅一遍，注意我希望从你那里学到的，既包括你对于基督徒的\OriginalTerm{退修生活}{Christian retirement}作为获取并讨论基督徒智慧真理的一种方法有何看法，也包括那种我原以为你在其中找到闲暇、但有人告诉我说你正不可思议地忙碌于事务的退修生活。\par

愿圣洁的天主因你赐给了我们极大的喜乐和安慰，也愿你在真正的幸福中，常记念我们而生活。\EditorialNote{此句为另一人所加。}\par

\SourceRecord{Translated by J.G. Cunningham.；Nicene and Post-Nicene Fathers, First Series；Vol. 1.；Edited by Philip Schaff.；Buffalo, NY: Christian Literature Publishing Co.,；1887.；<http://www.newadvent.org/fathers/1102095.htm>.}

% structure: p0001 source=h1 target=section reason=page-title

\section{第 96 封书信（主后 408 年）}

\Emph{致 \Person{奥林匹乌斯}——我至爱的主，在 \Person{基督} 内配得尊崇与礼遇的儿子，\Person{奥古斯丁} 致候。}\par

1. 无论你在今世的进程中拥有何等地位，我都满怀信心地向你，我深爱且真诚的基督徒同仆 \Person{奥林匹乌斯} 致函。因为我知道，在你看来，这个名号远胜一切光荣崇高的头衔。确实有消息传到我这里，说你在世俗尊荣上得到了擢升，但直到我抓住这次给你写信的机会时，还未接到证实此传闻的讯息。然而，既然我知道你从主那里学会了\BibleQuote{不可心高妄想，却要俯就卑微的人}，那么，无论你被提升到何等高位，我至爱的主，我堪当尊敬和重视、为主 \Person{基督} 的肢体的儿子，我们确信你仍会一如既往地乐意接受我的来信。至于你世间的昌达，我毫不怀疑你会明智地运用它，为获享永恒的利益；以致你在世上国家中的影响力愈大，你便愈献身于那个你因 \Person{基督} 而得诞生的天上之城的利益，因为这将要在活人之地、在那赋予永福的真平安中更丰厚地赏报给你。\par

2. 我再次恳请你善意考虑我的弟兄兼同工 \Person{波尼法爵} 的请求，希望以前不能办到的事，现在在你的权力下能够实现。他也许确实可以依法毫无困难地，保留他的前任以他人名义获得、并以教会名目开始占有的那些东西；但我们不希望，由于他的前任拖欠国库税款，而使我们的良心背负重担。因为那种欺诈行为，并不因其损害了公共收入而减少其真正的欺诈性。那位前任 \Person{保禄}，在受任命为主教时，因累积欠缴国库的债务而即将交出他所有财产，他取得了一笔到期债券的支付，便以当时一个极有权势的家族的名义，仿佛是为教会购买了这片田地，靠其出产维持生计，以便按照他的老习惯，在这些事上也能逃避税吏的麻烦，尽管他根本不纳税。然而，当 \Person{波尼法爵} 在同一教会受祝圣就任后，犹豫是否该接受 \Person{保禄} 如此拥有的田地；虽然他满可以仅向皇帝请求减免那位前任在此小额财产上所欠的税款，但他宁肯毫无保留地承认：\Person{保禄} 是以自己的资金在拍卖时买下这产业，而当时 \Person{保禄} 已因拖欠国库而破产。这样一来，教会如今若有可能获得这产业，便不是通过主教暗中的欺诈，而是经由基督徒皇帝公开的慷慨恩赐。倘若此事办不到，天主的仆人们情愿忍受匮乏之苦，也不愿在良心受不义交易谴责之下获得所需供给。\par

3. 我恳请你俯允支持这请求，因为他已决定不提出先前获得的有利判决，以免妨碍他再度提出申请的自由；因当时得到的答复未能满足他的愿望。如今，既然你仍怀有往昔的仁厚心肠，且拥有更大的影响力，我不失望：凭着你与皇帝的关系，靠主的助佑，这事会容易获准；即使你以你自己的名分请求赐予这产业，然后赠予我所说的那个教会，谁又能挑剔你的请求呢？不，反而谁不称赞这并非出于私欲，而是出于基督徒虔敬的善举呢？愿主、我们天主的仁慈荫庇你，并使你在 \Person{基督} 内日益喜乐，我主我儿。\par

\SourceRecord{Translated by J.G. Cunningham.；Nicene and Post-Nicene Fathers, First Series；Vol. 1.；Edited by Philip Schaff.；Buffalo, NY: Christian Literature Publishing Co.,；1887.；<http://www.newadvent.org/fathers/1102096.htm>.}

% structure: p0001 source=h1 target=section reason=page-title

\section{第97封信（主历408年）}

\Emph{致\Person{奥林皮乌斯}，我杰出而公正受尊敬的阁下，我在基督内极应受尊荣的儿子，\Person{奥古斯丁}在主内致问候。}\par

1. 虽然我们最近听到你荣获晋升至最高品级，尽管我们对这消息的确实性多少有些不确定，但仍确信你对教会——我们欣喜得知你真是教会的儿子——没有别的心情，如同你在信中向我们清楚表明的那样；然而，现在读了你的信，你在信中惠然主动向充满迟疑和缺乏信心的我们，发出最仁慈的劝勉，要我们用卑微的努力指出，主——藉其恩赐你才如此有权力——可如何借着你的虔敬顺从，时时援助他的教会；我们便因这更大的把握而写信给你，我卓越而公正受尊敬的阁下，我在基督内极应受尊荣的儿子。\par

2. 许多弟兄，即与我同僚的圣洁之人，因这里教会的困扰，已——我几乎可说是如同逃亡者——前往皇帝最堂皇的朝廷；这些弟兄你或已见到，或可能已从\Place{罗马}收到他们因各自上诉事由所写的信。我写信前无法与他们商议；然而，我不愿错过由带信人——我的弟兄和同僚司铎——带信的机会；他因迫不容缓的急需，为救一位同胞的性命，不得不在隆冬季节赶赴那方。所以我写信问候你，并凭你在主耶稣基督内所有的爱嘱咐你，务必以最大的勤奋迅速推进你的善工，好叫教会的敌人知道，那些关于拆毁偶像和纠正异端者的法律——在\Person{斯提里科}尚在生时便发到\Place{阿非利加}——是出于我们最虔敬忠信的皇帝的意愿；因为他们或狡猾地吹嘘，或不由自主地以为这些事是在皇帝不知情或违背他意愿下进行的，如此他们使无知者的心中充满煽动性的暴力，并挑动他们对我们的危险而激烈的敌意。\par

3. 我的确不怀疑，我在向尊驾呈递这请愿或恭敬的提议时，是与\Place{阿非利加}所有同僚的愿望一致的；并且我认为，贵职有责任采取行动，这事很容易办，只要一有机会，便让这些虚妄的人（我们虽寻求他们的救恩，他们却抗拒我们）明白，那些为保卫基督教会而发至\Place{阿非利加}的法律，其颁布并非出于\Person{斯提里科}的关注，而是出于\Person{狄奥多西}之子的关注。因此，我已提到的这位司铎，即本信的带信人，来自\Place{米莱维}地区，是奉他的主教——可敬的\Person{塞维鲁}之命——路过我所在的\Place{希波}；塞维鲁同我一起向你衷心问候，我们珍视你最真诚的爱。因为当我们在教会遭受严重磨难和痛苦时偶然相遇，我们寻找致书尊前的机会，却未找到。我曾确实为我们的圣善兄弟和同僚\Place{卡塔夸}主教\Person{博尼法斯}的事务送过一封信；但那时更沉重的灾难尚未临到我们，这些灾难注定要使我们更加不安；关于这些事，以及如何用基督的方法以最佳策略预防或惩罚，那些为此任务从本处乘船而去的主教们将能更方便地与你商谈，你在他们心中的善意令我们欣喜，因为他们能向你报告因时间有限而经周密共同商议的结果。但至于另一事，即让全省知晓我们至为慈祥、虔敬的皇帝对教会的态度，我建议，不，我请求、恳求并哀求足下，留心不要耽误时间，而要在你见到从我们这里出发的主教们之前，运用你对目前处于极大危险之中的基督的肢体所拥有的最卓越警惕，尽快促成此事；因为主在这些试炼中给了我们不小的安慰，因他乐于使你现在较从前拥有更大的权力，虽然我们早已因你善工的数目和重大而充满喜悦。\par

4. 我们为某些人的坚定稳固的信仰大大欢喜，这些人并非少数，他们借着这些法律归向了基督教，或从\OriginalTerm{裂教}{schism}归向了\OriginalTerm{公教}{Catholic}的和平，为了他们永恒的福祉，我们甘冒失落暂时福祉的风险。因为特别是在这一方面，我们现在不得不忍受来自极其顽固悖逆之人的更猛烈的敌意攻击，他们中的一些人同我们一起耐心忍受；但我们因他们的软弱而非常惧怕，直到他们学会，并靠着主至仁慈的恩宠帮助，能以更丰富的灵魂力量轻看今世和人的短暂岁月。若我所料不差，我的主教弟兄们若尚未到你那里，当他们抵达时，请惠然将我已送去的训令书交给他们。因为我们对你心中的真诚热心如此信任，以致我们渴望借着主的帮助，你不仅援助我们，而且参与我们的商议。\par

\SourceRecord{Translated by J.G. Cunningham.；Nicene and Post-Nicene Fathers, First Series；Vol. 1.；Edited by Philip Schaff.；Buffalo, NY: Christian Literature Publishing Co.,；1887.；<http://www.newadvent.org/fathers/1102097.htm>.}

% structure: p0001 source=h1 target=section reason=page-title

\section{第九十八封信（主后408年）}

\Emph{致他的主教同僚\Person{博尼法斯}，\Person{奥古斯丁}在主内致以问候。}\par

1. 你问我：\Quote{当父母在子女患病时，试图通过祭献异教伪神来救治他们，是否会对已受洗的子女造成伤害？} 又，\Quote{如果他们这样做对子女无害，那么这些子女受洗时，因着父母信德的功劳所得的任何益处，又怎能与父母背离信仰时对子女毫无损害这件事协调一致呢？} 对此，我的答复是：在基督奥体各肢体的神圣结合中，那带来救恩的圣事，即圣洗的效能如此伟大，以至于一旦那因他人（基于自然本能行事）而获得第一次出生的人，经由他人（基于属灵渴望行事）而分享第二次出生，他便不能自此被那另一个人所犯之罪的束缚所捆绑，因为他并未以自己的意志同意那罪。主说：\BibleQuote{父亲的生命和儿子的生命都属于我：谁犯罪，谁必死亡。}\BibleRef{则 18:4} 但如果父母或任何其他人在他不知情的情况下，为了他而陷入崇拜异教神明的亵渎行为，这孩子本人并没有犯罪。他由亚当所承受的那罪债的束缚，本应藉此圣事的恩宠予以消除；之所以如此，是因为在亚当犯罪时，他尚未是一个拥有独立生命的灵魂，\Emph{即}另一个可称\InnerQuote{父亲的生命和儿子的生命都属于我}的灵魂。因此，如今当人拥有个人的、独立的存在，从而与父母区分开来时，他便不再为那未经他同意而由他人所犯的罪负责。在前一种情况下，他由他人承受罪债，因为在他所承受的罪债被招致之时，他与那承受来源之人原是一体，且在他内。但当各人拥有各自独立的生命时，那句\Quote{犯罪的生命必死亡}的话便同样适用于两人，因此，一人不再由他人承受罪债。\par

2. 然而，当孩子被带来领受这神圣礼仪时，藉由他人意愿所履行的职务而得以重生的可能性，完全是圣神的工程，因为那被带来的孩子正是借着这位圣神而重生的。因为经上不是记载：\Quote{人除非因父母的意愿，或因献上孩子者的信德，或因施行这礼仪者的信德而生}，而是：\BibleQuote{人除非由水和圣神而生}\BibleRef{若 3:5}。因此，借着水——它以外在形式展示恩宠的圣事，并借着圣神——祂以内在能力赐予恩宠的益处，消除罪债的束缚，恢复\OriginalTerm{本性的善}{bonum naturæ}，那原本只由亚当获得第一次出生的人，便唯独在基督内得以重生。如今，那位使人重生的圣神既为献上孩子的父母所共享，也为那被献上并重生的婴儿所共享；因此，因着这同一圣神的分享，献上婴儿者的意愿便对该婴儿有益。但当父母因将孩子献于异教伪神，并试图将其置于这些伪神的不虔束缚之下而犯罪得罪孩子时，父母与孩子之间并不存在那样的灵魂相通，以致一方的罪债能归双方共同承担。因为我们不是通过他人的意志而成为罪债的分享者，就像我们通过圣神的合一与他人同享恩宠那样；因为那唯一的圣神能够临在于两个不同的人内，而他们彼此却不知晓因着祂，恩宠是他们共有的财产；但人的神魂却不能那样属于两个个体，以致在二人中一者犯罪，另一者不犯罪的情况下，使罪责归双方共同承担。因此，一个孩子一旦通过父母获得了自然的出生，便能由天主的圣神成为第二次（即属灵的）出生的分享者，从而使那由父母继承来的罪债束缚得以消除；但一个已经由天主的圣神获得这第二次出生的人，却不能再次通过父母成为自然出生的分享者，以致那已被消除的束缚重新捆绑他。因此，一旦领受了基督的恩宠，孩子除非因自己的不虔（若长大成人后变得如此败坏）才会失掉它。因为到那时，他便会开始有自己的罪，这些罪不能由重生来消除，而必须由其他的补救措施来医治。\par

3. 然而，那些有着更深忧惧的人，无论是带着自己孩子的父母，还是带着其他幼儿的人，若试图使已受洗者承担向异教神明进行亵渎崇拜的义务，便犯下了\OriginalTerm{灵魂的谋杀}{spiritual homicide}。的确，他们并未实际杀害孩子们的灵魂，但却尽其所能地将他们推向死亡。警告\Quote{不要杀害你们的小孩子}，可以恰如其分地对他们说；因为宗徒说：\BibleQuote{不要消灭圣神；}\BibleRef{得前 5:19}并非圣神能被消灭，而是那些行事如同想要消灭祂的人，理应被称为\OriginalTerm{消灭圣神者}{quenchers of the Spirit}。在这个意义上，也能正确理解至福者\Person{西彼廉}在其论\OriginalTerm{背教者}{the lapsed}的书信中所写的话，当他斥责那些在迫害时期向偶像献祭的人时，他说：\Quote{为了使他们的罪行达到满盈，他们的婴孩，或被抱在怀中，或被父母牵着手领到那里，在襁褓之中便失去了他们刚一出生就领受的恩宠。}他的意思是，他们失去了它，至少从那些迫使他们遭受此损失之人的罪咎而言：也就是说，是在那些对他们犯下如此错误之人的意图和愿望中失去了它。因为假如他们真的亲自失去了它，他们必然在天主的定罪判决下无可申诉；但如果圣\Person{西彼廉}持有这种观点，他就不会在紧接的上下文中加上为他们的辩护，说：\Quote{难道这些人在审判之日到来时不会说：\InnerQuote{我们什么也没有做；我们没有自愿匆忙参与不洁的仪式，离弃主的饼和杯；别人的背信导致了我们的毁灭；我们发现我们的父母是杀人犯，因为他们剥夺了我们的慈母教会和我们的天主圣父，以致因着别人的过错，我们陷入罗网，因为我们尚年幼，不能自己思考，我们因别人的行为，且对这种罪行完全无知，而成为他们罪行的同伙}吗？}如果他不是相信这些辩护是完全正当的，并且在天主审判的法庭上对这些婴儿有益，他就不会附加上这些辩护。因为如果他们确实所说为真：\Quote{我们什么也没有做，}那么\BibleQuote{那犯罪的灵魂，必死亡；}而在天主公义的审判安排下，那些其父母（至少就他们在这件事上的罪咎而言）已将其灵魂引向毁灭的人，不会被判定沉沦。\par

4. 关于同一封信中提到的事件：一个女婴被留给奶妈照料，她的父母因突然逃亡而失踪，奶妈迫使她参与了偶像崇拜的不洁仪式；后来在教会中，当给她\OriginalTerm{圣体}{Eucharist}时，她以奇妙的动作将圣体从口中吐出——在我看来，这事是出于天主的干预，为了让成年人不会认为在这罪中他们没有亏待孩子，反而能通过那些不会说话者具有明显意义的身体动作，明白一个奇迹般的警告赐给了他们，警告他们：在如此大罪之后，应该采取怎样的行动才相称——他们原本理当以各种方式，为证明补赎，远离那些圣事，然而他们却冒失地冲向了它们。当天主眷顾借婴儿成就此类事时，我们不可相信他们是在知识和理性的引导下行事；正如我们不必因天主曾乐意借驴的声音斥责先知的愚妄\BibleRef{户 22:28}，就去赞赏驴的智慧。因此，如果一个无理性的动物发出了完全像人声的声音，且这当归因于天主的奇迹，而非驴本身的能力，那么全能者同样能借婴儿的精神（其中理性并非不在，只是沉睡未发展），通过其身体的运动显明某事，使那些既得罪自己灵魂又得罪自己孩子的人应当留意。但由于孩子不能重新成为其生身父母的一部分，与他合一且在他内，而是一个完全独立的个体，拥有自己的身体和灵魂，\BibleQuote{那犯罪的灵魂，必死亡。}\par

5. 有些人确实带他们的幼小者来受洗，不是出于信德的期待，期望他们借着\OriginalTerm{灵性恩宠}{spiritual grace}重生而获得永生，而是因为他们以为借此作为疗法，可使孩子恢复或保持身体健康；但不要让这使你们心神不安，因为他们的重生并不因那些带他们来受洗的人无意于此恩福而被阻止。因为由这些人，完成了所必需的\OriginalTerm{职任行动}{ministerial actions}，并宣念了圣事经文，没有这些，婴儿不能祝圣于天主。但居住在圣徒内的圣神——即那被炽热的爱火融合成那只翅膀镀银的鸽子——甚至借着奴仆的牧职，亦即那些有时不仅因无知而糊涂、甚至因可咎的缺失而不配为祂所用的人，完成祂的工程。呈现幼小者以领受灵性恩宠，与其说是那些亲手怀抱他们之人的行为（虽然若他们自己是好信徒，这也部分地是他们的行为），不如说是整个圣徒和信徒的团体。因为理应视婴儿为所有喜悦他们受洗的人所呈献，并借着这些人圣洁而完全合一的爱情，他们获得协助，得领受圣神的共融。因此，这是由整个母教会所行的，这母教会存在于圣徒内，因为整个教会是所有圣徒的母亲，整个教会是每一个圣徒的母亲。因为如果基督徒圣洗圣事常是同一的，即使由异端者施行时也有价值，且虽然在这种情况下不足以保证受洗者分享永生，却足以印证他对天主的祝圣；而如果这祝圣使那有主印记却留在主的羊栈外的人，成为异端者而负有罪责，但同时也提醒我们，他应由健全的教义纠正，而不应再次藉重复该仪式而重新祝圣——如果这在异端者的洗礼中尚且如此，那么在天主教会内，那只是糠秕的东西，竟能有助于承载麦子到将要簸扬的禾场上，并借此使之准备好被加入好麦子的堆中，岂不更加可信！\par

6. 此外，我愿你们不要继续错误地以为，继承自亚当的罪债，除非由父母亲自把他们的小孩带来领受基督的恩宠，否则无法被免除；因为你们写道：\Quote{正如父母是那使他们应受惩罚的生命的制造者，子女也应由同一渠道，即同一父母的信德，获得成义；}而你们却看到许多人不是由父母呈献，而是由任何陌生人呈献，正如有时奴隶的幼年子女由其主人呈献。有时，父母去世后，幼小的孤儿由那些有能力以此方式表达怜悯的人呈献而受洗。再者，有时无情的父母抛弃的弃婴，本为让任何路过的人照顾，却被圣洁贞女捡起，并由这些既没有也不愿有自己子女的人呈献受洗：在这事上，你正看到了福音中所提到的那件事，即那人被强盗打伤，半死地被抛在路上，主曾问谁是他的近人，得到的回答是：\BibleQuote{是怜悯他的那人。}\BibleRef{路 10:37}\par

7. 你把放在你那系列问题末尾的问题，判断为最困难的，因为你惯于小心避免任何虚假；你这样说：\Quote{如果我把一个婴儿放在你面前，问：\InnerQuote{这孩子长大后会不会贞洁？}或\InnerQuote{他会不会成为盗贼？}你会回答：\InnerQuote{我不知道。}如果我问：\InnerQuote{他目前在婴儿状态下，是想着善还是想着恶？}你会回答：\InnerQuote{我不知道。}因此，如果你不敢就他日后的行为或此刻的思想承担任何明确陈述的责任，那么父母们所做的，在他们为子女受洗，作为\OriginalTerm{保证人（或代父母）}{sureties (or sponsors)}代为答复，并声称他们做了在那个年龄甚至不能理解的事，或者至少，关于这些事，他们的思想（若他们能思考）对我们来说是隐藏的，这又算什么呢？因为我们问那些呈献孩子的人：\InnerQuote{他信天主吗？}而在那个年龄，孩子甚至不知道有天主，代父母却回答：\InnerQuote{他信；}对其他每个问题，他们也同样作答。现在我惊奇的是，在这些事上，父母竟能如此自信地代表孩子说，当他们在回答施行洗礼者的问题时，那婴儿正在做如此值得称赏、如此卓越的事；然而，如果同一时刻我加上这样的问题，如：\InnerQuote{现在正在受洗的孩子长大后会不会贞洁？他会不会成为盗贼？}恐怕没有人敢回答：\InnerQuote{他会的？}或\InnerQuote{他不会，}尽管对于孩子信天主并转向天主，却毫不迟疑地做出回答。} 此后，你在结尾加上这句话：\Quote{我恳请你对这些问题屈尊给我一个简短答复，不要用习俗的传统权威来封我的口，而是用针对我理性的论证来满足我。}\par

8. 当我反复阅读你的这封信，并在有限的时间允许下思索其内容时，记忆让我想起了我的朋友\Person{内布里迪乌斯}。他是一位极其勤勉而热切钻研困难问题的人，尤其在基督信仰教义领域，但对于对一个重大问题给出简短回答，他却极其反感。如果有人坚持要求这样，他会非常不高兴；而且，倘若不是碍于对方的年龄或身份，他会以严厉的目光和言辞愤然责备这样的提问者；因为他认为，如果一个人不知道在一个重大问题上既能说许多话、也必须说许多话，就不配探究这类问题。但我不会像他被人要求简短答复时惯常的那样对你失去耐心；因为你与我一样，是一位被众多事务缠身的主教，因此既没有闲暇阅读冗长的通信，也没有闲暇写这样的信。他那时还是个年轻人，不满足于对此类问题的简短陈述，而且他本人那时有闲暇，便向一个也有闲暇的人就我们谈话中讨论的许多话题发问；而你，考虑到你自己作为提问者、以及我作为你要答复者的处境，坚持要我为你提出的这个问题给出一个简短的回答。好吧，我会尽力满足你；愿主帮助我达成你的要求。\par

9. 你知道，在日常用语中，当复活节临近时，我们常说：\Quote{明天或后天是主的受难日}，尽管他在多年前受难，且他的受难是一次而永远地承受了。同样，在复活主日，我们说：\Quote{今日主从死者中复活了}，尽管从他的复活以来已过了许多年。但没有人如此愚蠢，在我们使用这些说法时指责我们虚假，因为我们给这些日子这样的名称，是基于它们与那些事件实际发生之日的相似性，那一天被称为那事件的日子，虽然并非事件实际发生的那一日，而是通过每年同一时间的循环而与其对应的一日，而那事件本身也被说成在那一天发生，因为，虽然它确实发生在很久以前，却是在那一天以圣事的方式庆祝的。难道基督不是一次而永远地以他自己的身份被献为祭献吗？然而，他不也在圣事中被献为祭献，不但在复活节的特殊庆典中，而且每日在我们的信友聚会中；以致若有人被问及，回答说他在那圣事中被献为祭献，就是说出了完全真实的话吗？因为如果圣事与它们所是圣事的事物没有一些真正的相似点，他们就根本不是圣事。而且，在大多数情况下，借着这种相似性，它们确实具有它们所相似之现实的名称。因此，以某种方式，基督身体的圣事就是基督的身体，基督宝血的圣事就是基督的宝血，同样，信德的圣事就是信德。而相信无非就是具有信德；因此，当为一个尚不能行使信德的婴儿回答说他信时，这个回答的意思是他因信德的圣事而有信德，同样，也为着他因皈依的圣事而转向天主而作回答，因为这回答本身就属于这圣事的举行。所以，关于这圣洗圣事，宗徒说：\BibleQuote{我们借着圣洗与基督同葬于死。}\BibleRef{罗 6:4}他没有说：\BibleQuote{我们表明了与他同葬。}而是说：\BibleQuote{我们已与他同葬。}因此，他给予那关乎如此伟大事迹的圣事的，不是别的名称，正是描述这事迹本身的言辞。\par

10. 因此，一个婴儿，虽然他尚未是拥有那包含行使信德者同意之意志的信德的信徒，然而通过那信德的圣事，他成为了信徒。因为就像回答说他相信一样，他也被称为信徒，不是因为他凭自己判断的行为同意真理，而是因为他领受了那真理的圣事。然而，当他开始具有成人的分辨能力时，他不会重复这圣事，而是理解它的意义，并意合于心，随同他的意志也同意，而相似于它所包含的真理。在他因年幼尚不能做到此事的期间，这圣事将保护他抵抗敌对势力，并对他大有裨益；倘若他到达理性运用的年龄之前离开此世，他借着基督徒的援助，即教会借着这圣事将他托付于天主的爱，得以从那由一人进入世界的定罪中解救出来。\BibleRef{罗 5:12} 谁若不信这事，以为不可能，那确实是不信的人，即使他可能已领受了信德的圣事；而在功劳上远超过他的是那婴儿，虽尚未拥有由理智帮助的信德，却未以任何理智上的反对阻碍信德，因此有益地领受了信德的圣事。\par

我已回复了你的问题，就我看来，若我是与能力较弱、倾向反驳的人打交道，这种方式也许有所欠缺，但或许足以让平和而明理的人满意了。而且，我并未仅仅以这项习惯已牢固确立的事实来为自己辩护，而是已尽我所能，提出了论点，以支持这充满丰富祝福的作法。\par

\SourceRecord{Translated by J.G. Cunningham.；Nicene and Post-Nicene Fathers, First Series；Vol. 1.；Edited by Philip Schaff.；Buffalo, NY: Christian Literature Publishing Co.,；1887.；<http://www.newadvent.org/fathers/1102098.htm>.}

% structure: p0001 source=h1 target=section reason=page-title

\section{书信九十九（公元408年或409年初）}

\Emph{致极虔敬并受基督肢体公义虔诚赞颂的天主婢女\Person{意塔利卡}，奥古斯丁在主内致候。}\par

1. 在我写这封回信时，我已收到尊驾的三封来信，其中第一封急切请求我写一封信，第二封告知我写的回信已送达您处，而第三封则表达了您极仁慈地为我们的利益而关心属于那位极其显赫杰出的青年\Person{犹利安}的房屋之事，这房屋紧邻我们教会墙壁。刚收到这最后一封信，我立即不失时机地回信，因为尊驾的代理人写信给我说，他可以毫无延迟地将我的信送往\Place{罗马}。他的信使我们极为忧苦，因为他特地让我们了解城里（\Place{罗马}）或城墙周围正在发生的事情，从而就我们之前不愿相信的模糊传闻之可靠性，提供了确实的信息。在先前弟兄们寄给我们的信中，带给我们的消息虽确实是令人烦恼而伤心的，但远不及我们现在所听到的那般灾难深重。我无法形容地感到诧异，我的主教弟兄们在有如此良机通过您的信使寄信时竟未写信给我，而且您本人的信也未向我们传递任何有关降临到你们身上的痛苦患难的消息——这患难，因着基督徒爱德的柔情同感，也是我们的患难，如同是你们的一样。不过我想，你们认为最好不提这些伤心事，因为你们觉得这没有什么好处，或是因为你们不愿用信使我们忧伤。但在我看来，即使将诸如此类的事告知我们，也确有某种益处：首先，因为不宜准备\BibleQuote{与喜乐的一同喜乐}，却拒绝\BibleQuote{与哭泣的一同哭泣}；其次，因为\BibleQuote{患难生忍耐，忍耐生老练，老练生望德；望德不叫人蒙羞，因为天主的爱藉着赐给我们的圣神，已倾注在我们心中了}。\par

2. 因此，我们绝不可拒绝听闻临于我们至亲身上那些痛苦忧伤的事！因为，以我无法说明的方式，当一个肢体受苦时，如果所有其他肢体都一同受苦，那肢体的痛苦就会减轻。\BibleRef{格前 12:26} 这种减轻并非借着实际分担灾祸，而是藉由爱的安慰力量；因为虽然只有一些人实际承受苦难的负担，其他人是透过知晓他们所承受的而分担其痛苦，然而患难是由他们全体共同承担的，因为他们共有着同样的经历、望德和爱，以及同一的天主圣神。此外，主为我们所有人提供安慰，因为祂既已预先警戒我们这些现世的患难，又应许我们在此之后永生之福；而且那渴望在战斗结束时领受冠冕的士兵，不应在战斗持续时丧胆，因为那为得胜者预备无可言喻之赏报的祂，在他们于沙场战斗时，亲自赐予他们力量。\par

3. 不要让现在所写的这些打消你写信给我的信心，尤其因为你为避免我们的恐惧而可能提出的理由，并无不可非议之处。我们回致问候你的幼小儿女，并切愿他们能留给你，且在基督内成长，因为即使在目前的稚弱年龄，他们已经看出此世的贪爱是多么危险有害。愿天主使这些幼小仍然柔韧的植物，在此大树与硬木正被摇撼的时期，得以朝向正确的方向弯曲。至于你所说的那房子，除了对你的极仁慈关切表示感谢之外，我还能说什么呢？因为我们所能给出的房子，他们不想要；而他们想要的房子，我们不能给，因为它并非如他们被虚假告知的那样，是由我的前任留给教会的，而是属于教会的古老产业之一，且它附属于那座古老教堂，正如引起这疑问的房子附属于另一座古老教堂一样。\par

\SourceRecord{Translated by J.G. Cunningham.；Nicene and Post-Nicene Fathers, First Series；Vol. 1.；Edited by Philip Schaff.；Buffalo, NY: Christian Literature Publishing Co.,；1887.；<http://www.newadvent.org/fathers/1102099.htm>.}

% structure: p0001 source=h1 target=section reason=page-title

\section{书信一百（主历409年）}

\Emph{\Person{奥思定}谨向尊贵而备受敬仰的主、卓绝可赞的儿子\Person{多纳图}，在主内致候。}\par

1. 我实愿\Place{非洲教会}不必处于这般艰难的境遇，以致需要任何尘世权势的援助。然而，正如宗徒所说：\BibleQuote{没有权柄不是从天主来的}，\BibleRef{罗 13:1}，这是无可置疑的：当您——您那天主教慈母的真诚儿女——给予她援助时，我们的救助是在于主的名，\BibleQuote{祂创造了天地}。啊，尊贵而备受敬仰的主、卓绝可赞的儿子，谁不察觉，在此浩劫之中，天主赐予我们的安慰实非微小，因为您这样一位如此热心于基督之名的人，被擢升至\OriginalTerm{总督}{proconsul}之尊位，以便权柄与您的善意结合，能遏制教会的仇敌陷于邪恶亵圣的图谋？事实上，对于您执行正义，我们深为担忧的只有一件事：即深恐您因见不虔诚且忘恩负义之徒加于基督徒团体的每桩伤害，都比加于他人者更为严重邪恶，便据此认为当以与罪行之重大相称的严厉加以惩治，而非以适合基督徒忍耐的温和处之。我们因耶稣基督之名恳求您，不要如此行事。因为我们并不求在此世报复；我们所受的苦难，也不应使我们忧愁至极，以致心中毫无余地存记我们从祂所领受的有关此事之诫命——我们正为祂的真理、因祂的名而受苦；我们\BibleQuote{爱我们的仇敌}，并\BibleQuote{为他们祈祷}。\BibleRef{玛 5:44} 我们借法官与法律之威吓以求达成的，并非置他们于死地，乃是救他们脱离错谬，藉此可保全他们不致陷于永罚；我们既不愿见到对他们疏于管教，亦不愿见到他们遭受应得的严刑。因此，请您如此遏制他们的罪过，好使罪人得以存留，俾能悔改己罪。\par

2. 因此，我们恳求您，在审理有关教会案件而宣判之时，无论您查证到教会遭受或意欲遭受的伤害何等邪恶，都请忘记您拥有处以极刑之权，而勿忘我们的恳求。我最亲爱尊崇的儿子，我们请您饶恕那些我们为其祈求天主施予悔改之恩者的性命，此事切勿以为无足轻重而不屑垂顾。即使我们永不该偏离以善胜恶的坚定决心，也请您的智慧斟酌此事：除属于教会者外，无任何人会费心将有关她利益之案件呈于您前。因此，您若认为犯有此等罪行之人必处以死刑，则您将阻止我们竭诚将此类事项诉诸您的法庭；此事一旦泄露，他们将更为肆无忌惮地迅速完成对我们的毁灭，而我们则被迫并受命宁可死于其手，也不愿因在您庭前指控他们而使其致死。我恳求您，万勿轻弃我的此项建议、请求与哀恳。因为我并不认为您不记得，我即便未为主教，且您的品级远高于现时，我亦可有极大胆量向您进言。与此同时，请借您阁下的谕令，让这些\OriginalTerm{多纳图派}{Donatist}异端者立刻知晓，那些针对其谬误而通过、他们以为并夸耀宣布已废除的法律，依然有效，纵使他们得知后仍无法丝毫停止加害我们。然而，您若能留意使帝国为遏制其充满自负与不虔骄傲之教派而颁的法律，得以如此运用，以致他们自己或他人均不觉得是因真理与正义而遭受任何形式的苦难；而是当他们向您请求时，允其于阁下或下级法官面前，通过公开程序，以无可辩驳的确凿事实，说服并教导他们，使那些由您下令逮捕之人，甘心移其顽固意志于更善之途，并能有裨益地向其同党宣读这些事理。因为若仅以强制而非劝导训诲，使人弃绝大恶而把握大善，所受的辛劳便较为累赘，而非真有实益。\par

\SourceRecord{Translated by J.G. Cunningham.；Nicene and Post-Nicene Fathers, First Series；Vol. 1.；Edited by Philip Schaff.；Buffalo, NY: Christian Literature Publishing Co.,；1887.；<http://www.newadvent.org/fathers/1102100.htm>.}

% structure: p0001 source=h1 target=section reason=page-title

\section{第101封信（公元409年）}

\Emph{致\Person{Memor}，我最可敬的主、以全部敬意最亲爱的、衷心渴念的弟兄与同工，\Person{奥思定}在主内致意。}\par

1. 我不该仅给您的圣爱德写信，却不同时寄上那些您以神圣之爱不可抗拒的理由向我索要的书籍，以致至少可通过这一服从的行动，答复您那些既给了我极高荣誉，也给我压上沉重担子的来信。然而，当我因这重担而俯身时，也因您的爱而被扶起。因为爱着我、扶起我、使我立直的，并非一个寻常之人，而是一位主的司铎，我知道他在主前如此蒙悦纳，以至当你把你那颗怀着我的善心高举向主时，你便将我连同你自己一起举向了祂。因此，我本应在此刻寄上那些我曾允诺修订的书籍。我没有寄出的原因是我尚未修订，这并非出于不愿，而是由于不能，因我被许多极为紧迫的事务缠身。然而，如果让带信人——我那圣善的同工与弟兄\Person{波西迪乌斯}，您会发现他与我自己非常相似——既未能结识如此深爱我的您，又未曾带着我的信便来见您，那就显示出不可原谅的忘恩负义和心硬。因为他是一位由我的劳苦喂养长大的人，不是以那些被各种私欲的奴隶们称为自由的学问，而是以上主之粮，尽我菲薄所能供应给他的限度。\par

2. 对于虽不义又不虔敬，却自以为在自由学科方面受过良好教育的人，我们除了从那些真正配得上自由之名的著作中读到的话之外，还能对他们说什么呢？——\BibleQuote{如果圣子使你们自由，你们就确实是自由的。}\BibleRef{若 8:36} 因为正是通过祂，即使是在那些没有被召叫获得这种真自由的人所称的自由学问中，人们也可认识到其中任何配得上这个名称的东西。因为它们没有与自由相合的东西，除了其中与真理相合的东西；为此，圣子自己说：\BibleQuote{真理必会使你们自由。}\BibleRef{若 8:38} 因此，我们所享有的自由，与那些愚蠢诗人的诗句所充斥的无数不虔敬的寓言，与那些雄辩家浮夸而精雕细琢的谎言，最后，与那些哲人们纷繁芜杂的诡辩都毫无共同之处——这些哲人或根本不认识天主，或虽然认识了天主，却没有以祂为天主而予以光荣或感谢，\BibleQuote{他们虽然认识了天主，却没有以他为天主而予以光荣或感谢，反而所思所想的成了荒谬的，他们冥顽不灵的心陷入了黑暗；他们自负为智者，反而成为愚蠢，将不可朽坏的天主的光荣，改归为可朽坏的人、飞禽、走兽和爬虫形状的偶像，}\BibleRef{罗 1:21} 又或者，虽然没有完全或根本沉迷于偶像崇拜，却仍然敬拜事奉受造物，胜过了造物主。所以，我们绝不承认，那些不配的、不幸福的人们的虚妄幻想和错觉、空洞的琐屑和自负的谬误，配得上自由这一美称；他们不认识在我主基督耶稣内的天主的恩宠，只有藉着这恩宠，我们才\BibleQuote{得救脱离这死亡的肉身，}\BibleRef{罗 7:24} 而且他们甚至未能察觉自己所知事物中的那一点真理的分量。至于他们的历史著作，其作者声称最关心的乃是准确记述事件，对此，我或许可以承认，其中包含一些值得自由之人知道的东西，因为记述是真实的，无论所述对象是人类经历中的善或恶。但与此同时，我实在看不出，那些未得圣神之助而认识，且不得不因人性软弱的限制而搜集飘泊传闻的人，如何能在如此众多的事上避免被误导；不过，如果他们没有欺骗的意图，并且除了他们也因人性软弱而误入迷途之外，并没有以其他方式误导他人，那么，在这类著作中仍有一种接近自由的东西。\par

3. 然而，由于在所有运动中呈现的数字能力，在声音中最易于研究，而且这项研究提供了一条逐步向上攀登，通往更高真理奥秘的路径，正如经上所说，\BibleQuote{智慧在其中愉快地显示自己，并在眷顾的每一步与爱慕她的人相遇，}\BibleRef{智 6:17} 因此，当我开始有闲暇从事研究，心智尚未被更重大更紧要的挂虑所占据时，就立意着手撰写您索要的那些书籍，以此锻炼自己。那时我仅论节奏便写了六卷书，并且打算，还可以加上，再写六卷论音乐，因我那时预期会有闲暇。然而从教会的挂虑重担加于我身的那一刻起，所有这些消遣便全然脱手而去，以至于如今我既不能不尊重您的愿望和命令——因这已不止是请求——竟连我过去所写的都难以找到了。不过，即使我有能力将那论著寄给您，它所带来的遗憾，不是对我而言惋惜我听从了您的命令，而是为您惋惜，您竟如此急切地坚持要我把它寄出。因为这部论著的前五卷几乎无法理解，除非身边有人，在诵读时不仅能分辨讨论对话各方所讲的话，还能用声音标出音节应占的时值，以便表达出它们特有的节拍，并传入耳中，尤其是因为在某些地方会出现有确定长度的停顿，这些停顿自然只有诵读的人用其间静默的间隙向听者提示，否则便会被忽略。\par

然而，第六卷，我发现其已经修订完毕，其中包含了其余五卷的成果；我已毫不迟延地将其呈给您的爱德；它或许并非完全不适用于您这可敬的年纪。至于其余五卷，在我看来，几乎不值得我们的儿子、如今也是我们的同僚的\Person{犹利安}知晓和阅读，因为他身为执事，正与我们一起投身于同样的战斗。关于他，我不敢说我爱他胜过爱您，因为那并非事实；但我可以说，我对他更加思念。当我对两人同样爱慕时，却对其中一人更为渴念，这或许显得奇怪；但造成这一区别的原因是，我见到他的希望更大；因为我想，若他受您派遣或命令来到我们这里，他既是在做与他的年纪相称的事，尤其是他尚未被更重大的责任所阻碍，也会更快地将您亲自带到我面前。\par

我未曾在这篇论文中说明\WorkTitle{圣咏集}的诗句是以何种韵律写成的，因为我对此并不知晓。因为任何人在从希伯来文（我对这语言一无所知）翻译这些诗篇时，都不可能同时保留其韵律，否则他就会因格律的要求而被迫偏离准确的翻译，以至于不符合句子的原意。尽管如此，根据那些通晓该语言之人的见证，我相信它们的确是以若干不同的韵律写成的；因为那位圣人热爱圣乐，他比任何人都更在我心中燃起了对圣乐研究的热情。\par

\BibleQuote{愿至高者的翼荫永远作你们众人的居所！}\BibleRef{咏 91:4}你们，父亲与母亲，同心合意同住一屋，与你们的儿子们结为同一弟兄团体，全是唯一圣父的子女。请纪念我们。\par

\SourceRecord{Translated by J.G. Cunningham.；Nicene and Post-Nicene Fathers, First Series；Vol. 1.；Edited by Philip Schaff.；Buffalo, NY: Christian Literature Publishing Co.,；1887.；<http://www.newadvent.org/fathers/1102101.htm>.}

% structure: p0001 source=h1 target=section reason=page-title

\section{书信102（公元409年）}

\Emph{致我至诚的兄弟、同为长老的\Person{德奥格拉提亚斯}，\Person{奥古斯丁}以主内问候。}\par

1. 你选择将那些原本向你求问解答的问题转交给我，我认为你并非出于怠惰，而是因为你对我的爱超过了我的应得，所以宁愿通过我来听闻那些你已经十分清楚的事。然而，我更希望由你亲自给出答复，因为那位提出这些问题的朋友似乎不愿接受我的劝告——如果我可以从他没有回复我的信这一点来判断的话，至于他何以如此，他自己最清楚。不过，我这样猜疑是合理的，既无恶意，亦不荒谬；因为你也深知我多么爱他，也多么因他尚未成为基督徒而痛心；而若我看见他不愿回复我的信，便认为他不愿我为他写任何东西，这也是合乎情理的。因此，我恳求你应允我的请求，既然我已顺从了你，而且尽管有最繁重的职务，仍不敢辜负你所珍视的心意，而拒绝你的要求。我所求的是：请你不要拒绝亲自回答他所有的问题，因为正如你所告诉我的，他向你恳求过此事；而这项任务，在你收到这封信之前，你便已能够胜任；因为当你读了这封信后，你会看见，我所说的几乎没有什么是你所不知的，或即使我保持沉默，你也能知道的。因此，我恳请你为我保存这份工作，供你自己以及那些你视为适宜满足其求知渴望的人使用。至于我向你索求的、由你亲自撰写的论文，请交给最适合接受的人，不但给他，也给所有对你关于这些事所能作的陈述极为赞赏的人，而我本人也在此列。愿你永活在基督内，并记念我。\par

2. \Emph{问题一}。关于复活。此问题困扰着一些人，他们问道：这两种复活，哪一种与我们被应许的相符？是基督的复活，还是\Person{拉匝禄}的复活？他们说：\Quote{如果是前者，这怎能与那些经由寻常生育而出生之人的复活相符呢？因为他并非如此出生。如果反过来，说\Person{拉匝禄}的复活与我们相符，这里也似乎有出入，因为\Person{拉匝禄}的复活是在身体尚未分解时完成的，就是那具他以\Person{拉匝禄}之名而为人所知的身体；而我们的身体则要在许多世纪之后，从已与其他物体混同难辨的质料中被解救出来。再者，倘若我们复活后的境界是有福的，在此境界中身体免除一切创伤，也不再有饥饿的痛苦，那么基督复活后取食并显示他的创伤，这记载又意味着什么呢？因为如果他这样做是为了使怀疑者信服，而创伤并非真实的，他就对他们行了欺骗；然而，如果他显示的是真实的创伤，那么由此可知，身体所受的创伤将留存在复活后的境界中。}\par

3. 对于这一点，我的回答是：被应许给我们的复活，与基督的复活相符，而非与\Person{拉匝禄}的复活相符，因为\Person{拉匝禄}复活后仍要再死一次，而关于基督的记载则是：\BibleQuote{基督既从死者中复活，就不再死；死亡不再统治他了。} \BibleRef{罗 6:9} 同样的应许也赐给那些将在世界终穷时复活，并与基督一同永远为王的人。至于基督的诞世方式与他人不同，这并不影响他复活的性质，正如这也不影响他死亡的性质，使之与我们的死亡有所不同。他的死亡，并不因为他不曾由世上的父亲所生，就变得不那么真实；正如那直接用尘土造成的第一个人的身体起源方式，与我们由父母所生的身体起源方式不同，但并未因此使他的死亡与我们的死亡属于不同种类。因此，正如诞世方式的不同不会造成死亡性质的不同，同样，这也不会造成复活性质的不同。\par

4. 但为了不让对此怀疑的人，以同样的怀疑态度，拒绝相信关于第一个人受造的记载是真实的，就让他们至少就他们所能相信的，去探究或观察一下：有多少类别的动物是从土地生出，而非从父母获得生命，却通过寻常繁殖，产生与其自身相似的后代；在这些动物中，尽管起源方式不同，但从土而生的父母与其所生的后代，本性却是相同的；因为它们虽然诞生的方式不同，却同样地生活，同样地死亡。因此，说身体在起源上有所不同，但在复活上却相同，这并无荒谬之处。然而，这类人不能分辨事物之间的差异在何种情形下有影响，在何种情形下无影响，因此一旦发现事物在最初形成上有所不同，便坚持认为在此后一切事上，那不同也必继续存在。这样的人，也可以同样有理地推断：用油脂制成的油，应不能像橄榄油那样浮在水面，因为两种油的来源如此不同：一种来自树的果实，另一种来自动物的肉。\par

再次，关于所谓的基督身体与我们身体的复活之间的差异：祂的身体第三天复活，并未因腐朽而解体；而我们的身体则要在很长时间以后，并且要从那已分解得无法辨认的聚合体中重新形成。这两件事对人来说都不可能做到，但对天主的能力来说，两者都轻而易举。因为正如目光注视近处的物体并不比远处的更快，而是以同等速度射向任何距离，同样，当死人复活在\BibleQuote{眨眼之间}完成\BibleRef{格前 15:52}时，对天主的全能与祂意志的不可言喻之表达来说，使许久时间以来已分解的身体复活，如同使刚死去不久的身体复活，是同样容易的。这些事对某些人而言难以置信，因为它们超越其经验；虽说整个自然界充满如此多的奇事，以至于它们对我们并不显为奇妙，因此被认为不值得仔细研究或探究，并非因为我们的官能轻易能理解它们，而是因为我们已如此习惯于看见它们。至于我自己，以及所有与我一同努力通过受造之物\BibleRef{罗 1:20}理解天主不可见之事的人，我可以说，我们对如此微小的一粒种子里仿佛蕴含着我们在参天大树中所赞美的一切，所感到的惊奇，丝毫不亚于（或许更多）对如此广阔的这大地的怀抱将在未来复活时，归还它所接纳的、现已分解的人体全部完整元素这一事实所感到的惊奇。\par

再者，基督复活后进食这一事实，与所应许的复活状态中无需食物的教义，二者之间有何矛盾？因为我们读到，天使也曾以同样的方式吃了同样的食物，不是出于虚妄且骗人的模拟，而是无可置疑的真实；然而并非出于必需的压力，而是自由运用其能力。因为水以一种方式被干渴的土地吸收，以另一种方式被炽热的阳光吸收；在前者我们看到匮乏的后果，在后者我们看到能力的表现。那么，那未来复活状态的身体，若不能进食，其幸福便不完美；反之，若依赖食物，同样也不完美。我本可在此更充分地讨论物体性质可能的变化，以及高层物体对低层物体所具有的统辖权；但我已决定简短答复，并且我为那些心思如此聪慧的人写下这些，对他们而言真理的简单提示便已足够。\par

提出这些问题的人务必要知道，基督在复活后，向疑惑的门徒展示的是祂伤口的疤痕，而非伤口本身；也是为他们，祂乐意多次进食，免得他们以为祂的身体并非真实，而祂是幽灵，以幻影而非实体的形式显现。这些疤痕若先前没有伤口存在，确实就成了单纯的假象；然而，若祂有意，连这些疤痕也本不会留下。但祂乐意以明确的目的保留它们，也就是说，为要向他们——祂正在不虚伪的信德中培育的人——表明，并非一个身体替代了另一个，而是他们曾见被钉在十字架上的那个身体复活了。那么，有什么理由说：\Quote{如果祂这样做是为了说服怀疑者，祂就在行骗}呢？假设一位勇士，在为祖国作战时身负多处创伤，他向一位技术非凡、能使这些伤口愈合得不留任何可见疤痕的医师说，他宁愿以这样的方式得到治疗：让伤口痕迹留在身上，作为他所获荣誉的标记；在此情形下，你会说这位医师行骗吗？因为他虽然能用技术使疤痕完全消失，却出于明确的原因，反而用同样的技术让它们保持原样？正如我已经说过的，证明这些疤痕是骗局的唯一根据，便是根本没有伤口曾在这些疤痕所在之处得到医治。\par

8. \Emph{问题二} 至于基督宗教的时代，他们还提出了一些其他的论点，可称之选取了\Person{波菲利}反对基督徒的若干较为有力的论证：\Quote{如果基督，}他们说，\Quote{宣称自己是救恩的道路、恩宠与真理，并断言唯有在他内、且仅对信靠他的灵魂，才有回归天主的道路，\BibleRef{若 14:6}那么，在基督来临之前许多世纪生活的人，又如何了呢？暂且不提}他补充道，\Quote{\Place{拉丁姆}王国建立之前的时间，我们就以那个政权的开端作为人类历史之始吧。早在\Place{阿尔巴}建立之前，\Place{拉丁姆}本地就有对神明的崇拜；在\Place{阿尔巴}，神庙中的宗教仪式和敬拜形式也得以维持。\Place{罗马}本身，在不少于同样长的一段时期内，甚至在连续许多世纪中，对基督教的教义一无所知。那么，如此不可胜数的灵魂，既毫无罪过，且唯有相信他才能获得救恩的那一位尚未降临人间施恩予人，他们的结局又是如何？此外，整个世界在神庙中对诸神的崇拜，其热忱并不亚于\Place{罗马}。那么为甚么}他问道，\Quote{那位被称为救主者，竟让自己在世间隐藏如此之多的世纪呢？并且，不要推说}他又补充道，\Quote{已由旧犹太法律为人类作了预备。犹太法律是在很久之后才出现，并在\Place{叙利亚}狭窄的地域内昌盛，此后才逐渐扩展到\Place{意大利}沿海；但这不早于\Person{卡利古拉}统治的末期，或至早在他登基之时。那么，生活在\Person{凯撒}时代之前的\Place{罗马}与\Place{拉丁姆}之人，既缺乏基督的恩宠——因为他那时尚未降临——他们的灵魂又当如何呢？}\par

9. 对于这些说法，我们的回答是：首先要求那些提出这类问题的人告诉我们，我们知道在可以查考的各时期中被引入其神明崇拜的那些宗教仪式，是否曾对人有益？如果他们说，这些仪式对人的救恩毫无益处，那么他们便与我们一同否定了这些仪式，并承认它们是无用的。我们固然证明了它们是有害的；但至少他们承认这些仪式无用，这已是一个重要的让步。反之，如果他们为这些仪式辩护，并坚持它们是明智而有益的规条，那末，请问，在它们被设立之前去世的人，又如何了呢？因为他们被剥夺了这些仪式所拥有的救恩与效益。然而，若说他们能藉另一途径同样有效地涤除罪过，那末，为甚么这同一途径不为其后代继续有效呢？在崇拜中设立新规又有何用？\par

10. 倘若对此他们回答说，那些神明本身确实始终存在，并且在各处同样有能力赐予其崇拜者自由，但神明乐于依照时间、地点和方式的状况，按属世与属地事物中的种种变化，以他们知道最为适合某些时代和国家的方式，来规定人敬拜他们；那末，为甚么他们竟以这问题攻击基督宗教呢？如果这问题被用来质问他们自己的神明，他们要么无法回答，要么，即使能回答，也必然以同样方式为我们的宗教、并不亚于为他们的宗教作答。因为他们还能说甚么呢？无非是说：适应不同时代与地点的\OriginalTerm{圣事}{sacraments}，其间的差异无关紧要，只要在这一切中所被崇拜的对象是圣善的；正如属于不同语言、适应不同听众的语音之间的差异无关紧要，只要所言说的是真理；虽然在这一点上有区别：人可以彼此约定，安排传达思想的语言之音声，但那些辨识正路的人，在关乎合宜于神圣存在的宗教仪式上，只受到天主圣意的指引。这神圣的旨意，为了救助世人，从未吝于给予有死之人的正义与虔诚；而且，在同一宗教相联系的各民族中，崇拜礼仪纵有多样变化，最须关注的是：一方面，人的软弱因此如何在多大程度上得到鼓励而精进，或得以承受；另一方面，天主的权威又未受侵犯。\par

11. 因此，既然我们确认基督是天主的\OriginalTerm{圣言}{Word}，万物是藉他而造成的，并且他是圣子，因为他是圣言，不是说出即逝的言语，而是永恒不变地与不变的圣父同在，他自身也不变；在他掌管之下，整个宇宙，无论精神的或物质的，都以最适合不同时代与地点的方式被安排；并且他拥有完全的智慧与知识，知道在掌控与安排宇宙时，何事当作，以及何时何地作这一切——那么，至为确定的是：无论是在他使希伯来民族存在之前（他乐于通过适合当时代的圣事，以该民族预示他将在降临时显示自己），或是在犹太政体的时期，还是此后，当他以由童贞女所取的人形，以有死者的样式向有死的人显示自己之时，又从那时起直至我们今天，在这进程中他正在实现他昔日藉先知所预言的一切，并且从此刻起直到世界末日，那时他将把圣者与恶人分开，给予各人应得的报偿——在所有这些相继的时代中，他同是天主之子，与圣父同为永恒，是不可改变的\OriginalTerm{智慧}{Wisdom}；借着他，整个受造界得以存在，并且一切有理性的灵魂，都因分享他而成为有福的。\par

12. 因此，从人类的一开始，凡信靠祂的，并以任何方式认识祂，按照祂的诫命度虔敬而正义的生活的人，无疑都借着祂得救，无论他生活在何时何地。因为就如我们相信祂既与\OriginalTerm{圣父}{Father}同在，又已降生成人，同样，古代的人也相信祂既与圣父同在，且预定要降生成人。\OriginalTerm{信德}{faith}的本质并未改变，\OriginalTerm{救恩}{salvation}在我们的时代也没有变得不同，因为时代的差异，那从前预言为将来的事，现在被宣告为已过去的事。而且，我们不必认为，当同一件事藉着不同的神圣言语和敬拜仪式，在一个情形下宣布为已完成，在另一个情形下宣布为未来时，所指的事物和救恩种类就有不同。至于那关系到所有信者和虔敬者共同的同一救恩，其实现的方式和时间，让我们将智慧归于天主，而我们则服从祂的旨意。因此，真宗教虽然在古代以与现在不同的名号和象征性仪式来宣讲和实行，并且在古代比现在更为隐晦地启示，只为少数人所知，如今则光芒更亮，传播更广，但在两个时期却是同一且不变的。\par

13. 此外，我们不因为\Person{努马·庞皮利乌斯}为罗马人敬拜神明所制定的制度，与在那之前\Place{罗马}或\Place{意大利}其他地区所实行的制度之间的差异，而反对他们的宗教；也不因为\Person{毕达哥拉斯}时代那种哲学体系开始被普遍接受，而直到那时这种体系尚不存在，或可能只隐藏在极少数人中，他们的观点相同，但宗教实践与崇拜却不同，就加以反对：我们与他们较量的要点是：这些神是真神，还是值得崇拜？那种哲学是否适合促进人类灵魂的\OriginalTerm{救恩}{salvation}？这就是我们所坚持讨论的；在讨论中，我们要将他们的诡辩连根拔起。因此，他们不要再用那些对任一宗派、对一切名号的宗教都同样有力的诘难来反对我们。因为，正如他们所承认的，世界的时代并非在偶然的支配下运行，而是受\OriginalTerm{天主眷顾}{Providence}的控制，每个时代中合宜而有利的事，超越了人类的理解范围，是由那眷顾宇宙的同一智慧所决定的。\par

14. 因为，如果他们断言\Person{毕达哥拉斯}的学说没有在任何时代和任何地方普遍盛行的原因是，\Person{毕达哥拉斯}只是一个人，没有能力确保这一点，那么，他们能否也断言，在他的哲学盛行的时代和地区，所有有机会听到他的人都愿意相信并跟随他？因此，更确定的是，如果\Person{毕达哥拉斯}拥有在他所喜欢的时间和地点发表他学说的能力，并且同时拥有对事件的完美预知，他就会只在他预知人们会相信他教导的那些地方和时候出现。因此，既然他们不以\Person{基督}的教义未被普遍接受为由反对\Person{基督}——因为他们觉得，这种反对无论是针对哲学家的教导，还是针对他们自己神明的威严，都是徒劳的——那么，我要问，如果我们暂且不论天主智慧与知识的深邃，其中可能隐藏着更为深奥的某种神圣目的，也不预先判断可能存在的其他原因（这些原因值得智者耐心研究），而是为了在此讨论中简短起见，仅限于提出这一点：即\Person{基督}乐意按照祂对于人们会相信祂的时间与地点的预知，来决定祂将显现的时间和向哪些人宣讲祂的教义，那么他们能给出什么回答呢？因为祂预先知道，关于那些祂的\OriginalTerm{福音}{gospel}未曾宣讲的时代和地点，在那里，如果福音被宣讲，它将会遭到所有人的对待，无一例外，正如祂在地上亲身临在时，福音所遭遇的那样——并非从所有人，而是从许多人，他们即使看见祂让死人复活，也不信祂；又如我们在今日所见的，许多人虽然关于祂的\OriginalTerm{先知}{prophets}预言如此明显地应验，却仍拒绝相信，他们被人心乖僻的诡诈引入歧途，面对如此清晰明显、如此光辉且辉煌传扬的神圣权威，宁可抗拒而不顺从。只要人的心智在能力和力量上有限，他就应当顺服神圣的真理。那么，我们为什么要奇怪，如果\Person{基督}知道在古代世界充满了如此多的不信者，以致祂公义地拒绝向那些祂预知无论对祂的话语还是奇迹都不会相信的人显现自己或传道呢？因为，令人惊讶的是，从祂降临的时代直到现在，甚至直到今日，仍有许多人如此，那么当时所有人也可能都是这样，这并非不可信。\par

15. 然而，从人类之初，祂从未停止藉着祂的先知说话，有时比较隐晦，有时比较明白，这是依天主的上智看为最适合时机的；而且，从\Person{亚当}到\Person{梅瑟}，在\Place{以色列}民族本身中——它因一个特殊且奥妙的任命而成了一个先知性的民族——以及在祂降生成人之前的其他民族中，从不缺乏信祂的人。因为在神圣的希伯来经卷中提到一些人，甚至从\Person{亚巴郎}时代起，他们既不属于他的后裔，也不属于\Place{以色列}民族，也不是归附那民族的人，但他们却有份于这\OriginalTerm{神圣奥迹}{holy mystery}；那么，我们为何不能相信，在其他民族中，这里或那里，也有更多这样的人，尽管我们在这些权威记录中读不到他们的名字呢？因此，由这宗教所预备的救恩——唯有这宗教是真正的，也只有它真正应许了真正的救恩——从不曾缺乏任何配得的人，而那缺乏的人原是不配得的。而且，从人类大家庭的开始，直到时间的终结，这救恩被宣讲，为一些人是得益处，为一些人是受审判。所以，那些从未听到宣讲的人，原是\OriginalTerm{预知}{foreknown}不会信的人；那些虽确知不会信，却仍被宣报了救恩的人，是作为不信者的模本而被设立；那些作为将会相信的人，而被宣报真理的人，正在为天国和圣天使的团体作准备。\par

\Emph{问题三}现在让我们依序来看下一个问题。他说：\Quote{他们挑剔}，\Quote{神圣的仪式、祭牲、焚香，以及我们庙宇中其他一切崇拜部分；然而，同样的崇拜形式，在他们自己中间，或从他们所崇拜的神那里，自古就有，因为据他们描述，那神也需要初熟之果。}\par

17. 这个问题显然是依据我们圣经中的一段记载：\Person{加音}把田间的出产作礼物献给天主，而\Person{亚伯尔}却把羊群中\OriginalTerm{首生牲畜}{firstlings}作礼物献上。\BibleRef{创 4:3}因此，我们的答复是：从这段经文可得出更合适的推论，即祭献的礼仪何其古老，无谬误而神圣的经卷宣告，这礼仪只应归于独一无二的真天主；并非因为天主需要我们的奉献——因为在同一圣经中，最清楚地写着：\BibleQuote{我对主说：祢是我的主，因为祢不需要我的善，}——而是因为，在接受、拒绝或收纳这些奉献时，祂考虑的是人的益处，且只考虑人的益处。因为，在崇拜天主时，我们是使自己获益，而非使祂获益。所以，当祂赐下神启的教导，并指示人当如何崇拜祂时，祂这样做，不仅不是出于祂自己有任何需要，而是为了我们最高的益处。因为所有这些祭献都是有意义的，它们是某些事物的标记，我们应当藉此被激发去寻求、认识或记起它们所象征的事物。要圆满地讨论这个主题，需要比我们这次决定作出的简短答复更长的论述，尤其因为我们在其他论著中已经详细谈过。在我们之前讲解神圣启示的人们，也已大量论述过\OriginalTerm{旧约}{Old Testament}的祭献标记，乃是将来事物的影子和预像。\par

18. 然而，尽管我们力求简短，但有一件事我们决不可不提：那就是，假神——亦即\OriginalTerm{邪魔}{demons}，也就是\OriginalTerm{撒谎的天使}{lying angels}——若非知道这一切原本都当归于独一真天主，就决不会要求他们欺骗的敬拜者为之提供庙宇、司祭、祭献以及与此相关的一切事物。因此，当这些事物按照天主的启示和教导而献给天主时，便是真正的宗教；但当它们应邪魔傲慢的骄傲而归于邪魔时，便是邪恶的迷信。因此，认识基督徒的旧约与新约圣经的人，并不因异教徒建造庙宇、任命司祭、奉献祭献而责备其亵渎的礼仪，而是因他们为偶像和邪魔做这一切而责备。至于偶像，的确，谁怀疑它们是全然没有知觉的呢？然而，当它们被安置在这些庙宇中，高坐在尊位上，好让恳求者和敬拜者祈求并奉献祭献来伺候它们时，这些偶像，虽全无知觉和生命，却仅凭那有生命之物的肢体和感官的形象，就如此影响软弱的心灵，以致它们仿佛在生活和呼吸，尤其在群众出于深切的崇敬而白白付出如此昂贵的服事时，这种影响更加强烈。\par

19. 针对这些病态而有害的心灵倾向，圣经提供了一剂良药，它以有益劝诫的感染力重复一个熟知的事实，说：\BibleQuote{有眼不能看，有耳不能听，}等等。因为这些话语如此明白，且为众人承认是真实的，它们更能有效地使那些人产生有益的羞耻感：这些人带着宗教敬畏，在这些偶像前呈现敬拜，看着它们如同活物，在崇敬和朝拜它们时，向它们祈求，祭献，并在它们面前履行誓愿，仿佛它们就在眼前，他们完全被征服了，以至于不敢想象这些偶像毫无知觉。此外，为了防止这些崇拜者认为我们的圣经只意在说明，这种人心倾向自然源于偶像崇拜，圣经又用最明确的话语写道：\BibleQuote{万民的神祇都是魔鬼。}因此，宗徒的教导也不仅如我们在\Person{若望}书中读到的：\BibleQuote{孩子们，你们要远避偶像！}\BibleRef{若一 5:21} 还借\Person{保禄}的口说：\BibleQuote{那么我说什么呢？是说祭偶像的肉算得什么吗？或是说偶像算得什么吗？不是的。外邦人所献的祭，是献给魔鬼，而不是献给天主；我不愿意你们与魔鬼有份。}\BibleRef{格前 10:19} 由此可以清楚明白，真正宗教在异教迷信中所谴责的，并非只是祭献的行为（因为古代的圣人们也曾向真天主祭献），而是向虚假的神祇和不虔敬的魔鬼祭献。因为真理劝导人寻求与圣天使的共融，同样，不虔敬将人转离，去与邪恶的天使共融，为这些邪恶天使的同伴预备了永火，正如为圣天使的同伴预备了永恒的国度。\par

20. 异教徒为他们亵渎的仪式和偶像找到一个借口，即他们能机巧地解释每个偶像所象征的意义，但这个借口毫无用处。因为所有这些解释都指向受造物，而非造物主；惟有造物主才应受那在希腊语中以\OriginalTerm{钦崇}{\GreekText{λατρεία}}一词所指的敬礼。我们也并不说大地、海洋、天空、太阳、月亮、星辰，以及任何可能超出我们感知的天体力量是魔鬼；但因为一切受造物分为物质的和非物质的——后者我们也称为精神的——显然，我们基于虔敬和宗教所行的一切，乃出自我们灵魂中称为意志的官能，这官能属于精神的受造界，因此应优先于一切物质的受造物。由此推论，不可向任何物质的受造物祭献。剩下的就是精神性的受造部分，它要么虔敬，要么不虔敬：虔敬的包括正义的人与天使，他们恰当地事奉天主；不虔敬的包括邪恶的人和天使，我们也称他们为魔鬼。现在，即使向正义的精神受造物也不可祭献，这从以下考虑中显而易见：任何受造物越是虔敬和顺服天主，就越不敢觊觎那惟独应归于天主的尊荣。那么，向魔鬼——即向邪恶的精神受造物祭献，岂不是更糟糕？它们居住在相比之下较黑暗的、离地最近的那个天中，如同在空气中被指定的监狱里，注定要受永罚。因此，即使人们说他们是向更高层的天体力量祭献——这些力量并非魔鬼——并且以为我们与他们之间的区别仅仅在于名称，因为他们称这些为神，我们称这些为天使，实际上，这些被交付于各种欺骗戏弄的人们所真正面对的，只是那些以人类的谬误为乐，并在某种意义上以此为食粮的魔鬼。因为圣天使不赞许任何祭献，除非是按真正智慧和真正宗教的教导，向那唯一的真天主，即他们在神圣共融中事奉的天主所奉献的祭献。所以，正如不虔敬的僭越，无论是在人还是在天使中，命令或贪求将那当归于天主的尊荣归于自己；反之，虔敬的谦逊，无论是在人还是在圣天使中，都会拒绝这些尊荣，并明示这些尊荣唯独应归于谁。我们圣经中有极显著的例子显明这一点。\par

21. 在天主圣言所规定的祭献中，按照所实行的时代不同，有不同制度的多样性。有些祭献是在那新盟约的实际彰显之前奉献的，这新盟约的益处由那独一司铎的一次真实奉献，即由基督所流的宝血所提供；而另一种祭献则适应此彰显，在此时代由我们这些按已彰显者之名而被称为基督徒的人所奉献，这祭献不仅在福音中，在先知书中也已摆在我们面前。因为，受敬拜的天主未改变，宗教本身也未改变，惟有祭献和圣事的改变，若没有在先前的安排中预言过，如今宣告出来似乎便没有根据。正如同一人在早晨向天主献上一种祭品，傍晚献上另一种，按每日的时刻而定，他既不由此改变他的天主，也不改变他的宗教，正如他用一种方式在早晨请安，用另一种方式在傍晚请安，也不改变请安的性质：同样，在整个世代的完整循环中，当知道古代的圣人曾奉献一种祭献，而现时的圣人则呈现另一种时，这只表明这些神圣奥秘并非按照人的僭越，而是按照天主的权威，以最适合时代的方式进行庆祝。这里，在天主或宗教上都没有改变。\par

\Emph{问题四。}接下来，让我们思考他关于罪与罚之间比例所提出的观点，当时他歪曲福音说：\BibleQuote{基督威胁那些不信祂的人将受永罚；} \BibleRef{若 3:18} 然而祂在另一处又说，\BibleQuote{你们用什么量器量给人，也用什么量器量给你们。} \BibleRef{玛 7:2} \Quote{这里，}他评论道，\Quote{是相当荒谬和矛盾的；因为如果祂要按照量器赏罚，而所有的量器都受时间的限制，那么这些永罚的威胁是什么意思呢？}\par

很难相信，这个问题竟由一个自称为某种哲学家的人以反驳的形式提出；因为他说，\Quote{所有的量器都受时间限制}，好像人们除了时间的量器之外不习惯其他量器，例如小时、日、年，或者就像我们说短音节的时长是长音节的一半那样。因为我认为蒲式耳、小桶、瓮和\OriginalTerm{双耳瓶}{amphoræ}并不是时间的量器。那么，所有的量器如何都受时间限制呢？教外人自己不是断言太阳是永恒的吗？然而他们却敢于根据几何测量来计算并断定太阳与地球之间的比例。无论这种计算是在他们的能力范围之内还是之外，尽管如此，太阳确实具有确定尺寸的圆盘。因为如果他们确实确定太阳有多大，他们就知道它的尺寸；如果他们未能探究成功，他们就不知道这些尺寸；但是人们无法发现这些尺寸，并不能证明它们不存在。因此，某物有可能是永恒的，同时却具有确定的比例量度。我在这里是根据他们自己对太阳永恒持续的观点来说的，为的是让他们被自己的原则所说服，并不得不承认某物可以是永恒的，同时又是可量度的。因此，让他们不要因基督也说过\BibleQuote{你们用什么量器量给人，也用什么量器量给你们}，就认为祂关于永罚的威胁不可相信。\par

因为即使祂说过：\Quote{你们所量过的，也必量给你们}，即便如此，也不必将这话理解为指在各方面都相同的事物。我们可以正确地说，你们所栽种的，你们要收割，尽管人栽种的不是果实而是树木，收割的不是树木而是果实。然而，我们这样说的时候，是指树的种类而言；因为一个人栽种一棵无花果树，不会期望从它收获坚果。同样地，也可以说，你们所行的，你们要受报；这不是说，例如有人若犯奸淫，他就要受同样的报应，而是指他在那罪中对法律所行的事，法律也要对他施行\Emph{也就是说，}因为他从自己的生活中除去了禁止这类事的法律，法律就要报应他，将他从那法律所掌管的人类生活中除去。再者，如果祂说过：\Quote{你们量了多少，也要量给你们多少}，即使从这句话，也未必推出我们必须将惩罚理解为在各方面都与所罚的罪相等。例如，大麦和小麦在质量上不相等，但也可以说\Quote{你们量了多少，也量给你们多少}，意思是给这么多小麦，给这么多大麦。或者，如果所涉及的是痛苦，也可以说\Quote{你们加给别人的痛苦有多重，你们要受的痛苦也有多重}；这可能意味着痛苦的严重程度相等，但时间更长，因此由于持续而更重。假设我说两盏灯：\Quote{这盏灯的火焰和那盏灯的火焰一样热}，这并不虚假，尽管也许其中一盏比另一盏更早熄灭。因此，如果事物在某一方面同样大，但在其他方面不同，那么它们在所有方面不完全相同的事实，并不使它们在某一方面同样大的说法失效。\par

然而，鉴于基督的话是：\BibleQuote{你们用什么量器量给人，也用什么量器量给你们}，而且毫无疑问，用以衡量事物的量器是一回事，而其中所量的事物是另一回事，显然有可能的是，人们用同样的量器，比如说，量了一蒲式耳小麦，而量给他们的却是成千蒲式耳，这样，在量器没有丝毫差别的情况下，在数量上却可能有那么大的差异，更不用说所量之物的质量差异了；因为不但可能一个人用同样的量器量给人大麦，却得到量给他的小麦，甚至他用量器量谷物，却得到量给他的金子，而且谷物可能只有一蒲式耳，而金子却可能有很多。因此，尽管在种类和数量上都有差异，但对于如此不同的事物，仍可以真实地说：\BibleQuote{他用什么量器量给人，也用什么量器量给他。}\par

再者，基督之所以说这话，从紧接的上下文便十分清楚。\BibleQuote{你们不要判断，} 祂说，\BibleQuote{免得你们受判断；因为你们用什么判断来判断，你们也要受什么判断。} 这是否意味着，若他们以不义判断了别人，他们自己也会受到不义的判断呢？当然不是；因为天主毫无不义。但这话如此表达：\BibleQuote{因你们用什么判断来判断，你们也要受什么判断，} 仿佛是说，\Quote{在你们以善意待人的意志中，你们将得自由；或在你们以恶待人的意志中，你们将受惩罚。} 例如，若有人为满足卑贱的欲望而使用眼睛，被命失明，他所听到的公正判决就是：\Quote{你因那些眼睛犯了罪，就该在它们身上受惩罚。} 因为每个人都按自己心智的判断，或善或恶，去行善或作恶。因此，他在自己施行判断的事上受审判，并不不义；这就是说，当他被迫承受那些由他心智罪恶的判断所导致的邪恶时，他就是在心智的判断官能中承受惩罚。\par

26. 虽然那些为将来所预备的、将施加的痛苦是可见的，但由同一个根本原因——即败坏的意志——所引起的痛苦，同样也发生在心智本身之内；在心智中，意志的欲望是一切人类行为的尺度，犯罪之后立刻伴随惩罚，而且这惩罚往往因犯罪者感觉不到它而变得更重，因其盲目程度更深。因此，当祂说过：\BibleQuote{你们用什么判断来判断，你们也要受什么判断（或者更确切地说，如\Person{奥古斯丁}所译，\InnerQuote{你们在什么判断里判断，你们也要受什么判断}）}，祂接着又补充道：\BibleQuote{你们用什么尺度量给人，也要用什么尺度量给你们。}也就是说，善人会在自己的意志中量出善行，也必在同样的意志中得享真福；同样，恶人也会在自己的意志中量出恶行，也必在同样的意志中遭受痛苦；因为一个人，当其意志倾慕善时是善的，同样，当其意志倾慕恶时便是恶的。所以，他也是在这一点上经历真福或痛苦，即他自身意志所感受到的，这意志既是一切行为、也是一切行为报应的尺度。因为我们衡量行为的善与恶，是根据产生这些行为的意志的性质，而非根据行为所占的时间长短。若非如此，砍倒一棵树就会被视为比杀害一个人更大的罪行了。因为前者需要很长时间和许多次砍击，后者却可能在一瞬间一击致命；然而，若一个人因这瞬间犯下的大罪，仅被处以终身流放，人们会说对他的处罚过于宽大，尽管从时间长度来看，漫长的惩罚与瞬间的谋杀行为完全无法相比。那么，如果各种惩罚同样漫长，甚至同样永恒，只是相对的温和与严厉程度不同，所反驳的那句话又有何矛盾之处呢？时间长度相同，但所施加的痛苦程度不同，因为构成诸罪本身尺度的，并不在于它们所占的时间长度，而在于犯罪者的意志。\par

27. 诚然，意志本身承受惩罚，无论痛苦是施加于心灵还是身体；如此，那由罪所满足的同一个东西，便受到惩罚的打击，并且\BibleQuote{那不行怜悯的，也要受无怜悯的审判}；因为在这句话中，尺度的标准也仅在这一点上相同，即他没有施予别人的，自己也被拒绝，因此，对他的审判将是永恒的，尽管他所发出的审判不可能是永恒的。所以，正是在罪人自己的尺度中，永恒的惩罚被量给了他，尽管这受罚的罪不是永恒的；因为正如他曾希望永恒地享受罪，他所获得的判决便是永恒地忍受苦难。\par

我在这答复中所力求的简洁，不容许我搜集我们圣书中所有关于罪及其惩罚的论述，或者至少尽我所能地搜集，并从中推出一条无可争辩的命题；而且，即使我得到必要的闲暇，或许也未必具备胜任此事的才能。然而，我认为在此期间，我已经证明了，惩罚的永恒性与罪将按人们所犯的同样尺度受报应这一原则，并无矛盾。\par

28. \Emph{第五个问题}。这位从\Person{波菲利}提出问题的反对者，接下来又补充了这一个问题：请指教我，\Person{撒罗满}是否真实地说过，天主没有\OriginalTerm{儿子}{Son}？\par

29. 答案很简单：\Person{撒罗满}不但没有如此说过，相反，他明确地指出\OriginalTerm{天主}{God}有一位\OriginalTerm{圣子}{Son}。在他的一部著作中，\OriginalTerm{智慧}{Wisdom}说道：\BibleQuote{群山未定，丘陵未立，我已生出。} 而\Person{基督}不是别的，正是\OriginalTerm{天主的智慧}{Wisdom of God}。又在\BibleBook{箴言}{}另一处说：\BibleQuote{天主教训我智慧，我学得了神圣的知识。谁曾上升高天又降下？谁聚拢风在掌握中？谁包扎水在衣裳里？谁奠定大地的四极？他名叫什么，他的儿子名叫什么？} \BibleRef{箴 30:3} 这段引文结尾的两个问题，一个指向\OriginalTerm{圣父}{Father}，即 \BibleQuote{他名叫什么？}——与前述 \BibleQuote{天主教训我智慧} 相呼应——另一个显然指向\OriginalTerm{圣子}{Son}，因为他问 \BibleQuote{他儿子名叫什么？}——与其余语句呼应，那些语句更适于理解为指\OriginalTerm{圣子}{Son}，即 \BibleQuote{谁曾上升高天又降下？}——\Person{保禄}的话可作此问的印证：\BibleQuote{那降下的就是同一位上升超乎诸天之上的} \BibleRef{弗 4:10} —— \BibleQuote{谁聚拢风在掌握中？} \Emph{即} 信徒的灵魂藏在一个隐秘的地方，因此对他们说：\BibleQuote{你们已经死了，你们的生命与\Person{基督}一同藏在天主内} \BibleRef{哥 3:3} —— \BibleQuote{谁包扎水在衣裳里？} 由此可说：\BibleQuote{你们凡是受洗归于\Person{基督}的，就是穿上了\Person{基督}} \BibleRef{迦 3:27} —— \BibleQuote{谁奠定大地的四极？} 正是那位对门徒说：\BibleQuote{你们要为我作证，从\Place{耶路撒冷}、全\Place{犹太}、\Place{撒玛黎雅}，直到地极} \BibleRef{宗 1:8}\par

30. \Emph{问题六。} 最后一个问题是关于\Person{约纳}的，其提出方式仿佛并非出自\Person{波尔菲里}，而是外教人中一个常在嘲弄的话题；因为他这样说道：\Quote{接下来，关于\Person{约纳}，我们该相信什么？据说他在鲸鱼腹中三日。这事完全不可信，难以置信，一个穿着衣服被吞下的人竟能在鱼腹中存活。然而，如果这故事是比喻性的，请解释之。再者，那故事中说\Person{约纳}被鱼吐出后，有株蓖麻长起高过他的头，这又是什么意思？这蓖麻生长的原因是什么？} 我曾见外教人在大声嘲笑和极度蔑视中讨论这样的问题。\par

31. 对此，我这样回答：要么所有由神圣能力所行的奇迹都可被视为不可信，否则便没有理由不信这个故事中的奇迹。\Person{基督}本人在第三天复活，若基督信仰害怕面对外教人的嘲笑，我们就不会相信了。然而，既然我们的朋友并未基于此质问是否应相信\Person{拉匝禄}在第四天复活，或\Person{基督}在第三天复活，我就很惊讶他认为\Person{约纳}的事迹不可信；除非他或许认为死人从坟墓中复活比活人存留在庞大海怪腹中更容易。且不提那些有识之人向我们报告的庞大海洋生物，我倒要问：\Place{迦太基}一处公共场所所固定那巨大肋骨所围的腹中，能容纳多少人？那里所有人都熟知这些肋骨。谁不能揣想那庞大洞穴入口——即那张开的巨口——该有多宽？除非，或许正如我们的朋友所说，\Person{约纳}的衣服阻碍了他毫无损伤地被吞下，仿佛他不得不挤过一条狭窄的通道，而非实际情况那样：他被从空中猛然抛下，于是被那海怪接住，在未遭其牙齿伤害之前便落入其腹中。同时，\WorkTitle{圣经}并未说明他被抛入那洞穴时是否穿着衣服，因此我们完全可以作这样的理解：他迅速落下时是赤裸的，如果可能的话，他的衣服必须被剥下，如同蛋壳被剥去，以便让他更易被吞下。其实，人们对这位先知衣着的关切程度，就好像经上说他爬过一个极小的窗户，或是进去沐浴一般；而即便如此，若在那种情况下必须穿戴整齐进去，那也仅仅是不便，而非奇迹。\par

32. 但是，或许我们的反对者认为，关于这一天主奇迹，不可能相信腹中用以分解食物的湿热气体，能被调节到如此温度，以致能保存一个人的生命。若真如此，他们又将以何等的力量宣称，那三位被不虔敬的君王投入火窑中的青年，竟能在火焰中行走而不受伤害，这是不可信的！因此，这些反对者若拒绝相信\Emph{任何}记载天主奇迹的叙述，就必须用另一种论据来反驳。他们的责任不是单单挑剔某一事迹，质疑其不可信，而是要谴责所有记载同类或更为奇特奇迹的叙述为不可信。然而，如果关于\Person{约纳}的这段记载，被说成是由\Place{马达乌拉}的\Person{阿普列乌斯}或\Place{提亚纳}的\Person{阿波罗尼乌斯}——他们夸耀这些人行了许多奇事（尽管魔鬼也做一些类似圣天使的工作，不是在真理中，而是在外表上，不是凭智慧，而是明显凭诡诈），即使没有可靠证据的支持——所做的，我说，倘若这些被他们谄媚地称为术士或哲学家的人行过此类事迹，我们就会从他们口中听到的，不是嘲笑，而是胜利的欢呼。那么，就这样吧；让他们嘲笑我们的\WorkTitle{圣经}；让他们尽情嘲笑吧，当他们看到自己的人数日渐减少，有的因死亡离去，有的因信奉了基督信仰，并且看到过去嘲笑他们、预言他们将徒然与真理对抗并徒然狂吠、且将逐渐归向我们这边的先知们所预言的一切正在兑现时；这些先知不仅把这些陈述传给我们后代学习，还许诺它们将在我们的经验中应验。\par

33. 探讨这些奇迹有何寓意，既非不合理，也非无裨益，这样，在解释了它们的意义之后，人们便不仅相信它们确实发生过，而且相信它们之所以被记载，正是因其具有象征含义。因此，凡欲探究为何先知\Person{约纳}在大鱼的巨腹中三日的人，首先应消除对事实本身的怀疑；因为这确实发生了，且不是徒然发生的。如果仅以言语而非行动表达的象征有助于我们的信德，那么那些不仅以言语、更以行动表达的象征，岂不更应有助于我们的信德！人们惯于用言语说话；但天主的能力不只用言语，也通过行动说话。正如新颖或略显生疏的词语，若以适度且审慎的方式散置于人间演讲中，能为其增添光彩，同样，天主的启示的雄辩，也可说因那些本身既神奇又巧妙设计以传授灵性教导的行动而倍增光辉。\par

34. 至于大海怪在第三日将吞下的先知活活吐出来，预表了什么？这个问题何必问我们，因为基督自己已经回答了，祂说：\BibleQuote{邪恶淫乱的世代要求征兆，除了约纳先知的征兆外，必不给它任何征兆。正如约纳曾在大鱼腹中三天三夜，人子也必在地心三天三夜。} \BibleRef{玛 12:39} 至于主基督受死亡辖制的三天，如何被算为整整三天，即三天三夜，因为第一天和第三天各自的一部分被当作整体来理解，这需要长篇解释；而且这一点在其他讲论中已屡次详述。因此，正如\Person{约纳}从船转入大鱼腹中，\Person{基督}也从十字架进入坟墓，或进入死亡的深渊。又如\Person{约纳}为了那些在风暴中遇险的人而承受这事，\Person{基督}也为了那些在这世界的波浪中颠簸的人而受苦。并且起初命令\Person{约纳}向\Place{尼尼微}人宣讲天主的言语，但\Person{约纳}的宣讲直到大鱼将他吐出后才临到他们；同样，先知的教导早早便面向外邦人，但实际临到外邦人却是在\Person{基督}从坟墓中复活之后。\par

35. 其次，关于\Person{约纳}为自己搭了一个棚，坐在\Place{尼尼微}对面，等着看那城将要遭遇何事，这位先知在这里以他本人作为另一事实的象征。他预表了那属血肉的以色列子民。因为他也对尼尼微人的得救，即外邦人的救赎和解救感到不悦，而\Person{基督}正是从外邦人中来，不是召叫义人，而是召叫罪人悔改。\BibleRef{路 5:32} 因此，遮在他头上的那棵蓖麻的阴影，预表了旧约的应许，或更确切地说，旧约中已享有的特权，正如宗徒所说，其中\BibleQuote{含有未来事物的影子}，\BibleRef{哥 2:17} 在应许之地，它宛如提供庇护，使人免受今世苦难的炎热。再者，那条以啃咬的牙齿使蓖麻枯萎的早晨之虫，再次预表了\Person{基督}，因为借着从祂口中所宣讲的福音，所有那些曾在以色列人中一时兴盛，或在先前安排中具有影像象征意义的事物，如今都失去了意义而枯萎了。如今那个民族，丧失了王国、司祭职以及在\Place{耶路撒冷}先前设立的祭献——这一切都是未来事物的影子——在散居和掳掠的境遇中，遭受着沉重苦难的灼热，正如按历史记载，\Person{约纳}被太阳晒得发昏，且痛苦不堪；尽管如此，外邦人及悔改者的得救，在天主眼中，比以色列的这份忧伤以及那曾令犹太民族如此欣喜的\Quote{影子}更为重要。\par

36. 再次，让异教徒嘲笑吧，让他们傲慢而愚昧地讥讽基督是蠕虫，讥讽这个对预言象征的解释，只要祂还是逐渐且确实地吞噬他们。因为关于这一切，依撒意亚先知曾预言，天主借着他向我们说话：\BibleQuote{认识正义、心中存有我的法律的百姓啊，你们要听我！不要怕世人的辱骂，也不要因他们的毁谤而恐惧。因为蛀虫必吃掉他们，像吃掉衣服；蠕虫必吃掉他们，像吃掉羊毛；但我的正义必永存。} \BibleRef{依 51:7} 因此，我们要承认基督就是那晨蠕虫，此外，在那篇题为\WorkTitle{早晨的牝鹿}的圣咏中，祂乐意用这个名称称呼自己：\BibleQuote{我是}，祂说，\BibleQuote{是蠕虫，而不是人，世人的羞辱，百姓的轻蔑。}这羞辱正是依撒意亚的话中命令我们不要恐惧的那种羞辱：\BibleQuote{不要怕世人的辱骂。}那些人正被那条蠕虫，如同被蛀虫一样吞噬；他们在祂福音的牙齿下，每日惊见自己的人数因脱党而减少。所以，我们要承认这是基督的象征；并且为了天主的救恩，要耐心忍受世人的羞辱。祂是蠕虫，因为祂所取肉体的卑微——或许也因为是童贞女所生；因为蠕虫通常不由交合而生，而是自发产生于肉中或任何植物产物中[\OriginalTerm{不经交合而生}{sine concubitu nascitur}]。祂是那\Emph{早晨的}蠕虫，因为祂在天亮以前从坟墓复活了。那棵蓖麻当然可能没有根部的蠕虫也枯槁；最后，如果天主认为蠕虫为完成这项工作所必需，那么为何要加上\Emph{早晨的}这个形容词，若不是为了叫人认出祂就是那条曾在圣咏中，在\OriginalTerm{为早晨的救助}{pro susceptione matutina}的标题下咏唱：\BibleQuote{我是蠕虫，不是人}的那位呢？\par

37. 那么，还有什么比所预言之事得以实现更能显明这预言的应验呢？那条蠕虫确实曾在祂悬在十字架上时受人轻蔑，正如同一圣咏所记载的：\BibleQuote{他们撇唇，摇头说：他信赖主，主就该救他；主既喜爱他，就该拯救他。}又当这圣咏所预言的应验时：\BibleQuote{他们穿透了我的手和我的脚，他们数遍了我的骨骼，他们却冷眼望着我。他们分开我的外衣，为我的内衣拈阄。}——这些事在古代经卷中，先知描述将来时，其清晰程度，正如如今在福音史中，它们发生后被记述下来一样。但若那条蠕虫在受辱时被人轻蔑，那么当我们看见同一圣咏的后半部分所预言的这些事实现时，祂还要被轻蔑吗：\BibleQuote{整个大地将记起，归向主；万民的家族都要在祂面前叩拜。因为王国属于主；祂必统管万国}？因此，尼尼微人\BibleQuote{记起了，归向了主。}外邦人悔改而获得的救恩，早在那时就如此被预表，而以色列，如约纳所代表的，那时心怀悲戚，如今他们的民族也悲戚，失去了荫庇，并因痛苦的煎熬而忧愁。任何人可以自由地以不同的解释去揭示先知约纳象征历史中隐藏的其他细节，只要那解释与信德的准则相符；但显然，关于他在大鱼腹中度过的三日，除了按照天上导师在福音中亲口揭示的，如上文所引，以其他方式解释是不合法的。\par

38. 我已尽我所能回答了所提的问题；但提出问题的人应立刻成为基督徒，免得他若耽延，直到解决了有关圣经的一切疑难问题后才结束讨论，他会在由死亡转入生命之前，已走到了此生的尽头。因为他在被准入领受基督徒的圣事之前，探究关于死人复活的道理，是合理的。或许还应允许他坚持预先讨论所提出的关于基督的问题——祂为何在世界历史中来得如此之晚，以及除此之外的少数几个重大问题，而其他问题都属于这些之下。然而，若想在他成为基督徒之前，就解决完所有诸如关于\BibleQuote{你们用什么尺度量给人，也要用什么尺度量给你们}这句话、以及关于约纳之类的问题，那便暴露出对人之有限能力，以及余生之短暂的极大忽视。因为有无数的疑问，其解答不该在我们相信之前就要求解答，免得我们在不信中结束了生命。然而，当基督徒的信仰被彻底领受后，这些疑问便当以最大的勤勉去探究，为的是虔敬地满足信徒们的心灵。这种探究所发现的任何东西，都当毫无虚荣自满地与人分享；若仍有什么奥秘隐藏，我们必须耐心忍受知识上的缺欠，而这并不危害救恩。\par

\SourceRecord{Translated by J.G. Cunningham.；Nicene and Post-Nicene Fathers, First Series；Vol. 1.；Edited by Philip Schaff.；Buffalo, NY: Christian Literature Publishing Co.,；1887.；<http://www.newadvent.org/fathers/1102102.htm>.}

% structure: p0001 source=h1 target=section reason=page-title

\section{书函一〇三（主历409年）}

\Emph{致我的主和弟兄\Person{奥斯定}，堪受一切敬意与全备尊崇的，\Person{纳克塔里乌斯}在主内致候。}\par

阅览阁下的大函，您在其中推翻了偶像崇拜及其庙宇仪式，我彷佛听到一位哲学家的声音；不是像\OriginalTerm{学园派}{academician}那种哲学家——据说，这种人既没有新学说要提出，也没有自己的旧有主张需要维护，惯于坐在阴暗角落的地上，沉浸在深思之中，双膝缩至额前，头埋在膝间，盘算着如何诋毁别人已有的发现，或对他人赢得的声誉妄加挑剔；不，在您雄辩的魅力下，浮现于我眼前的形像，更像是那位伟大的政治家\Person{西塞罗}。他曾因拯救众多同胞的生命而荣膺桂冠，带着他法庭上的胜利成果，进入令世人惊奇的希腊哲学学园；那时他彷佛停下来喘口气，放下了那支号角——他曾以义愤的呼号吹响它，发出嘹亮的声调与言辞，抨击那些违法乱纪、图谋危害共和国生命的人——并仿效希腊披风的式样，解开托加袍宽大的褶皱，向后搭在肩上。\par

因此，当您敦促我们敬拜并信奉唯一至高的天主时，我欣然聆听；当您劝勉我们仰望天上的父邦时，我喜悦地接受了这一劝谕。因为您显然不是对我谈及那些为城墙所限的城池，也不是哲学家著作中所提到并宣称以全人类为公民的尘世国度，而是谈及那座由伟大天主及那些从祂那里获得奖赏的灵体所居住和拥有的\OriginalTerm{圣城}{City}，所有宗教都经由不同的道路和途径向往它，——这座圣城我们无法用语言描述，或许可以通过思想略窥其貌。然而，尽管这座圣城因此应当被我们渴望和爱慕超越一切，但我依然认为，我们不应对自己的本土不忠——这片土地首先让我们得享天日，我们在其中被养育和教化（而且，在此特别相关的是，藉着为这片土地服务，人死后才能在天上获得为他们预备的住处）；因为，按照博学之士的意见，得以升入那座天上圣城的，是那些对自己出生城邦有功之人；而那些被证实以自己的谋略或劳苦为故土福祉作出了贡献的人，则会获得与天主交往的更高体验。\par

至于您欣然以妙语提及我们城镇的评语，说它之出名并非由于战士的功绩，而是由于纵火者的焚烧，又说它出产的是荆棘而非花朵，这并非最严厉的责备，因为我们知道，花朵大多长在带刺的灌木上。因为谁不知道，连玫瑰也生长在野蔷薇上，而谷物的芒须中，麦穗由尖刺护卫；总的说来，令人愉悦与令人痛苦之物常是如此交织在一起的？\par

阁下信中的最后一点是，为补偿对教会所行的不义，既不要求判处死刑，也不要求流血，只是必须剥夺冒犯者最怕失去的财产。然而，依我审慎的判断——虽然我当然也可能犯错——丧失财产比丧失性命更为悲痛。因为，如您所知，在整个文学领域频繁出现的观点是，死亡结束了一切罪恶的体验，而贫困的一生只会带给我们无尽的悲惨；因为悲惨地活着，比以死亡结束悲惨更糟。您整个工作的性质和方式也表明了这一点，您援助穷人，为患病者施行医治，为痛苦的人施用药剂，总之，您的职责在于防止受苦者感到痛苦的持续延长。\par

再者，关于某些过错与其他过错相比的过失程度，在请求宽恕的案件中，过错的性质显得如何并不重要。因为，首先，如果悔改能获得宽恕并补赎罪行——而那些恳求原谅、谦卑地抱住他所冒犯者双脚的人，确实就是悔改者——再者，如果如某些哲学家所认为的，一切过失都是等同的，那么宽恕就应当不加区别地赐予所有人。我们的一个市民也许说话有些粗鲁：这是一个过失；另一个可能施以侮辱或伤害：这同样是一个过失；另一个可能强行夺取不属于他的东西：这被视为罪行；另一个可能攻击用于世俗或神圣用途的建筑物：他不应因此罪行就被置于宽恕的范围之外。最后，如果没有先前的过失，也就没有宽恕的场合了。\par

4. 既已回复了你的来信——虽非按这封信所应得的，而是尽我所能、仅此而已——我恳切地哀求你（啊，但愿我在你面前，这样你也可以看见我的眼泪！），一再思考你是谁，你公开宣认的身份是什么，你终生投身的事业又是什么。思考一下那城镇的景象：被判酷刑的人被拖出来；想想母亲们、妻子们、儿子们和父亲们的哀号；想想那些可能获释回来的人——确实重获自由，却已遭受了酷刑——所感到的羞耻；想想看到他们的伤口和疤痕必定会重新引发的悲伤和呻吟。当你深思这一切之后，首先要想到天主，想到你在人间的美名；更确切地说，要想到友善的\OriginalTerm{爱德}{charity}和共同人性之纽带要求你做什么，并力求因宽恕罪犯而非惩罚他们而受到称赞。事实上，关于你对待那些基于自己的招供而被实际罪过定罪之人的方式，可以说这样的话：你出于对你的宗教信仰的尊重，已对他们施以了宽恕；为此，我会永远称赞你。但现在，几乎无法表达那种残酷之深重：它追逼无辜者，传召那些确知未参与所指控罪行的人来接受死罪审判。倘若他们恰巧被无罪开释，我恳请你想想，他们的控告者会怀着怎样的恶意看待这一开释——这些控告者自行放过了法庭上的有罪之人，却只在企图加害无辜者而失败时，才让无辜者离去。\par

愿至高的天主作你的守护者，并保全你，使你成为祂的宗教的堡垒和我们国家的荣耀。\par

\SourceRecord{Translated by J.G. Cunningham.；Nicene and Post-Nicene Fathers, First Series；Vol. 1.；Edited by Philip Schaff.；Buffalo, NY: Christian Literature Publishing Co.,；1887.；<http://www.newadvent.org/fathers/1102103.htm>.}

% structure: p0001 source=h1 target=section reason=page-title

\section{书信104（公元409年）}

\Emph{致\Person{内克塔里乌斯}，我尊贵的大人与兄弟，堪受一切尊崇与敬重者，\Person{奥古斯丁}在主内问候。}\par

% structure: p0003 source=h2 target=subsection reason=relative-heading-rank; base=h2

\subsection{第一章}

1. 我已读了你好意寄来的回信。相较于我寄出我的信给你，你的回覆经过了极长的时间间隔。因为当我的圣善兄弟与同僚\Person{波西迪乌斯}仍与我们同在、尚未启程远航时，我便已写了回信给你；但你托他带交给我的信，我于3月27日收到，那已是写信给你大约八个月之后了。究竟为何我的信如此迟才到你那里，或你的信如此迟才寄给我，我不得而知。或许是你的明智如今才下笔答覆，而这答覆你先前曾因骄傲而不屑一顾。倘若如此，我不禁想知道是什么导致了这一转变。莫非你偶尔听到了什么我们尚无知的传言，说我兄弟\Person{波西迪乌斯}已获授权，可对你的市民采取更严厉的行动——你当原谅我这么说——他比你更以能促进他们福祉的方式爱他们？因为，当你嘱咐我设想一下\Quote{那城市所呈现的景象：命定受刑的人被拖出来}，并\Quote{想想母亲们、妻子们、儿子们和父亲们的哀号；想想那些可能返回的人所感到的羞耻，他们虽获释放，却已遭受了酷刑；再想想那看见他们的伤口与伤疤时必然重新激起的悲伤与呻吟}时，你的信表明了你对此种情况的忧虑。我们绝不愿自己，也不愿任何人，将这样的苦难加于我们的任何仇敌！但是，如我所说，如果传言已将此类严酷措施带进你的耳中，就请你更清楚、更详细地陈述所传之事，使我们得以知道该做什么来阻止这些事发生，或该怎样回答，以消除信谣之人的疑虑。\par

2. 请更仔细地检视我那封信，你回覆得如此勉强的那封信，因为我在其中已充分表明了我的看法；但我猜想你已忘了我所写的，在你的回信中提及了一些与我的观点极为不同、毫不相似的意见。因为你仿佛凭记忆引用我所写的，将一些我从未说过的话插入了你的信中。你说我那封信的结尾句是：\Quote{为了补偿对教会的不公，既不需要死刑，也不需要流血，但那些犯罪者必须被剥夺他们最怕失去的东西}；然后，在说明这将导致多么大的灾难时，你加上并与我的话语联系起来，说你\Quote{慎重判断——尽管你也许错了——被剥夺财产比被剥夺生命更为痛苦。} 为了更清楚地说明你所指的财产的种类，你接着说，我必须知道，\Quote{在整个文学领域中，这观察经常出现：死亡终结了一切罪恶的体验，但贫穷的生活只会赐给我们永恒的不幸。} 由此你得出结论：\Quote{凄惨地活着比藉死亡结束我们的痛苦更糟。}\par

3. 就我而言，我不记得曾在任何地方读到过——无论是在我们的基督教文献中（我承认我热心于此的时间，较我心愿的更晚），还是在你们的异教文献中（我自幼研读）——有\Quote{贫穷的生活只会赐给我们永恒的不幸}这样的话。因为勤劳者的贫穷本身绝非罪恶；不，在某种程度上，它是使人远离并约束罪恶的一种方法。因此，人今生生活贫困这一事实，不足以让我们担忧这会在短暂生命之后带给他\Quote{永恒的不幸}；而且在我们度过的尘世生活中，任何悲惨都不可能永恒，因为此生不能永恒，不，即便是那些活到极老年龄的人，此生也不长久。在提及的著作中，就我而言，我所读到的，并不是说在此生中——如你所想，并如你所声称这些著作经常断言的那样——会有永恒的不幸，而是说我们在此享受的此生本身是短暂的。你们的有些作者，并非全部，曾说过死亡是一切罪恶的终结：那确实是伊壁鸠鲁派和其他相信灵魂会死之人的意见。但那些被西塞罗以一定意义称为\OriginalTerm{安慰者}{consulates}的哲学家们，因为他对他们的权威颇为重视，持守这样的意见：当我们尘世最后一刻来临时，灵魂并非消灭，而是离开它的居所，继续存在，处于福乐或悲惨之境，按照人所行善或恶应得的报偿。这符合我们圣经的教导，我渴望能更全面地熟悉这教导。因此，死亡是一切罪恶的终结——但这只针对那些生命纯全、虔诚、正直、无可指摘的人；并不针对那些被短暂琐事与虚妄的炽烈欲望所点燃，因其欲望的彻底败坏而被证明为悲惨的人，虽然同时他们自以为幸福，死后却不得不不仅承受，而且在经验中体认到远为巨大的悲惨。\par

4. 因此，这些看法既频繁地表达于你所尊崇的作者的一些作品中，也表达于我们的所有圣经中。啊，你这位地上祖国当之无愧的爱人，你理当为你的同胞恐惧那奢侈放纵的生活，而非贫困的生活；或者，如果你惧怕贫困的生活，那就警告他们：切需躲避的贫穷，乃是那种人——虽拥有丰裕的世俗财物，却因对之贪求的永不满足，常处于匮乏之中，用你们自己的作者们的话说，这匮乏是丰裕与匮乏都不能缓解的。然而，在你所回复的那封信中，我并未说，你那些与教会为敌的市民，应被逼至缺乏生活必需品的那种极端贫困，并且给此贫困带来援助的，正是你所认为必须向我指出的，我们会在一切劳苦中声称自己所行的慈悲——即我们所引的\Quote{支持穷人，施医予病者，给痛苦者之身施以救治}；尽管确实，即使是如此极端的匮乏，也比滥用一切丰裕以纵邪恶欲情更为有益。但是，让我绝不要想，我们所论及的那些人，应通过所提议的强制措施被陷于如此赤贫。\par

% structure: p0008 source=h2 target=subsection reason=relative-heading-rank; base=h2

\subsection{第二章}

5. 虽然在你看来，当你要回复时，不值得通读我的信，也许你至少如此看重它，以至于保存下来，为能在你任何想要并取用它的时候，可以拿到它；若果真如此，请再次细阅，仔细留意我的话：你必将在其中发现一件事，以我之见，你必然承认你对此并未作答。因为在那封信中，出现了我现在所引的话：\Quote{我们不希望通过对过去的报复性惩治来满足我们的愤怒，而是关切于以真正慈悲的精神为未来作好安排。如今恶人具有某些可受罚之处，基督徒能用慈悲的方式予以惩罚，以促进他们自身的利益与福祉。因为他们拥有三样东西——身体的生命和健康、维持生命的手段，以及过邪恶生活的手段和机会。对于悔改己罪的人，但愿前两者不受触动而保留给他们；这是我们渴望的，也是我们不遗余力去保证的。至于第三样，若天主乐意将它视为一个腐烂或患病的部分，必须用修剪刀切除，那么祂将在这种惩罚中证明祂慈悲的伟大。}如果你当初乐于写回信时，能再次读到我这些话，你就会把向我提出请愿——不仅求免死，而且求免拷打——视为非善意的含沙射影，而非出于必要的职责，而我曾说到我们希望让那些人身体生命和健康的拥有不受损害。也没有理由担心我们会使那些人陷于贫困生活，或要靠别人供给每日的食粮；因为我曾说过，我们渴望保证他们上面所提三样东西中的第二样，即维持生命的手段。至于他们的第三样东西，即过邪恶生活的手段和机会，也就是说——暂且不论其它——他们用来铸造其伪神像的银子，他们就是在这些伪神的庇护、崇拜或亵渎的敬拜中，甚至试图烧毁天主的教会，而本为救济极虔诚者贫穷的物资，竟落入可悲暴民之手成为掠物，并任由鲜血横流——我要问，你那颗爱国的心，为什么害怕斩断这等事物的那一击，以免那致命的胆大妄为，因免受惩罚，而在万事中得到助长和巩固？我恳请你充分讨论此事，用深思熟虑的论据，向我指出这样做有什么不对；请仔细留意我的话，以免在你就我所说之事提出恳求的幌子下，你显得是在对我们发出间接指控。\par

6. 愿你的同胞因品德高尚而非因过多财富而得美名；我们不希望他们因我们的缘故，通过强制手段而落到\Person{辛辛纳图斯}的犁头或\Person{法布里奇乌斯}的灶台那般境地。然而，正是由于如此极度的贫穷，罗马共和国的这些政治家不仅没有招致同胞的轻蔑，反而因此特别受到爱戴，并被认为更有资格管理国家的资源。我们无意也未曾试图削减你们富人的产业，使其所剩无几，就如\Person{鲁菲努斯}那般，他虽两次担任执政官，却仅拥有十磅银子，而当时严厉的监察官还值得称赞地要求将其进一步削减，因其数量之大实属该罚。我们深受堕落时代流行情绪的影响，以更为温和的方式对待那极其脆弱的心理，以至于在基督徒的仁慈看来，那时监察官所认为的公义措施都显得过分严厉；然而你看，这两种情况之间的差异何其大：一种情况是，仅仅拥有十磅银子是否该被视为应受惩罚的罪行；另一种情况则是，是否允许任何人在犯下其他极大罪行之后，仍可保留前述这笔钱。我们所求的只是，昔日自身即构成罪行的行为，在我们的时代被定为对罪行的惩罚。然而，有一件事可行且应行：一方面，严厉不应推至我所提及的如此极端；另一方面，人们不应因抱有免罚之念而沉溺于过度的狂欢与暴乱，从而为其他不幸的人树立榜样，效仿之后便会招致最严厉、最闻所未闻的刑罚。至少请你们应允这一点：让那些企图以火与剑摧毁我们生活必需品的人，因害怕失去那些他们有害地充裕拥有的奢侈品而恐惧。还请允许我们为我们的敌人带来这一益处：通过使他们害怕那些失去也无损于己的东西，来阻止他们尝试那些会伤害自身的事。因为这应称之为明智的预防，而非对罪行的惩罚；这不是施加刑罚，而是保护人们免于遭受刑罚。\par

7. 当有人采取会带来某些痛苦的措施，以避免一个轻率的人因习惯于徒劳无益的罪行而招致最可怕的刑罚时，他就如同一个人揪住男孩的头发，以防止他拍手挑逗蛇一样；在这两种情形中，爱的行动固然使对象感到烦恼，身体却未受伤害，而受阻止所做的那些事本会危及安全与生命。我们施惠于人，并非每次都应允所求，而只在我们所行不损害请求者之时。因为在大多数情况下，我们不给予反倒是最好地服务他人，而给予他们所渴望的反而会伤害他们。因此有谚语云：\Quote{不要将剑放在孩子手中。} \Person{西塞罗}则说：\Quote{不仅如此，连你的独生子也当拒绝给他。因为我们越爱一个人，就越应当避免把那些容易导致极大过犯的东西交托给他。}如果我没有记错，他发表这些见解时指的是财富。因此，当某些东西被从那些保管它们有危险的人手中拿走时，对他们自身往往是一种益处，以免他们滥用这些东西。当外科医生看到坏疽必须切除或烧灼时，他们常常出于怜悯，对许多哭喊充耳不闻。倘若我们年幼时，每次犯有过失便恳求宽恕，而父母和教师总是宽容免罚，我们中又有谁不会变得令人无法忍受呢？我们之中又有谁能学到任何有用的东西呢？这样的惩罚是出于明智的关爱，而非出于蓄意的残忍。我恳求你，在这件事上，不要只考虑如何成就你的同胞所请求你做的事，而应仔细思量事情的方方面面。如果你忽视已成过去、无法挽回之事，那就请考虑未来；要明智地留意，不是对向你提出请求之人的欲望，而是对他们的真正利益。如果我们所宣告爱的人们，我们唯一关心的是唯恐拒绝他们的请求会削弱他们对我们的爱，那么我们就确是对自己不忠。那么，即便你们的文献所赞颂的那种美德又如何呢？那种美德体现在那治理国家的统治者身上，他远非只考虑百姓的愿望，而是更考虑他们的福祉。\par

% structure: p0012 source=h2 target=subsection reason=relative-heading-rank; base=h2

\subsection{第三章}

8. 你说：\Quote{在恳求宽恕的任何情况下，过失的性质如何并不重要。} 如果所涉问题在于对人的惩罚而非纠正，你这话就是说出了真理。基督徒的心绝不应受复仇欲望驱使而惩罚任何人。基督徒在宽恕他人的过犯时，绝不应不先期预见或至少立即回应那恳求宽恕者的请求；但他这样做的目的，当是为了克服那憎恨冒犯他人、以恶报恶、并被怒火煽动——若实际上未施伤害，也至少希望看到法律规定的刑罚施加到对方身上的诱惑；而不应是为了使自己免于考虑犯错者的利益，为他深思远虑并阻止他行恶。因为一方面，因更为强烈的敌意所驱动，人可能疏于纠正他所深恨之人；另一方面，通过包含某种苦痛的纠正，人可能确保他所深爱之人的改善。\par

9. 我承认，如你所写，\Quote{补赎能求得宽恕，并涂抹过犯，}但那是真宗教影响下所实行的\OriginalTerm{补赎}{penitence}，并且着眼于天主将来的审判；而不是那种在人面前暂时公开表示或假装，不是为了永远洁净灵魂除去罪过，而只是为了将即将逝去的生命从眼前痛苦的恐惧中解救出来的补赎。正因如此，对于某些基督徒，他们承认自己的过错，并因卷入那罪行的罪责而祈求宽恕——无论是没有保护面临被焚危险的教会，还是私吞了一部分恶人抢走的财物——我们相信\OriginalTerm{悔改}{repentance}的痛苦已结出果实，并认为这足以纠正他们，因为在他们的心中存有信德，藉此他们能明白因自己的罪当从天主的审判中畏惧什么。但是，那些不仅不承认，甚至坚持嘲笑并亵渎那宽恕之源者，他们的悔改怎能有任何医治的功效呢？同时，我们对这些人心中不怀任何敌意，我们的心赤裸敞开在祂的眼前，我们畏惧祂在今生和来世的审判，并将我们的希望寄于祂的助佑。但我们认为自己甚至是在采取措施为这些人谋利益，如果看到他们不敬畏天主，我们就通过做一些事使他们心生畏惧，藉此惩戒他们的愚妄，而他们的真正利益却不受损害。这样，我们就防止他们所藐视的天主因他们更大的罪行而被更重地激怒，他们会因灾难性的\OriginalTerm{不受惩罚}{impunity}的保证而胆大妄为，我们也防止他们不受惩罚的保证被更有害地展示出来，作为鼓励他人效仿的榜样。最后，对于那些你向我们为其代求的人，我们在天主面前为他们代祷，恳求祂使他们归向自己，教导他们施行真诚而有救赎功效的悔改，藉信德洁净他们的心。\par

10. 那么，请看我们是怎样爱那些你猜想我们对其充满愤怒的人——请允许我说，我们以比你对他们所怀的更为明智、更有裨益的爱去爱他们；因为我们为他们恳求，愿他们能逃脱远更大的苦难，并获得远更大的祝福。如果你也以爱去爱这些人，不是以人纯粹属世的感情，而是以那种作为天主属天恩赐的爱，并且如果你写信给我说你曾愉快地聆听我向你推荐至高天主的敬礼和宗教是真心的，那么你不仅会希望你的同胞得到我们为他们寻求的祝福，更会亲身以榜样的力量引领他们去获得这些祝福。这样，你为我们求情的整个事务就会以丰盛且极其合理的喜乐而圆满结束。这样，你对于那天上的故乡——你说你曾欣然接受我的劝告，将目光注视于其上——的权利，就会通过你对自己出生国度的真实而虔诚的爱而获得，那时你会努力确保你的同胞得到的，不是现世幸福虚妄的梦想，也不是对其过错应有惩罚的最危险的免除，而是永恒真福的恩宠。\par

11. 在此，关于此事我心中的想法和愿望，我向你坦诚相告。至于隐藏在天主圣意中之事，我承认我并不知道；我只是一介凡人；但无论如何，祂的旨意极为坚定，其公平与智慧无可比拟地超越人所能想象的一切。在我们的圣经中真实地说：\BibleQuote{人心中有许多谋算；惟独主的筹算才能立定。}见：\BibleRef{箴 19:21}。因此，至于时间会带来什么，会出现什么使我们的措施简化或复杂化，简言之，会因害怕失去或希望保有其现有财产而突然激发什么欲望；至于天主是否会因他们的所作所为而如此不悦，以致以灾难性的不受惩罚的更沉重、更严厉的判决惩罚他们，或者会指定以我们所提议的方式慈悲地纠正他们，或者会扭转正在为他们预备的某种可怕厄运，并以更为严厉却更有救赎功效的纠正，使之化为喜乐，引领他们真诚地转而不向人而向祂寻求怜悯——这一切祂全知道；我们不知道。那么，为何阁下与我要在此事上提前徒然劳苦呢？让我们暂且放下那时刻尚未到来的忧虑，如果你愿意，让我们专注那始终紧迫之事。因为没有哪个时候不适合且不需要我们思考如何能使天主喜悦；因为人在今生完全达至毫无任何罪过的完美，要么是不可能的，要么（若万一有人达到）是极其困难的：因此我们应当毫不迟延地立即投奔到祂的恩宠之中，对祂我们可以完全真实地说出某位诗人向一位显赫人物所说的话，那位诗人宣称这些诗句借自一个\OriginalTerm{库迈神谕}{Cumæan oracle}，或先知灵感的颂歌：\Quote{以你为领袖，我们所有罪恶残余痕迹的涂除，将使大地免于永久的惶恐。}因为以祂为领袖，所有罪恶都被涂除并宽赦；通过祂的道路，我们被引向那天上的故乡，当我竭尽全力将对其作为居所的思念托付给你的喜爱和渴望时，这思想曾使你深为喜悦。\par

% structure: p0017 source=h2 target=subsection reason=relative-heading-rank; base=h2

\subsection{第四章}

12. 但是，既然你说一切宗教都经由不同的道路和途径而向往同一居所，我担心，在你以为你如今所走的道路正趋向那里时，你或许会有些不愿接受那条唯一领人到达天堂的道路。然而，我更仔细思考了你所用的词语，我想，我用略有不同的方式解释其意义，也不算冒昧；因为你并没有说一切宗教经由不同的道路和途径而到达天堂，或揭示、或找到、或进入、或获得那福地，而是通过一句刻意斟酌而选用的短语说一切宗教都向往它，这表明，一切宗教共有的是对天堂的渴望，而非享有。你在这些话中既没有排斥那唯一真实的宗教，也没有接纳其他虚假的宗教；因为固然，引我们达到目标的道路是向往那里的，但并非每条向往那里的道路都把人带到那里，而所有被带到那里的人必定是蒙福的。如今，我们全都愿意，即向往，成为蒙福的；但我们并不能全都实现我们所愿意的，即我们不能全都获得我们所向往的。因此，唯有那行走在不但向往那里、而且确实把他带到那里的道路上的人，才得到天堂，他使自己与那些固守虽向往天堂却最终未能到达天堂的道路的人分开。如果人满足于无所向往，或人所向往的真理既已得到，便不会有迷失。然而，如果你在说\Quote{不同的路}时，意非指相反的路，而是指不同的路，就如我们说不同的规诫，它们全都趋向于建立圣善的生活——一个规定贞洁，另一个规定忍耐、信德或仁慈等等——若道路和途径仅在这个意义上是不同的，那么那个国度便不仅是向往的，而且是确实找到的。因为在圣经中，我们既读到多条道路，也读到一条道路——关于多条道路，\Emph{例如}在经上说：\BibleQuote{我要给恶人教导你的道路，罪人们都要回头，向你奔赴；}关于一条道路，\Emph{例如}在祈祷中说：\BibleQuote{求你教训我你的途径，叫我照你的真理去行。}那些道路与这条道路并无不同；而是在一条道路中包含了圣经在另一处所说的一切道路：\BibleQuote{上主的一切行径是慈爱和忠信。}对这些道路的细心研究，可成为长篇论述的主题，也是极为愉悦的默想；但这我要留到别的时候，如果需要的话。\par

13. 然而，在此期间——我想，这在对阁下作本次回复时或许已足够——既然基督说过：\BibleQuote{我是道路}\BibleRef{若 14:6} 那么仁慈和真理就当在他内寻求：如果我们以任何其他道路寻求这些，我们必会迷失，走入一条虽向往真正目标、却不能领人到那里的路。例如，如果我们决定遵循你提到的那句格言所指示的道路：\Quote{一切罪皆相等，} 这岂不将我们领入绝望的流亡，远离那真理与福乐的故乡吗？还有什么能比说，那过分粗鲁地大笑的人，与那狂怒地放火烧了他城市的人，被判为犯了同等的罪，更为荒谬无理的呢？这个见解，并非是许多通往天上居所的不同的路之一，而是一条邪路，必然引向最致命的错误，你认为有必要引用某些哲学家的这话，并非因为你赞同这种看法，而是因为它有助于你为你的同胞恳求——让我们按同样条件宽恕那些发怒焚烧了我们教堂的人，如同宽恕那些可能用某种无礼的辱骂攻击了我们的人一样。\par

14. 但是，请与我一同重新思考你用以支持你立场的推论。你说：\Quote{如果，如某些哲学家的主张，一切过失皆相等，便应一律不予区分地宽恕一切。}之后，你似乎努力证明一切过失皆相等，接着说道：\Quote{我们的一个同胞可能说话有些粗鲁：这是一个过失；另一个人可能犯下了侮辱或伤害：这同样是一个过失。}这不是教导真理，而是在没有证据支持下提出对真理的颠倒。因为对于你的陈述\Quote{这同样是一个过失，}我们立刻予以直接的否认。你或许要求证明；但我回答，你为你的陈述提供了什么证明？我们是否应把你下一句话当作证据：\Quote{另一个人可能用暴力夺取了不属于他的东西：这被视为不法行为}？这里你自己承认对引用的这句格言感到羞愧；你没有把握说这同样是一个过失，而是说\Quote{这被视为不法行为。}但这里的问题不在于这是否也算作不法行为，而在于这次冒犯以及你提到的其他行为是否等值的过失，当然，除非因为它们都是冒犯，就必须宣布它们相等；如果那样，老鼠和大象就必须被宣布相等，因为两者都是动物，苍蝇与鹰也因为两者都有翅膀。\par

15. 你更进一步，提出这样的主张：\Quote{另一个人可能攻击了用于世俗或神圣目的的建筑：他不应因为这一罪行而被置于无法获得宽恕的境地。}在这个句子中，你确实触及了你的同胞们最令人发指的罪行，即提到伤害神圣建筑的行为；但即使是你，也没有断言这一罪行仅仅等同于说一句无礼的话。你只是满足于为那些犯下此罪的人请求宽恕，这种宽恕本应基于基督徒的广施怜悯而提出，而不是基于所谓的众罪相等。我已经引用过一句圣经经文：\BibleQuote{主的一切行径，都是慈爱和真诚。}因此，他们若是不恨真理，就必寻得怜悯。这种怜悯的赐予，并非因为所有人的过犯仅仅等同于说粗话之人的过犯，而是因为基督的法律要求为悔罪者施以宽恕，不论他们的罪行是多么不人道和邪恶。尊敬的先生，我恳请你不要将这些斯多葛学派的悖论作为行为准则强加给你的儿子\Person{巴拉多克苏斯}，我们都希望他能在虔诚和兴旺中成长，令你满意。因为对他自己来说，还有什么比这更糟糕？对你来说，还有什么比这更危险？难道要让你的天真的孩子沾染上一种错误，这种错误会使我姑且不论弑亲之罪，即使是对父亲的无礼，也仅仅等同于对一个陌生人无意中说出的某句粗话？\par

16. 因此，当你向我们要为你的同胞恳求时，明智的做法是强调基督徒的怜悯心，而不是斯多葛学派那严厉的教条，这些教条不但毫无助益，反而大大阻碍你所承担的事业。因为斯多葛学派宣称，怜悯之心——若我们可能因你的说情或他们的恳求而受感动就必须具备的——是一种可鄙的软弱，他们将其全然从智者的心中驱逐出去；在他们看来，智者的完美在于像铁一般无情且刚硬。因此，你本更应该想到引用你自己的\Person{西塞罗}在赞美\Person{凯撒}时所言的那句话：\Quote{在你所有的美德中，没有一样比你的宽仁更令人敬佩，更优雅动人。}何况在那些跟随那位曾说\BibleQuote{我是道路}并学习祂的话语\BibleQuote{主的一切行径，都是慈爱和真诚}的教会中，这种怜悯之心更当普及！那么，不要害怕我们会试图将无辜者处死，事实上，我们甚至不愿让有罪者遭受其应得的惩罚，因为我们被那与真理相结合的、我们在基督内所爱的怜悯所感动。然而，有些人因害怕痛苦地违逆犯罪者的意愿，便姑息纵容恶行，恶行由此益发壮大；这样的人，还不如那个因不忍听到小孩哭声而不从他手中夺走危险刀具的人仁慈，后者也并未因担心孩子可能因他的软弱而受伤或死亡（这后果将使他悲叹不已）而心软。所以，请把为那些人向我们说情的工作留待恰当的时机——在爱他们这件事上（恕我直言），你不仅没有超越我们，甚至至今仍拒绝跟随我们的脚步；不如在回信中写下是什么影响你，使你避开我们所追随的道路，而我们恳求你在这条路上与我们同行，一同走向那高天之上的家乡，我们欣然得知你在其中享受着极大的喜乐。\par

17. 至于那些生来便是你同胞的人，你确曾说过他们中有些人，虽非全部，是无辜的；但如果你重读我的另一封信，必会发现你并未为他们作出辩护。当你提到你希望离开一个繁荣的故乡时，我回答说，在你的同胞中我们感受到的是荆棘，而非寻得花朵，这时你以为我是在说笑。仿佛在如此巨大的灾祸中，我们还有心情嬉戏。绝对不是。当我们的教堂被火焚毁，废墟上浓烟升腾时，我们岂可能就此事开玩笑？尽管在我看来，你们城中确实无人清白，除了那些不在场的人、受害者、或无力也无权阻止骚乱的人，然而在我的回复中，我仍然区分了罪孽较重者与应受责备较轻者，并指明，那些因害怕冒犯教会之有权势的敌人而行动的人，与那些渴望这些暴行发生的人，情况有别；同样，那些亲手犯下暴行的人与那些唆使他人犯罪的人之间，也有区别。不过，我决意不对唆使者进行审讯，因为若不借助刑讯，或许就无法查明他们，而我们对刑讯深恶痛绝，认为这与我们的目标完全不符。然而，你的朋友们斯多葛学派，他们主张一切过犯皆同等，若由他们作审判官，必会判定所有人同样有罪；如果他们再持守那种不屈不挠的严厉——他们以此将宽仁贬为罪恶——那么他们的判决必定是，不是所有人同样获得宽恕，而是所有人同样遭受惩罚。因此，请让这些哲学家完全退出本案辩护人的位置，但愿我们作为基督徒行事，这样，正如我们所愿，我们可以为基督赢得我们所宽恕的人，而不以对他们有害的放纵姑息他们。愿那行径充满慈爱与真诚的天主乐意以真正的幸福使你富足！\par

\SourceRecord{Translated by J.G. Cunningham.；Nicene and Post-Nicene Fathers, First Series；Vol. 1.；Edited by Philip Schaff.；Buffalo, NY: Christian Literature Publishing Co.,；1887.；<http://www.newadvent.org/fathers/1102104.htm>.}

% structure: p0001 source=h1 target=section reason=page-title

\section{书信 111（公元409年11月）}

\Emph{\Person{奥古斯丁}向\Person{维克托里亚努斯}——他所爱慕的主内尊长、最为渴念的弟兄与同为司铎者——致以主内的问候。}\par

1. 我的心因你的信而充满大忧伤。你希望我在回复中详细讨论一些事情；然而你所叙述的如此灾难，需要的更多是呻吟和泪水，而不是冗长的论述。确实，全世界都遭受着如此可怕的厄运，以至于几乎没有一个地方不在发生并抱怨你所描述的这类事。不久以前，一些弟兄甚至在埃及的旷野中被野蛮人屠杀，而他们正是为了绝对的安全，才选择了那些远离一切纷扰的地方作为隐修院的所在。此外，我想他们在意大利和高卢地区所犯的暴行，你也已知晓；而现在，类似的事件开始从西班牙的许多省份传报给我们，而这些省份长久以来似乎免于这些灾祸。但为什么舍近求远呢？看！在我们自己的\Place{希波}地区——野蛮人尚未触及之处——\OriginalTerm{多纳图派}{Donatist}的神职人员和\OriginalTerm{希尔库姆切利翁}{Circumcelliones}对我们的教会造成的破坏如此之大，相比之下，或许野蛮人的残忍都算轻微的了。因为哪个野蛮人能想出这些人所做的事呢？即把石灰和醋撒进我们神职人员的眼里，此外还残酷地殴打和伤害他们身体的每一部分。他们还常常抢劫和焚烧房屋，掠夺谷仓，倾倒油和酒；并以如此对待该地区所有其他人相威胁，迫使许多人甚至\OriginalTerm{再次受洗}{re-baptized}。就在昨天，我得到消息，说在一个地方有四十八人，因害怕此类事情，被迫再次受洗了。\par

2. 这些事应使我们哭泣，而非惊讶；我们应向天主呼号，并非因我们的功劳，而是按照他的仁慈，愿他从如此大恶中拯救我们。人类还能期待什么呢？因为这些事早在先知书和福音中就被预言了。因此，我们不应如此反复无常，即当我们在诵读中相信这些圣经，却在其应验时抱怨；的确，甚至那些在诵读或听闻圣经中这些事时拒绝相信的人，理应在目睹这话应验时成为信徒；这样，在这巨大的压力下，仿佛在我们主天主的橄榄榨机中，虽然有不信者的抱怨和亵渎的渣滓，但信徒的宣认和祈祷中也稳定地流出纯净的油。对于那些不断责难基督信仰，亵渎地说在基督的教义传遍世界之前人类并未遭受如此严重灾难的人，可以很容易地用福音中主亲口的话来回答：\BibleQuote{那不知道主人的旨意，而做了该受拷打之事的仆人，要少受拷打；但那知道主人的旨意，却偏不准备，也不遵从他的旨意行事的仆人，必然要多受拷打。}\BibleRef{路 12:47} 如果在基督宗教的安排下，世界像那个仆人一样，知道了主的旨意却拒绝遵行，因而遭到多受拷打，这又有什么可惊奇的呢？这些人注意到福音迅速传开，却未注意到许多人悖逆地轻蔑它。然而，天主的温和而虔诚的仆人，既要忍受双倍的现世灾难——因为他们既在恶人手中受苦，也与恶人一同受苦——也有专属于他们的安慰，以及对来世的希望；为此，宗徒说：\BibleQuote{我实在以为现时的苦楚，与将来在我们身上要显示的光荣，是不能较量的。}\BibleRef{罗 8:18}\par

因此，我亲爱的，即使你们遇到那些你们说无法忍受其言语的人，因为他们说，\Quote{如果我们因自己的罪应受这些苦，那么天主的仆人为何也同样被蛮族的刀剑所杀，天主的使女也被掳去呢？}——你们要这样谦卑、真诚而虔敬地回答他们：无论我们如何谨慎地遵守正义之路，服从我们的主，我们岂能比那三位因遵守天主的法律而被投入火窑的人更好？然而，请读读\Person{阿匝黎雅}，那三人中的一个，在火焰中开口所说的话：\BibleQuote{主，我们祖先的天主！你是可赞美的，应受称扬，应受颂扬，直到永远！你的一切作为都是真实的，你的道路都是正义的，你的一切审判都是正确的。你对我们和我们祖先的圣城\Place{耶路撒冷}所行的一切审判都是公义的；由于我们的罪过，你按照真理和正义，使这一切事临于我们身上。我们犯罪作恶，违背了你，在一切事上犯了重罪，没有遵行你的诫命，没有遵守，也没有遵行你为我们的好处，吩咐我们的事。因此，你所加于我们的一切，和你对我们所行的一切，都是按照正义的审判而行的；你将我们交在无法无天的仇人手中，他们是极可憎的叛教者，并交在不义的君王，普世最坏的暴君手中。现在我们再不能开口；你的仆人和敬拜你的人，都遭受了羞辱和耻笑。但为了你的名，求你不要永久抛弃我们，不要废除你的盟约；不要使你的仁爱离开我们，为了你的朋友\Person{亚巴郎}，你的仆人\Person{依撒格}和你的圣者\Person{以色列}，因为你曾向他们应许过，要使他们的后裔繁多，有如天上的星辰，海边的沙粒。主啊！我们在万民中是最微小的，今天在全世界上因我们的罪恶，我们成了最微小的。}我的兄弟，在此你确实可以看到，像他们这样的人——圣者、在磨难中的勇者（然而他们却从其中被解救出来，连火焰本身都不敢烧毁他们）——对自己的罪并不沉默，反而承认了，他们深知因着这些罪，他们理当且公正地被贬抑。\par

4. 不，我们能比达尼尔本人更义吗？天主曾藉先知厄则克耳对提洛的王子论到达尼尔说：\BibleQuote{你比达尼尔更有智慧吗？}\BibleRef{则 28:3} 达尼尔也被列在那三位义人之中，惟独这三位，天主说祂要拯救他们——无疑，这是以他们指代三位代表的义人——天主宣布只拯救诺厄、达尼尔和约伯，他们连带自己的儿女都不能救，只能救自己的灵魂？然而，也请读一读达尼尔的祈祷，看他如何在被掳时，不仅承认自己同胞的罪，也承认自己的罪，并且承认正因这些罪，天主的正义才使他们遭受被掳的惩罚和羞辱。经上这样记载：\BibleQuote{我面向我主天主，恳切祈祷，禁食，身穿苦衣，头上撒灰。我哀求主我的天主，认罪说：\InnerQuote{哎！我主，伟大可畏的天主！你对那些爱你和遵守你诫命的人，必守盟约施恩。我们犯了罪，行了不义，作恶背叛了你，离弃了你的诫命和律例。我们也没有听从你的仆人先知们，他们曾因你的名而向我们的君王、首长和祖先，以及全国人民说话。我主，正义归于你！而满面羞愧归于我们，如今日一样：归于犹大人、耶路撒冷的居民，以及你因他们所犯的罪过而将他们充军各地，或远或近的所有以色列人。我主，满面羞愧归于我们、我们的君王、我们的首长和我们的祖先，因为我们犯罪得罪了你。然而，慈悲和宽恕应归于主我们的天主，因为我们实在都背叛了他，没有听从主我们天主的声音，没有遵行他藉自己的仆人众先知向我们宣示的法律。以色列人全都违背了你的法律，甚至偏离了，不肯听从你的声音；因此，你仆梅瑟的法律上已写明的诅咒和誓词，都倾注在我们身上，因为我们犯罪得罪了他。他对我们和审判我们的判官所说的话，完全应验了，给我们的灾难是十分重大的：天下任何地方所遭受的，没有像耶路撒冷所遭受的。在梅瑟法律上所记载的一切灾祸都降到了我们身上，但我们仍然没有向主我们的天主讨恩，没有离开自己的罪恶，去领悟你的真理。因此，主注意了这灾祸，并把它降在我们身上，因为主我们的天主在他的一切作为中都是正义的，我们还是没有听从他的声音。我主，我们的天主！你曾以强有力的手，将你的人民从埃及地领了出来，而获得了名声，如今天这样；我们却犯了罪，行了恶。我主！求你按照你所有的正义，使你的忿怒和愤恨，远离你的城耶路撒冷，你的圣山！因为，由于我们的罪和我们祖先的过犯，耶路撒冷和你的人民，已在四周的邻人中成了笑柄。如今，我们的天主！请你俯听你仆人的祈祷和哀恳！我主，求你为了你自己的缘故，以你的面容光照你荒凉的圣所。我的天主！求你侧耳俯听，睁眼垂视我们的废墟，和属于你名下的城市！我们所以在你面前哀求，不是仗着我们的正义，而是凭你的大慈悲。我主，求你俯听！我主，求你宽宥！我主，求你垂听并速行，为了你自己的缘故，我的天主！因为这城和你的人民都是属于你名下的。}我还在陈诉、祈祷，承认我罪和我人民以色列的罪过时……}\BibleRef{达 9:3} 请注意，他首先陈述自己的罪，然后才是同胞的罪。他颂扬天主的正义，并为此赞美天主，因为天主即使对祂的圣者们，也并非不公义地，而是因他们的罪过，才以鞭杖惩治他们。因此，如果这是那些因着卓越圣德，以至于连烈焰包围和猛狮之口都毫无伤损之人的言语，那么对于我们这些处于如此低微地位的人，又该说什么合适呢？因为，无论我们看似行有多么仁义，我们都远远不配与他们相比。\par

5. 然而，恐怕有人会认为，那些死于蛮族之手的、你所提及的天主的仆人，理应如同那三位青年从火中获救、达尼尔从狮子坑中获救一样，也获得解救。那么，要让这人明白，那些奇迹之所以施行，是为了使那些将他们判处刑罚的君王相信，他们所崇拜的是真天主。因为天主以祂隐秘的旨意和慈悲，正通过这种方式，为这些君王的救恩作安排。然而，对于那位残酷杀害了玛加伯人的安提约古王，天主却不愿作此安排；反藉他们最光荣的苦难，以更严厉的方式惩罚了这位顽固君王的心。不过，请读一读他们中的一位——那第六位受难者——所说的话：\BibleQuote{在他之后，又解来第六个，他临死时这样说：\InnerQuote{你们不要无故自欺；我们遭遇这些事，是因为我们得罪了自己的天主，因而令人可奇的事临到我们身上。但你不要以为，你既伸手攻打天主和祂的法律，便能逃脱惩罚。}}\BibleRef{加下 7:18} 你看，他们同样在谦卑与真诚中有了智慧，承认自己因罪受到天主的惩戒，经上论到天主说：\BibleQuote{因为主惩戒他所爱的}\BibleRef{箴 3:12}，又\BibleQuote{鞭打他所收纳的每一个儿子}\BibleRef{希 12:6}；因此宗徒也说：\BibleQuote{如果我们先省察自己，就不至于受审判；但我们受审判，是受主的惩戒，免得我们和这世界一同被定罪。}\BibleRef{格前 11:31}\par

6. 你们要忠信地阅读这些事，忠信地宣讲；并在这些考验和磨难中，极尽全力提防抱怨天主，也要教导别人提防。你们告诉我，善良、忠信、圣洁的天主的仆人已被\OriginalTerm{蛮族}{barbarians}的刀剑所杀。但他们是因疾病还是因刀剑脱离肉身，这有何区别呢？主关注的是祂的仆人从这个世界离去时的品德，而非他们去世的具体方式，除非这一点：缠绵的病痛比突然的死亡带来更多痛苦；然而我们读到，这同样漫长而可畏的病痛竟是那位\Person{约伯}的遭遇，而那位不能受欺骗的天主亲自为他的正义作证。\par

7. 最可悲且最应哀恸的，是贞洁而圣洁的妇女被掳；但她们的天主并不在掳掠者的权下，祂也不离弃那些祂确知属于祂自己的被掳者。至于那些圣洁的人，我从圣经中引用了他们受苦与宣认的记录，他们被掳他们的敌人带走，却说了那些话；这些话存于经卷，我们可以亲自阅读，好使我们明白，天主的仆人即使在被掳时，也不会被他们的主所离弃。不仅如此，我们岂知道全能而仁慈的天主会乐意借着这些被掳妇女，甚至在\OriginalTerm{蛮族}{barbarians}的土地上成就何等大能与恩宠的奇事？无论如何，不要停止在天主前为她们叹息代祷；也要在你们的能力和祂的上智安排允许的范围内，寻求为她们做任何能做的事，并在时机许可时，给予她们一切可能的安慰。几年前，一位修女，主教\Person{塞维鲁}的孙女，从\Place{西提法}附近被蛮族掳走，却借着天主奇妙的仁慈被极其荣耀地归还给了她的父母。因为在少女进入掳掠她的蛮族人家时，那家由于主人突然患病而陷入极大痛苦，所有蛮族——如果我没记错的话，是三兄弟或更多——染上了极其危险的疾病。他们的母亲观察到这位少女是奉献于天主的，并相信借着她的祈祷，她的儿子们能从迫在眉睫的死亡危险中被解救出来。她恳求她为他们代祷，许诺若他们得治愈，就将她归还给她的父母。她禁食祈祷，立即蒙了应允；因为，正如结果所显示的，此事早已预定这样发生。因此，他们因这来自天主的意外恩惠恢复了健康，便以钦佩和尊敬对待她，并履行了他们母亲所作的许诺。\par

8. 所以，要为她们向天主祈祷，并恳求祂使她们能说出像我们提到的圣\Person{阿匝黎雅}在祈祷和向天主宣认时与他所说的其他话语一样的话。因为在被掳之地，这些妇女的处境与那三位希伯来青年相似：在那地方，他们不能按规定的方式向主他们的天主献祭；她们既不能将供物带到天主的祭台上，也找不到一位司铎将她们的供物呈献给天主。因此，愿天主赐予她们恩宠，能向祂说出阿匝黎雅在以下祈祷句中所说的话：\BibleQuote{目前我们没有元首，没有先知，没有领袖，没有全燔祭，没有祭祀，没有供物，没有馨香火祭，没有地方可以给你献初果，好蒙受你的恩惠。但愿我们能借着忏悔的心和谦虚的精神蒙你悦纳，如献上公羊和公牛的全燔祭，又如献上了成千上万的肥羊；这样也希望我们今天在你面前所行的祭献，能当作满全了我们对你应尽的义务，因为凡信赖你的，决不会蒙受羞辱。现在我们全心随从你，敬畏你，寻求你的慈颜，望你不要使我们羞惭。但愿你按你的宽容对待我们，按你丰厚的仁慈拯救我们！主啊，求你照你的奇迹拯救我们，使你的名受光荣，愿那些迫害你仆人的人都蒙受羞辱，丧失他们的一切权势，使他们蒙受耻辱，粉碎他们的一切力量，使他们知道你是上主，唯一的天主，在普天下应受光荣。}\par

9. 当天主的仆人用这些话向天主热切祈祷时，祂必站在他们身边，如同祂惯常对待自己的子民一样；要么不许她们贞洁的身体因敌人的淫欲而遭受任何侵犯，或者如果祂允许这事发生，祂也不会将这事上的罪归于她们。因为当灵魂并未因同意这种恶行而受到任何不洁的玷污时，身体也因而免于分担任何罪责；既然受害的一方未因情欲而有任何行径或放任，全部过失归于施暴者，而施加于受害者的一切暴力，并不视为一种下贱的顺从，而是无辜承受的创伤。灵魂中无瑕纯净的价值便是如此：只要灵魂保持完整无玷，即使身体因暴力而被强行制伏，它本身的纯洁依然不受玷污。\par

我恳求您的\OriginalTerm{爱德}{Charity} 满意这封信，虽然相对于我其他的工作它已经很长（尽管仍远不能满足您的愿望），而且写得有些仓促，因为送信人急于动身。若你们专心阅读祂的圣言，主必将赐给你们更丰富的安慰。\par

\SourceRecord{Translated by J.G. Cunningham.；Nicene and Post-Nicene Fathers, First Series；Vol. 1.；Edited by Philip Schaff.；Buffalo, NY: Christian Literature Publishing Co.,；1887.；<http://www.newadvent.org/fathers/1102111.htm>.}

% structure: p0001 source=h1 target=section reason=page-title

\section{书信第115封（公元410年）}

\Emph{致\Person{福图纳图斯}——我在司铎职上的同僚、我至可祝福的阁下、我深为敬爱的弟兄，并致与你同在的弟兄们：\Person{奥古斯丁}在主内向你们致意。}\par

尊驾熟知\Person{法文修斯}，他是\Place{帕拉提安森林}庄园上的佃农。他因惧怕来自该庄园主的一些伤害，逃到\Place{希波}的教堂中寻求庇护，并如逃亡者通常所做的那样，在那里等待能通过我的调解使事情得到解决。然而，日复一日，正如常发生的那样，他变得越来越不警觉，实际上完全丧失了戒备，仿佛对手已罢休了仇恨。一次，他晚餐后从朋友家出来时，突然被伯爵手下的一名军官\Person{弗洛伦提努斯}劫走；此人在这暴力行径中使用了一大队足够的武装人员。当我得知此事，且尚不知他是奉谁的命令、被何人的手劫走时，尽管怀疑自然落在那人身上——即因害怕受其伤害而曾要求教会保护的那人——我立即与指挥海岸卫队的指挥官取得联系。他派出了士兵，却未能找到任何人。但早晨我们得知他在哪所房子里过的夜，并得知他在鸡鸣后便带着押解他的人离开了那里。我也派人前往他据报已被转移的地点：在那里找到了上述军官，但他拒绝让我派去的司铎见他囚禁的人一面。次日我寄去一封信，请求允许他享有皇帝为此类案件所规定的特权，即被传唤出庭受审者，应在市法庭被询问是否愿意在城内受充分监视下度过三十天，以便安排事务或筹措审判费用；我期望在此期间我们或许能友善地解决他的事务。然而，他已在\Person{弗洛伦提努斯}军官的押解下被带得更远了；但我担心，万一他被带到地方法官面前，可能会遭受不公正的对待。因为尽管那位法官的正直广受称赞，被视为不受贿赂，但\Person{法文修斯}的对手乃是一个极为富有之人。为确保金钱在法庭上不占上风，我恳请尊驾，我敬爱的阁下和可敬的弟兄，劳烦将这封附带的信交给我们十分爱戴的尊敬的地方法官，并也把这封信读给他听；因为我认为没有必要将同一案件的陈述写两遍。我深信他会推迟听审，因为我不知道那人是无辜还是有罪。我也深信他不会忽视这一事实：法律已遭违反，因为法文修斯被突然劫走，而没有按照皇帝的法令被带到市法庭，以询问他是否愿意接受三十天延缓期的恩惠；如此，我们便可通过这种方式解决他与对手之间的纠纷。\par

\SourceRecord{Translated by J.G. Cunningham.；Nicene and Post-Nicene Fathers, First Series；Vol. 1.；Edited by Philip Schaff.；Buffalo, NY: Christian Literature Publishing Co.,；1887.；<http://www.newadvent.org/fathers/1102115.htm>.}

% structure: p0001 source=h1 target=section reason=page-title

\section{第116封信（公元410年）}

\EditorialNote{附在前一封信中。}\par

\Emph{致\Person{格内罗苏斯}，我尊贵且卓越的主，我敬爱且深爱的儿子，\Person{奥古斯丁}在主内致意。}\par

虽然我因对您功绩和仁爱应有的爱，而对您治理的赞誉和好评以及您自己卓越的美名总是感到极大的喜乐，但我迄今从未以代求者的身份劳烦阁下，我深爱的主和当受尊敬的儿子，向您请求任何恩惠。然而，一旦阁下从我寄给我的可敬的兄弟兼同工\Person{福尔图纳图斯}的信中得知，在我服务天主的教会的那座城镇里发生了何事，您仁慈的心会立刻觉察到我所处的困境，正是这困境迫使我不得不以这请求打扰您已经非常繁忙的时间。我完全确信，怀着我们因基督之名完全有理由期待的、对我们的那种情感，您在这件事上的行事将不仅是一位正直的，更是一位基督徒长官应有的作为。\par

\SourceRecord{Translated by J.G. Cunningham.；Nicene and Post-Nicene Fathers, First Series；Vol. 1.；Edited by Philip Schaff.；Buffalo, NY: Christian Literature Publishing Co.,；1887.；<http://www.newadvent.org/fathers/1102116.htm>.}

% structure: p0001 source=h1 target=section reason=page-title

\section{第117封信（公元410年）}

\Emph{\Person{迪奥斯科鲁斯}致\Person{奥古斯丁}}\par

对你而言，重视的是内容而非辞藻的修饰，因此本信的任何正式开场白不仅不必要，而且令人生厌。因此，不再赘述，我便直接恳请留意。年迈的\Person{阿利比乌斯}曾屡次应我的请求，答应我他将在你协助下，就我有关\Person{西塞罗}的\WorkTitle{对话录}的少数几个简短问题作出答复；而据说他目前身在\Place{毛里塔尼亚}，故我请求并恳切央求你，能在没有他帮助的情况下屈尊赐予答复——即便令弟在场，无疑这些答复也理当由你来提供。我所求的不是金钱，也不是黄金；纵然你富有这些，我确信你也乐意为了任何正当目的而给予我。但我这一请求，你只需开口说话，便可毫不费力地应允。我本可长篇累牍地，并透过你许多挚友来强求于你；但我深知你的性情，你不喜别人再三恳求，却乐于向众人表示善意，只要所求之事并无不当：而我所求的绝无半点不当。然而无论如何，我恳请你施以援手，因为我即将登程远行。你知道，我何等不愿成为任何人的负担，尤其不愿成为像你这样坦诚之人的负担；但惟有天主知道，我提出此请求是出于何等不可抗拒的急需。因为，我将向你告别，并将自己托付于天主的护佑，动身远行；你知道世人的作风，他们极好挑剔，你也明白，当有人向一个人提出问题却得不到答复时，那人会被视为何等愚昧无知。因此，我恳求你，速速回答我所有的问题。不要让我失望而去。我恳求此事，为的是能见到我的双亲；因为正是为此差事，我才派遣\Person{凯尔多}到你那里，而我现在也只为等他回来而耽搁。家兄\Person{泽诺比乌斯}已被任命为御前记事官，并为我送来了一纸通行凭照及途中的给养。若我不配得到你的答复，那至少请你因担心我因延误而失去这些给养，从而动心回答我的小问题吧。愿至高天主长久保全你的健康，与我们同在！\Signature{帕帕斯}谨向尊驾致以最诚挚的敬意。\par

\SourceRecord{Translated by J.G. Cunningham.；Nicene and Post-Nicene Fathers, First Series；Vol. 1.；Edited by Philip Schaff.；Buffalo, NY: Christian Literature Publishing Co.,；1887.；<http://www.newadvent.org/fathers/1102117.htm>.}

% structure: p0001 source=h1 target=section reason=page-title

\section{第118封书信（公元410年）}

\Signature{奥思定致狄奥斯科若}\par

% structure: p0003 source=h2 target=subsection reason=relative-heading-rank; base=h2

\subsection{第一章}

1. 你突然向我投来无数问题，你必定是想借此将我团团围住，甚至将我彻底掩埋，即便你以为我别无事务、悠闲自在。因为纵使我完全空闲，又如何能为如此急切、而且正如你所写是即将启程的你，解答这许多问题呢？实际上，即使问题容易解答，单凭数量之多，也会使我无法应付。然而这些问题错综复杂，艰深难解，即便数量稀少，且碰上我完全闲暇之时，单是解答它们所需的时间，也会耗尽我的精力，令我精疲力竭。但我很想将你从那让你乐此不疲的探究中强拽出来，让你专注于我所忧心的那些事上，好让你学会不要无益地好奇，或是打消那渴望，即把满足并助长你的好奇心的任务，强加给那些其最大牵挂之一便是遏制并约束过分好奇者的人。因为，与其花时间和精力给你写信，倒不如将其用来努力剪除那些虚幻而危险的欲念（这些欲念伪装在德行与\OriginalTerm{自由学术}{liberal studies}名义下，隐晦而容易欺骗我们，我们必须格外警惕），而不是借我们的服侍，或者更确切地说是奉承，激使它们更加猛烈地确立你卓越心灵所遭受的专制。\par

2. 因为请你告诉我，你所读过的许多\OriginalTerm{对话录}{Dialogues}，如果根本没有帮助你发现并达成你一切行动的目的，它们又有何益处呢？因为通过你的来信，你已相当清楚地表明，你为自己如此热切地从事这些研究——这于你无益，于我却烦扰——所设定的目的是什么。当你在信中竭尽全力说服我回答你提出的问题时，你写下了这些话：\Quote{我本可更详尽地，或藉你诸多挚友来敦促你；但我知道你的性情，你不愿受人恳求，而是乐于向所有人施以仁慈，只要所求并无不当之事；而我所求的，绝无不妥之处。不过，无论如何，我恳请你赐我此恩，因我即将登船启程。}在你信中的这些话里，你对我的评价的确没错，即我乐于向所有人表达善意，只要所求并无不当之处；但我却不认为你所求的毫无不当。因为当我想到一位主教被四面涌来的牧职牵挂所分心、所劳累时，我觉得他如同失聪一般，突然撇下这一切，专心致志地向一个学生解释\Person{西塞罗}的\WorkTitle{对话录}中一些无足轻重的问题，这并不合适。这一点的不合宜，你自己也有所察觉，尽管你因追求学问的热忱而冲昏头脑，绝不肯留意。因为，除了下面这个解释，我还能怎样理解你先是说在这件事上绝无不妥，紧接着又说\Quote{不过，无论如何，我恳请你赐我此恩，因我即将登船启程}呢？这表明，在你看来，至少你的请求并无不妥，但不论其中有何不妥，你还是请求我满足你的要求，因为你即将远航。现在，这个补充的理由\Quote{我即将登船启程}究竟有何力量？难道你是说，除非你处于这种境况，否则我不该帮你做任何可能包含不妥之事？你大概以为，那些不妥能被海水洗净吧。但即便真是如此，至少我那份罪过无法消除，因为我不打算启程。\par

3. 你进一步写道，我知道对你来说，成为任何人的负担是多么痛苦，你郑重声明，唯有天主知道你提出这项申请的迫切需要是何等不可抗拒。当我读到你信中的这一陈述时，我急切地注意，想了解这需要的性质；看，你用这些话把它呈现在我面前：\Quote{你知道人的方式，他们多么倾向于指责，以及当有人向他提问时，他若不能回答，就会被视为无知和愚蠢。} 读到这句话，我燃起了回复你信件的强烈渴望；因为，通过这所显示的病态的心灵软弱，你刺穿了我内心最深处，强行闯入我的挂虑之中，以致我无法拒绝按照天主赐予我的能力来帮助你解脱——不是通过设计解决你困难的方法，而是通过割断你的幸福与它现在所不牢固地依附的那个可怜支撑——即人的意见——之间的联系，并将它固定在一个坚实而不可动摇的基点上。\Person{狄奥斯库若}啊，你难道不记得你喜爱的\Person{佩尔西乌斯}的一句巧妙诗句吗？他在诗中不仅责备你的愚蠢，而且如果你稍有感觉，还会给你幼稚的头脑以应得的纠正，用这些话来抑制你的虚荣：\Quote{在你的眼中，知识本身毫无价值，除非别人知道你拥有知识。} 如我之前所说，你读过那么多对话录，专注于那么多哲学家的讨论——请告诉我，他们中有谁把自己的行为最高目标置于庸人的掌声中，甚至置于善良智慧之人的意见中？但是你——而这更应使你感到羞愧——当你即将启航离开\Place{阿非利加}时，你确实表现出你在这里的学习有了显著的进步，因为你断言，你把讲解\Person{西塞罗}的任务强加给那些已经工作繁重、忙于完全不同性质事务的主教们的唯一理由，就是你担心如果当那些倾向于指责的人向你提问时，你不能回答，你就会被他们视为无知和愚蠢。哦，这理由实在值得占用主教们在别人睡觉时用于学习的时光！\par

4. 在我看来，驱使你日夜努力思想的，无非是想要人们称赞你的勤奋和学识的野心。虽然我一直认为这对于追求真实和正确的人充满危险，但现在因你的情况，我比以往更确信其危险。因为正是由于这同样有害的习惯，你才未能看出，我们可能因什么动机而被促使同意你的请求；因为，由于扭曲的判断，你自己追求了解所提问的事情，动机无非是为了得到人们的赞扬或避免指责，所以你想象我们也会因类似的判断扭曲而受你请求中提出的考虑所影响。但愿当我们向你声明，你对自己写下的这些事情使我们感动，不是要答应你的请求，而是要责备和纠正你时，我们也能让你彻底摆脱人的掌声这种毫无价值且骗人的恩惠的影响！\Quote{这是人的方式，}你说，\Quote{倾向于指责。} 那么呢？\Quote{任何不能回答提问的人，}你说，\Quote{会被视为无知和愚蠢。} 那么，看，我问你一个问题，不是关于\Person{西塞罗}著作中的某事，其意义他的读者或许无法找到，而是关于你自己的信和你自己用词的意义。我的问题是：你为什么不说\Quote{任何不能回答的人将被\Emph{证明}为无知和愚蠢，}而宁愿说，他会被\Emph{视为}无知和愚蠢？为什么呢，若不是因为你自己已经足够明白，不能回答这类问题的人并非实际上，而只是在某些人的意见中，才是无知和愚蠢的？但我警告你，那害怕被这类人的舌头修剪之刃所触及的人是一根无汁的木头，因此不仅被视为无知和愚蠢，而且实际上就是如此，并被证明如此。\par

5. 你或许会说：\Quote{但我并不愚笨，且尤为殷切地努力不成为愚笨之人，因此我连被视为愚笨都不愿意。}此言极是；但我问你，你不愿意的动机是什么？因为当你陈述为何毫不犹豫地用那些你希望得到解答和阐明的问题来加重我们的负担时，你说这就是原因，这就是目的，而且在你看来，这个目的极其必要，以至于你说它具有压倒性的紧迫性——唯恐，若有人拿这些问题质问你而你无言以对，你就会被那些好挑剔的人视为无知且愚笨。现在我问你，这\Emph{对自己名声的嫉妒}是否是你向我们恳求这一切的全部原因，还是由于某种更深层的目的，使你不愿意被认为无知且愚笨？如果这是全部原因，那么，在我看来，正如你所见，这一样东西\Emph{众人的称赞}就是你那强烈热忱所追求的目的，而你也承认，这热忱给我们加上了负担。但是，就\Person{狄奧斯科魯斯}而言，除了\Person{狄奧斯科魯斯}自己不自觉地背负的负担——这负担只有当他试图站起时才开始感受到——除了这个负担之外，对我们来说，还有什么能成为负担呢？这个负担，我情愿相信，并没有如此紧缚于他，以至他无法奋力摆脱。我这样说，不是因为这些问题占据了你的研究，而是因为你研究它们是为了这样的目的。因为到如今，你肯定已感受到，这个目的是琐碎的、不实的、轻如空气。它也容易在灵魂中产生类似危险的肿胀，其下潜伏着腐败的根源，并使心灵之眼变得模糊，以致无法辨识真理的丰富。相信我，我的\Person{狄奧斯科魯斯}，这是真的：这样，我才能在你对真理的真诚渴望中，并在那真理的本然尊严中——你曾因追逐其影子而被引离——得享你的陪伴。如果我如今所用的方法未能使你信服这一点，我不知还能用什么方法。因为你看不见这一点；只要你还将你的喜乐建立在人类赞赏那易碎的基础上，你就绝不可能看见。\par

6. 然而，如果这些行动和你的这份热情所追求的目的并非如此，而是另有其他更深层的原因使你不愿被视为无知和愚笨，那么我请问，那是什么原因。如果是为了消除获取世俗财富、娶得妻子、博得尊荣以及诸如此类事物的障碍——这些事物如急流般奔涌而过，并将落入其中者拖入深渊——那么协助你达成此目的，断然不是我们的责任；相反，我们应当使你远离它。因为，我们并非在禁止你将研究的目标定于不可靠的名声之上，以至于使你仿佛离开\Place{明乔河}的水流而进入\Place{埃里达努斯河}；其实，即使你自己不改变方向，\Place{明乔河}也可能将你带入\Place{埃里达努斯河}。因为当人类赞誉的虚幻不能满足灵魂时——因为它无法为灵魂提供任何真实而实在的滋养——这同样的热切欲望便会迫使心灵转向其他更富有、更丰产的事物；而倘若这事物也同属随时间流逝之物，这就如同一条河流将我们引入另一条河流一样，只要我们在履行职责时所追求的目标是置于那不稳定的事物之上，我们的不幸就永无止歇。因此，我们盼望你将自己最坚定的努力之归宿，以及你所有良善可敬之行动的完全安息之所，建立在某种坚固而永恒的美善之上。或许你的意图是：如果你借着顺风般的名声之气息，甚或张满风帆追逐其忽来忽去的阵风，而获得了我所说的那种尘世幸福，然后使之服务于获得另一种美善——那确实、真实且令人满足的美善吗？但在我看来，这似乎不可能——真理本身也不容许这种假设——即通过如此迂回的途径达到它，何况它近在咫尺；或是在它白白赐予时，却要付出这样的代价才能得到。\par

7. 或许你认为，我们应当善用人类赞誉本身，将其作为工具，使他人易于听取有关美善与有益之事的劝告；而且，你或许担心，如果人们视你为无知和愚笨，当你劝勉人行善，或责备作恶者的邪恶与恶意时，他们会认为你不配获得他们的认真或耐心倾听。如果在提出这些问题时，你心中所设想的是这正义且有益的目的，那么你的信函并未优先提及这一点，以之作为可能打动我们乐意应允你所求的理由，或者，若拒绝你的请求，也当以其他可能阻碍我们的原因为依据，而非因为我们耻于接受那种迎合、甚或纵容你虚荣心企图的角色。你这样写，确实使我们蒙受了委屈。因为，请你想想，你若能从我们这里以远远更为确凿且耗时更少的方式，获得那些真理的原则，该有多么好、多么有益啊！凭借这些原则，你便能自行驳斥一切虚假，并因此不致抱持一种如此虚妄可鄙的见解——即以为若你以一颗骄傲多于审慎的心，研习了过去世代许多作家那些陈腐的错误，你便是博学而聪慧的。不过，我猜想你现在已不抱持这种见解了，因为我对\Person{狄奧斯科魯斯}说的这些如此明显的真理，必定没有白说；既已明了此点，我便继续写这封信了。\par

% structure: p0011 source=h2 target=subsection reason=relative-heading-rank; base=h2

\subsection{第二章}

8. 所以，既然你认为一个人不会仅仅因为不知道这些事就被视为不学无术和愚笨，而只有当他不知道真理本身时才会如此；而且，既然如此，任何已经或可能写这些题材的人的意见，要么是真的，因此已为你所持守，要么是假的，因此你可以安然于不去知道它们，也不必因担心自己若不如此就会成为不学无术和愚笨之人，而徒然地渴求了解他人的各式意见——我既然说情况如此，那么让我们现在，如果你愿意，考虑一下：如果如你所说那些好指摘的人发现你对这些事无知，并因此认为你（尽管是错误地）是个不学无术和愚笨的人，他们的这个错误对你是否应该如此重要，以至于你向主教们请教这些事就并非不体面？我提出这一点，是基于我们现在相信你寻求这种教导，是为了藉此帮助向人推荐真理，并挽回那些人——他们若认为你对\Person{西塞罗}的那些著作不学无术和愚笨，就会视你为他们觉得不屑于从其领受任何有用或有益教导的人。相信我，你弄错了。\par

9. 因为，首先，我完全没有看出，在那些你如此害怕被人视为缺乏教养和敏锐的国家里，会有任何人就这些事向你提出一个问题。无论在你前来学习这些事的这个国家，还是在\Place{罗马}，你都通过经验知道这些事是多么不受重视，因此它们既无人讲授也无人学习；而在整个\Place{非洲}，你根本不会受到任何此类提问者的烦扰，以至于你找不到任何愿为你的问题烦心的人，并且由于缺乏这样的人，你不得不把你的问题送给主教们去解决：好像这些主教，尽管在青年时代，在同样那股——我还不如说是错误——将你卷走的热忱影响下，他们也曾费心学习这些事，视之为至关重要，如今虽已白发苍苍并担负着主教职分的重任，却仍让它们留在记忆中；或者好像，即使他们想要将这些事保留在记忆中，更重大更严肃的忧虑也不会违背他们的意愿把它们从心里驱逐；或者好像，即使有些事因长久习惯的力量而萦绕于回忆中，他们也不愿宁可将其完全遗忘，而不是回答那些毫无意义的问题——在这样的时候，即使在学校里和修辞学家的讲堂中的相对闲暇之中，他们似乎都已如此丧失了声音和活力，以至于为了就这些事获得教导，人们认为必须从\Place{迦太基}派人到\Place{希波}——一个所有这些事都如此罕见、如此完全陌生的地方，如果我为回答你的问题而费心写信，想查阅任何章节以发现需要解释的词语前后文中的思想顺序，我将根本找不到一份\Person{西塞罗}著作的手稿。然而，\Place{迦太基}的这些修辞学教师在这件事上未能满足你，他们不仅不受责备，反而，在我看来，受到称赞，如果如我所想，他们没有忘记这些竞技的舞台过去通常不是\Place{罗马广场}，而是希腊的体育馆。但当你将心思转向这些体育馆，却发现它们在这些事上甚至也是贫乏而冰冷的，希波基督徒的教会就浮现在你面前，作为你放下忧虑的地方，因为那位目前占据此主教座的主教曾一度收费教导少年这些事。可是，一方面，我不希望你还是个少年，另一方面，我现在处理这类儿戏之事，无论收费还是作为人情，都不适宜。因此，情况既是如此——也就是说，\Place{罗马}和\Place{迦太基}这两座伟大城市，拉丁文学的活跃中心，既不用你所提的那些问题考验你的耐心，也不愿耐心听你提出这些问题——我无比惊讶，像你这样有良好判断力的年轻人竟然担心在\Place{希腊}和\Place{东方}诸城市里会受到任何提问者就这些问题的折磨。你在\Place{非洲}更可能听到寒鸦叫，在那些地方则听不到这种谈话。\par

10. 然而，进一步假设我错了，或许有人会提出这样的问题——由于在那些地区特别荒谬，这种现象更加不受欢迎——难道你不更担心那些身为希腊人、发现你定居在\Place{希腊}、并以希腊语为母语的人，可能会在你所阅读的哲学家原著中，问一些\Person{西塞罗}可能未收入其论著中的问题吗？如果发生这种情况，你将如何作答？你会说你宁愿从拉丁文著作而非希腊文著作中学习这些东西吗？这样的回答首先是对\Place{希腊}的侮辱；你知道该国的人对此会何等愤恨。其次，他们既已受到伤害而愤怒，你将多么轻易地发现你急于避免的事情：一方面，他们会认为你愚蠢，因为你宁愿在拉丁文对话录中学习希腊哲学家的见解，或者更确切地说，学习他们哲学中某些孤立零散的学说，而不愿在希腊原著中研究他们见解的完整连贯的体系；另一方面，他们会认为你没文化，因为尽管你对用自己母语写成的许多东西一无所知，却费尽心力从外文著作中搜集一些片断，结果并不成功。或者你也许会回答说，你并非轻视论及这些主题的希腊文著作，而是先致力于拉丁文著作的学习，如今已精通这些，正开始探究希腊学问？如果这不会令你脸红，承认你身为希腊人，却在少年时学习了拉丁文，如今像个外国人似的渴望研读希腊文献，那么你肯定也不会羞于承认，在拉丁文学领域，你对某些事物无知，而你只要想一想，就会发现许多精通拉丁学问的人也同样无知——你只要想想，尽管你生活在\Place{迦太基}如此众多的博学之士中间，你却向我保证，是迫于无奈才把这负担加在我身上。\par

11. 最后，假设你被问及所有你向我提出的那些问题，并且能够全部回答。看哪！你现在被誉为博学而敏锐之人；看哪！现在这股希腊赞誉的微弱气息将你捧上天。如今愿你记住你的责任，以及你贪求这些赞誉的目的，即：通过这些琐事已轻易赢得敬佩，并如今最亲切热切地倾听你言谈的人们，你可以传授某些真正重要且有益的教训；我想知道你是否拥有，并能正确地传授给他人那真正最重要最有益的教训。因为假如你为了预备人们的耳朵去接受必要之事而学了许多不必要之事，却发现你并不拥有这些必要之事——你正是用这些不必要之事为接受它们预备了道路——这是荒谬的；假如你忙于学习那些可以赢得人们注意的东西，却拒绝学习在他们注意被吸引后可以倾注于他们心中的东西，这也是荒谬的。但如果你回答说，你已经学了这些，并说那最必要的真理是基督宗教教义，我知道你将其置于万物之上，唯独将其作为你永久救恩的希望，那么，若要劝人倾听这教义，肯定不需要你通晓\WorkTitle{西塞罗的对话录}，以及搜集他人种种可怜而分裂的意见。让你的品格和生活方式引起那些要从你领受此类教导之人的注意。我不愿你借由先教导日后必须弃绝的东西，来为教导真理开路。\par

12. 因为如果了解他人那些不调和且相互矛盾的意见，对于那些想要为基督的真理打开入口、推翻错谬反对的人有任何用处的话，那仅仅在于，阻止真理的攻击者得以自由地只专注于驳斥你的主张，却小心隐藏自己的主张。因为真理的知识本身足以察验并推翻一切错谬，即使是以往未曾听闻的错谬，只要它们被提出来。然而，如果为了不仅摧毁公开的错谬，也根除那些潜藏于暗处的错谬，你有必要熟悉他人提出的那些错误意见，那么我恳请你，要让眼睛和耳朵保持警醒——仔细察看和聆听，看看我们的攻击者中，是否有谁从\Person{阿那克西米尼}和\Person{阿那克萨哥拉}那里提出任何论证；因为尽管斯多亚派和伊壁鸠鲁派的哲学较近，且广为传授，但连它们的灰烬也不够温热，以致无法从中引出一点火星来反对基督信仰。如今在争论战场上轰鸣的喧嚣，来自无数宗派的小团体和派系，其中有些很容易击败，有些则胆敢顽强抵抗——比如这里拥护\Person{多纳图斯}、\Person{马克西米安}和\Person{摩尼}的人，或是你前往之地盛行的亚略派、欧诺弥派、马西顿派、卡塔弗里基派等无法无天的成群害虫。如果你因畏缩而不愿去了解所有这些宗派的错谬，那么我们在为基督宗教辩护时，有何理由去探究\Person{阿那克西米尼}的学说，并以无益的好奇心重新激起已沉寂数世纪的争论，而那时甚至某些自诩享有基督之名的异端者，如马西昂派和撒伯里乌派等等许多人的盘问与论证，都已经消声匿迹了？尽管如此，如果如我所说，有必要预先了解一些与真理作战的意见，并充分熟悉它们，那么我们有责任在这样的研究中优先考虑那些自称基督徒的异端者，而不是\Person{阿那克萨哥拉}和\Person{德谟克利特}。\par

% structure: p0017 source=h2 target=subsection reason=relative-heading-rank; base=h2

\subsection{第三章}

13. 再者，无论谁向你提出你曾向我们询问的那些问题，他都应当明白，你是在更深奥的学识和更大智慧的指引下，才对这类事情一无所知的。因为如果\Person{忒弥斯托克利}在宴会上拒绝弹奏里拉琴而被视为未受充分教育时，他将此事看作小事，并且当他在承认对此技艺一无所知之后，有人问他：\Quote{那么，你到底知道什么呢？}，他回答说：\Quote{使一个小共和国变伟大的技艺}——那么，当你有能力回答任何可能问\Quote{那么，你到底知道什么呢？}的人说\Quote{那使人即使没有这类知识也能享真福的秘诀}时，你又何必犹豫承认自己在这些琐事上的无知呢？如果你尚未拥有这秘诀，那么你在探究那些其他事物时的盲目偏执，就如同一个患有危险身体疾病的人，不是求医问药，而是急切地追求珍馐美味和华美衣饰一样。因为这一目标无论如何不应以任何借口推迟，特别是在我们这时代，任何其他知识都不应在学习中占有优先地位。现在请看看，如果你渴望这知识，将多么容易得到它。凡探求如何能获得真福生活的人，无疑正是探究这个问题：至高的善在哪里？换言之，人的至高善存在于何处——不是按照人们歪曲而草率的意见，而是按照确定不移的真理？这至高善的所在，除了在身体内、或在心神内、或在天主内、或在其中两者内、或在三者结合之内，没有人能发现。那么，如果你已经认识到，至高善或其任何部分都不在身体内，剩下的选择就是：它在心神内，或在天主内，或在两者结合之内。而如果你现在也已认识到，在这点上身体所是的同样适用于心神，那么除了天主自己作为人的至高善所在的那一位以外，还剩什么呢？——并不是说没有其他的善，而是说至高善是其他一切善所关联的那个善。因为每个人享有真福，是当他享有那为它之故他渴望拥有一切别的事物的事物之时，因为那事物是为了它自身而被爱，而不是为了别的事物。而说至高善就在那里，是因为在这一点上找不到任何至高善能趋向或与之关联的事物。在这至高善中，是欲望的安息之所；在这至高善中，是确实的享有；在这至高善中，是一个道德上完善的意志最宁静的满足。\par

14. 请给我找一个人，他立即看出身体不是心神的善，心神反倒是身体的善：对于这样的人，我们当然无需探问我们所谈及的那至高善或其任何部分是否在身体内。因为心神优于身体，这是一个否认它就是极端愚蠢的真理。同样荒谬的是否认那带来真福生活或其任何部分的事物，优于那接受这恩赐的事物。因此，心神并不从身体接受至高善或至高善的任何部分。看不清这一点的人被肉欲的甘甜弄瞎了眼，而他们没有看出这肉欲的甘甜乃是健康不完满的结果。然而，身体完美的健康将是全人永生的完成。因为天主赋予灵魂一种如此强大的本性，以至于从许诺给末世圣人们的那圆满的喜悦中，有一部分满溢到我们本性的低级部分，即身体——并不是那属于享受和领悟的部分所特有的真福，而是健康的完满，即不朽的活力。正如我所说，那些看不清这一点的人们，在令人不满的辩论中彼此争论，各自坚持自己可能喜欢的观点，但全都将人的至高善置于身体内，从而煽动起一群群混乱的肉欲之徒，在这类人中，\Person{伊壁鸠鲁派}在未受教育的群众中享有卓越的声望。\par

15. 此外，请给我指出这样一个人：他能立刻看到，当心灵快乐时，其快乐并非源于属于自身的善，否则心灵便永不会不快乐。对于这样一个人，我们自然不必再去追问，那至高的、可说是赐予真福的善，或其任何部分，是否在心灵之中。因为当心灵因自身而得意，仿佛在属于自己的善中快乐时，这便是骄傲。但是，当心灵察觉到自身是可变的——即使没有其他证据，单从心灵可由愚钝变为智慧这一事实便可得知——并领悟到智慧是不变的，它就必然同时领悟到，智慧超越其自身的本性，并且在智慧的交流与光照中，要比在自身中找到更丰盈、更持久的喜乐。于是，心灵放弃并平息了炫耀与自负，努力依恋天主，并藉着那不变者而得到更新与再造；此时它认识到，天主不仅是它通过身体感官或理智官能所接触到的一切万物的各类形式的创造者，而且也是在任何形式被接受之前，那接受形式之能力本身的创造者，因为无形者正定义为能接受形式者。因此，心灵越不依恋那存在完美的天主，便越感到自身的不稳定性；它同时也辨明，天主的存在的完美是圆满的，因为祂是不变的，因而既无所得亦无所失；而对于心灵自身，凡使它获得更完美依恋天主之能力的改变，都是有益的，凡使它失去此种能力的改变，则是有害的；并且，一切的损失都趋向毁灭；虽然某物最终是否被彻底毁灭并不显而易见，但人人皆知，损失所带来的毁灭至少使那物不再成其原本所是。由此心灵推断，事物之所以遭受损失或可能遭受损失，唯一的原因是它们是从虚无中受造的；因此，它们的存在与恒定性，以及它们虽不完美却能在复杂整体中各安其位的秩序，全依赖于那存在完美、且能从虚无中不仅创造某物、更能创造伟大之物的创造者的善与全能；并且第一个罪，\Emph{即}第一个自愿的损失，便是因自己的权力而欢喜：因为若因天主的权力而欢喜，那无疑是更大的喜乐之源，而它却因较小的而欢喜。有些人没有感知到这一点，只着眼于人类心灵的能力及其在言行方面成就的大美，他们虽羞于将人的\OriginalTerm{至高善}{supreme good}置于身体中，却将其置于心灵内，无疑赋予了它一个比纯真理性所赋予的更低的范围。在持这些观点的希腊哲学家中，无论就信徒数目还是论证的精妙而言，斯多噶派都占据首要地位；然而，由于他们主张自然界中一切都是物质的，他们成功地使心灵脱离肉欲的对象，而非物质的对象。\par

16. 同样，在那些宣称我们\OriginalTerm{至高且唯一的善}{supreme and only good}是享受天主——我们自身及万物皆由祂所造——的人中，最杰出的当属柏拉图学派。他们并非没有理由地认为，驳斥斯多噶派和伊壁鸠鲁派——尤其是后者，且几乎是专门针对后者——乃是他们的职责。学园派与柏拉图学派本是一体，这从学派间不间断的传承脉络便可清楚看出。因为若问 \Person{阿尔克西拉乌斯} 的先驱是谁——他是第一位不宣说自己学说、而专心驳斥斯多噶派和伊壁鸠鲁派的人——你会发现那是 \Person{波莱蒙}；再问 \Person{波莱蒙} 的先驱，那是 \Person{色诺克拉底}；而 \Person{色诺克拉底} 是 \Person{柏拉图} 的门徒，并被 \Person{柏拉图} 任命为学园的继承人。因此，对于这个关于至高善的问题，若我们搁置各种对立观点的代表人物，仅考虑问题本身，你立即会发现两种错误针锋相对：一种主张身体，另一种主张心灵是人类至高善的所在。你也会发现，真正受启示的理性——藉此理性我们认识到天主是我们的至高善——与这两种错误都相对立，但它并不直接传授真理的知识，除非首先使人忘却虚假。若你联系各派观点的拥护者来考虑这个问题，你会发现伊壁鸠鲁派与斯多噶派激烈地相互论争，而柏拉图学派则试图仲裁他们之间的争论，隐藏自己所持有的真理，只致力于证明并推翻其他学派执着于错误的自负。\par

17. 然而，\Person{柏拉图}主义者并没有能力像其他人支持自己的谬误那样，有效地支持由真理所光照的理性一方。因为那时，他们都尚未得到属神的谦卑的榜样；这榜样在时期一满，便由我主\Person{耶稣基督}提供了——在这独一无二的榜样面前，即便在最顽固傲慢之人的心中，一切骄傲也都会屈服、破碎并消亡。因此，\Person{柏拉图}主义者无法凭他们的权威引导那被贪爱世俗之物所蒙蔽的大众，去信仰不可见的事物——尽管他们看到大众特别被\Person{伊壁鸠鲁}派的论证所动摇，不仅畅饮他们本性所趋向的身体欢愉之杯，甚至为这些欢愉辩护，声称它们构成人的至善；此外，他们也看到那些因对美德的赞美而致力于节制这些欢愉的人们，更容易将快乐视为真住所在于灵魂之中，那些良善的行为，他们多少还能有些看法，便是由此而生——同时，他们也看到，如果他们试图在人心中引入某种属神、至高不变者的观念——这观念不能由任何身体感官触及，只能由理性领悟，但理性在本性上却又超越心灵本身；并且教导说这便是天主，祂被置于人灵面前，当人灵从一切属人欲望的污秽中得到净化后便可享有祂，只有在祂内，所有对幸福的渴求才能得到安息，也只有在祂内，我们才应找到一切美善的完满——那么，人们将无法理解他们，而会更乐意把胜利的棕榈枝授予他们的对手，无论是\Person{伊壁鸠鲁}派还是\Person{斯多葛}派；其结果将是人类最为灾难性的，即真实而有益的教义，将因未受教育大众的轻蔑而受到玷污。关于伦理学的问题就谈这些。\par

18. 至于物理学，如果\Person{柏拉图}主义者教导说，一切本性的初始原因是非物质的智慧，而另一方面，与之竞争的哲学派别从未超越物质事物，他们将万物的起始归因于原子，或归因于四元素，并认为火在万物的构造中具有特殊的力量——那么，当绝大多数不思考的人们被物质事物所束缚，且根本无法理解一种非物质的德能如何能形成宇宙时，谁会看不出人们会对哪种意见作出有利的裁决呢？\par

19. 辩证法问题这一分支仍有待讨论；因为如你所知，一切追求智慧的议题都被分类为三个部分——伦理学、物理学和辩证法。因此，当\Person{伊壁鸠鲁}派说感官从不欺骗，而\Person{斯多葛}派虽然承认它们有时会出错，但双方都把感官当作理解真理的标准时，在这两派都反对他们的情况下，谁还会去听\Person{柏拉图}主义者的呢？如果他们毫不迟疑且毫无保留地断言：不仅存在着某种无法被触觉、嗅觉、味觉、听觉或视觉所辨别，也无法借由感官所认识的事物形成的任何形象来构想的东西，而且唯有这东西才真正存在，也唯有它才能被领悟，因为它独一不变且永恒，但只能由理性来领悟——而这理性正是我们赖以发现真理（就其能被我们发现而言）的唯一官能——那么，谁还会将他们视为配得上被当作人来看待，更不要说当作智士了呢？\par

20. 因此，看到\Person{柏拉图}主义者所持的见解，他们无法传给那些受肉身束缚的人；也看到他们在普通人中间并无如此权威，以致能说服人们接受那些他们本应相信的，直到心灵被训练到能理解这些事的程度——于是他们选择把自己的见解隐藏起来，而只满足于反驳那些人，那些人虽然宣称真理的发现是通过身体感官，却夸口说他们已经找到了真理。确实，我们有何必要去探究他们学说的本质呢？我们知道那并不是属神的，也不具有任何属神的权威。但有一点却值得我们注意：尽管\Person{西塞罗}在许多方面非常清楚地证明，\Person{柏拉图}将至善、事物的原因，以及理性过程的确实性，都置于智慧之中——不是人的智慧，而是属神的智慧，从这智慧中以某种方式衍生出人类智慧之光——这智慧是完全不变的，是始终自相一致的真理；并且我们也从\Person{西塞罗}那里得知，\Person{柏拉图}的追随者们努力去推翻那些被称为\Person{伊壁鸠鲁}派和\Person{斯多葛}派的哲学家，这些哲学家将至善、事物的原因和理性过程的确实性，要么置于身体的本性中，要么置于心灵的本性中——这场争论持续经历了数个世纪，以至在基督宗教肇始时期，当那关于不可见与永恒之事的信仰，正借着可见的奇迹，带着救恩的力量，被宣讲给那些除了物质事物之外既不能看见也不能想象的众人时，我们发现这些同样的\Person{伊壁鸠鲁}派和\Person{斯多葛}派，在\BibleBook{宗徒}{大事录}中竟然反对圣宗徒\Person{保禄}，而\Person{保禄}当时正开始在外邦人中播撒这信仰的种子。\par

21. 我觉得通过此事已充分证明，外邦人在伦理学、物理学和寻求真理的方式上所犯的许多各式各样的错误——这些错误显著地体现在这两个哲学学派中——一直延续到基督宗教时代，尽管有学问的人以最激烈的方式抨击它们，并运用了非凡的技巧和大量的辛劳来推翻它们。然而，我们看到这些错误在我们的时代已经完全销声匿迹，以至于现在在我们的修辞学校里，几乎从不提及他们的观点是什么的问题；这些争论甚至在以爱好讨论著称的希腊体育学校里也已被如此彻底地根除或压制，以至于现在每当任何错误的学派抬起头来反对真理，\Emph{i.e.}反对基督的教会时，它若不以基督徒的名号为掩护，就不敢进入竞技场。由此可见，柏拉图学派的哲学家们感到有必要——在他们改变了自己观点中受到基督教义谴责的少数内容之后——怀着虔诚的敬意顺服于基督，那位唯一不可战胜的君王，并领会那\OriginalTerm{降生成人的天主圣言}{Incarnate Word of God}，因着祂的命令，他们甚至曾害怕发表的真理立即就被相信了。\par

22. 我的\Person{狄奥斯库若斯}啊，我渴望你以毫无保留的虔敬顺服于祂，并且愿你除了祂所预备的道路之外，不为自己预备任何其他把握并持守真理的道路；祂作为天主，看到了我们步履的软弱。在这条道路上，第一部分是\OriginalTerm{谦卑}{humility}；第二部分是谦卑；第三部分是谦卑：无论你多少次询问指引，我都会继续这样重复，这并非因为不能给出其他训导，而是因为，除非谦卑在我们所做的每一件善行之前、之中、之后，同时作为我们眼前的目标、我们所依附的支撑、以及约束我们的告诫者，否则骄傲就会从我们的手中夺走我们正为之沾沾自喜的任何善工。所有其他恶习是在我们作恶时应当提防的；但是骄傲即使在行善时也当惧怕，免得那些以值得称赞的方式所做的事被对赞扬的渴望本身所败坏。因此，正如那位最杰出的演说家，当被问及在他看来演讲艺术中首要遵守的是什么时，据说回答说，表达；当被问及第二点是什么时，又回答说，表达；当被问及第三点是什么时，仍然只给出了这个回答，表达；同样地，如果你问我，无论你重复多少次这个问题，基督宗教的训导是什么，我总会倾向于回答，并且只回答，\Quote{谦卑}，尽管也许有需要时，我也会被迫谈论其他事情。\par

% structure: p0028 source=h2 target=subsection reason=relative-heading-rank; base=h2

\subsection{第四章}

23. 对这种最有益健康的谦卑——我们的主耶稣基督在其中是我们的老师，祂屈尊降卑，为要在谦卑上教导我们——我说，对这种谦卑，最可怕的敌人是某种最无见识的知识（如果我可以这样称呼它的话），这种知识使我们以知道\Person{阿那克西美尼}、\Person{阿那克萨哥拉}、\Person{毕达哥拉斯}、\Person{德谟克利特}以及其他同类人的观点为满足，自以为这样就成了有学问的人和学者，尽管这些造诣与真正的学问和博学相去甚远。因为一个人若明白了天主并不伸展或弥漫于空间（无论是有限的还是无限的），以致在一部分较大而在另一部分较小，而是完全地临在于每一处，正如真理一样——任何心智健全的人都不会说真理一部分在这里，一部分在那里——而真理就是天主自己——这样的人就不会被任何哲学家的观点所动摇，这些人（比如\Person{阿那克西美尼}）相信我们周围的无限空气就是真天主。这样的人即使不知道他们所说的形是指什么形体，又有什么关系呢？因为他们说的是一种在各方面受限制的形体。对他来说，\Person{西塞罗}反对\Person{阿那克西美尼}（他曾说天主是一种物质存在——因为空气是物质）时说天主必须有形式和美，这究竟是仅仅作为一名学院派哲学家，并且只是为反驳\Person{阿那克西美尼}而提出的，还是他自己认识到了真理具有非物质的形式和美——心智本身正是由此被塑造，我们也以此判断智者的一切行为都是美的——因而断言天主必定具有最完美的美，这不仅仅是为了反驳一个对手，而是基于深刻的洞察力，认识到没有什么比真理本身更美，真理只能由理智认识，并且是不变的？这两者对他又有何区别呢？再者，至于\Person{阿那克西美尼}的观点，他主张空气是被生成的，同时又相信它是天主，这一点丝毫不能动摇这样的人，他明白：既然空气肯定不是天主，那么空气被生成的方式（即由某种原因所产生）与祂被生出的方式（除了通过天主的启示，无人能理解）之间毫无相似之处——祂就是\OriginalTerm{天主的圣言}{the Word of God}，与天主同在的天主。此外，谁看不到，即使在物质事物上，他肯定空气被生成，同时又是天主，这是最愚蠢的说法，因为他拒绝将天主之名给予那产生空气者——因为空气不可能不由任何能力而被生成？再者，他说空气总是在运动，这个说法并不能作为空气是天主的证据来困扰这样的人，他知道一切物体的运动都低于灵魂的运动，但即使灵魂的运动，与那至高无上、不变的\OriginalTerm{智慧}{Wisdom}相比，也是无限缓慢的。\par

24. 同样，如果\Person{安那克萨哥拉}或任何其他人断言心灵是本质的真理和智慧，我有什么必要就一个词与人争论呢？因为很明显，心灵赋予万物的秩序和模式以存在，并且它可以恰当地被称为无限的，并非就其在空间中的广延而言，而是就它的能力而言，其能力范围超越一切人类思想。我也不会[就他的断言]争论说这种本质的智慧是无形式的；因为这是物质事物的属性，凡是无限的物体也都是无形式的。然而，\Person{西塞罗}出于反驳这类意见的愿望，据我猜想，他在与那些不相信任何非物质东西的对手争辩时，否认任何东西能够附加于无限者，因为在物质事物中，任何东西附加的部分都必须有一个边界。因此他说，\Person{安那克萨哥拉}\Quote{没有看到，与感觉相连并与它相连的运动}（\Emph{即}，在不可中断的联系中与之相连）\Quote{在无限者中是不可能的}（也就是在无限的本体中），仿佛在论述物质本体，对于物质本体，除了在它们的边界处，什么也不能连接。此外，在接下来的话中——\Quote{并且，整个自然系统在被触动时没有感觉的那种感觉是不可能的}——\Person{西塞罗}说话的方式仿佛\Person{安那克萨哥拉}说过，心灵——他把安排和塑造万物的能力归属于它——具有像灵魂通过身体所具有的那种感觉。因为很明显，当灵魂通过身体感觉到任何东西时，整个灵魂都有感觉；因为被感觉领悟的任何东西都不会向整个灵魂隐藏。现在，\Person{西塞罗}说整个自然系统必然意识到每一种感觉，其意图是，他仿佛要从那位哲学家那里夺走他断言为无限的心灵。因为如果自然是无限的，整个自然如何经历感觉呢？身体的感觉从某一点开始，并不遍及任何本体的全部，除非它能在其中到达一个尽头；但是，当然，对于无限者不能这么说。然而，\Person{安那克萨哥拉}并未说过任何关于身体感觉的话。此外，当我们谈论非物质的东西时，\Quote{整个}这个词的用法不同，因为它被理解为在空间中没有边界，所以它可以被称为整体的，同时又是无限的——前者是因为它的完整性，后者是因为它不受空间边界的限制。\par

25. \Quote{此外，}\Person{西塞罗}说，\Quote{如果他将断言心灵本身，可以说，是某种动物，那么必有某种内在的东西，它由此得到\InnerQuote{动物}的名称，}——这样，按照\Person{安那克萨哥拉}，心灵是一种物体，并且内部具有赋予生命的东西，因此被称为\Quote{动物}。请注意，他使用的是我们习惯于用于有形事物的语言——动物通常是可见的实体——据我猜想，他使自己适应那些与他争论的人迟钝的理解力；然而他说的某件事，如果他们能醒悟过来领悟它，可能足以教导他们，凡是向我们心灵呈现为活的身体的东西，都必须被理解为并非灵魂本身，而是一个具有灵魂的动物。因为他已经说过，\Quote{必须有某种内在的东西，它由此得到动物的名称，}他又补充说，\Quote{但有什么比心灵更内在呢？}因此，心灵不可能有任何内在的灵魂，因拥有它而成为动物；因为它本身就是最内在的。那么，如果它是动物，就让它有某种外在的身体，相对于这身体它可以在其内；因为这就是他说\Quote{因此它被一个外在的身体包围}的意思，仿佛\Person{安那克萨哥拉}说过，心灵不能不是属于某个动物的。然而\Person{安那克萨哥拉}认为那本质的至高智慧是心灵，尽管它不是任何特定生物的特有属性，可以这么说，因为真理同样接近所有能够享有它的灵魂。因此，注意他如何巧妙地总结论证：\Quote{既然这不是\Person{安那克萨哥拉}的意见}（\Emph{即}，鉴于他不认为他称为\Emph{天主}的心灵被一个外在的身体包围，因它与身体的关系而能成为动物），\Quote{我们必须说，纯粹而简单的心灵，不附加任何东西}（\Emph{即}任何身体）\Quote{使它能够行使感觉，似乎超出了我们理智的范围和构想。}\par

26. 没有什么比这更确定：即这超越了\Person{斯多葛}派和\Person{伊壁鸠鲁}派理智的范围和构想，他们不能思考任何非物质的东西。但是，所谓\Quote{我们的}理智，他是指\Quote{人类的}理智；并且他很恰当地没有说\Quote{它超越了我们的理智，}而是说\Quote{它似乎超越了。}因为他们的意见是，这超越了所有人的理解，因此他们认为任何这类东西都不可能存在。但是有些人，他们的理智领悟到，只要人被赋予这种能力，确实存在着纯粹而单一的智慧和真理，它不是任何特定生物的特有属性，而是将智慧和真理赋予所有同样能受其影响的灵魂。如果\Person{安那克萨哥拉}察觉到这至高智慧的存在，并领悟到它是\Emph{天主}，称之为心灵，那么，我们成为有学问和智慧的人，并非仅仅借着这位哲学家的名字——由于他在古代学术的远古地位，所有文学新手（采用军事用语）都喜欢炫耀自己认识他——也并非借着我们具有使他认识这真理的知识。因为真理对我来说应该是宝贵的，并非仅仅因为\Person{安那克萨哥拉}并非不知道它，而是因为，即使所有这些哲学家都不知道它，它仍然是真理。\par

27. 因此，若我们因那或许领悟了真理之人的知识——凭借这知识，我们仿佛获得了学识的表象——甚或因对真理本身的实在拥有——借此我们获得学识上的真实收获——而心高气傲，尚且不合宜，那么，那些陷于错谬中的人的名字和学说，又岂能对基督信仰的学习和彰显隐晦之事有所帮助呢？因为作为人，更为合宜的是，当我们听闻如此众多著名人物陷于错谬时，我们应为之悲伤，而非刻意探究这些错谬，以便在那些对此无知的人面前，享受一种最可鄙的、以知识自恃的满足。我从未听闻\Person{德谟克利特}之名，该有多好，胜于此刻悲伤地思索：一个在他那个时代备受推崇的人，竟认为众神是从固体中流溢出来的\OriginalTerm{影像}{images}，而非实体本身；这些影像凭借自身的运动四处流转，滑入人们的心灵，使神圣能力进入他们思想的领域；尽管，当然，那产生影像的物体，在卓越性和比例上，无疑应被判定为超越这影像，正如在坚实性上超越它一样。因此，如人们所说，他的观点摇曳不定，以至于他有时说\OriginalTerm{神性}{deity}是某种自然，影像由此流溢而出，然而神性本身只能通过其所倾泻和发出的影像而被思考；也就是说，这些影像从那种自然（他认为它是一种物质且永恒的事物，正因如此才是神圣的）中被携带着，如同一种蒸发或持续的\OriginalTerm{流溢}{emanation}，进入并闯入我们的心灵，从而使我们能够形成对神或众神的思想。因为这些哲学家认为，除了来自我们思想对象的物体的影像进入心灵之外，我们心灵中别无任何思想的原因；仿佛确实没有许多事物——是的，多于我们所能计数的——它们没有任何物质形式，但却是可理解的，并由那些知道如何理解这类事物的人所领会。以本质智慧与真理为例，如果他们对它不能形成任何观念，我倒奇怪他们为何竟要为此辩论；然而，如果他们确能在思想中对它形成某种观念，我希望他们能告诉我，真理的影像究竟从何种物体进入他们的心灵，或者它属于何种类型。\par

28. 然而，据说\Person{德谟克利特}在自然哲学学说上也与\Person{伊壁鸠鲁}不同；因为他主张，在\OriginalTerm{原子}{atoms}的聚合中有某种有生命和呼吸的力量，凭借这力量（我相信）他断言影像本身（不是所有事物的所有影像，而是众神的影像）被赋予了神圣的属性，并且心灵的最初开端在于那些被赋予神性的普遍元素之中，影像拥有生命，因为它们惯于要么有益于我们，要么伤害我们。然而，\Person{伊壁鸠鲁}在万物的最初开端中不设想任何东西，除了原子，即某些极其微小、无法分割且不能为视觉或触觉所感知的\OriginalTerm{微粒}{corpuscles}；他的学说认为，通过这些原子的偶然聚合（碰撞），便产生了无数世界、生物、赋予其生命的灵魂，以及众神——他赋予众神人形，却不置于任何世界之内，而是在世界之外、分隔诸世界的空间中；而且他不允许任何超出物质事物之外的思想对象。但为了使这些成为思想的对象，他说，从那些他描述为由原子构成的事物中，有比进入我们眼睛的影像更精微的影像流下并进入心灵。因为按照他的说法，我们视觉的原因在于某些极其巨大的影像，它们包裹了整个外部世界。但我料想你已了解他们对这些影像的看法。\par

29. 令我惊讶的是，\Person{德谟克利特}竟未能因如下事实而信服其哲学的错误：如此巨大的影像进入我们的心灵——心灵如此微小（如果如他们所声称的，灵魂是物质的，局限于身体的维度之内）——这些影像就其整个大小而言，不可能与心灵接触。因为当一个小物体与一个大物体接触时，它无论如何不可能在同一时刻被大物体的所有点触到。那么，这些影像如何在整体上同时成为思想的对象呢？如果它们仅在进入心灵并与其接触时才成为思想对象，鉴于它们就其整体而言既不可能进入如此小的身体，也不可能与如此小的心灵接触？请务必记住，我现在是按照他们的方式来说话；因为我并不认为心灵像他们声称的那样。确实，如果\Person{德谟克利特}认为心灵是非物质的，则这一论证仅能用于攻击\Person{伊壁鸠鲁}；但我们可以转而问他，为什么他没有认识到，对于非物质的灵魂而言，通过物质影像的接近和接触来思想，既是多余的，也是不可能的。两位哲学家肯定都同样被视觉的事实所驳倒；因为如此巨大的影像，就其整体而言，不可能触及如此小的眼睛。\par

30. 此外，当有人向他们提出，一个物体如何会产生一个影像，尽管从该物体放射出的影像数不胜数，他们的回答是，正因为影像如此大量地放射和消逝，它们拥挤聚集在一起，结果从众多影像中只看到一个。\Person{西塞罗}揭露其荒谬，他说，他们的神祇不能被看作是永恒的，正是由于这个原因：他是通过那些数不胜数、不断流出并消逝的影像而被思想的。当他们说，神明的形体是因无数\OriginalTerm{原子}{atoms}不断提供接替的支援而成为永恒的，因此其他的\OriginalTerm{微粒}{corpuscles}会立即取代那些离开神圣实体的微粒，并通过同样的接替防止神明的本性消散时，西塞罗回答说：\Quote{按此理，万物都将与神明一样永恒。}因为没有一样东西不具有同样无限量的原子来修复其持续不断的朽坏。他再次问道，他们的神怎么可能不害怕毁灭，因为他无时无刻不被原子不间断的冲击所打击和摇动——打击，是因为他被冲向他而来的原子所撞击；摇动，是因为他被穿过他而去的原子所穿透。不但如此，看到从他自己身上不断发出影像（关于这些我们已说得够了），他有何根据相信自己的不朽呢？\par

31. 至于那些持有这些见解之人的一切胡言乱语，尤为可叹的是，仅仅陈述它们，尚不足以使人在不加讨论反驳的情况下便予以摒弃；相反，那些心智最为聪慧的人却承担起详尽驳斥这些观点的任务，而这些观点一经说出，哪怕最愚钝的理智也应立即轻蔑地弃绝。因为即使我们承认有\OriginalTerm{原子}{atoms}存在，并且它们随意向撞击、相互击打和摇动，我们难道也可以承认，这样偶然相遇、组合在一起的原子，就能制造出什么东西，以致塑造其独特的形态、决定其形状、磨光其表面、赋予其色彩、或通过赋予生气使其活跃吗？所有这一切，只要人不仅用眼睛观察，更乐意用心灵观看，并从他的存在之创造者那里求得理智认知的能力，每一个人都会看到，除了通过\OriginalTerm{天主眷顾}{providence of God}以外，别无他途。不但如此，我们甚至没有自由承认原子本身的存在，因为，暂且不讨论学者们关于物质可分性的微妙理论，且看如何从他们自己的意见中轻易地证明原子的荒谬。众所周知，他们断言自然界中除了物体和\OriginalTerm{虚空}{empty space}，以及这些物体的\OriginalTerm{偶性}{accidents}外，别无他物；所谓偶性，我相信他们指的是运动、撞击以及由此产生的形式。那么，请他们告诉我们，他们将那些他们认为从较坚固的物体中流出的影像归入哪一类呢？这些影像，如果确实是物体，其坚固性如此之小，以至于只有当它们接触我们的眼睛时我们才能看到，接触我们的心灵时我们才能想到它们。因为这些哲学家的见解是：这些影像能从物质对象出发，到达眼睛或心灵，而他们却断言心灵也是物质的。现在我问：这些影像是否从原子本身流出？如果是，那么有些物体微粒在这个过程中被分离出来，这些原子又怎能是原子呢？如果不是，那么要么某种东西可以在没有这种影像的情况下成为思想的对象（对此他们极力否认），要么我们就要问，他们是从哪里获得关于原子的知识，既然原子绝不能成为我们思想的对象？但我羞于已经驳斥这些观点到如此地步，尽管他们毫不羞耻地持有这些观点。然而，当我想到他们竟敢为之辩护时，我羞耻的不是因为他们，而是因为人类本身的耳朵竟能容忍这样的胡言乱语。\par

% structure: p0038 source=h2 target=subsection reason=relative-heading-rank; base=h2

\subsection{第五章}

32. 因此，看到人心因罪污与肉欲而如此盲目，以致连这些荒谬的错误也能让博学之人花费闲暇去讨论它们，\Person{狄奥斯库若}，或者任何具有仆人之心的人，还会犹豫不确认，为了人类追求真理，最好的安排莫过于：有一个人被真理本身以不可言喻而奇妙的结合所摄取，并在世上彰显祂的位格，借着完美的教导和神圣的作为，激发人们相信那些尚不能凭理智领悟之事，从而获得救恩的信德？对于已成就此事的祂，我们献上侍奉；并且我们劝告你们，坚定而不移地信靠祂，因着祂，并非少数被选之人，而是全体子民，虽不能凭理性分辨这些事，却以信德接受了它们，直到他们因救恩真理的教导而被扶持，从这些困惑中逃脱，进入那完全纯洁而简朴的真理的氛围之中。此外，我们更应当毫无保留地顺服祂的权威，因为我们看到，在我们的时代，没有任何错误敢于高举自己，并聚集无知的群众围绕它，除非它寻求基督徒名号的庇护；而且，在所有属于更早时代、如今仍在基督徒名号之外的人当中，只有那些保留圣经的人，在他们的隐秘集会中继续拥有相当多的听从者，而圣经正是预言宣告主耶稣基督本人者，无论他们如何假装看不见或不理解这一点。此外，那些虽不在天主教的合一与共融之内，却以基督徒之名为夸耀的人，不得不反对信徒，并妄图以诉诸理性的伪饰来误导无知者，因为主带来了一种超越一切的治疗方法，即祂将信德的义务托付给万民。但如我所说，他们被迫采取这种策略，因为他们感到，若将自己的权威与天主教的权威相比较，便会极其可悲地被推翻。因此，他们企图借着所谓诉诸理性的名义和承诺，胜过那稳固不移的教会所坚立的权威。我们可以说，这种厚颜无耻是异端者的特征。然而，那位最仁慈的信德之主，既借着最著名的大公会议和宗座本身，将教会置于权威的堡垒之中，又借着少数具有虔诚博学与真诚灵修的人，赐予她同样不可战胜的丰富理性武装。培训门徒方法的完美之处在于，软弱者被极力鼓励进入权威的堡垒，以便他们安全置身其中后，能以最奋力的理性运用来维持为保卫自己所需的战斗。\par

33. 然而，柏拉图主义者们，当时在各方面受到种种错误哲学的围攻，他们宁愿隐藏自己的学说以待人探索，而不愿将其公开以遭受贬损；他们因为没有一位神圣人物来命令人信从，便在基督之名已广为人知，令这个世界诸国惊奇且不安之后，才开始展示和阐发\Person{柏拉图}的学说。那时，\Person{普罗提诺}学派在\Place{罗马}兴盛起来，该学派有许多非常敏锐和能干的人士为学者。但他们中有些人因对魔术的猎奇探究而堕落，另一些人则在主耶稣基督身上认出他们努力追求的、那本质且不变之真理与智慧的化身，转而事奉祂。这样，为重生和革新人类所需的一切权威的至高地位与理性之光，都已寓于那唯一的救恩圣名，以及祂的唯一教会之中。\par

34. 我丝毫不后悔在这封信中长篇大论地陈述了这些事，虽然你或许宁愿我采取另一种方式；因为你在真理中进步越大，就越会赞许我所写的内容，届时你也会赞许我的劝告，尽管你现在认为它对你的学业无益。同时，我已尽我所能回答了你的问题——其中一些在此信中，其余的几乎都通过简短批注写在你送来那些问题的羊皮纸上。如果你在这些答复中认为我做得太少，或者做了你意料之外的事，我的\Person{狄奥斯库若}，那你可没有适当考虑到你所问问题的对象是谁。关于那位演说家以及\Person{西塞罗}的\WorkTitle{论演说家}各书的所有问题，我略而未答。如果我着手阐述这些论题，在我自己看来，我就像是一个可鄙的琐屑之徒。因为，关于所有其他主题，如果有人希望我处理和阐明它们，而不与\Person{西塞罗}的著作相连，只是就其本身而言，我可能恰当地被提问；但在这些问题上，论题本身与我现在的身份不相调和。然而，若不是你的信使来到时，我正受病痛之苦，并在病后暂时离开\Place{希波}一段时间，我就不会在这封信里做这所有的一切。即使在近日，我仍遭受健康不佳和发烧的侵袭，因此寄送这些答复给你，可能比原本要有更多延迟。我恳切请求你写信，让我知道你是如何收到它们的。\par

\SourceRecord{Translated by J.G. Cunningham.；Nicene and Post-Nicene Fathers, First Series；Vol. 1.；Edited by Philip Schaff.；Buffalo, NY: Christian Literature Publishing Co.,；1887.；<http://www.newadvent.org/fathers/1102118.htm>.}

% structure: p0001 source=h1 target=section reason=page-title

\section{第122封信（主历410年）}

\Emph{致我所亲爱的弟兄众圣职人员，并\Place{希坡}全体子民，\Signature{奥思定}在主内致候。}\par

1. 首先，朋友们，我为基督的缘故恳切请求并哀恳你们，不要因我身体不在而忧伤。因为我想你们不会以为，我能在精神上与无伪的爱中与你们分离，尽管也许我比你们更感忧伤，因为我的软弱，使我难以担当基督的肢体加给我的一切挂虑；为了服事他们，我既出于对天主的敬畏，也出于对他们的爱，甘愿献身。亲爱的诸位，你们知道，我从未因贪图安逸而离开你们，只是当某些职责迫使我时，我才不得不如此；这些职责使我的几位圣善的同工与弟兄，不得不在海上及海外各地，承受我因身体的欠佳而非精神上的懈怠，得以免除的辛劳。因此，我最亲爱的弟兄们，你们行事当如宗徒所说：\BibleQuote{无论我来见你们，或不在你们那里，我可以听到你们的事，知道你们在同一精神内，同心协力为福音的\OriginalTerm{信德}{faith}奋斗。} \BibleRef{斐 1:27} 若是任何属世的苦恼令你们忧愁，这本身更应提醒你们，该当如何将心思专注于那可使你们无任何重担而生活的生命，所逃避的，不是这短暂生命中令人烦恼的艰辛，而是永火的可怕烈焰。因为你们若以如此多的焦虑、如此大的热切、如此多的劳苦，力求避免陷于此世的短暂痛苦，那么，你们该当何等急切地逃避永苦！如果那结束尘世劳苦的死亡尚且如此令人畏惧，那么，那引人进入永苦的死亡，我们该如何畏惧！如果此世短暂而鄙陋的欢乐尚且如此被爱，那么，我们该当以何等更大的热切，去寻求来世那纯洁无垠的喜乐！默想这些事，你们在善工上不可懈怠，好能按时候收获你们所播种的。\par

2. 我听闻你们已忘记了为穷人提供衣物的惯例，这\OriginalTerm{爱德}{charity}的工作，是我与你们同在时曾经劝勉你们的；如今我再次劝勉你们，不要被此世的患难所压倒而变得懈怠，你们现在所见的这些灾祸，正是那不能说谎的我们的主和救赎主所预言的。在此情形下，你们在\OriginalTerm{爱德}{charity}工作上不应懈怠，反而应比以往更加丰厚。正如人看见墙壁摇动，房屋即将倒塌时，会更快地逃往更安全的地方一样，基督徒越是从日益频繁的苦难中，察觉到此世的毁灭已近，就越应更迅速、更积极地将他们原先打算在地上积蓄的财物，转移到天国的宝库里去；这样，倘若任何人所难免的意外发生，那逃离了注定毁灭的居所的人，便可以欢欣；反之，若这类事并未发生，那么，那些无论何时死去、已将财产交托于那永生之主——即他即将前往投奔的那位——保守的人，便可以免去忧苦的焦虑。因此，我最亲爱的弟兄们，你们各人要按照自己的能力，各人自己就是自己能力最好的判断者，从自己的财物中取出一份，如你们素常所做的；并且，要比以往更加甘心乐意地去做；在此生的一切愁苦中，你们心中要常怀宗徒的劝勉：\BibleQuote{主已临近，你们什么也不要挂虑。} \BibleRef{斐 4:5} 愿有人向我汇报关于你们的情况，使我能晓得，你们遵循那美好的习惯，并非因为我有在与你们同在，而是出于服从那位永不远离的天主的诫命；这习惯，多年以来当我和你们在一起时，以及我离开后一些时候，你们一直遵守的。\par

愿主保守你们平安！亲爱的弟兄们，请为我们祈祷。\par

\SourceRecord{Translated by J.G. Cunningham.；Nicene and Post-Nicene Fathers, First Series；Vol. 1.；Edited by Philip Schaff.；Buffalo, NY: Christian Literature Publishing Co.,；1887.；<http://www.newadvent.org/fathers/1102122.htm>.}

% structure: p0001 source=h1 target=section reason=page-title

\section{\Person{热罗尼莫}致\Person{奥斯定}（主后410年）}

\EditorialNote{热罗尼莫致奥斯定。}\par

有许多人两脚跛行，即使颈项折断也不肯低头，固执于先前的错误，纵使已没有先前宣扬错误的自由，仍坚持不改。\par

与您卑微的仆人同在的圣善弟兄们，尤其您的虔诚而可敬的女儿们，向您致以恭敬的问候。我恳请尊驾以我的名义，问候您的弟兄——我的主\Person{阿利皮乌斯}和我的主\Person{埃沃迪乌斯}。\Place{耶路撒冷}被\Person{拿步高}掳获，不肯听从\Person{耶肋米亚}的劝告，宁愿切望\Place{埃及}，好死在\Place{塔黑培乃斯}，在那里永受奴役而灭亡。\par

\SourceRecord{Translated by J.G. Cunningham.；Nicene and Post-Nicene Fathers, First Series；Vol. 1.；Edited by Philip Schaff.；Buffalo, NY: Christian Literature Publishing Co.,；1887.；<http://www.newadvent.org/fathers/1102123.htm>.}

% structure: p0001 source=h1 target=section reason=page-title

\section{书信一二四（主后四一一年）}

\Emph{致\Person{阿尔比娜}、\Person{皮尼亚努斯}和\Person{梅拉尼娅}}，\Emph{主内可敬、圣洁可爱、以兄弟之情所思念的，\Person{奥古斯丁}在主内致候。}\par

1. 我，无论是出于目前的病弱还是天生的体质，极易感寒；然而，在这异常可怕的冬季，我所受的炎热之苦却不可能更甚，因为我因忧苦而发烧，以致无法——我不说赶往，而是飞往——你们那里（去拜访你们，我本该飞越大海才是），在你们既已定居于我近旁，又从如此遥远的地方来看我之后。也许你们还会以为，今冬严酷的天气是我遭受这一失望的唯一原因；亲爱的各位，我恳求你们，勿作此想。因为这些倾盆大雨所能带来的不便、艰辛甚至危险，哪一样我不会迎受和忍受，以便前往你们那里？你们在我们巨大的灾难中给予了我们如此大的安慰，你们在这弯曲悖谬的世代中，是由至高之光点燃的强烈火焰的明灯，因心灵的谦卑而被高举，又因你们所轻看的荣耀而散发出更辉煌的光芒。此外，我本可享受那赐给我尘世故乡的灵性福乐，因它获允有你们在场，而你们的声名在你们缺席时已为同乡人所耳闻——所闻如此之多，以致他们虽因爱德而相信了关于你们本性及因恩宠所成为的报道，却或许不敢向他人复述，以免不被相信。\par

2. 因此，我将告诉你们我未能前来的原因，以及使我被阻不得享有如此殊恩的试炼，好使我不仅能获得你们的宽恕，也能经由你们的祈祷，获得那位在你们内工作使你们为祂而生活者的慈悲。\Place{希波}的会众，主立定我服侍的，大部分且几乎全都是如此软弱的体质，即使遭遇相对轻微的磨难也可能严重危及它的安宁；然而目前，它正遭受如此沉重的苦难打击，即使它坚强，也几乎难以承受这重担的加诸。而且，当我最近返回时，我发现它因我的离开而被冒犯到了极危险的程度；你们，我们在主内因你们的灵性力量而欢欣的，能凭着健康的品味喜爱并赞同\Person{保禄}的话：\BibleQuote{谁软弱，我不软弱？谁跌倒，我不心焦？}\BibleRef{格后 11:29} 我对此感受尤深，因为这里有许多人通过贬损我们，试图煽动那些看似爱我们的人来反对我们，好为魔鬼在他们中腾出地方。但当那些我们顾念其救恩的人生我们的气时，他们对我们报复的强烈决心，不过是一种为他们自己带来死亡的不理性欲望——不是肉身的死亡，而是灵魂的死亡，其中死亡的事实借着败坏的臭气神秘地显露出来，远在我们的关心能够预见和防备之前。\par

无疑，你们会欣然原谅我这份焦虑，尤其是因为，如果你们不高兴并想惩罚我，你们或许想不出比我未能与你们在\Place{塔加斯特}见面更严厉的痛苦了。但我相信，借着你们的祈祷相助，当目前的阻碍尽速消除后，我可得允前来你们这里，无论你们在\Place{阿非利加}的哪一部分，如果我所劳苦的这城镇不配得（我也并不自以为它能配得）与我们一同因你们的临在而欢欣的话。\par

\SourceRecord{Translated by J.G. Cunningham.；Nicene and Post-Nicene Fathers, First Series；Vol. 1.；Edited by Philip Schaff.；Buffalo, NY: Christian Literature Publishing Co.,；1887.；<http://www.newadvent.org/fathers/1102124.htm>.}

% structure: p0001 source=h1 target=section reason=page-title

\section{书信一二五（公元四一一年）}

\Emph{致\Person{阿利比乌斯}，至可颂扬的主，以全敬畏最亲爱的弟兄，我\OriginalTerm{司铎职}{Priestly Office}的同伴，以及与他同在的弟兄们，\Person{奥思定}及与他同在的弟兄们，在主内请安。}\par

1. 我们深感悲痛，绝无法视为小事，\Place{希波}的民众竟大声说了那么多贬损您圣德的话；但是，良善的弟兄，他们大声说出这些事，并不像我们虽未受公开指责，却被判定犯了类似的事那么令人悲痛。因为当人们相信我们之所以在天主的众仆人中将他们中某些人留下来，不是出于对义德的爱，而是出于贪爱金钱时，难道我们不该期望那些如此相信我们的人，用他们的口说出隐藏在他们心中的话，从而尽可能获得对症的大剂量良药，而不是在种种致命猜疑的毒害下默然丧亡吗？因此，我们应当更加用心（而且正是为此，我们在此事发生前就已经商议过），设法使那些我们奉命在善工上作其榜样的人，确信他们所怀的种种猜疑毫无根据；更甚于设法责斥那些用言语明确表达出自己猜疑的人。\par

2. 因此，我并没有对虔诚的\Person{阿耳比纳}发怒，也不认为她该受责斥；但我认为她需要从这类猜疑中治愈。诚然，她并没有把那些我指的话指向我，而是在抱怨\Place{希波}的民众，好像是说：这些人的\OriginalTerm{贪吝}{covetousness}已被暴露；他们渴望把一个众人皆知轻看金钱、乐意施舍的富人留在他们中间，并不是出于他堪当此职，而是看中了他丰厚的财产；尽管如此，她几乎公开地说了她对我本人也怀有同样的猜疑；不仅她，就连她虔诚的女婿和女儿，也在那一天，在教堂的后殿里，说了同样的话。依我看，消除这些人的猜疑，比责斥他们说出猜疑更为必要。若这类荆棘甚至在我们最亲密的朋友心中，就是在如此虔诚、如此为我们所爱的人心中，也能生长起来攻击我们，我们怎能获得并得到免受这类烦恼的安息呢？正是无知的大众那样为你着想，而我却蒙受了来自教会之光的相似猜疑；因此你看，我们中谁的悲痛更大。我认为，这两种情形都不需要抨击，而需要救治；因为他们是人，他们的猜疑是人的，因此他们所猜疑的事，即使可能虚假，也并非难以置信。当然，这样的人还不至于如此愚蠢，竟相信民众贪求他们的金钱，特别是他们已体验到\Place{塔加斯忒}的民众未曾从他们那里得到分文，\Place{希波}的民众必然也将什么也得不到。不，这憎恨的一切毒火，全是针对\OriginalTerm{神职}{clergy}人员，尤其是针对主教，因为他们的权力显着地高高在上，人以为他们就像占有者和主上那样，在使用和享用教会的财产。我亲爱的\Person{阿利比乌斯}，切莫让软弱的人因我们的榜样而滋生这种有害且致命的贪吝。回想起在此次诱惑发生之前我们互相说过的那些话，这使得此事更加急迫。但愿我们凭着天主的助佑，通过友好的商谈，设法消除这困难；我们不应以为单凭我们的良心便能安然引导自己；因为此事并非仅凭良心之声便足以给我们提供正确指引的事例。如果我们不是我们天主的无用的仆人，如果我们内还存着一丝不求自己益处的\OriginalTerm{爱德}{charity}火花，我们就必定要竭尽所能，在天主面前不但在天主眼中行事端正，而且在人面前也行事端正；免得我们虽然在自身良心内畅饮不受干扰的水，却因用不小心的脚踏入，而使主的羊群畅饮混浊的溪流。\par

3. 对于你信中的提议——即我们应共同讨论以暴力强行索取之誓言的责任，我恳求你，莫让我们的讨论方式将完全清晰之事陷入晦暗。因为，若以求不可避之死，来强迫天主的仆人发誓去做人律与神律皆禁止之事，那么他理应宁可选择死亡而不发此誓，以免在履行誓言时，犯下罪行。但在这一事例中，强加于此人的，仅是民众坚决的叫嚣，而非罪恶，乃是那若行出来便是合法之事；此外，诚然有人担忧，恐怕有些鲁莽之徒——正如混杂于众多善人中的那样——若寻得骚乱的口实，及貌似有理可据的愤怒，便会因喜好暴动，而做出某些凶暴的恶行，然而，这一恐惧并无实现的把握——那么，谁会断言，为了逃避不确定的后果——我且不说损失或身体伤害，甚至死亡本身——而故意犯下发虚誓之罪是合法的呢？\Person{雷古卢斯}未曾从圣经中听过任何关于发虚誓之不虔敬的事，他从未听说过\Person{匝加利亚}的飞卷，他用以向\Place{迦太基}人确认其誓言的，不是基督的圣事，而是邪神的可憎之物；然而，面对无可避免的酷刑，及史无前例的恐怖死亡，他并未因恐惧而被迫发誓，而是由于他既已发誓，便自由地自愿承受这些，以免犯下发虚誓的罪。在那个时代，罗马监察官也拒绝将有些人录入名册——那并非继承天国荣耀的圣人之册，而是录入\Place{罗马元老院会堂}的元老之册——其中不仅有那些因畏惧死亡及残忍酷刑，而宁愿选择明显的发虚誓，也不愿回到无情的敌人那里的人；还包括一个自以为没有发虚誓之罪的人，因为他发誓之后，借着所谓必需的借口，以返回来破坏了誓言；由此可见，那些将他逐出元老院的人所考虑的，并不是他发誓时自己心中所想，而是那些接受他誓言的人所期待于他的。然而，他们从未读过我们在\BibleBook{圣咏}{}中不断歌唱的话：\BibleQuote{他虽起誓损己，也不变更。}我们惯常以最高度的钦佩谈论这些德行的实例，虽然它们见于那些对基督的恩宠与圣名陌生的人身上；然而，我们竟当真以为，发虚誓是否有时合法这个问题，是一个我们应该查阅圣经来寻求答案的问题吗？在这圣经中，为阻止我们因轻率发誓而犯此罪，已写下了这条禁令：\BibleQuote{总不可发誓}？\par

4. 我绝不怀疑这一格言的绝对正确性，即：诚信要求遵守誓言，并非仅依据发誓者的言辞，而是依据发誓者所知晓的、那接受其誓言者的期望。因为，要以言语——尤其是以寥寥数语——来界定那要求发誓者作出保证的许诺，是极为困难的。因此，有些人虽遵守了许诺的字面，却辜负了接受其誓言者已知的期望，他们便犯下发虚誓的罪；而另一些人，即便偏离了许诺的字面，却履行了当初发誓时对方所期望的，他们便没有犯下发虚誓的罪。因此，既见\Place{希波}人民之愿望，是要将圣善的\Person{丕尼亚努斯}留为居民，而非作为一名失去自由的囚犯，而是作为城内备受爱戴的居民；那么，他们对他所期望的限度，虽不能完全用他许诺之言词来涵盖，却是如此明显，以至他在发誓要留在他们中间之后，此刻却离开了，这事实并不会令任何人不安——只要他们听说他是为某一确定目的而离去，且心怀着归来的意愿。因此，除非他辜负了他们的期望，他便不会犯下发虚誓的罪，也不会被他们视为违背誓言；而他不会辜负他们的期望，除非他要么放弃了居住在他们中间的计划，要么在将来某个时候离开，且无意重返。愿天主禁止他如此背离他对基督与教会所应尽的圣洁与忠诚！因为，且不提天主对发虚誓者那可畏的审判——对此你同我一样清楚——我完全确信，倘若我们发现，像\Person{丕尼亚努斯}这样的人所犯的发虚誓，不仅被容忍而不激起义愤，甚至还被辩护，那么从今往后，我们便再无理由对任何拒绝相信我们凭誓言所证实之事的人表示不悦了。惟愿那拯救信靠祂者脱离诱惑的天主，以其仁慈拯救我们免于此事！因此，愿\Person{丕尼亚努斯}，如同你在来信中所写的那样，履行他所承诺的不离开\Place{希波}的许诺，正如我本人及本城其他居民不离开此地一般；当然，我们拥有随时出入的完全自由；唯一的区别在于：那些不受任何居住于此的誓言约束的人，也有权在任何时候，决意不再归来而离开，却不会负发虚誓的罪责。\par

5. 至于我们的圣职人员，以及定居在我们隐修院中的弟兄们，我不知道是否可以证明，他们曾对你们所遭受的指责给予帮助或教唆。因为当我调查此事时，我得知只有一位来自我们隐修院的人，一位\Place{迦太基}人，参与了民众的叫嚣；且那并不是在民众对你们口出侮辱之言时，而是在他们要求\Person{丕尼亚努斯}当司铎的时候。\par

我已在本信后附上一份给予他的许诺的副本，该副本乃取自他亲自在我监督下签署并校订的那份文件。\par

\SourceRecord{Translated by J.G. Cunningham.；Nicene and Post-Nicene Fathers, First Series；Vol. 1.；Edited by Philip Schaff.；Buffalo, NY: Christian Literature Publishing Co.,；1887.；<http://www.newadvent.org/fathers/1102125.htm>.}

% structure: p0001 source=h1 target=section reason=page-title

\section{第126封书信（主历411年）}

\Emph{致天主的圣洁女士及可敬的婢女\Person{阿尔比纳}，\Person{奥斯定}在主内向您致意。}\par

至于您所描述的、难以言喻的心灵忧伤，我理当抚平而非加深其苦涩，尽可能消除您的猜疑，而不是因我在此事上所遭受的，而发泄愤慨，致使一位如此可敬且献身于天主的人更加不安。\Place{希波}的人民并未对我们的圣洁弟兄、您的女婿\Person{彼尼亚努斯}做出任何足以合理激起他对死亡恐惧的事，尽管他本人或许怀有这样的恐惧。事实上，我们也曾担忧，人群中那些常为作恶而暗中勾结的亡命之徒，可能会借某个煽动激情的似是借口，寻隙发起暴乱，从而爆发大胆的暴行。然而，据我事后有机会查知，无人说过或企图做过任何此类的事情；但人民却对我的弟兄\Person{阿利比乌斯}高声叫嚷，发出许多辱骂和不配的指责，为此大罪，我切愿他们因他的祈祷而获得宽恕。就我而言，在他们的叫嚷开始后，我告诉了他们，我因承诺而不能违背他的意愿授予他圣秩，并补充说，如果我背信弃义，使他们得以拥有他作他们的\OriginalTerm{司铎}{presbyter}，他们就无法保留我作他们的主教；之后我便离开人群，回到我的座位上。于是，他们因我这出乎意料的答复而一度停顿动摇，如同火焰被风暂时吹退，但随后却开始变得更为激烈，他们想象着或许能从我这里强行逼出对承诺的违背，或是，若我信守诺言，他可能被另一位主教授予圣秩。对所有我能交谈的人，即那些在后殿中来到我身边的较为可敬和年长的人，我声明，我不会动摇而违背我的诺言，并且，在托付给我照管的教会中，除非征得并获得了我的同意，任何其他主教都不能授予他圣秩；而在表示同意这件事上，我同样会犯失信之罪。我接着说，如果他是在违背自己意愿的情况下被授予圣秩，那么，人民其实只是希望他一被授予圣秩就离开我们。他们不相信这会是可能的。然而，聚集在台阶前的人群，以可怕而不断的喧嚷和叫喊，坚持同样的决心，使得他们犹豫不决，不知所措。那时，人们高声辱骂我的弟兄\Person{阿利比乌斯}；那时，我们也忧虑会有更为严重的后果。\par

然而，尽管我因人民如此巨大的骚动，以及教会神职人员如此的惶恐而不安，我对那群暴民所说的，无非是我不能违背他的意愿授予他圣秩；在经历了这一切之后，我也没有受到影响，去做我同样已承诺不做的事，那就是以任何方式劝他接受\OriginalTerm{司铎}{presbyter}之职；倘若我能够说服他这样做，那么他受圣秩时便是自愿的。我信守了所作出的两个承诺：不仅是不久前向人民表明的那个，还有那个，就人而言，我只向一个人作过证的承诺。我守信，我说，不是对誓言，而是对我赤裸的诺言，即使面对如此的危险。的确，后来我们查明，对危险的恐惧是毫无根据的；但无论危险可能是什么，我们所有人都同样分担了它。恐惧也为众人所共有；我自己也曾想过要退避，主要是为着我们聚集所在的建筑物之安全而担忧。然而，有理由担心，如果我缺席，某种灾祸便更可能发生，因为那时人民会因失望而更加激怒，且较少为敬畏之情所约束。再者，如果我和\Person{阿利比乌斯}一同穿过稠密的暴民，我则有理由担心，会有人敢向他施暴；反之，如果我不带他离去，那么，倘若\Person{阿利比乌斯}遭遇了什么不测，而我似乎抛弃了他，将他交到了愤怒民众的手中，人们又会是多么自然地形成怎样的看法呢？\par

在极度的亢奋与悲痛之中，当我们智穷力竭、几乎无法喘息之际，看，我们虔诚的儿子\Person{皮尼亚努斯}突然且意外地，打发了一位天主之仆来见我，告诉我他愿意向民众起誓：假如他被迫接受祝圣，他将完全离开\Place{阿非利加}——我想，他认为民众若知道他当然不能违背誓言，就不会再继续喧嚷，明白坚持己见非但一无所获，反而只会迫使这个人离开我们，而即使他不能留下，我们至少也该把他当作邻居那样保留在身边。然而在我看来，若他这样起誓，恐怕民众的悲痛会更加强烈；对此，我保持沉默。他也通过那位信使恳请我去他那里，于是我立刻去了。他把先前托信使说的话又对我说了一遍，立刻又在那誓言上加了一条——他派另一位信使在我赶往他那里时转告我的——那就是：若无人强迫他违反自己的意愿接受圣职的重担，他愿意居住在\Place{希波}。对此，我就像在窒息之险中得到一阵清风似的从困窘中得着安慰，一言不发，只是加快脚步到了我的弟兄\Person{阿利皮乌斯}那里，把\Person{皮尼亚努斯}的话告诉了他。可是，他——我猜是小心谨慎，不想因自己的认可而做出任何可能冒犯你们的事——说道：\Quote{在这个问题上，不要问我的意见。} 听了这话，我赶到喧嚷的民众前，获得安静后，把所允诺的事连同那位起誓的保证，都宣告给了他们。然而，民众除了要他作自己的司铎以外，别无他想，因此并不像我预料的那样接受这个提议，而是交头接耳窃窃私语了一番后，请求在这允诺和誓约外加上一项条款：假如他日后乐意接受圣职，只能在\Place{希波}的教会，而非其他教会受职。我把这要求报告给他，他毫不迟疑地答应了。我带着他的答复回到他们那里，他们满心欢喜，随即要求按所允诺的起誓。\par

我回到你们女婿那里，发现他正为如何措辞表达他那由誓言确认的允诺而犯难，因为有多种可能发生的必要性，可能迫使他不得不离开\Place{希波}。他当时说出了自己的担忧，即可能会发生蛮族的入侵，为了避免这种危险就不得不离开此地。圣妇\Person{梅拉尼亚}还想加上气候的不利于健康作为可能离开的理由，但因他的回答而作罢。不过我指出，他所提出的理由是重大且值得考虑的，而且这个理由一旦发生，连市民自己也会被迫弃城而去；可是，假如将这个理由向民众说明，我们很有理由担心他们以为我们在预言灾祸，另一方面，如果把一个用以撤销允诺的借口笼统地冠以必要性的名义，人们可能会认为这必要性只是掩盖一个欺骗的意图。因此，他觉得我们应就此事试探民众的反应，结果正如我所预料。当执事宣读了他所口授的文字，并被众人欣然接受后，一听到有可能妨碍履行允诺的必要性条款传入耳中，立刻响起一阵反对的呼喊，允诺也被拒绝了；喧嚣再次爆发，民众以为这些交涉无非是要欺骗他们。当我们虔诚的儿子看见这情形，便命令将关于必要性的条款删除，民众这才又欢欣鼓舞起来。\par

我本乐意借疲倦为借口脱身，但他不愿在没有我陪同的情况下到民众那里去；于是我们一同前往。他告诉民众，他亲口授给执事宣读的那些话就是他自己的意思，他已用誓言确认了那个允诺，并将履行所允之事，然后他立刻用他口授的话把这一切复述了一遍。民众的回应是：\Quote{感谢天主！} 他们请求把所写的一切都予以签署。我们遣散了慕道者们，他立即在文件上签了名。然后，我们（\Person{阿利皮乌斯}和我）开始被催促——不是群众鼓噪，而是由德高望重的信友作为代表——催促我们这些作主教的也该在文件上署名。可是当我正要这样做的时候，虔诚的\Person{梅拉尼亚}表示反对。我奇怪她何以这么晚才提出反对，仿佛若我们因省略而不署名，真能叫他的允诺和誓言归于无效似的；不过，我听从了，因此我的签名并未完成，也没有人认为有必要再坚持要我们署名。\par

6. 我费心向你的圣善传达，据我看来已足够说明，民众在次日听说他已离开该城时的感受，或更确切地说，他们的议论。因此，若有人向你报告任何与我所述相矛盾的事，他或是故意，或是由于自己的错误，在误导你。因为我知道，我略过了一些在我看来无关紧要的事，但我晓得我所说的全是真话。所以，事确是真的：我们圣善的儿子\Person{Pinianus}在我面前并得到我的允许后发了誓，但并非如传言所说，他是奉了我的任何命令而这样做的。他自己知道这一点；他所派遣给我的天主仆人也知道此事，第一位是虔诚的\Person{Barnabas}，第二位是\Person{Timasius}，他也是通过他们向我传送了他将留在\Place{Hippo}的承诺。至于民众本身，他们当时正在用呼喊催促他接受\OriginalTerm{司铎}{presbyter}的职务。他们没有要求他发誓，但当誓言被提出时，他们也没有拒绝，因为他们希望，如果他留在我们当中，可能在他内产生同意接受\OriginalTerm{圣秩}{ordination}的意愿，而他们担心，如果违背他的意愿祝圣他为司铎，他可能会依照他的誓言离开\Place{Africa}。因此，他们在喧嚷的举动中，也是出于对天主事工的重视（因为\OriginalTerm{祝圣司铎}{consecration of a presbyter}确是天主的事工）；而且，既然他们不仅对他留在\Place{Hippo}的承诺感到不满足，除非他也承诺，若他日后接受圣职，他绝不在别处，只在他们当中接受，那就极其明显地看出，他们盼望他居住在他们当中，以及他们并没有放弃对天主事工的热诚。\par

7. 那么，你凭什么声称民众这样做是出于对金钱的卑鄙贪欲呢？首先，那些喧嚷的民众根本不可能从中获得任何此类的利益；因为\Place{Thagaste}的民众从你赐予他们教会的捐赠中，除了因看到你的善工而喜乐外，别无所得，\Place{Hippo}的民众，或任何其他地方，只要你已经或将要遵守你主关于\BibleQuote{不义的钱财}的律法，情况也是如此。因此，民众极其激烈地坚持要指导他们的教会在处理如此伟大人物时的行事方式，并非向你要求金钱上的利益，而是表明他们对你轻视金钱的赞赏。因为就我自己的情形而言，由于他们听说我轻看祖产——那不过是几小块田地——而将自己奉献于服事天主的自由，他们便喜爱这种无私的精神，并不嫉妒将这份礼物给予我的出生地教会\Place{Thagaste}，反而，当\Place{Thagaste}教会没有将圣职加给我时，他们便强使我，可以说，成为他们自己的人。那么，对于我们的\Person{Pinianus}，他以如此非凡的决心克服并踏碎了如此巨大的财富、如此光明的希望以及享乐此世的强劲本性，他们岂不会更加热切地爱慕他吗？的确，在许多人看来——那些以自己衡量自己——我似乎与其说是抛弃了财富，不如说是找到了财富。因为我的祖产几乎够不上我现在身为管理者所拥有的教会财产的二十分之一。但无论在哪个教会，尤其是在\Place{Africa}，我们的\Person{Pinianus}若被祝圣为（我不说\OriginalTerm{司铎}{presbyter}，而是）\OriginalTerm{主教}{bishop}，与他从前的富裕相比，他仍然处于极度的贫穷中，即使他怀着凌驾于天主基业之上的心态来使用教会的收入。基督徒的贫穷，在那些没有理由使人怀疑其意在增加财富的人身上，是更为清晰和确实地受人爱戴的。正是这种赞赏点燃了民众的心灵，并激起他们如此猛烈而持久的喧嚷。因此，我们不可无缘无故地指责他们卑鄙贪财，而不如，在不加诸不良动机的前提下，至少允许他们去爱慕别人身上自己所没有的那种优点。因为尽管人群中可能混杂着一些穷人或乞丐，他们助长了喧嚷，并出于希望从你那受人敬重的富裕中得到一些救济以解困乏的动机，即便如此，在我看来，这也并非卑鄙贪财。\par

8. 因此，所剩下的便是，这种可耻贪欲的指责必然间接地指向\OriginalTerm{神职班}{clergy}，尤其是\OriginalTerm{主教}{bishop}。因为人们设想我们作为教会财产的主人行事；我们被设想为享用它的收入。总之，我们为教会所收受的任何金钱，要么仍在我们手中，要么已按我们的判断花费掉了；并且，除了\OriginalTerm{神职班}{clergy}和\OriginalTerm{隐修院}{monastery}中的弟兄们，以及极少数贫困者之外，我们没有将其中任何部分分给任何民众。我这样说，并非意味着你所说的事必定是特别针对我们的，而是说，如果这些话是针对我们之外的任何人所说，那必是难以置信的。那么，我们该怎么办呢？如果不可能在敌人面前澄清自己，我们至少该用什么方法在你面前澄清自己呢？这事关乎灵魂；它在我们内，隐密不为人眼所见，惟有天主知晓。那么，除了呼求明知此事的天主作证之外，我们还剩下什么可做的呢？因此，当你对我们怀有这些猜疑时，你不是命令，而是完全强迫我们立誓——这远比命令人立誓更为严重，而你在信中认为命令立誓是极其可责的，并以此谴责我；我说，你强迫我们，不是以死亡威胁身体，如同\Place{Hippo}的民众所为的那样，而是以死亡威胁我们的美名，这美名对我们而言应当比生命本身更为宝贵，这是为了那些软弱的弟兄们的缘故，我们努力在一切事上为他们树立善行的榜样。\par

9. 然而，我们不怨恨强迫我们宣誓的你们，正如你们怨恨\Place{希波}的人民一样。因为你们相信，就像人判断别人一样，那些虽然实际上不存在于我们身上，却可能本来存在的事情。你们的猜疑，我们必须努力消除，而不只是谴责；并且既然我们在天主眼中良心清白，我们也必须努力使我们的品格在你们眼中清白。可能，正如\Person{阿利比乌斯}和我在这次考验发生前彼此说过的，天主会恩赐，不仅你们——我们深爱的\OriginalTerm{基督奥体}{Christ's body}同伴——而且甚至我们最不可化解的敌人，也会完全满意地看到，我们在管理教会事务上并未被任何贪财之心所玷污。在这事成就之前（如果主应允我们的祈祷，允许它成就），在此期间请听我们被迫要做的事，而不要长时间拖延你们心灵的治愈。天主为我作证，关于那些我们被以为喜欢行使管辖权的教会财产的全部管理，我只是将其视为一种负担，是因着对弟兄们的爱和敬畏天主而加于我的：我并不喜欢它；不，如果我可以，只要不违背我的职分，我宁愿摆脱它。天主也为我作证，我相信\Person{阿利比乌斯}在这件事上的想法与我相同。然而，一方面，人民，尤其是更糟糕的\Place{希波}人民，因对他的品格持有另一种看法，轻率地大大冤枉了\Person{阿利比乌斯}；另一方面，你们作为天主的圣人，充满真挚的怜悯，由于相信了关于我们的这些事，认为适合在名义上责备那同一\Place{希波}人民的同时，触动和劝诫我们，而\Place{希波}人民在所谓的贪婪之罪上毫无份。你们毫无疑问渴望纠正我们，而且不带恨恶（这远离你们吧！）；因此我不应向你们发怒，反而应感谢你们，因为你们不可能比这更恰当地将谦恭与自由结合在一起了：你们不将你们的想法作为对主教的冒犯性指控陈述出来，而是留待人们通过间接推论发现它们。\par

10. 不要让我认为必须这样用誓言来确认我的陈述这件事，使你们因觉得受到严厉对待而烦恼。\Person{保禄}对他写信的对象并没有心硬或缺乏仁慈之情：\BibleQuote{我们从来没有用过谄媚的话，这是你们知道的；也没有托词贪心，天主可以作证。} \BibleRef{得前 2:5}在人们能观察的事上，他诉诸他们自己的见证，但关于隐藏的事，除了天主，他还能诉诸谁呢？因此，如果就连\Person{保禄}——正如众人所知，他的劳苦是这样的：除非极度必要，他从不从他施予基督恩宠的团体中为自己取任何东西养身，在一切其他情况下，都是亲手劳碌来获取所需——都合理地害怕人们的无知会使他们对他产生这样的想法，那么，我们这些在圣德的功绩和心灵的坚强上都远远落后于他的人，而且我们不仅不能亲手劳作来供养自己，即使我们能劳作，也因着积累的大量职分而被禁止（我相信宗徒们是免于这些职分的），必须费多大的努力才能令人相信我们的无私啊！因此，不要让最卑鄙的贪婪指控在这件事上再被加于基督子民——即基督的教会。因为，这指控加于我们，是比加于那人民更可容忍的：对我们来说，怀疑虽然无根，却可能并非全然不可能地落下；而关于那人民，众所周知，他们既不可能怀有贪婪的欲望，也不可能被合理地怀疑有此欲望。\par

11. 因为，对于拥有任何信德的人——更何况是拥有基督信仰的人！——违背誓言，我不说是做相反誓言的事，而是对履行誓言的任何犹豫都算，是完全错误的。我对这个问题的判断，在我寄给我兄弟\Person{阿利比乌斯}的信中已经足够充分地阐明了。你们的圣善写信问我：\Quote{是否我或\Place{希波}人民认为，任何人都有义务履行被暴力强迫所发的誓言。}但你们的意见如何？你们是否认为，即使死亡——在这件事上，是毫无理由地恐惧的——确实迫近，一个基督徒就可以使用他主的名来证实谎言，并呼号他的天主为一个虚假见证？因为确实地，一个基督徒，如果被即死威胁所迫，要通过假见证发虚誓，就应该宁可失去生命，也不可失去荣誉。互相敌对的军队在战场上以死亡相互威胁，对此没有不确定性；然而，当他们彼此以誓言保证时，我们赞扬那些信守承诺的人，我们也公正地憎恶那些不守信的人。那么，是什么动机导致他们彼此发誓呢，不就是双方都害怕被杀或被俘吗？而通过这承诺，即使这样的人也认为自己受约束，以免犯下渎圣和虚誓之罪，如果他们不履行因怕死或怕被俘而被迫所发的誓，并且违背了在这种情形下做出的承诺：他们更害怕违背誓言，胜过夺取一个人的生命。而我们难道打算讨论这样一个可争议的问题：是否必须履行天主的仆人——他们对圣洁负有卓越的义务——在怕受伤害下发过的誓言，以及那些按基督命令分施财产、奔向基督徒成全的隐修士们所发的誓言？\par

12. 请告诉我，我恳求你，在他应许居住在这里的承诺中，有什么值得被称为流放、流徙或放逐的艰难呢？我认为司铎职并非流放。难道我们的\Person{品尼安努斯}会宁愿选择流放而不愿接受这个职位吗？我们万万不可如此为一个天主的圣人且为我们所深爱之人辩护：我说，天主不许，绝不可说他宁愿选择流放也不愿担任司铎职，宁愿\OriginalTerm{背弃誓言}{perjure himself}也不愿忍受流放。即使他是被迫向我们起誓应许住在我们中间，我也会说这话；但事实却是，这誓言并非在他拒绝时被强迫发出，而是在他自己主动提出后被接受的。正如我已说过的，这誓言被接受，更是因为希望——他的留下鼓励了这希望——他也可能同意遵行我们的愿望而接受圣职。总之，无论人们对\Place{希波}的民众持有何种看法，那些可能强迫他起誓的人，与那些可能——我不是说强迫，但至少是——劝他违背誓言的人，两者的情形大不相同。我也相信，\Person{品尼安努斯}本人不会不愿认真思考：在恐惧压力下发誓，无论那压力有多大，抑或在全无恐慌之际故意\OriginalTerm{故意背誓}{deliberate perjury}，哪一样更为恶劣。\par

13. 感谢天主，\Place{希波}的民众认为，只要他带着将此城当作自己家园的意图而来，并在必要时外出时也怀着再回到我们这里的意向而行，那么他居住这里的承诺就算是充分履行了。因为如果他们要求逐字照做承诺的话，那么一个天主的仆人理当严守每一条文句，而不是\OriginalTerm{背弃誓言}{forswear himself}。但正如他们如此捆绑任何人——更何况是他这样的人——是一种罪过，同样，他们自己也证明了他们并无如此不合理的期望；因为一听说他离去时是打算回来的，他们就表示满意了；而对誓言的信守，所要求的不过是履行接受该誓言之人所期望的。此外，容我问一下，说他在亲口起誓时，曾提到有某种必要可能会妨碍他履行承诺，这究竟是何意？事实是，他亲口命令将\OriginalTerm{限定条款}{qualifying clause}删除了。如果他要加上，那应该是在他自己对民众说话时；但若他那样做了，他们肯定不会回答说，\Quote{感谢天主}，而会重新发出先前执事宣读了带有限定条款的誓言时他们所作的那些抗议。其实，这种免于遵守承诺的必要性理由是否被插入，究竟有何不同？人们对他所期望的，无非是我们上面所说的；而一个使接受其誓言之人所明知期望落空的人，只能是一个\OriginalTerm{背誓者}{perjured person}。\par

14. 因此，愿他的承诺得以履行，愿软弱者的心得到医治，以免一方面赞同它的人因这个显著的榜样而学去模仿\OriginalTerm{背誓}{perjury}的行为，另一方面谴责它的人有正当的理由说，我们当中没有一个人配得信任，不仅在我们许下承诺时，甚至在我们起誓时也是如此。我们尤其要提防在这件事上给敌人的口舌留下把柄，这些口舌被\OriginalTerm{大敌}{the great Enemy}用作射杀软弱者的利箭。然而，天主不许我们指望像\Person{品尼安努斯}这样的人，所行会不同于天主敬畏所激发的，以及他自己超凡虔诚所认可的。至于我本人，你们责备我没有干涉以拦阻他的誓言，我承认我无法让自己相信，在如此混乱和可耻的情况下，我应当宁可让我所服务的教会被颠覆，也不接受这样一个人赐给我们的解救。\par

\SourceRecord{Translated by J.G. Cunningham.；Nicene and Post-Nicene Fathers, First Series；Vol. 1.；Edited by Philip Schaff.；Buffalo, NY: Christian Literature Publishing Co.,；1887.；<http://www.newadvent.org/fathers/1102126.htm>.}

% structure: p0001 source=h1 target=section reason=page-title

\section{第130号书信（公元412年）}

\Emph{致\Person{普罗巴}，天主虔诚的婢女，基督的仆人及基督众仆之仆\Person{奥思定}主教，奉万主之主的名致以问候。}\par

% structure: p0003 source=h2 target=subsection reason=relative-heading-rank; base=h2

\subsection{第一章}

想起你的请求和我的承诺，即一旦我们向之祈求的那位赐予时间和机会，我就会就向天主祈祷的主题写一些东西给你；我现在感到有责任清偿这笔债，并在基督的爱中满足你虔诚的愿望。我无法用言语表达我因这请求而多么喜乐，因为我从中看出你对这件至关重要的事是多么关切。因为，根据宗徒的告诫：\BibleQuote{那真正做寡妇的，孤独无依，已仰望天主，昼夜不断祈求}\BibleRef{弟前 5:5}，有什么事比你在寡居中日夜继续祈祷更相宜呢？如果仅就现世而言，你高贵富有，又是那样显赫家庭的母亲，虽然寡居，却并非孤独，但你却明智地理解，在今世和今生中，灵魂没有稳妥的份；那么，对祈祷的关切竟占据了你的心，并在其中居于首位，这看起来确实令人惊奇。\par

因此，那启发你这种思想的主，确实正在做祂向门徒所应许的事；那时门徒们忧闷，不是为自己，而是为全人类，当他们从祂那里听到骆驼穿过针眼比富人进入天国还容易之后，便对任何人的得救感到绝望。祂给了他们这个奇妙而仁慈的回答：\BibleQuote{在人不可能的事，在天主是可能的。}\BibleRef{玛 19:21} 因此，那位能使富人也进入天国的，启发你怀有这种虔诚的焦虑，使你想到有必要就如何祈祷的问题来请教我。因为当祂还在地上时，就曾接纳富有的\Person{匝凯}\BibleRef{路 19:9} 进入天国；并且在祂复活升天、受到光荣之后，祂使许多富人鄙视这个现世，并通过赐予他们圣神，消除他们对财富的贪欲，使他们成为更真正的富裕。因为，如果你不信任祂，你又怎能如此渴望向天主祈祷呢？如果你寄望于无常的财富，而忽视宗徒的健全劝告，你又怎能信任祂呢？宗徒的劝告是：\BibleQuote{你要嘱咐那些今世富有的人，不要心高气傲，也不要寄望于无常的财富，而要寄望于那将万物丰富地赐给我们享用的永生天主；又要劝他们行善，在善工上致富，甘心施舍，乐意供应，为自己积存美好的根基，为要掌握永生。}\BibleRef{弟前 6:17}\par

% structure: p0006 source=h2 target=subsection reason=relative-heading-rank; base=h2

\subsection{第二章}

因此，你出于对真正生命的爱，应该自视为在此世是\Quote{孤独}的，无论你的境况多么繁荣。因为，与那真正的生命相比，这深受喜爱的现世生命，无论如何幸福和长久，都不配称为生命；同样，主在\Person{依撒意亚}的话中所应许的：\BibleQuote{我要赐给他真正的安慰，平安上加平安}，才是真正的安慰；没有这种安慰，人们在任何世俗的慰藉中，与其说是得到安慰，不如说是更感孤独。因为至于财富和高位，以及人们不认识真正幸福时所想象的幸福所在的其他一切事物，它们带来什么安慰呢？其实，不依赖这些事物比充分享受它们更好，因为当拥有它们时，由于害怕失去它们，所引起的烦恼，比强烈渴望拥有它们时所感受到的欲望更甚。人并不因为拥有这些所谓的\Quote{好}事物而成为善的；相反，如果人以其他方式成为善的，他才能善用这些事物，使之真正成为好的。因此，真正的安慰不在于这些事物，而在于那些可以找到真正生命的事物。因为使人有福的，正是那使人成为善的同一能力。\par

4. 的确，善人在此世常常是他人不小的安慰之源。若我们为贫困所扰，为丧失亲人而悲伤，为身体的痛苦所折磨，或在放逐中憔悴，或被任何灾祸所困扰，若有善人来看望我们——这些人不仅能与喜乐的一同喜乐，也能与哭泣的一同哭泣\BibleRef{罗 12:15}，且懂得给予有益的劝告，引导我们在交谈中吐露心声——结果便是：崎岖变为平坦，重担负于减轻，困难奇妙地得以克服。但这乃是借那位以圣神使他们成为善人的天主，在他们内并通过他们而成就的。反之，即便财富丰裕，没有丧亲之痛，身体康健，安居本土，平安稳妥，若邻人中恶徒充斥，无一可信赖，无人不是我们既疑又实遭其奸诈、欺骗、暴怒、纷争与陷阱，那么上述种种岂不尽成苦涩与伤悲，其中再无甘饴与愉悦？因此，无论我们在此世境遇如何，若无朋友，便无真乐可言。然而在此生中，能遇到一个在心灵与行为上均可全然信赖的真朋友，是何等稀罕！因为无人能像深知自己那般深知他人，而即使对自己，也无人能确知明日自身行为；因此，虽有许多人可由其果实而被认识，有些人以其善行令邻人欢欣，另有些人以其恶行令邻人忧伤，然而人心如此难测，如此易变，宗徒的劝诫便极具智慧：\BibleQuote{所以时候未到以前，你们什么也不要判断，直到主来，祂要照出暗中的隐情，显露人心的计谋；那时各人才可由天主那里获得称誉。}\BibleRef{格前 4:5}\par

5. 故此，身处此世的黑暗之中，我们犹如离家而寄居尘世的旅客，尚不在主前，因为\BibleQuote{我们现在只是凭信德往来，并非亲眼看见。}\BibleRef{格后 5:6}基督徒的灵魂应自觉孤寂，恒心祈祷，并学着将信德的目光专注于天主神圣圣经之言，犹如\BibleQuote{在暗处发光的灯，直到天亮，晨星在你们的心中升起的时候。}\BibleRef{伯后 1:19}因为这盏灯所借以照耀的那不可言喻的根源，乃是在黑暗中照耀的光，而黑暗却未能胜过它——这光，若要看见，我们的心必须借信德得以净化；因为\BibleQuote{心里洁净的人是有福的，因为他们要看见天主。}\BibleRef{玛 5:8}且\BibleQuote{我们知道，一旦祂显现出来，我们必要相似祂，因为我们要看见祂实在怎样。}\BibleRef{若一 3:2}死亡之后，才是真实的生命；孤寂之后，才是真正的安慰。那生命将\BibleQuote{救我们灵魂免于死亡}，那安慰将\BibleQuote{拭去我们眼中的泪}，并且，如圣咏所续，我们的脚将免于颠仆；因为那里再无诱惑。再者，若无诱惑，亦无须祈祷；因为在那里，我们不再期盼所许之恩，而是瞻仰已赐之福；故圣咏作者又加说：\BibleQuote{我要在上主的领域，在生命的地域里行走。}那时我们将生活其处——而非如今所处的死者旷野，正如宗徒所说：\BibleQuote{因为你们已经死了，你们的生命已与基督一同藏在天主内了；当基督，我们的生命显现时，那时，你们也要与祂一同出现在光荣之中。}\BibleRef{哥 3:3}因为这才是真实的生命，宗徒曾劝勉富有的人在此生命中要行善致富；在这真生命中才有真实的安慰，正因欠缺这安慰，一位寡妇即使有儿孙，并以虔诚治理家政，敦促所有亲人仰望天主，她依然是\Quote{孤寂}的：她会在祈祷中说：\BibleQuote{我的灵魂渴慕你，我的肉身切望你，在这干旱涸竭无水的地域。}这奄奄一息的生命岂非正是如此一片地域，无论我们有再多有限的慰藉，有再多旅途中的友伴，有再多的财产丰盈。你深知这一切是何等无常；即便它们恒常不变，与来生所应许的幸福相比，又算得了什么呢！\par

6. 我对你说这些话，你是一位富有的贵妇，身为名门之母，向我请求一篇关于祈祷的讲论，我的用意是让你感受到：即使你的家人都健在，且一如你心愿地生活，但只要你尚未获得那拥有真实和常存安慰的生命，你便是孤寂的。在那生命里，圣经的预言将得以应验：\BibleQuote{清晨我们就为你的慈爱所满足，我们一生一世欢欣鼓舞；你按照使我们受苦的时日，和我们遭遇祸患的年岁，使我们欢乐。}\par

% structure: p0011 source=h2 target=subsection reason=relative-heading-rank; base=h2

\subsection{第三章}

因此，直到那安慰来临，你要记着，为了你\BibleQuote{昼夜恒在祈祷和恳求}，无论你周遭的世俗繁荣多么丰盛，你是孤独无靠的。因为宗徒并非将这份恩赐归于每一个寡妇，而是归于那真正为寡妇、又孤独无靠的，她\BibleQuote{信赖天主，昼夜恒在恳求。}然而，要极其警醒地留意接下来的警告：\BibleQuote{但那好宴乐的寡妇，虽生犹死；}\BibleRef{弟前 5:5} 因为一个人活在他所爱、所极欲、并自认为有福的事物中。因此，圣经论及财富说：\BibleQuote{若财富增加，勿要系心于其上。}我论及娱乐也如此说：\Quote{若娱乐增加，勿要系心于其上。}所以不要因为这些事物不缺乏，且你大量拥有，仿佛从世俗幸福最丰沛的泉源涌流而出，就高看自己。务要以淡泊和轻看的心对待这些拥有，除了身体的健康以外，不要从中寻求什么。因为这是不可轻视的祝福，因为它是此生工作所需，直到\BibleQuote{这可死的人穿上了不可死}\BibleRef{格前 15:54}；换言之，那真正的、完善的、永恒的健康，既不会被今世的软弱削弱，也不会被可朽的欢愉修复，而是以天上的活力持久，被永恒不朽的生命所激发。因为宗徒自己说：\BibleQuote{不要为肉体操心，以满足其私欲，}\BibleRef{罗 13:14} 因为我们必须照顾肉体，但只限于健康所必需的；\BibleQuote{因为从来没有人恨过自己的肉身，}\BibleRef{弗 5:39} 正如他也同样说过。因此，他也劝勉\Person{弟茂德}——他看起来对身体过于严苛——要\BibleQuote{因你的胃口和屡次的疾病，要稍微用一点酒。}\BibleRef{弟前 5:23}\par

许多圣洁的男女，对那使人若活在其中的娱乐——即那些若依附并居住其中、以之为心的喜悦、虽生犹死的人——极力防范，他们已将那如同娱乐之母的东西抛弃，将财富分施给穷人，如此储藏在更安全的天堂宝库里。如果你因对家庭的某种义务顾虑而\Emph{受阻}不能这样做，你自己知道，你怎样向天主交代你对财富的使用。因为除了在人内的神，没有人知道在人内的事。\BibleRef{格前 2:11} 我们不应该在\BibleQuote{时候未到以前，什么也不要判断，只等主来，祂要揭发暗中的隐情，且要显露人心的计谋：那时，各人才由天主那里获得称誉。}\BibleRef{格前 4:5} 因此，作为寡妇，你的职责在于留心：若娱乐增加，勿要系心于其上，免得那应当升起以便活着的，却因与它们腐朽影响接触而死。你要将自己算作经上所载的人之一：\BibleQuote{他们的心灵要永生。}\par

% structure: p0014 source=h2 target=subsection reason=relative-heading-rank; base=h2

\subsection{第四章}

你们已听到，若要祈祷，应当是怎样的人；接下来，请听应当为什么祈祷。这是你认为最需要征询我意见的问题，因为你因宗徒的话感到不安：\BibleQuote{我们不知道如何祈求才对}\BibleRef{罗 8:26}；并且你开始担心，若不按当有的方式祈祷，比完全不祈祷，为你更有害。你的困难可简短解答如下：\Quote{为幸福的生活祈祷吧。}这是所有人都愿意获得的；因为即使那些品性最坏、最被遗弃的人，也绝不愿那样生活，除非他们认为，通过这种方式，他们已经或能够真正获得幸福。那么，除了那恶人与善人都渴望，但只有善人才能获得的事以外，我们还应该为什么祈祷呢？\par

% structure: p0016 source=h2 target=subsection reason=relative-heading-rank; base=h2

\subsection{第五章}

你也许要问：这幸福生活是什么？在这个问题上，耗费了许多哲学家的才智和闲暇，然而，他们对它的探究失败，恰恰与他们未曾尊崇那从中产生幸福者，且对祂不感恩的程度相当。那么，首先，请考虑我们是否应接受那些哲学家的意见，他们宣称，人若随心所欲地生活，便是幸福的。我们断不可相信此言；因为如果一个人的愿望驱使他度邪恶的生活，那岂不是说，在堕落意志得以顺利实现时，他反而被证实为悲惨的人吗？就连未曾敬拜天主的哲学家们，也对此观点深恶痛绝，弃之如敝屣。其中一位最雄辩者说：\Quote{然而，看哪，还有另一些人，他们并非真正的哲学家，而是善于辩论的人，他们断言，凡随心所欲生活的人都是幸福的。但这肯定是不真实的，因为希求不相称的事物本身就是最悲惨的事；未能如愿获得所想要的，并不像想要其实不当贪求的事物那样悲惨。}你的意见如何？这些话，不论出于何人之口，岂非源于真理本身吗？因此，我们在此可以说宗徒论及一位\Place{克里特}诗人时所言，那位诗人的观点曾使他喜悦：\BibleQuote{这见证是真的。}\BibleRef{铎 1:13}\par

11. 因此，凡拥有一切他所愿有的，并且不愿有任何他所不该愿的事物的人，才是真正幸福的。明白了这一点，现在我们来观察一下，人们可以正当地希望拥有哪些事物。有人渴望婚姻；另有人，成为鳏夫后，选择此后过有\OriginalTerm{节欲}{continence}的生活；第三个人虽已结婚，仍选择实践节欲。尽管这三种状况中，一种可能优于另一种，但我们不能说这三个人中的任何一个人，所愿的是他不该愿的：对于渴望儿女作为婚姻的果实，以及渴望已得的儿女享有生命与健康——这些渴望中，后者是那些不再婚的寡妇们通常所牵挂的；因为尽管她们拒绝第二次婚姻，现在不想再有儿女，但她们希望已有的儿女能健康生活。那些保守完整童贞的人，则免于此类的挂虑。然而，所有人都有一些所爱的人，他们也正当地渴望这些人获得现世福利。但当人们为自己以及他们所爱的人获得了这种健康时，我们就可以说他们现在是幸福的吗？他们固然拥有了某种完全可以渴望的事物；但倘若他们没有其他更大、更美、更具效用与美善的事物，他们依然远未拥有幸福的生活。\par

% structure: p0019 source=h2 target=subsection reason=relative-heading-rank; base=h2

\subsection{第六章}

12. 那么，我们是否可以说，除了身体的这种健康之外，人们还可为自己及所爱的人渴望荣誉与权力呢？全然可以，只要他们渴望这些，是为了借此促进那些依靠他们之人的利益。如果他们追求这些事物，不是为了事物本身，而是为了借此达成某种善事，那么这种愿望便是正当的；但若只是为了满足骄傲与狂妄的空虚快感，以及虚荣之多余的、有害的胜利，这愿望便不正当。因此，人们为自己及亲属渴望必需事物的适度份额，这并没有错；宗徒论及此事时说：\BibleQuote{的确，虔敬是一个获利的富源，但应有知足的心，因为我们没有带什么到世界上，同样也不能带走什么，只要我们有吃有穿，就当知足。至于那些想望致富的人，却陷于诱惑，堕入罗网和许多不理智及有害的欲望中，这欲望使人沉溺于败坏和灭亡之中，因为贪爱钱财乃万恶的根源；有些人曾因贪求钱财而离弃了信德，使自己受了许多刺心的痛苦。}\BibleRef{弟前 6:6} 人若只渴望这适度的份额，而不奢求更多，这渴望便是正当的；因为倘若他的渴望超出了此限，他便不再渴望这适度的份额，因而其愿望便不正当。那名曾如此渴望并祈求的人说：\BibleQuote{愿你使虚伪和欺诈远离我，赐我不贫乏也不富裕，只供予我必需的食粮；免得我饱饫而背叛你说：\InnerQuote{谁是上主？}或过于贫乏而行窃，加辱我天主的名。}\BibleRef{箴 30:8} 你确实看见，这种知足并非为了其自身而被渴望，而是为了确保身体的健康，以及相称于人处境所需的居所与衣着的供应，使他能在所当生活的群体中显得得体，从而维持他们的尊重，并履行其职分。\par

13. 在这一切事物中，我们自身的福祉以及友谊促使我们为他人所求的益处，都是为其本身的缘故而可渴望；但对生活必需品的足够所需，通常却不是为了其本身，而是为了那两项更高益处而寻求的（若以正当方式寻求）。福祉在于拥有生命本身，以及心灵和身体的健康与健全。再者，友谊的要求，不应局限于\Emph{过于}狭窄的范围，因为它涵盖了一切应受爱戴与善意之情的人，尽管心向其中一些人更为自由地敞开，对另一些人则更为谨慎；是的，它甚至延及我们的仇敌，因为我们也被命令为他们祈祷。因此，在整个人类大家庭中，没有一个人不因共有人性之纽带而应受善意之情，儘管这情谊可能并非基于相互的爱而应得；\par

% structure: p0022 source=h2 target=subsection reason=relative-heading-rank; base=h2

\subsection{第七章}

因此，为这些事物，我们理当祈祷：若我们已拥有它们，是为能持守；若我们尚未拥有，是为能获得。\par

14. 难道这是全部吗？难道幸福生活单单在于这些好事吗？抑或真理教导我们，另有某事物比这一切更应优先？我们知道，必需品的足够所需，以及我们自身和朋友的安康，只要它们仅关涉此现世，与获得永生相比，都应被视为渣滓而弃置；因为即使身体健康，若心灵不将永恒之事置于现世之事之上，便不能被视为健全；而且，我们在时间中所度的生活，若不用于获取那使我们堪得永生的事物，便是虚度了。因此，一切作为有益且合宜的渴望之对象的事物，无疑都应参照那与天主同度、并从祂而来的唯一生活来看待。如此而行时，我们若爱天主，便是爱我们自己；同样，按照第二条大诫命，我们若在能力范围内，劝导他们同样爱天主，我们便是真正爱邻如己了。因此，我们爱天主，是因祂本身之所是；我们爱自己和邻人，则是为了祂的缘故。即使如此生活，我们也不要以为在那幸福生活中已稳固建立，仿佛再没有需要祈求的了。因为，当那唯一作为正确导向之生活的目标，仍为我们所欠缺时，我们怎能说现在正过着幸福生活呢？\par

% structure: p0025 source=h2 target=subsection reason=relative-heading-rank; base=h2

\subsection{第八章}

15. 那么，为何我们的欲望分散于许多事物，为何我们因害怕未按应祈祷的方式祈祷，而询问我们应该祈求什么，却不与圣咏作者一同说：\BibleQuote{有一事我曾求主，我仍要寻求：就是一生一世住在主的殿里，瞻仰主的美善，在祂的殿里求问}？因为在主的殿里，\BibleQuote{一生一世}的日子并非以相继来去而区分的日子：前一天的开始并非后一天的终结；而是在那地方，它们全都是同样无终的，因为由它们构成的生命本身没有终结。为了使我们获得这真福生命，那本身就是真福生命的主教导我们祈祷时，不要多言，仿佛我们是否被俯听取决于我们说话的流利，因为我们祈祷的对象是一位主所告诉我们的，\BibleQuote{在你们求祂以前，已知道你们需要什么}的主。\BibleRef{玛 6:7} 由此似乎令人惊讶，虽然祂禁止了\BibleQuote{多言}，但在我们求以前已知道我们需要什么的那一位，却用这样的话语勉励我们祈祷：\BibleQuote{人应当时常祈祷，不要灰心}；祂给我们讲了一个寡妇的比喻，这寡妇渴望在不义的法官面前为她对付仇人求得正义，并藉着恒切的恳求说服了这个不义的法官听取了她，这法官既不是因为顾及正义，也不是因为怜悯，而是被她令人疲惫的纠缠所胜；主这样教导我们，是要让我们明白，主天主，慈悲而正义的，是多么更加肯定地俯听我们不断地向祂祈祷，既然这位寡妇以她不间断的恳求胜过了不义邪恶法官的冷漠，而祂又是多么乐意和宽仁地满全那些祂知道已宽恕别人过犯之人的善愿，既然这位恳求者虽寻求向仇人报仇，却也得到了她所求的。\BibleRef{路 18:1} 主在另一个比喻中给了类似的教导，说到一个人有朋友行路来到他那里，他没有东西摆在他面前，想要向另一个朋友借三个饼（这或许象征同一性体的\OriginalTerm{天主圣三}{Trinity}的三个位格），发现他已经和家人一同睡着了，便藉着非常急切和纠缠的恳求把他叫醒，以至于他给了他所需数量的饼，这样做是出于避免进一步的搅扰，而非善意的想法；主愿意我们从这比喻明白，如果连一个睡着了的人，在被求他的人打扰睡眠后，也不得不给出，即使心不甘情不愿，那么那位从不睡眠，且唤醒我们使我们得以向祂祈求的主，会多么更慷慨地给予呢。\BibleRef{路 11:5}\par

16. 抱着同样的用意，祂接着说：\BibleQuote{你们求，必要给你们；你们找，必要找着；你们敲，必要给你们开。因为凡求的，就必得到；找的，就必找到；敲的，就必给他开。你们中间哪有为父亲的，儿子向他求饼，反而给他石头呢？或是求鱼，反而给他蛇呢？或者求鸡蛋，反而给他蝎子呢？你们纵然不善，尚且知道把好东西给你们的儿女，何况你们在天之父，岂不更将好东西赐与求祂的人吗？} 这里我们有与宗徒所称赞的那三样相对应的事物：\Emph{信德}由鱼来象征，要么是因为圣洗圣事所用的水元素，要么是因为它在这世界的惊涛骇浪中不受伤害——与鱼相对的是蛇，蛇以有毒的欺骗说服人不信天主；\Emph{望德}由蛋来象征，因为雏鸟的生命尚不在其中，但将要存在——看不见，却盼望得到，因为\BibleQuote{所希望的若已看见，就不是希望了}\BibleRef{罗 8:24}——与蛋相对的是蝎子，因为盼望永生的人忘尽背后的，伸手向前面的，因为他回头看是危险的；但需要提防蝎子，因为它尾巴上有东西，即尖刺和毒钩；\Emph{爱德}由饼来象征，因为\BibleQuote{其中最大的，是爱德}，而饼在一切食物中最有益处——与饼相对的是石头，因为铁石心肠拒绝实践爱德。不论这些象征是否有这意思，还是能找到其他更合适的意思，至少确定的是，那知道如何把好东西赐给子女的那位，催促我们要\BibleQuote{求、找、敲}。\par

17. 那位\BibleQuote{在我们求祂以前，已知道我们需要什么}的主，为何要我们这样做，这可能会困扰我们的心灵，如果我们不明白主我们的天主并非为了把我们的愿望告诉祂而要求我们祈求，因为对祂而言这不可能是未知的，而是为了藉着祈祷，在我们内通过恳求操练那渴望，使我们藉以领受祂所预备的恩赐。祂的恩赐非常伟大，但我们微小而狭隘，接受的能力有限。因此有话说：\BibleQuote{扩展你们的心胸，不要与不信者共负一轭}。\BibleRef{格后 6:13} 因为，按照我们信德的单纯、望德的确切、渴望的热切，我们将更丰盛地领受那无限伟大的事物；这事物\BibleQuote{眼所未见}，因为它不是色彩；\BibleQuote{耳所未闻}，因为它不是声音；\BibleQuote{人心所未想到的}，因为人心必须上升才能达到它。\BibleRef{格前 2:9}\par

% structure: p0029 source=h2 target=subsection reason=relative-heading-rank; base=h2

\subsection{第九章}

18. 当我们借着实践\OriginalTerm{信德}{faith}、\OriginalTerm{望德}{hope}和\OriginalTerm{爱德}{charity}而怀有不间断的渴望时，我们便是在\Quote{常常祈祷}了。但我们也按时按刻用言语向天主祈祷，好借着这些事物记号提醒自己，并使我们了解自己在这渴望上进到了什么程度，也更热切地激励自己获得其增长的力量。因为祈祷的效果将与说出祷词之前渴望的热切程度成正比。因此，\OriginalTerm{宗徒}{apostle}的话：\BibleQuote{不断祈祷}\BibleRef{得前 5:17}，除了\Quote{从不间断地从唯一能赐予福乐者那里渴望幸福的生命，而这种生命只能是\OriginalTerm{永生}{eternal life}}之外，还指别的什么呢？所以，让我们不断向主我们的天主渴望这永生吧；这样我们便是在不断祈祷。但我们在某些时刻，要从其他挂虑和俗务中收回心神——因为渴望本身在这些俗务中或多或少会冷却下来——转向祈祷之务，借我们的祷词提醒自己将注意力集中在所渴望之事上，以免那开始变凉的渴望变得完全冷却，并最终熄灭——若这火焰不更频繁地被煽动的话。因此，当同一位宗徒说：\BibleQuote{将你们的请求呈于天主前}\BibleRef{斐 4:6}，这话不应被理解为此举能让天主知道这些请求，因为祂无疑在祈求之前便已知道；而应理解为，这些请求应在天主面前借着耐心等待祂而向我们自己显明，而不是在人面前借着炫耀的敬拜。或者，也许它们也向那些在天主面前的天使显明，好让这些存在以某种方式将它们呈献给天主，为这些事求问祂，并可将他们听从祂命令而得知是祂旨意、且必须由他们按他们在那里所学到的职责来履行的事情，或明显或隐秘地带给我们；因为天使对\Person{多俾亚}说：\BibleQuote{现在，当你和你的儿媳妇撒辣祈祷时，我上达了你们的祈祷的记念于上主之前。}\BibleRef{多 12:12}\par

% structure: p0031 source=h2 target=subsection reason=relative-heading-rank; base=h2

\subsection{第十章}

19. 所以，若暇时甚多，且不阻碍本分所要求的其他善工和必要工作，那么长时间祈祷既非错误，也非无益。尽管即便在做这些工作时，正如我所说，我们也应通过怀有神圣的渴望而不停地祈祷。因为长时间祈祷并不像某些人所想的那样等同于\Quote{唠唠叨叨}地祈祷。词汇增多是一回事，渴望久燃是另一回事。因为论及主自己，经上记载祂彻夜祈祷\BibleRef{路 6:12}，并且当祂在极度痛苦中时，祂的祈祷更为长久；在这事上，那位在时空中作我们所需的代祷者，并永恒与圣父同在、俯听祈祷的那一位，岂不是给我们立了榜样吗？\par

20. 据说在\Place{埃及}的弟兄们有非常频繁的祈祷，但这些祈祷非常简短，并且是突如其来的、短促的，以免祈祷时必不可少的警醒振奋的注意力因长时间的修炼而消失或失去其敏锐。他们在这一点上已经很清楚地表明：正如注意力若不能持续很久就不应让它耗竭，同样，若注意力能持续，它也不应骤然中断。我们切不可在祈祷中\Quote{唠唠叨叨}，也不该在灵魂的热切注意力持续时抑制长久的祈祷。在祈祷中唠唠叨叨，就是为求一件必要的事而用大量多余的言词；而延长祈祷，是让内心对那位我们向其祈祷者保持持续的虔诚情感跳动。因为在大多数情况下，祈祷更多在于呼号而非言谈，在于泪水而非言辞。祂把我们的泪水放在祂眼前，我们的呼号对那位以圣言创造万物、并不需要人类言语的祂来说，并非隐藏。\par

% structure: p0034 source=h2 target=subsection reason=relative-heading-rank; base=h2

\subsection{第十一章}

21. 因此，对我们来说，言语是必要的，借助言语我们可以帮助自己思考并留意我们所祈求的，而不是把它们当作我们指望天主由此获得信息或被说服而顺从的手段。所以，当我们说：\BibleQuote{愿你的名受显扬}，我们是在提醒自己切望祂的名——它永远是圣的——也能在人间被尊为圣，也就是说，不被轻慢；这带来的利益不是给天主的，而是给人的。当我们说：\BibleQuote{愿你的国来临}，这国度无论我们愿意与否都必定来临，我们借这些话激起我们对那国度的渴望，好使它临到我们，我们也堪当在其中为王。当我们说：\BibleQuote{愿你的旨意奉行在人间，如同在天上}，我们是在为自己祈求：求祂赐给我们服从的恩宠，使祂的旨意能借着我们而承行，如同在天上由祂的天使们承行一样。当我们说：\BibleQuote{求你今天赏给我们日用的食粮}，\Quote{今天}一词指的是现在的时候；在这时候我们祈求的，或者是我们在前面所讲过的现世恩惠方面的充足供应（\Quote{食粮}用来代表全部那些恩惠，因为它是其中如此重要的一部分），或者是信徒的圣事，它在现在的时候是必需的，但必需是为了获得不是今世而是永世的幸福。当我们说：\BibleQuote{求你宽恕我们的罪债，如同我们也宽恕得罪我们的人}，我们提醒自己应当祈求什么，以及应当做什么，才能堪当领受所祈求的。当我们说：\BibleQuote{不要让我们陷于诱惑}，我们提醒自己寻求的是：不要因失去天主的助佑，或则被诱入圈套而同意，或则被迫屈服于诱惑。当我们说：\BibleQuote{但救我们免于凶恶}，我们提醒自己，我们还没有享有那不再经历任何凶恶的幸福境界。这个求恩，位于天主经的末位，是如此的包罗万象，以至于一个基督徒，无论处于何种苦难中，都可以用它来发出悲叹、倾流眼泪——可以以这个求恩开始祈祷，继续它，并用它结束祈祷。因为我们必须借着使用这些话，使它们所表示的意思常存于我们的记忆之中。\par

% structure: p0036 source=h2 target=subsection reason=relative-heading-rank; base=h2

\subsection{第十二章}

22. 因为无论我们可能说的其他任何话语——无论是祈祷者的愿望在话语之先，并用话语来明确表达其请求，或是愿望在话语之后，并专注于话语以便增长热忱——如果我们正确地祈祷，并按照我们的需要，我们就无非是在说天主经里已经包含的内容。任何人在祈祷中所说的，若在那福音的祈祷中找不到一席之地，他的祈祷即使不算不合法的，至少也不是属神的；我不知道属血肉的祈祷怎能是合法的，因为那由圣神而重生的人，只应以属神的方式祈祷。例如，当一个人祈求：\BibleQuote{愿你在万民中受光荣，如同你已在我们中受了光荣}，以及\BibleQuote{愿你的先知们被显为忠信}，\BibleRef{德 36:4}，他所求的岂不是：\BibleQuote{愿你的名受显扬}？当一个人说：\BibleQuote{万军的上主，求你使我们回头，显示你的慈颜，好拯救我们}，他岂不是在说：\BibleQuote{愿你的国来临}？当一个人说：\BibleQuote{求你引导我的脚步履行你的训令，不要允许任何邪恶主宰我}，他岂不是在说：\BibleQuote{愿你的旨意奉行在人间，如同在天上}？当一个人说：\BibleQuote{求你使我也不穷也不富}，\BibleRef{箴 30:8}，这岂不是：\BibleQuote{求你今天赏给我们日用的食粮}？当一个人说：\BibleQuote{上主，求你记起达味和他所有的善行}，或：\BibleQuote{上主，如果我作了这事，如果我的双手沾了不义，如果我以恶报善，就请……}，这岂不是：\BibleQuote{求你宽恕我们的罪债，如同我们也宽恕得罪我们的人}？当一个人说：\BibleQuote{求你除去我肚腹的贪欲，不要让我陷于肉情的欲望}，\BibleRef{德 23:6}，这岂不是：\BibleQuote{不要让我们陷于诱惑}？当一个人说：\BibleQuote{我的天主，求你救我脱离我的仇敌，求你使我摆脱那起来攻击我的人}，这岂不是：\BibleQuote{但救我们免于凶恶}？你若遍阅所有圣善祈祷的言辞，我相信，你找不到任何不能包括并归结在天主经的各求恩中的话语。因此，在祈祷时，我们可以自由使用不同的话语，范围可以任意宽广，但我们必须祈求相同的事物；在这点上，我们没有选择。\par

23. 这些事，我们理当毫无犹豫地为自己、为朋友、为陌生人——是的，甚至为仇敌祈求；虽然祈祷者心中对某人和某人的渴望，根据关系的近远，在性质上或强度上会有差别。但是，谁若在祈祷中说这样的话，比如：\Quote{主啊，增多我的财富}；或者，\Quote{求你给我像你赐给这人或那人的那么多家产}；或者，\Quote{增加我的尊荣，让我在这世界上权势显赫、名声出众}，或类似这样的其他话，而且仅仅是出于对这些事物的贪欲而求，并不是为了通过它们按照天主的旨意有益于人，我不认为他能找到天主经的哪一部分，可以正好把这些请求放入其中。因此，我们至少应该羞于祈求这些事，如果我们会以渴望它们为耻的话。然而，如果我们连渴望它们也感到羞耻，却感到自己被渴望所胜，那么，向那位我们向祂说\BibleQuote{但救我们免于凶恶}的，祈求从这渴望的灾病中解救出来，岂不是好得多吗！\par

% structure: p0039 source=h2 target=subsection reason=relative-heading-rank; base=h2

\subsection{第十三章}

24. 你若我没弄错，现在已得到两个问题的答案——如果你要祈祷，你应当是怎样的人，以及在祈祷中你应求什么；而这答案并非来自我的教导，而是出于那位俯就教导我们众人的。应追求幸福的生活，并且这应从主天主那里求得。许多人在讨论真正的幸福在于何处时给出了许多不同的答案；但我们为何要求助于许多师傅，或考虑这问题的许多答案呢？圣经中已简短而真实地陈述了：\BibleQuote{那奉主为天主的百姓，是有福的。}为使我们得列入这百姓中，且能达到看见他并永远与他同住，经上说：\BibleQuote{这训令的目的就是爱德，由纯洁的心、善良的良心和真诚的信德所发出的爱德。}\BibleRef{弟前 1:5}在这三者中，用望德取代了善良的良心。因此，信德、望德和爱德引领祈祷的人，\Emph{即}那相信、盼望和渴望的人，归向天主，并借着学习\WorkTitle{天主经}而知道该向主求什么。守斋、戒绝肉欲的享乐而不损害健康，尤其是时常施舍，对于祈祷是大有裨益的；这样我们便能说：\BibleQuote{在我患难之日我寻求主，夜间向主举手，我并未受骗。}因为天主是神，是不能触摸的，若不借着善工，怎能用双手寻找他呢？\par

% structure: p0041 source=h2 target=subsection reason=relative-heading-rank; base=h2

\subsection{第14章}

25. 也许你还要问，为什么宗徒说：\BibleQuote{我们不知道我们如何祈求才对，}\BibleRef{罗 8:26}因为如果说他或他写信的对象都不认识天主经，那简直是不可信的。他不可能轻率地或虚假地说这话；那么，我们猜想他这样说的理由是什么呢？难道不是因为，世上的磨难和困苦大多有益，或是为医治骄傲的膨胀，或是为考验和锻练忍耐，在如此的考验和训练之后，将有更大的赏报存留，或是为惩罚和消除某些罪；而我们却不知道这些磨难可能有什么益处，只愿脱离一切患难。对于这种无知，宗徒指出，连他自己也不是外人（除非，他明知应如何祈求，却仍这样说），当时，为免他因启示的高超而过于高举自己，就有一根刺加在他身上，就是撒殚的使者来打击他；为此事，他既确实不知该如何祈求，就三次求主使这刺离开他。最后他得到天主的答复，说明了为什么如此伟大的人所求的没被应允，以及为什么这事成就更为有益：\BibleQuote{我的恩宠为你足够了，因为我的德能在软弱中才全显出来。}\BibleRef{格后 12:7}\par

26. 因此，对于磨难，我们不知道如何祈求才对，因为它们可能对我们有益或有害；可是，由于它们是艰难痛苦的，且与我们软弱本性的自然感觉相悖，我们便怀着人类共有的愿望，祈求它们离开我们。但我们应当这样服从主天主的旨意：即使他不除去那些磨难，我们也不要想自己被他忽视，反而要在忍耐邪恶时，希望分享更大的善，因为他的德能正是在我们的软弱中才得完全。有时，天主在愤怒中应允了急躁祈求者的请求，正如他出于慈悲拒绝了宗徒的请求。因为我们读过以色列子民所求的，以及他们怎样求而获得了所求的事；但当他们的愿望得到满足时，他们的急躁却受到了严厉的处分。他又按他们所求的，给他们立了一位合他们心意的君王，如经上记载的，却非合他心意的。他也应允了魔鬼所求的，让他的仆人，那要受考验的，受试探。他也应允了不洁之神的请求，当他们求他把那群污鬼赶入大猪群中时。\BibleRef{路 8:32}这些事被记载下来，为的是不使任何人因自己急切祈求的某种东西得到应允而自视过高，因为那东西不获得对他更为有益；或是为不使任何人，在祈祷未蒙垂听时，感到沮丧，对天主的慈爱绝望，因为或许他所求的，一旦获得，反而会使他更加痛苦，或是被顺境的腐蚀影响而彻底毁灭。因此，对于这类事情，我们不知道如何祈求才对。所以，如果某事的发生与我们的祈祷相违，我们就该忍耐地承受失望，在一切事上感谢天主，丝毫不怀疑天主的旨意而非我们的旨意应当承行才是正确的。关于这一点，中保已给我们立了榜样，因为他在说了：\BibleQuote{父啊！若是可能，就让这杯离开我罢！但不要照我，而要照你所愿意的。}\BibleRef{玛 26:39}因此，因一人的服从，众人都成了义人，这并非没有理由的。\BibleRef{罗 5:19}\par

27. 但是，无论谁向主渴望那\BibleQuote{一件事}，并寻求它，就是以确定和信心祈求，并且不怕得到后对自己有害，因为若没有它，通过正确方式祈求而得到的任何其他事物都对他毫无益处。那件事就是唯一真实且唯一有福的生命，在其中，身体和灵魂不死不灭，我们能够永远瞻仰主的喜乐。所有其他的事物都是为了这一件事而渴望，并毫无不当之处地祈求。因为凡拥有它的人，就会拥有他所希望的一切，并且绝不可能希望与它一起得到任何不相称的事物。因为在它里面有生命的泉源，我们现在必须在祈祷中渴望着它，只要我们在望德中，还没有看见我们所期望的，在祂翅膀的荫庇下投靠，在祂面前有我们一切的渴望，好让我们能以祂殿里的肥甘得以饱足，喝祂乐河的水；因为与祂同在的是生命的泉源，在祂的光明中我们必得见光明，当我们的渴望因美好事物而满足时，当不再有什么需要叹息寻求，而是我们在欢欣中拥有一切时。与此同时，因为这一恩福不是别的，正是\BibleQuote{超越一切理解的平安}\BibleRef{斐 4:7}，即使我们在祈祷中祈求它，我们也不知道该如何以宜于的方式祈求。因为实际上我们不能将它凭其真实状况呈现在心思中，我们不知道它，但我们心思中可能出现的任何形象，我们都拒绝、否认并谴责；我们知道它并非我们所寻求的，尽管我们还不足以知道如何界定我们所寻求的。\par

% structure: p0045 source=h2 target=subsection reason=relative-heading-rank; base=h2

\subsection{第十五章}

28. 因此，在我们内有一种可说是\OriginalTerm{有所知的不知}{learned ignorance}——这是我们从那扶助我们软弱的圣神所学得的一种不知。因为宗徒说过：\BibleQuote{如果我们希望那我们所未见的，我们就耐心等待它}，在同一段中他又加上：\BibleQuote{同样，圣神也扶助我们的软弱：因为我们不知道我们该如何以宜于的方式祈求，但圣神本身却亲自以无可言喻的叹息，代我们转求。那洞悉人心的，知道圣神的思念是什么，因为祂是按照天主的旨意代圣徒转求。}\BibleRef{罗 8:25}这不应被理解为：天主圣神，就是在天主圣三中不变的天主，与圣父及圣子同为一体，自己像一位非神圣者那样为圣徒转求；因为经上说\BibleQuote{祂代圣徒转求}意思是祂使圣徒能够转求，如同经上另一处说：\BibleQuote{上主你们的天主试探你们，为要知道你们是否爱祂}\BibleRef{申 12:3}，即：为使你们知道。所以，祂使圣徒以无可言喻的叹息转求，就是当祂激发他们对那尚属未知的伟大恩福的渴望时，我们为此耐心等待。因为所渴望的若是未知，何以用言语表达？因为若是完全未知，就不会渴望；另一方面，若是已见到，就不会渴望，也不会叹息寻求了。\par

% structure: p0047 source=h2 target=subsection reason=relative-heading-rank; base=h2

\subsection{第十六章}

29. 思考了这一切，以及主在此事上可能启示给你的、我或未想到或在此陈述过于冗长的其他事情之后，你要在祈祷中努力战胜这个世界：怀着望德祈祷，怀着信德祈祷，怀着爱德祈祷，恳切而忍耐地祈祷，如同属于基督的寡妇那样祈祷。因为尽管祈祷，如祂所教导的，是所有肢体——即所有信祂并与祂身体结合的人——的职责，但圣经似乎特别嘱咐寡妇更勤奋地专心祈祷。圣经中尊荣地提到两个名叫亚纳的女子——一位是厄耳卡纳的妻子，圣撒慕尔的母亲；另一位则是那位在圣婴尚为襁褓时认出至圣者来的寡妇。前者虽已婚，却因没有儿子而满心忧伤地祈祷，破碎心灵；她获得了撒慕尔，并将他献于上主，因为她在为他祈求时曾许下此愿。\EditorialNote{撒慕尔纪上一章}不过，她的祈求很难归入\WorkTitle{天主经}中的哪一项祈求，除非是最后一项：\BibleQuote{救我们免于凶恶}，因为当时被认为是凶恶的，是已婚而无后嗣作为婚姻的果实。然而，请看关于另一位寡妇亚纳的记载：她\BibleQuote{总不离开圣殿，以斋戒和祈祷，日夜事奉天主。}\BibleRef{路 2:36}同样，宗徒也在前面引用过的话中说：\BibleQuote{那真正作寡妇的，孤独无靠，仰望天主，日夜在哀求和祈祷中坚持不懈；}\BibleRef{弟前 5:5}而主在劝人应当不断祈祷、不要灰心时，也提到一位寡妇，她以坚持不懈的缠扰，使一个不义且邪恶、既不敬畏天主也不尊重人的判官受理了她的案件。寡妇更加殷勤祈祷的责任有多么重大，从以下事实就可以清楚看出：即用于劝勉所有人恳切祈祷的榜样，正是取自她们中间。\par

30. 那么，究竟是什么使这工作特别适合寡妇，岂不是她们丧偶且孤寂的状况吗？所以，谁若明白自己在这世上，只要身为朝圣者远离他的主，便是丧偶且孤寂的，他必用心将自己的守寡，可以这么说，交托与他的天主作为盾牌，在持续不断而极为热切的祈祷中。所以，你要以基督的寡妇的身份祈祷，尚未得见那位你恳求帮助的主。虽然你非常富有，仍要以贫穷人的身份祈祷，因为你尚未拥有将来世界的真正财富，在那里你无需害怕损失。即使你有儿孙和众多家人，仍要像我已说过的那样，以孤寂者的身份祈祷，因为对于所有现世的恩惠，我们并无把握它们必存留为我们的安慰，直到今生终结。如果你寻求并体味天上的事，你便渴望永恒而确实的事物；而只要你尚未拥有它们，就当视自己为孤寂的，纵然所有家人都蒙保全，并按你所愿地生活。而你若如此行动，你的极为虔诚的儿媳，以及在你的照管下安居的其他圣洁的寡妇和贞女，必定会效法你的榜样；因为你越是以虔诚的方式治理家庭，便越有义务坚持热切祈祷，除了信仰所要求的之外，不可让自己卷入世俗的事务。\par

31. 务必记得为我恳切祈祷。我不愿你们因我肩负的危险职务而如此恭敬，以致于不给予我知道自己所需的帮助。基督的家庭曾为\Person{伯多禄}和\Person{保禄}祈祷。我欢喜你们在祂的家庭中；而我，比\Person{伯多禄}和\Person{保禄}更无可比拟地需要弟兄们祈祷的帮助。你们要以圣洁的竞争在祈祷上彼此效仿，同心合意，因为你们不是与彼此角力，而是与魔鬼角力，他是所有圣人的公敌。\BibleQuote{用斋戒、守夜和一切克苦身体，祈祷大受助益。} \BibleRef{多 12:8} 各人尽己所能；一人所不能做的，若她爱在另一人身上那仅因个人无能而无法做的，便由那能做的人代做；因此，能力较小者不可拦阻能力较大者，能力较大者也不可过分催迫能力较小者。因为你们的良心对天主负责；彼此之间只应有彼此相爱。愿主，祂能成就远超我们所求所想的事，俯听你们的祈祷。 \BibleRef{弗 3:20}\par

\SourceRecord{Translated by J.G. Cunningham.；Nicene and Post-Nicene Fathers, First Series；Vol. 1.；Edited by Philip Schaff.；Buffalo, NY: Christian Literature Publishing Co.,；1887.；<http://www.newadvent.org/fathers/1102130.htm>.}

% structure: p0001 source=h1 target=section reason=page-title

\section{第131封信（公元412年）}

\Emph{致他最杰出的女儿、高贵而理应显赫的\Person{普罗巴}女士，\Person{奥斯定}在主内致以问候。}\par

你说得对，灵魂寓居于可朽坏的身体内，因与尘世接触的程度而受牵制，且不知何故被这重担压弯压碎，以致其欲望和思想更容易向下奔向繁多事物，而非向上归向那一位。因为圣经也这样说：\BibleQuote{这必朽坏的肉身，重压着灵魂；这属土的帐幕，使多虑的心神滞重。}\BibleRef{智 9:15} 但我们的救主，曾以他医治的圣言使福音中那个被病魔缠身十八年的女人直起身来\BibleRef{路 13:11}（她的情况或许正是灵性疾病的一个预像），他正是为此而来，为使基督徒不致徒然听到\Quote{请举心向上}的呼唤，而能真实地回答：\Quote{我们举心向主}。有鉴于此，你做得很好：因着对来世的盼望，便视此世的邪恶为易于忍受。因为如此，这些邪恶若善加利用，便成为祝福，因为它们不激起我们对今世的贪恋，反而锻炼我们的忍耐；关于这点，宗徒说：\BibleQuote{我们晓得万事都互相效力，叫爱天主的人得益处。}\BibleRef{罗 8:28}他说\Emph{一切}——因此，不仅包括那些因愉悦而渴望的，也包括那些因痛苦而躲避的；因为我们对前者既加以接受而不被其冲昏头脑，对后者则加以承受而不被其压垮，且在一切事上，按着天主的命令，向那一位感恩，我们曾对他说：\BibleQuote{我必要时时赞美主，对他的赞颂常在我口；}以及\BibleQuote{为叫我能学习你的法度，受苦遭难于我确有好处。} 事实上，最高贵的女士，倘若这不可靠的顺境之平静总是向我们微笑，人的灵魂便永远不会驶向那真正安全的港湾。因此，在向阁下致上应有的敬意问候，并对你如此虔敬地关心我的福祉表示感谢之际，我恳求主，愿他赐予你来世的赏报，并在此生中给予安慰；我也将自己托付于你们所有人心中的爱和祈祷，因着信德，基督寓居在你们心中。\par

（\EditorialNote{另一个手笔}）愿真实而信实的天主真正安慰你的心灵并保守你的健康，我最杰出的女儿、高贵的女士，理应显赫。\par

\SourceRecord{Translated by J.G. Cunningham.；Nicene and Post-Nicene Fathers, First Series；Vol. 1.；Edited by Philip Schaff.；Buffalo, NY: Christian Literature Publishing Co.,；1887.；<http://www.newadvent.org/fathers/1102131.htm>.}

% structure: p0001 source=h1 target=section reason=page-title

\section{书信132（主历412年）}

\Emph{致\Person{沃卢西亚努斯}，我尊贵的主上、至为公义杰出的儿子，主教\Person{奥斯定}在主内致以问候。}\par

在我为你在此世与在基督内的福祉所怀的渴望上，我或许甚至不亚于你虔诚母亲的祈祷。因此，我在以与你身份相称的敬意回复你的问候时，恳切地劝勉你，不要吝于专心研读那些真实无疑乃神圣的\OriginalTerm{著作}{Writings}。因为它们是真实而坚固的真理，不以矫揉的文辞赢取人心，也不发出谄媚之声，徒有虚浮而不可靠的声响。它们深深吸引那饥渴于事物而非言辞的人；它们首先使那将要被它们从恐惧中确保安全的人大为惊惶。我特别劝你阅读宗徒们的著作，因为从它们你将受到激励，去熟悉众先知，宗徒们正是引用了他们的证言。如果你在阅读或默想所读内容时，有任何问题，其解答似乎需要我，请写信给我，我也好写信回复。因为，有主相助，我或许能以此方式比亲自与你交谈这些事更好地服务你，部分原因在于，由于我们事务不同，你不可能恰好在我有空时也有空闲，更尤其是因为那些人大都不适于此类讨论，且更乐于言辞之战，而不喜爱知识的清晰明了，他们令人难以推却地闯入；而书写下来的东西则随时听候读者在自己闲暇时的差遣，且与那随己意拿起或放下的东西相伴，绝不会成为重担。\par

\SourceRecord{Translated by J.G. Cunningham.；Nicene and Post-Nicene Fathers, First Series；Vol. 1.；Edited by Philip Schaff.；Buffalo, NY: Christian Literature Publishing Co.,；1887.；<http://www.newadvent.org/fathers/1102132.htm>.}

% structure: p0001 source=h1 target=section reason=page-title

\section{第133封信（公元412年）}

\Emph{致\Person{马凯利努斯}，我尊贵的主、公正杰出者、我极亲爱的儿子，\Person{奥斯定}在主内向您致意。}\par

1. 我得知，在\Place{希波}地区属于多纳图斯派的切尔库姆凯利翁派及其圣职人员，因他们的行为而被公共秩序的守护者送审，并由阁下进行了审讯；他们中的大多数人已承认，司铎\Person{雷斯提图图斯}遭受的暴力致死，以及他们殴打另一位天主教司铎\Person{依诺森爵}，并挖出\Person{依诺森爵}的眼睛、切断他的手指，他们皆有参与。这消息使我陷入极深的焦虑之中，恐怕阁下会依照法律，判定他们应受的惩罚不轻于让他们亲身遭受他们对别人所施加的同样伤害。因此，我写这封信，凭您对基督的信德，并藉着主基督本人的仁慈，恳求您切勿如此行，也勿容许此事发生。因为尽管我们或许可以对那些罪犯的处决保持沉默——他们并非由我们提出控告，而是由那些负责维护公共安宁的人们提起公诉而受审——但我们不希望天主的仆人所受的苦难，通过完全相同的伤害以报复的方式得到报复。当然，我们并不反对剥夺这些恶人继续作恶的自由；但我们所渴望的是，正义得以满足，却不用夺去他们的生命或伤残他们身体的任何部分，并且，通过符合法律的强制措施，使他们从狂乱的疯狂转向心智健全之人的宁静，或者迫使他们放弃有害的暴力，转而从事某种有益的劳作。这确被称为刑罚判决；但有谁看不出，当对野蛮暴行的肆无忌惮加以约束，并保留那些适合引发悔改的补救方法时，这种管教更应称为一种益处，而非报复性的惩罚呢？\par

2. 基督徒法官啊，您当履行慈父的职责；让您对他们罪行的义愤受到人性考量的节制；不要因他们罪恶行为的残暴而被激怒去满足复仇的激情，而应被这些行为在他们自己灵魂上所造成的创伤所触动，去行使治愈他们的愿望。不要现在失去您在审讯过程中所保持的那种慈父般的关怀，您当初正是借此使他们供认了如此骇人听闻的罪行，而所用的方式并非将他们绑在刑架上，并非用铁爪撕裂他们的皮肉，并非用火焰灼烧他们，而是用棍棒击打——这是教师常用于管教的方式，也是父母责罚子女时所用的方式，主教们在其判决中亦常使用。因此，对于您当初以宽仁查明的罪行，如今切勿施以极刑。查究时严苛的必要性，大于施罚时；因为即使是最温和的人，在搜寻隐藏的罪行时也会用尽勤勉和严苛，以便找到他们可以施予怜悯的对象。因此，在侦查时通常需要更为严厉，以便等到罪行被揭露后，能有宽仁施展的空间。因为一切善工都喜爱被置于光中，并非为了从人获得荣耀，而是如主所说：\BibleQuote{你们的光也当在人前照耀，好使他们看见你们的善行，光荣你们在天之父。}\BibleRef{玛 5:16} 并且，为了同样的缘故，宗徒不满足于仅仅劝勉我们实行温和，还命令我们要将其彰显出来：\BibleQuote{你们的宽仁}，他说，\BibleQuote{当使众人知道；}\BibleRef{斐 4:5} 并在另一处说：\BibleQuote{对众人表示十分温和。}\BibleRef{铎 3:2} 因此，圣\Person{达味}那极其显著的宽容——当他仁慈地饶恕了落入他手中的仇敌时——若没有他能够采取相反行动的明显能力，就不会如此昭彰了。\BibleRef{撒上 24:7} 所以，不要让施行报复的权力激发您的严苛，既然审讯罪犯的必要性未能使您放下宽仁。如今罪行既已查明，切莫召唤刽子手，当初您在查寻真相时既已克制未曾召来刑讯者。\par

3. 总之，您被派遣至此是为了教会的益处。我郑重声明，我所建议的，对于天主教会的利益是有益的，或者，为了不让人以为我越出了自己职权的界限，我申明，这对于属于\Place{希波}教区的教会是有益的。倘若您不听从我作为朋友恳求这一恩惠，那么请听从我作为主教提出的建议；不过，既然我是在对一位基督徒说话，尤其是在此类事情上，我这样说并不算冒昧——我尊贵、公正杰出、极亲爱的儿子啊，您理当听从我作为主教以权威所发的命令。我知道，与教会事务有关的法律案件的主要职责已托付于阁下，但我相信这一特殊案件属于最显赫可敬的总督大人管辖，因此我也给他写了一封信，恳请您不要拒绝将此信转交给他，或者必要时，向他推荐此信；我恳请你们两位都不要因为我们以切切之情代祷、建议或忧虑而见怪。请不要让天主教的天主仆人们的苦难——这些苦难本应为软弱者的灵性建树带来助益——因对施害者的同等报复而受玷污，而是要调和正义的严苛，作为教会的子女，你们应致力于彰显你们自己的信德以及你们慈母的宽仁。\par

愿全能的天主以一切美善丰盈地赐予阁下，我尊贵、公正杰出、极亲爱的主和儿子！\par

\SourceRecord{Translated by J.G. Cunningham.；Nicene and Post-Nicene Fathers, First Series；Vol. 1.；Edited by Philip Schaff.；Buffalo, NY: Christian Literature Publishing Co.,；1887.；<http://www.newadvent.org/fathers/1102133.htm>.}

% structure: p0001 source=h1 target=section reason=page-title

\section{第135封信（公元412年）}

\Emph{致主教\Person{奥古斯丁}，我真正圣善的主教大人，及应受尊敬的父亲，\Signature{伏鲁西亚努斯敬启}}\par

1. 良善正直的典范啊，您要求我向您请教，就我阅读中遇到的一些晦涩经文请求指导。 我遵命领受这一善意，甘愿受您的教导，并承认古人箴言的权威：\Quote{活到老，学到老。} 这则箴言的作者理当不设任何界限或终点来限制我们对智慧的追求；因为真理，隐匿于其原始原理之中，对于接近它的人，从不完全揭示，使之全然知晓。 因此，我真正圣善的主教大人，及应受尊敬的父亲，我认为值得向您转述一场最近在我们之间开展的对话的要点。 我出席了一次朋友们的聚会，会上提出了许多意见，各随其性情与兴趣。 然而，我们的谈话主要围绕着修辞学中讨论合宜布局的那部分。 我与熟谙此道者交谈，因为您不久前正是这些学问的教师。 随后讨论了修辞学中的\Quote{创意}，其本质、所需何等胆识、系统布局所涉何等巨大辛劳、隐喻的魅力、例证的优美，以及运用与所论题材之特性与本质相合的修饰语的能力。 另一些人则偏爱地盛赞诗人的艺术。 您对雄辩术的这一部分也未曾忽视或轻视。 我们可以恰如其分地将诗人的诗句应用于您：\Quote{常春藤与酬报您凯旋的月桂枝交织。} 因此，我们谈论了巧妙布局赋予诗篇的润饰，隐喻之美，以及精良比喻的崇高；接着我们谈论了流畅顺滑的诗律，以及，容我如此说，诗行中停顿的和谐变化。 谈话接着转向了一个您极为熟知的主题，即您自己曾以\Person{亚里士多德}和\Person{伊索克拉底}的方式所珍爱的哲学。 我们问到，吕克昂学园的哲学家、学园派各式各样无休止的怀疑、柱廊学派的论辩者、自然哲学家的发现、伊壁鸠鲁学派的自我放纵，各自取得了什么成就；他们之间无限热情地互相争辩的结果又是什么，以及在他们认定真理可知之后，真理又如何比以往更加隐晦不明。\par

2. 正当我们继续这些话题时，众人中的一位开口说：\Quote{我们中间有谁如此通晓基督教所教导的智慧，足以解开纠缠我的疑惑，并借着真理或或然性所要求的论证，使我犹豫不决地接受其教导变得坚定？} 我们全都惊愕得哑口无言。 然后，他自发地说出这些话：\Quote{我很想知道，那位主宰与统御世界的主，是否确实充满了童贞女的子宫——他的母亲是否承受了长达十个月的疲惫，并且，她虽为童贞女，如期分娩了她的孩子，之后仍保持其童贞无玷？} 他还补充了其他陈述：\Quote{那位宇宙几乎无法容纳的主，竟隐藏于一个号哭的婴孩的弱小身体内；他承受童年岁月，成长，被确立于成年男子的刚毅中；这位治御者如此长久地放逐于自己的居所之外，而整个世界的照料全交托在一个尺寸微小的身体上。 再者，他入睡，进食以养身，经历凡人一切感受。 如此伟大的威严也未曾以足够充分的证据彰显；因为若你想到别的曾行这些奇事的人，驱逐魔鬼、治愈病者、复活死者，对天主而言不过些小之事。} 我们制止了他继续提出此类问题，聚会随即解散，我们将此事诉诸我们本身经验以外的宝贵裁决，以免因过于轻率地窥探隐秘之事，使本属无意的错误变得应受责备。\par

您已然听闻，一切尊荣堪当之人啊，我们对无知的坦言；您明了我们向您请求的是什么。 您的声望关系着我们能否获得这些问题的答案。 对于别的司铎，无知或许可以在不损害信仰的前提下被容忍；但当谈及\Person{奥古斯丁}主教时，凡我们发现他所不知的，便不属于基督信仰的体系。 愿至高天主护佑您可敬的阁下，我真正圣善且应受尊敬的主教大人！\par

\SourceRecord{Translated by J.G. Cunningham.；Nicene and Post-Nicene Fathers, First Series；Vol. 1.；Edited by Philip Schaff.；Buffalo, NY: Christian Literature Publishing Co.,；1887.；<http://www.newadvent.org/fathers/1102135.htm>.}

% structure: p0001 source=h1 target=section reason=page-title

\section{第136封信（公元412年）}

\Emph{致\Person{奥思定}，我最可敬的主，并堪受我一切事奉的绝无仅有之父台，\Person{玛尔凯利努斯}谨致问候。}\par

1. 高贵的\Person{沃卢西安}读了我圣善主教的来信，在我的恳切请求下，他向更多人朗读了令我钦赏的句子，因为如同出自您笔下的其他一切，它们确实值得赞赏。这封信洋溢着谦逊的精神，并充满了神圣口才的恩宠，轻易地讨得了读者的欢心。尤其令我高兴的是您竭力扶持一位犹豫不定者的脚步，通过劝导他立下善志。因为我每天都与这同一个人进行讨论，只要我的能力——更确切地说，我的才拙学浅——还能允许。受他虔诚母亲的热切恳求所感动，我费心经常去探望他，而他也不时好心回访。但接到尊贵的您写来的这封信后，虽然他因这座城市中颇多的一类人的影响而一直未能坚定信仰真天主，他却深受触动，正如他亲口告诉我的，他惟恐信函过于冗长，才没有向您一一陈述他对基督信仰的所有疑虑。然而他非常恳切地请您解释某些问题，以洗练精准的风格，清晰而辉煌地展现了罗马雄辩的才华，相信您自己也会认为值得称许。他提交给您的问题，在争论中确已陈滥，而来自同一源头的狡诈、以非难主的\OriginalTerm{降生成人}{incarnation}来攻击的技俩，也是众所周知的。但我深信您无论写下什么答复，都将对极多的人有益，因我向您提出请求，即使在对这个问题的答复中，也请屈尊针对他们的谬说——声称主在祂的工程中没有做过任何超越凡人之所能的事——给予严密彻底的回应。他们惯于搬出他们的\Person{阿波罗纽斯}和\Person{阿普列尤斯}，以及其他一些自称精于魔法的人，声称他们的奇迹比主的更大。\par

2. 上述高贵的\Person{沃卢西安}还在众人面前宣称，如果不是如我已经说过的，他写信时特意追求简短，还有许多其他问题原可合理地附加到他已提出的问题上。然而，尽管他没有写下来，他并未在沉默中放过这些问题。因为他惯常说，即使有人现在给他讲明了主降生成人的合理依据，仍然很难给出令人满意的理由，来说明为什么这位天主——据断言也是\OriginalTerm{旧约}{Old Testament}的天主——在拒绝了旧时的祭献后，又喜爱新的祭献。因为他声称，只有先前被证明行事不正的，才能纠正；而先前行事正直的，则丝毫不得更改。他说，正直而行的事，不能更改而不犯错误，尤其是因为这种改变可能会使天主蒙受反复无常的指责。他提出的另一点反对意见是，基督徒的教义和宣讲，绝不符合公民的义务和权利；因为，举一个常被引用的例子来说，在其诫律中我们看到：\BibleQuote{不要以恶报恶}\BibleRef{罗 12:17}，以及\BibleQuote{有人打你面颊的一边，也把另一边转给他；有人拿去你的内衣，连外衣也让他拿去；有人强迫你走一里路，你就同他走两里。}\BibleRef{玛 5:39}——所有这些他断言都违背公民的义务和权利。因为谁肯容忍被敌人夺去任何东西，或者不去报复那劫掠罗马行省的入侵者所施加的战争恶行？正如您所理解的，其他诫命也面临类似的反对。\Person{沃卢西安}认为，所有这些问题都可加在前述问题上，尤其是在（虽然他对此点保持沉默）大部分信奉基督宗教的皇帝治理下，共和国遭受了极重的大灾难。\par

3. 因此，正如您的恩宠屈尊同我一起承认的，应该以一份全面、彻底且明晰的答复来回应所有这些疑难（因为您圣善的答复必受欢迎，无疑将被许多人传阅）；尤其是因为，在这场讨论进行之际，有一位在\Place{希波}地区的显赫领主和业主在场，他以讽刺的口吻说了一些恭维您圣善的话，并断言他自己在探究这些问题时根本未曾感到满意。\par

因此，我并未忘记您的许诺，而是坚持您实现它，恳请您就所提出的问题撰文，这些著作将对教会产生难以置信的裨益，尤其在当前的时候。\par

\SourceRecord{Translated by J.G. Cunningham.；Nicene and Post-Nicene Fathers, First Series；Vol. 1.；Edited by Philip Schaff.；Buffalo, NY: Christian Literature Publishing Co.,；1887.；<http://www.newadvent.org/fathers/1102136.htm>.}

% structure: p0001 source=h1 target=section reason=page-title

\section{书信一三七（主后412年）}

\Emph{致我最杰出的儿子、尊贵且公义卓越的\Person{沃鲁西亚努斯}大人，\Signature{奥古斯丁}在主内致候。}\par

% structure: p0003 source=h2 target=subsection reason=relative-heading-rank; base=h2

\subsection{第一章}

我读了你的来信，信中扼要记述了一次可称道的谈话，值得赞赏的简洁。我自感有责任回复，并避免以任何迟延为借口；因我恰好有一段短暂空闲，可免于处理他人事务。在此期间，我也推迟了向我的书记员口述那些原本打算利用这段闲暇来完成的文字，因为我认为，若延迟回答那些问题——这些问题本是我自己敦促提问者提出的——那将是一种严重的不公。我们中间那些尽力施行基督恩宠的人，谁愿意看到你对基督宗教教义的了解仅限于足以保证你自己得救的程度呢？——不是此世的救恩，正如天主圣言不厌其烦地提醒我们的，此世生命不过是出现片刻便消失的云雾，而是我们基督徒为之追求并获得永恒拥有的那种救恩。在我们看来，你只接受足以使自己获救的教导，未免太少了。因为你的天赋才智，以及你那非凡清晰有力的口才，理应为周围其他人所用，对于那些或因迟钝或因固执而需要以适当的方式为之辩护的人，必须为之捍卫如此丰沛的恩宠的救恩计划；傲慢的狭隘之辈轻视这计划，自夸能行大事，实则对于医治甚或约束自己的恶习，他们一无所能。\par

你问道：\Quote{世界的主与统治者果真充满了童贞女的胎吗？祂的母亲是否忍受了长达十个月的孕期劳累，而且虽仍为童贞，却在适当时期生下她的孩子，并在其后依然保持童贞无损？那位据认为宇宙几乎难以容纳的祂，竟被隐藏在一个啼哭婴儿的弱小身体内吗？祂历经童年岁月，逐渐长大，确立了成年男子的刚毅吗？这位统治者如此长久地流亡于自己的居所之外，而整个世界的管理竟被交托给一个尺寸如此微小的身体吗？祂睡觉，祂进食以获滋养，且经受凡人一切感觉吗？} 你接着说：\Quote{祂伟大威能的证据并未以足够充分的证据显明出来；因为驱魔、治病、复活死人——若我们想到其他也行过这些奇迹的人——对天主而言，不过是小事一桩。} 你说，这个问题是在一次朋友的聚会中由其中一人提出的，但你们其余的人阻止了他进一步提出别的问题，并解散了聚会，把对此事的考虑推迟到你们能获得超出自身经验之外的知识之时，以免因过于冒失地探究隐密之事，而使原本无辜的错误变得应受谴责。\par

于是你向我恳请，并要求我留意，在你这般承认无知之后，对我有何期望。你补充说，我的声誉取决于你能否得到对这些问题的答复，因为尽管在其他司铎那里，无知可被宽容而不损害信仰，但当你向我这身为主教的人提问时，任何我所不知的，都必定不属于基督教体系。\par

因此，我首先请你放下那过于轻率对我形成的看法，并抛弃那些情怀——尽管它们是你善意的愉快证据——而如果你回报我的情感，就请相信我对自己的见证，胜过他人对我的评断。因为圣经的深度如此之深，以至于即使我从童年早期直到衰迈老年，用最充分的闲暇、最不懈的热忱和比我所有的更优越的才能，去钻研它而不及其他，我仍会每日在发掘其宝藏中不断进步；这并非说，通过圣经获知得救所必需的事理有何等巨大的困难，而是当任何人以那为虔诚正直生活所不可或缺的基础之信德接受了这些真理之后，仍有许多被重重奥秘之影所遮蔽的事情，有待那些在研究中进阶的人去探究；不仅在表达这些事理的话语中，而且在事理本身中，都蕴含着如此之深的智慧，以致最年长、最有能力、最热忱的圣经研究者的经验印证了圣经自己所说的话：\BibleQuote{当人完成时，他才开始。}\BibleRef{德 18:6}\par

% structure: p0008 source=h2 target=subsection reason=relative-heading-rank; base=h2

\subsection{第二章}

4. 但对此何必再多言呢？我更应当针对你提出的问题发言。首先，我愿你明白，基督宗教的道理并不认为，\OriginalTerm{天主性}{Godhead} 与祂由童贞女所生的人性混合到了如此地步，以致祂放弃或丧失了对宇宙的管理，或将这管理转移给了那作为一个微小而受限的物质实体的肉身。这种看法只被那些除了物质实体——无论是较稠密的，如水与土，或较精微的，如空气与光——便不能设想任何事物的人所持有；然而所有物质实体都同样具有这一特点：它们中间没有一个能整体地遍在一切处，因为，由于其众多部分之故，它不能不具有一部分在此处，另一部分在彼处，并且无论该物体或大或小，它都必须占据某个位置，且如此填充位置，以致其整体不在所占据空间的任何一个部分。因此，物质身体能够被凝聚和稀薄、收缩和膨胀、压成碎片并扩成巨大体积，便是其独特的特性。灵魂的本性与身体的本性大不相同；而既是灵魂与身体二者之造物主的天主，其本性更当是如何地不同啊！不能说天主填充世界的方式如同水、空气，乃至光占据空间那样，以致用自己较大或较小的部分去占据世界较大或较小的部分。祂能以自己整个的存在遍在一切处：祂不能被限制在任何地方：祂能来临却不离开原先所在之处：祂能离去却不弃舍所来到之处。\par

5. 人的心灵对此感到惊奇，并因无法理解，便或许拒绝相信。然而，心灵不可在对自身内的奥秘同样感到惊奇之前，便继续不信地惊奇于天主的属性；可能的话，让它稍稍提升到身体之上，超越那些它惯于通过身体器官知觉的事物，并让它观想那使用身体作为工具的东西是什么。也许它做不到这点，因为正如有人所说，需要有极大的心灵能力才能使心灵离开感官，并使思想脱离其惯常的轨道。那么，让心灵以这种颇为不寻常的方式，并带着最大的专注，来考察身体的感官。有五种不同的身体感官，它们既不能没有身体，也不能没有灵魂而存在；因为，一方面，只有在人活着时，感官知觉才可能，而身体从灵魂获得生命；另一方面，只有借助于身体的管道与器官，我们才能行使视觉、听觉及其他三种感官。让理性灵魂集中注意于此题目，并不要通过这些感官本身，而是用自己的理智与理性来判断身体的感官。一个人除非活着，当然不能通过这些感官知觉；但在灵魂与身体因死亡分离之前，他生活在身体内。那么，他那只生活在他身体内的灵魂，又是如何知觉到他身体表面之外的事物的呢？天上的星辰岂不是远离他的身体吗？然而他岂不看见那边的太阳吗？视觉岂不正是身体感官的一种运用——岂不正是其中最尊贵的一种吗？那么，又如何呢？他是否既生活在他的身体内，又生活在天上，因为他用自己的一种感官知觉到天上的事物，而感官知觉不能在知觉者没有生命的地方存在？或是即使他并非生活在那里，他也能知觉——因为尽管他只生活在他自己的身体内，他的知觉感官也在那些位于他身体之外、远离他身体、包含着他用视觉所接触的对象的地方活动？你看到了吗，甚至在我们如此显而易见的感官，即我们称之为视觉的感官中，也蕴含着多么大的奥秘？再考察一下听觉，并说说灵魂是否以某种方式扩散到了身体之外。因为，除非我们在发出敲门声的地方运用了听觉，我们又怎能说\Quote{有人在敲门}呢？因此，在这情形中，我们也生活在我们身体的界限之外。或者，我们能够在一个我们并非生活其中的地方，用感官去知觉吗？然而我们知道，在没有生命的地方，感官是无法运作的。\par

6. 另外三种感官是通过与其自身器官的直接接触而运作的。也许在嗅觉方面，这一点可以合理地争论；但对于味觉和触觉，则没有争议：我们所尝到和触到的事物，我们除了通过与自己身体的接触之外，不可能在别处感知到。因此，把这三感官暂且搁置不论。视觉和听觉给我们提出了一个奇妙的问题，要求我们解释：灵魂如何能通过这些感官在它所不居住的地方感知，或者，它如何能居住在它不在的地方。因为灵魂除了在自己的身体中，不会在任何地方，但它却通过这些感官在身体之外的地方感知到事物。因为灵魂在任何地方看见任何事物，就在那地方运用其感知能力，因为看是一种感知行为；灵魂在任何地方听见任何事物，就在那地方运用其感知能力，因为听是一种感知行为。因此，灵魂要么生活在它看见或听见的地方，因而自己就在那地方，要么在它不生活的地方运用其感知能力，要么它生活在一个地方，却同时不在那里。所有这些都是令人惊奇的；其中任何一个说法都难免显得荒谬；而我们还只是在谈论可朽的感官。那么，那超越了肉体感官的灵魂本身，即那居住于理智中、用以思考这些奥秘的灵魂，又是什么呢？因为它并不是藉着感官来对感官本身作出判断的。那么，当我们断言\OriginalTerm{圣言}{Verbum Dei}——万物都是藉着祂而造成的——如此由童贞玛利亚取了肉身，并以可朽的感官将自己显示出来，既没有破坏自己的不朽，也没有改变自己的永恒，既没有减弱自己的权能，也没有放弃掌管世界，更没有离开圣父的怀抱（即祂与圣父同在、在圣父内的隐秘之处）时，我们是否就认为，我们听到了某种不可信的事呢？\par

7. 你要理解\OriginalTerm{圣言}{Verbum Dei}的本质，万物都是藉着祂造成的，祂的本质使得你不能设想圣言的任何部分会消逝，由将来变成过去。祂保持自己的所是，并整个地无所不在。当祂显示时，祂就来临；当祂隐藏时，祂就离去。然而，无论隐藏还是显示，祂都与我们同在，正如光既临在于能看见的人眼前，也临在于盲人眼前；但能看见的人感受到光的临在，而盲人则感受不到。同样，声音既接近能听见的人，也接近聋人，但它让前者知晓其临在，而对后者则隐藏。然而，我们声音和话语的声响所发生的事——一件确实转瞬即逝的事——岂不比这更奇妙吗？因为当我们说话时，在前一个音节的声音停止之后，下一个音节才能发出；然而，如果有一个听者在场，他听到我们所说的一切；如果有两个听者在场，两人听到的是同样的，并且对每一人来说，都是整个的；如果有许多人静静地聆听，他们并不像分饼那样把声音分割开来，分给各人，而是把所说出的一切完整地分施给所有的人，也分施给每一个人。请想想这一点，然后说说，难道永恒的天主的圣言在宇宙中竟不能成就那会过去的世人的言语在听者耳中所成就的事，这难道不是更不可信的吗？\par

8. 因此，没有理由害怕主在婴孩时的弱小身体，会使得\OriginalTerm{天主性}{Deitas}似乎在它里面受到局限。因为天主的伟大不在于体积，而在于能力：祂以自己的眷顾，赐给了蚂蚁和蜜蜂比驴和骆驼更优越的感官；祂从最微小的种子里，形成无花果树的巨大身躯，而许多较小的植物却从大得多的种子里生长出来；祂也赋予了小小的瞳孔这样的能力，使其在一瞥之间，瞬间扫过几乎半个天空；祂将整个神经系统的五重体系，从大脑中的一个中心和一点，扩散到全身；祂从心脏这一相对较小的肢体，将生命运动分配至全身；通过这些以及其他类似的事，祂——在小事上为大的那位——神秘地从极其微小的事物中产生出伟大的事物。祂的能力如此伟大，以至于祂在困难的事上觉察不到困难。正是这同一能力，不是从外而来，而是从内，在童贞女的胎中，肇始了一个婴孩的成孕：这同一能力，将一个人类的灵魂，并通过灵魂也将一个人类的身体——简言之，整个人性，使其在与祂的结合中被提升——与祂自己联合在一起，而祂自己却丝毫没有因此而被贬低；祂公义地从人性接受了人类的名字，同时丰沛地将\OriginalTerm{天主性}{Deitas}的名字赐予人性。婴孩耶稣的身体，是由祂那仍是童贞的母亲腹中诞生的，正是由于这同一能力；后来，当祂是成人时，也是这同一能力使祂的身体经由关闭的门进入了楼上的房间。\BibleRef{若 20:26} 在这里，如果追问此事的缘由，它就不再是奇迹了；如果要求一个完全相似的事例，它也就不再是独一无二的了。让我们承认，天主能做某件我们必须承认超出我们理解的事。在这些奇妙的事中，整个工作的解释就在于那位行事者的权能。\par

% structure: p0014 source=h2 target=subsection reason=relative-heading-rank; base=h2

\subsection{第3章}

9. 祂在睡眠中休息、以食物滋养、经历人性的一切感受，这些事实向人证明，祂所取的人性真实不虚，而非加以摧毁。看，这就是事实；然而一些\OriginalTerm{异端者}{heretic}，因对祂能力的扭曲仰慕和赞美，完全拒绝承认祂人性的真实，而在这人性中，包藏着祂拯救信祂之人的一切恩宠，其中蕴藏着智慧与知识的深邃宝藏，并将其信德赋予那些心灵，使它们上升到对不变真理的永恒默观。倘若全能者创造基督的人性，不是让祂由母亲所生，而是通过其他方式，并且突然将祂呈现于世人眼前，那会如何？倘若主未曾经历从婴儿到成年的成长阶段，既未进食也未睡眠，那又会如何？这岂不会印证上述错误的印象，使人完全无法相信祂果真取了真实的人性，并在保留奇迹的同时，从祂的行动中剔除慈悲的元素吗？但如今祂如此显现，作为\OriginalTerm{中保}{Mediator}，在一个位格内结合两性，既以超凡的使平凡显得尊贵，又以平凡调和祂内的超凡。\par

10. 然而，在受造界一切各异的运动中，有哪一件天主的化工不奇妙呢？若非因习以为常，这些奇事竟在我们眼中变得渺小。其实，有多少寻常事物被践踏于脚下，若细加审察，却足以唤醒我们的惊叹！以种子的特性为例：谁能理解或说明其种类的繁多，那生命力、活力与隐秘的能力，使得它们从微小的容积中演变出宏大之物？祂所取的人之灵魂与身体，乃是由那位在自然界中原本无需先在种子而创造种子者，无需种子而受造。在如此成为祂的身体内，那自身毫无可变，按其旨意编织了过去诸世纪变迁的祂，也顺服于时序接替与人生命的一般阶段。因为祂的身体，既始于时间的一点，便随时间流逝而发育成长。然而，那在起初就有的、万世因之而存在的\OriginalTerm{天主圣言}{Word of God}，并非屈从于时间，仿佛\OriginalTerm{降生成人}{Incarnation}的事未经祂同意而临到；祂选择了那个时刻，自由地摄取了我们的本性。人性与天主性结合为一；天主并未撤去自身。\par

11. 一些人坚决要求，对于天主性如何与人灵肉结合以构成基督的唯一位格——此事在世界历史中只需发生一次——给予解释，其大胆程度，仿佛他们自己能够解释灵魂如何与身体结合以构成人的唯一位格——这事每天都在发生——一样。因为正如灵魂在一个位格内与身体结合而构成人，同样，天主在一个位格内与人结合而构成基督。在前者的位格中，是灵魂与身体的结合；在后者的位格中，是天主性与人的结合。然而，我的读者应警惕，不要从物质物体的属性中借用\Quote{结合}的概念，当两种流体结合时，它们混合得如此彻底，以至于两者都不保持其原有的特性；尽管即使在物质物体中也有例外，比如光，当它与空气结合时，并未发生改变。因此，在人的位格中，有灵魂与身体的结合；在基督的位格中，有天主性与人的结合；因为当天主圣言与一个具有身体的灵魂结合时，祂将灵魂与身体一同摄取到与自身的合一中。前一件事在人类个体生命之始每日发生；后一件事为人类的救恩而一次成就。在这两件事中，两种非物质实体的结合，理应比一种非物质与一种物质实体的结合更易被相信。因为，如果灵魂对自己的本性没有误解，它就明白自己是非物质的。这一属性更确定地属于天主圣言；因此，圣言与人类灵魂的结合，应当比灵魂与身体的结合更值得相信。后者我们在自身中体验得到；前者我们受命相信已在基督内成就。但若这两者都与我们的经验相异，我们受命相信两者都已发生，那么哪个我们会更愿意相信已经发生呢？我们岂不承认两种非物质实体能比一种非物质和一种物质实体更容易结合吗？除非，也许用\Quote{结合}一词形容这些事物并不恰当，因为它们的本性，无论在自身内还是在我们所知中，都与物质物体的本性不同。\par

12. 因此，天主圣言，祂也是\OriginalTerm{圣子}{Son}，与\Person{圣父}同为永恒，是天主的能力与智慧，\BibleRef{格前 1:24}，强有力地渗透并和谐地安排万物，从理智受造物的最高界限到物质受造物的最低界限，\BibleRef{智 8:1}，显而又隐，无处受限，无处分离，无处伸展，却无量度，以其整体无所不在地临在——我说，这位天主圣言，以完全不同于其临在于其他受造物的方式，摄取了一个人的灵魂与身体，并藉由此结合，构成了唯一位格的\Person{耶稣基督}，天主与人之间的中保，\BibleRef{弟前 2:5}，在其天主性上，与\Person{圣父}同等，在其肉身上，\Emph{即}在其人性上，次于\Person{圣父}——就天主性而言，祂与\Person{圣父}同等，因而永恒不朽；就脆弱而言，这是祂因与我们的本性有份而取的，因而可变化且有死。\par

在这基督内，祂在自己知道最为适宜的时刻，也就是在世界创立之前就已预定的时刻，给人类带来了为获得永恒救恩所必需的\Emph{教导}和\Emph{援助}。教导借由祂而来，因为那些在祂之前、为人类的益处而在世上所宣讲的真理——不仅由言语皆真的圣先知们，而且由哲学家们、甚至诗人和各类文学作者（毫无疑问，他们将许多真知与虚假混合）——都通过祂那具人性权威的实际展示，被证实为真理，这是为了那些无法在本质真理的光照下领悟和辨别它们的人；而这真理甚至早在祂摄取人性之前，就已临在于所有能够接受真理的人。此外，借着祂降生成人的事实，祂为我们的益处教导了超越一切的事——即，那些渴望神圣存在的人，出于骄傲而非虔诚，曾以为必须不直接接近祂，而是要经由他们视作神明的天上权能，以及各种神圣却亵渎的禁忌礼仪（在这些敬拜中，魔鬼通过骄傲在人类与他们之间形成纽带，取代圣天使的位置而成功），如今人们可以明白，那位他们曾以为遥不可及、并借中介之力而非直接接近的天主，实际上如此亲近于人们对祂的虔敬渴望，以至于祂屈尊摄取人的灵魂与肉身，与自己结合得如此紧密，以至于这完整的人与祂的结合，就如同人体内灵魂与肉身的结合一样，只是暂且不论灵魂与肉身皆有共同的逐步发展，而祂并不参与此一成长，因为成长意味着可变性，而这属性是天主所不能具有的。再次，在这基督内，人类获得了得救所必需的\Emph{援助}，因为没有源于祂的信德的恩宠，无人能克制邪恶的欲望，或借着宽恕从任何他所未能完全战胜的罪恶欲望之力量带来的罪咎中获得洁净。至于祂教导所产生的效果，如今难道还有一个无论多么软弱的痴愚者，或一个无论多么卑微的愚蠢妇人不相信灵魂的不朽和死后生命的实在性吗？然而这些真理，当亚述人\Person{菲雷西德}首次在古代希腊人的讨论中主张它们时，其新颖性深深打动了萨摩斯的\Person{毕达哥拉斯}，使他从运动员的锻炼转向哲学家的研究。但现在，我们全都见证了\Person{维吉尔}所说的话：\Quote{亚述的香膏在各处生长。}至于通过基督的恩宠所给予的援助，在祂内确实应验了同一位诗人的话：\Quote{以你为领袖，我们罪过的一切痕迹的消除，将使大地免于永恒的恐惧。}\par

% structure: p0020 source=h2 target=subsection reason=relative-heading-rank; base=h2

\subsection{第四章}

13. \Quote{但是，\InnerQuote{他们说，}如此伟大的尊威的证据并没有以充分明晰的方式显现出来；因为驱逐魔鬼、医治病人、使死者复活，这些事对天主来说不过是小事一桩，如果我们想到其他行过类似奇迹的人的话。}我们自己承认，先知们曾行过一些与基督所行类似的奇迹。因为在那些奇迹中，有什么比使死者复活更奇妙呢？然而\Person{厄里亚}和\Person{厄里叟}都行了这事。\BibleRef{列上 17:22}至于术士们的奇迹，以及他们是否也能使死者复活的问题，就让那些人去发表意见吧，他们不是作为控告者，而是作为颂赞者，竭力想证明\Person{阿普列尤斯}犯下了他本人曾竭力为自己辩白的行巫术的罪行。我们读到，埃及的术士们，即精通这些法术的人，当他们借着不虔敬的咒语行奇事时，却被天主的仆人\Person{梅瑟}所挫败，\Person{梅瑟}只是简单地呼求天主就推翻了他们的一切诡计。然而这\Person{梅瑟}本人和其他所有真先知都预言了主基督，并给予祂极大的荣耀；他们预告祂来临，不是作为在行奇迹的能力上与他们同等或稍胜一筹的一位，而是作为确是真天主、万有之主，为人类的益处而成为人的一位。祂也乐意行一些他们曾行过的奇迹，以避免祂没有亲自行那些祂曾藉他们行的奇迹所造成的不协调。不过，祂也要行一些唯独属于祂自己的事，即由童贞女诞生、从死者中复活、升天。我不知道，那些认为这些事对于天主来说太微不足道的人，还能期待什么更大的事。\par

14. 因为我想，这些反对者所要求的天主大能的标记，并不适合祂在取了人性时施行。\BibleQuote{在起初已有圣言，圣言与天主同在，圣言就是天主。万物是藉着祂而造成的。}\BibleRef{若 1:1,3} 那么，当圣言成为血肉时，祂必须创造另一个世界，我们才能相信祂就是那创造世界的吗？但在这个世界之内，不可能创造一个比它更大或与之相等的世界。然而，如果祂创造一个比现有世界更小的世界，这也将被认为是一件太小而不能证明祂天主性的工程。因此，既然祂没有必要创造一个新世界，祂便在世间行了新事。因为一个人由童贞女诞生，从死者中复活得永生，并被高举于诸天之上，这或许是一项比创造世界更显其大能的事工。在此，反对者可能回答说，他们不相信这些事发生过。那么，对于那些轻视较小的证据，认为其不足为凭，而拒绝较大的证据，认为其难以置信的人，又能做什么呢？让死者复活已被人相信，因为别人也曾行过此事，而这事太小，不足以证明施行者就是天主；但是，一个真实的身体在童贞女内受造，并从死亡中复活得永生，被提升天，却没有人相信，因为无人做过这事，而这是唯独天主才能做的。按照这个原则，每个人都会泰然接受他自己认为容易——不是自己去做，而是容易设想的事，而将一切超越这种界限的事拒绝为虚假和虚构。我恳求你，不要像这些人一样。\par

15. 这些主题在别处有更详尽的讨论，而在有关教义的基本问题上，每一个疑难之点都通过彻底的探讨和辩论得以阐明；但信德使人领悟这些事物，不信则关闭门户。有谁看到从起初以来那奇妙的一系列事件，以及世界各时代相互衔接的方式——如此，我们对现今事物的信德便得到过去所发生之事的辅助，而早期和远古的记载则得到后来、较近之事的证实——能不归向对基督教导的信德呢？从迦勒底人中拣选出一人，一位以极卓越的虔敬和信德著称的人，为要赐他天主的许诺，这许可是要在世界末期，经过如此诸多世纪之后才得应验的；并预先宣告，地上万民都要因他的后裔而得蒙祝福。\EditorialNote{创 12} 这人敬拜唯一真天主，宇宙的创造者，在他年迈之时生了一个儿子，那时他的妻子因不育和年事已高，早已放弃成为母亲的希望。这儿子的后裔成了一个非常繁盛的支派，在埃及人数大增，他们是从东方被迁到那里去的，天主的上智安排随着时间的推移，使得赐给他们的许诺和为他们行的大事都不断增多。他们从埃及出来，成为一个强大的民族，是藉着可怕的征兆和奇事被领出来的；应许之地的邪恶民族从他们面前被赶出去，他们被领进去定居在那里，并被提升到王国的地位。此后，他们因屡次以盛行的罪恶和崇拜偶像的不虔敬行为，激怒那赐予他们如此众多恩惠的真天主，并交替经历灾祸的惩罚和复兴的安慰，这个民族的历史一直延续到基督的降生成人和公开显现。所有关于这基督的预言——祂是天主的圣言，天主的圣子，而且是天主自己，祂要降生成人，受难，复活，升天，藉着祂名字的权能，吸引万民中众多的人归向祂，并且藉着祂，凡信祂的人都将获得罪过的赦免和永恒的救恩——这些预言，我说，已经通过赐给那个民族的诸多许诺，通过所有的预言，司祭职的建立，祭献，圣殿，简言之，通过他们所有神圣的奥迹，公之于世了。\par

16. 因此，基督来了：祂的诞生、生活、言语、行动、苦难、死亡、复活、升天，所有先知所预言的都得以应验。\BibleRef{玛 1:22} 祂派遣圣神；当信徒们聚集在一间屋里，以祈祷和热切渴望期待这所应许的恩赐时，祂就以这圣神充满他们。既被圣神充满，他们便立刻用万邦的方言讲话，勇敢地驳斥谬误，宣讲对人类最有益的真理，劝人悔改过去应受谴责的生活，并许下因天主的白白恩宠而得宽赦。适合作为确证的标记与奇迹伴随着他们对虔敬和真宗教的宣讲。不信的残酷仇恨被煽动起来反对他们；他们承受着预言中的试炼，盼望所应许的福祉，并教导他们受命要传扬的道理。起初人数稀少，他们却像种子般分散普世；他们以奇妙的速度使万民归化；在仇敌中人数日益增长；藉着迫害反而增多；在严峻的艰难困苦中，不但未被遏制，反而将影响力扩展到地极。从极其无知、受人轻视、人数稀少，他们变为受到光照、出类拔萃、人数众多，成为才能卓越、辩才优雅的人；他们也使那些以天才、口才和博学著称、成就非凡的人，归服于基督的轭下，并吸引他们参与宣讲圣善与救恩之道的工程。在顺逆交替之中，他们警醒地操练忍耐与节制；当世界之日临近终穷，耗尽其精力的灾祸预示着终结将临时，他们从中看到预言的应验，便以加增的信心只期待着天上之城的永恒福乐。而且，在这所有的变动之中，异教民族的不信继续向基督的教会发怒；教会却藉着坚忍的忍耐，并在面对仇敌残酷时保持不动摇的信德，而获得胜利。那一位的祭献——在祂身上，久藏于奥秘预许之下的真理得以启示——既已奉献，那些预表此祭的祭献便因犹太圣殿的彻底毁灭而最终废止。犹太民族本身因不信而被弃绝，如今从本土被根除，分散到世界各地，为要把\WorkTitle{圣经}带到各处，这样一来，我们的仇敌自己便可将关于基督及其教会的预言所提供的见证带到人类面前，从而排除了人们认为这些预言是我们为顺应时事而伪造的可能；在这些预言中，也预告了这些犹太人的不信。异教神祇的庙宇、偶像和渎神的崇拜，依照先知的预示，逐渐而有序地被摧毁。各种异端，虽假托基督之名，却正如预言所说，起来反对基督之名，藉着这些异端，圣教的道理受到考验并发展。所有这些事，如今都已应验，完全符合我们在\WorkTitle{圣经}中所读的预言；从这些众多而重大的已应验的预言，我们便满怀信心地期待其余预言的成就。那么，凡是向往永恒并为此现世生命的短暂所驱动的心灵，又怎能抗拒这些证明我们信德具有神圣起源的清楚而完备的证据呢？\par

% structure: p0025 source=h2 target=subsection reason=relative-heading-rank; base=h2

\subsection{第5章}

17. 有什么哲学家的言论或著作，有何地何代任何国家的法律，能有一刻值得与那两条诫命相比？基督说，全部法律和先知都系于这两条诫命：\BibleQuote{你应全心、全灵、全意，爱上主你的天主，并应爱近人如你自己}？\BibleRef{玛 22:37} 这里包含了全部哲学——物理学、伦理学、逻辑学：\Emph{第一}，因为在造物主天主内存在着自然界一切存在的所有原因；\Emph{第二}，因为善良正直的生活，只有通过以应有的方式爱那当爱的对象——即天主和邻人——才能产生；\Emph{第三}，因为唯有天主是理性灵魂的真理和光明。这里也包含了国家福祉与声誉的保障；因为没有任何国家能完美地建立和保存，除非建立在信德与坚固和谐的基础上并借此维系，那时，最高最真的公共福祉——即天主——为所有人所爱，人们也在祂内毫无虚伪地彼此相爱，因为他们是为了祂而彼此相爱，无法向祂隐藏其爱的真实品质。\par

18. 你还要思考\WorkTitle{圣经}的文体——它如何为所有人所理解，尽管其中更深的奥秘只有极少数人能洞悉。它以亲切友谊般的质朴语言，将其中所包含的浅显真理宣告给无学问和有学问之人的心；但即使是那些用象征遮盖的真理，它也并非用僵硬威严的句子表达，使迟钝无知的心灵望而却步，如同穷人回避富人一般；而是藉着文体的俯就，它邀请众人不但以显明的真理为食粮，也以隐藏的真理作为操练，无论在简易还是晦涩的部分，都有着同一真理。为了避免容易理解的内容使读者生厌，同一真理在别处以较为晦涩的方式表达时，便重新成为所愿，而一旦成为所愿，就奇妙地获得新的吸引力，从而愉快地被接纳入心。通过这种方式，悖逆的心得到纠正，软弱的心得到滋养，坚强的心充满喜乐，而这一切对所有人都大有裨益。这教义没有敌人，除非那人陷于错谬，不知其无比的益处，或是灵性患病，厌恶其医治的能力。\par

19. 您看，我写了这样一封长信。因此，若有任何事困扰您，而您认为此事足够重要，值得我们彼此讨论，那么请勿让自己因拘守寻常信函的篇幅而受限制；因为您同任何人一样清楚，古人论述他们无法简短解释的题目时，曾写下多么长的信函。而且，即使文学其他领域作者的习惯有所不同，基督徒作者们的权威——他们的芳表更值得我们效法——也可摆在我们面前。所以，请看看宗徒书信的长度，以及对这些神圣谕旨所作的注释的篇幅，便不要迟疑：若您有许多疑难，便多多发问；或更充分地论述您所提出的问题，好使我们力所能及之处，不再有疑惑的阴云遮蔽真理之光。\par

20. 因为我知道，阁下您必须面对某些人极其坚决的反对；这些人以为——或但愿他人以为——基督教义与国家的福祉不相容，因为他们希望国家不是靠坚定地践行德行，而是靠容许恶行不受惩罚来建立。但在天主面前，许多人结伙所犯的罪行，不会像在一位国王或任一仅仅是人间的官长面前那样，往往不遭受惩罚。此外，祂的慈悲和恩宠，由基督——祂自己是人——向人宣告，并由同一位基督——祂也是天主和圣子——分施于人；这慈悲和恩宠，从不离弃那些凭信德在祂内生活、虔诚朝拜祂的人：在逆境中，他们耐心而勇敢地承受此世的考验；在顺境中，他们有节制地、怀着对他人怜悯之心，善用此生的美物；注定要在那神圣的天上的城邑中，为这两种景况下的忠信，领受永恒的赏报；在那里，再没有灾难需要痛苦忍受，也没有混乱的欲望需要用辛劳的顾虑去控制，我们唯一的工作只是毫无困难地、以完全的自由，保存我们对天主和近人的爱。\par

愿天主无限慈悲的全能护佑您平安，并增进您的幸福，我尊贵而卓越的阁下，我最卓越的儿子。怀着应给予您尊位的深切敬意，我向您虔诚且真正可敬的母亲致意，愿天主俯听她为您所做的祈祷！我虔诚的弟兄兼同为主教的\Person{波西迪乌斯}，热诚地向尊驾致意。\par

\SourceRecord{Translated by J.G. Cunningham.；Nicene and Post-Nicene Fathers, First Series；Vol. 1.；Edited by Philip Schaff.；Buffalo, NY: Christian Literature Publishing Co.,；1887.；<http://www.newadvent.org/fathers/1102137.htm>.}

% structure: p0001 source=h1 target=section reason=page-title

\section{书信一三八（公元412年）}

\Emph{致\Person{马塞林}，我尊贵而著名的、公正的大人，我最亲爱、最想念的儿子，\Person{奥古斯丁}在主内致意。}\par

% structure: p0003 source=h2 target=subsection reason=relative-heading-rank; base=h2

\subsection{第一章}

1. 当我写信给那位我们两人都真诚地敬爱的、既杰出又极具口才的\Person{沃鲁西亚努斯}时，我认为只应局限于回答他自己认为适宜提出的问题；至于你在来信中提交给我讨论和解答的问题，无论是沃鲁西亚努斯本人还是其他人提起或提出的，我将针对这些给予尽可能的答复，寄给你更为妥当。我将试着去做，不是以那种需要写成正式论著的方式，而是以书信的闲谈亲近适合的方式，以便若你，因日常的讨论了解他们的心境，认为恰当，这封信或许也可以读给你的朋友们听。但若这通书信并不适于他们，因他们未有信仰的虔敬而预备聆听，那么，你认为适合他们的内容，先在我们自己之间准备，然后再将如此预备好的传给他们。因为有许多事，他们的心灵目前还可能退缩躲避，或许再过些时候，他们能逐渐被说服，接受它们为真，或是通过更详尽的、更巧妙的论证，或是诉诸于依他们看来不当拒斥的权威。\par

2. 在来信中你提到，有人对此问题感到困惑：\Quote{为何这位天主，就是也被证明是\WorkTitle{旧约}中天主的，在拒绝了古老的祭献后，又喜欢新的祭献。因为他们声称，除了那些被证实先前所行不当的事外，没有什么是可纠正的，或者说，一旦曾行得正确，便丝毫也不应更改：他们说，那行得正确的事，若无错误便不可更改。} 我引用你信中的这些话。如果我要详尽地答复这个异议，即使时间耗尽，我也无法穷举自然界本身的过程和人的工作都随着时间的情境发生变化的实例，而与此同时，管理这些变化的原则或计划本身却是毫不变动的。我可以提及少数几个例子，受它们激发，你警觉的观察力便犹如由它们导向更多同类事例。夏天不是跟随着冬天，气温逐渐变暖吗？黑夜和白昼不是轮流接替吗？我们的人生经历多少次变化啊！一去不返的童年，给青年让位；壮年，它自己也注定只持续一段时期，又轮替着取代了青年；而老年，结束壮年之期，自身又因死亡而终结。所有这一切都在变化，可是那安排这些相继变化的天主眷顾计划却不变。我认为，同样地，当农夫指定夏天要做的工作与冬天所吩咐的不同时，耕作的原则并未改变。清晨起身，经过夜间休息，并不被视为更改了生活的计划。老师给成人指定的作业，不同于他惯常给少儿学生规定的；他的教导始终如一，只是当功课改变时，教学随之而变，自身却并未改变。\par

3. 我们时代的那位杰出医师\Person{温迪齐亚努斯}，曾应一病人之请，为其病症开具了当时在他看来相宜的药方；用之而健康得以恢复。若干年后，这位病人发现又为同一病症所困扰，便以为应施用同一药方；但施用后却使病情恶化。他惊愕之下，再度返回医师那里，告知所发生之事；而这位思维极其敏达的医师回答：\Quote{你因此施用而受害的原因，在于我并不曾吩咐这样做；} 此时，所有听了此言而又不了解此人的，都以为他所依赖的不是医术，而是某种被禁止的超自然力量。他后来受到一些因他的话而惊异的人询问时，便解释了他们所不解的：那就是，他不会给已达到如今这个年龄的那位病人开具同一药方。因此，即使医术的原则和方法毫无改变，但按照它们，依照时代差异而可能必需的改变却是巨大的。\par

4. 因此，断定某件曾经做得正确的事无论如何也不可更改，是在肯定虚妄之事。因为倘若促成某事的时间环境发生了改变，真实的理性几乎在所有情况下都要求，先前那些环境下做得正确的事，如今要如此改变，以至于尽管他们声称若改变便做得不对了，真理却反而明言，若不改变，便做得不对；因为在两个时期，只要按照时间的差异来调节差异，就都将做得正确。因为正如在不同的人那里，或许发生同一时刻这个人可以毫无不利地去做的事，另一个人却不能，其原因不在于所做之事不同，而在于做事之人不同，同样，在同一个人而时间不同时，以前是义务的事，现在却不是了，并非因为此人不同于昔日之我，而是因为做此事的时间不同了。\par

5. 凡是能胜任且细心观察\OriginalTerm{美}{pulchrum}与\OriginalTerm{适宜}{aptum}之间对比的人，都能看出这一问题所开辟的广阔领域；我们可以说，这种对比的实例散布在宇宙各处。因为美，其对立面是丑陋与畸形，是根据其自身来评估和赞扬的。但适宜，其对立面是不协调，则依赖于与之相连的其他事物，不是根据其自身，而是根据它与别物的关联来评估的：\OriginalTerm{合宜}{decens}与\OriginalTerm{不合宜}{indecens}之间的对比，要么相同，要么至少被视为相同。现在，将我们所说的应用到当前主题上。祭献的神圣制度在先前的\OriginalTerm{安排}{dispensatio}中是适宜的，但现在却不适宜了。因为适合现今时代的改变，是由天主所命令的，祂比人无限地更知道什么对每个时代是合宜的；祂，无论是赐予或增添、废除或删减、增加或减少，都是不可改变的\OriginalTerm{治理者}{gubernator}，犹如祂是不可改变的\OriginalTerm{创造者}{creator}之于可改变的事物，在祂的\OriginalTerm{天主眷顾}{providentia}中安排万有，直到时间进程完成的美——其组成部分乃是适应各个相继时代的安排——得以实现，就像某位不可言喻地睿智的歌曲大师所创作的宏伟乐章，那时，那些在此世虽然处在\OriginalTerm{信德}{fides}、而非亲眼看见的时代，却以悦纳的方式敬拜祂的人，将进入天主永恒的当面\OriginalTerm{默观}{contemplatio}。\par

6. 此外，那些认为天主为了自己的益处或快乐而规定这些法则的人，是错误的；难怪他们由于这种错误而困惑，仿佛祂是因为心情变化，才在一个时代命令给自己献上一样东西，而现在又命令献上别样东西。但事实并非如此。天主绝不为了自己的益处而命令什么，而是为了接受命令者的益处。因此，祂是真正的主，因为祂不需要仆役，而仆役却需要祂。在那些旧约圣经中，在祭献仍被献上而如今已被废除的时代，经上说：\BibleQuote{我对上主说：你是我的天主，因为你不需要我的好事。}因此，天主不需要那些祭献，也永远不需要任何东西；但有一些行动，作为这些神圣恩赐的象征，灵魂藉以获得现今的\OriginalTerm{恩宠}{gratia}或\OriginalTerm{永恒光荣}{gloria aeterna}，在举行和履行这些行动时，我们实行了虔诚的功夫，不是对天主有用，而是对我们自己有益。\par

7. 然而，要足够充分地讨论古代和现今的象征性行动的差异，恐怕过于冗长，这些行动由于涉及神圣事物而被称为\OriginalTerm{圣事}{sacramenta}。因为正如一个人在早晨做一件事、晚上做另一件事，这个月做一件事、下个月做另一件事，今年做一件事、明年做另一件事，并非反复无常，天主也没有变化，尽管在世界历史的前一阶段，祂曾命令一种祭献，而在后一阶段命令另一种祭献，从而在与各个时代的变化相协调的情况下，安排了与真宗教的神圣教导相关的象征性行动，而祂自身却毫无变化。为了让那些对此感到困惑的人明白，变化早已在天主的计划中，而且当新的法则被指定时，并非因为旧的法则由于祂旨意的不恒常而突然丧失了天主的悦纳，而是这早已由那位天主的智慧所确定和决定——圣经上关于更大的变化对祂说：你变换它们，它们便被变换；但你是常存的——所以，有必要说服他们，以旧约圣事换成新约圣事的这种交换，早已由先知们的声音预言了。因为这样，如果他们能看出什么的话，他们就会看出，在时间上是新的事，对于那位确定了各个时代，并且没有时间先后而拥有祂按照各种类分配给各个时代的一切的天主来说，并不新。因为在上面我引用了\BibleQuote{我对上主说：你是我的天主，因为你不需要我的好事}这句话的圣咏中，为证明天主不需要我们的祭献，圣咏作者以基督的名义稍后又说：\BibleQuote{我不会聚集他们的血祭集会；}意思是指，为了从他们的牲畜中献上动物，犹太人习惯聚集的集会；在另一处他又说：\BibleQuote{我不从你的家中取公牛，也不从你的圈中取公山羊；}另一位先知还说：\BibleQuote{看，日子将到——上主的断语——我必要与以色列家和犹大家订立新约，不像我昔日——握住他们的手，引他们出离埃及地时——与他们的祖先订立的盟约。}\BibleRef{耶 31:32}此外，关于这个主题，还有许多其他见证，预言天主将要行祂如今所行的事；但提及它们会耗时太长。\par

8. 如果现已确定，在一个时代正当地规定的，可以在另一个时代正当地改变——这种改变表明实行改变者的工程有变化，但计划没有变化，这计划是由祂的理智所构想的，在不受时间先后影响的理智中，那些由于各个时代彼此接续而不能同时实际进行的事是同时呈现的——那么，此时也许有人会期望听我讲明这一改变的原因。你知道，要完全讨论这些原因会需要多久。这件事可以用以下话语概括性地、但对于一个判断敏锐的人来说足够地说明：基督将来的\OriginalTerm{来临}{adventus}应由某些\OriginalTerm{圣事}{sacramenta}预言，而祂来临之后应有其他圣事来宣扬；正如事实的差异迫使我们改变我们谈论来临时所用的语词，或指将来或指过去：预言是一回事，宣扬是另一回事，将要来临是一回事，已经来临是另一回事。\par

% structure: p0012 source=h2 target=subsection reason=relative-heading-rank; base=h2

\subsection{第二章}

9. 我们现在来观察第二点，就是你的信中所继续述说的。你补充说，他们声称\OriginalTerm{基督宗教的教义}{Christian doctrine}和宣讲与\OriginalTerm{公民的义务与权利}{duties and rights of citizens}绝不相容，因为在其规诫中我们发现：\BibleQuote{不可以恶报恶}，\BibleRef{罗 12:17}以及\BibleQuote{若有人掌击你的右颊，你把另一面也转给他；那愿与你争讼，拿你内衣的，你连外衣也让给他；若有人强迫你走一千步，你就同他走两千步}，\BibleRef{玛 5:39}——所有这些都被认为与公民的义务与权利相悖；因为谁肯任凭敌人夺走他的东西，或不对蹂躏\Place{罗马}行省的入侵者进行战争报复呢？对于这些以及类似的轻蔑言论，或许我该说，是那些寻求真理之人的探问，我本可以作出更详尽答复，若非与我辩论的那些人都是受过自由教育的人。对这样的人，何必延长辩论，而不反过来首先自己探究：\Place{罗马}共和国如何可能由那些在遭受伤害时宁愿宽恕而不惩罚冒犯者的人，从微不足道与贫穷治理并扩张至伟大与富庶；或者，\Person{西塞罗}在对当时最伟大的政治家\Person{凯撒}致辞时，如何称赞其品性，说他素来只忘记别人对他的伤害？因为西塞罗在此所说，要么是赞美，要么是奉承：如果是赞美，那是因为他知道凯撒果如其言；如果是奉承，他就表明了\OriginalTerm{共和体}{commonwealth}的最高长官应当做那些他虚假地颂扬凯撒所为的事。然而，\Quote{不以恶报恶}，岂不是抑制\OriginalTerm{复仇的欲情}{passion of revenge}——换言之，就是选择在遭受伤害时，宁愿宽恕而不惩罚冒犯者，并只忘记别人对我们的伤害吗？\par

10. 当这些事迹在他们的作家著作中读到，人们报以热烈掌声；它们被视为美德的记载与推崇，\Place{罗马}共和国正是因实践这些美德而应得统治如此多民族，因为其公民宁愿宽恕而不惩罚那些冒犯他们的人。可是，当\BibleQuote{不可以恶报恶}这条规诫作为\OriginalTerm{天主的权威}{divine authority}所赐被宣读，当这条有益的劝告从我们教堂的讲道台在会众中间公布，或者我们可以说，在向一切性别、年龄和阶层开放的训导之所公布时，我们的宗教却被指责为共和国的敌人！然而，若我们的宗教得到应有的听从，它必会以一种超越\Person{罗穆卢斯}、\Person{努马}、\Person{布鲁图斯}和罗马历史上所有其他著名人物所成就的方式，去建立、祝圣、巩固和扩展这个\OriginalTerm{共和体}{commonwealth}。因为共和国岂不就是一个共和体？因此，它的利益是众人共有的；它就是国家的利益。而国家岂不就是由某种和睦纽带联结起来的众人？在他们自己一位作家的著作中我们读到：\Quote{一群散乱无定的人，藉着和睦在短时间内成为一个国家。}可他们曾在自己的庙宇中规定要宣读哪些劝睦之言呢？远非如此，他们不幸地不得不设计如何能以不得罪任何神祇的方式去敬拜，而这些神祇彼此间如此纷争，以至于若其敬拜者效法它们的争吵，国家必定因缺乏和睦的纽带而分崩离析，正如不久之后民众道德沦丧败坏时，内战始作，国家便沦落至此。\par

11. 但是，即使有人对我们的宗教陌生，又怎能如此耳聋，以致不知道多少规劝和睦的诫命，并非由人的讨论所发明，而是凭着天主的权威写就，在基督的教会中不断宣读呢？就连那些他们宁愿辩论而不愿遵行的规诫，其旨趣也在于此：\Quote{凡掌击我们一颊的，我们应把另一颊也转给对方；凡拿我们内衣的，我们也应把外衣给他；并且，凡强迫我们走一千步的，我们就同他走两千步。}因为做这些事，只为着\OriginalTerm{以善胜恶}{overcome evil by good}，使恶人被善胜，或更确切地说，使恶人里面的恶为善所胜，并使那人脱离恶——不是脱离任何外在、外在于他的恶，而是脱离那内在、属于他自己的恶，这种恶使他遭受比任何外来仇敌的残暴所能加予的更为严重而致命的损失。所以，那以善胜恶的人，耐心承受现世利益的损失，为要显明，那些因过度贪爱而使人变得邪恶的事物，相比于\OriginalTerm{信德与义德}{faith and righteousness}，理应受到鄙视；好使那伤害人的人，能从他所伤害的人那里学到他犯下伤害所为之物的真正本质，并带着对自己罪过的痛悔，被赢回到\OriginalTerm{和睦}{concord}之中——对国家来说，再没有比这更裨益的了——他不是被忿怒报复的力量所胜，而是被坚忍受苦的善性所胜。因为，当如此行似乎有益于那为此而施为的人，使他改过迁善并与人和睦时，才算是行得正确。无论如何，即便结果与预期不同，即使那本欲藉此施以这所谓治愈、\OriginalTerm{救命的良药}{salutary medicine}的人，拒绝改正与和解，也当抱持这样的意向而行。\par

12. 此外，如果我们注意这一诫命的字句，并认为自己受字面解释的束缚，那么当左颊被打时，我们就不应把右颊递上去。经上说：\BibleQuote{凡掌击你的右颊，你把另一面也转给他；}\BibleRef{玛 5:39}然而左颊更易被打，因为攻击者的右手打左颊比打另一侧更方便。但通常这些话被理解为，仿佛我们的主在说：若有人在你们所拥有的较高的财物上伤害了你们，那么也把较低的财物给他；免得你们更关心报复而非宽容，致使在暂时的事物面前轻视永恒的事物，其实暂时的事物相对于永恒的事物，应当如同左相对于右一样被轻视。这始终是殉道者的目标；因为最终的报复惟有当不再有改过迁善的余地时——即在最后的伟大审判中——才能公义地要求。但在此期间，我们必须提防，免得因渴望报复而失去忍耐本身——这一德行比敌人所能从我们这里夺走的一切都更有价值，即使我们抵抗。另一位圣史在记述同一诫命时，并没有提及右颊，而只是说这一颊或那一颊；\BibleRef{路 6:29}因此，虽然从\BibleBook{玛窦}{福音}可以稍更清楚地学得这一义务，但他却只是赞扬了同样的忍耐工夫。所以，一个正义而虔诚的人应当准备以忍耐去承受他想使之成为善人的那些人对他的伤害，这样好人的数目得以增加，而非他自己因报复伤害而加入恶人的数目。\par

13. 总之，这些诫命更关乎内心的意向，而非在人眼前所行的行动；要求我们在心灵深处，怀有忍耐与仁爱，但在外在行动上，去做那些看起来最有益于我们应寻求其善的人的事。这一点从以下事实显明：我们的主\Person{耶稣基督}本人，就是忍耐的完美典范，当祂被打脸时，祂回答说：\BibleQuote{我若说得不对，你指证那里不对；若对，你为什么打我？}\BibleRef{若 18:23}如果我们只看字句，祂在这件事上并没有遵守自己的诫命，因为祂没有把另一面脸转给那打祂的人，反而阻止了那行恶的人再施恶；然而祂来，却已准备不仅被打脸，甚至为那些对祂施以钉刑的人被钉死在十字架上；当祂悬在十字架上时，祂为这些人祈祷说：\BibleQuote{父啊，宽赦他们吧！因为他们不知道他们做的是什么。}\BibleRef{路 23:34}同样，宗徒\Person{保禄}似乎也没有遵守他的主和师傅的诫命；当他像祂一样被打脸时，他对大司祭说：\BibleQuote{天主将要打击你，你这粉白的墙！你坐下审判我，应按照法律，你竟违反法律，下令打我吗？}旁边站着的人说：\BibleQuote{你竟敢辱骂天主的大司祭吗？}他便费心以讽刺的方式表明他的话意指什么，好让那些明辨的人能够明白，那\OriginalTerm{粉白的墙}{whited wall}，\Emph{即}犹太司祭制的虚伪，已注定要因基督的来临而被推倒；因为他说：\BibleQuote{弟兄们，我原不知道他是大司祭，因为经上记载：不可诅咒你百姓的首长。}\BibleRef{宗 23:3}然而，完全确定的是，他既在那个民族中长大，并在那里受过法律的训练，就不可能不知道审判他的是大司祭；他也不能藉着声称自己对此无知，而欺骗那些对他非常熟悉的人。\par

14. 这些关于忍耐的训诫，应当时常存留于内心习惯性的纪律中；而那种防止以恶报恶的仁爱，也必须始终充分地怀持于性情之内。同时，在纠正他人时，即使违背其意愿，也有许多事需要以一种怀有仁爱之心的严厉去完成；对于这些人，我们应当顾念的是他们的福祉，而非他们的意愿；而圣经，已极明确地称赞了官长身上的这种德行。因为，在纠正儿子时，即便带着某些严厉，父亲的爱也决不会减少；然而，在那纠正的过程中，所行之事，却必为那一位似乎有必要藉痛苦而获医治的人所勉强而痛苦地接受。基于同一原理，若国家遵守基督宗教的训诫，那么甚至其战争本身，也不会不带有这样的仁爱意图：在那些顽抗的民族被征服之后，可以更易为他们预备条件，使他们在和平中享受那由虔敬与正义所结成的相互联系。因为，那被剥夺了他用以做恶之自由的人，他的被战胜其实于他有益；因为，再没有比犯罪者的那种好运更为真正的不幸了，藉此好运，有害的免罚得以维持，而邪恶的倾向，如同人内心的敌人，便得以加强。然而，人心乖谬而刚愎，当人们关心华美的宅邸，对灵魂的毁灭却漠不关心时；当巨大的剧场被建造起来，而美德的根基却被暗中破坏时；当奢侈挥霍的狂乱备受尊崇，而慈悲的工作反遭轻蔑时；当从富人的财富与充裕中，为演员们备办奢华的供给，而穷人却连生活的必需品也吝于施予时；当那位以其教义的公开宣告谴责公共恶德的天主，却被不虔敬的社群所亵渎，而这些社群所要求的，正是那样一种神祇，以至于那些使身体和灵魂都蒙受耻辱的戏剧表演，竟可以恰如其分地用于尊崇它们——这时，人心便以为世事顺遂。若天主容许这些事得势，祂便是在那容许中显示了更为严厉的不悦：若祂让这些罪行不遭惩罚，这种免罚本身就是一种更可怕的审判。反之，当祂推倒恶习的支柱，并使那些由富足滋养的情欲沦为贫困时，祂乃是出于仁慈而加以磨难。并且，若这种事是可能的，甚至善良者也能出于仁慈而发动战争，为要藉着将人们放纵的情欲置于轭下，而废除那些在公义的政府之下理应被根除或抑制的恶德。\par

15. 因为，若基督宗教谴责一切战争，那么福音中对那些寻求救恩之道的士兵所给的命令，就该是要求他们丢弃武器，并完全退出军役；然而，对这样的人所说的话却是：\BibleQuote{不要强暴待人，也不要讹诈人，你们应当满足于你们的粮饷。}\BibleRef{路 3:14}——这诫命要求他们满足于粮饷，便显然暗示并不禁止他们继续从军。因此，让那些说基督的教义与国家的福祉不相容的人，给我们一支由基督的教义所要求的那样的士兵组成的军队吧；让他们给我们这样的臣民、这样的丈夫和妻子、这样的父母和子女、这样的主人和仆人、这样的君王、这样的法官——总而言之，甚至这样的纳税人和税吏吧，即基督宗教所教导人应当成为的样子，然后让他们敢于说这教义与国家的福祉相悖吧；是的，更确切地说，让他们不要再犹豫承认：这教义若被遵守，便是国家的救恩。\par

% structure: p0020 source=h2 target=subsection reason=relative-heading-rank; base=h2

\subsection{第三章}

16. 然而，对于有人断言，\Place{罗马帝国}因某些基督徒皇帝而遭遇许多灾难，我该如何回答呢？这种笼统的指控是一种诽谤。因为，如果他们能从过去诸皇帝的历史中更清楚地引用某些无可争辩的事实来支持这种说法，那我也可以提及在其他非基督徒皇帝统治下发生的类似、甚或更大的灾难；从而使人们明白，这些灾难要么是由于人的过错，而非宗教之过，要么不能归咎于皇帝本身，而应归咎于那些皇帝离开他们便一事无成的人。至于\Place{罗马共和国}开始衰落的年代，有充分的证据；他们自己的文献对此说得明明白白。远在基督之名尚未照耀大地之前，就有人这样论及\Place{罗马}：\Quote{啊，可买卖的城，若它能找到买主，就注定速速灭亡！}那位最杰出的罗马史家，在其写于基督来临之前的论\WorkTitle{喀提林阴谋}的书中，清楚地说明了\Place{罗马}人民的军队从何时开始变得放纵和酗酒；开始高度估价雕像、绘画和浮雕花瓶；从个人和国家那里强夺这些东西；劫掠庙宇并玷污一切神圣与世俗的事。因此，当那腐败而放纵之风尚的贪婪与攫取暴力，既不放过人，也不放过他们所奉为神明者时，国家那著名的荣耀与安定便开始衰落了。从那时起，最恶劣的罪恶发展到何种地步，帝国的邪恶以何等大的危害继续损害人类的利益，说来话长，无法一一重述。让他们听听他们自己的讽刺诗人以戏谑却真实的笔调这样说道：——\par

从前贫穷，因而贞洁，我们的主妇，\newline{}在往昔岁月中，奢华无立足之处，\newline{}在低矮的房屋和裸露的土墙里；\newline{}她们的双手白天劳作不歇，\newline{}清贫的睡眠供给了宁静的夜；\newline{}当匮乏紧逼，饥饿使她们束紧腰带，\newline{}正值\Person{汉尼拔}兵临城下之时；\newline{}而今却放荡，慵懒地安享逸乐，\newline{}我们饱受和平带来的陈年痼疾\newline{}和穷奢极欲，其毁灭性的魔力\newline{}在向我们胜利的武装所征服的世界复仇。\newline{}任何罪恶，任何放荡的姿态都不再陌生，\newline{}自从贫穷，我们的守护神离去。\newline{}\par

那么，你为什么期待我列举更多由因繁荣而嚣张的邪恶所引入的灾祸的例子呢？因为在他们当中，那些比较仔细地观察事件的人看出，\Place{罗马}更有理由惋惜其贫穷的离去，而非其富足的离去；因为在贫穷中，其德行的完整得以保障，但通过其富足，可怕的腐化，比任何入侵者更可怖，已暴烈地占据了——不是城市的城墙，而是国家的心灵。\par

17. 感谢我们的主天主，祂给我们送来了前所未有的援助，以抵抗这些邪恶。因为若非\Person{基督}的十字架——可以说——建立在如此坚固的权威磐石上，被树立得如此崇高而强大，以致洪水无法冲走它，人类那骇人听闻的邪恶洪流岂不会将人卷向何方？它会饶过谁？又有什么深渊不会吞没其牺牲品？因此，藉着抓住它的力量，我们可以坚定不移，而不至于在邪恶计谋和邪恶冲动的巨大漩涡中失足并被吞没。当帝国正在彻底败坏的风俗和衰朽的古代宗教的邪恶深渊中沉沦时，特别重要的是，天上的权威应来救援，说服人们实践自愿的贫穷、节制、仁爱、正义和彼此之间的和谐，以及对天主的真正虔诚，和人生所有其他光辉而纯全的德行——不仅是为了以最高尚的方式度过今生，也不仅是为了确保地上国家中最完美的和谐纽带，而且也是为了获得永恒的救恩，并在永远长存的人民的神圣、天上的共和国中获得一席之地——一个共和国，信德、望德和爱德使我们成为其公民；因而，在此世旅居期间远离它时，我们可以耐心容忍那些希望通过对恶行不加惩罚来使那个共和国稳固的人，如果我们不能纠正他们的话；那个共和国是早期\Place{罗马}人凭借其德行创立并扩张的，那时，虽然他们没有对真天主的真正虔诚，不能通过具有拯救能力的宗教将他们带入永恒的公共福祉，但他们确实遵守了一种自身的正直，足以创立、扩张和维持一个地上的国家。因为在最富有、最辉煌的\Place{罗马}帝国中，天主已显示出，没有真正宗教的公民德行本身具有多大的影响，以便使人明白，当真正宗教加入这些德行时，人类就成为另一个国家的公民，那个国家的君王是真理，法律是爱，存续时间是永恒。\par

% structure: p0025 source=h2 target=subsection reason=relative-heading-rank; base=h2

\subsection{第四章}

18. 谁能不觉得他们试图将\Person{阿波罗尼乌斯}和\Person{阿普列乌斯}，以及其他最擅长魔法技艺的人，与\Person{基督}相比，甚至置于更高地位，这简直是荒谬可笑的呢？然而，这尚可较不恼怒地容忍，因为他们将这些人而非他们自己的神祇与祂相比；因为我们必须承认，\Person{阿波罗尼乌斯}比他们称为\Person{朱庇特}的那位无数粗俗不道德行为的始作俑者和实施者，品格要高尚得多。他们回答说，\Quote{关于我们神祇的这些传说，}\Quote{都是寓言。} 那么，他们为什么继续赞美共和国那种奢靡、放荡且明显亵渎的繁荣呢？这种繁荣发明了神祇的这些卑劣罪行，不仅让它们作为寓言传到人们耳中，还在剧院里展示在人们眼前；在剧院中，神祇们乐于看到公开进行的罪行比他们的神祇更多，而这些罪行本是为尊崇他们而公开犯下的，他们却本应惩罚那些甚至容忍这类表演的崇拜者？他们又回答说，\Quote{但是，}\Quote{那些不是他们自己的神祇，它们的神礼是根据这类寓言的虚假编造来庆祝的。} 那么，通过实践这类可憎之事以求取悦的，到底是谁呢？因为确实，基督信仰揭露了那些魔鬼的邪恶和狡诈，它们的魔力也通过魔法技艺欺骗了人们的心灵，并因此将这一点昭告于世，且分辨出圣天使与这些恶毒敌对者之间的区别，警告人们提防它们，同时指示了如何做到这一点——基督信仰因此被称为共和国的敌人，好像即使可以通过它们的帮助获得一时的繁荣，任何逆境都不如通过这种手段获得的繁荣可取似的。然而天主愿意防止人们对此困惑；因为在相对黑暗的\WorkTitle{旧约}时代——其中隐藏着\WorkTitle{新约}——祂以如此显著的此世繁荣使那敬拜真天主而弃绝假神的第一个民族卓然出众，以致任何人都能从他们的情况看出，繁荣并非由魔鬼支配，而只属于那位天使事奉、魔鬼畏惧的天主。\par

19. \Person{阿普列乌斯}（我宁愿谈论他，因为作为我们的同胞，他更为我们非洲人所知）虽然出生于一个有些知名的地方，且是一位受过优越教育且极具口才之人，尽管拥有所有\OriginalTerm{巫术}{magical arts}，却从未成功获得，我且不说至高权力，甚至帝国中的任何下级官职。他作为一个哲学家，自愿轻视了这些事物似乎可能吗？他身为一省的祭司，如此野心勃勃追求伟大，以至于他举办角斗士表演，为那些在竞技场与野兽搏斗的人提供服装，并且为了在\Place{科亚}镇——他妻子出生地——树立自己的雕像，竟与一些反对此提议的公民对簿公堂，然后，为了不让后人遗忘此事，他发表了那次演讲并将其出版。因此，就世俗繁荣而言，那位巫术师为成功竭尽了全力；由此可见，他之失败并非因为不愿，而是因为他不能做得更多。同时我们承认，他以出众的口才为自己辩护，反驳那些将施行巫术的罪行归咎于他的人；这使得我对他的赞美者感到惊奇，他们一面断言他通过这些巫术行了一些奇迹，一面却试图提出与他为自己所作的无罪辩护相矛盾的证据。那么，让他们检查一下，究竟是他们提出了真实的证词，而他对自己明知是虚假的事作了辩护。对于那些仅为追求世俗繁荣或出于可咒的好奇心而施行巫术的人，以及那些虽未行此类技艺却仍以危险的赞赏赞美它们的人，我愿劝诫他们，如果他们明智的话，要留心观察：没有任何此类技艺，牧羊人\Person{达味}是如何被提升到君王的尊位的；圣经上忠实地记载了他的罪过和功德行为，使我们既知道如何避免冒犯天主，也知道在冒犯了祂之后，如何平息祂的义怒。\par

20. 至于那些为激发人的惊奇而施行的奇迹，那些将异教的巫术师与\Emph{圣先知}相比的人大错特错了，因为圣先知凭其伟大奇迹的盛名完全超越他们。如果他们再将他们与\Person{基督}相比，错误就更大了；那些先知们——他们无与伦比地胜过一切名号的巫术师——曾预言祂将以由\Person{童贞玛利亚}所生而取得的人性，以及祂从未与\Person{圣父}分离的天主性来临！\par

我看出我写了一封极长的信，却还没有说尽所有关于\Person{基督}的事，这些事本可应对那些因心智迟钝而无法领悟神圣事物的人，或那些虽禀赋敏锐，却因喜好争辩和心灵长期固守谬误的偏见而无法辨认真理的人。然而，请留意任何使他们反对我们教义的事，再写信给我，以便如果主帮助我们，我们可以藉书信或论著答复他们所有的异议。愿你藉着主的恩宠与仁慈，在祂内得享幸福，我尊贵而堪受尊敬的阁下，我深爱而切望的儿子！\par

\SourceRecord{Translated by J.G. Cunningham.；Nicene and Post-Nicene Fathers, First Series；Vol. 1.；Edited by Philip Schaff.；Buffalo, NY: Christian Literature Publishing Co.,；1887.；<http://www.newadvent.org/fathers/1102138.htm>.}

% structure: p0001 source=h1 target=section reason=page-title

\section{书信第一三九封（公元412年）}

\Emph{致\Person{马塞利努斯}，我公义卓越的主，我至爱且切切思念的儿子，\Person{奥古斯丁}在主内致候。}\par

1. 我急切期待阁下应允寄送的卷宗，并渴望尽快能在\Place{希波}的教会宣读它们，若可能，也要在教区内所有已建立的教会宣读，使所有人都能听到并彻底了解那些认罪的人；他们认罪不是因为敬畏天主使之悔改，而是因为审判官的严厉击碎了他们极其残暴之心的刚硬——其中一些人承认杀害了一位司铎\Person{雷斯提图特}，并弄瞎并伤残了另一位\Person{依诺森提乌斯}；另一些人不敢否认他们可能知道这些暴行，尽管他们说他们不赞同，并坚持\OriginalTerm{裂教}{schism}的邪恶不虔，与如此众多的凶残恶棍为伍，却以不愿被他人罪行玷污为借口，离开了天主教会的平安；还有一些人宣称，即使天主教真理和\OriginalTerm{多纳特派}{Donatists}的错谬已向他们证明了，他们也不会弃绝\OriginalTerm{裂教者}{schismatics}。天主乐意托付给阁下勤勉的这项工作，极为重要。我心中的愿望是，许多类似的多纳特派案件能由你像这些一样审理并判决，从而这些人的罪行和疯狂顽固能经常被揭露；并且记录这些审理的卷宗得以公布，让所有人都知晓。\par

至于阁下信中所说，您不确定是否应命令将上述卷宗在\Place{特奥普雷皮亚}公布，我的答复是：若那里能聚集大量听众，就照此执行；若不能，则须提供另一处更常有人聚集的地方；然而，无论如何不可省略此事。\par

2. 至于对这些人的惩罚，我恳求你使他们所受的不至于死刑，尽管他们已亲口承认犯下如此重大的罪行。我这样请求，是出于对我们自己良心的关切，也为由此而显示的天主教宽仁作证。因为他们的认罪给了我们一个特别的益处，就是天主教会找到了一个机会，得以维持并向她最暴烈的仇敌展现宽容；因为在如此残酷的罪行已发生的情况下，凡不处以死刑的惩罚，都将被众人视为源于极大的仁慈。尽管在我们团体内的一些人看来，他们的心被这些暴行所扰动，认为这样的处置与罪行不相称，甚至显得有些软弱和失职，然而，当那些在事件刚发生时通常被过度激起的情绪随着时间平息之后，向有罪者显示的仁惠将发出最显赫的光辉，人们将更乐意阅读这些卷宗并向他人展示，我公义卓越的主，至爱且切切思念的儿子。\par

我的圣洁弟兄兼同仁主教\Person{波尼法爵}就在那里，我已通过与他同行的执事\Person{佩雷格里努斯}转交了一封指示信；请接受这封信，视如我在场。凡经你们共同判断认为有益于教会的事，靠着主的帮助，就行吧；主能在如此众多的大恶中慈悲地援助你们。他们的一位主教\Person{马克罗比乌斯}，此时正带领成群结队的可怜男女四处奔走，为自己打开了那些[多纳特派的]教堂——这些教堂此前因些许恐惧而曾一度关闭。然而，由于我向你的爱所推举并再次衷心推举的\Person{斯彭德乌斯}——显赫的\Person{塞勒}的代理——的到场，他们的放肆确实有所收敛；但现在，自从他前往\Place{迦太基}之后，马克罗比乌斯甚至在他的地产范围内打开了[多纳特派的]教堂，并在其中聚集会众崇拜。此外，他身边还有\Person{多纳图斯}，一名执事，曾被他们\OriginalTerm{重洗}{rebaptized}，甚至在他还是教会地产的一个佃户时就已是如此；此人作为头目牵连于那次[对\Person{依诺森提乌斯}]的暴行。当他与之为伍时，谁能知道他的随从会是何等货色？若对这些人的判决将由总督大人，或由你们二人共同宣布，且万一他坚持施以死刑——尽管他是基督徒，且据我们有机会观察，他并非倾向于如此严厉——如果他的决心迫使如此，那就命令将我认为有责任就此主题分别写给你们的那些信，在审理仍在继续时呈递上来；因为我惯常听说，法官有权减轻刑罚，并判处比法律规定的更轻的处罚。然而，如果尽管有我这些信，他仍拒绝答应这一请求，至少让他允许将这些人暂缓羁押；我们将努力从皇帝们的仁慈中获得这一许可，以便殉道者们的苦难——这些苦难本应为教会增添光辉荣耀——不致被其仇敌的鲜血玷污；因为我得知，在\Place{阿瑙尼亚}谷，圣职人员被异教徒杀害，如今被尊为殉道者，皇帝欣然批准了一项请求，即不对已经发现并被监禁的凶手处以死刑。\par

关于那些论述\OriginalTerm{婴孩圣洗}{baptism of infants}的书籍，我曾将其原稿寄给阁下；我已忘记究竟是何缘故又从您那里收回了它们；或许是因为，经过检查，我发现了其中的缺陷，并希望作某些修正，但由于特殊的阻碍，至今尚未来得及完成。我也必须承认，那封打算附在这些著作后致您的信，是我与您在一起时开始口授的，至今仍未完成，自那以后几乎没有再添加什么。不过，如果我能向您陈述我日夜不得不分心于其他事务的辛劳，您定会深深同情我，并惊讶于那不容拖延的众多事务，它们使我心神不宁，阻碍我完成您以祈求和劝勉所敦促我的那些事；这些恳求正是向一个极愿效力、却又因力不从心而无法言喻地忧伤的人发出的。因为当我从那些迫切必要的事务中获得片刻闲暇时——那些人如此催迫我为他们服务，使我既无法逃避，也不能忽视——总有必须优先口授给书记员的主题，因为它们与当下时刻紧密相连，不容拖延。这类事务之一是\WorkTitle{会议录摘要}，这是一项极其费力但必要的工作，因为我看到，没有人会试图阅读如此卷帙浩繁的记录；另一项是致\OriginalTerm{多纳图派平信徒}{Donatist laity}论及上述会议的信，我为此劳作了数夜，刚刚完成；另一项是撰写两封长信，一封致您，我所爱的朋友，另一封致杰出的\Person{沃鲁西亚努斯}，我想你们二人都已收到；另一项是一部正在撰写的书，致我们的朋友\Person{奥诺拉托斯}，答复他在一封信中向我提出的五个问题——您看得出，我对他责无旁贷，须迅速回复。因为爱对待子女，如同保姆对待孩童，并非按照对其情感深浅的顺序来关注他们，而是依照每一事件的急迫程度；她优先照顾软弱的，意欲将更强壮者所具有的力量分施给他们，对强壮者她暂且略过，并非轻视，而是因对他们，心中安稳。这类紧急情况，迫使我动用书记员去书写那些使我无法用他们来从事更契合我内心热望之工作的主题，总是不可避免，因为即便在事务堆积中——这些事务或是出于他人的意愿，或是迫于需要，都违背我的本意将我卷入——我极难获得哪怕一点点闲暇；而我究竟该怎么办，我实在不知道。\par

您已听到我的重负，我愿您以祈祷与我一同寻求解脱；但同时，我不希望您停止像以往那样，如此殷切而频繁地劝诫我；您的言语并非毫无效果。同时，我向阁下推荐一座建立在\Place{努米底亚}的教会；为了它目前的急需，我的圣善兄弟与同主教\Person{德尔斐努斯}，已受那些在该地区分担劳苦与危险的弟兄及同主教们的差派。关于此事，我不再多写，因为他来见您时，您将从他的口中听到一切。其余所有必要细节，您将在指令信中读到，这些信或由我现在交给\OriginalTerm{司铎}{presbyter}，或由执事\Person{佩雷格里努斯}转交，因此我无需在此重复。\par

愿您的心在基督内永远坚强，当受尊崇的阁下，极受敬爱与渴念的儿子！\par

我向阁下推荐我们的儿子\Person{鲁菲努斯}，\Place{锡尔塔}的\OriginalTerm{教长}{Provost}。\par

\SourceRecord{Translated by J.G. Cunningham.；Nicene and Post-Nicene Fathers, First Series；Vol. 1.；Edited by Philip Schaff.；Buffalo, NY: Christian Literature Publishing Co.,；1887.；<http://www.newadvent.org/fathers/1102139.htm>.}

% structure: p0001 source=h1 target=section reason=page-title

\section{第143封信（主后412年）}

致我尊贵的主、公正卓越的、我深爱的儿子\Person{玛尔策利努斯}：\Person{奥思定}在主内致意。\par

1. 我渴望回复我通过我们圣善的弟兄、我的同僚主教\Person{玻尼法爵}从你那里收到的信，我寻找过它，但没有找到。然而，我记起你在信中问过我，\Person{法郎}的术士们在\Place{埃及}所有的水都变成血之后，怎能找到任何水来模仿这个奇迹。通常有两种方式来回答这个问题：要么可能从海里取水，要么，更可信的是，这些灾祸并未临到以色列子民所住的地区；因为\WorkTitle{圣经}叙述中某些部分对此有清楚、明确的陈述，这使得我们可以在陈述被省略的地方假定这一点。\par

2. 在由司铎\Person{乌尔巴努斯}带给我的你的另一封信中，提出了一个问题，它不是出自\WorkTitle{圣经}，而是出自我自己的一本书，即我所写的\WorkTitle{论自由意志}。然而，对于这类问题，我并不花费太多精力；因为即使受到反对的陈述不容无可辩驳的辩护，那也仅是我个人的；它并非那位作者的言语，拒绝祂的话乃是不虔敬的，即使我们由于误解其含义而对我们对它们的解释加以拒绝。因此我坦然承认，我努力成为那些有所进步而写作，并藉着写作而获得更大进步的人之一。所以，若因疏忽或无知，我说过任何可以有充分理由受谴责的话，不仅别人能发现这一点，我自己也同样能（因为如果我正在进步，我至少应当在被指出之后看到它），这样的错误不应被看作令人惊讶或悲哀，反而应被原谅，并成为向我道贺的时机，当然不是庆贺我犯了错，而是庆贺我弃绝了错误。因为一个渴望别人犯错，以便他自己犯错的事实不被发现的人，其自爱中存在一种极度的偏执。多么更好和更有益的是，在他犯错的地方别人不犯错，这样他就可以藉他们的劝告而从错误中被拯救出来，或者，如果他拒绝这样，至少可以在他的错误中没有追随者。因为，如我所渴望的，如果天主准许我，在一部为此明确目的而写的著作中，汇集并指出我的书中所有最令我不悦的一切，那么，人们将会看到我在判断自己的事情上是多么远离偏袒。\par

3. 至于你，你这般热切地爱我，如果，在反对那些出于恶意、无知或卓越的智慧而责难我的人时，你坚持认为我在著作中没有任何错误，你就是在从事一项无望的事业——你承担了一个糟糕的辩护，在这件事上，即使我自己作裁决，你也会轻易被击败；因为，让最亲爱的朋友们把我看作并非我实际的样子，这对我并不是一件乐事，因为毫无疑问，如果他们爱的不是真正的我，而是非我的样子，那么他们爱的就不是我，而是以我的名义爱着另一个人；因为他们了解我，或相信关于我的真实的事，我才为他们所爱；但当他们将我不知道自己拥有的东西归于我时，他们就是爱着另一个他们凭空想象的人。罗马演说家之首\Person{西塞罗}论到某人说：\Quote{他从未说过任何他想要收回的话。}这一赞辞，虽然看似最高的赞誉，然而更可能适用于一个十足的愚人，而非一个完全智慧的人；因为这正是白痴的真实写照：他们越是荒谬、愚蠢，越偏离众所公认的观点，就越不可能说出任何他们想要收回的话；因为对于一个智慧的人来说，后悔一句邪恶、愚蠢或不智的言语是惯常之事。然而，如果那句引语被善意地理解为要我们相信，有人因凡事都说得很智慧，以至于他从未说过任何想要收回的话——那么按照纯正的虔诚，我们更应当相信那些受圣神感动而说话的属天主的人，而非\Person{西塞罗}所称许的人。就我而言，我与这种卓越相去甚远，如果我从未说过任何想要收回的话，那必定是因为我更类似于白痴而非智慧人。其著作最有资格享有最高权威的人，是那从未说过任何——我不说他想收回，而是他有义务收回的话的人。那尚未达到这一境界的人，由于无法藉智慧占据首位，就当藉着谦卑安于次位。既然他已尽所有谨慎仍无法使所有用语都免于可公义地后悔之虞，那么当他如今可能已发现什么不当说的话时，就当承认自己的懊悔。\par

4. 因此，既然我所说的那些——如果可能的话我希望收回的——话语，并非如我最亲爱的朋友们所设想的那样很少或者根本没有，反而可能比我的敌人们想象的还要多；我并非因受称赞而高兴，如\Person{西塞罗}所言：\Quote{他从未说过一句他希望能收回的话}，而是因\Person{贺拉斯}的话而深感痛心：\Quote{话一出口，无法收回}。这就是为什么我将我所写的那些关于\BibleBook{}{创世纪}和\WorkTitle{论天主圣三}这些困难而重要问题的著作，长久地留在身边，甚至超出你们所愿或所耐心承受的时间，希望的是，如果说有些内容难免会被合理地挑出毛病，那么这些毛病至少会比那因急躁催促而未经深思熟虑便出版时更少一些；因为你们的信函（我们的圣洁弟兄兼同任主教\Person{弗罗伦提乌斯}已就此事给我来信）表明，你们现在急于这些著作的出版，以便能在我的有生之年亲自为之辩护，那时，或许它们会在某些方面开始受到攻击，无论是来自敌人的吹毛求疵还是朋友们的误解。你们这样说，无疑是认为书中没有任何可公正地受到指责的地方，否则你们就不会催促我出版它们，而是劝我更仔细地修订它们。然而我更注目于那些真正公正而严厉的审判者，我渴望首先在自身与这些审判者之间确立稳固的根基，好让他们唯一能责难的，就是那些虽经我极为勤勉的审察，却仍无法注意到的问题。\par

5. 尽管我刚才说了这些，我仍准备为我的\WorkTitle{论自由意志}第三卷中的一句话进行辩护。在那句话中，论及理性的实体时，我这样表达了我的观点：\Quote{灵魂被指定在罪进入之后占据那本性低于自身的身体，它并非按照其自由意志绝对地管理自身的身体，而只是就宇宙法则所允许的范围内管理。}我特别提请那些认为我在此将人的灵魂确定为通过父母繁衍而来，或认为灵魂因在更高天界生活中犯下罪过而被罚禁锢于可朽坏身体中的人的注意。我请他们注意，我是如何仔细地斟酌了这些措辞，使之在确定我所确信的这件事上——即原祖犯罪之后，所有其他的人都在那罪恶的肉身中出生并继续出生，为了治愈这肉身，\BibleQuote{罪恶肉身的形状}以主的位格来临——的同时，也措辞如此，避免对我在后文中阐述并区分的四种意见中的任何一种作出定论——既不试图认定某种意见优于其他，而只是目前先处理所讨论的问题，保留对这些意见的考量，以便无论其中哪一种为真，都应毫不犹豫地赞美天主。\par

6. 因为，无论所有灵魂是由原祖繁衍而来，还是在每个个体中特别受造，或是受造于身体之外而后被送入身体，或自己自愿进入身体，毫无疑问，这赋有理智的受造物，即人的灵魂——被指定占据一个低级的，即属土的身体——在罪进入之后，并非按照其自由意志绝对地管理自身的身体。因为我没有说\Quote{在他犯罪之后}或\Quote{在他犯了罪之后}，而是说\Quote{在罪进入之后}，这样，无论后来若是可能，理性如何确定这罪是他自身的，还是人类原祖的罪，人们都可以看出，我说\Quote{灵魂被指定在罪进入之后占据一个低级身体，并非按照其自由意志绝对地管理自身的身体}这句话，是真实的；因为\BibleQuote{肉身的贪恋相反圣神}，\BibleQuote{我们在这肉身内叹息，被重压}，而\BibleQuote{可朽坏的身体重压着灵魂}；总之，谁能历数肉身所引发的一切邪恶呢？这一切必然结束，当\BibleQuote{这可朽坏的穿上了不可朽坏的}，使得\BibleQuote{那必死的被生命所吞灭}的时候。因此，在那将来的状态中，灵魂将以绝对自由的意志管理其属神的身体；但在此期间，它的自由并非绝对，而是受到宇宙法则的制约，根据这些法则，凡经历出生的身体必经历死亡，成长至成熟则衰退于衰老。因为原祖的灵魂，在罪进入之前，确实以完全自由的意志管理他的身体，尽管那身体还不是属神的，而是属魂的；但在罪进入之后——即在那肉身中犯罪之后，由此罪身将繁衍——有理性的灵魂被如此指定占据低级身体，以至于它并非以绝对自由的意志管理自身的身体。在我看来，即使婴儿在犯下任何自身之罪以前，他们也分受了罪恶的肉身，这由他们也需要在圣洗中应用那以\BibleQuote{罪恶肉身的形状}来临的主所提供的良药而得到治愈这件事，得到证明。但即使那些不赞成此观点的人，也没有正当理由对我书中引用的那句话感到冒犯；因为，如果我没有弄错的话，确定无疑的是，即使那软弱不是来自罪而是来自本性，无论如何，在罪进入之后，具有这种软弱的身体才开始产生；因为\Person{亚当就是这样被造的，而且他在犯罪之前没有生育任何后代。}\par

7. 所以，让我的批评者们去寻找其他段落来责备吧，不仅是在我其他仓促发表的著作中，也包括在我这些\WorkTitle{论自由意志}的书卷中。因为我绝不否认，他们可能在寻找中发现有机会给我带来益处；因为，如果这些书已经流传甚广，现在无法订正，我自己还在世，倒还可以。然而，我精心选择那些词语，是为了避免在关于灵魂起源的四种看法或理论中承诺任何一方；只有那些认为我对如此隐晦之事迟疑不定是可指责的人，才会去责备这些话。对于这些人，我不为自己辩护说我认为不对这个主题发表任何意见是正确的，因为我毫不怀疑：灵魂是不死的——不是如同唯独不死不灭的天主那样的不死，而是以其独特的方式——灵魂是受造物而非造物主实体的一部分，以及其他我认为关于其本性最为确实的事情。但是，既然这个最为神秘的主题——灵魂起源——的隐晦迫使我如此行事，就让他们向我伸出友谊之手，承认我的无知，并渴望知道有关这主题的任何真理吧；让他们，如果能够，教导或向我证明他们通过健全理性的运用所学到的，或是根据圣经中无可置疑的明白见证所相信的。因为，如果发现理性与圣经的权威相矛盾，那么理性只是以真理的假象欺骗人，无论它多么敏锐，因为其推演在这种情况下不可能是真的。另一方面，如果有人提出某种东西，声称具有圣经的权威，却与理性最明显、最可靠的见证相悖，那么这样做的人是误解了他所读的内容；他将自己对圣经的解释，而非未能发现的圣经真义，立为真理的对立面；他所宣称的，不是他从圣经中发现的，而是他作为解释者在自身中发现的。\par

8. 让我举一个例子，我恳请你认真注意。在\BibleBook{训道}{篇}结尾附近，作者论及人因死亡而分解，灵魂与身体分离时，写道：\BibleQuote{其实尘土要归于他的故土，灵魂要归于天主，因为原是天主所赐。}如此权威所依据的陈述无疑是真实的，绝无意欺骗任何人；然而，如果有人愿意如此解释它，以帮助他试图支持\OriginalTerm{灵魂传衍}{propagation of souls}的理论，即所有其他灵魂都源于天主赐予第一个人的那一个灵魂，那么经文中以\Quote{尘土}之名所说的关于身体的话（因为显然，\Quote{尘土}和\Quote{灵魂}在此段中指的就是身体和灵魂），似乎对他的观点有利；因为他可以断言，经文说灵魂要归于天主，是因为它源于天主赐予第一个人的那个灵魂的原始根苗，正如说身体要归于尘土，是因为它源于第一个人的那个由尘土造成的身体的原始根苗一样；因此他可以辩称，从我们完全确知的身体方面，我们应当相信关于灵魂那隐秘不可见的事；因为关于身体的原始根苗没有异议，但关于灵魂的原始根苗却有。在这样被引为证据的经文中，对二者都做了陈述，仿佛各自回归其本原的方式完全相似：一方面，身体归于原来的土中，因为它是在第一个人受造时从那里取出的；另一方面，灵魂归于天主，因为当天主将生命的气息吹进祂所造的人的鼻孔，人就成了有灵的生物\BibleRef{创 2:7}，天主赐予了它。从此以后，每一部分都从父母相应的部分繁衍下去。\par

9. 然而，如果灵魂起源的真实情况是，天主赐予每个人一个灵魂，不是从那第一个灵魂传衍而来，而是以其他方式受造，那么\BibleQuote{灵魂要归于天主，因为原是天主所赐。}这句话，与这种观点同样一致。那么，关于灵魂起源的另外两种看法，似乎是这节经文唯一排除的。因为，首先，关于每个人的灵魂在其受造之时各自在其内造成的看法，人们认为，若是这样，灵魂应该说归于创造它的天主，而不是赐予它的天主；因为\Quote{赐予}一词似乎意味着，能被赐予的东西已经独立存在。有些人进一步强调\Quote{归于天主}这句话，他们说：它怎能回到它从未去过的地方呢？因此他们主张，如果相信灵魂以前从未与天主同在，那么经文应当说\Quote{它往}，或\Quote{它去}，或\Quote{它离开}，而不应说它\Quote{归于}天主。同样，就每个灵魂自行潜入其身体的看法而言，难以解释这种理论如何与天主赐予它的陈述相协调。因此，这段经文的措辞，多少不利于这两种看法：一种是假定每个灵魂在其自身身体内受造，另一种是假定每个灵魂自行进入其身体。但是，要显示这经文与另外两种看法之一是一致的，却毫无困难。这两种看法是：所有灵魂都由第一个受造的灵魂传衍而来；或者，灵魂在天主那里受造并预备好，按需要赐给每个身体。\par

尽管如此，即使每个灵魂在其各自身体内受造的理论可能并未被这节经文完全排除——因为如果其支持者断言，经上说天主赋予人精神（或灵魂），与经上说天主赋予人眼睛、耳朵、手或其他此类肢体是同一方式，而这些肢体并非天主在别处造好并存留以供赋予，即附加并连接到我们的身体上，而是由天主在那被赋予的身体中造成的——我看不出能回应什么，除非或许可以用圣经的其他章节或有效的推理来反驳此观点。同样，那些认为每个灵魂自行流入其身体的人，将\BibleQuote{天主赋予它}这句话理解为经上所说\BibleQuote{天主任凭他们随从心中的情欲，行污秽的事}的那种意义。因此，只有一个词似乎与每个灵魂在其身体内受造的理论明显无法协调，即\Quote{返回}一词，在\BibleQuote{返回天主}的表达中；因为灵魂如何能返回它从前未曾与之一同的那位呢？仅仅这一个词就使四种观点中这一观点的支持者感到困惑。但我并不认为应当仅仅因这一个词就认为此观点已被驳倒，因为有可能表明，按照圣经语言的通常用法，使用\Quote{返回}一词可能完全正确，以表示由天主所造的精神返回天主，并非因为它在与身体结合之前曾与天主同在，而是因为它从天主的创造大能获得了存在。\par

我写下这些，是为了表明，任何愿意坚持并捍卫关于灵魂起源这四种理论中任何一种的人，必须提出，要么来自为教会权威所接受的圣经中那些不能作任何其他解释的章节——正如\BibleQuote{天主造了人}这一陈述——要么建立在前提上的推理，这些前提如此明显真实，以至于对其质疑便是疯狂，例如\Quote{只有活着者才能认知或犯错}这一陈述；因为像这样的陈述不需要圣经的权威来证明其真实性，仿佛人类的常识本身没有以其透明而有力的理性自动宣告其真实性，以至于凡反驳它的人必被视为无可救药的疯子。若有人在探讨\Emph{灵魂起源这一极为隐晦的问题}时能够提出这样的论据，就请帮助我这无知的人；但若他不能，就请克制自己，不要指责我在这一问题上的犹豫。\par

至于圣玛利亚的童贞，如果我关于此题目所写的不足以证明其可能性，那么我们必须拒绝相信关于任何身体中发生奇迹的任何记载。然而，如果对这一奇迹的异议在于它只发生过一次，请询问对此仍感困惑的朋友，是否从世俗文献中可以举出一些事件的例子，它们像这一事件一样是唯一的，却仍被相信——不是像头脑简单的人相信寓言那样，而是以接受事实历史的那种信仰——我恳请你问他这个问题。因为如果他说这类事件在这些著作中找不到，那么应当有人向他指出这类例子；如果他承认这一点，问题就因他的承认而解决了。\par

\SourceRecord{Translated by J.G. Cunningham.；Nicene and Post-Nicene Fathers, First Series；Vol. 1.；Edited by Philip Schaff.；Buffalo, NY: Christian Literature Publishing Co.,；1887.；<http://www.newadvent.org/fathers/1102143.htm>.}

% structure: p0001 source=h1 target=section reason=page-title

\section{第144封书信}

致我可敬且公正受尊敬的诸位大人，即\Place{希尔塔}的全体居民，各阶层；亲爱的且深切思念的弟兄们，主教\Person{奥古斯丁}致以问候。\par

1. 如果说曾使我极为忧伤的，你们城中的那件事现在已经消除了；如果说那些抵抗最明显且（我们可称之为）众所周知的真理的顽固心肠，已经被真理的力量所征服；如果说和平的甜蜜已被品味，并且那导向合一的爱情，对于患病的眼睛不再是痛苦的根由，而是对于复原的眼睛成为光明与活力——那么，这乃是天主的工作，不是我们的；即使当我在你们中间时，伴随着我自己的讲道和劝勉，你们全体的如此显著的归化曾发生过，我也无论如何不会将这些成果归于人的努力。这行动与成功都属于祂，祂借着自己的仆人，藉着事物的标记向外召唤人的注意，并亲自藉着事物本身向内教导人。然而，你们中间发生的任何值得称赞的改变，都当归于我们，不归于那独行奇事的祂，这一事实难道不正是我们更应乐意应邀去拜访你们的原因吗？因为我们更应当急切地去观看天主的工作，而非我们自己的工作，因为我们自己若在任何工作中有所贡献，也不当归功于人，而当归功于祂；故此宗徒说：\BibleQuote{栽种的算不得什么，浇灌的也算不得什么，只在那使之生长的天主。}\par

2. 你们在信中提及了一个我也从古典文献中记得的事：\Person{克塞诺克拉特}因谈论节制的益处，竟突然使\Person{波勒莫}从放荡的生活转变为节制的生活，尽管此人不但一向放纵，当时且已酩酊大醉。如今，虽然此事如你们明智且真确地所理解的，并非归向天主，而是从自我放纵的奴役中获得解脱，然而即使在他身上发生的这点改善，我也不归于人的工作，而归于天主的工作。因为即使在我们本性中最低的部分——身体上，一切美好的事物，如美丽、活力、健康等，都是天主的工作，自然受造及其完善都归于祂；那么，灵魂的卓越品质更非他人所能赋予，这岂不更为确然！人的狂妄想像中，有什么能比这种念头更充满骄傲和忘恩负义的呢：即认为，虽然只有天主能赋予身体美貌，但赋予灵魂纯洁却属于人？在\BibleBook{智慧}{篇}中写道：\BibleQuote{我深知，除非天主赐予，无人能有节制，并且知道这是谁的恩赐，乃是真智慧的一部分。}因此，如果\Person{波勒莫}在改变放荡为节制的生活时，明白了这恩赐来自何方，从而弃绝异教的迷信，崇拜那位神圣的赐予者，那么他不但成为了有节制的人，更成为了真正智慧和有救恩的信徒，这不仅使他在此生修德，更将在来世获得永生。那么，我岂敢将你们的归化，或你们人民中你们现在向我报告的归化，归功于我自己呢？那些归化，当我不在向你们讲话，甚至不在场时，无可置疑地是由天主的德能在所有真正经历归化的人身上完成的。因此，你们应首先知道这一点，并以虔敬的谦逊加以默想：我的弟兄们，你们要感谢天主！敬畏祂，免得你们倒退；爱慕祂，好使你们前进。\par

3. 然而，如果对人的畏惧仍然使一些人暗中与基督的羊群分离，而对别人的畏惧则迫使他们伪装和好，那么我警告所有这些人都要思想：在天主台前，人的良心毫无遮掩，他们既不能欺骗祂作见证者，也不能逃避祂作审判者。但如果由于对自己得救的焦虑，关于基督羊群的合一问题使他们困惑，就让他们对自己提出这个要求——在我看来这是最公正的要求——即对于\OriginalTerm{天主教会}{Catholic Church}，也就是遍布全世界的教会，他们要更相信圣经的话语，而不是人舌头的诽谤。此外，关于在人中间产生的分裂（那些人，无论他们是怎样的人，绝不会使天主对\Person{亚巴郎}的应许落空：\BibleQuote{地上万民都要因你的后裔蒙受祝福}——这些应许，当他们听在耳中时，被当作预言相信了；但当摆在眼前，作为已成就的事实时，他们却不承认了），让他们暂且思索下面这个非常简短，但若我没有弄错的话，无可辩驳的论证：引致争端的那个问题，要么曾在海外的教会法庭审理过，要么没有；如果没有在这些法庭审理过，那么此事上的罪责丝毫也不能归咎于海外万民中基督的整个羊群，而我们与这些羊群在同道共融中喜乐欢欣，因此，他们与这些无罪的团体分离，就是一种不虔敬的分裂行为；反之，如果这问题已经在这些教会的法庭前审理过了，谁还不明白、不感觉到，不，谁还看不见，那些现在与这些教会断绝共融的人，就是在诉讼中败诉的一方呢？因此，让他们选择究竟该相信谁：是决定此案的教会审判官们，还是败诉者们的怨言？请明智地观察，他们要对这个简短而极其明晰的两难推理做出合理的答复，是多么不可能的事；然而，要使\Person{波勒莫}脱离放纵的生活，比之将他们从根深蒂固的错谬疯狂中驱逐出来，还要容易些。\par

请宽恕我，我高贵而可敬的诸位大人，我最亲爱的、深切思念的弟兄们，我给你们写了一封冗长胜过令人愉快的信，但依我想来，适于使你们受益而非奉承你们。至于来到你们这里，愿天主满全我们双方同怀的愿望！因为我无法用言语表达，但我确信你们会乐意相信，我怀着何等炽热的爱火切望见到你们。\par

\SourceRecord{Translated by J.G. Cunningham.；Nicene and Post-Nicene Fathers, First Series；Vol. 1.；Edited by Philip Schaff.；Buffalo, NY: Christian Literature Publishing Co.,；1887.；<http://www.newadvent.org/fathers/1102144.htm>.}

% structure: p0001 source=h1 target=section reason=page-title

\section{书信一四五（公元412年或413年）}

致\Person{阿纳斯塔修斯}，我圣洁而亲爱的弟兄，\Person{奥古斯丁}在主内向您致意。\par

我们的弟兄\Person{卢皮基努斯}和\Person{孔科尔迪亚利斯}——天主可敬的仆人——提供了一个最令人满意的向您致意您真诚尊贵的机会，即使我不写信，您也能从他们那里了解我们这里的一切。但我知道您在基督内多么爱我们，正如您知道我们在祂内多么热切地回报您的爱；深信，如果您见到他们，并且不可能不知道他们直接来自我们，与我们有着最亲密的友谊，却未曾从我这里收到一封信，那一定会令您失望。此外，我还欠您一封回信，因为据我所知，自从收到您上次来信后，我一直没有给您写过信；身陷如此繁重的挂虑，让我都不知道以前是否曾给您写过信。\par

我们渴望知道您的近况，以及主是否在此世所能赐予的范围内，给了您一些安息；因为\BibleQuote{若一个肢体受光荣，所有的肢体都一同欢乐}；所以根据我们几乎一贯的经验，当我们在焦虑中想到那些处于相对安宁的弟兄时，便如同在那种安宁中，我们自己也享受到了更平安、更宁静的生活，从而获得了不小的振奋。同时，当这无常生活中的烦恼挂虑倍增时，便迫使我们渴慕那永恒的安息。因为这个世界在顺境中比在逆境中对我们更为危险，当它诱使我们爱它时，比当它警告并迫使我们轻视它时，更须倍加防范。因为虽然\BibleQuote{世界上的一切}就是\BibleQuote{肉身的贪欲，眼目的贪欲，以及人生的骄奢}，然而，即使那些宁愿选择属灵的、看不见的、永恒的事物的人，尘世事物的甜蜜也会渗入我们的情感，并以诱人的魅力伴随我们履行本分的脚步。因为现世事物对我们软弱之人的控制力，恰恰与未来事物命令我们爱的更高价值成正比。但愿那些学会觉察并哀叹这一点的人，能够战胜并摆脱这尘世事物的势力！这样的胜利与解脱，若无天主的恩宠，单凭人的意志是绝不能获得的，因为人的意志，只要还受制于那占上风并奴役人的私欲偏情，就根本不能称为自由的；\BibleQuote{因为人被谁制胜，就是谁的奴隶}。圣子也亲自说过：\BibleQuote{如果天主子使你们自由了，你们的确是自由的}。\par

因此，法律通过教导和命令那些没有恩宠就无法完成的事，向人指出他的软弱，好让这被证实的软弱投奔救主，借祂的医治，意志才能做到在软弱中所不能的事。这样，法律便领我们到信德那里，信德使人更充足地获得圣神，圣神将爱德倾注在我们心中，而爱德成全法律。为此，法律被称为\BibleQuote{启蒙师}，在其威吓和严苛之下，\BibleQuote{凡呼号主名的人，必然获救}。但\BibleQuote{人若不信祂，又怎能呼号祂呢}？因此，对于信从并呼号祂的人，就赐予那赋予生命的圣神，免得那没有圣神的文字使他们死亡。而借着那赐给我们的圣神，天主的爱已倾注在我们心中了。所以，同一位宗徒的话：\BibleQuote{爱德就是法律的满全}便实现了。如此，法律对于依法使用的人是好的；所谓依法使用，就是明白法律为何而赐，并在其威吓之下投奔那使人自由的恩宠。凡是不知感恩地轻视这使不虔敬者成义的恩宠，却信赖自己的力量，仿佛凭此就能满全法律，不认识天主的正义而企图建立自己的正义的人，是不顺从天主的正义的；这样，法律对他来说，就不是赦罪的助佑，而是系紧他罪过的枷锁。并不是法律不好，而是罪借着那好的事物，在这些人的身上造成了死亡。因为借着诫命，那因诫命而知道自己所犯的罪是多么邪恶的人，便犯罪更重了。\par

然而，若有人仅仅出于对惩罚的恐惧而戒避犯罪，便自认为已战胜了罪恶，那是枉然的；因为，虽然那由恶欲所激发的行为未付诸实行，但那在人心内的恶欲本身仍是一个未制服的仇敌。谁若只要除去所畏惧的惩罚就愿意去犯那被禁止的罪，这样的人又能在天主眼前被视为无辜呢？因此，就在这意愿本身，凡渴望做非法之事，只是因为不能不受惩罚而止步的人，就已经犯下了罪；因为就他个人而言，他宁愿没有那禁止并惩罚罪恶的正义。而确实，如果他宁愿没有正义，谁能怀疑，他若能完全废除正义，就必会付诸实行呢？那么，一个如此敌视正义的人——若有权能，他便会废除其权威，以免受其威吓或惩罚——又怎能被称为义人呢？\par

因此，那仅因害怕惩罚而戒避犯罪的人，是正义的仇敌；但若因爱正义而不犯罪，他便将成为正义的朋友，因为那时他才是真正地害怕犯罪。因为只畏惧地狱之火的人，怕的不是犯罪，而是被焚烧；但憎恨罪恶如同憎恨地狱的人，才是害怕犯罪。这就是\BibleQuote{敬畏上主}，这敬畏是\BibleQuote{纯洁的，存留直至永远}。因为对惩罚的惧怕含有痛苦，并不在于爱内；圆满的爱把恐惧驱逐于外。\par

5. 再者，每一个人恨恶罪的程度，恰如他爱慕正义的程度；而他能做到这一点，不是借着法律以字面的禁令使他恐惧，而是借着圣神以恩宠医治他。那时，宗徒在劝勉中所吩咐的就实现了：\BibleQuote{我因你们肉身的软弱，就按常情对你们说：你们从前怎样将你们的肢体当作奴隶，献于不洁和不法，行不法的事；如今也要怎样将你们的肢体当作奴隶，献于正义，行圣善的事。} 那么，\Quote{正如}和\Quote{照样}这两个词的力量是什么，岂不就是：\Quote{正如没有恐惧强迫你们犯罪，而是对罪的欲望和罪中之乐驱使你们，照样，也不要让对惩罚的恐惧驱迫你们去度正义的生活；而要让在正义中寻得的喜悦和你们对正义的爱吸引你们去实践它}？在我看来，这甚至可以说是一种较为成熟的正义，但仍非完美的正义。因为他若没有想到，到他们那时能承受时，还有更当说的话，他就不会用\BibleQuote{我因你们肉身的软弱，就按常情对你们说}来作为这劝勉的开场白。因为对正义，理当有比人们通常向罪所献的更专心的事奉。肉身的痛苦阻止人，即便不能阻止犯罪的欲望，至少能阻止犯罪的行为；我们不易找到任何人，若确知刑罚的痛苦会立刻随罪行而来，还会公然犯那带给他不洁且非法满足的罪。但正义却应受到如此的爱慕，以致连肉身的痛苦也不能阻止我们实践正义的工程，反而，即便我们在残酷仇敌的手中，我们的善工也应在人前如此照耀，好使那些能从中获得喜乐的人，光荣我们在天之父。\par

6. 因此，那位最热诚爱慕正义的宗徒呼喊说：\BibleQuote{谁能使我们与基督的爱隔绝？是困苦吗？是窘迫吗？是迫害吗？是饥饿吗？是赤贫吗？是危险吗？是刀剑吗？正如经上所载：\InnerQuote{为了你，我们整日被置于死地，人看我们像待宰的羊。}然而，靠着那爱我们的主，我们在这一切事上，大获全胜，因为我深信：无论是死亡，是生活，是天使，是掌权者，是现存的或将来的事物，是有权能者，是崇高或深远的势力，或其他任何受造之物，都不能使我们与天主的爱相隔绝，即是与我们的主基督耶稣之内的爱相隔绝。} 请留意他并没有简单地说，\Quote{谁能使我们与基督隔绝？}而是指出我们与基督相连的凭借，说：\Quote{谁能使我们与基督的爱隔绝？}我们与基督相连，是借着爱，而非借着对惩罚的畏惧。再者，在列举了那些看似猛烈，却不足以造成隔绝的事物之后，他在结论中称那先前称为基督之爱的为天主的爱。而这\Quote{基督的爱}是什么，难道不就是对正义的爱吗？因为论到祂，经上说：\BibleQuote{祂由天主成了我们的智慧、正义、圣化者和救赎者，正如经上所记载的：\InnerQuote{凡要夸耀的，应因主而夸耀。}} 因此，正如那连肉身的痛苦也不能阻止他从事污秽情欲的可耻行为的人，是极端的邪恶；照样，那因畏惧肉身的痛苦也不能阻止他实践最光荣的爱的神圣工程的人，是极端的正义。\par

7. 这种对天主的爱，必须借不懈的虔敬默想来保持，\BibleQuote{是借着所赐与我们的圣神，已倾注在我们心中了，}好使\BibleQuote{凡要夸耀的，应因主而夸耀。} 因此，我们既然感觉自己贫乏，缺少那能真正满全法律的爱，我们就不应从自己的匮乏中去期待和要求这爱的富足，而应在祈祷中祈求、寻找、叩门，好使那位拥有\BibleQuote{生命的泉源}的，\BibleQuote{以祂殿中的盛筵饱饫我们，使我们畅饮祂乐河的水，} 如此，我们既被这满溢的洪流灌溉而复苏，不但能免于被忧伤所吞没，甚至还能\BibleQuote{在磨难中也欢跃，因为我们知道：磨难生忍耐，忍耐生老练，老练生望德，望德不叫人蒙羞；}——这并不是我们凭自己能做得到的，而是\BibleQuote{因为天主的爱，借着所赐与我们的圣神，已倾注在我们心中了。}\par

8. 我感到很高兴，至少能借这封信，说出当你在场时我未能说出的话。我写这些，不是针对你，因为你并不好高骛远，而是满足于卑微的事？而是针对某些过于归功于人的意志的人，他们设想，法律既已颁布，人的意志凭自身的力量就足以满全那法律，即使除了法律的教导之外，没有受到圣神所赋予的任何恩宠的协助；他们凭自己的推论，说服人可怜而又贫乏的软弱，使之相信我们无须祈祷以免陷于诱惑。\par

这并不是他们敢于公开这样说；但不论他们承认与否，这是他们学说的必然结论。因为倘使这完全属于人的意志的能力，而不需要恩宠的扶助就能实现，那么，为什么对我们说：\BibleQuote{醒寤祈祷，免得陷于诱惑；} 又为什么当祂教导我们祈祷时，依照这劝诫，规定了\BibleQuote{不要让我们陷于诱惑}的祈求呢？\par

我还有什么可说的呢？请向与你在一起的弟兄们致意，并为我们祈祷，使我们能得救，就是那所说过的救恩：\BibleQuote{不是健康的人需要医生，而是有病的人；我不是来召义人，而是召罪人。} 所以请为我们祈祷，使我们成为义人——这种成就是人完全无法凭自己达到的，除非他认识正义并愿意实践正义，而当人完全愿意时，就能立刻实现；但这意志的完全同意，除非他受圣神的恩宠医治和扶助，就永不能在他内存在。\par

\SourceRecord{Translated by J.G. Cunningham.；Nicene and Post-Nicene Fathers, First Series；Vol. 1.；Edited by Philip Schaff.；Buffalo, NY: Christian Literature Publishing Co.,；1887.；<http://www.newadvent.org/fathers/1102145.htm>.}

% structure: p0001 source=h1 target=section reason=page-title

\section{第146封书信（公元413年）}

致\Person{\OriginalTerm{白拉奇}{Pelagius}}，我极爱慕的主公，且是我极渴望的弟兄，\Person{奥古斯丁}在主内致以问候。\par

我非常感谢你体贴地通过一封信使我欢欣，并告知你的近况。愿主以那些福分回报你，使你因拥有这些福分而永远良善，并能与那位永恒者永远生活在一起。我极爱慕的主公，我极渴望的弟兄，尽管我不承认我内有任何配得上你仁爱之信中对我的那些颂词，但我仍不能不为你对如此微不足道之人所显示的善意而感恩，同时建议你更应为我祈求，求主使我成为你已经认为我是的那种人。\par

\EditorialNote{另一人笔迹}愿你得享平安与主的恩惠，并记念我们！\par

\SourceRecord{Translated by J.G. Cunningham.；Nicene and Post-Nicene Fathers, First Series；Vol. 1.；Edited by Philip Schaff.；Buffalo, NY: Christian Literature Publishing Co.,；1887.；<http://www.newadvent.org/fathers/1102146.htm>.}

% structure: p0001 source=h1 target=section reason=page-title

\section{第148封信（公元413年）}

一封指导信\OriginalTerm{训导书}{commonitorium}致圣善的弟兄\Person{福图纳提安}。\par

% structure: p0003 source=h2 target=subsection reason=relative-heading-rank; base=h2

\subsection{第一章}

1. 我写这封信是为了提醒你，当我与你在一起时曾提出的请求，请你费心拜访我们在谈话中提及的那位弟兄，以便请求他宽恕我，如果他把我最近的一封信（我并不后悔写了这封信）中的任何言论解释为对他本人严厉而不友善的攻击，断言这身体的眼睛不能看见天主，并且永远无法看见祂。我紧接着补充了我做此断言的理由，即，为防止人们相信天主本身是物质的、可见的，以为祂占据一个由尺寸和与我们之间的距离确定的位置（因为这身体的眼睛除非在这些条件下，否则什么也看不见），并防止人们将\Quote{面对面}这一表述理解为仿佛天主被局限在身体的肢体之内。因此，我并不后悔做了这一断言，作为反对我们形成如此不配且亵渎的观念的一种抗议，即认为祂并非以其整体无处不在，而是可被分割，并散布于空间的各处；因为只有这样的对象才能被我们这些眼睛所认识。\par

2. 然而，若有人虽未对天主持有此类意见，却相信祂是神体，不变易，无形体，以其全部存有遍在处处，而他认为我们身体的这种改变（当它从本性之身体变为属神之身体时）将会如此之大，以致在这样一个身体中，我们将有可能看见一个属神实体，此实体既不按空间距离或尺寸分割，也不局限在身体肢体的界限内，而是以其整体处处呈现，我希望他教导我此事，若他所发现的属实；但若他在这意见上有误，将不属于身体的东西归于身体，比夺去属于天主的东西更为可接受。即使那个意见正确，也不与我在那封信中的话相矛盾；因为我说过这身体的眼睛不能看见天主，意思是，我们这身体的眼睛只能看见那些与它们有空间间隔的物体，因为如果没有间隔，连物体本身也不能通过眼睛被我们看到。\par

3. 再者，如果我们的身体将被改变成为与现在迥然不同的某物，以致具有这样的眼睛：藉着它们，能看见一种不散布于空间之中、不囿限于范围之内——不在空间上有一部分在此、另一部分在彼，较小者在较窄处、较大者在较广处——而是以其整体属神地遍在一切的实体——这些身体就将是与现在的身体极为不同的，且不再是其本身了，它们将不仅摆脱死亡、腐朽和重负，而且将以某种方式被转化成心灵本身的特质，只要它们能以那种届时将赋予心灵、而目前却连心灵本身也未获得的方式去看。因为，如果当一个人的习性改变时，我们说他已经不是原来那个人了——如果在我们的年岁改变时，我们说身体已不是原来那样了，那么当身体经历了如此巨大的变化，不仅获得了不朽的生命，而且有能力看见那不可见者时，我们更可以说身体已不再是同一个身体了。因此，如果他们将如此看见天主，那也不是藉这身体的眼睛看见祂，因为此时就连这身体也不再是同一个身体了，它在能力和权能上已被改变到如此巨大的程度；因而此意见并不与我信中的话相矛盾。然而，如果身体只改变到这种程度：即现今是可死的，那时成为不死的；现今拖累灵魂，那时却毫无重负，对各种运动都极为敏捷，但观看的能力却未改变，仍然只能看到那些依其尺寸和距离而辨别的事物，那么它仍然完全不可能看见一个无形体且整体遍在的实体。因此，不论前一种假定还是后一种假定正确，在这两种情况下，这身体的眼睛不能看见天主这一点都依然为真；或者说，如果要看见祂，那将不再是这身体的眼睛，因为在如此巨大的改变之后，它们将是一个与此时截然不同的身体的眼睛。\par

4. 但是，如果这位弟兄能够就这个问题提出更好的见解，我乐意向他或他的导师学习。如果我这话出于讥讽，我或许也会说，我准备接受关于天主的这样一种道理，即祂拥有一个具有肢体的身体，并且在空间中可分割到不同的位置；但我没有这样说，因为我并非在讥讽，我完全确信天主在任何方面都不具备这样的性质；我写那封信，正是为了防止人们如此相信祂。在那封信中，我被自己警戒错谬的热忱所驱使，又因我没有指名道姓地批评那个人的观点而写得更为随意，以致过于激烈，不够谨慎，没有尽到我作为一位弟兄和主教对另一位的尊位所应有的尊重：对此，我不为自己辩护，反而加以责备；我谴责而不为自己开脱，恳求得到宽恕。我恳请他记念我们昔日的友谊，忘掉我近来的冒犯。让他去做他因我未做到而责怪我未做的事吧；让他以宽恕表现温和，正如我在写那封信时未能表现温和一样。因此，我借着您的善意调解，请求他，如果我曾有机会，我早就决定亲自向他请求了。我的确曾努力争取与他会面（我以我的名义写信请求这位可敬的、配得我们所有人尊敬的人），但他谢绝前来，我猜想，他怀疑有人要对他设下圈套，这种事在人与人之间是屡见不鲜的。我恳请您尽最大努力向他保证，我绝无此等诡诈，这事您亲自见他会更容易做到。请告诉他，我是怀着何等深切真实的忧伤与您谈及我伤害他感情一事。让他知道，我绝非轻视他，我在他身上何等敬畏天主，且记念我们的元首，我们同是祂身体上的弟兄。我所以认为最好不去他居住的地方，是为了避免我们成为教会范围之外那些人的笑柄，使朋友忧伤，自己蒙羞。这一切，借着您的圣德与爱德的斡旋，都能圆满解决；更进一步说，圆满的结局掌握在那位借着祂所赐的信德住在您心里的主手中，我深信我们的弟兄不会拒绝在您身上尊敬祂，因为他在自己内经验到基督的内住。\par

5. 无论如何，我不知这件事上除了向那位抱怨因我信中的严厉而受伤害的弟兄请求宽恕之外，还能做什么更好的事。我希望，他会做那位借着宗徒说话的主所命令他的事：\BibleQuote{你们要彼此宽恕，如果有人对某人有什么过犯，就如天主在基督内宽恕了你们一样}\BibleRef{哥 3:13}；\BibleQuote{所以你们该效法天主，如同蒙宠爱的儿女一样；又要在爱中行事，就如基督爱了我们}\BibleRef{弗 5:1-2}。在这爱中行事，让我们同心合意，而且，如果可能，以比以往更大的勤奋，探讨我们复活后将要拥有的那个\OriginalTerm{属神的身体}{spiritual body}的性体。\BibleQuote{如果我们有什么不同的见解，天主也必将这事启示给我们}\BibleRef{斐 3:15}，只要我们住在祂内。那住在爱内的，就住在天主内，因为\BibleQuote{天主是爱}\BibleRef{若一 4:8}——无论是指那不可言喻的本质中的爱的泉源，还是指那借着祂的圣神白白赐给我们的爱的泉源。那么，如果能够证明爱在任何时候都能成为我们肉眼可见的事物，那么我们就可以承认天主也可能如此；但如果爱永远不能成为可见的，那么就更不用说祂本身是爱的泉源，或无论用什么更卓越更适当的象征性名称来论及如此伟大的一位，能够成为可见的了。\par

% structure: p0009 source=h2 target=subsection reason=relative-heading-rank; base=h2

\subsection{第二章}

6. 有些具有伟大恩赐、深通圣经的人，当他们遇到机会时，通过著作大大造福了教会，促进了信徒的教导；他们曾说，那不可见的天主是以不可见的方式被看见的，也就是通过我们内那同样不可见的本性，即一颗纯洁的心灵或内心。圣\Person{盎博罗削}在论及基督是圣言时说：\Quote{耶稣不是被肉体的眼睛所看见，而是被属神的眼目所看见；}稍后他又补充说：\Quote{犹太人看不见祂，因为他们愚昧的心被蒙蔽了，}以此说明基督如何被看见。还有，当他在论及圣神时，引用了主的话说：\BibleQuote{我要祈求圣父，祂将赐给你们另一位护慰者，使祂与你们永远同在，即真理之神；世界不能领受祂，因为见不到祂，也不认识祂；}\BibleRef{若 14:16-17}接着又说：\Quote{所以，祂在身体内显现自己，是有充分理由的，因为在祂天主性的本体中，祂是不可见的。我们看见了圣神，不过是借着有形的形象；让我们也看见圣父；可因为我们不能看见祂，就让我们听从祂吧。}过了一会儿，他又说：\Quote{那么，让我们听从圣父，因为圣父是不可见的；但圣子就其天主性而言也是不可见的，因为从来没有人看见过天主；}\BibleRef{若 1:18}\Quote{既然圣子是天主，在祂是天主的方面，当然不能被看见。}\par

7. 圣\Person{热罗尼莫}也说：\Quote{人的眼睛不能看见天主在其自身性体中的本来面目；这点不仅对人如此；就是天使，无论是座天使、掌权者、统权者，或任何可以称呼的名号，都不能看见天主，因为没有一个受造物能看见其造物主。}通过这几句话，这位学识渊博的人充分表明了他对这个问题的看法，不仅涉及今生，也涉及来世。因为无论我们肉身的眼睛会变得多么美好，也只不过等同于天使的眼睛。然而，在这里，热罗尼莫确认了造物主的性体即使对天使，对天堂里的一切受造物，无一例外，都是不可见的。不过，如果有人对此提出疑问，怀疑我们将来是否会高于天使，主自己的心意是清楚的，祂在论及那些复活进入天国的人时所用的话表明了这一点：\BibleQuote{他们将相似天使}\BibleRef{路 20:36}。\par

因此，同一圣\Person{热罗尼莫}在另一处如此表达：\Quote{因此，人不能看见\OriginalTerm{天主的面容}{face of God}，但教会中最小的天使却常看见天主的面容。\BibleQuote{现在我们是在镜子中观看，模糊不清，如同猜谜，到那时就要面对面的观看了}；当我们从人提升到天使的等级时，我们就能同宗徒一起说：\BibleQuote{我们众人以揭开的面容，在镜子中观看主的光荣，就变成了与主同样的肖像，光荣上加光荣，如同由主圣神所变化的}；尽管就\OriginalTerm{性体的本质属性}{essential properties of His nature}而言，没有任何受造物能看见天主的面容，但在这些情况下，祂是由心灵看见的，因为祂被认为是不可见的。}\par

8. 在这位天主的人的话语中，有许多值得我们深思之处：首先，根据主极清晰明确的声明，他也认为，当我们将提升到天使的等级，即与天使同等时——无疑是在死人复活时——我们才会看见\OriginalTerm{天主的面容}{face of God}。其次，他充分运用宗徒的证言解释，当说到我们将\BibleQuote{面对面的观看}时，这里所说的面容不应理解为外在的人，而是内心的人；因为宗徒当时所论的是\OriginalTerm{心灵的面容}{face of the heart}，当他说出那些在此被圣\Person{热罗尼莫}引用的话时：\Quote{\BibleQuote{我们众人以揭开的面容，在镜子中观看主的光荣，就变成了与主同样的肖像}}。若有人对此怀疑，请他再细读这段经文，并留意宗徒所论的是何事，即那帕子，在诵读\WorkTitle{旧约}时仍蒙在众人心上的帕子，直到他们归向基督，帕子才被揭去。因为他在那里说：\Quote{\BibleQuote{我们众人也以揭开的面容，在镜子中观看主的光荣}}——而犹太人的面容却未被揭开，关于他们，宗徒说，\Quote{\BibleQuote{帕子仍蒙在他们的心上}}——为要显示，当帕子揭去后，在我们内揭开的面容正是那心灵的面容。\par

最后，为避免有人因观察这些事不够仔细而误解其意，竟以为天主如今或将来，当人达到与天使同等时，要为天使或人所目睹，他便极其清晰地表明了自己的观点，肯定地说：\Quote{就\OriginalTerm{性体的本质属性}{essential properties of His nature}而言，没有任何受造物能看见\OriginalTerm{天主的面容}{face of God}}，并且\Quote{在这些情况下，祂是由心灵看见的，因为祂被认为是不可见的}。从这些陈述中，他充分说明，当天主通过肉身的眼目，以仿佛有肉身的形式显现于人时，祂并非根据其性体的本质属性而被看见，因为就那属性而言，祂是由心灵看见的，因为祂被认为是不可见的——即对天上的存在，也就是\Person{热罗尼莫}上文中论及的天使、有权能者、宰制者，祂也是不可见的。那么，对地上的存在，祂更是何等不可见啊！\par

9. 因此，在另一处，\Person{热罗尼莫}以更加清晰的言辞说，的确，不但圣父的天主性，而且圣子的天主性及圣神的天主性，构成了天主圣三内同一的性体，都不能被肉身的眼目看见，而只能被\OriginalTerm{心灵的眼目}{eyes of the mind}看见；关于这心灵的眼目，救主自己曾说：\Quote{\BibleQuote{心里洁净的人是有福的，因为他们要看见天主}}。\par

还有什么比这陈述更清晰呢？因为如果他仅仅说，圣父的天主性、或圣子的天主性、或圣神的天主性，不可能被肉身的眼目看见，而没有加上\Quote{而只能被心灵的眼目看见}这句话，也许有人会说，当肉身成为属神的之后，便不能再称为\Quote{肉身}了；但通过加上\Quote{而只能被心灵的眼目看见}这句话，他就把看见天主的可能性，从任何形体的范畴中排除了。然而，为避免任何人以为，他说的只指现世的状态，请留意，他还附加了主的一句证言，特意为界定他所论的心灵的眼目；在这句证言中，所许下的不是现今的，而是未来的看见：\Quote{\BibleQuote{心里洁净的人是有福的，因为他们要看见天主}}。\par

10. 同样，至福的圣人\Person{亚大纳削}，\Place{亚历山大里亚}的主教，在驳斥亚略派——他们肯定唯独圣父是不可见的，却认为圣子和圣神是可见的——时，断言天主圣三的所有位格同等不可见，他引证\WorkTitle{圣经}的许多章节来证明，并以他惯有的争辩之细致加以论证，竭力说服他的对手：天主从未被人看见，除非是借取了受造物的形象；而在其本质的天主性中，天主是不可见的，也就是说，圣父、圣子、圣神都是不可见的，除非就天主圣三的位格能由心灵和神魂所认识而言。东方的一位圣主教\Person{额我略}，也非常明确地说，天主本性不可见，但在祂被古圣祖看见（如\Person{梅瑟}，天主曾与他面对面谈话）的那些场合，祂使自己能被看见，是借取了某种物质和可见的形象。我们的圣人\Person{安博罗削}也说了同样的话：\Quote{\InnerQuote{父、子及圣神，在显现时，是选取了随意选取的形象，而非由天主性所规定的形象}}；这样，便澄清了\Quote{\BibleQuote{从来没有人见过天主}}这句话的真理，这是主基督自己的话，也澄清了另一句话：\Quote{\BibleQuote{没有人看见过，也不能看见}}的真理，这是宗徒，更恰当地说，是基督借着宗徒所说的话；同时，也维护了\WorkTitle{圣经}中那些提及天主被人看见的章节的一致性，因为天主在天主性的本质性体中是不可见的，但祂能随意取一个受造的形象，使自己成为可见的，只要祂愿意。\par

% structure: p0018 source=h2 target=subsection reason=relative-heading-rank; base=h2

\subsection{第三章}

11. 再者，如果不可见性是天主本性的一个属性，就如不朽性一样，那么这本性在来世决不会经受这样的改变，以至于不再是不可见的而变为可见的；因为它决不可能不再是不可朽的而变为可朽坏，因为在这两个属性上它同样是不变的。宗徒的确宣告了这天主本性的卓越，将这两项并列，说：\BibleQuote{那万世的君王、不可见的、不朽的、唯一的天主，愿尊荣和荣耀归于祂，直到永远。}因此，我不敢做这样的区分说：不朽的，固然是直到永远，而不可见的——却不是直到永远，只在今世。同时，既然我们接下来要引用的证据不可能是假的，\par

\BibleQuote{心里洁净的人是有福的，因为他们要看见天主，}以及\BibleQuote{我们晓得，当祂显现的时候，我们必要像祂，因为我们要看见祂实在怎样。}——我们不能否认天主的子女将要看见天主；但他们看见天主，是如同不可见的事物被看见一样，正像那曾以血肉之躯显现、为人所见的，向人预许祂要显示自己时，在门徒面前说话、为他们的肉眼所见时所说的：\BibleQuote{我要爱他，并要向他显现我自己。}除了用心灵的眼目，还有什么别的方式能看见不可见的事物呢？对此，作为我们看见天主的工具，我先前曾简短引述了热罗尼莫的意见。\par

12. 因此，我以前引述过的米兰主教也说过，即使在复活的时候，除了那些心里洁净的人，谁也不能轻易地看见天主，所以经上写道：\BibleQuote{心里洁净的人是有福的，因为他们要看见天主。}他说道：\Quote{祂已经列举了多少有福的人，却还没有应许他们能看见天主；}然后他推论说：\Quote{所以，如果心里洁净的人要看见天主，那么显然别的人就不会看见天主；}并且为免我们以为他是指主所说\InnerQuote{神贫的人是有福的，温良的人是有福的}中那些别的人，他立刻补充说：\Quote{因为那些不相配的人，将不会看见天主，}他的意思是，那些不相配的人就是那些虽然会复活，却不能看见天主的人，因为他们的复活是被定罪的，因为他们拒绝借着那\InnerQuote{以爱德行事的}真信德来洁净自己的心。为此他接着说：\Quote{凡不愿看见天主的，就不能看见祂。}然后，由于他想到在某种意义上连所有的恶人也渴望看见天主，他随即解释说，他之所以说\Quote{凡不愿看见天主的}，是因为恶人不愿洁净自己的心（唯有借心才能看见天主），这表明他们并不渴望看见天主，接着说了下面的话：\Quote{天主不是在空间中被人看见，而是在洁净的心里；也不是用肉体的眼睛去寻觅；也不是用我们的视觉能力去界定祂的形像；也不能用手触摸去把握；祂的声音不落于耳中；祂的行踪也不会被感官所觉察。}通过这些话语，圣安博罗修愿意教导人，如果他们希望看见天主，应该做的准备就是：借着以爱德行事的信德，靠着圣神的恩赐来洁净自己的心，我们已经从祂领受了那保证，借此我们被教导去渴望那亲眼看见天主的神视。\par

% structure: p0022 source=h2 target=subsection reason=relative-heading-rank; base=h2

\subsection{第四章}

13. 至于圣经时常提及的天主的肢体，免有人按照身体的形状外貌设想我们与天主相似，同一部圣经还说天主也有翅膀，这我们当然没有。所以，当我们听到天主的\Quote{翅膀}时，便理解是天主的护佑；同样，天主的\Quote{手}应理解为祂的作为；\Quote{脚}是指祂的临在；\Quote{眼}是指祂看见并知晓万有的能力；\Quote{面}是指祂借以将自己启示给我们认识的那个标记。并且我相信，圣经中任何其他类似的表述，都应当从灵性的意义上去理解。这个主张并非我一人独有，也不是我第一个提出的。所有以这种灵性解释理解圣经这类语言的人，都持此意见，并抵抗那些被称为拟人论者的人。为了不多占时间而大量引述这些人的著作，我在此只引出虔诚的热罗尼莫的一段话，好叫我们的弟兄知道，如果有什么感动他持守相反的意见，他不但要与我，也要与在我之先的前辈们进行辩论。\par

14. 在圣经那位最博学的学者对那首圣咏所作的注释中，这首圣咏里有这样的话：\BibleQuote{你们愚昧的百姓，应当知晓！愚顽的人，你们几时才能明白？装备耳朵的，难道自己听不见？形成眼睛的，难道自己看不见？}他说道：\Quote{这段经文有力地反对那些拟人论者，他们说天主拥有和我们一样的肢体。例如，他们说天主有眼，因为\InnerQuote{主的眼目遍察一切}；同样，他们按字面接受\InnerQuote{主的手成就一切}以及\InnerQuote{亚当听见主在乐园中行走的脚步声}这类话，从而将人的软弱归于天主的尊威。但我断言，天主全为眼，全为手，全为脚：全为眼，因为祂看见一切；全为手，因为祂运作一切；全为脚，因为祂处处临在。所以，请看圣咏作者怎么说：\BibleQuote{装备耳朵的，难道自己听不见？形成眼睛的，难道自己看不见？}他没有说：装备耳朵的，难道自己没有耳朵？形成眼睛的，难道自己没有眼睛？他怎么说呢？\BibleQuote{装备耳朵的，难道自己听不见？形成眼睛的，难道自己看不见？}圣咏作者把看和听的能力归于天主，却没有将肢体归于祂。}\par

15. 我自觉有责任从拉丁和希腊作家们的著作中引用所有这些段落，他们先于我们的时代在天主教会内，曾为\OriginalTerm{神圣神谕}{divine oracles}撰写注释，以便我们的弟兄，如果他持有任何与他们不同的意见，可以知道，他应当放下一切争论的苦毒，并完全保持或恢复弟兄之爱的温柔，以勤勉和平静的心去探究，或是他必须向他人学习什么，或是他人必须向他学习什么。因为任何人的推理，即使他们是天主教徒，且享有崇高声誉，也不可按照我们对待\OriginalTerm{圣经正典}{canonical Scriptures}的方式对待。我们完全可以自由地，在不损害这些人士所应受的尊敬的情况下，摒弃和拒绝他们著作中的任何内容，倘若我们发现他们所持的观点，与其他人或我们自己，依靠天主的助佑，所发现的真理有所不同。我如此对待他人的著作，并希望我明智的读者也如此对待我的著作。总之，我靠着主的助佑，最坚定地相信，并且在祂恩赐我的范围内，我理解在一切我所引用的、出自神圣而博学的\Person{圣安博罗削}、\Person{圣热罗尼莫}、\Person{圣亚大纳削}和\Person{圣额我略}著作中的陈述，以及在其他我所读过的、但为简洁起见而未引用的作家著作中类似的陈述中所教导的，即：天主并非一个形体，祂没有人体肢体的构造，祂不因空间而分割，祂是不可见且永恒不变的，祂并非以自己的本质和实体显现，而是以祂所愿意的可见形式，向那些祂显现的人显现，正如\WorkTitle{圣经}记载祂曾被圣人们以肉眼看见的情况。\par

% structure: p0026 source=h2 target=subsection reason=relative-heading-rank; base=h2

\subsection{第五章}

16. 关于我们在复活时将拥有的\OriginalTerm{属神的身体}{spiritual body}，它将经历何等巨大的改善——它是变成\OriginalTerm{纯粹的精神体}{pure spirit}，使整个人那时成为一个精神体，还是（我更倾向于这种观点，但尚不能确信地坚持）成为一个\OriginalTerm{属神的身体}{spiritual body}，以致因为其运动具有某种不可言喻的轻灵而被称为属神的，但同时仍保留其物质实体，这实体不能靠自身生活和感觉，而只能通过运用它的精神体（因为我们目前的状态也是如此，虽然身体被称为\OriginalTerm{有生命的}{animal}，但赋予生命之源的特性却不同于身体的特性）；又，那时不死而不朽的身体的属性是否将保持不变，它是否将在某种程度上帮助精神体观看可见的，即物质的事物，正如我们目前若不通过肉眼就无法看见任何这类事物，或者我们的精神体那时将能够，即使处于更高状态，也不需要身体的辅助而认识物质事物（因为天主自己就不是通过身体感官认识这些事物的）——对于这些和许多其他可能在讨论这主题时困扰我们的问题，我承认，我尚未在任何地方读到过任何我认为足够确定、值得人们学习或教导的东西。\par

17. 因此，如果我们的弟兄愿意耐心忍受我方面的任何犹豫，让我们在此期间，因为所记载的：\BibleQuote{我们必要看见祂实在怎样}，就借着天主自己的助佑，尽我们所能，为那异象预备我们净化了的心。让我们同时更平静而仔细地探究关于\OriginalTerm{属神的身体}{spiritual body}的事，因为或许天主，如果祂知道这为我们是获益的，就会屈尊，按照祂书写的话语，向我们显示关于这主题的某种明确而清晰的观点。因为如果更仔细的研究将发现，身体的变化如此之大，以致能看见不可见的事物，那么，我想，赋予身体的这种能力，不会剥夺心灵观看的能力，从而使外在的人看见天主，而内在的人却不能——仿佛与\WorkTitle{圣经}的明确话语相矛盾，即\BibleQuote{天主成为万物中的万有}，天主只是在那人旁边——在他之外，而不是在他内，在他的内心深处；或者仿佛那位处处整个地临在、不受空间限制的天主，居然在人内，只能被外在的人从外面看见，而不能被内在的人从里面看见。如果这些观点显而易见是荒谬的——因为，相反，圣人们将充满天主；他们不会在内心空乏时，被祂从外部包围；他们也不会因为内心盲目，看不见他们所充满的那位，而只有为外面的事物所设的眼睛，才能看见那位他们将被包围的——我认为，如果这些事是荒谬的，我们就只能暂时安然确信内在的人看见天主。但是，如果身体通过某种奇妙的变化而被赋予这种能力，那么新增的是另一种新的官能；而原先拥有的官能并不会被取走。\par

18. 那么，最好肯定我们所无疑的事——天主将由那\OriginalTerm{内在的人}{inward man}看见，唯有这内在的人，在我们当前状态下，能看见那爱，宗徒为赞扬这爱曾说：\BibleQuote{天主是爱；}唯有这内在的人能看见\BibleQuote{和平与圣德，没有这两样，谁也见不到主。}因为如今没有肉体的眼睛能看见爱、和平、圣德，以及这类事物；然而这一切，就它们能被看见的程度而言，都是通过心灵之眼看见的，而且心灵之眼越纯净，就看得越清晰。因此我们可以毫不犹豫地相信，我们将看见天主，无论我们关于将来身体的本性的探讨是成功还是失败——尽管同时我们确信，身体必将复活，成为不死不灭的，因为关于这一点，我们有圣经最清楚、最有力的证据：然而，如果我们的弟兄现在断言，他已就那\OriginalTerm{属神的身体}{spiritual body}获得了确切知识——对此我不过是在探究——那么若我拒绝心平气和地聆听他的教导，便生我的气，他确有正当的理由，只要他也心平气和地聆听我的疑问。然而现在，我恳求你，看在基督的份上，为我获致他对我的宽恕，因为我信中的严厉，据我所知，并非无故地冒犯了他；愿你靠天主的助佑，以你的回信来安慰我的心神。\par

\SourceRecord{Translated by J.G. Cunningham.；Nicene and Post-Nicene Fathers, First Series；Vol. 1.；Edited by Philip Schaff.；Buffalo, NY: Christian Literature Publishing Co.,；1887.；<http://www.newadvent.org/fathers/1102148.htm>.}

% structure: p0001 source=h1 target=section reason=page-title

\section{第150封信（公元413年）}

\Emph{致\Person{普罗芭}与\Person{尤莉安娜}，最值得尊敬的夫人，享誉四方且最为卓越的女儿们，奥古斯丁在主内向你们致候。}\par

你们使我们心中充满了极为愉悦的欣喜：这欣喜因着我们对你们的爱而令人快慰，又因消息迅速传来而令人欣然。当你们家女儿的献身童贞生活之祝圣，正借着最繁忙的声誉在你们所知的各处——即世界各地——传扬之际，你们却以一封你们的信函，以更确实可靠的消息，超过了名声的飞驰，使我们在还来不及质疑任何有关如此蒙福且非凡事件的传闻之真实性之前，就已确知喜讯。谁能用言语述说，或以相当的赞词阐明，你们的家族由于献给基督一些献身事奉祂的妇女，而获得的荣耀与益处，比起献给世界一些蒙召获得执政官荣誉的男子，要超越得多么无可比拟呢？因为，如在这流转的尘世岁月中留下荣耀名字的印记是伟大而高贵的事，那么借着心灵与身体未受玷污的贞洁而超越其上，该是何等更伟大更高贵的事啊！因此，这位女郎，门第显赫，而因虔诚更为显赫，让她在与她神圣的主缔结婚约而获致天上的卓越荣耀中，寻得更大的喜乐，胜过她若与人间配偶缔结婚约，成为显赫子孙之母所可能有的喜乐。这位阿尼基乌斯家族的女儿，表现了更高尚的胸怀：她选择了借着不婚而为那高贵家族带来祝福，而不是增添其后裔；宁愿在身体的纯洁中，已相似天使，而不愿借着生育果实增添凡人的数目。因为这真福的境界是更丰盛、更富饶的：不是拥有怀孕的子宫，而是发展灵魂的崇高能力；不是拥有满溢乳汁的乳房，而是拥有纯洁如雪的心灵；不是在分娩的阵痛中孕育属世之物，而是以恒切的祈祷孕育属天之物。我的女儿们，你们最堪当地位的尊荣，愿你们在她身上享有你们自己所缺少的；愿她坚持到底，持守那永无尽头的婚姻结合。愿许多婢女们效法她们主母的表样；愿卑微者效法这位出身高贵的夫人，愿在此无常世上享有显赫者，向往那谦逊赋予她的更可贵的显赫。愿那些羡慕阿尼基乌斯家族荣耀的贞女们，转而热衷于效法其虔诚；因为前者是她们无论怎样热切渴求也无法企及的，而后者若她们全心寻求，立刻就能获致。愿至高者的右手护佑你们，赐给你们平安与更大的幸福，最值得尊敬的夫人，最卓越的女儿们！在主的爱内，并以一切应有的敬意，我们问候你们圣善的子女们，尤其是那位在圣德上超群出众的。我们非常欣喜地收到了为纪念她献身修道而寄来的礼物。\par

\SourceRecord{Translated by J.G. Cunningham.；Nicene and Post-Nicene Fathers, First Series；Vol. 1.；Edited by Philip Schaff.；Buffalo, NY: Christian Literature Publishing Co.,；1887.；<http://www.newadvent.org/fathers/1102150.htm>.}

% structure: p0001 source=h1 target=section reason=page-title

\section{第151封信（公元413或414年）}

\Emph{致\Person{塞西利安}，我公义闻名的大人，并因他的地位而最堪当受我敬重的儿子，奥古斯丁在主内致意。}\par

1. 您信中向我提出的劝诫，因其所显示的爱而令我喜悦。因此，若我试图就我的沉默为自己开脱，我必须证明的是，您没有正当理由对我不悦。但既然您屈尊因我的沉默——我原以为在您许多挂虑中这无关紧要——而不悦，这比任何事都更让我欢喜，那么我若如此试图自辩，便无异于控告自己。因为如果您因我不给您写信而对我表示不悦是错误的，那么必定是由于您对我评价不高，以致您对我说话或缄默全不在意。不，您因我的沉默而感到痛苦，由此产生的不悦并非真的不悦。因此，我对于自己的沉默，与其哀痛，不如欢喜，因您渴望我的音讯。因为，能被一位老朋友，一位（虽然您也许不说，但我们有责任承认）具有如此卓著的价值与伟大、身在异邦、肩负公共责任而身居高位的人记念，这对我是一种荣誉，而非烦扰。所以，请原谅我表达感谢，因您并未将我视为一个其沉默也不值得您恼怒的人。因为，藉着那比您的高位更能显示您的仁慈，我现在相信，在您众多且重大的公务——非私事，而是涉及全民利益——之中，我的一封信会被您看作不添负担，而受欢迎。\par

2. 当我收到因卓越功德而可敬的圣教宗\Person{依诺森}的信函，并由弟兄们交付给我的时候——那信函明显是由您的卓越转达给我的——我形成的意见是，没有您的书信相伴，是因为您忙于更重大的事务，不愿被通信的麻烦所困。因为那时我似乎并非不合理地期望，当您屈尊将一位圣者的著作寄给我时，我会收到一些您自己的文字。因此，我本已决定不写信打扰您，除非必要为向您推荐某人，而我又不能拒绝为之效劳代祷——这是我们惯常对所有的人施予的恩惠，这种习惯虽带来许多麻烦，却不应全然责难。我照此做了，向您的善意推荐了我的一位朋友，我现在收到了他表示感谢的信，我也在此加上我的感谢，因您的帮助。\par

3. 然而，如果我曾对您形成过任何不利的印象，特别是关于那件事——虽未明言，却仿佛有一股微妙的气味弥漫在您整封信中——我决不会给您写那样的短笺，为自己或他人请求任何恩惠。在那种情形下，我要么保持沉默，等待有机会亲自见您；要么，如果我认为有责任就此事写信，我会将它放在信的首要位置，并以一种几乎使您无法表示不悦的方式来处理。因为，当我们在与您共享的忧虑下，向他提出劝诫——热烈恳求他，却是徒劳——求他不要以如此巨大的悲伤刺透我们的心，并以如此重大的罪使自己的良心受致命伤时，他却犯下了那不虔、野蛮、背信的罪行；为此，我立即秘密地离开\Place{迦太基}，原因是：那些畏惧他的刀剑而在教堂内寻求庇护的众多有权势的人，若想象我的同在对他们有所帮助，就会以他们痛哭和呻吟阻止我，使我为保全他们身体的缘故，不得不向一个我无法为了其灵魂的福祉而以他的罪行当受的严厉加以斥责的人恳求恩惠。至于他们的人身安全，我知道教堂的围墙足以保护他们。但对我自己[如果我留在那里为他们向他求情]，那只会陷入痛苦尴尬的处境，因为他不会容忍我以我所应当的方式对待他，而我也会被迫采取一种于我不宜的行动。同时，我确实为那位可敬的\OriginalTerm{同侪主教}{co-bishop}，如此重要教会的主管，遭遇的不幸感到悲伤；人们竟期望他，即使在此人犯下如此可耻的背叛罪行之后，仍以顺服恭敬来对待他，为使他人的性命得以保全。我承认我离开的原因：我将无法以必要的坚强面对如此巨大的灾难。\par

4. 当初使我离去的那种考虑，同样会使我对您保持沉默，如果我曾相信您利用对他的影响力来报复如此邪恶的伤害的话。只有那些不知道您如何、多少次、以何种措辞向我们表达您想法的人，才会这样相信您。当时我们怀着焦虑的心情，竭力确保，因为他与您如此亲密，您又经常拜访他，并常常与他单独交谈，他更该谨慎地维护您的名誉，免得人们以为您未尽力阻止他对那些据说是您仇敌的人处以那种死刑。事实上，无论是我，还是那些听过您谈话、并从您的言语和一举一动中看到您对遭处决者内心善意的我的弟兄们，都不这样相信您。但是，我恳求您，宽恕那些这样相信的人；因为他们是人，人心有那样的隐蔽之处和那么深的深渊，以至于一切多疑之人都该受责备，他们自己却以为他们的谨慎甚至值得称赞。归罪于您的行为是有原因的：我们知道，您曾受到他突然下令逮捕的那些人中一位极其严重的伤害。据说，他的兄弟——他特别借此人迫害教会——曾答复您，言辞中似乎含有某种严厉的控诉。这两人都被认为受到您的猜疑。当他们被传唤后离开时，您仍留在原地，据说与\Person{马里努斯}进行了一次比平时更私密的交谈，随后他们突然被下令扣留。人们纷纷议论您与他之间的友谊并非新近，而是由来已久。您与他交往的密切，以及私下交谈的频繁，证实了这种传闻。那时他权势很大。对任何人进行诬告何等容易，是众所周知的。要找到某个人，在得到自身安全的许诺后，说出他命令说出的任何话，并非难事。当时的一切，使任何人都能使一个人不经下令处决者任何审讯就被处死，只要有一个证人提出一项看起来既可憎又可信的控诉。\par

5. 与此同时，有传言说教会的权力可以解救他们，我们便被一些虚假的承诺所愚弄，以至于不仅得到\Person{马里努斯}的同意，而且似乎是在他的迫切要求下，一位主教被派往朝廷为他们说情，同时还让主教们听到这样一个承诺：在朝廷听取为囚犯们提出的某种申诉之前，不会对他们的案件进行任何审讯。最后，就在他们被处死的前一天，阁下您来到我们这里；您给了我们前所未有的鼓励，说他在您（前往\Place{罗马}）启程之前，可能会将他们的性命作为一份人情赐给您，因为您曾郑重而谨慎地对他说，他这样屈尊，让您如此经常地与他进行熟稔而私密的交谈，带给您的将不是荣誉而是耻辱，而且当这些人的死亡已经成为你们之间交谈和商议的主题之后，这会让每个人都毫不怀疑这些会谈将导致什么结果。当您告诉我们您已向他说了这些话时，您一边说，一边伸手指向信友举行圣事之处，就在我们惊讶地聆听时，您以极其庄重的誓言证实您确实说过这些话，以致不仅当时，甚至在囚犯们遭遇这可怕而出乎意料的死亡之后，如今我回忆起您当时整个举止，都觉得如果我对您有任何不良的想法，那将是一种加倍的侮辱。您还说到，他听了您这番话深受感动，打算在您启程之前，将这几个人的性命作为友谊的信物赠予您。\par

6. 因此，我郑重向阁下保证，当次日——即那桩恶名昭彰的阴谋实施之日——消息出人意料地带给我们，得知他们已被提出监狱，站在他面前受审，尽管我们有些惊恐，然而回想起您前一天对我们所说的话，并想起次日正值圣\Person{西彼廉}殉道的周年纪念，我便推想，他甚至是特意选择了这样一个日子：在这一天，他不仅能够允准您的请求，而且还能因突然将喜乐带给整个基督的教会，而渴求效法如此伟大的一位殉道者的德行，证明自己真正更伟大之处，在于运用仁慈保全性命，而非拥有置人于死的权力。我正这样想着，看！一个报信者闯入我们中间，我们还没来得及询问他审判进行得如何，就已得知他们被斩首了。原来已经特意安排，就在紧邻的地方作为行刑地点，那个地方并不是用来处决罪犯的，而是供市民休憩的场所；就在几天前，他曾命令在那里处决过一些人，其用意（人们有理由相信）是为避免首次将这些人的处决用于该地而招致憎恨，他原希望能将这些人从代他们求情的教会手中偷偷夺走，因此不仅下令立即处决，而且下令在最近便的地点执行。因此，他充分表明，他并不惧怕使那一位他畏惧其干预的慈母——即圣教会——遭受残酷的伤痛；因为我们知道，他本人也算在她怀中受洗的忠信子女之列。因此，在这样一个重大阴谋实施之后——其中煞费苦心地与我们周旋，以致我们，甚至也因您，虽然是无意地，几乎完全免于忧虑，并在前一天几乎确信他们的安全——谁能不按照常人的判断方式看此事，而认为这绝无疑问：通过您，我们得到了话语，而他们被夺去了生命呢？那么，正如我说过的，请宽恕那些相信这些攻击您的流言的人，尽管我们自己并不相信，卓绝的人啊。\par

7. 然而，无论在心中还是在行为上，即便我在许多事上有缺失，也绝不敢向您为任何人求情，或向您为任何人请求恩惠，假若我相信您应对这骇人的冤屈、这邪恶的残忍负有责任的话。但我坦白向您承认，如果您在那事件之后，依然与他保持着先前那样亲密的友谊，那么请您务必谅解我之感到悲伤；因为这样一来，您便会迫使我们相信许多我们宁愿不相信的事情。不过，恰当的是，正如我不相信您犯有某些人归咎于您的其他罪行，我也不应相信此事。您的这位朋友，在骤然得势之时出人意料地得逞，对您的名誉造成的伤害，不下于对这些人的生命造成的伤害。我写这些话的目的，并非要在您心中燃起怨恨；那样做便有违我的情感与信仰告白。但我劝您更忠实向他施行爱德。因为那种以善待恶人、使其痛改前非之人，乃是懂得如何对他们发怒，却仍为他们的福祉着想的人；因为坏伙伴因顺从而损害人的幸福，同样，好朋友因反对恶行而予人帮助。他以骄横滥用权力夺去他人性命的同一件凶器，在他自己的灵魂上造成了更深更重的创伤；如果他不以悔改补救，善用天主的忍耐，那么当生命终结时，他将不得不发现并感受到这一点。再者，天主眷顾常常准许善人在今世的生命被恶人剥夺，为的是避免人们相信遭受这样的事便是一场灾祸。因为对于注定终有一死的人来说，肉身的死亡能造成什么损害呢？畏惧死亡的人，其殚精竭虑所成就的，除了将死亡的时刻短暂推迟之外，又是什么呢？凡有死之人所遭遇的一切灾祸，并非来自死亡，而是来自生命；倘若他们在临终时，灵魂得基督恩宠的支撑，死亡对他们而言，便不是善生终结的黑暗之夜，而是更美好生命起始的黎明。\par

8. 两兄弟中兄长的生活和言谈，看来确实更随从世俗而非基督，尽管他婚后也大幅纠正了早年不信宗教时的过失。然而，我们仁慈的天主命定他成为兄弟在死亡中的同伴，这未尝不是出于仁慈。至于那位弟弟，他度着虔诚的生活，无论在内心还是行为上都是一位卓越的基督徒。当他受命服务教会时，关于他将证明自己如此的佳誉，在他抵达非洲之前就已传来，且在他抵达后仍随之不绝。论他的品行，何等纯真！论他的友谊，何等恒常！论他钻研基督徒真理，何等热忱！论他的宗教信仰，何等诚恳！论他的家庭生活，何等纯洁！论他的公务职责，何等正直！他对仇敌表现出何等的忍耐，对朋友何等和蔼，对虔敬者何等谦卑，对众人何等爱德！他施予恩惠时何等爽快，请求恩惠时何等羞涩！他对人的善行何等真诚地满意，对人的过失何等哀伤！他身上闪耀着何等无玷的荣誉、高贵的风度与严谨的虔诚！在提供帮助时，他何等怜悯！在宽恕伤害时，他何等慷慨！在祈祷时，他何等信赖！当他熟悉任何事物时，总是以何等的谦虚传达有用的知识！当他意识到自己无知时，又以何等的勤奋努力探究以弥补不足！他对现世事物何等轻看！他对永恒福分的希望与渴望何等炽烈！他本会放下一切世俗事务，佩上基督徒战斗的徽章，若非他已进入婚姻状态而受阻；因为他在陷入这些束缚之前尚未开始渴望更美善之事，而一旦已经进入，尽管那些事较为低劣，他若将其撕毁便是不合法的了。\par

9. 有一天，当他们一同被关在监牢里时，他的兄弟对他说：\Quote{如果我因自己的罪恶而遭受这些作为公义的惩罚，那么是什么恶行使你遭遇同样的命运呢？因为我们知道你的生活是最严格热心的基督徒生活。}他回答说：\Quote{即使你对我生活的见证属实，你以为天主所赐予我的，不是一份微小的恩惠吗？祂命定我在这些苦难中受罚——即使至死——而非留待那将来的审判中面对我，这岂不是大恩？}这些话也许会让有些人以为他意识到自己有一些隐密的恶行。因此，我将提及上主天主乐意命定我从他口中听到，并确切知道、使我大得安慰的事。正因我对此事感到挂虑，因为人性容易陷入这样的邪恶，在他被囚禁之后，我单独与他一起时，便问他是否有任何罪，使他应藉着某种更严厉且特别的补赎来寻求与天主的和好。他带着特有的谦虚，对于我毫无根据的猜疑仅仅提及便脸红了，却极其热忱地感谢我的警告，并带着庄重而谦虚的微笑，用双手抓住我的右手，说：\Quote{我凭藉这手所分施给我的圣事起誓，无论在婚前还是婚后，我从未犯过不道德的自我放纵之罪。}\par

10. 那么，死亡给他带来了什么恶呢？不，相反，这岂不是为他带来了最大可能的善吗？因为，拥有这些恩赐，他离世归向基督，唯独在祂内才真正拥有它们。如果我相信你因我称赞他而见怪，我就不会向你提及这些事。然而，既然我不信你会如此，我同样不信他的被杀是出于你的意愿或希望，更非应你的请求。因此，以与你在这事上的无辜相称的诚恳，你无疑会与我们一同认为，那杀害他的人，对自己的灵魂所造成的伤害，远比对他受苦的身体更为残酷；因为他不仅不顾我们，不顾自己所许的诺言，不顾你那么多恳求和警告，最后，还不顾基督的教会（并在教会内不顾基督自己），竟卑劣地谋划并杀害了这人。那个有权势者的高位，与那囚徒的境遇，岂能相提并论？他虽为囚犯，但当掌权者因愤怒而发狂时，受苦者在监牢中却充满喜乐。在这个世界的一切地牢里——不，甚至在地狱中——都没有什么比得上一个恶棍因良心懊悔而陷入的黑暗可怖命运。即使对你，他又造成了什么伤害？他并没有摧毁你的清白，尽管严重损害了你的名声；然而，这名声在那些比我们更了解你的人的判断中，以及在我们这亲眼目睹你为阻止如此骇人的罪行而表露的焦虑（与我们一样）的人看来，仍是不受损伤的；因为你的焦虑如此明显地流露，我们几乎能用眼看见你心中无形的活动。因此，无论他造成了什么伤害，他都只伤害了自己；他刺透了自己的灵魂，自己的生命，自己的良心；最后，通过那盲目的残忍行为，他甚至摧毁了自己的好名声——这是连最坏的人也通常不愿丢弃的。因为对一切善人而言，他越是竭力获取恶人的赞许，或越是乐于接受这种赞许，就越是令人厌恶。\par

11. 有什么能比那件事更清楚地证明他并非处于他所伪装的不得已——他声称自己作为别无选择的善人而做了这件恶事——呢？他竟敢以那人的命令充当托辞，而那人却反对如此行事。我们派遣此信的虔诚执事曾与那位我们派去为他们求情的主教同行；因此，请他向大人陈述，皇帝认为最好连正式赦免都不颁发，免得由此将罪行的污名多少加之于他们，而只是发出一道命令，要求他们立即获释，免受一切进一步的烦扰。他出于纯粹无缘无故的残忍，绝非迫于无奈（尽管，或许还有其他我们怀疑的原因，但不必形诸笔墨），竟肆意骚扰教会——正是这教会，他的兄弟曾因恐惧死亡而投奔其庇护怀抱，却因保护他的性命而得到的回报，是发现他积极谋划犯下这罪行——正是这教会，他本人曾因开罪恩主，而寻求那不可被拒绝的荫庇之所。你若爱这个人，就当恨恶他的罪行；你若不愿他陷入永罚，就当厌恶与他为伍。你理应采取此类措施，既为你自己的好名声，也为他的生命；因为谁若爱这个人为天主所憎恶的东西，他其实不仅恨这个人，也恨自己的灵魂。\par

12. 事既如此，我深知你的仁厚，绝不相信你是这罪行的策划者，或参与作恶的同谋，也不会相信你以恶意残忍欺骗我们：这等行径与你的生活和为人相去甚远！同时，我也不愿你的友谊竟具有如此性质，以致促使他自取沉沦，以罪行为荣，并证实人们对你的自然猜疑；反而，但愿那友谊能感动他悔改，且是那种为了医治如此可怖的创伤，所要求的在质与量上都相称的悔改。因为你越憎恨他的罪行，便越真实地做他本人的朋友。我们极愿得知，藉由大人的回信，罪行发生之日你身在何处，你如何得知消息，其后你做了什么，以及你于最近见到他时对他说了什么，又从他听到什么；因为我于次日骤然离去后，对此事与你有关的一切，未能闻悉。\par

13. 至于你信中所说，你现在不得不信，我拒绝访问\Place{迦太基}是因为怕我在那里见到你，你这话反倒逼我明言缺席的理由。一个理由是，我被迫在那城中承受的劳苦——我若详述，势必使这封信再加长数倍——如今已非我能担负；因为除了我自己特有的、我较亲近的朋友皆知的病弱外，我还受着人类家庭共同的软弱，即老年的衰微。另一理由是，在教会所迫切要求，且我之职分特别责成我服务的工作之余暇中，我已决定，若主愿意，将全部时间致力于教会学识的研究；倘蒙天主仁慈，我想借此或能对后世亦有所裨益。\par

14. 既然你愿听实话，你身上确有一事令我极为忧愁：那就是，你虽在年龄、生活及品行上都有资格另作选择，却仍宁愿做个\OriginalTerm{慕道者}{catechumen}；仿佛信者无法在基督信仰及善行中进步，从而在公务管理中成为更忠信有用的人。诚然，促进人类福祉乃你一切重大忧劳的唯一目标。而且，倘若你的公共服务并非归结于此，那么你即便日夜安眠，也胜于牺牲休息，去做那些对同胞毫无裨益的工作。我毫不怀疑大人……\par

（\Emph{其余部分缺失。}）\par

\SourceRecord{Translated by J.G. Cunningham.；Nicene and Post-Nicene Fathers, First Series；Vol. 1.；Edited by Philip Schaff.；Buffalo, NY: Christian Literature Publishing Co.,；1887.；<http://www.newadvent.org/fathers/1102151.htm>.}

% structure: p0001 source=h1 target=section reason=page-title

\section{第158封信（公元414年）}

\Emph{致我主\Person{奥斯定}，我在司铎职上的兄弟兼同工，至诚至爱，满怀敬意，并向与他同在的弟兄们，\Person{埃沃迪乌斯}} \Emph{及与他同在的弟兄们，在主内致候。}\par

1. 我急切地恳求你，回复我上一封信当得的答复。事实上，我本更愿意先得知我那时所问的，然后再提出我现在向你提出的这些问题。请你留意听我讲述一件事，你必会乐于关注，而这件事使我急不可待，只要能在此生中，若有可能，获得我所渴望的知识，便不愿再耽误任何时间。我曾有一个年轻人作我的文书，是\Place{梅洛尼塔}的司铎\Person{阿尔梅努斯}之子，藉着我这卑微器皿的服事，天主在他几乎已深陷世俗事务时救拔了他，因为他当时受雇于总督的法律顾问，担任速记员。那时，他确实如一般少年那样，敏捷而有些浮躁，但随着年龄渐长（因他去世时年仅二十二岁），举止的端庄和一生审慎的正直使他如此卓然，以致追忆他实为一种乐事。此外，他还是一位聪明的速记员，笔耕不倦；他也开始热衷于阅读，甚至敦促我超出自己怠惰所愿，将夜晚的时辰用于读书，因为他每晚在万籁俱寂后，曾有一段时间为我朗读；而在朗读时，除非明白了，他绝不放过任何一句话，定要重复三遍甚至四遍，直到他所想明白的得以澄清为止。我已开始不单视他为一个少年和文书，而看作一位比较亲密而令人愉快的朋友，因为他的言谈给我很多快乐。\par

2. 他也渴望\Quote{离开，与基督同在}，这渴望已经实现。他在父亲家中病了十六天，凭着记忆的力气，几乎在整个患病期间不停地背诵圣经的篇章。但当他生命临近终了时，他大声歌唱，使所有人都可听到：\Quote{我的灵魂渴慕并奔向主的庭院}，随后他又唱道：\Quote{你以油傅抹了我的头，你的杯爵何其甘美，令我神魂颠倒！} 他全神贯注于这些事；在它们所带来的慰藉中，他感到满足。最后，当死亡即将来临时，他开始在额上画十字圣号，画完时手正移到口边，也想以同一圣号标记那里，但在此之前，那内在的人（他实在是日日更新的）\BibleRef{格后 4:16} 已经离弃了泥土的帐幕。我本人被赐予了如此大的欣喜若狂，以至于我想，在他离开自己的身体后，他已进入了我的精神，并在那里从他亲临中赋予我某种圆满的光明，因为我因他的解脱与安然——实在无法言喻——而感受到超乎一切的喜乐。我为他曾感到不小的焦虑，害怕他这年龄特有的危险。因为我不辞辛苦地亲自询问他，是否或许曾为与女人交合所玷污；他庄严地向我们保证自己未染此污点，这声明更增添了我们的喜乐。他就这样去世了。我们以相称的殡葬礼仪纪念他，如此卓越之人理应如此，因为我们一连三天在他的墓前以赞美诗颂扬主，并在第三日献上了救赎的圣事。\par

3. 可是，看，两天之后，一位来自\Place{菲根特斯}、可敬的寡妇，天主的婢女，她说自己已守寡十二年，在梦中看见了以下神视。她看见一位四年前去世的执事，正在天主的男女仆人（贞女和寡妇）协助下，预备一座宫殿。那宫殿被装饰得如此华丽，以致那地方光辉灿烂，似乎全由白银造成。当她急切地询问这座宫殿是为谁预备的，前述执事回答说：\Quote{是为那个年轻人，司铎的儿子，昨天被夺去的。} 在同一宫殿中，出现一位身穿白衣的老人，他命令另外两位也穿着白衣的人，去将遗体从坟墓中启出，与他们一同带上天堂。她又补充说，当遗体一从坟墓中被取去带到天上，同一墓穴中便生出玫瑰枝条，那玫瑰因花瓣紧叠而被称为童贞玫瑰。\par

4. 我讲述了事件：现在，若你愿意，请听我的问题，并教导我所问的，因为那位年轻人灵魂的离去迫使我发出这样的问题。当我们在此肉身内时，我们拥有一种内在的知觉官能，其警觉程度与我们的专注活跃成正比，且我们越专心致志，它就越清醒急切：在我们看来，即使在其最高活动时，它也可能受肉身的拖累而迟缓，因为谁能完全描述心灵通过肉身所遭受的一切！在由肉身所引发的暗示、诱惑、需求及各种苦恼所带来的纷扰与烦恼之中，心灵并未放弃其力量，它抵抗并得胜。有时它被击败；然而，铭记自身的本性，它在这些劳苦的激励下变得更加活跃更加警惕，冲破邪恶的罗网，从而向着更善之事前进。可敬者你必会仁慈地明白我的意思。因此，当我们在此生中时，我们被这些缺欠所阻碍，然而，如经上所载，\BibleQuote{靠着那爱我们的主，我们在这一切事上，大获全胜}。\BibleRef{罗 8:37} 当我们离开这肉身，摆脱一切重担，并摆脱那不停活动的罪时，我们将会如何？\par

首先，我要问，是否可能有某种身体（也许是由四元素之一，即气或以太构成），当灵魂离弃此世之肉身时，它并不脱离\OriginalTerm{无形体原则}{incorporeal principle}，即那严格称为灵魂的实体。因为灵魂的本性是无形的，如果它因死亡而完全脱去形体，那么，所有离世者的灵魂便合成唯一灵魂了。如此一来，那身着紫红袍的富家人与浑身疮痍的拉匝禄，如今会在何处呢？再者，倘若所有无形体的灵魂结合为唯一的灵魂，他们又如何能按各自的功过被区分，一个受苦，一个享福呢？除非，当然，这些事须以象征意义去理解。无论如何，毫无疑问，那些被置于特定处所的灵魂（如那富家人在火焰中，那贫穷人在亚巴郎的怀抱中），是被置于形体之中的。既有不同的处所，便有形体，而灵魂就寓居在这些形体中；纵使赏罚只在良心内被感知，那感知赏罚的灵魂仍处于形体之内。不管那由众多灵魂合成的单一灵魂性质如何，若它真要证明自己确是由众魂聚而为一的实体，它就必须在那不可分割的统一体中，能于同一时刻既悲又喜。然而，若说此灵魂被称为一，仅仅如同无形体的心神被称为一那样——虽然它内含记忆、意志和理智——且若宣称所有这些均为各自独立、无形体的原因或官能，各有其特定的职能与运作方式，彼此毫无妨碍，那么，我认为，在某种程度上可以回答说：在由众魂聚而为一的这一实体中，必然可能有一部分灵魂受罚，同时另一部分灵魂享福。\par

或者，倘若并非如此（即，若无任何形体在灵魂脱离此世肉身后继续与无形体原则结合），又有何能阻止每个灵魂，在离开其在此居住的坚实肉身后，拥有另一形体，使得灵魂始终赋予某种形体以生命呢？又或，灵魂是借着何种形体迁至某处——若确有那迫使其前往的处所呢？就连天使本身，倘若他们不是因所具有的某种形体而被计数，便不能被称作众多，正如真理之主在福音中所说：\BibleQuote{我能祈求圣父，祂如今必为我派遣十二军天使。} \BibleRef{玛 26:53}再者，撒慕尔应撒乌耳之请被召上来时，确是以形体显现的；\BibleRef{撒上 28:14}至于梅瑟，他的肉身虽已埋葬，但福音叙事清楚表明，他曾在肉身中出现在主面前，在那主与门徒退隐的山上。\BibleRef{玛 17:3}在伪经中，尤其是在毫无权威的\WorkTitle{梅瑟奥秘书}里，确曾提到：当他上山赴死之时，因其肉身所具之能，有一个身体被交于土中，而另一个身体则与陪同他的天使结合；但我认为自己并无义务偏重伪经之言，胜于上文所引的确切陈述。因此，我们必须关注此事，并借助启示的权威或理性的光照，探求我们所追问的事理。然而，有人称，将来肉身的复活正证明灵魂死后完全脱去形体。但这一反对并非无法回应，因为天使——他们如同我们的灵魂，本是不可见的——有时愿意显为形体，被人看见；而且（无论这些神魂适于采取何种形体），他们确曾显现，例如，显现给亚巴郎，\BibleRef{创 18:6}也显现给多俾亚。\BibleRef{多 12:16}因此，很有可能的是，肉身的复活——正如我们所深信——确要发生，而灵魂也可与之重新结合，却不必发现它在任一时刻完全不具备某种形体。如今，灵魂在此所居的肉身，由四元素构成，其中一种，即热，似乎与灵魂同时离开此身。因为死后，由土构成的部分依旧存留，水分并不匮乏，冷性元素也未消失；唯有热已消逝，而灵魂若由此迁至彼处，或许便带去了这热。这就是我目前关于形体所要说的。\par

在我看来也是如此：若灵魂居于有生命的肉身时，能如我所说，专注于严格的思维运用，那么，若在此身中已学会喜爱德行，它将会多么无拘无束、活泼有力、恳切坚决、坚忍不拔，其能力何等扩大，品性何等改善！因为，放下这肉身，或更确切地说，在将这云雾扫除之后，灵魂将摆脱一切纷扰，享受宁静，免于诱惑，见到它所渴慕的一切，拥抱它所爱的一切。那时，它也将能够记起并认出朋友，无论是先它离世者，或是它留在尘世者。或许这是真的。我虽不知，却渴望得知。然而，若想到灵魂死后陷入一种麻木状态，如同被埋葬一般，犹如它在肉身之中沉睡之时，只在希望中存活，却一无所有，一无所知，尤其若在沉睡中甚至不被任何梦境所触动，这念头便叫我极度恐惧，并似乎表明灵魂的生命在死亡时便熄灭了。\par

8. 我还要问：假如发现灵魂具有我们所说的这样一种身体，那身体是否缺少任何感官？当然，如果像我所设想的那样，它不可能有嗅觉、味觉或触觉的需要，那么这些感官将会缺失；但对于视觉和听觉，我却犹豫不决。因为，据说魔鬼不是能听到吗（当然，并非在他们所困扰的所有人身上，因为关于那些人还有疑问），即使他们带着自己的身体显现时？至于视觉能力，如果他们有身体却没有视力来指引其运动，他们又如何能从一个地方移动到另一个地方呢？你认为，当人的灵魂离开肉身后，情况不是这样吗——他们仍然具有某种身体，并且至少不缺少某些适合这种身体的感官？否则，我们如何解释这样的事实：非常多已故的人曾在白天被人看见，或者被清醒且在外行走的人在夜间看见，他们进入房屋，就像他们生前习惯做的那样？这件事我听过不止一次，而是经常听到；我还听人说，在埋葬尸体的地方，尤其是在教堂里，有喧闹声和祈祷声，多半在夜间的某个特定时刻被听到。我清楚地记得从不止一个人那里听说过这件事：因为有一位圣洁的司铎亲眼目睹过这样的显现，他看到一大群这样的形貌从圣洗堂里出来，带着充满光的身体，之后他听到他们在教堂中间的祈祷声。所有这类事要么是真的，从而有助于我们正在进行的探究，要么只是无稽之谈，那样的话，它们被编造出来这件事本身就令人惊奇；然而，我还是想从它们会来探访人，并在梦以外被看见这个事实中得到一些信息。\par

9. 这些梦引出了另一个问题。我不在此刻操心那些纯粹由幻念所造的东西，它们是由未受教育者的情绪所形成的。我说的是睡梦中的探访，就像在梦中对\Person{若瑟}的显现\BibleRef{玛 1:20}那样，以大多数这类情况所经历的方式。因此，同样地，在我们之前离世的朋友们有时也会在梦上来向我们显现，并对我们说话。因为我自己记得，\Person{普罗福图鲁斯}、\Person{普里瓦图斯}和\Person{塞尔维留斯}，在我记忆中是从我们隐修院被死亡带走的圣洁之人，曾对我说话，而他们所谈及的事果然如其言成就了。或者，如果是另外某种更高的神灵取了他们的形象并访问我们的心灵，我把这留给那洞悉一切、从最高到最低都显露无遗者的全见之眼。因此，如果天主乐意通过理性向你的圣善讲论所有这些事，我恳求你发发慈心，让我分享你所领受的知识。还有一件事我决意不可略去不提，因为它或许与我们当前所查究的问题有关：\par

10. 正是这个年轻人，因他而引发了这些问题，他在临终时可以说是领受了召唤后离世了。因为他的一个同伴，曾同为学子，担任读经员和我的口述速记员，于八个月前去世，有人在梦中看见他来接近这年轻人。当那个清楚看到他的人问他为什么来时，他说：\Quote{我来是要带走这位朋友。}结果果真如此。因为在这家屋里，也有一位几乎醒着的老人看见一个手里拿着一条月桂枝的人显现，枝上写着些东西。不仅如此，当这个形貌被看见后，据进一步报导，在这年轻人死后，他那位司铎父亲开始与年迈的\Person{提亚修斯}一起住在隐修院里，以寻得安慰，但看！在他死后第三天，年轻人被看见进入隐修院，某位弟兄在某种梦中问他是否知道自己已死。他回答知道。那人问他是否已蒙天主悦纳，他极其欢喜地作了肯定的回答。当被问及为什么来时，他回答说：\Quote{我被派来召唤我的父亲。}蒙示这些事的人醒来，述说所发生的。这事传到\Person{提亚修斯}主教耳中。他吃惊了，严厉告诫告知他的人，以免这事很容易地传到司铎本人耳中，使他被这样的消息搅扰。但何必拖长叙述呢？在这次探访后大约四天之内，他（指那位司铎，因他曾发轻度热症）说他已经脱离危险，医生已停止给他看病，向他保证没有任何可忧虑的原因；但就在那一天，这位司铎在他躺下后便断气了。我也不应不提，在年轻人离世那天，他三次请求父亲宽恕他可能有所冒犯的任何事，每一次亲吻父亲时都对他说：\Quote{父啊，让我们感谢天主，}并坚持要父亲和他一起说这话，仿佛他在劝勉一位将与他一同离开此世的同伴。事实上，两次死亡之间只隔了七天。对于如此奇妙的事，我们能说什么呢？谁能成为这些奥秘安排完全可靠的导师呢？在困惑的时刻，我激动的心向你倾吐。年轻人及其父亲的死亡由天主所命定是毫无疑问的，因为没有我们天父的旨意，两只麻雀也不会掉在地上。\BibleRef{玛 10:29}\par

11. 在我看来，灵魂不能绝对脱离某种身体而存在，这一事实被证明：无身体的存在只属于天主。但我认为，在离开此世的过程中，卸下身体如此巨大的重担，证明了灵魂那时将比现在更加清醒；因为那时，我认为灵魂不再受如此大的障碍所累，在行动和知识上都显得远为高贵，而全然的神魂安息证明它摆脱了所有扰乱和错误的原因，但这并不会使它变得无力、迟缓、麻木和困窘，因为灵魂已充分享有它因脱离世界和身体而获得的自由，这便足够；因为，正如您智慧地所言，在那幸福而蒙福的境界中，理智得享饱饫，并以神魂的双唇贴近生命之泉，在自身官能的无可争议的主权中；因为我离开隐修院之前，曾在梦中见到已故的\Person{塞维利乌斯}兄弟，他说我们正努力通过理性的运用达到对真理的理解，而他与那些与他同在一处的人，则始终安息于纯然静观的喜乐中。\par

12. 我也恳请您向我解释\OriginalTerm{智慧}{wisdom}一词有几种用法：比如，天主是智慧，睿智的心灵也是智慧（在这种情况下，它被称为如同光明）；我们也读到\Person{贝匝肋耳}的智慧，他建造会幕或调配制香膏，以及\Person{撒罗满}的智慧，或任何其他智慧（如果存在的话），并且它们彼此间有何不同；以及那与圣父同在的唯一永恒智慧，是否应在这些不同的层次上来理解，正如它们被称为圣神的不同恩赐，圣神随自己的心意，个别地分配与各人。或者，除了那非受造的智慧本身之外，这些智慧都是受造的吗？它们是否为独立的存在？还是说它们只是效果，并从其工作定义的名称而得名？我问了许多问题。愿主赐您恩宠，既发现所寻求的真理，又有足够的智慧将其付诸笔墨，并尽快传递给我。我是在极度无知中，以朴素的文笔写成的；但既然您认为值得了解我所询问的事，我因主基督之名恳求您，在我有误之处纠正我，并将您知道而我渴望学习的教导我。\par

\SourceRecord{Translated by J.G. Cunningham.；Nicene and Post-Nicene Fathers, First Series；Vol. 1.；Edited by Philip Schaff.；Buffalo, NY: Christian Literature Publishing Co.,；1887.；<http://www.newadvent.org/fathers/1102158.htm>.}

% structure: p0001 source=h1 target=section reason=page-title

\section{第159封信（主历415年）}

\Emph{致\Person{埃沃迪乌斯}，我至为有福的主，我尊敬亲爱的弟兄和司铎职的同工，并致与他同在的弟兄们：\Person{奥古斯丁}和与他同在的弟兄们在主内致意。}\par

1. 我们的弟兄\Person{巴尔巴鲁斯}，这封信的携带者，是天主的仆人，他已在\Place{希波}定居许久，并一直是天主圣言的热切而勤勉的聆听者。他请求我们写这封信给您的圣德，借此我们在主内向您推荐他，并通过他向您转达我们有责任奉上的致意。要回复您那些信中交织著极其困难的问题，即使对于有闲暇且拥有远胜于我们的讨论能力和敏锐理解力的人，也是一项极繁重的任务。那两封信中，其中一封提出了许多重大问题，却不知何故遗失了，虽经长期寻找仍未能找到；另一封已找到的，则包含了一个关于天主仆人的非常令人愉悦的记述，讲述一位善良而贞洁的青年如何离世，以及如您所说，通过弟兄们的异象所传达的见证，他的功德便向您显明了。借这个青年的案例，您提出并讨论了一个关于灵魂的极隐晦的问题――灵魂离开此身时，是否与另一种身体相结合，藉此能被搬运或拘禁于有物质界限的地方？如果像我们这样存在的人，确实能对此问题进行满意的研究，那需要最勤勉的专注与辛劳，因而也要求一颗完全摆脱那等耗费我时间的俗务、处于闲暇的心灵。然而，如果您愿意聆听，我概括为一句话的看法是：我绝不相信灵魂离开身体时，有任何种类的另一身体伴随。\par

2. 至于这些异象和未来事件的预知是如何产生的这个问题，就让那明白是何能力使我们在人思绪繁忙时，心中产生那等巨大奇妙现象的人，去尝试解释它们吧。因为我们都看到并清楚觉知，心灵中会产生无数由眼目或其他感官可辨识的众多事物的影像――它们是有序生成还是杂乱出现，目前与我们无关：我们所要说的只是，既然这些影像的产生是无可置疑的，那么，唯有那能够阐明是以何种能力、以何种方式解释这每日持续经验的种种现象的人，才可适切地冒险去推测或提出这些极罕见异象的任何解释。至于我，越是在一生中、醒时或睡时，更清楚地发现自己无力解释我们经验的寻常事实，便越不敢贸然解释那非凡之事。因为就在我向您口授这封信时，我一直在心中默想著您的形象――当然，您始终不在场，对我的思想一无所知――并且我根据对您内心的了解，想象您会如何感受我的话；而我既不能通过观察，也不能通过探究，明白这过程是如何在我心中完成的。但有一点我可以肯定：虽然那心理影像非常类似某种物质的东西，但它既不是由物质的团块、也不是由其性质产生的。请暂且接受这封来自一个在其它职责压力下匆促写信的人的信。在我论\BibleBook{创世}{纪}的第十二卷书中，对此问题有极详细的讨论，那篇论述充满了由实际经验或可靠叙述所得的大量例证。至于我有多大能力处理这问题，以及我完成了什么，您读了那部作品后自会判断；如果主仁慈地恩准我现在将那些已经尽我所能系统修订的书卷出版，以回应众多弟兄的期待，而不是继续延迟他们的希望，继续讨论那已占用我很长时间的主题的话。\par

3. 然而，我要简短叙述一件事，供您默想。您认识我们的弟兄\Person{革纳迪乌斯}，一位几乎人人皆知的医生，我们非常亲爱，他现今住在\Place{迦太基}，早年在\Place{罗马}曾是一位杰出的医师。您知道他是一个虔敬且极仁慈的人，对穷人照顾积极、充满同情、时刻慷慨。然而，甚至他，在年轻时，且最热心于这些爱德行动时，如他自己告诉我的，有时也会怀疑死后是否有生命。因此，既然天主绝不会离弃一个性情和行为如此慈悲的人，他在睡梦中见到一位外貌非凡、威严堂堂的青年，对他说：\Quote{跟随我。}跟随他，来到一座城，在那里他开始听到右边传来极其悦耳的旋律，甜美得超越他以往所听过的一切。当他问那是什么时，向导说：\Quote{这是蒙福者和圣者的歌咏。}据他报告，他在左边看到的，我已记不清了。他醒过来；梦消失，他以为那只是一个梦。\par

4. 但是，在第二夜，同一个青年显现给\Person{Gennadius}，问他是否认得他，他回答说完全认得他，毫无疑惑。于是青年问\Person{Gennadius}是在哪里认识他的。在那里，他的记忆也没有辜负合适的回答：他叙述了整个异象，以及在他引导下听到的圣人们的赞美诗，带着回忆最近经历的一切从容。青年就问是在睡时还是醒时看见了他所叙述的。\Person{Gennadius}回答：\Quote{在梦中。}青年接着说：\Quote{你记得很清楚；确实你在梦中看到了这些事，但我愿意你知道，甚至现在你也是在梦中观看。}听到此，\Person{Gennadius}信服了其真实，在回答中声明他相信。然后他的导师继续说：\Quote{你的身体现在在哪里？}他答：\Quote{在我的床上。}\Quote{你知道，}青年说，\Quote{你身体的这双眼睛现在被绑着、闭着、休息着，你用这双眼睛什么也看不见吗？}他答：\Quote{我知道。}\Quote{那么，}青年说，\Quote{你用什么样的眼睛看见我？}他答不上来，默然无语。在他犹豫时，青年向他揭示了通过这些问话要教导他的道理，立即说：\Quote{正如你睡着躺在床上的时候，你身体的这双眼睛现在闲置无用，而你却有眼睛能看见我，享受这异象；同样，死后，你身体的眼睛完全无用时，你里面将有一个生命使你仍然活着，有一个知觉能力使你仍然感知。所以，从此以后你要小心，不要再怀疑人死后生命是否继续。}这位信徒说，藉此方式，他对此事的一切疑惑都被消除了。是谁这样教导了他，岂不是经由天主仁慈的眷顾吗？\par

5. 有人可能说，藉此叙述我非但未解决问题，反而使问题复杂化了。然而，虽然每个人都可以自由相信或不相信这些叙述，但每个人都有他自己的意识作为导师在近旁，可以藉其帮助而致力于这个最深奥的问题。人每天醒寤、睡眠、思想；那么，让任何人回答：那些虽非物体的东西，却类似物体的形状、属性和动作，是从哪里来的呢？——让他回答这个，若是能够。如果他不能做到这一点，为什么他竟这样急于对那些非常罕见或超出他经验范围的事情下断语，而不能解释每日恒见的事情呢？至于我，虽然我完全无法用言语解释那些似乎有形质却没有任何真实形体的相似物是如何产生的，但我可以说，我愿意，如同我确信这些事并非由身体产生那样，也能确知那些偶尔被灵魂见到、且被认为由身体感官见到的东西，是怎样被感知的；或者，那些被迷惑，或更常见的是被不虔所误导的人们所见的异象，有什么独特标记，因为这些异象的例子如此众多，类似于虔诚和圣洁人们所见的异象，以致我若想引用它们，时间会不够，而非例子不足。\par

愿您，藉着主的仁慈，在恩宠中成长，最可称颂的阁下，可敬而亲爱的兄弟！\par

\SourceRecord{Translated by J.G. Cunningham.；Nicene and Post-Nicene Fathers, First Series；Vol. 1.；Edited by Philip Schaff.；Buffalo, NY: Christian Literature Publishing Co.,；1887.；<http://www.newadvent.org/fathers/1102159.htm>.}

% structure: p0001 source=h1 target=section reason=page-title

\section{书信 163（主后 414 年）}

\Emph{埃伏第乌斯主教向奥斯定主教问安。}\par

前一段时间我向主教阁下提出两个问题；第一个，我想是由隐修院仆人\Person{乔比努斯}送去的，涉及天主和理性，第二个是关于一个观点，即救主的身体是否能看见\OriginalTerm{天主的本体}{substance of the Deity}。我现在提出第三个问题：我们救主随同其身体所取的理性灵魂，是否属于在讨论灵魂起源时通常提出的任何理论（如果确有某种理论能够确定的话）——或者说，祂的灵魂虽然是理性的，却并不属于生灵灵魂所划分的任何一类，而是属于另一类？\par

我也问第四个问题：宗徒伯多禄用以下的话为天主作证时，所指的那些灵魂是谁：\BibleQuote{就肉身说，祂固然被处死了；但就神魂说，祂却复活了。祂借这神魂，也曾去给那些在狱中的灵魂宣讲过。} 这使我们理解到，他们是在地狱里，基督下降地狱，向他们所有人宣讲了福音，并借着恩宠将他们所有人从黑暗和惩罚中解救出来，以至从主复活之时起，人们等待的是审判，而那时地狱已被完全清空。\par

我恳切希望知道主教阁下对此事持有怎样的信念。\par

\SourceRecord{Translated by J.G. Cunningham.；Nicene and Post-Nicene Fathers, First Series；Vol. 1.；Edited by Philip Schaff.；Buffalo, NY: Christian Literature Publishing Co.,；1887.；<http://www.newadvent.org/fathers/1102163.htm>.}

% structure: p0001 source=h1 target=section reason=page-title

\section{信函一六四（公元四一四年）}

\Emph{致我至福之主\Person{埃沃迪乌斯}，我在主内共主教职的兄弟与同工，\Person{奥古思定}致以主内问候。}\par

% structure: p0003 source=h2 target=subsection reason=relative-heading-rank; base=h2

\subsection{第一章}

1. 你从宗徒\Person{伯多禄}的书信中所提出的问题，正如我想你意识到的那样，是最常使我深感困惑的一个问题，即：倘若认为你所引用的那些话是指\OriginalTerm{地狱}{hell}说的，那么这些话应如何理解？因此，我把这问题再交还给你，好让你自己，或者能找到其他能解决的人，来消除并了结我在这个问题上的困惑。如果主在我之先赐给我明白这些话的能力，而我又能将我从他领受的传授予你，我必不向如此深爱的朋友隐瞒。当下，我将那书信中使我感到困难之处告知于你，好让你留意我对宗徒这些话的评注，去反复思考，或者去请教任何你认为有资格发表见解的人。\par

2. 宗徒在说过\BibleQuote{就肉身说，基督被置于死地；就神魂说，他复活了，}这句话后，紧接着说：\BibleQuote{他藉这神魂，曾去给那些在狱中的神魂宣讲；这些神魂从前在诺厄建造方舟的时日，天主耐心期待之时，原是不信的人；当时赖方舟经由水而得救的不多，只有八个生灵。}此后他又加上一句：\BibleQuote{这水所预表的圣洗，如今也拯救你们。}\BibleRef{伯前 3:18}因此，这使我感到困难。如果主在死后下到地狱，给那些在狱中的神魂宣讲，那么为什么只有那些在诺厄建造方舟时仍不信的人，才被算为堪当此恩惠，即主降临地狱？因为在诺厄的时代与基督受难之间，有无数国家成千上万的人死亡，主在地狱中本可以找到他们。当然，我这里不是指那些在那段时期内信了天主的人，例如\Person{亚巴郎}家族的众先知和圣祖，或者更早的\Person{诺厄}本人和他的一家，他们藉着水得救（或许除开那后来被弃绝的一个儿子），此外还有\Person{雅各伯}后裔之外所有信天主的人，如\Person{约伯}、\Place{尼尼微}的居民，以及任何其他不论圣经中提及的，还是在广大人类家庭中为我们所不知的任一时代的人。我只是指那些从\Person{诺厄}到基督受难期间，不认识天主，专务敬拜魔鬼或偶像，从今生逝去的无数万人。基督在地狱中找到了他们，为什么不向他们宣讲，却只向那些在\Person{诺厄}建造方舟时不信的人宣讲？或者，如果他向所有人宣讲了，为什么\Person{伯多禄}只提到这些人，而无视了其他无可计数的众人？\par

% structure: p0006 source=h2 target=subsection reason=relative-heading-rank; base=h2

\subsection{第二章}

3. 毫无疑问，主在就肉身被处死之后，\Quote{下降地狱}了；因为谁也不能否认那预言论及的：\BibleQuote{你绝不会将我的灵魂遗弃在地狱}——\Person{伯多禄}自己在\WorkTitle{宗徒大事录}中解释过这句话，以免有人敢对它另作解释——以及同一位宗徒肯定主的话：\BibleQuote{他解除了地狱的痛苦，因为他不能受地狱的控制。}所以除了不信者，谁会否认基督曾去过地狱呢？至于如何调和他解除了地狱的痛苦，而他自己从未被这些痛苦束缚过，他解除它们并非他挣脱了捆绑他的锁链，这困难是容易解决的，只要我们明白那些痛苦被解除的方式，如同猎人的圈套被解开以防止它们套住猎物，不是因为它们已经套住了猎物。这也可以理解为教导我们相信他解除了那些不可能拘禁他的痛苦，而是拘禁着那些他决意要解救的人。\par

4. 但是这些人是谁，由我们来界定是僭越的。因为我们若说那里所有的人都毫无例外地得到了解救，如果我们能证明这一点，谁不欢喜呢？人们尤其会为某些因文学著作而为我们熟知的人欢喜，我们赞赏他们的口才与才华——不仅那些诗人和演说家，他们在著作的许多部分中鄙视讥讽这些列国的假神，甚至偶尔宣认唯一真天主，尽管他们与其他人一样遵守迷信礼仪，而且也包括那些不以诗歌或修辞，而是以哲学家身份说出同样道理的人；并且也为更多没有文献遗存的人欢喜，但我们从这些人的著作中得知，他们的生活在一定程度上是值得称赞的，以至于（除了他们的事奉天主方面，他们在其中犯了错，崇拜那些被立为公开敬拜对象的虚幻之物，事奉受造物而非造物主）他们可以被公正地奉为其他各种德行的模范：节俭、克己、贞洁、节制、为保卫国家视死如归，以及不仅对同胞，甚至对敌人忠信不渝。然而，这一切若并非为了纯正真虔敬的伟大目标，而是为了人的赞美和荣耀的空虚骄傲去实践时，在某种意义上就变得无价值且无益的；尽管如此，作为某种心性倾向的标示，它们如此使我们喜悦，以致我们希望具有这些品质的人——无论是特别优先还是与其他人一起——从地狱的痛苦中得到解救，若非人的感觉之判断不同于造物主公义之判断。\par

5. 事既如此，倘若救主从那里解救了所有人，并引用你信中所提问题的话语，\Quote{清空了地狱，以致从那时起只待最后的审判，} 那么以下各点对此问题引起并非不合理的困惑，每当我思考时，这些困惑便常出现在我面前。首先，有何权威论断可证实此见解？因为圣经所言基督的死使\BibleQuote{地狱的痛苦得以解除}，并不能确立此见解，因为这句话可理解为指基督自己，意即祂解除了（即使之无效）地狱的痛苦，以致祂自己不被它们拘禁，尤其因为经上又补充说，祂\BibleQuote{不可能被它们拘禁}。或者，若有人[反对此解释]问及祂为何选择降到地狱，那里有痛苦，却不能拘禁祂那位——如圣经所说，\BibleQuote{在死者中自由者}——死亡的王子和统领在祂身上找不到任何应受惩罚之处，那么\BibleQuote{地狱的痛苦得以解除}这句话可理解为不是指所有人，而只是指祂判定堪当此解救的某些人；这样，既不能认为祂降到那里是徒劳无益的，无意为那些被囚禁在那里的人带来好处，也不能必然推断天主仁慈和正义赐予某些人的，必定也要赐予所有人。\par

% structure: p0010 source=h2 target=subsection reason=relative-heading-rank; base=h2

\subsection{第三章}

6. 至于第一人，即人类之父，几乎整个教会都同意主把他从那牢狱中解救出来；这一信念应当认为并非毫无理由地被接受——无论它是从何种来源传至教会的——尽管正典圣经的权威不能被引用来明确支持此说，但\WorkTitle{智慧篇}中的这些话语似乎最有力地支持了这一意见。\BibleRef{智 10:1} 有些人还[传统上]加上，古代的圣人们——\Person{亚伯尔}、\Person{舍特}、\Person{诺厄}及其一家、\Person{亚巴郎}、\Person{依撒格}、\Person{雅各伯}，以及其他圣祖和先知们——也得到了同样的恩惠，当主降到地狱时，他们也同样从那些痛苦中被解救出来。\par

7. 但是，依我看来，我无法理解\Person{亚巴郎}——比喻中那位虔敬的乞丐也被接收到他怀里——怎能被认为曾处于那些痛苦之中；能者或许可以解释这一点。但我想，每个人都必看出，若想像在\Person{亚巴郎}那奇妙安息的怀抱中，在主降到地狱之前，只有两人，即\Person{亚巴郎}和\Person{拉匝禄}，并认为只为这两个人，才向那富人说了话：\BibleQuote{在我们与你们之间，隔着一个巨大的鸿沟，以致人即便愿意，从这边到你们那边也不能，从那边到我们这边也不能。}\BibleRef{路 16:26} 这显然是荒谬的。更何况，如果那里不止有两人，谁敢说圣祖和先知们不在那里呢？他们的义德和虔敬，在天主的话语中受到了如此卓越的见证。在那情形下，那位解除地狱痛苦的主——他们并未被这些痛苦拘禁——究竟给他们带来了什么益处，我至今还不明白，尤其因为我在圣经中找不到任何地方将地狱一词用于好的意义。而若在圣经中找不到这种用法，那么\Person{亚巴郎}的怀抱，即某一宁静隐退之所，就决不应被视为地狱的一部分。不，从那位伟大导师亲口所说\Person{亚巴郎}讲的：\BibleQuote{在我们与你们之间，隔着一个巨大的鸿沟}这句话中，我认为已足够明显地看出，那光荣真福的怀抱并非地狱的任何组成部分。因为那巨大的鸿沟，若非一道完全分隔两侧的深渊，又是什么呢？它不仅存在，而且是固定的。因此，倘若圣经没有提及地狱及其痛苦，就说基督死时去了\Person{亚巴郎}的怀抱，我怀疑是否有人敢说祂\Quote{降到地狱}。\par

8. 但是，鉴于明显的圣经见证提到了地狱及其痛苦，没有理由相信祂，即救主，去了那里，除非祂要拯救人脱离其痛苦；但祂是救了所有祂发现陷于痛苦中的人，还只是救了祂判定堪当此恩惠的某些人，我仍在问：然而，祂确实曾在地狱，并给那些遭受这些痛苦的人带来了这种益处，这我毫不怀疑。因此，我尚未发现当祂降到地狱时，给那些在\Person{亚巴郎}怀抱中的义人带来了什么益处，因为我看到，就祂天主性的真福临在而言，祂从未离开过他们；因为正是在祂死了的那一天，祂应许那盗贼当天就与祂同在乐园，而那时祂正要降到地狱去解除地狱的痛苦。所以，最确定的是，在那之前，祂以祂的真福智慧已在乐园和\Person{亚巴郎}的怀抱中，同时以祂的惩罚之能也在地狱；因为天主性不受任何界限的限制，何处没有祂的临在呢？然而，就祂在一特定时刻所取的受造本性而言——祂在继续是天主的同时，成为了人——就是说，就祂的灵魂而言，祂曾在地狱：这在圣经的这些话中明确宣告，这些话既已在预言中预先说出，也由解释充分说明：\BibleQuote{你不会将我的灵魂遗弃在地狱里。}\par

9. 我知道有些人认为，在基督死亡之时，一次如同在世界末日应许给我们的复活，已赐给了义人；他们的根据是圣经的记载：在祂死亡的那一刻，发生地震，岩石崩裂，坟墓洞开，许多圣人的身体复活了，在祂复活后，与祂一同显现于圣城。诚然，如果这些圣人没有再次安眠，即他们的身体没有被再次葬入坟墓，那么就必须考察，如果这么多人在祂之前复活了，基督又如何能被称为\BibleQuote{死者中的首生者}\BibleRef{默 1:5}呢。如果有人对此回答说，经文是提前叙述的，即地震在基督被钉十字架时确实使坟墓洞开了，但圣人的身体当时并未复活，而是在基督先复活之后才复活的——即便按此经文提前叙述的假设，加上的这些话也不会妨害我们仍然相信，一方面，基督毫无疑问是\BibleQuote{死者中的首生者}，另一方面，这些圣人蒙准许，在祂先复活之后，也复活进入一种永恒的不朽不灭状态；但仍然存在一个困难，即，在这种情况下，\Person{伯多禄}怎么能那样说，并且他所说的无疑是绝对真实的，当他断言，在上面引用的预言中，\BibleQuote{祂的肉身不见腐朽}这句话，不是指\Person{达味}，而是指基督，并且他补充关于\Person{达味}的话，\BibleQuote{他已被埋葬，他的坟墓至今还在我们这里}\BibleRef{宗 2:28}——如果\Person{达味}的遗体在坟墓中并非安然无恙，这句话作为论据便毫无力量；因为当然，即便圣人的遗体在死后立即复活了，并且因此没有见到腐朽，坟墓可能仍然存在。但如果说\Person{达味}没有包括在这些圣人的复活之内，而圣人们却获得了永生，这似乎困难，因为经上如此频繁、如此清楚，并且以如此光荣的方式提及\Person{达味}的名字，宣告基督将出自\Person{达味}的后裔。此外，\BibleBook{希伯来人}{书}中关于古时信者的这些话：\BibleQuote{天主为我们预备了更美好的事，叫他们若不与我们同得，就不能完全}\BibleRef{希 11:40}，也将受到威胁，如果这些信者已被建立在那永不朽坏的复活状态中，那是应许给我们，在世界终结我们得以完全之时所要得到的。\par

% structure: p0015 source=h2 target=subsection reason=relative-heading-rank; base=h2

\subsection{第四章}

10. 你因此看出，\Person{伯多禄}为何选择只提及那些在\Person{诺厄}建造方舟的日子不信从的人，作为在监狱中被囚禁时得到福音传报的对象——这个问题是何等复杂，以及种种困难如何阻止我对此发表任何定见。我犹豫不决的另一个原因是，在宗徒说了以下的话之后：\BibleQuote{这圣洗如今藉着耶稣基督的复活，也拯救了你们：它并非涤除肉体的污秽，而是向天主求有一颗纯洁的良心；这位基督已升了天，坐在天主的右边，众天使、掌权者和异能都屈伏于祂，因为祂已吞噬了死亡，使我们成为永生的承继者。}他又补充说：\BibleQuote{所以，基督既然在肉身内为你们受了苦难，你们就该用同样的心志装备自己，因为在肉身内受苦的，就戒断了罪恶；为能在今后，不再随从人的情欲，只随从天主的旨意，在肉身内度其余生。}然后他继续说：\BibleQuote{过去度日足以使我们随从外邦人的欲望，生活在放荡、情欲、酗酒、宴乐、狂欢和违法的偶像崇拜中；由于你们不再同他们狂奔于纵欲败坏的漩涡，他们便觉得奇怪，遂毁谤你们；但他们要向那已准备好审判生者死者的主交账。}在说了这些话之后，他又附加说：\BibleQuote{也正是为此，福音也曾传报给死者，使他们虽在肉身中按人受审判，但在神魂中却按天主生活。}\par

11. 面对如此深奥的话语，谁能不感到困惑呢？他说：\BibleQuote{给死者宣讲了这福音}；如果我们将\Quote{死者}理解为已经离开肉身的人，我想他指的必定是上文所说的\BibleQuote{在诺厄时代不信的}，也必定是基督在\OriginalTerm{地狱}{hell}中所遇见的一切人。那么，\BibleQuote{使他们虽然肉身如人一样受审判，但神魂按天主而生活}是什么意思呢？他们怎么可能在肉身上受审判呢？如果他们在地狱中，便不再有肉身；如果他们已经从地狱的刑罚中得到解脱，那他们也尚未复得肉身。因为即使如你在问题中所提出的，\Quote{地狱被清空了}，也不能相信当时在那里的所有人都已经肉身复活了，或者那些复活后与主一同显现的人是为了这个目的而重新获得肉身，好使他们在肉身内按着人受审判；但我看不出，对那些在\Person{诺厄}时代不信从的人来说，这怎么可能成立，因为圣经并未确认他们被赋予肉身的生命，也不能相信他们得脱地狱的刑罚，是为了让这些被解救者重新获得肉身以便承受惩罚。那么，\BibleQuote{使他们虽然肉身如人一样受审判，但神魂按天主而生活}到底是什么意思呢？难道是说，基督在地狱中所遇见的人得到了这样的恩赐，即他们借着福音在\OriginalTerm{神魂}{spirit}上获得生命，尽管在未来的复活中，他们必须在肉身内受审判，好让他们通过肉身内的一些惩罚而得以进入天主的国？如果这就是这节经文的意思，那么为什么只有\Person{诺厄}时代的那些不信者（而不是基督降入地狱时在那里所遇见的所有其他人）借着福音的宣讲而在神魂上获得生命，然后再在肉身内受一个有限期的惩罚呢？然而，如果我们认为这适用于所有人，问题仍然存在：为什么\Person{伯多禄}只提到那些在\Person{诺厄}时代不信从的人，而没有提到其他人呢？\par

12. 此外，我发现，那些试图对这件事提出解释的人所提出的理由也有困难。他们说，当基督降到地狱时，在那里发现的所有人从未听过福音，那个刑罚或监禁之地因为福音在他们生前没有传遍天下，他们因未曾听见宣报而有足够的理由不信，因此所有的人都从那里被释放了；但是从那时起，既然福音已传遍各国，到处宣扬，人若轻慢福音便无可推诿，所以自那监牢被清空以后，仍留有公正的审判，其中顽固不信的人要遭受永火的惩罚。持这种意见的人没有考虑到，同样的借口也同样适用于所有在基督复活之后、福音传给他们之前离世的人。因为即使主已从地狱回来了，并不是说从此以后未经听过福音的人就可以不下地狱，事实上，世界上众多的人在福音佳音传给他们之前就已死去，所有这些人都有权提出那个借口，而这个借口据说是从那些未听过福音的人那里被剥夺的，因为据说主下降到地狱时曾给他们宣讲了福音。\par

13. 或许有人会这样答复这种异议：那些在主复活之后去世，或现今在离世之前未闻福音宣讲的人，也可能在他们在的地方（即地狱里）听到过或正在听到福音，从而在那里他们相信应该相信的有关基督真理的事，并获得基督曾宣讲过的人所获得的宽恕和救恩；因为基督从地狱复活升天，并不等于关于他的消息在那里就从记忆中消失了，因为在这地上他也是升到了天上，然而借着福音的宣扬，凡相信他的人必将得救；而且，\BibleQuote{天主极其举扬他，赐给了他一个名字，超越其他所有的名字，致使上天、地上和地下的一切，一听到耶稣的名字，无不屈膝叩拜。}但是，如果我们接受这种意见，认为生前不信的人可以在地狱里相信基督，那么随之而来的理性与信仰的矛盾谁能承受呢？首先，如果这是真的，我们似乎就没有理由为那些未蒙恩宠而离开肉身的人哀悼了，也没有理由急切地敦促人应在死前接受天主的恩宠，以免他们遭受永死。再者，如果说地狱里只有那些在世上拒绝接受已宣讲给他们的福音的人，他们的相信是徒劳无功的，而相信对那些从未轻视过福音的人——因为他们根本没有能力听到福音——则是有益的，那就会引出另一个更荒谬的结论：既然所有人都必定死去，而且他们都应完全摆脱轻视福音的罪责进入地狱；否则他们在地狱里的相信将毫无用处，那么福音就不该在地上宣讲了，这种想法既不智而又亵圣。\par

% structure: p0020 source=h2 target=subsection reason=relative-heading-rank; base=h2

\subsection{第五章}

14. 因此，让我们最坚定地持守\OriginalTerm{信德}{faith}所宣称的——这信德建立在毫无疑义的权威之上——即：\BibleQuote{基督照圣经记载死了}，且\BibleQuote{祂被埋葬了}，且\BibleQuote{祂第三天照圣经记载复活了}，以及所有在已被充分证实为真实的文献中记载关于祂的其他事。在这些教义中，我们包括祂下降\OriginalTerm{地狱}{hell}，并解除了地狱的痛苦——那痛苦原不可能拘禁祂——并合理地相信祂也由那里解救和释放了祂所愿的人，然后祂重新取回那曾在十字架上留下、并安放在坟墓中的身体。我说，让我们坚定持守这些事；但至于你从宗徒\Person{伯多禄}的话中提出的问题，既然你现在看到了我在其中遇到的困难，并且如果更仔细研究这个问题，可能还会发现其他困难，那么让我们继续探讨它，无论是通过我们自己对该问题的思考，还是通过询问任何适合且可能咨询的人的意见。\par

15. 然而，我恳请你思考，宗徒\Person{伯多禄}关于那些被囚禁在监狱里的灵魂所说的一切——他们在\Person{诺厄}的日子里不信——是否可能根本不是指地狱，而是指那些时代，他将其预像意义转移到了现今。因为那件事原是未来事件的预像，好使那些在我们这时代不信福音的人——此时教会正在万民中被建立——可以被理解为如同那些在方舟被建造的时代不信的人；同时，那些相信并借着圣洗得救的人，可以比作那些当时在方舟里借着水得救的人；为此他说：\BibleQuote{这水所预表的圣洗，如今也拯救你们。}因此，让我们这样解释其余关于不信之人的话，使之与这预像的类比相和谐，并拒绝这样的想法：福音曾在地狱里被宣讲，或者甚至如今还在那里宣讲，为使人相信并脱离其痛苦，好像在那里也如同在地上建立了一个教会一样。\par

16. 那些从\BibleQuote{祂去向那些在监狱里的灵魂宣讲}这话推断出\Person{伯多禄}持有那使你困惑的见解的人，在我看来，是被这样的想象所误导：他们认为\Quote{灵魂}一词不能用来指那些当时仍在人身体内的灵魂，这些灵魂因被封闭在无知的黑暗中，可以说是\Quote{在监狱里}——这样的监狱正是圣咏作者在祈祷中所寻求解脱的：\BibleQuote{求你领我出离这监狱，好让我赞美你的名}；这在别处被称为\BibleQuote{死亡的阴影}，从那阴影中得到解救的，当然不是在地狱里，而是在这个世界上，就是那些经上指着说：\BibleQuote{在死亡的阴影之地居住的人，有光照射他们。}\BibleRef{依 9:2}但在\Person{诺厄}时代，福音是徒然宣讲了，因为他们不相信，当天主的恒忍在方舟被建造的许多年里等待他们的时候（因为方舟的建造本身在某种意义上就是一种仁慈的宣讲）；正如现今那些与他们相似的不信者，用同样的预像来说，他们被封闭在无知的黑暗中如同在监狱里，徒然观看正在全世界被建造的教会，而审判即将来临，就像那时洪水淹死了所有不信者一样；因为主说：\BibleQuote{在诺厄的日子里怎样，在人子的日子里也要怎样：那时人们吃喝婚嫁，直到诺厄进入方舟的那一天，洪水来了，把他们全都消灭了。}\BibleRef{路 17:26}但因为那件事也是一个未来事件的预像，那洪水对信徒来说是圣洗的预像，对不信者来说是毁灭的预像，正如在那个预像中，不是借着事件而是借着言语，预言了关于基督的两件事：在圣经中祂被描述为一块石头，注定对不信者是绊脚石，对信徒是基石。然而，有时在同一个预像中，无论是典型事件的形式还是比喻的形式，两样东西被用来代表一样，正如信徒既由建造方舟的木材代表，又由在方舟中得救的八个灵魂代表，又如\BibleBook{若望}{福音}中羊栈的比喻，\Person{基督}既是牧人又是门。\BibleRef{若 10:1}\par

% structure: p0024 source=h2 target=subsection reason=relative-heading-rank; base=h2

\subsection{第六章}

17. 不要让以下这一点被视为对上述解释的反对：宗徒\Person{伯多禄}说，\Person{基督}亲自向那些在\Person{诺厄}日子不信而被囚禁在监狱中的人宣讲过，好像我们必须认为这个解释与基督那时尚未来临的事实不符。因为尽管祂还没有在肉身内来临，如同祂后来\BibleQuote{在地上出现，与世人往来}那样来临\BibleRef{巴 3:37}，然而从人类之初，祂确实常常来到这世上，或是为责备恶人，如\Person{加音}，以及在那之前，\Person{亚当}和他的妻子，当他们犯罪时，或是为安慰善人，或是为劝告双方，好使一些人相信而得救，另一些人拒绝相信而受审判——那时祂来临不是在肉身内，而是在精神内，借着合宜的自我显现，以祂所喜悦的方式和对象说话。至于\Quote{祂在精神内来临}这个说法，当然，作为\OriginalTerm{圣子}{Son of God}的祂，在其天主性的本质中就是神，因为那不是形体的；但圣子何时离了圣神或圣父行事呢？因为\OriginalTerm{天主圣三}{Trinity}的一切工程都是不可分的。\par

18. 在我看来，那些正在思考的圣经的话，本身足以向细心留意的人充分说明这道理：\BibleQuote{因为基督也曾一次为罪而死，且是义人代替不义的人，为将我们领到天主面前；就肉身说，他固然被处死了；但就神魂说，他却复活了；他藉这神魂，曾去给那些在狱中的灵魂宣讲过；他们从前在诺厄建造方舟的时日，天主耐心期待时，原是不信的人}。现在，我设想你们已细心留意到这些词的顺序：\BibleQuote{基督就肉身说，固然被处死了，但就神魂说，却复活了}；他藉这圣神，曾去给那些在诺厄的日子里一度拒绝相信他话的灵魂宣讲过；因为在他以肉身来为我们而死之前——这他只做了一次——他常藉圣神，按自己的意愿，以神视指教他愿意的人，确实藉着同一圣神，就是他在受难时肉身被处死后、使他复活的那圣神，来到他们那里。那么，他\BibleQuote{就神魂说却复活了}是什么意思呢？无非是说，那唯一经历死亡的肉身，因赋予生命的圣神而从死者中复活了。\par

% structure: p0027 source=h2 target=subsection reason=relative-heading-rank; base=h2

\subsection{第七章}

19. 因为谁敢说耶稣在祂的灵魂，\Emph{即}属于祂为人的神魂中，被处死了呢？因为灵魂能经历的死亡只是罪，而祂在为我们肉身被处死时，是完全无罪无辜的。如果所有人的灵魂都源于天主将生气吹入第一人的那一个灵魂，而藉着他\BibleQuote{罪恶进入了世界，死亡藉罪恶进入了世界，这样死亡就殃及了众人，} \BibleRef{罗 5:12}那么要么基督的灵魂不与其他灵魂同源，因为祂丝毫没有原罪或本罪，因此死亡在祂身上没有可以理应收到的惩罚，因为祂替我们偿还了并非由祂自己而负的债；这世界的首领——即掌握死亡权势的——在祂身上找不到任何可指摘之处（\BibleRef{若 14:30}）——而且设想那为第一人创造灵魂的，为自己创造一个灵魂，也并非不合理；要么基督的灵魂若源于亚当的灵魂祂在取这灵魂归于自身时，洁净了它，以致当祂来到这世界，由童贞女所生时，完全没有任何实际的或遗传的罪。然而，如果人的灵魂并非源于那一个灵魂，而原罪只是通过肉身由亚当遗传，那么天主子为自己创造了一个灵魂，正如祂为其他所有人创造灵魂一样，但祂将这灵魂与\BibleQuote{罪恶肉身的形状。} \BibleRef{罗 8:3} 相结合，而非与罪恶的肉身结合。因为祂确实从童贞女取了肉身的真实本质，然而不是\BibleQuote{罪恶的肉身}，因为这肉身既非因肉欲而生，亦非因肉欲而成孕，而是有死的，且在生命各阶段可变化的，与罪恶的肉身在各个方面相似，只是没有罪。\par

20. 因此，无论关于灵魂起源的真实理论是什么——对此，我认为我目前不宜贸然发表任何意见，除非是彻底拒绝那种断言每个灵魂被推入它所居住的身体如同推入监狱，在那里为自己某些我无从知晓的前世行为作补赎的理论——关于基督的灵魂，我们可以确定的是，不仅按照所有灵魂的本性是不死不灭的，而且它既未因罪而被处死，也未因惩罚而被定罪，这是灵魂可能经历死亡仅有的两种方式；因此，就不能说基督在灵魂方面\BibleQuote{就神魂说却复活了}。因为祂复活是在于祂被处死的那方面；所以这是就祂的肉身说的，因为祂的肉身在灵魂回来时重获生命，正如它在灵魂离去时死亡一样。因此，祂被说成\BibleQuote{就肉身说被处死了}，因为祂只在肉身中经历死亡，但\BibleQuote{就神魂说却复活了}，因为藉着那祂惯于来临并向祂所愿意的人宣讲的圣神的运作，祂用来临于人的那同一肉身复活了，并从坟墓中起来。\par

21. 所以，现在转而论到稍后关于不信者的经文：\BibleQuote{他们要向那准备审判活人死人的主交账}，我们不必把这里的\BibleQuote{死人}理解为已经离开身体的人。因为宗徒可能用\BibleQuote{死人}一词指不信者，即灵魂上死了的人，正如论到那些人说：\BibleQuote{任凭死人去埋葬他们的死人} \BibleRef{玛 8:22}，而用\BibleQuote{活人}指那些信祂的人，他们不会白白听到这呼唤：\BibleQuote{你这睡眠的，醒来，从死者中起来，基督必要光照你} \BibleRef{弗 5:14}；论到这些人，主也说：\BibleQuote{时候将到，且现在就是，死者要听见天主子的声音，听见的人，必得生存。} \BibleRef{5:25} 按同样的解释原则，也没有任何理由迫使我们把紧接着的伯多禄的话——\BibleQuote{为此，福音也曾给死者宣讲了，为使他们虽然在肉身方面像人一样受审判，但靠天主在神魂方面生活} \BibleRef{伯4:6}——理解为描述在地狱发生的事。\BibleQuote{福音曾宣讲}在今生中给\BibleQuote{死，即不信的恶人，为}们相信后，\BibleQuote{在肉身像人一样受审判}——即藉各种苦难和肉身本身的死亡；为此，同一宗徒在另一处说：\BibleQuote{时候已经来到，审判必从天主的家开始} \BibleRef{伯前 4:17}—\BibleQuote{但靠天主在圣神里面生活}，因为就在那同一神魂中，当他们被囚禁在不信和邪恶的死亡中时，他们原是死的。\par

22. 如果这关于\Person{伯多禄}的话的解释冒犯了任何人，或者即便没有冒犯，至少未能使任何人满意，那么让他尝试在假定它们指的是地狱的前提下解释它们；如果他能够解决我上面提到的那些困难，消除它们带来的困惑，就让他把他的解释告诉我；这样，这些话或许原本就可能意在从两方面去理解，但我所提出的观点并不因此就被证明为错误。\par

我托执事\Person{阿塞鲁斯}带去一封信，我想你已经收到了，信中我尽可能回答了你先前提出的问题，只有一个问题除外：即关于藉肉身感官得到的天主的异象，对于这个问题，必须尝试一部更长的论著。在你最近一封便函（即我现在回复的这封）中，你提出了两个问题，分别关于宗徒\Person{伯多禄}的某些话，以及主的灵魂；这两个问题我都讨论过了——前者较详细，后者较简短。我恳求你不要吝于费心再寄另一份包含以下问题的信函副本给我：即天主的实体能否以局限于空间的形体被看见；因为这封信不知怎么在这里遗失了，尽管我找了很久，却仍未找到。\par

\SourceRecord{Translated by J.G. Cunningham.；Nicene and Post-Nicene Fathers, First Series；Vol. 1.；Edited by Philip Schaff.；Buffalo, NY: Christian Literature Publishing Co.,；1887.；<http://www.newadvent.org/fathers/1102164.htm>.}

% structure: p0001 source=h1 target=section reason=page-title

\section{\Person{热罗尼莫}致\Person{玛尔切利诺}与\Person{阿纳普西基亚}（主后410年）}

\Emph{致我真正虔诚的主人\Person{玛尔切利诺}与\Person{阿纳普西基亚}，你们是配受一切与其地位相称之爱的儿子，\Person{热罗尼莫}在基督内致候。}\par

% structure: p0003 source=h2 target=subsection reason=relative-heading-rank; base=h2

\subsection{第一章}

1. 我终于收到了你们从非洲寄来的联名信，并不后悔我当初的固执——尽管你们沉默，我仍坚持给你们写信——为能得着答复，不藉他人的传闻，而由你们亲口最可喜的陈述，得知你们身体健康。我未曾忘记你们提出的那个简短问题，或更恰当地说，那个极为重要的神学问题，即关于灵魂的起源：灵魂是从天降下的，如同哲学家\Person{毕达哥拉斯}、所有柏拉图主义者以及\Person{奥力振}所认为的吗？还是如斯多噶派、\Person{摩尼}及西班牙的普里西利安派所臆想的那样，是神性本质的一部分？抑或灵魂被保存在一个神圣的库房里，乃是从前积蓄在那里的，如同某些愚蠢受骗的教会人士所相信的？又或者如福音所载：\BibleQuote{我父到现在一直工作，我也工作}\BibleRef{若 5:17}，它们是天天由天主所造并送入肉身内的？还是如\Person{德尔图良}、\Person{阿波利那里}及大多数西方神学家所推测的，灵魂真的是藉繁衍而生的，即如同身体是身体的后代，灵魂也是灵魂的后代，并以类似的方式存在于较下等的动物中？我知道，我曾在一篇短文中反驳\Person{鲁菲诺}时发表了我对这个问题的看法，当时是为回应他致具有神圣记忆的罗马教会主教\Person{亚纳大削}的一篇论文。在那篇论文中，他试图以油滑而狡诈却又愚拙的表白欺哄头脑简单的读者，却使之鄙弃他自己的信德，或更确切地说，他的背信。我想，这些书卷在你们的神圣亲属\Person{奥切安}手中，因为它们是很久以前发表的，为回应\Person{鲁菲诺}众多著作中的诽谤。无论如何，你们在非洲有那位圣人、博学的主教\Person{奥斯定}，他能像俗话所说的，以\OriginalTerm{口头}{viva voce}教导你们这个问题，并向你们阐述他的意见——或者我应该说，是以他的言辞陈述的我的意见。\par

% structure: p0005 source=h2 target=subsection reason=relative-heading-rank; base=h2

\subsection{第二章}

2. 我心已久愿开始撰写\WorkTitle{厄则克耳书注释}，以兑现多次向好学的读者许下的诺言；但正当我着手口授所拟的释义时，由于帝国西部行省的沦陷，尤其是在野蛮人手中的\Place{罗马}城被劫掠，我的心情如此激荡，以至用一句常言道，我几乎连自己的名字都不知道了；我缄默了许久，深知这是一段流泪的时期。再者，当我在这一年中预备好\WorkTitle{厄则克耳书注释}的三卷书时，突然，被你们的\Person{维吉尔}称为\Quote{广泛游荡的巴尔凯人}，并被圣经在论及\Person{依市玛耳}时所说\BibleQuote{他必要住在众兄弟的对面}\BibleRef{创 16:12}的那些蛮族部落，以一股洪流般猛烈之势席卷了整个\Place{埃及}、\Place{巴勒斯坦}、\Place{腓尼基}及\Place{叙利亚}，其势如狂澜，我们仅是靠着基督的仁慈才极艰难地从他们手中逃脱。若如一位著名的演说家所言：\Quote{在刀兵碰撞声中法律沉默}，那么对于圣经研究，这话就更适用了，因为圣经研究需要众多书籍与静默，需要抄写员不间断的勤奋，尤其需要口授者享有安宁闲适！因此，我已给我的圣善女儿\Person{法比奥拉}寄去了两卷书，如果你们需要抄本，可以向她借阅。由于时间不足，我未能抄录其余部分；当你们读了这两卷，并如同见了门廊一般，便可轻易设想房屋本身的构造。但我信赖天主的仁慈，祂已经在上述著作那极为艰难的开端帮助了我，祂也将帮助我（完成）有关\Person{哥格}与\Person{玛哥格}战争的预言，这预言占据了先知书倒数第二部分，以及最后描述那座至圣而玄妙的圣殿之建筑、细节与比例的部分。\par

% structure: p0007 source=h2 target=subsection reason=relative-heading-rank; base=h2

\subsection{第三章}

3. 你们希望我提及的我们那位圣善兄弟\Person{奥切安}，是一位禀赋与品格如此优越、且深明主的法律的人，他或许无需我的任何请求就能教导你们，并能依照我们共同的能力限度，在所有圣经问题上传达我们已形成的意见。\par

愿我们全能的天主基督，保护你们，我真正虔诚的主人，平安顺利，直到寿高年迈！\par

\SourceRecord{Translated by J.G. Cunningham.；Nicene and Post-Nicene Fathers, First Series；Vol. 1.；Edited by Philip Schaff.；Buffalo, NY: Christian Literature Publishing Co.,；1887.；<http://www.newadvent.org/fathers/1102165.htm>.}

% structure: p0001 source=h1 target=section reason=page-title

\section{\Person{奥古斯丁}致\Person{热罗尼莫}：论灵魂的起源（主后415年）}

\WorkTitle{论人类灵魂的起源，致热罗尼莫}\par

% structure: p0003 source=h2 target=subsection reason=relative-heading-rank; base=h2

\subsection{第一章}

1. 我曾向我们的天主祈祷，就是召叫我们进入他的国度和光荣的那一位，\BibleRef{得前 2:12} 如今仍在祈祷，愿我写给你的，圣洁的弟兄\Person{热罗尼莫}，为求教我所不知之事而询问你的意见，能赖他的喜悦，于我们双方都有裨益。因为虽然我向你请教，是询问一位比我年长得多的人，然而我自己也日渐衰老；但我并不认为，在人生的任何阶段，学习我们所需要知道的事就会太迟，因为虽然老年人更适宜施教而非受教，但他们学习，总比对于应教导他人之事依旧蒙昧无知更为适宜。我向你保证，当我专心研究极其困难的问题时，不得不忍受诸多不便，其中最令我感到痛苦的，莫过于与你的爱德因相隔遥远而分离；这距离如此之大，使我几乎无法寄信给你，也几乎无法收到你的来信，即便偶尔能通音信，也不是隔几天或几个月，而是相隔数年；然而我所渴望的，是倘若可能，每日有你在我身边，如同一位我可以与之谈论任何话题的挚友。尽管如此，虽然我未能做到所有我想做之事，但我仍须竭尽所能。\par

2. 看，有一位虔诚的青年到我这里来，名叫\Person{奥罗修斯}，在天主教和平的联结中，他是我的弟兄，在年龄上，他如我的儿子，在尊荣上，他是我的同僚司铎——他悟性敏速，言语便给，热心如焚，渴望在主的家中成为有用的器皿，驳斥那些虚假而有害的教义，\Place{西班牙}人的灵魂因这些教义所受的创伤，比他们身体所受蛮族刀剑的创伤更为惨重。因为他从\Place{西班牙}遥远的西部海岸，热切匆忙地赶到我们这里，促使他如此行的，是传闻说，凡他想要明白的事，都可以向我学习。他的到来并非全然徒劳。首先，他学会了不要相信传闻所说关于我的一切；其次，我已经将我所能教导的，全都教给了他，至于那些我不能教导他的，我已告诉他可以向谁去学习，并劝勉他到你那里去。他欣然接受了我的这项建议，更恰当地说是我的训令，并有意遵照执行，我便请他由你那里回来时，在返回\Place{西班牙}的途中，再来看望我们。他如此应许了，我便相信，主许给了我一个机会，为那些我愿通过你而学习的事写信给你。因为我曾寻找一位可以差遣到你那里去的人，而要想找到一位既可靠又乐于承担任务，并且富有旅行经验的人，实属不易。因此，当我认识这位青年时，我毫不怀疑，他正是我向主所求的那等人。\par

% structure: p0006 source=h2 target=subsection reason=relative-heading-rank; base=h2

\subsection{第二章}

3. 因此，请容许我向你提出一个问题，恳求你不要拒绝为我剖析并与我讨论。许多人因有关灵魂的问题而困惑，我承认我自己也是其中之一。在这封信中，我将首先明确陈述关于灵魂我所绝对确信之事，然后再提出我仍渴望获得解释之事。\par

人的灵魂，在某种意义上，本身是不死不灭的。但它并非绝对不死不灭，如同天主那样，经上论到他写道：\BibleQuote{惟有他是不死不灭的，}\BibleRef{弟前 6:16} 因为圣经也提及灵魂会遭受的死亡——如所说的：\BibleQuote{让死人去埋葬他们的死人！}\BibleRef{玛 8:22} 然而，当灵魂与天主的生命隔绝时，它的死并非完全丧失其本性中的生命，故而可发现，它由于某种原因之所以能死，却又不无理由地被同时称为不死不灭的。\par

灵魂并非天主的一部分。因为若是，它便绝对不变且不朽，如此它既不能堕落变得更坏，也不能精进变得更好；它也不能在自身中开始拥有先前所无之物，也不能停止拥有其自身经验范围内所曾有之物。然而事实真相是何等不同，无需外来证据，凡返观内照的人都会承认。再者，那些主张灵魂是天主一部分的人，其辩称我们所见于最恶之人身上的败坏与卑劣，以及见于所有人身上的软弱与瑕疵，并非来自灵魂本身，而是来自身体，这是徒劳无益的；因为，如果灵魂确实不变，那么衰弱源于何处又有何关系？无论出于何种原因，它都不可能变得衰弱。因为真正不变且不朽者，不会受任何外来影响而改变或朽坏，否则，寓言中阿喀琉斯之肉所具之不可伤害，便非其独有，而是人人皆具，只要不遭意外。凡能以任何方式、因任何原因、在任何部分上被改变者，便非本性不变；但若以为天主并非真正且至极不变，便是亵渎：因此，灵魂并非天主的一部分。\par

4. 灵魂是\OriginalTerm{非物质的}{immaterial}，这是我自认完全信服的事实，尽管悟性迟钝的人很难被说服相信这一点。然而，为避免因用词不当而给他人带来不必要的麻烦，或给自己招致无谓的争论，我要说，既然事实本身毋庸置疑，就无需在单纯的术语上争执。如果\OriginalTerm{物质}{matter}一词用来指在任何形式下都有独立存在的一切，无论称之为\OriginalTerm{本质}{essence}、\OriginalTerm{实体}{substance}，还是其他名称，那么灵魂就是物质的。再者，如果你选择只将非物质的这一修饰语用于那绝对不变、全然处处临在的本性，那么灵魂就是物质的，因为它根本不具有这类特质。但如果\OriginalTerm{物质}{matter}一词仅仅用来指那些无论静止或运动，都具有长度、宽度和高度的东西，以至于其较大部分占据较大空间，较小部分占据较小空间，且每一部分都小于整体，那么灵魂就不是物质的。因为灵魂渗透它所激活的整个身体，不是通过局部分布，而是凭借某种\OriginalTerm{生命力的影响}{vital influence}，在每一时刻整体地临在于身体的所有部分，并非在较小部分中较少，在较大部分中较多，而是此处活力较强，彼处活力较弱；它整体地临在于整个身体，也临在于身体的每一部分。因为即便心灵在身体的一个部分所感知的，也绝非仅由部分心灵感知，而是由整个心灵感知；因为当活生生的肉体任一部分被尖细的工具触碰时，尽管受触之处不仅不是整个身体，甚至在其表面也几乎难以辨认，整个心灵仍不会忽略这一接触，然而这种接触感并不遍布全身，而只发生在触点所在的那一点。那么，为何仅发生于身体一部分的事立刻为整个心灵所知，除非整个心灵临在于那部分，同时又没有撇下身体所有其他部分，为要整体地临在于这一处？因为身体中未发生此类接触的所有其他部分，仍因灵魂的临在而活着。如果类似的接触发生在其他部分，且两处同时发生接触，整个心灵也会在同一时刻同样地感知两者。这种心灵在同一时刻临在于身体所有部分，以致在身体的每一部分，整个心灵都同时临在，若心灵像我们所见物质在空间中分布那样，以较小部分占据较小空间，较大部分占据较大空间，那么上述临在便不可能。因此，若心灵要被称为物质的，它也不像土、水、气、以太那样是物质的。因为一切由这些元素构成的事物，在较大空间中便较大，在较小空间中便较小，它们中没有一个是整体地临在于自身的任何部分，物质实体的尺寸与所占空间的尺寸相应。由此可知，灵魂无论被称为物质或非物质，都具有某种自身的本性，由一种超越此世元素的更高级\OriginalTerm{实体}{substance}所造——这种实体无法通过身体感官所知觉的物质形象的任何表象来真实构想，而是由理智领悟，并经由其生命活力向我们的意识显明。我陈述这些，并非意在教导你们已经知道的事，而是为要明确宣示我关于灵魂所坚信不疑的立场，免得当我提出我所寻求解答的问题时，有人以为我对于我们从科学或启示所学得有关灵魂的教义，竟一无所持。\par

5. 此外，我完全信服：灵魂陷于罪恶，不是由于天主的过错，也不是由于天主本性或灵魂自身本性中的任何必然，而是由于灵魂自己的\OriginalTerm{自由意志}{free will}；灵魂能脱离这必死的肉身，既不是靠自身意志的力量——仿佛这力足以成就此事，也不是靠肉身的死亡本身，而是唯独靠天主的\OriginalTerm{恩宠}{grace}，通过我们的主耶稣基督；\BibleRef{罗 7:24} 在人类大家庭中，没有一个人的灵魂的得救，不需要那在天主与人之间的\OriginalTerm{唯一中保}{one Mediator}，即人基督耶稣。再者，任何年龄的灵魂，若未领受中保的\OriginalTerm{恩宠}{grace}和这恩宠的\OriginalTerm{圣事}{sacrament}而离开此生，都将去受未来的惩罚，并在公审判时重新得回自己的身体，共享惩罚。但灵魂若在出自亚当的自然生育之后，在基督内\OriginalTerm{重生}{regenerated}，便属于祂的团体，不仅肉身死后将享安息，还要重新得回自己的身体共享光荣。这些是我关于灵魂所持的最坚定的真理。\par

% structure: p0012 source=h2 target=subsection reason=relative-heading-rank; base=h2

\subsection{第三章}

6. 因此，现在请允许我向你提出我切望得到解答的问题，不要拒绝我；如此，那为我们的缘故屈尊被拒绝的祂，才不会拒绝你！\par

我问：即使是那刚出生便被死亡夺走的婴儿的灵魂，又能从哪里染上罪债？除非基督的恩宠借着连婴儿也领受的圣洗圣事前来救援，这罪债就会使之陷入永罚。我知道你不是那些近来开始散布某些新颖荒谬见解的人，他们声称婴儿领受圣洗并不能清除从亚当而来的罪债。如果我知道你持这种观点——或者更确切地说，如果我不知道你拒绝这种观点，我肯定既不会向你提出这个问题，也根本不会认为应当向你提出。然而，关于这个问题，我们所坚持的，是与不可动摇的天主教信仰一致的见解；你自己在驳斥\Person{约维尼安}的荒谬言论时，引用\BibleBook{约伯}{传}中的那句话：\BibleQuote{在你眼中，没有人是洁净的，就连在世仅有一天的婴儿也不例外，}并且补充说：\Quote{因为我们因亚当背命的相似而负有罪债，}——你的那本关于\WorkTitle{约纳书注释}中，更突出且清晰地表明了这一主张，你在那里确认，\Place{尼尼微}的幼儿之所以理所当然地与民众一同禁食，纯粹是因为他们的原罪——因此，向你提出这个问题并非不妥：灵魂究竟从哪里沾染了罪债，即使在那样的年纪，也必须借着基督恩宠的圣事获得解脱？\par

7. 几年前，我撰写了一些关于\WorkTitle{论自由意志}的书，这些书流传到了许多人手中，如今也为极多读者所拥有。在提及关于灵魂如何降生的四种见解之后——（1）所有其他灵魂都源自赋予第一人的那个灵魂；（2）每个个体都有一个新的灵魂被创造；（3）灵魂已存在于某处，并藉着天主的作为被送入身体；或（4）灵魂自行溜入身体——我当时认为，必须这样处理它们，即无论其中哪一种可能为真，都不应妨碍我力争的目标，那时我正全力驳斥那些企图将罪责归于具有自身恶之本原的那类人，即摩尼教徒；因为那时我尚未听说过普里西利安派，他们的亵渎之言与这些相差不远。至于第五种见解，即灵魂是天主的一部分——你致\Person{马切利努斯}（一个虔诚纪念并因基督的恩宠而为我们极亲爱的）的信中，为不遗漏任何见解，将此见解与其他一同提及，他当时曾就这一问题向你请教——我未将其列入其他之中，原因有二：首先，在考察此见解时，我们讨论的不是灵魂的降生，而是它的本性；其次，因为这是我正在驳斥的那些人所持的观点，而我论证的主要目的是要证明：造物主那无可指摘且不可侵犯的本性，与受造物的过失和瑕疵毫无关系，而他们却主张，善良的天主的实体本身，就其被掳获、被败坏、被压迫、并因恶的实体而不可避免地犯罪而言，是被恶所左右的，他们将统治权与掌权者归于恶。因此，且将这异端错误搁置不论，我渴望知道，我们应从其他四种见解中选哪一种。因为，无论哪一种能公正地要求我们优先选择，我们都绝不可攻击这一我们所确信的信德道理，即：每个灵魂，即使是婴儿的灵魂，都必须从罪债的束缚中获得解脱，并且除了通过耶稣基督并他钉十字架以外，别无解脱。\par

% structure: p0016 source=h2 target=subsection reason=relative-heading-rank; base=h2

\subsection{第四章}

为避免冗长，因此，且让我谈一下你据我所了解所持的看法，即天主至今仍在每个人出生时为其分别创造灵魂。为了应对从\Quote{天主在第六日完成了全部创造工程，第七日安息了}这一事实可能引出的反对意见，你引用福音里的话作证：\BibleQuote{我父到现在还在工作，我也工作。}\BibleRef{若 5:17}。这写在你致\Person{马切利努斯}的信中，信中你还很亲切地屈尊提及我的名字，说他在\Place{非洲}有我在这里，可以更容易地为他解释你所持的见解。但是，假如我当时能做到这一点，他就不会向远在他方的你求教了，尽管他也许并非是从\Place{非洲}写信给你。因为我不知道他何时写的信；我只知道他很清楚我对于这一问题迟迟不肯接受任何明确的观点，因此他宁愿写信给你，而不咨询我。然而，即使他来咨询我，我也会鼓励他写信给你，并因我们大家可能获得的益处而向你表示感谢，只是你宁愿寄来一封简短的信函，而不是充分的答复；我猜想，你这是为了避免在我——你估计我完全熟悉他所请教的问题——就近的地方作无谓的劳神。看，我愿意将你所持的见解也变成我的见解；但我明确告诉你，我至今尚未接受它。\par

你差遣了一些学者到我这里来，你希望我将自己尚不知道的学问传授给他们。因此，请教我，我该教导他们什么；因为许多人极力催促我就此课题作教师，而我向他们承认，我对这事，正如对其他许多事一样，是无知的；或许，他们虽然在我面前保持恭敬的态度，私下里却说：\BibleQuote{你是\Place{以色列}的师傅，连这些事也不知道吗？}\BibleRef{若 3:10} 主曾以此斥责那些喜欢被人称为\OriginalTerm{辣彼}{Rabbi}的人；这也是他夜间来到真导师跟前的原因，因为这位已习惯了教导别人，大概羞于充当学生。然而，就我本人而言，我更喜欢接受别人的教导，远胜于自己被人聆听为师傅。因为我记得祂对祂特别拣选的人所说的话：\BibleQuote{你们不要被称为辣彼，因为你们的师傅只有一位，就是\Person{基督}。}\BibleRef{玛 23:8} 教导\Person{梅瑟}的，不是别人，正是\Person{耶特洛}；\BibleRef{出 18:14} 教导\Person{科尔乃略}的，是先前的宗徒\Person{伯多禄}；\BibleRef{宗 10:25} 而教导\Person{伯多禄}的，是后来的宗徒\Person{保禄}；\BibleRef{迦 2:11} 因为无论谁宣讲真理，都是靠着那位真理之主的恩赐而宣讲。倘若我们至今仍对这些事无知，即便在祈祷、阅读、思考和推论之后仍未能发现，其原因或许正是：不仅要充分证明我们教导无知者时的爱，更要证明我们接受博学者教导时的谦逊？\par

所以，我恳求你教导我，我该将什么教训传给别人；教导我，我该持守什么作为自己的主张；并请告诉我：如果灵魂是每天为每一个人出生时分别受造的，那么，就婴孩来说，这些灵魂犯了什么罪，以致他们需要在\Person{基督}的圣洗中获得罪的赦免，因为他们是在\Person{亚当}内犯了罪，而有罪的肉身正源于\Person{亚当}？或者，如果他们并未犯罪，那么，由于他们与源自他人的、有死的肢体结合，便被置于那人的罪愆之下，以致若非教会拯救他们，丧亡便临到他们，而他们自己却无力获得圣洗的恩宠以便被拯救，这又如何能与造物主的公义相容呢？如此多的灵魂——成千上万——在婴孩死亡时未得\Person{基督}圣事的福泽便离开此世，如果这些灵魂是全新受造的，并非由于自己先前的任何罪过，而是由于造物主的旨意，各自与不同的身体结合，为赋予它们生命，而这些灵魂正是造物主为了这个目的而创造并赐予的，造物主当然知道他们每一个都注定，不是由于自身的过失，而是在没有领受\Person{基督}圣洗的情况下离开肉身；那么，对这些灵魂的定罪，公义又在哪里呢？因此，我们既不能说天主强迫他们成为罪人，也不能说祂惩罚无辜的灵魂；另一方面，我们也不被允许否认，那些没有\Person{基督}圣洗就脱离肉身的灵魂，甚至是小孩的灵魂，其命运只能是丧亡——既然如此，我恳请你告诉我，在哪里能找到支持以下意见的证据：灵魂并非全都出自那第一个人的一个灵魂，而是为每一个人分别受造，正如\Person{亚当}的灵魂是为他而造的一样？\par

% structure: p0020 source=h2 target=subsection reason=relative-heading-rank; base=h2

\subsection{第五章}

至于有人提出的其他一些反对这见解的论据，我想我能轻易答辩。例如，有些人认为他们提出了一个决定性的论据：他们问道，如果天主仍在创造新灵魂，那么他怎能在第六天完成了他的一切工程，并在第七天安息呢？\BibleRef{创 2:2} 如果我们以（你在致\Person{玛尔凯利努斯}的信中已引用的）\BibleBook{若望}{福音}的引文——\BibleQuote{我父直到如今工作}——来回应他们，他们便会说，祂\Quote{工作}是指维持祂已创造的本性，而非创造新的本性；否则，这说法就会与\BibleBook{创}{世纪}中圣经之言相矛盾，因为那里最明确地声明，天主完成了他的一切工程。再者，圣经中说他安息的话，无疑应理解为：他停止了创造新的受造物，而不是不去掌管他已创造的万物；因为在那时，他创造了先前并不存在的事物，并从造物的工作中安息，因为他已经完成了他所预见在存在之前必须创造的一切受造物，所以从那以后，他不再创造和制造先前并不存在的事物，而是从已有之物中制造和形成他所做的一切。因此，所谓\BibleQuote{他从自己的工程安息了}和\BibleQuote{他直到如今工作}，这两种说法，都是真实的，因为\BibleBook{}{福音}不可能与\BibleBook{创}{世纪}相矛盾。\par

12. 然而，当有些人提出这些事，作为反对天主如今如同创造第一个人的灵魂那样创造新灵魂这一意见的决定性论据，并且主张祂或是由那个在祂停止创造之前就已存在的唯一灵魂形成它们，或是如今从某个祂当时所创造的灵魂的储库或仓库中将它们注入肉身时，我们很容易驳斥他们的论证，回答道：即使在最初的六天里，天主也是从祂已经创造的本性中形成了许多东西，比如，飞鸟和鱼类由水中形成，树木、青草和走兽由大地形成，但不可否认的是，祂当时确实在创造先前尚未存在的东西。因为那时先前并没有鸟，没有鱼，没有树，没有走兽，并且显然易见，祂停止了创造那些先前不存在、那时才被造的东西，也就是说，在此意义上，之后祂就不再造任何先前不存在的东西了。但是，如果我们在摈弃了所有人的那些意见之后——无论是认为天主从某种不可思议的储库中将灵魂送入人内，或认为祂使灵魂像露珠一样从祂自身滴下，作为祂本体的微粒，或认为祂从第一个人的那个灵魂中引出它们，或认为祂因为灵魂在先前的生存状态下犯下的罪恶而将它们束缚在肉身的桎梏中——如果，我说，摈弃了这些之后，我们肯定说，祂为每一个出生的人单独地创造一个崭新的灵魂，那么，我们在此并不是肯定说祂造了任何祂先前未曾造过的东西。因为祂早已在第六日按自己的肖像造了人；而祂的这项工程无疑应被理解为人所具的\OriginalTerm{理性灵魂}{rational soul}。祂现今仍在进行这同一工程，不是创造先前不存在之物，而是增多那已然存在之物。因此，一方面，祂停止了创造先前不存在之物，这是真的；另一方面，祂继续工作，不仅掌管祂所造的，而且制造（不是先前不存在之物，而是）祂早已受造的受造物的更多数量，这也是真的。因此，或以此解说，或以任何似乎更好的解说，我们就能避开那些人所提出的异议，他们试图将天主停止工程的这一事实，作为一个决定性的论据，来反对我们相信新灵魂仍在每日受造，并非出自第一个灵魂，而是按照第一个灵魂受造的同一方式受造。\par

13. 再者，对于问题中所述的另一项异议：\Quote{为何祂为那些祂预知将早逝的人创造灵魂？}我们可以回答说，孩童的死亡乃是为谴责或惩罚父母的罪过。此外，我们完全可以恰当地将这类事情留给万有之主安排，因为我们知道，祂为事件在时间中的演替，从而也为包含其中的生灵之生死，指定了一条在一切部分的安排上都极致美丽、完美无缺的路径；而我们却无法领悟那些事，如果能够领悟，我们就会因一种无以言喻的宁静喜乐而心满意足。因为先知并非徒然地蒙受天主的启示，论及天主说，\Quote{祂以和谐的节拍引领时间的进程。}为此缘故，音乐，即正确和声的学问或能力，也由天主的仁慈赐予具有\OriginalTerm{理性灵魂}{reasonable souls}的世人，为使他们牢记这伟大的真理。因为，如果人在为一首特定的旋律谱写歌曲时，懂得如何分配每个字所允许的时间长度，使歌曲在极美地适应旋律不断变化的音符中流利地进展，那么，天主的智慧，既当被尊为无限超越人的技艺，岂不更能万无一失地安排，使赋予有生有灭之本性的一切时段——这些时段就像是时间进程里连续纪元的字与音节——在我们可称之为此世变迁的崇高圣咏中，不使其任何一段时间短于或长于祂那预知并预定的和谐所要求的呢！因为，既然关于每一树的叶子、我们头上的头发数量，我尚且可以如此说，那么关于人之生死，我岂不更可如此说，因为每个人在世的生命所延续的时间，既不长于也不短于天主照祂治理宇宙的计划所知道与之相和谐的时间。\par

14. 至于那种断言一切在时间内开始存在的东西都不能不死，因为凡受生的都会死，一切生长的都会因年龄而衰败，以及他们认为由此必然得出的意见，即人的灵魂之不死必定归因于其创造是在时间开始之前——这一切并不扰乱我的信德；因为，暂且不提其他能最终驳倒此断言的事例，我只需指出\Person{基督}的身体，现在\BibleQuote{不再死；死亡不再统治他了。}\BibleRef{罗 6:9}\par

15. 此外，关于你在你的\WorkTitle{反卢非诺论}一书中所说的，有人提出反对灵魂为每一单独个体在出生时受造的见解，他们反对的理由是：天主将灵魂赐给奸淫者的后代，似乎于天主不配；他们并试图以此为基础，构建一种理论，认为灵魂可能被\Emph{仿佛}禁锢在这样的身体中，以受前世生活所行之恶事的苦。这个反对意见并不困扰我，因为当我思考时，许多可以反驳它的答案涌现于心。你本人所给出的回答，即：在被偷窃的麦子一事上，麦粒本身并无过错，过犯只在于偷窃者；而大地并没有义务因播种者可能以不诚实所玷污的手撒下种子，就拒绝滋育那种子——这是一个最恰当的比喻。但在我读到这回答之前，当我普遍地观察到一个事实，即天主甚至从我们的恶行与罪恶中，也引出许多善事时，我就感到，从奸淫者的后代所引出的反对意见，并不构成严重的困难。任何活的受造物的创造，都迫使每一个以虔诚和智慧来观之的人，将无法以言词表达的赞美归于造物主；倘若任何受造物的创造都能唤起这种赞美，那么人的创造更当唤起何等的赞美！因此，如果寻找任何创造行为的理由，那么没有比\Quote{天主的每一受造物都是美好的}更简短或更好的回答了。而且，[这样的行为不但不于天主不配，]还有什么比这更与他相配呢？——他既为善者，就创造了那些除他以外，无一能造的美好事物？\par

% structure: p0026 source=h2 target=subsection reason=relative-heading-rank; base=h2

\subsection{第六章}

16. 这些事，以及我能提出的其他论点，我习惯于竭尽所能地陈述，以反驳那些试图用这些反对意见推翻这一见解的人；这见解是：灵魂是为每一个体单独受造的，正如第一个人的灵魂是为他受造的一样。\par

但论到婴儿所受的惩罚之苦时，请相信我，我陷入了极大的困难，完全茫然，不知如何回答来解决它们；我在此不仅论及来生中的惩罚，即他们若未领受\OriginalTerm{基督恩宠的圣事}{sacrament of Christian grace}而离世，便必然被拖入的那种\OriginalTerm{丧亡}{perdition}；也论及我们忧伤地看见他们在今生、在我们眼前所遭受的那些痛苦，这些痛苦如此多样，如果我要试图枚举它们，我恐怕会耗尽时间而非找不到实例。他们易罹消耗性的疾病、剧痛、饥渴的煎熬、肢体的软弱、身体感官的缺失，以及\OriginalTerm{不洁之灵}{unclean spirits}的恼人攻击。显然，我们有责任指明，婴儿并无自身的恶行作为起因，却承受这一切，这如何能符合正义。因为若说这些事发生而天主不知晓，或说他不能阻止那些导致这些事的人，或说他行这些事或允许这些事发生是不公义——这些说法都是不虔之词。我们可以有把握地说，天主将无理性的动物赐给具有更高本性的受造物，即使这些受造物是邪恶的，正如我们在福音中极清楚地看到，\Place{革辣撒}人的猪曾应魔鬼军团的要求而被交给它们；但我们能有把握地以此论及人吗？断然不能。人的确是一种动物，却是一种有理性、虽然会死的动物。在他肢体中寓居着一个理性的灵魂，这灵魂在如此沉重的痛苦中正在承受惩罚。但天主是美善的，天主是公义的，天主是全能的——只有疯子才会怀疑他是这样；因此，应让婴孩所遭受的巨大痛苦，用某种符合公义的理由来解释。当成年人遭遇这样的试炼时，我们通常的确会说，要么他们的价值正在受考验，如\Person{约伯}的例子；要么他们的邪恶正在受惩罚，如\Person{黑落德}的例子；而且从一些天主乐意完全澄清的例子中，人们得以推测其他更隐晦之事物的性质；但这都是针对成年人的。那么请告诉我，论到婴孩，我们必须作何回答呢？难道说，他们虽然承受如此重大的惩罚，身上却没有应受惩罚的罪过吗？因为，显然在他们那个年龄，并没有任何需要通过考验来证明的正义。\par

17. 至于[不同灵魂间才能的差异，尤其是某些灵魂完全缺乏理性，这一困难在个体出生时灵魂一一单独受造的理论上所带来的困境，]我还有什么可说的呢？这差异在婴儿早期并不明显，但从生命之初持续发展，在孩童身上便显而易见：其中有的如此迟钝、记忆缺损，连字母都无法学会；而有些（通常被称为痴呆者）则如此愚拙，几乎无异于田野间的走兽。对于这一点，也许有人会回答说是身体导致了这些缺陷。但我们希望看到被辩护的见解，并不要求我们肯定说，灵魂为自己选择了使其受损的身体，由于选择错误而犯了错；或者说，灵魂由于出生必然被迫进入某个身体，因在它之前已有众多灵魂占据了其他身体而受阻找不到别的，于是就像一个人在戏院里随便找一个空位坐下那样，灵魂之占有那个不完美的身体，并非出于其选择，而是由其处境所导引。我们当然不能这么说，也不该相信这类事。那么，请告诉我们，我们应当相信和说些什么，以便从这一困难中辩明每个身体都特创一个新灵魂的理论。\par

% structure: p0030 source=h2 target=subsection reason=relative-heading-rank; base=h2

\subsection{第七章}

18. 在我业已提及的\WorkTitle{论自由意志}诸书中，我并非论及不同灵魂在能力上的差异，而是至少论及了婴孩在此生所遭受的苦痛。我将在那里所表意见的性质，及何以它不足以满足我们眼下探究之故，现呈于你，并将我引用的第三卷中那段话誊录于此信。其文如下：——\Quote{关于那些因年幼稚嫩而无罪的幼童所经受的肉体之苦——倘若赋予他们生命的灵魂在他们降生于人类家庭之前并不存在——一种更为沉痛且近乎悲悯的怨诉常见于这样的言辞：\InnerQuote{他们行了什么恶，竟要受这些苦呢？}仿佛一个人在可能行恶之前，就能拥有可称功绩的无辜！再者，倘若天主藉着父母所爱的子女受苦或死亡，以某种方式惩戒他们，从而纠正他们，那么，这些事为何不能发生呢？因为，这些苦难过后，对经历者而言，将如同从未发生过一般；而那些因苦难而被惩戒的人，或会因暂时的痛苦之杖而悔改，选择更善的生活方式；或者，尽管有此生的忧愁，若他们仍拒绝将愿望转向永生，则在将来的审判中，受罚时便无可推诿。而且，有谁知道天主会藉着这些幼童的苦难，赐给他们什么呢？这些苦难使父母的顽梗之心得以软化，或使他们的信德得以锻炼，或使他们的怜悯之心得以验证。有谁知道天主在祂隐秘的审判中，为这些幼童预备了什么美善的酬报呢？因为他们虽未行善，却未犯下自己的过犯，而承受了这些试炼。教会将那在黑落德搜求我主耶稣基督欲杀害祂时，被屠杀的幼童尊为殉道者之列，实非无故。}\par

19. 这些话是我当时写作的，那时我正试图捍卫如今讨论的观点。因为，正如我稍前提及，我努力要证明：无论关于灵魂的\OriginalTerm{降生}{incarnation}这四种观点中哪一种是正确的，造物主的本体都完全无可指责，且绝不染指我们的罪过。因此，无论这些观点中的哪一种会在真理中得以确立或被推翻，都无碍于我那时所从事作品的宗旨；因为无论哪一种观点能在更彻底的讨论之后胜过其余，这都不会引起我任何不安，因为我那时的目标是要证明，即使承认所有这些观点，我所持的学说仍屹立不摇。但我现在的目标是，尽可能藉着健全的推理，从这四种观点中选出一种；因此，当我仔细考量那书中引述的言辞时，我并不认为那些为辩护我们现在所讨论的观点而提出的论据是有效且确定的。\par

20. 我辩护的所谓主要支柱就在于这句话：\Quote{而且，有谁知道天主会藉着这些幼童的苦难，赐给他们什么呢？这些苦难使父母的顽梗之心得以软化，或使他们的信德得以锻炼，或使他们的怜悯之心得以验证。有谁知道天主在祂隐秘的审判中，为这些幼童预备了什么美善的酬报呢？}我看到，这对于那些以任何方式为基督并为真正宗教之故而受苦（虽然他们自己并不知道）的婴孩，以及那些已经领受了\OriginalTerm{基督徒团契圣事}{sacrament of Christian fellowship}的婴孩而言，并非毫无依据的推测；因为，若不与那位唯一中保结合，他们就不能脱离永罚，从而也就无法被置于一种甚至有可能为他们在此世种种苦难中所受之磨难获得补偿的地位。然而，这问题若要完全解决，答案就必须也涵盖那些未曾领受\OriginalTerm{基督徒团契圣事}{sacrament of Christian fellowship}、在忍受了最剧烈的苦难后于婴儿期死去的人，那么，他们的情况又能设想有何补偿呢？因为，除了在此生所受的一切痛苦之外，等待着他们的还有来世的丧亡。关于婴儿的圣洗，我在同一本书中并非详尽，但就当时所从事的工作而言，已尽可能给予了回答，证明了圣洗对婴儿有益，即便他们不知其为何，且尚无自己的信德；但关于那些未及领受圣洗而离世的婴儿之丧亡一题，我当时没有认为有必要说什么，因为那本书所讨论的问题与我们现在的论题不同。\par

21. 然而，如果我们忽略不计那些短暂且一旦承受便不会重复的痛苦，我们当然也不能同样地忽略这样的事实，即\BibleQuote{死亡既因一人而来，死者的复活也因一人而来；就如在亚当内，众人都死了，照样在基督内，众人都要复活。}\BibleRef{格前 15:21}因为根据这神圣而清晰的宣言，足够清楚的是，没有人不是通过亚当走向死亡，也没有人不是通过基督走向永生。这就是这句话中\Emph{所有}一词在两处的力量；正如所有人，通过他们的第一次诞生，即本性的诞生，属于亚当，同样，所有来到基督面前的人，无论他们是谁，都来到了第二次诞生，即属神的诞生。因此，\Emph{所有}一词在两个分句中都使用，因为正如所有死亡的人无一不是死于亚当内，所有将要复活的人也无一不是活于基督内。所以，凡告诉我们任何人能在死者的复活中活于基督以外的人，都应被视为共同信仰的瘟疫般敌人，应予憎恶。同样，凡说那些离开此生而未领受那\OriginalTerm{圣事}{sacrament}的孩童能在基督内复活的人，无疑在反对宗徒的宣言，并谴责普世教会，因为在教会中，惯常毫不迟延地急速为婴儿施行\OriginalTerm{圣洗}{baptism}，因为人们相信，作为不容置疑的真理，若非如此，他们便不能在基督内复活。如今，那不在基督内复活的人，必然仍处于\OriginalTerm{定罪}{condemnation}之下，关于这定罪，宗徒说：\BibleQuote{因一人的过犯，审判就临到众人，使他们被定罪。}\BibleRef{罗 5:18}婴儿生来就负有这过犯的罪债，这是整个教会所相信的。这也是您极其忠实地阐述的道理，既在您反对\Person{约维尼安}的论著中，也在您对《\WorkTitle{约纳}》的解释中，如上文所述，而且，如果我没记错的话，也在您其他我尚未读过或目前忘记的著作中。因此我问，这些未受洗婴儿被定罪的理由何在？因为，如果每个人的灵魂是在出生时个别新造的，一方面，我看不出他们在婴儿期能有什么罪，另一方面，我也不相信天主会定任何一个祂视为无罪的灵魂的罪。\par

% structure: p0035 source=h2 target=subsection reason=relative-heading-rank; base=h2

\subsection{第八章}

22. 难道我们要这样回答：在婴儿身上，只有身体是罪的原因；但对于每个身体，一个新灵魂被造，如果这灵魂依从天主的诫命而生活，藉着基督\OriginalTerm{恩宠}{grace}的助佑，当身体被制服并置于轭下时，得到不朽坏的赏报便可为这身体本身获得；并且，由于婴儿的灵魂还不能这样做，除非它领受基督的\OriginalTerm{圣事}{sacrament}，那原本不能由灵魂的圣德为身体获得的，便由此圣事的恩宠为它获得；但是，如果婴儿的灵魂未领受此\OriginalTerm{圣事}{sacrament}而去世，它自己将居于永生之中，因为它没有罪，不会与永生隔绝，然而它所居的身体却不复活于基督内，因为它在死前未领受这\OriginalTerm{圣事}{sacrament}？\par

23. 这种观点我从未在何处听说或读到过。然而，我确实听到并相信了那个引领我如此说的声明，即\BibleQuote{时候要到，那时凡在坟墓里的，都要听见他的声音，而出来；行过善的，复活进入生命，}——即那经上所说\BibleQuote{因一人死亡而来的，是死者的复活，}并在其中\BibleQuote{众人都要在基督内复活}的复活——\BibleQuote{而作过恶的，复活而受审判。}\BibleRef{若 5:29}现在，关于那些在能行善作恶之前，未受\OriginalTerm{圣洗}{baptism}而离世的婴儿，该怎么理解？这里没有提及他们。但是，如果这些婴儿的身体不复活，因为他们从未行善作恶，那么那些领受圣洗的\OriginalTerm{恩宠}{grace}后死去的婴儿的身体也将没有复活，因为他们在此生中也不能行善作恶。然而，如果后者要在圣人中复活，\Emph{即}在那些行善的人中复活，那么前者又将在谁中复活呢？岂不是在那些作恶的人中复活——除非我们相信有些人的灵魂，无论是在生命的复活还是在受审的复活中，都不会得回他们死时失去的身体？这种观点，即使未经正式驳斥，也因其绝对的新奇而被定罪；除此之外，谁又能忍受这样的想法：那些带着婴儿跑去受洗的人，是出于对身体的关心，而非对灵魂的关心？\Person{圣西彼廉}确实说过，为纠正那些认为婴儿在第八天之前不应受洗的人，他说需要从丧亡中拯救的是灵魂而非身体——在这句话中，他并没有发明任何新道理，而是在维护教会坚立的信仰；并且，他与一些主教同僚一起主张，婴儿在出生后立即受洗是合宜的。\par

24. 然而，各人尽可相信任何他认为合理的事，即使这与\Person{西彼廉}的某些观点相左，因他或许未在这件事上看出应该看出的东西。但是，没有人可以相信任何与那完全明确的宣言相左的事，即因一人的过犯，众人都被定罪，而且从这定罪中，除了天主的\OriginalTerm{恩宠}{grace}，藉着我们的主耶稣基督，无人能使人得救，只有在祂内，生命才被赐予一切复活的人。并且，没有人可以相信任何与教会牢固的实践相左的事，在那实践中，如果加速施行\OriginalTerm{圣洗}{baptism}的惟一理由是为了拯救孩童，那么死者和活人都会被带受洗。\par

25. 既然如此，如果灵魂是为每个人在出生时新造的，我们仍须探究并阐明，为何那些在婴儿时期未领受\OriginalTerm{基督的圣事}{sacrament of Christ}而死去的人注定要丧亡；因为如果他们这样离开身体便注定如此，这是由圣经和圣教会所共同见证的。因此，关于你认为灵魂是新造的那个意见，如果它不抵触这一稳固的信理，也让它成为我的意见；但若抵触，就让它不再是你的。\par

26. 不要对我说我们应当接受在\BibleBook{匝加利亚}{}中的话作为支持这个意见的根据：\Quote{\BibleQuote{祂在人内形成人的气息，}}\BibleRef{匝 12:1} 以及在\BibleBook{圣咏}{集}中的话：\Quote{\BibleQuote{祂塑造他们每个人的心。}}我们必须寻求最有力且最无可争议的证据，免得我们不得不相信天主是一位判官，祂惩罚任何无罪的灵魂。因为\Quote{创造}的意义等于，或很可能大于\Quote{塑造}（\OriginalTerm{塑造}{fingere}）；尽管如此，经上记载：\Quote{\BibleQuote{天主，求祢给我再造一颗纯洁的心，}}然而不能认为这里的灵魂在它开始存在之前表达一种被造的渴望。因此，正如那已经存在的灵魂因在正义中得到更新而被创造，同样那已经存在的灵魂因教义的塑造能力而被塑造。你在\BibleBook{训道}{篇}中的那句话，即\Quote{\BibleQuote{那时灰尘将归于原来的土中，生气将归于赐予它的天主。}}\BibleRef{训 12:7}，也并不支持你的意见——而我本乐意将它作为我自己的意见。不，它反而有利于那些认为所有灵魂都源于同一个灵魂的人；因为他们说，正如灰尘归于原来的土中，然而说这话时所指的身体并不是归于它所出自的那个人，而是归于第一个人的被造的土，同样，灵魂虽源于第一个人的灵魂，但并不归于他，而是归于主，因为是主将它赐给了我们的原祖。然而，由于这段经文对他们有利的证据并不那么确定，以至于显得与那个我乐于见到被辩护的意见完全对立，我认为最好将这些评论提交于你的判断，免得你想用这些经文来解除我的困惑。因为虽然任何人的愿望都不能使不真实的事成为真实，然而，如果可能的话，我会愿望这个意见是真的，正如我愿望，如果它是真的，它应得到你最清晰和无可辩驳的辩护。\par

% structure: p0041 source=h2 target=subsection reason=relative-heading-rank; base=h2

\subsection{第九章}

27. 同样的困难也困扰着那些认为灵魂已存于别处，并从天主创世之初就预备好，被祂派遣进入身体的人。因为对这些人也可以提出同样的问题：如果这些灵魂没有任何过错，顺从地进入被派往的身体，为什么在婴儿的情况中，如果他们在这一生结束时仍未受洗，就要遭受惩罚？同样的困难无疑地伴随着这两种意见。那些断言每个灵魂按照其先前状态中的行为功过，与今生指派给它的身体结合的人，自以为更容易逃脱这一困难。因为他们认为\Quote{\OriginalTerm{在亚当内}{in Adam}死亡}意指在那源自亚当的血肉中受惩罚，他们说，基督的恩宠解救年轻的和年老的脱离这一罪咎的状态。就此而言，他们教导的固然是正确、真实、极好的，当他们说基督的恩宠解救年轻的和年老的脱离罪恶的罪咎时。但我绝不相信灵魂在另一个先前的生命犯罪，并且因此被投入身体如同监狱：我拒绝并抗议这样的意见。我这样做，首先，因为他们声称这是借着某种不可理解的轮回完成的，以至于我不知道经过多少个周期后，灵魂必须再次回到同样的可朽血肉的重负和惩罚中——我想不出有什么比这更恐怖的意见了。其次，如果他们所说是真的，对于任何离世的义人，我们岂不要担心，生怕他在\OriginalTerm{亚巴郎的怀抱}{Abraham's bosom}中犯罪，而被投入那比喻中煎熬富人的火焰里吗？\BibleRef{路 16:22} 因为灵魂若能在进入身体之前犯罪，为什么不能在离开身体之后犯罪呢？最后，\OriginalTerm{在亚当内}{in Adam}犯了罪（关于这一点，宗徒说众人都因他犯了罪）是一回事，而在亚当之外、我不知道什么地方犯了罪，然后因此被推入亚当——也就是那源自亚当的身体，如同推入监牢——那完全是另一回事。至于上面提到的另一种意见，即所有灵魂都源于一个灵魂，除非有必要，我不会开始讨论；我的愿望是，如果我们现在讨论的这个意见是真实的，愿你如此为之辩护，以至于不再有任何必要去考察另一种意见。\par

28. 然而，虽然我愿意、祈求，并以热切的祈祷希望和盼望，借着你，主能消除我在此事上的无知；但如果最终我发现自己不堪当获得这知识，我便会恳求主我们的天主的忍耐之恩宠；我们对他怀有如此信德，以至于即使有些事，当我们敲门时他并未给我们开启，我们也知道，哪怕对他有丝毫的抱怨，都是错误的。我记得他亲自对宗徒们说过的话：\BibleQuote{我还有许多事要告诉你们，但你们现在不能承受。} \BibleRef{若 16:12} 至少就我而言，在这些事中，让我仍把这件事算进去；并让我提防，勿因自认为不堪当拥有这知识而发怒，免得这发怒反而更清楚地显明我确是不堪当的。我同样对许多其他事也一无所知，甚至多得无法尽述或数算；我本来会更耐心地忍受对此事的无知，但我害怕这些意见中的某一个，涉及与我们十分确信的真理相矛盾的内容，会悄悄潜入那些不警醒之人的心中。同时，虽然我尚未知道这些意见中哪一个更可取，但我明确宣认我深思熟虑后的坚信：真实的意见不会与基督教会信德中那最坚定、最有根据的信条相冲突，即婴儿，即使是刚出生的婴儿，除了借着基督之名的恩宠——这名已由他在自己的圣事中赐予——之外，别无他法可免于丧亡。\par

\SourceRecord{Translated by J.G. Cunningham.；Nicene and Post-Nicene Fathers, First Series；Vol. 1.；Edited by Philip Schaff.；Buffalo, NY: Christian Literature Publishing Co.,；1887.；<http://www.newadvent.org/fathers/1102166.htm>.}

% structure: p0001 source=h1 target=section reason=page-title

\section{第167封信（\Person{奥斯定}）或第132封（\Person{热罗尼莫}）}

\EditorialNote{奥斯定致热罗尼莫，论雅各伯书2:10（公元415年）}\par

% structure: p0003 source=h2 target=subsection reason=relative-heading-rank; base=h2

\subsection{第一章}

我敬爱的主内弟兄热罗尼莫，我由衷地在基督内尊敬你：当我给你写信提出这个关于人类灵魂的问题——现今每个人出生时是否有一个新的灵魂被造，如果是，灵魂从何处沾染了\OriginalTerm{罪愆}{bond of guilt}，而我们坚信这罪愆能藉着\OriginalTerm{基督恩宠的圣事}{sacrament of the grace of Christ}，即使施行于新生儿，也得以除去？——因为这封信论及这主题时，篇幅已达厚厚一卷，那时我不愿再给你增添其他问题的负担；但有一个问题更迫切地要求我关注，因为它更沉重地压在我心头。因此，我请求你，并因天主的名恳求你，做一件我相信会对许多人极有益处的事，就是为我解释（或指给我看你或任何其他注释者已有的著作中曾解释过）我们在\WorkTitle{雅各伯书}中应如何理解这句话：\BibleQuote{因为谁若遵守全部法律，但只触犯了一条，就算是全犯了。}\BibleRef{雅 2:10}这个问题如此重要，我非常后悔没有早就为此事写信给你。\par

2. 因为在我认为必须向你请教有关灵魂的那个问题中，我们的探讨涉及一段全然过去并沉入遗忘之中的生命；而在这个问题中，我们研究的是我们当下的生命，以及若要获得永生，这生命当如何度过。为说明这一点，请允许我引述一段有趣的轶事。有一个人掉进一口井里，井里的水足以缓冲他的坠落，使他免于死，但水不深，不足以淹没他的口而令他无法说话。另一个人走近，看到他时惊呼：\Quote{你是怎么掉进去的？}他答道：\Quote{我恳求你谋划如何救我出来，而不是问我怎么掉进去的。}同样，既然我们承认并持守天主教信仰的一项信条，即连一个婴儿的灵魂也需要藉着基督的恩宠，从罪愆的深坑中被解救出来，那么对这样一个灵魂而言，我们只要知道它得救的途径便足够了，即使我们永远无法知道它如何落入那悲惨的境地。但我曾认为我们有责任探讨这个问题，以免我们不加谨慎地持有某种关于灵魂如何与肉身结合的看法，这种看法可能否认幼小婴儿的灵魂也需要被解救，从而否定了它们受制于罪愆。那么，我们非常坚定地持守着一点：每个婴儿的灵魂都需要从罪愆中被解救出来，并且除了藉着我们的主耶稣基督、经由天主的恩宠之外，别无他途可得解救。如果我们能查明那邪恶本身的原因和起源，我们就能更好地装备自己去抵抗那些对手；这些人空洞的言论，我不称之为推理，而称之为狡辩。然而，倘若我们不能查明原因，这悲惨境地的起源隐而未显，也不应成为我们在慈悲所要求的工作上松懈的理由。不过，在与那些自以为知道他们所不知道的人交锋时，我们还有一个额外的力量和安全保障，那就是承认自己在这方面的无知，我们并不以此为耻。因为有些事情，不知道就是恶；有些事情则无法知道，或者不必知道，或者与我们所寻求的生命无关；但我现在从\Person{雅各伯}宗徒的著作中向你提出的这个问题，则与我们生活中的行为方式紧密相连，我们正是藉此生活，以获永生为目标，努力取悦于天主。\par

3. 那么，我恳请教我，应如何理解这句话：\BibleQuote{因为谁若遵守全部法律，但只触犯了一条，就算是全犯了}？这是不是断言，那犯了偷窃，甚至只对富人说\Quote{请坐在这里}而对穷人说\Quote{你站在那里}的人，也犯了杀人、奸淫和亵圣的罪？如果不是这样，怎能说人在一条上跌倒，就成了全犯呢？或者，宗徒论及富人和穷人时所说的那些事，不应算在那些只要触犯其中任何一条就算全犯的事情之中？但我们必须记住，这段话是从何处引来的，其前文是什么，并是在什么关联下说出的。宗徒说：\BibleQuote{我的弟兄们，你们既信仰我们的主耶稣基督，光荣的主，就不该按外貌待人。若有人戴着金戒指，穿着华美的衣服，进入你们的会堂，同时有一个衣服肮脏的穷人也进来，你们就专看那穿华美衣服的人，且对他说：请坐在这边好位上！而对那穷人说：你站在那里！或说：坐在我的脚凳下边！这岂不是你们自己立定区别，而按偏邪的心思判断人吗？我亲爱的弟兄们，请听！天主不是选了世俗视为贫穷的人，使他们富于信德，并继承他向爱他的人所预许的国吗？可是你们，竟侮辱穷人！}\BibleRef{雅 2:1}——正是因为你们对那穷人说\Quote{你站在那里}，而对戴金戒指的人却说\Quote{请坐在这边好位上}。接着，又有一段话解释并详论这同一结论：\BibleQuote{岂不是富贵人仗势欺压你们，亲自拉你们上法庭吗？岂不是他们亵渎你们被称的美名吗？的确，如果你们按照经书所说\InnerQuote{你应当爱你的近人如你自己}的话，满了这最高贵的法律，你们便作的对了；但若你们按外貌待人，那就是犯罪，就被法律指证为犯法者。}\BibleRef{雅 2:6}请看，宗徒如何称那些对富人说\Quote{请坐在这里}，而对穷人说\Quote{你站在那里}的人，是法律的犯法者。请看，他为免他们认为在这一件事上犯法是小罪，便继续补充说：\BibleQuote{因为谁若遵守全部法律，但只触犯了一条，就算是全犯了，因为那说了\InnerQuote{不可行奸淫}的，也说了\InnerQuote{不可杀人}。纵然你不杀人，但你行奸淫，你仍成了犯法的人。}这正是照他先前所说的：\Quote{就被法律指证为犯法者。}既然这些事如此，那么，除非能证明我们应以另一种方式理解，否则似乎可以推论说：凡对富人说\Quote{请坐在这里}，而对穷人说\Quote{你站在那里}，不按同样的尊重对待两人者，就该被判定为崇拜偶像者、亵渎者、奸淫者和杀人者——简言之，不必一一列举，以免冗长——犯了所有罪，因为他在一条上跌倒，就有了全部罪。\par

% structure: p0007 source=h2 target=subsection reason=relative-heading-rank; base=h2

\subsection{第二章}

4. 可是，难道一个人有一项德行，就拥有一切德行吗？难道缺少一项德行的人，就毫无德行吗？如果这是真的，那么宗徒的话便因此得到证实。但我所盼望的，是解释那句话，而不是证实它，因为那句话本身在我们的心目中，比一切哲学家的权威更确定。而且，即使刚才关于德行与恶习所讲的是真的，也不因此能推出一切罪都相等。至于德行不可分离地共存，这一学说，如果我没有记错的话——事实上我几乎已经忘记（虽然也许我记错了）——所有主张德行是正当生活所必需的哲学家们都是同意的。然而，罪相等的学说，只有斯多葛派\OriginalTerm{斯多葛派}{Stoics}竟敢提出，与人类一致的意见相悖：这是个荒谬的论点，你在驳斥\Person{约维尼安}（他在这个观点上是斯多葛派，但在追求并维护享乐上却是伊壁鸠鲁派\OriginalTerm{伊壁鸠鲁派}{Epicurean}）时，已从圣经中非常清楚地驳斥了它。在那篇极令人愉快且卓越的论文中，你已充分说明，这并不是我们那些作者们的教义，更确切地说，不是那借他们发言的真理本身的教义，认为一切罪都是相等的。我现在要尽最大努力，在天主的帮助下，试图表明：虽然哲学家关于德行的学说是真的，我们却因此仍不必承认斯多葛派一切罪相等的学说。如果我成功了，我将期待你的嘉许；若我在任何方面有不足之处，恳请你补足我的缺欠。\par

5. 那些主张一个人有一项德行就有一切德行，缺少一项就缺少一切德行的人，其推理是正确的，因为明智不能是怯懦的、不义的或不节制的；如果它是这些中的任何一种，就不再是明智了。此外，只有当明智是勇敢的、公义的、节制的，它才是明智，那么，凡是明智所在之处，必然具有其他德行。同样，勇敢也不能是不明智的、不节制的或不义的；节制必须明智、勇敢、公义；公义若不具有明智、勇敢、节制，便不存在。因此，这些德行中任何一项真正存在的地方，其他德行也同样存在；而缺失某些德行的地方，那看来有点像德行的东西，其实并不是真正的德行。\par

6. 如你所知，有些恶习与德性截然对立，如轻率是智德的对立面。但有些恶习与德性对立，仅因它们是恶习，却以欺骗性的外表相似于德性；例如，在与智德这德性相关时，不是轻率，而是机警。我这里说的机警，是常被理解并用于形容奸邪之人的那种，并非我们的圣经中通常所取的意义；在圣经中，它常用于好的意义，如：\BibleQuote{你们要机警如同蛇，} \BibleRef{玛 10:16}，又如：\BibleQuote{使愚人获得机警} \BibleRef{箴 1:4}。的确，在外教作家中，有一位极富才华的拉丁作家论及\Person{喀提林}时曾说：\Quote{他也不乏机警以防范危险，}将\OriginalTerm{机警}{astutia}用于好的意义；但这种用法在他们当中极少见，而在我们当中却极普遍。同样，对于节德之下的诸德性也是如此。奢侈浪费显然与节俭的德性对立；但常人所谓的吝啬，固然是一种恶习，却并非就其本质，而是以极似的外表窃取了节俭之名。同样，不义与义德截然对立；但报复的欲望常冒充义德，却是恶习。勇德与怯懦之间存在着明显的对立；但莽勇虽在本质上异于勇德，却因其相似而欺骗我们。坚定是德性的一部分；反复无常是一种与之远隔且无疑对立的恶习；但顽固却僭用坚定之名，实则全异，因为坚定是德性，而顽固是恶习。\par

7. 为避免再次重复同一论点，我们仅取一事为例，由此可明其余。据记载他的人所知，\Person{喀提林}能忍受不可思议的寒冷、干渴、饥饿，且极能忍耐禁食、寒冷与不眠，远超出任何人的相信；因此，他自视并为其追随者所视，为一个极具勇德之人。然而这勇德并非智德的，因为他舍善取恶；并非节德的，因为他生活放荡至极，污名累累；并非义德的，因为他阴谋叛国；故这并非勇德，而是莽勇，僭用勇德之名以欺骗愚人；因为若是勇德，便不会是恶习而会是德性；若是德性，便绝不会被其他与它不可分离的德性所弃舍。\par

8. 因此，当人被问及恶习时，若问是否有一恶习存在则所有恶习同样存在，或一恶习不存在则所有恶习都不存在，这实难加以阐明，因为常有两种恶习与同一德性对立：一种明显对立，另一种表面相似。因此，\Person{喀提林}的莽勇更易见其并非勇德，因其未偕同其他德性；但难使人信服他的莽勇即是怯懦，因为他惯于忍受并坚忍最大艰难，几近难以置信。然而，若更细察，或许这莽勇本身正显示为怯懦，因为他逃避磨练真正勇德所需的自由学术之辛劳。尽管如此，既有鲁莽而不犯怯懦之人，亦有怯懦而不犯鲁莽之人，且二者均为恶习，因为真正勇德者既不鲁莽行事，也不无端畏惧；我们不得不承认恶习比德性更多。\par

9. 故此，有时一恶习被另一恶习取代，如贪财被好名取代。偶有一恶习离去，以便更多恶习占其位置，如醉汉节饮后，可能受制于吝啬与野心。因此，恶习可能让位于恶习，而非德性为其继，故恶习更多。然当一德性进入时，必然（因其携同所有其他德性）原在人内的一切恶习尽皆退去；盖非所有恶习皆在其人之内，然有时众多，有时或多或少居于其位。\par

% structure: p0014 source=h2 target=subsection reason=relative-heading-rank; base=h2

\subsection{第三章}

10. 我们必须更加仔细地考察这些事是否如此；因为所谓\Quote{有一德行即有诸德行，而缺少一德则诸德尽失}的说法，并非出于默感，而是出于许多人的见解，这些人固然机敏勤勉，但毕竟是人。然而我必须承认，我不说那些据称\InnerQuote{德行}一词由其名字派生而来的人，只说一位忠于丈夫的妇女，她出于对天主的诫命和恩许的尊重，且首先是忠于天主——对于她，我不知如何能说她是不贞的，或者说贞洁不是德行或微不足道。对于一位忠于妻子的丈夫，我也有同感；但这样的人很多，我却不能断言他们中没有任何罪，毫无疑问，他们身上的罪，无论是什么，都源于某种恶习。由此可知，尽管虔诚男女的婚姻忠诚无疑是一种德行，因为它既非虚无亦非恶习，但它并不带来一切德行，因为如果一切德行都在那里，就不会有恶习，如果没有恶习，就不会有罪；然而，哪里有完全无罪的人？哪里有一个没有任何恶习，亦即所谓的罪的燃料或根源的人，当那位曾靠在主胸前的\Person{若望}说：\BibleQuote{如果我们说自己没有罪，便是自欺，真理也不在我们内}\BibleRef{若一 1:8}？我们无须在给你的信中更长地讨论这一点，但我之所以作此陈述，是为着或许会读这封信的其他人。事实上，你在那篇同样出色的\WorkTitle{驳若维尼安}中，已经仔细地从圣经证明了这一点；在那篇著作中，你也引用了\BibleQuote{我们在许多事上都犯过失}\BibleRef{雅 3:2}这句话，它正是出自我们正在探究其含义的这封书信。因为尽管说话的是基督的宗徒，他并没有说\Quote{你们犯过失}，而是说\Quote{我们犯过失}；而且，尽管在正讨论的这节经文中他说：\BibleQuote{谁若遵守全部法律，只犯了一条，就算是全犯了}\BibleRef{雅 2:10}，但在刚才引用的这话中，他断言我们犯过失不是在一件事上，而是在\Emph{许多}事上，并且不是有些人犯过失，而是我们\Emph{众人}都犯过失。\par

11. 任何信徒都绝不应认为，那成千上万为了不自欺、使真理不在他们内而诚实地承认自己有罪的基督的仆人，竟毫无德行！因为智慧是一项大德行，智慧本身曾对人说：\BibleQuote{看，敬畏主就是智慧。} 那么，我们绝不能说，如此众多、如此伟大的信从和虔诚之人没有敬畏主之心，希腊人称之为\OriginalTerm{虔诚}{\GreekText{εὐσέβεια}}，或更直译、更完整地，\OriginalTerm{敬畏天主}{\GreekText{θεοσέβεια}}。而敬畏主岂不是对主的敬拜吗？而真正的敬拜岂不源于爱德吗？因此，由纯洁的心、光明磊落的良心和真诚的信心所发出的爱德，是伟大而真实的德行，因为它是\BibleQuote{这训令的目的}\BibleRef{弟前 1:5}。爱德当之无愧地被说成\BibleQuote{爱情猛如死亡，}\BibleRef{歌 8:6}因为它像死亡一样，不为任何人所战胜；或者因为此生的爱德，其尺度直至死亡，正如主说：\BibleQuote{人若为自己的朋友舍掉性命，再没有比这更大的爱情了}\BibleRef{若 15:13}；又或者因为，正如死亡强行使灵魂与身体的感官分离，爱德也使灵魂与肉欲分离。知识，若是正确的那种，是爱德的婢女，因为没有爱德，\BibleQuote{知识只会使人傲慢自大}\BibleRef{格前 8:1}，但当爱德通过建树而充满了心灵，知识在那里就找不到任何可以自大的空虚了。此外，\Person{约伯}指明，什么是有益的知识，他定其义说：在说了\BibleQuote{敬畏主就是智慧}之后，又补充道\BibleQuote{远离邪恶就是明智。}\BibleRef{约 28:28}那么，我们为什么不说，有这德行的人就有诸德，因为\BibleQuote{爱德是法律的满全}？\BibleRef{罗 13:10}难道不是这样吗：一个人越有爱德，便越富有德行，爱德越少，德行也越少，因为爱德本身就是德行；一个人德行越少，恶习就越多。因此，当爱德圆满无缺时，便再没有恶习存留。\par

12. 因此，在我看来，斯多葛派的错误在于，拒绝承认一个在智慧上进步的人有任何智慧，而断言只有那些在智慧上达到完全完美的人才有智慧。他们确实不否认他取得了进步，但说他在任何程度上都不配称为智者，除非他可以说从深渊中一跃而出，突然升入智慧的自由空气。因为，正如一个人在溺水时，不论他头顶的水深是几斯他狄亚或只有一掌一手指之宽，都无关紧要，他们在论及那些走向智慧的人的进步时也持同样的说法，他们认为这等人就像人从漩涡底部上升趋向空气，但除非他们不断前进，完全从愚昧中逃脱出来，如同从没顶洪水中逃脱，否则他们就没有德行，就不是智者；然而，一旦他们这样逃脱了，他们立刻就拥有完美的智慧，没有任何愚昧的痕迹存留，以至于不会产生任何罪过。\par

13. 這個比喻將愚昧比作水，智慧比作空氣，彷彿心靈從愚昧的窒息影響中浮現，突然呼吸到智慧的自由空氣；但在我看來，這個比喻與我們聖經的權威陳述並不十分吻合。至少就可用物質事物說明屬靈事物而言，更好的比喻是將惡習或愚昧比作黑暗，而將德行或智慧比作光明。因此，通往智慧的路不像一個人從水中升到空中，在浮出水面的瞬間突然自由呼吸；而是像一個人從黑暗走進光明，在前進時逐漸受到更多光照。因此，只要這過程尚未完全成就，我們就稱這人為從巨大洞穴的黑暗深處走向洞口的人，他越靠近洞口，就越受光的接近的影響；所以他所有的光都來自他正在接近的光，而一切存留的黑暗都來自他從中擺脫的黑暗。因此，確實，在天主面前，\BibleQuote{凡活著的人沒有一個能成義}，然而，\BibleQuote{義人必因他的信德而生活}。 \BibleRef{哈 2:4} 一方面，\BibleQuote{聖人們以正義為衣}，\BibleRef{約 29:14} 有的更多，有的更少；另一方面，在此塵世，沒有人完全無罪——這人犯罪多，那人犯罪少，而最好的人是最少犯罪的人。\par

% structure: p0019 source=h2 target=subsection reason=relative-heading-rank; base=h2

\subsection{第四章}

14. 但是，我為何彷彿忘記了自己是在對誰說話，竟以教師的口氣陳述那個我希望得到你指教的問題呢？不過，既然我已決定將我對諸罪相等的看法（這個問題與我正在撰寫的事情有密切關係）提交給你審閱，那麼現在就讓我最終把我的陳述做個總結吧。即使一個人有一種\OriginalTerm{德行}{virtue}便有所有德行，而缺乏一種德行便毫無德行這句話是真的，這也不意味著所有罪都是相等的；因為，雖然沒有德行就沒有正確的事物，但絕不是說，在惡行中，一個不能比另一個更壞，或偏離正確並不存在程度之別。然而，我認為，更符合真理、也與聖經更一致的說法，是靈魂的不同傾向（它們雖然不能見到佔據不同位置，但通過各自獨特的運作，與身體的不同肢體一樣明顯地被感知）就像身體的肢體那樣：在同一身體中，一個肢體受到光的照耀更完全，另一個較少照耀，第三個完全沒有光，並隱藏在某些不透光的覆蓋之下，靈魂的各種傾向也發生類似的情況。如果這是正確的，那麼根據每個人受神聖之愛的光照方式，可以說他擁有一種德行而缺乏另一種德行，或擁有一種德行的分量較大而另一種分量較小。因為，對於那稱為敬畏天主的愛，我們可以正確地說，它在一個人身上比在另一個人身上更大，也可以說它在一個人身上有一點，在另一個人身上完全沒有；我們也可以正確地說，就一個人而言，他的貞潔比忍耐更大，如果他在進步，他今天擁有的任何一種德行比昨天更高級，或者他仍然缺乏節制，但同時擁有大量的憐憫。\par

15. 簡而言之，就聖潔生活而言，我對德行的觀點可以總結為：德行就是愛，就是愛那應當愛的對象。這種愛，在某些人身上較大，在某些人身上較小，在有些人身上則完全不存在；但是，在絕對充滿、不能增加的意義上，它不存在於任何生活在世上的人身上；然而，只要它還能增加，毫無疑問，它越是不及所應是的程度，這缺失便是出於\OriginalTerm{惡習}{vice}。由於這種惡習，\BibleQuote{世上沒有只行善而不犯罪的義人；} \BibleRef{訓 5:7} 由於這種惡習，\BibleQuote{在天主面前，凡活著的人沒有一個能成義。} 因著這惡習，\BibleQuote{如果我們說我們沒有罪過，就是欺騙自己，真理也不在我們內。} \BibleRef{若一 1:8} 也因著這惡習，無論我們取得了多大的進步，我們都必須說，\BibleQuote{寬免我們的罪債；} \BibleRef{瑪 6:12} 儘管一切在言語、行為和思想上的罪債都已在聖洗中洗淨。因此，正確去看的人，便看出他應當從何處、在何時、以及在何種境況下，希望達到那不能再增加的完美境界。此外，如果沒有誡命，人就沒有確實的途徑來省察自己，看出自己應當避開什麼，應當追求什麼，以及在什麼事上應當感謝或祈禱。所以，誡命大有裨益，只要給予自由意志足夠的自由，好使天主的恩寵得到更豐厚的尊崇。\par

% structure: p0022 source=h2 target=subsection reason=relative-heading-rank; base=h2

\subsection{第五章}

16. 若果真如此，那么一个人要是遵守了全部法律，却在一个点上犯了罪，怎么就算是全犯了呢？岂不是因为，法律的满全就是我们对天主和近人的爱，而在这爱的诫命上，\BibleQuote{全部法律和先知都系于其上}\BibleRef{玛 22:40}，因此，凡违背那一切所系之爱的，就理所当然地被判决为全犯了吗？然而，没有人犯罪而不触犯这爱；\BibleQuote{因为这\InnerQuote{不可奸淫，不可杀人，不可偷盗，不可贪恋}，以及其他任何诫命，都包含在这句话里：\InnerQuote{应当爱你的近人如你自己。}爱不加害于人，所以爱就是法律的满全。}\BibleRef{罗 13:9} 可是，没有人能爱近人，除非他出于对天主的爱，尽己所能，也将他所爱如己的近人引向对天主的爱；他若不爱天主，就既不爱自己，也不爱近人。因此，如果一个人遵守了全部法律，却在一个点上犯了罪，他就成了全犯，因为他所做的，违背了那一切法律所系之爱。所以，人因违背那一切所系的爱，就成了全犯。\par

17. 那么，为何不能说所有罪都是相等的呢？岂不是因为，犯更重大罪的人，对爱之法律的违背就更严重，而犯更轻微罪的人，对爱之法律的违背就较轻吗？而且，正因为他犯了任何一个罪，就成了全犯；但他若犯了更重大的罪，或在多方面犯了罪，他的罪责就更重；若犯了更轻微的罪，或在较少方面犯了罪，他的罪责就较轻——这样，罪责的轻重与他犯罪的大小多寡成正比。然而，即使只在一点上触犯，他仍是全犯，因为违背了那一切法律所系的爱。如果这些是真的，这就解释了那位富有宗徒恩宠的人所说的话：\BibleQuote{我们众人都犯许多过失。}\BibleRef{雅 3:2} 因为我们众人都犯错，但有的更严重，有的更轻微，这取决于各人所犯罪的重与轻；每个人在犯罪上的程度，与他缺乏对天主和近人的爱的程度相当；相反，随着他对天主和近人爱的增长，犯罪就减少。因此，一个人爱越少，罪就越多。爱的圆满，乃是在我们内不再存有任何罪恶的脆弱时达到的。\par

18. 而且，依我看来，我们不可将\Quote{怀着对容貌的偏私，而信仰我们的主耶稣基督}视为轻微的罪过，如果我们把上下文提及的坐着和站着的区别，是指教会内的尊位而言；因为，谁能忍受看到有钱人在教会中被选授予尊位，而更有资格、更圣善的穷人却被轻视呢？然而，如果宗徒在那里说的是我们日常的聚会，谁在这事上没有犯过呢？但同时，只有在那些心里存着按财富估量人价值的念头的人，才真正在这事上犯罪；因为，这似乎是那句\Quote{你们自己心里不存成见，反而成了怀邪恶思想的判官了吗}的意思。\par

19. 因此，自由的法律，即爱的法律，就是宗徒所说的：\BibleQuote{如果你们按照圣经所说，履行了\InnerQuote{应当爱你的近人如你自己}这条至尊的法律，你们便做得不错；但你们若按外貌待人，那就是犯罪，法律就要指证你们为犯法者。}\BibleRef{雅 2:8} 接着（在\Quote{谁若遵守全部法律，但只触犯了一条，就算全犯了}这段难解的话之后——关于这点，我已作了足够充分的陈述），他再提到同样的自由的法律说：\BibleQuote{你们要怎样说话，怎样行事，按照那使人自由的法律受审判的人一样。}并且，由于他经验到他在不久之前所说的话——\BibleQuote{我们众人都犯许多过失}，他便提出了一剂万灵的药方，好像要天天敷在我们灵魂天天遭受的、虽较轻微却真实的创伤上，因为他说：\BibleQuote{对不行怜悯的人，审判也没有怜悯。}\BibleRef{雅 2:13} 因为上主怀着同一目的说：\BibleQuote{你们要赦免，也就蒙赦免；你们给，也就给你们。}\BibleRef{路 6:37} 之后宗徒又说：\Quote{怜悯必得胜审判。}这并不是说怜悯胜过审判（因为怜悯并非审判的对手），而是怜悯\OriginalTerm{夸胜}{rejoices}于审判，因为有更多人因怜悯而被聚集；但这些人就是那些施行怜悯的人，因为，\BibleQuote{怜悯人的人是有福的，因为他们要受怜悯。}\BibleRef{玛 5:7}\par

20. 因此，他们得到宽恕，实在是非常公正的，因为他们宽恕了别人；他们所需求的，也赐予他们，因为他们将所有的施予了别人。因为天主在审判时运用仁慈，在施予仁慈时也运用审判。因此，圣咏作者说：\BibleQuote{我要歌颂仁慈与审判，向祢，主。}因为若有人自以为过于正义而不需要仁慈，便安然自恃，好像无忧无虑，等待那没有仁慈的审判，他便激起了那最公义的义怒，圣咏作者正是出于对此的敬畏而说：\BibleQuote{求祢不要审判祢的仆人。}为此，主对悖逆的子民说：\BibleQuote{你们为何在审判中与我争辩？}因为当正义的君王坐在祂的宝座上时，谁能夸耀说自己有纯洁的心？谁能夸耀说自己无罪呢？那么，除了仁慈在审判中\BibleQuote{得胜}审判之外，还有什么希望呢？但这只有在那些已施予仁慈，并且诚恳地说：\BibleQuote{宽免我们的罪债，如同我们也宽免别人一样}，和那些毫无怨言地施予的人身上，才会发生，因为\BibleQuote{主爱乐捐的人。}\BibleRef{格后 9:7}。结论是，\Person{圣雅各伯}在这一段中谈及慈悲的善工，正是为了安慰那些因前文所述而可能深感不安的人，他的目的是要劝诫我们，我们今生永远无法摆脱的日常小罪，也当如何用每日的良药来补赎；免得任何人在一点上跌倒就成了犯全法的，因在许多事上跌倒（因为\BibleQuote{我们在许多事上都有过失}），而在极度累积的罪债下，被带到至高审判者的法庭前，在那法庭上得不到丝毫仁慈，因为他未曾向别人施仁慈善；如果相反，他通过对别人施予宽恕和慷慨，反而配得自己罪过的赦免，并享受圣经所许诺的恩赐。\par

21. 我写得非常冗长，或许令你感到厌烦，因为你虽然赞同上面所说的话，却并不期望被当作一个刚刚学习这些你已经惯于教导别人的真理的人来对待。然而，如果这些话中——不是在阐述时所用的语言风格上，因我对辞藻并不十分在意，而是在所说内容的实质方面——有你博学的判断认为不对的地方，我恳请你在回复中向我指出，并毫不犹豫地纠正我的错误。我实在可怜那一个人，他鉴于你不知疲倦的辛劳和工作的神圣性质，却不因此既献上你当得的尊敬，又感谢我们的主天主，因祂的恩宠使你成为今日的你。因此，既然我更当乐意向任何一位老师学习那些我所不知而对我有害的事，而非急于教导别人我所知道的事，那么，我更有理由向你要求这爱的还报，因为借着你的学问，在主的圣名和祂的助佑下，拉丁语教会文学已达到了前所未有的高度。然而，我特别请你注意这句话：\BibleQuote{谁若遵守全部法律，但只触犯了一条，就算是全犯了。}我亲爱的弟兄，如果你知道有更好的解释方式，我恳切地请求你，靠着主，屈尊将你的解释传授给我们。\par

\SourceRecord{Translated by J.G. Cunningham.；Nicene and Post-Nicene Fathers, First Series；Vol. 1.；Edited by Philip Schaff.；Buffalo, NY: Christian Literature Publishing Co.,；1887.；<http://www.newadvent.org/fathers/1102167.htm>.}

% structure: p0001 source=h1 target=section reason=page-title

\section{第169封信（公元415年）}

\Emph{\Person{奥斯定}主教致\Person{埃伏迪乌斯}主教。}\par

% structure: p0003 source=h2 target=subsection reason=relative-heading-rank; base=h2

\subsection{第一章}

1. 如果尊驾如此看重那些特别占据我心神、且我不愿被任何其他事打扰而中断的著作，那么请派人来为您抄写它们吧。因为我已在今年复活节前，将近四旬期开始之时，完成了今年着手撰写的几部书中好几部。在已完成的\WorkTitle{天主之城}前三卷——为反对其仇敌，即魔鬼的崇拜者——之后，我又增添了两卷；在这五卷书中，我认为已充分驳斥了那些人：他们主张，必须敬拜［异教的］神明，才能获取今生今世的繁荣，并且他们因认为这种繁荣因我们而受阻，故而敌视基督徒之名。接下来，我必须如我在第一卷中所承诺的，回答那些人，他们认为敬拜他们的神明是获得死后生命的唯一途径，而我们之所以成为基督徒，正是为了获得这生命。我还口授了篇幅可观的注释，讲解三篇圣咏，即第68篇、第72篇和第78篇。而对其他圣咏的注释——既未口授，甚至尚未动笔——人们正殷切地期待并要求于我。我不愿从这些研究工作中被任何从别处突如其来的问题拖走并阻扰；是的，我如此不情愿，以至于目前甚至不愿转向那些\WorkTitle{论天主圣三}的著作，这些著作我已着手多年，至今尚未完成，因为它们需要极大的辛劳，而且我相信它们性质深奥，只有少数人能理解；因此，它们远不及那些我希望将有益于极多人的著作，能更迫切地唤起我的注意。\par

2. 因为宗徒所说的那句话：\BibleQuote{那无知的，就由他无知吧。}\BibleRef{格前 14:38}，并非如您的信中所断言，是用以论及这个主题的；仿佛那无法运用其理智辨别出\OriginalTerm{天主圣三}{Trinity}不可言喻的一体性的人——正如人灵魂中的记忆、悟性与意志的一体性可被辨别一样——就要受到这惩罚似的。宗徒说这些话，原是为了完全不同的目的。查阅该段经文，您便会看出，他所谈论的，乃是有助于多数人在信德与圣德上获得建树的事，而非那些只有少数人才能艰难领悟的事，且即便是这些人，在今生也只能在极微小的程度上领悟如此伟大的主题。他所定的原则是——说预言应胜过说方言；这些恩赐不应无序地运用，仿佛预言之灵强迫他们违背己意说话；妇女应在教会中保持缄默；凡事都应端庄且按秩序而行。在谈论这些事时，他说：\BibleQuote{若有人自以为先知，或属灵的，就该知道我所写给你们的事，因为那是主的命令。若有人无知，就由他无知吧。}他藉这些话，意在约束并整顿那些在教会中特别容易制造混乱的人，因为他们自以为在属灵恩赐上出类拔萃，实则却以僭越的行为扰乱一切。\BibleQuote{若有人自以为先知，或属灵的，}他说，\BibleQuote{就该知道我所写给你们的事，因为那是主的命令。}如果有人自以为为先知，而实际上并非如此——因为真正的先知无疑知道，且无需提醒和劝勉，因为\BibleQuote{他洞察一切，却不受任何人的审断。}\BibleRef{格前 2:15}因此，那些在教会中自以为是什么，实则并非如此的人，便引起了混乱和麻烦。他教导这些人认识主的命令，因为祂不是\BibleQuote{混乱的天主，而是和平的天主。}\BibleRef{格前 14:33}但\BibleQuote{若有人无知，就由他无知吧，}意即他必被弃绝；因为天主并非不知道——就单纯的知识而言——那些将来祂要对他们说：\BibleQuote{我不认识你们，}\BibleRef{路 13:27}的人，但这句话所意指的，正是他们被弃绝。\par

3. 再者，既然主说：\BibleQuote{心里洁净的人是有福的，因为他们要看见天主，}\BibleRef{玛 5:8}并且那直观被应许给我们，作为最终的至高赏报，我们就无须担心：若我们现在还未能清晰地看见我们关于天主本性的所信的事，这种有缺陷的理解会使我们遭受那\BibleQuote{无知的，就由他无知吧}的判决。因为当\BibleQuote{世人凭自己的智慧不认识天主，天主就乐用愚妄的道理，拯救那些信的人}的时候——这传道的愚妄和\BibleQuote{天主的愚妄总比人智慧}——将许多人引向救恩，以致不仅那些尚不能以清晰的理智领悟自己在信德中所持守的天主本性的人，而且甚至那些尚未这样认识自己灵魂的本性，以致能以相同的确定性区分其无形无像的本质与整个身体——正如他们觉察到自己活着、理解、愿意那样——的人，都并不因此而被排斥在那种因传道的愚妄而赋予信者的救恩之外。\par

4. 因为，如果\Person{基督}只为那些能以清晰的理智辨识这些事的人受死，我们在教会中的劳苦几乎就徒然了。但如果事实上，大批没有大智慧的普通群众，凭着信仰奔向那位医师，结果因着\Person{基督}并祂的十字架而得医治，就如\BibleQuote{罪恶在哪里越多，恩宠也格外丰盛}\BibleRef{罗5:20}，这事便以奇妙的方式发生，借着天主深奥的智慧与知识的丰富，以及祂不可测度的判断：一方面，一些在自身本性中辨别物质与精神的人，以此成就自夸，轻看那使信者得救的愚拙道理，就远远偏离那通向永生的唯一道路；另一方面，因为\Person{基督}为之而死的人无一丧亡\BibleRef{若17:12}，许多以\Person{基督}的十字架为荣，且不离弃这同一道路的人，虽然对某些人以极深奥精微探究之事一无所知，却达到那永恒、真理与爱——即那持久、清晰而圆满的真福——在那里，凡恒居、看见并爱慕的人，一切事都明了。\par

% structure: p0008 source=h2 target=subsection reason=relative-heading-rank; base=h2

\subsection{第二章}

5. 因此，我们应以坚定的虔信相信一个天主，即圣父、圣子和圣神；同时我们相信圣子不是圣父，圣父不是圣子，圣父和圣子也都不是圣神。不要以为在天主圣三内有时间或空间上的分离，而是这三位同等、同永恒，且完全同一本性：并且，受造物并非一部分由圣父造，一部分由圣子造，一部分由圣神造，而是已造和正在受造的一切，每一位和全部，都存在于创造者天主圣三之内；没有人单由圣父而无圣子和圣神得救，或单由圣子而无圣父和圣神，或单由圣神而无圣父和圣子得救，而是由圣父、圣子和圣神，这唯一、真实、真正不朽（即绝对不变）的天主得救。同时，我们相信圣经中许多事分别论及三位中的每一位，为教导我们：虽然他们是不可分离的天主圣三，然而他们是三位。因为，正如他们的名字在人类语言中被说出时，不能同时被称呼，虽然他们不可分离的结合在任何时刻都是同时的，同样，在圣经的一些地方，他们也借着某些受造物，分别而彼此关联地呈现给我们：例如，在\Person{基督}受洗时，圣父在声音中被听到说：\BibleQuote{你是我的儿子；}圣子在由童贞女所生而取得的人性中被看到；圣神在鸽子的形体中被看到\BibleRef{路3:22}——这些事物分别呈现出三位，但绝非分离。\par

6. 为了以一种理智能够理解的方式呈现这一点，我们从\OriginalTerm{记忆}{Memory}、\OriginalTerm{理解力}{Understanding}与\OriginalTerm{意志}{Will}借用一个比方。因为，虽然我们可以按各自的次序并分别地谈论这些官能中的每一个，但我们若不经由其他两个，就既不能运用，也不能提及其中任何一个。然而，尽管我们使用这三个官能与天主圣三作比较，绝不可认为相比之物在各方面都一致，因为在任何推理过程中，哪里能找到完全对应、在每一点上都可应用于所要说明之物的比喻呢？因此，首先，这种相似性在这一方面是不完全的：记忆、理解力和意志并非灵魂，而只存在于灵魂之内；但天主圣三并不存在于天主之内，而就是天主。所以，在天主圣三中显出一种令我们惊叹的\Emph{\OriginalTerm{单纯性}{simplicitas}}，因为在天主圣三中，存在是一回事，理解或做任何其他归给天主本性的行为，并不是另一回事；但在灵魂中，存在是一回事，理解是另一回事，因为即使不使用理解力，灵魂也依然存在。其次，谁敢说圣父不是凭自己，而是凭借圣子来理解，就如记忆不是凭自己，而是凭理解力来理解，或者更准确地说，拥有这些官能的灵魂，除了借理解力便不能理解，正如它只借记忆来记忆，只借意志来行使意志呢？因此，这个比喻所要说明的要点是：无论我们以何种方式理解，就灵魂中的这三个官能而言，当那分别代表它们的各个名称被说出来时，每一个名称的说出，却总是在三个官能的共同运作中完成的，因为它是借着记忆、理解力和意志的行动才被说出——同样，我们也这样理解圣父、圣子与圣神：任何有时只向我们心灵呈现天主圣三中一位的受造物，其产生没有不是通过天主圣三同时的、本质上不可分离的运作的；因此，圣父的声音、圣子的身体与灵魂、圣神的鸽子，也都不是以任何其他方式产生，而是通过天主圣三的共同运作。\par

7. 此外，那声音确实没有与圣父的\OriginalTerm{位格}{person}不可分解地合一，因为它一经发出便停止了。鸽子的形状也没有与圣神的位格不可分解地合一，因为它也像那在山上遮盖救主和祂三个门徒的光明云彩\BibleRef{玛 17:5}，或者更像那曾代表同一圣神的火焰之舌，一旦完成了作为\OriginalTerm{象征}{symbol}的使命便不复存在了。但降生成人的圣子所显示的身体与灵魂则不然：鉴于人的得救乃是一切这些事物的目的，祂所显现的\OriginalTerm{人性}{human nature}，以一种奇妙而独特的方式，被真实地\OriginalTerm{摄取}{assumed}与圣言——亦即天主的独生子——的位格合一，而圣言则不变地存留于其自身\OriginalTerm{性体}{nature}中，在祂内不可能有任何复合元素，能与只是幻象的人的灵魂共存。我们确实读到：\BibleQuote{智慧之神是多样的；}\BibleRef{智 7:22}但祂也被恰当地称为单纯的。祂是多样的，因为祂拥有许多事物；但又是单纯的，因为祂与祂所拥有的并非不同，正如子被说成在自身内拥有生命，而祂自己就是那生命。人性来到了圣言处；圣言并非带着可变性进入人性；因此，在与祂所摄取的人性结合中，祂仍被恰当地称为圣子；正因为如此，同一位既是永恒不变、与圣父同为永恒的圣子，又是被安放在坟墓里的圣子——前者仅适用于作为圣言的祂，后者仅适用于作为人的祂。\par

8. 因此，我们在读到关于圣子的任何陈述时，必须留意这些陈述是就祂的两个\OriginalTerm{性体}{nature}中的哪一个而言的。因为祂\OriginalTerm{摄取}{assumption}了一个人的灵魂与身体，并未增加位格的数目：天主圣三仍如从前一样。正如在每个人内——除了祂所唯一摄取而与之位格合一的那一位外——灵魂与身体构成一个位格，同样，在基督内，圣言及其人的灵魂和身体构成一个位格。又如，\Quote{哲学家}之名，无疑仅指人的灵魂而言，但我们说哲学家被杀、哲学家死去、哲学家被埋葬，尽管这一切事件发生于他的身体而非他作为哲学家的那一部分，这并非荒诞，反倒是语言上最为恰当且普通的用法；同样地，天主，或圣子，或荣耀之主，或其它任何类似名称，是指作为圣言的基督，然而，说天主被钉十字架，也是正确的，因为毫无疑问，祂是在人性而非在祂作为荣耀之主的那一部分承受了此死亡。\BibleRef{格前 2:8}\par

9. 然而，那声音、鸽子的形体、以及分开有如火焰的舌头停留在每人头上的形状，这些都像在西乃山发生的可畏奇事\BibleRef{出 19:18}，又像日间的云柱和夜间的火柱\BibleRef{出 13:21}，只是作为象征而产生，并在完成其目的后便消失了。我们对此尤其要防范的是，免得有人以为天主的性体——无论是圣父、或子、或圣神的性体——是可变的或可转化的。我们也不可因以下事实而不安：\OriginalTerm{标记}{sign}有时接受了\OriginalTerm{所表示的事物}{thing signified}的名称，例如圣经说圣神以鸽子的形体降下并停留在祂身上；因为同样地，被击打的磐石被称为基督\BibleRef{格前 10:4}，因为它原是基督的象征。\par

% structure: p0014 source=h2 target=subsection reason=relative-heading-rank; base=h2

\subsection{第三章}

10. 然而，我奇怪的是，虽然你相信那说\BibleQuote{你是我的爱子}的声音，能够通过天主的行动，无需\OriginalTerm{灵魂}{soul}的媒介，由某种物质性的东西产生，但你却不相信，任何真实生物的形体与动作，完全相似地，也能以同样的方式，即通过天主的行动，无需赋予生命的灵魂作为媒介而产生。因为如果无生命的物质，无需生命灵魂的工具，就能服从天主，发出如同有生命物体惯于发出的声音，以将清晰的话语传达给人的耳朵，那么为什么它不能同样服从天主，以同样的创造主的德能，无需任何生命灵魂为工具，而将鸟的形状和动作呈现给人的眼睛呢？视觉和听觉的对象——那冲击耳朵的声音和呈现在眼前的外观，声音的吐字和肢体的轮廓，一切可听可见的运动——都同样来自邻近我们的物质；那么，难道只准许听觉，而不准许视觉，告诉我们关于这由身体感官所知觉的身体，它既是真实的身体，而且除了身体感官所知觉的之外，什么也不是吗？因为每个生物的灵魂，当然不是任何身体感官所能知觉的。因此，我们无需探究鸽子的形体如何呈现在眼前，正如我们无需探究一个能说话的形体，其声音如何传入耳中一样。因为如果在一个声音被产生的情况中，能够无需灵魂的媒介（而这声音并非类似的声音），那么，在记载圣神\BibleQuote{\Emph{如同}鸽子}降下的情况中，岂不是更容易可能吗？这句话表明，呈现在眼前的只是一个形体，并非断言看到了一只真实的生物！同样，在五旬节那天，记载说：\BibleQuote{忽然从天上来了一阵响声，好像暴风颳来，又有如同火的舌头显现给他们}，\BibleRef{宗 2:2} 其中提到像风和像火的东西，\Emph{即}类似这些常见而熟悉的自然现象，被认为是被感知到了，但似乎并没有暗示这些常见而熟悉的自然现象真的产生了。\par

11. 然而，如果更精微的推理或更彻底的研究证明，那在时间和空间中自然不具备运动的东西（\Emph{即}\OriginalTerm{物质}{matter}），除非通过那只能在时间、不能在空间中运动的东西（\Emph{即}\OriginalTerm{精神体}{spirit}）的中介作用，否则不能运动，那么由此将得出，所有那些事必然是通过生物的工具性完成的，就如天使所行的事一样——关于这个主题，更详尽的讨论会显得冗长，且无必要。此外还要加上，有些异象不仅出现在睡眠或神志不清时，也出现在清醒而理智健全的人面前，它们如此清晰地显现于心神，如同显现于身体感官。这些异象并非由于魔鬼嘲弄人的欺骗，而是借着某种精神启示，以类似物体的非物质形式完成，并且根本无法与真实对象区分，除非借着天主的助佑，由理智更充分地启示和分辨，有时（但很困难）能在当时做到，但大多在它们消失之后。既然关于这些异象——不论其本性真的是物质的，或仅是表面物质而实际是精神的——似乎如身体感官所知觉到的那样，呈现于我们的心神，那么，当圣经中记载这类事情时，我们不应仓促断定它们应归于这两类中的哪一类，或者，如果它们属于前者，是否通过精神体的中介作用而产生；同时，至于创造主，即至高而不可言喻的天主圣三的无形而不可变的本性，我们或单纯地、毫无疑问地相信，或除此以外，再以某种程度的理性领悟去理解，知道这本性完全远离并有别于血肉之躯的感官，以及一切可能变坏或变好的易变性，或任何其它可变的事物的性质。\par

% structure: p0017 source=h2 target=subsection reason=relative-heading-rank; base=h2

\subsection{第4章}

12. 关于你的两个问题，我将这些寄给你——一个涉及天主圣三，另一个涉及鸽子，圣神在其中并非以祂自己的本性，而是以象征的形象显现；同样，天主子并非以祂永恒的圣子身份（关于此，圣父说：\BibleQuote{在晓明之前，我已生了你}），而是以祂从童贞玛利亚胎中所取的人性，被犹太人钉在十字架上。请留意，你正有闲暇，我尽管事务极其繁忙，仍能写这么多。然而，我并不认为有必要讨论你信中所提的一切；但对于你希望我解决的这两个问题，我想我已按基督徒的爱德要求，写了足够多的内容，尽管可能仍未满足你强烈的愿望。\par

13. 除了前述加在\WorkTitle{天主之城}前三卷上的两卷书，以及三部圣咏的注释外，我还给圣洁的司铎\Person{热罗尼莫}写了一篇关于灵魂起源的论文，向他请教：当他在写给\Person{马尔切利努斯}（一位虔诚纪念中的人）时，他承认自己主张每个个体在出生时都造有一个新的灵魂；这种观点怎能与教会信仰中那个最确定无疑的信条并行不悖呢？根据那个信条，我们坚信众人都死于\Person{亚当}内\BibleRef{格前 15:22}，并堕入永罚，除非他们因\Person{基督}的\OriginalTerm{恩宠}{grace}而得救——这恩宠借着祂的\OriginalTerm{圣事}{sacrament}，甚至在婴孩身上也发生效力。此外，我还写信给同一个人，询问他对\Person{雅各伯}所说\BibleQuote{谁若遵守全部法律，但只触犯了一条，就算全犯了}这句话应如何理解。在这封信中，我也陈述了自己的意见；而在那封关于灵魂起源的信中，我只询问了他的意见，将此问题交由他判断，同时进行了一定程度的讨论。我给\Person{热罗尼莫}写这些信，是因为我不想错过一个通信的机会，这个机会是由一位非常虔诚且好学的年轻司铎\Person{奥洛西乌斯}提供的。他纯粹出于对圣经的炽热热忱，从\Place{西班牙}最遥远的地方，也就是大洋之滨，来到我们这里，而我劝他从我们这里再前往\Person{热罗尼莫}那里。为了回答\Person{奥洛西乌斯}本人就困扰他的\OriginalTerm{普里西利安派异端}{heresy of the Priscillianists}和教会未曾接受的\Person{奥力振}的某些观点所提出的问题，我尽己所能，以简洁清晰的方式写了一篇中等篇幅的论文。我还写了一本相当有分量的书，驳斥\OriginalTerm{白拉奇异端}{heresy of Pelagius}，我不得不这样做，因为有些弟兄被他引入歧途，接受了他那否认\Person{基督}恩宠的致命错误。如果你想要所有这些作品，就派人来全部抄写给你。然而，请允许我专心研究和口授我的书记员那些许多人所迫切需要的、在我看来比你的那些只有极少数人感兴趣的问题更为优先的事情。\par

\SourceRecord{Translated by J.G. Cunningham.；Nicene and Post-Nicene Fathers, First Series；Vol. 1.；Edited by Philip Schaff.；Buffalo, NY: Christian Literature Publishing Co.,；1887.；<http://www.newadvent.org/fathers/1102169.htm>.}

% structure: p0001 source=h1 target=section reason=page-title

\section{\Person{热罗尼莫}致\Person{奥斯定}书（公元四一六年）}

\Emph{致\Person{奥斯定}，我至虔敬的主、父亲，堪当我至深情与崇敬者，\Person{热罗尼莫}在基督内致候。}\par

1. 那位可敬的人，我的弟兄，也是尊驾的儿子，\OriginalTerm{司铎}{presbyter} \Person{奥罗西乌斯}，我一方面因他本人，另一方面因遵从你的请求，已经接纳了他。但这个最为艰难的时刻已然临到我们身上，于此期间，我发觉自己宁可缄默，也不愿发言，这样我们的学业遂告中止，以免激起\Person{阿皮乌斯}所谓的\Quote{狗辩之才}。因此之故，我目前还未能对那两部题献给我的、博大精深且闪烁着华丽雄辩光彩的著作——我原本是会作出答复的——作出答复；并非我认为书中所述有任何需要纠正之处，而是因为我记念真福宗徒论及人们判断之不同时所说的话：\BibleQuote{各人对自己的心意应有充分的确信}。的确，凡能就那些主题加以讨论的，以及凡能以高超天赋从圣经的泉源中汲取的有关它们的一切，都已在你的这些信函中，通过你的立场得到了陈述，并借由你的论证得到了阐明。然而我恳求尊座，容我略略赞赏你的天赋。因为在你们彼此之间的任何讨论，我们双方所追求的目标都是学问上的进步。可我们的对手，尤其是异端者，倘若看到我们持守不同的意见，必将以诽谤之词攻击我们，说我们的分歧源于彼此嫉妒。但就我而言，我已下定决心要爱你，仰慕你，尊敬你，钦佩你，并将你的意见当作我自己的意见加以维护。此外，我在新近发表的一篇对话录中，也曾以恰当的言辞提及你的真福。故此，我们更当致力于为教会除去那最为邪恶的异端；这异端总是佯装悔改，为要在我们的教会中享有施教的自由，而不致被逐出与铲除——如果它在光天化日之下显露其真实面目，这种结局便是在所难免的了。\par

2. 你虔敬而可敬的女儿们，\Person{欧斯托基}和\Person{保拉}，继续过着配得上自己的出身与你的劝勉的生活，她们特别问候你的真福；与我们一同劳苦服事主、我们救主的全体弟兄也一同问候。至于\OriginalTerm{圣司铎}{holy presbyter} \Person{菲尔穆斯}，去年我们已派遣他前往处理\Person{欧斯托基}和\Person{保拉}的事务，先往\Place{拉文纳}，随后又赴\Place{阿非利加}与\Place{西西里}，我们猜想他如今正淹留于\Place{阿非利加}的某地。我恳请你代我向与你交往的圣者们致以敬礼。我亦已将一封我致\OriginalTerm{圣司铎}{holy presbyter} \Person{菲尔穆斯}的信交由你转递；若此信寄达你处，我恳请你费心转交于他。愿主基督保佑你平安，并记念我，我至虔敬的主、至可称颂的父亲。\par

\Emph{（附言。）}我们在此省中深受通晓拉丁语的圣职人员严重匮乏之苦；这就是我们无法遵从你的指示的原因，尤其是关于那附有区分星号与横线号的\OriginalTerm{七十贤士译本}{Septuagint}的版本；因为由于某人的不诚实，我们此前劳苦所得的大部分成果已经遗失了。\par

\SourceRecord{Translated by J.G. Cunningham.；Nicene and Post-Nicene Fathers, First Series；Vol. 1.；Edited by Philip Schaff.；Buffalo, NY: Christian Literature Publishing Co.,；1887.；<http://www.newadvent.org/fathers/1102172.htm>.}

% structure: p0001 source=h1 target=section reason=page-title

\section{第173封信（公元416年）}

\Emph{致\OriginalTerm{多纳徒派}{Donatist Party}\OriginalTerm{司铎}{Presbyter}\Person{多纳图斯}，\OriginalTerm{天主教会}{Catholic Church}\OriginalTerm{主教}{Bishop}\Person{奥古斯丁}致候。}\par

假若你能看见我内心的忧伤，以及我对你救恩的关切，或许你会怜悯你自己的灵魂，做中悦天主的事，听从那不属于我们而是属于祂的话；不再让\WorkTitle{圣经}只停留在你的记忆中，却将它拒于你的心灵之外。你愤怒，是因为你正被拉向救恩，尽管你已把许多我们的基督徒同道引向毁灭。除了命人将你逮捕，带到官长面前并看管起来，以免你丧亡之外，我们何曾命令过别的呢？至于你身体受伤，这你应该责怪你自己，因为你不愿骑上那即刻为你带来的马，反而猛烈地摔到地上；正如你清楚知道的，与你一同被带来的同伴，并没有像你那样伤害自己，而是安然无恙地抵达。\par

然而，你认为连我们所做的那些事也不该做，因为在你看来，无人应被强迫去行善。因此，留心宗徒的话：\BibleQuote{谁若想望主教的职分，便是想望一件善事。}然而，为使主教职分为许多人接受，人们违背他们的意愿被抓住，受到强求规劝，被关押拘禁，被迫忍受许多他们不悦的事，直到在他们身上找到甘愿承担这善工的心意。那么，更应当如何强拉你们脱离那有害的错谬——在其中你们还是自己灵魂的仇敌——并使你们认识真理，或在认识后选择真理，不仅为让你们稳妥而有益地持守你们职位的尊荣，更为使你们免于悲惨地丧亡！你说天主赐给了我们\OriginalTerm{自由意志}{free will}，因此，即使为善，人也不应受强迫。那么，我上面提到的人，又为何被强迫行善呢？所以，你要留意你所不愿思考的一些事。善的意志以怜悯之心努力追求的目标，乃是确保恶的意志得到正确的引导。因为谁不知道，一个人受罚，没有别的原因，只因他恶的意志应得此罚，而没有一个不具有善的意志的人能得救？然而，这并不意味著，被爱的人就应当被残忍地弃之不顾，任由他们放纵恶的意志而不受惩罚；反之，只要力所能及，就应当阻止他们作恶，并强迫他们行善。\par

假若恶的意志总该任其自由，那么为何背逆怨言的以色列人，被如此严厉的惩戒所约束，不得作恶，并被迫进入福地呢？假若恶的意志总该任其自由，为何\Person{保禄}未被听任于那曾用以迫害教会的最悖逆的意志之自由运用呢？为何他被击倒在地，以使他眼瞎，眼瞎以使他改变，改变以使他被差为\OriginalTerm{宗徒}{apostle}，被差以使他为真理的缘故，遭受他曾在错谬中加于他人的同样迫害呢？假若恶的意志总该任其自由，为何\WorkTitle{圣经}教导作父亲的，不仅要藉责备的话纠正执拗的儿子，还要责打他的腰，好叫他受强迫而归顺，能被引向善行呢？\BibleRef{德 30:12}为此缘故，\Person{撒罗满}也说：\BibleQuote{你应用杖责打他，才能救他的灵魂免下地狱。}\BibleRef{箴 23:14}假若恶的意志总该任其自由，为何失职的牧者受到责备？为何对他们说：\BibleQuote{你们没有领回迷路的羊，也没有寻找丧亡的羊。}\BibleRef{则 34:4}你也是属于基督的羊，在你所领受的\OriginalTerm{圣事}{sacrament}中，你带有主的印记，但你却在迷路丧亡。所以，我们不要惹你动怒，因为我们领回迷路的，寻找丧亡的；因为我们宁愿遵行主的旨意——祂命我们强迫你回到祂的羊栈——而不愿顺从迷羊的私意，任由你丧亡。因此，不要再说我听说你常常说的话：\Quote{我宁愿如此迷路；我宁愿如此丧亡；}因为我们最好尽己所能，绝对不容许你迷路丧亡。\par

那天你投井自尽，无疑你这样做是出于自己的\OriginalTerm{自由意志}{free will}。但天主的仆人若任由你这恶的意志结出果实，而不从死中救出你，那将是多么残忍！谁不理所当然地责备他们？谁不理所当然地谴责他们毫无人性？然而你，凭自己的\OriginalTerm{自由意志}{free will}，投水自溺，他们却违背你的意愿，将你从水中拉出，免得你溺死。你的行为随从己意，却是趋向毁灭；他们对待你违背你意，却是为保全你。因此，如果单是肉身的安危就应如此谨慎地守护，以致爱邻人者有责任即使违其己愿也要保其身于无害，那么，对于那灵魂的健康——失去它之可怕的后果乃是永死——这责任该何等更大！同时我应指出，你若在试图自杀时死去，你不是暂死，而是永死，因为即使有人强迫你——不是叫你接受救恩，不是叫你进入教会的和平、基督身体的合一、神圣不可分割的\OriginalTerm{爱德}{charity}，而是——叫你遭受某些恶事，你也绝不可夺去自己的生命。\par

5. 请思考神圣的圣经，尽你所能地加以考查，看看是否有任何一位义人和信徒曾这样做过，即便是遭受了最沉重的凶恶，而这些凶恶是有人试图迫使他们不去往永生——我们正迫使你走向的永生——而是去往永死。我听说你声称宗徒\Person{保禄}暗示自杀是合法的，当他说\BibleQuote{虽然我舍身被焚烧}时，\BibleRef{格前 13:3}你以为因为他正在列举所有没有爱德就毫无益处的善事，诸如说人和天使的万语，以及一切奥秘、一切知识、一切预言，并将自己的所有分施给穷人，他就打算在这些善事中包括自行取死的行为。但是请仔细留意并学习，圣经说任何人都可以舍身被焚烧，是在何种意义上说的。这肯定不是说任何人被追杀的敌人困扰时，可以投火自尽，而是说当他被迫在选择行恶与受苦之间做出抉择时，他应拒绝行恶而非拒绝受苦，从而将自己的身体交予施刑者，正如那三个人所行的一样，他们被迫朝拜金像，而逼迫他们的人威胁说，若不服从就将他们投入烈火窑中。他们拒绝朝拜那像：他们并未自己投入火中，然而圣经上却记载说他们\BibleQuote{交付了自己的身体，宁死也不愿侍奉或敬拜除自己天主以外的任何神明。}，\BibleRef{达 3:28}宗徒所说的\BibleQuote{如果我舍身被焚烧}，正是这个意义。\par

6. 也要留意接下来所说的：\BibleQuote{我若没有爱德，为我毫无益处。}你正是蒙召去获得这爱德；正是这爱德阻止你丧亡。然而，你竟然以为，通过自己的行为将自己一头栽入毁灭之中，在某种程度上对你有利；即便你是死于他人之手，只要你与爱德为敌，那就对你毫无益处。不仅如此，你若被排除在教会之外，与合一的奥体及爱德的纽带隔绝，即便你为基督之名被活活烧死，也必遭受永罚；因为这就是宗徒的声明：\BibleQuote{我若舍身被焚烧，却没有爱德，为我毫无益处。}因此，让你的心神回归理性的反思与清醒的思想吧；仔细想一想，你所蒙召走向的是否是谬误与不敬虔；若你仍这样想，就为了真理的缘故，耐心忍受一切艰辛吧。然而，如果事实却是你生活在谬误与不敬虔之中，而你所蒙召走向的教会中却存有真理与虔敬，因为在那里有基督徒的合一，和圣神的\OriginalTerm{爱德}{charitas}，你为何还要继续做自己的仇敌呢？\par

7. 为此目的，主的仁慈安排使我们和你们的主教们在\Place{迦太基}举行了一次会议，这次会议举行了多次集会，参加人数众多，并以最有序的方式就我们彼此分离的缘由共同进行了讨论。那次会议的记录被记了下来；我们的署名也附在了记录上：请阅读它，或让别人读给你听，然后选择你所倾向的一方。我听说你曾说过，如果我们在讨论中略去你们主教们所说的这番话：\Quote{任何案件都不妨碍对另一案件的审理，任何人也不牵累他人的立场。}那么你可以在某种程度上与我们讨论那份记录中的陈述。你希望我们略去这些话，但在这些话中，尽管他们自己不知道，真理却借着他们说了出来。你确实会说，他们在这一点上犯了错误，因考虑不周而陷入错误的见解。但我们断言，他们在这里说了真实的话，并且我们很容易就能从你身上证明这一点。因为关于你们自己的这些主教，他们是整个\Person{多纳图斯}派推选出来的，条件是他们应充当代表，且其余所有人都应认定他们所做的一切都是可接受且令人满意的，然而，你却拒绝让他们因你所认为的轻率、错误的言论而牵累你的立场，在这种拒绝中，你证实了他们所说的真理：\Quote{任何案件都不妨碍对另一案件的审理，任何人也不牵累他人的立场。}同时，你应当承认，如果你拒绝让代表你们众多主教的这七个人的联合权威来牵累\Place{穆图根纳}的司铎\Person{多纳图斯}，那么，即便在\Person{凯西利安}身上发现了某些邪恶，一个人便牵累基督的整个合一奥体——教会——的立场，这是无比不合理的；教会并非局限于\Place{穆图根纳}一个村庄，而是遍布整个世界。\par

8. 但是，请看，我们正按你所愿望的做；我们对待你们，就好像你们的主教从未说过：\Quote{彼此的案件互不妨碍审理，个人也不牵连他人的立场。} 如果你能，你就去发现他们应当如何回应，当有人针对他们提出\Person{普里米安}的案件和此人本身时，他们原本该怎样回答，而不是像现在这样回应。因为\Person{普里米安}虽然曾与其他主教一同宣判，谴责了那些曾谴责他的人，但他后来仍将那些由他谴责和痛斥的人恢复到原有的尊位，并选择承认和接受这些人在他们\Quote{已死}时所施行的圣洗，而非鄙视和拒绝（因为在那著名的\Place{巴盖}大公会议法令中论到他们说：\Quote{海岸上满是死人}），他这样做便推翻了你们惯常依据对经文：\Quote{\OriginalTerm{由死人施洗的人，他的圣洗对他有何益处？}{Qui baptizatur a mortuo quid ei prodest lavacrum ejus?}} 的曲解所建立的论点。因此，如果你们的主教没有说：\Quote{彼此的案件互不妨碍审理，个人也不牵连他人的立场。}，他们便不得不在\Person{普里米安}的事上认罪；但是，他们说了这话，就宣布了天主教会（正如我们提到的）在\Person{凯奇利安}的事上是无罪的。\par

9. 然而，请你阅读其余的部分，仔细审查。留意他们可曾成功地证明对\Person{凯奇利安}本人的任何恶的指控，正是他们试图通过此人来牵连教会的立场。留意他们难道不是反倒提出了许多对他有利的证据，并且用他们提出并宣读的大量摘录，驳斥了自身的主张，从而证实了他的案件是正直的吗？请你阅读这些记录，或者让人读给你听。仔细思考整件事，慎重掂量，然后选择你应当遵从哪一边：是在基督的平安中、在天主教会的合一中、在弟兄的爱中分享我们的喜乐；还是在邪恶的纷争、多纳徒派和亵渎的裂教中，继续忍受那因我们爱你而不得不采取的措施带给你的烦恼。\par

10. 我听说，你曾提及并经常引用福音中记载的事实：七十门徒离弃主而去，他们被任由在这邪恶不虔的背叛中自行选择，而主对那唯独留下的十二人说：\BibleQuote{难道你们也愿走吗？} \BibleRef{若 6:67} 但你却忽略了一件事：那时教会只是从刚播下的种子中才开始生长，在她身上尚未应验这预言：\BibleQuote{众王都要崇拜祂，万民都要事奉祂；} 而正是随着这预言更充分的应验，教会才拥有了更大的权力，使她不仅能邀请，甚至能强迫人接受美善。我们的主当时意在说明这一点，因为尽管祂拥有巨大的权能，祂却宁愿彰显自己的谦卑。祂也在宴席的比喻中以足够清晰的方式教导了这一点。在这比喻中，家主派人去邀请客人，被邀的人拒绝前来，他便对仆人说：\BibleQuote{你快出去，到城中的大街小巷，把那些贫穷的、残废的、瞎眼的、瘸腿的，都领到这里来。仆人说：主，已经照你的吩咐办了，可是还有空位子。主人对仆人说：你出去，到大道上以及篱笆边，强拉人进来，好坐满我的屋子。} \BibleRef{路 14:21} 现在，请注意，论到那些先来的人，是说\BibleQuote{领他们进来}；却没有说\BibleQuote{强拉他们进来}——这指的是教会的起初阶段，那时她正在成长，直至获得力量能强迫人进来。因此，当教会的力量和范围都加强后，理当强迫人进入永生救恩的宴席，所以比喻中随后又说：\BibleQuote{仆人说：主，已经照你的吩咐办了，可是还有空位子。主人对仆人说：你出去，到大道上以及篱笆边，强拉人进来。} 因此，如果你安然行走，却远离这永生救恩和教会神圣合一的宴席，我们就会如同在\BibleQuote{大道}上遇见你；但因为你对我们人民施加种种伤害和残暴，你就像满身荆棘和粗糙，我们如同在\BibleQuote{篱笆}中遇见你，于是我们强迫你进来。被强迫的羊被赶往它不愿意去的地方，但进入之后，它便自愿地在被带到的牧场上进食。所以，约束你悖逆和顽抗的心神吧，好使你在基督的真教会中找到救恩的宴席。\par

\SourceRecord{Translated by J.G. Cunningham.；Nicene and Post-Nicene Fathers, First Series；Vol. 1.；Edited by Philip Schaff.；Buffalo, NY: Christian Literature Publishing Co.,；1887.；<http://www.newadvent.org/fathers/1102173.htm>.}

% structure: p0001 source=h1 target=section reason=page-title

\section{第180封信（公元416年）}

\Emph{致\Person{奥切安}，他在基督内当受敬爱的我主与兄弟，在基督奥体中备受尊敬者，\Person{奥斯定}致候。}\par

1. 我同时收到了你的两封信，在其中一封信中，你提到了第三封信，并说你已在那两封信之前寄出了它。我不记得曾收到过这封信，或者更确切地说，我认为我记忆的见证是：我并未收到它；但对于已收到的信，我深深感谢你对我的厚意。若非一直被接连不断、急需处理的其他事务所缠累，我本应立即回信。现在既然从这些事务中偷得片刻闲暇，我选择了寄出一些回复，尽管不完美，也好过对一位如此真诚友善的朋友长久沉默，而一言不发比说太多更令你烦恼。\par

2. 我早已知道\Person{圣热罗尼莫}关于灵魂起源的看法，并且曾在你的信中读到你所引述的他的书中的话。对于某些人因这个问题而感到困惑的难题——\Quote{天主怎能公平地将灵魂赋予通奸者所生的后代？}——并不会使我为难，因为即使是他们自身的罪，更不用说他们父母的罪，都不能损害那些度圣善生活、归向天主、并在信德与虔诚中生活的人。真正困难的问题是：如果确实是在每个孩子出生时，一个新的、从无中创造的灵魂被赋予他，那么，对于那些天主确切知道在达到理智年龄之前、在能够知晓或理解是非之前就将脱离肉身、且未受洗的无数婴孩的灵魂，怎能被那位\BibleQuote{绝无不义}\BibleRef{罗 9:14}者公正地交付永死呢！关于这个主题，无需再多说，因为你已知道我想说什么，或更确切地说，我现在不想说什么。我想，对于一个明智的人来说，我所说的已经足够了。不过，如果你曾读过、或从\Person{圣热罗尼莫}的口中听到、或在默想这个难题时从主那里领受了任何能解决它的见解，我恳求你，请告诉我，以便我承认我欠你更多的情谊。\par

3. 至于说谎在任何情况下是否可视为正当且权宜的问题，你认为应当以主论及世界末日那日子和时辰的话——\BibleQuote{子也不知道。}\BibleRef{谷 13:32}——为例来解决。当我读到这解释时，不禁为你的精巧构思所吸引；但我绝不认为一种比喻的表达方式可以正确地被称为谎言。因为，称某日为\Quote{喜乐的}，因为它使人喜乐；或说羽扇豆为\Quote{苦涩的}，因其苦味使尝者面容现出苦涩；或者说天主知道某件事，其实是指祂使人知道（你自己也引用了天主对\Person{亚巴郎}说的这句话：\BibleQuote{现在我知道你敬畏天主。}\BibleRef{创 22:12}）——这些都绝不是谎言，正如你自己容易看到的。因此，当\Person{圣依拉略}通过这种隐晦的比喻性语言来解释主这一隐晦的陈述，说我们应该将基督宣称祂不知道那日子的话理解为：祂只是通过隐藏它，使别人不知道它；他这样解释并非为这陈述辩解为情有可原的谎言，而是指出它根本不是谎言，这不仅是藉着与这些常见的修辞格比较来证明，更是借与日常谈话中人人熟悉的隐喻相比较来证明。有谁会责备说\Quote{收割的田亩\Emph{波浪滚滚}}和\Quote{儿童\Emph{花朵般绽放}}的人是在说谎呢？只因他并未在这些事物中看到这些词语字面上所指的波浪和花朵。\par

4. 此外，以您的才华和学识，不难看出这些隐喻表达与宗徒的话何其不同：\BibleQuote{我一见他们的行为，不合福音的真理，就在众人前对\Person{伯多禄}说：\InnerQuote{你既是犹太人，竟按外邦人的方式，而不按犹太人的方式生活，怎么敢强迫外邦人犹太化呢？}}\BibleRef{迦 2:14}。这里没有比喻语言的隐晦；这些都是明确陈述的字面话语。显然，这位外邦人的大宗徒，在向那些\BibleQuote{他愿为他们再受产痛，直到基督在他们内形成}\BibleRef{迦 4:19}的人说话，并且曾郑重呼求天主为他的话作证说：\BibleQuote{我给你们写的，在天主前保证，我并没有说谎}\BibleRef{迦 1:20}，在以上所引的话语中，他所肯定的要么是真，要么是假；如果他说的是假话——愿天主禁止这样想——您会看到随之而来的后果；而\Person{保禄}亲自对他自己诚实的肯定，以及宗徒\Person{伯多禄}所树立的奇妙谦卑的榜样，都当警示您远离这种想法。\par

5. 但何须多言？可敬的教父\Person{热罗尼莫}与我曾在互通的书信中，详尽讨论了这个问题；而在他以\Emph{Critobulus}之名发表、反驳\Person{白拉奇}的最新著作中，他关于那项讨论以及宗徒的话所持的见解，与我本人遵循\Person{圣西彼廉}的观点所持的看法相同。关于灵魂的起源问题，我认为有合理的探究空间——不是说奸淫父母所生子如何被赋予灵魂，而是指无辜者竟被定罪（但愿天主不许）。您若从像\Person{热罗尼莫}这样品格崇高、卓越的人那里学到过什么，足以解答对此困惑的人，我恳请您不要拒绝告知我。在您的书信往来中，您表现得如此博学且和蔼可亲，以致能与您通信交流实在是一种殊荣。我请求您不要迟延，将同一\OriginalTerm{天主之人}{man of God}的那本书寄来，即司铎\Person{奥罗修斯}带来并交给您抄写的那本；书中论述肉身复活的方式，据说值得极高度的赞美。先前我们没有请求，是因知道您既要抄写又要校订；但我想现在这两项工作您已有充足时间完成。愿您在天主内生活，并记念我们。\par

\EditorialNote{关于致伯爵\Person{卜尼法斯}的第\Emph{一八五}号信函的翻译，其中包含对多纳特派分裂的详尽历史，请见\WorkTitle{反多纳特派著作}。}\par

\SourceRecord{Translated by J.G. Cunningham.；Nicene and Post-Nicene Fathers, First Series；Vol. 1.；Edited by Philip Schaff.；Buffalo, NY: Christian Literature Publishing Co.,；1887.；<http://www.newadvent.org/fathers/1102180.htm>.}

% structure: p0001 source=h1 target=section reason=page-title

\section{书信 185}

\BibleRef{路 24:46}圣经中的证词不可胜数，我无需将它们全都辑录于此书中。在这一切证词中，正如主基督彰显出来了——无论是按祂的天主性，在这方面祂与圣父同等，正如\BibleQuote{在起初已有圣言，圣言与天主同在，圣言就是天主}；还是按祂所取的人性的卑微，因此\BibleQuote{圣言成了血肉，寄居在我们中间}——同样，祂的教会也彰显出来了，不仅限于\Place{非洲}，如同他们极其无耻地在自己虚妄的疯狂中斗胆断言的那样，而是遍布全世界。\par

4. 因为他们宁要自己的争辩，而不要圣经的证词，因为，关于\Place{迦太基}教会的前任主教\Person{塞西里安}，他们虽对他提出指控，却过去不能、现在也不能证实这些指控，于是他们就从天主教会——即从万民的合一中——分裂出去。虽然，即使他们所提出的针对\Person{塞西里安}的指控是真实的，并且最终能向我们证明，那么，尽管我们甚至可以在他死后宣布对他的绝罚，我们仍然一定不能为了任何人的缘故而离开那建立在神圣证言之上、而非好争讼意见所捏造的教会，因为\BibleQuote{倚赖上主胜过信赖世人}。因为，我们不能允许，如果\Person{塞西里安}犯了错——我作此假设，并无损于他的正直——这样，基督就因此丧失祂的产业。一个人很容易相信关于他人的事情，不论是真是假；但是，这标志着一种肆无忌惮的无耻行径：由于针对某一个人所提出的指控——而这些指控你们并不能在全世界面前证实其真实性——竟希望谴责整个世界的共融。\par

5. \Person{塞西里安}是否是由那些交出过圣经的人所祝圣的，我不知道。我没有看见，只是从他的敌人那里听说的。这并没有在天主的法律中，或在先知的宣讲中，或在\BibleBook{}{圣咏}的神圣诗歌中，或在任何一位基督宗徒的著作中，或在基督自己的话语中向我宣告。然而，所有个别的圣经证据都一致宣告，与多纳图派不相通共融的教会，是遍布全世界的。天主的法律宣告：\BibleQuote{地上万民都要因你的后裔蒙受祝福}\BibleRef{创 26:4}。上主藉先知的口说：\BibleQuote{从日出到日落，必有纯洁的祭献献于我的名，因为我的名在列邦中必受尊崇}\BibleRef{拉 1:11}。上主藉\BibleBook{}{圣咏}的作者说：\BibleQuote{祂的权柄由这海到那海，从大河直到地极}。上主藉祂的宗徒说：\BibleQuote{福音已传至你们，就如它在全世界一样，并且结出果实}\BibleRef{哥 1:6}。天主子亲口说：\BibleQuote{你们要在\Place{耶路撒冷}、\Place{犹太}全境、\Place{撒玛黎雅}，直到地极，为我作证人}\BibleRef{宗 1:8}。\Place{迦太基}教会的主教\Person{塞西里安}受到人的争讼指控；而建立在万民中的基督教会，则受到天主圣言的推荐。纯正的虔敬、真理与爱德禁止我们接纳那些不在具有天主证言的教会中的人作证反对\Person{塞西里安}；因为凡不遵循天主证言的人，已丧失了作为人所应有的证言的份量。\par

% structure: p0005 source=h2 target=subsection reason=relative-heading-rank; base=h2

\subsection{第二章}

6. 我还要补充说，他们自己，正是通过将这事件作为诉讼，将\Person{塞西里安}的案件呈交\Person{君士坦丁}皇帝裁决；并且在主教们宣布判决后，发现无法击垮\Person{塞西里安}，便以最坚决的迫害精神，将他亲自带到上述皇帝面前受审。这样，他们自己率先做了他们指责我们所做的事，以便欺骗无学识的人，说基督徒不应当向基督的敌人请求基督徒皇帝的任何援助。而这正是他们在我们同时在\Place{迦太基}举行的会议上不敢否认的：不，他们甚至胆敢自豪地宣称，他们的祖辈曾向皇帝提出过对\Person{塞西里安}的刑事起诉；此外还附加一个谎言，说他们在那里击败了他，并使他被判了罪。那么，他们怎么能不是迫害者呢？既然他们藉着指控迫害\Person{塞西里安}，被他击败后，却以一种最无耻的谎言为自己谋求虚假的荣耀；不仅不以为耻，反而引以为荣，以为只要能证明他们的祖辈指控\Person{塞西里安}导致了定罪，便能增加他们的声望。但是，关于在会议本身中他们如何节节败退，由于记录极其冗长，而你们正忙于\Place{罗马}和平所必需的其他事务，向你们宣读这些记录恐怕是一件繁重的事；也许可以向你们宣读一份摘要，我相信这份摘要在我弟兄及同为主教的\Person{奥普塔特}手中；或者，如果他没有副本，他可以很容易地从\Place{西提法}教会获得一份；因为我可以相信，就连那一卷书，由于其冗长，在你们众多事务的重压下，也足以令你们感到厌倦了。\par

7. 因为\OriginalTerm{多纳徒派}{Donatists}遭遇了与控告圣\Person{达尼尔}者相同的命运。\BibleRef{达 6:24} 正如狮子转过来攻击他们，同样，多纳徒派想用来压垮无辜受害者的法律，也转过来攻击他们自己；所不同者，由于基督的仁慈，那些看似反对他们的法律，其实是他们最真实的朋友；因为藉着这些法律的运作，他们中许多人已经并正在日日改过，并为得以改过自新、脱离毁灭性的疯狂而归感谢于天主。那些从前怀恨的人，如今充满了爱；如今他们恢复了神智清醒，就以他们在疯狂时憎恨这些法律的同样热忱，庆幸这些最有益的法律曾加于他们；并且他们对那些尚停留于我们从前光景中的人，怀有与我们一样的炽热爱情，切望我们同样努力，拯救那些他们曾与之同趋丧亡的人。因为医生对于狂暴的疯子是恼人的，父亲对于不守规矩的儿子也是恼人——前者因约束而恼人，后者因所施的惩戒而恼人；然而二者都是出于爱而行。但如果他们忽略自己的责任，任由他们丧亡，这种错误的仁慈，更真是残忍。因为马和骡子没有理智，却用尽脚踢牙咬的力气，抗拒那些为治愈它们而处理伤口的人；然而，人们虽然常冒着被它们的牙齿和蹄子伤害的危险，有时也确实受伤，却仍不放弃它们，直到通过治疗过程所带来的苦痛和烦扰，使它们恢复健康——那么，人对同伴，或兄弟对兄弟，更不该放弃，以免他们永远丧亡；因为人现在已能理解，在自己改过时所得恩赐何其浩大，而那时他却正在抱怨遭受迫害。\par

8. 因此，正如宗徒所说：\BibleQuote{所以，我们一有机会，就应向众人行善，尤其应向有同样信德的家人。}\BibleRef{迦 6:9} 照样，所有的人都应蒙召获得救恩，所有的人都应被领回，脱离丧亡之路——能藉天主教宣讲者的讲道而领回的，就藉讲道；能藉天主教君主的谕令而领回的，就藉谕令；有些人通过服从天主的警告者，有些人通过服从皇帝的命令者。再者，当皇帝们颁行支持谬误、反对真理的恶法时，那些保持正统信仰的人便蒙悦纳，若坚忍到底，必获得荣冠；但当皇帝们颁行维护真理、反对谬误的善法时，那些反抗这些法律的人便被恐惧所慑，而明理的人则改过迁善。因此，凡拒绝服从那些反对天主真理的皇帝法律的人，为自己赢得丰厚的赏报；但拒绝服从那些维护真理的皇帝法律的人，为自己赢得重大的惩罚。因为在先知时代，那些在处理天主子民时，既不禁止也不废除违犯天主命令的法规的君王，都受到谴责；而那些禁止并废除这些法规的君王，便被誉为超越常人。\Person{拿步高}王，当他还是一个偶像的奴仆时，曾颁下不虔敬的法律，命令人朝拜某一个偶像；但那些拒绝服从他不虔敬的命令的人，却行了虔敬而忠信的事。而正是这同一位君王，因天主的奇迹而归正后，却为真理颁下了虔敬且堪受赞美的法律，命令凡毁谤真天主，即\Person{沙得辣客}、\Person{默沙客}和\Person{阿贝得乃哥}的天主的人，本人及其全家都该全然灭亡。\BibleRef{达 3:97-98} 如果有人不服从这条法律，并且受了应得的惩罚，他们原可以说这些人所说的话，说自己因受君王所颁法律的迫害而成为义人；如果他们像这些分裂基督肢体、轻贱基督圣事、因皇帝们为维护基督的合一而颁行的法律阻止他们行这些事，便以受迫害自夸，虚假地夸耀自己的无罪，并想从人得到主不会赐给他们的殉道光荣的人一样疯狂，他们必定也会这样说。\par

9. 然而，真正的殉道者，乃是主所说的那些人：\BibleQuote{为义而受迫害的人是有福的。}\BibleRef{玛 5:10} 所以，不是因为自己的不义，并因不虔敬地在基督徒的合一中制造分裂而受迫害的人，而是为义而受迫害的人，才是真正的殉道者。因为\Person{哈加尔}也曾受\Person{撒辣}迫害；\BibleRef{创 16:6} 而在那种情形下，迫害人的是义人，受迫害的则是不义的人。难道我们要将哈加尔所受的这种迫害，与圣\Person{达味}遭受不义的\Person{撒乌耳}迫害相比拟吗？显然，本质的差别不在于他所受的苦，而在于他是为义而受苦。主自己也与两个强盗同钉在十字架上；\BibleRef{路 23:33} 然而，那些同受苦难的人，却因动机的不同而被区分开来。因此，在\WorkTitle{圣咏}中，我们必须把这样一节解释为指着那些切愿与假殉道者区分的真殉道者所说的：\BibleQuote{天主，求你为我伸冤，向不虔诚的国为我辩护。} 他不说\InnerQuote{求你辨明我的刑罚}，而是说\InnerQuote{求你辨明我的理}。因为不虔敬者的刑罚可能相同；但殉道者的理却总是不同的。还有这话适合出自他们的口：\BibleQuote{他们无故地迫害我，求你扶助我；} 圣咏作者在这话中声称，他有权利获得正义的援助，因为他的仇敌无故迫害他；因为如果他们对他的迫害是正当的，他所应得的就不是援助，而是纠正。\par

10. 然而，如果他们以为没有一个人能以使用暴力为义——正如他们在会议中所说的，真教会必须是那遭受迫害的，而不是那施加迫害的——在这种情况下，我就不再坚持我上面所论述的；因为如果事情确如他们所主张的那样，那么\Person{凯基利安}必属于真教会，因为他们的祖先曾迫害他，甚至将控告诉诸皇帝的法庭。因为我们坚持他属于真教会，不仅因为他遭受了迫害，而且因为他为义而受了迫害；而他们背离了教会，不仅因为他们施行了迫害，而且因为他们在不义中施行了迫害。这便是我们的立场。但是，如果他们不追究各人施行迫害或遭受迫害的原因，反而认为一个真基督徒的充足标记就是他不行使迫害，而是遭受迫害，那么毫无疑问，他们便将\Person{凯基利安}归入此定义之下，因为他不曾施行迫害，而是遭受了迫害；同样，他们也把自己的祖先排除在这定义之外，因为祖先们施行了迫害，却没有遭受迫害。\par

11. 然而，我说过，我不再强调此点。但有一点我必须坚持：如果真教会是那真正受迫害的，而非施加迫害的，那么让他们去问宗徒，当\Person{撒辣}迫害她的婢女时，她是哪一个教会的预象。因为宗徒宣称，我们众人的自由之母，\Place{天上的耶路撒冷}，即天主的真教会，已由那位苦待婢女的女人预表了。\BibleRef{迦 4:22} 但是，如果我们进一步探究这故事，就会发现，婢女以她的傲慢迫害\Person{撒辣}，更甚于\Person{撒辣}以严厉对待婢女：因为婢女是对主母行不义在先；主母只是对她的傲慢施加了适当的管教。我再次发问：如果善良圣洁的人从不向任何人施加迫害，而只遭受迫害，那么他们认为圣咏中我们读到的这段话是谁的话呢：\BibleQuote{我追捕我的仇敌且把他们擒获，直到他们被歼灭，我才回还。} 因此，如果我们愿意宣扬或认识真理，那么就有一种不义的迫害，是亵圣者施加于基督的教会；还有一种义的迫害，是基督的教会施加于亵圣者的。所以，她为义而受苦，是有福的；而他们为不义而受苦，是悲惨的。再者，她以爱德的精神施行迫害，他们以愤怒的精神施行迫害；她是为了纠正，他们是为了倾覆；她是为了从错误中召回，他们是为了驱入错误。最后，她迫害并逮捕她的敌人，直到他们厌倦他们虚妄的意见，从而使他们在真理上进步；但他们对我们的美意以恶相报，因为我们为他们的益处，为保障他们的永恒救恩而采取行动，他们甚至图谋剥去我们暂时的安全，如此嗜好杀害，以至于当他们找不到其他牺牲品时，便对自己下手。因为教会的基督徒爱德越是努力把他们从那种毁灭中解救出来，以免他们中任何一人死亡，他们的疯狂就越是设法要么杀害我们，以满足他们残忍的欲望，要么甚至自杀，以免显得他们丧失了置人于死地的能力。\par

% structure: p0012 source=h2 target=subsection reason=relative-heading-rank; base=h2

\subsection{第3章}

12. 但是，那些不熟悉他们的习性的人以为，他们现在之所以自杀，是因为全体民众借着为维护合一而通过的法律，摆脱了他们篡夺统治权的可怕疯狂。然而，那些知道他们在法律颁布之前惯常行径的人，并不对他们的死亡感到惊讶，而是想起他们的品行；特别是他们曾经常成群结队地游行至异教徒最盛大的崇拜仪式上，那时偶像崇拜依然存在——不是为了捣毁偶像，而是为了被偶像崇拜者处死。因为他们若是寻求在合法权威的许可下捣毁偶像，那么万一遭遇不测，还可能有一些理由被视为殉道者；但他们前来的唯一目的，就是在偶像完好无损的情况下，让他们自己遭遇死亡。因为偶像崇拜者中的最强壮青年通常的习惯是，将他们杀死的任何牺牲者祭献给偶像。有些人甚至到了这样的地步：将自己交给他们遇到的任何携带武器的旅客去宰杀，并使用暴力威胁，如果旅客不杀死他们，他们就要杀死旅客。有时，他们还用暴力强迫任何路过的法官，要求由刽子手或其法庭的官吏将他们处死。因此，流传着这样一个故事：有一位法官对他们施了一个诡计，下令将他们捆绑带走，像是要处决一样，从而摆脱了他们的暴力，既未伤害自己，也未伤害他们。再者，他们日常的消遣便是自杀，从悬崖上跳下，或投水，或投火。因为魔鬼教会了他们这三种自杀方式，使得当他们想死，却找不到任何人能威吓用刀杀死他们时，他们就跳下岩石，或投身烈火或漩涡之中。然而，有谁能被认为教给他们这些事，并占据了他们的心呢？难道不是那位曾经向我们的救主本人建议，作为法律所许可的义务，从圣殿的顶端跳下去的吗？\BibleRef{路 4:9} 如果他们心中怀着\Person{基督}作为他们的师傅，他们早就远远摒弃他的建议了。但既然他们反而在心中给魔鬼留了地步，他们要么像那群被魔鬼从山坡赶入海中的猪一样灭亡，\BibleRef{谷 5:13}，要么，若从那种毁灭中获救，并被聚集在我们慈母教会慈爱的怀抱中，他们便如同\Person{主}所拯救的那个孩子一样获得释放；那孩子的父亲带他来给魔鬼附身的孩子治病，说：\BibleQuote{他屡次跌入火中，屡次跌入水中。}\BibleRef{玛 17:14}\par

13. 由此可见，向他们显示了极大的仁慈：借着那些帝国法令的强制力，他们先是被迫违背自己的意愿，从那个教派中被救拔出来——正是在那个教派中，他们因着撒谎的魔鬼的教导，学到了那些邪恶的教义——以便后来在天主教会内得以痊愈，逐渐习惯教会中良好的教导和例子。因为我们现在因着他们信德的热忱和爱德，在\Person{基督}的合一中所赞美的许多人，都满心喜乐地感谢天主，因为他们不再处于那种使他们错把恶事当作善事的错误之中；如果他们不是先被迫，甚至违背自己的意愿，脱离那个不敬虔的团体，他们如今就不会自愿地献上这样的感谢。那么对于那些向我们承认——正如每天都有人这样做——即使在往昔的日子，他们早就渴望成为天主教徒的人，我们又该说什么呢？但他们生活在这样的人群中：在那里，想要成为天主教徒的人因胆怯懦弱而无法如愿，因为如果有人在那个地方为天主教会说一句好话，他和他的全家就会立刻被彻底毁灭。谁会疯狂到否认，通过帝国法令给予援助是正当的，好使他们能从如此大的邪恶中解救出来，同时让那些他们从前所畏惧的人反过来畏惧，并使他们要么因同样的恐惧而改过自新，要么至少，在他们假装改过时，不再迫害那些真正改过的人，而这些人从前是他们不断恐惧的对象？\par

但若他们选择自我毁灭，为阻止那些理应获救的人得解救，并以此试图恐吓解救者虔诚的心，使解救者因担心少数几个自暴自弃者灭亡，便任由其他人丧失脱离毁灭的机会——这些人要么本不愿灭亡，要么本可借强迫手段得救——那么，在此情形下，基督徒爱德的作用是什么？尤其是，那些以暴力和自愿自杀相威胁的人，与有待解救的众多民族相比，数目极少。那么，兄弟之爱的作用又是什么？难道它因惧怕少数人短暂的窑火，便舍弃所有人于地狱的永火？难道它因防备少数人亲手自杀，便任由如此众多的人永远灭亡，而这少数人活着只为阻碍他人得救，他们不许他人按基督的教义生活，期望有朝一日也教导他人像自己现在这样，以其魔鬼般信条所灌输的方式，亲手寻死，使邻居恐惧？或者，兄弟之爱宁可尽力拯救一切能救之人，即使无法拯救之人迷失在自己的愚妄中灭亡？因为它热切盼望所有人都能活，但它更努力避免所有人死亡。但感谢主，在我们中间——虽非处处都如此，但在绝大多数地方——并在非洲其他地区，天主教会的平安正在且已经扩展，这些疯狂之徒无一被杀。然而，那些可悲的暴行发生在一些地方，那里有一群极端狂暴而无用的人，甚至在古时便惯于此行径。\par

% structure: p0016 source=h2 target=subsection reason=relative-heading-rank; base=h2

\subsection{第四章}

实际上，在那些信奉天主教信仰的皇帝们实施这些法律之前，基督的和平与合一的教义已开始逐渐传播，人们甚至从\OriginalTerm{多纳图派}{Donatists}中归向它，按照各人学习增多、更为愿意、更能自己作主的程度；然而，与此同时，在多纳图派中，成群自暴自弃的人正以最肆无忌惮的疯狂精神，出于种种理由扰乱无辜者的和平。哪个家主不被迫提心吊胆地惧怕自己的仆人，如果他投靠了多纳图派的庇护？谁敢以惩罚威胁那图谋自己毁灭的人？谁敢向那些寻求他们援助或保护的挥霍财物者或任何债务者追索债务？在鞭打、火烧和立即处死的威胁下，所有能使最恶的奴隶定罪的文件都被销毁，使他们得以自由离去。从债务人勒索到的借据被退还。任何蔑视他们严厉言辞的人都被更严厉的殴打强迫去做他们所愿的。那些触怒他们的无辜者的房屋要么被夷平，要么被焚毁。有些家庭出身尊贵、受过良好教育的家长，在他们暴力之后，被半死不活地拖走，或被绑在磨上，像最下贱的驮畜一样被鞭打着转动磨盘。因为国家权力施行的法律对他们有何效力呢？哪个官员胆敢在他们面前喘口气？哪个执行官曾向他们不愿清偿债务的人追索债务？谁曾试图为他们屠杀中杀死的人报仇？除非，的确，他们自己的疯狂向他们报复，当一些人激起他人对自己动剑，迫使他人在立即死亡的恐惧下杀害他们时；另一些人投身各处悬崖；又一些人投水；还有一些人纵火，在种种情形下自取死亡，将自己的生命作为献给死者的祭品，亲手对自己处以刑罚。\par

这些暴行被许多深陷同一迷信异端的人视为恐怖；因此，当他们以为仅对自己不满于这些行为便足证自己的清白时，天主教徒便以此驳斥他们：如果这些恶行不玷污你的清白，那么你又怎能声称整个基督宗教世界被\OriginalTerm{凯西里安}{Cæcilianus}的所谓罪过玷污了？这些指控要么纯属诽谤，要么至少未能证实。你如何以极不虔敬的行为，将自己与天主教会的合一分离，犹如离开主的禾场，这禾场必需容纳麦子和糠秕，直到最后的簸扬之时，麦子要收在仓里，糠秕却要用火烧尽？\BibleRef{玛 3:12}。这样，有些人被道理说服，归向了天主教会的合一，甚至预备面对那些自暴自弃者的敌意；而人数更多者，虽同样被说服，且渴望同样行，却不敢与这些暴力横行无忌的人为敌，因为他们看到，一些归向我们的人遭受了他们极其残酷的对待。\par

17. 此外，我们还可以补充一点：在\Place{迦太基}本城，同一派别的一些主教，在他们中间制造分裂，并在\Place{迦太基}下层民众中瓜分\Person{多纳图}派的势力，祝圣了一位名叫\Person{马西米安}的执事为主教，以对抗主教，此人无法忍受自己本主教的管理。由于此事使大部分人不悦，他们绝罚了前述的\Person{马西米努斯}，以及曾参与他祝圣的另外十二人，却给了参与同一分裂的其余人一个机会，在指定的日子重返他们的共融。但后来，这十二人中的一些人，以及那些被给予宽限期、却在指定日期之后才回头的其他人，被他们接纳，并且没有撤销他们的神品；他们竟不敢给那些受绝罚的圣职人员在他们的共融之外所施洗的人再次施洗。他们的这一举动，立刻有力地表明了他们自相矛盾，有利于天主教会一方，以至于他们完全哑口无言。当此事被理所当然地四处传扬，以治愈人们灵魂上分裂的毒害，并当天主教会的学者们通过讲道和辩论从各方面指陈：他们为了维护\Person{多纳图}的和平，不仅重新接纳了他们所绝罚的人，且完全承认其神品，甚至竟不敢宣布：那些被他们绝罚或甚至停职的人，在他们的教会之外所施行的圣洗是无效的；然而，他们却违背基督的和平，指责全世界的（教会）因与某些罪人接触而沾染了污点（至于那些罪人是谁，却无关紧要），并因此宣布，就连福音传入非洲之本源的那些教会中所施行的圣洗也是无效的；——看到这一切，许多人开始困惑，并在他们所见大部分明显真理面前羞愧，于是比以往更多的人接受纠正；人们开始从他们的残酷中感受到一种较为自由的呼吸，而且这种趋势在各个方向上都在大幅度扩展。\par

18. 那时，他们确实爆发出如此狂怒，并在仇恨的刺激下如此激动，以至于我们共融下的教会几乎无一能免遭他们的背叛、暴力和毫不掩饰的抢劫；几乎没有一条道路是安全的，能让人们前往宣讲天主教会的和平以对抗他们的疯狂，并通过清晰陈述真理来揭露他们轻率的愚昧。此外，他们还在和解条件上提出严苛要求，不仅对平信徒或任何圣职人员如此，甚至在某种程度上也对某些天主教会的主教如此。因为提出的唯一选择就是要么对真理缄默不语，要么忍受他们野蛮的暴怒。然而，如果他们不讲真理，那么不仅任何人不可能因他们的沉默而得救，更有许多人必定会因误入歧途而遭毁灭；而如果通过他们宣讲真理，再次激起多纳图派的怒气发泄其疯狂，虽然有些人会得救，那些已站在我们一方的人会得到坚固，但软弱者却会因恐惧而再次不敢追随真理。因此，当教会在苦难中被逼到如此困境时，任何人若认为应忍受一切，而非寻求天主的援助——这援助将藉基督徒皇帝的手来施行——他便是没有充分注意到，为忽略这一预防措施，绝无法交出一份好的账目。\par

% structure: p0021 source=h2 target=subsection reason=relative-heading-rank; base=h2

\subsection{第五章}

19. 至于那些不愿自己的不虔敬行为受到正义法律制裁的人所提出的论点，他们说宗徒们从未向地上的君王寻求过此类措施，却未曾考虑到那个时代的不同特征，以及万事各有其时。因为，当先知的话还只是在实现的过程中时——\BibleQuote{外邦人为什么嚣张，万民为什么妄想？地上的君王起来，首领们聚集在一起，反对上主，反对祂的受傅者；}——那时尚未有迹象表明同一圣咏后面稍后所说的话：\BibleQuote{众王！现在你们应当自觉；地上的判官！你们应受教：应以敬畏之情事奉上主，战战兢兢地欢庆。}那么，君王如何以敬畏之情事奉上主呢？除非通过以宗教的热忱，阻止并惩罚一切干犯上主诫命的行为。因为一个人以人的身份事奉天主是一回事，以君王身份事奉天主又是另一回事。他以人的身份事奉天主，是通过度忠信的生活；但他以君王身份事奉天主，则是通过以适当的严厉执行那些规定正义、惩罚悖逆的法律。正如\Person{希则克雅}事奉祂，摧毁了木偶、偶像的庙宇，以及那些违反天主诫命而建立的高丘；\BibleRef{列下 18:4}再如\Person{约史雅}事奉祂，他也做了同样的事；\BibleRef{列下 23:4}又如\Place{尼尼微}王事奉祂，强迫全城的人向上主悔改补过；\BibleRef{纳 3:6}又如\Person{达理阿}事奉祂，将偶像交给\Person{达内尔}毁掉，并将他的敌人投入狮子圈；又如\Person{拿步高}事奉祂（我之前提到过他），他发布了严厉的法律，禁止臣民中任何人亵渎天主。\BibleRef{达 3:29}因此，君王们正是以这种方式，即使在身为君王的职分上，也能事奉上主，当他们为祂服务时所行的，是他们若非君王便不能行的事。\par

20. 由此可见，在宗徒时代，地上的列王尚未侍奉主，仍旧图谋虚妄的事，抗拒主，抗拒祂的\OriginalTerm{受傅者}{Anointed}，为使先知们所说的一切得以应验，我们必须承认，那时不虔敬的行为不可能藉法律防止，反而是在法律的认可下进行的。因为事态的进展如此展开，甚至犹太人杀害宣讲基督的人，还以为是在尽恭敬天主的义务，正如基督所预言的\BibleRef{若 16:2}，而外邦人也在狂怒攻击基督徒，殉道者的忍耐却胜过一切。但一旦开始应验较晚诗篇所记的话：\BibleQuote{众王都要向祂俯伏敬拜，万民都要侍奉祂}，有哪个清醒的人能对君王们说：\Quote{不要让任何忧虑扰乱你国内的心，关于谁抑制或攻击你主的教会；不要认为你臣民中谁选择虔诚或\OriginalTerm{亵圣}{sacrilege}与你有关}，因为你不能对他们说：\Quote{你臣民中谁选择贞洁或不贞洁与你无关}？既然天主赋予人\OriginalTerm{自由意志}{free-will}，为何法律惩罚奸淫，却允许\OriginalTerm{亵圣}{sacrilege}？难道灵魂对天主不守信，比妻子对丈夫不忠是更轻的事吗？或者，如果是出于无知而非藐视所犯的违背宗教真理的过错应受较轻惩罚，难道就可以完全忽视吗？\par

% structure: p0024 source=h2 target=subsection reason=relative-heading-rank; base=h2

\subsection{第六章}

21. 的确，以教导引导人敬拜天主，胜于以惩罚或痛苦的恐惧驱使人，这是无人能否认的；但这并不必然导致因为前者造就更好的人，就可以忽视那些不服从者。因为许多人是先因恐惧或痛苦而被迫，后来才受到影响教导，或得以实践他们已在言语上学习的内容，这一点我们已充分证明并日复一日验证着，因此他们获得了益处。有些人确实向我们提出某位世俗作者的情感，他说：\par

\Quote{我认为以羞耻训练年轻人，\newline{}畏惧卑贱，胜于以痛苦，是好的。}\newline{}这是不争的事实。但是，虽然那些由爱正确引导的人更好，由恐惧纠正的人肯定更多。因为，用他们自己的作者来回答这些人，我们发现他在另一处说，\newline{}\Quote{除非你因痛苦和磨难受教导，\newline{}你不能在任何事上正确地引导自己。}\par

但是，圣经论及前一类更好的人说：\BibleQuote{爱内没有恐惧；圆满的爱把恐惧驱逐出去；}\BibleRef{若一 4:18}，也论及提供多数的后一类较低的人说：\BibleQuote{仆人不会因言语受纠正；因为他虽明白，却不答理。}\BibleRef{箴 29:19}。他说\BibleQuote{他不会因言语受纠正，}并不是命令放任他，而是暗示应如何纠正他；否则他不会说\BibleQuote{他不会因言语受纠正，}而是毫无限制地说\BibleQuote{他不会受纠正。}因为在另一处他说，不仅仆人，就连不顺服的儿子也必须受鞭打纠正，且带来重大成果；因为他说：\BibleQuote{你要用杖打他，就能救他的灵魂免下地狱；}\BibleRef{箴 23:14}；在别处又说：\BibleQuote{不忍用杖打儿子的，是恨他。}\BibleRef{箴 13:24}。因为，若给我们一个人，他凭着正确的信德和真实的领悟，能倾心全力说：\BibleQuote{我的灵魂渴慕天主，生活的天主：我几时才能前来，出现在天主面前？}对这样的人，地狱的恐怖便不需要，更不谈暂时的惩罚或帝国的法律，因为对他来说，亲近主是绝对不可或缺的恩福，他不仅怕失去这福乐，以为重罚，连迟延获得都几乎难以忍受。然而，在好的儿女们能说他们有\BibleQuote{渴望死去，与基督同在，}\BibleRef{斐 1:23}，许多人必须先藉现世的鞭打，像恶仆那样，在某种程度上像无用的逃亡者那样，被召回他们的主那里。\par

22. 谁能爱我们胜过基督？祂为自己的羊舍弃了性命。\BibleRef{若 10:15}，然而，当祂仅凭言语召叫了\Person{伯多禄}和其他宗徒之后，到召叫\Person{保禄}（先前叫\Person{扫禄}），就是后来教会的大建筑者，但起初却是残忍的迫害者时，祂不仅以声音约束他，更以能力击倒他于地；为了强使一个在不信仰的黑暗中发狂的人渴望心灵的光明，祂先以肉眼的失明打击他。若未施加那惩罚，他后来就不会因此得医治；因为他睁眼时本什么也看不见，若眼睛未受损，圣经就不会告诉我们，当\Person{阿纳尼雅}覆手为使复明时，眼上仿佛有鳞片掉下，这鳞片曾遮蔽视线\BibleRef{宗 9:1}。\OriginalTerm{多纳图派}{Donatists}惯常呼喊的话在哪里：\Quote{人有自由选择相信或不相信？}基督向谁施暴？祂强迫了谁？这里他们有宗徒\Person{保禄}。让他们在他的事例中认识基督先强迫，后教导；先打击，后安慰。因为奇妙的是，那在最初因身体之罚的强迫下进入福音服务的人，后来在福音上比所有仅凭言语被召的人更劳苦；\BibleRef{格前 15:10}，而那被更大的恐惧强迫去爱的人，彰显了那驱逐恐惧的圆满的爱。\par

23. 所以，如果迷失的儿子强迫别人走向毁灭，教会为什么不能使用强力强迫她的迷失的儿子回来呢？虽然那些并非受强迫、只是被引入歧途的人，如果通过严厉但救恩的法律得以回归她的怀抱，他们的慈爱母亲会以更大的亲情接纳他们，他们成为远比那些从未迷失过的人更深值得庆祝的对象。当羊只离开羊群时，即使它们不是被强力驱离，而是被柔声蜜语诱惑而误入歧途，牧人也有责任在找到它们时，若它们显出抵抗迹象，就以鞭子的恐吓甚至痛楚将它们带回主人的羊栈，难道这不是牧人的本分吗？尤其当它们在这些逃亡奴隶和强盗中繁衍增多时，牧人更有权这样做，因为在他们身上认出了主人的印记，这印记在那些我们接纳却不重新施洗的人身上并未受到冒犯。因为羊的迷路固然应当纠正，但救主的印记却不应在它身上被毁灭。因为即使一个人由一个自己带有王印的背弃者盖了王印，两人都获得宽恕，那人回去服役，另一人也开始参与以前没有参与的服役，那么印记在两人之中都没有被抹去，反而在两人身上都被认出，并因为它是王的印记而受到应有的尊崇与认可。既然他们不能证明强迫他们去的归宿是坏的，他们却坚持认为即使对善事也应用强力。但我们已表明，保禄是被基督强迫的；因此，教会在试图强迫\OriginalTerm{多纳图斯派}{Donatists}时，是在效法她的主的榜样，尽管一开始她曾等待，希望不需要强迫任何人，为使先知关于君王和万民信仰的预言得以应验。\par

24. 在这个意义上，我们也不荒谬地解释圣保禄宗徒的话，他说：\BibleQuote{已经准备妥当，在你们完全服从之后，要责罚一切的不服从。}\BibleRef{格后 10:6} 因此，主自己也曾吩咐先邀请客人赴他的盛筵，之后再强迫他们来；因为当他的仆人回答说：\BibleQuote{主人，你吩咐的事已办妥，但还有空位，}他就向他们说：\BibleQuote{你们去到大道和篱笆那里，强迫他们进来。}\BibleRef{路 14:22} 这样，那些先被温和地领进来的人，先前的服从得以完成；而那些被强迫的人，背命受到了惩罚。因为在先前说\BibleQuote{领进来}，仆人回答说\BibleQuote{主人，你吩咐的事已办妥，但还有空位}之后，又说\BibleQuote{强迫他们进来}，这除了报复背命之外，还能有什么别的意思呢？如果他希望我们理解这强迫是借着神迹的恐怖力量，那么，许多神圣的神迹反而更在先被呼召的人眼前行过，尤其是那些犹太人，如经上说：\BibleQuote{犹太人要求的是神迹；}\BibleRef{格前 1:22} 而且，在宗徒时代，福音在外邦人中也是借着神迹如此被宣扬，以至于如果神迹就是主所命令的强迫方式，那么我们更有理由认为，如我以前所说，反而是先前的客人被强迫了。因此，如果教会在适当的时期，借着宗教特性和君王的信德，从天主规定接受了权力，作为工具，用来强迫那些在大道和篱笆处——即在异端与分裂中——的人进来，那么，他们就不要抱怨受强迫，反而该反省他们是否如此受强迫。主的晚餐是基督身体的合一，不仅在于祭台上的圣事，也在于和平的连系。至于\OriginalTerm{多纳图斯派}{Donatists}自己，我们固然可以说，他们不以任何好事强迫任何人；因为无论他们强迫谁，都只是强迫人走向邪恶。\par

% structure: p0031 source=h2 target=subsection reason=relative-heading-rank; base=h2

\subsection{第7章}

25. 然而，在那些强迫人参加圣体圣事的法律被送往\Place{非洲}之前，某些弟兄——我本人也是其中之一——认为，尽管\OriginalTerm{多纳图斯派}{Donatists}的疯狂在各地蔓延，我们仍不应请求皇帝们下令，透过对凡愿生活其中者施加惩罚，使异端绝对不复存在；相反，他们应满足于下令：那些或用声音宣讲天主教真理，或藉研究确立天主教真理的人，不再遭受异端者的激烈暴行。此外，他们认为，这或许可以透过某种方式达成，如果他们能采用由\Person{狄奥多西}——怀着虔敬的记忆——为普遍对付各类异端者而颁布的法律，该法律规定，在任何地方发现的任何异端主教或圣职人员，应被罚款十磅黄金，并以更明确的措辞将其针对否认自己是异端者的多纳图斯派予以确认；但附带以下保留：罚款不应施加于他们全体，而只应施加于那些天主教会遭受其圣职人员、\OriginalTerm{西康塞里奥}{Circumcelliones}或其任何会众暴行的地区；这样，在遭受暴行的天主教徒提出正式申诉之后，主教或其他圣职人员应依据授予官员的委任，立即被迫缴纳罚款。因为我们想，以此方式，如果他们受到恐吓而不敢再行此类之事，天主教真理或许可以在这样的条件下被自由地教导和持守：即没有人被迫接受，任何愿意的人可以毫无畏惧地追随它，这样我们便不会有虚假和假冒的天主教徒。尽管其他弟兄持有不同看法——他们或是年资较深，或是对许多地区和地点有经验，我们在那些地方看到真正的天主教会已稳固建立，然而，这教会是在先前皇帝们的法律迫使人们进入天主教共融时，由天主的大仁慈所植立和坚固的——但我们仍贯彻了我们的主张，即应优先向皇帝们寻求上述措施：在我们的会议中做出了决定，\BibleRef{若 16:2} 并派遣使者前往伯爵的宫廷。\par

26. 但天主以祂的大仁慈，深知这些法律所激起的恐惧，以及对许多人心冷酷邪恶及那种不能以言语软化、却可藉由某种轻微纪律的严厉而得以软化之硬心而言，犹如一种药石般的困扰，是多么必要，遂使我们的使者未能获得他们所请愿的事。因为来自其他地区的若干主教已先于我们抵达，他们提出了严重的申诉，述说自己曾遭受多纳图斯派本身的许多虐待，并被逐出自己的主教座堂；尤其是，企图以难以置信的残暴方式杀害\Place{巴盖}教会天主教主教\Person{马克西米安}的事件，导致了各种措施的采取，使我们的代表团无计可施。因为一项法律已经颁布：多纳图斯派的异端如此凶残，对之仁慈实比其疯狂本身造成更大的残忍，故此今后不仅要防止其施行暴力，而且要使其根本不能安然存在；然而，并未对其处以死刑，以便即使对待不配之人，也能遵守基督徒的温和，但规定了罚款，并对其主教或圣职人员宣判流放。\par

27. 至于前述巴盖主教，在其诉求于普通法庭获准，且双方依次听审后，在一座多纳图斯派占据的、本属天主教徒的教堂内，他们趁他站在祭台前，以可怕的暴力和残忍的狂暴向他扑去，用棍棒和各种武器野蛮地殴打他，最后甚至用破碎祭台的木板打他。他们还用匕首在他腹股沟处严重刺伤，以至若非他们进一步的残暴反倒有助于保全其生命，大出血很快就会使之丧命；因为当他们将他重伤倒地拖行时，灰尘灌进喷血的血管，堵住了那原本急速流向死亡的血流。随后，当他们终于弃他而去，我们的人试图唱着圣咏将他抬走；但他的敌人怒火更炽，从抬他的人手中将他夺去，严惩那些天主教徒，将他们击溃，因他们在人数上远超，且轻易以其暴力引起恐慌。最后，他们将他投入一座高塔，以为他此时已经死去，其实他仍在呼吸。他幸运地落在一堆松软的泥土上，夜间有路人藉灯光发现，认出并救起，将他带到一所修院，经悉心照料，在数日内从几乎无望的危险中恢复了。然而谣言甚至越过大海传来，说他已被多纳图斯派的暴力杀害；当后来他自己出来，且出乎意料地被看见活着时，他藉着伤痕的数量、严重和新鲜程度，充分表明谣言带来他的死讯是何等有理。\par

因此，他向那位基督徒皇帝寻求协助，并非多么渴望为自己复仇，而是为了保卫他所受托的教会。倘若他未这样做，就不应因其容忍而受称赞，反而该因疏忽而受责备。因为\OriginalTerm{宗徒}{Apostle} \Person{保禄}并非为自己短暂的生命采取预防措施，而是为了天主的教会，当他使那些密谋杀害他的人的计谋被罗马长官知晓时，结果是他由一队武装士兵护送到他们计划送他去的地方，从而得以逃脱敌人的埋伏。\BibleRef{宗 23:17} 他也毫不迟疑地援引罗马法律的保护，宣称自己是罗马公民，当时不可受鞭打；\BibleRef{宗 22:25} 并且，为了不被交给那些想要杀害他的犹太人，他上诉于\Person{凯撒}，\BibleRef{宗 25:11} ——那是一位罗马皇帝，但并非基督徒。他借此足够清楚地表明，基督的仆人后来在教会的危险中，发现皇帝们是基督徒时，应尽的责任是什么。因此，随后发生的是，一位虔诚而虔敬的皇帝，在得知此类事务时，认为通过制定最为虔敬的法律，完全纠正这极不虔敬的错误，甚至以恐吓和强制的手段，将那些举着基督的旗帜反对基督事业的人纳入\OriginalTerm{天主教会}{Catholic Church}的合一之中，比仅仅除去他们施暴的能力，而任其迷失于歧途并在错误中丧亡要好。\par

29. 不久，当这些法律本身抵达\Place{非洲}时，首先是那些已经在寻找机会、或惧怕\OriginalTerm{多纳特派}{Donatists}狂怒、或先前因不愿得罪朋友而迟疑的人，立刻归向了教会。也有许多人，只是由于家庭传统习惯的束缚，来自父母，但从未思考过异端本身的根基是什么——实际上，从未愿意探究和思考其本质——现在开始运用他们的观察力，发现其中没有任何东西能够补偿他们所被要求遭受的如此重大的损失，便毫不困难地成为了\OriginalTerm{天主教徒}{Catholics}；因为他们先前因安全感而漫不经心，现在因焦虑而受到教导。然而，当所有这些人都树立了榜样之后，许多自己不太能理解多纳特派的错误与天主教真理之间差异的人也纷纷效仿。\par

30. 因此，当广大的人民群众被真正的母亲欣然地接纳到她的怀抱中时，外面仍然留下了残酷的团伙，他们以不幸的敌意坚持着那种疯狂。即便这些人中，大多数也在假装的修和中进行了共融，另一些人则因人数稀少而未受注意。然而那些假意顺从的人，逐渐习惯了我们的共融，并聆听真理的宣讲，尤其是在我们与他们的\OriginalTerm{主教}{bishops}在\Place{迦太基}举行的会谈和辩论之后，在很大程度上被引向了正确的信仰。但在某些地方，一个更为顽固和不肯妥协的群体占了优势，那些对与我们共融持有更好看法的少数人无法抵抗，或者群众受到少数更有权势的领导者的影响，追随他们走入歧途，我们的困难持续得更久一些。在这些地方中，少数地方至今仍有麻烦，其间天主教徒，特别是\OriginalTerm{主教}{bishops}和\OriginalTerm{圣职人员}{clergy}，遭受了许多可怕的苦难，要详细叙述过于冗长，其中有些人被挖去了眼睛，一位主教被砍去双手和舌头，有些人实际上被杀害了。更不用提那些最残忍的屠杀和趁夜入室抢劫，以及烧毁不仅是私人住宅，甚至是教堂——有人竟如此堕落，将\OriginalTerm{圣书}{sacred books}投入火中。\par

31. 然而，我们因这些恶事所遭受的痛苦，却由从中结出的果实而得了安慰。因为不信者无论在何处行这些事，那里基督徒的合一便以更大的热忱和完善而进展，主也受到更热切的赞颂，因为他竟肯恩赐他的仆人们能以自身的苦难赢得他们的弟兄，并通过他的血，聚集那些散布在致命错误中的羊群，进入永远\OriginalTerm{救恩}{salvation}的平安之中。主是大能和满有怜悯的，我们每日向他祈祷，求他赐予其余的人悔改之心，使他们能脱离魔鬼的网罗，他们曾被魔鬼活捉，随从他的意志，\BibleRef{弟后 2:26} 尽管他们现在只寻求诽谤我们的材料，以恶报善；因为他们没有知识能理解我们对他们持续怀有的感情和爱，以及我们是何等焦心地，按照主藉着\Person{厄则克耳}先知的口赐给他牧者的命令，要把那被赶走的领回，寻找那失丧的。\BibleRef{则 34:4}\par

% structure: p0039 source=h2 target=subsection reason=relative-heading-rank; base=h2

\subsection{第八章}

32. 但他们，正如我们有时在其他地方说过的，并不把他们对我们所做的事归咎于自己；相反地，他们却把他们对自己所做的事归咎于我们。我们这边有谁会希望——我且不说他们中一人丧亡，就连任何财产损失都不希望呢？但是，如果\Person{达味}家族不能以其他条件换得和平，除非他的儿子\Person{阿贝沙隆}在那场他攻打父亲的战争中被杀，尽管他极为谨慎地严令跟随他的人务必尽全力保全他性命安全，好使他的父爱能在阿贝沙隆忏悔时宽恕他，那么他除了为失去的儿子哀哭，并藉着想到国度获得了和平来安慰自己的悲伤之外，还能做什么呢？同样的情况适用于我们的母亲——天主教会；因为当她受到人们——他们当然不同于儿子们——的攻击时，必须承认从这棵枝繁叶茂蔓延全世界的巨树上，这根在\Place{非洲}的小枝被折断了，尽管她出于爱愿意让他们重生，好使他们能回归那无之则不能有真生命的根，同时，如果她通过损失某些人而聚集起如此众多的其余人，她便以这些强大民族的得救之思来抚慰并治愈慈母心怀的悲痛；尤其她想到那些丧亡者之死是由他们自己导致的，而非如阿贝沙隆那样死于战争中的厄运。如果你能看到那些在\Person{基督}的平安中得到拯救的人的喜乐，他们拥挤的集会，热切的渴望，他们成群结队前来听和唱赞美诗，并接受天主的话教导时的欢欣；他们中许多人回忆起从前的错误时的大忧伤，他们得知所学真理时的喜乐，以及他们发现他们的说谎教师关于我们圣事所散布的种种谎话之后，对这些人感到的愤慨和憎恶；此外，多少人承认他们很久以前就渴望成为天主教徒，但在如此暴力的人中间却不敢迈出这一步——我再说，如果你看到这些从如此毁灭中得救的民族的集会，那么你就会说，那是极端的残忍，倘若因害怕某些人数远不能与被救者群众相比的绝望之人，可能会死于他们自愿点起的烈火之中，而任由其他人永远丧亡，在不灭的火中受煎熬。\par

33. 假若有两个人同住一屋，我们确知即将倒塌，而我们警告他们有危险时，他们却不愿相信，仍执意留在屋内；假若我们有能力救出他们，甚至违背他们的意愿，并且事后能向他们展示那威胁着房屋的毁灭，使他们不敢再回到危险所及之处，我认为，如果我们放弃这样做，就理当被责为残忍。而且，假如其中一人对我们说：\Quote{你们既已进入房屋来救我们的性命，我立刻自杀；}而另一人虽不愿出屋，也不愿被救，却也无胆自杀：我们应该选择哪一种做法——是让两人一起被压在废墟下，还是，无论如何，藉着我们仁慈的努力救出一人，而另一人丧亡，不是由于我们的过错，而是咎由自取？没有人会糊涂到连在这种情况下该怎么做都弄不清楚。我已经提出两个人的问题——也就是说，一个人丧亡，一个人得救；那么，对于少数人丧亡，而无数民族得救的情形，我们又该作何想法呢？因为实际上，出于\OriginalTerm{自由意志}{free will}而如此丧亡的人，没有那些因所论及的法律而从致命且永恒的毁灭中得救的庄园、村庄、街道、堡垒、市镇和城市那么多。\par

34. 但是，如果我们更仔细地思考正在讨论的这件事，我想，倘若那将要倒塌的屋子里有许多人，其中任何一人都可能得救，而当我们努力救他时，其他人却跳窗自杀，那么，我们虽因其余人的损失而哀伤，却可因那一个人的安全而得安慰；我们不会因担心剩下的人会自杀，就任由所有人无一获救地丧亡。那么，对于那项我们应当致力于使人类获得\OriginalTerm{永生}{eternal life}、逃避永罚的仁慈工作，我们又当作何想法呢？既然真正的理性和仁爱迫使我们给予人们这样的援助，以确保他们的安全——不仅仅是暂时的，而且是极其短暂的，仅仅为他们尘世生命的短暂时光？\par

% structure: p0043 source=h2 target=subsection reason=relative-heading-rank; base=h2

\subsection{第九章}

35. 关于他们对我们提出的指控，即我们觊觎并掠夺他们的财产，我愿他们成为\OriginalTerm{天主教徒}{Catholics}，与我们在平安与爱中共同拥有，不仅是他们所称属于他们的，还包括公认属于我们的。但他们因渴望诽谤而如此盲目，以致没有注意到他们的陈述彼此多么不一致。无论如何，他们声称，并似乎在他们中间将其作为最令人反感的抱怨主题，即我们通过法律的强制权力强迫他们进入我们的共融——如果我们想获取他们的财产，我们绝不会这样做。有哪个贪婪的人曾希望别人分享他的财产？谁被权力欲所煽动，或因拥有权力而骄傲，曾希望有个伙伴？让他们至少看看那些曾属于他们、但如今因兄弟情谊与我们结合在一起的弟兄，看看他们如何不仅持有他们过去所有的，还持有原本属于我们、他们过去所没有的；然而，如果我们像穷人一样与穷人共融生活，这些财产就同属于我们和他们；而如果我们个人拥有的财富足以满足我们的需求，那便不再是我们自己的，因为我们不会犯下如此可耻的霸占行为，将穷人的财产据为己有，在某种意义上我们是这些穷人的受托管理人。\par

36. 因此，凡以多纳特派各教会名义持有的财产，连同建筑物本身的拥有权，基督徒皇帝们以其虔敬的法律，下令转移给天主教会。那么，既然我们中间有那些教会的贫穷成员，他们过去靠这些微薄的财产维持生活，让他们自己停止贪图属于别人的东西，而仍留在外面，不如让他们进入合一的纽带，好让我们大家一同管理，不仅是他们称之为自己的财产，还有据称属于我们的财产。因为经上记载：\BibleQuote{一切都是你们的；你们却是基督的，基督是天主的。}\BibleRef{格前3:22} 在作我们元首的祂之下，让我们在祂的一个身体内合而为一；并且，在你们所提及的所有这类事上，让我们效法在 \BibleBook{宗徒}{大事录} 中记载的榜样：\BibleQuote{他们一心一魂，没有人说他所拥有的任何东西是自己的；他们一切所有都是共同的。}\BibleRef{宗4:32} 让我们喜爱我们所唱的：\BibleQuote{看，弟兄们同居，是多么美善，多么愉快！}好使他们藉着自己的经验，明白他们的母亲——天主教会，以何种完全的真实向他们呼喊那位真福宗徒写给格林多人的话：\BibleQuote{我所求的不是你们的财物，而是你们自己。}\BibleRef{格后12:14}\par

37. 然而，如果我们想一想 \BibleBook{}{智慧篇} 中所说的：\BibleQuote{因此义人掠夺了不敬虔的人；}\BibleRef{智10:20} 以及 \BibleBook{}{箴言} 中所说的：\BibleQuote{罪人的财富是为义人积存的；}\BibleRef{箴13:22} 那么我们将会看到，问题不在于谁占有了异端者的财产？而是谁处于义人的团体中？我们确实知道，多纳特派将如此大量的\OriginalTerm{义德}{righteousness}归于自己，他们不但夸耀自己拥有它，而且还声称将它赐予别人。因为他们说，任何被他们施洗的人就是被他们\OriginalTerm{成义}{justified}了，此后他们只需对由他们施洗的人说，他必须相信那位施行圣事的人；因为既然宗徒说：\BibleQuote{相信那使不敬虔者成义者的人，这人的信德就算为他的义德}\BibleRef{罗4:5}，他为什么不应如此呢？因此，如果使人成义的别无他人，就让他相信为他施洗的人，好让他的信德得算为义德。但我认为，即使他们自己，倘若有一刻胆敢怀有这样的想法，也会对自己感到恐惧。因为除了天主，没有谁是义人并能使人成义。但对他们所说的话，正如同宗徒论及犹太人所说的：\BibleQuote{他们因为不知道天主的义德，想要立自己的义德，就不顺从天主的义德。}\BibleRef{罗10:3}\par

38. 然而，我们当中绝不可有人自称如此义德，以致他竟想要立自己的义德，仿佛这义德是由他自己赋予自己的，却不对他说：\BibleQuote{你有什么不是领受的呢？}\BibleRef{格前4:7} 或者胆敢夸耀自己在此世没有罪，如同多纳特派在我们会议中宣称的那样，他们是一个已经没有瑕疵，没有皱纹，或任何这类事物的教会的成员\BibleRef{弗5:27}——他们不知道，这只有在那些\OriginalTerm{圣洗}{baptism}后立即离开此身，或在祈祷中祈求罪恶赦免之后离世的人身上才得以实现；但对于整个教会而言，完全没有瑕疵、皱纹或任何这类事物的时刻不会来临，直到我们听见那话：\BibleQuote{死亡！你的胜利在那里？死亡！你的刺在那里？死亡的刺就是罪。}\BibleRef{格前15:55}\par

39. 但在此生，当这败坏的身体重压着灵魂时\BibleRef{智 9:15}，如果他们的教会已经具有他们所主张的那种性质，他们就不会向天主献上主教导我们使用的祈祷：\BibleQuote{求祢宽恕我们的罪债。}\BibleRef{玛 6:12}因为既然在圣洗中一切罪都已被赦免，为什么教会还做这样的祈求，如果已经在此生就没有斑点、皱纹或这类的事了？他们还会轻视宗徒若望的警告，当他在书信中大声说：\BibleQuote{如果我们说我们没有罪，就是欺骗自己，真理也不在我们内。但若我们明认我们的罪，天主是忠信正义的，必赦免我们的罪，并洗净我们的各种不义。}\BibleRef{若一 1:8}因为这一希望，普世教会献上祈求：\BibleQuote{求祢宽恕我们的罪债，}好使主看到我们并非虚荣，而是乐意承认己罪，洗净我们的各种不义，这样主耶稣基督可以在那日将荣耀的教会展示给自己，没有斑点或皱纹，或这类的事，现在祂藉着水洗，藉着圣言来净化她：因为一方面，在圣洗中没有什么留下阻挡一切过去之罪的赦免（只要圣洗不是在教会之外无效地领受，而是在教会内施行，或者至少，如果已在外部施行，领受人不是带着它留在外面）；另一方面，那些在此生圣洗后生活的人，因人性的软弱而沾染的任何罪污，无论何种，都藉着同一洗濯的功效得以洗净。因为对一个未受洗的人来说，说\BibleQuote{求祢宽恕我们的罪债。}毫无益处。\par

40. 因此，祂现今如此藉着水洗，藉着圣言净化自己的教会，以便将来将她展示给自己，作为没有斑点、皱纹或这类的事的——也就是说，全然美丽，并在绝对完美中，那时死亡要被\BibleQuote{胜利吞灭}\BibleRef{格前 15:54}。所以现在，在我们内，凡由天主所生、凭信德而生活的生命兴旺到什么程度，我们也就在什么程度上是义人；但我们也拖着来自亚当的必死性情的痕迹到什么程度，我们就不能无罪。因为以下两句话都是确实的：\BibleQuote{凡由天主生的，就不犯罪。}\BibleRef{若一 3:9}和前面的话：\BibleQuote{如果我们说我们没有罪，就是欺骗自己，真理也不在我们内。}\BibleRef{若一 1:8}因此，主耶稣既是义者，又能使人成义；但我们得称为义，是白白地，不是凭别的恩宠，而是凭祂的恩宠\BibleRef{罗 3:24}。因为除了祂的身体——即教会——之外，没有什么使人成义；所以如果基督的身体夺得不义之人的战利品，不义之人的财富被储存起来作为基督身体的宝藏，那么不义之人不应因此留在外面，反而应进入其内，好使他们得以成义。\par

41. 因此我们也可以肯定，关于审判之日所写的：\BibleQuote{那时义人在磨难他的人面前，并对他所轻视的劳苦之人，将大无畏地站立，}\BibleRef{智 5:1}不应被解释为：客纳罕人将站在以色列人面前，尽管以色列人轻视客纳罕人的劳苦；而应解释为：纳波特将站在阿哈布面前，因为阿哈布轻视纳波特的劳苦；因为客纳罕人是不义的，而纳波特是义人。同样，外邦人不会站在基督徒面前，基督徒曾轻视他的劳苦，当偶像的庙宇被劫掠摧毁时；但基督徒将站在外邦人面前，外邦人曾轻视他的劳苦，当殉道者的身体被置于死地时。因此，同样地，异端者不会站在公教徒面前，公教徒轻视他的劳苦，当公教皇帝的法律得到执行时；但公教徒将站在异端者面前，异端者曾轻视他的劳苦，当不敬虔的\OriginalTerm{游荡派}{Circumcelliones}的疯狂被允许横行时。因为圣经的这段话本身就解决了问题，它不是说：那时人们将站立，而是说：\BibleQuote{那时义人将站立；}并且他们\BibleQuote{大无畏地站立}，因为他们站在良心的力量中。\par

42. 但是在这个世界上，没有人凭自己的义成为义人——也就是说，义似乎是由他自己并为他自己的行为造成的；但正如宗徒所说：\BibleQuote{天主按每人的信德分量分配了。}然后他继续补充：\BibleQuote{就如我们在一个身体上有许多肢体，但每个肢体的作用不同；同样，我们众人在基督内，都是一个身体。}\BibleRef{罗 12:3}根据这一教导，只要一个人与这身体的合一分离，就不能成为义人。因为正如一个肢体如果从一个活人的身体上被割下，就不能再保有生命的气息；同样那从义者基督的身体上被割下的人，绝无法保有正义的精神，即使他保有在身体中时所领受的肢体形状。因此让他们进入这身体的架构内，这样他们才能拥有自己的劳苦，不是出于统治的欲望，而是出于善用它们的虔敬。但我们，如前所说，即使在你乐意指定的任何仇敌的判断中，也洁净我们的意志免受这贪欲的污染，因为我们竭尽全力恳求我们享用其劳苦的那些人，在天主教会的共融内，与我们一同享用他们和我们劳苦的果实。\par

% structure: p0052 source=h2 target=subsection reason=relative-heading-rank; base=h2

\subsection{第10章}

43. 但他们说，这正是令我们不安的事——如果我们是不义的，你们为什么还寻求与我们结交？对此我们回答：我们寻求你们这不义之人的结交，是叫你们不再是不义的；我们寻找你们这些失丧的人，是为了一找到你们，就因你们欢喜，说：\BibleQuote{我这个兄弟是死而复生，失而复得的。} \BibleRef{路 15:32}。那么，他问：你为什么不给我施洗，好洗去我的罪呢？我答复：因为当我在纠正一个叛逃者的恶行时，我并不藐视君王的印记。他又问：我为什么不能在你们中间做补赎呢？其实，你若不先做\OriginalTerm{补赎}{penance}，就不能得救；因为你若不先悔恨自己曾经迷失，又怎能因改邪归正而喜乐呢？那么，他再问：当我们转到你们这边时，我们从你们这里得到什么呢？我回答：你们确实没有领受\OriginalTerm{圣洗}{Baptism}，因为圣洗虽能在基督身体的构架之外存在于你们身上，却不能使你们获益；但你们领受的是圣神的合一，在以和平彼此联系中，\BibleRef{弗 4:3}。若没有这合一，谁也不能看见天主；你们还领受了\OriginalTerm{爱德}{charity}，正如经上写的：\BibleQuote{爱德遮盖许多罪过。}\BibleRef{伯前 4:8}。对于这极大的恩赐，我们有宗徒的见证，若没有它，即使有能说人或天使的话，能明白一切奥秘，有先知之恩，有足以移山的信德，舍得把全部家产施舍给穷人，甚至舍身被火焚烧，也无益处。\BibleRef{格前 13:1}。我告诉你们，如果你们认为这浩大恩赐是无用的或微不足道的，你们就理所当然地、可怜地迷失了；而且你们必然要丧亡，除非你们归向\OriginalTerm{天主教会}{Catholic Church}的合一。\par

44. 那么，他们又问：如果我们为了获得救恩，有必要悔改我们曾置身于教会之外并与教会为敌的过错，那为什么在你们要求我们做了悔改之后，我们仍然可以在你们中间担任\OriginalTerm{圣职}{clergy}，甚至可能成为主教呢？事实上，老实说，我们必须承认，这种情况本不该发生，若不是因为和平自身的补偿力量医治了这恶。但让他们自己受教吧，尤其是那些陷入与教会分离的深沉死亡之中的人，要心里悲伤，因为正是通过在我们\OriginalTerm{天主教慈母教会}{Catholic mother Church}身上施以这种创伤，他们才得以恢复生命。因为当被砍下的树枝被嫁接时，树上就产生一个新伤口，为了接纳它，好把生命赐给那个因缺乏来自根部的生命而快要枯死的枝条。当新接上的枝条与它所接上的树干长成一体时，结果就是活力与果效；但如果它们不能长成一体，嫁接的枝条就会枯萎，而树的生机丝毫不损。此外，还有一种嫁接方式，无需切除树内的任何原有树枝，而是将不属于这树的枝条插入，虽不免造成伤害，但这是对树最轻微的伤害。同样，当他们来到存在于\OriginalTerm{天主教会}{Catholic Church}内的根部，虽对他们的错误做了无比深刻的悔改，却无需放弃任何属于他们作为圣职人员或主教的地位时，那就好比在慈母教会的树皮上留下某种创伤，打破了她纪律的严格；但是，正如经上说：\BibleQuote{栽植的不算什么，浇灌的也不算什么，}\BibleRef{格前 3:7}，一旦借着向天主的仁慈倾注的祈祷，通过所嫁接的枝条与母株的合一而获得平安，爱德就会遮盖许多罪过。\par

45. 虽然教会曾规定，凡因过错而被要求做补赎者，不得准许进入\OriginalTerm{圣秩}{holy orders}，亦不可重返或继续留在圣职团体中，但这规定的目的，并不是要使人对得蒙任何宽容感到绝望，而仅是为了维持严格的纪律；否则，就会构成反对那交予教会的钥匙的证据，关于这点我们有圣经的见证：\BibleQuote{凡你在地上所释放的，在天上也要被释放。}\BibleRef{玛 16:19}。然而，为了避免这样的情况发生：即在过犯被揭露之后，一个因渴望教会高位而心怀骄矜的人，以骄傲的态度去做补赎，于是极为严格地决定：凡犯了任何\OriginalTerm{大罪}{mortal sin}而做了补赎后，任何人都不得再被接纳进入圣职行列，这样，当所有对现世高升的希望都断绝以后，谦卑的良药便具有更大的力量和真实。因为连圣\Person{达味}也为死罪做了补赎，却并未丧失其尊位。我们也知道，真福\Person{伯多禄}流了最痛苦的眼泪后，懊悔自己否认了主，却仍然留在\OriginalTerm{宗徒}{apostle}之列。但我们绝不能因此就认为后世人的谨慎是多余的：他们——依我看，必是经验到某些人因贪图职权而假装悔改——在不危及救恩的前提下，对谦卑又加上了一些约束，好使救恩得到更彻底的保障。因为许多疾病的经验，必然会带来许多救治方法的发明。然而在此类事件中，当教会因严重的分裂破裂，不再是某一个个人的危险，而是整个民族都处于亡塌之中时，就有理由从我们的严格中稍作让步，以便真爱德能发挥其辅助作用，来医治更严重的弊害。\par

46. 因此，让他们为他们过去可憎的错误深感悲痛，如同\Person{伯多禄}因恐惧而说谎后所做的那样；让他们归向基督的真教会，也就是我们的母亲\OriginalTerm{天主教会}{Catholic Church}；让他们在其中作圣职人员，让他们作主教以造福教会，正如他们迄今一直敌对教会那样。我们并不嫉妒他们，不，我们接纳他们；我们愿意、我们劝导、我们甚至勉强那些我们在路上和篱笆边找到的人进来，尽管我们尚未说服其中一些人，我们所求的不是他们的财物，而是他们自身。宗徒\Person{伯多禄}，当他不认他的救主而哭泣，且仍不失为宗徒时，尚未来领受所应许的\OriginalTerm{圣神}{Holy Spirit}；但这些人更未领受圣神，因为他们割裂了身体的结构——唯有圣神赋予这身体生命——并在教会外、敌对教会的情况下，篡夺了教会的\OriginalTerm{圣事}{sacraments}，还用我们自己的武器和旗帜反对我们，仿佛在进行一场内战。让他们来吧；让和平在\Place{耶路撒冷}的力量中得以缔结，这力量就是\OriginalTerm{基督的爱德}{Christian charity}，——对这圣城有话说：\BibleQuote{愿你城垣内永享和平，愿你堡垒中常保安全！}让他们不要反对母亲的挂虑而高举自己，她过去和现在都怀着这挂虑，为将他们和所有强大的民族——那些他们正在或最近仍在欺骗的民族——聚集在自己怀中；让他们不要因她这般接纳他们而骄傲自满；让他们不要将她为了缔造和平而行的善事，归因于他们自己高举的邪恶。\par

47. 教会向来习惯援助因分裂和异端而沦丧的大众。当这种做法被用来接纳并医治那些曾陷于\OriginalTerm{亚略异端}{Arian heresy}毒素而沦丧的人时，这令\Person{路西法}不悦；因对此不悦，他堕入了分裂的黑暗，丧失了\OriginalTerm{基督的爱德}{Christian charity}的光明。根据这一原则，\Place{非洲}教会从一开始就承认了多纳徒派，在此事上服从了主教们的判决，这些主教曾在\Place{罗马}教会中就\Person{凯基利安}与\Person{多纳图斯}一党作出裁决；在谴责了被证实为分裂肇始者的名为\Person{多纳图斯}的主教后，他们决定，即便其他人曾在教会外受祝圣，经矫正后也应获接纳，并完全承认其\OriginalTerm{圣秩}{orders}——并非因为他们即便在基督身体合一之外也能拥有圣神，而是首先为了那些他们留在教会外时可能欺骗、使之无法获得此恩赐的人；其次，也为了他们自己的软弱能在内部蒙怜悯收纳，从而可获得治愈，不再有固执阻碍他们闭眼不见真理的明证。他们自己接纳了\Person{马克西米安}派——他们曾以亵渎的分裂罪名谴责他们，正如其会议所表明的，并且为了填补其空缺已祝圣了其他人——当他们看到民众并未离开他们的团体，以免所有人都卷入灭亡时，若非为此意图，又何以解释他们的行为？他们又基于何种理由，既不非议也不质疑那些被他们谴责的人在教会外所施行的\OriginalTerm{圣洗}{baptism}的有效性？那么，他们为何惊讶、为何抱怨，并以此诽谤我们，说我们如此接纳他们是为促进基督的真和平，而他们却不记得自己曾为促进与基督对立的、\Person{多纳图斯}的假和平所做之事？因为若铭记他们的这种行为，并在辩论中理智地用以反驳他们，他们将无言以对。\par

% structure: p0058 source=h2 target=subsection reason=relative-heading-rank; base=h2

\subsection{第十一章}

48. 但他们如此争辩说：\Quote{如果我们得罪了圣神，就是因为我们轻看了你们的圣洗，那么你们为何还寻找我们呢？因为正如主所说：\BibleQuote{凡出言干犯圣神的，在今世及来世，都不得赦免？}} \BibleRef{玛 12:32} ——他们却没有认识到，按照他们对这节经文的解释，无人能得救。因为谁不出言干犯圣神、不得罪他呢？无论是对于一个尚未成为基督徒的人，还是对于一个信奉\Person{亚略}、或\Person{欧诺米}、或\Person{马其顿纽斯}异端的人——这些人都说他是受造物；或是对于一个信奉\Person{福蒂努斯}的人——他完全否认他具有任何实体，说只有\OriginalTerm{天主圣父}{the Father}是唯一的天主；或是对于任何其他的异端者，要一一细说将太费时间。那么，难道这些人中没有一个能得救吗？或者，如果那些主曾亲自指斥的\OriginalTerm{犹太人}{Jews}相信了他，难道就不许他们受圣洗吗？因为救主并没有说：在圣洗中得赦免；而是说：\BibleQuote{在今世及来世，都不得赦免。}\par

49. 因此，他们应当明白，不是所有的罪，而只是某些亵渎圣神的罪，是不能获得赦免的。正如我们的主说：\BibleQuote{如果我没有来，没有向他们说话，他们就没有罪，} \BibleRef{若 15:22}显然，他并不是要说他们会完全没有罪，因为他们充满了许多重罪，而是要说他们会没有某种特殊的罪，这种罪的缺失将使他们能够获得他们内仍存的所有其他罪的赦免，即不信从来到他们那里的他的罪；因为他们若没有他来临，就不可能犯此罪。同样，当他说：\BibleQuote{无论谁亵渎圣神}或\BibleQuote{无论谁说话攻击圣神}时，显然他不是指任何在言语或行为上反对圣神的罪，而是要我们理解为某种特殊而个别的罪。这就是直到生命终结时的心硬，使人拒绝在基督身体的合一内接受罪的赦免；圣神赋予了这身体生命。因为当他对门徒们说\BibleQuote{你们领受圣神罢！}之后，立即补充说：\BibleQuote{你们赦免谁的罪，就给谁赦免；你们存留谁的，就给谁存留。} \BibleRef{若 20:22}因此，任何人若是抗拒或对抗这天主恩宠的赐予，或以任何方式与之疏离，直到这必死生命的终结，那么无论是在今世还是来世，他都不能获得那罪的赦免，因为这罪如此重大，以至于包含了所有的罪；然而，除非一个人离世，否则无法证实他犯有这罪。但只要他还活在世上，正如宗徒所说，\BibleQuote{天主的慈爱要引领他悔改}；但如果他故意地，并在极度的固执于邪恶中，如宗徒在接下来的章节中所补充的，\BibleQuote{随从自己顽固不悔改的心，为自己在忿怒和显示天主正义审判的日子积蓄忿怒} \BibleRef{罗 2:4}那么无论是在今世还是来世，他都不能获得赦免。\par

50. 但那些我们与之争论的，或是我们为之争论的人，不可对他们失望，因为他们仍在肉身内；但他们只能在基督的身体内寻求圣神，他们在教会之外拥有这身体的外在标记，却没有在教会内拥有那外在标记所指向的实质本身，因此他们吃喝自己的罪罚。 \BibleRef{格前 11:29}因为只有一个饼，它是合一的圣事，正如宗徒所说：\BibleQuote{我们虽多，只是一个饼，一个身体。} \BibleRef{格前 10:17}此外，只有天主教会是基督的身体，他是这身体的头，他身体的救主。 \BibleRef{弗 5:23}在这身体之外，圣神不给任何人生命，正如宗徒自己所说：\BibleQuote{天主的爱，藉着所赐与我们的圣神，已倾注在我们心中了；} \BibleRef{罗 5:5}但仇恨合一的人，不能分享这天主的爱。因此，那些在教会之外的人，没有圣神；因为经上论及他们说：\BibleQuote{他们分化人，属于血肉，没有圣神。} \BibleRef{犹 1:19}但那些在教会内却虚伪不实的人，同样不能领受圣神，因为经上也是为此而写：\BibleQuote{因为智慧的圣神远避欺诈。} \BibleRef{智 1:5}因此，若有人愿意领受圣神，就要谨防继续与教会隔绝，谨防以伪善的心进入教会；或者，如果他已经这样进入了，就要谨防坚持这样的伪善，好能确实而真正地与生命树结合为一。\par

51. 我已寄给你一封稍长的信函，或许会在你繁多的事务中成为负担。因此，如果这封信能分段读给你听，主会赐你理解，使你能够有所回答，以纠正和医治那些由我们慈母教会托付给你这位忠信儿子的人；求你藉主的助佑，无论在哪里，无论用什么方式，都尽力纠正和医治他们，或是亲自向他们说话和答复，或是让他们与教会的导师们沟通。\par

\SourceRecord{Translated by J.R. King.；Nicene and Post-Nicene Fathers, First Series；Vol. 1.；Edited by Philip Schaff.；Buffalo, NY: Christian Literature Publishing Co.,；1887.；<http://www.newadvent.org/fathers/1102185.htm>.}

% structure: p0001 source=h1 target=section reason=page-title

\section{第188封信（公元416年）}

\Emph{致\Person{儒利亚纳}夫人——在基督内堪当受与其位份相称的尊敬、我们当之无愧的卓越爱女，\Person{阿利比乌斯}和\Person{奥斯定}在主内致候。}\par

% structure: p0003 source=h2 target=subsection reason=relative-heading-rank; base=h2

\subsection{第一章}

夫人，您堪当在基督内受与其位份相称的尊敬，是我们当之无愧的卓越爱女；您的信在\Place{希坡}同时送达我们，这使我们非常高兴和欣慰，因此我们能一同给您回信，表达我们对您安康的喜悦，并以回应您真诚的爱意来让您知道我们的安康——我们确信您对此深为关切。我们清楚地知道，您并非不知道我们认为在基督内对您应尽的爱心有多大，以及我们在天主面前和人面前对您有多么关切。因为，尽管我们起初通过书信、后来通过个人交往而知道您是虔诚而天主教的，即是基督身体的真实肢体，然而我们卑微的职务也对您有益，因为当您从我们领受了天主的话时，如宗徒所说的：\BibleQuote{你们接受了它，并不以为是人的话，而确实是天主的话。}\BibleRef{得前 2:13}借救主的恩宠和仁慈，我们在您家中这项职务结出了如此丰硕的果实，以至当她的婚事准备就绪时，圣\Person{德梅特里亚}宁愿选择那位比世人更美的净配的精神拥抱；童贞女与祂结合而保持其贞洁，并获得更丰盛的精神果实。然而，我们至今仍不知道这位忠诚而高贵的少女是如何接受我们这一劝勉的，因为在她誓守贞洁之愿前不久我们就离开了；若不是从您信中的欣然宣告和可靠见证中得知，天主这份大恩——其实是由祂的仆人栽种和浇灌，但生长却全在于祂——已赐予我们，作为在祂葡萄园中劳作的工人。\par

既然情况如此，谁也不能指责我们冒昧，在这种更亲密的神缘基础上，我们表示对您的关怀，劝告您避免那些与天主的恩宠相反的意见。因为尽管宗徒命令我们在宣讲圣道时必须\BibleQuote{不论顺境逆境，务要宣讲}\BibleRef{弟后 3:2}，但我们仍不认为您会在那些认为我们劝您谨慎避开与健全道理相左之事的信或言为\Quote{不合时宜}的人之列。因此，您曾如此亲切地接受了我们的规劝，以至于您在我们现在回复的信中说：\Quote{我衷心感谢尊驾给我的虔诚劝告，叫我不要听从那些人，他们用有害的著作常败坏我们的神圣信仰。}\par

在这封信中您继续说：\Quote{但尊驾知道，我和我全家完全与这类人无涉；我们全家都严格遵循天主教信仰，从未在任何时候偏离过它，或陷于任何异端——我不是说那些几乎无法弥补地错误了的派别的异端，而是说那些看似微不足道的谬误。}这个说明使得我们更有必要在信中对您不保持沉默，谈论那些企图腐蚀甚至信仰纯正者的人的行为。我们认为您的家庭并非微不足道的基督的教会，那些认为我们内心所有的任何义德、节制、虔诚、贞洁都是出于我们自己，人如此受造，以至于除了已给我们的启示外，天主不再给我们任何进一步的帮助，使我们凭自己的抉择去完成通过学习所知道的义务；他们宣称本性与知识就是天主的恩宠，也是过正直公义生活的唯一助力。他们主张，我们拥有向善的意愿，由此产生正直的生活和那远超其他一切恩赐的爱德——天主本身甚至被称为爱，也只有借此才能在我们内满全我们所能满全的天主的法律和诫命——他们声称我们之拥有这样的意愿，并不归功于天主的助佑，而主张我们自己意志足以达成这些事。请不要让您觉得这只是个微不足道的错误：竟有人想自称为基督徒，却不愿听从基督的宗徒。宗徒说：\BibleQuote{天主的爱已倾注在我们心中}，免得任何人以为这爱是由于自己的自由意志，便立即补充说：\BibleQuote{借着赐给我们的圣神。}\BibleRef{罗 5:5}因此您应明白，那人不承认这是\BibleQuote{救主的大恩赐}\BibleRef{弗 4:7}是何等巨大、何等致命的错误；救主升上高天时，\Quote{带领俘虏，将恩赐给予人。}\par

% structure: p0007 source=h2 target=subsection reason=relative-heading-rank; base=h2

\subsection{第二章}

那么，我们怎能如此隐藏真实感受，而不警告您——您对我们如此关切——提防这些教义呢？我们读过一本写给圣\Person{德梅特里亚}的某本书之后，更觉得必须这样做。这本书是否已到达您手中，作者是谁，我们希望在您的回信中得知。在这本书中，若允许这样的人读它，基督的贞女将读到，她的圣德和一切属神财富只能从她自己而来，这样，在她达到真福的完美之前，她将学到——但愿天主不容！——对天主忘恩负义。因为在那本书中对她说的话是：\Quote{那么，您在此拥有了那些使您堪当、甚至更应优先于他人的东西；因为您属世的地位和财富，人们认为来自您的亲属，而非您自己，但您的属神财富却无人能给您，只能出于您自己；因此，您理当为此受称赞，理当为此而优先于他人，因为这些只能来自您自己，且在您自己内。}\par

5. 无疑你看到，这些话里的道理是何等危险，你必须提防。因为断言这些灵性财富只能存在于你内，这确实说得很好且真实：那显然是食粮；但断言它们只能从你而出，却是纯粹的毒药。任何基督的贞女都决不可甘心听信此类言论。每位基督的贞女都明白人心生来的贫穷，因此不愿用她净配恩赐之外的任何事物来装饰它。她宁可听从宗徒的话，他说：\BibleQuote{我曾把你们许配给一个丈夫，为把你们当作贞洁的童女献给基督。但我怕，就像那蛇以诡诈诱惑了\Person{厄娃}，你们的心也会失去在基督内的纯朴。}\BibleRef{格后 11:2}因此，关于这些灵性财富，她不应听从那说：\Quote{无人能将它们分施给你，只除了你自己，它们只能从你而出并在你内；}的人，却应听从那说：\BibleQuote{我们是在瓦器中存有这宝贝，为彰显那卓著的力量是出于天主，并非出于我们。}\BibleRef{格后 4:7}的那一位。\par

6. 关于那种神圣的童贞贞洁，她也应知道，这并不是出于她自己，而是天主的恩赐，虽则赐予那信且愿意的人。愿她听从那位真实而热诚的导师，他在论及这主题时说：\BibleQuote{我本来愿意众人都如同我一样，可是每人都有他各自从天主得来的恩宠：有人这样，有人那样。}\BibleRef{格前 4:7}愿她也听从那唯一净配，不单是她，也是整个教会的净配，论及这贞洁和纯洁时说：\BibleQuote{这话并不是人人都能领悟的，只有那些得了恩赐的人，才能领悟；}\BibleRef{玛 19:11}好使她明白，为获得这如此大而卓越的恩赐，她理当更该感谢我们的天主、我们的主，而不要听从任何说她凭自己拥有它的人的话——我们不可称这话为讨好人的谄媚者之言，免得我们似是轻率判断人的隐密思念，但这肯定是误导人的颂辞。因为，正如宗徒\Person{雅各伯}所说，\BibleQuote{一切美好的赠与，一切完善的恩赐，都是从上面，从光明之父降下来的；}\BibleRef{雅 1:17}所以，这圣洁的童贞就是从那里来的，你既然喜爱它并以此为乐，你的女儿在这方面却超越了你，她后你而出生，却在品行上走在你前面；从你而出，却在荣耀上超越你；年龄在你之后，却在圣德上超过你；因此，那在你自身不能有的事，在她身上却开始为你所有了。因为她没有缔结婚姻，为使自己不仅在灵性上富裕，也使你富裕，甚至比你更富裕，因为即使加上她，你仍比她不如，因为你所缔结的婚约产生了她。这些事都是天主的恩赐，确实是你的，却不是从你来的；因为你拥有这宝藏在脆弱的身体内，这些身体仍像陶匠的瓦器，为彰显那卓著的力量是出于天主，并非出于你。你不要惊奇我们说这些事是你的，却不是从你来的，因为我们称\Quote{日用粮}是我们的，但接着却说：\Quote{求祢赐给我们，}\BibleRef{路 11:3}免得以为它是从我们自己来的。\par

7. 因此，你要服从圣经的命令：\Quote{不断祈祷，事事感谢；}因为你祈祷，是为能时常且日益拥有这些恩赐，你感谢，是因为你拥有它们并非靠你自己。因为谁把你从那出自\Person{亚当}的死亡与丧亡之团中分离出来？岂不是那\BibleQuote{来寻找并拯救迷失者}\BibleRef{路 19:10}的？那么，人听见宗徒的问题：\Quote{谁使你异于别人？}\Quote{我自己的善意、信德、义德，}却不顾那紧接着的话：\BibleQuote{你有什么不是领受的呢？既然是领受的，为什么你还夸耀，好像不是领受的呢？}\BibleRef{格前 4:7}如此回答，是对的吗？因此，我们不愿意，绝对不愿意，当一个献身贞女听到或读到这些话：\Quote{你的灵性财富无人能分施给你；为此你当受赞美，为此你当受称扬，因为它们只能从你并在你内，}就如此夸耀财富，好像不是领受的一样。她固然可以说：\Quote{天主，我向祢所许的愿在我内，我要向祢献上赞颂；}但既然这些是在她内，不是从她来，她也应记得说：\Quote{主，因祢的旨意，祢赐予力量给我的美丽，}因为，虽然这也从她而来，即由于她意志的行为（没有它，我们不能行善），但我们不能像他那样说，这\Quote{只从她而来。}因为我们自己的意志，若非有天主的恩宠相助，就连名义上的善意也不可能是，因为宗徒说：\BibleQuote{是天主在你们内工作，使你们愿意，并使你们力行，为成就祂的善意，}\BibleRef{斐 2:13}——不是像这些人所想的，仅仅藉着启示知识，使我们知道该做什么，而是也藉着激发基督徒的爱德，使我们也能藉着选择实行我们所学来的事。\par

8. 因为，那说过的人无疑知道\OriginalTerm{节制}{continence}恩赐的价值：\Quote{我得知，除非天主赐与，没有人能节制。}他不仅知道那是多么大的益处，以及应当多么热切地渴求它，而且也知道，除非天主赐与，它便不能存在；因为智慧教导了他这一点，因为他说：这也是智慧的一点，知道这是谁的恩赐；但这知识对他还不够，他接着说：\Quote{我便去到上主那里，向他恳切祈求。}\BibleRef{智 8:21} 因此，天主在这事上援助我们，不仅让我们知道该做什么，还让我们因着爱慕而践行我们通过学识所已知的。所以，没有人能拥有，不仅是知识，就连节制，除非天主赐与他。正因如此，他有了知识后，便祈求能获得节制，好让这节制在他内，因为他知道这不是出于他自己；或者，若因他意志的自由，在某种意义上这是出于他自己，但也不单是出于他自己，因为除非天主赐与这恩赐，没有人能节制。然而，我所驳斥其意见的人，在谈论灵性财富——其中无疑包括那光彩而美好的节制恩赐——时，并没有说它们可以存在于你内、且由你而来，而是说它们\Emph{唯独}从你而来、并在你内，其方式是：正如基督的贞女除了在她自己内，无处拥有这些事，所以可以相信，她只可能从她自己获得这些事，且以这种方式（但愿仁慈的天主使她的心远离此事！）她自夸，仿佛她不是领受的！\par

% structure: p0013 source=h2 target=subsection reason=relative-heading-rank; base=h2

\subsection{第三章}

9. 我们确实对这位圣洁贞女的培育，以及她在其中被教养长大的基督徒谦逊，持有这样的看法，以至于确信：当她读到这些话时——如果她确实读了——她会放声痛哭，谦卑地捶胸，或许泪如泉涌，并怀着信德向那她献身服事、且蒙其圣化的主祈祷，恳求祂，这些不是她自己的话，而是别人的，并祈求她的信德不至于如此，竟相信她有什么可在自己内夸耀，而不在主内夸耀。因为，正如宗徒所说，她的光荣是在她自己内，而不是在别人的话里：\BibleQuote{各人应当考验自己的行为，这样，他所夸的（欢欣），就只在自己身上，而不在别人身上。}\BibleRef{迦 6:4} 但天主不许她的光荣在自己内，而不在那位圣咏作者对之说的那一位内：\BibleQuote{你是我的光荣，是你使我昂首。}因为，当内在于她的天主自己就是她的光荣时，她在自己身上的光荣才是有益的；她从祂获得一切美善，借此她成为善的，并必将获得一切使她变得更好的事物，按照她在此生能变得更好的程度，也借着天主的恩宠而非人的赞美成为完全的时，获享一切。\BibleQuote{愿她的灵魂因上主而受赞美，} \BibleQuote{以美物满足她心愿的那位，}因为正是祂启示了这渴望，为使祂的贞女不会因任何善而自夸，仿佛她不是领受的。\par

10. 那么，请你们回信告知，我们如此判断，认为这些是令嫒的看法，是否正确。因为我们清楚地知道，您和您全家过去和现在都是\OriginalTerm{不可分割的天主圣三}{indivisible Trinity}的朝拜者。然而，人的错误除了在关于不可分割的天主圣三的错误意见之外，还会以其他形式潜入。还有其他主题，人们在这些主题上陷入非常危险的错误。正如我们在本信中更详细地谈论的那个主题，或许对于像您这样坚定而纯洁的智慧，已经说得太多。然而，我们不知道，那否认由天主而来的善是源自天主的人，是在伤害谁，除非是伤害天主，因此也就是伤害圣三；愿天主使你们免于这种邪恶，正如我们所相信的上主确是如此！愿天主绝对禁止那本书——我们认为有责任从中摘录一些话，以便它们能被更容易地理解——产生这样的印象，我们不说是在您的思想上，或在圣洁的贞女令嫒的思想上，而说是在您家中最末等的男女仆人的思想上。\par

11. 但即使你更仔细地研读那些作者似乎为恩宠或天主的助佑辩护的话语，你也会发现它们模棱两可，既可以指本性，可以指知识，也可以指罪过的赦免。因为，甚至在他们不得不承认的那一点——即我们应当祈求不要陷于诱惑——上，他们可能会认为，这话的意思是我们受到帮助，通过祈祷和叩门，真理的知识如此启示给我们，使我们能学到我们应尽的本分，但我们的意志并未因此获得力量，从而能践行我们所学的本分；至于他们说，是由于天主的恩宠或助佑，主基督才被立为圣洁生活的榜样，他们对此的解释，无非是教导同样的教义，即肯定地说，我们通过祂的榜样学习应该如何生活，但否认我们受到如此的助佑，以致能因爱慕而践行从学习所知的。\par

12. 请在这本书中找一找，如果你可以的话，除了本性、意志的自由（这属于同一本性）、罪的赦免和教义的启示以外，是否承认有任何天主的神助，是如同那位说以下这话的人所承认的：\BibleQuote{我觉知，除非天主赐下恩赐，无人能自制；明白这恩赐属谁，也是一种智慧，我便转向主，向祂全心祈求。}\BibleRef{智 8:21} 因为他并未通过祈祷求得他受造的本性；也并未急切想要获得他受造时拥有的意志的自然自由；也并未渴望罪的赦免，因为他祈求自制以免犯罪；也并未渴望知道应做何事，因为他已承认他知道这自制之恩是谁的恩赐；但他渴望从\OriginalTerm{智慧的圣神}{Spirit of wisdom}那里获得这样的意志力量、这样的爱火，足以完全实践这伟大的自制之德。因此，如果你成功在那本书中找到任何这样的陈述，你若在回复中屈尊告知我们，我们将衷心感谢。\par

13. 我们无法形容我们是多么渴望在这些人的著作中找到对那恩宠的公开宣认，这恩宠是宗徒竭力赞美的，他说：\BibleQuote{天主把信德的度量分给了每个人，}\BibleRef{罗 12:3}\BibleQuote{缺少它，不可能中悦于天主，}\BibleRef{希 11:6}\BibleQuote{义人因之而生活，}\BibleRef{罗 1:17}\BibleQuote{它藉着爱德行事，}\BibleRef{迦 5:6}在它之前和没有它的情况下，任何人的任何行为都不能被算为善，因为\BibleQuote{凡不出于信德的行为，就是罪恶。}\BibleRef{罗 14:23}他肯定天主把信德分给每个人，\BibleRef{罗 12:3}并且我们领受神圣的助佑以虔诚和正义地生活，不仅是借着那知识的启示——这知识若无爱德会\BibleQuote{叫人自高自大，}\BibleRef{格前 8:1}——而且是借着我们被灌输以那种\BibleQuote{爱，这爱就是法律的满全，}\BibleRef{罗 13:10}这爱如此建树我们的心，使知识不致叫人自高自大。但迄今为止，我未能在这些人的著作中找到任何这样的陈述。\par

14. 但我们尤其希望这些观点应该在那本书中找到，我们曾从那本书中引用过那些话，作者在其中赞扬基督的一位贞女，仿佛除她自己之外没有人能将属灵的财富赋予她，又仿佛这些财富除她之外不能存在，他不愿她因主而夸耀，倒愿她像没有领受过这些财富似的夸耀。在这本书中，尽管既没有提他的名字，也没有提您尊敬的名字，他却提到贞女的母亲曾请求他给她写信。然而，在他的一封书信中，他公开署上了自己的名字，并且没有隐瞒那位神圣贞女的名字，同一位\Person{白拉奇}说他曾给她写过信，并试图引用该作品来证明，他极其公开地宣认了天主的恩宠，而这恩宠据称他曾沉默不谈或加以否认。但我们恳请您屈尊在您的答复中告知我们，那是否就是那本他插入这些话、论及属灵财富的书，以及这本书是否已达于您的\OriginalTerm{圣德}{Holiness}。\par

\SourceRecord{Translated by J.G. Cunningham.；Nicene and Post-Nicene Fathers, First Series；Vol. 1.；Edited by Philip Schaff.；Buffalo, NY: Christian Literature Publishing Co.,；1887.；<http://www.newadvent.org/fathers/1102188.htm>.}

% structure: p0001 source=h1 target=section reason=page-title

\section{第189号书信（主历418年）}

\Emph{致\Person{博尼法斯}，我尊贵的主、公正卓越且可敬的儿子，\Person{奥古斯丁}在主内致以问候。}\par

1. 我已给你的\OriginalTerm{爱德}{Charity}写了回信，但当我等待机会转交信件时，我亲爱的儿子\Person{福斯图斯}在前往你卓越之处的路上来到这里。他收到我本打算由他带给你的仁心的信之后，向我说，你非常希望我写些能在我们的主耶稣基督内、你寄望于祂的那永恒救恩中建树你的东西。虽然我当时正忙于事务，但他以你所知他对你的爱之恳切，坚持要我不迟延地完成。因此，为了他的方便，因他急欲离开，我觉得即便没有多少时间思考，也最好写一封，而不是推迟满足你的虔诚愿望，我尊贵的主、公正卓越且可敬的儿子。\par

2. 所有这一切都包含在这简短的几句话中：\BibleQuote{你当全心、全灵、全力爱上主你的天主；并爱近人如你自己。}\BibleRef{玛 22:37}因为这是主在世时对宗教的总括，祂在福音中说：\BibleQuote{全部法律和先知，都系于这两条诫命。}因此，你要藉着祈祷和行善，在这爱上日日前进，好能借着那命令你且赐予你这爱的祂的助佑，使这爱得到滋养和增长，直至达到成全，使你成为成全的。\BibleQuote{因为这就是爱，}正如宗徒所说，\BibleQuote{借着所赐与我们的圣神，已倾注在我们心中了。}\BibleRef{罗 5:5}这就是\BibleQuote{法律的满全；}\BibleRef{罗 13:10}这就是以爱德行事的信德，关于这信德他又说：\BibleQuote{割损或不割损都算不得什么，惟有以爱德行事的信德才有价值。}\BibleRef{迦 5:6}\par

3. 因着这爱，我们所有的\OriginalTerm{圣者}{holy fathers}、\OriginalTerm{圣祖}{patriarchs}、先知和宗徒们都得到天主的喜悦。因着这爱，所有真正的殉道者都力战魔鬼，甚至流血，而因为它（这爱）在他们内既不冷淡也不失败，他们便成为战胜者。因着这爱，所有真正的信徒日复一日地前进，力求获得的不是地上的国，而是天上的国；不是暂时的，而是永恒的产业；不是金银，而是天使们不朽的财富；不是今世那些带着战栗享用、无人死时能带走的世物，而是得见天主，祂的恩宠和赐福能力超越地上乃至天上一切形体的美好，超越人间一切无论多么公义圣洁的灵性优美，超越上界众天使和掌权者的一切荣耀，不但超越言语所能表达的，也超越思想能够想像的关于祂的一切。我们也不要因如此伟大的许诺而绝望于其实现，因为它是极其伟大的，而更该相信我们将得到它，因为那许诺者是极其伟大的，正如圣宗徒若望所说：\BibleQuote{亲爱的，现在我们是天主的子女；我们将来如何，还没有显明；但我们知道，当祂显现时，我们要相似祂，因为我们要看见祂实在怎样。}\BibleRef{若一 3:2}\par

4. 不要以为任何人在从事现役军务时便不能取悦天主。在这些人中，有圣\Person{达味}，天主对他给予了极大的赞许；在他们中，也有那时代的许多义人；还有那位百夫长，他对主说：\BibleQuote{主！我不堪当你到我的舍下来，但只要你一句话，我的仆人就会好的。因为我也是属人权下的，有士兵在我下；我对这个说：去，他便去；对那个说：来，他便来；对我的仆人说：做这个，他便做。}主论及他说：\BibleQuote{我实在告诉你们，在以色列中，我没有遇到过这样大的信德。}\BibleRef{玛 8:8}在他们中，还有那\Person{科尔乃略}，一位天使对他说：\BibleQuote{科尔乃略，你的施舍已蒙记念，你的祈祷已蒙听闻，}\BibleRef{宗 10:4}并指示他派人去见圣宗徒\Person{伯多禄}，听从他指教你该做什么；为此他派遣一位虔诚士兵去请伯多禄来。在他们中，还有那些来受洗于若翰的士兵——若翰是主的神圣前驱，新郎的朋友，主论他说：\BibleQuote{在妇女所生者中，没有兴起一位比洗者若翰更大的，}\BibleRef{玛 11:11}——他们曾问若翰他们该做什么，得到的回答是：\BibleQuote{不要对任何人施暴，也不要诬告任何人，并要满足于你们的粮饷。}\BibleRef{路 3:14}确实，当他命令他们满足于服役的粮饷时，并未禁止他们当兵。\par

5. 那些放弃所有世俗职务、以最严格的贞洁事奉天主的人，确实在天主前拥有更高的地位；但正如宗徒所说：\BibleQuote{每人都有他各自从天主得来的恩宠：有人这样，有人那样。}\BibleRef{格前 7:7}因此，有些人在为你们祈祷时，是与你们不可见的仇敌作战；你们在为他们作战时，是与蛮族，即可见的仇敌争战。但愿众人同有一个信德，因为那样，疲乏的挣扎便会减少，而魔鬼及其使者也会更容易被战胜了；但既然在此生中，天国的子民必须处于迷途者和不虔敬者中间经受试探，好使他们得到锻炼，并\BibleQuote{像在熔炉中的黄金一样受到考验，}\BibleRef{智 3:6}我们就不该在预定的时间之前渴望只与那些圣善和义人生活，而要通过忍耐，使我们配在适当的时机领受这真福。\par

6. 那么，当你准备征战的时候，首先要想到：即使你身体的力量也是天主的恩赐；想到了这一点，你便不会用天主的恩赐来反对天主。因为，既然信义已经作出承诺，即使是对自己与之交战的敌人，也应当遵守，更何况是对自己为之战斗的朋友呢！和平应当是你渴望的目标；战争只应出于不得已而进行，且进行战争只是为了天主能借以将人从不得已的状态中解救出来，并在和平中保全他们。因为寻求和平不是为了燃起战火，发动战争却是为了求得和平。所以，即使在作战时，也要怀有缔造和平的精神，好使你们因战胜你们所攻击的人，而引领他们回归和平的益处；因为我们的主说：\BibleQuote{缔造和平的人是有福的，因为他们要称为天主的子女。} \BibleRef{玛 5:9} 如果说人与人之间的和平如此甘饴，因为它能带来现世的安全，那么与天主之间的和平就更加甘饴，因为它为人类带来天使的永恒福乐！因此，致死与你交战的敌人，要出于不得已，而非出于你的意愿。正如对反叛与抵抗的人加以武力，对战败或被俘的人却应施行怜悯，尤其是在无需担心将来和平再受扰乱的状况下。\par

7. 你的生活仪表要以贞洁、节制和适度为装饰；因为若一个人在人们眼中是不可战胜的，却被贪欲所制伏；刀剑无法伤害他，却被美酒所击败，这实为极可耻的事。至于世间的财富，你若没有，不可藉作恶在世上追求；你若有，便当通过善工积蓄在天上。大丈夫气概和基督徒的精神，既不该因得着这世界的财宝而趾高气扬，也不该因失去它们而一蹶不振。我们更当思念主所说的话：\BibleQuote{你的财宝在那里，你的心也必在那里；} \BibleRef{玛 6:21} 而且当然，每当我们听到举心向上的劝导时，我们有责任诚恳地作出那你知道我们惯常所作的回应。\par

8. 的确，在这些事上，我知道你非常谨慎；我听到有关你的良好报告，使我满心欢喜，并促使我因你而在主内向你祝贺。所以，这封信与其说是指导你当如何行的指南，不如说是一面镜子，让你从中看到自己的本相；然而，无论你从这封信或从圣经中发现，自己在圣洁生活上还有所欠缺，你都要恒心，既通过努力，也通过祈祷，去恳切寻求；对于你已经拥有的，你要感谢天主，他是美善的泉源，你由他领受了这一切；在一切善工中，应将光荣归于天主，并要躬行谦卑，因为经上记载：\BibleQuote{一切美好的赠与，一切完善的恩赐，都是从上，从光明之父降下来的。} \BibleRef{雅 1:17} 但是，无论你在爱天主及爱你的近人以及真正热心方面有多大的进步，只要你还在此生，切莫想象自己没有罪过，因为关于这事，我们在圣经上读到：\BibleQuote{人生在世，岂非受试探的生活吗？} 因此，既然只要你还在这肉躯内，就常必须在祈祷中像主教导我们的那样说：\BibleQuote{宽免我们的罪债，如同我们也宽免得罪我们的人；} \BibleRef{玛 6:12} 那么，如果有人亏负了你且求你宽恕，你就要赶快宽免，好使你能真诚地祈祷，并能有效地为自己求得罪过的宽恕。\par

这些事，我亲爱的朋友，我匆忙写给你，因为送信的人急于动身，催我勿要耽搁；但我感谢天主，我总算多少满全了你热诚的愿望。愿天主的仁慈常护佑你，我高贵的主，我理应杰出的孩子。\par

\SourceRecord{Translated by J.G. Cunningham.；Nicene and Post-Nicene Fathers, First Series；Vol. 1.；Edited by Philip Schaff.；Buffalo, NY: Christian Literature Publishing Co.,；1887.；<http://www.newadvent.org/fathers/1102189.htm>.}

% structure: p0001 source=h1 target=section reason=page-title

\section{书信191（主历418年）}

\Emph{致我可敬的主和圣善的弟兄与同司铎\Person{西克斯图斯}}，\Emph{堪当在基督的爱中被接纳者，\Person{奥古斯丁}在主内向您致意。}\par

1. 自从我不在时，您通过我们圣善的弟兄司铎\Person{斐尔慕斯}转来的那封信到达后，我在返回\Place{希波}后才读到，但那时送信人已经离去；因此，现在是我首次有机会给您任何回信，我极其乐意将这信托付给我们至爱的儿子、\OriginalTerm{辅祭员}{acolyte}\Person{阿尔比努斯}。您的信是寄给\Person{阿利比乌斯}和我共同收的，那时我们不在一起，这就是为什么您现在会收到我们每人一封回信，而不是一封我们两人共同写的回信。因为这封信的送信人刚刚从我这里离开，去探望我可敬的弟兄和同主教\Person{阿利比乌斯}，他会亲自给您的圣座写一封回信，他带去了您的信，我已经细读过。至於那封信使我心中充满的极大喜乐，人岂能试图言说那无法表达的呢？诚然，我不认为您自己充分了解您寄来那封信给我们所带来的益处有多大；但请相信我们的话，因为正如您见证了您的感受，我们也见证我们的感受，表明我们如何深深地被那封信中所阐明的见解之全然透明正确的确凿性所感动。因为，当您通过辅祭员\Person{良}将一封关于同一主题的很短的信寄给至福的老\Person{奥勒留}时，我们曾喜乐敏捷地抄录，并满怀热情地读给所有能接触到的人听，作为您的思想的阐释，既关于那个最致命的\OriginalTerm{教义}{dogma} [贝拉基的]，也关于天主白白赐予大小众人的\OriginalTerm{恩宠}{grace}，而那个教义正与这恩宠截然相反；那么，您认为我们读到您笔下这更详尽的阐述时，是何等的喜乐，我们又是何等热切地关注著，要让它被所有我们已能或将来能让其知晓的人所阅读！因为，还有什么能比读到或听到一个人如此清晰地捍卫天主的恩宠，反对其仇敌，更能令人满足呢？而这个人以前曾被这些仇敌骄傲地宣称是其主张的有力支持者！或者，还有什么事比这更应使我们向天主献上更丰富的感谢呢？即：祂赐予恩宠的这些人，如此出色地捍卫这恩宠，反对那些未获赐恩宠的人，或反对那些虽获赐恩宠却不接受的人，因为按照天主隐秘而公义的审判，那接受恩宠的心情并未赐给他们。\par

2. 因此，我可敬的主和圣善的弟兄，堪当在基督的爱中被接纳者，虽然您这样写信给那些弟兄们时——那些对手们习惯于在他们面前夸口说您是他们朋友——您提供了最卓越的服务，然而，您还有更重大的责任，不仅要以有益的严厉惩罚那些敢于更公开地宣讲那对基督徒之名最具危险敌意的错谬的人，而且还要以牧人的警觉，为那些更软弱、更纯朴的主的羊群，极其谨慎地防范那些更为隐蔽、确实、胆怯但坚持不懈、仿佛轻声低语般教导这错谬的人，正如宗徒所说的\BibleQuote{潜入人家}，又熟练地以不虔敬之心做著那经文中紧接着提及的其他一切事。\BibleRef{弟后 3:6} 那些因恐惧束缚而在最深沉的沉默中隐藏其思想的人，也不应被忽视，他们却仍然像以前一样未停止怀揣那些错谬的见解。因为，在\OriginalTerm{宗座}{apostolic see}最明确的谴责宣布那瘟疫之前，您可能认识他们中的一些人，现在突然发现他们沉默无声了；也无法确定他们是否真的痊愈了，除非他们不仅不再讲说那些错误的教义，而且也以他们以往支持另一方的同样热情，来捍卫那些与他们以往的错谬相反的真理。但是，这些人应当更温和地对待；因为对于那些其沉默本身就证明已足够惊惧的人，何必再引起他们更深的恐惧呢？同时，他们虽不该受惊骇，却该受教导；依我之见，既然他们对於严厉的恐惧有助于真理的教导者的工作，他们就会更容易地，靠着主的助佑，在他们学会理解并爱慕他的恩宠之后，大声地说出反对他们目前还不敢明认的错谬。\par

\SourceRecord{Translated by J.G. Cunningham.；Nicene and Post-Nicene Fathers, First Series；Vol. 1.；Edited by Philip Schaff.；Buffalo, NY: Christian Literature Publishing Co.,；1887.；<http://www.newadvent.org/fathers/1102191.htm>.}

% structure: p0001 source=h1 target=section reason=page-title

\section{第192封信（公元418年）}

\Emph{致我可敬的主教、最尊贵且圣洁的弟兄\Person{策肋定}}，\Emph{\Person{奥斯定}在主内致候}。\par

1. 当宗座的信函经由书记\Person{普罗耶克托斯}之手送达位于\Place{希波}的我时，我正离家甚远。当我返回家中，读了你的信，感到我应该给你回信时，我仍在等待与你通信的某种方式，恰在此时，看！一个极好的机会随着我们最亲爱的弟兄、\OriginalTerm{辅祭}{acolyte}\Person{阿尔比努斯}的离去而出现，他即将立刻离开我们。因此，我因你的健康而欢欣——你的健康是我最恳切盼望的——我向宗座回致我欠下的问候。但是我总是欠你爱德，这是唯一的债，即使已经偿还，仍使偿还者成为负债者。因为当它偿还时便是给予了，但它即使在给予之后仍被亏欠，因为没有一刻它不再被亏欠。当它给予时也没有失去，反而是通过给予而增多；因为是通过拥有它，而非舍弃它，才给予它。既然除非拥有便不能给予，同样除非给予也不能拥有；不仅如此，当一个人给予它时，它就在那人内增长，并根据给予之人的数量，所获得的份量变得更大。此外，甚至对敌人所亏欠的，怎能拒绝给朋友呢？然而，对敌人付这份债是慎重的，对朋友则是怀着信赖偿还。尽管如此，它尽一切努力确保它收回它所给予的，即使对那些以善报恶的人。因为我们希望那作为敌人而真正去爱的那个人成为朋友，除非我们希望他成为善人，否则我们并非真诚地爱他，而除非他从怀恨的罪恶中被解脱出来，他便不能成为善人。\par

2. 因此，爱德与金钱偿还的方式不同；因为金钱在付出时减少，爱德却因付出而增长。它们也有此区别——我们若不想索回金钱，便表示对那人有更大的善意；但没有人能成为爱德的真正施予者，除非他以爱来要求偿还。因为金钱，当它被收下时，归于受予者所有，却远离赐予者；爱德则相反，即使没有被偿还，它仍随着要求其所爱之人偿还的人而增长；不但如此，那偿还爱德之人，除非他再付出，否则不会开始拥有它。\par

因此，我的主教与弟兄，我欣然给予你，并喜乐地接受你我彼此亏欠的爱德。我所接受的爱德我仍要求，我所给予的爱德我仍亏欠。因为我们当以顺服之心遵守独一师尊的诫命——我们都自认是祂的门徒——祂借着祂的宗徒对我们说：\BibleQuote{你们除了彼此相爱之外，不可再欠人什么。}\BibleRef{罗 13:8}\par

\SourceRecord{Translated by J.G. Cunningham.；Nicene and Post-Nicene Fathers, First Series；Vol. 1.；Edited by Philip Schaff.；Buffalo, NY: Christian Literature Publishing Co.,；1887.；<http://www.newadvent.org/fathers/1102192.htm>.}

% structure: p0001 source=h1 target=section reason=page-title

\section{热罗尼莫致奥斯定（主后418年）}

\Emph{致他的圣主及最蒙福之父}, \Emph{\Person{奥斯定}，\Person{热罗尼莫}谨致问候。}\par

一直以来，我都以相称的尊敬与荣耀敬重您的真福，并爱那居于您内的主和救主。但现在，如若可能，我们要为那已臻极点的再增添些许，并为那已经满盈的再堆聚些许，致使我们不容一个时辰逝去而不提及您的名字，因为您怀着坚定不移的\OriginalTerm{信德}{faith}之热忱，在敌对的暴风雨中坚守阵地，并宁愿——只要还在您能力范围内——从\Place{索多玛}被解救出来，哪怕您独自出来，也胜过逗留在那些注定要灭亡的人后面。您的智慧领悟了我意欲表达的事。继续前进，蒸蒸日上吧！您闻名于普世；天主教徒尊崇您，仰视您为古\OriginalTerm{信德}{faith}的恢复者，并且——这更是赫赫荣耀的标志——所有\OriginalTerm{异端者}{heretics}都憎恨您。他们以同样的仇恨迫害我，试图用诅咒夺去他们无法用刀剑触及的生命。愿基督我等主的仁慈保全您安康，并念及我，我可敬的主及最蒙福之父。\par

\SourceRecord{Translated by J.G. Cunningham.；Nicene and Post-Nicene Fathers, First Series；Vol. 1.；Edited by Philip Schaff.；Buffalo, NY: Christian Literature Publishing Co.,；1887.；<http://www.newadvent.org/fathers/1102195.htm>.}

% structure: p0001 source=h1 target=section reason=page-title

\section{书信201（公元419年）}

\Emph{皇帝\Person{奥诺留}·奥古斯都和皇帝\Person{狄奥多西}·奥古斯都向主教\Person{奥勒留}致意。}\par

1. 其实，早已颁令，邪恶异端的始作俑者\Person{白拉奇}与\Person{切莱斯提乌斯}，既为天主教真理的祸害腐化者，应被逐出\Place{罗马}城，以免他们以毒害影响败坏无知者的心灵。吾等仁慈之举在此秉承了尊座之判断；依据此判断，毫无疑问，他们在其主张经公正审查后遭一致谴责。然而，他们顽固不改，使敕令不得不再次颁布；吾等又藉近旨进一步规定：若有人明知他们在各行省某处藏匿，却迟延不将他们驱逐或告发，则此人以共犯论，应受规定之处罚。\par

2. 然而，为集合众人基督热忱之共力，以铲除此荒谬异端，至爱的父亲，最宜者，乃由尊座之权威施于若干主教之纠正，此辈或以缄默认同而支持彼等邪说，或避免公开反对而不加抨击。故此，请尊驾以适当文书，令所有主教受劝诫（一旦他们藉尊座之命令，得知此项指令加于彼等）：凡有人因不虔之顽固，怠于签署上述诸人之谴责以表白其教义之纯正者，将先受褫夺主教职位之罚，复遭绝罚，并被终生放逐出离其主教座。盖因吾等自身，凭着真诚之宣信，遵从\OriginalTerm{尼西亚大公会议}{Council of Nicaea}，敬拜天主为万有之创造者，兼吾皇权之源泉；尊座必不容忍此可憎派系之成员，发明新奇怪异之邪说以损害宗教，将公权所谴责之谬误，隐于私下流传之书卷中。盖因，至爱可亲的父亲啊，凡以伪善之默然暗助其谬误，或因缄默不斥责而予以致命赞同者，其异端之罪愆，绝不在轻减之列。\par

\Emph{（另一方笔迹。）}愿天主保佑你平安多年！\par

颁于\Place{拉文纳}，六月九日，\Person{摩纳西乌斯}与\Person{普林塔}执政之年。\par

一封相同措辞的信亦寄给了圣\Person{奥古斯丁}主教。\par

\SourceRecord{Translated by J.G. Cunningham.；Nicene and Post-Nicene Fathers, First Series；Vol. 1.；Edited by Philip Schaff.；Buffalo, NY: Christian Literature Publishing Co.,；1887.；<http://www.newadvent.org/fathers/1102201.htm>.}

% structure: p0001 source=h1 target=section reason=page-title

\section{\Person{热罗尼莫}致\Person{阿利比乌斯}与\Person{奥古斯丁}书（主后419年）}

\Emph{致\Person{阿利比乌斯}与\Person{奥古斯丁}主教，我至圣而堪受敬爱尊崇的主教们，\Person{热罗尼莫}在基督内致候。}\par

1. 圣司铎\Person{英诺森提乌斯}是这封信的携带者，他去年没有带我的信给你们尊驾，因为他没指望返回非洲。我们感谢天主，这反而给你们机会，\BibleQuote{以善胜恶}，用你们的书信弥补了我们的沉默。任何写信给崇敬的教父们的机会，对我都是最乐意的。我呼求天主作证，若可能，我愿取\BibleQuote{鸽子般的翅膀}，飞去投入你们的怀抱。我一向因你们可敬的品格而深爱你们，而如今尤其如此，因为你们的合作与领导成功地扼制了\Person{塞莱斯提乌斯}的异端，这异端毒害了许多人的心，以致他们虽感到被击败和被谴责，却仍未放下其毒恶思想，并因他们唯一尚存的能力，而憎恨我们，因他们以为是我们的缘故，他们才失去了教导异端道理的自由。\par

2. 至于你们询问我是否已著书反对\Person{阿尼安努斯}的书，此人自称是\Place{塞莱德阿}的执事，享受着充裕的供给，只为了让他去散布别人异端思想的轻浮记述。你们当知，我收到这些书不久，是由我们圣洁的弟兄、司铎\Person{优西比乌斯}以散页形式寄来的。自那以来，因病患的攻击，以及你们杰出而圣洁的女儿\Person{欧斯托基雍}的安息，我备受煎熬，几乎想以轻蔑的沉默放过这些作品。因为他从头至尾在同样的泥潭中挣扎，除了一些并非他原创的叮当作响的词句外，全是老生常谈。然而，我大大得益，因他试图回应我的信时，更不保留地表白了他的观点，并将其亵渎之言公之于众；因为他在\Place{狄奥斯波利斯}那场可悲的主教会议中所否认的每一错误，在这篇论文中都公开承认了。回答他那极其愚蠢的幼稚言论，实不足道；但若主保全我，又有足够多的助手为我抄写，我将以几篇简短的夜作回答他，不是为了驳斥一个已死的异端，而是以论证封住他的无知与亵渎：此事您的尊驾远比我做得更好，因为您可免去我向异端者写作时自我夸赞的必要。我们圣洁的女儿\Person{阿尔比娜}和\Person{梅拉尼娅}，以及我们的儿子\Person{皮尼安努斯}，衷心地问候你们。我从圣地\Place{白冷}，请圣司铎\Person{英诺森提乌斯}把这封短笺带给你们。您的侄女\Person{保拉}恳切地求您记念她，并向您热烈致意。愿我们的主耶稣基督的慈悲保守你们安康并记念我，我的真正圣洁的主教，可敬可爱的教父们。\par

\SourceRecord{Translated by J.G. Cunningham.；Nicene and Post-Nicene Fathers, First Series；Vol. 1.；Edited by Philip Schaff.；Buffalo, NY: Christian Literature Publishing Co.,；1887.；<http://www.newadvent.org/fathers/1102202.htm>.}

% structure: p0001 source=h1 target=section reason=page-title

\section{书信二〇三（主后四二〇年）}

\Emph{致我尊贵的阁下、最卓越且亲爱的儿子\Person{拉尔古斯}，\Person{奥古斯丁}在主内致意。}\par

我收到了阁下您的来信，您在其中请求我给您写信。您必定不会这样做，除非您认为我所能写给您的，是可接受的、令人愉悦的。换言之，我假定，当您尚未尝试过此生的虚幻时，曾渴望它们；如今，经过尝试之后，您已鄙视它们，因为在其中，快乐是欺骗的，劳苦是无益的，焦虑是不断的，高升是危险的。人们起初因轻率而寻求它们，最终却带着失望和懊悔而放弃它们。对于在此尘世生活的忧虑中，世人以其不安的欲求，不是以智慧，而是以更大的急切所贪求的一切，这话都是真实的。但虔敬者的\OriginalTerm{望德}{hope}则全然不同：他们劳苦的果实大不相同，他们危险的赏报大不相同。只要我们活在此世，恐惧、忧伤、劳苦和危险就是不可避免的；然而重大的问题是，一个人因何缘故、带着何种期望、以何种目标而忍受这些。当我默观此世的贪恋者时，我不知道在什么时刻，智慧能最适宜地尝试改善他们的道德；因为当他们拥有表面的顺境时，他们轻蔑地拒绝她有益的规劝，视之为老妇的无稽之谈；当他们处于逆境时，他们想的多半只是逃避当前的痛苦，而非获得真实的救药，借此得以治愈，并抵达那免除一切痛苦的地方。然而，偶尔却有人向真理敞开双耳和心灵——在顺境中罕见，在逆境中较常。这些人确实是少数，因为经上预言他们将是少数。我希望您也在这些人之中，因为我真诚地爱您，我尊贵的阁下、最卓越且亲爱的儿子。但愿这番劝言作为我回信的答复，因为虽然我不愿您从此再忍受您所曾忍受的一切，然而我更为伤痛的是，若您经历了这一切，而未使您的生活有任何改善。\par

\SourceRecord{Translated by J.G. Cunningham.；Nicene and Post-Nicene Fathers, First Series；Vol. 1.；Edited by Philip Schaff.；Buffalo, NY: Christian Literature Publishing Co.,；1887.；<http://www.newadvent.org/fathers/1102203.htm>.}

% structure: p0001 source=h1 target=section reason=page-title

\section{第208封信（公元423年）}

\Emph{致\Person{斐莉齐亚}女士，他在信德内的女儿，在基督的肢体中堪受尊敬者，\Person{奥斯定}在主内致候。}\par

1. 我不怀疑，当我想到你的信德，以及他人的软弱或邪恶时，你的心已被扰乱，因为即使是一位充满怜悯之爱的圣宗徒，也承认类似的经验，说：\BibleQuote{谁软弱，我不软弱？谁跌倒，我不心焦？}\BibleRef{格后 11:29} 因此，既然我分担你的痛苦，并在基督内挂念你的福祉，我认为有责任以此信部分给予安慰、部分给予劝勉，致书于你的圣德，因为在我们主耶稣基督的身体内，祂的所有肢体都合而为一，你在这身体内与我们紧密相连，作为那身体中可敬的肢体被爱着，并藉着祂的圣神与我们共享生命。\par

2. 我劝你，不要过于因那些恶表而烦扰，这些恶表正是为此目的被预言必要来临，好叫它们来到时，我们能记起它们曾被预言，而不致过于惊惶。因为主自己在福音中预言了它们，说：\BibleQuote{世界因了恶表是有祸的！恶表固然免不了要来，但立恶表的那个人是有祸的！}\BibleRef{玛 18:7} 这些人就是宗徒所说的：\BibleQuote{他们只求自己的事，而不求耶稣基督的事。}\BibleRef{斐 2:21} 因此，有些人担任牧人的尊贵职务，是为了照顾基督的羊群；另一些人占据那职位，是为了享受与职务相连的世俗尊荣和利益。这两种牧人，一些去世，另一些继任，必将继续存在于天主教会内，直到世末及主的审判。如果在宗徒时代已有这样的人，\Person{保禄}因他们的行为而忧伤，在列举自己的磨难时提到：\BibleQuote{假弟兄中的危险，}\BibleRef{格前 11:26} 然而他并没有傲慢地将他们驱逐，而是耐心地容忍他们；那么在我们的时代，又当有多少这样的人兴起呢？因为主极其明确地论及这个接近终结的时代说：\BibleQuote{由于罪恶增多，许多人的爱德将要冷淡。}\BibleRef{玛 24:12} 然而，随后的话应给我们安慰和劝勉，因为祂又加上一句：\BibleQuote{唯独坚持到底的，才可得救。}\par

3. 再者，既有好牧人和坏牧人，羊群中也有好有坏。好的以\Quote{羊}的名称代表，坏的则被称为\Quote{山羊}。然而，它们在同一个草场上并肩吃草，直到首席牧者——即被称为唯一牧者的那位——来临，按祂的应许\Quote{如同牧人分开绵羊和山羊一般}，将它们彼此分开。祂将聚集羊群的责任放在我们身上，而把最终分开的工作保留给自己，因为这事只属于那位不会错误的；因为那些僭越的仆人，竟敢在主留在自己权下的时刻之前就进行分开，结果不但没有分开别人，自己反而与公教的合一分离；因为那些因分裂而成为不洁的人，怎能拥有洁净的羊群呢？\par

4. 因此，为使我们能保持在信德的合一中，不因糠秕所招致的恶表而在主的打谷场上绊倒，反要如同麦子坚持到最后的扬净，\BibleRef{玛 3:12} 并藉着那使我们稳定的爱德容忍被打的草秸，我们的大牧者便在福音中论及好牧人时告诫我们，不要因他们的善工而寄望于他们，却要因天父使他们成为如此而归光荣于天父；论及坏牧人（祂用经师和法利塞人的名称予以暗指），则提醒我们：他们讲论的是善事，所行的却是恶事。\BibleRef{玛 3:12}\par

5. 论及好牧人，祂这样说道：\BibleQuote{你们是世界的光；建在山上的城，是不能隐藏的。人点灯，并不放在斗底下，而是放在灯台上，照耀屋中所有的人。照样，你们的光也当在人前照耀，好使他们看见你们的善行，光荣你们在天之父。} 论及坏牧人，祂用这些话劝诫羊群：\BibleQuote{经师和法利塞人坐在\Person{梅瑟}的讲座上：凡他们对你们所说的，你们要行要守；但不要照他们的行为去做，因为他们只说不做。}\BibleRef{玛 23:2} 当这些话被聆听时，基督的羊即使透过邪恶的教师，也听见祂的声音，而没有离弃祂羊群的合一，因为他们听到他们所教导的善，并不属于牧人，而是属于祂，因此羊群得到安全的喂养，因为即使在坏牧人手下，他们仍在上主的牧场上得到滋养。然而，他们并不效法坏牧人的行为，因为那些行为不属于世界，而是属于牧人自己。至于那些他们视为好牧人的人，他们不仅听从他们所教导的善事，也效法他们所行的善工。宗徒\Person{保禄}就属于这一类，他说：\BibleQuote{你们该效法我，如同我效法了基督一样。}\BibleRef{格前 11:1} 他是被永恒之光，即主耶稣基督自己点燃的光，并被放在灯台上，因为他以十字架夸耀，论及这十字架他说：\BibleQuote{至于我，我只以我们的主耶稣基督的十字架来夸耀。}\BibleRef{迦 6:14} 再者，既然他不求自己的事，而只求耶稣基督的事，当他劝勉那些他\BibleQuote{藉着福音所生的}\BibleRef{格前 4:15}人效法自己的生活时，他仍然严厉责斥那些以宗徒的名义制造分裂的人，并叱责那些说\BibleQuote{我是属\Person{保禄}的，难道\Person{保禄}为你们被钉在十字架上吗？或者你们受洗是归于\Person{保禄}的名下吗？}\BibleRef{格前 1:12}的人。\par

6. 因此，我们明白，善牧们是那些不求私利，只求耶稣基督之事的人；而好羊虽然效法那些善牧的善工——他们也靠这些牧者的服务得以聚集——却不把希望寄托在他们身上，而是寄托在用自己的血救赎了他们的主身上。这样，当他们偶然处于宣扬基督教导却自己行恶的恶牧之下时，他们会实行他们所教导的，却不效法他们的恶行，也不会因为这些邪恶之子而离开唯一真教会的牧场。因为在天主教会内既有善人也有恶人，这教会不像\OriginalTerm{多纳图斯派}{Donatist sect}，它不只限于非洲，而是扩展及天下万国，正如宗徒所说：\Quote{在全世界不断结果、增长。} 但那些与教会分离的人，只要他们与教会对立，就不可能是善的。尽管他们中有些人因看似可称赞的品行，似乎证明他们为善，但他们与教会本身分离这一事实，已使他们成为恶的，正如主的话说：\BibleQuote{不与我相合的，就是敌我的；不同我收集的，就是分散的。}\BibleRef{玛 12:30}\par

7. 因此，我的女儿，在基督的肢体中堪受一切欢迎和尊敬的，我劝你要忠信地持守主所托付给你的，并全心爱慕祂和祂的教会，因祂不容你与迷失者联合，以致丧失你\OriginalTerm{童贞}{virginity}的赏报，或与他们一同丧亡。因为如果你离开此世时，与基督身体的合一分离，那么即使你保全了童贞的完整，也毫无益处。但富于仁慈的天主，为你成就了福音上所记载的：当被邀的宾客向家主推辞时，他对仆人说：\BibleQuote{你们出去，到大路上和篱笆边，凡遇到的，都勉强他们进来。}\BibleRef{玛 22:9} 然而，虽然你应当对那些藉他们之手把你勉强领进来的祂的忠仆怀有最诚挚的感激，但你的责任仍是要把希望寄托在那预备了宴席的祂身上，也因祂的劝导你才来赴永生和真福的生命。将你的心、你的圣愿、你的神圣童贞，以及你的信德、望德和爱德都交托给祂，你便不会被那直到终结必将增多的恶表所动摇；反而，凭着不可动摇的虔敬之力，在主内获得安全，并在主内凯旋，坚守祂身体的合一，直到终结。请回信告诉我，你如何看待我对你的挂虑，我已在信中用尽己力表达了。愿天主的恩宠和仁慈永远护佑你！\par

\SourceRecord{Translated by J.G. Cunningham.；Nicene and Post-Nicene Fathers, First Series；Vol. 1.；Edited by Philip Schaff.；Buffalo, NY: Christian Literature Publishing Co.,；1887.；<http://www.newadvent.org/fathers/1102208.htm>.}

% structure: p0001 source=h1 target=section reason=page-title

\section{书信 209（公元 423 年）}

\Emph{致\Person{塞莱斯廷}}，\Emph{我最蒙福之主，并以一切应有之敬爱受尊崇的圣父，\Person{奥古斯丁}在主内问候。}\par

1. 首先，我庆贺你，因为我们听说，我们的主天主已立你于那显赫的宗座之上，且在祂的子民中毫无分裂。接下来，我将我们这里的情况陈明于圣座前，使你不仅能以祈祷，更能以建议和援助来帮助我们。因为我此刻是在深深的痛苦中给你写信，由于我本想造福于我们邻近地区基督的某些肢体，却因我行事缺乏审慎与警惕，给他们带来了一场大灾难。\par

2. 与\Place{希波}地区接壤处，有一个小镇，名叫\Place{福萨拉}：以前那里没有主教，但它连同邻近地区一起被包括在\Place{希波}的堂区内。该地区天主教徒很少；\OriginalTerm{多纳特派}{Donatists}的谬误以其可悲的影响力控制着所有其他位于人口众多之处的会众，以致在\Place{福萨拉}镇本身竟没有一个天主教徒。仰赖天主的仁慈，所有这些地方都被引领归附于教会的合一；我们原先任命去聚集这些民众进入羊栈的\OriginalTerm{司铎们}{presbyters}曾遭受抢劫、殴打、伤残、失明，甚至被杀；若要细述其间经历了多少辛劳与多少危险，说来话长。然而，他们的苦难并非徒劳无果，因为藉由他们，合一的恢复得以成就。但由于\Place{福萨拉}距离\Place{希波}有四十英里之遥，而我看出，在管理其民众，并聚集那分散四处的、无论其人数如何稀少、并非威胁他人而是为自身安全而逃难的男女人等中的残余者时，我的工作范围将超出应有的限度，且我无法给予我清楚看到所必需的关注，于是我便安排在那里祝圣并任命一位主教。\par

3. 为实施此事，我寻找一位适合该地方与民众、同时又通晓\OriginalTerm{布匿语}{Punic}的人；我心里已有一位适合此职的司铎。我致函当时担任\Place{努米底亚}\OriginalTerm{首席主教}{Primate}的年高德劭的主教，获得他的同意，远道而来祝圣这位司铎。他到来之后，正当我们所有人的心思都专注于如此重大的事务时，到了最后时刻，那位我以为已准备好受祝圣的人却断然拒绝接受此职，令我们大失所望。那时，我自己——正如事态所表明的，我本该推迟而非草率从事如此危险的事——由于不愿让那位极其可敬而圣洁的老人，在如此疲惫地来到我们这里之后，不完成他远道而来的使命就返回家园，便在没有民众主动请求的情况下，向他们推荐了一位当时与我在一起的年轻人，\Person{安东尼乌斯}。他自幼在我们的\OriginalTerm{隐修院}{monastery}长大，但除了担任\OriginalTerm{读经员}{reader}外，对于教会职务中各级别所涉及的劳苦毫无经验。那些不幸的民众，不知随后将发生何事，出于对我的信任而顺从地接受了我的建议。我还需多言吗？事情就这样做了；他上任成为他们的主教。\par

4. 我该怎么办？我不愿在可敬的圣座前指控那由我引入羊栈、精心抚养的人；我也不愿抛弃那些我曾在恐惧焦虑之中、为将他们收归教会而辛劳过的羊群；而我发现无法对双方都做到公正。事情已发展到如此痛苦的危机地步：那些当初听从我的意愿、以为是在为自己谋取利益而选他为主教的人，现在竟到我面前控告他。其中最为严重的指控，即那些不是由他主教的子民、而是由某些其他人提出的严重不道德的指控，被发现全无证据支持，且在我们看来，他已完全摆脱了那些最无端归咎于他的罪行；因此，他不仅被我们自己，也被他人如此同情地看待，以致\Place{福萨拉}及周边地区民众所抱怨的事情——诸如不可容忍的专横、掠夺、勒索和各种压迫——决不显得如此严重，以致我们有必要因其中一项、甚或因所有这些项合在一起，就认为应当剥夺他的主教职务；在我们看来，只要他归还那些可被证明不公取走之物就足够了。\par

5. 最终，在我们的判决中，我们将仁慈与严厉如此结合，以致在为他保留主教职务的同时，我们并未让那些既不应当由他自己重犯、也不应被树为他人的榜样的过犯完全免受惩罚。因此，在纠正他时，我们为这年轻人保留了其职务的尊位不受减损，但同时，作为惩罚，我们剥夺了他的权力，规定他不得再管辖那些他曾以如此方式对待的人，因为他们正当的愤恨使他们无法服从他的权威，并且或许会因无法忍耐的愤怒而爆发某些暴力行为，使双方都陷于危险。这种情绪状态在主教们就\Person{安东尼乌斯}之事与他们交涉时显然表现出来，尽管目前那位显赫人士\Person{塞莱尔}——他曾抱怨此人对他施加了有力的干涉——无论在\Place{阿非利加}或其他地方都不再拥有权力。\par

6. 但我何必用更多细节耽搁您呢？我恳求您在这件费力的事上协助我们，真福阁下，神圣教宗，因您的虔敬而受尊敬，因应有的爱戴而被崇敬；请吩咐所有已转呈的文件向您朗读。请注意\Person{安东尼}以何种方式履行其主教职责；当为了\Place{弗萨拉}人得到完全赔偿之前他被禁止共融时，他如何服从了我们的判决，并现已留出一笔款项，以便在调查后支付被认为公正的赔偿，为使共融的特权得以恢复；他如何以狡诈的说辞说服了我们年迈的首席主教，一位极可敬的人，相信他的一切陈述，并向可敬的教宗\Person{波尼法爵}推荐他为全然无可指摘。但我何必复述其余一切，既然前述的可敬老人必定已将全部事件呈报给陛下您了呢？\par

7. 在记录我们对他判决的众多会议记录中，我本会担心我们在贵处看来可能作出了较应有的更为宽松的判决，若非我知道您如此倾向怜悯，以致您会认为您有责任不仅宽恕我们，因为我们宽恕了他，也要宽恕他本人。然而我们所做的，无论是出于慈爱还是宽纵，他却试图利用，并作为我们判决的法律上的反对理由。他大胆地抗议：\Quote{要么我应坐在我自己的主教座上，要么根本不应是主教，}好像他现在正在坐的并非他自己的座位。因为，正是为此缘故，那些他曾任主教的地方被划分并指定给他，为免他被指责为非法调至另一教区，有违教父法规。抑或，人们应坚持如此严格的辩护，无论是为严苛还是为宽大，以致要求，对那些似乎不应被免去主教职的人，不得施加任何惩罚；或是对所有似乎应受任何惩罚的人，一律宣判免职？\par

8. 有案可查，宗座曾或亲自判决，或确认他人的判决，认可某些决定，使某些犯罪者既不被免去主教职，也不被完全不加惩罚。我不必列举那些距今很久远的案例；我只提近期的事。让\Place{凯撒勒亚}省的\Person{普里斯库}主教大胆抗议：\Quote{要么首席主教职应向我开放，如对其他主教一般，要么我不应继续作主教。}让同省的\Person{维克托}主教，他与\Person{普里斯库}处于同一判决，除自己教区外，无任何主教与他共融，让他，我说，以同样的信心抗议：\Quote{要么我应在各处享有共融，要么我不应在我自己的区域内享有。}让同省的第三位主教\Person{劳伦提乌斯}说话，并能以与此人完全相同的话喊出：\Quote{要么我应坐在我被祝圣的座位上，要么我不应是主教。}然而，谁又能指责这些判决呢，除非是不考虑，一方面，并非所有罪过都应不加惩罚，另一方面，也并非所有罪过都应以同一方式惩罚的人。\par

9. 既然，至福教宗\Person{波尼法爵}在论及\Person{安东尼}主教时，在其书信中以牧者应具的警觉审慎，在他的判决中加入了附加条款：\Quote{倘若他已将事实如实向我们报告，}如今请接受那些他向你陈述时略过不提的事实，以及在那位永福之人书信在\Place{阿非利加}被宣读之后所发生的事，并在基督的怜悯中，向那些比他——他们渴望脱离其骚扰——更恳切地哀求您的人，伸出您的援手。因为，或者是从他本人，或者至少从甚嚣尘上的传言中，他们受到威胁，说司法法庭、政府当局和军队暴力，将要执行宗座的决定；其后果是，这些不幸的人，如今是公教信友，对公教主教的恐惧，竟超过了他们过去身为异端者时对公教皇帝法律的畏惧。我恳求你，借着基督的血，借着宗徒\Person{伯多禄}的记忆——他曾警告那些管理基督徒子民的人，不要暴力地\BibleQuote{主宰他们的兄弟}。\BibleRef{伯前 5:3}——不要让这些事发生。我将\Place{弗萨拉}的公教信友，我在基督内的孩子们，以及\Person{安东尼}主教，我在基督内的儿子，都托付于陛下您亲切的爱中，因为我爱他们双方，并将两者都托付于您。我不责怪\Place{弗萨拉}人民向您陈情，对我强加一个我未曾测试过、且至少年龄尚不成熟，却使他们如此受苦的人给他们而提出的正当控诉；我也不愿对\Person{安东尼}做出任何错待，我以与我对他的真诚爱意相称的决心，对抗他的邪恶贪婪。愿您的怜悯扩展至双方——对他们，使恶事不致加诸其身；对他，使他不致作恶：对他们，为使如果他们在公教主教们，尤其是宗座本身那里，寻不着庇护以抵御一位公教主教的侵犯，他们便不会憎恨\OriginalTerm{天主教会}{Catholic Church}；对他，则是为使他不要陷入如此严重的邪恶，以致将那些他企图违背其意愿而据为己有的人，与基督隔绝。\par

10. 至于我本人，我必须向\OriginalTerm{圣座}{Sanctitas}承认，在这威胁到双方的危难中，我焦虑悲伤至极，以至于想到要辞去主教职务的责任，并任由自己表现出与我的错误之严重程度相应的悲痛，如果我看到（由于我轻率投票支持其当选主教的那人的行为）天主的教会遭到破坏，甚至（愿天主禁止）灭亡，并使造成破坏者卷入其毁灭之中。想起宗徒所说的话：\BibleQuote{如果我们省察自己，我们就不至于受审判。} \BibleRef{格前 11:31} 我将省察自己，好使那将来审判生者死者的能宽容我。然而，如果您能成功地使那个地区的\OriginalTerm{基督的肢体}{membra Christi}脱离致命的恐惧和忧伤，并以仁慈调和正义的治理来安慰我的晚年，那位在此苦难中通过您赐予我们解脱、并立您于您所在之位者，必将在此生及来生\OriginalTerm{以善报善}{bonum pro bono}，赏报于您。\par

\SourceRecord{Translated by J.G. Cunningham.；Nicene and Post-Nicene Fathers, First Series；Vol. 1.；Edited by Philip Schaff.；Buffalo, NY: Christian Literature Publishing Co.,；1887.；<http://www.newadvent.org/fathers/1102209.htm>.}

% structure: p0001 source=h1 target=section reason=page-title

\section{书信210（公元423年）}

\Emph{致最亲爱的至圣母亲\Person{费利齐塔斯}、\Person{鲁斯蒂库斯}弟兄，以及与她们在一起的姊妹们，\Person{奥古斯丁}及与他在一起的人，在主内致候。}\par

1. 主是美善的，祂的仁慈遍及各处，这仁慈藉着你们在祂内对我们的爱，安慰了我们。祂多么爱那些对祂怀有\OriginalTerm{信德}{faith}和\OriginalTerm{望德}{hope}，且既爱祂又彼此相爱的人，并为他们将来存留了什么福分，祂以最显著的方式证明：对于不信、堕落、乖戾的人，祂虽以永火警告他们，如果他们坚持恶念到底，祂却在今生赐给他们如此多的恩惠，使\BibleQuote{祂使太阳上升，光照恶人，也光照善人；降雨给义人，也给不义的人}\BibleRef{玛 5:45}，这些话简洁地说出一些例子，好使人联想到更多；因为谁能数清恶人在今生从他们轻视的天主那里得到多少恩惠呢？其中有一个极大的恩惠，就是通过时而临到的痛苦——祂如同良医将之与今生的欢乐调剂——告诫他们，只要他们愿意留意，就要逃避将来的义怒，趁着还在路上，即还在今生，与主的话和解，因为他们邪恶的生活已使自己与主的话为敌。那么，如果天主所赐的痛苦本身也是一种祝福，那么祂岂不是更仁慈地赐予人们什么呢？因为顺利是天主安慰时的恩赐，逆境是天主警告时的恩赐；既然祂如我所说，甚至将这些赐给恶人，那么祂为那些彼此担待的人预备了什么呢？你们因祂的\OriginalTerm{恩宠}{grace}而被聚集在这数目中，为此欢喜吧，\BibleQuote{以爱德彼此担待，尽力以和平的联繫，保持圣神的合一}\BibleRef{弗 4:2}。因为不会缺乏你们彼此担待的机会，直到天主如此净化了你们，以致死亡\BibleQuote{被胜利吞灭了}\BibleRef{格前 15:24}，\BibleQuote{天主成为万物之中的万有}\BibleRef{格前 15:28}。\par

2. 我们绝不应喜欢争吵；但无论我们多么不愿意，争吵有时或出于爱，或考验爱。因为很难找到一个愿意受责备的人；经上说的那个智慧人在哪里，\BibleQuote{责备明智的人，他必会爱你}\BibleRef{箴 9:8}？但我们因此就不责备规劝弟兄，以免他因虚假的安全而走向死亡吗？因为这是常见而时常的经历：当一位弟兄受责备时，他开始时感到难堪，抗拒反驳朋友，但后来在独处中与主一起反省，在那里他不怕因顺服纠正而得罪人，却怕因拒绝改正而得罪主，从此便不再做他受正当责备的事；并且，他越是恨恶自己的罪，就越是爱那位他觉得是自己罪的仇敌的弟兄。但他若属于经上所说\BibleQuote{不要责斥傲慢的人，免得他恨你}\BibleRef{箴 9:8}的那类人，争吵并非出于被责者的爱，而是操练并考验责者的爱；因为他不以恨还恨，但那迫使他规劝的爱却坚定不移，即使受责的人以恨相报。但如果责者以恶报恶对待那因受责备而见怪的人，他就不配责备别人，而显然自己应受责备。你们要按这些原则行事，好使争吵不发生，或者若发生，也能迅速以和平终结。要更努力于和谐相处，胜过在争论中压倒对方。因为醋若久留则腐蚀器皿，同样，愤怒若怀抱到明天，就会腐蚀心灵。因此，你们要遵行这些事，和平的天主必与你们同在。也要联合祈祷，为我们代祷，好使我们能愉快地实践我们给你们的善劝。\par

\SourceRecord{Translated by J.G. Cunningham.；Nicene and Post-Nicene Fathers, First Series；Vol. 1.；Edited by Philip Schaff.；Buffalo, NY: Christian Literature Publishing Co.,；1887.；<http://www.newadvent.org/fathers/1102210.htm>.}

% structure: p0001 source=h1 target=section reason=page-title

\section{书信211（公元423年）}

\Emph{在这封信中，\Person{奥古斯丁}谴责了他姐姐曾担任院长的那个\OriginalTerm{隐修院}{monastery}的修女们，因为她们对继任者表现出某些骚动不满，并为她们的指引制定了一般规则。}\par

1. 严厉随时准备惩罚它可能发现的过错，而\OriginalTerm{爱德}{charity}却不愿发现它必须惩罚的过错。这就是我没有同意你们要求我前来访问的原因，因为那时如果我来，我必定不是来为你们的和睦而欢欣，而是来加剧你们的纷争。因为，如果在我不在场时，就发生了如此大的骚乱，以至于你们喧嚷的声音冲击我的耳朵，尽管我的眼睛没有目睹你们的混乱，那么在我面前，我怎能对你们的行为无动于衷，或允许它不受惩罚呢？因为在我面前，你们反抗权威的行为可能会更加激烈，因为我必定拒绝你们所要求的那些让步——这些让步会给你们自己带来不利，开一些非常危险的先例，破坏健全的纪律；这样，我就必定看到你们不是我所希望的样子，而你们也必定看到我不是你们所希望的样子。\par

2. 宗徒致书格林多人说：\BibleQuote{我呼求天主给我的灵魂作证：我没有再到格林多，是为了顾惜你们。我们并不是管制你们的\OriginalTerm{信德}{faith}，而是作你们喜乐的合作者。}\BibleRef{格后 1:23} 我也同样对你们说：为了顾惜你们，我没有到你们那里去。我也顾惜了自己，免得我愁上加愁，并且决定不再见到你们的面，而是为你们在天主面前倾吐我的心，并为你们可怕的危险处境在天主面前辩护，不是用言语对你们说，而是用眼泪向天主祈求；恳求祂不要将使我对你们常感到的喜乐变为悲伤，而在这个世界上处处充满严重罪恶之际，我有时可以因想到你们的数目、你们的纯洁爱德、你们的圣善品行，以及天主赐给你们的丰富\OriginalTerm{恩宠}{grace}而得到安慰，以致你们不仅放弃了婚姻，而且选择同心合意地居住在同一屋檐下，共同生活，好使你们在天主内成为一个灵魂，一个心灵。\par

3. 当我默想你们这些美好的品质，这些天主在你们内的恩赐时，我的心灵虽然在别处的邪恶所搅动的重重风暴中，常常能获得完全的安宁。\BibleQuote{你们跑得好，有谁拦阻了你们，使你们不再服从真理呢？这种说服，决不是出于那召叫你们的。}\BibleRef{迦 5:7} \BibleQuote{一点酵母}——\BibleRef{格前 5:6} 我不愿把这句话说完，因为我更乐意希望、恳求和劝勉这酵母本身转变为更好的东西，免得它败坏整个面团，就像它几乎已经造成的情况那样。因此，如果你们已经重新开始对你们的责任表现出明智的辨别力的幼芽，那么请祈祷不要陷入诱惑，也不要再陷入纷争、嫉妒、怨恨、分裂、恶言、扰乱和私语。因为我们不是为了从你们那里收获这些荆棘，才在你们中间栽种并浇灌上主的园子的。然而，如果你们的软弱仍然被骚动所困扰，那么请祈祷能摆脱这个诱惑。至于那些扰乱你们平安的人，如果你们中间还有这样的人，他们若不改正行为，无论是谁，都将承担他们的审判。\par

4. 想想这是多么邪恶的事：就在我们为多纳徒派回归我们教会合一而欢欣的时候，我们却不得不在我们的隐修院内部哀叹不和。你们要坚守你们所发的善愿，这样你们就不会想换掉那位多年来不辞辛劳地照料隐修院的女院长，在她领导下，你们的人数增加了，年岁增长了，她也将你们在心中放在母亲对亲生子女的位置上。你们所有人来到隐修院时，她就在那里，要么忠实地履行已故圣善女院长——我姐姐的助手职责，要么在她自己继任该职位后，欢迎你们加入修女团体。在她手下，你们度过了初学阶段，在她手下，你们领受了面纱，在她手下，你们的人数增多了，而你们却吵吵闹闹地要求用另一个人替换她，而如果这个别人取代她的提议是由我们提出的，你们理应为此哀悼。因为她是你们所熟知的人；你们最初就是到她这里来的，并且在她手下许多年来年龄增长、人数增多。除了院长之外，没有任命你们以前不认识的其他职员；如果你们是因为他的缘故而寻求更换，如果你们因厌恶他而这样反抗你们的母亲，为什么不请求撤换他呢？然而，如果你们避而不谈这个建议，因为我知道你们在基督内多么尊敬和爱戴他，那么你们为什么不更多地为了他的缘故而尊敬和爱戴她呢？因为这位最近任命的院长在管理你们时，一开始的举措就受到你们混乱行为的阻碍，以至于他自己也打算离开你们，免得因为你们的缘故而蒙受耻辱和憎恶，这种耻辱和憎恶必然源于传出这样的消息：除非你们开始有他作你们的院长，否则你们不会寻求另一位女院长。因此，愿天主平静和安定你们的心：不要让魔鬼的作为在你们内得胜，但愿基督的和平在你们心中赢得胜利；不要因为所渴望的被拒绝而心神恼怒，或因渴求了不该渴求的事而羞愧，就冒然冲向死亡，反要藉着悔改，重新开始良心的责任；不要效法叛徒\Person{犹达斯}的悔改，而要效法牧者\Person{伯多禄}的眼泪。\par

5. 我们为你们这些定居于隐修院的人所制定的应遵守的规则如下：首先，为达成你们聚集在一个团体内的目的，你们要以心神合一住在院内，使你们的心思念虑在天主内合而为一。也不要把任何东西称为私人的财产，却要让一切都成为公有，并由\OriginalTerm{女隐修院长}{prioress}给你们每个人分配食物和衣物——不是一律均等，因为你们并非人人都同样强健，而是按照每人的需要分配。因为你们在\BibleBook{宗徒}{大事录}中读到：\BibleQuote{他们凡物公用，并照每人所需要的分配。}凡进入隐修院时拥有世俗财物的人，应欣然愿意将这些财物归于公用。凡原本没有世俗财物的人，在隐修院内不可要求那些她们在院墙外时无法获得的奢侈品；然而，若有人因体弱而需用舒适之物，即便她们在进入隐修院前极度贫穷，甚至无法为自己获得生活必需品，也应供给她们；但在这种情况下，她们要小心，不要将她们现今在隐修院内找到了衣食视为最大的幸福，而这些衣食是她们在别处无法获得的。\par

6. 再者，你们不要因为与那些在世俗中你们不敢接近的人处于平等地位而昂首自傲；却要举心向上，不寻求世上的财富，以免在我们的隐修院中，富者谦卑而贫者因虚荣自大，致使这些修院只对富者有益，对贫者有害。然而，另一方面，那些看似在世俗中有些地位的人，不可轻视她们的姊妹，这些姊妹为加入这神圣的团体而舍弃了贫穷的境况；她们要更注意因与贫苦的姊妹共处而自豪，不要因自己富裕的出身而自夸。也不要因为或许曾从自己的财产中拿出一些来维持团体，就自视高于他人，免得她们因财富在隐修院内与别人共享，反而比当初在世俗中自己享用财富时，更加容易骄傲。因为其他各种罪都是在恶行中施展，借此得以实行；而骄傲却潜伏在善工中，借此使善工败坏。如果一个人向穷人挥霍钱财而使自己贫穷，却因轻视财富而使不幸的灵魂变得比拥有财富时更为骄傲，这又有什么益处呢？所以，你们大家要同心合意，和睦相处，彼此尊敬，使那使你们成为其殿宇的天主受光荣。\par

7. 在规定的时刻和时间，要\Emph{\OriginalTerm{勤行}{instate}}祈祷。在\OriginalTerm{经堂}{oratory}内，除了该场所原定且因此得名的本分之外，任何人不得做其他事；这样，若你们中有人空闲时想在规定时间以外祈祷，就不致因他人将经堂用作他途而受打扰。在你们向天主祈祷时所用的圣咏和赞美诗中，要以心中默思口中宣唱的；只可咏唱规定要唱的；凡未如此规定的，则不可咏唱。\par

8. 依健康所容许的程度，以守斋和节制饮食来克制肉身。若有人不能守斋，除非生病，否则只可在惯常的用膳时刻进食。从你们来到餐桌直至离开，要在安静无争辩中聆听给你们诵读的内容；不要只动口进食，也要用耳领受天主的话。\par

9. 那些因早年教养而体质较弱者，如果在食物方面受到一些不同的待遇，这对那些因不同的教养而体魄更健壮者，不应感到苦恼或不公平。而且，她们不要认为这些较弱者比自己更受偏爱，因为她们所得到的膳食不如自己的清苦；反而要庆幸自己享有他人所没有的强健体魄。如果对那些在家中受过更娇惯教养后进入隐修院的人，给予一些食物、衣服、床铺或遮盖物——这些是对其他更为强健，因而所处条件更优者所不给的——那么，那些未获得这些宽待的姊妹应想到，这些人从世俗的生活方式下降到她们现在的生活，是多么大的退步，尽管她们或许还未能完全降下到那些体质更强健者所持的艰辛简约的水平。当那些原本较为富有的人，看到别人——并非作为更高荣誉的表示，而是出于对体弱的体恤——所获得的比她们自己多时，她们不应因恐惧隐修纪律可能出现如此可憎的颠倒——即在一座隐修院里，富者在能忍受的范围内受训于刻苦，而贫者受训于奢侈——而感到不安。因为，理所当然的，正如患病者必须少吃食物，否则会加重病情，同样，在病后，正处于恢复期的人，必须为了更快康复而得到如此照顾——即使她们是来自最贫穷的状态进入隐修院的——就好像她们近来的疾病赋予了她们特殊待遇的同等权利，如同她们以前的富裕生活在进入隐修院之前所赋予富者一样。然而，一旦她们恢复了惯常的健康，就应回到她们自己那种更幸福的生活方式，这种生活方式因需求更少而更适合于天主的仆人；不要让私欲在她们健壮时，将她们滞留在那种因软弱而由需要所提升的舒适水平上。那些被赋予更大力量以忍受艰辛的人，应视自己为真正更富有。因为需求少比拥有多更好。\par

10. 你们的衣着切勿引人注目；要渴望以品行而非服饰取悦他人。你们头巾不可薄得让下面的网子显露。外出隐修院时，头发要完全遮盖，既不可蓬乱不修，也不可过于刻意梳理。无论去何处，要结伴同行；到达目的地后，要站在一起。行走、站立、举止和一切动作，都不可有任何可能勾起他人不当欲望之处，而应一切符合你们的圣洁身份。即使目光偶然瞥向男子，也不可凝视任何人；因为在行走时并非不许看见男子，但你们不可对之动情，也不可希望成为他们的情欲对象。因为女人在男子心中唤起或助长这等欲望的，不仅是接触，内在的情感和眼神也能做到。不要说自己心怀贞洁而目光放荡，因为放荡的眼就是放荡的心的表记。当放荡的心彼此以眼神交流，舌头虽缄默，却因感官情欲之力而互相激发欲火，以此为乐时，即使身体未受任何污秽行为沾染，其品格的纯洁也已丧失。那定睛注视异性、并因对方也注视自己而欣喜的女子，不要以为自己的行为未被旁人察觉；她绝对已被她以为不会注意她的人看在眼里。但即使她瞒过众目，无人眼目得见，在天上的那位见证者面前，她又能如何呢？万物在他面前都无法隐藏。难道因他以相应的智慧恒忍鉴察万物，我们就当他是看不见吗？愿每位圣洁的女子，因怀着害怕得罪天主的心，而谨慎提防以罪恶的方式取悦于人；愿她因记念天主的眼目鉴察万物，而克制以罪恶的方式注视男子的欲念。因为正是在此事上，圣经的话劝勉我们敬畏天主：——\BibleQuote{定睛注视的人，为主所憎恶。}所以，当你们聚集在圣堂，或任何有男子在场的地方时，要彼此看管守护你们的贞洁，如此，住在你们内的天主便会借着你们自己来护佑你们。\par

11. 若你们中有人察觉某位姐妹眼神如此轻佻，当立即警诫她，使已萌的恶念不再蔓延，而得以及时遏止。然而，若经劝诫后，你见她重犯此过，或日后又行同样之事，那么无论是谁有机会看见，都必须将她当作一个受伤而需医治的人来报告，但不可只告诉另一人，也许还要告知第三位姐妹，使她凭两三个人的见证而被定罪，\BibleRef{玛 18:16}，并受应有的严责。不要以为这样彼此告发便是心怀恶意。其实真相是，若你们保持缄默，任由姐妹沉沦，而本可通过报告其过失去纠正她们，你们便不是无罪的。比如你的姐妹身上有处伤口，她因害怕外科医生的刀而想隐瞒，你保持沉默，岂非残忍？你若告知他人，岂非真怜悯？那么，你们更当如何将她的罪过告知他人，以免她因忽视灵性的创伤而遭受更致命的痛苦呢？但在指出其他作为见证的人、以便她若否认时可被定罪之前，若经劝诫后她仍拒绝改正，那么应先将她带到\OriginalTerm{女院长}{prioress}面前，使她能这样私下受责，而她的过错不致为众人所知。然而，若她此时否认指控，则须安排别人在她否认后观察其行为，这样，如今在全姐妹团体面前，她就不是被一个人控告，而是被两三个人定罪了。一旦罪证确凿，她就必须服从女院长或\OriginalTerm{院长}{prior}所定的纠正纪律。若她拒绝服从，又不主动离开你们，就当将她从你们的团体中驱逐出去。因为这样做并非残酷，而是出于慈悲，为保护众人不受那染有瘟疫者的传染而灭亡。再者，我方才关于戒绝放荡眼神所说的话，在对待其他一切过犯时，也应怀着对人的合宜之爱和对罪的憎恶，谨慎遵行，无论是观察、禁止、报告、证罪还是惩处，都当如此。但若你们中有人陷入如此大罪，竟暗中接受任何男子的信件或任何形式的礼物，若她主动坦白，就当饶恕她并为之祈祷。然而，若她被查出或证实有此行为，则应按女院长、院长或甚至主教的判决，予以更严厉的惩罚。\par

12. 你们要把衣服存放在一个地方，由一两个人照管，或按需要由几个人负责抖开衣服，以免虫蛀；你们的食物既来自同一个储藏室，你们的衣服也应由同一个衣物间供应。无论按季节配发给你们的衣着是什么，倘若可能，就不要计较每人是否领回自己原先搁置的那一件，或者领到别人穿过的，只要任何人都不缺少必需品。然而，如果因此而在你们中间引起争执和抱怨，有人埋怨说自己领到的衣服不如从前穿过的，认为这样穿着比不上另一位姊妹，觉得有失身分，那么，你们就当借此省察自己，并判断你们若为身体的衣服互相争吵，你们内心那圣洁的衣饰该是多么缺乏。不过，倘若因你们的软弱而容许你们领回原先搁置的那件衣服，那么，你们所搁置的衣服仍应集中存放，并由衣物间的固定管理员负责保管；当然，要明白谁也不得为自己私自享用而缝制任何衣服、床榻用品、腰带、头巾或头饰，你们的一切工作都应为众人的公益而做，要比各自为私人利益而做时更热心、更愉快地坚持。因为经上论这\OriginalTerm{爱德}{charity}说：\BibleQuote{不求己益。}\BibleRef{格前 13:5}，这应理解为爱德是谋求公益而非私益，不是私益先于公益。所以，你们越是把公益置于私人和个别的利益之上，就越能体认到进步，在那些供应转眼即逝的需求的事物上，那永存的爱德将会占有显赫而有力的地位。由这些规定必然得出的结论是：当男女信友给他们在隐修院里的女儿，或任何关系的住院者，带来衣物或其他视为必需品的礼物时，这些礼物不得私下收受，而必须交由院长管理，加进公用物品，以供任何有需要的住院者使用。若有人隐藏别人给她的任何礼物，应定她如同犯了偷窃的罪。\par

13. 你们的衣服，无论由自己或由洗衣妇来洗，间隔多久应由院长批准，以免过度追求衣服的洁白无瑕而使心灵染上污点。沐浴洁身及使用澡堂不可频繁，应按惯常的间隔，\Emph{即}每月一次。但若是因病而需沐浴洁身，不可过分耽延；应遵照医生的建议去行，不可抱怨；即使病人不情愿，也须服从院长的命令，做健康所要求的事。然而，病人若想洗澡，而那对她并非有益，就不可随从她的愿望，因为有时因带来快感而被认为有益，实际上却可能有害。最后，若天主的婢女身体有隐痛，她对病痛的陈述应毫不犹豫地相信；但如果她所喜欢的方法是否真能治愈疼痛尚有疑问，就该请教医生。去澡堂或任何必须去的地方，每次不得少于三人同行。而且，需要外出的姊妹不可自己挑选同伴，而应由院长指定。对病人、需要照料的康复者，以及虽无发热症状但身体虚弱的人，应指定你们中的一人负责照料，她应从储藏室为各人领取她认为必需的物品。此外，管理储藏室、衣物间或图书馆的人，为姊妹们服务时不可抱怨。手抄本应在每天固定的时刻申请，在别的时刻索取的人不得领取。但若有人急需衣服、鞋履，无论何时，负责部门的人都不可迟延供应。\par

14. 你们中间不应有争吵；至少，一旦发生，就当尽快平息，免得忿怒化为仇恨，把\BibleQuote{木屑变成大梁，}\BibleRef{玛 7:3} 并使灵魂负上谋杀的罪责。因为圣经的话：\BibleQuote{凡恼恨自己弟兄的，便是杀人的，}\BibleRef{若一 3:15} 并不只针对男人，女人也受这条法律的约束，因这法律原是颁给在创造次序上居先的另一性别。无论谁，若以辱骂、恶言或诬告伤害了他人，务必尽快道歉补救；受伤害的一方则当宽恕，不再争论。若伤害是相互的，由于你们祈祷频繁，而这些祈祷在你们眼中本该更神圣，双方就更应彼此宽恕。但那爽快地请求她承认被她伤害的人宽恕的姐妹，虽然可能因性急而多次犯错，却比那虽不易发怒、却很难被说服去求恕的人更好。拒绝宽恕姐妹的，不要指望自己的祈祷蒙应允；至于那从不求恕、或并非出于真心求恕的姐妹，即便留在隐修院内，对她毫无益处，即使她或许不会被逐出。因此，要避免刻薄的话；但若它们已脱口而出，就不要迟延，要用那造成创伤的同一张嘴，说出医治的话语。然而，当纪律的需要迫使你们以严词管教年幼者时，即便你们觉得自己已说得太过，也无需向她们求恕，以免你们对那本该服从的人过度谦卑，损及管理她们所必需的权威；但仍须向万有的主求得宽恕，祂知道你们是怀着多少善意去爱，甚至爱那些你们或许责之过严的人。你们彼此间的爱不应是肉性的，而应是属神的：因为那些不端庄的女子在嬉戏狎昵中所行的事，即便已出嫁或准备出嫁的女子尚且不该做，更何况是寡妇，或因圣愿而献身作基督婢女的贞女，则更不该行这些事了。\par

15. 要如同顺从母亲一样顺从院长，给予她一切应有的尊重，免得你们因忘记对她当尽的本分而得罪天主；你们更要顺从那位负责照料你们全体的司铎。院长负有特别的责任，要确保这一切规矩都得遵守，若有任何规矩被人忽略，不可姑息，须仔细纠正并处罚；当然，若遇超过她权限或能力的事，她可转呈司铎处理。但她要称自己为有福，不是凭借行使统治的权力，而是在于实践服事的爱。在人前的尊荣上，她须被高举于你们之上；但在天主前的敬畏中，她要如同在你们脚下。她要在众人前显出\BibleQuote{善行的模范。}\BibleRef{铎 2:7} \BibleQuote{劝戒不守规矩的，安慰怯懦的，扶持软弱的，对众人忍耐。}\BibleRef{得前 5:14} 她要欣然遵守规矩，谨慎推行规矩。尽管两者都必要，她却要更切愿为你们所爱，而非为人所惧；要时常想到她须为你们向天主交帐。为此，你们顺从她，不但要出于对自身的怜悯，也要出于对她的怜悯，因为她在你们中间地位愈高，她的危险也相应比你们更大。\par

16. 愿主赐予你们，能以爱慕属神的华美，借着良好的品行散发基督的芬芳，不是像法律下婢女那样，而是像在恩宠下享有自由的人那样，甘心服从这一切规矩。为使你们能借这篇文字如同镜子般省察自己，不至因遗忘而疏忽任何事，当每周为你们诵读一遍；你们若发现自己在实行所记载的事，就要为此感谢诸恩之源的天主；若有人察觉自己在某些方面有所欠缺，就要为过去哀伤，并在将来警醒，祈求既蒙宽赦罪债，又不被引入诱惑。\par

\SourceRecord{Translated by J.G. Cunningham.；Nicene and Post-Nicene Fathers, First Series；Vol. 1.；Edited by Philip Schaff.；Buffalo, NY: Christian Literature Publishing Co.,；1887.；<http://www.newadvent.org/fathers/1102211.htm>.}

% structure: p0001 source=h1 target=section reason=page-title

\section{第212封信（公元423年）}

\Emph{致\Person{昆体良努斯}，我最为有福的主、兄弟及同为主教，当受尊敬者，\Person{奥思定}谨致问候。}\par

可敬的父亲，我在基督的爱内将这几位可敬的天主仆婢与基督的宝贵肢体托付于您：寡妇\Person{加拉}（她已誓发圣愿）和她的女儿\Person{辛普利西亚}，一位献身于主的贞女，她因年龄而服从母亲，但因其圣德而高于母亲。我们已尽所能以天主的圣言滋养了她们；并藉此信函，如同亲手将她们交托给您，请您在她们所需或急难中，以各种方式安慰并扶助她们。您圣善的主教即使没有我的推荐，也无疑会承担此责任；因为倘若由于那在天上的\OriginalTerm{耶路撒冷}{Jerusalem above}（我们都是其公民，而她们切愿在其中获得卓越圣善的位份），我们当对她们怀有不仅公民间应有的情谊，甚至弟兄般的友爱，那么她们在您身上的要求岂不更为迫切？您居住在这尘世的同一国度，而这两位女士为了基督的爱已弃绝了此世的荣华。我也请您不吝以我献上此函时的同等爱德，收受我作为主教身份的致意，并在您的祈祷中记念我。这几位女士随身携有最可颂及光荣的殉道者\Person{斯德望}的\OriginalTerm{圣髑}{relics}：您圣善的主教知道如何给予它们应有的敬礼，一如我们所行。\par

\SourceRecord{Translated by J.G. Cunningham.；Nicene and Post-Nicene Fathers, First Series；Vol. 1.；Edited by Philip Schaff.；Buffalo, NY: Christian Literature Publishing Co.,；1887.；<http://www.newadvent.org/fathers/1102212.htm>.}

% structure: p0001 source=h1 target=section reason=page-title

\section{第213封信（公元426年9月26日）}

\Emph{由圣\Person{奥斯定}编写的记录，记载了他指定\Person{埃拉克利乌斯}接替其主教职位的经过，并使\Person{埃拉克利乌斯}在其年老时代替他分担部分职责。}\par

\Emph{在\Place{希波}地区内的和平教会，于最著名的\Person{狄奥多西}第十二任执政官、\Person{瓦伦提尼安·奥古斯都}第二任执政官之年的9月26日——\Person{奥斯定}主教与他的同伴主教\Person{瑞利吉亚努斯}和\Person{马尔提尼亚努斯}一同就座，在场的\OriginalTerm{司铎}{presbyters}有\Person{萨图尔尼努斯}、\Person{勒波里乌斯}、\Person{巴尔纳伯}、\Person{福尔图纳提亚努斯}、\Person{鲁斯提库斯}、\Person{拉匝禄}和\Person{埃拉克利乌斯}，同时有圣职人员及众多平信徒站立一旁——\Person{奥斯定}主教说：——}\par

\Quote{亲爱的诸位，我昨天向你们提出的那件事——为此我请你们来参加，正如我所见，你们比往常来得更多——必须立即处理。因为当你们的心情焦虑地专注于此事时，若我讲论其他主题，你们几乎不会听我。我们所有人都有一死，而世上生命最后的一天，对每个人而言，在任何时候都是不确定的；但在婴孩期，有进入少年期的希望，如此我们的盼望继续，从少年期展望青年期，从青年期展望成年期，从成年期展望老年：这些盼望能否实现是不确定的，但在每一阶段都有某些可盼望的事。然而老年却没有任何生命阶段可展望：老年人能延长多久是不确定的，但确定的是，没有任何别的阶段注定在老年之后来临。我来到这个城镇——因为这是天主的旨意——那时我正值壮年。那时我年轻，但现在我已年老。我知道教会习惯在主教逝世后，因有野心或好争论的派系而受扰乱；我觉得我有责任采取措施，避免此团体因我的离世而遭受那种我在别处经常目睹且为之哀叹的情况。亲爱的诸位，你们都知道，我最近拜访了\Place{米勒维}教会；因为那里的某些弟兄，特别是天主的仆人们，请我去，因为担心在我的弟兄兼同工、已故的主教\Person{塞维鲁斯}去世后会有某些骚动。因此我去了，主仁慈地乐意帮助我们，使他们和谐地接纳了那位由他们前主教生前指定的人为主教；因为当这一指定已为他们所知时，他们欣然默许了他所作的选择。然而，有个疏漏发生，令某些人不悦；因为我们的弟兄\Person{塞维鲁斯}相信，只要将继任者的名字告知圣职人员就够了，没有对民众提及此事，这使得有些人心中不悦。但我何必多说呢？按天主的喜悦，不满消除了，喜乐充满众人的心，而\Person{塞维鲁斯}所提名的那人被祝圣为主教。为避免本教会在此事上发生任何此类的怨言，我现在向你们在座的各位宣告我的心愿，我相信这也是天主的旨意：我希望司铎\Person{埃拉克利乌斯}做我的继任者。}\par

民众喊道，\Quote{感谢天主！赞美基督！}（重复了二十三次）。\Quote{基督，求你垂听我们；愿\Person{奥斯定}长寿！}（重复了十六次）。\Quote{我们要你作我们的父亲，作我们的主教}（重复了八次）。\par

\Emph{2. 在获得安静后，\Person{奥斯定}主教说：——}\par

我不必说什么赞扬\Person{埃拉克利乌斯}的话；我赞赏他的智慧，也顾全他的谦逊；你们认识他就够了。而我所宣布的、我所愿望于他的，我知道也是你们所愿望的；如果我以前不知道，今天我便得到了证据。因此，我渴望此事；我以祈祷，祈求主我们的天主，这些祈祷的热忱不因年迈的寒冷而减退；我也劝勉、告诫并请求你们与我一同祈祷——祈求天主能坚定那已在我们内成就的事，藉着基督的平安，将我们众人的心灵融合在一起。愿派遣他到我这里来的那位保护他！保护他安全，保护他无可指摘，使他在我活着时给我喜乐，在我离世时接替我的位置。\par

\Quote{教会的书记员，正如你们所见，正在记录我所说的和你们所说的；我的讲话和你们的欢呼都不会被忽略。更明白地说，我们正在为今天的事项制定一份教会记录；因为我希望这些事情能以这种方式，在涉及人的方面得到确认。}\par

民众呼喊了三十六次，\Quote{感谢天主！赞美基督！}基督，求你垂听我们；愿\Person{奥斯定}长寿！说了十三次。\Quote{你，我们的父亲！你，我们的主教！}说了八次。\Quote{他相称且正义，}说了二十次。\Quote{堪当此任，十分相称！}说了五次。\Quote{他相称且正义！}说了六次。\par

\Emph{3. 在获得安静后，\Person{奥斯定}主教说：——}\par

\Quote{正如我刚才说的，我希望我的愿望和你们的愿望，在关乎人的方面，通过列入教会记录而得到确认；但在关乎全能者的旨意方面，让我们都祈祷，正如我先前所说，愿天主坚定那在我们内所成就的事。}\par

民众喊叫，说了十六次\Quote{我们感谢你的决定：}然后十二次\Quote{同意！同意！}然后六次\Quote{你，我们的父亲！\Person{埃拉克利乌斯}，我们的主教！}\par

\Emph{4. 在获得安静后，\Person{奥斯定}主教说：——}\par

\Quote{我赞同你们也表示赞同的事；但我不希望在他身上重演曾发生在我身上的那件事。你们许多人都知道所发生的事；事实上，你们所有人都知道，除非那些当时尚未出生，或未达懂事年龄的人。当我那位安息主怀的年迈父亲主教 \Person{Valerius} 尚在世时，我被祝圣为主教，并与他共同占据主教之位，而我当时并不知道这为 \Place{Nicaea}公会议所禁止；他同样对此禁令一无所知。我不愿让我在此的儿子受到与我当年相同的责难。}\par

人群呼喊，说了十三遍：\Quote{感谢天主！赞美基督！}\par

5. \Emph{待众人安静后，奥古斯丁主教说：——}\par

\Quote{他将暂且保持现状，仍是一位 \OriginalTerm{长老}{presbyter}；但之后，当天主乐意之时，成为你们的主教。但现在，我确实将开始借助基督的怜悯，着手做我迄今未做的事。你们知道，如果你们曾允许的话，我多年来本会每周用五天时间，不让人打扰，以便履行我的弟兄和父亲们，即共教的主教们，在 \Place{Numidia} 和 \Place{Carthage} 的两次公会议上乐于托付给我的 \WorkTitle{圣经} 研习之责。该协议已正式记录，你们表示了同意，并通过欢呼予以确认。你们同意和欢呼的记录，已当众向你们宣读。有短暂时间，你们遵循了协议；之后，却毫无顾忌地违反了它，使我无法获得空闲，从事我渴望的工作：上午和下午，我都深陷于他人需要我关注的事务之中。我现在恳求你们，并因基督的缘故郑重地要求你们，容我担起将这部分劳苦的重担转给这位年轻人，我指的是 \Person{Eraclius} 长老，我今天以基督之名指定他为我主教职位的继任者。}\par

人群呼喊，说了二十六遍：\Quote{我们感谢你的决定。}\par

6. \Emph{待众人安静后，奥古斯丁主教说：——}\par

我在主我们的天主面前为你们的爱和善意感谢；是的，我为这些感谢天主。因此，弟兄们，从今以后，所有你们惯常带到我这里来的事，都带到他那里去。在任何他认为需要我的建议的情况下，我绝不推辞；我绝不会收回这一点：不过，所有惯常带到我这里的事，都带到他那里去。倘若 \Person{Eraclius} 本人，在任何情况下，或许对应该做的事感到困惑，他可以咨询我，或求助于我这位被他视为父亲的人。通过这一安排，你们一方面不会遭受任何损失，而我，终于可以在天主延续我生命的短暂时期内，将我的余生，无论多寡，不是用于闲散或贪图安逸，而是在 \Person{Eraclius} 善意允许的范围内，致力于研习 \WorkTitle{圣经}：这也有益于他，并通过他同样有益于你们。因此，愿没有人不愿给我这一空闲，因为我提出这个要求只是为了做重要的工作。\par

\Quote{我看，为召集你们的目的，所有必需的事务我已与你们处理完毕。最后一件事是，请你们中尽可能多的人，乐意为这份记录署名。此刻我需要你们的回应。请让我听到你们的回应：通过欢呼表示你们的赞同。}\par

人群呼喊，说了二十五遍：\Quote{同意！同意！} 然后二十八遍：\Quote{当之无愧，理所当然！} 然后十四遍：\Quote{同意！同意！} 然后二十五遍：\Quote{他久已堪当，他久已当之无愧！} 然后十三遍：\Quote{我们感谢你的决定！} 然后十八遍：\Quote{基督，请俯听我们；护佑 \Person{Eraclius}！}\par

7. \Emph{待众人安静后，奥古斯丁主教说：——}\par

\Quote{我们能够围绕祂的祭献来处理那些属于天主的事，这是美好的；在这个我们指定的祈祷时刻，我特别劝勉你们，诸位亲爱的，放下你们一切的事务和营生，为这个教会，为我，并为 \OriginalTerm{长老}{presbyter} \Person{Eraclius}，在主面前倾吐你们的祈求。}\par

\SourceRecord{Translated by J.G. Cunningham.；Nicene and Post-Nicene Fathers, First Series；Vol. 1.；Edited by Philip Schaff.；Buffalo, NY: Christian Literature Publishing Co.,；1887.；<http://www.newadvent.org/fathers/1102213.htm>.}

% structure: p0001 source=h1 target=section reason=page-title

\section{书信二一四}

致我最亲爱的主、在基督的肢体中最受尊敬的弟兄\Person{瓦伦提努斯}，及与你同在的弟兄们，\Person{奥思定}在主内致候。\par

1. 有两个年轻人，\Person{克雷斯孔尼乌斯}和\Person{费利克斯}，来到了我们这里，并自称属于你们的团体，告诉我们说，你们的隐修院因某些人宣讲\OriginalTerm{恩宠}{grace}的方式而大受扰乱，因为他们竟否认人有自由意志；更严重的是，他们主张在审判的日子，天主不按各人的行为施行报应。同时，他们向我们指出，你们中许多人不赞成这种看法，反而承认自由意志靠天主的恩宠得到助佑，使我们能思想善事、行为正直；这样，当主来按各人的行为施行报应时，必看见我们的那些善工，就是天主所预备，叫我们行在其中的。持这种想法的人，想得对。\par

2. 因此，弟兄们，正如宗徒劝导格林多人那样，\BibleQuote{我因我们的主耶稣基督之名，求你们众人言谈一致，在你们中不要有分裂}\BibleRef{格前 1:10}。首先，主耶稣，正如宗徒\Person{若望}的\BibleBook{若望}{福音}中所记载的，\BibleQuote{来到世界上，不是为审判世界，而是为叫世界借着祂而获救}\BibleRef{若 3:17}。其次，正如宗徒\Person{保禄}所写的，\BibleQuote{当祂来临时，天主必要审判世界}\BibleRef{格前 4:5}，一如整个教会在信经中所宣认的，\Quote{审判生者死者}。那麽我要问，如果没有天主的恩宠，祂怎能拯救世界？如果没有自由意志，祂怎能审判世界？因此，我希望你们按照这信仰理解我的那本书或那封书信，就是上述弟兄带给你们的；这样，你们既不要否认天主的恩宠，也不要如此主张自由意志，以致将自由意志与天主的恩宠分开，好像没有天主的恩宠，我们便无法按天主的意思思想或行动一样——这完全超出我们的能力。的确，正是因此，主论及义德的果实时说：\BibleQuote{离了我，你们什么也不能做。}\BibleRef{若 15:5}\par

3. 由此你们可以明白，我为何写给\Place{罗马}教会的司铎\Person{西克斯图斯}那封提及的书信，反对新兴的\OriginalTerm{培拉基派}{Pelagian}异端者，他们说，天主的恩宠是按我们自己的功劳赋予的，因此，凡夸耀的，不可在主内夸耀，而要在自己内夸耀——就是说，在人内，而不是在主内。然而，宗徒用以下的话禁止了这一点：\BibleQuote{谁也不要拿人来夸耀}\BibleRef{格前 3:21}；在另一处他又说：\BibleQuote{凡要夸耀的，应当因主而夸耀。}\BibleRef{格前 1:31}。但这些异端者，以为自己是由自己成义的，好像天主并没有把这恩赐赋予他们，而是他们自己凭自己获得的，自然在自己内夸耀，而不是在主内。如今，宗徒对这样的人说：\BibleQuote{谁使你异于别人呢？}\BibleRef{格前 4:7}，他这样说，原是鉴于那由\Person{亚当}而出的\OriginalTerm{丧亡之众}{mass of perdition}中，除了天主将一个人分别出来，使他成为\OriginalTerm{尊贵的器皿}{vessel to honour}，而非\OriginalTerm{卑贱的器皿}{vessel to dishonour}之外，别无其他分别。\par

然而，为使属血肉的人不致因愚妄的骄傲，在听到\BibleQuote{谁使你异于别人呢？}\BibleRef{格前 4:7}这问题时，或在思想上或在言语上回答说：是我的\OriginalTerm{信德}{faith}，或我的祈祷，或我的义德使我异于别人——宗徒立刻在问题之后加上了这些话，以回应所有这类想法，说：\BibleQuote{你有什么不是领受的呢？既然是领受的，为什么你还夸耀，好像不是领受的呢？}\BibleRef{格前 4:7}。如今，凡夸耀的人，正好像他们的恩赐不是由\OriginalTerm{恩宠}{grace}领受的；他们以为自己是凭自己成义的，并因此在自己内夸耀，而不在主内。\par

4. 所以，在这封已送达你们的书信中，我引用了圣经的章节，你们可以自行查考，证明了我们的善工、虔诚的祈祷和正确的信德，若非从祂那里领受了一切，便不可能存在于我们内；关于祂，宗徒\Person{雅各伯}说：\BibleQuote{一切美好的赠与，一切完善的恩赐，都是从上，从光明之父降下来的。}\BibleRef{雅 1:17}。因此，没有人能说天主的恩宠是凭他自己善工的功劳、或凭他祈祷的功劳、或凭他信德的功劳而赐给他的；也不能以为那些异端者所持的道理是真实的，即天主的恩宠是按照我们自己的功劳而赋予我们的。这完全是一种极其错误的意见；其实，这并非因为在虔敬的人中没有善功，在不虔敬的人中没有恶行（否则天主将如何审判世界呢？），而是因为一个人的皈依是由于天主的仁慈和恩宠，正如圣咏作者所说：\BibleQuote{我的天主，愿你的慈爱作我的前导}\BibleRef{咏 59:10}；这样，不义的人得以成义，即从邪恶变为正义，并开始拥有那善功，当天主审判世界时，祂将予以赏报。\par

5. 我本来有许多事情想要寄给你们，你们读了之后，便能更精确而充分地了解主教们在公会议上为反对这些培拉基派异端者所行的一切。但来自你们团体的弟兄们行色匆忙。我们已托他们带给你们这封书信；然而，这并非对任何来函的答复，因为事实上，他们并没有从你们挚爱的本人那里带给我们任何书信。尽管如此，我们仍毫不犹豫地接待了他们；因为他们纯朴的举止，充分向我们证明，他们此次来访绝无虚假或欺骗。不过，他们非常匆忙，因为希望能与你们一起在家里过复活佳节；我恳切祈求，愿这神圣的日子，赖主的助佑，为你们带来和平，而非分裂。\par

6. 你确实会采取更好的做法（我恳切地请求你），如果你不拒绝把那个他们说是使他们受到扰乱的人送到我这里来。因为他要么不明白我的书，要么当他自己试图解决和说明一个极为困难且仅有少数人能明白的问题时，他或许自己也被误解了。因为这正是天主的\OriginalTerm{恩宠}{grace}的问题，它使一些不明智的人以为宗徒\Person{保禄}给我们规定了一条规则：\BibleQuote{让我们作恶以成善吧。}\BibleRef{罗 3:8} 正是针对这些人，宗徒\Person{伯多禄}在他的第二封书信中写道：\BibleQuote{为此，可爱的诸位，你们既等候这些事，就当努力，使祂在平安中看见你们没有玷污，无可指摘。并应以我们主的容忍当作得救的机会；就如我们亲爱的兄弟保禄，按着赐给他的智慧，给你们写的；也正如他在论及这些事的所有书信上所写过的；在这些书信内，有些难懂的地方，不学无术和站立不稳的人，便加以曲解，一如曲解其它圣经一样，而自取灭亡。}\BibleRef{伯后 3:14-16}\par

7. 因此，你要留意这位伟大宗徒的这些可怕的话；当你感到不理解时，同时要信赖天主的默示言语，要相信人的意志是自由的，也有天主的\OriginalTerm{恩宠}{grace}，没有\OriginalTerm{恩宠}{grace}的帮助，人的\OriginalTerm{自由意志}{free will}既不能转向天主，也不能在天主内取得任何进步。你虔诚信仰的事，你要祈祷，为能明智理解。而实际上，正是为了这个目的——即我们能明智理解——才有了\OriginalTerm{自由意志}{free will}。因为除非我们以\OriginalTerm{自由意志}{free will}理解并变得明智，否则圣经中就不会以这样的话吩咐我们：\BibleQuote{愚昧的人群，如今你们要理解；愚顽的人，你们何时才明智？}\BibleRef{咏 94:8} 正是这呼唤我们要明智和理解的诫命与训令，也要求我们服从；而没有\OriginalTerm{自由意志}{free will}，我们就无法履行服从。但是，如果我们能够凭\OriginalTerm{自由意志}{free will}服从这条要我们明智和理解的诫命，而不需要天主\OriginalTerm{恩宠}{grace}的帮助，那么向天主说：\BibleQuote{赐我理解，使我学习你的诫命；}\BibleRef{咏 119:73} 就毫无必要；也不会在福音中记载：\BibleQuote{祂遂开启他们的明悟，使他们理解圣经；}\BibleRef{路 24:45} 宗徒\Person{雅各伯}也不至于对我们说：\BibleQuote{你们中谁若缺乏智慧，应当向那慷慨施恩于众人，且不责备的天主祈求，就必获得。}\BibleRef{雅 1:5} 但主能够赐予你们和我们，使我们可以因很快听到你们平安和虔诚合一的消息而欢欣。我问候你们，不仅以我自己的名义，也以与我同在的弟兄们的名义；并请求你们同心合意、孜孜不倦地为我们祈祷。愿主与你们同在。\par

\SourceRecord{Translated by J.G. Cunningham.；Nicene and Post-Nicene Fathers, First Series；Vol. 1.；Edited by Philip Schaff.；Buffalo, NY: Christian Literature Publishing Co.,；1887.；<http://www.newadvent.org/fathers/1102214.htm>.}

% structure: p0001 source=h1 target=section reason=page-title

\section{书信215}

致我在主内极亲爱且最尊敬的兄弟、基督肢体中的\Person{瓦伦提努斯}，以及与你同在的弟兄们，\Person{奥古斯丁}在主内致以问候。\par

1. \Person{克雷斯科尼乌斯}和\Person{菲利克斯}，以及另一位\Person{菲利克斯}，天主的仆人，从你们的兄弟团体来到我们这里，与我们一同庆祝了复活节，此事你们爱德已得知。我们留他们稍久，为使他们能更好地受训，以对抗新的白拉奇异端者；凡以为天主的\OriginalTerm{恩宠}{grace}是按人的任何功绩赐予我们的，都陷入这种异端的错误，而这恩宠唯独通过我们的主耶稣基督拯救人。但同样错误的是，有人以为，当主来审判时，一个毕生能运用\OriginalTerm{自由意志}{free will}的人不会按他的行为受审判。因为只有婴孩，他们尚未做过任何自己的、或好或坏的行为，若未在重生的圣洗中借着救主的恩宠获得拯救，就将只因\OriginalTerm{原罪}{original sin}而受罚。至于所有其他的人，他们在运用自由意志时，在自己所犯的罪上增加了原罪，但未蒙天主的恩宠从黑暗的权势中被拯救出来而迁入基督的国，他们将按功过受审判，不仅根据原罪，也根据他们自己意志的行为。善人固然将按自己善意的功绩得到赏报，但他们之所以有此善意，原是藉着天主的恩宠获得的；于是圣经上的那句话就应验了：\BibleQuote{忿怒与愤恨、患难和困苦，必加于一切作恶的人，先是犹太人，后是外邦人；光荣、尊贵与平安，加于一切行善的人，先是犹太人，后是外邦人。}\par

2. 关于意志与恩宠这个极为困难的问题，我觉得在本信中无需进一步讨论，因为在他们将更急速返回时，我已给了他们另一封信。我也为你们写了一本书，如果你们靠主的助佑研读并透彻领悟它，我想你们中间就不会再起分歧了。他们还带了其他文件，我们以为应当送给你们，为使你们能从中得知，天主教会靠天主的仁慈采用了什么方法来驱逐白拉奇异端的毒害。因为致\Place{罗马}主教教宗\Person{依诺增爵}的信——来自\Place{迦太基}省会议和\Place{努米底亚}会议的信，以及一份由五位主教极其谨慎撰写的信，还有他对这三封信的回复；我们为非洲会议致教宗\Person{佐西穆斯}的信，和他写给全世界主教们的答复；还有我们在后来全\Place{非洲}全体会议上针对这错误拟定的一份简短决议；以及我上面提及的刚刚为你们写成的那本书——所有这些，在他们与我们在一起时，我们一同向他们宣读过了，现在已托他们带给你们。\par

3. 此外，我们向他们宣读了至福的殉道者圣\Person{西彼廉}的\WorkTitle{论主祷文}著作，并向他们指出，他教导说，一切有关我们品行、构成正直生活的事，都必须向我们在天上的父祈求，免得我们因妄用自由意志而失落天主的恩宠。从同一篇论著中，我们还向他们指出，同一荣耀的殉道者如何教导我们：甚至要为那些尚不信基督的仇敌祈祷，使他们能信；这样祈祷当然全心信赖，除非教会相信，人的邪恶和不信的意志也能藉天主的恩宠转变为善。不过，这本圣\Person{西彼廉}的书我们并没有送给你们，因为他们告诉我们你们手中早已有了。还有我曾寄给\Place{罗马}教会司铎\Person{西斯多}的那封信，他们也带来给我们了，我们一同读过，并指出它是为反对那些说天主的恩宠是按我们的功绩赐予的人而写的——那就是反对同样的白拉奇派。\par

4. 那么，我们已尽我们的能力，像对待你们的兄弟和我们自己的兄弟一样，劝勉了他们，为使他们坚守\OriginalTerm{公教信德}{Catholic faith}的纯正，这信德既不否认自由意志可用于为恶或为善的生活，也不赋予它如此大的能力，以致没有天主的恩宠，它便不能有所成就，无论是使之由恶变善，还是使之坚持行善，或是使之达到永恒的美善——那时再没有跌倒的恐惧了。同样地，我最亲爱的你们，我在这封信中也给你们同一劝告，就是\OriginalTerm{宗徒}{Apostle}向我们众人所发的：\BibleQuote{不要把自己看得太高，超过所当看的；但要怀着合理的谦恭，按照天主分给各人的信德的尺度来估价自己。}\par

5. 请仔细留意圣神借\Person{撒罗满}给我们的劝言：\BibleQuote{你应修平你脚下的行径，坚定你的路途。不可偏左偏右，要使你的脚远离邪恶；因为主认识右边的道路，至于左边的道路则是邪恶的。他必使你的行径平坦，引导你的脚步得享平安。}现在，弟兄们，请想想，圣经的这些话语中，若没有自由意志，便不会说：\BibleQuote{你应修平你脚下的行径，坚定你的路途；不可偏左偏右。}同样，若没有天主的恩宠，我们便无法做到此事，那么随后也不会加上：\BibleQuote{他必使你的行径平坦，引导你的脚步得享平安。}\par

6. 因此，你们既不要偏向右，也不要偏向左，尽管右边的道路受称赞，左边的道路受责备。为此他加上一句：\BibleQuote{使你的脚远离邪恶的道路，}——那就是指左边的道路。这一点他在以下的话中显明出来，他说：\BibleQuote{因为主认识右边的道路；至于左边的道路却是邪恶的。}我们本当行走在主所认识的道路中；在\BibleBook{圣咏}{}中我们读到的正是这些道路：\BibleQuote{主认识义人的道路，不虔敬者的道路必趋灭亡；}因为这条路，就是左边的路，主不认识。正如他在审判之日，也必最终对站在他左边的人说：\BibleQuote{我不认识你们。}那么，那洞悉人内一切善恶者所不认识的是什么？但\BibleQuote{我不认识你们}这句话，岂不意指你们如今成了我从未创造过的样子？正如经上论及主耶稣基督所说：\BibleQuote{他不认识罪。}他怎能不认识呢？除非他从未造过罪。因此，\BibleQuote{主认识右边的道路}这句话，除了指他自己亲手造了那些道路——即\Quote{义人的途径}之外，还能怎样理解呢？而这些\Quote{义人的途径}无疑正是\BibleQuote{\OriginalTerm{善工}{good works}，即天主——正如宗徒告诉我们的——所预先安排，叫我们行于其中的}。至于左边的道路——那些不义者的邪恶途径——他确实一无所知，因为他从未为人造过这些，却是人自己为自己造的。所以他说：\BibleQuote{邪恶者的邪路我深恶痛绝；它们都在左边。}\par

7. 但有人会问：为什么他说：\BibleQuote{不可偏向右，也不可偏向左，} 而他显然更应该说：持守右边的道路，不要转向左边，既然右边的道路是好的？我们想，原因岂不是：右边的道路如此美善，以致偏离它们——即使偏向右——也非善事？因为那选择将本身属于右边道路的善工归于自己、而不归于天主的人，确实须被理解为偏向了右。因此，在说了\BibleQuote{因为主认识右边的道路，至于左边的道路却是邪恶的}之后，仿佛有人向他提出异议：那么，你为何不愿我们转向右呢？他立刻又加上：\BibleQuote{他必亲自修直你的途径，使你的道路得享平安。}因此，你要如此理解这诫命——\BibleQuote{为你自己的脚修直途径，又要为你自己的道路平顺}——就是要知道：无论何时你做了这一切，都是主天主使你能够做到的。这样，即使你行走在右边的道路上，也不致偏向右，因为你不仗恃自己的力量；他自己将成为你的力量，他必亲自修直你的途径，使你的道路得享平安。\par

8. 因此，最亲爱的诸位，凡说\Quote{我的意志足以为我成就\OriginalTerm{善工}{good works}}的人，便是偏向了右。但另一方面，有人一听到天主的\OriginalTerm{恩宠}{grace}被宣讲，而致使人设想并相信：恩宠本身能使人的意志由恶变善，甚至恩宠本身也能保存它已造成的状态，便认为应抛弃善的生活方式；并由于这种见解，进而说：\BibleQuote{让我们作恶以成善吧}——这些人便是偏向了左。这就是为什么他对你们说：\BibleQuote{不可偏向右，也不可偏向左；}换句话说，既不要如此高举\OriginalTerm{自由意志}{free will}，以致将善工归功于它而不归因于天主的恩宠，也不要如此维护并坚持恩宠，以致好像因着恩宠，你就可以安稳无忧地喜爱恶行——愿天主的恩宠亲自使你们远避这事！宗徒所针对的，正是这类人的话，当他说：\BibleQuote{那么，我们可说什么呢？我们要常留在罪恶中，好叫恩宠洋溢吗？}对于这谬误之人所发的怪论——他们根本不懂天主的恩宠——他作了应当作的答复：\BibleQuote{断乎不可！我们这些\OriginalTerm{死于罪恶}{dead to sin}的人，如何还能在罪恶中生活呢？}没有比这说得更简要、也更中肯的了。在这邪恶的现世，天主的恩宠赐予我们什么比死于罪恶更有益的恩惠呢？因此，谁若选择并凭那使我们死于罪恶的恩宠而活在罪恶中，便显明自己对恩宠本身毫无感恩之心。但愿富于\OriginalTerm{慈悲}{mercy}的天主，赏赐你们健全而明智的思想，并使你们在各种善志与决意中坚忍不懈、日进于善，直至终结。请为你们自己、为我们、为一切爱你们的人，也为恨你们的人祈祷，使能获得这项恩赐——在兄弟般的和平中，恳切而警醒地祈祷吧。为天主而生活。若我在你们眼中蒙受些许恩待，请让\Person{弗洛路斯}兄弟到我这里来。\par

\SourceRecord{Translated by J.G. Cunningham.；Nicene and Post-Nicene Fathers, First Series；Vol. 1.；Edited by Philip Schaff.；Buffalo, NY: Christian Literature Publishing Co.,；1887.；<http://www.newadvent.org/fathers/1102215.htm>.}

% structure: p0001 source=h1 target=section reason=page-title

\section{书信二一八（主历四二六年）}

\Emph{\Person{奥古斯丁}致\Person{帕拉提努斯}，我至爱的大人与儿子，至为深情渴念的问候。}\par

1. 你向着主我们的天主所度的卓越坚毅且多结果实的生活，带给我们极大的喜乐。因为 \BibleQuote{你自幼便选择了训诲，到白发时仍可获得智慧。}\BibleRef{德 6:18} 因为 \BibleQuote{智慧就是人的白发，无瑕的生活便是老年。}\BibleRef{智 4:9} 愿那知道把好恩赐分施给自己儿女的主\BibleRef{玛 7:11}，将之赐给祈求、寻求、叩门的你。尽管你已有许多谋士、许多计策引导你走向永生的光荣之路，并且首先拥有基督的恩宠，这恩宠已在你心中以救恩的德能如此有效地发言，但我们自己，出于对你所欠的责任，借着在此回复你的问候，也向你提供一些劝言，目的不在于唤醒你这懒惰或贪睡的人，而是要激励并加快你这已在跑着的赛跑。\par

2. 我的儿子，你要在这赛跑中保持坚忍，就需要智慧，正如你起初是在智慧的影响下踏上这赛程的。那么，就让这成为 \BibleQuote{你智慧的一部分：知道智慧是谁的恩赐。}\BibleRef{智 8:20} \Quote{将你的行径委托给主，并要信赖他，他必成就；他必使你的义德如光出现，使你的判断如日中天。} \Quote{他必使你的路途平坦，使你的脚步安定平安。} 正如你从前轻看了此世的尊荣，免得你效法这世代的子女，因你开始贪恋的丰厚财富而自夸，如今，当你负起主的轭和他的担子时，不可依赖自己的力量；如此，\BibleQuote{他的轭是柔和的，他的担子是轻省的。}\BibleRef{玛 11:30} 因为在\BibleBook{}{圣咏集}中，\Quote{凡信赖自己力量的，}和\Quote{凡夸耀自己钱财丰厚的，}同受谴责。所以，正如你从前不在财富中寻求荣耀，并且极其智慧地轻视了你曾开始贪恋的，如今你也要谨防以自己的力量为信赖这阴险的诱惑；因为你不过是人，而\BibleQuote{凡信赖人的，是可咒骂的。}\BibleRef{耶 17:5} 但你务要全心信赖天主，他自己将成为你的力量，使你以虔诚和感恩之心信赖，并以谦卑和胆识向他说：主，我的力量，我爱慕你；因为即使天主的爱，当它完全时，\BibleQuote{驱逐了恐惧，}\BibleRef{若一 4:18} 这爱倾注在我们心中，不是靠我们的力量，即任何人的能力，而是如宗徒所说：\BibleQuote{借着赐给我们的圣神。}\BibleRef{罗 5:5}\par

3. \BibleQuote{所以，你们要警醒祈祷，为免陷于诱惑。}\BibleRef{谷 14:38} 这样的祈祷本身就劝诫你，你需要主的帮助，并且你不可将度圣善生活的希望寄托于自己。因为如今你祈祷，不是为了获得今世的财富和尊荣，或任何虚幻的人世财物，而是为免陷于诱惑，这件事若人凭自己的意志就能为自己成就，便无须在祈祷中祈求了。因此，如果我们的意志足以保护我们，我们就不会祈求免陷于诱惑；然而，如果我们缺乏躲避诱惑的意志，我们就不能如此祈祷。所以，也许我们有愿意的心，\BibleRef{罗 7:18} 这是由于我们因他的恩赐而成为智慧的人，但我们必须祈求能践行我们所愿意的。因你已开始实践这真正的智慧，你就有理由感谢。\BibleQuote{你有什么不是领受的呢？既然是领受的，为什么你还夸耀，好像不是领受的呢？}\BibleRef{格前 4:7} 就是说，好像你本可从自己拥有它似的。然而，你知道你是从哪里领受的，你当祈求那位以恩赐开启你的，也恩赐你得以成全。\BibleQuote{你们要怀着恐惧战栗，努力成就你们得救的事，因为是天主在你们内工作，使你们愿意，并使你们力行，为成就他的善意。}\BibleRef{斐 2:12} 因为\Quote{意志是由主预备的，}\Quote{善人的脚步由主安排，他喜爱他的道路。} 对这些事的圣善默想将保全你，使你的智慧成为虔诚，就是说，使你因天主的恩赐成为善人，而不忘恩负义于基督的恩宠。\par

4. 你的父母，真诚地与你同乐，因你在主内开始怀抱的更美好的希望，正切切渴望你的在场。但无论我们身体上与你分离还是同在，我们渴望在将爱倾注于我们心中的同一圣神内与你同在，这样，无论我们的身体寄居何处，我们的心灵绝不至彼此分离。\par

我们极其感谢地收到了你寄给我们的山羊鬃毛衣，在这礼物中，你已先我而提醒了应常祈祷，并在祈求中实践谦卑的责任。\par

\SourceRecord{Translated by J.G. Cunningham.；Nicene and Post-Nicene Fathers, First Series；Vol. 1.；Edited by Philip Schaff.；Buffalo, NY: Christian Literature Publishing Co.,；1887.；<http://www.newadvent.org/fathers/1102218.htm>.}

% structure: p0001 source=h1 target=section reason=page-title

\section{书信219（公元436年）}

\Emph{致最亲爱的可敬的弟兄、司铎职的同工普罗库鲁斯与西利努斯：奥古斯丁、弗洛伦提乌斯与塞昆迪努斯，在主内向你们致意。}\par

1. 当我们的儿子勒波里乌斯——你们因其执迷于错误而公正且恰当地责斥了他——被你们驱逐后，来到我们这里时，我们收留了他，视他为一个为使他得益而受苦的人；我们应当尽可能把他从错误中解救出来，并恢复他灵魂的健康。因为，你们在对待他时，遵守了宗徒的规诫：\BibleQuote{劝戒不守规矩的}，而我们的责任则是遵守那紧接着的命令：\BibleQuote{安慰忧心的人，扶持软弱的人}。\BibleRef{得前 5:14} 他的错误实在不小，因为他在某些关于天主独生子的道理上，既不赞同正确的事，也不体认真实的事。关于这位独生子，经上记载说：\BibleQuote{在起初已有圣言，圣言与天主同在，圣言就是天主}，又说：到了时候满足，\BibleQuote{圣言成了血肉，寄居在我们中间}；因为他否认天主成为人，认为若承认这一教理，就必然推论出：祂在那与圣父同等的\OriginalTerm{天主性体}{divine substance}受了某种不配的改变或朽坏，他却没有看出，如此一来，他便在\OriginalTerm{天主圣三}{Trinity}内引入了第四位，这完全违背了\OriginalTerm{信经}{Creed}的纯正教理和天主教的真理。然而，最亲爱的可敬的弟兄们，他作为一个容易犯错的人，陷入这一错误；我们在主的帮助下，尽力\BibleQuote{以温和的精神}教导他，尤其记得那位特选的器皿在给予我们所引用的这条命令时，还加上一句：\BibleQuote{你自己也要留心，免得也受试探}——免得有人因自己的灵修进步而沾沾自喜，以为像别人一样不会受试探了——并且紧接着他又加上一句有益且促进和平的话：\BibleQuote{你们应彼此分担重担，这样，你们就满全了基督的法律。因为人若自以为是什么，自己其实什么也不是，便是欺骗自己。}\par

2. 勒波里乌斯的这种恢复，若非你们先前对他应予纠正之处加以责备并予以惩罚，我们或许绝不可能实现。因此，同一位主，我们的神圣医生，运用祂自己的工具及仆人，曾借着你们击伤骄傲的他，又借着我们医治忏悔的他，正如祂亲口所说：\BibleQuote{我击伤，也加以治疗}。\BibleRef{申 32:39} 这位神圣的管理者和祂自己房屋的监督，借着你们拆毁了建筑中残缺的部分，又借我们按序重建了祂所拆毁的。这位神圣的农夫，以其细心勤勉，借着你们拔除祂田地中荒芜有害之物，又借着我们栽种有益丰产之物。所以，我们不要将荣耀归于自己，而应归于祂的仁慈，因为祂手中掌握着我们和我们的一切言语。正如我们谦卑地赞扬你们身为祂的仆人在我们上述儿子的事上所做的工作，也请你们以圣洁的喜乐为我们在工作中所完成的而欢欣。因此，以父爱和兄弟之爱接纳那位我们以仁慈的严厉所纠正的人吧。因为虽然工作的一部分是你们做的，另一部分是我们做的，但对于我们弟兄的救恩而言，这两部分不可或缺，都是同一的爱所做的。所以，是同一天主在两处工作，因为\BibleQuote{天主是爱}。\BibleRef{若一 4:8}\par

3. 所以，既然他因悔改而被我们接纳进入共融，也让他因这封信而被你们接纳；我们以为理当在这封信上附上我们的签名，以证明其真实性。我们毫不怀疑，你们在实践基督徒爱德时，不仅会欣然听闻他的改过，更会将这事告知那些曾被他错误绊倒的人。因为同他一起来到我们这里的人也与他一同得到纠正和恢复，正如他们当着我们的面在信上签了名所表明的。余下的就是，你们既因一位弟兄的得救而喜乐，便请屈尊回信给我们，使我们亦因此而喜乐。愿你们在主内平安，最亲爱的可敬的弟兄们；这是我们对你们的愿望：请记念我们。\par

\SourceRecord{Translated by J.G. Cunningham.；Nicene and Post-Nicene Fathers, First Series；Vol. 1.；Edited by Philip Schaff.；Buffalo, NY: Christian Literature Publishing Co.,；1887.；<http://www.newadvent.org/fathers/1102219.htm>.}

% structure: p0001 source=h1 target=section reason=page-title

\section{书信 220（公元 427 年）}

\Emph{致\Person{博尼法斯}阁下}，\Emph{我儿，被托付于天主慈悲的守护与引导，为今生与永世的救恩，\Person{奥思定}致意。}\par

1. 我从未能找到比他更值得信赖、也更能随时将我的信呈送于你耳旁的人，他就是这位基督的仆人与职事，执事\Person{保禄}，我们两人都极为亲爱。如今，主将他带到我面前，好使我有机会向你进言——不是关乎你在\Emph{这邪恶世界}中所拥有的权力和尊荣，也不是为着你可朽坏、有死的肉身的保全——因为这也终将逝去，且何时逝去无人可知——而是关乎那由基督许给我们的救恩；基督曾在此世受轻慢、被钉死，为教导我们宁愿轻看而非贪恋今世的美好事物，并将我们的情感与希望寄托于祂藉复活所启示的那个世界。因为祂已从死者中复活，如今\BibleQuote{不再死亡；死亡不再统治祂。}\BibleRef{罗 6:9}\par

2. 我知道你不乏朋友，他们就此生而言爱你，并为此给予你忠告，有时有益，有时相反；因为他们是人，所以尽管他们尽其智慧针对当下建言，却不知明日将发生何事。然而，为着天主而给你忠告，以阻止你灵魂的丧亡，却非易事；这并非因你缺少愿行此事的朋友，而是因为他们难以找到与你谈论这些事的机会。我自己就曾屡次渴望如此，却从未觅得处所与时机，能够按当有的方式对待一个我在基督内热切爱慕的人。此外，你知晓我在\Place{希坡}的情形，当你惠临探望我时，我因肉身虚弱几近不能言语。那么，我儿，现在听我言，至少我得以用一封信向你进言——这真是个难得的机会，在你们处于危险时，我本无力将这样的信件送达于你，既因担心信使遇险，又怕信件落入我所不愿之人手中。因此，你若认为我过于惧怕，我恳求你宽宥；无论如何，我已陈明我所忧惧的。\par

3. 因此，请听我言；不，更该听我们的主、天主藉我——祂卑微的仆人——所说的话。请忆念：在你的前妻——那位备受怀念者——仍然在世时，你是何等人；在她的死亡打击下，在事件还记忆犹新时，今世的虚幻使你退缩远离它；你又曾多么热切地渴望投身事奉天主。我们知道，且能作证，你在\Place{图布奈}与我们交谈时，你所说的心志与愿望。我的弟兄\Person{阿里比}和我，当时只与你二人相处。[因此，我恳求你忆念那次交谈]因为我不认为，如今你全神贯注的世俗挂虑能有如此力量，以致将其从你的记忆中全然抹去。那时，你渴望放弃一切你所从事的公务，退入神圣的静修，度如天主仆人的隐修生活。那么，是什么阻止你按这些愿望行事呢？岂非是因了我们向你说明此事后，你考虑到，若你以唯一目标——即保护教会免受蛮族骚扰，使它们能如宗徒所说：\BibleQuote{以全备的虔敬与端庄，度平静安宁的生活}\BibleRef{弟前 2:2}——去从事当时占据你的工作，这工作可对基督的众教会大有裨益；同时你也下定决心，不再从这世界寻求超出自身及眷属生活所需的，腰系全然贞洁自制的\OriginalTerm{束带}{cincture}，并且在军士装束之下，穿着那更为确实、更为坚强的神性铠甲。\par

4. 就在我们为你立下此志而满心喜乐时，你却登船起航，并娶了第二位妻子。你的起航，如宗徒所教导的，是对\BibleQuote{在上掌权者}\BibleRef{罗 13:1}的服从行为；然而，若非你放弃了先前献身的自制，被\OriginalTerm{贪欲}{concupiscence}所胜，你是不会再次结婚的。当我获知此事，我必须承认，我惊愕得哑口无言；但在悲痛中，我稍得安慰的是，听人说你拒绝娶她，除非她在结婚前成为公教徒；然而，那拒绝相信天主真子的异端，在你的家中竟如此猖獗，以致你的女儿由这些异端者施洗。如今，如果传说是真的（但愿天主使它是假的！），甚至有些献身于天主，作为祂婢女的，也被这些异端者重洗，那么我们应当以何等泪河痛哭这一巨大灾祸！此外，人们还说——或许是无稽之谈——一个妻子不能满足你欲念，你更与别的女人姘居，玷污了自己。\par

关于这些恶事——大家都看得分明，而且既多且大——自你结婚以来，你直接或间接成了它们的起因，我还能说些什么呢？你是一个基督徒，你有良心，你敬畏天主；那么，你自己想想那些我宁愿不说的事，你便会发现你应当为多么重大的恶事做补赎；而且我相信，主现在宽容你，救你脱离一切危险，正是要给你机会，让你以当行之道去行补赎。可是，如果你愿意听从圣经的劝言，我恳求你：\BibleQuote{归向主，不要迟缓，不要一天一天地拖延。}\BibleRef{德 5:8} 确实，你声称你所做的是有充分理由的，而我不能判断那理由是否充分，因为我无法听到辩论双方的陈词；然而，无论你的理由是什么——其性质目前既无需探究，也无需讨论——你能在天主面前断言，如果你不爱这世界上的好东西，你绝不会陷入你目前的困境吗？作为一个天主的仆人，正如我们过去所知道的你那样，你的本分是彻底轻视这些好东西，视它们毫无价值——固然，接受提供给你的东西，以便将其用于虔敬的用途，但不应贪求那拒绝给你的，或交托你照管的，以致因此陷入你目前的困境；在这种困境中，一方面爱好善，另一方面却做出恶事——固然，你亲自做的很少，但有许多是因你而做的；一方面惧怕那些即使有害也只伤害一时的事，另一方面却做出那些真正有害于永恒的事？\par

只提其中一件事——谁看不见有许多人追随你，是为了保护你的权力或安全？这些人虽然可能全都对你忠诚，而且不必担心他们中有任何背信，但他们渴望通过你获得某些好处，这些好处他们自己也贪求，并非出于虔诚的心，而是出于世俗的动机。这样一来，你——你的本分是约束并克制自己的激情——却被迫去满足别人的激情。为了做到这一点，必须做出许多令天主不悦的事；然而，终究，这些激情还是不能如此满足，因为在爱慕天主的人身上，这些激情最终更容易被治死，而在爱慕世界的人身上，即使暂时满足也做不到。因此，圣经说：\BibleQuote{你们不要爱世界，也不要爱世界上的事；谁若爱世界，圣父的爱就不在他内。原来世界上的一切：肉身的贪欲，眼目的贪欲，以及人生的骄奢，都不是出于圣父，而是出于世界。这世界和它的贪欲都要过去；但那履行天主旨意的，却永远存在，如同天主永远存在一样。}\BibleRef{若一 2:15} 所以，你与如此众多的武装人员为伍，他们的激情必须迁就，他们的残忍令人惧怕，如何能够——我且不说满足，甚至说部分满足——这些贪爱世界之人的欲望呢？你又急于防止更广大的恶事，除非你做出天主所禁止的事，从而招致那悬在头上的审判。为了放纵他们的贪欲，所造成的破坏已经如此彻底，如今连掠夺者也几乎无可掠夺了。\par

然而，此刻\Place{非洲}正遭受成群非洲蛮族的蹂躏，无人抵抗，而你却全神贯注于自身的困境，不采取任何措施避免这场灾难，对此我还能说什么呢？有谁会相信，有谁会担心，在\Person{博尼法斯}成为帝国和\Place{非洲}的伯爵，统率如此庞大的军队，拥有如此大的权柄之后，那位从前仅任保民官时，凭着一小队勇敢的盟军，便攻下蛮族的堡垒，威吓他们，从而令所有这些蛮族都得享太平的同一个人，如今竟让蛮族如此胆大妄为，深入侵袭，破坏掠夺如此之广，把以往人口稠密的广大地区变为荒漠？难道不是每个人都曾预言，一旦你取得伯爵的权柄，非洲的蛮族不但会被遏制，而且会成为罗马帝国的附庸吗？而现在，事态如何完全辜负了人们的希望，你自己也看到了；其实，在这件事上我无需多言，因为你自己的反思必定会提示出比我能用言语表达的多得多的东西。\par

8. 或许你替自己辩护说，这过错理当归于那些伤害你，且不按你职分所尽的功劳给予公道回报，反而以恶报善的人。这些事我无法审断。我恳求你反躬自省，探究此事，你知道，你在此事上面对的根本不是人，而是天主；作为在基督内的信徒，你务要敬畏，不要触怒祂。因为我更关注的是更高的原因，相信人应把\Place{阿非利加}的巨大灾难归于他们自己的罪过。然而，我不愿你成为天主用来施加现世惩罚的工具——那些邪恶不义的人中的一员；因为天主公义地利用他们的恶意向他人施加现世的审判，却为那不肯悔改的不义者自己预备了永罚。你要定志仰望天主，注视基督，祂赐予你如此巨大的恩惠，并为你忍受如此巨大的痛苦。凡渴望属于祂的国度，永远幸福地与祂同在、受祂统治的人，都爱自己的仇敌，善待憎恨自己的人，为迫害自己的人祈祷；\BibleQuote{爱你们的仇敌，善待恼恨你们的人，为迫害你们的人祈祷；}\BibleRef{玛 5:44} 并且，就算为了惩戒，他们有时会采用令人生厌的严厉手段，但他们从未放弃至诚的爱。如果这些现世且短暂的恩惠是\Place{罗马帝国}给予你的——因为那帝国本身是现世的，而非天上的，它不能赐予自己无权拥有的东西——假如，我是说，你蒙受了恩惠，就不要以恶报善；假如你遭遇恶待，就不要以恶报恶。这两种情况中哪一种发生在你身上，我不愿讨论，也无法判断。我对一位基督徒说——既不要以恶报善，也不要以恶报恶。\par

9. 或许你对我说：在如此艰难的环境中，你希望我做什么呢？如果你向我征询世俗的观点，即如何在这暂世的生命中保全你的安全，保全甚至增加你目前拥有的权力和财富，我不知道该怎么回答你，因为对于这些如此不确定的事物，任何建议都必定同样不确定。但如果你就你与天主的关系及你灵魂的救恩来征询我，并且如果你敬畏这真理之言，就是经上说的：\BibleQuote{人纵然赚得了全世界，却赔上了自己的灵魂，为他有什么益处？}\BibleRef{玛 16:26} 那么我可以给你一个明确的答复。我已准备好忠告，你理应听从。可是，我又何必再说我已经说过的话呢。\BibleQuote{不要爱世界，也不要爱世界上的事；谁若爱世界，圣父的爱就不在他内。因为凡世界上的事，肉身的贪欲，眼目的贪欲，以及人生的骄奢，都不是出于圣父，而是出于世界。这世界和它的贪欲都要过去；但那履行天主旨意的，却永存不朽。}\BibleRef{若一 2:15} 这就是忠告！抓住它并付诸行动。显示出你是个勇毅的人。战胜你爱恋世俗的欲望。为过去生活中的恶事做补赎，那时你被私欲偏情战胜，被罪愆的欲望牵引而去。如果你接受这忠告，持守它并行出来，你就既能获得那确定的福祉，又能在这些不确定的事物中度日，而不失去你灵魂的救恩。\par

10. 但也许你又要问我，你深陷如此重大的世俗困难中，怎样能做到这些呢。恳切祈祷，并引用圣咏的话对天主说：\BibleQuote{求祢救我脱离我的忧患}，因为这些忧患在产生它们的私欲偏情被战胜时就终止了。那位俯听你的祈祷，以及我们为你代祷的祈求，使你从有形可见战争的众多大危险中解救出来的天主——在这场战争中，身体有失去迟早必亡生命的危险，但灵魂除非被邪恶的私欲俘获，否则不致丧亡——我说，那位天主也会俯听你的祈祷，使你在不可见的精神斗争中，战胜你内在不可见的仇敌，就是你的私欲偏情本身，并能如此享用世界，而不滥用世界，以致用世上的美善行善，不致因拥有它们而变坏。因为这些事物本身是好的，除非掌管天地万有的那位赐予，便不会给予人。祂的恩赐不应被视为恶，因此它们也赐予善人；同时，它们也不应被视为伟大或至高的善，因此也赐予恶人。此外，这些事物从善人手中被夺去，是为试探他们；从恶人手中被夺去，是为惩罚他们。\par

11. 有谁如此无知，如此愚昧，竟然看不出，这能死身体的健康、可朽肢体的力量、战胜敌人的胜利、现世的尊荣和权力，以及其他一切纯属世间的利益，既赐给善人也赐给恶人，同样也从善人和恶人手中被夺去？但灵魂的救恩，以及肉身的永生、义德的力量、战胜敌对私欲的凯旋、光荣、尊荣和永恒的平安，却只赐给善人。所以，要爱慕这些，渴求这些，尽你所能寻求它们。为了获得并持守这些，你要施舍，要热切祈祷，要在身体不受伤害的情况下尽力禁食。但不要贪爱这些世间财物，不论它们多么丰盛。要这样使用它们：借着它们行许多善事，却绝不为它们做一件恶事。因为所有这类事物都将消逝；但善工，甚至那些借着世间可朽之物完成的善工，却永远不会消逝。\par

12. 如果你现在没有妻子，我就会对你说我们在\Place{图布奈}所说的话，即你应该度圣洁的贞洁生活；我还要补充说，你现在应当去做我们当时阻止你去做的事，那就是：在不损害公共福祉的前提下，尽可能从军旅劳苦中抽身，并且享受你当时所渴望的那种闲暇，去度圣人们的团体中的生活；在那团体里，基督的战士在静默中战斗，不是为了杀人，而是为了\BibleQuote{与率领者和掌权者，以及天界里邪恶的鬼神交战}\BibleRef{弗 6:12}，也就是魔鬼及其使者。因为圣人们战胜了不可见的仇敌，而他们之所以能战胜这些不可见的仇敌，正是因为他们战胜了感官所触及的事物。然而，我之所以不能劝你过那种生活，是因为你有妻子，因为若没有她的同意，你度贞洁的誓愿的生活是不合法的；因为，尽管你在图布奈作出声明之后又再婚，犯了错误，但她对此一无所知，因此她是清白无罪、不受任何约束地成为你的妻子的。但愿你能说服她同意贞洁的誓愿，这样你就能毫无阻碍地将你知道应该归于天主的东西归还给祂！然而，如果你不能与她达成这个协议，就要用一切方法谨慎地保守婚姻的贞洁，并祈求天主，祂将救你脱离困境，使你将来能做到目前尚不可能的事。不过，这并不影响你应尽的义务：爱天主，不爱世界，即使仍不得不投身战事，也要在战事的忧虑中坚守信德，寻求和平；使此世的美物在善工中为人服务，而不在追求这些世物时作恶——在所有这些责任上，你的妻子不阻碍你，即使有所阻碍，也不应该是阻碍。\par

我亲爱的儿子，我写这些话，是出于我对你的爱——这爱不是着眼于这个世界，而是着眼于天主；而且，因为我想起了圣经的话：\BibleQuote{你不要责斥轻狂的人，免得他恨你；却要责斥智慧者，他必爱你}\BibleRef{箴 9:8}，所以我理应认为你当然不是愚昧人，而是智慧人。\par

\SourceRecord{Translated by J.G. Cunningham.；Nicene and Post-Nicene Fathers, First Series；Vol. 1.；Edited by Philip Schaff.；Buffalo, NY: Christian Literature Publishing Co.,；1887.；<http://www.newadvent.org/fathers/1102220.htm>.}

% structure: p0001 source=h1 target=section reason=page-title

\section{信函227（公元428或429年）}

\Emph{致年迈的\Person{阿利比乌斯}，\Person{奥古斯丁}致以问候。}\par

\Person{保禄}弟兄已平安抵达此处：他报告说，他为所从事的事务付出的辛劳已获成功；主会准许他在那件事上的麻烦就此结束。他向你致以热烈的问候，并向我们讲述了关于\Person{加比尼安努斯}的消息，这消息令我们喜乐，即：他因天主的仁慈而在其案子中获得圆满解决后，如今不仅名义上是基督徒，而且真诚地成为信仰的极好皈依者，并在最近的复活节受洗，心中与口唇上都拥有他所领受的恩宠。我多么渴望他，无法用言语表达；但你知道我爱他。\par

医学系主任\Person{狄奥斯科鲁斯}也宣认了基督信仰，同时获得了恩宠。听我讲述他皈依的经过，因为若非某种奇迹，他那固执的颈项与狂妄的口舌是无法被制服的。他的女儿，他生命中唯一的安慰，病了，而且病情变得如此严重，以致据她父亲自己承认，她的生命已无希望。据报告——而且此报告的真实性无可置疑，因为在\Person{保禄}弟兄回来之前，此事已由\Person{佩雷格里努斯}伯爵向我提及，他是一位极为可敬且真正虔诚的基督徒，与\Person{狄奥斯科鲁斯}和\Person{加比尼安努斯}同时受洗——我再说，据报告，这位老人最后感到不得不恳求基督的怜悯，便立下誓言：若他看见女儿恢复健康，就成为一名基督徒。女儿康复了，但他背信弃义地收回了誓愿。然而，主的手仍然伸着，因为他突然被击打而失明，随即这灾祸的原因便印入他的心中。他高声认错，再次立誓：若视力恢复，必履行所许的愿。他恢复视力，履行了誓愿，但天主的手仍伸着。他未曾将信经牢记在心，或许拒绝背诵，并借口无力为之。天主已看到这一切。在他完成所有入门礼仪之后，立刻患上瘫痪，累及多处，几乎遍及四肢，甚至舌头。随后，在一次梦境的警示下，他书面承认，此事之发生乃因他未曾背诵信经。认错之后，所有肢体的功能均得恢复，唯有舌头例外；然而，尽管如此折磨，他仍通过文字表明，他已学习信经，并仍铭记于心；这样就完全摧毁了他那轻浮的多言（如你所知，这多言本损伤了他天性的和善，并使他在嘲弄基督徒时极其不敬）。我能说什么呢，只有\BibleQuote{请赞美主，歌颂称扬他，直到永远！阿们。}\par

\SourceRecord{Translated by J.G. Cunningham.；Nicene and Post-Nicene Fathers, First Series；Vol. 1.；Edited by Philip Schaff.；Buffalo, NY: Christian Literature Publishing Co.,；1887.；<http://www.newadvent.org/fathers/1102227.htm>.}

% structure: p0001 source=h1 target=section reason=page-title

\section{书信二二八（主后428或429年）}

\Emph{致他的圣善弟兄、同为主教的\Person{奥诺拉图斯}}，\Emph{\Person{奥斯定}在主内谨致问候。}\par

1. 我本以为，将我给我们的弟兄及同工主教\Person{郭德福尔德乌斯}的信抄送一份给尊驾，就可以免去您加给我的重担——您询问我，在这些时日临到我们的危险中，我当作什么。因为，虽然我写得简短，但我认为我没有遗漏任何在这问题上通信时我必须说或提问者必须听的必要内容：因为我说过，一方面，那些愿意迁往坚固之地的人，若能做到，不应被禁止；另一方面，我们也不应扯断那由基督的爱所系结的联系，这联系使我们身为\OriginalTerm{圣职人员}{ministri}，不可离弃我们有责任服务的教会。我在那封信中用到的话是：\Quote{因此，不论我们所在的天主子民的会众多么微小，只要我们的服务为他们如此必要，明显有义务不可撤去，我们就只好对主说：\InnerQuote{求你作我们避难的堡垒，作坚固的保障。}}\par

2. 但这建议并未获得您的认同，因为正如您在信中说的，我们试图违背主的诫命或榜样行事，是不合宜的，主曾劝勉我们，应在一座城受迫害时，逃到另一座城去。我们确实记得主的话：\BibleQuote{人若在这城迫害你们，你们就逃到另一城去；}\BibleRef{玛 10:23}但谁能相信，主愿意我们在这种情况下这样做，以致祂用自己宝血所赎回的羊群，因牧者的离开而缺乏那为生命所不可少的事奉呢？基督自己曾这样做吗？当祂在婴孩时，由父母带着逃往\Place{埃及}？没有，因为那时祂还没有聚集我们可以说被祂离弃的教会。或者，当宗徒\Person{保禄}\BibleQuote{被人从窗口用篮子缒下，}以逃脱敌人的捉拿，因而脱身时，\BibleRef{格后 11:33}\Place{大马士革}的教会就缺乏基督仆人所必须的劳苦吗？在他离开之后，其他定居那城的弟兄岂不就供应了一切所需的事奉吗？因为宗徒这样做是在他们的要求下，为保全自己的生命以益于教会，而那迫害者当时特别想要毁灭他的生命。所以，那些是基督的仆人、祂圣言和\OriginalTerm{圣事}{sacramentum}的圣职人员，当遵照祂所命令或允许的去做。当有圣职人员特别被迫害者追索时，他无论如何应从一城逃到另一城去，只要教会不因此而无人，而其他未被特别追索的人仍留下来，为他们的同工提供神粮，因为他们知道若不然这些同工便无法维持灵性生命。然而，当主教、圣职人员和平信徒众人的危险都相同时，就不应让那些依靠别人帮助的人，被他们所依靠的人离弃。在这种情况下，或是一同迁往坚固之地，或是那些必须留下的人不可被那些必须为他们提供教会所需的人离弃；这样，他们或一同生活，或一同承受那一家之父所命定他们受的苦。\par

3. 但若遇到所有人一同受苦，无论有些人受苦较轻，有些人较重，或所有人都受同样的苦，就不难看出他们中间谁是为他人而受苦：显然是那些本可通过逃走而免于这些灾难，却选择留下来，而不愿离弃那些需要他们的人。这样的行为特别证明了宗徒\Person{若望}所赞许的爱，他说：\BibleQuote{基督为我们舍了自己的性命：我们也应当为弟兄们舍命。}\BibleRef{若一 3:16}因为那些投身于逃亡，或只因环境阻碍其计划而未能逃脱的人，如果被捉拿而受苦，他们只是为自己受这苦，而不是为弟兄；但那些因决意不离弃他人——他人的基督徒福祉有赖于他们——而受牵连受苦的人，无疑正是\BibleQuote{为弟兄们舍命。}\par

4. 因此，我们听到有人将某位主教的话说成：\Quote{如果主命令我们在那能收获殉道果实的迫害中逃走，那么，在蛮族入侵造成的无益苦难面前，若能逃走，我们就更当逃走了！}这话诚然真实而且值得接受，但只适用于那些不受教会职务约束的人。因为，一个人若本有能力逃脱敌人的暴力，却选择不逃，以免因此离弃基督的事奉——没有了这事奉，人既不能成为基督徒，也不能过基督徒的生活——这样的人，必定得到更大的爱的赏报，比那并非为弟兄，而是为自己而逃，却被迫害者捉拿，因拒绝背弃基督而殉道的人，所得的赏报更大。\par

5. 那么，对于你在前信中如此陈述的立场，我们该说什么呢：——\Quote{我看不出继续留在教堂里能为自己或民众带来什么好处，无非是眼睁睁看着男人被杀、妇女受辱、教堂被焚，而我们在酷刑中奄奄一息，为要逼我们交出并不拥有的东西}？天主有能力垂听祂子女的祈祷，并消除他们所恐惧的那些事；我们也不应该因为不确定的恶事，就断然决定放弃那牧职，因为没有这牧职，民众必定遭受毁灭，不是在今生的事务上，而是在那另一生命上，那生命本该以不可比拟地更大的殷勤和关切去照料。因为如果那些我们所担心的、可能临到我们所在之地的恶事是确然的，那么所有我们为了他们才应留守的人，就会在我们之前逃走，从而免去我们留守的必要；因为没有人说圣职人员有义务留在任何地方，如果那里已无人需要他们的牧职。这样，有些圣洁的主教从\Place{西班牙}逃走，那时他们的信友团体在他们逃走之前已被消灭，成员要么逃散，要么死于刀剑，要么在城镇被围困时丧生，要么被掳去：但那个国家有更多的主教留在这些蔓延的危险之中，因为他们为之留下的人仍留在那里。如果有人抛弃了自己的羊群，我们说这是不应该做的，因为他们这样做不是由神圣权威教导，而是出于人性的软弱，或因错误受骗，或因恐惧而屈服。\par

6. [我们坚持，作为另一种可能，他们是被错误所骗，] 否则，他们为什么认为必须不加区别地遵从那条诫命，就是他们读到要从一城逃到另一城的那条，却不厌恶那股\Quote{佣工}的品性，他\BibleQuote{看见狼来，就逃走，因为他不关心羊}？\BibleRef{若 10:12} 他们为什么不同样尊重主的这两条真实的话，那一条允许或命令逃跑，这一条却谴责和责备逃跑，并努力理解它们，从而发现它们并非彼此对立，正如事实那样？又怎能找到它们的协调呢，除非牢记我上面所证明的，即在迫害压力下，我们这些基督的仆役只应在两种情形之一时逃离我们所在的地方：要么那里没有我们要服务的会众，要么虽然有会众，但必要的牧职可以由其他人提供，而这些人没有我们那样迫切的逃跑理由？关于这点，我们已经提到一个例子，就是宗徒\Person{保禄}被人用筐子从城墙上缒下逃脱，当时他正被追捕者亲自搜寻，而当地有其他没有同样逃跑必要的人，他们的留下可以防止教会缺少圣职人员的服务。另一个例子是圣\Person{亚大纳削}，\Place{亚历山大里亚}的主教，他在\Person{君士坦提乌斯}皇帝特别想捉拿他时逃跑了，而留在\Place{亚历山大里亚}的公教徒并没有被其他天主的仆人抛弃。但是，当民众留下而天主的仆人逃跑，他们的服务被撤去，这岂不就是那不关心羊的\Quote{佣工}的可耻逃跑吗？因为狼要来——不是人，而是魔鬼，他常常使那些缺少主身体日常牧职的信友背教堕落；因此，不是\BibleQuote{因你的知识}，而是因你的无知，\BibleQuote{那软弱的弟兄为基督而死的，沉沦了}。\BibleRef{格前 8:9}\par

7. 然而，对于那些不是因为受错误欺骗，而是因为被恐惧所胜而逃跑的人，他们为什么不靠着主所赐的怜悯和帮助，勇敢地与恐惧作战，免得有无比更重、更可怕的恶事临到他们呢？这种对恐惧的胜利，无论在何处，只要爱天主的火焰在心中燃烧，没有世俗的烟雾，就能赢得。因为爱说：\BibleQuote{谁软弱，我不软弱呢？谁跌倒，我不焦急呢？}\BibleRef{格后 11:29}但爱来自天主。因此，让我们恳求那要求我们爱祂的天主，将这爱赐给我们，并在爱的感召下，让我们更害怕基督的羊被那能刺入人心的属灵邪恶之剑所杀，胜过害怕他们倒在只能伤害身体的剑下，这身体人无论如何，或早或晚，总会以某种方式死去。让我们更害怕信德的纯洁因内心腐败的污损而丧失，胜过害怕我们的妇女遭受暴力凌辱；因为如果贞洁在心灵中得以保全，它就不会被这种暴力所毁，既然当受害者没有卑劣地同意犯罪，而只是意志未曾同意而屈服于他人的行为时，贞洁甚至在身体上也不会被毁。让我们更害怕\OriginalTerm{活石}{living stones}中的生命火花因我们的缺席而熄灭，胜过害怕我们地上建筑的石头和木料在我们眼前被焚。让我们更害怕基督身体的肢体因缺乏灵性食粮而死亡，胜过害怕我们自己身体的肢体被敌人的暴力制服，遭受酷刑折磨。这并不是说，当我们有能力时不应该躲避这些事，而是说，当这些事无法避免而不陷入不虔时，我们应宁愿受苦，除非或许有人主张，那在特别需要虔诚服务之时，撤去虔诚所必需之服务的仆人，不算犯了不虔之罪。\par

8. 我们岂忘记了，当这些危险达到极点，无法以逃避摆脱时，一大群不同性别、各种年龄的人，惯常聚集在教堂里——有的急切要求\OriginalTerm{圣洗}{baptism}，有的要求\OriginalTerm{修和}{reconciliation}，有的甚至要求行\OriginalTerm{补赎}{penance}，众人莫不呼求藉施行\OriginalTerm{圣事}{sacraments}而获得安慰与坚固？在此关头，若天主之仆不在其岗位，那些未经\OriginalTerm{重生}{regenerated}或未获\OriginalTerm{罪赦}{loosed from sins}而离世者，将遭逢何等浩劫！他们信仰的亲属又该何等忧伤，因他们无法在\OriginalTerm{永生}{eternal life}安息之福中与这些沉沦者同聚！总之，由于缺乏礼仪和无人施行，众人的哀号与不少人的愤怒咒骂何其响亮！看看，对现世灾难的恐惧能产生什么，它又可能成为导致何等大量永远灾难的原因。但若仆人在其岗位，藉天主所赐的力量，所有人都得到帮助——有的领受\OriginalTerm{圣洗}{baptism}，有的与教会\OriginalTerm{修和}{reconciled}。无人被亏缺\OriginalTerm{主圣体}{Lord's body}的共融；所有人都得安慰、得建树，并被劝勉祈求那能成就一切的天主，转离所畏惧的一切——做好两种准备，以致\BibleQuote{如果这杯不能离开}他们，愿祂的旨意成就\BibleRef{玛 26:42}，祂绝不能愿意任何恶事。\par

9. 无疑，你现在看清了（据你的信，你以前未能看清），若在灾难临头时，基督仆人的临在不从他们中撤走，基督徒民众可获得多大的益处。你也看清了，当\BibleQuote{他们求自己的事，而不求耶稣基督的事}\BibleRef{斐 2:21}，当他们在\OriginalTerm{爱德}{charity}上匮乏——经上论这爱德说：\BibleQuote{不求自己的益处}\BibleRef{格前 13:5}——，且未效法那位说过\BibleQuote{我不求自己的利益，只求大众的利益，为使他们得救}\BibleRef{格前 10:33}的人；并且，倘若不是为了那些需要他的人而保存自己，他本不会逃避帝国迫害者的阴险攻击——正因此他说：\BibleQuote{我正夹在两者之间：我渴望求解脱而与基督同在一起：这实在是再好没有了；但存留在肉身内，对你们却更为重要。}\BibleRef{斐 1:23}，他们的缺席造成了多大的损害。\par

10. 在此，也许有人会说，当如此邪恶即将临头时，天主的仆人们应该以逃避保全性命，好使自己留待更和平的时期为教会谋益。某些人如此行是对的，但前提是不乏他人可供应教会职务，且工作并未被所有人遗弃，正如我们上文所陈述的\Person{亚大纳削}所做的那样；因为整个\OriginalTerm{公教世界}{Catholic world}知道，他为教会如此行是何等必要，这位以言语和爱心服侍她、抵抗\OriginalTerm{亚略异端者}{Arian heretics}的人，其长久生命又是何等有益。但当危险为所有人共有时，却绝不该如此行；并且最应畏惧的，是恐怕有人认为这样做不是出于保障他人福祉的愿望，而是出于害怕丧失自己的生命，从而因其擅离职守的榜样而造成更大的损害，远超过他保全性命以备将来服务所能行的一切善。最后，请看\Person{圣达味}怎样听从了民众的迫切请求，不让自己暴露于战阵的危险，正如叙述中所说：\BibleQuote{免得熄灭以色列的明灯}\BibleRef{撒下 21:17}，但他自己并非首先提出；若他首先提出，就会使许多人效法那可能被他们归咎于他的怯懦，以为他如此行并非顾及别人的益处，而是出于对自己生命的恐惧激动。\par

11. 这里提示了另一个我们不可视为不值一提的问题。因为，若教会的利益不可忽视，且若此利益使然，当大灾难临头时，某些牧者应当逃亡，以便幸存下来，去牧养那些灾难过后仍存留的人——问题便产生了：当看来若不逃亡，所有人必将一同灭亡时，当作什么？倘若毁灭者的怒火仅仅针对教会的牧者，又当作什么？我们该怎么回答？教会是否要因惧怕其牧者死于其职守中，而更不幸地被剥夺他们的服务，从而因逃避职守而丧失其牧者的服务？当然，若\OriginalTerm{平信徒}{laity}免于迫害，他们便有能力以某种方式庇护和隐藏他们的\OriginalTerm{主教}{bishops}和\OriginalTerm{圣职人员}{clergy}，正如统辖万有的那位将帮助他们，祂以其奇妙能力，甚至能保全那不逃避危险的人。然而，我们探询在此情形下何为义务之路的理由，正是为了避免因每遇事便期望天主的神奇干预，而负上试探主的罪责。\par

的确，当危险为平信徒与圣职人员所共有，犹如同船上的商人与水手共担风暴之险时，灾难风暴的严酷程度有所区别。但我们决不可轻视我们的这艘船，以至于认为船员，尤其是舵手，在危难时刻可以弃船逃生，即使他们有能力跳上小船甚至游到岸上。因为对于那些我们担忧若我们离弃职守便会丧亡的人，我们所惧怕的恶并非终必来临的\OriginalTerm{暂世死亡}{temporal death}，而是\OriginalTerm{永死}{eternal death}——而永死之来与不来，全在于我们是疏忽还是采取规避之策。再者，当平信徒与圣职人员的生命同遭共同危险时，我们有何理由认为敌人所到之处，一切圣职人员都将被害，而所有平信徒却不会死？在那种情况下，圣职人员及他们所服侍的人将同时离世。又为何不能希望，正如一些平信徒可能幸存，一些圣职人员也可能被留下，借以向他们施行必要的圣事？\par

12. 但愿在此境况中，天主的仆人们所争论的问题是：哪些人应当留下，以免教会因所有人躲避危险而陷于荒废；哪些人应当逃离，以免教会因所有人丧身于危险而同样陷于荒废。在所有同样炽热爱德并满足爱德要求的弟兄之间，将会兴起这样的争论。若在任何情形下，都无法以其它方式结束辩论，我认为应当以抽签决定谁留下、谁逃离。因为那些声称自己理应优先于他人逃离的人，要么显得胆怯，如同不愿面对逼近的危险，要么显得傲慢，仿佛认为自己的生命对于教会的利益比他人的生命更需保全。再者，或许那些更卓越的人会率先选择为弟兄们舍命；如此，得以借逃离而保全的，将是那些生命较为无足轻重的人，因为他们在辅导和治理教会方面的能力较弱；然而，这些人生若虔诚明智，便会抗拒那些人的愿望——他们看到，一方面，那些人活着对教会更为重要，另一方面，那些人宁肯舍命也不愿逃离危险。在这种情况下，正如经上所载：\BibleQuote{抽签平息争端，也能解散强权}\BibleRef{箴 18:18}，因为，在这类困境中，天主比人判断得更好：祂或乐于召叫祂仆人中较卓越者去领受苦难的赏报，而宽免软弱的；或使软弱者坚强以承受试炼，然后让他们离开此世，因为他们的生命对教会所能做的服务，不如那些比他们坚强者的生命。我承认，若采用这样的抽签方法，将是一件不寻常的事；但若真的在某事上这样做了，谁敢指责呢？除了无知或有偏见者，谁又会迟疑不给予它应得的赞许呢？然而，若因无先例可循而不采用抽签，那就务必确保教会不因任何人的逃离而失去在这等巨大危险中对它来说尤为必要和应有的牧职。勿让任何人因自诩在某样恩宠上明显优越，便认为自己的生命比他人更宝贵，因而更有权寻求逃离以自保。因为凡如此以为者，太过自满；而凡如此言说者，必令所有人对他不满。\par

13. 有些人以为，主教和圣职人员在如此危险中不逃离而留下，可能会误导民众，因为民众见到管理他们的人留下，便也不逃离危险。然而，他们很容易反驳这种异议和误导他人的指责，只要向会众说明，说：\Quote{我们并非因自己的缘故，而是因你们的缘故才留在这里，好能继续向你们施予我们知道为你们的救恩——即在基督内的救恩——所必需的一切；因此，你们若选择逃离，便从我们被束缚于留下的义务中，将我们释放了。} 我认为，当确实觉得迁往更安全之处更为有益时，就当这样说。如果在他们公开说了这些话之后，全体或部分人会说：\Quote{我们全在于祂的掌握，无论逃到何处，无人能逃脱祂的义怒；也无论命运将我们抛于何处，都能寻获祂的仁慈；那些或因已知无法克服的困难受阻，或因不愿追寻未知的避难所——在那里危险可能只是变换而非终结——而宁可不往别处去的人，}——那么，那些决意留下的人绝对不应缺少基督徒司牧的服务。反之，若在听了他们的主教和圣职人员如上所说之后，人民宁愿离开该地，那么留下的人就不再负有留下的义务，因为他们原本只是为了那些人才留下的，既然已无人因他们的缘故而需留下，他们便无需再留。\par

14. 因此，凡在教会不因其逃离而失去必要服务的情况下逃离危险的人，是在做主所命令或许可的事。但那圣职人员，若其逃离导致基督的羊群失去维持其灵性生命的滋养，他便是那\BibleQuote{佣工看见狼来，便弃羊逃跑，因为他并不关心羊。}\par

怀着我所确知的真诚之爱，我现已就这个问题写下我所信为真实之事，因为我至爱的弟兄你征询了我的意见；但我并未强令你遵从我的建议，如果你能找到比我更好的见解。无论如何，在这些危险中，我们所能行的，莫过于不断恳求主我们的天主怜悯我们。至于我所写的那件事，即\OriginalTerm{牧者}{ministers}不应离弃天主的教会，有些智慧和圣善的人，藉着天主的\OriginalTerm{恩赐}{gift}，既得以愿意，也得以践行此事，并且在坚定地追求其目的的过程中，丝毫没有动摇，即使他们遭到了\OriginalTerm{诽谤者}{slanderers}的攻击。\par

\SourceRecord{Translated by J.G. Cunningham.；Nicene and Post-Nicene Fathers, First Series；Vol. 1.；Edited by Philip Schaff.；Buffalo, NY: Christian Literature Publishing Co.,；1887.；<http://www.newadvent.org/fathers/1102228.htm>.}

% structure: p0001 source=h1 target=section reason=page-title

\section{书信229（公元429年）}

\Emph{致\Person{达里乌斯}，当受尊崇的、显赫且大有能力的领主，在主内亲爱的儿子，\Person{奥斯定}谨在主内致意。}\par

1. 我从我的圣善弟兄及同为主教的\Person{乌尔巴努斯}与\Person{诺瓦图斯}那里，得知了你的品格与地位。前者在\Place{迦太基}附近的\Place{希拉里}城，及更近时在\Place{西加}城，与你相识；后者则在\Place{西提非斯}与你相识。因着他们，我便不能视你为与我素不相识之人。因为尽管我肉身的软弱与年迈的衰微不容我亲与交谈，却不能因此说我未曾见过你；因为乌尔巴努斯友善来访时的言谈，及诺瓦图斯的书信，向我形容的，并非你面孔的特征，而是你心灵的容貌，所以我已看见了你，并且以更大的喜乐看见了你，因为我看的不是外在的仪表，而是\OriginalTerm{内在的人}{inner man}。你品格的这些特征，不但我们，而且赖天主的仁慈，你自己也喜乐地看见，正如在圣福音这面镜子中观看一样；在圣福音中，那位即真理者亲口所说的话记载着：\BibleQuote{缔造和平的人是有福的，因为他们要称为天主的子女。}\BibleRef{玛 5:9}\par

2. 那些战士诚然伟大，配受特别的尊崇，不但因其至极的勇武，且（这是更高的称赞）因其卓越的忠诚；藉着他们的辛劳与危险，并赖天主护佑与助佑的降福，先前未平服的敌人得以征服，国家获得和平，行省得以归顺。然而，以言语平息战争，胜于用刀剑杀人；藉和平促成或维持和平，而非藉战争——这是更崇高的荣耀。因为那些作战的人，如果是善人，无疑也寻求和平；但却是通过流血达成。而你的使命，正是阻止流血的灾祸。因此，你的特权在于免除那别人不得不造成的灾难。所以，我当受尊崇的、显赫且大有能力的领主，在主内最亲爱的儿子，你当为天主所赐予你的这独一无二、伟大而真实的福分而欢欣，并在天主内享受这福分；你之所以成为你今日的样子，并得以承担如此的事工，全当归功于祂。愿天主\Quote{坚固祂藉你为我们所成就的。} 请收纳我们这致意，并惠然赐复。我从我的弟兄诺瓦图斯的来信中看到，他曾费心使博学的尊驾，通过我的著作，也能认识我。若你已读过他给你的这些著作，那么，在你的内心知觉中，我也已为你所认识。依我所能判断，你若以亲爱之情而非挑剔之心去阅读，它们当不至于大大地令你不悦。若你为此信及我的著作惠赐一封回信，这虽非过分之求，却将是一大厚惠。我以应有的爱心向和平的保证致意，这保证你赖我们的主和天主的恩宠，已欣然领受。\par

\SourceRecord{Translated by J.G. Cunningham.；Nicene and Post-Nicene Fathers, First Series；Vol. 1.；Edited by Philip Schaff.；Buffalo, NY: Christian Literature Publishing Co.,；1887.；<http://www.newadvent.org/fathers/1102229.htm>.}

% structure: p0001 source=h1 target=section reason=page-title

\section{第231封信（公元429年）}

\Emph{致\Person{达里乌斯}，我的儿子及基督的肢体，\Person{奥古斯丁}，基督与基督众肢体的仆人，在主内致意。}\par

1. 你曾请求我回信作为我已欣然收悉你信的证据。那么，请看，我再次写信；然而我既不能以此回信，也不能以任何其他方式，无论是简略地写还是尽情挥洒，来表达我所感受到的喜悦，因为无论言辞多少，我都不可能向你表达那言语无法表达之事。我虽然能言善道，却并不雄辩；但也决不会允许任何人，无论他多么雄辩，即使他能像我一样洞悉我的内心，去做我力所不及的事，即在一封信中，无论其文笔如何优美、篇幅如何长，来描述你的信在我心中所产生的效果。那么，我只有这样向你表达你所愿了解的事，好使你从我未能表达的话语中领会到那未能表达的情意。那么，我该说什么呢？我因你的信而欣喜，极其欣喜——这个词的重复不是简单的重复，而仿佛是一种恒久的确认；因为不可能一直这样说，所以至少重复一次，这样或许我的感受才得以表达。\par

2. 如果此时有人问，究竟你信中什么使我如此欣喜——\Quote{是它的雄辩吗？}我将回答：不是；而他或许会再问：\Quote{那么，是那些给予你的赞美吗？}但我仍将回答：不是；我之所以如此回答，并不是因为这些不在这封信中，因为信中雄辩如此出众，分明可见你天生具有最高才华，并且受过极为精心的教育；而且你的信无疑充满了对我的赞美。那么有人可能会说：\Quote{这些不令你欣喜吗？}是的，真的欣喜，因为正如诗人所说，\Quote{我的心并非角质的}，所以我不会注意不到这些，或注意到却不会不喜悦。这些的确令我欣喜；但它们与我所说我极其欣喜的事有何关系呢？你的雄辩令我欣喜，因为它既文思温雅，又措辞庄重；虽然我确实不会因来自所有人的各种赞美而欣喜，而只是因那些你认为我配得的赞美，以及来自像你那样的人的赞美——即那些为基督的缘故爱他众仆人的人——我无法否认，我对你信中的赞美感到欣喜。\par

3. 善于思考且富有经验的人，对于他们应当对\Person{地米斯托克利}（如果我还记对了名字的话）形成什么看法，是不会感到困惑的；据说，在一次宴会上，他拒绝弹奏里拉琴，而这原是\Place{希腊}那些杰出而有学问的人惯常做的，并因此被视为没有教养；当他表示对这种消遣的鄙弃时，有人问他：\Quote{那么，听什么令你高兴呢？}据传他回答说：\Quote{我自己的赞美。}善于思考且富有经验的人会立刻看出，如果他们要相信他确实说了这些话，那么他是以什么意图、在什么意义上说出这些话的，或者他们应如何理解这些话；因为他在这个世界的事务中是一位极其杰出的人物，可从他在进一步被追问\Quote{那么，你知道什么？}时所作的回答中得见一斑，\Quote{我知道，}他回答说，\Quote{如何使一个小共和国变得伟大。}至于\Person{恩尼乌斯}所说的对赞美的渴求——\Quote{所有人都渴望受到赞美}，我认为对这种渴求既应部分地予以肯定，也应部分地加以提防。因为一方面，我们应热切渴望真理，真理无疑值得热切追求，作为唯一值得赞美之物，即使它未受赞美；另一方面，我们必须小心避免那种随着人的赞美而容易悄然渗入的虚荣：这种虚荣存在于心中，当那些值得赞美的事物除非因之得到他人的赞美才被视为值得拥有时，或是当一个人因拥有某些他因此渴望受到大力赞美的事物，而这些事物实际上只配得到微小的赞美，甚或可能应受严厉的谴责时。因此，比\Person{恩尼乌斯}观察更为细致的\Person{贺拉斯}说：\Quote{你渴求名望吗？智慧的强大魅力，若诵读三遍，便能使名望的威力解除。}\par

4. 因此，那位诗人认为，由喜爱人的赞美所产生的病患，既已受到他讽刺作品的彻底抨击，就当以具有治愈能力的话语来驱除。那位伟大的导师已藉着祂的宗徒教导我们，我们行善，不应着眼于得人的赞美，也就是说，我们不应以人的赞美作为我们行善的动机；然而，为了人本身的缘故，祂教导我们要寻求他们的赞许。因为当善人受到赞美时，这赞美并不使受赞美的人受益，而是使赞美的人受益。因为就善人本身而言，他们是善的，这已足够；但那些因善人受到他们赞美而乐于效法善人的人，是值得庆贺的，因为他们藉此表明，他们所真诚赞美的人，正是他们所赏识其品行的人。宗徒在某处说：\BibleQuote{如果我仍求人的欢心，我就不是基督的仆役。}\BibleRef{迦 1:10}同一宗徒在另一处说：\BibleQuote{在一切事上使众人喜欢，}并加上理由说：\BibleQuote{不求自己的利益，只求大众的利益，为使他们得救。}\BibleRef{格前 10:33}请看，他为在人的赞美中寻求什么，正如下面这些话所表明的：\BibleQuote{此外，弟兄们！凡是真实的，凡是高尚的，凡是正义的，凡是纯洁的，凡是可爱的，凡是荣誉的，不管是美德，不管是称誉：这一切你们都该思念；凡你们在我身上所学得的，所领受的，所听见的，所看到的：这一切你们都该实行：这样，赐平安的天主必与你们同在。}\BibleRef{斐 4:8}以上我所列举的其他各事，他都用\BibleQuote{美德}一词概括了，说：\BibleQuote{不管是美德；}但他附加上去的定义，\BibleQuote{凡是荣誉的，}则以另一个合适的词\BibleQuote{不管是称誉}来补充。因此，宗徒在这些经文中第一处所说的话：\BibleQuote{如果我仍求人的欢心，我就不是基督的仆役，}应当如此理解，就如他说：如果我行善是以人的赞美为动机，如果我因喜爱赞美而骄傲自满，那我就不是基督的仆役。所以，宗徒愿意使众人喜欢，并因取悦他们而欢喜，并非要自己因他们的赞美而自高，而是要藉着受赞美，在基督内建树他们。那么，我怎会不因你的赞美而喜悦呢？因为你是一个至诚的人，不会言不由衷，而且你将赞美加于你所喜爱的事物上，而这些事物即使并不在我身上，也是有益而值得喜爱的。而且，这不仅对你有益，也对我有益。因为如果它们不在我身上，那么我因羞愧而脸红，并燃起渴望拥有它们的热忱，这对我也是好事。至于你的赞美中那些我认作已拥有的品质，我因拥有它们而喜乐，也因你爱慕这些品质，并因它们的缘故而爱我而喜乐。至于那些我不承认属于我的品质，我不仅渴望获得它们，好使自己拥有，也切望那些真诚爱我的人，不致因这些赞美而一直误会我。\par

5. 请看，我说了这么多，却仍未提及你信中那比你的口才、远比你对我的赞美更令我喜悦的事。卓越的人啊，你以为这会是什么？那就是：我获得了你这样一位如此卓越之人的友谊，甚至尚未与你晤面；的确，如果说我看见了一个人的灵魂，却没看见他的身体，那么我该称他为未见面的人，因为我在他的信函中见到了他的灵魂，却未见到他的身体。在这些信中，我基于自己的了解，而非像从前那样基于弟兄们的见证，形成了对你的看法。因为我早已听闻你的品格如何，但直到现在才知你对我的态度。由此，从你对我的这份友谊，我深信，就连你对我的赞美——我已就其原因说得够多了，它也令我喜乐——必将更加丰盛地裨益基督的教会，因为你拥有、研读、喜爱并推崇我那些为捍卫福音、反对残余不虔敬的偶像崇拜者而写的著作，这使得我的这些作品因你所居的高位而获得更广的影响；因为，你本人既显赫，便自然而然地给它们增添了光彩。你因闻名遐迩，而使它们声名远播；无论你在何处看到这些作品的流传可能有所裨益，你都不会让它们完全湮没无闻。如果你问我如何知道这些，我的回答是，这是通过阅读你的信函，你在我心中产生的印象。于此，你现在将看到，你的信能带给我多么大的喜悦，因为倘若你对我是器重的，你便知道，基督事业的利益能给予我多大的喜悦。此外，当你告诉我有关你自己的事，即如你所说，你出身于一个不止一代两代，甚至远溯祖先，就以能接受基督的教义而著称的家庭，你却仍因我那些反对外邦人礼仪的著作而得到帮助，对这些礼仪有了前所未有的领悟，那么，我岂能认为这是小事一桩：我们的著作，经过你的推荐与传播，可能给别人带来多大的益处，而你的见证又可能使这些著作到达多少且多么显贵的人手中，并通过他们，又如何轻易而有益地到达其他人？或者，想到这些，我心中洋溢的喜悦，岂会是微小或平常的呢？\par

6. 既然我无法用言语表达我从你的信函中获得了何等大的喜乐，我已说明了它令我喜乐的理由；至于我在此事上未能充分言说的，就留给你去揣度了。那么，请接受，我的儿子——接受吧，卓越的人啊，你不仅是由于外在的宣认，更是凭着基督徒的爱而成为基督徒——我再说，请接受包含我的\WorkTitle{忏悔录}的书卷，这正是你所切望得到的。在其中注视我吧，好使你对我的称扬不超越我的本来面目；在其中相信我所讲述的关于我自己的事，不是听信他人，而是听信我自己；在其中默观我，看看我凭自己、在自身上曾是怎样的人；倘若我身上有什么令你喜欢，便为此与我一同赞颂祂，因为我渴望赞颂归于祂，而非归于我自己。因为\BibleQuote{祂造成了我们，我们并非自己造成的}；的确，我们曾自取灭亡，但那造成我们的祂已重造了我们。然而，当你在这些书中发现我的时候，请为我祈祷，使我不致失败，而得成全。祈祷吧，我的儿子，祈祷吧。我感受到我所说的话；我知道我所祈求的。不要让你觉得此事不合适，仿佛超出你的功劳。你若不做，便是剥夺了我的一大助佑。不仅你自己，连所有因你的见证而喜爱我的人，也都要为我祈祷。告诉他们我已恳求了此事，若你看重我们，就当视我们所求的为命令；无论如何，无论是作为允诺请求还是服从命令，都请为我们祈祷。请阅读\WorkTitle{圣经}，你就会发现，宗徒们自己——基督羊群的领袖——也曾向他们的子弟请求此事，或是对他们的听众如此吩咐。至于我，既然你向我请求，我定会尽力为你这样做。那位是祈祷的聆听者看见了这事，祂在我祈求之前便已看见了我为你祈祷；但愿这爱的明证能由你来回报。我们被安置在你们之上；你们是天主的羊群。想想看，我们的危险比你们更大，并为你们祈祷，因为这对我们双方而言都是合宜的，好使我们对你们能向总牧者和万有之首交出一份好的账目，并且能逃避此世的考验及其诱惑，这些诱惑更为危险，除非此世的平安能产生宗徒吩咐我们祈求的那效果：\BibleQuote{为叫我们能以全事虔敬和端庄，度宁静平安的生活}\BibleRef{弟前 2:2}。因为若缺少了虔敬和端庄，那么宁静且平安地免于世间的邪恶又算什么呢？无非是诱人进入或助人追随放纵与沦亡之途的时机罢了。那么，你为我们所求的，就是我们为你们所求的：叫我们能以全事虔敬和端庄，度宁静平安的生活。无论你们在哪里，我们也在哪里，让我们为彼此祈求这事，因为我们属于祂，祂无处不在。\par

7. 我还寄给了你其他一些你未曾索要的书籍，以免我过于拘泥你的请求：我的\WorkTitle{论未见之事}、\WorkTitle{论忍耐}、\WorkTitle{论节制}、\WorkTitle{论天意}，以及一部大部头的\WorkTitle{论信德、望德和爱德}。当你身在非洲时，若读了所有这些书，请将你对它们的意见寄给我，或者寄到某处，以便我的主和弟兄\Person{奥勒留}能转寄给我们；不过，无论你在何处，我们都希望能收到你的信函；也请你期待我们的信函，只要我们还能够。我极感恩地收到了你寄给我的东西，在其中，你惠然援助我，既顾及我的身体健康——因你希望我能在疾病不受妨碍的情况下委身于天主，又顾及我的藏书，使我能购买新书、修补旧书。愿天主在今生和来世，以祂为那些祂所愿意你成为的人预备的恩惠报答你。现在我请求你，如前一样，代我问候那托付给你的和平的抵押，也为我和你至亲至爱者。\par

\SourceRecord{Translated by J.G. Cunningham.；Nicene and Post-Nicene Fathers, First Series；Vol. 1.；Edited by Philip Schaff.；Buffalo, NY: Christian Literature Publishing Co.,；1887.；<http://www.newadvent.org/fathers/1102231.htm>.}

% structure: p0001 source=h1 target=section reason=page-title

\section{书信 232}

\Emph{致\Place{马达乌拉}人民，我应受赞颂的主人们，最亲爱的弟兄们，\Person{奥斯定}问候，答复由\Person{弗洛伦提努斯}弟兄转交的信函。}\par

1. 若偶然间，我所收到的这样一封信是由你们中那些天主教基督徒寄给我的，那么唯一令我惊讶的是，这封信是以自治市的名义寄出的，而非以他们自己的名义。倘若却是你们中全部或几乎全部有地位的人乐意给我写信，那么令我惊讶的是你们在信中用\Quote{父亲}的称号和\Quote{在主内的致意}来称呼我；我对你们确切地知道，并深感遗憾，你们竟以迷信的崇敬去尊崇那些偶像，而你们对之关闭的庙宇远不如你们的心那样容易关闭；或者不如说，那些偶像并非更真实地被关在你们的庙宇里，而是被关在你们的心中。难道你们终于在明智的反省后，认真思考了那在主内的救恩吗？你们正是以主的名义选择向我致意的。因为若非如此，那么我请问你们，我全然应受赞颂的主人们，最亲爱的弟兄们，我究竟在哪件事上伤害了你们，在哪件事上冒犯了你们的善意，竟使你们认为在信首的致意中戏弄我，而非尊重我，是合宜的呢？\par

2. 当我读到\Quote{致\Person{奥斯定}神父，在主内的永恒救恩}这些话时，我骤然满怀希望，以至我相信你们或已归依了主自己，并归依了那由祂创始的永恒救恩，或渴望通过我们的职务来归依。但当我读了信的其他部分，我的心就凉了。于是，我询问了那位送信人，你们是否已经是基督徒，或渴望成为基督徒。从他的答复中得知你们丝毫未变之后，我深为痛心，你们竟认为不仅要弃绝基督的名——你们已经看见全世界都臣服于祂——还要在我这个人身上侮辱祂的名；因为我无法想象除主基督以外的任何主，一位主教能被你们称作父亲，如果对你们用词的含义有任何疑窦，那也会被你们信末的句子所驱散，你们在那里明说：\Quote{我们期望，您的尊驾在圣职团中长久喜乐于天主和祂的基督内。}在读了并思量了这一切之后，我（或者任何人）又能怎么想，除了认为这些话要么是写信者内心的真实表达，要么是出于欺骗的意图？如果你们把这作为内心真实写照来写，那么是谁阻挡了你们通往真理的道路？谁在上面撒了荆棘？什么仇敌将巨石横在你们的路途中？说到底，如果你们渴望进来，是谁将我们圣堂的大门向你们关闭，以致你们不愿意在同一的主内与我们一同享受同一的救恩，而你们正是以同一主的名向我们致意的？但如果你们欺骗地并嘲弄地写这些事，那么，就在你们将你们的事务托付给我之时，竟敢用虚伪奉承的言语侮辱那一位的名，唯有藉着祂我才能做任何事，而不是以应有的崇敬尊崇祂吗？\par

3. 最亲爱的弟兄们，请确信，我是怀着对你们难以言喻的内心战栗写这封信的，因为我知道，如果我对你们徒劳地说了这些话，你们在天主的审判中将有的罪责和惩罚更重。关于人类历史中我们祖先所记下并传给给我们的一切，以及关于寻求和持守真宗教的一切——这些我们现在看见并记录下来传给后代——\BibleQuote{圣经}绝非缄默；远非如此，万事正是按照圣经所预言的实现。你们不能否认你们看见犹太人被从他们祖先的居所逐出，分散到几乎每一个国家：如今，这民族的起源、渐长、王国的丧失、分散全世界，都正如所预言的发生了。你们不能否认你们看见主的话和法律从这个民族经由基督传出，基督在他们的民族中奇迹般诞生，并且这话和法律取得了所有民族的信仰并予以保存：我们现在读到的，早已预先宣告了，正如我们所见。你们不能否认你们看见我们称之为\OriginalTerm{异端}{heresies}和\OriginalTerm{裂教}{schisms}的事，即许多从基督徒团体的根割裂出去的人，而那个团体借着\OriginalTerm{宗座}{Apostolic Sees}和主教们的继承，以无可争辩的世界性扩展广传，他们自称基督徒，如同枯枝夸耀仅仅外表上来自\OriginalTerm{真葡萄树}{true vine}：这一切圣经都已预见、预言并描述过了。你们不能否认你们看见一些偶像的庙宇因荒废而倒塌，另一些被暴力推倒，另一些被关闭，一些被改作他用；你们看见偶像本身或被打碎，或被焚烧，或被锁起，或被毁坏，而同一世界的权势，曾为保卫偶像迫害基督徒，现已被基督徒战胜和征服——基督徒并非为真理而战，而是为之而死——他们转而将攻击和法律指向那些他们曾为保卫之而处死基督徒的偶像；而最崇高帝国中最高显贵，也脱下冠冕，跪伏在渔夫\Person{伯多禄}的墓前恳求。\par

4. 圣经，如今已传到众人手中，很早以前就证实这一切事都必发生。我们看到这一切事都已应验，便欢欣鼓舞，因为我们藉此发现，这些事在圣经中所宣告的权威是何等伟大，我们的信德也就越发坚固。既然这一切都按预言应验了，我问你们：我们难道要以为，在同一圣经中所载的天主的审判——那审判被定为最终将信者与不信者分别开来——是预言中唯一不会应验的事吗？是的，如同这一切事都已来临，审判也必来临。在那日子，我们这个时代的人，没有一个能为自己的不信找到任何辩护的借口。因为基督之名挂在每个人的嘴边：义人履行正义时呼求它，发假誓者行骗时呼求它，君王为巩固统治呼求它，士兵为激励自己作战呼求它，丈夫为确立自己的权威呼求它，妻子为表明自己的顺服呼求它，父亲为加强命令呼求它，儿子为表明服从呼求它，主人为支持治理之权呼求它，奴隶为履行职责呼求它，谦卑者为激起虔敬呼求它，骄傲者为煽动野心呼求它，富人施舍时呼求它，穷人领受施舍时呼求它，醉汉在酒杯前呼求它，乞丐在门口呼求它，好人守信时呼求它，恶人背约时呼求它：所有的人都频繁使用基督之名，基督徒出于真诚的敬畏，外教人出于假装的尊敬；他们将来无疑都要为他们呼求的那一位，就他们提及祂名时所怀的心态和所说的话语，向祂交账。\par

5. 有一位不可见者，我们所见的一切都是由祂——作为造物主和第一因——获得存在：祂是至高、永恒、不变的，除祂自己以外，无人能领会。有一位，至高者借以发言并启示自己，即圣言，祂不次于生祂并发出祂的那一位，借着这圣言，那生祂者得以彰显。有一位是圣，是使一切成圣者的圣化者，祂是那不可分离且不可分割的相互共融，存在于那启示第一因的不变之圣言，与那借其同等之圣言启示自己的第一因之间。然而，有谁能够以完全平静而纯洁的心灵，去默观这整个本质（我努力描述祂，却不提祂的名字，而非仅仅提出名字而不加描述），并从这默观中汲取真福，仿佛沉浸于这种冥想的神往之中，忘却自我，奋力奔赴那超越我们知觉范围的视见；换言之，穿上不朽，并获得那永恒的救恩，就是你们在问候中乐意为我祈求的呢？我要说，谁能做到这一切呢？只有那承认己罪，将骄傲一切虚妄的膨胀都践踏于尘土，并以温良和谦卑俯伏在地，接受天主作他导师的人，才能做到。\par

6. 因此，既然我们必须先从虚幻的自满自足中被贬抑到心灵的谦卑，然后由此上升，达到真正的崇高，那么，要使我们产生这种精神，除了以下方法之外，不可能有更荣耀、更温柔（用劝服而非暴力来制服我们的高傲）的途径了：那就是，圣父借以向天使启示自己、作为祂的德能与智慧、在人心被爱慕有形之物所蒙蔽时无法认出的那位圣言，竟屈尊就卑，摄取我们的本性，并在降生成人时如此行事和显示自己的位格，使人更惧怕因人而骄傲自大，而不是惧怕效法天主而受贬抑。因此，传遍整个世界的基督，不是头戴尘世冠冕的基督，不是富有尘世财宝的基督，不是以世俗昌盛而显赫的基督，而是被钉十字架的基督。这在起初，曾受整个高傲民族的讥笑，至今仍受各民族中少数残余者的讥笑；但起初对少数人，如今对全体民族，它却是信德的对象。因为当那位被钉十字架的基督在当初被宣讲时，尽管受各民族讥笑，对少数相信的人，瘸子得到行走的能力，哑巴开口说话，聋子得以听见，瞎子得以看见，死人得以复活。这样，世俗的骄傲终于被折服，它被说服：即使在今世的事物中，也没有什么比天主的谦卑更强大的了，见：\BibleRef{格前 1:23}，以至于在那神圣榜样的庇护下，那种对人最有益的谦卑，得以抵御骄傲的轻蔑攻击。\par

7. \Place{马道拉}的人啊，我的弟兄们，不，我的父老们，我恳求你们最终醒悟过来：天主赐给了我这个给你们写信的机会。我已尽力为\Person{弗洛伦提努斯}弟兄的事务提供了服务和帮助；你们按天主的旨意，通过他给我写了信。但那事务本不需要我的协助也能轻易办妥，因为几乎所有住在\Place{希波}的他的家人，都认识\Person{弗洛伦提努斯}，并且深切哀悼他的丧亲之痛。然而，你们把信寄给我，是要我在回复时，借你们提供的机会，对崇拜偶像的人说一些有关基督的话，这样就不至于显得我过于冒昧。但我恳求你们，如果你们在那信中并非妄用祂的名，就不要让我写给你们的这些话落空；但如果你们提及祂的名是想嘲弄我，那么，你们当敬畏祂：祂曾为世俗的骄傲所藐视，被当作判了刑的罪犯，而如今，这同一个世界臣服于祂的统治之下，正等待祂作审判者。因为我心中对你们所怀的渴望——我已尽力在这封信中表达——将在祂的审判座前作证反对你们；祂将永远坚立那些相信祂的人，并使不信的人蒙羞。愿唯一真天主完全救你们脱离这世界的虚幻，并使你们归向祂自己，我极可称赞的主人们，极亲爱的弟兄们。\par

\SourceRecord{Translated by J.G. Cunningham.；Nicene and Post-Nicene Fathers, First Series；Vol. 1.；Edited by Philip Schaff.；Buffalo, NY: Christian Literature Publishing Co.,；1887.；<http://www.newadvent.org/fathers/1102232.htm>.}

% structure: p0001 source=h1 target=section reason=page-title

\section{\WorkTitle{书信 245}\newline{}}

\Emph{致\Person{普希迪乌斯}，我最爱的主教、可敬的弟兄及司铎圣职的同工，以及与他同在的弟兄们，\Person{奥古斯丁}和与他同在的弟兄们在主内致以问候。}\par

1. 决定如何处置那些拒绝服从你的人，比查明如何证明他们所行之事为非法，需要更多的思考。然而，在此期间，你的圣善的信函在我极为事务繁忙之际送达，而信差的匆匆离去使我不得不给你回信，却没有时间像应当的那样，就你咨询我的事情作出答复。不过，关于黄金装饰和昂贵服饰的问题，我要说的是，我不希望你轻率地作出禁止使用的决定，除非是那些既未结婚也不打算结婚的人，他们只应考虑如何取悦于天主。但属于世界的人还必须考虑如何在这方面取悦于自己的配偶：丈夫当使妻子喜悦，妻子当使丈夫喜悦；\BibleRef{格前 7:32} 不过有限制，就是已婚妇女也不宜露发，因为宗徒命令妇女要蒙头。\BibleRef{格前 11:5} 至于妇女用脂粉涂面，以求肤色红润或白皙，这是一种不诚实的诡计，我确信连她们自己的丈夫也不愿受其欺骗；而妇女仅为了自己的丈夫才应被允许修饰，这是圣经的宽容，而非命令。因为真正的装饰，尤其是基督徒男女的装饰，不仅在于完全抛弃一切面部涂画的欺骗，更不在于拥有华丽的黄金首饰或贵重服饰，而在于无可指摘的生活。\par

2. 至于佩戴护身符的可咒的迷信（其中包括男子单侧耳朵上部所戴的耳环），这种做法不是为了取悦于人，而是为了敬拜魔鬼。但有谁能指望在圣经中找到对每一种邪恶迷信的明确禁止呢？因为宗徒概括地说：\BibleQuote{我不愿意你们与魔鬼有分，}\BibleRef{格前 10:20} 又说：\BibleQuote{基督与贝里雅耳有什么协调？}\BibleRef{格后 6:15} 除非或许因为他提了\Person{贝里雅耳}，同时概括地禁止与魔鬼有份，就留给基督徒去祭祀\Person{乃普顿}的开放性，因为我们没有在哪处读到明确禁止敬拜\Person{乃普顿}！同时，要告诫这些不幸的人，若他们固执地不服从有益的训令，他们至少必须停止为这些不虔敬的行为辩护，以免陷入更大的罪过。但是我们又何必与他们争辩呢，如果他们害怕摘掉耳环，却戴着魔鬼的标记去领受基督的身体而不害怕？\par

关于祝圣一个在\OriginalTerm{陶纳派}{Donatist}中受洗的人，我不能承担劝你这样做的责任；你在别无选择的情况下这样做是一回事，我建议你这样做则是另一回事。\par

\SourceRecord{Translated by J.G. Cunningham.；Nicene and Post-Nicene Fathers, First Series；Vol. 1.；Edited by Philip Schaff.；Buffalo, NY: Christian Literature Publishing Co.,；1887.；<http://www.newadvent.org/fathers/1102245.htm>.}

% structure: p0001 source=h1 target=section reason=page-title

\section{书信 246}

\Emph{\Person{奥古斯丁}致\Person{兰帕迪乌斯}的问候}\par

1. 关于命运与机缘这个主题，我在与你在一起时就已经察觉，如今更从你的来信中满心欣慰地得知，你的心思因此严重不安；我本当为你写下相当篇幅的论述；主将使我能够以祂所知最适宜保守你信德的方式来阐述。这可不是小恶：当人们接受错谬意见时，他们不仅被享乐的诱惑引诱去犯罪，还转而为自己的罪辩护，而不是通过承认自己做错了来寻求医治。\par

2. 因此，请允许我简要提醒你一件与这问题相关、你肯定知道的事：一切法律、一切管束手段——褒奖、谴责、劝诫、恐吓、赏赐、惩罚，以及治理统御人类的其他所有事物——若意志不是人所犯罪恶的原因，便全然被颠覆推翻，且处于绝对没有正义的境地。因此，对我们来说，拒斥\OriginalTerm{占星家}{mathematici}的荒谬之说，比接受不得不谴责并拒斥那源于神圣权威的法律、甚或治理我们自己家庭所必需的手段这样的后果，要正当合理得多！就在这一点上，占星家自己也无视他们关于命运与机缘的学说；因为他们中任何一人向有钱的愚人售卖了关于命运的愚蠢预言之后，将心思从象牙牌板收回，转去管理和照料自己的家室，一旦发现妻子——我且不说轻浮嬉笑，即便只是向窗外张望过多——便不仅用言语，更用拳打脚踢来责罚她。然而，若是妻子在这样的情形下争辩说：\Quote{为什么要打\Emph{我}？你若能够，不如去打维纳斯吧，因为正是在那颗星的影响下，我才被迫做了你所抱怨的事。}——他必定会全力以赴，不是去编造那些蒙骗公众的荒谬行话，而是施行一些公正的纠正，以维持他在家中的权威。\par

3. 因此，当有人受责备时断言命运是行为的原因，并坚称自己不该受责备，因为他是在命运的强迫下做了那受谴责的事，那么让他回到他自己的事上应用这个原则：叫他不要责罚不忠的仆人，不要抱怨无礼的儿子，不要向惹事生非的邻舍发出威胁。因为，如果他遭受所有这些人的伤害都不是出于他们自身的过错，而是命运驱使所致，那么他做这些事中的哪一件是正义的呢？然而，如果他凭着内心固有的责任感，以及作为家长对眼下受他管辖的众人所负的义务，劝勉人行善，警戒人避恶，命令人服从他的旨意，奖赏绝对顺服的人，惩罚轻慢他的人，感谢善待他的人，憎恨忘恩负义的人——当我发现他用行动（而非言语）如此斩钉截铁地驳斥自己的虚张声势，似乎亲手拔光了占星家满头的每一根头发时，还需要我去证明占星家对命运的推算何其荒谬吗？\par

倘若你热切的渴求不满足于这寥寥数语，而是需要一本需花更长时间诵读的论著来探讨这主题，你必须耐心等待，直到我从其他事务中获得稍许喘息；并且你要向天主祈祷，求祂乐意赐予写作所需的闲暇与能力，好使你的心灵在这事上得享安宁。然而，若你的爱德不吝于频繁写信提醒我，并在回信中告诉我你对这封信的看法，我将更欣然乐意地去做这件事。\par

\SourceRecord{Translated by J.G. Cunningham.；Nicene and Post-Nicene Fathers, First Series；Vol. 1.；Edited by Philip Schaff.；Buffalo, NY: Christian Literature Publishing Co.,；1887.；<http://www.newadvent.org/fathers/1102246.htm>.}

% structure: p0001 source=h1 target=section reason=page-title

\section{书信 250}

\Emph{致他所爱的尊长、可敬的弟兄及司铎职务的同工\Person{奥克西留斯}，\Person{奥思定}在主内致候。}\par

1. 我们的儿子\Person{克拉西亚努斯}，一位有地位的人，寄给我一封信，痛苦地抱怨他受到了你的圣善不公正的绝罚。他对此事的叙述是：他带着一小队适合他官职权力的随从来到教会，恳求你不要损害他们灵性上的利益，不要将圣所的特权扩展给那些人；这些人曾违背他们在福音前所发的誓愿，又在信仰之家内寻求他们背信罪行的帮助和保护；随后，这些人反省自己所犯的罪，自愿地，而不是被强行逼迫，离开了教会；而因为这件事，你的圣善对他非常不悦，以教会的惯常程序，将他和他全家处以绝罚。\par

读到他这封信，我内心十分不安，思绪如汹涌海浪般翻腾，感到不能不写信给你，恳请你在彻底审查了你对此事的判断，并用无可辩驳的推理或圣经证据证明之后，也能善意地教导我，究竟有何根据，竟使儿子因父亲的罪、妻子因丈夫的罪、仆人因主人的罪而受绝罚是正义的，或者，甚至未出生的婴儿，如果出生在一个全家都处于绝罚禁令之下的家庭时，即使死亡临近，也被剥夺通过\OriginalTerm{重生之洗}{laver of regeneration}获得的救援，这又如何是正义的。因为这不是那些只影响身体的审判，在那些审判中，如我们在圣经中所读到的，一些蔑视天主的人连同他们全家都被处死，尽管这些家人并未参与他们的不敬。在这些情况下，确实，作为对幸存者的警告，死亡施加于那些本来终有一死的\Emph{身体}；但灵性的审判，基于所记载的，\BibleQuote{凡你在地上所束缚的，在天上也要被束缚}\BibleRef{玛 16:19}——约束的是\Emph{灵魂}，关于灵魂，经上说：\BibleQuote{父亲的生命属于我，儿子的生命也属于我：谁犯罪，谁丧亡。}\BibleRef{则 18:14}\par

2. 你可能听说过，其他享有盛誉的司铎在某些情况下，将罪人的全家也包括在对他所宣布的绝罚之中；但如果需要，他们或许能够给出这样做的充分理由。至于我自己，尽管对于某些人极端残酷地骚扰教会的行为，我曾感到极度忧虑，但我从未敢像你那样做，因为如果有人要求我为此举辩护，我无法做出令人满意的答复。不过，如果主启示你这可能是正义之举，我绝不会轻视你的年轻和缺乏经验，因为你最近才被擢升到教会的高位。请看，尽管我已年迈，却愿向年轻人学习；尽管我已任主教多年，却愿向任同一圣职尚不足一年的你学习，如果你能教导我，我们如何能在人或天主面前，为我们在无辜灵魂上施加灵性惩罚的行为辩护，而这些灵魂并未以同样的方式参与他人的罪行，如同他们参与\Person{亚当}的原罪那样，因为在他内\Quote{众人都犯了罪}。虽然\Person{克拉西亚努斯}的儿子通过他的父亲，从我们的始祖那里继承了需要藉由\OriginalTerm{重生之洗}{laver of regeneration}洗净的罪愆，但谁会有片刻犹豫地说，他绝不为父亲可能犯下的任何罪行负责，因为他出生时并未参与？对于他的妻子又该怎么说？对于全家那么多灵魂又该怎么说？——如果因着这包括全家在内的绝罚的严厉，其中哪怕只有一个灵魂，在未受圣洗的情况下离开肉身而丧亡，那么所造成的损失，将比无数即使是无辜者也被从圣所庭院拖出并被杀害的肉体死亡，更为不可比拟的巨大灾难。因此，如果你能为此给出充分的理由，我相信你会在回信中令我知晓，使我也能如此行；但如果你不能，你又有何权利，在轻率激动的驱使下，做出你若被要求提供令人满意的理由时却无从找到的事情呢？\par

3. 我至今所说的话，适用于这种情况：即便假设我们的孩子\Person{Classicianus}做了某事，而这事似乎完全有理由使你对他施以\OriginalTerm{绝罚}{excommunication}的惩罚。但若他寄给我的信中所言为真，那么就连他本人（即使他的家人得以免于此罚）也不应受到这样的惩罚。然而，关于此事，我并不干预圣善的您；我只恳求您，当他请求宽恕且承认己过时，就原谅他；另一方面，若您经过反省，承认他并未做错什么，因为事实上公理更在他这一边：他恳切要求在信仰之家内，信德应被神圣地持守，且不应在每日教导背信之罪恶的地方遭破坏——那么，在这种情形下，就做一位虔诚者所当做的事吧——也就是说，若您作为人，遭遇了如同那位真正属天主的人在圣咏中所宣认的事：\BibleQuote{我的眼睛因忿怒而昏花}，那么您也不可不效法他向主呼求：\BibleQuote{主，求你可怜我，因为我软弱}，好使祂向您伸出右手，斥责您情绪的风暴，使您心灵平静，从而能看清并践行正义；因为，正如经上所说：\BibleQuote{人的忿怒，并不成全天主的正义。} \BibleRef{雅 1:20}不要以为因我们是主教，不公的激愤就不可能暗中侵袭我们；我们当更记住，正因为我们是人，处于诱惑的罗网中，生命就满布极大的危险。因此，取消那或许是在异常激动下作出的教会判决；让你们彼此之间那从您是\OriginalTerm{慕道者}{catechumen}时就已联合的友爱，重新恢复；驱逐纷争，迎回和平，免得您失去这位朋友，而憎恨你们的魔鬼却在你们双方身上幸灾乐祸。我们天主的仁慈是伟大的；祂的怜悯或许会俯听我的祈祷——我恳求祂使我对您的忧伤不再加深，反而将我已开始承受的痛苦移除；愿您的青年时代，不轻视我的老年，藉着祂的恩宠，得到鼓励并充满喜乐！再会。\par

\EditorialNote{附在这封信后面的，是一封同时写给\Person{Classicianus}的信的残篇；内容如下：}\par

[为了制止那些因一人之罪便将违犯者的全家——即众多灵魂——置于同一条\OriginalTerm{绝罚}{excommunication}判决之下，特别是为了避免有人因这种\OriginalTerm{绝罚}{anathema}而在未领受圣洗的情况下离开此生；也为了决定一个问题：对那些逃入教堂，为要背弃向保证人所许之信德的人，是否不应将其逐出教堂——我希望靠着主的帮助，在我们的会议中采取必要措施，并在必要时致书\OriginalTerm{宗座}{Apostolic See}；以便通过全体一致、权威的决定，确定并规定在这些情况下应遵循的途径。我特意指出一件事，作为不容置疑的真理：若有信徒被不公正地施以绝罚，则这判决对宣告者的伤害，大于对受此冤屈者的伤害。因为任何人被释放或束缚，都是藉着住在圣者内的圣神，而祂绝不将不应得的惩罚加于任何人；因为藉着祂，那不作恶的爱德已倾注在我们心中。]\par

\SourceRecord{Translated by J.G. Cunningham.；Nicene and Post-Nicene Fathers, First Series；Vol. 1.；Edited by Philip Schaff.；Buffalo, NY: Christian Literature Publishing Co.,；1887.；<http://www.newadvent.org/fathers/1102250.htm>.}

% structure: p0001 source=h1 target=section reason=page-title

\section{书信二五四}

\Emph{致\Person{贝内纳图斯}，我至福的主教，我所尊敬而亲爱的弟兄及\OriginalTerm{司铎职务}{priestly office}上的同工，以及与他在一起的弟兄们；\Person{奥思定}及与他在一起的弟兄们，在主内致候。}\par

关于主教阁下写信给我提及的那位少女，她目前的想法是，如果她已成年，她将拒绝一切婚配的提议。然而她还如此年轻，即使她倾向于婚姻，现在也还不该将她许配或订婚给任何人。此外，\Person{贝内纳图斯}主教，我所尊敬和亲爱的弟兄，必须记得天主将她置于祂的教会的监护之下，旨在保护她免遭恶人侵害；因此将她托付给我照顾，并非可以由我随意将她给与我所选中的任何人，而是使任何不适合的人不能违我意愿将她夺走。您乐于提及的那一婚议，如果她愿意并准备结婚的话，并不会令我不悦；但至于她是否会嫁给任何人——虽然就我本人而言，我更希望她实现她现在所说的愿望——我暂时还不得而知，因为在她这个年纪，她说希望成为修女的声明，与其说是发下庄重誓愿者的深思熟虑的宣示，不如说是一个随口乱说者的轻率之言。此外，她有一位姨母，嫁给了我们可敬的弟兄\Person{斐理斯}，我已就此事与他商议——因为我既不能、也确实不该避免与他商量——他并未不愿接受此婚议，反而表示满意；但他不无理由地表达了不满，因为尽管他作为亲属有权获知此事，却没有人就此事写信给他。因为，或许少女的母亲也会出面，尽管目前她尚不现身，而在我看来，在嫁女之事上，母亲的心愿自然应优先于其他所有人，除非少女本人已足够年长，得以合法地有更强的主张权选择自己所愿。我也愿阁下明了，倘若她在婚姻上的最终和全权都托付给我，而她本人业已成年并愿意结婚，在天主作为我的审判者的注视下将自己托付给我，由我将她许配给我认为最佳之人选——我声明，并且是确确实实地声明，您提及的婚议暂时令我满意，但因为天主是我的审判者，我不能保证在她有更好婚议出现时为她而拒绝；但将来是否有任何这样的婚议出现，全然无法确定。因此，主教阁下明见，有多少重要因素共同使她目前不可能确定地许给任何人。\par

\SourceRecord{Translated by J.G. Cunningham.；Nicene and Post-Nicene Fathers, First Series；Vol. 1.；Edited by Philip Schaff.；Buffalo, NY: Christian Literature Publishing Co.,；1887.；<http://www.newadvent.org/fathers/1102254.htm>.}

% structure: p0001 source=h1 target=section reason=page-title

\section{书信263}

\Emph{向极其虔诚的女士、圣善的女儿\Person{萨皮达}，\Person{奥古斯丁}在主内致以问候。}\par

我接受了那件由你公正而虔诚的巧手预备、并蒙你亲切赠予我的礼物，免得加增那——我看出——更需要我安慰的人的忧伤；尤其因为你表示，若我穿上这件你为那圣善的天主之仆、你的弟兄所缝制的袍子，对你来说是不小的安慰，因为他已离开将死之地，不再需要那些会因使用而朽坏之物。因此，我顺从了你的愿望，无论你由此感受到何种安慰，我都不拒绝你对弟兄的挚爱。因此，我已收下你寄来的袍子，并在写信给你之前，就已开始穿它。所以，你当心怀喜乐；但我恳求你，致力于远更美好、远更伟大的安慰，好让那因人性的软弱而密布你心的阴云，被神圣权威的话语驱散；并且，时刻如此生活，好能与你弟兄一同生活，因为他已如此死，却仍活着。\par

确实让人流泪的是，你那爱你的弟兄，并以你的虔诚生活和\OriginalTerm{献身贞女}{consecrated virgin}的身分特别敬重你的弟兄，不再像从前那样在你的眼前，勤勉地执行\Place{迦太基}教会\OriginalTerm{执事}{deacon}的教会职务，进进出出；你再也不能从他口中听到，他惯常用亲切、虔诚而合宜的挚爱之情，对他如此亲爱的姊妹的圣德，所作的尊崇见证。当这些事被反复思量，并以积久弥深的爱全部热切地追念时，心灵被刺透，泪水便如刺透之心的血，流淌不止。但让你的心举目向天，你的眼泪就会停止。你哀悼失去的那些事，确实已经消逝，因为它们本质上是暂时的，但失去它们并不消灭\Person{弟茂德}对[姊妹]\Person{萨皮达}的爱，他仍然爱着她：这爱存留在自己的宝库中，与基督一同藏在天主内。守财奴将黄金藏于隐秘之处，难道就失去了黄金吗？当他将黄金保存在某个更安全的藏匿处，甚至让他眼目不见，难道他岂不因此，在他力所能及的范围内，更确知黄金为他所有吗？世俗的贪恋，在不再看到它所爱的宝物时，相信自己找到了更稳妥的守护；而天上的爱，难道要将那不过是提前送往天上粮仓的，当作永远失去而悲伤吗？\Person{萨皮达}啊，你当将自己完全献于崇高的召叫，思念天上的事，在那里，基督坐在天主的右边，祂曾为我们屈尊就死，好使我们虽死犹生，并确保没有人会害怕死亡，仿佛死亡注定要毁灭他；也没有人会在那生命为之而死的人死后，为他哀哭，仿佛他已失去生命。你当以此等以及类似的属天安慰为你的慰藉，让世俗的忧伤为之赧颜而逃。\par

凡人因挚爱之人的死亡而忧伤，其中并无当受责备之处；但信徒的忧伤不应长久。因此，倘若你至今一直忧伤，愿这忧伤止息，不要像那些没有希望的\OriginalTerm{外教人}{heathen}一样忧伤，\BibleQuote{没有望德的人}。\BibleRef{得前 4:13} 因为\Person{保禄}宗徒说这话时，并非完全禁止忧伤，而只是禁止那些没有希望的外教人所表露的忧伤。就连\Person{玛尔大}和\Person{玛利亚}，两位虔诚的姊妹和信徒，也为自己知道将会复活的弟兄\Person{拉匝禄}哀哭，尽管她们不知道他当时就会恢复生命；而主自己也为同一个拉匝禄流泪，祂即将使他从死中复活；\BibleRef{若 11:19} 在其中，祂无疑以榜样允许了，而不是以任何命令规诫了，对那些我们相信会复活获得真生命的人，在坟墓前流泪。而且，无怪乎\BibleBook{德训}{篇}上说：\BibleQuote{让眼泪流下，在死者身上哀悼，就如遭受大难似的，开始悲伤吧}；但稍后又补充了这劝告：\BibleQuote{然后安慰自己，平息忧苦。因为忧伤致死，心中忧闷，摧残生命。} \BibleRef{德 38:16}\par

4. 你的兄弟，我的女儿，在灵魂上是活着的，在身体上是安眠的：\Quote{那睡眠的，岂不也要醒来吗？} 天主，既已接纳了他的灵魂，也必将他的身体归还给他；天主并没有将他的身体拿走以毁灭，而仅仅是暂时取去，为要使他重新拥有。因此，没有理由长久悲伤，却有更坚强的理由永远喜乐。因为即使你兄弟那已埋葬于地下的可朽部分，也必不会永远失去——正是借着这一部分，他曾显然地与你同住，向你说话，与你交谈，借这声音发言，而他的声音之为你耳朵所熟悉，并不亚于他的面容之为你眼睛所见；因此，无论何处听到他说话的声音，即使看不见他，你也能立刻辨认出他来。这些东西确实已被收去，不再为活人的感官所觉知，以致死者的离去使尚存的朋友哀悼他们。然而，连死人的身体也必不灭亡（正如连一根头发也不会失落，\BibleRef{路 21:18}），它们被暂时搁置后，将重新被取回，不再被搁置，最终固定在已被改变的更卓越的生命状态中；那么，相对于为短暂片时忧伤，岂不更有理由为无可限量的永恒满怀感恩吗？这希望外邦人没有，因为他们不认识圣经，也不认识天主的能力\BibleRef{玛 22:29}；天主能恢复那失去的，使死亡的复生，使受腐朽的更新，使分开者重聚，并从此永远保全那原先是可朽而短命的事物。这些事，天主已许下，且借着对其他许诺的实践，已给予我们的信德以充分的根据，以相信这些事也必将实现。你的信德如今当常常向你论述这些事，因为你的望德必不会落空，尽管你的爱德现在或许暂时无法施展；默想这些事吧，在它们之中寻得更坚实而丰富的安慰。因为倘若我现在穿上（因他未能穿上的）那件你为你兄弟所织的外衣这一事实，能给你带来些许安慰，那么，当你想到这衣服原为他所预备，而他那时并不需要这不会朽坏的衣服，但他将穿上不朽和不死之时，你所得的安慰将是何等充实而满足啊！\par

\SourceRecord{Translated by J.G. Cunningham.；Nicene and Post-Nicene Fathers, First Series；Vol. 1.；Edited by Philip Schaff.；Buffalo, NY: Christian Literature Publishing Co.,；1887.；<http://www.newadvent.org/fathers/1102263.htm>.}

% structure: p0001 source=h1 target=section reason=page-title

\section{信函269}

\Emph{致\Person{诺比利乌斯}，我最蒙福且可敬的兄弟、司铎职务中的伙伴，\Person{奥古斯丁}致以问候。}\par

您兄弟般的盛情邀请我参加的庆典如此重要，以致若非体弱使我无法成行，我内心的渴望定会带我这可怜的身体到您那里去。若此时非冬日，我本可成行；若我尚年少，我本可无畏寒冬：因为若是年少，青春的温暖会毫无怨言地承受季节的寒冷；若是夏日，夏季的温暖会温和地面对老年的寒气与倦怠。目前，我最蒙福的阁下，我圣善而可敬的司铎职务伙伴，我无法在冬日进行如此漫长的旅程，因为我必须带着因极多年岁而来的冰冷虚弱之躯。我以您应得的问候回敬，为自身的福祉，我请求您在祈祷中纪念我，我自己也恳求主天主，使如此宏伟的圣堂为敬礼祂而举行的奉献礼带来和平的繁荣。\par

\SourceRecord{Translated by J.G. Cunningham.；Nicene and Post-Nicene Fathers, First Series；Vol. 1.；Edited by Philip Schaff.；Buffalo, NY: Christian Literature Publishing Co.,；1887.；<http://www.newadvent.org/fathers/1102269.htm>.}
